\documentclass[sigconf]{acmart}
\usepackage{pgf}
\usepackage[T1]{fontenc}
\usepackage{libertine}
\settopmatter{printacmref=false}
\pagestyle{plain}
\begin{document}
\title{Paper figure gallery (placeholder data)}
\author{fm-switching}
\maketitle
\noindent Each figure is shown at its true size. Status tags: OLD DATA, DEV DATA, or NO DATA YET where a figure
has no committed data. Captions start with the status placeholder.
\begin{figure}[t]
\centering
{\fontfamily{LinuxBiolinumT}\selectfont
%% Creator: Matplotlib, PGF backend
%%
%% To include the figure in your LaTeX document, write
%%   \input{<filename>.pgf}
%%
%% Make sure the required packages are loaded in your preamble
%%   \usepackage{pgf}
%%
%% Also ensure that all the required font packages are loaded; for instance,
%% the lmodern package is sometimes necessary when using math font.
%%   \usepackage{lmodern}
%%
%% Figures using additional raster images can only be included by \input if
%% they are in the same directory as the main LaTeX file. For loading figures
%% from other directories you can use the `import` package
%%   \usepackage{import}
%%
%% and then include the figures with
%%   \import{<path to file>}{<filename>.pgf}
%%
%% Matplotlib used the following preamble
%%
\begingroup%
\makeatletter%
\begin{pgfpicture}%
\pgfpathrectangle{\pgfpointorigin}{\pgfqpoint{3.330000in}{1.950000in}}%
\pgfusepath{use as bounding box, clip}%
\begin{pgfscope}%
\pgfsetbuttcap%
\pgfsetmiterjoin%
\definecolor{currentfill}{rgb}{1.000000,1.000000,1.000000}%
\pgfsetfillcolor{currentfill}%
\pgfsetlinewidth{0.000000pt}%
\definecolor{currentstroke}{rgb}{1.000000,1.000000,1.000000}%
\pgfsetstrokecolor{currentstroke}%
\pgfsetdash{}{0pt}%
\pgfpathmoveto{\pgfqpoint{0.000000in}{0.000000in}}%
\pgfpathlineto{\pgfqpoint{3.330000in}{0.000000in}}%
\pgfpathlineto{\pgfqpoint{3.330000in}{1.950000in}}%
\pgfpathlineto{\pgfqpoint{0.000000in}{1.950000in}}%
\pgfpathlineto{\pgfqpoint{0.000000in}{0.000000in}}%
\pgfpathclose%
\pgfusepath{fill}%
\end{pgfscope}%
\begin{pgfscope}%
\pgfsetbuttcap%
\pgfsetmiterjoin%
\definecolor{currentfill}{rgb}{1.000000,1.000000,1.000000}%
\pgfsetfillcolor{currentfill}%
\pgfsetlinewidth{0.000000pt}%
\definecolor{currentstroke}{rgb}{0.000000,0.000000,0.000000}%
\pgfsetstrokecolor{currentstroke}%
\pgfsetstrokeopacity{0.000000}%
\pgfsetdash{}{0pt}%
\pgfpathmoveto{\pgfqpoint{1.282050in}{0.321750in}}%
\pgfpathlineto{\pgfqpoint{3.230100in}{0.321750in}}%
\pgfpathlineto{\pgfqpoint{3.230100in}{1.686750in}}%
\pgfpathlineto{\pgfqpoint{1.282050in}{1.686750in}}%
\pgfpathlineto{\pgfqpoint{1.282050in}{0.321750in}}%
\pgfpathclose%
\pgfusepath{fill}%
\end{pgfscope}%
\begin{pgfscope}%
\pgfpathrectangle{\pgfqpoint{1.282050in}{0.321750in}}{\pgfqpoint{1.948050in}{1.365000in}}%
\pgfusepath{clip}%
\pgfsetrectcap%
\pgfsetroundjoin%
\pgfsetlinewidth{0.401500pt}%
\definecolor{currentstroke}{rgb}{0.901961,0.901961,0.901961}%
\pgfsetstrokecolor{currentstroke}%
\pgfsetdash{}{0pt}%
\pgfpathmoveto{\pgfqpoint{1.498500in}{0.321750in}}%
\pgfpathlineto{\pgfqpoint{1.498500in}{1.686750in}}%
\pgfusepath{stroke}%
\end{pgfscope}%
\begin{pgfscope}%
\pgfsetbuttcap%
\pgfsetroundjoin%
\definecolor{currentfill}{rgb}{0.101961,0.101961,0.101961}%
\pgfsetfillcolor{currentfill}%
\pgfsetlinewidth{0.803000pt}%
\definecolor{currentstroke}{rgb}{0.101961,0.101961,0.101961}%
\pgfsetstrokecolor{currentstroke}%
\pgfsetdash{}{0pt}%
\pgfsys@defobject{currentmarker}{\pgfqpoint{0.000000in}{-0.027778in}}{\pgfqpoint{0.000000in}{0.000000in}}{%
\pgfpathmoveto{\pgfqpoint{0.000000in}{0.000000in}}%
\pgfpathlineto{\pgfqpoint{0.000000in}{-0.027778in}}%
\pgfusepath{stroke,fill}%
}%
\begin{pgfscope}%
\pgfsys@transformshift{1.498500in}{0.321750in}%
\pgfsys@useobject{currentmarker}{}%
\end{pgfscope}%
\end{pgfscope}%
\begin{pgfscope}%
\definecolor{textcolor}{rgb}{0.101961,0.101961,0.101961}%
\pgfsetstrokecolor{textcolor}%
\pgfsetfillcolor{textcolor}%
\pgftext[x=1.498500in,y=0.245361in,,top]{\color{textcolor}\sffamily\fontsize{7.000000}{8.400000}\selectfont 0}%
\end{pgfscope}%
\begin{pgfscope}%
\pgfpathrectangle{\pgfqpoint{1.282050in}{0.321750in}}{\pgfqpoint{1.948050in}{1.365000in}}%
\pgfusepath{clip}%
\pgfsetrectcap%
\pgfsetroundjoin%
\pgfsetlinewidth{0.401500pt}%
\definecolor{currentstroke}{rgb}{0.901961,0.901961,0.901961}%
\pgfsetstrokecolor{currentstroke}%
\pgfsetdash{}{0pt}%
\pgfpathmoveto{\pgfqpoint{1.979500in}{0.321750in}}%
\pgfpathlineto{\pgfqpoint{1.979500in}{1.686750in}}%
\pgfusepath{stroke}%
\end{pgfscope}%
\begin{pgfscope}%
\pgfsetbuttcap%
\pgfsetroundjoin%
\definecolor{currentfill}{rgb}{0.101961,0.101961,0.101961}%
\pgfsetfillcolor{currentfill}%
\pgfsetlinewidth{0.803000pt}%
\definecolor{currentstroke}{rgb}{0.101961,0.101961,0.101961}%
\pgfsetstrokecolor{currentstroke}%
\pgfsetdash{}{0pt}%
\pgfsys@defobject{currentmarker}{\pgfqpoint{0.000000in}{-0.027778in}}{\pgfqpoint{0.000000in}{0.000000in}}{%
\pgfpathmoveto{\pgfqpoint{0.000000in}{0.000000in}}%
\pgfpathlineto{\pgfqpoint{0.000000in}{-0.027778in}}%
\pgfusepath{stroke,fill}%
}%
\begin{pgfscope}%
\pgfsys@transformshift{1.979500in}{0.321750in}%
\pgfsys@useobject{currentmarker}{}%
\end{pgfscope}%
\end{pgfscope}%
\begin{pgfscope}%
\definecolor{textcolor}{rgb}{0.101961,0.101961,0.101961}%
\pgfsetstrokecolor{textcolor}%
\pgfsetfillcolor{textcolor}%
\pgftext[x=1.979500in,y=0.245361in,,top]{\color{textcolor}\sffamily\fontsize{7.000000}{8.400000}\selectfont 100}%
\end{pgfscope}%
\begin{pgfscope}%
\pgfpathrectangle{\pgfqpoint{1.282050in}{0.321750in}}{\pgfqpoint{1.948050in}{1.365000in}}%
\pgfusepath{clip}%
\pgfsetrectcap%
\pgfsetroundjoin%
\pgfsetlinewidth{0.401500pt}%
\definecolor{currentstroke}{rgb}{0.901961,0.901961,0.901961}%
\pgfsetstrokecolor{currentstroke}%
\pgfsetdash{}{0pt}%
\pgfpathmoveto{\pgfqpoint{2.460500in}{0.321750in}}%
\pgfpathlineto{\pgfqpoint{2.460500in}{1.686750in}}%
\pgfusepath{stroke}%
\end{pgfscope}%
\begin{pgfscope}%
\pgfsetbuttcap%
\pgfsetroundjoin%
\definecolor{currentfill}{rgb}{0.101961,0.101961,0.101961}%
\pgfsetfillcolor{currentfill}%
\pgfsetlinewidth{0.803000pt}%
\definecolor{currentstroke}{rgb}{0.101961,0.101961,0.101961}%
\pgfsetstrokecolor{currentstroke}%
\pgfsetdash{}{0pt}%
\pgfsys@defobject{currentmarker}{\pgfqpoint{0.000000in}{-0.027778in}}{\pgfqpoint{0.000000in}{0.000000in}}{%
\pgfpathmoveto{\pgfqpoint{0.000000in}{0.000000in}}%
\pgfpathlineto{\pgfqpoint{0.000000in}{-0.027778in}}%
\pgfusepath{stroke,fill}%
}%
\begin{pgfscope}%
\pgfsys@transformshift{2.460500in}{0.321750in}%
\pgfsys@useobject{currentmarker}{}%
\end{pgfscope}%
\end{pgfscope}%
\begin{pgfscope}%
\definecolor{textcolor}{rgb}{0.101961,0.101961,0.101961}%
\pgfsetstrokecolor{textcolor}%
\pgfsetfillcolor{textcolor}%
\pgftext[x=2.460500in,y=0.245361in,,top]{\color{textcolor}\sffamily\fontsize{7.000000}{8.400000}\selectfont 200}%
\end{pgfscope}%
\begin{pgfscope}%
\pgfpathrectangle{\pgfqpoint{1.282050in}{0.321750in}}{\pgfqpoint{1.948050in}{1.365000in}}%
\pgfusepath{clip}%
\pgfsetrectcap%
\pgfsetroundjoin%
\pgfsetlinewidth{0.401500pt}%
\definecolor{currentstroke}{rgb}{0.901961,0.901961,0.901961}%
\pgfsetstrokecolor{currentstroke}%
\pgfsetdash{}{0pt}%
\pgfpathmoveto{\pgfqpoint{2.941500in}{0.321750in}}%
\pgfpathlineto{\pgfqpoint{2.941500in}{1.686750in}}%
\pgfusepath{stroke}%
\end{pgfscope}%
\begin{pgfscope}%
\pgfsetbuttcap%
\pgfsetroundjoin%
\definecolor{currentfill}{rgb}{0.101961,0.101961,0.101961}%
\pgfsetfillcolor{currentfill}%
\pgfsetlinewidth{0.803000pt}%
\definecolor{currentstroke}{rgb}{0.101961,0.101961,0.101961}%
\pgfsetstrokecolor{currentstroke}%
\pgfsetdash{}{0pt}%
\pgfsys@defobject{currentmarker}{\pgfqpoint{0.000000in}{-0.027778in}}{\pgfqpoint{0.000000in}{0.000000in}}{%
\pgfpathmoveto{\pgfqpoint{0.000000in}{0.000000in}}%
\pgfpathlineto{\pgfqpoint{0.000000in}{-0.027778in}}%
\pgfusepath{stroke,fill}%
}%
\begin{pgfscope}%
\pgfsys@transformshift{2.941500in}{0.321750in}%
\pgfsys@useobject{currentmarker}{}%
\end{pgfscope}%
\end{pgfscope}%
\begin{pgfscope}%
\definecolor{textcolor}{rgb}{0.101961,0.101961,0.101961}%
\pgfsetstrokecolor{textcolor}%
\pgfsetfillcolor{textcolor}%
\pgftext[x=2.941500in,y=0.245361in,,top]{\color{textcolor}\sffamily\fontsize{7.000000}{8.400000}\selectfont 300}%
\end{pgfscope}%
\begin{pgfscope}%
\definecolor{textcolor}{rgb}{0.101961,0.101961,0.101961}%
\pgfsetstrokecolor{textcolor}%
\pgfsetfillcolor{textcolor}%
\pgftext[x=0.581831in, y=1.523118in, left, base]{\color{textcolor}\sffamily\fontsize{7.000000}{8.400000}\bfseries\selectfont \textbf{COVEX (ours)}}%
\end{pgfscope}%
\begin{pgfscope}%
\definecolor{textcolor}{rgb}{0.101961,0.101961,0.101961}%
\pgfsetstrokecolor{textcolor}%
\pgfsetfillcolor{textcolor}%
\pgftext[x=0.526318in, y=1.438335in, left, base]{\color{textcolor}\sffamily\fontsize{7.000000}{8.400000}\selectfont StreamScout~\cite{streamscout2026}}%
\end{pgfscope}%
\begin{pgfscope}%
\definecolor{textcolor}{rgb}{0.101961,0.101961,0.101961}%
\pgfsetstrokecolor{textcolor}%
\pgfsetfillcolor{textcolor}%
\pgftext[x=0.244336in, y=1.353553in, left, base]{\color{textcolor}\sffamily\fontsize{7.000000}{8.400000}\selectfont Confidence cascade~\cite{chen2024frugalgpt}}%
\end{pgfscope}%
\begin{pgfscope}%
\definecolor{textcolor}{rgb}{0.101961,0.101961,0.101961}%
\pgfsetstrokecolor{textcolor}%
\pgfsetfillcolor{textcolor}%
\pgftext[x=0.483951in, y=1.268770in, left, base]{\color{textcolor}\sffamily\fontsize{7.000000}{8.400000}\selectfont SimpleStream~\cite{simplestream2026}}%
\end{pgfscope}%
\begin{pgfscope}%
\definecolor{textcolor}{rgb}{0.101961,0.101961,0.101961}%
\pgfsetstrokecolor{textcolor}%
\pgfsetfillcolor{textcolor}%
\pgftext[x=0.907953in, y=1.183987in, left, base]{\color{textcolor}\sffamily\fontsize{7.000000}{8.400000}\selectfont LLF~\cite{mok1983llf}}%
\end{pgfscope}%
\begin{pgfscope}%
\definecolor{textcolor}{rgb}{0.101961,0.101961,0.101961}%
\pgfsetstrokecolor{textcolor}%
\pgfsetfillcolor{textcolor}%
\pgftext[x=0.638226in, y=1.099205in, left, base]{\color{textcolor}\sffamily\fontsize{7.000000}{8.400000}\selectfont ECO-LLM~\cite{ecollm2025}}%
\end{pgfscope}%
\begin{pgfscope}%
\definecolor{textcolor}{rgb}{0.101961,0.101961,0.101961}%
\pgfsetstrokecolor{textcolor}%
\pgfsetfillcolor{textcolor}%
\pgftext[x=0.640758in, y=1.014422in, left, base]{\color{textcolor}\sffamily\fontsize{7.000000}{8.400000}\selectfont Mooncake~\cite{mooncake2025}}%
\end{pgfscope}%
\begin{pgfscope}%
\definecolor{textcolor}{rgb}{0.101961,0.101961,0.101961}%
\pgfsetstrokecolor{textcolor}%
\pgfsetfillcolor{textcolor}%
\pgftext[x=0.333964in, y=0.929639in, left, base]{\color{textcolor}\sffamily\fontsize{7.000000}{8.400000}\selectfont EDF + feasibility~\cite{liu1973edf}}%
\end{pgfscope}%
\begin{pgfscope}%
\definecolor{textcolor}{rgb}{0.101961,0.101961,0.101961}%
\pgfsetstrokecolor{textcolor}%
\pgfsetfillcolor{textcolor}%
\pgftext[x=0.883816in, y=0.844857in, left, base]{\color{textcolor}\sffamily\fontsize{7.000000}{8.400000}\selectfont EDF~\cite{liu1973edf}}%
\end{pgfscope}%
\begin{pgfscope}%
\definecolor{textcolor}{rgb}{0.101961,0.101961,0.101961}%
\pgfsetstrokecolor{textcolor}%
\pgfsetfillcolor{textcolor}%
\pgftext[x=0.735787in, y=0.760074in, left, base]{\color{textcolor}\sffamily\fontsize{7.000000}{8.400000}\selectfont KVFlow~\cite{kvflow2025}}%
\end{pgfscope}%
\begin{pgfscope}%
\definecolor{textcolor}{rgb}{0.101961,0.101961,0.101961}%
\pgfsetstrokecolor{textcolor}%
\pgfsetfillcolor{textcolor}%
\pgftext[x=0.401875in, y=0.677992in, left, base]{\color{textcolor}\sffamily\fontsize{7.000000}{8.400000}\selectfont Always-cloud FCFS}%
\end{pgfscope}%
\begin{pgfscope}%
\definecolor{textcolor}{rgb}{0.101961,0.101961,0.101961}%
\pgfsetstrokecolor{textcolor}%
\pgfsetfillcolor{textcolor}%
\pgftext[x=0.612964in, y=0.593210in, left, base]{\color{textcolor}\sffamily\fontsize{7.000000}{8.400000}\selectfont Warm-all LRU}%
\end{pgfscope}%
\begin{pgfscope}%
\definecolor{textcolor}{rgb}{0.101961,0.101961,0.101961}%
\pgfsetstrokecolor{textcolor}%
\pgfsetfillcolor{textcolor}%
\pgftext[x=0.513490in, y=0.505726in, left, base]{\color{textcolor}\sffamily\fontsize{7.000000}{8.400000}\selectfont QueryStream~\cite{querystream2026}}%
\end{pgfscope}%
\begin{pgfscope}%
\definecolor{textcolor}{rgb}{0.101961,0.101961,0.101961}%
\pgfsetstrokecolor{textcolor}%
\pgfsetfillcolor{textcolor}%
\pgftext[x=0.447999in, y=0.420944in, left, base]{\color{textcolor}\sffamily\fontsize{7.000000}{8.400000}\selectfont AttentionStore~\cite{cachedattention2024}}%
\end{pgfscope}%
\begin{pgfscope}%
\definecolor{textcolor}{rgb}{0.101961,0.101961,0.101961}%
\pgfsetstrokecolor{textcolor}%
\pgfsetfillcolor{textcolor}%
\pgftext[x=0.672322in, y=0.336161in, left, base]{\color{textcolor}\sffamily\fontsize{7.000000}{8.400000}\selectfont LMCache~\cite{lmcache2025}}%
\end{pgfscope}%
\begin{pgfscope}%
\pgfpathrectangle{\pgfqpoint{1.282050in}{0.321750in}}{\pgfqpoint{1.948050in}{1.365000in}}%
\pgfusepath{clip}%
\pgfsetbuttcap%
\pgfsetmiterjoin%
\definecolor{currentfill}{rgb}{0.156863,0.356863,0.568627}%
\pgfsetfillcolor{currentfill}%
\pgfsetlinewidth{0.000000pt}%
\definecolor{currentstroke}{rgb}{0.000000,0.000000,0.000000}%
\pgfsetstrokecolor{currentstroke}%
\pgfsetstrokeopacity{0.000000}%
\pgfsetdash{}{0pt}%
\pgfpathmoveto{\pgfqpoint{1.498500in}{1.534141in}}%
\pgfpathlineto{\pgfqpoint{2.523030in}{1.534141in}}%
\pgfpathlineto{\pgfqpoint{2.523030in}{1.585011in}}%
\pgfpathlineto{\pgfqpoint{1.498500in}{1.585011in}}%
\pgfpathlineto{\pgfqpoint{1.498500in}{1.534141in}}%
\pgfpathclose%
\pgfusepath{fill}%
\end{pgfscope}%
\begin{pgfscope}%
\pgfpathrectangle{\pgfqpoint{1.282050in}{0.321750in}}{\pgfqpoint{1.948050in}{1.365000in}}%
\pgfusepath{clip}%
\pgfsetbuttcap%
\pgfsetmiterjoin%
\definecolor{currentfill}{rgb}{0.784314,0.784314,0.784314}%
\pgfsetfillcolor{currentfill}%
\pgfsetlinewidth{0.000000pt}%
\definecolor{currentstroke}{rgb}{0.000000,0.000000,0.000000}%
\pgfsetstrokecolor{currentstroke}%
\pgfsetstrokeopacity{0.000000}%
\pgfsetdash{}{0pt}%
\pgfpathmoveto{\pgfqpoint{1.498500in}{1.449359in}}%
\pgfpathlineto{\pgfqpoint{1.840010in}{1.449359in}}%
\pgfpathlineto{\pgfqpoint{1.840010in}{1.500228in}}%
\pgfpathlineto{\pgfqpoint{1.498500in}{1.500228in}}%
\pgfpathlineto{\pgfqpoint{1.498500in}{1.449359in}}%
\pgfpathclose%
\pgfusepath{fill}%
\end{pgfscope}%
\begin{pgfscope}%
\pgfpathrectangle{\pgfqpoint{1.282050in}{0.321750in}}{\pgfqpoint{1.948050in}{1.365000in}}%
\pgfusepath{clip}%
\pgfsetbuttcap%
\pgfsetmiterjoin%
\definecolor{currentfill}{rgb}{0.784314,0.784314,0.784314}%
\pgfsetfillcolor{currentfill}%
\pgfsetlinewidth{0.000000pt}%
\definecolor{currentstroke}{rgb}{0.000000,0.000000,0.000000}%
\pgfsetstrokecolor{currentstroke}%
\pgfsetstrokeopacity{0.000000}%
\pgfsetdash{}{0pt}%
\pgfpathmoveto{\pgfqpoint{1.498500in}{1.364576in}}%
\pgfpathlineto{\pgfqpoint{1.840010in}{1.364576in}}%
\pgfpathlineto{\pgfqpoint{1.840010in}{1.415446in}}%
\pgfpathlineto{\pgfqpoint{1.498500in}{1.415446in}}%
\pgfpathlineto{\pgfqpoint{1.498500in}{1.364576in}}%
\pgfpathclose%
\pgfusepath{fill}%
\end{pgfscope}%
\begin{pgfscope}%
\pgfpathrectangle{\pgfqpoint{1.282050in}{0.321750in}}{\pgfqpoint{1.948050in}{1.365000in}}%
\pgfusepath{clip}%
\pgfsetbuttcap%
\pgfsetmiterjoin%
\definecolor{currentfill}{rgb}{0.784314,0.784314,0.784314}%
\pgfsetfillcolor{currentfill}%
\pgfsetlinewidth{0.000000pt}%
\definecolor{currentstroke}{rgb}{0.000000,0.000000,0.000000}%
\pgfsetstrokecolor{currentstroke}%
\pgfsetstrokeopacity{0.000000}%
\pgfsetdash{}{0pt}%
\pgfpathmoveto{\pgfqpoint{1.498500in}{1.279793in}}%
\pgfpathlineto{\pgfqpoint{1.815960in}{1.279793in}}%
\pgfpathlineto{\pgfqpoint{1.815960in}{1.330663in}}%
\pgfpathlineto{\pgfqpoint{1.498500in}{1.330663in}}%
\pgfpathlineto{\pgfqpoint{1.498500in}{1.279793in}}%
\pgfpathclose%
\pgfusepath{fill}%
\end{pgfscope}%
\begin{pgfscope}%
\pgfpathrectangle{\pgfqpoint{1.282050in}{0.321750in}}{\pgfqpoint{1.948050in}{1.365000in}}%
\pgfusepath{clip}%
\pgfsetbuttcap%
\pgfsetmiterjoin%
\definecolor{currentfill}{rgb}{0.784314,0.784314,0.784314}%
\pgfsetfillcolor{currentfill}%
\pgfsetlinewidth{0.000000pt}%
\definecolor{currentstroke}{rgb}{0.000000,0.000000,0.000000}%
\pgfsetstrokecolor{currentstroke}%
\pgfsetstrokeopacity{0.000000}%
\pgfsetdash{}{0pt}%
\pgfpathmoveto{\pgfqpoint{1.498500in}{1.195011in}}%
\pgfpathlineto{\pgfqpoint{1.705330in}{1.195011in}}%
\pgfpathlineto{\pgfqpoint{1.705330in}{1.245880in}}%
\pgfpathlineto{\pgfqpoint{1.498500in}{1.245880in}}%
\pgfpathlineto{\pgfqpoint{1.498500in}{1.195011in}}%
\pgfpathclose%
\pgfusepath{fill}%
\end{pgfscope}%
\begin{pgfscope}%
\pgfpathrectangle{\pgfqpoint{1.282050in}{0.321750in}}{\pgfqpoint{1.948050in}{1.365000in}}%
\pgfusepath{clip}%
\pgfsetbuttcap%
\pgfsetmiterjoin%
\definecolor{currentfill}{rgb}{0.784314,0.784314,0.784314}%
\pgfsetfillcolor{currentfill}%
\pgfsetlinewidth{0.000000pt}%
\definecolor{currentstroke}{rgb}{0.000000,0.000000,0.000000}%
\pgfsetstrokecolor{currentstroke}%
\pgfsetstrokeopacity{0.000000}%
\pgfsetdash{}{0pt}%
\pgfpathmoveto{\pgfqpoint{1.498500in}{1.110228in}}%
\pgfpathlineto{\pgfqpoint{1.695710in}{1.110228in}}%
\pgfpathlineto{\pgfqpoint{1.695710in}{1.161098in}}%
\pgfpathlineto{\pgfqpoint{1.498500in}{1.161098in}}%
\pgfpathlineto{\pgfqpoint{1.498500in}{1.110228in}}%
\pgfpathclose%
\pgfusepath{fill}%
\end{pgfscope}%
\begin{pgfscope}%
\pgfpathrectangle{\pgfqpoint{1.282050in}{0.321750in}}{\pgfqpoint{1.948050in}{1.365000in}}%
\pgfusepath{clip}%
\pgfsetbuttcap%
\pgfsetmiterjoin%
\definecolor{currentfill}{rgb}{0.478431,0.478431,0.478431}%
\pgfsetfillcolor{currentfill}%
\pgfsetlinewidth{0.000000pt}%
\definecolor{currentstroke}{rgb}{0.000000,0.000000,0.000000}%
\pgfsetstrokecolor{currentstroke}%
\pgfsetstrokeopacity{0.000000}%
\pgfsetdash{}{0pt}%
\pgfpathmoveto{\pgfqpoint{1.498500in}{1.025446in}}%
\pgfpathlineto{\pgfqpoint{1.686090in}{1.025446in}}%
\pgfpathlineto{\pgfqpoint{1.686090in}{1.076315in}}%
\pgfpathlineto{\pgfqpoint{1.498500in}{1.076315in}}%
\pgfpathlineto{\pgfqpoint{1.498500in}{1.025446in}}%
\pgfpathclose%
\pgfusepath{fill}%
\end{pgfscope}%
\begin{pgfscope}%
\pgfpathrectangle{\pgfqpoint{1.282050in}{0.321750in}}{\pgfqpoint{1.948050in}{1.365000in}}%
\pgfusepath{clip}%
\pgfsetbuttcap%
\pgfsetmiterjoin%
\definecolor{currentfill}{rgb}{0.784314,0.784314,0.784314}%
\pgfsetfillcolor{currentfill}%
\pgfsetlinewidth{0.000000pt}%
\definecolor{currentstroke}{rgb}{0.000000,0.000000,0.000000}%
\pgfsetstrokecolor{currentstroke}%
\pgfsetstrokeopacity{0.000000}%
\pgfsetdash{}{0pt}%
\pgfpathmoveto{\pgfqpoint{1.498500in}{0.940663in}}%
\pgfpathlineto{\pgfqpoint{1.666850in}{0.940663in}}%
\pgfpathlineto{\pgfqpoint{1.666850in}{0.991533in}}%
\pgfpathlineto{\pgfqpoint{1.498500in}{0.991533in}}%
\pgfpathlineto{\pgfqpoint{1.498500in}{0.940663in}}%
\pgfpathclose%
\pgfusepath{fill}%
\end{pgfscope}%
\begin{pgfscope}%
\pgfpathrectangle{\pgfqpoint{1.282050in}{0.321750in}}{\pgfqpoint{1.948050in}{1.365000in}}%
\pgfusepath{clip}%
\pgfsetbuttcap%
\pgfsetmiterjoin%
\definecolor{currentfill}{rgb}{0.784314,0.784314,0.784314}%
\pgfsetfillcolor{currentfill}%
\pgfsetlinewidth{0.000000pt}%
\definecolor{currentstroke}{rgb}{0.000000,0.000000,0.000000}%
\pgfsetstrokecolor{currentstroke}%
\pgfsetstrokeopacity{0.000000}%
\pgfsetdash{}{0pt}%
\pgfpathmoveto{\pgfqpoint{1.498500in}{0.855880in}}%
\pgfpathlineto{\pgfqpoint{1.551410in}{0.855880in}}%
\pgfpathlineto{\pgfqpoint{1.551410in}{0.906750in}}%
\pgfpathlineto{\pgfqpoint{1.498500in}{0.906750in}}%
\pgfpathlineto{\pgfqpoint{1.498500in}{0.855880in}}%
\pgfpathclose%
\pgfusepath{fill}%
\end{pgfscope}%
\begin{pgfscope}%
\pgfpathrectangle{\pgfqpoint{1.282050in}{0.321750in}}{\pgfqpoint{1.948050in}{1.365000in}}%
\pgfusepath{clip}%
\pgfsetbuttcap%
\pgfsetmiterjoin%
\definecolor{currentfill}{rgb}{0.478431,0.478431,0.478431}%
\pgfsetfillcolor{currentfill}%
\pgfsetlinewidth{0.000000pt}%
\definecolor{currentstroke}{rgb}{0.000000,0.000000,0.000000}%
\pgfsetstrokecolor{currentstroke}%
\pgfsetstrokeopacity{0.000000}%
\pgfsetdash{}{0pt}%
\pgfpathmoveto{\pgfqpoint{1.498500in}{0.771098in}}%
\pgfpathlineto{\pgfqpoint{1.512930in}{0.771098in}}%
\pgfpathlineto{\pgfqpoint{1.512930in}{0.821967in}}%
\pgfpathlineto{\pgfqpoint{1.498500in}{0.821967in}}%
\pgfpathlineto{\pgfqpoint{1.498500in}{0.771098in}}%
\pgfpathclose%
\pgfusepath{fill}%
\end{pgfscope}%
\begin{pgfscope}%
\pgfpathrectangle{\pgfqpoint{1.282050in}{0.321750in}}{\pgfqpoint{1.948050in}{1.365000in}}%
\pgfusepath{clip}%
\pgfsetbuttcap%
\pgfsetmiterjoin%
\definecolor{currentfill}{rgb}{0.784314,0.784314,0.784314}%
\pgfsetfillcolor{currentfill}%
\pgfsetlinewidth{0.000000pt}%
\definecolor{currentstroke}{rgb}{0.000000,0.000000,0.000000}%
\pgfsetstrokecolor{currentstroke}%
\pgfsetstrokeopacity{0.000000}%
\pgfsetdash{}{0pt}%
\pgfpathmoveto{\pgfqpoint{1.498500in}{0.686315in}}%
\pgfpathlineto{\pgfqpoint{1.508120in}{0.686315in}}%
\pgfpathlineto{\pgfqpoint{1.508120in}{0.737185in}}%
\pgfpathlineto{\pgfqpoint{1.498500in}{0.737185in}}%
\pgfpathlineto{\pgfqpoint{1.498500in}{0.686315in}}%
\pgfpathclose%
\pgfusepath{fill}%
\end{pgfscope}%
\begin{pgfscope}%
\pgfpathrectangle{\pgfqpoint{1.282050in}{0.321750in}}{\pgfqpoint{1.948050in}{1.365000in}}%
\pgfusepath{clip}%
\pgfsetbuttcap%
\pgfsetmiterjoin%
\definecolor{currentfill}{rgb}{0.478431,0.478431,0.478431}%
\pgfsetfillcolor{currentfill}%
\pgfsetlinewidth{0.000000pt}%
\definecolor{currentstroke}{rgb}{0.000000,0.000000,0.000000}%
\pgfsetstrokecolor{currentstroke}%
\pgfsetstrokeopacity{0.000000}%
\pgfsetdash{}{0pt}%
\pgfpathmoveto{\pgfqpoint{1.498500in}{0.601533in}}%
\pgfpathlineto{\pgfqpoint{1.493690in}{0.601533in}}%
\pgfpathlineto{\pgfqpoint{1.493690in}{0.652402in}}%
\pgfpathlineto{\pgfqpoint{1.498500in}{0.652402in}}%
\pgfpathlineto{\pgfqpoint{1.498500in}{0.601533in}}%
\pgfpathclose%
\pgfusepath{fill}%
\end{pgfscope}%
\begin{pgfscope}%
\pgfpathrectangle{\pgfqpoint{1.282050in}{0.321750in}}{\pgfqpoint{1.948050in}{1.365000in}}%
\pgfusepath{clip}%
\pgfsetbuttcap%
\pgfsetmiterjoin%
\definecolor{currentfill}{rgb}{0.784314,0.784314,0.784314}%
\pgfsetfillcolor{currentfill}%
\pgfsetlinewidth{0.000000pt}%
\definecolor{currentstroke}{rgb}{0.000000,0.000000,0.000000}%
\pgfsetstrokecolor{currentstroke}%
\pgfsetstrokeopacity{0.000000}%
\pgfsetdash{}{0pt}%
\pgfpathmoveto{\pgfqpoint{1.498500in}{0.516750in}}%
\pgfpathlineto{\pgfqpoint{1.479260in}{0.516750in}}%
\pgfpathlineto{\pgfqpoint{1.479260in}{0.567620in}}%
\pgfpathlineto{\pgfqpoint{1.498500in}{0.567620in}}%
\pgfpathlineto{\pgfqpoint{1.498500in}{0.516750in}}%
\pgfpathclose%
\pgfusepath{fill}%
\end{pgfscope}%
\begin{pgfscope}%
\pgfpathrectangle{\pgfqpoint{1.282050in}{0.321750in}}{\pgfqpoint{1.948050in}{1.365000in}}%
\pgfusepath{clip}%
\pgfsetbuttcap%
\pgfsetmiterjoin%
\definecolor{currentfill}{rgb}{0.478431,0.478431,0.478431}%
\pgfsetfillcolor{currentfill}%
\pgfsetlinewidth{0.000000pt}%
\definecolor{currentstroke}{rgb}{0.000000,0.000000,0.000000}%
\pgfsetstrokecolor{currentstroke}%
\pgfsetstrokeopacity{0.000000}%
\pgfsetdash{}{0pt}%
\pgfpathmoveto{\pgfqpoint{1.498500in}{0.431967in}}%
\pgfpathlineto{\pgfqpoint{1.479260in}{0.431967in}}%
\pgfpathlineto{\pgfqpoint{1.479260in}{0.482837in}}%
\pgfpathlineto{\pgfqpoint{1.498500in}{0.482837in}}%
\pgfpathlineto{\pgfqpoint{1.498500in}{0.431967in}}%
\pgfpathclose%
\pgfusepath{fill}%
\end{pgfscope}%
\begin{pgfscope}%
\pgfpathrectangle{\pgfqpoint{1.282050in}{0.321750in}}{\pgfqpoint{1.948050in}{1.365000in}}%
\pgfusepath{clip}%
\pgfsetbuttcap%
\pgfsetmiterjoin%
\definecolor{currentfill}{rgb}{0.478431,0.478431,0.478431}%
\pgfsetfillcolor{currentfill}%
\pgfsetlinewidth{0.000000pt}%
\definecolor{currentstroke}{rgb}{0.000000,0.000000,0.000000}%
\pgfsetstrokecolor{currentstroke}%
\pgfsetstrokeopacity{0.000000}%
\pgfsetdash{}{0pt}%
\pgfpathmoveto{\pgfqpoint{1.498500in}{0.347185in}}%
\pgfpathlineto{\pgfqpoint{1.469640in}{0.347185in}}%
\pgfpathlineto{\pgfqpoint{1.469640in}{0.398054in}}%
\pgfpathlineto{\pgfqpoint{1.498500in}{0.398054in}}%
\pgfpathlineto{\pgfqpoint{1.498500in}{0.347185in}}%
\pgfpathclose%
\pgfusepath{fill}%
\end{pgfscope}%
\begin{pgfscope}%
\pgfpathrectangle{\pgfqpoint{1.282050in}{0.321750in}}{\pgfqpoint{1.948050in}{1.365000in}}%
\pgfusepath{clip}%
\pgfsetrectcap%
\pgfsetroundjoin%
\pgfsetlinewidth{0.602250pt}%
\definecolor{currentstroke}{rgb}{0.101961,0.101961,0.101961}%
\pgfsetstrokecolor{currentstroke}%
\pgfsetdash{}{0pt}%
\pgfpathmoveto{\pgfqpoint{1.498500in}{0.321750in}}%
\pgfpathlineto{\pgfqpoint{1.498500in}{1.686750in}}%
\pgfusepath{stroke}%
\end{pgfscope}%
\begin{pgfscope}%
\pgfsetrectcap%
\pgfsetmiterjoin%
\pgfsetlinewidth{0.602250pt}%
\definecolor{currentstroke}{rgb}{0.101961,0.101961,0.101961}%
\pgfsetstrokecolor{currentstroke}%
\pgfsetdash{}{0pt}%
\pgfpathmoveto{\pgfqpoint{1.282050in}{0.321750in}}%
\pgfpathlineto{\pgfqpoint{3.230100in}{0.321750in}}%
\pgfusepath{stroke}%
\end{pgfscope}%
\begin{pgfscope}%
\definecolor{textcolor}{rgb}{0.101961,0.101961,0.101961}%
\pgfsetstrokecolor{textcolor}%
\pgfsetfillcolor{textcolor}%
\pgftext[x=2.561510in,y=1.559576in,left,]{\color{textcolor}\sffamily\fontsize{7.000000}{8.400000}\bfseries\selectfont \textbf{+213}}%
\end{pgfscope}%
\begin{pgfscope}%
\definecolor{textcolor}{rgb}{0.101961,0.101961,0.101961}%
\pgfsetstrokecolor{textcolor}%
\pgfsetfillcolor{textcolor}%
\pgftext[x=1.878490in,y=1.474793in,left,]{\color{textcolor}\sffamily\fontsize{7.000000}{8.400000}\selectfont +71}%
\end{pgfscope}%
\begin{pgfscope}%
\definecolor{textcolor}{rgb}{0.101961,0.101961,0.101961}%
\pgfsetstrokecolor{textcolor}%
\pgfsetfillcolor{textcolor}%
\pgftext[x=1.878490in,y=1.390011in,left,]{\color{textcolor}\sffamily\fontsize{7.000000}{8.400000}\selectfont +71}%
\end{pgfscope}%
\begin{pgfscope}%
\definecolor{textcolor}{rgb}{0.101961,0.101961,0.101961}%
\pgfsetstrokecolor{textcolor}%
\pgfsetfillcolor{textcolor}%
\pgftext[x=1.854440in,y=1.305228in,left,]{\color{textcolor}\sffamily\fontsize{7.000000}{8.400000}\selectfont +66}%
\end{pgfscope}%
\begin{pgfscope}%
\definecolor{textcolor}{rgb}{0.101961,0.101961,0.101961}%
\pgfsetstrokecolor{textcolor}%
\pgfsetfillcolor{textcolor}%
\pgftext[x=1.743810in,y=1.220446in,left,]{\color{textcolor}\sffamily\fontsize{7.000000}{8.400000}\selectfont +43}%
\end{pgfscope}%
\begin{pgfscope}%
\definecolor{textcolor}{rgb}{0.101961,0.101961,0.101961}%
\pgfsetstrokecolor{textcolor}%
\pgfsetfillcolor{textcolor}%
\pgftext[x=1.734190in,y=1.135663in,left,]{\color{textcolor}\sffamily\fontsize{7.000000}{8.400000}\selectfont +41}%
\end{pgfscope}%
\begin{pgfscope}%
\definecolor{textcolor}{rgb}{0.101961,0.101961,0.101961}%
\pgfsetstrokecolor{textcolor}%
\pgfsetfillcolor{textcolor}%
\pgftext[x=1.724570in,y=1.050880in,left,]{\color{textcolor}\sffamily\fontsize{7.000000}{8.400000}\selectfont +39}%
\end{pgfscope}%
\begin{pgfscope}%
\definecolor{textcolor}{rgb}{0.101961,0.101961,0.101961}%
\pgfsetstrokecolor{textcolor}%
\pgfsetfillcolor{textcolor}%
\pgftext[x=1.705330in,y=0.966098in,left,]{\color{textcolor}\sffamily\fontsize{7.000000}{8.400000}\selectfont +35}%
\end{pgfscope}%
\begin{pgfscope}%
\definecolor{textcolor}{rgb}{0.101961,0.101961,0.101961}%
\pgfsetstrokecolor{textcolor}%
\pgfsetfillcolor{textcolor}%
\pgftext[x=1.589890in,y=0.881315in,left,]{\color{textcolor}\sffamily\fontsize{7.000000}{8.400000}\selectfont +11}%
\end{pgfscope}%
\begin{pgfscope}%
\definecolor{textcolor}{rgb}{0.101961,0.101961,0.101961}%
\pgfsetstrokecolor{textcolor}%
\pgfsetfillcolor{textcolor}%
\pgftext[x=1.551410in,y=0.796533in,left,]{\color{textcolor}\sffamily\fontsize{7.000000}{8.400000}\selectfont +3}%
\end{pgfscope}%
\begin{pgfscope}%
\definecolor{textcolor}{rgb}{0.101961,0.101961,0.101961}%
\pgfsetstrokecolor{textcolor}%
\pgfsetfillcolor{textcolor}%
\pgftext[x=1.546600in,y=0.711750in,left,]{\color{textcolor}\sffamily\fontsize{7.000000}{8.400000}\selectfont +2}%
\end{pgfscope}%
\begin{pgfscope}%
\definecolor{textcolor}{rgb}{0.101961,0.101961,0.101961}%
\pgfsetstrokecolor{textcolor}%
\pgfsetfillcolor{textcolor}%
\pgftext[x=1.532170in,y=0.626967in,left,]{\color{textcolor}\sffamily\fontsize{7.000000}{8.400000}\selectfont \(\displaystyle -\)1}%
\end{pgfscope}%
\begin{pgfscope}%
\definecolor{textcolor}{rgb}{0.101961,0.101961,0.101961}%
\pgfsetstrokecolor{textcolor}%
\pgfsetfillcolor{textcolor}%
\pgftext[x=1.517740in,y=0.542185in,left,]{\color{textcolor}\sffamily\fontsize{7.000000}{8.400000}\selectfont \(\displaystyle -\)4}%
\end{pgfscope}%
\begin{pgfscope}%
\definecolor{textcolor}{rgb}{0.101961,0.101961,0.101961}%
\pgfsetstrokecolor{textcolor}%
\pgfsetfillcolor{textcolor}%
\pgftext[x=1.517740in,y=0.457402in,left,]{\color{textcolor}\sffamily\fontsize{7.000000}{8.400000}\selectfont \(\displaystyle -\)4}%
\end{pgfscope}%
\begin{pgfscope}%
\definecolor{textcolor}{rgb}{0.101961,0.101961,0.101961}%
\pgfsetstrokecolor{textcolor}%
\pgfsetfillcolor{textcolor}%
\pgftext[x=1.508120in,y=0.372620in,left,]{\color{textcolor}\sffamily\fontsize{7.000000}{8.400000}\selectfont \(\displaystyle -\)6}%
\end{pgfscope}%
\begin{pgfscope}%
\pgfsetroundcap%
\pgfsetroundjoin%
\pgfsetlinewidth{0.602250pt}%
\definecolor{currentstroke}{rgb}{0.101961,0.101961,0.101961}%
\pgfsetstrokecolor{currentstroke}%
\pgfsetdash{}{0pt}%
\pgfpathmoveto{\pgfqpoint{1.849327in}{1.620620in}}%
\pgfpathquadraticcurveto{\pgfqpoint{2.181520in}{1.620620in}}{\pgfqpoint{2.513713in}{1.620620in}}%
\pgfusepath{stroke}%
\end{pgfscope}%
\begin{pgfscope}%
\pgfsetroundcap%
\pgfsetroundjoin%
\pgfsetlinewidth{0.602250pt}%
\definecolor{currentstroke}{rgb}{0.101961,0.101961,0.101961}%
\pgfsetstrokecolor{currentstroke}%
\pgfsetdash{}{0pt}%
\pgfpathmoveto{\pgfqpoint{1.888216in}{1.601175in}}%
\pgfpathlineto{\pgfqpoint{1.849327in}{1.620620in}}%
\pgfpathlineto{\pgfqpoint{1.888216in}{1.640064in}}%
\pgfusepath{stroke}%
\end{pgfscope}%
\begin{pgfscope}%
\pgfsetroundcap%
\pgfsetroundjoin%
\pgfsetlinewidth{0.602250pt}%
\definecolor{currentstroke}{rgb}{0.101961,0.101961,0.101961}%
\pgfsetstrokecolor{currentstroke}%
\pgfsetdash{}{0pt}%
\pgfpathmoveto{\pgfqpoint{2.474824in}{1.640064in}}%
\pgfpathlineto{\pgfqpoint{2.513713in}{1.620620in}}%
\pgfpathlineto{\pgfqpoint{2.474824in}{1.601175in}}%
\pgfusepath{stroke}%
\end{pgfscope}%
\begin{pgfscope}%
\definecolor{textcolor}{rgb}{0.101961,0.101961,0.101961}%
\pgfsetstrokecolor{textcolor}%
\pgfsetfillcolor{textcolor}%
\pgftext[x=2.181520in,y=1.630793in,,bottom]{\color{textcolor}\sffamily\fontsize{7.000000}{8.400000}\bfseries\selectfont 3.0\(\displaystyle \times\)}%
\end{pgfscope}%
\begin{pgfscope}%
\definecolor{textcolor}{rgb}{0.101961,0.101961,0.101961}%
\pgfsetstrokecolor{textcolor}%
\pgfsetfillcolor{textcolor}%
\pgftext[x=1.665000in,y=0.023400in,,bottom]{\color{textcolor}\sffamily\fontsize{7.500000}{9.000000}\selectfont Additional correct answers over the on-device model}%
\end{pgfscope}%
\begin{pgfscope}%
\pgfsetbuttcap%
\pgfsetmiterjoin%
\definecolor{currentfill}{rgb}{0.156863,0.356863,0.568627}%
\pgfsetfillcolor{currentfill}%
\pgfsetlinewidth{0.000000pt}%
\definecolor{currentstroke}{rgb}{0.000000,0.000000,0.000000}%
\pgfsetstrokecolor{currentstroke}%
\pgfsetstrokeopacity{0.000000}%
\pgfsetdash{}{0pt}%
\pgfpathmoveto{\pgfqpoint{1.122327in}{1.795806in}}%
\pgfpathlineto{\pgfqpoint{1.212605in}{1.795806in}}%
\pgfpathlineto{\pgfqpoint{1.212605in}{1.859000in}}%
\pgfpathlineto{\pgfqpoint{1.122327in}{1.859000in}}%
\pgfpathlineto{\pgfqpoint{1.122327in}{1.795806in}}%
\pgfpathclose%
\pgfusepath{fill}%
\end{pgfscope}%
\begin{pgfscope}%
\definecolor{textcolor}{rgb}{0.101961,0.101961,0.101961}%
\pgfsetstrokecolor{textcolor}%
\pgfsetfillcolor{textcolor}%
\pgftext[x=1.248716in,y=1.795806in,left,base]{\color{textcolor}\sffamily\fontsize{6.500000}{7.800000}\selectfont COVEX (ours)}%
\end{pgfscope}%
\begin{pgfscope}%
\pgfsetbuttcap%
\pgfsetmiterjoin%
\definecolor{currentfill}{rgb}{0.478431,0.478431,0.478431}%
\pgfsetfillcolor{currentfill}%
\pgfsetlinewidth{0.000000pt}%
\definecolor{currentstroke}{rgb}{0.000000,0.000000,0.000000}%
\pgfsetstrokecolor{currentstroke}%
\pgfsetstrokeopacity{0.000000}%
\pgfsetdash{}{0pt}%
\pgfpathmoveto{\pgfqpoint{1.900257in}{1.795806in}}%
\pgfpathlineto{\pgfqpoint{1.990535in}{1.795806in}}%
\pgfpathlineto{\pgfqpoint{1.990535in}{1.859000in}}%
\pgfpathlineto{\pgfqpoint{1.900257in}{1.859000in}}%
\pgfpathlineto{\pgfqpoint{1.900257in}{1.795806in}}%
\pgfpathclose%
\pgfusepath{fill}%
\end{pgfscope}%
\begin{pgfscope}%
\definecolor{textcolor}{rgb}{0.101961,0.101961,0.101961}%
\pgfsetstrokecolor{textcolor}%
\pgfsetfillcolor{textcolor}%
\pgftext[x=2.026646in,y=1.795806in,left,base]{\color{textcolor}\sffamily\fontsize{6.500000}{7.800000}\selectfont Stateful}%
\end{pgfscope}%
\begin{pgfscope}%
\pgfsetbuttcap%
\pgfsetmiterjoin%
\definecolor{currentfill}{rgb}{0.784314,0.784314,0.784314}%
\pgfsetfillcolor{currentfill}%
\pgfsetlinewidth{0.000000pt}%
\definecolor{currentstroke}{rgb}{0.000000,0.000000,0.000000}%
\pgfsetstrokecolor{currentstroke}%
\pgfsetstrokeopacity{0.000000}%
\pgfsetdash{}{0pt}%
\pgfpathmoveto{\pgfqpoint{2.423900in}{1.795806in}}%
\pgfpathlineto{\pgfqpoint{2.514178in}{1.795806in}}%
\pgfpathlineto{\pgfqpoint{2.514178in}{1.859000in}}%
\pgfpathlineto{\pgfqpoint{2.423900in}{1.859000in}}%
\pgfpathlineto{\pgfqpoint{2.423900in}{1.795806in}}%
\pgfpathclose%
\pgfusepath{fill}%
\end{pgfscope}%
\begin{pgfscope}%
\definecolor{textcolor}{rgb}{0.101961,0.101961,0.101961}%
\pgfsetstrokecolor{textcolor}%
\pgfsetfillcolor{textcolor}%
\pgftext[x=2.550289in,y=1.795806in,left,base]{\color{textcolor}\sffamily\fontsize{6.500000}{7.800000}\selectfont Stateless}%
\end{pgfscope}%
\end{pgfpicture}%
\makeatother%
\endgroup%
}
\caption{[PLACEHOLDER: OLD DATA, rung 3 stand-in for COVEX, seeds 4-7] Correct, on-time answers gained over the on-device model. COVEX extracts 3$\times$ more of the cloud's benefit than any baseline.}
\label{fig:F01_intro_teaser}
\end{figure}
\clearpage
\begin{figure*}[t]
\centering
{\fontfamily{LinuxBiolinumT}\selectfont
%% Creator: Matplotlib, PGF backend
%%
%% To include the figure in your LaTeX document, write
%%   \input{<filename>.pgf}
%%
%% Make sure the required packages are loaded in your preamble
%%   \usepackage{pgf}
%%
%% Also ensure that all the required font packages are loaded; for instance,
%% the lmodern package is sometimes necessary when using math font.
%%   \usepackage{lmodern}
%%
%% Figures using additional raster images can only be included by \input if
%% they are in the same directory as the main LaTeX file. For loading figures
%% from other directories you can use the `import` package
%%   \usepackage{import}
%%
%% and then include the figures with
%%   \import{<path to file>}{<filename>.pgf}
%%
%% Matplotlib used the following preamble
%%
\begingroup%
\makeatletter%
\begin{pgfpicture}%
\pgfpathrectangle{\pgfpointorigin}{\pgfqpoint{7.000000in}{4.600000in}}%
\pgfusepath{use as bounding box, clip}%
\begin{pgfscope}%
\pgfsetbuttcap%
\pgfsetmiterjoin%
\definecolor{currentfill}{rgb}{1.000000,1.000000,1.000000}%
\pgfsetfillcolor{currentfill}%
\pgfsetlinewidth{0.000000pt}%
\definecolor{currentstroke}{rgb}{1.000000,1.000000,1.000000}%
\pgfsetstrokecolor{currentstroke}%
\pgfsetdash{}{0pt}%
\pgfpathmoveto{\pgfqpoint{0.000000in}{0.000000in}}%
\pgfpathlineto{\pgfqpoint{7.000000in}{0.000000in}}%
\pgfpathlineto{\pgfqpoint{7.000000in}{4.600000in}}%
\pgfpathlineto{\pgfqpoint{0.000000in}{4.600000in}}%
\pgfpathlineto{\pgfqpoint{0.000000in}{0.000000in}}%
\pgfpathclose%
\pgfusepath{fill}%
\end{pgfscope}%
\begin{pgfscope}%
\pgfsetbuttcap%
\pgfsetmiterjoin%
\definecolor{currentfill}{rgb}{1.000000,1.000000,1.000000}%
\pgfsetfillcolor{currentfill}%
\pgfsetlinewidth{0.000000pt}%
\definecolor{currentstroke}{rgb}{0.000000,0.000000,0.000000}%
\pgfsetstrokecolor{currentstroke}%
\pgfsetstrokeopacity{0.000000}%
\pgfsetdash{}{0pt}%
\pgfpathmoveto{\pgfqpoint{2.695000in}{0.598000in}}%
\pgfpathlineto{\pgfqpoint{3.850000in}{0.598000in}}%
\pgfpathlineto{\pgfqpoint{3.850000in}{4.462000in}}%
\pgfpathlineto{\pgfqpoint{2.695000in}{4.462000in}}%
\pgfpathlineto{\pgfqpoint{2.695000in}{0.598000in}}%
\pgfpathclose%
\pgfusepath{fill}%
\end{pgfscope}%
\begin{pgfscope}%
\pgfpathrectangle{\pgfqpoint{2.695000in}{0.598000in}}{\pgfqpoint{1.155000in}{3.864000in}}%
\pgfusepath{clip}%
\pgfsetrectcap%
\pgfsetroundjoin%
\pgfsetlinewidth{0.501875pt}%
\definecolor{currentstroke}{rgb}{0.901961,0.901961,0.901961}%
\pgfsetstrokecolor{currentstroke}%
\pgfsetdash{}{0pt}%
\pgfpathmoveto{\pgfqpoint{2.695000in}{4.398656in}}%
\pgfpathlineto{\pgfqpoint{3.850000in}{4.398656in}}%
\pgfusepath{stroke}%
\end{pgfscope}%
\begin{pgfscope}%
\pgfpathrectangle{\pgfqpoint{2.695000in}{0.598000in}}{\pgfqpoint{1.155000in}{3.864000in}}%
\pgfusepath{clip}%
\pgfsetrectcap%
\pgfsetroundjoin%
\pgfsetlinewidth{0.501875pt}%
\definecolor{currentstroke}{rgb}{0.901961,0.901961,0.901961}%
\pgfsetstrokecolor{currentstroke}%
\pgfsetdash{}{0pt}%
\pgfpathmoveto{\pgfqpoint{2.695000in}{2.427066in}}%
\pgfpathlineto{\pgfqpoint{3.850000in}{2.427066in}}%
\pgfusepath{stroke}%
\end{pgfscope}%
\begin{pgfscope}%
\pgfpathrectangle{\pgfqpoint{2.695000in}{0.598000in}}{\pgfqpoint{1.155000in}{3.864000in}}%
\pgfusepath{clip}%
\pgfsetrectcap%
\pgfsetroundjoin%
\pgfsetlinewidth{0.501875pt}%
\definecolor{currentstroke}{rgb}{0.901961,0.901961,0.901961}%
\pgfsetstrokecolor{currentstroke}%
\pgfsetdash{}{0pt}%
\pgfpathmoveto{\pgfqpoint{2.695000in}{1.564000in}}%
\pgfpathlineto{\pgfqpoint{3.850000in}{1.564000in}}%
\pgfusepath{stroke}%
\end{pgfscope}%
\begin{pgfscope}%
\pgfpathrectangle{\pgfqpoint{2.695000in}{0.598000in}}{\pgfqpoint{1.155000in}{3.864000in}}%
\pgfusepath{clip}%
\pgfsetrectcap%
\pgfsetroundjoin%
\pgfsetlinewidth{0.401500pt}%
\definecolor{currentstroke}{rgb}{0.901961,0.901961,0.901961}%
\pgfsetstrokecolor{currentstroke}%
\pgfsetdash{}{0pt}%
\pgfpathmoveto{\pgfqpoint{2.926000in}{0.598000in}}%
\pgfpathlineto{\pgfqpoint{2.926000in}{4.462000in}}%
\pgfusepath{stroke}%
\end{pgfscope}%
\begin{pgfscope}%
\definecolor{textcolor}{rgb}{0.101961,0.101961,0.101961}%
\pgfsetstrokecolor{textcolor}%
\pgfsetfillcolor{textcolor}%
\pgftext[x=2.926000in,y=0.549389in,,top]{\color{textcolor}\sffamily\fontsize{7.000000}{8.400000}\selectfont 0.4}%
\end{pgfscope}%
\begin{pgfscope}%
\pgfpathrectangle{\pgfqpoint{2.695000in}{0.598000in}}{\pgfqpoint{1.155000in}{3.864000in}}%
\pgfusepath{clip}%
\pgfsetrectcap%
\pgfsetroundjoin%
\pgfsetlinewidth{0.401500pt}%
\definecolor{currentstroke}{rgb}{0.901961,0.901961,0.901961}%
\pgfsetstrokecolor{currentstroke}%
\pgfsetdash{}{0pt}%
\pgfpathmoveto{\pgfqpoint{3.388000in}{0.598000in}}%
\pgfpathlineto{\pgfqpoint{3.388000in}{4.462000in}}%
\pgfusepath{stroke}%
\end{pgfscope}%
\begin{pgfscope}%
\definecolor{textcolor}{rgb}{0.101961,0.101961,0.101961}%
\pgfsetstrokecolor{textcolor}%
\pgfsetfillcolor{textcolor}%
\pgftext[x=3.388000in,y=0.549389in,,top]{\color{textcolor}\sffamily\fontsize{7.000000}{8.400000}\selectfont 0.6}%
\end{pgfscope}%
\begin{pgfscope}%
\definecolor{textcolor}{rgb}{0.000000,0.000000,0.000000}%
\pgfsetstrokecolor{textcolor}%
\pgfsetfillcolor{textcolor}%
\pgftext[x=2.976272in, y=0.367676in, left, base]{\color{textcolor}\sffamily\fontsize{7.500000}{9.000000}\selectfont Robot correct}%
\end{pgfscope}%
\begin{pgfscope}%
\definecolor{textcolor}{rgb}{0.000000,0.000000,0.000000}%
\pgfsetstrokecolor{textcolor}%
\pgfsetfillcolor{textcolor}%
\pgftext[x=3.030876in, y=0.267754in, left, base]{\color{textcolor}\sffamily\fontsize{7.500000}{9.000000}\selectfont by deadline}%
\end{pgfscope}%
\begin{pgfscope}%
\pgfpathrectangle{\pgfqpoint{2.695000in}{0.598000in}}{\pgfqpoint{1.155000in}{3.864000in}}%
\pgfusepath{clip}%
\pgfsetrectcap%
\pgfsetroundjoin%
\pgfsetlinewidth{0.602250pt}%
\definecolor{currentstroke}{rgb}{0.156863,0.356863,0.568627}%
\pgfsetstrokecolor{currentstroke}%
\pgfsetdash{}{0pt}%
\pgfpathmoveto{\pgfqpoint{3.454500in}{3.369311in}}%
\pgfpathlineto{\pgfqpoint{3.489500in}{3.210951in}}%
\pgfpathlineto{\pgfqpoint{3.384500in}{3.052590in}}%
\pgfpathlineto{\pgfqpoint{3.454500in}{2.894230in}}%
\pgfpathlineto{\pgfqpoint{3.419500in}{2.735869in}}%
\pgfpathlineto{\pgfqpoint{3.472000in}{2.577508in}}%
\pgfusepath{stroke}%
\end{pgfscope}%
\begin{pgfscope}%
\pgfsetrectcap%
\pgfsetmiterjoin%
\pgfsetlinewidth{0.803000pt}%
\definecolor{currentstroke}{rgb}{0.101961,0.101961,0.101961}%
\pgfsetstrokecolor{currentstroke}%
\pgfsetdash{}{0pt}%
\pgfpathmoveto{\pgfqpoint{2.695000in}{0.598000in}}%
\pgfpathlineto{\pgfqpoint{3.850000in}{0.598000in}}%
\pgfusepath{stroke}%
\end{pgfscope}%
\begin{pgfscope}%
\pgfpathrectangle{\pgfqpoint{2.695000in}{0.598000in}}{\pgfqpoint{1.155000in}{3.864000in}}%
\pgfusepath{clip}%
\pgfsetbuttcap%
\pgfsetroundjoin%
\definecolor{currentfill}{rgb}{0.435294,0.560784,0.709804}%
\pgfsetfillcolor{currentfill}%
\pgfsetlinewidth{0.000000pt}%
\definecolor{currentstroke}{rgb}{0.435294,0.560784,0.709804}%
\pgfsetstrokecolor{currentstroke}%
\pgfsetdash{}{0pt}%
\pgfsys@defobject{currentmarker}{\pgfqpoint{-0.024056in}{-0.024056in}}{\pgfqpoint{0.024056in}{0.024056in}}{%
\pgfpathmoveto{\pgfqpoint{0.000000in}{-0.024056in}}%
\pgfpathcurveto{\pgfqpoint{0.006380in}{-0.024056in}}{\pgfqpoint{0.012499in}{-0.021522in}}{\pgfqpoint{0.017010in}{-0.017010in}}%
\pgfpathcurveto{\pgfqpoint{0.021522in}{-0.012499in}}{\pgfqpoint{0.024056in}{-0.006380in}}{\pgfqpoint{0.024056in}{0.000000in}}%
\pgfpathcurveto{\pgfqpoint{0.024056in}{0.006380in}}{\pgfqpoint{0.021522in}{0.012499in}}{\pgfqpoint{0.017010in}{0.017010in}}%
\pgfpathcurveto{\pgfqpoint{0.012499in}{0.021522in}}{\pgfqpoint{0.006380in}{0.024056in}}{\pgfqpoint{0.000000in}{0.024056in}}%
\pgfpathcurveto{\pgfqpoint{-0.006380in}{0.024056in}}{\pgfqpoint{-0.012499in}{0.021522in}}{\pgfqpoint{-0.017010in}{0.017010in}}%
\pgfpathcurveto{\pgfqpoint{-0.021522in}{0.012499in}}{\pgfqpoint{-0.024056in}{0.006380in}}{\pgfqpoint{-0.024056in}{0.000000in}}%
\pgfpathcurveto{\pgfqpoint{-0.024056in}{-0.006380in}}{\pgfqpoint{-0.021522in}{-0.012499in}}{\pgfqpoint{-0.017010in}{-0.017010in}}%
\pgfpathcurveto{\pgfqpoint{-0.012499in}{-0.021522in}}{\pgfqpoint{-0.006380in}{-0.024056in}}{\pgfqpoint{0.000000in}{-0.024056in}}%
\pgfpathlineto{\pgfqpoint{0.000000in}{-0.024056in}}%
\pgfpathclose%
\pgfusepath{fill}%
}%
\begin{pgfscope}%
\pgfsys@transformshift{3.087000in}{4.161115in}%
\pgfsys@useobject{currentmarker}{}%
\end{pgfscope}%
\end{pgfscope}%
\begin{pgfscope}%
\definecolor{textcolor}{rgb}{0.549020,0.549020,0.549020}%
\pgfsetstrokecolor{textcolor}%
\pgfsetfillcolor{textcolor}%
\pgftext[x=3.127425in,y=4.161115in,left,]{\color{textcolor}\sffamily\fontsize{7.000000}{8.400000}\selectfont 0.47}%
\end{pgfscope}%
\begin{pgfscope}%
\pgfpathrectangle{\pgfqpoint{2.695000in}{0.598000in}}{\pgfqpoint{1.155000in}{3.864000in}}%
\pgfusepath{clip}%
\pgfsetbuttcap%
\pgfsetroundjoin%
\definecolor{currentfill}{rgb}{0.435294,0.560784,0.709804}%
\pgfsetfillcolor{currentfill}%
\pgfsetlinewidth{0.000000pt}%
\definecolor{currentstroke}{rgb}{0.435294,0.560784,0.709804}%
\pgfsetstrokecolor{currentstroke}%
\pgfsetdash{}{0pt}%
\pgfsys@defobject{currentmarker}{\pgfqpoint{-0.024056in}{-0.024056in}}{\pgfqpoint{0.024056in}{0.024056in}}{%
\pgfpathmoveto{\pgfqpoint{0.000000in}{-0.024056in}}%
\pgfpathcurveto{\pgfqpoint{0.006380in}{-0.024056in}}{\pgfqpoint{0.012499in}{-0.021522in}}{\pgfqpoint{0.017010in}{-0.017010in}}%
\pgfpathcurveto{\pgfqpoint{0.021522in}{-0.012499in}}{\pgfqpoint{0.024056in}{-0.006380in}}{\pgfqpoint{0.024056in}{0.000000in}}%
\pgfpathcurveto{\pgfqpoint{0.024056in}{0.006380in}}{\pgfqpoint{0.021522in}{0.012499in}}{\pgfqpoint{0.017010in}{0.017010in}}%
\pgfpathcurveto{\pgfqpoint{0.012499in}{0.021522in}}{\pgfqpoint{0.006380in}{0.024056in}}{\pgfqpoint{0.000000in}{0.024056in}}%
\pgfpathcurveto{\pgfqpoint{-0.006380in}{0.024056in}}{\pgfqpoint{-0.012499in}{0.021522in}}{\pgfqpoint{-0.017010in}{0.017010in}}%
\pgfpathcurveto{\pgfqpoint{-0.021522in}{0.012499in}}{\pgfqpoint{-0.024056in}{0.006380in}}{\pgfqpoint{-0.024056in}{0.000000in}}%
\pgfpathcurveto{\pgfqpoint{-0.024056in}{-0.006380in}}{\pgfqpoint{-0.021522in}{-0.012499in}}{\pgfqpoint{-0.017010in}{-0.017010in}}%
\pgfpathcurveto{\pgfqpoint{-0.012499in}{-0.021522in}}{\pgfqpoint{-0.006380in}{-0.024056in}}{\pgfqpoint{0.000000in}{-0.024056in}}%
\pgfpathlineto{\pgfqpoint{0.000000in}{-0.024056in}}%
\pgfpathclose%
\pgfusepath{fill}%
}%
\begin{pgfscope}%
\pgfsys@transformshift{3.122000in}{4.002754in}%
\pgfsys@useobject{currentmarker}{}%
\end{pgfscope}%
\end{pgfscope}%
\begin{pgfscope}%
\definecolor{textcolor}{rgb}{0.549020,0.549020,0.549020}%
\pgfsetstrokecolor{textcolor}%
\pgfsetfillcolor{textcolor}%
\pgftext[x=3.162425in,y=4.002754in,left,]{\color{textcolor}\sffamily\fontsize{7.000000}{8.400000}\selectfont 0.48}%
\end{pgfscope}%
\begin{pgfscope}%
\pgfpathrectangle{\pgfqpoint{2.695000in}{0.598000in}}{\pgfqpoint{1.155000in}{3.864000in}}%
\pgfusepath{clip}%
\pgfsetbuttcap%
\pgfsetroundjoin%
\definecolor{currentfill}{rgb}{0.435294,0.560784,0.709804}%
\pgfsetfillcolor{currentfill}%
\pgfsetlinewidth{0.000000pt}%
\definecolor{currentstroke}{rgb}{0.435294,0.560784,0.709804}%
\pgfsetstrokecolor{currentstroke}%
\pgfsetdash{}{0pt}%
\pgfsys@defobject{currentmarker}{\pgfqpoint{-0.024056in}{-0.024056in}}{\pgfqpoint{0.024056in}{0.024056in}}{%
\pgfpathmoveto{\pgfqpoint{0.000000in}{-0.024056in}}%
\pgfpathcurveto{\pgfqpoint{0.006380in}{-0.024056in}}{\pgfqpoint{0.012499in}{-0.021522in}}{\pgfqpoint{0.017010in}{-0.017010in}}%
\pgfpathcurveto{\pgfqpoint{0.021522in}{-0.012499in}}{\pgfqpoint{0.024056in}{-0.006380in}}{\pgfqpoint{0.024056in}{0.000000in}}%
\pgfpathcurveto{\pgfqpoint{0.024056in}{0.006380in}}{\pgfqpoint{0.021522in}{0.012499in}}{\pgfqpoint{0.017010in}{0.017010in}}%
\pgfpathcurveto{\pgfqpoint{0.012499in}{0.021522in}}{\pgfqpoint{0.006380in}{0.024056in}}{\pgfqpoint{0.000000in}{0.024056in}}%
\pgfpathcurveto{\pgfqpoint{-0.006380in}{0.024056in}}{\pgfqpoint{-0.012499in}{0.021522in}}{\pgfqpoint{-0.017010in}{0.017010in}}%
\pgfpathcurveto{\pgfqpoint{-0.021522in}{0.012499in}}{\pgfqpoint{-0.024056in}{0.006380in}}{\pgfqpoint{-0.024056in}{0.000000in}}%
\pgfpathcurveto{\pgfqpoint{-0.024056in}{-0.006380in}}{\pgfqpoint{-0.021522in}{-0.012499in}}{\pgfqpoint{-0.017010in}{-0.017010in}}%
\pgfpathcurveto{\pgfqpoint{-0.012499in}{-0.021522in}}{\pgfqpoint{-0.006380in}{-0.024056in}}{\pgfqpoint{0.000000in}{-0.024056in}}%
\pgfpathlineto{\pgfqpoint{0.000000in}{-0.024056in}}%
\pgfpathclose%
\pgfusepath{fill}%
}%
\begin{pgfscope}%
\pgfsys@transformshift{3.384500in}{3.844393in}%
\pgfsys@useobject{currentmarker}{}%
\end{pgfscope}%
\end{pgfscope}%
\begin{pgfscope}%
\definecolor{textcolor}{rgb}{0.549020,0.549020,0.549020}%
\pgfsetstrokecolor{textcolor}%
\pgfsetfillcolor{textcolor}%
\pgftext[x=3.424925in,y=3.844393in,left,]{\color{textcolor}\sffamily\fontsize{7.000000}{8.400000}\selectfont 0.60}%
\end{pgfscope}%
\begin{pgfscope}%
\pgfpathrectangle{\pgfqpoint{2.695000in}{0.598000in}}{\pgfqpoint{1.155000in}{3.864000in}}%
\pgfusepath{clip}%
\pgfsetbuttcap%
\pgfsetroundjoin%
\definecolor{currentfill}{rgb}{0.435294,0.560784,0.709804}%
\pgfsetfillcolor{currentfill}%
\pgfsetlinewidth{0.000000pt}%
\definecolor{currentstroke}{rgb}{0.435294,0.560784,0.709804}%
\pgfsetstrokecolor{currentstroke}%
\pgfsetdash{}{0pt}%
\pgfsys@defobject{currentmarker}{\pgfqpoint{-0.024056in}{-0.024056in}}{\pgfqpoint{0.024056in}{0.024056in}}{%
\pgfpathmoveto{\pgfqpoint{0.000000in}{-0.024056in}}%
\pgfpathcurveto{\pgfqpoint{0.006380in}{-0.024056in}}{\pgfqpoint{0.012499in}{-0.021522in}}{\pgfqpoint{0.017010in}{-0.017010in}}%
\pgfpathcurveto{\pgfqpoint{0.021522in}{-0.012499in}}{\pgfqpoint{0.024056in}{-0.006380in}}{\pgfqpoint{0.024056in}{0.000000in}}%
\pgfpathcurveto{\pgfqpoint{0.024056in}{0.006380in}}{\pgfqpoint{0.021522in}{0.012499in}}{\pgfqpoint{0.017010in}{0.017010in}}%
\pgfpathcurveto{\pgfqpoint{0.012499in}{0.021522in}}{\pgfqpoint{0.006380in}{0.024056in}}{\pgfqpoint{0.000000in}{0.024056in}}%
\pgfpathcurveto{\pgfqpoint{-0.006380in}{0.024056in}}{\pgfqpoint{-0.012499in}{0.021522in}}{\pgfqpoint{-0.017010in}{0.017010in}}%
\pgfpathcurveto{\pgfqpoint{-0.021522in}{0.012499in}}{\pgfqpoint{-0.024056in}{0.006380in}}{\pgfqpoint{-0.024056in}{0.000000in}}%
\pgfpathcurveto{\pgfqpoint{-0.024056in}{-0.006380in}}{\pgfqpoint{-0.021522in}{-0.012499in}}{\pgfqpoint{-0.017010in}{-0.017010in}}%
\pgfpathcurveto{\pgfqpoint{-0.012499in}{-0.021522in}}{\pgfqpoint{-0.006380in}{-0.024056in}}{\pgfqpoint{0.000000in}{-0.024056in}}%
\pgfpathlineto{\pgfqpoint{0.000000in}{-0.024056in}}%
\pgfpathclose%
\pgfusepath{fill}%
}%
\begin{pgfscope}%
\pgfsys@transformshift{3.454500in}{3.686033in}%
\pgfsys@useobject{currentmarker}{}%
\end{pgfscope}%
\end{pgfscope}%
\begin{pgfscope}%
\definecolor{textcolor}{rgb}{0.549020,0.549020,0.549020}%
\pgfsetstrokecolor{textcolor}%
\pgfsetfillcolor{textcolor}%
\pgftext[x=3.494925in,y=3.686033in,left,]{\color{textcolor}\sffamily\fontsize{7.000000}{8.400000}\selectfont 0.63}%
\end{pgfscope}%
\begin{pgfscope}%
\pgfpathrectangle{\pgfqpoint{2.695000in}{0.598000in}}{\pgfqpoint{1.155000in}{3.864000in}}%
\pgfusepath{clip}%
\pgfsetbuttcap%
\pgfsetroundjoin%
\definecolor{currentfill}{rgb}{0.435294,0.560784,0.709804}%
\pgfsetfillcolor{currentfill}%
\pgfsetlinewidth{0.000000pt}%
\definecolor{currentstroke}{rgb}{0.435294,0.560784,0.709804}%
\pgfsetstrokecolor{currentstroke}%
\pgfsetdash{}{0pt}%
\pgfsys@defobject{currentmarker}{\pgfqpoint{-0.024056in}{-0.024056in}}{\pgfqpoint{0.024056in}{0.024056in}}{%
\pgfpathmoveto{\pgfqpoint{0.000000in}{-0.024056in}}%
\pgfpathcurveto{\pgfqpoint{0.006380in}{-0.024056in}}{\pgfqpoint{0.012499in}{-0.021522in}}{\pgfqpoint{0.017010in}{-0.017010in}}%
\pgfpathcurveto{\pgfqpoint{0.021522in}{-0.012499in}}{\pgfqpoint{0.024056in}{-0.006380in}}{\pgfqpoint{0.024056in}{0.000000in}}%
\pgfpathcurveto{\pgfqpoint{0.024056in}{0.006380in}}{\pgfqpoint{0.021522in}{0.012499in}}{\pgfqpoint{0.017010in}{0.017010in}}%
\pgfpathcurveto{\pgfqpoint{0.012499in}{0.021522in}}{\pgfqpoint{0.006380in}{0.024056in}}{\pgfqpoint{0.000000in}{0.024056in}}%
\pgfpathcurveto{\pgfqpoint{-0.006380in}{0.024056in}}{\pgfqpoint{-0.012499in}{0.021522in}}{\pgfqpoint{-0.017010in}{0.017010in}}%
\pgfpathcurveto{\pgfqpoint{-0.021522in}{0.012499in}}{\pgfqpoint{-0.024056in}{0.006380in}}{\pgfqpoint{-0.024056in}{0.000000in}}%
\pgfpathcurveto{\pgfqpoint{-0.024056in}{-0.006380in}}{\pgfqpoint{-0.021522in}{-0.012499in}}{\pgfqpoint{-0.017010in}{-0.017010in}}%
\pgfpathcurveto{\pgfqpoint{-0.012499in}{-0.021522in}}{\pgfqpoint{-0.006380in}{-0.024056in}}{\pgfqpoint{0.000000in}{-0.024056in}}%
\pgfpathlineto{\pgfqpoint{0.000000in}{-0.024056in}}%
\pgfpathclose%
\pgfusepath{fill}%
}%
\begin{pgfscope}%
\pgfsys@transformshift{3.472000in}{3.527672in}%
\pgfsys@useobject{currentmarker}{}%
\end{pgfscope}%
\end{pgfscope}%
\begin{pgfscope}%
\definecolor{textcolor}{rgb}{0.549020,0.549020,0.549020}%
\pgfsetstrokecolor{textcolor}%
\pgfsetfillcolor{textcolor}%
\pgftext[x=3.512425in,y=3.527672in,left,]{\color{textcolor}\sffamily\fontsize{7.000000}{8.400000}\selectfont 0.64}%
\end{pgfscope}%
\begin{pgfscope}%
\pgfpathrectangle{\pgfqpoint{2.695000in}{0.598000in}}{\pgfqpoint{1.155000in}{3.864000in}}%
\pgfusepath{clip}%
\pgfsetbuttcap%
\pgfsetroundjoin%
\definecolor{currentfill}{rgb}{0.156863,0.356863,0.568627}%
\pgfsetfillcolor{currentfill}%
\pgfsetlinewidth{0.000000pt}%
\definecolor{currentstroke}{rgb}{0.156863,0.356863,0.568627}%
\pgfsetstrokecolor{currentstroke}%
\pgfsetdash{}{0pt}%
\pgfsys@defobject{currentmarker}{\pgfqpoint{-0.024056in}{-0.024056in}}{\pgfqpoint{0.024056in}{0.024056in}}{%
\pgfpathmoveto{\pgfqpoint{0.000000in}{-0.024056in}}%
\pgfpathcurveto{\pgfqpoint{0.006380in}{-0.024056in}}{\pgfqpoint{0.012499in}{-0.021522in}}{\pgfqpoint{0.017010in}{-0.017010in}}%
\pgfpathcurveto{\pgfqpoint{0.021522in}{-0.012499in}}{\pgfqpoint{0.024056in}{-0.006380in}}{\pgfqpoint{0.024056in}{0.000000in}}%
\pgfpathcurveto{\pgfqpoint{0.024056in}{0.006380in}}{\pgfqpoint{0.021522in}{0.012499in}}{\pgfqpoint{0.017010in}{0.017010in}}%
\pgfpathcurveto{\pgfqpoint{0.012499in}{0.021522in}}{\pgfqpoint{0.006380in}{0.024056in}}{\pgfqpoint{0.000000in}{0.024056in}}%
\pgfpathcurveto{\pgfqpoint{-0.006380in}{0.024056in}}{\pgfqpoint{-0.012499in}{0.021522in}}{\pgfqpoint{-0.017010in}{0.017010in}}%
\pgfpathcurveto{\pgfqpoint{-0.021522in}{0.012499in}}{\pgfqpoint{-0.024056in}{0.006380in}}{\pgfqpoint{-0.024056in}{0.000000in}}%
\pgfpathcurveto{\pgfqpoint{-0.024056in}{-0.006380in}}{\pgfqpoint{-0.021522in}{-0.012499in}}{\pgfqpoint{-0.017010in}{-0.017010in}}%
\pgfpathcurveto{\pgfqpoint{-0.012499in}{-0.021522in}}{\pgfqpoint{-0.006380in}{-0.024056in}}{\pgfqpoint{0.000000in}{-0.024056in}}%
\pgfpathlineto{\pgfqpoint{0.000000in}{-0.024056in}}%
\pgfpathclose%
\pgfusepath{fill}%
}%
\begin{pgfscope}%
\pgfsys@transformshift{3.454500in}{3.369311in}%
\pgfsys@useobject{currentmarker}{}%
\end{pgfscope}%
\end{pgfscope}%
\begin{pgfscope}%
\definecolor{textcolor}{rgb}{0.549020,0.549020,0.549020}%
\pgfsetstrokecolor{textcolor}%
\pgfsetfillcolor{textcolor}%
\pgftext[x=3.494925in,y=3.369311in,left,]{\color{textcolor}\sffamily\fontsize{7.000000}{8.400000}\selectfont 0.63}%
\end{pgfscope}%
\begin{pgfscope}%
\pgfpathrectangle{\pgfqpoint{2.695000in}{0.598000in}}{\pgfqpoint{1.155000in}{3.864000in}}%
\pgfusepath{clip}%
\pgfsetbuttcap%
\pgfsetroundjoin%
\definecolor{currentfill}{rgb}{0.435294,0.560784,0.709804}%
\pgfsetfillcolor{currentfill}%
\pgfsetlinewidth{0.000000pt}%
\definecolor{currentstroke}{rgb}{0.435294,0.560784,0.709804}%
\pgfsetstrokecolor{currentstroke}%
\pgfsetdash{}{0pt}%
\pgfsys@defobject{currentmarker}{\pgfqpoint{-0.024056in}{-0.024056in}}{\pgfqpoint{0.024056in}{0.024056in}}{%
\pgfpathmoveto{\pgfqpoint{0.000000in}{-0.024056in}}%
\pgfpathcurveto{\pgfqpoint{0.006380in}{-0.024056in}}{\pgfqpoint{0.012499in}{-0.021522in}}{\pgfqpoint{0.017010in}{-0.017010in}}%
\pgfpathcurveto{\pgfqpoint{0.021522in}{-0.012499in}}{\pgfqpoint{0.024056in}{-0.006380in}}{\pgfqpoint{0.024056in}{0.000000in}}%
\pgfpathcurveto{\pgfqpoint{0.024056in}{0.006380in}}{\pgfqpoint{0.021522in}{0.012499in}}{\pgfqpoint{0.017010in}{0.017010in}}%
\pgfpathcurveto{\pgfqpoint{0.012499in}{0.021522in}}{\pgfqpoint{0.006380in}{0.024056in}}{\pgfqpoint{0.000000in}{0.024056in}}%
\pgfpathcurveto{\pgfqpoint{-0.006380in}{0.024056in}}{\pgfqpoint{-0.012499in}{0.021522in}}{\pgfqpoint{-0.017010in}{0.017010in}}%
\pgfpathcurveto{\pgfqpoint{-0.021522in}{0.012499in}}{\pgfqpoint{-0.024056in}{0.006380in}}{\pgfqpoint{-0.024056in}{0.000000in}}%
\pgfpathcurveto{\pgfqpoint{-0.024056in}{-0.006380in}}{\pgfqpoint{-0.021522in}{-0.012499in}}{\pgfqpoint{-0.017010in}{-0.017010in}}%
\pgfpathcurveto{\pgfqpoint{-0.012499in}{-0.021522in}}{\pgfqpoint{-0.006380in}{-0.024056in}}{\pgfqpoint{0.000000in}{-0.024056in}}%
\pgfpathlineto{\pgfqpoint{0.000000in}{-0.024056in}}%
\pgfpathclose%
\pgfusepath{fill}%
}%
\begin{pgfscope}%
\pgfsys@transformshift{3.489500in}{3.210951in}%
\pgfsys@useobject{currentmarker}{}%
\end{pgfscope}%
\end{pgfscope}%
\begin{pgfscope}%
\definecolor{textcolor}{rgb}{0.549020,0.549020,0.549020}%
\pgfsetstrokecolor{textcolor}%
\pgfsetfillcolor{textcolor}%
\pgftext[x=3.529925in,y=3.210951in,left,]{\color{textcolor}\sffamily\fontsize{7.000000}{8.400000}\selectfont 0.64}%
\end{pgfscope}%
\begin{pgfscope}%
\pgfpathrectangle{\pgfqpoint{2.695000in}{0.598000in}}{\pgfqpoint{1.155000in}{3.864000in}}%
\pgfusepath{clip}%
\pgfsetbuttcap%
\pgfsetroundjoin%
\definecolor{currentfill}{rgb}{0.435294,0.560784,0.709804}%
\pgfsetfillcolor{currentfill}%
\pgfsetlinewidth{0.000000pt}%
\definecolor{currentstroke}{rgb}{0.435294,0.560784,0.709804}%
\pgfsetstrokecolor{currentstroke}%
\pgfsetdash{}{0pt}%
\pgfsys@defobject{currentmarker}{\pgfqpoint{-0.024056in}{-0.024056in}}{\pgfqpoint{0.024056in}{0.024056in}}{%
\pgfpathmoveto{\pgfqpoint{0.000000in}{-0.024056in}}%
\pgfpathcurveto{\pgfqpoint{0.006380in}{-0.024056in}}{\pgfqpoint{0.012499in}{-0.021522in}}{\pgfqpoint{0.017010in}{-0.017010in}}%
\pgfpathcurveto{\pgfqpoint{0.021522in}{-0.012499in}}{\pgfqpoint{0.024056in}{-0.006380in}}{\pgfqpoint{0.024056in}{0.000000in}}%
\pgfpathcurveto{\pgfqpoint{0.024056in}{0.006380in}}{\pgfqpoint{0.021522in}{0.012499in}}{\pgfqpoint{0.017010in}{0.017010in}}%
\pgfpathcurveto{\pgfqpoint{0.012499in}{0.021522in}}{\pgfqpoint{0.006380in}{0.024056in}}{\pgfqpoint{0.000000in}{0.024056in}}%
\pgfpathcurveto{\pgfqpoint{-0.006380in}{0.024056in}}{\pgfqpoint{-0.012499in}{0.021522in}}{\pgfqpoint{-0.017010in}{0.017010in}}%
\pgfpathcurveto{\pgfqpoint{-0.021522in}{0.012499in}}{\pgfqpoint{-0.024056in}{0.006380in}}{\pgfqpoint{-0.024056in}{0.000000in}}%
\pgfpathcurveto{\pgfqpoint{-0.024056in}{-0.006380in}}{\pgfqpoint{-0.021522in}{-0.012499in}}{\pgfqpoint{-0.017010in}{-0.017010in}}%
\pgfpathcurveto{\pgfqpoint{-0.012499in}{-0.021522in}}{\pgfqpoint{-0.006380in}{-0.024056in}}{\pgfqpoint{0.000000in}{-0.024056in}}%
\pgfpathlineto{\pgfqpoint{0.000000in}{-0.024056in}}%
\pgfpathclose%
\pgfusepath{fill}%
}%
\begin{pgfscope}%
\pgfsys@transformshift{3.384500in}{3.052590in}%
\pgfsys@useobject{currentmarker}{}%
\end{pgfscope}%
\end{pgfscope}%
\begin{pgfscope}%
\definecolor{textcolor}{rgb}{0.549020,0.549020,0.549020}%
\pgfsetstrokecolor{textcolor}%
\pgfsetfillcolor{textcolor}%
\pgftext[x=3.424925in,y=3.052590in,left,]{\color{textcolor}\sffamily\fontsize{7.000000}{8.400000}\selectfont 0.60}%
\end{pgfscope}%
\begin{pgfscope}%
\pgfpathrectangle{\pgfqpoint{2.695000in}{0.598000in}}{\pgfqpoint{1.155000in}{3.864000in}}%
\pgfusepath{clip}%
\pgfsetbuttcap%
\pgfsetroundjoin%
\definecolor{currentfill}{rgb}{0.435294,0.560784,0.709804}%
\pgfsetfillcolor{currentfill}%
\pgfsetlinewidth{0.000000pt}%
\definecolor{currentstroke}{rgb}{0.435294,0.560784,0.709804}%
\pgfsetstrokecolor{currentstroke}%
\pgfsetdash{}{0pt}%
\pgfsys@defobject{currentmarker}{\pgfqpoint{-0.024056in}{-0.024056in}}{\pgfqpoint{0.024056in}{0.024056in}}{%
\pgfpathmoveto{\pgfqpoint{0.000000in}{-0.024056in}}%
\pgfpathcurveto{\pgfqpoint{0.006380in}{-0.024056in}}{\pgfqpoint{0.012499in}{-0.021522in}}{\pgfqpoint{0.017010in}{-0.017010in}}%
\pgfpathcurveto{\pgfqpoint{0.021522in}{-0.012499in}}{\pgfqpoint{0.024056in}{-0.006380in}}{\pgfqpoint{0.024056in}{0.000000in}}%
\pgfpathcurveto{\pgfqpoint{0.024056in}{0.006380in}}{\pgfqpoint{0.021522in}{0.012499in}}{\pgfqpoint{0.017010in}{0.017010in}}%
\pgfpathcurveto{\pgfqpoint{0.012499in}{0.021522in}}{\pgfqpoint{0.006380in}{0.024056in}}{\pgfqpoint{0.000000in}{0.024056in}}%
\pgfpathcurveto{\pgfqpoint{-0.006380in}{0.024056in}}{\pgfqpoint{-0.012499in}{0.021522in}}{\pgfqpoint{-0.017010in}{0.017010in}}%
\pgfpathcurveto{\pgfqpoint{-0.021522in}{0.012499in}}{\pgfqpoint{-0.024056in}{0.006380in}}{\pgfqpoint{-0.024056in}{0.000000in}}%
\pgfpathcurveto{\pgfqpoint{-0.024056in}{-0.006380in}}{\pgfqpoint{-0.021522in}{-0.012499in}}{\pgfqpoint{-0.017010in}{-0.017010in}}%
\pgfpathcurveto{\pgfqpoint{-0.012499in}{-0.021522in}}{\pgfqpoint{-0.006380in}{-0.024056in}}{\pgfqpoint{0.000000in}{-0.024056in}}%
\pgfpathlineto{\pgfqpoint{0.000000in}{-0.024056in}}%
\pgfpathclose%
\pgfusepath{fill}%
}%
\begin{pgfscope}%
\pgfsys@transformshift{3.454500in}{2.894230in}%
\pgfsys@useobject{currentmarker}{}%
\end{pgfscope}%
\end{pgfscope}%
\begin{pgfscope}%
\definecolor{textcolor}{rgb}{0.549020,0.549020,0.549020}%
\pgfsetstrokecolor{textcolor}%
\pgfsetfillcolor{textcolor}%
\pgftext[x=3.494925in,y=2.894230in,left,]{\color{textcolor}\sffamily\fontsize{7.000000}{8.400000}\selectfont 0.63}%
\end{pgfscope}%
\begin{pgfscope}%
\pgfpathrectangle{\pgfqpoint{2.695000in}{0.598000in}}{\pgfqpoint{1.155000in}{3.864000in}}%
\pgfusepath{clip}%
\pgfsetbuttcap%
\pgfsetroundjoin%
\definecolor{currentfill}{rgb}{0.435294,0.560784,0.709804}%
\pgfsetfillcolor{currentfill}%
\pgfsetlinewidth{0.000000pt}%
\definecolor{currentstroke}{rgb}{0.435294,0.560784,0.709804}%
\pgfsetstrokecolor{currentstroke}%
\pgfsetdash{}{0pt}%
\pgfsys@defobject{currentmarker}{\pgfqpoint{-0.024056in}{-0.024056in}}{\pgfqpoint{0.024056in}{0.024056in}}{%
\pgfpathmoveto{\pgfqpoint{0.000000in}{-0.024056in}}%
\pgfpathcurveto{\pgfqpoint{0.006380in}{-0.024056in}}{\pgfqpoint{0.012499in}{-0.021522in}}{\pgfqpoint{0.017010in}{-0.017010in}}%
\pgfpathcurveto{\pgfqpoint{0.021522in}{-0.012499in}}{\pgfqpoint{0.024056in}{-0.006380in}}{\pgfqpoint{0.024056in}{0.000000in}}%
\pgfpathcurveto{\pgfqpoint{0.024056in}{0.006380in}}{\pgfqpoint{0.021522in}{0.012499in}}{\pgfqpoint{0.017010in}{0.017010in}}%
\pgfpathcurveto{\pgfqpoint{0.012499in}{0.021522in}}{\pgfqpoint{0.006380in}{0.024056in}}{\pgfqpoint{0.000000in}{0.024056in}}%
\pgfpathcurveto{\pgfqpoint{-0.006380in}{0.024056in}}{\pgfqpoint{-0.012499in}{0.021522in}}{\pgfqpoint{-0.017010in}{0.017010in}}%
\pgfpathcurveto{\pgfqpoint{-0.021522in}{0.012499in}}{\pgfqpoint{-0.024056in}{0.006380in}}{\pgfqpoint{-0.024056in}{0.000000in}}%
\pgfpathcurveto{\pgfqpoint{-0.024056in}{-0.006380in}}{\pgfqpoint{-0.021522in}{-0.012499in}}{\pgfqpoint{-0.017010in}{-0.017010in}}%
\pgfpathcurveto{\pgfqpoint{-0.012499in}{-0.021522in}}{\pgfqpoint{-0.006380in}{-0.024056in}}{\pgfqpoint{0.000000in}{-0.024056in}}%
\pgfpathlineto{\pgfqpoint{0.000000in}{-0.024056in}}%
\pgfpathclose%
\pgfusepath{fill}%
}%
\begin{pgfscope}%
\pgfsys@transformshift{3.419500in}{2.735869in}%
\pgfsys@useobject{currentmarker}{}%
\end{pgfscope}%
\end{pgfscope}%
\begin{pgfscope}%
\definecolor{textcolor}{rgb}{0.549020,0.549020,0.549020}%
\pgfsetstrokecolor{textcolor}%
\pgfsetfillcolor{textcolor}%
\pgftext[x=3.459925in,y=2.735869in,left,]{\color{textcolor}\sffamily\fontsize{7.000000}{8.400000}\selectfont 0.61}%
\end{pgfscope}%
\begin{pgfscope}%
\pgfpathrectangle{\pgfqpoint{2.695000in}{0.598000in}}{\pgfqpoint{1.155000in}{3.864000in}}%
\pgfusepath{clip}%
\pgfsetbuttcap%
\pgfsetroundjoin%
\definecolor{currentfill}{rgb}{0.156863,0.356863,0.568627}%
\pgfsetfillcolor{currentfill}%
\pgfsetlinewidth{0.000000pt}%
\definecolor{currentstroke}{rgb}{0.156863,0.356863,0.568627}%
\pgfsetstrokecolor{currentstroke}%
\pgfsetdash{}{0pt}%
\pgfsys@defobject{currentmarker}{\pgfqpoint{-0.024056in}{-0.024056in}}{\pgfqpoint{0.024056in}{0.024056in}}{%
\pgfpathmoveto{\pgfqpoint{0.000000in}{-0.024056in}}%
\pgfpathcurveto{\pgfqpoint{0.006380in}{-0.024056in}}{\pgfqpoint{0.012499in}{-0.021522in}}{\pgfqpoint{0.017010in}{-0.017010in}}%
\pgfpathcurveto{\pgfqpoint{0.021522in}{-0.012499in}}{\pgfqpoint{0.024056in}{-0.006380in}}{\pgfqpoint{0.024056in}{0.000000in}}%
\pgfpathcurveto{\pgfqpoint{0.024056in}{0.006380in}}{\pgfqpoint{0.021522in}{0.012499in}}{\pgfqpoint{0.017010in}{0.017010in}}%
\pgfpathcurveto{\pgfqpoint{0.012499in}{0.021522in}}{\pgfqpoint{0.006380in}{0.024056in}}{\pgfqpoint{0.000000in}{0.024056in}}%
\pgfpathcurveto{\pgfqpoint{-0.006380in}{0.024056in}}{\pgfqpoint{-0.012499in}{0.021522in}}{\pgfqpoint{-0.017010in}{0.017010in}}%
\pgfpathcurveto{\pgfqpoint{-0.021522in}{0.012499in}}{\pgfqpoint{-0.024056in}{0.006380in}}{\pgfqpoint{-0.024056in}{0.000000in}}%
\pgfpathcurveto{\pgfqpoint{-0.024056in}{-0.006380in}}{\pgfqpoint{-0.021522in}{-0.012499in}}{\pgfqpoint{-0.017010in}{-0.017010in}}%
\pgfpathcurveto{\pgfqpoint{-0.012499in}{-0.021522in}}{\pgfqpoint{-0.006380in}{-0.024056in}}{\pgfqpoint{0.000000in}{-0.024056in}}%
\pgfpathlineto{\pgfqpoint{0.000000in}{-0.024056in}}%
\pgfpathclose%
\pgfusepath{fill}%
}%
\begin{pgfscope}%
\pgfsys@transformshift{3.472000in}{2.577508in}%
\pgfsys@useobject{currentmarker}{}%
\end{pgfscope}%
\end{pgfscope}%
\begin{pgfscope}%
\definecolor{textcolor}{rgb}{0.549020,0.549020,0.549020}%
\pgfsetstrokecolor{textcolor}%
\pgfsetfillcolor{textcolor}%
\pgftext[x=3.512425in,y=2.577508in,left,]{\color{textcolor}\sffamily\fontsize{7.000000}{8.400000}\selectfont 0.64}%
\end{pgfscope}%
\begin{pgfscope}%
\pgfpathrectangle{\pgfqpoint{2.695000in}{0.598000in}}{\pgfqpoint{1.155000in}{3.864000in}}%
\pgfusepath{clip}%
\pgfsetbuttcap%
\pgfsetroundjoin%
\definecolor{currentfill}{rgb}{0.156863,0.356863,0.568627}%
\pgfsetfillcolor{currentfill}%
\pgfsetlinewidth{0.000000pt}%
\definecolor{currentstroke}{rgb}{0.156863,0.356863,0.568627}%
\pgfsetstrokecolor{currentstroke}%
\pgfsetdash{}{0pt}%
\pgfsys@defobject{currentmarker}{\pgfqpoint{-0.024056in}{-0.024056in}}{\pgfqpoint{0.024056in}{0.024056in}}{%
\pgfpathmoveto{\pgfqpoint{0.000000in}{-0.024056in}}%
\pgfpathcurveto{\pgfqpoint{0.006380in}{-0.024056in}}{\pgfqpoint{0.012499in}{-0.021522in}}{\pgfqpoint{0.017010in}{-0.017010in}}%
\pgfpathcurveto{\pgfqpoint{0.021522in}{-0.012499in}}{\pgfqpoint{0.024056in}{-0.006380in}}{\pgfqpoint{0.024056in}{0.000000in}}%
\pgfpathcurveto{\pgfqpoint{0.024056in}{0.006380in}}{\pgfqpoint{0.021522in}{0.012499in}}{\pgfqpoint{0.017010in}{0.017010in}}%
\pgfpathcurveto{\pgfqpoint{0.012499in}{0.021522in}}{\pgfqpoint{0.006380in}{0.024056in}}{\pgfqpoint{0.000000in}{0.024056in}}%
\pgfpathcurveto{\pgfqpoint{-0.006380in}{0.024056in}}{\pgfqpoint{-0.012499in}{0.021522in}}{\pgfqpoint{-0.017010in}{0.017010in}}%
\pgfpathcurveto{\pgfqpoint{-0.021522in}{0.012499in}}{\pgfqpoint{-0.024056in}{0.006380in}}{\pgfqpoint{-0.024056in}{0.000000in}}%
\pgfpathcurveto{\pgfqpoint{-0.024056in}{-0.006380in}}{\pgfqpoint{-0.021522in}{-0.012499in}}{\pgfqpoint{-0.017010in}{-0.017010in}}%
\pgfpathcurveto{\pgfqpoint{-0.012499in}{-0.021522in}}{\pgfqpoint{-0.006380in}{-0.024056in}}{\pgfqpoint{0.000000in}{-0.024056in}}%
\pgfpathlineto{\pgfqpoint{0.000000in}{-0.024056in}}%
\pgfpathclose%
\pgfusepath{fill}%
}%
\begin{pgfscope}%
\pgfsys@transformshift{3.071293in}{2.189525in}%
\pgfsys@useobject{currentmarker}{}%
\end{pgfscope}%
\end{pgfscope}%
\begin{pgfscope}%
\definecolor{textcolor}{rgb}{0.549020,0.549020,0.549020}%
\pgfsetstrokecolor{textcolor}%
\pgfsetfillcolor{textcolor}%
\pgftext[x=3.111718in,y=2.189525in,left,]{\color{textcolor}\sffamily\fontsize{7.000000}{8.400000}\selectfont 0.46}%
\end{pgfscope}%
\begin{pgfscope}%
\pgfpathrectangle{\pgfqpoint{2.695000in}{0.598000in}}{\pgfqpoint{1.155000in}{3.864000in}}%
\pgfusepath{clip}%
\pgfsetbuttcap%
\pgfsetroundjoin%
\definecolor{currentfill}{rgb}{0.156863,0.356863,0.568627}%
\pgfsetfillcolor{currentfill}%
\pgfsetlinewidth{0.000000pt}%
\definecolor{currentstroke}{rgb}{0.156863,0.356863,0.568627}%
\pgfsetstrokecolor{currentstroke}%
\pgfsetdash{}{0pt}%
\pgfsys@defobject{currentmarker}{\pgfqpoint{-0.024056in}{-0.024056in}}{\pgfqpoint{0.024056in}{0.024056in}}{%
\pgfpathmoveto{\pgfqpoint{0.000000in}{-0.024056in}}%
\pgfpathcurveto{\pgfqpoint{0.006380in}{-0.024056in}}{\pgfqpoint{0.012499in}{-0.021522in}}{\pgfqpoint{0.017010in}{-0.017010in}}%
\pgfpathcurveto{\pgfqpoint{0.021522in}{-0.012499in}}{\pgfqpoint{0.024056in}{-0.006380in}}{\pgfqpoint{0.024056in}{0.000000in}}%
\pgfpathcurveto{\pgfqpoint{0.024056in}{0.006380in}}{\pgfqpoint{0.021522in}{0.012499in}}{\pgfqpoint{0.017010in}{0.017010in}}%
\pgfpathcurveto{\pgfqpoint{0.012499in}{0.021522in}}{\pgfqpoint{0.006380in}{0.024056in}}{\pgfqpoint{0.000000in}{0.024056in}}%
\pgfpathcurveto{\pgfqpoint{-0.006380in}{0.024056in}}{\pgfqpoint{-0.012499in}{0.021522in}}{\pgfqpoint{-0.017010in}{0.017010in}}%
\pgfpathcurveto{\pgfqpoint{-0.021522in}{0.012499in}}{\pgfqpoint{-0.024056in}{0.006380in}}{\pgfqpoint{-0.024056in}{0.000000in}}%
\pgfpathcurveto{\pgfqpoint{-0.024056in}{-0.006380in}}{\pgfqpoint{-0.021522in}{-0.012499in}}{\pgfqpoint{-0.017010in}{-0.017010in}}%
\pgfpathcurveto{\pgfqpoint{-0.012499in}{-0.021522in}}{\pgfqpoint{-0.006380in}{-0.024056in}}{\pgfqpoint{0.000000in}{-0.024056in}}%
\pgfpathlineto{\pgfqpoint{0.000000in}{-0.024056in}}%
\pgfpathclose%
\pgfusepath{fill}%
}%
\begin{pgfscope}%
\pgfsys@transformshift{3.259032in}{2.031164in}%
\pgfsys@useobject{currentmarker}{}%
\end{pgfscope}%
\end{pgfscope}%
\begin{pgfscope}%
\definecolor{textcolor}{rgb}{0.549020,0.549020,0.549020}%
\pgfsetstrokecolor{textcolor}%
\pgfsetfillcolor{textcolor}%
\pgftext[x=3.299457in,y=2.031164in,left,]{\color{textcolor}\sffamily\fontsize{7.000000}{8.400000}\selectfont 0.54}%
\end{pgfscope}%
\begin{pgfscope}%
\pgfpathrectangle{\pgfqpoint{2.695000in}{0.598000in}}{\pgfqpoint{1.155000in}{3.864000in}}%
\pgfusepath{clip}%
\pgfsetbuttcap%
\pgfsetroundjoin%
\definecolor{currentfill}{rgb}{0.101961,0.101961,0.101961}%
\pgfsetfillcolor{currentfill}%
\pgfsetlinewidth{0.000000pt}%
\definecolor{currentstroke}{rgb}{0.101961,0.101961,0.101961}%
\pgfsetstrokecolor{currentstroke}%
\pgfsetdash{}{0pt}%
\pgfsys@defobject{currentmarker}{\pgfqpoint{-0.032572in}{-0.032572in}}{\pgfqpoint{0.032572in}{0.032572in}}{%
\pgfpathmoveto{\pgfqpoint{-0.016286in}{-0.032572in}}%
\pgfpathlineto{\pgfqpoint{0.000000in}{-0.016286in}}%
\pgfpathlineto{\pgfqpoint{0.016286in}{-0.032572in}}%
\pgfpathlineto{\pgfqpoint{0.032572in}{-0.016286in}}%
\pgfpathlineto{\pgfqpoint{0.016286in}{0.000000in}}%
\pgfpathlineto{\pgfqpoint{0.032572in}{0.016286in}}%
\pgfpathlineto{\pgfqpoint{0.016286in}{0.032572in}}%
\pgfpathlineto{\pgfqpoint{0.000000in}{0.016286in}}%
\pgfpathlineto{\pgfqpoint{-0.016286in}{0.032572in}}%
\pgfpathlineto{\pgfqpoint{-0.032572in}{0.016286in}}%
\pgfpathlineto{\pgfqpoint{-0.016286in}{0.000000in}}%
\pgfpathlineto{\pgfqpoint{-0.032572in}{-0.016286in}}%
\pgfpathlineto{\pgfqpoint{-0.016286in}{-0.032572in}}%
\pgfpathclose%
\pgfusepath{fill}%
}%
\begin{pgfscope}%
\pgfsys@transformshift{3.079456in}{1.872803in}%
\pgfsys@useobject{currentmarker}{}%
\end{pgfscope}%
\end{pgfscope}%
\begin{pgfscope}%
\definecolor{textcolor}{rgb}{0.549020,0.549020,0.549020}%
\pgfsetstrokecolor{textcolor}%
\pgfsetfillcolor{textcolor}%
\pgftext[x=3.119881in,y=1.872803in,left,]{\color{textcolor}\sffamily\fontsize{7.000000}{8.400000}\selectfont 0.47}%
\end{pgfscope}%
\begin{pgfscope}%
\pgfpathrectangle{\pgfqpoint{2.695000in}{0.598000in}}{\pgfqpoint{1.155000in}{3.864000in}}%
\pgfusepath{clip}%
\pgfsetbuttcap%
\pgfsetroundjoin%
\definecolor{currentfill}{rgb}{0.101961,0.101961,0.101961}%
\pgfsetfillcolor{currentfill}%
\pgfsetlinewidth{0.000000pt}%
\definecolor{currentstroke}{rgb}{0.101961,0.101961,0.101961}%
\pgfsetstrokecolor{currentstroke}%
\pgfsetdash{}{0pt}%
\pgfsys@defobject{currentmarker}{\pgfqpoint{-0.032572in}{-0.032572in}}{\pgfqpoint{0.032572in}{0.032572in}}{%
\pgfpathmoveto{\pgfqpoint{-0.016286in}{-0.032572in}}%
\pgfpathlineto{\pgfqpoint{0.000000in}{-0.016286in}}%
\pgfpathlineto{\pgfqpoint{0.016286in}{-0.032572in}}%
\pgfpathlineto{\pgfqpoint{0.032572in}{-0.016286in}}%
\pgfpathlineto{\pgfqpoint{0.016286in}{0.000000in}}%
\pgfpathlineto{\pgfqpoint{0.032572in}{0.016286in}}%
\pgfpathlineto{\pgfqpoint{0.016286in}{0.032572in}}%
\pgfpathlineto{\pgfqpoint{0.000000in}{0.016286in}}%
\pgfpathlineto{\pgfqpoint{-0.016286in}{0.032572in}}%
\pgfpathlineto{\pgfqpoint{-0.032572in}{0.016286in}}%
\pgfpathlineto{\pgfqpoint{-0.016286in}{0.000000in}}%
\pgfpathlineto{\pgfqpoint{-0.032572in}{-0.016286in}}%
\pgfpathlineto{\pgfqpoint{-0.016286in}{-0.032572in}}%
\pgfpathclose%
\pgfusepath{fill}%
}%
\begin{pgfscope}%
\pgfsys@transformshift{3.079456in}{1.714443in}%
\pgfsys@useobject{currentmarker}{}%
\end{pgfscope}%
\end{pgfscope}%
\begin{pgfscope}%
\definecolor{textcolor}{rgb}{0.549020,0.549020,0.549020}%
\pgfsetstrokecolor{textcolor}%
\pgfsetfillcolor{textcolor}%
\pgftext[x=3.119881in,y=1.714443in,left,]{\color{textcolor}\sffamily\fontsize{7.000000}{8.400000}\selectfont 0.47}%
\end{pgfscope}%
\begin{pgfscope}%
\pgfpathrectangle{\pgfqpoint{2.695000in}{0.598000in}}{\pgfqpoint{1.155000in}{3.864000in}}%
\pgfusepath{clip}%
\pgfsetbuttcap%
\pgfsetroundjoin%
\definecolor{currentfill}{rgb}{0.156863,0.356863,0.568627}%
\pgfsetfillcolor{currentfill}%
\pgfsetlinewidth{0.000000pt}%
\definecolor{currentstroke}{rgb}{0.156863,0.356863,0.568627}%
\pgfsetstrokecolor{currentstroke}%
\pgfsetdash{}{0pt}%
\pgfsys@defobject{currentmarker}{\pgfqpoint{-0.024056in}{-0.024056in}}{\pgfqpoint{0.024056in}{0.024056in}}{%
\pgfpathmoveto{\pgfqpoint{0.000000in}{-0.024056in}}%
\pgfpathcurveto{\pgfqpoint{0.006380in}{-0.024056in}}{\pgfqpoint{0.012499in}{-0.021522in}}{\pgfqpoint{0.017010in}{-0.017010in}}%
\pgfpathcurveto{\pgfqpoint{0.021522in}{-0.012499in}}{\pgfqpoint{0.024056in}{-0.006380in}}{\pgfqpoint{0.024056in}{0.000000in}}%
\pgfpathcurveto{\pgfqpoint{0.024056in}{0.006380in}}{\pgfqpoint{0.021522in}{0.012499in}}{\pgfqpoint{0.017010in}{0.017010in}}%
\pgfpathcurveto{\pgfqpoint{0.012499in}{0.021522in}}{\pgfqpoint{0.006380in}{0.024056in}}{\pgfqpoint{0.000000in}{0.024056in}}%
\pgfpathcurveto{\pgfqpoint{-0.006380in}{0.024056in}}{\pgfqpoint{-0.012499in}{0.021522in}}{\pgfqpoint{-0.017010in}{0.017010in}}%
\pgfpathcurveto{\pgfqpoint{-0.021522in}{0.012499in}}{\pgfqpoint{-0.024056in}{0.006380in}}{\pgfqpoint{-0.024056in}{0.000000in}}%
\pgfpathcurveto{\pgfqpoint{-0.024056in}{-0.006380in}}{\pgfqpoint{-0.021522in}{-0.012499in}}{\pgfqpoint{-0.017010in}{-0.017010in}}%
\pgfpathcurveto{\pgfqpoint{-0.012499in}{-0.021522in}}{\pgfqpoint{-0.006380in}{-0.024056in}}{\pgfqpoint{0.000000in}{-0.024056in}}%
\pgfpathlineto{\pgfqpoint{0.000000in}{-0.024056in}}%
\pgfpathclose%
\pgfusepath{fill}%
}%
\begin{pgfscope}%
\pgfsys@transformshift{3.005993in}{1.326459in}%
\pgfsys@useobject{currentmarker}{}%
\end{pgfscope}%
\end{pgfscope}%
\begin{pgfscope}%
\definecolor{textcolor}{rgb}{0.549020,0.549020,0.549020}%
\pgfsetstrokecolor{textcolor}%
\pgfsetfillcolor{textcolor}%
\pgftext[x=3.046418in,y=1.326459in,left,]{\color{textcolor}\sffamily\fontsize{7.000000}{8.400000}\selectfont 0.43}%
\end{pgfscope}%
\begin{pgfscope}%
\pgfpathrectangle{\pgfqpoint{2.695000in}{0.598000in}}{\pgfqpoint{1.155000in}{3.864000in}}%
\pgfusepath{clip}%
\pgfsetbuttcap%
\pgfsetroundjoin%
\definecolor{currentfill}{rgb}{0.156863,0.356863,0.568627}%
\pgfsetfillcolor{currentfill}%
\pgfsetlinewidth{0.000000pt}%
\definecolor{currentstroke}{rgb}{0.156863,0.356863,0.568627}%
\pgfsetstrokecolor{currentstroke}%
\pgfsetdash{}{0pt}%
\pgfsys@defobject{currentmarker}{\pgfqpoint{-0.024056in}{-0.024056in}}{\pgfqpoint{0.024056in}{0.024056in}}{%
\pgfpathmoveto{\pgfqpoint{0.000000in}{-0.024056in}}%
\pgfpathcurveto{\pgfqpoint{0.006380in}{-0.024056in}}{\pgfqpoint{0.012499in}{-0.021522in}}{\pgfqpoint{0.017010in}{-0.017010in}}%
\pgfpathcurveto{\pgfqpoint{0.021522in}{-0.012499in}}{\pgfqpoint{0.024056in}{-0.006380in}}{\pgfqpoint{0.024056in}{0.000000in}}%
\pgfpathcurveto{\pgfqpoint{0.024056in}{0.006380in}}{\pgfqpoint{0.021522in}{0.012499in}}{\pgfqpoint{0.017010in}{0.017010in}}%
\pgfpathcurveto{\pgfqpoint{0.012499in}{0.021522in}}{\pgfqpoint{0.006380in}{0.024056in}}{\pgfqpoint{0.000000in}{0.024056in}}%
\pgfpathcurveto{\pgfqpoint{-0.006380in}{0.024056in}}{\pgfqpoint{-0.012499in}{0.021522in}}{\pgfqpoint{-0.017010in}{0.017010in}}%
\pgfpathcurveto{\pgfqpoint{-0.021522in}{0.012499in}}{\pgfqpoint{-0.024056in}{0.006380in}}{\pgfqpoint{-0.024056in}{0.000000in}}%
\pgfpathcurveto{\pgfqpoint{-0.024056in}{-0.006380in}}{\pgfqpoint{-0.021522in}{-0.012499in}}{\pgfqpoint{-0.017010in}{-0.017010in}}%
\pgfpathcurveto{\pgfqpoint{-0.012499in}{-0.021522in}}{\pgfqpoint{-0.006380in}{-0.024056in}}{\pgfqpoint{0.000000in}{-0.024056in}}%
\pgfpathlineto{\pgfqpoint{0.000000in}{-0.024056in}}%
\pgfpathclose%
\pgfusepath{fill}%
}%
\begin{pgfscope}%
\pgfsys@transformshift{3.087618in}{1.168098in}%
\pgfsys@useobject{currentmarker}{}%
\end{pgfscope}%
\end{pgfscope}%
\begin{pgfscope}%
\definecolor{textcolor}{rgb}{0.549020,0.549020,0.549020}%
\pgfsetstrokecolor{textcolor}%
\pgfsetfillcolor{textcolor}%
\pgftext[x=3.128043in,y=1.168098in,left,]{\color{textcolor}\sffamily\fontsize{7.000000}{8.400000}\selectfont 0.47}%
\end{pgfscope}%
\begin{pgfscope}%
\pgfpathrectangle{\pgfqpoint{2.695000in}{0.598000in}}{\pgfqpoint{1.155000in}{3.864000in}}%
\pgfusepath{clip}%
\pgfsetbuttcap%
\pgfsetroundjoin%
\definecolor{currentfill}{rgb}{0.156863,0.356863,0.568627}%
\pgfsetfillcolor{currentfill}%
\pgfsetlinewidth{0.000000pt}%
\definecolor{currentstroke}{rgb}{0.156863,0.356863,0.568627}%
\pgfsetstrokecolor{currentstroke}%
\pgfsetdash{}{0pt}%
\pgfsys@defobject{currentmarker}{\pgfqpoint{-0.024056in}{-0.024056in}}{\pgfqpoint{0.024056in}{0.024056in}}{%
\pgfpathmoveto{\pgfqpoint{0.000000in}{-0.024056in}}%
\pgfpathcurveto{\pgfqpoint{0.006380in}{-0.024056in}}{\pgfqpoint{0.012499in}{-0.021522in}}{\pgfqpoint{0.017010in}{-0.017010in}}%
\pgfpathcurveto{\pgfqpoint{0.021522in}{-0.012499in}}{\pgfqpoint{0.024056in}{-0.006380in}}{\pgfqpoint{0.024056in}{0.000000in}}%
\pgfpathcurveto{\pgfqpoint{0.024056in}{0.006380in}}{\pgfqpoint{0.021522in}{0.012499in}}{\pgfqpoint{0.017010in}{0.017010in}}%
\pgfpathcurveto{\pgfqpoint{0.012499in}{0.021522in}}{\pgfqpoint{0.006380in}{0.024056in}}{\pgfqpoint{0.000000in}{0.024056in}}%
\pgfpathcurveto{\pgfqpoint{-0.006380in}{0.024056in}}{\pgfqpoint{-0.012499in}{0.021522in}}{\pgfqpoint{-0.017010in}{0.017010in}}%
\pgfpathcurveto{\pgfqpoint{-0.021522in}{0.012499in}}{\pgfqpoint{-0.024056in}{0.006380in}}{\pgfqpoint{-0.024056in}{0.000000in}}%
\pgfpathcurveto{\pgfqpoint{-0.024056in}{-0.006380in}}{\pgfqpoint{-0.021522in}{-0.012499in}}{\pgfqpoint{-0.017010in}{-0.017010in}}%
\pgfpathcurveto{\pgfqpoint{-0.012499in}{-0.021522in}}{\pgfqpoint{-0.006380in}{-0.024056in}}{\pgfqpoint{0.000000in}{-0.024056in}}%
\pgfpathlineto{\pgfqpoint{0.000000in}{-0.024056in}}%
\pgfpathclose%
\pgfusepath{fill}%
}%
\begin{pgfscope}%
\pgfsys@transformshift{3.095781in}{1.009738in}%
\pgfsys@useobject{currentmarker}{}%
\end{pgfscope}%
\end{pgfscope}%
\begin{pgfscope}%
\definecolor{textcolor}{rgb}{0.549020,0.549020,0.549020}%
\pgfsetstrokecolor{textcolor}%
\pgfsetfillcolor{textcolor}%
\pgftext[x=3.136206in,y=1.009738in,left,]{\color{textcolor}\sffamily\fontsize{7.000000}{8.400000}\selectfont 0.47}%
\end{pgfscope}%
\begin{pgfscope}%
\pgfpathrectangle{\pgfqpoint{2.695000in}{0.598000in}}{\pgfqpoint{1.155000in}{3.864000in}}%
\pgfusepath{clip}%
\pgfsetbuttcap%
\pgfsetroundjoin%
\definecolor{currentfill}{rgb}{0.101961,0.101961,0.101961}%
\pgfsetfillcolor{currentfill}%
\pgfsetlinewidth{0.000000pt}%
\definecolor{currentstroke}{rgb}{0.101961,0.101961,0.101961}%
\pgfsetstrokecolor{currentstroke}%
\pgfsetdash{}{0pt}%
\pgfsys@defobject{currentmarker}{\pgfqpoint{-0.032572in}{-0.032572in}}{\pgfqpoint{0.032572in}{0.032572in}}{%
\pgfpathmoveto{\pgfqpoint{-0.016286in}{-0.032572in}}%
\pgfpathlineto{\pgfqpoint{0.000000in}{-0.016286in}}%
\pgfpathlineto{\pgfqpoint{0.016286in}{-0.032572in}}%
\pgfpathlineto{\pgfqpoint{0.032572in}{-0.016286in}}%
\pgfpathlineto{\pgfqpoint{0.016286in}{0.000000in}}%
\pgfpathlineto{\pgfqpoint{0.032572in}{0.016286in}}%
\pgfpathlineto{\pgfqpoint{0.016286in}{0.032572in}}%
\pgfpathlineto{\pgfqpoint{0.000000in}{0.016286in}}%
\pgfpathlineto{\pgfqpoint{-0.016286in}{0.032572in}}%
\pgfpathlineto{\pgfqpoint{-0.032572in}{0.016286in}}%
\pgfpathlineto{\pgfqpoint{-0.016286in}{0.000000in}}%
\pgfpathlineto{\pgfqpoint{-0.032572in}{-0.016286in}}%
\pgfpathlineto{\pgfqpoint{-0.016286in}{-0.032572in}}%
\pgfpathclose%
\pgfusepath{fill}%
}%
\begin{pgfscope}%
\pgfsys@transformshift{3.095781in}{0.851377in}%
\pgfsys@useobject{currentmarker}{}%
\end{pgfscope}%
\end{pgfscope}%
\begin{pgfscope}%
\definecolor{textcolor}{rgb}{0.549020,0.549020,0.549020}%
\pgfsetstrokecolor{textcolor}%
\pgfsetfillcolor{textcolor}%
\pgftext[x=3.136206in,y=0.851377in,left,]{\color{textcolor}\sffamily\fontsize{7.000000}{8.400000}\selectfont 0.47}%
\end{pgfscope}%
\begin{pgfscope}%
\pgfpathrectangle{\pgfqpoint{2.695000in}{0.598000in}}{\pgfqpoint{1.155000in}{3.864000in}}%
\pgfusepath{clip}%
\pgfsetbuttcap%
\pgfsetroundjoin%
\definecolor{currentfill}{rgb}{0.101961,0.101961,0.101961}%
\pgfsetfillcolor{currentfill}%
\pgfsetlinewidth{0.000000pt}%
\definecolor{currentstroke}{rgb}{0.101961,0.101961,0.101961}%
\pgfsetstrokecolor{currentstroke}%
\pgfsetdash{}{0pt}%
\pgfsys@defobject{currentmarker}{\pgfqpoint{-0.032572in}{-0.032572in}}{\pgfqpoint{0.032572in}{0.032572in}}{%
\pgfpathmoveto{\pgfqpoint{-0.016286in}{-0.032572in}}%
\pgfpathlineto{\pgfqpoint{0.000000in}{-0.016286in}}%
\pgfpathlineto{\pgfqpoint{0.016286in}{-0.032572in}}%
\pgfpathlineto{\pgfqpoint{0.032572in}{-0.016286in}}%
\pgfpathlineto{\pgfqpoint{0.016286in}{0.000000in}}%
\pgfpathlineto{\pgfqpoint{0.032572in}{0.016286in}}%
\pgfpathlineto{\pgfqpoint{0.016286in}{0.032572in}}%
\pgfpathlineto{\pgfqpoint{0.000000in}{0.016286in}}%
\pgfpathlineto{\pgfqpoint{-0.016286in}{0.032572in}}%
\pgfpathlineto{\pgfqpoint{-0.032572in}{0.016286in}}%
\pgfpathlineto{\pgfqpoint{-0.016286in}{0.000000in}}%
\pgfpathlineto{\pgfqpoint{-0.032572in}{-0.016286in}}%
\pgfpathlineto{\pgfqpoint{-0.016286in}{-0.032572in}}%
\pgfpathclose%
\pgfusepath{fill}%
}%
\begin{pgfscope}%
\pgfsys@transformshift{3.014155in}{0.693016in}%
\pgfsys@useobject{currentmarker}{}%
\end{pgfscope}%
\end{pgfscope}%
\begin{pgfscope}%
\definecolor{textcolor}{rgb}{0.549020,0.549020,0.549020}%
\pgfsetstrokecolor{textcolor}%
\pgfsetfillcolor{textcolor}%
\pgftext[x=3.054580in,y=0.693016in,left,]{\color{textcolor}\sffamily\fontsize{7.000000}{8.400000}\selectfont 0.44}%
\end{pgfscope}%
\begin{pgfscope}%
\definecolor{textcolor}{rgb}{0.549020,0.549020,0.549020}%
\pgfsetstrokecolor{textcolor}%
\pgfsetfillcolor{textcolor}%
\pgftext[x=2.660350in,y=4.319475in,left,]{\makebox[0pt][r]{\color{textcolor}\sffamily\fontsize{7.000000}{8.400000}\itshape\selectfont n16, 1x, standard mix}}%
\end{pgfscope}%
\begin{pgfscope}%
\definecolor{textcolor}{rgb}{0.549020,0.549020,0.549020}%
\pgfsetstrokecolor{textcolor}%
\pgfsetfillcolor{textcolor}%
\pgftext[x=2.660350in,y=2.347885in,left,]{\makebox[0pt][r]{\color{textcolor}\sffamily\fontsize{7.000000}{8.400000}\itshape\selectfont n20, 3:11:2 mix, no outage}}%
\end{pgfscope}%
\begin{pgfscope}%
\definecolor{textcolor}{rgb}{0.549020,0.549020,0.549020}%
\pgfsetstrokecolor{textcolor}%
\pgfsetfillcolor{textcolor}%
\pgftext[x=2.660350in,y=1.484820in,left,]{\makebox[0pt][r]{\color{textcolor}\sffamily\fontsize{7.000000}{8.400000}\itshape\selectfont n20, 3:11:2 mix, 14-robot outage}}%
\end{pgfscope}%
\begin{pgfscope}%
\definecolor{textcolor}{rgb}{0.101961,0.101961,0.101961}%
\pgfsetstrokecolor{textcolor}%
\pgfsetfillcolor{textcolor}%
\pgftext[x=2.660350in,y=4.161115in,left,]{\makebox[0pt][r]{\color{textcolor}\sffamily\fontsize{7.000000}{8.400000}\selectfont rung 0 base}}%
\end{pgfscope}%
\begin{pgfscope}%
\definecolor{textcolor}{rgb}{0.101961,0.101961,0.101961}%
\pgfsetstrokecolor{textcolor}%
\pgfsetfillcolor{textcolor}%
\pgftext[x=2.660350in,y=4.002754in,left,]{\makebox[0pt][r]{\color{textcolor}\sffamily\fontsize{7.000000}{8.400000}\selectfont rung 1 rules}}%
\end{pgfscope}%
\begin{pgfscope}%
\definecolor{textcolor}{rgb}{0.101961,0.101961,0.101961}%
\pgfsetstrokecolor{textcolor}%
\pgfsetfillcolor{textcolor}%
\pgftext[x=2.660350in,y=3.844393in,left,]{\makebox[0pt][r]{\color{textcolor}\sffamily\fontsize{7.000000}{8.400000}\selectfont rung 2 state}}%
\end{pgfscope}%
\begin{pgfscope}%
\definecolor{textcolor}{rgb}{0.101961,0.101961,0.101961}%
\pgfsetstrokecolor{textcolor}%
\pgfsetfillcolor{textcolor}%
\pgftext[x=2.660350in,y=3.686033in,left,]{\makebox[0pt][r]{\color{textcolor}\sffamily\fontsize{7.000000}{8.400000}\selectfont rung 3 fair}}%
\end{pgfscope}%
\begin{pgfscope}%
\definecolor{textcolor}{rgb}{0.101961,0.101961,0.101961}%
\pgfsetstrokecolor{textcolor}%
\pgfsetfillcolor{textcolor}%
\pgftext[x=2.660350in,y=3.527672in,left,]{\makebox[0pt][r]{\color{textcolor}\sffamily\fontsize{7.000000}{8.400000}\selectfont rung 3 fair + rr}}%
\end{pgfscope}%
\begin{pgfscope}%
\definecolor{textcolor}{rgb}{0.101961,0.101961,0.101961}%
\pgfsetstrokecolor{textcolor}%
\pgfsetfillcolor{textcolor}%
\pgftext[x=2.660350in,y=3.369311in,left,]{\makebox[0pt][r]{\color{textcolor}\sffamily\fontsize{7.000000}{8.400000}\selectfont \textbf{rung 3 (COVEX, old)}}}%
\end{pgfscope}%
\begin{pgfscope}%
\definecolor{textcolor}{rgb}{0.101961,0.101961,0.101961}%
\pgfsetstrokecolor{textcolor}%
\pgfsetfillcolor{textcolor}%
\pgftext[x=2.660350in,y=3.210951in,left,]{\makebox[0pt][r]{\color{textcolor}\sffamily\fontsize{7.000000}{8.400000}\selectfont +0b, +P (resync)}}%
\end{pgfscope}%
\begin{pgfscope}%
\definecolor{textcolor}{rgb}{0.101961,0.101961,0.101961}%
\pgfsetstrokecolor{textcolor}%
\pgfsetfillcolor{textcolor}%
\pgftext[x=2.660350in,y=3.052590in,left,]{\makebox[0pt][r]{\color{textcolor}\sffamily\fontsize{7.000000}{8.400000}\selectfont +B}}%
\end{pgfscope}%
\begin{pgfscope}%
\definecolor{textcolor}{rgb}{0.101961,0.101961,0.101961}%
\pgfsetstrokecolor{textcolor}%
\pgfsetfillcolor{textcolor}%
\pgftext[x=2.660350in,y=2.894230in,left,]{\makebox[0pt][r]{\color{textcolor}\sffamily\fontsize{7.000000}{8.400000}\selectfont +P (wire, host model)}}%
\end{pgfscope}%
\begin{pgfscope}%
\definecolor{textcolor}{rgb}{0.101961,0.101961,0.101961}%
\pgfsetstrokecolor{textcolor}%
\pgfsetfillcolor{textcolor}%
\pgftext[x=2.660350in,y=2.735869in,left,]{\makebox[0pt][r]{\color{textcolor}\sffamily\fontsize{7.000000}{8.400000}\selectfont +question priority}}%
\end{pgfscope}%
\begin{pgfscope}%
\definecolor{textcolor}{rgb}{0.101961,0.101961,0.101961}%
\pgfsetstrokecolor{textcolor}%
\pgfsetfillcolor{textcolor}%
\pgftext[x=2.660350in,y=2.577508in,left,]{\makebox[0pt][r]{\color{textcolor}\sffamily\fontsize{7.000000}{8.400000}\selectfont \textbf{+latch fix (dev)}}}%
\end{pgfscope}%
\begin{pgfscope}%
\definecolor{textcolor}{rgb}{0.101961,0.101961,0.101961}%
\pgfsetstrokecolor{textcolor}%
\pgfsetfillcolor{textcolor}%
\pgftext[x=2.660350in,y=2.189525in,left,]{\makebox[0pt][r]{\color{textcolor}\sffamily\fontsize{7.000000}{8.400000}\selectfont \textbf{rung 3 (COVEX, old)}}}%
\end{pgfscope}%
\begin{pgfscope}%
\definecolor{textcolor}{rgb}{0.101961,0.101961,0.101961}%
\pgfsetstrokecolor{textcolor}%
\pgfsetfillcolor{textcolor}%
\pgftext[x=2.660350in,y=2.031164in,left,]{\makebox[0pt][r]{\color{textcolor}\sffamily\fontsize{7.000000}{8.400000}\selectfont \textbf{+latch fix (dev)}}}%
\end{pgfscope}%
\begin{pgfscope}%
\definecolor{textcolor}{rgb}{0.101961,0.101961,0.101961}%
\pgfsetstrokecolor{textcolor}%
\pgfsetfillcolor{textcolor}%
\pgftext[x=2.660350in,y=1.872803in,left,]{\makebox[0pt][r]{\color{textcolor}\sffamily\fontsize{7.000000}{8.400000}\selectfont counterfactual admission}}%
\end{pgfscope}%
\begin{pgfscope}%
\definecolor{textcolor}{rgb}{0.101961,0.101961,0.101961}%
\pgfsetstrokecolor{textcolor}%
\pgfsetfillcolor{textcolor}%
\pgftext[x=2.660350in,y=1.714443in,left,]{\makebox[0pt][r]{\color{textcolor}\sffamily\fontsize{7.000000}{8.400000}\selectfont counterfactual admission, admission only}}%
\end{pgfscope}%
\begin{pgfscope}%
\definecolor{textcolor}{rgb}{0.101961,0.101961,0.101961}%
\pgfsetstrokecolor{textcolor}%
\pgfsetfillcolor{textcolor}%
\pgftext[x=2.660350in,y=1.326459in,left,]{\makebox[0pt][r]{\color{textcolor}\sffamily\fontsize{7.000000}{8.400000}\selectfont \textbf{rung 3 (COVEX, old)}}}%
\end{pgfscope}%
\begin{pgfscope}%
\definecolor{textcolor}{rgb}{0.101961,0.101961,0.101961}%
\pgfsetstrokecolor{textcolor}%
\pgfsetfillcolor{textcolor}%
\pgftext[x=2.660350in,y=1.168098in,left,]{\makebox[0pt][r]{\color{textcolor}\sffamily\fontsize{7.000000}{8.400000}\selectfont \textbf{+latch fix (dev)}}}%
\end{pgfscope}%
\begin{pgfscope}%
\definecolor{textcolor}{rgb}{0.101961,0.101961,0.101961}%
\pgfsetstrokecolor{textcolor}%
\pgfsetfillcolor{textcolor}%
\pgftext[x=2.660350in,y=1.009738in,left,]{\makebox[0pt][r]{\color{textcolor}\sffamily\fontsize{7.000000}{8.400000}\selectfont \textbf{+S2 (dev)}}}%
\end{pgfscope}%
\begin{pgfscope}%
\definecolor{textcolor}{rgb}{0.101961,0.101961,0.101961}%
\pgfsetstrokecolor{textcolor}%
\pgfsetfillcolor{textcolor}%
\pgftext[x=2.660350in,y=0.851377in,left,]{\makebox[0pt][r]{\color{textcolor}\sffamily\fontsize{7.000000}{8.400000}\selectfont eager catch-up}}%
\end{pgfscope}%
\begin{pgfscope}%
\definecolor{textcolor}{rgb}{0.101961,0.101961,0.101961}%
\pgfsetstrokecolor{textcolor}%
\pgfsetfillcolor{textcolor}%
\pgftext[x=2.660350in,y=0.693016in,left,]{\makebox[0pt][r]{\color{textcolor}\sffamily\fontsize{7.000000}{8.400000}\selectfont counterfactual admission}}%
\end{pgfscope}%
\begin{pgfscope}%
\pgfsetbuttcap%
\pgfsetmiterjoin%
\definecolor{currentfill}{rgb}{1.000000,1.000000,1.000000}%
\pgfsetfillcolor{currentfill}%
\pgfsetlinewidth{0.000000pt}%
\definecolor{currentstroke}{rgb}{0.000000,0.000000,0.000000}%
\pgfsetstrokecolor{currentstroke}%
\pgfsetstrokeopacity{0.000000}%
\pgfsetdash{}{0pt}%
\pgfpathmoveto{\pgfqpoint{4.165000in}{0.598000in}}%
\pgfpathlineto{\pgfqpoint{5.320000in}{0.598000in}}%
\pgfpathlineto{\pgfqpoint{5.320000in}{4.462000in}}%
\pgfpathlineto{\pgfqpoint{4.165000in}{4.462000in}}%
\pgfpathlineto{\pgfqpoint{4.165000in}{0.598000in}}%
\pgfpathclose%
\pgfusepath{fill}%
\end{pgfscope}%
\begin{pgfscope}%
\pgfpathrectangle{\pgfqpoint{4.165000in}{0.598000in}}{\pgfqpoint{1.155000in}{3.864000in}}%
\pgfusepath{clip}%
\pgfsetrectcap%
\pgfsetroundjoin%
\pgfsetlinewidth{0.501875pt}%
\definecolor{currentstroke}{rgb}{0.901961,0.901961,0.901961}%
\pgfsetstrokecolor{currentstroke}%
\pgfsetdash{}{0pt}%
\pgfpathmoveto{\pgfqpoint{4.165000in}{4.398656in}}%
\pgfpathlineto{\pgfqpoint{5.320000in}{4.398656in}}%
\pgfusepath{stroke}%
\end{pgfscope}%
\begin{pgfscope}%
\pgfpathrectangle{\pgfqpoint{4.165000in}{0.598000in}}{\pgfqpoint{1.155000in}{3.864000in}}%
\pgfusepath{clip}%
\pgfsetrectcap%
\pgfsetroundjoin%
\pgfsetlinewidth{0.501875pt}%
\definecolor{currentstroke}{rgb}{0.901961,0.901961,0.901961}%
\pgfsetstrokecolor{currentstroke}%
\pgfsetdash{}{0pt}%
\pgfpathmoveto{\pgfqpoint{4.165000in}{2.427066in}}%
\pgfpathlineto{\pgfqpoint{5.320000in}{2.427066in}}%
\pgfusepath{stroke}%
\end{pgfscope}%
\begin{pgfscope}%
\pgfpathrectangle{\pgfqpoint{4.165000in}{0.598000in}}{\pgfqpoint{1.155000in}{3.864000in}}%
\pgfusepath{clip}%
\pgfsetrectcap%
\pgfsetroundjoin%
\pgfsetlinewidth{0.501875pt}%
\definecolor{currentstroke}{rgb}{0.901961,0.901961,0.901961}%
\pgfsetstrokecolor{currentstroke}%
\pgfsetdash{}{0pt}%
\pgfpathmoveto{\pgfqpoint{4.165000in}{1.564000in}}%
\pgfpathlineto{\pgfqpoint{5.320000in}{1.564000in}}%
\pgfusepath{stroke}%
\end{pgfscope}%
\begin{pgfscope}%
\pgfpathrectangle{\pgfqpoint{4.165000in}{0.598000in}}{\pgfqpoint{1.155000in}{3.864000in}}%
\pgfusepath{clip}%
\pgfsetrectcap%
\pgfsetroundjoin%
\pgfsetlinewidth{0.401500pt}%
\definecolor{currentstroke}{rgb}{0.901961,0.901961,0.901961}%
\pgfsetstrokecolor{currentstroke}%
\pgfsetdash{}{0pt}%
\pgfpathmoveto{\pgfqpoint{4.165000in}{0.598000in}}%
\pgfpathlineto{\pgfqpoint{4.165000in}{4.462000in}}%
\pgfusepath{stroke}%
\end{pgfscope}%
\begin{pgfscope}%
\definecolor{textcolor}{rgb}{0.101961,0.101961,0.101961}%
\pgfsetstrokecolor{textcolor}%
\pgfsetfillcolor{textcolor}%
\pgftext[x=4.165000in,y=0.549389in,,top]{\color{textcolor}\sffamily\fontsize{7.000000}{8.400000}\selectfont 0.4}%
\end{pgfscope}%
\begin{pgfscope}%
\pgfpathrectangle{\pgfqpoint{4.165000in}{0.598000in}}{\pgfqpoint{1.155000in}{3.864000in}}%
\pgfusepath{clip}%
\pgfsetrectcap%
\pgfsetroundjoin%
\pgfsetlinewidth{0.401500pt}%
\definecolor{currentstroke}{rgb}{0.901961,0.901961,0.901961}%
\pgfsetstrokecolor{currentstroke}%
\pgfsetdash{}{0pt}%
\pgfpathmoveto{\pgfqpoint{4.550000in}{0.598000in}}%
\pgfpathlineto{\pgfqpoint{4.550000in}{4.462000in}}%
\pgfusepath{stroke}%
\end{pgfscope}%
\begin{pgfscope}%
\definecolor{textcolor}{rgb}{0.101961,0.101961,0.101961}%
\pgfsetstrokecolor{textcolor}%
\pgfsetfillcolor{textcolor}%
\pgftext[x=4.550000in,y=0.549389in,,top]{\color{textcolor}\sffamily\fontsize{7.000000}{8.400000}\selectfont 0.5}%
\end{pgfscope}%
\begin{pgfscope}%
\pgfpathrectangle{\pgfqpoint{4.165000in}{0.598000in}}{\pgfqpoint{1.155000in}{3.864000in}}%
\pgfusepath{clip}%
\pgfsetrectcap%
\pgfsetroundjoin%
\pgfsetlinewidth{0.401500pt}%
\definecolor{currentstroke}{rgb}{0.901961,0.901961,0.901961}%
\pgfsetstrokecolor{currentstroke}%
\pgfsetdash{}{0pt}%
\pgfpathmoveto{\pgfqpoint{4.935000in}{0.598000in}}%
\pgfpathlineto{\pgfqpoint{4.935000in}{4.462000in}}%
\pgfusepath{stroke}%
\end{pgfscope}%
\begin{pgfscope}%
\definecolor{textcolor}{rgb}{0.101961,0.101961,0.101961}%
\pgfsetstrokecolor{textcolor}%
\pgfsetfillcolor{textcolor}%
\pgftext[x=4.935000in,y=0.549389in,,top]{\color{textcolor}\sffamily\fontsize{7.000000}{8.400000}\selectfont 0.6}%
\end{pgfscope}%
\begin{pgfscope}%
\definecolor{textcolor}{rgb}{0.000000,0.000000,0.000000}%
\pgfsetstrokecolor{textcolor}%
\pgfsetfillcolor{textcolor}%
\pgftext[x=4.505552in, y=0.367676in, left, base]{\color{textcolor}\sffamily\fontsize{7.500000}{9.000000}\selectfont Single-shot}%
\end{pgfscope}%
\begin{pgfscope}%
\definecolor{textcolor}{rgb}{0.000000,0.000000,0.000000}%
\pgfsetstrokecolor{textcolor}%
\pgfsetfillcolor{textcolor}%
\pgftext[x=4.426946in, y=0.267754in, left, base]{\color{textcolor}\sffamily\fontsize{7.500000}{9.000000}\selectfont correct in time}%
\end{pgfscope}%
\begin{pgfscope}%
\pgfpathrectangle{\pgfqpoint{4.165000in}{0.598000in}}{\pgfqpoint{1.155000in}{3.864000in}}%
\pgfusepath{clip}%
\pgfsetrectcap%
\pgfsetroundjoin%
\pgfsetlinewidth{0.602250pt}%
\definecolor{currentstroke}{rgb}{0.156863,0.356863,0.568627}%
\pgfsetstrokecolor{currentstroke}%
\pgfsetdash{}{0pt}%
\pgfpathmoveto{\pgfqpoint{4.922059in}{3.369311in}}%
\pgfpathlineto{\pgfqpoint{4.922059in}{3.210951in}}%
\pgfpathlineto{\pgfqpoint{4.922059in}{3.052590in}}%
\pgfpathlineto{\pgfqpoint{4.932843in}{2.894230in}}%
\pgfpathlineto{\pgfqpoint{4.932843in}{2.735869in}}%
\pgfpathlineto{\pgfqpoint{4.932843in}{2.577508in}}%
\pgfusepath{stroke}%
\end{pgfscope}%
\begin{pgfscope}%
\pgfsetrectcap%
\pgfsetmiterjoin%
\pgfsetlinewidth{0.803000pt}%
\definecolor{currentstroke}{rgb}{0.101961,0.101961,0.101961}%
\pgfsetstrokecolor{currentstroke}%
\pgfsetdash{}{0pt}%
\pgfpathmoveto{\pgfqpoint{4.165000in}{0.598000in}}%
\pgfpathlineto{\pgfqpoint{5.320000in}{0.598000in}}%
\pgfusepath{stroke}%
\end{pgfscope}%
\begin{pgfscope}%
\pgfpathrectangle{\pgfqpoint{4.165000in}{0.598000in}}{\pgfqpoint{1.155000in}{3.864000in}}%
\pgfusepath{clip}%
\pgfsetbuttcap%
\pgfsetroundjoin%
\definecolor{currentfill}{rgb}{0.435294,0.560784,0.709804}%
\pgfsetfillcolor{currentfill}%
\pgfsetlinewidth{0.000000pt}%
\definecolor{currentstroke}{rgb}{0.435294,0.560784,0.709804}%
\pgfsetstrokecolor{currentstroke}%
\pgfsetdash{}{0pt}%
\pgfsys@defobject{currentmarker}{\pgfqpoint{-0.024056in}{-0.024056in}}{\pgfqpoint{0.024056in}{0.024056in}}{%
\pgfpathmoveto{\pgfqpoint{0.000000in}{-0.024056in}}%
\pgfpathcurveto{\pgfqpoint{0.006380in}{-0.024056in}}{\pgfqpoint{0.012499in}{-0.021522in}}{\pgfqpoint{0.017010in}{-0.017010in}}%
\pgfpathcurveto{\pgfqpoint{0.021522in}{-0.012499in}}{\pgfqpoint{0.024056in}{-0.006380in}}{\pgfqpoint{0.024056in}{0.000000in}}%
\pgfpathcurveto{\pgfqpoint{0.024056in}{0.006380in}}{\pgfqpoint{0.021522in}{0.012499in}}{\pgfqpoint{0.017010in}{0.017010in}}%
\pgfpathcurveto{\pgfqpoint{0.012499in}{0.021522in}}{\pgfqpoint{0.006380in}{0.024056in}}{\pgfqpoint{0.000000in}{0.024056in}}%
\pgfpathcurveto{\pgfqpoint{-0.006380in}{0.024056in}}{\pgfqpoint{-0.012499in}{0.021522in}}{\pgfqpoint{-0.017010in}{0.017010in}}%
\pgfpathcurveto{\pgfqpoint{-0.021522in}{0.012499in}}{\pgfqpoint{-0.024056in}{0.006380in}}{\pgfqpoint{-0.024056in}{0.000000in}}%
\pgfpathcurveto{\pgfqpoint{-0.024056in}{-0.006380in}}{\pgfqpoint{-0.021522in}{-0.012499in}}{\pgfqpoint{-0.017010in}{-0.017010in}}%
\pgfpathcurveto{\pgfqpoint{-0.012499in}{-0.021522in}}{\pgfqpoint{-0.006380in}{-0.024056in}}{\pgfqpoint{0.000000in}{-0.024056in}}%
\pgfpathlineto{\pgfqpoint{0.000000in}{-0.024056in}}%
\pgfpathclose%
\pgfusepath{fill}%
}%
\begin{pgfscope}%
\pgfsys@transformshift{4.922059in}{4.161115in}%
\pgfsys@useobject{currentmarker}{}%
\end{pgfscope}%
\end{pgfscope}%
\begin{pgfscope}%
\definecolor{textcolor}{rgb}{0.549020,0.549020,0.549020}%
\pgfsetstrokecolor{textcolor}%
\pgfsetfillcolor{textcolor}%
\pgftext[x=4.962484in,y=4.161115in,left,]{\color{textcolor}\sffamily\fontsize{7.000000}{8.400000}\selectfont 0.60}%
\end{pgfscope}%
\begin{pgfscope}%
\pgfpathrectangle{\pgfqpoint{4.165000in}{0.598000in}}{\pgfqpoint{1.155000in}{3.864000in}}%
\pgfusepath{clip}%
\pgfsetbuttcap%
\pgfsetroundjoin%
\definecolor{currentfill}{rgb}{0.435294,0.560784,0.709804}%
\pgfsetfillcolor{currentfill}%
\pgfsetlinewidth{0.000000pt}%
\definecolor{currentstroke}{rgb}{0.435294,0.560784,0.709804}%
\pgfsetstrokecolor{currentstroke}%
\pgfsetdash{}{0pt}%
\pgfsys@defobject{currentmarker}{\pgfqpoint{-0.024056in}{-0.024056in}}{\pgfqpoint{0.024056in}{0.024056in}}{%
\pgfpathmoveto{\pgfqpoint{0.000000in}{-0.024056in}}%
\pgfpathcurveto{\pgfqpoint{0.006380in}{-0.024056in}}{\pgfqpoint{0.012499in}{-0.021522in}}{\pgfqpoint{0.017010in}{-0.017010in}}%
\pgfpathcurveto{\pgfqpoint{0.021522in}{-0.012499in}}{\pgfqpoint{0.024056in}{-0.006380in}}{\pgfqpoint{0.024056in}{0.000000in}}%
\pgfpathcurveto{\pgfqpoint{0.024056in}{0.006380in}}{\pgfqpoint{0.021522in}{0.012499in}}{\pgfqpoint{0.017010in}{0.017010in}}%
\pgfpathcurveto{\pgfqpoint{0.012499in}{0.021522in}}{\pgfqpoint{0.006380in}{0.024056in}}{\pgfqpoint{0.000000in}{0.024056in}}%
\pgfpathcurveto{\pgfqpoint{-0.006380in}{0.024056in}}{\pgfqpoint{-0.012499in}{0.021522in}}{\pgfqpoint{-0.017010in}{0.017010in}}%
\pgfpathcurveto{\pgfqpoint{-0.021522in}{0.012499in}}{\pgfqpoint{-0.024056in}{0.006380in}}{\pgfqpoint{-0.024056in}{0.000000in}}%
\pgfpathcurveto{\pgfqpoint{-0.024056in}{-0.006380in}}{\pgfqpoint{-0.021522in}{-0.012499in}}{\pgfqpoint{-0.017010in}{-0.017010in}}%
\pgfpathcurveto{\pgfqpoint{-0.012499in}{-0.021522in}}{\pgfqpoint{-0.006380in}{-0.024056in}}{\pgfqpoint{0.000000in}{-0.024056in}}%
\pgfpathlineto{\pgfqpoint{0.000000in}{-0.024056in}}%
\pgfpathclose%
\pgfusepath{fill}%
}%
\begin{pgfscope}%
\pgfsys@transformshift{4.846569in}{4.002754in}%
\pgfsys@useobject{currentmarker}{}%
\end{pgfscope}%
\end{pgfscope}%
\begin{pgfscope}%
\definecolor{textcolor}{rgb}{0.549020,0.549020,0.549020}%
\pgfsetstrokecolor{textcolor}%
\pgfsetfillcolor{textcolor}%
\pgftext[x=4.886994in,y=4.002754in,left,]{\color{textcolor}\sffamily\fontsize{7.000000}{8.400000}\selectfont 0.58}%
\end{pgfscope}%
\begin{pgfscope}%
\pgfpathrectangle{\pgfqpoint{4.165000in}{0.598000in}}{\pgfqpoint{1.155000in}{3.864000in}}%
\pgfusepath{clip}%
\pgfsetbuttcap%
\pgfsetroundjoin%
\definecolor{currentfill}{rgb}{0.435294,0.560784,0.709804}%
\pgfsetfillcolor{currentfill}%
\pgfsetlinewidth{0.000000pt}%
\definecolor{currentstroke}{rgb}{0.435294,0.560784,0.709804}%
\pgfsetstrokecolor{currentstroke}%
\pgfsetdash{}{0pt}%
\pgfsys@defobject{currentmarker}{\pgfqpoint{-0.024056in}{-0.024056in}}{\pgfqpoint{0.024056in}{0.024056in}}{%
\pgfpathmoveto{\pgfqpoint{0.000000in}{-0.024056in}}%
\pgfpathcurveto{\pgfqpoint{0.006380in}{-0.024056in}}{\pgfqpoint{0.012499in}{-0.021522in}}{\pgfqpoint{0.017010in}{-0.017010in}}%
\pgfpathcurveto{\pgfqpoint{0.021522in}{-0.012499in}}{\pgfqpoint{0.024056in}{-0.006380in}}{\pgfqpoint{0.024056in}{0.000000in}}%
\pgfpathcurveto{\pgfqpoint{0.024056in}{0.006380in}}{\pgfqpoint{0.021522in}{0.012499in}}{\pgfqpoint{0.017010in}{0.017010in}}%
\pgfpathcurveto{\pgfqpoint{0.012499in}{0.021522in}}{\pgfqpoint{0.006380in}{0.024056in}}{\pgfqpoint{0.000000in}{0.024056in}}%
\pgfpathcurveto{\pgfqpoint{-0.006380in}{0.024056in}}{\pgfqpoint{-0.012499in}{0.021522in}}{\pgfqpoint{-0.017010in}{0.017010in}}%
\pgfpathcurveto{\pgfqpoint{-0.021522in}{0.012499in}}{\pgfqpoint{-0.024056in}{0.006380in}}{\pgfqpoint{-0.024056in}{0.000000in}}%
\pgfpathcurveto{\pgfqpoint{-0.024056in}{-0.006380in}}{\pgfqpoint{-0.021522in}{-0.012499in}}{\pgfqpoint{-0.017010in}{-0.017010in}}%
\pgfpathcurveto{\pgfqpoint{-0.012499in}{-0.021522in}}{\pgfqpoint{-0.006380in}{-0.024056in}}{\pgfqpoint{0.000000in}{-0.024056in}}%
\pgfpathlineto{\pgfqpoint{0.000000in}{-0.024056in}}%
\pgfpathclose%
\pgfusepath{fill}%
}%
\begin{pgfscope}%
\pgfsys@transformshift{4.911275in}{3.844393in}%
\pgfsys@useobject{currentmarker}{}%
\end{pgfscope}%
\end{pgfscope}%
\begin{pgfscope}%
\definecolor{textcolor}{rgb}{0.549020,0.549020,0.549020}%
\pgfsetstrokecolor{textcolor}%
\pgfsetfillcolor{textcolor}%
\pgftext[x=4.951700in,y=3.844393in,left,]{\color{textcolor}\sffamily\fontsize{7.000000}{8.400000}\selectfont 0.59}%
\end{pgfscope}%
\begin{pgfscope}%
\pgfpathrectangle{\pgfqpoint{4.165000in}{0.598000in}}{\pgfqpoint{1.155000in}{3.864000in}}%
\pgfusepath{clip}%
\pgfsetbuttcap%
\pgfsetroundjoin%
\definecolor{currentfill}{rgb}{0.435294,0.560784,0.709804}%
\pgfsetfillcolor{currentfill}%
\pgfsetlinewidth{0.000000pt}%
\definecolor{currentstroke}{rgb}{0.435294,0.560784,0.709804}%
\pgfsetstrokecolor{currentstroke}%
\pgfsetdash{}{0pt}%
\pgfsys@defobject{currentmarker}{\pgfqpoint{-0.024056in}{-0.024056in}}{\pgfqpoint{0.024056in}{0.024056in}}{%
\pgfpathmoveto{\pgfqpoint{0.000000in}{-0.024056in}}%
\pgfpathcurveto{\pgfqpoint{0.006380in}{-0.024056in}}{\pgfqpoint{0.012499in}{-0.021522in}}{\pgfqpoint{0.017010in}{-0.017010in}}%
\pgfpathcurveto{\pgfqpoint{0.021522in}{-0.012499in}}{\pgfqpoint{0.024056in}{-0.006380in}}{\pgfqpoint{0.024056in}{0.000000in}}%
\pgfpathcurveto{\pgfqpoint{0.024056in}{0.006380in}}{\pgfqpoint{0.021522in}{0.012499in}}{\pgfqpoint{0.017010in}{0.017010in}}%
\pgfpathcurveto{\pgfqpoint{0.012499in}{0.021522in}}{\pgfqpoint{0.006380in}{0.024056in}}{\pgfqpoint{0.000000in}{0.024056in}}%
\pgfpathcurveto{\pgfqpoint{-0.006380in}{0.024056in}}{\pgfqpoint{-0.012499in}{0.021522in}}{\pgfqpoint{-0.017010in}{0.017010in}}%
\pgfpathcurveto{\pgfqpoint{-0.021522in}{0.012499in}}{\pgfqpoint{-0.024056in}{0.006380in}}{\pgfqpoint{-0.024056in}{0.000000in}}%
\pgfpathcurveto{\pgfqpoint{-0.024056in}{-0.006380in}}{\pgfqpoint{-0.021522in}{-0.012499in}}{\pgfqpoint{-0.017010in}{-0.017010in}}%
\pgfpathcurveto{\pgfqpoint{-0.012499in}{-0.021522in}}{\pgfqpoint{-0.006380in}{-0.024056in}}{\pgfqpoint{0.000000in}{-0.024056in}}%
\pgfpathlineto{\pgfqpoint{0.000000in}{-0.024056in}}%
\pgfpathclose%
\pgfusepath{fill}%
}%
\begin{pgfscope}%
\pgfsys@transformshift{4.922059in}{3.686033in}%
\pgfsys@useobject{currentmarker}{}%
\end{pgfscope}%
\end{pgfscope}%
\begin{pgfscope}%
\definecolor{textcolor}{rgb}{0.549020,0.549020,0.549020}%
\pgfsetstrokecolor{textcolor}%
\pgfsetfillcolor{textcolor}%
\pgftext[x=4.962484in,y=3.686033in,left,]{\color{textcolor}\sffamily\fontsize{7.000000}{8.400000}\selectfont 0.60}%
\end{pgfscope}%
\begin{pgfscope}%
\pgfpathrectangle{\pgfqpoint{4.165000in}{0.598000in}}{\pgfqpoint{1.155000in}{3.864000in}}%
\pgfusepath{clip}%
\pgfsetbuttcap%
\pgfsetroundjoin%
\definecolor{currentfill}{rgb}{0.435294,0.560784,0.709804}%
\pgfsetfillcolor{currentfill}%
\pgfsetlinewidth{0.000000pt}%
\definecolor{currentstroke}{rgb}{0.435294,0.560784,0.709804}%
\pgfsetstrokecolor{currentstroke}%
\pgfsetdash{}{0pt}%
\pgfsys@defobject{currentmarker}{\pgfqpoint{-0.024056in}{-0.024056in}}{\pgfqpoint{0.024056in}{0.024056in}}{%
\pgfpathmoveto{\pgfqpoint{0.000000in}{-0.024056in}}%
\pgfpathcurveto{\pgfqpoint{0.006380in}{-0.024056in}}{\pgfqpoint{0.012499in}{-0.021522in}}{\pgfqpoint{0.017010in}{-0.017010in}}%
\pgfpathcurveto{\pgfqpoint{0.021522in}{-0.012499in}}{\pgfqpoint{0.024056in}{-0.006380in}}{\pgfqpoint{0.024056in}{0.000000in}}%
\pgfpathcurveto{\pgfqpoint{0.024056in}{0.006380in}}{\pgfqpoint{0.021522in}{0.012499in}}{\pgfqpoint{0.017010in}{0.017010in}}%
\pgfpathcurveto{\pgfqpoint{0.012499in}{0.021522in}}{\pgfqpoint{0.006380in}{0.024056in}}{\pgfqpoint{0.000000in}{0.024056in}}%
\pgfpathcurveto{\pgfqpoint{-0.006380in}{0.024056in}}{\pgfqpoint{-0.012499in}{0.021522in}}{\pgfqpoint{-0.017010in}{0.017010in}}%
\pgfpathcurveto{\pgfqpoint{-0.021522in}{0.012499in}}{\pgfqpoint{-0.024056in}{0.006380in}}{\pgfqpoint{-0.024056in}{0.000000in}}%
\pgfpathcurveto{\pgfqpoint{-0.024056in}{-0.006380in}}{\pgfqpoint{-0.021522in}{-0.012499in}}{\pgfqpoint{-0.017010in}{-0.017010in}}%
\pgfpathcurveto{\pgfqpoint{-0.012499in}{-0.021522in}}{\pgfqpoint{-0.006380in}{-0.024056in}}{\pgfqpoint{0.000000in}{-0.024056in}}%
\pgfpathlineto{\pgfqpoint{0.000000in}{-0.024056in}}%
\pgfpathclose%
\pgfusepath{fill}%
}%
\begin{pgfscope}%
\pgfsys@transformshift{4.932843in}{3.527672in}%
\pgfsys@useobject{currentmarker}{}%
\end{pgfscope}%
\end{pgfscope}%
\begin{pgfscope}%
\definecolor{textcolor}{rgb}{0.549020,0.549020,0.549020}%
\pgfsetstrokecolor{textcolor}%
\pgfsetfillcolor{textcolor}%
\pgftext[x=4.973268in,y=3.527672in,left,]{\color{textcolor}\sffamily\fontsize{7.000000}{8.400000}\selectfont 0.60}%
\end{pgfscope}%
\begin{pgfscope}%
\pgfpathrectangle{\pgfqpoint{4.165000in}{0.598000in}}{\pgfqpoint{1.155000in}{3.864000in}}%
\pgfusepath{clip}%
\pgfsetbuttcap%
\pgfsetroundjoin%
\definecolor{currentfill}{rgb}{0.156863,0.356863,0.568627}%
\pgfsetfillcolor{currentfill}%
\pgfsetlinewidth{0.000000pt}%
\definecolor{currentstroke}{rgb}{0.156863,0.356863,0.568627}%
\pgfsetstrokecolor{currentstroke}%
\pgfsetdash{}{0pt}%
\pgfsys@defobject{currentmarker}{\pgfqpoint{-0.024056in}{-0.024056in}}{\pgfqpoint{0.024056in}{0.024056in}}{%
\pgfpathmoveto{\pgfqpoint{0.000000in}{-0.024056in}}%
\pgfpathcurveto{\pgfqpoint{0.006380in}{-0.024056in}}{\pgfqpoint{0.012499in}{-0.021522in}}{\pgfqpoint{0.017010in}{-0.017010in}}%
\pgfpathcurveto{\pgfqpoint{0.021522in}{-0.012499in}}{\pgfqpoint{0.024056in}{-0.006380in}}{\pgfqpoint{0.024056in}{0.000000in}}%
\pgfpathcurveto{\pgfqpoint{0.024056in}{0.006380in}}{\pgfqpoint{0.021522in}{0.012499in}}{\pgfqpoint{0.017010in}{0.017010in}}%
\pgfpathcurveto{\pgfqpoint{0.012499in}{0.021522in}}{\pgfqpoint{0.006380in}{0.024056in}}{\pgfqpoint{0.000000in}{0.024056in}}%
\pgfpathcurveto{\pgfqpoint{-0.006380in}{0.024056in}}{\pgfqpoint{-0.012499in}{0.021522in}}{\pgfqpoint{-0.017010in}{0.017010in}}%
\pgfpathcurveto{\pgfqpoint{-0.021522in}{0.012499in}}{\pgfqpoint{-0.024056in}{0.006380in}}{\pgfqpoint{-0.024056in}{0.000000in}}%
\pgfpathcurveto{\pgfqpoint{-0.024056in}{-0.006380in}}{\pgfqpoint{-0.021522in}{-0.012499in}}{\pgfqpoint{-0.017010in}{-0.017010in}}%
\pgfpathcurveto{\pgfqpoint{-0.012499in}{-0.021522in}}{\pgfqpoint{-0.006380in}{-0.024056in}}{\pgfqpoint{0.000000in}{-0.024056in}}%
\pgfpathlineto{\pgfqpoint{0.000000in}{-0.024056in}}%
\pgfpathclose%
\pgfusepath{fill}%
}%
\begin{pgfscope}%
\pgfsys@transformshift{4.922059in}{3.369311in}%
\pgfsys@useobject{currentmarker}{}%
\end{pgfscope}%
\end{pgfscope}%
\begin{pgfscope}%
\definecolor{textcolor}{rgb}{0.549020,0.549020,0.549020}%
\pgfsetstrokecolor{textcolor}%
\pgfsetfillcolor{textcolor}%
\pgftext[x=4.962484in,y=3.369311in,left,]{\color{textcolor}\sffamily\fontsize{7.000000}{8.400000}\selectfont 0.60}%
\end{pgfscope}%
\begin{pgfscope}%
\pgfpathrectangle{\pgfqpoint{4.165000in}{0.598000in}}{\pgfqpoint{1.155000in}{3.864000in}}%
\pgfusepath{clip}%
\pgfsetbuttcap%
\pgfsetroundjoin%
\definecolor{currentfill}{rgb}{0.435294,0.560784,0.709804}%
\pgfsetfillcolor{currentfill}%
\pgfsetlinewidth{0.000000pt}%
\definecolor{currentstroke}{rgb}{0.435294,0.560784,0.709804}%
\pgfsetstrokecolor{currentstroke}%
\pgfsetdash{}{0pt}%
\pgfsys@defobject{currentmarker}{\pgfqpoint{-0.024056in}{-0.024056in}}{\pgfqpoint{0.024056in}{0.024056in}}{%
\pgfpathmoveto{\pgfqpoint{0.000000in}{-0.024056in}}%
\pgfpathcurveto{\pgfqpoint{0.006380in}{-0.024056in}}{\pgfqpoint{0.012499in}{-0.021522in}}{\pgfqpoint{0.017010in}{-0.017010in}}%
\pgfpathcurveto{\pgfqpoint{0.021522in}{-0.012499in}}{\pgfqpoint{0.024056in}{-0.006380in}}{\pgfqpoint{0.024056in}{0.000000in}}%
\pgfpathcurveto{\pgfqpoint{0.024056in}{0.006380in}}{\pgfqpoint{0.021522in}{0.012499in}}{\pgfqpoint{0.017010in}{0.017010in}}%
\pgfpathcurveto{\pgfqpoint{0.012499in}{0.021522in}}{\pgfqpoint{0.006380in}{0.024056in}}{\pgfqpoint{0.000000in}{0.024056in}}%
\pgfpathcurveto{\pgfqpoint{-0.006380in}{0.024056in}}{\pgfqpoint{-0.012499in}{0.021522in}}{\pgfqpoint{-0.017010in}{0.017010in}}%
\pgfpathcurveto{\pgfqpoint{-0.021522in}{0.012499in}}{\pgfqpoint{-0.024056in}{0.006380in}}{\pgfqpoint{-0.024056in}{0.000000in}}%
\pgfpathcurveto{\pgfqpoint{-0.024056in}{-0.006380in}}{\pgfqpoint{-0.021522in}{-0.012499in}}{\pgfqpoint{-0.017010in}{-0.017010in}}%
\pgfpathcurveto{\pgfqpoint{-0.012499in}{-0.021522in}}{\pgfqpoint{-0.006380in}{-0.024056in}}{\pgfqpoint{0.000000in}{-0.024056in}}%
\pgfpathlineto{\pgfqpoint{0.000000in}{-0.024056in}}%
\pgfpathclose%
\pgfusepath{fill}%
}%
\begin{pgfscope}%
\pgfsys@transformshift{4.922059in}{3.210951in}%
\pgfsys@useobject{currentmarker}{}%
\end{pgfscope}%
\end{pgfscope}%
\begin{pgfscope}%
\definecolor{textcolor}{rgb}{0.549020,0.549020,0.549020}%
\pgfsetstrokecolor{textcolor}%
\pgfsetfillcolor{textcolor}%
\pgftext[x=4.962484in,y=3.210951in,left,]{\color{textcolor}\sffamily\fontsize{7.000000}{8.400000}\selectfont 0.60}%
\end{pgfscope}%
\begin{pgfscope}%
\pgfpathrectangle{\pgfqpoint{4.165000in}{0.598000in}}{\pgfqpoint{1.155000in}{3.864000in}}%
\pgfusepath{clip}%
\pgfsetbuttcap%
\pgfsetroundjoin%
\definecolor{currentfill}{rgb}{0.435294,0.560784,0.709804}%
\pgfsetfillcolor{currentfill}%
\pgfsetlinewidth{0.000000pt}%
\definecolor{currentstroke}{rgb}{0.435294,0.560784,0.709804}%
\pgfsetstrokecolor{currentstroke}%
\pgfsetdash{}{0pt}%
\pgfsys@defobject{currentmarker}{\pgfqpoint{-0.024056in}{-0.024056in}}{\pgfqpoint{0.024056in}{0.024056in}}{%
\pgfpathmoveto{\pgfqpoint{0.000000in}{-0.024056in}}%
\pgfpathcurveto{\pgfqpoint{0.006380in}{-0.024056in}}{\pgfqpoint{0.012499in}{-0.021522in}}{\pgfqpoint{0.017010in}{-0.017010in}}%
\pgfpathcurveto{\pgfqpoint{0.021522in}{-0.012499in}}{\pgfqpoint{0.024056in}{-0.006380in}}{\pgfqpoint{0.024056in}{0.000000in}}%
\pgfpathcurveto{\pgfqpoint{0.024056in}{0.006380in}}{\pgfqpoint{0.021522in}{0.012499in}}{\pgfqpoint{0.017010in}{0.017010in}}%
\pgfpathcurveto{\pgfqpoint{0.012499in}{0.021522in}}{\pgfqpoint{0.006380in}{0.024056in}}{\pgfqpoint{0.000000in}{0.024056in}}%
\pgfpathcurveto{\pgfqpoint{-0.006380in}{0.024056in}}{\pgfqpoint{-0.012499in}{0.021522in}}{\pgfqpoint{-0.017010in}{0.017010in}}%
\pgfpathcurveto{\pgfqpoint{-0.021522in}{0.012499in}}{\pgfqpoint{-0.024056in}{0.006380in}}{\pgfqpoint{-0.024056in}{0.000000in}}%
\pgfpathcurveto{\pgfqpoint{-0.024056in}{-0.006380in}}{\pgfqpoint{-0.021522in}{-0.012499in}}{\pgfqpoint{-0.017010in}{-0.017010in}}%
\pgfpathcurveto{\pgfqpoint{-0.012499in}{-0.021522in}}{\pgfqpoint{-0.006380in}{-0.024056in}}{\pgfqpoint{0.000000in}{-0.024056in}}%
\pgfpathlineto{\pgfqpoint{0.000000in}{-0.024056in}}%
\pgfpathclose%
\pgfusepath{fill}%
}%
\begin{pgfscope}%
\pgfsys@transformshift{4.922059in}{3.052590in}%
\pgfsys@useobject{currentmarker}{}%
\end{pgfscope}%
\end{pgfscope}%
\begin{pgfscope}%
\definecolor{textcolor}{rgb}{0.549020,0.549020,0.549020}%
\pgfsetstrokecolor{textcolor}%
\pgfsetfillcolor{textcolor}%
\pgftext[x=4.962484in,y=3.052590in,left,]{\color{textcolor}\sffamily\fontsize{7.000000}{8.400000}\selectfont 0.60}%
\end{pgfscope}%
\begin{pgfscope}%
\pgfpathrectangle{\pgfqpoint{4.165000in}{0.598000in}}{\pgfqpoint{1.155000in}{3.864000in}}%
\pgfusepath{clip}%
\pgfsetbuttcap%
\pgfsetroundjoin%
\definecolor{currentfill}{rgb}{0.435294,0.560784,0.709804}%
\pgfsetfillcolor{currentfill}%
\pgfsetlinewidth{0.000000pt}%
\definecolor{currentstroke}{rgb}{0.435294,0.560784,0.709804}%
\pgfsetstrokecolor{currentstroke}%
\pgfsetdash{}{0pt}%
\pgfsys@defobject{currentmarker}{\pgfqpoint{-0.024056in}{-0.024056in}}{\pgfqpoint{0.024056in}{0.024056in}}{%
\pgfpathmoveto{\pgfqpoint{0.000000in}{-0.024056in}}%
\pgfpathcurveto{\pgfqpoint{0.006380in}{-0.024056in}}{\pgfqpoint{0.012499in}{-0.021522in}}{\pgfqpoint{0.017010in}{-0.017010in}}%
\pgfpathcurveto{\pgfqpoint{0.021522in}{-0.012499in}}{\pgfqpoint{0.024056in}{-0.006380in}}{\pgfqpoint{0.024056in}{0.000000in}}%
\pgfpathcurveto{\pgfqpoint{0.024056in}{0.006380in}}{\pgfqpoint{0.021522in}{0.012499in}}{\pgfqpoint{0.017010in}{0.017010in}}%
\pgfpathcurveto{\pgfqpoint{0.012499in}{0.021522in}}{\pgfqpoint{0.006380in}{0.024056in}}{\pgfqpoint{0.000000in}{0.024056in}}%
\pgfpathcurveto{\pgfqpoint{-0.006380in}{0.024056in}}{\pgfqpoint{-0.012499in}{0.021522in}}{\pgfqpoint{-0.017010in}{0.017010in}}%
\pgfpathcurveto{\pgfqpoint{-0.021522in}{0.012499in}}{\pgfqpoint{-0.024056in}{0.006380in}}{\pgfqpoint{-0.024056in}{0.000000in}}%
\pgfpathcurveto{\pgfqpoint{-0.024056in}{-0.006380in}}{\pgfqpoint{-0.021522in}{-0.012499in}}{\pgfqpoint{-0.017010in}{-0.017010in}}%
\pgfpathcurveto{\pgfqpoint{-0.012499in}{-0.021522in}}{\pgfqpoint{-0.006380in}{-0.024056in}}{\pgfqpoint{0.000000in}{-0.024056in}}%
\pgfpathlineto{\pgfqpoint{0.000000in}{-0.024056in}}%
\pgfpathclose%
\pgfusepath{fill}%
}%
\begin{pgfscope}%
\pgfsys@transformshift{4.932843in}{2.894230in}%
\pgfsys@useobject{currentmarker}{}%
\end{pgfscope}%
\end{pgfscope}%
\begin{pgfscope}%
\definecolor{textcolor}{rgb}{0.549020,0.549020,0.549020}%
\pgfsetstrokecolor{textcolor}%
\pgfsetfillcolor{textcolor}%
\pgftext[x=4.973268in,y=2.894230in,left,]{\color{textcolor}\sffamily\fontsize{7.000000}{8.400000}\selectfont 0.60}%
\end{pgfscope}%
\begin{pgfscope}%
\pgfpathrectangle{\pgfqpoint{4.165000in}{0.598000in}}{\pgfqpoint{1.155000in}{3.864000in}}%
\pgfusepath{clip}%
\pgfsetbuttcap%
\pgfsetroundjoin%
\definecolor{currentfill}{rgb}{0.435294,0.560784,0.709804}%
\pgfsetfillcolor{currentfill}%
\pgfsetlinewidth{0.000000pt}%
\definecolor{currentstroke}{rgb}{0.435294,0.560784,0.709804}%
\pgfsetstrokecolor{currentstroke}%
\pgfsetdash{}{0pt}%
\pgfsys@defobject{currentmarker}{\pgfqpoint{-0.024056in}{-0.024056in}}{\pgfqpoint{0.024056in}{0.024056in}}{%
\pgfpathmoveto{\pgfqpoint{0.000000in}{-0.024056in}}%
\pgfpathcurveto{\pgfqpoint{0.006380in}{-0.024056in}}{\pgfqpoint{0.012499in}{-0.021522in}}{\pgfqpoint{0.017010in}{-0.017010in}}%
\pgfpathcurveto{\pgfqpoint{0.021522in}{-0.012499in}}{\pgfqpoint{0.024056in}{-0.006380in}}{\pgfqpoint{0.024056in}{0.000000in}}%
\pgfpathcurveto{\pgfqpoint{0.024056in}{0.006380in}}{\pgfqpoint{0.021522in}{0.012499in}}{\pgfqpoint{0.017010in}{0.017010in}}%
\pgfpathcurveto{\pgfqpoint{0.012499in}{0.021522in}}{\pgfqpoint{0.006380in}{0.024056in}}{\pgfqpoint{0.000000in}{0.024056in}}%
\pgfpathcurveto{\pgfqpoint{-0.006380in}{0.024056in}}{\pgfqpoint{-0.012499in}{0.021522in}}{\pgfqpoint{-0.017010in}{0.017010in}}%
\pgfpathcurveto{\pgfqpoint{-0.021522in}{0.012499in}}{\pgfqpoint{-0.024056in}{0.006380in}}{\pgfqpoint{-0.024056in}{0.000000in}}%
\pgfpathcurveto{\pgfqpoint{-0.024056in}{-0.006380in}}{\pgfqpoint{-0.021522in}{-0.012499in}}{\pgfqpoint{-0.017010in}{-0.017010in}}%
\pgfpathcurveto{\pgfqpoint{-0.012499in}{-0.021522in}}{\pgfqpoint{-0.006380in}{-0.024056in}}{\pgfqpoint{0.000000in}{-0.024056in}}%
\pgfpathlineto{\pgfqpoint{0.000000in}{-0.024056in}}%
\pgfpathclose%
\pgfusepath{fill}%
}%
\begin{pgfscope}%
\pgfsys@transformshift{4.932843in}{2.735869in}%
\pgfsys@useobject{currentmarker}{}%
\end{pgfscope}%
\end{pgfscope}%
\begin{pgfscope}%
\definecolor{textcolor}{rgb}{0.549020,0.549020,0.549020}%
\pgfsetstrokecolor{textcolor}%
\pgfsetfillcolor{textcolor}%
\pgftext[x=4.973268in,y=2.735869in,left,]{\color{textcolor}\sffamily\fontsize{7.000000}{8.400000}\selectfont 0.60}%
\end{pgfscope}%
\begin{pgfscope}%
\pgfpathrectangle{\pgfqpoint{4.165000in}{0.598000in}}{\pgfqpoint{1.155000in}{3.864000in}}%
\pgfusepath{clip}%
\pgfsetbuttcap%
\pgfsetroundjoin%
\definecolor{currentfill}{rgb}{0.156863,0.356863,0.568627}%
\pgfsetfillcolor{currentfill}%
\pgfsetlinewidth{0.000000pt}%
\definecolor{currentstroke}{rgb}{0.156863,0.356863,0.568627}%
\pgfsetstrokecolor{currentstroke}%
\pgfsetdash{}{0pt}%
\pgfsys@defobject{currentmarker}{\pgfqpoint{-0.024056in}{-0.024056in}}{\pgfqpoint{0.024056in}{0.024056in}}{%
\pgfpathmoveto{\pgfqpoint{0.000000in}{-0.024056in}}%
\pgfpathcurveto{\pgfqpoint{0.006380in}{-0.024056in}}{\pgfqpoint{0.012499in}{-0.021522in}}{\pgfqpoint{0.017010in}{-0.017010in}}%
\pgfpathcurveto{\pgfqpoint{0.021522in}{-0.012499in}}{\pgfqpoint{0.024056in}{-0.006380in}}{\pgfqpoint{0.024056in}{0.000000in}}%
\pgfpathcurveto{\pgfqpoint{0.024056in}{0.006380in}}{\pgfqpoint{0.021522in}{0.012499in}}{\pgfqpoint{0.017010in}{0.017010in}}%
\pgfpathcurveto{\pgfqpoint{0.012499in}{0.021522in}}{\pgfqpoint{0.006380in}{0.024056in}}{\pgfqpoint{0.000000in}{0.024056in}}%
\pgfpathcurveto{\pgfqpoint{-0.006380in}{0.024056in}}{\pgfqpoint{-0.012499in}{0.021522in}}{\pgfqpoint{-0.017010in}{0.017010in}}%
\pgfpathcurveto{\pgfqpoint{-0.021522in}{0.012499in}}{\pgfqpoint{-0.024056in}{0.006380in}}{\pgfqpoint{-0.024056in}{0.000000in}}%
\pgfpathcurveto{\pgfqpoint{-0.024056in}{-0.006380in}}{\pgfqpoint{-0.021522in}{-0.012499in}}{\pgfqpoint{-0.017010in}{-0.017010in}}%
\pgfpathcurveto{\pgfqpoint{-0.012499in}{-0.021522in}}{\pgfqpoint{-0.006380in}{-0.024056in}}{\pgfqpoint{0.000000in}{-0.024056in}}%
\pgfpathlineto{\pgfqpoint{0.000000in}{-0.024056in}}%
\pgfpathclose%
\pgfusepath{fill}%
}%
\begin{pgfscope}%
\pgfsys@transformshift{4.932843in}{2.577508in}%
\pgfsys@useobject{currentmarker}{}%
\end{pgfscope}%
\end{pgfscope}%
\begin{pgfscope}%
\definecolor{textcolor}{rgb}{0.549020,0.549020,0.549020}%
\pgfsetstrokecolor{textcolor}%
\pgfsetfillcolor{textcolor}%
\pgftext[x=4.973268in,y=2.577508in,left,]{\color{textcolor}\sffamily\fontsize{7.000000}{8.400000}\selectfont 0.60}%
\end{pgfscope}%
\begin{pgfscope}%
\pgfpathrectangle{\pgfqpoint{4.165000in}{0.598000in}}{\pgfqpoint{1.155000in}{3.864000in}}%
\pgfusepath{clip}%
\pgfsetbuttcap%
\pgfsetroundjoin%
\definecolor{currentfill}{rgb}{0.156863,0.356863,0.568627}%
\pgfsetfillcolor{currentfill}%
\pgfsetlinewidth{0.000000pt}%
\definecolor{currentstroke}{rgb}{0.156863,0.356863,0.568627}%
\pgfsetstrokecolor{currentstroke}%
\pgfsetdash{}{0pt}%
\pgfsys@defobject{currentmarker}{\pgfqpoint{-0.024056in}{-0.024056in}}{\pgfqpoint{0.024056in}{0.024056in}}{%
\pgfpathmoveto{\pgfqpoint{0.000000in}{-0.024056in}}%
\pgfpathcurveto{\pgfqpoint{0.006380in}{-0.024056in}}{\pgfqpoint{0.012499in}{-0.021522in}}{\pgfqpoint{0.017010in}{-0.017010in}}%
\pgfpathcurveto{\pgfqpoint{0.021522in}{-0.012499in}}{\pgfqpoint{0.024056in}{-0.006380in}}{\pgfqpoint{0.024056in}{0.000000in}}%
\pgfpathcurveto{\pgfqpoint{0.024056in}{0.006380in}}{\pgfqpoint{0.021522in}{0.012499in}}{\pgfqpoint{0.017010in}{0.017010in}}%
\pgfpathcurveto{\pgfqpoint{0.012499in}{0.021522in}}{\pgfqpoint{0.006380in}{0.024056in}}{\pgfqpoint{0.000000in}{0.024056in}}%
\pgfpathcurveto{\pgfqpoint{-0.006380in}{0.024056in}}{\pgfqpoint{-0.012499in}{0.021522in}}{\pgfqpoint{-0.017010in}{0.017010in}}%
\pgfpathcurveto{\pgfqpoint{-0.021522in}{0.012499in}}{\pgfqpoint{-0.024056in}{0.006380in}}{\pgfqpoint{-0.024056in}{0.000000in}}%
\pgfpathcurveto{\pgfqpoint{-0.024056in}{-0.006380in}}{\pgfqpoint{-0.021522in}{-0.012499in}}{\pgfqpoint{-0.017010in}{-0.017010in}}%
\pgfpathcurveto{\pgfqpoint{-0.012499in}{-0.021522in}}{\pgfqpoint{-0.006380in}{-0.024056in}}{\pgfqpoint{0.000000in}{-0.024056in}}%
\pgfpathlineto{\pgfqpoint{0.000000in}{-0.024056in}}%
\pgfpathclose%
\pgfusepath{fill}%
}%
\begin{pgfscope}%
\pgfsys@transformshift{4.558122in}{2.189525in}%
\pgfsys@useobject{currentmarker}{}%
\end{pgfscope}%
\end{pgfscope}%
\begin{pgfscope}%
\definecolor{textcolor}{rgb}{0.549020,0.549020,0.549020}%
\pgfsetstrokecolor{textcolor}%
\pgfsetfillcolor{textcolor}%
\pgftext[x=4.598547in,y=2.189525in,left,]{\color{textcolor}\sffamily\fontsize{7.000000}{8.400000}\selectfont 0.50}%
\end{pgfscope}%
\begin{pgfscope}%
\pgfpathrectangle{\pgfqpoint{4.165000in}{0.598000in}}{\pgfqpoint{1.155000in}{3.864000in}}%
\pgfusepath{clip}%
\pgfsetbuttcap%
\pgfsetroundjoin%
\definecolor{currentfill}{rgb}{0.156863,0.356863,0.568627}%
\pgfsetfillcolor{currentfill}%
\pgfsetlinewidth{0.000000pt}%
\definecolor{currentstroke}{rgb}{0.156863,0.356863,0.568627}%
\pgfsetstrokecolor{currentstroke}%
\pgfsetdash{}{0pt}%
\pgfsys@defobject{currentmarker}{\pgfqpoint{-0.024056in}{-0.024056in}}{\pgfqpoint{0.024056in}{0.024056in}}{%
\pgfpathmoveto{\pgfqpoint{0.000000in}{-0.024056in}}%
\pgfpathcurveto{\pgfqpoint{0.006380in}{-0.024056in}}{\pgfqpoint{0.012499in}{-0.021522in}}{\pgfqpoint{0.017010in}{-0.017010in}}%
\pgfpathcurveto{\pgfqpoint{0.021522in}{-0.012499in}}{\pgfqpoint{0.024056in}{-0.006380in}}{\pgfqpoint{0.024056in}{0.000000in}}%
\pgfpathcurveto{\pgfqpoint{0.024056in}{0.006380in}}{\pgfqpoint{0.021522in}{0.012499in}}{\pgfqpoint{0.017010in}{0.017010in}}%
\pgfpathcurveto{\pgfqpoint{0.012499in}{0.021522in}}{\pgfqpoint{0.006380in}{0.024056in}}{\pgfqpoint{0.000000in}{0.024056in}}%
\pgfpathcurveto{\pgfqpoint{-0.006380in}{0.024056in}}{\pgfqpoint{-0.012499in}{0.021522in}}{\pgfqpoint{-0.017010in}{0.017010in}}%
\pgfpathcurveto{\pgfqpoint{-0.021522in}{0.012499in}}{\pgfqpoint{-0.024056in}{0.006380in}}{\pgfqpoint{-0.024056in}{0.000000in}}%
\pgfpathcurveto{\pgfqpoint{-0.024056in}{-0.006380in}}{\pgfqpoint{-0.021522in}{-0.012499in}}{\pgfqpoint{-0.017010in}{-0.017010in}}%
\pgfpathcurveto{\pgfqpoint{-0.012499in}{-0.021522in}}{\pgfqpoint{-0.006380in}{-0.024056in}}{\pgfqpoint{0.000000in}{-0.024056in}}%
\pgfpathlineto{\pgfqpoint{0.000000in}{-0.024056in}}%
\pgfpathclose%
\pgfusepath{fill}%
}%
\begin{pgfscope}%
\pgfsys@transformshift{4.704325in}{2.031164in}%
\pgfsys@useobject{currentmarker}{}%
\end{pgfscope}%
\end{pgfscope}%
\begin{pgfscope}%
\definecolor{textcolor}{rgb}{0.549020,0.549020,0.549020}%
\pgfsetstrokecolor{textcolor}%
\pgfsetfillcolor{textcolor}%
\pgftext[x=4.744750in,y=2.031164in,left,]{\color{textcolor}\sffamily\fontsize{7.000000}{8.400000}\selectfont 0.54}%
\end{pgfscope}%
\begin{pgfscope}%
\pgfpathrectangle{\pgfqpoint{4.165000in}{0.598000in}}{\pgfqpoint{1.155000in}{3.864000in}}%
\pgfusepath{clip}%
\pgfsetbuttcap%
\pgfsetroundjoin%
\definecolor{currentfill}{rgb}{0.101961,0.101961,0.101961}%
\pgfsetfillcolor{currentfill}%
\pgfsetlinewidth{0.000000pt}%
\definecolor{currentstroke}{rgb}{0.101961,0.101961,0.101961}%
\pgfsetstrokecolor{currentstroke}%
\pgfsetdash{}{0pt}%
\pgfsys@defobject{currentmarker}{\pgfqpoint{-0.032572in}{-0.032572in}}{\pgfqpoint{0.032572in}{0.032572in}}{%
\pgfpathmoveto{\pgfqpoint{-0.016286in}{-0.032572in}}%
\pgfpathlineto{\pgfqpoint{0.000000in}{-0.016286in}}%
\pgfpathlineto{\pgfqpoint{0.016286in}{-0.032572in}}%
\pgfpathlineto{\pgfqpoint{0.032572in}{-0.016286in}}%
\pgfpathlineto{\pgfqpoint{0.016286in}{0.000000in}}%
\pgfpathlineto{\pgfqpoint{0.032572in}{0.016286in}}%
\pgfpathlineto{\pgfqpoint{0.016286in}{0.032572in}}%
\pgfpathlineto{\pgfqpoint{0.000000in}{0.016286in}}%
\pgfpathlineto{\pgfqpoint{-0.016286in}{0.032572in}}%
\pgfpathlineto{\pgfqpoint{-0.032572in}{0.016286in}}%
\pgfpathlineto{\pgfqpoint{-0.016286in}{0.000000in}}%
\pgfpathlineto{\pgfqpoint{-0.032572in}{-0.016286in}}%
\pgfpathlineto{\pgfqpoint{-0.016286in}{-0.032572in}}%
\pgfpathclose%
\pgfusepath{fill}%
}%
\begin{pgfscope}%
\pgfsys@transformshift{4.623101in}{1.872803in}%
\pgfsys@useobject{currentmarker}{}%
\end{pgfscope}%
\end{pgfscope}%
\begin{pgfscope}%
\definecolor{textcolor}{rgb}{0.549020,0.549020,0.549020}%
\pgfsetstrokecolor{textcolor}%
\pgfsetfillcolor{textcolor}%
\pgftext[x=4.663526in,y=1.872803in,left,]{\color{textcolor}\sffamily\fontsize{7.000000}{8.400000}\selectfont 0.52}%
\end{pgfscope}%
\begin{pgfscope}%
\pgfpathrectangle{\pgfqpoint{4.165000in}{0.598000in}}{\pgfqpoint{1.155000in}{3.864000in}}%
\pgfusepath{clip}%
\pgfsetbuttcap%
\pgfsetroundjoin%
\definecolor{currentfill}{rgb}{0.101961,0.101961,0.101961}%
\pgfsetfillcolor{currentfill}%
\pgfsetlinewidth{0.000000pt}%
\definecolor{currentstroke}{rgb}{0.101961,0.101961,0.101961}%
\pgfsetstrokecolor{currentstroke}%
\pgfsetdash{}{0pt}%
\pgfsys@defobject{currentmarker}{\pgfqpoint{-0.032572in}{-0.032572in}}{\pgfqpoint{0.032572in}{0.032572in}}{%
\pgfpathmoveto{\pgfqpoint{-0.016286in}{-0.032572in}}%
\pgfpathlineto{\pgfqpoint{0.000000in}{-0.016286in}}%
\pgfpathlineto{\pgfqpoint{0.016286in}{-0.032572in}}%
\pgfpathlineto{\pgfqpoint{0.032572in}{-0.016286in}}%
\pgfpathlineto{\pgfqpoint{0.016286in}{0.000000in}}%
\pgfpathlineto{\pgfqpoint{0.032572in}{0.016286in}}%
\pgfpathlineto{\pgfqpoint{0.016286in}{0.032572in}}%
\pgfpathlineto{\pgfqpoint{0.000000in}{0.016286in}}%
\pgfpathlineto{\pgfqpoint{-0.016286in}{0.032572in}}%
\pgfpathlineto{\pgfqpoint{-0.032572in}{0.016286in}}%
\pgfpathlineto{\pgfqpoint{-0.016286in}{0.000000in}}%
\pgfpathlineto{\pgfqpoint{-0.032572in}{-0.016286in}}%
\pgfpathlineto{\pgfqpoint{-0.016286in}{-0.032572in}}%
\pgfpathclose%
\pgfusepath{fill}%
}%
\begin{pgfscope}%
\pgfsys@transformshift{4.671835in}{1.714443in}%
\pgfsys@useobject{currentmarker}{}%
\end{pgfscope}%
\end{pgfscope}%
\begin{pgfscope}%
\definecolor{textcolor}{rgb}{0.549020,0.549020,0.549020}%
\pgfsetstrokecolor{textcolor}%
\pgfsetfillcolor{textcolor}%
\pgftext[x=4.712260in,y=1.714443in,left,]{\color{textcolor}\sffamily\fontsize{7.000000}{8.400000}\selectfont 0.53}%
\end{pgfscope}%
\begin{pgfscope}%
\pgfpathrectangle{\pgfqpoint{4.165000in}{0.598000in}}{\pgfqpoint{1.155000in}{3.864000in}}%
\pgfusepath{clip}%
\pgfsetbuttcap%
\pgfsetroundjoin%
\definecolor{currentfill}{rgb}{0.156863,0.356863,0.568627}%
\pgfsetfillcolor{currentfill}%
\pgfsetlinewidth{0.000000pt}%
\definecolor{currentstroke}{rgb}{0.156863,0.356863,0.568627}%
\pgfsetstrokecolor{currentstroke}%
\pgfsetdash{}{0pt}%
\pgfsys@defobject{currentmarker}{\pgfqpoint{-0.024056in}{-0.024056in}}{\pgfqpoint{0.024056in}{0.024056in}}{%
\pgfpathmoveto{\pgfqpoint{0.000000in}{-0.024056in}}%
\pgfpathcurveto{\pgfqpoint{0.006380in}{-0.024056in}}{\pgfqpoint{0.012499in}{-0.021522in}}{\pgfqpoint{0.017010in}{-0.017010in}}%
\pgfpathcurveto{\pgfqpoint{0.021522in}{-0.012499in}}{\pgfqpoint{0.024056in}{-0.006380in}}{\pgfqpoint{0.024056in}{0.000000in}}%
\pgfpathcurveto{\pgfqpoint{0.024056in}{0.006380in}}{\pgfqpoint{0.021522in}{0.012499in}}{\pgfqpoint{0.017010in}{0.017010in}}%
\pgfpathcurveto{\pgfqpoint{0.012499in}{0.021522in}}{\pgfqpoint{0.006380in}{0.024056in}}{\pgfqpoint{0.000000in}{0.024056in}}%
\pgfpathcurveto{\pgfqpoint{-0.006380in}{0.024056in}}{\pgfqpoint{-0.012499in}{0.021522in}}{\pgfqpoint{-0.017010in}{0.017010in}}%
\pgfpathcurveto{\pgfqpoint{-0.021522in}{0.012499in}}{\pgfqpoint{-0.024056in}{0.006380in}}{\pgfqpoint{-0.024056in}{0.000000in}}%
\pgfpathcurveto{\pgfqpoint{-0.024056in}{-0.006380in}}{\pgfqpoint{-0.021522in}{-0.012499in}}{\pgfqpoint{-0.017010in}{-0.017010in}}%
\pgfpathcurveto{\pgfqpoint{-0.012499in}{-0.021522in}}{\pgfqpoint{-0.006380in}{-0.024056in}}{\pgfqpoint{0.000000in}{-0.024056in}}%
\pgfpathlineto{\pgfqpoint{0.000000in}{-0.024056in}}%
\pgfpathclose%
\pgfusepath{fill}%
}%
\begin{pgfscope}%
\pgfsys@transformshift{4.736814in}{1.326459in}%
\pgfsys@useobject{currentmarker}{}%
\end{pgfscope}%
\end{pgfscope}%
\begin{pgfscope}%
\definecolor{textcolor}{rgb}{0.549020,0.549020,0.549020}%
\pgfsetstrokecolor{textcolor}%
\pgfsetfillcolor{textcolor}%
\pgftext[x=4.777239in,y=1.326459in,left,]{\color{textcolor}\sffamily\fontsize{7.000000}{8.400000}\selectfont 0.55}%
\end{pgfscope}%
\begin{pgfscope}%
\pgfpathrectangle{\pgfqpoint{4.165000in}{0.598000in}}{\pgfqpoint{1.155000in}{3.864000in}}%
\pgfusepath{clip}%
\pgfsetbuttcap%
\pgfsetroundjoin%
\definecolor{currentfill}{rgb}{0.156863,0.356863,0.568627}%
\pgfsetfillcolor{currentfill}%
\pgfsetlinewidth{0.000000pt}%
\definecolor{currentstroke}{rgb}{0.156863,0.356863,0.568627}%
\pgfsetstrokecolor{currentstroke}%
\pgfsetdash{}{0pt}%
\pgfsys@defobject{currentmarker}{\pgfqpoint{-0.024056in}{-0.024056in}}{\pgfqpoint{0.024056in}{0.024056in}}{%
\pgfpathmoveto{\pgfqpoint{0.000000in}{-0.024056in}}%
\pgfpathcurveto{\pgfqpoint{0.006380in}{-0.024056in}}{\pgfqpoint{0.012499in}{-0.021522in}}{\pgfqpoint{0.017010in}{-0.017010in}}%
\pgfpathcurveto{\pgfqpoint{0.021522in}{-0.012499in}}{\pgfqpoint{0.024056in}{-0.006380in}}{\pgfqpoint{0.024056in}{0.000000in}}%
\pgfpathcurveto{\pgfqpoint{0.024056in}{0.006380in}}{\pgfqpoint{0.021522in}{0.012499in}}{\pgfqpoint{0.017010in}{0.017010in}}%
\pgfpathcurveto{\pgfqpoint{0.012499in}{0.021522in}}{\pgfqpoint{0.006380in}{0.024056in}}{\pgfqpoint{0.000000in}{0.024056in}}%
\pgfpathcurveto{\pgfqpoint{-0.006380in}{0.024056in}}{\pgfqpoint{-0.012499in}{0.021522in}}{\pgfqpoint{-0.017010in}{0.017010in}}%
\pgfpathcurveto{\pgfqpoint{-0.021522in}{0.012499in}}{\pgfqpoint{-0.024056in}{0.006380in}}{\pgfqpoint{-0.024056in}{0.000000in}}%
\pgfpathcurveto{\pgfqpoint{-0.024056in}{-0.006380in}}{\pgfqpoint{-0.021522in}{-0.012499in}}{\pgfqpoint{-0.017010in}{-0.017010in}}%
\pgfpathcurveto{\pgfqpoint{-0.012499in}{-0.021522in}}{\pgfqpoint{-0.006380in}{-0.024056in}}{\pgfqpoint{0.000000in}{-0.024056in}}%
\pgfpathlineto{\pgfqpoint{0.000000in}{-0.024056in}}%
\pgfpathclose%
\pgfusepath{fill}%
}%
\begin{pgfscope}%
\pgfsys@transformshift{4.769304in}{1.168098in}%
\pgfsys@useobject{currentmarker}{}%
\end{pgfscope}%
\end{pgfscope}%
\begin{pgfscope}%
\definecolor{textcolor}{rgb}{0.549020,0.549020,0.549020}%
\pgfsetstrokecolor{textcolor}%
\pgfsetfillcolor{textcolor}%
\pgftext[x=4.809729in,y=1.168098in,left,]{\color{textcolor}\sffamily\fontsize{7.000000}{8.400000}\selectfont 0.56}%
\end{pgfscope}%
\begin{pgfscope}%
\pgfpathrectangle{\pgfqpoint{4.165000in}{0.598000in}}{\pgfqpoint{1.155000in}{3.864000in}}%
\pgfusepath{clip}%
\pgfsetbuttcap%
\pgfsetroundjoin%
\definecolor{currentfill}{rgb}{0.156863,0.356863,0.568627}%
\pgfsetfillcolor{currentfill}%
\pgfsetlinewidth{0.000000pt}%
\definecolor{currentstroke}{rgb}{0.156863,0.356863,0.568627}%
\pgfsetstrokecolor{currentstroke}%
\pgfsetdash{}{0pt}%
\pgfsys@defobject{currentmarker}{\pgfqpoint{-0.024056in}{-0.024056in}}{\pgfqpoint{0.024056in}{0.024056in}}{%
\pgfpathmoveto{\pgfqpoint{0.000000in}{-0.024056in}}%
\pgfpathcurveto{\pgfqpoint{0.006380in}{-0.024056in}}{\pgfqpoint{0.012499in}{-0.021522in}}{\pgfqpoint{0.017010in}{-0.017010in}}%
\pgfpathcurveto{\pgfqpoint{0.021522in}{-0.012499in}}{\pgfqpoint{0.024056in}{-0.006380in}}{\pgfqpoint{0.024056in}{0.000000in}}%
\pgfpathcurveto{\pgfqpoint{0.024056in}{0.006380in}}{\pgfqpoint{0.021522in}{0.012499in}}{\pgfqpoint{0.017010in}{0.017010in}}%
\pgfpathcurveto{\pgfqpoint{0.012499in}{0.021522in}}{\pgfqpoint{0.006380in}{0.024056in}}{\pgfqpoint{0.000000in}{0.024056in}}%
\pgfpathcurveto{\pgfqpoint{-0.006380in}{0.024056in}}{\pgfqpoint{-0.012499in}{0.021522in}}{\pgfqpoint{-0.017010in}{0.017010in}}%
\pgfpathcurveto{\pgfqpoint{-0.021522in}{0.012499in}}{\pgfqpoint{-0.024056in}{0.006380in}}{\pgfqpoint{-0.024056in}{0.000000in}}%
\pgfpathcurveto{\pgfqpoint{-0.024056in}{-0.006380in}}{\pgfqpoint{-0.021522in}{-0.012499in}}{\pgfqpoint{-0.017010in}{-0.017010in}}%
\pgfpathcurveto{\pgfqpoint{-0.012499in}{-0.021522in}}{\pgfqpoint{-0.006380in}{-0.024056in}}{\pgfqpoint{0.000000in}{-0.024056in}}%
\pgfpathlineto{\pgfqpoint{0.000000in}{-0.024056in}}%
\pgfpathclose%
\pgfusepath{fill}%
}%
\begin{pgfscope}%
\pgfsys@transformshift{4.688080in}{1.009738in}%
\pgfsys@useobject{currentmarker}{}%
\end{pgfscope}%
\end{pgfscope}%
\begin{pgfscope}%
\definecolor{textcolor}{rgb}{0.549020,0.549020,0.549020}%
\pgfsetstrokecolor{textcolor}%
\pgfsetfillcolor{textcolor}%
\pgftext[x=4.728505in,y=1.009738in,left,]{\color{textcolor}\sffamily\fontsize{7.000000}{8.400000}\selectfont 0.54}%
\end{pgfscope}%
\begin{pgfscope}%
\pgfpathrectangle{\pgfqpoint{4.165000in}{0.598000in}}{\pgfqpoint{1.155000in}{3.864000in}}%
\pgfusepath{clip}%
\pgfsetbuttcap%
\pgfsetroundjoin%
\definecolor{currentfill}{rgb}{0.101961,0.101961,0.101961}%
\pgfsetfillcolor{currentfill}%
\pgfsetlinewidth{0.000000pt}%
\definecolor{currentstroke}{rgb}{0.101961,0.101961,0.101961}%
\pgfsetstrokecolor{currentstroke}%
\pgfsetdash{}{0pt}%
\pgfsys@defobject{currentmarker}{\pgfqpoint{-0.032572in}{-0.032572in}}{\pgfqpoint{0.032572in}{0.032572in}}{%
\pgfpathmoveto{\pgfqpoint{-0.016286in}{-0.032572in}}%
\pgfpathlineto{\pgfqpoint{0.000000in}{-0.016286in}}%
\pgfpathlineto{\pgfqpoint{0.016286in}{-0.032572in}}%
\pgfpathlineto{\pgfqpoint{0.032572in}{-0.016286in}}%
\pgfpathlineto{\pgfqpoint{0.016286in}{0.000000in}}%
\pgfpathlineto{\pgfqpoint{0.032572in}{0.016286in}}%
\pgfpathlineto{\pgfqpoint{0.016286in}{0.032572in}}%
\pgfpathlineto{\pgfqpoint{0.000000in}{0.016286in}}%
\pgfpathlineto{\pgfqpoint{-0.016286in}{0.032572in}}%
\pgfpathlineto{\pgfqpoint{-0.032572in}{0.016286in}}%
\pgfpathlineto{\pgfqpoint{-0.016286in}{0.000000in}}%
\pgfpathlineto{\pgfqpoint{-0.032572in}{-0.016286in}}%
\pgfpathlineto{\pgfqpoint{-0.016286in}{-0.032572in}}%
\pgfpathclose%
\pgfusepath{fill}%
}%
\begin{pgfscope}%
\pgfsys@transformshift{4.671835in}{0.851377in}%
\pgfsys@useobject{currentmarker}{}%
\end{pgfscope}%
\end{pgfscope}%
\begin{pgfscope}%
\definecolor{textcolor}{rgb}{0.549020,0.549020,0.549020}%
\pgfsetstrokecolor{textcolor}%
\pgfsetfillcolor{textcolor}%
\pgftext[x=4.712260in,y=0.851377in,left,]{\color{textcolor}\sffamily\fontsize{7.000000}{8.400000}\selectfont 0.53}%
\end{pgfscope}%
\begin{pgfscope}%
\pgfpathrectangle{\pgfqpoint{4.165000in}{0.598000in}}{\pgfqpoint{1.155000in}{3.864000in}}%
\pgfusepath{clip}%
\pgfsetbuttcap%
\pgfsetroundjoin%
\definecolor{currentfill}{rgb}{0.101961,0.101961,0.101961}%
\pgfsetfillcolor{currentfill}%
\pgfsetlinewidth{0.000000pt}%
\definecolor{currentstroke}{rgb}{0.101961,0.101961,0.101961}%
\pgfsetstrokecolor{currentstroke}%
\pgfsetdash{}{0pt}%
\pgfsys@defobject{currentmarker}{\pgfqpoint{-0.032572in}{-0.032572in}}{\pgfqpoint{0.032572in}{0.032572in}}{%
\pgfpathmoveto{\pgfqpoint{-0.016286in}{-0.032572in}}%
\pgfpathlineto{\pgfqpoint{0.000000in}{-0.016286in}}%
\pgfpathlineto{\pgfqpoint{0.016286in}{-0.032572in}}%
\pgfpathlineto{\pgfqpoint{0.032572in}{-0.016286in}}%
\pgfpathlineto{\pgfqpoint{0.016286in}{0.000000in}}%
\pgfpathlineto{\pgfqpoint{0.032572in}{0.016286in}}%
\pgfpathlineto{\pgfqpoint{0.016286in}{0.032572in}}%
\pgfpathlineto{\pgfqpoint{0.000000in}{0.016286in}}%
\pgfpathlineto{\pgfqpoint{-0.016286in}{0.032572in}}%
\pgfpathlineto{\pgfqpoint{-0.032572in}{0.016286in}}%
\pgfpathlineto{\pgfqpoint{-0.016286in}{0.000000in}}%
\pgfpathlineto{\pgfqpoint{-0.032572in}{-0.016286in}}%
\pgfpathlineto{\pgfqpoint{-0.016286in}{-0.032572in}}%
\pgfpathclose%
\pgfusepath{fill}%
}%
\begin{pgfscope}%
\pgfsys@transformshift{4.704325in}{0.693016in}%
\pgfsys@useobject{currentmarker}{}%
\end{pgfscope}%
\end{pgfscope}%
\begin{pgfscope}%
\definecolor{textcolor}{rgb}{0.549020,0.549020,0.549020}%
\pgfsetstrokecolor{textcolor}%
\pgfsetfillcolor{textcolor}%
\pgftext[x=4.744750in,y=0.693016in,left,]{\color{textcolor}\sffamily\fontsize{7.000000}{8.400000}\selectfont 0.54}%
\end{pgfscope}%
\begin{pgfscope}%
\pgfsetbuttcap%
\pgfsetmiterjoin%
\definecolor{currentfill}{rgb}{1.000000,1.000000,1.000000}%
\pgfsetfillcolor{currentfill}%
\pgfsetlinewidth{0.000000pt}%
\definecolor{currentstroke}{rgb}{0.000000,0.000000,0.000000}%
\pgfsetstrokecolor{currentstroke}%
\pgfsetstrokeopacity{0.000000}%
\pgfsetdash{}{0pt}%
\pgfpathmoveto{\pgfqpoint{5.635000in}{0.598000in}}%
\pgfpathlineto{\pgfqpoint{6.790000in}{0.598000in}}%
\pgfpathlineto{\pgfqpoint{6.790000in}{4.462000in}}%
\pgfpathlineto{\pgfqpoint{5.635000in}{4.462000in}}%
\pgfpathlineto{\pgfqpoint{5.635000in}{0.598000in}}%
\pgfpathclose%
\pgfusepath{fill}%
\end{pgfscope}%
\begin{pgfscope}%
\pgfpathrectangle{\pgfqpoint{5.635000in}{0.598000in}}{\pgfqpoint{1.155000in}{3.864000in}}%
\pgfusepath{clip}%
\pgfsetrectcap%
\pgfsetroundjoin%
\pgfsetlinewidth{0.501875pt}%
\definecolor{currentstroke}{rgb}{0.901961,0.901961,0.901961}%
\pgfsetstrokecolor{currentstroke}%
\pgfsetdash{}{0pt}%
\pgfpathmoveto{\pgfqpoint{5.635000in}{4.398656in}}%
\pgfpathlineto{\pgfqpoint{6.790000in}{4.398656in}}%
\pgfusepath{stroke}%
\end{pgfscope}%
\begin{pgfscope}%
\pgfpathrectangle{\pgfqpoint{5.635000in}{0.598000in}}{\pgfqpoint{1.155000in}{3.864000in}}%
\pgfusepath{clip}%
\pgfsetrectcap%
\pgfsetroundjoin%
\pgfsetlinewidth{0.501875pt}%
\definecolor{currentstroke}{rgb}{0.901961,0.901961,0.901961}%
\pgfsetstrokecolor{currentstroke}%
\pgfsetdash{}{0pt}%
\pgfpathmoveto{\pgfqpoint{5.635000in}{2.427066in}}%
\pgfpathlineto{\pgfqpoint{6.790000in}{2.427066in}}%
\pgfusepath{stroke}%
\end{pgfscope}%
\begin{pgfscope}%
\pgfpathrectangle{\pgfqpoint{5.635000in}{0.598000in}}{\pgfqpoint{1.155000in}{3.864000in}}%
\pgfusepath{clip}%
\pgfsetrectcap%
\pgfsetroundjoin%
\pgfsetlinewidth{0.501875pt}%
\definecolor{currentstroke}{rgb}{0.901961,0.901961,0.901961}%
\pgfsetstrokecolor{currentstroke}%
\pgfsetdash{}{0pt}%
\pgfpathmoveto{\pgfqpoint{5.635000in}{1.564000in}}%
\pgfpathlineto{\pgfqpoint{6.790000in}{1.564000in}}%
\pgfusepath{stroke}%
\end{pgfscope}%
\begin{pgfscope}%
\pgfpathrectangle{\pgfqpoint{5.635000in}{0.598000in}}{\pgfqpoint{1.155000in}{3.864000in}}%
\pgfusepath{clip}%
\pgfsetrectcap%
\pgfsetroundjoin%
\pgfsetlinewidth{0.401500pt}%
\definecolor{currentstroke}{rgb}{0.901961,0.901961,0.901961}%
\pgfsetstrokecolor{currentstroke}%
\pgfsetdash{}{0pt}%
\pgfpathmoveto{\pgfqpoint{5.635000in}{0.598000in}}%
\pgfpathlineto{\pgfqpoint{5.635000in}{4.462000in}}%
\pgfusepath{stroke}%
\end{pgfscope}%
\begin{pgfscope}%
\definecolor{textcolor}{rgb}{0.101961,0.101961,0.101961}%
\pgfsetstrokecolor{textcolor}%
\pgfsetfillcolor{textcolor}%
\pgftext[x=5.635000in,y=0.549389in,,top]{\color{textcolor}\sffamily\fontsize{7.000000}{8.400000}\selectfont 0.0}%
\end{pgfscope}%
\begin{pgfscope}%
\pgfpathrectangle{\pgfqpoint{5.635000in}{0.598000in}}{\pgfqpoint{1.155000in}{3.864000in}}%
\pgfusepath{clip}%
\pgfsetrectcap%
\pgfsetroundjoin%
\pgfsetlinewidth{0.401500pt}%
\definecolor{currentstroke}{rgb}{0.901961,0.901961,0.901961}%
\pgfsetstrokecolor{currentstroke}%
\pgfsetdash{}{0pt}%
\pgfpathmoveto{\pgfqpoint{6.212500in}{0.598000in}}%
\pgfpathlineto{\pgfqpoint{6.212500in}{4.462000in}}%
\pgfusepath{stroke}%
\end{pgfscope}%
\begin{pgfscope}%
\definecolor{textcolor}{rgb}{0.101961,0.101961,0.101961}%
\pgfsetstrokecolor{textcolor}%
\pgfsetfillcolor{textcolor}%
\pgftext[x=6.212500in,y=0.549389in,,top]{\color{textcolor}\sffamily\fontsize{7.000000}{8.400000}\selectfont 0.5}%
\end{pgfscope}%
\begin{pgfscope}%
\pgfpathrectangle{\pgfqpoint{5.635000in}{0.598000in}}{\pgfqpoint{1.155000in}{3.864000in}}%
\pgfusepath{clip}%
\pgfsetrectcap%
\pgfsetroundjoin%
\pgfsetlinewidth{0.401500pt}%
\definecolor{currentstroke}{rgb}{0.901961,0.901961,0.901961}%
\pgfsetstrokecolor{currentstroke}%
\pgfsetdash{}{0pt}%
\pgfpathmoveto{\pgfqpoint{6.790000in}{0.598000in}}%
\pgfpathlineto{\pgfqpoint{6.790000in}{4.462000in}}%
\pgfusepath{stroke}%
\end{pgfscope}%
\begin{pgfscope}%
\definecolor{textcolor}{rgb}{0.101961,0.101961,0.101961}%
\pgfsetstrokecolor{textcolor}%
\pgfsetfillcolor{textcolor}%
\pgftext[x=6.790000in,y=0.549389in,,top]{\color{textcolor}\sffamily\fontsize{7.000000}{8.400000}\selectfont 1.0}%
\end{pgfscope}%
\begin{pgfscope}%
\definecolor{textcolor}{rgb}{0.000000,0.000000,0.000000}%
\pgfsetstrokecolor{textcolor}%
\pgfsetfillcolor{textcolor}%
\pgftext[x=5.901537in, y=0.367676in, left, base]{\color{textcolor}\sffamily\fontsize{7.500000}{9.000000}\selectfont Dashcam alert}%
\end{pgfscope}%
\begin{pgfscope}%
\definecolor{textcolor}{rgb}{0.000000,0.000000,0.000000}%
\pgfsetstrokecolor{textcolor}%
\pgfsetfillcolor{textcolor}%
\pgftext[x=5.977291in, y=0.267754in, left, base]{\color{textcolor}\sffamily\fontsize{7.500000}{9.000000}\selectfont attainment}%
\end{pgfscope}%
\begin{pgfscope}%
\pgfpathrectangle{\pgfqpoint{5.635000in}{0.598000in}}{\pgfqpoint{1.155000in}{3.864000in}}%
\pgfusepath{clip}%
\pgfsetrectcap%
\pgfsetroundjoin%
\pgfsetlinewidth{0.602250pt}%
\definecolor{currentstroke}{rgb}{0.156863,0.356863,0.568627}%
\pgfsetstrokecolor{currentstroke}%
\pgfsetdash{}{0pt}%
\pgfpathmoveto{\pgfqpoint{6.456363in}{3.369311in}}%
\pgfpathlineto{\pgfqpoint{6.445789in}{3.210951in}}%
\pgfpathlineto{\pgfqpoint{6.460826in}{3.052590in}}%
\pgfpathlineto{\pgfqpoint{6.452999in}{2.894230in}}%
\pgfpathlineto{\pgfqpoint{6.467247in}{2.735869in}}%
\pgfpathlineto{\pgfqpoint{6.472882in}{2.577508in}}%
\pgfusepath{stroke}%
\end{pgfscope}%
\begin{pgfscope}%
\pgfsetrectcap%
\pgfsetmiterjoin%
\pgfsetlinewidth{0.803000pt}%
\definecolor{currentstroke}{rgb}{0.101961,0.101961,0.101961}%
\pgfsetstrokecolor{currentstroke}%
\pgfsetdash{}{0pt}%
\pgfpathmoveto{\pgfqpoint{5.635000in}{0.598000in}}%
\pgfpathlineto{\pgfqpoint{6.790000in}{0.598000in}}%
\pgfusepath{stroke}%
\end{pgfscope}%
\begin{pgfscope}%
\pgfpathrectangle{\pgfqpoint{5.635000in}{0.598000in}}{\pgfqpoint{1.155000in}{3.864000in}}%
\pgfusepath{clip}%
\pgfsetbuttcap%
\pgfsetroundjoin%
\definecolor{currentfill}{rgb}{0.435294,0.560784,0.709804}%
\pgfsetfillcolor{currentfill}%
\pgfsetlinewidth{0.000000pt}%
\definecolor{currentstroke}{rgb}{0.435294,0.560784,0.709804}%
\pgfsetstrokecolor{currentstroke}%
\pgfsetdash{}{0pt}%
\pgfsys@defobject{currentmarker}{\pgfqpoint{-0.024056in}{-0.024056in}}{\pgfqpoint{0.024056in}{0.024056in}}{%
\pgfpathmoveto{\pgfqpoint{0.000000in}{-0.024056in}}%
\pgfpathcurveto{\pgfqpoint{0.006380in}{-0.024056in}}{\pgfqpoint{0.012499in}{-0.021522in}}{\pgfqpoint{0.017010in}{-0.017010in}}%
\pgfpathcurveto{\pgfqpoint{0.021522in}{-0.012499in}}{\pgfqpoint{0.024056in}{-0.006380in}}{\pgfqpoint{0.024056in}{0.000000in}}%
\pgfpathcurveto{\pgfqpoint{0.024056in}{0.006380in}}{\pgfqpoint{0.021522in}{0.012499in}}{\pgfqpoint{0.017010in}{0.017010in}}%
\pgfpathcurveto{\pgfqpoint{0.012499in}{0.021522in}}{\pgfqpoint{0.006380in}{0.024056in}}{\pgfqpoint{0.000000in}{0.024056in}}%
\pgfpathcurveto{\pgfqpoint{-0.006380in}{0.024056in}}{\pgfqpoint{-0.012499in}{0.021522in}}{\pgfqpoint{-0.017010in}{0.017010in}}%
\pgfpathcurveto{\pgfqpoint{-0.021522in}{0.012499in}}{\pgfqpoint{-0.024056in}{0.006380in}}{\pgfqpoint{-0.024056in}{0.000000in}}%
\pgfpathcurveto{\pgfqpoint{-0.024056in}{-0.006380in}}{\pgfqpoint{-0.021522in}{-0.012499in}}{\pgfqpoint{-0.017010in}{-0.017010in}}%
\pgfpathcurveto{\pgfqpoint{-0.012499in}{-0.021522in}}{\pgfqpoint{-0.006380in}{-0.024056in}}{\pgfqpoint{0.000000in}{-0.024056in}}%
\pgfpathlineto{\pgfqpoint{0.000000in}{-0.024056in}}%
\pgfpathclose%
\pgfusepath{fill}%
}%
\begin{pgfscope}%
\pgfsys@transformshift{6.214187in}{4.161115in}%
\pgfsys@useobject{currentmarker}{}%
\end{pgfscope}%
\end{pgfscope}%
\begin{pgfscope}%
\definecolor{textcolor}{rgb}{0.549020,0.549020,0.549020}%
\pgfsetstrokecolor{textcolor}%
\pgfsetfillcolor{textcolor}%
\pgftext[x=6.254612in,y=4.161115in,left,]{\color{textcolor}\sffamily\fontsize{7.000000}{8.400000}\selectfont 0.50}%
\end{pgfscope}%
\begin{pgfscope}%
\pgfpathrectangle{\pgfqpoint{5.635000in}{0.598000in}}{\pgfqpoint{1.155000in}{3.864000in}}%
\pgfusepath{clip}%
\pgfsetbuttcap%
\pgfsetroundjoin%
\definecolor{currentfill}{rgb}{0.435294,0.560784,0.709804}%
\pgfsetfillcolor{currentfill}%
\pgfsetlinewidth{0.000000pt}%
\definecolor{currentstroke}{rgb}{0.435294,0.560784,0.709804}%
\pgfsetstrokecolor{currentstroke}%
\pgfsetdash{}{0pt}%
\pgfsys@defobject{currentmarker}{\pgfqpoint{-0.024056in}{-0.024056in}}{\pgfqpoint{0.024056in}{0.024056in}}{%
\pgfpathmoveto{\pgfqpoint{0.000000in}{-0.024056in}}%
\pgfpathcurveto{\pgfqpoint{0.006380in}{-0.024056in}}{\pgfqpoint{0.012499in}{-0.021522in}}{\pgfqpoint{0.017010in}{-0.017010in}}%
\pgfpathcurveto{\pgfqpoint{0.021522in}{-0.012499in}}{\pgfqpoint{0.024056in}{-0.006380in}}{\pgfqpoint{0.024056in}{0.000000in}}%
\pgfpathcurveto{\pgfqpoint{0.024056in}{0.006380in}}{\pgfqpoint{0.021522in}{0.012499in}}{\pgfqpoint{0.017010in}{0.017010in}}%
\pgfpathcurveto{\pgfqpoint{0.012499in}{0.021522in}}{\pgfqpoint{0.006380in}{0.024056in}}{\pgfqpoint{0.000000in}{0.024056in}}%
\pgfpathcurveto{\pgfqpoint{-0.006380in}{0.024056in}}{\pgfqpoint{-0.012499in}{0.021522in}}{\pgfqpoint{-0.017010in}{0.017010in}}%
\pgfpathcurveto{\pgfqpoint{-0.021522in}{0.012499in}}{\pgfqpoint{-0.024056in}{0.006380in}}{\pgfqpoint{-0.024056in}{0.000000in}}%
\pgfpathcurveto{\pgfqpoint{-0.024056in}{-0.006380in}}{\pgfqpoint{-0.021522in}{-0.012499in}}{\pgfqpoint{-0.017010in}{-0.017010in}}%
\pgfpathcurveto{\pgfqpoint{-0.012499in}{-0.021522in}}{\pgfqpoint{-0.006380in}{-0.024056in}}{\pgfqpoint{0.000000in}{-0.024056in}}%
\pgfpathlineto{\pgfqpoint{0.000000in}{-0.024056in}}%
\pgfpathclose%
\pgfusepath{fill}%
}%
\begin{pgfscope}%
\pgfsys@transformshift{5.702189in}{4.002754in}%
\pgfsys@useobject{currentmarker}{}%
\end{pgfscope}%
\end{pgfscope}%
\begin{pgfscope}%
\definecolor{textcolor}{rgb}{0.549020,0.549020,0.549020}%
\pgfsetstrokecolor{textcolor}%
\pgfsetfillcolor{textcolor}%
\pgftext[x=5.742614in,y=4.002754in,left,]{\color{textcolor}\sffamily\fontsize{7.000000}{8.400000}\selectfont 0.06}%
\end{pgfscope}%
\begin{pgfscope}%
\pgfpathrectangle{\pgfqpoint{5.635000in}{0.598000in}}{\pgfqpoint{1.155000in}{3.864000in}}%
\pgfusepath{clip}%
\pgfsetbuttcap%
\pgfsetroundjoin%
\definecolor{currentfill}{rgb}{0.435294,0.560784,0.709804}%
\pgfsetfillcolor{currentfill}%
\pgfsetlinewidth{0.000000pt}%
\definecolor{currentstroke}{rgb}{0.435294,0.560784,0.709804}%
\pgfsetstrokecolor{currentstroke}%
\pgfsetdash{}{0pt}%
\pgfsys@defobject{currentmarker}{\pgfqpoint{-0.024056in}{-0.024056in}}{\pgfqpoint{0.024056in}{0.024056in}}{%
\pgfpathmoveto{\pgfqpoint{0.000000in}{-0.024056in}}%
\pgfpathcurveto{\pgfqpoint{0.006380in}{-0.024056in}}{\pgfqpoint{0.012499in}{-0.021522in}}{\pgfqpoint{0.017010in}{-0.017010in}}%
\pgfpathcurveto{\pgfqpoint{0.021522in}{-0.012499in}}{\pgfqpoint{0.024056in}{-0.006380in}}{\pgfqpoint{0.024056in}{0.000000in}}%
\pgfpathcurveto{\pgfqpoint{0.024056in}{0.006380in}}{\pgfqpoint{0.021522in}{0.012499in}}{\pgfqpoint{0.017010in}{0.017010in}}%
\pgfpathcurveto{\pgfqpoint{0.012499in}{0.021522in}}{\pgfqpoint{0.006380in}{0.024056in}}{\pgfqpoint{0.000000in}{0.024056in}}%
\pgfpathcurveto{\pgfqpoint{-0.006380in}{0.024056in}}{\pgfqpoint{-0.012499in}{0.021522in}}{\pgfqpoint{-0.017010in}{0.017010in}}%
\pgfpathcurveto{\pgfqpoint{-0.021522in}{0.012499in}}{\pgfqpoint{-0.024056in}{0.006380in}}{\pgfqpoint{-0.024056in}{0.000000in}}%
\pgfpathcurveto{\pgfqpoint{-0.024056in}{-0.006380in}}{\pgfqpoint{-0.021522in}{-0.012499in}}{\pgfqpoint{-0.017010in}{-0.017010in}}%
\pgfpathcurveto{\pgfqpoint{-0.012499in}{-0.021522in}}{\pgfqpoint{-0.006380in}{-0.024056in}}{\pgfqpoint{0.000000in}{-0.024056in}}%
\pgfpathlineto{\pgfqpoint{0.000000in}{-0.024056in}}%
\pgfpathclose%
\pgfusepath{fill}%
}%
\begin{pgfscope}%
\pgfsys@transformshift{5.934520in}{3.844393in}%
\pgfsys@useobject{currentmarker}{}%
\end{pgfscope}%
\end{pgfscope}%
\begin{pgfscope}%
\definecolor{textcolor}{rgb}{0.549020,0.549020,0.549020}%
\pgfsetstrokecolor{textcolor}%
\pgfsetfillcolor{textcolor}%
\pgftext[x=5.974945in,y=3.844393in,left,]{\color{textcolor}\sffamily\fontsize{7.000000}{8.400000}\selectfont 0.26}%
\end{pgfscope}%
\begin{pgfscope}%
\pgfpathrectangle{\pgfqpoint{5.635000in}{0.598000in}}{\pgfqpoint{1.155000in}{3.864000in}}%
\pgfusepath{clip}%
\pgfsetbuttcap%
\pgfsetroundjoin%
\definecolor{currentfill}{rgb}{0.435294,0.560784,0.709804}%
\pgfsetfillcolor{currentfill}%
\pgfsetlinewidth{0.000000pt}%
\definecolor{currentstroke}{rgb}{0.435294,0.560784,0.709804}%
\pgfsetstrokecolor{currentstroke}%
\pgfsetdash{}{0pt}%
\pgfsys@defobject{currentmarker}{\pgfqpoint{-0.024056in}{-0.024056in}}{\pgfqpoint{0.024056in}{0.024056in}}{%
\pgfpathmoveto{\pgfqpoint{0.000000in}{-0.024056in}}%
\pgfpathcurveto{\pgfqpoint{0.006380in}{-0.024056in}}{\pgfqpoint{0.012499in}{-0.021522in}}{\pgfqpoint{0.017010in}{-0.017010in}}%
\pgfpathcurveto{\pgfqpoint{0.021522in}{-0.012499in}}{\pgfqpoint{0.024056in}{-0.006380in}}{\pgfqpoint{0.024056in}{0.000000in}}%
\pgfpathcurveto{\pgfqpoint{0.024056in}{0.006380in}}{\pgfqpoint{0.021522in}{0.012499in}}{\pgfqpoint{0.017010in}{0.017010in}}%
\pgfpathcurveto{\pgfqpoint{0.012499in}{0.021522in}}{\pgfqpoint{0.006380in}{0.024056in}}{\pgfqpoint{0.000000in}{0.024056in}}%
\pgfpathcurveto{\pgfqpoint{-0.006380in}{0.024056in}}{\pgfqpoint{-0.012499in}{0.021522in}}{\pgfqpoint{-0.017010in}{0.017010in}}%
\pgfpathcurveto{\pgfqpoint{-0.021522in}{0.012499in}}{\pgfqpoint{-0.024056in}{0.006380in}}{\pgfqpoint{-0.024056in}{0.000000in}}%
\pgfpathcurveto{\pgfqpoint{-0.024056in}{-0.006380in}}{\pgfqpoint{-0.021522in}{-0.012499in}}{\pgfqpoint{-0.017010in}{-0.017010in}}%
\pgfpathcurveto{\pgfqpoint{-0.012499in}{-0.021522in}}{\pgfqpoint{-0.006380in}{-0.024056in}}{\pgfqpoint{0.000000in}{-0.024056in}}%
\pgfpathlineto{\pgfqpoint{0.000000in}{-0.024056in}}%
\pgfpathclose%
\pgfusepath{fill}%
}%
\begin{pgfscope}%
\pgfsys@transformshift{6.210523in}{3.686033in}%
\pgfsys@useobject{currentmarker}{}%
\end{pgfscope}%
\end{pgfscope}%
\begin{pgfscope}%
\definecolor{textcolor}{rgb}{0.549020,0.549020,0.549020}%
\pgfsetstrokecolor{textcolor}%
\pgfsetfillcolor{textcolor}%
\pgftext[x=6.250948in,y=3.686033in,left,]{\color{textcolor}\sffamily\fontsize{7.000000}{8.400000}\selectfont 0.50}%
\end{pgfscope}%
\begin{pgfscope}%
\pgfpathrectangle{\pgfqpoint{5.635000in}{0.598000in}}{\pgfqpoint{1.155000in}{3.864000in}}%
\pgfusepath{clip}%
\pgfsetbuttcap%
\pgfsetroundjoin%
\definecolor{currentfill}{rgb}{0.435294,0.560784,0.709804}%
\pgfsetfillcolor{currentfill}%
\pgfsetlinewidth{0.000000pt}%
\definecolor{currentstroke}{rgb}{0.435294,0.560784,0.709804}%
\pgfsetstrokecolor{currentstroke}%
\pgfsetdash{}{0pt}%
\pgfsys@defobject{currentmarker}{\pgfqpoint{-0.024056in}{-0.024056in}}{\pgfqpoint{0.024056in}{0.024056in}}{%
\pgfpathmoveto{\pgfqpoint{0.000000in}{-0.024056in}}%
\pgfpathcurveto{\pgfqpoint{0.006380in}{-0.024056in}}{\pgfqpoint{0.012499in}{-0.021522in}}{\pgfqpoint{0.017010in}{-0.017010in}}%
\pgfpathcurveto{\pgfqpoint{0.021522in}{-0.012499in}}{\pgfqpoint{0.024056in}{-0.006380in}}{\pgfqpoint{0.024056in}{0.000000in}}%
\pgfpathcurveto{\pgfqpoint{0.024056in}{0.006380in}}{\pgfqpoint{0.021522in}{0.012499in}}{\pgfqpoint{0.017010in}{0.017010in}}%
\pgfpathcurveto{\pgfqpoint{0.012499in}{0.021522in}}{\pgfqpoint{0.006380in}{0.024056in}}{\pgfqpoint{0.000000in}{0.024056in}}%
\pgfpathcurveto{\pgfqpoint{-0.006380in}{0.024056in}}{\pgfqpoint{-0.012499in}{0.021522in}}{\pgfqpoint{-0.017010in}{0.017010in}}%
\pgfpathcurveto{\pgfqpoint{-0.021522in}{0.012499in}}{\pgfqpoint{-0.024056in}{0.006380in}}{\pgfqpoint{-0.024056in}{0.000000in}}%
\pgfpathcurveto{\pgfqpoint{-0.024056in}{-0.006380in}}{\pgfqpoint{-0.021522in}{-0.012499in}}{\pgfqpoint{-0.017010in}{-0.017010in}}%
\pgfpathcurveto{\pgfqpoint{-0.012499in}{-0.021522in}}{\pgfqpoint{-0.006380in}{-0.024056in}}{\pgfqpoint{0.000000in}{-0.024056in}}%
\pgfpathlineto{\pgfqpoint{0.000000in}{-0.024056in}}%
\pgfpathclose%
\pgfusepath{fill}%
}%
\begin{pgfscope}%
\pgfsys@transformshift{6.462372in}{3.527672in}%
\pgfsys@useobject{currentmarker}{}%
\end{pgfscope}%
\end{pgfscope}%
\begin{pgfscope}%
\definecolor{textcolor}{rgb}{0.549020,0.549020,0.549020}%
\pgfsetstrokecolor{textcolor}%
\pgfsetfillcolor{textcolor}%
\pgftext[x=6.502797in,y=3.527672in,left,]{\color{textcolor}\sffamily\fontsize{7.000000}{8.400000}\selectfont 0.72}%
\end{pgfscope}%
\begin{pgfscope}%
\pgfpathrectangle{\pgfqpoint{5.635000in}{0.598000in}}{\pgfqpoint{1.155000in}{3.864000in}}%
\pgfusepath{clip}%
\pgfsetbuttcap%
\pgfsetroundjoin%
\definecolor{currentfill}{rgb}{0.156863,0.356863,0.568627}%
\pgfsetfillcolor{currentfill}%
\pgfsetlinewidth{0.000000pt}%
\definecolor{currentstroke}{rgb}{0.156863,0.356863,0.568627}%
\pgfsetstrokecolor{currentstroke}%
\pgfsetdash{}{0pt}%
\pgfsys@defobject{currentmarker}{\pgfqpoint{-0.024056in}{-0.024056in}}{\pgfqpoint{0.024056in}{0.024056in}}{%
\pgfpathmoveto{\pgfqpoint{0.000000in}{-0.024056in}}%
\pgfpathcurveto{\pgfqpoint{0.006380in}{-0.024056in}}{\pgfqpoint{0.012499in}{-0.021522in}}{\pgfqpoint{0.017010in}{-0.017010in}}%
\pgfpathcurveto{\pgfqpoint{0.021522in}{-0.012499in}}{\pgfqpoint{0.024056in}{-0.006380in}}{\pgfqpoint{0.024056in}{0.000000in}}%
\pgfpathcurveto{\pgfqpoint{0.024056in}{0.006380in}}{\pgfqpoint{0.021522in}{0.012499in}}{\pgfqpoint{0.017010in}{0.017010in}}%
\pgfpathcurveto{\pgfqpoint{0.012499in}{0.021522in}}{\pgfqpoint{0.006380in}{0.024056in}}{\pgfqpoint{0.000000in}{0.024056in}}%
\pgfpathcurveto{\pgfqpoint{-0.006380in}{0.024056in}}{\pgfqpoint{-0.012499in}{0.021522in}}{\pgfqpoint{-0.017010in}{0.017010in}}%
\pgfpathcurveto{\pgfqpoint{-0.021522in}{0.012499in}}{\pgfqpoint{-0.024056in}{0.006380in}}{\pgfqpoint{-0.024056in}{0.000000in}}%
\pgfpathcurveto{\pgfqpoint{-0.024056in}{-0.006380in}}{\pgfqpoint{-0.021522in}{-0.012499in}}{\pgfqpoint{-0.017010in}{-0.017010in}}%
\pgfpathcurveto{\pgfqpoint{-0.012499in}{-0.021522in}}{\pgfqpoint{-0.006380in}{-0.024056in}}{\pgfqpoint{0.000000in}{-0.024056in}}%
\pgfpathlineto{\pgfqpoint{0.000000in}{-0.024056in}}%
\pgfpathclose%
\pgfusepath{fill}%
}%
\begin{pgfscope}%
\pgfsys@transformshift{6.456363in}{3.369311in}%
\pgfsys@useobject{currentmarker}{}%
\end{pgfscope}%
\end{pgfscope}%
\begin{pgfscope}%
\definecolor{textcolor}{rgb}{0.549020,0.549020,0.549020}%
\pgfsetstrokecolor{textcolor}%
\pgfsetfillcolor{textcolor}%
\pgftext[x=6.496788in,y=3.369311in,left,]{\color{textcolor}\sffamily\fontsize{7.000000}{8.400000}\selectfont 0.71}%
\end{pgfscope}%
\begin{pgfscope}%
\pgfpathrectangle{\pgfqpoint{5.635000in}{0.598000in}}{\pgfqpoint{1.155000in}{3.864000in}}%
\pgfusepath{clip}%
\pgfsetbuttcap%
\pgfsetroundjoin%
\definecolor{currentfill}{rgb}{0.435294,0.560784,0.709804}%
\pgfsetfillcolor{currentfill}%
\pgfsetlinewidth{0.000000pt}%
\definecolor{currentstroke}{rgb}{0.435294,0.560784,0.709804}%
\pgfsetstrokecolor{currentstroke}%
\pgfsetdash{}{0pt}%
\pgfsys@defobject{currentmarker}{\pgfqpoint{-0.024056in}{-0.024056in}}{\pgfqpoint{0.024056in}{0.024056in}}{%
\pgfpathmoveto{\pgfqpoint{0.000000in}{-0.024056in}}%
\pgfpathcurveto{\pgfqpoint{0.006380in}{-0.024056in}}{\pgfqpoint{0.012499in}{-0.021522in}}{\pgfqpoint{0.017010in}{-0.017010in}}%
\pgfpathcurveto{\pgfqpoint{0.021522in}{-0.012499in}}{\pgfqpoint{0.024056in}{-0.006380in}}{\pgfqpoint{0.024056in}{0.000000in}}%
\pgfpathcurveto{\pgfqpoint{0.024056in}{0.006380in}}{\pgfqpoint{0.021522in}{0.012499in}}{\pgfqpoint{0.017010in}{0.017010in}}%
\pgfpathcurveto{\pgfqpoint{0.012499in}{0.021522in}}{\pgfqpoint{0.006380in}{0.024056in}}{\pgfqpoint{0.000000in}{0.024056in}}%
\pgfpathcurveto{\pgfqpoint{-0.006380in}{0.024056in}}{\pgfqpoint{-0.012499in}{0.021522in}}{\pgfqpoint{-0.017010in}{0.017010in}}%
\pgfpathcurveto{\pgfqpoint{-0.021522in}{0.012499in}}{\pgfqpoint{-0.024056in}{0.006380in}}{\pgfqpoint{-0.024056in}{0.000000in}}%
\pgfpathcurveto{\pgfqpoint{-0.024056in}{-0.006380in}}{\pgfqpoint{-0.021522in}{-0.012499in}}{\pgfqpoint{-0.017010in}{-0.017010in}}%
\pgfpathcurveto{\pgfqpoint{-0.012499in}{-0.021522in}}{\pgfqpoint{-0.006380in}{-0.024056in}}{\pgfqpoint{0.000000in}{-0.024056in}}%
\pgfpathlineto{\pgfqpoint{0.000000in}{-0.024056in}}%
\pgfpathclose%
\pgfusepath{fill}%
}%
\begin{pgfscope}%
\pgfsys@transformshift{6.445789in}{3.210951in}%
\pgfsys@useobject{currentmarker}{}%
\end{pgfscope}%
\end{pgfscope}%
\begin{pgfscope}%
\definecolor{textcolor}{rgb}{0.549020,0.549020,0.549020}%
\pgfsetstrokecolor{textcolor}%
\pgfsetfillcolor{textcolor}%
\pgftext[x=6.486214in,y=3.210951in,left,]{\color{textcolor}\sffamily\fontsize{7.000000}{8.400000}\selectfont 0.70}%
\end{pgfscope}%
\begin{pgfscope}%
\pgfpathrectangle{\pgfqpoint{5.635000in}{0.598000in}}{\pgfqpoint{1.155000in}{3.864000in}}%
\pgfusepath{clip}%
\pgfsetbuttcap%
\pgfsetroundjoin%
\definecolor{currentfill}{rgb}{0.435294,0.560784,0.709804}%
\pgfsetfillcolor{currentfill}%
\pgfsetlinewidth{0.000000pt}%
\definecolor{currentstroke}{rgb}{0.435294,0.560784,0.709804}%
\pgfsetstrokecolor{currentstroke}%
\pgfsetdash{}{0pt}%
\pgfsys@defobject{currentmarker}{\pgfqpoint{-0.024056in}{-0.024056in}}{\pgfqpoint{0.024056in}{0.024056in}}{%
\pgfpathmoveto{\pgfqpoint{0.000000in}{-0.024056in}}%
\pgfpathcurveto{\pgfqpoint{0.006380in}{-0.024056in}}{\pgfqpoint{0.012499in}{-0.021522in}}{\pgfqpoint{0.017010in}{-0.017010in}}%
\pgfpathcurveto{\pgfqpoint{0.021522in}{-0.012499in}}{\pgfqpoint{0.024056in}{-0.006380in}}{\pgfqpoint{0.024056in}{0.000000in}}%
\pgfpathcurveto{\pgfqpoint{0.024056in}{0.006380in}}{\pgfqpoint{0.021522in}{0.012499in}}{\pgfqpoint{0.017010in}{0.017010in}}%
\pgfpathcurveto{\pgfqpoint{0.012499in}{0.021522in}}{\pgfqpoint{0.006380in}{0.024056in}}{\pgfqpoint{0.000000in}{0.024056in}}%
\pgfpathcurveto{\pgfqpoint{-0.006380in}{0.024056in}}{\pgfqpoint{-0.012499in}{0.021522in}}{\pgfqpoint{-0.017010in}{0.017010in}}%
\pgfpathcurveto{\pgfqpoint{-0.021522in}{0.012499in}}{\pgfqpoint{-0.024056in}{0.006380in}}{\pgfqpoint{-0.024056in}{0.000000in}}%
\pgfpathcurveto{\pgfqpoint{-0.024056in}{-0.006380in}}{\pgfqpoint{-0.021522in}{-0.012499in}}{\pgfqpoint{-0.017010in}{-0.017010in}}%
\pgfpathcurveto{\pgfqpoint{-0.012499in}{-0.021522in}}{\pgfqpoint{-0.006380in}{-0.024056in}}{\pgfqpoint{0.000000in}{-0.024056in}}%
\pgfpathlineto{\pgfqpoint{0.000000in}{-0.024056in}}%
\pgfpathclose%
\pgfusepath{fill}%
}%
\begin{pgfscope}%
\pgfsys@transformshift{6.460826in}{3.052590in}%
\pgfsys@useobject{currentmarker}{}%
\end{pgfscope}%
\end{pgfscope}%
\begin{pgfscope}%
\definecolor{textcolor}{rgb}{0.549020,0.549020,0.549020}%
\pgfsetstrokecolor{textcolor}%
\pgfsetfillcolor{textcolor}%
\pgftext[x=6.501251in,y=3.052590in,left,]{\color{textcolor}\sffamily\fontsize{7.000000}{8.400000}\selectfont 0.72}%
\end{pgfscope}%
\begin{pgfscope}%
\pgfpathrectangle{\pgfqpoint{5.635000in}{0.598000in}}{\pgfqpoint{1.155000in}{3.864000in}}%
\pgfusepath{clip}%
\pgfsetbuttcap%
\pgfsetroundjoin%
\definecolor{currentfill}{rgb}{0.435294,0.560784,0.709804}%
\pgfsetfillcolor{currentfill}%
\pgfsetlinewidth{0.000000pt}%
\definecolor{currentstroke}{rgb}{0.435294,0.560784,0.709804}%
\pgfsetstrokecolor{currentstroke}%
\pgfsetdash{}{0pt}%
\pgfsys@defobject{currentmarker}{\pgfqpoint{-0.024056in}{-0.024056in}}{\pgfqpoint{0.024056in}{0.024056in}}{%
\pgfpathmoveto{\pgfqpoint{0.000000in}{-0.024056in}}%
\pgfpathcurveto{\pgfqpoint{0.006380in}{-0.024056in}}{\pgfqpoint{0.012499in}{-0.021522in}}{\pgfqpoint{0.017010in}{-0.017010in}}%
\pgfpathcurveto{\pgfqpoint{0.021522in}{-0.012499in}}{\pgfqpoint{0.024056in}{-0.006380in}}{\pgfqpoint{0.024056in}{0.000000in}}%
\pgfpathcurveto{\pgfqpoint{0.024056in}{0.006380in}}{\pgfqpoint{0.021522in}{0.012499in}}{\pgfqpoint{0.017010in}{0.017010in}}%
\pgfpathcurveto{\pgfqpoint{0.012499in}{0.021522in}}{\pgfqpoint{0.006380in}{0.024056in}}{\pgfqpoint{0.000000in}{0.024056in}}%
\pgfpathcurveto{\pgfqpoint{-0.006380in}{0.024056in}}{\pgfqpoint{-0.012499in}{0.021522in}}{\pgfqpoint{-0.017010in}{0.017010in}}%
\pgfpathcurveto{\pgfqpoint{-0.021522in}{0.012499in}}{\pgfqpoint{-0.024056in}{0.006380in}}{\pgfqpoint{-0.024056in}{0.000000in}}%
\pgfpathcurveto{\pgfqpoint{-0.024056in}{-0.006380in}}{\pgfqpoint{-0.021522in}{-0.012499in}}{\pgfqpoint{-0.017010in}{-0.017010in}}%
\pgfpathcurveto{\pgfqpoint{-0.012499in}{-0.021522in}}{\pgfqpoint{-0.006380in}{-0.024056in}}{\pgfqpoint{0.000000in}{-0.024056in}}%
\pgfpathlineto{\pgfqpoint{0.000000in}{-0.024056in}}%
\pgfpathclose%
\pgfusepath{fill}%
}%
\begin{pgfscope}%
\pgfsys@transformshift{6.452999in}{2.894230in}%
\pgfsys@useobject{currentmarker}{}%
\end{pgfscope}%
\end{pgfscope}%
\begin{pgfscope}%
\definecolor{textcolor}{rgb}{0.549020,0.549020,0.549020}%
\pgfsetstrokecolor{textcolor}%
\pgfsetfillcolor{textcolor}%
\pgftext[x=6.493424in,y=2.894230in,left,]{\color{textcolor}\sffamily\fontsize{7.000000}{8.400000}\selectfont 0.71}%
\end{pgfscope}%
\begin{pgfscope}%
\pgfpathrectangle{\pgfqpoint{5.635000in}{0.598000in}}{\pgfqpoint{1.155000in}{3.864000in}}%
\pgfusepath{clip}%
\pgfsetbuttcap%
\pgfsetroundjoin%
\definecolor{currentfill}{rgb}{0.435294,0.560784,0.709804}%
\pgfsetfillcolor{currentfill}%
\pgfsetlinewidth{0.000000pt}%
\definecolor{currentstroke}{rgb}{0.435294,0.560784,0.709804}%
\pgfsetstrokecolor{currentstroke}%
\pgfsetdash{}{0pt}%
\pgfsys@defobject{currentmarker}{\pgfqpoint{-0.024056in}{-0.024056in}}{\pgfqpoint{0.024056in}{0.024056in}}{%
\pgfpathmoveto{\pgfqpoint{0.000000in}{-0.024056in}}%
\pgfpathcurveto{\pgfqpoint{0.006380in}{-0.024056in}}{\pgfqpoint{0.012499in}{-0.021522in}}{\pgfqpoint{0.017010in}{-0.017010in}}%
\pgfpathcurveto{\pgfqpoint{0.021522in}{-0.012499in}}{\pgfqpoint{0.024056in}{-0.006380in}}{\pgfqpoint{0.024056in}{0.000000in}}%
\pgfpathcurveto{\pgfqpoint{0.024056in}{0.006380in}}{\pgfqpoint{0.021522in}{0.012499in}}{\pgfqpoint{0.017010in}{0.017010in}}%
\pgfpathcurveto{\pgfqpoint{0.012499in}{0.021522in}}{\pgfqpoint{0.006380in}{0.024056in}}{\pgfqpoint{0.000000in}{0.024056in}}%
\pgfpathcurveto{\pgfqpoint{-0.006380in}{0.024056in}}{\pgfqpoint{-0.012499in}{0.021522in}}{\pgfqpoint{-0.017010in}{0.017010in}}%
\pgfpathcurveto{\pgfqpoint{-0.021522in}{0.012499in}}{\pgfqpoint{-0.024056in}{0.006380in}}{\pgfqpoint{-0.024056in}{0.000000in}}%
\pgfpathcurveto{\pgfqpoint{-0.024056in}{-0.006380in}}{\pgfqpoint{-0.021522in}{-0.012499in}}{\pgfqpoint{-0.017010in}{-0.017010in}}%
\pgfpathcurveto{\pgfqpoint{-0.012499in}{-0.021522in}}{\pgfqpoint{-0.006380in}{-0.024056in}}{\pgfqpoint{0.000000in}{-0.024056in}}%
\pgfpathlineto{\pgfqpoint{0.000000in}{-0.024056in}}%
\pgfpathclose%
\pgfusepath{fill}%
}%
\begin{pgfscope}%
\pgfsys@transformshift{6.467247in}{2.735869in}%
\pgfsys@useobject{currentmarker}{}%
\end{pgfscope}%
\end{pgfscope}%
\begin{pgfscope}%
\definecolor{textcolor}{rgb}{0.549020,0.549020,0.549020}%
\pgfsetstrokecolor{textcolor}%
\pgfsetfillcolor{textcolor}%
\pgftext[x=6.507672in,y=2.735869in,left,]{\color{textcolor}\sffamily\fontsize{7.000000}{8.400000}\selectfont 0.72}%
\end{pgfscope}%
\begin{pgfscope}%
\pgfpathrectangle{\pgfqpoint{5.635000in}{0.598000in}}{\pgfqpoint{1.155000in}{3.864000in}}%
\pgfusepath{clip}%
\pgfsetbuttcap%
\pgfsetroundjoin%
\definecolor{currentfill}{rgb}{0.156863,0.356863,0.568627}%
\pgfsetfillcolor{currentfill}%
\pgfsetlinewidth{0.000000pt}%
\definecolor{currentstroke}{rgb}{0.156863,0.356863,0.568627}%
\pgfsetstrokecolor{currentstroke}%
\pgfsetdash{}{0pt}%
\pgfsys@defobject{currentmarker}{\pgfqpoint{-0.024056in}{-0.024056in}}{\pgfqpoint{0.024056in}{0.024056in}}{%
\pgfpathmoveto{\pgfqpoint{0.000000in}{-0.024056in}}%
\pgfpathcurveto{\pgfqpoint{0.006380in}{-0.024056in}}{\pgfqpoint{0.012499in}{-0.021522in}}{\pgfqpoint{0.017010in}{-0.017010in}}%
\pgfpathcurveto{\pgfqpoint{0.021522in}{-0.012499in}}{\pgfqpoint{0.024056in}{-0.006380in}}{\pgfqpoint{0.024056in}{0.000000in}}%
\pgfpathcurveto{\pgfqpoint{0.024056in}{0.006380in}}{\pgfqpoint{0.021522in}{0.012499in}}{\pgfqpoint{0.017010in}{0.017010in}}%
\pgfpathcurveto{\pgfqpoint{0.012499in}{0.021522in}}{\pgfqpoint{0.006380in}{0.024056in}}{\pgfqpoint{0.000000in}{0.024056in}}%
\pgfpathcurveto{\pgfqpoint{-0.006380in}{0.024056in}}{\pgfqpoint{-0.012499in}{0.021522in}}{\pgfqpoint{-0.017010in}{0.017010in}}%
\pgfpathcurveto{\pgfqpoint{-0.021522in}{0.012499in}}{\pgfqpoint{-0.024056in}{0.006380in}}{\pgfqpoint{-0.024056in}{0.000000in}}%
\pgfpathcurveto{\pgfqpoint{-0.024056in}{-0.006380in}}{\pgfqpoint{-0.021522in}{-0.012499in}}{\pgfqpoint{-0.017010in}{-0.017010in}}%
\pgfpathcurveto{\pgfqpoint{-0.012499in}{-0.021522in}}{\pgfqpoint{-0.006380in}{-0.024056in}}{\pgfqpoint{0.000000in}{-0.024056in}}%
\pgfpathlineto{\pgfqpoint{0.000000in}{-0.024056in}}%
\pgfpathclose%
\pgfusepath{fill}%
}%
\begin{pgfscope}%
\pgfsys@transformshift{6.472882in}{2.577508in}%
\pgfsys@useobject{currentmarker}{}%
\end{pgfscope}%
\end{pgfscope}%
\begin{pgfscope}%
\definecolor{textcolor}{rgb}{0.549020,0.549020,0.549020}%
\pgfsetstrokecolor{textcolor}%
\pgfsetfillcolor{textcolor}%
\pgftext[x=6.513307in,y=2.577508in,left,]{\color{textcolor}\sffamily\fontsize{7.000000}{8.400000}\selectfont 0.73}%
\end{pgfscope}%
\begin{pgfscope}%
\pgfpathrectangle{\pgfqpoint{5.635000in}{0.598000in}}{\pgfqpoint{1.155000in}{3.864000in}}%
\pgfusepath{clip}%
\pgfsetbuttcap%
\pgfsetroundjoin%
\definecolor{currentfill}{rgb}{0.156863,0.356863,0.568627}%
\pgfsetfillcolor{currentfill}%
\pgfsetlinewidth{0.000000pt}%
\definecolor{currentstroke}{rgb}{0.156863,0.356863,0.568627}%
\pgfsetstrokecolor{currentstroke}%
\pgfsetdash{}{0pt}%
\pgfsys@defobject{currentmarker}{\pgfqpoint{-0.024056in}{-0.024056in}}{\pgfqpoint{0.024056in}{0.024056in}}{%
\pgfpathmoveto{\pgfqpoint{0.000000in}{-0.024056in}}%
\pgfpathcurveto{\pgfqpoint{0.006380in}{-0.024056in}}{\pgfqpoint{0.012499in}{-0.021522in}}{\pgfqpoint{0.017010in}{-0.017010in}}%
\pgfpathcurveto{\pgfqpoint{0.021522in}{-0.012499in}}{\pgfqpoint{0.024056in}{-0.006380in}}{\pgfqpoint{0.024056in}{0.000000in}}%
\pgfpathcurveto{\pgfqpoint{0.024056in}{0.006380in}}{\pgfqpoint{0.021522in}{0.012499in}}{\pgfqpoint{0.017010in}{0.017010in}}%
\pgfpathcurveto{\pgfqpoint{0.012499in}{0.021522in}}{\pgfqpoint{0.006380in}{0.024056in}}{\pgfqpoint{0.000000in}{0.024056in}}%
\pgfpathcurveto{\pgfqpoint{-0.006380in}{0.024056in}}{\pgfqpoint{-0.012499in}{0.021522in}}{\pgfqpoint{-0.017010in}{0.017010in}}%
\pgfpathcurveto{\pgfqpoint{-0.021522in}{0.012499in}}{\pgfqpoint{-0.024056in}{0.006380in}}{\pgfqpoint{-0.024056in}{0.000000in}}%
\pgfpathcurveto{\pgfqpoint{-0.024056in}{-0.006380in}}{\pgfqpoint{-0.021522in}{-0.012499in}}{\pgfqpoint{-0.017010in}{-0.017010in}}%
\pgfpathcurveto{\pgfqpoint{-0.012499in}{-0.021522in}}{\pgfqpoint{-0.006380in}{-0.024056in}}{\pgfqpoint{0.000000in}{-0.024056in}}%
\pgfpathlineto{\pgfqpoint{0.000000in}{-0.024056in}}%
\pgfpathclose%
\pgfusepath{fill}%
}%
\begin{pgfscope}%
\pgfsys@transformshift{6.198981in}{2.189525in}%
\pgfsys@useobject{currentmarker}{}%
\end{pgfscope}%
\end{pgfscope}%
\begin{pgfscope}%
\definecolor{textcolor}{rgb}{0.549020,0.549020,0.549020}%
\pgfsetstrokecolor{textcolor}%
\pgfsetfillcolor{textcolor}%
\pgftext[x=6.239406in,y=2.189525in,left,]{\color{textcolor}\sffamily\fontsize{7.000000}{8.400000}\selectfont 0.49}%
\end{pgfscope}%
\begin{pgfscope}%
\pgfpathrectangle{\pgfqpoint{5.635000in}{0.598000in}}{\pgfqpoint{1.155000in}{3.864000in}}%
\pgfusepath{clip}%
\pgfsetbuttcap%
\pgfsetroundjoin%
\definecolor{currentfill}{rgb}{0.156863,0.356863,0.568627}%
\pgfsetfillcolor{currentfill}%
\pgfsetlinewidth{0.000000pt}%
\definecolor{currentstroke}{rgb}{0.156863,0.356863,0.568627}%
\pgfsetstrokecolor{currentstroke}%
\pgfsetdash{}{0pt}%
\pgfsys@defobject{currentmarker}{\pgfqpoint{-0.024056in}{-0.024056in}}{\pgfqpoint{0.024056in}{0.024056in}}{%
\pgfpathmoveto{\pgfqpoint{0.000000in}{-0.024056in}}%
\pgfpathcurveto{\pgfqpoint{0.006380in}{-0.024056in}}{\pgfqpoint{0.012499in}{-0.021522in}}{\pgfqpoint{0.017010in}{-0.017010in}}%
\pgfpathcurveto{\pgfqpoint{0.021522in}{-0.012499in}}{\pgfqpoint{0.024056in}{-0.006380in}}{\pgfqpoint{0.024056in}{0.000000in}}%
\pgfpathcurveto{\pgfqpoint{0.024056in}{0.006380in}}{\pgfqpoint{0.021522in}{0.012499in}}{\pgfqpoint{0.017010in}{0.017010in}}%
\pgfpathcurveto{\pgfqpoint{0.012499in}{0.021522in}}{\pgfqpoint{0.006380in}{0.024056in}}{\pgfqpoint{0.000000in}{0.024056in}}%
\pgfpathcurveto{\pgfqpoint{-0.006380in}{0.024056in}}{\pgfqpoint{-0.012499in}{0.021522in}}{\pgfqpoint{-0.017010in}{0.017010in}}%
\pgfpathcurveto{\pgfqpoint{-0.021522in}{0.012499in}}{\pgfqpoint{-0.024056in}{0.006380in}}{\pgfqpoint{-0.024056in}{0.000000in}}%
\pgfpathcurveto{\pgfqpoint{-0.024056in}{-0.006380in}}{\pgfqpoint{-0.021522in}{-0.012499in}}{\pgfqpoint{-0.017010in}{-0.017010in}}%
\pgfpathcurveto{\pgfqpoint{-0.012499in}{-0.021522in}}{\pgfqpoint{-0.006380in}{-0.024056in}}{\pgfqpoint{0.000000in}{-0.024056in}}%
\pgfpathlineto{\pgfqpoint{0.000000in}{-0.024056in}}%
\pgfpathclose%
\pgfusepath{fill}%
}%
\begin{pgfscope}%
\pgfsys@transformshift{6.425820in}{2.031164in}%
\pgfsys@useobject{currentmarker}{}%
\end{pgfscope}%
\end{pgfscope}%
\begin{pgfscope}%
\definecolor{textcolor}{rgb}{0.549020,0.549020,0.549020}%
\pgfsetstrokecolor{textcolor}%
\pgfsetfillcolor{textcolor}%
\pgftext[x=6.466245in,y=2.031164in,left,]{\color{textcolor}\sffamily\fontsize{7.000000}{8.400000}\selectfont 0.68}%
\end{pgfscope}%
\begin{pgfscope}%
\pgfpathrectangle{\pgfqpoint{5.635000in}{0.598000in}}{\pgfqpoint{1.155000in}{3.864000in}}%
\pgfusepath{clip}%
\pgfsetbuttcap%
\pgfsetroundjoin%
\definecolor{currentfill}{rgb}{0.101961,0.101961,0.101961}%
\pgfsetfillcolor{currentfill}%
\pgfsetlinewidth{0.000000pt}%
\definecolor{currentstroke}{rgb}{0.101961,0.101961,0.101961}%
\pgfsetstrokecolor{currentstroke}%
\pgfsetdash{}{0pt}%
\pgfsys@defobject{currentmarker}{\pgfqpoint{-0.032572in}{-0.032572in}}{\pgfqpoint{0.032572in}{0.032572in}}{%
\pgfpathmoveto{\pgfqpoint{-0.016286in}{-0.032572in}}%
\pgfpathlineto{\pgfqpoint{0.000000in}{-0.016286in}}%
\pgfpathlineto{\pgfqpoint{0.016286in}{-0.032572in}}%
\pgfpathlineto{\pgfqpoint{0.032572in}{-0.016286in}}%
\pgfpathlineto{\pgfqpoint{0.016286in}{0.000000in}}%
\pgfpathlineto{\pgfqpoint{0.032572in}{0.016286in}}%
\pgfpathlineto{\pgfqpoint{0.016286in}{0.032572in}}%
\pgfpathlineto{\pgfqpoint{0.000000in}{0.016286in}}%
\pgfpathlineto{\pgfqpoint{-0.016286in}{0.032572in}}%
\pgfpathlineto{\pgfqpoint{-0.032572in}{0.016286in}}%
\pgfpathlineto{\pgfqpoint{-0.016286in}{0.000000in}}%
\pgfpathlineto{\pgfqpoint{-0.032572in}{-0.016286in}}%
\pgfpathlineto{\pgfqpoint{-0.016286in}{-0.032572in}}%
\pgfpathclose%
\pgfusepath{fill}%
}%
\begin{pgfscope}%
\pgfsys@transformshift{6.213448in}{1.872803in}%
\pgfsys@useobject{currentmarker}{}%
\end{pgfscope}%
\end{pgfscope}%
\begin{pgfscope}%
\definecolor{textcolor}{rgb}{0.549020,0.549020,0.549020}%
\pgfsetstrokecolor{textcolor}%
\pgfsetfillcolor{textcolor}%
\pgftext[x=6.253873in,y=1.872803in,left,]{\color{textcolor}\sffamily\fontsize{7.000000}{8.400000}\selectfont 0.50}%
\end{pgfscope}%
\begin{pgfscope}%
\pgfpathrectangle{\pgfqpoint{5.635000in}{0.598000in}}{\pgfqpoint{1.155000in}{3.864000in}}%
\pgfusepath{clip}%
\pgfsetbuttcap%
\pgfsetroundjoin%
\definecolor{currentfill}{rgb}{0.101961,0.101961,0.101961}%
\pgfsetfillcolor{currentfill}%
\pgfsetlinewidth{0.000000pt}%
\definecolor{currentstroke}{rgb}{0.101961,0.101961,0.101961}%
\pgfsetstrokecolor{currentstroke}%
\pgfsetdash{}{0pt}%
\pgfsys@defobject{currentmarker}{\pgfqpoint{-0.032572in}{-0.032572in}}{\pgfqpoint{0.032572in}{0.032572in}}{%
\pgfpathmoveto{\pgfqpoint{-0.016286in}{-0.032572in}}%
\pgfpathlineto{\pgfqpoint{0.000000in}{-0.016286in}}%
\pgfpathlineto{\pgfqpoint{0.016286in}{-0.032572in}}%
\pgfpathlineto{\pgfqpoint{0.032572in}{-0.016286in}}%
\pgfpathlineto{\pgfqpoint{0.016286in}{0.000000in}}%
\pgfpathlineto{\pgfqpoint{0.032572in}{0.016286in}}%
\pgfpathlineto{\pgfqpoint{0.016286in}{0.032572in}}%
\pgfpathlineto{\pgfqpoint{0.000000in}{0.016286in}}%
\pgfpathlineto{\pgfqpoint{-0.016286in}{0.032572in}}%
\pgfpathlineto{\pgfqpoint{-0.032572in}{0.016286in}}%
\pgfpathlineto{\pgfqpoint{-0.016286in}{0.000000in}}%
\pgfpathlineto{\pgfqpoint{-0.032572in}{-0.016286in}}%
\pgfpathlineto{\pgfqpoint{-0.016286in}{-0.032572in}}%
\pgfpathclose%
\pgfusepath{fill}%
}%
\begin{pgfscope}%
\pgfsys@transformshift{6.214082in}{1.714443in}%
\pgfsys@useobject{currentmarker}{}%
\end{pgfscope}%
\end{pgfscope}%
\begin{pgfscope}%
\definecolor{textcolor}{rgb}{0.549020,0.549020,0.549020}%
\pgfsetstrokecolor{textcolor}%
\pgfsetfillcolor{textcolor}%
\pgftext[x=6.254507in,y=1.714443in,left,]{\color{textcolor}\sffamily\fontsize{7.000000}{8.400000}\selectfont 0.50}%
\end{pgfscope}%
\begin{pgfscope}%
\pgfpathrectangle{\pgfqpoint{5.635000in}{0.598000in}}{\pgfqpoint{1.155000in}{3.864000in}}%
\pgfusepath{clip}%
\pgfsetbuttcap%
\pgfsetroundjoin%
\definecolor{currentfill}{rgb}{0.156863,0.356863,0.568627}%
\pgfsetfillcolor{currentfill}%
\pgfsetlinewidth{0.000000pt}%
\definecolor{currentstroke}{rgb}{0.156863,0.356863,0.568627}%
\pgfsetstrokecolor{currentstroke}%
\pgfsetdash{}{0pt}%
\pgfsys@defobject{currentmarker}{\pgfqpoint{-0.024056in}{-0.024056in}}{\pgfqpoint{0.024056in}{0.024056in}}{%
\pgfpathmoveto{\pgfqpoint{0.000000in}{-0.024056in}}%
\pgfpathcurveto{\pgfqpoint{0.006380in}{-0.024056in}}{\pgfqpoint{0.012499in}{-0.021522in}}{\pgfqpoint{0.017010in}{-0.017010in}}%
\pgfpathcurveto{\pgfqpoint{0.021522in}{-0.012499in}}{\pgfqpoint{0.024056in}{-0.006380in}}{\pgfqpoint{0.024056in}{0.000000in}}%
\pgfpathcurveto{\pgfqpoint{0.024056in}{0.006380in}}{\pgfqpoint{0.021522in}{0.012499in}}{\pgfqpoint{0.017010in}{0.017010in}}%
\pgfpathcurveto{\pgfqpoint{0.012499in}{0.021522in}}{\pgfqpoint{0.006380in}{0.024056in}}{\pgfqpoint{0.000000in}{0.024056in}}%
\pgfpathcurveto{\pgfqpoint{-0.006380in}{0.024056in}}{\pgfqpoint{-0.012499in}{0.021522in}}{\pgfqpoint{-0.017010in}{0.017010in}}%
\pgfpathcurveto{\pgfqpoint{-0.021522in}{0.012499in}}{\pgfqpoint{-0.024056in}{0.006380in}}{\pgfqpoint{-0.024056in}{0.000000in}}%
\pgfpathcurveto{\pgfqpoint{-0.024056in}{-0.006380in}}{\pgfqpoint{-0.021522in}{-0.012499in}}{\pgfqpoint{-0.017010in}{-0.017010in}}%
\pgfpathcurveto{\pgfqpoint{-0.012499in}{-0.021522in}}{\pgfqpoint{-0.006380in}{-0.024056in}}{\pgfqpoint{0.000000in}{-0.024056in}}%
\pgfpathlineto{\pgfqpoint{0.000000in}{-0.024056in}}%
\pgfpathclose%
\pgfusepath{fill}%
}%
\begin{pgfscope}%
\pgfsys@transformshift{6.298207in}{1.326459in}%
\pgfsys@useobject{currentmarker}{}%
\end{pgfscope}%
\end{pgfscope}%
\begin{pgfscope}%
\definecolor{textcolor}{rgb}{0.549020,0.549020,0.549020}%
\pgfsetstrokecolor{textcolor}%
\pgfsetfillcolor{textcolor}%
\pgftext[x=6.338632in,y=1.326459in,left,]{\color{textcolor}\sffamily\fontsize{7.000000}{8.400000}\selectfont 0.57}%
\end{pgfscope}%
\begin{pgfscope}%
\pgfpathrectangle{\pgfqpoint{5.635000in}{0.598000in}}{\pgfqpoint{1.155000in}{3.864000in}}%
\pgfusepath{clip}%
\pgfsetbuttcap%
\pgfsetroundjoin%
\definecolor{currentfill}{rgb}{0.156863,0.356863,0.568627}%
\pgfsetfillcolor{currentfill}%
\pgfsetlinewidth{0.000000pt}%
\definecolor{currentstroke}{rgb}{0.156863,0.356863,0.568627}%
\pgfsetstrokecolor{currentstroke}%
\pgfsetdash{}{0pt}%
\pgfsys@defobject{currentmarker}{\pgfqpoint{-0.024056in}{-0.024056in}}{\pgfqpoint{0.024056in}{0.024056in}}{%
\pgfpathmoveto{\pgfqpoint{0.000000in}{-0.024056in}}%
\pgfpathcurveto{\pgfqpoint{0.006380in}{-0.024056in}}{\pgfqpoint{0.012499in}{-0.021522in}}{\pgfqpoint{0.017010in}{-0.017010in}}%
\pgfpathcurveto{\pgfqpoint{0.021522in}{-0.012499in}}{\pgfqpoint{0.024056in}{-0.006380in}}{\pgfqpoint{0.024056in}{0.000000in}}%
\pgfpathcurveto{\pgfqpoint{0.024056in}{0.006380in}}{\pgfqpoint{0.021522in}{0.012499in}}{\pgfqpoint{0.017010in}{0.017010in}}%
\pgfpathcurveto{\pgfqpoint{0.012499in}{0.021522in}}{\pgfqpoint{0.006380in}{0.024056in}}{\pgfqpoint{0.000000in}{0.024056in}}%
\pgfpathcurveto{\pgfqpoint{-0.006380in}{0.024056in}}{\pgfqpoint{-0.012499in}{0.021522in}}{\pgfqpoint{-0.017010in}{0.017010in}}%
\pgfpathcurveto{\pgfqpoint{-0.021522in}{0.012499in}}{\pgfqpoint{-0.024056in}{0.006380in}}{\pgfqpoint{-0.024056in}{0.000000in}}%
\pgfpathcurveto{\pgfqpoint{-0.024056in}{-0.006380in}}{\pgfqpoint{-0.021522in}{-0.012499in}}{\pgfqpoint{-0.017010in}{-0.017010in}}%
\pgfpathcurveto{\pgfqpoint{-0.012499in}{-0.021522in}}{\pgfqpoint{-0.006380in}{-0.024056in}}{\pgfqpoint{0.000000in}{-0.024056in}}%
\pgfpathlineto{\pgfqpoint{0.000000in}{-0.024056in}}%
\pgfpathclose%
\pgfusepath{fill}%
}%
\begin{pgfscope}%
\pgfsys@transformshift{6.381723in}{1.168098in}%
\pgfsys@useobject{currentmarker}{}%
\end{pgfscope}%
\end{pgfscope}%
\begin{pgfscope}%
\definecolor{textcolor}{rgb}{0.549020,0.549020,0.549020}%
\pgfsetstrokecolor{textcolor}%
\pgfsetfillcolor{textcolor}%
\pgftext[x=6.422148in,y=1.168098in,left,]{\color{textcolor}\sffamily\fontsize{7.000000}{8.400000}\selectfont 0.65}%
\end{pgfscope}%
\begin{pgfscope}%
\pgfpathrectangle{\pgfqpoint{5.635000in}{0.598000in}}{\pgfqpoint{1.155000in}{3.864000in}}%
\pgfusepath{clip}%
\pgfsetbuttcap%
\pgfsetroundjoin%
\definecolor{currentfill}{rgb}{0.156863,0.356863,0.568627}%
\pgfsetfillcolor{currentfill}%
\pgfsetlinewidth{0.000000pt}%
\definecolor{currentstroke}{rgb}{0.156863,0.356863,0.568627}%
\pgfsetstrokecolor{currentstroke}%
\pgfsetdash{}{0pt}%
\pgfsys@defobject{currentmarker}{\pgfqpoint{-0.024056in}{-0.024056in}}{\pgfqpoint{0.024056in}{0.024056in}}{%
\pgfpathmoveto{\pgfqpoint{0.000000in}{-0.024056in}}%
\pgfpathcurveto{\pgfqpoint{0.006380in}{-0.024056in}}{\pgfqpoint{0.012499in}{-0.021522in}}{\pgfqpoint{0.017010in}{-0.017010in}}%
\pgfpathcurveto{\pgfqpoint{0.021522in}{-0.012499in}}{\pgfqpoint{0.024056in}{-0.006380in}}{\pgfqpoint{0.024056in}{0.000000in}}%
\pgfpathcurveto{\pgfqpoint{0.024056in}{0.006380in}}{\pgfqpoint{0.021522in}{0.012499in}}{\pgfqpoint{0.017010in}{0.017010in}}%
\pgfpathcurveto{\pgfqpoint{0.012499in}{0.021522in}}{\pgfqpoint{0.006380in}{0.024056in}}{\pgfqpoint{0.000000in}{0.024056in}}%
\pgfpathcurveto{\pgfqpoint{-0.006380in}{0.024056in}}{\pgfqpoint{-0.012499in}{0.021522in}}{\pgfqpoint{-0.017010in}{0.017010in}}%
\pgfpathcurveto{\pgfqpoint{-0.021522in}{0.012499in}}{\pgfqpoint{-0.024056in}{0.006380in}}{\pgfqpoint{-0.024056in}{0.000000in}}%
\pgfpathcurveto{\pgfqpoint{-0.024056in}{-0.006380in}}{\pgfqpoint{-0.021522in}{-0.012499in}}{\pgfqpoint{-0.017010in}{-0.017010in}}%
\pgfpathcurveto{\pgfqpoint{-0.012499in}{-0.021522in}}{\pgfqpoint{-0.006380in}{-0.024056in}}{\pgfqpoint{0.000000in}{-0.024056in}}%
\pgfpathlineto{\pgfqpoint{0.000000in}{-0.024056in}}%
\pgfpathclose%
\pgfusepath{fill}%
}%
\begin{pgfscope}%
\pgfsys@transformshift{6.305874in}{1.009738in}%
\pgfsys@useobject{currentmarker}{}%
\end{pgfscope}%
\end{pgfscope}%
\begin{pgfscope}%
\definecolor{textcolor}{rgb}{0.549020,0.549020,0.549020}%
\pgfsetstrokecolor{textcolor}%
\pgfsetfillcolor{textcolor}%
\pgftext[x=6.346299in,y=1.009738in,left,]{\color{textcolor}\sffamily\fontsize{7.000000}{8.400000}\selectfont 0.58}%
\end{pgfscope}%
\begin{pgfscope}%
\pgfpathrectangle{\pgfqpoint{5.635000in}{0.598000in}}{\pgfqpoint{1.155000in}{3.864000in}}%
\pgfusepath{clip}%
\pgfsetbuttcap%
\pgfsetroundjoin%
\definecolor{currentfill}{rgb}{0.101961,0.101961,0.101961}%
\pgfsetfillcolor{currentfill}%
\pgfsetlinewidth{0.000000pt}%
\definecolor{currentstroke}{rgb}{0.101961,0.101961,0.101961}%
\pgfsetstrokecolor{currentstroke}%
\pgfsetdash{}{0pt}%
\pgfsys@defobject{currentmarker}{\pgfqpoint{-0.032572in}{-0.032572in}}{\pgfqpoint{0.032572in}{0.032572in}}{%
\pgfpathmoveto{\pgfqpoint{-0.016286in}{-0.032572in}}%
\pgfpathlineto{\pgfqpoint{0.000000in}{-0.016286in}}%
\pgfpathlineto{\pgfqpoint{0.016286in}{-0.032572in}}%
\pgfpathlineto{\pgfqpoint{0.032572in}{-0.016286in}}%
\pgfpathlineto{\pgfqpoint{0.016286in}{0.000000in}}%
\pgfpathlineto{\pgfqpoint{0.032572in}{0.016286in}}%
\pgfpathlineto{\pgfqpoint{0.016286in}{0.032572in}}%
\pgfpathlineto{\pgfqpoint{0.000000in}{0.016286in}}%
\pgfpathlineto{\pgfqpoint{-0.016286in}{0.032572in}}%
\pgfpathlineto{\pgfqpoint{-0.032572in}{0.016286in}}%
\pgfpathlineto{\pgfqpoint{-0.016286in}{0.000000in}}%
\pgfpathlineto{\pgfqpoint{-0.032572in}{-0.016286in}}%
\pgfpathlineto{\pgfqpoint{-0.016286in}{-0.032572in}}%
\pgfpathclose%
\pgfusepath{fill}%
}%
\begin{pgfscope}%
\pgfsys@transformshift{6.280717in}{0.851377in}%
\pgfsys@useobject{currentmarker}{}%
\end{pgfscope}%
\end{pgfscope}%
\begin{pgfscope}%
\definecolor{textcolor}{rgb}{0.549020,0.549020,0.549020}%
\pgfsetstrokecolor{textcolor}%
\pgfsetfillcolor{textcolor}%
\pgftext[x=6.321142in,y=0.851377in,left,]{\color{textcolor}\sffamily\fontsize{7.000000}{8.400000}\selectfont 0.56}%
\end{pgfscope}%
\begin{pgfscope}%
\pgfpathrectangle{\pgfqpoint{5.635000in}{0.598000in}}{\pgfqpoint{1.155000in}{3.864000in}}%
\pgfusepath{clip}%
\pgfsetbuttcap%
\pgfsetroundjoin%
\definecolor{currentfill}{rgb}{0.101961,0.101961,0.101961}%
\pgfsetfillcolor{currentfill}%
\pgfsetlinewidth{0.000000pt}%
\definecolor{currentstroke}{rgb}{0.101961,0.101961,0.101961}%
\pgfsetstrokecolor{currentstroke}%
\pgfsetdash{}{0pt}%
\pgfsys@defobject{currentmarker}{\pgfqpoint{-0.032572in}{-0.032572in}}{\pgfqpoint{0.032572in}{0.032572in}}{%
\pgfpathmoveto{\pgfqpoint{-0.016286in}{-0.032572in}}%
\pgfpathlineto{\pgfqpoint{0.000000in}{-0.016286in}}%
\pgfpathlineto{\pgfqpoint{0.016286in}{-0.032572in}}%
\pgfpathlineto{\pgfqpoint{0.032572in}{-0.016286in}}%
\pgfpathlineto{\pgfqpoint{0.016286in}{0.000000in}}%
\pgfpathlineto{\pgfqpoint{0.032572in}{0.016286in}}%
\pgfpathlineto{\pgfqpoint{0.016286in}{0.032572in}}%
\pgfpathlineto{\pgfqpoint{0.000000in}{0.016286in}}%
\pgfpathlineto{\pgfqpoint{-0.016286in}{0.032572in}}%
\pgfpathlineto{\pgfqpoint{-0.032572in}{0.016286in}}%
\pgfpathlineto{\pgfqpoint{-0.016286in}{0.000000in}}%
\pgfpathlineto{\pgfqpoint{-0.032572in}{-0.016286in}}%
\pgfpathlineto{\pgfqpoint{-0.016286in}{-0.032572in}}%
\pgfpathclose%
\pgfusepath{fill}%
}%
\begin{pgfscope}%
\pgfsys@transformshift{6.273762in}{0.693016in}%
\pgfsys@useobject{currentmarker}{}%
\end{pgfscope}%
\end{pgfscope}%
\begin{pgfscope}%
\definecolor{textcolor}{rgb}{0.549020,0.549020,0.549020}%
\pgfsetstrokecolor{textcolor}%
\pgfsetfillcolor{textcolor}%
\pgftext[x=6.314187in,y=0.693016in,left,]{\color{textcolor}\sffamily\fontsize{7.000000}{8.400000}\selectfont 0.55}%
\end{pgfscope}%
\begin{pgfscope}%
\definecolor{textcolor}{rgb}{0.549020,0.549020,0.549020}%
\pgfsetstrokecolor{textcolor}%
\pgfsetfillcolor{textcolor}%
\pgftext[x=6.930000in,y=4.531000in,right,top]{\color{textcolor}\sffamily\fontsize{5.500000}{6.600000}\selectfont OLD DATA}%
\end{pgfscope}%
\end{pgfpicture}%
\makeatother%
\endgroup%
}
\caption{[PLACEHOLDER: OLD DATA, rung 3 stand-in and COVEX dev, seed 4 only] Per-type effect of each addition (cumulative, seed 4 only). Top block: standard 16-device mix without outage, rungs 0 to 3 then the candidate additions (0b with the engine-resync part of P, B, the rest of P, question priority, touch latch fix). Middle and bottom blocks: 20 devices, 3:11:2 mix, without and with a 14-robot outage; +S2 exists only with the outage. X marks failed alternatives: counterfactual admission and eager catch-up. P, 0b and B cannot be separated one by one in the committed cells.}
\label{fig:F05_ablation_ladder}
\end{figure*}
\clearpage
\begin{figure*}[t]
\centering
{\fontfamily{LinuxBiolinumT}\selectfont
%% Creator: Matplotlib, PGF backend
%%
%% To include the figure in your LaTeX document, write
%%   \input{<filename>.pgf}
%%
%% Make sure the required packages are loaded in your preamble
%%   \usepackage{pgf}
%%
%% Also ensure that all the required font packages are loaded; for instance,
%% the lmodern package is sometimes necessary when using math font.
%%   \usepackage{lmodern}
%%
%% Figures using additional raster images can only be included by \input if
%% they are in the same directory as the main LaTeX file. For loading figures
%% from other directories you can use the `import` package
%%   \usepackage{import}
%%
%% and then include the figures with
%%   \import{<path to file>}{<filename>.pgf}
%%
%% Matplotlib used the following preamble
%%
\begingroup%
\makeatletter%
\begin{pgfpicture}%
\pgfpathrectangle{\pgfpointorigin}{\pgfqpoint{7.000000in}{4.200000in}}%
\pgfusepath{use as bounding box, clip}%
\begin{pgfscope}%
\pgfsetbuttcap%
\pgfsetmiterjoin%
\definecolor{currentfill}{rgb}{1.000000,1.000000,1.000000}%
\pgfsetfillcolor{currentfill}%
\pgfsetlinewidth{0.000000pt}%
\definecolor{currentstroke}{rgb}{1.000000,1.000000,1.000000}%
\pgfsetstrokecolor{currentstroke}%
\pgfsetdash{}{0pt}%
\pgfpathmoveto{\pgfqpoint{0.000000in}{0.000000in}}%
\pgfpathlineto{\pgfqpoint{7.000000in}{0.000000in}}%
\pgfpathlineto{\pgfqpoint{7.000000in}{4.200000in}}%
\pgfpathlineto{\pgfqpoint{0.000000in}{4.200000in}}%
\pgfpathlineto{\pgfqpoint{0.000000in}{0.000000in}}%
\pgfpathclose%
\pgfusepath{fill}%
\end{pgfscope}%
\begin{pgfscope}%
\pgfsetbuttcap%
\pgfsetmiterjoin%
\definecolor{currentfill}{rgb}{1.000000,1.000000,1.000000}%
\pgfsetfillcolor{currentfill}%
\pgfsetlinewidth{0.000000pt}%
\definecolor{currentstroke}{rgb}{0.000000,0.000000,0.000000}%
\pgfsetstrokecolor{currentstroke}%
\pgfsetstrokeopacity{0.000000}%
\pgfsetdash{}{0pt}%
\pgfpathmoveto{\pgfqpoint{0.595000in}{3.192000in}}%
\pgfpathlineto{\pgfqpoint{3.430000in}{3.192000in}}%
\pgfpathlineto{\pgfqpoint{3.430000in}{3.822000in}}%
\pgfpathlineto{\pgfqpoint{0.595000in}{3.822000in}}%
\pgfpathlineto{\pgfqpoint{0.595000in}{3.192000in}}%
\pgfpathclose%
\pgfusepath{fill}%
\end{pgfscope}%
\begin{pgfscope}%
\pgfpathrectangle{\pgfqpoint{0.595000in}{3.192000in}}{\pgfqpoint{2.835000in}{0.630000in}}%
\pgfusepath{clip}%
\pgfsetbuttcap%
\pgfsetmiterjoin%
\pgfsetlinewidth{0.000000pt}%
\definecolor{currentstroke}{rgb}{0.000000,0.000000,0.000000}%
\pgfsetstrokecolor{currentstroke}%
\pgfsetstrokeopacity{0.000000}%
\pgfsetdash{}{0pt}%
\pgfpathmoveto{\pgfqpoint{1.557719in}{3.192000in}}%
\pgfpathlineto{\pgfqpoint{1.557719in}{3.822000in}}%
\pgfpathlineto{\pgfqpoint{2.001869in}{3.822000in}}%
\pgfpathlineto{\pgfqpoint{2.001869in}{3.192000in}}%
\pgfpathlineto{\pgfqpoint{1.557719in}{3.192000in}}%
\pgfpathclose%
\pgfusepath{}%
\end{pgfscope}%
\begin{pgfscope}%
\pgfpathrectangle{\pgfqpoint{0.595000in}{3.192000in}}{\pgfqpoint{2.835000in}{0.630000in}}%
\pgfusepath{clip}%
\pgfsetrectcap%
\pgfsetroundjoin%
\pgfsetlinewidth{0.401500pt}%
\definecolor{currentstroke}{rgb}{0.901961,0.901961,0.901961}%
\pgfsetstrokecolor{currentstroke}%
\pgfsetdash{}{0pt}%
\pgfpathmoveto{\pgfqpoint{0.595000in}{3.192000in}}%
\pgfpathlineto{\pgfqpoint{3.430000in}{3.192000in}}%
\pgfusepath{stroke}%
\end{pgfscope}%
\begin{pgfscope}%
\definecolor{textcolor}{rgb}{0.101961,0.101961,0.101961}%
\pgfsetstrokecolor{textcolor}%
\pgfsetfillcolor{textcolor}%
\pgftext[x=0.362745in, y=3.158242in, left, base]{\color{textcolor}\sffamily\fontsize{7.000000}{8.400000}\selectfont 0.00}%
\end{pgfscope}%
\begin{pgfscope}%
\pgfpathrectangle{\pgfqpoint{0.595000in}{3.192000in}}{\pgfqpoint{2.835000in}{0.630000in}}%
\pgfusepath{clip}%
\pgfsetrectcap%
\pgfsetroundjoin%
\pgfsetlinewidth{0.401500pt}%
\definecolor{currentstroke}{rgb}{0.901961,0.901961,0.901961}%
\pgfsetstrokecolor{currentstroke}%
\pgfsetdash{}{0pt}%
\pgfpathmoveto{\pgfqpoint{0.595000in}{3.402000in}}%
\pgfpathlineto{\pgfqpoint{3.430000in}{3.402000in}}%
\pgfusepath{stroke}%
\end{pgfscope}%
\begin{pgfscope}%
\definecolor{textcolor}{rgb}{0.101961,0.101961,0.101961}%
\pgfsetstrokecolor{textcolor}%
\pgfsetfillcolor{textcolor}%
\pgftext[x=0.362745in, y=3.368242in, left, base]{\color{textcolor}\sffamily\fontsize{7.000000}{8.400000}\selectfont 0.25}%
\end{pgfscope}%
\begin{pgfscope}%
\pgfpathrectangle{\pgfqpoint{0.595000in}{3.192000in}}{\pgfqpoint{2.835000in}{0.630000in}}%
\pgfusepath{clip}%
\pgfsetrectcap%
\pgfsetroundjoin%
\pgfsetlinewidth{0.401500pt}%
\definecolor{currentstroke}{rgb}{0.901961,0.901961,0.901961}%
\pgfsetstrokecolor{currentstroke}%
\pgfsetdash{}{0pt}%
\pgfpathmoveto{\pgfqpoint{0.595000in}{3.612000in}}%
\pgfpathlineto{\pgfqpoint{3.430000in}{3.612000in}}%
\pgfusepath{stroke}%
\end{pgfscope}%
\begin{pgfscope}%
\definecolor{textcolor}{rgb}{0.101961,0.101961,0.101961}%
\pgfsetstrokecolor{textcolor}%
\pgfsetfillcolor{textcolor}%
\pgftext[x=0.362745in, y=3.578242in, left, base]{\color{textcolor}\sffamily\fontsize{7.000000}{8.400000}\selectfont 0.50}%
\end{pgfscope}%
\begin{pgfscope}%
\pgfpathrectangle{\pgfqpoint{0.595000in}{3.192000in}}{\pgfqpoint{2.835000in}{0.630000in}}%
\pgfusepath{clip}%
\pgfsetrectcap%
\pgfsetroundjoin%
\pgfsetlinewidth{0.401500pt}%
\definecolor{currentstroke}{rgb}{0.901961,0.901961,0.901961}%
\pgfsetstrokecolor{currentstroke}%
\pgfsetdash{}{0pt}%
\pgfpathmoveto{\pgfqpoint{0.595000in}{3.822000in}}%
\pgfpathlineto{\pgfqpoint{3.430000in}{3.822000in}}%
\pgfusepath{stroke}%
\end{pgfscope}%
\begin{pgfscope}%
\definecolor{textcolor}{rgb}{0.101961,0.101961,0.101961}%
\pgfsetstrokecolor{textcolor}%
\pgfsetfillcolor{textcolor}%
\pgftext[x=0.362745in, y=3.788242in, left, base]{\color{textcolor}\sffamily\fontsize{7.000000}{8.400000}\selectfont 0.75}%
\end{pgfscope}%
\begin{pgfscope}%
\definecolor{textcolor}{rgb}{0.000000,0.000000,0.000000}%
\pgfsetstrokecolor{textcolor}%
\pgfsetfillcolor{textcolor}%
\pgftext[x=0.176479in, y=3.271368in, left, base,rotate=90.000000]{\color{textcolor}\sffamily\fontsize{7.500000}{9.000000}\selectfont Occupancy}%
\end{pgfscope}%
\begin{pgfscope}%
\definecolor{textcolor}{rgb}{0.000000,0.000000,0.000000}%
\pgfsetstrokecolor{textcolor}%
\pgfsetfillcolor{textcolor}%
\pgftext[x=0.282884in, y=3.302173in, left, base,rotate=90.000000]{\color{textcolor}\sffamily\fontsize{7.500000}{9.000000}\selectfont (fraction)}%
\end{pgfscope}%
\begin{pgfscope}%
\pgfpathrectangle{\pgfqpoint{0.595000in}{3.192000in}}{\pgfqpoint{2.835000in}{0.630000in}}%
\pgfusepath{clip}%
\pgfsetrectcap%
\pgfsetroundjoin%
\pgfsetlinewidth{0.903375pt}%
\definecolor{currentstroke}{rgb}{0.156863,0.356863,0.568627}%
\pgfsetstrokecolor{currentstroke}%
\pgfsetdash{}{0pt}%
\pgfpathmoveto{\pgfqpoint{0.597605in}{3.192000in}}%
\pgfpathlineto{\pgfqpoint{0.621686in}{3.192000in}}%
\pgfpathlineto{\pgfqpoint{0.645772in}{3.287469in}}%
\pgfpathlineto{\pgfqpoint{0.670231in}{3.192000in}}%
\pgfpathlineto{\pgfqpoint{0.694408in}{3.319719in}}%
\pgfpathlineto{\pgfqpoint{0.718655in}{3.192000in}}%
\pgfpathlineto{\pgfqpoint{0.742766in}{3.192000in}}%
\pgfpathlineto{\pgfqpoint{0.766928in}{3.247317in}}%
\pgfpathlineto{\pgfqpoint{0.791068in}{3.192000in}}%
\pgfpathlineto{\pgfqpoint{0.815279in}{3.231939in}}%
\pgfpathlineto{\pgfqpoint{0.839787in}{3.206310in}}%
\pgfpathlineto{\pgfqpoint{0.864012in}{3.215280in}}%
\pgfpathlineto{\pgfqpoint{0.888185in}{3.192000in}}%
\pgfpathlineto{\pgfqpoint{0.912408in}{3.256500in}}%
\pgfpathlineto{\pgfqpoint{0.936598in}{3.265684in}}%
\pgfpathlineto{\pgfqpoint{0.960851in}{3.192000in}}%
\pgfpathlineto{\pgfqpoint{0.985126in}{3.369910in}}%
\pgfpathlineto{\pgfqpoint{1.009230in}{3.355600in}}%
\pgfpathlineto{\pgfqpoint{1.033692in}{3.397675in}}%
\pgfpathlineto{\pgfqpoint{1.057865in}{3.192000in}}%
\pgfpathlineto{\pgfqpoint{1.082102in}{3.403655in}}%
\pgfpathlineto{\pgfqpoint{1.106270in}{3.445944in}}%
\pgfpathlineto{\pgfqpoint{1.130431in}{3.365425in}}%
\pgfpathlineto{\pgfqpoint{1.154612in}{3.192000in}}%
\pgfpathlineto{\pgfqpoint{1.178742in}{3.251375in}}%
\pgfpathlineto{\pgfqpoint{1.202932in}{3.192000in}}%
\pgfpathlineto{\pgfqpoint{1.227133in}{3.251375in}}%
\pgfpathlineto{\pgfqpoint{1.251358in}{3.284693in}}%
\pgfpathlineto{\pgfqpoint{1.275557in}{3.260986in}}%
\pgfpathlineto{\pgfqpoint{1.299799in}{3.376104in}}%
\pgfpathlineto{\pgfqpoint{1.323937in}{3.405791in}}%
\pgfpathlineto{\pgfqpoint{1.348122in}{3.434410in}}%
\pgfpathlineto{\pgfqpoint{1.372271in}{3.192000in}}%
\pgfpathlineto{\pgfqpoint{1.396503in}{3.192000in}}%
\pgfpathlineto{\pgfqpoint{1.420884in}{3.192000in}}%
\pgfpathlineto{\pgfqpoint{1.445105in}{3.402160in}}%
\pgfpathlineto{\pgfqpoint{1.469237in}{3.206096in}}%
\pgfpathlineto{\pgfqpoint{1.493437in}{3.370764in}}%
\pgfpathlineto{\pgfqpoint{1.517643in}{3.393404in}}%
\pgfpathlineto{\pgfqpoint{1.541776in}{3.192000in}}%
\pgfpathlineto{\pgfqpoint{1.566005in}{3.282984in}}%
\pgfpathlineto{\pgfqpoint{1.590476in}{3.328690in}}%
\pgfpathlineto{\pgfqpoint{1.614742in}{3.192000in}}%
\pgfpathlineto{\pgfqpoint{1.638877in}{3.192000in}}%
\pgfpathlineto{\pgfqpoint{1.663199in}{3.503823in}}%
\pgfpathlineto{\pgfqpoint{1.687429in}{3.192000in}}%
\pgfpathlineto{\pgfqpoint{1.711609in}{3.384433in}}%
\pgfpathlineto{\pgfqpoint{1.735732in}{3.192000in}}%
\pgfpathlineto{\pgfqpoint{1.759849in}{3.233220in}}%
\pgfpathlineto{\pgfqpoint{1.784128in}{3.192000in}}%
\pgfpathlineto{\pgfqpoint{1.808478in}{3.403014in}}%
\pgfpathlineto{\pgfqpoint{1.832899in}{3.192000in}}%
\pgfpathlineto{\pgfqpoint{1.857028in}{3.192000in}}%
\pgfpathlineto{\pgfqpoint{1.881213in}{3.242191in}}%
\pgfpathlineto{\pgfqpoint{1.905440in}{3.373968in}}%
\pgfpathlineto{\pgfqpoint{1.929644in}{3.287469in}}%
\pgfpathlineto{\pgfqpoint{1.953786in}{3.247530in}}%
\pgfpathlineto{\pgfqpoint{1.977959in}{3.279140in}}%
\pgfpathlineto{\pgfqpoint{2.002470in}{3.192000in}}%
\pgfpathlineto{\pgfqpoint{2.026731in}{3.251802in}}%
\pgfpathlineto{\pgfqpoint{2.050965in}{3.256500in}}%
\pgfpathlineto{\pgfqpoint{2.075096in}{3.349407in}}%
\pgfpathlineto{\pgfqpoint{2.099335in}{3.315234in}}%
\pgfpathlineto{\pgfqpoint{2.123478in}{3.346416in}}%
\pgfpathlineto{\pgfqpoint{2.147684in}{3.192000in}}%
\pgfpathlineto{\pgfqpoint{2.172173in}{3.271237in}}%
\pgfpathlineto{\pgfqpoint{2.196356in}{3.342999in}}%
\pgfpathlineto{\pgfqpoint{2.220579in}{3.224250in}}%
\pgfpathlineto{\pgfqpoint{2.244797in}{3.356241in}}%
\pgfpathlineto{\pgfqpoint{2.269003in}{3.206096in}}%
\pgfpathlineto{\pgfqpoint{2.293230in}{3.192000in}}%
\pgfpathlineto{\pgfqpoint{2.317432in}{3.301565in}}%
\pgfpathlineto{\pgfqpoint{2.341775in}{3.445730in}}%
\pgfpathlineto{\pgfqpoint{2.365962in}{3.255860in}}%
\pgfpathlineto{\pgfqpoint{2.390112in}{3.192000in}}%
\pgfpathlineto{\pgfqpoint{2.414271in}{3.192000in}}%
\pgfpathlineto{\pgfqpoint{2.438529in}{3.220406in}}%
\pgfpathlineto{\pgfqpoint{2.462662in}{3.613602in}}%
\pgfpathlineto{\pgfqpoint{2.486929in}{3.192000in}}%
\pgfpathlineto{\pgfqpoint{2.511353in}{3.338087in}}%
\pgfpathlineto{\pgfqpoint{2.535524in}{3.395112in}}%
\pgfpathlineto{\pgfqpoint{2.559687in}{3.192000in}}%
\pgfpathlineto{\pgfqpoint{2.583884in}{3.503609in}}%
\pgfpathlineto{\pgfqpoint{2.607970in}{3.421169in}}%
\pgfpathlineto{\pgfqpoint{2.632079in}{3.490368in}}%
\pgfpathlineto{\pgfqpoint{2.656243in}{3.362435in}}%
\pgfpathlineto{\pgfqpoint{2.680484in}{3.370978in}}%
\pgfpathlineto{\pgfqpoint{2.704660in}{3.215280in}}%
\pgfpathlineto{\pgfqpoint{2.728802in}{3.375463in}}%
\pgfpathlineto{\pgfqpoint{2.753211in}{3.306691in}}%
\pgfpathlineto{\pgfqpoint{2.777363in}{3.192000in}}%
\pgfpathlineto{\pgfqpoint{2.801539in}{3.614029in}}%
\pgfpathlineto{\pgfqpoint{2.825740in}{3.324845in}}%
\pgfpathlineto{\pgfqpoint{2.849927in}{3.312458in}}%
\pgfpathlineto{\pgfqpoint{2.874160in}{3.233220in}}%
\pgfpathlineto{\pgfqpoint{2.898378in}{3.486950in}}%
\pgfpathlineto{\pgfqpoint{2.922579in}{3.292595in}}%
\pgfpathlineto{\pgfqpoint{2.946821in}{3.694334in}}%
\pgfpathlineto{\pgfqpoint{2.971024in}{3.454700in}}%
\pgfpathlineto{\pgfqpoint{2.995318in}{3.192000in}}%
\pgfpathlineto{\pgfqpoint{3.019553in}{3.346630in}}%
\pgfpathlineto{\pgfqpoint{3.043723in}{3.479475in}}%
\pgfpathlineto{\pgfqpoint{3.067925in}{3.192000in}}%
\pgfpathlineto{\pgfqpoint{3.092299in}{3.259490in}}%
\pgfpathlineto{\pgfqpoint{3.116427in}{3.192000in}}%
\pgfpathlineto{\pgfqpoint{3.140586in}{3.360299in}}%
\pgfpathlineto{\pgfqpoint{3.164922in}{3.286828in}}%
\pgfpathlineto{\pgfqpoint{3.189086in}{3.292595in}}%
\pgfpathlineto{\pgfqpoint{3.213311in}{3.389346in}}%
\pgfpathlineto{\pgfqpoint{3.237576in}{3.306264in}}%
\pgfpathlineto{\pgfqpoint{3.262061in}{3.433983in}}%
\pgfpathlineto{\pgfqpoint{3.286260in}{3.260345in}}%
\pgfpathlineto{\pgfqpoint{3.310483in}{3.609971in}}%
\pgfpathlineto{\pgfqpoint{3.334738in}{3.337873in}}%
\pgfpathlineto{\pgfqpoint{3.359015in}{3.283625in}}%
\pgfpathlineto{\pgfqpoint{3.383186in}{3.192000in}}%
\pgfpathlineto{\pgfqpoint{3.407423in}{3.224250in}}%
\pgfpathlineto{\pgfqpoint{3.431854in}{3.410703in}}%
\pgfusepath{stroke}%
\end{pgfscope}%
\begin{pgfscope}%
\pgfpathrectangle{\pgfqpoint{0.595000in}{3.192000in}}{\pgfqpoint{2.835000in}{0.630000in}}%
\pgfusepath{clip}%
\pgfsetrectcap%
\pgfsetroundjoin%
\pgfsetlinewidth{0.903375pt}%
\definecolor{currentstroke}{rgb}{0.478431,0.478431,0.478431}%
\pgfsetstrokecolor{currentstroke}%
\pgfsetdash{}{0pt}%
\pgfpathmoveto{\pgfqpoint{0.621686in}{3.192000in}}%
\pgfpathlineto{\pgfqpoint{0.645772in}{3.192000in}}%
\pgfpathlineto{\pgfqpoint{0.670231in}{3.192000in}}%
\pgfpathlineto{\pgfqpoint{0.694408in}{3.192000in}}%
\pgfpathlineto{\pgfqpoint{0.718655in}{3.192000in}}%
\pgfpathlineto{\pgfqpoint{0.742766in}{3.192000in}}%
\pgfpathlineto{\pgfqpoint{0.766928in}{3.192000in}}%
\pgfpathlineto{\pgfqpoint{0.791068in}{3.192000in}}%
\pgfpathlineto{\pgfqpoint{0.815279in}{3.229674in}}%
\pgfpathlineto{\pgfqpoint{0.839787in}{3.192000in}}%
\pgfpathlineto{\pgfqpoint{0.864012in}{3.192000in}}%
\pgfpathlineto{\pgfqpoint{0.888185in}{3.192000in}}%
\pgfpathlineto{\pgfqpoint{0.912408in}{3.192000in}}%
\pgfpathlineto{\pgfqpoint{0.936598in}{3.192000in}}%
\pgfpathlineto{\pgfqpoint{0.960851in}{3.192000in}}%
\pgfpathlineto{\pgfqpoint{0.985126in}{3.192000in}}%
\pgfpathlineto{\pgfqpoint{1.009230in}{3.198198in}}%
\pgfpathlineto{\pgfqpoint{1.033692in}{3.231497in}}%
\pgfpathlineto{\pgfqpoint{1.057865in}{3.192000in}}%
\pgfpathlineto{\pgfqpoint{1.082102in}{3.192000in}}%
\pgfpathlineto{\pgfqpoint{1.106270in}{3.222139in}}%
\pgfpathlineto{\pgfqpoint{1.130431in}{3.192000in}}%
\pgfpathlineto{\pgfqpoint{1.154612in}{3.192000in}}%
\pgfpathlineto{\pgfqpoint{1.178742in}{3.192000in}}%
\pgfpathlineto{\pgfqpoint{1.202932in}{3.192000in}}%
\pgfpathlineto{\pgfqpoint{1.227133in}{3.192000in}}%
\pgfpathlineto{\pgfqpoint{1.251358in}{3.295785in}}%
\pgfpathlineto{\pgfqpoint{1.275557in}{3.192000in}}%
\pgfpathlineto{\pgfqpoint{1.299799in}{3.322399in}}%
\pgfpathlineto{\pgfqpoint{1.323937in}{3.313649in}}%
\pgfpathlineto{\pgfqpoint{1.348122in}{3.308910in}}%
\pgfpathlineto{\pgfqpoint{1.372271in}{3.192000in}}%
\pgfpathlineto{\pgfqpoint{1.396503in}{3.192000in}}%
\pgfpathlineto{\pgfqpoint{1.420884in}{3.192000in}}%
\pgfpathlineto{\pgfqpoint{1.445105in}{3.192000in}}%
\pgfpathlineto{\pgfqpoint{1.469237in}{3.192000in}}%
\pgfpathlineto{\pgfqpoint{1.493437in}{3.252156in}}%
\pgfpathlineto{\pgfqpoint{1.517643in}{3.192000in}}%
\pgfpathlineto{\pgfqpoint{1.541776in}{3.192000in}}%
\pgfpathlineto{\pgfqpoint{1.566005in}{3.199535in}}%
\pgfpathlineto{\pgfqpoint{1.590476in}{3.192000in}}%
\pgfpathlineto{\pgfqpoint{1.614742in}{3.192000in}}%
\pgfpathlineto{\pgfqpoint{1.638877in}{3.192000in}}%
\pgfpathlineto{\pgfqpoint{1.663199in}{3.301618in}}%
\pgfpathlineto{\pgfqpoint{1.687429in}{3.199535in}}%
\pgfpathlineto{\pgfqpoint{1.711609in}{3.192000in}}%
\pgfpathlineto{\pgfqpoint{1.735732in}{3.192000in}}%
\pgfpathlineto{\pgfqpoint{1.759849in}{3.192000in}}%
\pgfpathlineto{\pgfqpoint{1.784128in}{3.192000in}}%
\pgfpathlineto{\pgfqpoint{1.808478in}{3.192729in}}%
\pgfpathlineto{\pgfqpoint{1.832899in}{3.192000in}}%
\pgfpathlineto{\pgfqpoint{1.857028in}{3.192000in}}%
\pgfpathlineto{\pgfqpoint{1.881213in}{3.192000in}}%
\pgfpathlineto{\pgfqpoint{1.905440in}{3.215212in}}%
\pgfpathlineto{\pgfqpoint{1.929644in}{3.192000in}}%
\pgfpathlineto{\pgfqpoint{1.953786in}{3.192000in}}%
\pgfpathlineto{\pgfqpoint{1.977959in}{3.192000in}}%
\pgfpathlineto{\pgfqpoint{2.002470in}{3.192000in}}%
\pgfpathlineto{\pgfqpoint{2.026731in}{3.207677in}}%
\pgfpathlineto{\pgfqpoint{2.050965in}{3.192000in}}%
\pgfpathlineto{\pgfqpoint{2.075096in}{3.199656in}}%
\pgfpathlineto{\pgfqpoint{2.099335in}{3.192000in}}%
\pgfpathlineto{\pgfqpoint{2.123478in}{3.192000in}}%
\pgfpathlineto{\pgfqpoint{2.147684in}{3.192000in}}%
\pgfpathlineto{\pgfqpoint{2.172173in}{3.192729in}}%
\pgfpathlineto{\pgfqpoint{2.196356in}{3.237451in}}%
\pgfpathlineto{\pgfqpoint{2.220579in}{3.192000in}}%
\pgfpathlineto{\pgfqpoint{2.244797in}{3.222747in}}%
\pgfpathlineto{\pgfqpoint{2.269003in}{3.192000in}}%
\pgfpathlineto{\pgfqpoint{2.293230in}{3.192000in}}%
\pgfpathlineto{\pgfqpoint{2.317432in}{3.192000in}}%
\pgfpathlineto{\pgfqpoint{2.341775in}{3.211931in}}%
\pgfpathlineto{\pgfqpoint{2.365962in}{3.192000in}}%
\pgfpathlineto{\pgfqpoint{2.390112in}{3.192000in}}%
\pgfpathlineto{\pgfqpoint{2.414271in}{3.192000in}}%
\pgfpathlineto{\pgfqpoint{2.438529in}{3.192000in}}%
\pgfpathlineto{\pgfqpoint{2.462662in}{3.245351in}}%
\pgfpathlineto{\pgfqpoint{2.486929in}{3.192000in}}%
\pgfpathlineto{\pgfqpoint{2.511353in}{3.236722in}}%
\pgfpathlineto{\pgfqpoint{2.535524in}{3.254222in}}%
\pgfpathlineto{\pgfqpoint{2.559687in}{3.192000in}}%
\pgfpathlineto{\pgfqpoint{2.583884in}{3.192000in}}%
\pgfpathlineto{\pgfqpoint{2.607970in}{3.237816in}}%
\pgfpathlineto{\pgfqpoint{2.632079in}{3.192000in}}%
\pgfpathlineto{\pgfqpoint{2.656243in}{3.192000in}}%
\pgfpathlineto{\pgfqpoint{2.680484in}{3.192000in}}%
\pgfpathlineto{\pgfqpoint{2.704660in}{3.192000in}}%
\pgfpathlineto{\pgfqpoint{2.728802in}{3.192000in}}%
\pgfpathlineto{\pgfqpoint{2.753211in}{3.192000in}}%
\pgfpathlineto{\pgfqpoint{2.777363in}{3.192000in}}%
\pgfpathlineto{\pgfqpoint{2.801539in}{3.225785in}}%
\pgfpathlineto{\pgfqpoint{2.825740in}{3.254222in}}%
\pgfpathlineto{\pgfqpoint{2.849927in}{3.192000in}}%
\pgfpathlineto{\pgfqpoint{2.874160in}{3.192000in}}%
\pgfpathlineto{\pgfqpoint{2.898378in}{3.255559in}}%
\pgfpathlineto{\pgfqpoint{2.922579in}{3.192000in}}%
\pgfpathlineto{\pgfqpoint{2.946821in}{3.360194in}}%
\pgfpathlineto{\pgfqpoint{2.971024in}{3.235385in}}%
\pgfpathlineto{\pgfqpoint{2.995318in}{3.192000in}}%
\pgfpathlineto{\pgfqpoint{3.019553in}{3.192000in}}%
\pgfpathlineto{\pgfqpoint{3.043723in}{3.231497in}}%
\pgfpathlineto{\pgfqpoint{3.067925in}{3.192000in}}%
\pgfpathlineto{\pgfqpoint{3.092299in}{3.240733in}}%
\pgfpathlineto{\pgfqpoint{3.116427in}{3.192000in}}%
\pgfpathlineto{\pgfqpoint{3.140586in}{3.192000in}}%
\pgfpathlineto{\pgfqpoint{3.164922in}{3.192729in}}%
\pgfpathlineto{\pgfqpoint{3.189086in}{3.192000in}}%
\pgfpathlineto{\pgfqpoint{3.213311in}{3.192000in}}%
\pgfpathlineto{\pgfqpoint{3.237576in}{3.192000in}}%
\pgfpathlineto{\pgfqpoint{3.262061in}{3.260420in}}%
\pgfpathlineto{\pgfqpoint{3.286260in}{3.192000in}}%
\pgfpathlineto{\pgfqpoint{3.310483in}{3.194917in}}%
\pgfpathlineto{\pgfqpoint{3.334738in}{3.192000in}}%
\pgfpathlineto{\pgfqpoint{3.359015in}{3.192000in}}%
\pgfpathlineto{\pgfqpoint{3.383186in}{3.192000in}}%
\pgfpathlineto{\pgfqpoint{3.407423in}{3.192000in}}%
\pgfpathlineto{\pgfqpoint{3.431854in}{3.316323in}}%
\pgfusepath{stroke}%
\end{pgfscope}%
\begin{pgfscope}%
\pgfsetrectcap%
\pgfsetmiterjoin%
\pgfsetlinewidth{0.803000pt}%
\definecolor{currentstroke}{rgb}{0.101961,0.101961,0.101961}%
\pgfsetstrokecolor{currentstroke}%
\pgfsetdash{}{0pt}%
\pgfpathmoveto{\pgfqpoint{0.595000in}{3.192000in}}%
\pgfpathlineto{\pgfqpoint{0.595000in}{3.822000in}}%
\pgfusepath{stroke}%
\end{pgfscope}%
\begin{pgfscope}%
\pgfsetrectcap%
\pgfsetmiterjoin%
\pgfsetlinewidth{0.803000pt}%
\definecolor{currentstroke}{rgb}{0.101961,0.101961,0.101961}%
\pgfsetstrokecolor{currentstroke}%
\pgfsetdash{}{0pt}%
\pgfpathmoveto{\pgfqpoint{0.595000in}{3.192000in}}%
\pgfpathlineto{\pgfqpoint{3.430000in}{3.192000in}}%
\pgfusepath{stroke}%
\end{pgfscope}%
\begin{pgfscope}%
\definecolor{textcolor}{rgb}{0.101961,0.101961,0.101961}%
\pgfsetstrokecolor{textcolor}%
\pgfsetfillcolor{textcolor}%
\pgftext[x=0.623350in,y=3.840900in,left,bottom]{\color{textcolor}\sffamily\fontsize{7.000000}{8.400000}\selectfont (a) normal 1x run}%
\end{pgfscope}%
\begin{pgfscope}%
\pgfsetrectcap%
\pgfsetroundjoin%
\pgfsetlinewidth{0.903375pt}%
\definecolor{currentstroke}{rgb}{0.156863,0.356863,0.568627}%
\pgfsetstrokecolor{currentstroke}%
\pgfsetdash{}{0pt}%
\pgfpathmoveto{\pgfqpoint{0.682500in}{3.687872in}}%
\pgfpathlineto{\pgfqpoint{0.740833in}{3.687872in}}%
\pgfpathlineto{\pgfqpoint{0.799167in}{3.687872in}}%
\pgfusepath{stroke}%
\end{pgfscope}%
\begin{pgfscope}%
\definecolor{textcolor}{rgb}{0.101961,0.101961,0.101961}%
\pgfsetstrokecolor{textcolor}%
\pgfsetfillcolor{textcolor}%
\pgftext[x=0.876944in,y=3.653844in,left,base]{\color{textcolor}\sffamily\fontsize{7.000000}{8.400000}\selectfont KV pool}%
\end{pgfscope}%
\begin{pgfscope}%
\pgfsetrectcap%
\pgfsetroundjoin%
\pgfsetlinewidth{0.903375pt}%
\definecolor{currentstroke}{rgb}{0.478431,0.478431,0.478431}%
\pgfsetstrokecolor{currentstroke}%
\pgfsetdash{}{0pt}%
\pgfpathmoveto{\pgfqpoint{1.316778in}{3.687872in}}%
\pgfpathlineto{\pgfqpoint{1.375111in}{3.687872in}}%
\pgfpathlineto{\pgfqpoint{1.433445in}{3.687872in}}%
\pgfusepath{stroke}%
\end{pgfscope}%
\begin{pgfscope}%
\definecolor{textcolor}{rgb}{0.101961,0.101961,0.101961}%
\pgfsetstrokecolor{textcolor}%
\pgfsetfillcolor{textcolor}%
\pgftext[x=1.511222in,y=3.653844in,left,base]{\color{textcolor}\sffamily\fontsize{7.000000}{8.400000}\selectfont host tier}%
\end{pgfscope}%
\begin{pgfscope}%
\pgfsetbuttcap%
\pgfsetmiterjoin%
\definecolor{currentfill}{rgb}{1.000000,1.000000,1.000000}%
\pgfsetfillcolor{currentfill}%
\pgfsetlinewidth{0.000000pt}%
\definecolor{currentstroke}{rgb}{0.000000,0.000000,0.000000}%
\pgfsetstrokecolor{currentstroke}%
\pgfsetstrokeopacity{0.000000}%
\pgfsetdash{}{0pt}%
\pgfpathmoveto{\pgfqpoint{0.595000in}{2.289000in}}%
\pgfpathlineto{\pgfqpoint{3.430000in}{2.289000in}}%
\pgfpathlineto{\pgfqpoint{3.430000in}{2.919000in}}%
\pgfpathlineto{\pgfqpoint{0.595000in}{2.919000in}}%
\pgfpathlineto{\pgfqpoint{0.595000in}{2.289000in}}%
\pgfpathclose%
\pgfusepath{fill}%
\end{pgfscope}%
\begin{pgfscope}%
\pgfpathrectangle{\pgfqpoint{0.595000in}{2.289000in}}{\pgfqpoint{2.835000in}{0.630000in}}%
\pgfusepath{clip}%
\pgfsetbuttcap%
\pgfsetmiterjoin%
\pgfsetlinewidth{0.000000pt}%
\definecolor{currentstroke}{rgb}{0.000000,0.000000,0.000000}%
\pgfsetstrokecolor{currentstroke}%
\pgfsetstrokeopacity{0.000000}%
\pgfsetdash{}{0pt}%
\pgfpathmoveto{\pgfqpoint{1.557719in}{2.289000in}}%
\pgfpathlineto{\pgfqpoint{1.557719in}{2.919000in}}%
\pgfpathlineto{\pgfqpoint{2.001869in}{2.919000in}}%
\pgfpathlineto{\pgfqpoint{2.001869in}{2.289000in}}%
\pgfpathlineto{\pgfqpoint{1.557719in}{2.289000in}}%
\pgfpathclose%
\pgfusepath{}%
\end{pgfscope}%
\begin{pgfscope}%
\pgfpathrectangle{\pgfqpoint{0.595000in}{2.289000in}}{\pgfqpoint{2.835000in}{0.630000in}}%
\pgfusepath{clip}%
\pgfsetrectcap%
\pgfsetroundjoin%
\pgfsetlinewidth{0.401500pt}%
\definecolor{currentstroke}{rgb}{0.901961,0.901961,0.901961}%
\pgfsetstrokecolor{currentstroke}%
\pgfsetdash{}{0pt}%
\pgfpathmoveto{\pgfqpoint{0.595000in}{2.289000in}}%
\pgfpathlineto{\pgfqpoint{3.430000in}{2.289000in}}%
\pgfusepath{stroke}%
\end{pgfscope}%
\begin{pgfscope}%
\definecolor{textcolor}{rgb}{0.101961,0.101961,0.101961}%
\pgfsetstrokecolor{textcolor}%
\pgfsetfillcolor{textcolor}%
\pgftext[x=0.494739in, y=2.255242in, left, base]{\color{textcolor}\sffamily\fontsize{7.000000}{8.400000}\selectfont 0}%
\end{pgfscope}%
\begin{pgfscope}%
\pgfpathrectangle{\pgfqpoint{0.595000in}{2.289000in}}{\pgfqpoint{2.835000in}{0.630000in}}%
\pgfusepath{clip}%
\pgfsetrectcap%
\pgfsetroundjoin%
\pgfsetlinewidth{0.401500pt}%
\definecolor{currentstroke}{rgb}{0.901961,0.901961,0.901961}%
\pgfsetstrokecolor{currentstroke}%
\pgfsetdash{}{0pt}%
\pgfpathmoveto{\pgfqpoint{0.595000in}{2.604000in}}%
\pgfpathlineto{\pgfqpoint{3.430000in}{2.604000in}}%
\pgfusepath{stroke}%
\end{pgfscope}%
\begin{pgfscope}%
\definecolor{textcolor}{rgb}{0.101961,0.101961,0.101961}%
\pgfsetstrokecolor{textcolor}%
\pgfsetfillcolor{textcolor}%
\pgftext[x=0.391439in, y=2.570242in, left, base]{\color{textcolor}\sffamily\fontsize{7.000000}{8.400000}\selectfont 100}%
\end{pgfscope}%
\begin{pgfscope}%
\pgfpathrectangle{\pgfqpoint{0.595000in}{2.289000in}}{\pgfqpoint{2.835000in}{0.630000in}}%
\pgfusepath{clip}%
\pgfsetrectcap%
\pgfsetroundjoin%
\pgfsetlinewidth{0.401500pt}%
\definecolor{currentstroke}{rgb}{0.901961,0.901961,0.901961}%
\pgfsetstrokecolor{currentstroke}%
\pgfsetdash{}{0pt}%
\pgfpathmoveto{\pgfqpoint{0.595000in}{2.919000in}}%
\pgfpathlineto{\pgfqpoint{3.430000in}{2.919000in}}%
\pgfusepath{stroke}%
\end{pgfscope}%
\begin{pgfscope}%
\definecolor{textcolor}{rgb}{0.101961,0.101961,0.101961}%
\pgfsetstrokecolor{textcolor}%
\pgfsetfillcolor{textcolor}%
\pgftext[x=0.391439in, y=2.885242in, left, base]{\color{textcolor}\sffamily\fontsize{7.000000}{8.400000}\selectfont 200}%
\end{pgfscope}%
\begin{pgfscope}%
\definecolor{textcolor}{rgb}{0.000000,0.000000,0.000000}%
\pgfsetstrokecolor{textcolor}%
\pgfsetfillcolor{textcolor}%
\pgftext[x=0.217057in, y=2.404051in, left, base,rotate=90.000000]{\color{textcolor}\sffamily\fontsize{7.500000}{9.000000}\selectfont Decisions}%
\end{pgfscope}%
\begin{pgfscope}%
\definecolor{textcolor}{rgb}{0.000000,0.000000,0.000000}%
\pgfsetstrokecolor{textcolor}%
\pgfsetfillcolor{textcolor}%
\pgftext[x=0.316979in, y=2.429420in, left, base,rotate=90.000000]{\color{textcolor}\sffamily\fontsize{7.500000}{9.000000}\selectfont per 10 s}%
\end{pgfscope}%
\begin{pgfscope}%
\pgfpathrectangle{\pgfqpoint{0.595000in}{2.289000in}}{\pgfqpoint{2.835000in}{0.630000in}}%
\pgfusepath{clip}%
\pgfsetrectcap%
\pgfsetroundjoin%
\pgfsetlinewidth{0.903375pt}%
\definecolor{currentstroke}{rgb}{0.156863,0.356863,0.568627}%
\pgfsetstrokecolor{currentstroke}%
\pgfsetdash{}{0pt}%
\pgfpathmoveto{\pgfqpoint{0.606813in}{2.289000in}}%
\pgfpathlineto{\pgfqpoint{2.437750in}{2.289000in}}%
\pgfpathlineto{\pgfqpoint{2.437750in}{2.408700in}}%
\pgfpathlineto{\pgfqpoint{2.461375in}{2.408700in}}%
\pgfpathlineto{\pgfqpoint{2.461375in}{2.604000in}}%
\pgfpathlineto{\pgfqpoint{2.485000in}{2.604000in}}%
\pgfpathlineto{\pgfqpoint{2.485000in}{2.600850in}}%
\pgfpathlineto{\pgfqpoint{2.508625in}{2.600850in}}%
\pgfpathlineto{\pgfqpoint{2.508625in}{2.604000in}}%
\pgfpathlineto{\pgfqpoint{2.555875in}{2.604000in}}%
\pgfpathlineto{\pgfqpoint{2.555875in}{2.600850in}}%
\pgfpathlineto{\pgfqpoint{2.579500in}{2.600850in}}%
\pgfpathlineto{\pgfqpoint{2.579500in}{2.604000in}}%
\pgfpathlineto{\pgfqpoint{2.650375in}{2.604000in}}%
\pgfpathlineto{\pgfqpoint{2.650375in}{2.600850in}}%
\pgfpathlineto{\pgfqpoint{2.674000in}{2.600850in}}%
\pgfpathlineto{\pgfqpoint{2.674000in}{2.597700in}}%
\pgfpathlineto{\pgfqpoint{2.697625in}{2.597700in}}%
\pgfpathlineto{\pgfqpoint{2.697625in}{2.600850in}}%
\pgfpathlineto{\pgfqpoint{2.721250in}{2.600850in}}%
\pgfpathlineto{\pgfqpoint{2.721250in}{2.566200in}}%
\pgfpathlineto{\pgfqpoint{2.744875in}{2.566200in}}%
\pgfpathlineto{\pgfqpoint{2.744875in}{2.518950in}}%
\pgfpathlineto{\pgfqpoint{2.768500in}{2.518950in}}%
\pgfpathlineto{\pgfqpoint{2.768500in}{2.871750in}}%
\pgfpathlineto{\pgfqpoint{2.792125in}{2.871750in}}%
\pgfpathlineto{\pgfqpoint{2.792125in}{2.862300in}}%
\pgfpathlineto{\pgfqpoint{2.815750in}{2.862300in}}%
\pgfpathlineto{\pgfqpoint{2.815750in}{2.541000in}}%
\pgfpathlineto{\pgfqpoint{2.839375in}{2.541000in}}%
\pgfpathlineto{\pgfqpoint{2.839375in}{2.481150in}}%
\pgfpathlineto{\pgfqpoint{2.863000in}{2.481150in}}%
\pgfpathlineto{\pgfqpoint{2.863000in}{2.490600in}}%
\pgfpathlineto{\pgfqpoint{2.886625in}{2.490600in}}%
\pgfpathlineto{\pgfqpoint{2.886625in}{2.289000in}}%
\pgfpathlineto{\pgfqpoint{3.193750in}{2.289000in}}%
\pgfpathlineto{\pgfqpoint{3.193750in}{2.553600in}}%
\pgfpathlineto{\pgfqpoint{3.217375in}{2.553600in}}%
\pgfpathlineto{\pgfqpoint{3.217375in}{2.604000in}}%
\pgfpathlineto{\pgfqpoint{3.264625in}{2.604000in}}%
\pgfpathlineto{\pgfqpoint{3.264625in}{2.597700in}}%
\pgfpathlineto{\pgfqpoint{3.288250in}{2.597700in}}%
\pgfpathlineto{\pgfqpoint{3.288250in}{2.600850in}}%
\pgfpathlineto{\pgfqpoint{3.311875in}{2.600850in}}%
\pgfpathlineto{\pgfqpoint{3.311875in}{2.430750in}}%
\pgfpathlineto{\pgfqpoint{3.335500in}{2.430750in}}%
\pgfpathlineto{\pgfqpoint{3.335500in}{2.289000in}}%
\pgfpathlineto{\pgfqpoint{3.418188in}{2.289000in}}%
\pgfpathlineto{\pgfqpoint{3.418188in}{2.289000in}}%
\pgfusepath{stroke}%
\end{pgfscope}%
\begin{pgfscope}%
\pgfsetrectcap%
\pgfsetmiterjoin%
\pgfsetlinewidth{0.803000pt}%
\definecolor{currentstroke}{rgb}{0.101961,0.101961,0.101961}%
\pgfsetstrokecolor{currentstroke}%
\pgfsetdash{}{0pt}%
\pgfpathmoveto{\pgfqpoint{0.595000in}{2.289000in}}%
\pgfpathlineto{\pgfqpoint{0.595000in}{2.919000in}}%
\pgfusepath{stroke}%
\end{pgfscope}%
\begin{pgfscope}%
\pgfsetrectcap%
\pgfsetmiterjoin%
\pgfsetlinewidth{0.803000pt}%
\definecolor{currentstroke}{rgb}{0.101961,0.101961,0.101961}%
\pgfsetstrokecolor{currentstroke}%
\pgfsetdash{}{0pt}%
\pgfpathmoveto{\pgfqpoint{0.595000in}{2.289000in}}%
\pgfpathlineto{\pgfqpoint{3.430000in}{2.289000in}}%
\pgfusepath{stroke}%
\end{pgfscope}%
\begin{pgfscope}%
\definecolor{textcolor}{rgb}{0.156863,0.356863,0.568627}%
\pgfsetstrokecolor{textcolor}%
\pgfsetfillcolor{textcolor}%
\pgftext[x=0.623350in,y=2.900100in,left,top]{\color{textcolor}\sffamily\fontsize{7.000000}{8.400000}\selectfont warms (each paired with a touch)}%
\end{pgfscope}%
\begin{pgfscope}%
\pgfsetbuttcap%
\pgfsetmiterjoin%
\definecolor{currentfill}{rgb}{1.000000,1.000000,1.000000}%
\pgfsetfillcolor{currentfill}%
\pgfsetlinewidth{0.000000pt}%
\definecolor{currentstroke}{rgb}{0.000000,0.000000,0.000000}%
\pgfsetstrokecolor{currentstroke}%
\pgfsetstrokeopacity{0.000000}%
\pgfsetdash{}{0pt}%
\pgfpathmoveto{\pgfqpoint{0.595000in}{1.386000in}}%
\pgfpathlineto{\pgfqpoint{3.430000in}{1.386000in}}%
\pgfpathlineto{\pgfqpoint{3.430000in}{2.016000in}}%
\pgfpathlineto{\pgfqpoint{0.595000in}{2.016000in}}%
\pgfpathlineto{\pgfqpoint{0.595000in}{1.386000in}}%
\pgfpathclose%
\pgfusepath{fill}%
\end{pgfscope}%
\begin{pgfscope}%
\pgfpathrectangle{\pgfqpoint{0.595000in}{1.386000in}}{\pgfqpoint{2.835000in}{0.630000in}}%
\pgfusepath{clip}%
\pgfsetbuttcap%
\pgfsetmiterjoin%
\pgfsetlinewidth{0.000000pt}%
\definecolor{currentstroke}{rgb}{0.000000,0.000000,0.000000}%
\pgfsetstrokecolor{currentstroke}%
\pgfsetstrokeopacity{0.000000}%
\pgfsetdash{}{0pt}%
\pgfpathmoveto{\pgfqpoint{1.557719in}{1.386000in}}%
\pgfpathlineto{\pgfqpoint{1.557719in}{2.016000in}}%
\pgfpathlineto{\pgfqpoint{2.001869in}{2.016000in}}%
\pgfpathlineto{\pgfqpoint{2.001869in}{1.386000in}}%
\pgfpathlineto{\pgfqpoint{1.557719in}{1.386000in}}%
\pgfpathclose%
\pgfusepath{}%
\end{pgfscope}%
\begin{pgfscope}%
\pgfpathrectangle{\pgfqpoint{0.595000in}{1.386000in}}{\pgfqpoint{2.835000in}{0.630000in}}%
\pgfusepath{clip}%
\pgfsetrectcap%
\pgfsetroundjoin%
\pgfsetlinewidth{0.401500pt}%
\definecolor{currentstroke}{rgb}{0.901961,0.901961,0.901961}%
\pgfsetstrokecolor{currentstroke}%
\pgfsetdash{}{0pt}%
\pgfpathmoveto{\pgfqpoint{0.595000in}{1.386000in}}%
\pgfpathlineto{\pgfqpoint{3.430000in}{1.386000in}}%
\pgfusepath{stroke}%
\end{pgfscope}%
\begin{pgfscope}%
\definecolor{textcolor}{rgb}{0.101961,0.101961,0.101961}%
\pgfsetstrokecolor{textcolor}%
\pgfsetfillcolor{textcolor}%
\pgftext[x=0.494739in, y=1.352242in, left, base]{\color{textcolor}\sffamily\fontsize{7.000000}{8.400000}\selectfont 0}%
\end{pgfscope}%
\begin{pgfscope}%
\pgfpathrectangle{\pgfqpoint{0.595000in}{1.386000in}}{\pgfqpoint{2.835000in}{0.630000in}}%
\pgfusepath{clip}%
\pgfsetrectcap%
\pgfsetroundjoin%
\pgfsetlinewidth{0.401500pt}%
\definecolor{currentstroke}{rgb}{0.901961,0.901961,0.901961}%
\pgfsetstrokecolor{currentstroke}%
\pgfsetdash{}{0pt}%
\pgfpathmoveto{\pgfqpoint{0.595000in}{1.701000in}}%
\pgfpathlineto{\pgfqpoint{3.430000in}{1.701000in}}%
\pgfusepath{stroke}%
\end{pgfscope}%
\begin{pgfscope}%
\definecolor{textcolor}{rgb}{0.101961,0.101961,0.101961}%
\pgfsetstrokecolor{textcolor}%
\pgfsetfillcolor{textcolor}%
\pgftext[x=0.494739in, y=1.667242in, left, base]{\color{textcolor}\sffamily\fontsize{7.000000}{8.400000}\selectfont 2}%
\end{pgfscope}%
\begin{pgfscope}%
\pgfpathrectangle{\pgfqpoint{0.595000in}{1.386000in}}{\pgfqpoint{2.835000in}{0.630000in}}%
\pgfusepath{clip}%
\pgfsetrectcap%
\pgfsetroundjoin%
\pgfsetlinewidth{0.401500pt}%
\definecolor{currentstroke}{rgb}{0.901961,0.901961,0.901961}%
\pgfsetstrokecolor{currentstroke}%
\pgfsetdash{}{0pt}%
\pgfpathmoveto{\pgfqpoint{0.595000in}{2.016000in}}%
\pgfpathlineto{\pgfqpoint{3.430000in}{2.016000in}}%
\pgfusepath{stroke}%
\end{pgfscope}%
\begin{pgfscope}%
\definecolor{textcolor}{rgb}{0.101961,0.101961,0.101961}%
\pgfsetstrokecolor{textcolor}%
\pgfsetfillcolor{textcolor}%
\pgftext[x=0.494739in, y=1.982242in, left, base]{\color{textcolor}\sffamily\fontsize{7.000000}{8.400000}\selectfont 4}%
\end{pgfscope}%
\begin{pgfscope}%
\definecolor{textcolor}{rgb}{0.000000,0.000000,0.000000}%
\pgfsetstrokecolor{textcolor}%
\pgfsetfillcolor{textcolor}%
\pgftext[x=0.320357in, y=1.509794in, left, base,rotate=90.000000]{\color{textcolor}\sffamily\fontsize{7.500000}{9.000000}\selectfont Deferrals}%
\end{pgfscope}%
\begin{pgfscope}%
\definecolor{textcolor}{rgb}{0.000000,0.000000,0.000000}%
\pgfsetstrokecolor{textcolor}%
\pgfsetfillcolor{textcolor}%
\pgftext[x=0.420279in, y=1.526420in, left, base,rotate=90.000000]{\color{textcolor}\sffamily\fontsize{7.500000}{9.000000}\selectfont per 10 s}%
\end{pgfscope}%
\begin{pgfscope}%
\pgfpathrectangle{\pgfqpoint{0.595000in}{1.386000in}}{\pgfqpoint{2.835000in}{0.630000in}}%
\pgfusepath{clip}%
\pgfsetbuttcap%
\pgfsetmiterjoin%
\definecolor{currentfill}{rgb}{0.101961,0.101961,0.101961}%
\pgfsetfillcolor{currentfill}%
\pgfsetlinewidth{0.000000pt}%
\definecolor{currentstroke}{rgb}{0.000000,0.000000,0.000000}%
\pgfsetstrokecolor{currentstroke}%
\pgfsetstrokeopacity{0.000000}%
\pgfsetdash{}{0pt}%
\pgfpathmoveto{\pgfqpoint{0.596181in}{1.386000in}}%
\pgfpathlineto{\pgfqpoint{0.617444in}{1.386000in}}%
\pgfpathlineto{\pgfqpoint{0.617444in}{1.386000in}}%
\pgfpathlineto{\pgfqpoint{0.596181in}{1.386000in}}%
\pgfpathlineto{\pgfqpoint{0.596181in}{1.386000in}}%
\pgfpathclose%
\pgfusepath{fill}%
\end{pgfscope}%
\begin{pgfscope}%
\pgfpathrectangle{\pgfqpoint{0.595000in}{1.386000in}}{\pgfqpoint{2.835000in}{0.630000in}}%
\pgfusepath{clip}%
\pgfsetbuttcap%
\pgfsetmiterjoin%
\definecolor{currentfill}{rgb}{0.101961,0.101961,0.101961}%
\pgfsetfillcolor{currentfill}%
\pgfsetlinewidth{0.000000pt}%
\definecolor{currentstroke}{rgb}{0.000000,0.000000,0.000000}%
\pgfsetstrokecolor{currentstroke}%
\pgfsetstrokeopacity{0.000000}%
\pgfsetdash{}{0pt}%
\pgfpathmoveto{\pgfqpoint{0.619806in}{1.386000in}}%
\pgfpathlineto{\pgfqpoint{0.641069in}{1.386000in}}%
\pgfpathlineto{\pgfqpoint{0.641069in}{1.386000in}}%
\pgfpathlineto{\pgfqpoint{0.619806in}{1.386000in}}%
\pgfpathlineto{\pgfqpoint{0.619806in}{1.386000in}}%
\pgfpathclose%
\pgfusepath{fill}%
\end{pgfscope}%
\begin{pgfscope}%
\pgfpathrectangle{\pgfqpoint{0.595000in}{1.386000in}}{\pgfqpoint{2.835000in}{0.630000in}}%
\pgfusepath{clip}%
\pgfsetbuttcap%
\pgfsetmiterjoin%
\definecolor{currentfill}{rgb}{0.101961,0.101961,0.101961}%
\pgfsetfillcolor{currentfill}%
\pgfsetlinewidth{0.000000pt}%
\definecolor{currentstroke}{rgb}{0.000000,0.000000,0.000000}%
\pgfsetstrokecolor{currentstroke}%
\pgfsetstrokeopacity{0.000000}%
\pgfsetdash{}{0pt}%
\pgfpathmoveto{\pgfqpoint{0.643431in}{1.386000in}}%
\pgfpathlineto{\pgfqpoint{0.664694in}{1.386000in}}%
\pgfpathlineto{\pgfqpoint{0.664694in}{1.386000in}}%
\pgfpathlineto{\pgfqpoint{0.643431in}{1.386000in}}%
\pgfpathlineto{\pgfqpoint{0.643431in}{1.386000in}}%
\pgfpathclose%
\pgfusepath{fill}%
\end{pgfscope}%
\begin{pgfscope}%
\pgfpathrectangle{\pgfqpoint{0.595000in}{1.386000in}}{\pgfqpoint{2.835000in}{0.630000in}}%
\pgfusepath{clip}%
\pgfsetbuttcap%
\pgfsetmiterjoin%
\definecolor{currentfill}{rgb}{0.101961,0.101961,0.101961}%
\pgfsetfillcolor{currentfill}%
\pgfsetlinewidth{0.000000pt}%
\definecolor{currentstroke}{rgb}{0.000000,0.000000,0.000000}%
\pgfsetstrokecolor{currentstroke}%
\pgfsetstrokeopacity{0.000000}%
\pgfsetdash{}{0pt}%
\pgfpathmoveto{\pgfqpoint{0.667056in}{1.386000in}}%
\pgfpathlineto{\pgfqpoint{0.688319in}{1.386000in}}%
\pgfpathlineto{\pgfqpoint{0.688319in}{1.386000in}}%
\pgfpathlineto{\pgfqpoint{0.667056in}{1.386000in}}%
\pgfpathlineto{\pgfqpoint{0.667056in}{1.386000in}}%
\pgfpathclose%
\pgfusepath{fill}%
\end{pgfscope}%
\begin{pgfscope}%
\pgfpathrectangle{\pgfqpoint{0.595000in}{1.386000in}}{\pgfqpoint{2.835000in}{0.630000in}}%
\pgfusepath{clip}%
\pgfsetbuttcap%
\pgfsetmiterjoin%
\definecolor{currentfill}{rgb}{0.101961,0.101961,0.101961}%
\pgfsetfillcolor{currentfill}%
\pgfsetlinewidth{0.000000pt}%
\definecolor{currentstroke}{rgb}{0.000000,0.000000,0.000000}%
\pgfsetstrokecolor{currentstroke}%
\pgfsetstrokeopacity{0.000000}%
\pgfsetdash{}{0pt}%
\pgfpathmoveto{\pgfqpoint{0.690681in}{1.386000in}}%
\pgfpathlineto{\pgfqpoint{0.711944in}{1.386000in}}%
\pgfpathlineto{\pgfqpoint{0.711944in}{1.386000in}}%
\pgfpathlineto{\pgfqpoint{0.690681in}{1.386000in}}%
\pgfpathlineto{\pgfqpoint{0.690681in}{1.386000in}}%
\pgfpathclose%
\pgfusepath{fill}%
\end{pgfscope}%
\begin{pgfscope}%
\pgfpathrectangle{\pgfqpoint{0.595000in}{1.386000in}}{\pgfqpoint{2.835000in}{0.630000in}}%
\pgfusepath{clip}%
\pgfsetbuttcap%
\pgfsetmiterjoin%
\definecolor{currentfill}{rgb}{0.101961,0.101961,0.101961}%
\pgfsetfillcolor{currentfill}%
\pgfsetlinewidth{0.000000pt}%
\definecolor{currentstroke}{rgb}{0.000000,0.000000,0.000000}%
\pgfsetstrokecolor{currentstroke}%
\pgfsetstrokeopacity{0.000000}%
\pgfsetdash{}{0pt}%
\pgfpathmoveto{\pgfqpoint{0.714306in}{1.386000in}}%
\pgfpathlineto{\pgfqpoint{0.735569in}{1.386000in}}%
\pgfpathlineto{\pgfqpoint{0.735569in}{1.386000in}}%
\pgfpathlineto{\pgfqpoint{0.714306in}{1.386000in}}%
\pgfpathlineto{\pgfqpoint{0.714306in}{1.386000in}}%
\pgfpathclose%
\pgfusepath{fill}%
\end{pgfscope}%
\begin{pgfscope}%
\pgfpathrectangle{\pgfqpoint{0.595000in}{1.386000in}}{\pgfqpoint{2.835000in}{0.630000in}}%
\pgfusepath{clip}%
\pgfsetbuttcap%
\pgfsetmiterjoin%
\definecolor{currentfill}{rgb}{0.101961,0.101961,0.101961}%
\pgfsetfillcolor{currentfill}%
\pgfsetlinewidth{0.000000pt}%
\definecolor{currentstroke}{rgb}{0.000000,0.000000,0.000000}%
\pgfsetstrokecolor{currentstroke}%
\pgfsetstrokeopacity{0.000000}%
\pgfsetdash{}{0pt}%
\pgfpathmoveto{\pgfqpoint{0.737931in}{1.386000in}}%
\pgfpathlineto{\pgfqpoint{0.759194in}{1.386000in}}%
\pgfpathlineto{\pgfqpoint{0.759194in}{1.386000in}}%
\pgfpathlineto{\pgfqpoint{0.737931in}{1.386000in}}%
\pgfpathlineto{\pgfqpoint{0.737931in}{1.386000in}}%
\pgfpathclose%
\pgfusepath{fill}%
\end{pgfscope}%
\begin{pgfscope}%
\pgfpathrectangle{\pgfqpoint{0.595000in}{1.386000in}}{\pgfqpoint{2.835000in}{0.630000in}}%
\pgfusepath{clip}%
\pgfsetbuttcap%
\pgfsetmiterjoin%
\definecolor{currentfill}{rgb}{0.101961,0.101961,0.101961}%
\pgfsetfillcolor{currentfill}%
\pgfsetlinewidth{0.000000pt}%
\definecolor{currentstroke}{rgb}{0.000000,0.000000,0.000000}%
\pgfsetstrokecolor{currentstroke}%
\pgfsetstrokeopacity{0.000000}%
\pgfsetdash{}{0pt}%
\pgfpathmoveto{\pgfqpoint{0.761556in}{1.386000in}}%
\pgfpathlineto{\pgfqpoint{0.782819in}{1.386000in}}%
\pgfpathlineto{\pgfqpoint{0.782819in}{1.386000in}}%
\pgfpathlineto{\pgfqpoint{0.761556in}{1.386000in}}%
\pgfpathlineto{\pgfqpoint{0.761556in}{1.386000in}}%
\pgfpathclose%
\pgfusepath{fill}%
\end{pgfscope}%
\begin{pgfscope}%
\pgfpathrectangle{\pgfqpoint{0.595000in}{1.386000in}}{\pgfqpoint{2.835000in}{0.630000in}}%
\pgfusepath{clip}%
\pgfsetbuttcap%
\pgfsetmiterjoin%
\definecolor{currentfill}{rgb}{0.101961,0.101961,0.101961}%
\pgfsetfillcolor{currentfill}%
\pgfsetlinewidth{0.000000pt}%
\definecolor{currentstroke}{rgb}{0.000000,0.000000,0.000000}%
\pgfsetstrokecolor{currentstroke}%
\pgfsetstrokeopacity{0.000000}%
\pgfsetdash{}{0pt}%
\pgfpathmoveto{\pgfqpoint{0.785181in}{1.386000in}}%
\pgfpathlineto{\pgfqpoint{0.806444in}{1.386000in}}%
\pgfpathlineto{\pgfqpoint{0.806444in}{1.386000in}}%
\pgfpathlineto{\pgfqpoint{0.785181in}{1.386000in}}%
\pgfpathlineto{\pgfqpoint{0.785181in}{1.386000in}}%
\pgfpathclose%
\pgfusepath{fill}%
\end{pgfscope}%
\begin{pgfscope}%
\pgfpathrectangle{\pgfqpoint{0.595000in}{1.386000in}}{\pgfqpoint{2.835000in}{0.630000in}}%
\pgfusepath{clip}%
\pgfsetbuttcap%
\pgfsetmiterjoin%
\definecolor{currentfill}{rgb}{0.101961,0.101961,0.101961}%
\pgfsetfillcolor{currentfill}%
\pgfsetlinewidth{0.000000pt}%
\definecolor{currentstroke}{rgb}{0.000000,0.000000,0.000000}%
\pgfsetstrokecolor{currentstroke}%
\pgfsetstrokeopacity{0.000000}%
\pgfsetdash{}{0pt}%
\pgfpathmoveto{\pgfqpoint{0.808806in}{1.386000in}}%
\pgfpathlineto{\pgfqpoint{0.830069in}{1.386000in}}%
\pgfpathlineto{\pgfqpoint{0.830069in}{1.386000in}}%
\pgfpathlineto{\pgfqpoint{0.808806in}{1.386000in}}%
\pgfpathlineto{\pgfqpoint{0.808806in}{1.386000in}}%
\pgfpathclose%
\pgfusepath{fill}%
\end{pgfscope}%
\begin{pgfscope}%
\pgfpathrectangle{\pgfqpoint{0.595000in}{1.386000in}}{\pgfqpoint{2.835000in}{0.630000in}}%
\pgfusepath{clip}%
\pgfsetbuttcap%
\pgfsetmiterjoin%
\definecolor{currentfill}{rgb}{0.101961,0.101961,0.101961}%
\pgfsetfillcolor{currentfill}%
\pgfsetlinewidth{0.000000pt}%
\definecolor{currentstroke}{rgb}{0.000000,0.000000,0.000000}%
\pgfsetstrokecolor{currentstroke}%
\pgfsetstrokeopacity{0.000000}%
\pgfsetdash{}{0pt}%
\pgfpathmoveto{\pgfqpoint{0.832431in}{1.386000in}}%
\pgfpathlineto{\pgfqpoint{0.853694in}{1.386000in}}%
\pgfpathlineto{\pgfqpoint{0.853694in}{1.386000in}}%
\pgfpathlineto{\pgfqpoint{0.832431in}{1.386000in}}%
\pgfpathlineto{\pgfqpoint{0.832431in}{1.386000in}}%
\pgfpathclose%
\pgfusepath{fill}%
\end{pgfscope}%
\begin{pgfscope}%
\pgfpathrectangle{\pgfqpoint{0.595000in}{1.386000in}}{\pgfqpoint{2.835000in}{0.630000in}}%
\pgfusepath{clip}%
\pgfsetbuttcap%
\pgfsetmiterjoin%
\definecolor{currentfill}{rgb}{0.101961,0.101961,0.101961}%
\pgfsetfillcolor{currentfill}%
\pgfsetlinewidth{0.000000pt}%
\definecolor{currentstroke}{rgb}{0.000000,0.000000,0.000000}%
\pgfsetstrokecolor{currentstroke}%
\pgfsetstrokeopacity{0.000000}%
\pgfsetdash{}{0pt}%
\pgfpathmoveto{\pgfqpoint{0.856056in}{1.386000in}}%
\pgfpathlineto{\pgfqpoint{0.877319in}{1.386000in}}%
\pgfpathlineto{\pgfqpoint{0.877319in}{1.386000in}}%
\pgfpathlineto{\pgfqpoint{0.856056in}{1.386000in}}%
\pgfpathlineto{\pgfqpoint{0.856056in}{1.386000in}}%
\pgfpathclose%
\pgfusepath{fill}%
\end{pgfscope}%
\begin{pgfscope}%
\pgfpathrectangle{\pgfqpoint{0.595000in}{1.386000in}}{\pgfqpoint{2.835000in}{0.630000in}}%
\pgfusepath{clip}%
\pgfsetbuttcap%
\pgfsetmiterjoin%
\definecolor{currentfill}{rgb}{0.101961,0.101961,0.101961}%
\pgfsetfillcolor{currentfill}%
\pgfsetlinewidth{0.000000pt}%
\definecolor{currentstroke}{rgb}{0.000000,0.000000,0.000000}%
\pgfsetstrokecolor{currentstroke}%
\pgfsetstrokeopacity{0.000000}%
\pgfsetdash{}{0pt}%
\pgfpathmoveto{\pgfqpoint{0.879681in}{1.386000in}}%
\pgfpathlineto{\pgfqpoint{0.900944in}{1.386000in}}%
\pgfpathlineto{\pgfqpoint{0.900944in}{1.386000in}}%
\pgfpathlineto{\pgfqpoint{0.879681in}{1.386000in}}%
\pgfpathlineto{\pgfqpoint{0.879681in}{1.386000in}}%
\pgfpathclose%
\pgfusepath{fill}%
\end{pgfscope}%
\begin{pgfscope}%
\pgfpathrectangle{\pgfqpoint{0.595000in}{1.386000in}}{\pgfqpoint{2.835000in}{0.630000in}}%
\pgfusepath{clip}%
\pgfsetbuttcap%
\pgfsetmiterjoin%
\definecolor{currentfill}{rgb}{0.101961,0.101961,0.101961}%
\pgfsetfillcolor{currentfill}%
\pgfsetlinewidth{0.000000pt}%
\definecolor{currentstroke}{rgb}{0.000000,0.000000,0.000000}%
\pgfsetstrokecolor{currentstroke}%
\pgfsetstrokeopacity{0.000000}%
\pgfsetdash{}{0pt}%
\pgfpathmoveto{\pgfqpoint{0.903306in}{1.386000in}}%
\pgfpathlineto{\pgfqpoint{0.924569in}{1.386000in}}%
\pgfpathlineto{\pgfqpoint{0.924569in}{1.386000in}}%
\pgfpathlineto{\pgfqpoint{0.903306in}{1.386000in}}%
\pgfpathlineto{\pgfqpoint{0.903306in}{1.386000in}}%
\pgfpathclose%
\pgfusepath{fill}%
\end{pgfscope}%
\begin{pgfscope}%
\pgfpathrectangle{\pgfqpoint{0.595000in}{1.386000in}}{\pgfqpoint{2.835000in}{0.630000in}}%
\pgfusepath{clip}%
\pgfsetbuttcap%
\pgfsetmiterjoin%
\definecolor{currentfill}{rgb}{0.101961,0.101961,0.101961}%
\pgfsetfillcolor{currentfill}%
\pgfsetlinewidth{0.000000pt}%
\definecolor{currentstroke}{rgb}{0.000000,0.000000,0.000000}%
\pgfsetstrokecolor{currentstroke}%
\pgfsetstrokeopacity{0.000000}%
\pgfsetdash{}{0pt}%
\pgfpathmoveto{\pgfqpoint{0.926931in}{1.386000in}}%
\pgfpathlineto{\pgfqpoint{0.948194in}{1.386000in}}%
\pgfpathlineto{\pgfqpoint{0.948194in}{1.386000in}}%
\pgfpathlineto{\pgfqpoint{0.926931in}{1.386000in}}%
\pgfpathlineto{\pgfqpoint{0.926931in}{1.386000in}}%
\pgfpathclose%
\pgfusepath{fill}%
\end{pgfscope}%
\begin{pgfscope}%
\pgfpathrectangle{\pgfqpoint{0.595000in}{1.386000in}}{\pgfqpoint{2.835000in}{0.630000in}}%
\pgfusepath{clip}%
\pgfsetbuttcap%
\pgfsetmiterjoin%
\definecolor{currentfill}{rgb}{0.101961,0.101961,0.101961}%
\pgfsetfillcolor{currentfill}%
\pgfsetlinewidth{0.000000pt}%
\definecolor{currentstroke}{rgb}{0.000000,0.000000,0.000000}%
\pgfsetstrokecolor{currentstroke}%
\pgfsetstrokeopacity{0.000000}%
\pgfsetdash{}{0pt}%
\pgfpathmoveto{\pgfqpoint{0.950556in}{1.386000in}}%
\pgfpathlineto{\pgfqpoint{0.971819in}{1.386000in}}%
\pgfpathlineto{\pgfqpoint{0.971819in}{1.386000in}}%
\pgfpathlineto{\pgfqpoint{0.950556in}{1.386000in}}%
\pgfpathlineto{\pgfqpoint{0.950556in}{1.386000in}}%
\pgfpathclose%
\pgfusepath{fill}%
\end{pgfscope}%
\begin{pgfscope}%
\pgfpathrectangle{\pgfqpoint{0.595000in}{1.386000in}}{\pgfqpoint{2.835000in}{0.630000in}}%
\pgfusepath{clip}%
\pgfsetbuttcap%
\pgfsetmiterjoin%
\definecolor{currentfill}{rgb}{0.101961,0.101961,0.101961}%
\pgfsetfillcolor{currentfill}%
\pgfsetlinewidth{0.000000pt}%
\definecolor{currentstroke}{rgb}{0.000000,0.000000,0.000000}%
\pgfsetstrokecolor{currentstroke}%
\pgfsetstrokeopacity{0.000000}%
\pgfsetdash{}{0pt}%
\pgfpathmoveto{\pgfqpoint{0.974181in}{1.386000in}}%
\pgfpathlineto{\pgfqpoint{0.995444in}{1.386000in}}%
\pgfpathlineto{\pgfqpoint{0.995444in}{1.386000in}}%
\pgfpathlineto{\pgfqpoint{0.974181in}{1.386000in}}%
\pgfpathlineto{\pgfqpoint{0.974181in}{1.386000in}}%
\pgfpathclose%
\pgfusepath{fill}%
\end{pgfscope}%
\begin{pgfscope}%
\pgfpathrectangle{\pgfqpoint{0.595000in}{1.386000in}}{\pgfqpoint{2.835000in}{0.630000in}}%
\pgfusepath{clip}%
\pgfsetbuttcap%
\pgfsetmiterjoin%
\definecolor{currentfill}{rgb}{0.101961,0.101961,0.101961}%
\pgfsetfillcolor{currentfill}%
\pgfsetlinewidth{0.000000pt}%
\definecolor{currentstroke}{rgb}{0.000000,0.000000,0.000000}%
\pgfsetstrokecolor{currentstroke}%
\pgfsetstrokeopacity{0.000000}%
\pgfsetdash{}{0pt}%
\pgfpathmoveto{\pgfqpoint{0.997806in}{1.386000in}}%
\pgfpathlineto{\pgfqpoint{1.019069in}{1.386000in}}%
\pgfpathlineto{\pgfqpoint{1.019069in}{1.386000in}}%
\pgfpathlineto{\pgfqpoint{0.997806in}{1.386000in}}%
\pgfpathlineto{\pgfqpoint{0.997806in}{1.386000in}}%
\pgfpathclose%
\pgfusepath{fill}%
\end{pgfscope}%
\begin{pgfscope}%
\pgfpathrectangle{\pgfqpoint{0.595000in}{1.386000in}}{\pgfqpoint{2.835000in}{0.630000in}}%
\pgfusepath{clip}%
\pgfsetbuttcap%
\pgfsetmiterjoin%
\definecolor{currentfill}{rgb}{0.101961,0.101961,0.101961}%
\pgfsetfillcolor{currentfill}%
\pgfsetlinewidth{0.000000pt}%
\definecolor{currentstroke}{rgb}{0.000000,0.000000,0.000000}%
\pgfsetstrokecolor{currentstroke}%
\pgfsetstrokeopacity{0.000000}%
\pgfsetdash{}{0pt}%
\pgfpathmoveto{\pgfqpoint{1.021431in}{1.386000in}}%
\pgfpathlineto{\pgfqpoint{1.042694in}{1.386000in}}%
\pgfpathlineto{\pgfqpoint{1.042694in}{1.386000in}}%
\pgfpathlineto{\pgfqpoint{1.021431in}{1.386000in}}%
\pgfpathlineto{\pgfqpoint{1.021431in}{1.386000in}}%
\pgfpathclose%
\pgfusepath{fill}%
\end{pgfscope}%
\begin{pgfscope}%
\pgfpathrectangle{\pgfqpoint{0.595000in}{1.386000in}}{\pgfqpoint{2.835000in}{0.630000in}}%
\pgfusepath{clip}%
\pgfsetbuttcap%
\pgfsetmiterjoin%
\definecolor{currentfill}{rgb}{0.101961,0.101961,0.101961}%
\pgfsetfillcolor{currentfill}%
\pgfsetlinewidth{0.000000pt}%
\definecolor{currentstroke}{rgb}{0.000000,0.000000,0.000000}%
\pgfsetstrokecolor{currentstroke}%
\pgfsetstrokeopacity{0.000000}%
\pgfsetdash{}{0pt}%
\pgfpathmoveto{\pgfqpoint{1.045056in}{1.386000in}}%
\pgfpathlineto{\pgfqpoint{1.066319in}{1.386000in}}%
\pgfpathlineto{\pgfqpoint{1.066319in}{1.386000in}}%
\pgfpathlineto{\pgfqpoint{1.045056in}{1.386000in}}%
\pgfpathlineto{\pgfqpoint{1.045056in}{1.386000in}}%
\pgfpathclose%
\pgfusepath{fill}%
\end{pgfscope}%
\begin{pgfscope}%
\pgfpathrectangle{\pgfqpoint{0.595000in}{1.386000in}}{\pgfqpoint{2.835000in}{0.630000in}}%
\pgfusepath{clip}%
\pgfsetbuttcap%
\pgfsetmiterjoin%
\definecolor{currentfill}{rgb}{0.101961,0.101961,0.101961}%
\pgfsetfillcolor{currentfill}%
\pgfsetlinewidth{0.000000pt}%
\definecolor{currentstroke}{rgb}{0.000000,0.000000,0.000000}%
\pgfsetstrokecolor{currentstroke}%
\pgfsetstrokeopacity{0.000000}%
\pgfsetdash{}{0pt}%
\pgfpathmoveto{\pgfqpoint{1.068681in}{1.386000in}}%
\pgfpathlineto{\pgfqpoint{1.089944in}{1.386000in}}%
\pgfpathlineto{\pgfqpoint{1.089944in}{1.386000in}}%
\pgfpathlineto{\pgfqpoint{1.068681in}{1.386000in}}%
\pgfpathlineto{\pgfqpoint{1.068681in}{1.386000in}}%
\pgfpathclose%
\pgfusepath{fill}%
\end{pgfscope}%
\begin{pgfscope}%
\pgfpathrectangle{\pgfqpoint{0.595000in}{1.386000in}}{\pgfqpoint{2.835000in}{0.630000in}}%
\pgfusepath{clip}%
\pgfsetbuttcap%
\pgfsetmiterjoin%
\definecolor{currentfill}{rgb}{0.101961,0.101961,0.101961}%
\pgfsetfillcolor{currentfill}%
\pgfsetlinewidth{0.000000pt}%
\definecolor{currentstroke}{rgb}{0.000000,0.000000,0.000000}%
\pgfsetstrokecolor{currentstroke}%
\pgfsetstrokeopacity{0.000000}%
\pgfsetdash{}{0pt}%
\pgfpathmoveto{\pgfqpoint{1.092306in}{1.386000in}}%
\pgfpathlineto{\pgfqpoint{1.113569in}{1.386000in}}%
\pgfpathlineto{\pgfqpoint{1.113569in}{1.386000in}}%
\pgfpathlineto{\pgfqpoint{1.092306in}{1.386000in}}%
\pgfpathlineto{\pgfqpoint{1.092306in}{1.386000in}}%
\pgfpathclose%
\pgfusepath{fill}%
\end{pgfscope}%
\begin{pgfscope}%
\pgfpathrectangle{\pgfqpoint{0.595000in}{1.386000in}}{\pgfqpoint{2.835000in}{0.630000in}}%
\pgfusepath{clip}%
\pgfsetbuttcap%
\pgfsetmiterjoin%
\definecolor{currentfill}{rgb}{0.101961,0.101961,0.101961}%
\pgfsetfillcolor{currentfill}%
\pgfsetlinewidth{0.000000pt}%
\definecolor{currentstroke}{rgb}{0.000000,0.000000,0.000000}%
\pgfsetstrokecolor{currentstroke}%
\pgfsetstrokeopacity{0.000000}%
\pgfsetdash{}{0pt}%
\pgfpathmoveto{\pgfqpoint{1.115931in}{1.386000in}}%
\pgfpathlineto{\pgfqpoint{1.137194in}{1.386000in}}%
\pgfpathlineto{\pgfqpoint{1.137194in}{1.386000in}}%
\pgfpathlineto{\pgfqpoint{1.115931in}{1.386000in}}%
\pgfpathlineto{\pgfqpoint{1.115931in}{1.386000in}}%
\pgfpathclose%
\pgfusepath{fill}%
\end{pgfscope}%
\begin{pgfscope}%
\pgfpathrectangle{\pgfqpoint{0.595000in}{1.386000in}}{\pgfqpoint{2.835000in}{0.630000in}}%
\pgfusepath{clip}%
\pgfsetbuttcap%
\pgfsetmiterjoin%
\definecolor{currentfill}{rgb}{0.101961,0.101961,0.101961}%
\pgfsetfillcolor{currentfill}%
\pgfsetlinewidth{0.000000pt}%
\definecolor{currentstroke}{rgb}{0.000000,0.000000,0.000000}%
\pgfsetstrokecolor{currentstroke}%
\pgfsetstrokeopacity{0.000000}%
\pgfsetdash{}{0pt}%
\pgfpathmoveto{\pgfqpoint{1.139556in}{1.386000in}}%
\pgfpathlineto{\pgfqpoint{1.160819in}{1.386000in}}%
\pgfpathlineto{\pgfqpoint{1.160819in}{1.386000in}}%
\pgfpathlineto{\pgfqpoint{1.139556in}{1.386000in}}%
\pgfpathlineto{\pgfqpoint{1.139556in}{1.386000in}}%
\pgfpathclose%
\pgfusepath{fill}%
\end{pgfscope}%
\begin{pgfscope}%
\pgfpathrectangle{\pgfqpoint{0.595000in}{1.386000in}}{\pgfqpoint{2.835000in}{0.630000in}}%
\pgfusepath{clip}%
\pgfsetbuttcap%
\pgfsetmiterjoin%
\definecolor{currentfill}{rgb}{0.101961,0.101961,0.101961}%
\pgfsetfillcolor{currentfill}%
\pgfsetlinewidth{0.000000pt}%
\definecolor{currentstroke}{rgb}{0.000000,0.000000,0.000000}%
\pgfsetstrokecolor{currentstroke}%
\pgfsetstrokeopacity{0.000000}%
\pgfsetdash{}{0pt}%
\pgfpathmoveto{\pgfqpoint{1.163181in}{1.386000in}}%
\pgfpathlineto{\pgfqpoint{1.184444in}{1.386000in}}%
\pgfpathlineto{\pgfqpoint{1.184444in}{1.386000in}}%
\pgfpathlineto{\pgfqpoint{1.163181in}{1.386000in}}%
\pgfpathlineto{\pgfqpoint{1.163181in}{1.386000in}}%
\pgfpathclose%
\pgfusepath{fill}%
\end{pgfscope}%
\begin{pgfscope}%
\pgfpathrectangle{\pgfqpoint{0.595000in}{1.386000in}}{\pgfqpoint{2.835000in}{0.630000in}}%
\pgfusepath{clip}%
\pgfsetbuttcap%
\pgfsetmiterjoin%
\definecolor{currentfill}{rgb}{0.101961,0.101961,0.101961}%
\pgfsetfillcolor{currentfill}%
\pgfsetlinewidth{0.000000pt}%
\definecolor{currentstroke}{rgb}{0.000000,0.000000,0.000000}%
\pgfsetstrokecolor{currentstroke}%
\pgfsetstrokeopacity{0.000000}%
\pgfsetdash{}{0pt}%
\pgfpathmoveto{\pgfqpoint{1.186806in}{1.386000in}}%
\pgfpathlineto{\pgfqpoint{1.208069in}{1.386000in}}%
\pgfpathlineto{\pgfqpoint{1.208069in}{1.386000in}}%
\pgfpathlineto{\pgfqpoint{1.186806in}{1.386000in}}%
\pgfpathlineto{\pgfqpoint{1.186806in}{1.386000in}}%
\pgfpathclose%
\pgfusepath{fill}%
\end{pgfscope}%
\begin{pgfscope}%
\pgfpathrectangle{\pgfqpoint{0.595000in}{1.386000in}}{\pgfqpoint{2.835000in}{0.630000in}}%
\pgfusepath{clip}%
\pgfsetbuttcap%
\pgfsetmiterjoin%
\definecolor{currentfill}{rgb}{0.101961,0.101961,0.101961}%
\pgfsetfillcolor{currentfill}%
\pgfsetlinewidth{0.000000pt}%
\definecolor{currentstroke}{rgb}{0.000000,0.000000,0.000000}%
\pgfsetstrokecolor{currentstroke}%
\pgfsetstrokeopacity{0.000000}%
\pgfsetdash{}{0pt}%
\pgfpathmoveto{\pgfqpoint{1.210431in}{1.386000in}}%
\pgfpathlineto{\pgfqpoint{1.231694in}{1.386000in}}%
\pgfpathlineto{\pgfqpoint{1.231694in}{1.386000in}}%
\pgfpathlineto{\pgfqpoint{1.210431in}{1.386000in}}%
\pgfpathlineto{\pgfqpoint{1.210431in}{1.386000in}}%
\pgfpathclose%
\pgfusepath{fill}%
\end{pgfscope}%
\begin{pgfscope}%
\pgfpathrectangle{\pgfqpoint{0.595000in}{1.386000in}}{\pgfqpoint{2.835000in}{0.630000in}}%
\pgfusepath{clip}%
\pgfsetbuttcap%
\pgfsetmiterjoin%
\definecolor{currentfill}{rgb}{0.101961,0.101961,0.101961}%
\pgfsetfillcolor{currentfill}%
\pgfsetlinewidth{0.000000pt}%
\definecolor{currentstroke}{rgb}{0.000000,0.000000,0.000000}%
\pgfsetstrokecolor{currentstroke}%
\pgfsetstrokeopacity{0.000000}%
\pgfsetdash{}{0pt}%
\pgfpathmoveto{\pgfqpoint{1.234056in}{1.386000in}}%
\pgfpathlineto{\pgfqpoint{1.255319in}{1.386000in}}%
\pgfpathlineto{\pgfqpoint{1.255319in}{1.386000in}}%
\pgfpathlineto{\pgfqpoint{1.234056in}{1.386000in}}%
\pgfpathlineto{\pgfqpoint{1.234056in}{1.386000in}}%
\pgfpathclose%
\pgfusepath{fill}%
\end{pgfscope}%
\begin{pgfscope}%
\pgfpathrectangle{\pgfqpoint{0.595000in}{1.386000in}}{\pgfqpoint{2.835000in}{0.630000in}}%
\pgfusepath{clip}%
\pgfsetbuttcap%
\pgfsetmiterjoin%
\definecolor{currentfill}{rgb}{0.101961,0.101961,0.101961}%
\pgfsetfillcolor{currentfill}%
\pgfsetlinewidth{0.000000pt}%
\definecolor{currentstroke}{rgb}{0.000000,0.000000,0.000000}%
\pgfsetstrokecolor{currentstroke}%
\pgfsetstrokeopacity{0.000000}%
\pgfsetdash{}{0pt}%
\pgfpathmoveto{\pgfqpoint{1.257681in}{1.386000in}}%
\pgfpathlineto{\pgfqpoint{1.278944in}{1.386000in}}%
\pgfpathlineto{\pgfqpoint{1.278944in}{1.386000in}}%
\pgfpathlineto{\pgfqpoint{1.257681in}{1.386000in}}%
\pgfpathlineto{\pgfqpoint{1.257681in}{1.386000in}}%
\pgfpathclose%
\pgfusepath{fill}%
\end{pgfscope}%
\begin{pgfscope}%
\pgfpathrectangle{\pgfqpoint{0.595000in}{1.386000in}}{\pgfqpoint{2.835000in}{0.630000in}}%
\pgfusepath{clip}%
\pgfsetbuttcap%
\pgfsetmiterjoin%
\definecolor{currentfill}{rgb}{0.101961,0.101961,0.101961}%
\pgfsetfillcolor{currentfill}%
\pgfsetlinewidth{0.000000pt}%
\definecolor{currentstroke}{rgb}{0.000000,0.000000,0.000000}%
\pgfsetstrokecolor{currentstroke}%
\pgfsetstrokeopacity{0.000000}%
\pgfsetdash{}{0pt}%
\pgfpathmoveto{\pgfqpoint{1.281306in}{1.386000in}}%
\pgfpathlineto{\pgfqpoint{1.302569in}{1.386000in}}%
\pgfpathlineto{\pgfqpoint{1.302569in}{1.701000in}}%
\pgfpathlineto{\pgfqpoint{1.281306in}{1.701000in}}%
\pgfpathlineto{\pgfqpoint{1.281306in}{1.386000in}}%
\pgfpathclose%
\pgfusepath{fill}%
\end{pgfscope}%
\begin{pgfscope}%
\pgfpathrectangle{\pgfqpoint{0.595000in}{1.386000in}}{\pgfqpoint{2.835000in}{0.630000in}}%
\pgfusepath{clip}%
\pgfsetbuttcap%
\pgfsetmiterjoin%
\definecolor{currentfill}{rgb}{0.101961,0.101961,0.101961}%
\pgfsetfillcolor{currentfill}%
\pgfsetlinewidth{0.000000pt}%
\definecolor{currentstroke}{rgb}{0.000000,0.000000,0.000000}%
\pgfsetstrokecolor{currentstroke}%
\pgfsetstrokeopacity{0.000000}%
\pgfsetdash{}{0pt}%
\pgfpathmoveto{\pgfqpoint{1.304931in}{1.386000in}}%
\pgfpathlineto{\pgfqpoint{1.326194in}{1.386000in}}%
\pgfpathlineto{\pgfqpoint{1.326194in}{1.701000in}}%
\pgfpathlineto{\pgfqpoint{1.304931in}{1.701000in}}%
\pgfpathlineto{\pgfqpoint{1.304931in}{1.386000in}}%
\pgfpathclose%
\pgfusepath{fill}%
\end{pgfscope}%
\begin{pgfscope}%
\pgfpathrectangle{\pgfqpoint{0.595000in}{1.386000in}}{\pgfqpoint{2.835000in}{0.630000in}}%
\pgfusepath{clip}%
\pgfsetbuttcap%
\pgfsetmiterjoin%
\definecolor{currentfill}{rgb}{0.101961,0.101961,0.101961}%
\pgfsetfillcolor{currentfill}%
\pgfsetlinewidth{0.000000pt}%
\definecolor{currentstroke}{rgb}{0.000000,0.000000,0.000000}%
\pgfsetstrokecolor{currentstroke}%
\pgfsetstrokeopacity{0.000000}%
\pgfsetdash{}{0pt}%
\pgfpathmoveto{\pgfqpoint{1.328556in}{1.386000in}}%
\pgfpathlineto{\pgfqpoint{1.349819in}{1.386000in}}%
\pgfpathlineto{\pgfqpoint{1.349819in}{1.386000in}}%
\pgfpathlineto{\pgfqpoint{1.328556in}{1.386000in}}%
\pgfpathlineto{\pgfqpoint{1.328556in}{1.386000in}}%
\pgfpathclose%
\pgfusepath{fill}%
\end{pgfscope}%
\begin{pgfscope}%
\pgfpathrectangle{\pgfqpoint{0.595000in}{1.386000in}}{\pgfqpoint{2.835000in}{0.630000in}}%
\pgfusepath{clip}%
\pgfsetbuttcap%
\pgfsetmiterjoin%
\definecolor{currentfill}{rgb}{0.101961,0.101961,0.101961}%
\pgfsetfillcolor{currentfill}%
\pgfsetlinewidth{0.000000pt}%
\definecolor{currentstroke}{rgb}{0.000000,0.000000,0.000000}%
\pgfsetstrokecolor{currentstroke}%
\pgfsetstrokeopacity{0.000000}%
\pgfsetdash{}{0pt}%
\pgfpathmoveto{\pgfqpoint{1.352181in}{1.386000in}}%
\pgfpathlineto{\pgfqpoint{1.373444in}{1.386000in}}%
\pgfpathlineto{\pgfqpoint{1.373444in}{1.543500in}}%
\pgfpathlineto{\pgfqpoint{1.352181in}{1.543500in}}%
\pgfpathlineto{\pgfqpoint{1.352181in}{1.386000in}}%
\pgfpathclose%
\pgfusepath{fill}%
\end{pgfscope}%
\begin{pgfscope}%
\pgfpathrectangle{\pgfqpoint{0.595000in}{1.386000in}}{\pgfqpoint{2.835000in}{0.630000in}}%
\pgfusepath{clip}%
\pgfsetbuttcap%
\pgfsetmiterjoin%
\definecolor{currentfill}{rgb}{0.101961,0.101961,0.101961}%
\pgfsetfillcolor{currentfill}%
\pgfsetlinewidth{0.000000pt}%
\definecolor{currentstroke}{rgb}{0.000000,0.000000,0.000000}%
\pgfsetstrokecolor{currentstroke}%
\pgfsetstrokeopacity{0.000000}%
\pgfsetdash{}{0pt}%
\pgfpathmoveto{\pgfqpoint{1.375806in}{1.386000in}}%
\pgfpathlineto{\pgfqpoint{1.397069in}{1.386000in}}%
\pgfpathlineto{\pgfqpoint{1.397069in}{1.543500in}}%
\pgfpathlineto{\pgfqpoint{1.375806in}{1.543500in}}%
\pgfpathlineto{\pgfqpoint{1.375806in}{1.386000in}}%
\pgfpathclose%
\pgfusepath{fill}%
\end{pgfscope}%
\begin{pgfscope}%
\pgfpathrectangle{\pgfqpoint{0.595000in}{1.386000in}}{\pgfqpoint{2.835000in}{0.630000in}}%
\pgfusepath{clip}%
\pgfsetbuttcap%
\pgfsetmiterjoin%
\definecolor{currentfill}{rgb}{0.101961,0.101961,0.101961}%
\pgfsetfillcolor{currentfill}%
\pgfsetlinewidth{0.000000pt}%
\definecolor{currentstroke}{rgb}{0.000000,0.000000,0.000000}%
\pgfsetstrokecolor{currentstroke}%
\pgfsetstrokeopacity{0.000000}%
\pgfsetdash{}{0pt}%
\pgfpathmoveto{\pgfqpoint{1.399431in}{1.386000in}}%
\pgfpathlineto{\pgfqpoint{1.420694in}{1.386000in}}%
\pgfpathlineto{\pgfqpoint{1.420694in}{1.386000in}}%
\pgfpathlineto{\pgfqpoint{1.399431in}{1.386000in}}%
\pgfpathlineto{\pgfqpoint{1.399431in}{1.386000in}}%
\pgfpathclose%
\pgfusepath{fill}%
\end{pgfscope}%
\begin{pgfscope}%
\pgfpathrectangle{\pgfqpoint{0.595000in}{1.386000in}}{\pgfqpoint{2.835000in}{0.630000in}}%
\pgfusepath{clip}%
\pgfsetbuttcap%
\pgfsetmiterjoin%
\definecolor{currentfill}{rgb}{0.101961,0.101961,0.101961}%
\pgfsetfillcolor{currentfill}%
\pgfsetlinewidth{0.000000pt}%
\definecolor{currentstroke}{rgb}{0.000000,0.000000,0.000000}%
\pgfsetstrokecolor{currentstroke}%
\pgfsetstrokeopacity{0.000000}%
\pgfsetdash{}{0pt}%
\pgfpathmoveto{\pgfqpoint{1.423056in}{1.386000in}}%
\pgfpathlineto{\pgfqpoint{1.444319in}{1.386000in}}%
\pgfpathlineto{\pgfqpoint{1.444319in}{1.386000in}}%
\pgfpathlineto{\pgfqpoint{1.423056in}{1.386000in}}%
\pgfpathlineto{\pgfqpoint{1.423056in}{1.386000in}}%
\pgfpathclose%
\pgfusepath{fill}%
\end{pgfscope}%
\begin{pgfscope}%
\pgfpathrectangle{\pgfqpoint{0.595000in}{1.386000in}}{\pgfqpoint{2.835000in}{0.630000in}}%
\pgfusepath{clip}%
\pgfsetbuttcap%
\pgfsetmiterjoin%
\definecolor{currentfill}{rgb}{0.101961,0.101961,0.101961}%
\pgfsetfillcolor{currentfill}%
\pgfsetlinewidth{0.000000pt}%
\definecolor{currentstroke}{rgb}{0.000000,0.000000,0.000000}%
\pgfsetstrokecolor{currentstroke}%
\pgfsetstrokeopacity{0.000000}%
\pgfsetdash{}{0pt}%
\pgfpathmoveto{\pgfqpoint{1.446681in}{1.386000in}}%
\pgfpathlineto{\pgfqpoint{1.467944in}{1.386000in}}%
\pgfpathlineto{\pgfqpoint{1.467944in}{1.701000in}}%
\pgfpathlineto{\pgfqpoint{1.446681in}{1.701000in}}%
\pgfpathlineto{\pgfqpoint{1.446681in}{1.386000in}}%
\pgfpathclose%
\pgfusepath{fill}%
\end{pgfscope}%
\begin{pgfscope}%
\pgfpathrectangle{\pgfqpoint{0.595000in}{1.386000in}}{\pgfqpoint{2.835000in}{0.630000in}}%
\pgfusepath{clip}%
\pgfsetbuttcap%
\pgfsetmiterjoin%
\definecolor{currentfill}{rgb}{0.101961,0.101961,0.101961}%
\pgfsetfillcolor{currentfill}%
\pgfsetlinewidth{0.000000pt}%
\definecolor{currentstroke}{rgb}{0.000000,0.000000,0.000000}%
\pgfsetstrokecolor{currentstroke}%
\pgfsetstrokeopacity{0.000000}%
\pgfsetdash{}{0pt}%
\pgfpathmoveto{\pgfqpoint{1.470306in}{1.386000in}}%
\pgfpathlineto{\pgfqpoint{1.491569in}{1.386000in}}%
\pgfpathlineto{\pgfqpoint{1.491569in}{1.386000in}}%
\pgfpathlineto{\pgfqpoint{1.470306in}{1.386000in}}%
\pgfpathlineto{\pgfqpoint{1.470306in}{1.386000in}}%
\pgfpathclose%
\pgfusepath{fill}%
\end{pgfscope}%
\begin{pgfscope}%
\pgfpathrectangle{\pgfqpoint{0.595000in}{1.386000in}}{\pgfqpoint{2.835000in}{0.630000in}}%
\pgfusepath{clip}%
\pgfsetbuttcap%
\pgfsetmiterjoin%
\definecolor{currentfill}{rgb}{0.101961,0.101961,0.101961}%
\pgfsetfillcolor{currentfill}%
\pgfsetlinewidth{0.000000pt}%
\definecolor{currentstroke}{rgb}{0.000000,0.000000,0.000000}%
\pgfsetstrokecolor{currentstroke}%
\pgfsetstrokeopacity{0.000000}%
\pgfsetdash{}{0pt}%
\pgfpathmoveto{\pgfqpoint{1.493931in}{1.386000in}}%
\pgfpathlineto{\pgfqpoint{1.515194in}{1.386000in}}%
\pgfpathlineto{\pgfqpoint{1.515194in}{1.386000in}}%
\pgfpathlineto{\pgfqpoint{1.493931in}{1.386000in}}%
\pgfpathlineto{\pgfqpoint{1.493931in}{1.386000in}}%
\pgfpathclose%
\pgfusepath{fill}%
\end{pgfscope}%
\begin{pgfscope}%
\pgfpathrectangle{\pgfqpoint{0.595000in}{1.386000in}}{\pgfqpoint{2.835000in}{0.630000in}}%
\pgfusepath{clip}%
\pgfsetbuttcap%
\pgfsetmiterjoin%
\definecolor{currentfill}{rgb}{0.101961,0.101961,0.101961}%
\pgfsetfillcolor{currentfill}%
\pgfsetlinewidth{0.000000pt}%
\definecolor{currentstroke}{rgb}{0.000000,0.000000,0.000000}%
\pgfsetstrokecolor{currentstroke}%
\pgfsetstrokeopacity{0.000000}%
\pgfsetdash{}{0pt}%
\pgfpathmoveto{\pgfqpoint{1.517556in}{1.386000in}}%
\pgfpathlineto{\pgfqpoint{1.538819in}{1.386000in}}%
\pgfpathlineto{\pgfqpoint{1.538819in}{1.386000in}}%
\pgfpathlineto{\pgfqpoint{1.517556in}{1.386000in}}%
\pgfpathlineto{\pgfqpoint{1.517556in}{1.386000in}}%
\pgfpathclose%
\pgfusepath{fill}%
\end{pgfscope}%
\begin{pgfscope}%
\pgfpathrectangle{\pgfqpoint{0.595000in}{1.386000in}}{\pgfqpoint{2.835000in}{0.630000in}}%
\pgfusepath{clip}%
\pgfsetbuttcap%
\pgfsetmiterjoin%
\definecolor{currentfill}{rgb}{0.101961,0.101961,0.101961}%
\pgfsetfillcolor{currentfill}%
\pgfsetlinewidth{0.000000pt}%
\definecolor{currentstroke}{rgb}{0.000000,0.000000,0.000000}%
\pgfsetstrokecolor{currentstroke}%
\pgfsetstrokeopacity{0.000000}%
\pgfsetdash{}{0pt}%
\pgfpathmoveto{\pgfqpoint{1.541181in}{1.386000in}}%
\pgfpathlineto{\pgfqpoint{1.562444in}{1.386000in}}%
\pgfpathlineto{\pgfqpoint{1.562444in}{1.386000in}}%
\pgfpathlineto{\pgfqpoint{1.541181in}{1.386000in}}%
\pgfpathlineto{\pgfqpoint{1.541181in}{1.386000in}}%
\pgfpathclose%
\pgfusepath{fill}%
\end{pgfscope}%
\begin{pgfscope}%
\pgfpathrectangle{\pgfqpoint{0.595000in}{1.386000in}}{\pgfqpoint{2.835000in}{0.630000in}}%
\pgfusepath{clip}%
\pgfsetbuttcap%
\pgfsetmiterjoin%
\definecolor{currentfill}{rgb}{0.101961,0.101961,0.101961}%
\pgfsetfillcolor{currentfill}%
\pgfsetlinewidth{0.000000pt}%
\definecolor{currentstroke}{rgb}{0.000000,0.000000,0.000000}%
\pgfsetstrokecolor{currentstroke}%
\pgfsetstrokeopacity{0.000000}%
\pgfsetdash{}{0pt}%
\pgfpathmoveto{\pgfqpoint{1.564806in}{1.386000in}}%
\pgfpathlineto{\pgfqpoint{1.586069in}{1.386000in}}%
\pgfpathlineto{\pgfqpoint{1.586069in}{1.543500in}}%
\pgfpathlineto{\pgfqpoint{1.564806in}{1.543500in}}%
\pgfpathlineto{\pgfqpoint{1.564806in}{1.386000in}}%
\pgfpathclose%
\pgfusepath{fill}%
\end{pgfscope}%
\begin{pgfscope}%
\pgfpathrectangle{\pgfqpoint{0.595000in}{1.386000in}}{\pgfqpoint{2.835000in}{0.630000in}}%
\pgfusepath{clip}%
\pgfsetbuttcap%
\pgfsetmiterjoin%
\definecolor{currentfill}{rgb}{0.101961,0.101961,0.101961}%
\pgfsetfillcolor{currentfill}%
\pgfsetlinewidth{0.000000pt}%
\definecolor{currentstroke}{rgb}{0.000000,0.000000,0.000000}%
\pgfsetstrokecolor{currentstroke}%
\pgfsetstrokeopacity{0.000000}%
\pgfsetdash{}{0pt}%
\pgfpathmoveto{\pgfqpoint{1.588431in}{1.386000in}}%
\pgfpathlineto{\pgfqpoint{1.609694in}{1.386000in}}%
\pgfpathlineto{\pgfqpoint{1.609694in}{1.386000in}}%
\pgfpathlineto{\pgfqpoint{1.588431in}{1.386000in}}%
\pgfpathlineto{\pgfqpoint{1.588431in}{1.386000in}}%
\pgfpathclose%
\pgfusepath{fill}%
\end{pgfscope}%
\begin{pgfscope}%
\pgfpathrectangle{\pgfqpoint{0.595000in}{1.386000in}}{\pgfqpoint{2.835000in}{0.630000in}}%
\pgfusepath{clip}%
\pgfsetbuttcap%
\pgfsetmiterjoin%
\definecolor{currentfill}{rgb}{0.101961,0.101961,0.101961}%
\pgfsetfillcolor{currentfill}%
\pgfsetlinewidth{0.000000pt}%
\definecolor{currentstroke}{rgb}{0.000000,0.000000,0.000000}%
\pgfsetstrokecolor{currentstroke}%
\pgfsetstrokeopacity{0.000000}%
\pgfsetdash{}{0pt}%
\pgfpathmoveto{\pgfqpoint{1.612056in}{1.386000in}}%
\pgfpathlineto{\pgfqpoint{1.633319in}{1.386000in}}%
\pgfpathlineto{\pgfqpoint{1.633319in}{1.543500in}}%
\pgfpathlineto{\pgfqpoint{1.612056in}{1.543500in}}%
\pgfpathlineto{\pgfqpoint{1.612056in}{1.386000in}}%
\pgfpathclose%
\pgfusepath{fill}%
\end{pgfscope}%
\begin{pgfscope}%
\pgfpathrectangle{\pgfqpoint{0.595000in}{1.386000in}}{\pgfqpoint{2.835000in}{0.630000in}}%
\pgfusepath{clip}%
\pgfsetbuttcap%
\pgfsetmiterjoin%
\definecolor{currentfill}{rgb}{0.101961,0.101961,0.101961}%
\pgfsetfillcolor{currentfill}%
\pgfsetlinewidth{0.000000pt}%
\definecolor{currentstroke}{rgb}{0.000000,0.000000,0.000000}%
\pgfsetstrokecolor{currentstroke}%
\pgfsetstrokeopacity{0.000000}%
\pgfsetdash{}{0pt}%
\pgfpathmoveto{\pgfqpoint{1.635681in}{1.386000in}}%
\pgfpathlineto{\pgfqpoint{1.656944in}{1.386000in}}%
\pgfpathlineto{\pgfqpoint{1.656944in}{1.858500in}}%
\pgfpathlineto{\pgfqpoint{1.635681in}{1.858500in}}%
\pgfpathlineto{\pgfqpoint{1.635681in}{1.386000in}}%
\pgfpathclose%
\pgfusepath{fill}%
\end{pgfscope}%
\begin{pgfscope}%
\pgfpathrectangle{\pgfqpoint{0.595000in}{1.386000in}}{\pgfqpoint{2.835000in}{0.630000in}}%
\pgfusepath{clip}%
\pgfsetbuttcap%
\pgfsetmiterjoin%
\definecolor{currentfill}{rgb}{0.101961,0.101961,0.101961}%
\pgfsetfillcolor{currentfill}%
\pgfsetlinewidth{0.000000pt}%
\definecolor{currentstroke}{rgb}{0.000000,0.000000,0.000000}%
\pgfsetstrokecolor{currentstroke}%
\pgfsetstrokeopacity{0.000000}%
\pgfsetdash{}{0pt}%
\pgfpathmoveto{\pgfqpoint{1.659306in}{1.386000in}}%
\pgfpathlineto{\pgfqpoint{1.680569in}{1.386000in}}%
\pgfpathlineto{\pgfqpoint{1.680569in}{1.386000in}}%
\pgfpathlineto{\pgfqpoint{1.659306in}{1.386000in}}%
\pgfpathlineto{\pgfqpoint{1.659306in}{1.386000in}}%
\pgfpathclose%
\pgfusepath{fill}%
\end{pgfscope}%
\begin{pgfscope}%
\pgfpathrectangle{\pgfqpoint{0.595000in}{1.386000in}}{\pgfqpoint{2.835000in}{0.630000in}}%
\pgfusepath{clip}%
\pgfsetbuttcap%
\pgfsetmiterjoin%
\definecolor{currentfill}{rgb}{0.101961,0.101961,0.101961}%
\pgfsetfillcolor{currentfill}%
\pgfsetlinewidth{0.000000pt}%
\definecolor{currentstroke}{rgb}{0.000000,0.000000,0.000000}%
\pgfsetstrokecolor{currentstroke}%
\pgfsetstrokeopacity{0.000000}%
\pgfsetdash{}{0pt}%
\pgfpathmoveto{\pgfqpoint{1.682931in}{1.386000in}}%
\pgfpathlineto{\pgfqpoint{1.704194in}{1.386000in}}%
\pgfpathlineto{\pgfqpoint{1.704194in}{1.386000in}}%
\pgfpathlineto{\pgfqpoint{1.682931in}{1.386000in}}%
\pgfpathlineto{\pgfqpoint{1.682931in}{1.386000in}}%
\pgfpathclose%
\pgfusepath{fill}%
\end{pgfscope}%
\begin{pgfscope}%
\pgfpathrectangle{\pgfqpoint{0.595000in}{1.386000in}}{\pgfqpoint{2.835000in}{0.630000in}}%
\pgfusepath{clip}%
\pgfsetbuttcap%
\pgfsetmiterjoin%
\definecolor{currentfill}{rgb}{0.101961,0.101961,0.101961}%
\pgfsetfillcolor{currentfill}%
\pgfsetlinewidth{0.000000pt}%
\definecolor{currentstroke}{rgb}{0.000000,0.000000,0.000000}%
\pgfsetstrokecolor{currentstroke}%
\pgfsetstrokeopacity{0.000000}%
\pgfsetdash{}{0pt}%
\pgfpathmoveto{\pgfqpoint{1.706556in}{1.386000in}}%
\pgfpathlineto{\pgfqpoint{1.727819in}{1.386000in}}%
\pgfpathlineto{\pgfqpoint{1.727819in}{1.386000in}}%
\pgfpathlineto{\pgfqpoint{1.706556in}{1.386000in}}%
\pgfpathlineto{\pgfqpoint{1.706556in}{1.386000in}}%
\pgfpathclose%
\pgfusepath{fill}%
\end{pgfscope}%
\begin{pgfscope}%
\pgfpathrectangle{\pgfqpoint{0.595000in}{1.386000in}}{\pgfqpoint{2.835000in}{0.630000in}}%
\pgfusepath{clip}%
\pgfsetbuttcap%
\pgfsetmiterjoin%
\definecolor{currentfill}{rgb}{0.101961,0.101961,0.101961}%
\pgfsetfillcolor{currentfill}%
\pgfsetlinewidth{0.000000pt}%
\definecolor{currentstroke}{rgb}{0.000000,0.000000,0.000000}%
\pgfsetstrokecolor{currentstroke}%
\pgfsetstrokeopacity{0.000000}%
\pgfsetdash{}{0pt}%
\pgfpathmoveto{\pgfqpoint{1.730181in}{1.386000in}}%
\pgfpathlineto{\pgfqpoint{1.751444in}{1.386000in}}%
\pgfpathlineto{\pgfqpoint{1.751444in}{1.386000in}}%
\pgfpathlineto{\pgfqpoint{1.730181in}{1.386000in}}%
\pgfpathlineto{\pgfqpoint{1.730181in}{1.386000in}}%
\pgfpathclose%
\pgfusepath{fill}%
\end{pgfscope}%
\begin{pgfscope}%
\pgfpathrectangle{\pgfqpoint{0.595000in}{1.386000in}}{\pgfqpoint{2.835000in}{0.630000in}}%
\pgfusepath{clip}%
\pgfsetbuttcap%
\pgfsetmiterjoin%
\definecolor{currentfill}{rgb}{0.101961,0.101961,0.101961}%
\pgfsetfillcolor{currentfill}%
\pgfsetlinewidth{0.000000pt}%
\definecolor{currentstroke}{rgb}{0.000000,0.000000,0.000000}%
\pgfsetstrokecolor{currentstroke}%
\pgfsetstrokeopacity{0.000000}%
\pgfsetdash{}{0pt}%
\pgfpathmoveto{\pgfqpoint{1.753806in}{1.386000in}}%
\pgfpathlineto{\pgfqpoint{1.775069in}{1.386000in}}%
\pgfpathlineto{\pgfqpoint{1.775069in}{1.386000in}}%
\pgfpathlineto{\pgfqpoint{1.753806in}{1.386000in}}%
\pgfpathlineto{\pgfqpoint{1.753806in}{1.386000in}}%
\pgfpathclose%
\pgfusepath{fill}%
\end{pgfscope}%
\begin{pgfscope}%
\pgfpathrectangle{\pgfqpoint{0.595000in}{1.386000in}}{\pgfqpoint{2.835000in}{0.630000in}}%
\pgfusepath{clip}%
\pgfsetbuttcap%
\pgfsetmiterjoin%
\definecolor{currentfill}{rgb}{0.101961,0.101961,0.101961}%
\pgfsetfillcolor{currentfill}%
\pgfsetlinewidth{0.000000pt}%
\definecolor{currentstroke}{rgb}{0.000000,0.000000,0.000000}%
\pgfsetstrokecolor{currentstroke}%
\pgfsetstrokeopacity{0.000000}%
\pgfsetdash{}{0pt}%
\pgfpathmoveto{\pgfqpoint{1.777431in}{1.386000in}}%
\pgfpathlineto{\pgfqpoint{1.798694in}{1.386000in}}%
\pgfpathlineto{\pgfqpoint{1.798694in}{1.386000in}}%
\pgfpathlineto{\pgfqpoint{1.777431in}{1.386000in}}%
\pgfpathlineto{\pgfqpoint{1.777431in}{1.386000in}}%
\pgfpathclose%
\pgfusepath{fill}%
\end{pgfscope}%
\begin{pgfscope}%
\pgfpathrectangle{\pgfqpoint{0.595000in}{1.386000in}}{\pgfqpoint{2.835000in}{0.630000in}}%
\pgfusepath{clip}%
\pgfsetbuttcap%
\pgfsetmiterjoin%
\definecolor{currentfill}{rgb}{0.101961,0.101961,0.101961}%
\pgfsetfillcolor{currentfill}%
\pgfsetlinewidth{0.000000pt}%
\definecolor{currentstroke}{rgb}{0.000000,0.000000,0.000000}%
\pgfsetstrokecolor{currentstroke}%
\pgfsetstrokeopacity{0.000000}%
\pgfsetdash{}{0pt}%
\pgfpathmoveto{\pgfqpoint{1.801056in}{1.386000in}}%
\pgfpathlineto{\pgfqpoint{1.822319in}{1.386000in}}%
\pgfpathlineto{\pgfqpoint{1.822319in}{1.386000in}}%
\pgfpathlineto{\pgfqpoint{1.801056in}{1.386000in}}%
\pgfpathlineto{\pgfqpoint{1.801056in}{1.386000in}}%
\pgfpathclose%
\pgfusepath{fill}%
\end{pgfscope}%
\begin{pgfscope}%
\pgfpathrectangle{\pgfqpoint{0.595000in}{1.386000in}}{\pgfqpoint{2.835000in}{0.630000in}}%
\pgfusepath{clip}%
\pgfsetbuttcap%
\pgfsetmiterjoin%
\definecolor{currentfill}{rgb}{0.101961,0.101961,0.101961}%
\pgfsetfillcolor{currentfill}%
\pgfsetlinewidth{0.000000pt}%
\definecolor{currentstroke}{rgb}{0.000000,0.000000,0.000000}%
\pgfsetstrokecolor{currentstroke}%
\pgfsetstrokeopacity{0.000000}%
\pgfsetdash{}{0pt}%
\pgfpathmoveto{\pgfqpoint{1.824681in}{1.386000in}}%
\pgfpathlineto{\pgfqpoint{1.845944in}{1.386000in}}%
\pgfpathlineto{\pgfqpoint{1.845944in}{1.386000in}}%
\pgfpathlineto{\pgfqpoint{1.824681in}{1.386000in}}%
\pgfpathlineto{\pgfqpoint{1.824681in}{1.386000in}}%
\pgfpathclose%
\pgfusepath{fill}%
\end{pgfscope}%
\begin{pgfscope}%
\pgfpathrectangle{\pgfqpoint{0.595000in}{1.386000in}}{\pgfqpoint{2.835000in}{0.630000in}}%
\pgfusepath{clip}%
\pgfsetbuttcap%
\pgfsetmiterjoin%
\definecolor{currentfill}{rgb}{0.101961,0.101961,0.101961}%
\pgfsetfillcolor{currentfill}%
\pgfsetlinewidth{0.000000pt}%
\definecolor{currentstroke}{rgb}{0.000000,0.000000,0.000000}%
\pgfsetstrokecolor{currentstroke}%
\pgfsetstrokeopacity{0.000000}%
\pgfsetdash{}{0pt}%
\pgfpathmoveto{\pgfqpoint{1.848306in}{1.386000in}}%
\pgfpathlineto{\pgfqpoint{1.869569in}{1.386000in}}%
\pgfpathlineto{\pgfqpoint{1.869569in}{1.386000in}}%
\pgfpathlineto{\pgfqpoint{1.848306in}{1.386000in}}%
\pgfpathlineto{\pgfqpoint{1.848306in}{1.386000in}}%
\pgfpathclose%
\pgfusepath{fill}%
\end{pgfscope}%
\begin{pgfscope}%
\pgfpathrectangle{\pgfqpoint{0.595000in}{1.386000in}}{\pgfqpoint{2.835000in}{0.630000in}}%
\pgfusepath{clip}%
\pgfsetbuttcap%
\pgfsetmiterjoin%
\definecolor{currentfill}{rgb}{0.101961,0.101961,0.101961}%
\pgfsetfillcolor{currentfill}%
\pgfsetlinewidth{0.000000pt}%
\definecolor{currentstroke}{rgb}{0.000000,0.000000,0.000000}%
\pgfsetstrokecolor{currentstroke}%
\pgfsetstrokeopacity{0.000000}%
\pgfsetdash{}{0pt}%
\pgfpathmoveto{\pgfqpoint{1.871931in}{1.386000in}}%
\pgfpathlineto{\pgfqpoint{1.893194in}{1.386000in}}%
\pgfpathlineto{\pgfqpoint{1.893194in}{1.386000in}}%
\pgfpathlineto{\pgfqpoint{1.871931in}{1.386000in}}%
\pgfpathlineto{\pgfqpoint{1.871931in}{1.386000in}}%
\pgfpathclose%
\pgfusepath{fill}%
\end{pgfscope}%
\begin{pgfscope}%
\pgfpathrectangle{\pgfqpoint{0.595000in}{1.386000in}}{\pgfqpoint{2.835000in}{0.630000in}}%
\pgfusepath{clip}%
\pgfsetbuttcap%
\pgfsetmiterjoin%
\definecolor{currentfill}{rgb}{0.101961,0.101961,0.101961}%
\pgfsetfillcolor{currentfill}%
\pgfsetlinewidth{0.000000pt}%
\definecolor{currentstroke}{rgb}{0.000000,0.000000,0.000000}%
\pgfsetstrokecolor{currentstroke}%
\pgfsetstrokeopacity{0.000000}%
\pgfsetdash{}{0pt}%
\pgfpathmoveto{\pgfqpoint{1.895556in}{1.386000in}}%
\pgfpathlineto{\pgfqpoint{1.916819in}{1.386000in}}%
\pgfpathlineto{\pgfqpoint{1.916819in}{1.386000in}}%
\pgfpathlineto{\pgfqpoint{1.895556in}{1.386000in}}%
\pgfpathlineto{\pgfqpoint{1.895556in}{1.386000in}}%
\pgfpathclose%
\pgfusepath{fill}%
\end{pgfscope}%
\begin{pgfscope}%
\pgfpathrectangle{\pgfqpoint{0.595000in}{1.386000in}}{\pgfqpoint{2.835000in}{0.630000in}}%
\pgfusepath{clip}%
\pgfsetbuttcap%
\pgfsetmiterjoin%
\definecolor{currentfill}{rgb}{0.101961,0.101961,0.101961}%
\pgfsetfillcolor{currentfill}%
\pgfsetlinewidth{0.000000pt}%
\definecolor{currentstroke}{rgb}{0.000000,0.000000,0.000000}%
\pgfsetstrokecolor{currentstroke}%
\pgfsetstrokeopacity{0.000000}%
\pgfsetdash{}{0pt}%
\pgfpathmoveto{\pgfqpoint{1.919181in}{1.386000in}}%
\pgfpathlineto{\pgfqpoint{1.940444in}{1.386000in}}%
\pgfpathlineto{\pgfqpoint{1.940444in}{1.386000in}}%
\pgfpathlineto{\pgfqpoint{1.919181in}{1.386000in}}%
\pgfpathlineto{\pgfqpoint{1.919181in}{1.386000in}}%
\pgfpathclose%
\pgfusepath{fill}%
\end{pgfscope}%
\begin{pgfscope}%
\pgfpathrectangle{\pgfqpoint{0.595000in}{1.386000in}}{\pgfqpoint{2.835000in}{0.630000in}}%
\pgfusepath{clip}%
\pgfsetbuttcap%
\pgfsetmiterjoin%
\definecolor{currentfill}{rgb}{0.101961,0.101961,0.101961}%
\pgfsetfillcolor{currentfill}%
\pgfsetlinewidth{0.000000pt}%
\definecolor{currentstroke}{rgb}{0.000000,0.000000,0.000000}%
\pgfsetstrokecolor{currentstroke}%
\pgfsetstrokeopacity{0.000000}%
\pgfsetdash{}{0pt}%
\pgfpathmoveto{\pgfqpoint{1.942806in}{1.386000in}}%
\pgfpathlineto{\pgfqpoint{1.964069in}{1.386000in}}%
\pgfpathlineto{\pgfqpoint{1.964069in}{1.386000in}}%
\pgfpathlineto{\pgfqpoint{1.942806in}{1.386000in}}%
\pgfpathlineto{\pgfqpoint{1.942806in}{1.386000in}}%
\pgfpathclose%
\pgfusepath{fill}%
\end{pgfscope}%
\begin{pgfscope}%
\pgfpathrectangle{\pgfqpoint{0.595000in}{1.386000in}}{\pgfqpoint{2.835000in}{0.630000in}}%
\pgfusepath{clip}%
\pgfsetbuttcap%
\pgfsetmiterjoin%
\definecolor{currentfill}{rgb}{0.101961,0.101961,0.101961}%
\pgfsetfillcolor{currentfill}%
\pgfsetlinewidth{0.000000pt}%
\definecolor{currentstroke}{rgb}{0.000000,0.000000,0.000000}%
\pgfsetstrokecolor{currentstroke}%
\pgfsetstrokeopacity{0.000000}%
\pgfsetdash{}{0pt}%
\pgfpathmoveto{\pgfqpoint{1.966431in}{1.386000in}}%
\pgfpathlineto{\pgfqpoint{1.987694in}{1.386000in}}%
\pgfpathlineto{\pgfqpoint{1.987694in}{1.386000in}}%
\pgfpathlineto{\pgfqpoint{1.966431in}{1.386000in}}%
\pgfpathlineto{\pgfqpoint{1.966431in}{1.386000in}}%
\pgfpathclose%
\pgfusepath{fill}%
\end{pgfscope}%
\begin{pgfscope}%
\pgfpathrectangle{\pgfqpoint{0.595000in}{1.386000in}}{\pgfqpoint{2.835000in}{0.630000in}}%
\pgfusepath{clip}%
\pgfsetbuttcap%
\pgfsetmiterjoin%
\definecolor{currentfill}{rgb}{0.101961,0.101961,0.101961}%
\pgfsetfillcolor{currentfill}%
\pgfsetlinewidth{0.000000pt}%
\definecolor{currentstroke}{rgb}{0.000000,0.000000,0.000000}%
\pgfsetstrokecolor{currentstroke}%
\pgfsetstrokeopacity{0.000000}%
\pgfsetdash{}{0pt}%
\pgfpathmoveto{\pgfqpoint{1.990056in}{1.386000in}}%
\pgfpathlineto{\pgfqpoint{2.011319in}{1.386000in}}%
\pgfpathlineto{\pgfqpoint{2.011319in}{1.386000in}}%
\pgfpathlineto{\pgfqpoint{1.990056in}{1.386000in}}%
\pgfpathlineto{\pgfqpoint{1.990056in}{1.386000in}}%
\pgfpathclose%
\pgfusepath{fill}%
\end{pgfscope}%
\begin{pgfscope}%
\pgfpathrectangle{\pgfqpoint{0.595000in}{1.386000in}}{\pgfqpoint{2.835000in}{0.630000in}}%
\pgfusepath{clip}%
\pgfsetbuttcap%
\pgfsetmiterjoin%
\definecolor{currentfill}{rgb}{0.101961,0.101961,0.101961}%
\pgfsetfillcolor{currentfill}%
\pgfsetlinewidth{0.000000pt}%
\definecolor{currentstroke}{rgb}{0.000000,0.000000,0.000000}%
\pgfsetstrokecolor{currentstroke}%
\pgfsetstrokeopacity{0.000000}%
\pgfsetdash{}{0pt}%
\pgfpathmoveto{\pgfqpoint{2.013681in}{1.386000in}}%
\pgfpathlineto{\pgfqpoint{2.034944in}{1.386000in}}%
\pgfpathlineto{\pgfqpoint{2.034944in}{1.386000in}}%
\pgfpathlineto{\pgfqpoint{2.013681in}{1.386000in}}%
\pgfpathlineto{\pgfqpoint{2.013681in}{1.386000in}}%
\pgfpathclose%
\pgfusepath{fill}%
\end{pgfscope}%
\begin{pgfscope}%
\pgfpathrectangle{\pgfqpoint{0.595000in}{1.386000in}}{\pgfqpoint{2.835000in}{0.630000in}}%
\pgfusepath{clip}%
\pgfsetbuttcap%
\pgfsetmiterjoin%
\definecolor{currentfill}{rgb}{0.101961,0.101961,0.101961}%
\pgfsetfillcolor{currentfill}%
\pgfsetlinewidth{0.000000pt}%
\definecolor{currentstroke}{rgb}{0.000000,0.000000,0.000000}%
\pgfsetstrokecolor{currentstroke}%
\pgfsetstrokeopacity{0.000000}%
\pgfsetdash{}{0pt}%
\pgfpathmoveto{\pgfqpoint{2.037306in}{1.386000in}}%
\pgfpathlineto{\pgfqpoint{2.058569in}{1.386000in}}%
\pgfpathlineto{\pgfqpoint{2.058569in}{1.386000in}}%
\pgfpathlineto{\pgfqpoint{2.037306in}{1.386000in}}%
\pgfpathlineto{\pgfqpoint{2.037306in}{1.386000in}}%
\pgfpathclose%
\pgfusepath{fill}%
\end{pgfscope}%
\begin{pgfscope}%
\pgfpathrectangle{\pgfqpoint{0.595000in}{1.386000in}}{\pgfqpoint{2.835000in}{0.630000in}}%
\pgfusepath{clip}%
\pgfsetbuttcap%
\pgfsetmiterjoin%
\definecolor{currentfill}{rgb}{0.101961,0.101961,0.101961}%
\pgfsetfillcolor{currentfill}%
\pgfsetlinewidth{0.000000pt}%
\definecolor{currentstroke}{rgb}{0.000000,0.000000,0.000000}%
\pgfsetstrokecolor{currentstroke}%
\pgfsetstrokeopacity{0.000000}%
\pgfsetdash{}{0pt}%
\pgfpathmoveto{\pgfqpoint{2.060931in}{1.386000in}}%
\pgfpathlineto{\pgfqpoint{2.082194in}{1.386000in}}%
\pgfpathlineto{\pgfqpoint{2.082194in}{1.386000in}}%
\pgfpathlineto{\pgfqpoint{2.060931in}{1.386000in}}%
\pgfpathlineto{\pgfqpoint{2.060931in}{1.386000in}}%
\pgfpathclose%
\pgfusepath{fill}%
\end{pgfscope}%
\begin{pgfscope}%
\pgfpathrectangle{\pgfqpoint{0.595000in}{1.386000in}}{\pgfqpoint{2.835000in}{0.630000in}}%
\pgfusepath{clip}%
\pgfsetbuttcap%
\pgfsetmiterjoin%
\definecolor{currentfill}{rgb}{0.101961,0.101961,0.101961}%
\pgfsetfillcolor{currentfill}%
\pgfsetlinewidth{0.000000pt}%
\definecolor{currentstroke}{rgb}{0.000000,0.000000,0.000000}%
\pgfsetstrokecolor{currentstroke}%
\pgfsetstrokeopacity{0.000000}%
\pgfsetdash{}{0pt}%
\pgfpathmoveto{\pgfqpoint{2.084556in}{1.386000in}}%
\pgfpathlineto{\pgfqpoint{2.105819in}{1.386000in}}%
\pgfpathlineto{\pgfqpoint{2.105819in}{1.386000in}}%
\pgfpathlineto{\pgfqpoint{2.084556in}{1.386000in}}%
\pgfpathlineto{\pgfqpoint{2.084556in}{1.386000in}}%
\pgfpathclose%
\pgfusepath{fill}%
\end{pgfscope}%
\begin{pgfscope}%
\pgfpathrectangle{\pgfqpoint{0.595000in}{1.386000in}}{\pgfqpoint{2.835000in}{0.630000in}}%
\pgfusepath{clip}%
\pgfsetbuttcap%
\pgfsetmiterjoin%
\definecolor{currentfill}{rgb}{0.101961,0.101961,0.101961}%
\pgfsetfillcolor{currentfill}%
\pgfsetlinewidth{0.000000pt}%
\definecolor{currentstroke}{rgb}{0.000000,0.000000,0.000000}%
\pgfsetstrokecolor{currentstroke}%
\pgfsetstrokeopacity{0.000000}%
\pgfsetdash{}{0pt}%
\pgfpathmoveto{\pgfqpoint{2.108181in}{1.386000in}}%
\pgfpathlineto{\pgfqpoint{2.129444in}{1.386000in}}%
\pgfpathlineto{\pgfqpoint{2.129444in}{1.386000in}}%
\pgfpathlineto{\pgfqpoint{2.108181in}{1.386000in}}%
\pgfpathlineto{\pgfqpoint{2.108181in}{1.386000in}}%
\pgfpathclose%
\pgfusepath{fill}%
\end{pgfscope}%
\begin{pgfscope}%
\pgfpathrectangle{\pgfqpoint{0.595000in}{1.386000in}}{\pgfqpoint{2.835000in}{0.630000in}}%
\pgfusepath{clip}%
\pgfsetbuttcap%
\pgfsetmiterjoin%
\definecolor{currentfill}{rgb}{0.101961,0.101961,0.101961}%
\pgfsetfillcolor{currentfill}%
\pgfsetlinewidth{0.000000pt}%
\definecolor{currentstroke}{rgb}{0.000000,0.000000,0.000000}%
\pgfsetstrokecolor{currentstroke}%
\pgfsetstrokeopacity{0.000000}%
\pgfsetdash{}{0pt}%
\pgfpathmoveto{\pgfqpoint{2.131806in}{1.386000in}}%
\pgfpathlineto{\pgfqpoint{2.153069in}{1.386000in}}%
\pgfpathlineto{\pgfqpoint{2.153069in}{1.386000in}}%
\pgfpathlineto{\pgfqpoint{2.131806in}{1.386000in}}%
\pgfpathlineto{\pgfqpoint{2.131806in}{1.386000in}}%
\pgfpathclose%
\pgfusepath{fill}%
\end{pgfscope}%
\begin{pgfscope}%
\pgfpathrectangle{\pgfqpoint{0.595000in}{1.386000in}}{\pgfqpoint{2.835000in}{0.630000in}}%
\pgfusepath{clip}%
\pgfsetbuttcap%
\pgfsetmiterjoin%
\definecolor{currentfill}{rgb}{0.101961,0.101961,0.101961}%
\pgfsetfillcolor{currentfill}%
\pgfsetlinewidth{0.000000pt}%
\definecolor{currentstroke}{rgb}{0.000000,0.000000,0.000000}%
\pgfsetstrokecolor{currentstroke}%
\pgfsetstrokeopacity{0.000000}%
\pgfsetdash{}{0pt}%
\pgfpathmoveto{\pgfqpoint{2.155431in}{1.386000in}}%
\pgfpathlineto{\pgfqpoint{2.176694in}{1.386000in}}%
\pgfpathlineto{\pgfqpoint{2.176694in}{1.386000in}}%
\pgfpathlineto{\pgfqpoint{2.155431in}{1.386000in}}%
\pgfpathlineto{\pgfqpoint{2.155431in}{1.386000in}}%
\pgfpathclose%
\pgfusepath{fill}%
\end{pgfscope}%
\begin{pgfscope}%
\pgfpathrectangle{\pgfqpoint{0.595000in}{1.386000in}}{\pgfqpoint{2.835000in}{0.630000in}}%
\pgfusepath{clip}%
\pgfsetbuttcap%
\pgfsetmiterjoin%
\definecolor{currentfill}{rgb}{0.101961,0.101961,0.101961}%
\pgfsetfillcolor{currentfill}%
\pgfsetlinewidth{0.000000pt}%
\definecolor{currentstroke}{rgb}{0.000000,0.000000,0.000000}%
\pgfsetstrokecolor{currentstroke}%
\pgfsetstrokeopacity{0.000000}%
\pgfsetdash{}{0pt}%
\pgfpathmoveto{\pgfqpoint{2.179056in}{1.386000in}}%
\pgfpathlineto{\pgfqpoint{2.200319in}{1.386000in}}%
\pgfpathlineto{\pgfqpoint{2.200319in}{1.386000in}}%
\pgfpathlineto{\pgfqpoint{2.179056in}{1.386000in}}%
\pgfpathlineto{\pgfqpoint{2.179056in}{1.386000in}}%
\pgfpathclose%
\pgfusepath{fill}%
\end{pgfscope}%
\begin{pgfscope}%
\pgfpathrectangle{\pgfqpoint{0.595000in}{1.386000in}}{\pgfqpoint{2.835000in}{0.630000in}}%
\pgfusepath{clip}%
\pgfsetbuttcap%
\pgfsetmiterjoin%
\definecolor{currentfill}{rgb}{0.101961,0.101961,0.101961}%
\pgfsetfillcolor{currentfill}%
\pgfsetlinewidth{0.000000pt}%
\definecolor{currentstroke}{rgb}{0.000000,0.000000,0.000000}%
\pgfsetstrokecolor{currentstroke}%
\pgfsetstrokeopacity{0.000000}%
\pgfsetdash{}{0pt}%
\pgfpathmoveto{\pgfqpoint{2.202681in}{1.386000in}}%
\pgfpathlineto{\pgfqpoint{2.223944in}{1.386000in}}%
\pgfpathlineto{\pgfqpoint{2.223944in}{1.386000in}}%
\pgfpathlineto{\pgfqpoint{2.202681in}{1.386000in}}%
\pgfpathlineto{\pgfqpoint{2.202681in}{1.386000in}}%
\pgfpathclose%
\pgfusepath{fill}%
\end{pgfscope}%
\begin{pgfscope}%
\pgfpathrectangle{\pgfqpoint{0.595000in}{1.386000in}}{\pgfqpoint{2.835000in}{0.630000in}}%
\pgfusepath{clip}%
\pgfsetbuttcap%
\pgfsetmiterjoin%
\definecolor{currentfill}{rgb}{0.101961,0.101961,0.101961}%
\pgfsetfillcolor{currentfill}%
\pgfsetlinewidth{0.000000pt}%
\definecolor{currentstroke}{rgb}{0.000000,0.000000,0.000000}%
\pgfsetstrokecolor{currentstroke}%
\pgfsetstrokeopacity{0.000000}%
\pgfsetdash{}{0pt}%
\pgfpathmoveto{\pgfqpoint{2.226306in}{1.386000in}}%
\pgfpathlineto{\pgfqpoint{2.247569in}{1.386000in}}%
\pgfpathlineto{\pgfqpoint{2.247569in}{1.386000in}}%
\pgfpathlineto{\pgfqpoint{2.226306in}{1.386000in}}%
\pgfpathlineto{\pgfqpoint{2.226306in}{1.386000in}}%
\pgfpathclose%
\pgfusepath{fill}%
\end{pgfscope}%
\begin{pgfscope}%
\pgfpathrectangle{\pgfqpoint{0.595000in}{1.386000in}}{\pgfqpoint{2.835000in}{0.630000in}}%
\pgfusepath{clip}%
\pgfsetbuttcap%
\pgfsetmiterjoin%
\definecolor{currentfill}{rgb}{0.101961,0.101961,0.101961}%
\pgfsetfillcolor{currentfill}%
\pgfsetlinewidth{0.000000pt}%
\definecolor{currentstroke}{rgb}{0.000000,0.000000,0.000000}%
\pgfsetstrokecolor{currentstroke}%
\pgfsetstrokeopacity{0.000000}%
\pgfsetdash{}{0pt}%
\pgfpathmoveto{\pgfqpoint{2.249931in}{1.386000in}}%
\pgfpathlineto{\pgfqpoint{2.271194in}{1.386000in}}%
\pgfpathlineto{\pgfqpoint{2.271194in}{1.386000in}}%
\pgfpathlineto{\pgfqpoint{2.249931in}{1.386000in}}%
\pgfpathlineto{\pgfqpoint{2.249931in}{1.386000in}}%
\pgfpathclose%
\pgfusepath{fill}%
\end{pgfscope}%
\begin{pgfscope}%
\pgfpathrectangle{\pgfqpoint{0.595000in}{1.386000in}}{\pgfqpoint{2.835000in}{0.630000in}}%
\pgfusepath{clip}%
\pgfsetbuttcap%
\pgfsetmiterjoin%
\definecolor{currentfill}{rgb}{0.101961,0.101961,0.101961}%
\pgfsetfillcolor{currentfill}%
\pgfsetlinewidth{0.000000pt}%
\definecolor{currentstroke}{rgb}{0.000000,0.000000,0.000000}%
\pgfsetstrokecolor{currentstroke}%
\pgfsetstrokeopacity{0.000000}%
\pgfsetdash{}{0pt}%
\pgfpathmoveto{\pgfqpoint{2.273556in}{1.386000in}}%
\pgfpathlineto{\pgfqpoint{2.294819in}{1.386000in}}%
\pgfpathlineto{\pgfqpoint{2.294819in}{1.386000in}}%
\pgfpathlineto{\pgfqpoint{2.273556in}{1.386000in}}%
\pgfpathlineto{\pgfqpoint{2.273556in}{1.386000in}}%
\pgfpathclose%
\pgfusepath{fill}%
\end{pgfscope}%
\begin{pgfscope}%
\pgfpathrectangle{\pgfqpoint{0.595000in}{1.386000in}}{\pgfqpoint{2.835000in}{0.630000in}}%
\pgfusepath{clip}%
\pgfsetbuttcap%
\pgfsetmiterjoin%
\definecolor{currentfill}{rgb}{0.101961,0.101961,0.101961}%
\pgfsetfillcolor{currentfill}%
\pgfsetlinewidth{0.000000pt}%
\definecolor{currentstroke}{rgb}{0.000000,0.000000,0.000000}%
\pgfsetstrokecolor{currentstroke}%
\pgfsetstrokeopacity{0.000000}%
\pgfsetdash{}{0pt}%
\pgfpathmoveto{\pgfqpoint{2.297181in}{1.386000in}}%
\pgfpathlineto{\pgfqpoint{2.318444in}{1.386000in}}%
\pgfpathlineto{\pgfqpoint{2.318444in}{1.386000in}}%
\pgfpathlineto{\pgfqpoint{2.297181in}{1.386000in}}%
\pgfpathlineto{\pgfqpoint{2.297181in}{1.386000in}}%
\pgfpathclose%
\pgfusepath{fill}%
\end{pgfscope}%
\begin{pgfscope}%
\pgfpathrectangle{\pgfqpoint{0.595000in}{1.386000in}}{\pgfqpoint{2.835000in}{0.630000in}}%
\pgfusepath{clip}%
\pgfsetbuttcap%
\pgfsetmiterjoin%
\definecolor{currentfill}{rgb}{0.101961,0.101961,0.101961}%
\pgfsetfillcolor{currentfill}%
\pgfsetlinewidth{0.000000pt}%
\definecolor{currentstroke}{rgb}{0.000000,0.000000,0.000000}%
\pgfsetstrokecolor{currentstroke}%
\pgfsetstrokeopacity{0.000000}%
\pgfsetdash{}{0pt}%
\pgfpathmoveto{\pgfqpoint{2.320806in}{1.386000in}}%
\pgfpathlineto{\pgfqpoint{2.342069in}{1.386000in}}%
\pgfpathlineto{\pgfqpoint{2.342069in}{1.386000in}}%
\pgfpathlineto{\pgfqpoint{2.320806in}{1.386000in}}%
\pgfpathlineto{\pgfqpoint{2.320806in}{1.386000in}}%
\pgfpathclose%
\pgfusepath{fill}%
\end{pgfscope}%
\begin{pgfscope}%
\pgfpathrectangle{\pgfqpoint{0.595000in}{1.386000in}}{\pgfqpoint{2.835000in}{0.630000in}}%
\pgfusepath{clip}%
\pgfsetbuttcap%
\pgfsetmiterjoin%
\definecolor{currentfill}{rgb}{0.101961,0.101961,0.101961}%
\pgfsetfillcolor{currentfill}%
\pgfsetlinewidth{0.000000pt}%
\definecolor{currentstroke}{rgb}{0.000000,0.000000,0.000000}%
\pgfsetstrokecolor{currentstroke}%
\pgfsetstrokeopacity{0.000000}%
\pgfsetdash{}{0pt}%
\pgfpathmoveto{\pgfqpoint{2.344431in}{1.386000in}}%
\pgfpathlineto{\pgfqpoint{2.365694in}{1.386000in}}%
\pgfpathlineto{\pgfqpoint{2.365694in}{1.386000in}}%
\pgfpathlineto{\pgfqpoint{2.344431in}{1.386000in}}%
\pgfpathlineto{\pgfqpoint{2.344431in}{1.386000in}}%
\pgfpathclose%
\pgfusepath{fill}%
\end{pgfscope}%
\begin{pgfscope}%
\pgfpathrectangle{\pgfqpoint{0.595000in}{1.386000in}}{\pgfqpoint{2.835000in}{0.630000in}}%
\pgfusepath{clip}%
\pgfsetbuttcap%
\pgfsetmiterjoin%
\definecolor{currentfill}{rgb}{0.101961,0.101961,0.101961}%
\pgfsetfillcolor{currentfill}%
\pgfsetlinewidth{0.000000pt}%
\definecolor{currentstroke}{rgb}{0.000000,0.000000,0.000000}%
\pgfsetstrokecolor{currentstroke}%
\pgfsetstrokeopacity{0.000000}%
\pgfsetdash{}{0pt}%
\pgfpathmoveto{\pgfqpoint{2.368056in}{1.386000in}}%
\pgfpathlineto{\pgfqpoint{2.389319in}{1.386000in}}%
\pgfpathlineto{\pgfqpoint{2.389319in}{1.701000in}}%
\pgfpathlineto{\pgfqpoint{2.368056in}{1.701000in}}%
\pgfpathlineto{\pgfqpoint{2.368056in}{1.386000in}}%
\pgfpathclose%
\pgfusepath{fill}%
\end{pgfscope}%
\begin{pgfscope}%
\pgfpathrectangle{\pgfqpoint{0.595000in}{1.386000in}}{\pgfqpoint{2.835000in}{0.630000in}}%
\pgfusepath{clip}%
\pgfsetbuttcap%
\pgfsetmiterjoin%
\definecolor{currentfill}{rgb}{0.101961,0.101961,0.101961}%
\pgfsetfillcolor{currentfill}%
\pgfsetlinewidth{0.000000pt}%
\definecolor{currentstroke}{rgb}{0.000000,0.000000,0.000000}%
\pgfsetstrokecolor{currentstroke}%
\pgfsetstrokeopacity{0.000000}%
\pgfsetdash{}{0pt}%
\pgfpathmoveto{\pgfqpoint{2.391681in}{1.386000in}}%
\pgfpathlineto{\pgfqpoint{2.412944in}{1.386000in}}%
\pgfpathlineto{\pgfqpoint{2.412944in}{1.386000in}}%
\pgfpathlineto{\pgfqpoint{2.391681in}{1.386000in}}%
\pgfpathlineto{\pgfqpoint{2.391681in}{1.386000in}}%
\pgfpathclose%
\pgfusepath{fill}%
\end{pgfscope}%
\begin{pgfscope}%
\pgfpathrectangle{\pgfqpoint{0.595000in}{1.386000in}}{\pgfqpoint{2.835000in}{0.630000in}}%
\pgfusepath{clip}%
\pgfsetbuttcap%
\pgfsetmiterjoin%
\definecolor{currentfill}{rgb}{0.101961,0.101961,0.101961}%
\pgfsetfillcolor{currentfill}%
\pgfsetlinewidth{0.000000pt}%
\definecolor{currentstroke}{rgb}{0.000000,0.000000,0.000000}%
\pgfsetstrokecolor{currentstroke}%
\pgfsetstrokeopacity{0.000000}%
\pgfsetdash{}{0pt}%
\pgfpathmoveto{\pgfqpoint{2.415306in}{1.386000in}}%
\pgfpathlineto{\pgfqpoint{2.436569in}{1.386000in}}%
\pgfpathlineto{\pgfqpoint{2.436569in}{1.701000in}}%
\pgfpathlineto{\pgfqpoint{2.415306in}{1.701000in}}%
\pgfpathlineto{\pgfqpoint{2.415306in}{1.386000in}}%
\pgfpathclose%
\pgfusepath{fill}%
\end{pgfscope}%
\begin{pgfscope}%
\pgfpathrectangle{\pgfqpoint{0.595000in}{1.386000in}}{\pgfqpoint{2.835000in}{0.630000in}}%
\pgfusepath{clip}%
\pgfsetbuttcap%
\pgfsetmiterjoin%
\definecolor{currentfill}{rgb}{0.101961,0.101961,0.101961}%
\pgfsetfillcolor{currentfill}%
\pgfsetlinewidth{0.000000pt}%
\definecolor{currentstroke}{rgb}{0.000000,0.000000,0.000000}%
\pgfsetstrokecolor{currentstroke}%
\pgfsetstrokeopacity{0.000000}%
\pgfsetdash{}{0pt}%
\pgfpathmoveto{\pgfqpoint{2.438931in}{1.386000in}}%
\pgfpathlineto{\pgfqpoint{2.460194in}{1.386000in}}%
\pgfpathlineto{\pgfqpoint{2.460194in}{1.386000in}}%
\pgfpathlineto{\pgfqpoint{2.438931in}{1.386000in}}%
\pgfpathlineto{\pgfqpoint{2.438931in}{1.386000in}}%
\pgfpathclose%
\pgfusepath{fill}%
\end{pgfscope}%
\begin{pgfscope}%
\pgfpathrectangle{\pgfqpoint{0.595000in}{1.386000in}}{\pgfqpoint{2.835000in}{0.630000in}}%
\pgfusepath{clip}%
\pgfsetbuttcap%
\pgfsetmiterjoin%
\definecolor{currentfill}{rgb}{0.101961,0.101961,0.101961}%
\pgfsetfillcolor{currentfill}%
\pgfsetlinewidth{0.000000pt}%
\definecolor{currentstroke}{rgb}{0.000000,0.000000,0.000000}%
\pgfsetstrokecolor{currentstroke}%
\pgfsetstrokeopacity{0.000000}%
\pgfsetdash{}{0pt}%
\pgfpathmoveto{\pgfqpoint{2.462556in}{1.386000in}}%
\pgfpathlineto{\pgfqpoint{2.483819in}{1.386000in}}%
\pgfpathlineto{\pgfqpoint{2.483819in}{1.386000in}}%
\pgfpathlineto{\pgfqpoint{2.462556in}{1.386000in}}%
\pgfpathlineto{\pgfqpoint{2.462556in}{1.386000in}}%
\pgfpathclose%
\pgfusepath{fill}%
\end{pgfscope}%
\begin{pgfscope}%
\pgfpathrectangle{\pgfqpoint{0.595000in}{1.386000in}}{\pgfqpoint{2.835000in}{0.630000in}}%
\pgfusepath{clip}%
\pgfsetbuttcap%
\pgfsetmiterjoin%
\definecolor{currentfill}{rgb}{0.101961,0.101961,0.101961}%
\pgfsetfillcolor{currentfill}%
\pgfsetlinewidth{0.000000pt}%
\definecolor{currentstroke}{rgb}{0.000000,0.000000,0.000000}%
\pgfsetstrokecolor{currentstroke}%
\pgfsetstrokeopacity{0.000000}%
\pgfsetdash{}{0pt}%
\pgfpathmoveto{\pgfqpoint{2.486181in}{1.386000in}}%
\pgfpathlineto{\pgfqpoint{2.507444in}{1.386000in}}%
\pgfpathlineto{\pgfqpoint{2.507444in}{1.701000in}}%
\pgfpathlineto{\pgfqpoint{2.486181in}{1.701000in}}%
\pgfpathlineto{\pgfqpoint{2.486181in}{1.386000in}}%
\pgfpathclose%
\pgfusepath{fill}%
\end{pgfscope}%
\begin{pgfscope}%
\pgfpathrectangle{\pgfqpoint{0.595000in}{1.386000in}}{\pgfqpoint{2.835000in}{0.630000in}}%
\pgfusepath{clip}%
\pgfsetbuttcap%
\pgfsetmiterjoin%
\definecolor{currentfill}{rgb}{0.101961,0.101961,0.101961}%
\pgfsetfillcolor{currentfill}%
\pgfsetlinewidth{0.000000pt}%
\definecolor{currentstroke}{rgb}{0.000000,0.000000,0.000000}%
\pgfsetstrokecolor{currentstroke}%
\pgfsetstrokeopacity{0.000000}%
\pgfsetdash{}{0pt}%
\pgfpathmoveto{\pgfqpoint{2.509806in}{1.386000in}}%
\pgfpathlineto{\pgfqpoint{2.531069in}{1.386000in}}%
\pgfpathlineto{\pgfqpoint{2.531069in}{1.386000in}}%
\pgfpathlineto{\pgfqpoint{2.509806in}{1.386000in}}%
\pgfpathlineto{\pgfqpoint{2.509806in}{1.386000in}}%
\pgfpathclose%
\pgfusepath{fill}%
\end{pgfscope}%
\begin{pgfscope}%
\pgfpathrectangle{\pgfqpoint{0.595000in}{1.386000in}}{\pgfqpoint{2.835000in}{0.630000in}}%
\pgfusepath{clip}%
\pgfsetbuttcap%
\pgfsetmiterjoin%
\definecolor{currentfill}{rgb}{0.101961,0.101961,0.101961}%
\pgfsetfillcolor{currentfill}%
\pgfsetlinewidth{0.000000pt}%
\definecolor{currentstroke}{rgb}{0.000000,0.000000,0.000000}%
\pgfsetstrokecolor{currentstroke}%
\pgfsetstrokeopacity{0.000000}%
\pgfsetdash{}{0pt}%
\pgfpathmoveto{\pgfqpoint{2.533431in}{1.386000in}}%
\pgfpathlineto{\pgfqpoint{2.554694in}{1.386000in}}%
\pgfpathlineto{\pgfqpoint{2.554694in}{1.386000in}}%
\pgfpathlineto{\pgfqpoint{2.533431in}{1.386000in}}%
\pgfpathlineto{\pgfqpoint{2.533431in}{1.386000in}}%
\pgfpathclose%
\pgfusepath{fill}%
\end{pgfscope}%
\begin{pgfscope}%
\pgfpathrectangle{\pgfqpoint{0.595000in}{1.386000in}}{\pgfqpoint{2.835000in}{0.630000in}}%
\pgfusepath{clip}%
\pgfsetbuttcap%
\pgfsetmiterjoin%
\definecolor{currentfill}{rgb}{0.101961,0.101961,0.101961}%
\pgfsetfillcolor{currentfill}%
\pgfsetlinewidth{0.000000pt}%
\definecolor{currentstroke}{rgb}{0.000000,0.000000,0.000000}%
\pgfsetstrokecolor{currentstroke}%
\pgfsetstrokeopacity{0.000000}%
\pgfsetdash{}{0pt}%
\pgfpathmoveto{\pgfqpoint{2.557056in}{1.386000in}}%
\pgfpathlineto{\pgfqpoint{2.578319in}{1.386000in}}%
\pgfpathlineto{\pgfqpoint{2.578319in}{1.543500in}}%
\pgfpathlineto{\pgfqpoint{2.557056in}{1.543500in}}%
\pgfpathlineto{\pgfqpoint{2.557056in}{1.386000in}}%
\pgfpathclose%
\pgfusepath{fill}%
\end{pgfscope}%
\begin{pgfscope}%
\pgfpathrectangle{\pgfqpoint{0.595000in}{1.386000in}}{\pgfqpoint{2.835000in}{0.630000in}}%
\pgfusepath{clip}%
\pgfsetbuttcap%
\pgfsetmiterjoin%
\definecolor{currentfill}{rgb}{0.101961,0.101961,0.101961}%
\pgfsetfillcolor{currentfill}%
\pgfsetlinewidth{0.000000pt}%
\definecolor{currentstroke}{rgb}{0.000000,0.000000,0.000000}%
\pgfsetstrokecolor{currentstroke}%
\pgfsetstrokeopacity{0.000000}%
\pgfsetdash{}{0pt}%
\pgfpathmoveto{\pgfqpoint{2.580681in}{1.386000in}}%
\pgfpathlineto{\pgfqpoint{2.601944in}{1.386000in}}%
\pgfpathlineto{\pgfqpoint{2.601944in}{1.386000in}}%
\pgfpathlineto{\pgfqpoint{2.580681in}{1.386000in}}%
\pgfpathlineto{\pgfqpoint{2.580681in}{1.386000in}}%
\pgfpathclose%
\pgfusepath{fill}%
\end{pgfscope}%
\begin{pgfscope}%
\pgfpathrectangle{\pgfqpoint{0.595000in}{1.386000in}}{\pgfqpoint{2.835000in}{0.630000in}}%
\pgfusepath{clip}%
\pgfsetbuttcap%
\pgfsetmiterjoin%
\definecolor{currentfill}{rgb}{0.101961,0.101961,0.101961}%
\pgfsetfillcolor{currentfill}%
\pgfsetlinewidth{0.000000pt}%
\definecolor{currentstroke}{rgb}{0.000000,0.000000,0.000000}%
\pgfsetstrokecolor{currentstroke}%
\pgfsetstrokeopacity{0.000000}%
\pgfsetdash{}{0pt}%
\pgfpathmoveto{\pgfqpoint{2.604306in}{1.386000in}}%
\pgfpathlineto{\pgfqpoint{2.625569in}{1.386000in}}%
\pgfpathlineto{\pgfqpoint{2.625569in}{1.386000in}}%
\pgfpathlineto{\pgfqpoint{2.604306in}{1.386000in}}%
\pgfpathlineto{\pgfqpoint{2.604306in}{1.386000in}}%
\pgfpathclose%
\pgfusepath{fill}%
\end{pgfscope}%
\begin{pgfscope}%
\pgfpathrectangle{\pgfqpoint{0.595000in}{1.386000in}}{\pgfqpoint{2.835000in}{0.630000in}}%
\pgfusepath{clip}%
\pgfsetbuttcap%
\pgfsetmiterjoin%
\definecolor{currentfill}{rgb}{0.101961,0.101961,0.101961}%
\pgfsetfillcolor{currentfill}%
\pgfsetlinewidth{0.000000pt}%
\definecolor{currentstroke}{rgb}{0.000000,0.000000,0.000000}%
\pgfsetstrokecolor{currentstroke}%
\pgfsetstrokeopacity{0.000000}%
\pgfsetdash{}{0pt}%
\pgfpathmoveto{\pgfqpoint{2.627931in}{1.386000in}}%
\pgfpathlineto{\pgfqpoint{2.649194in}{1.386000in}}%
\pgfpathlineto{\pgfqpoint{2.649194in}{1.386000in}}%
\pgfpathlineto{\pgfqpoint{2.627931in}{1.386000in}}%
\pgfpathlineto{\pgfqpoint{2.627931in}{1.386000in}}%
\pgfpathclose%
\pgfusepath{fill}%
\end{pgfscope}%
\begin{pgfscope}%
\pgfpathrectangle{\pgfqpoint{0.595000in}{1.386000in}}{\pgfqpoint{2.835000in}{0.630000in}}%
\pgfusepath{clip}%
\pgfsetbuttcap%
\pgfsetmiterjoin%
\definecolor{currentfill}{rgb}{0.101961,0.101961,0.101961}%
\pgfsetfillcolor{currentfill}%
\pgfsetlinewidth{0.000000pt}%
\definecolor{currentstroke}{rgb}{0.000000,0.000000,0.000000}%
\pgfsetstrokecolor{currentstroke}%
\pgfsetstrokeopacity{0.000000}%
\pgfsetdash{}{0pt}%
\pgfpathmoveto{\pgfqpoint{2.651556in}{1.386000in}}%
\pgfpathlineto{\pgfqpoint{2.672819in}{1.386000in}}%
\pgfpathlineto{\pgfqpoint{2.672819in}{1.386000in}}%
\pgfpathlineto{\pgfqpoint{2.651556in}{1.386000in}}%
\pgfpathlineto{\pgfqpoint{2.651556in}{1.386000in}}%
\pgfpathclose%
\pgfusepath{fill}%
\end{pgfscope}%
\begin{pgfscope}%
\pgfpathrectangle{\pgfqpoint{0.595000in}{1.386000in}}{\pgfqpoint{2.835000in}{0.630000in}}%
\pgfusepath{clip}%
\pgfsetbuttcap%
\pgfsetmiterjoin%
\definecolor{currentfill}{rgb}{0.101961,0.101961,0.101961}%
\pgfsetfillcolor{currentfill}%
\pgfsetlinewidth{0.000000pt}%
\definecolor{currentstroke}{rgb}{0.000000,0.000000,0.000000}%
\pgfsetstrokecolor{currentstroke}%
\pgfsetstrokeopacity{0.000000}%
\pgfsetdash{}{0pt}%
\pgfpathmoveto{\pgfqpoint{2.675181in}{1.386000in}}%
\pgfpathlineto{\pgfqpoint{2.696444in}{1.386000in}}%
\pgfpathlineto{\pgfqpoint{2.696444in}{1.858500in}}%
\pgfpathlineto{\pgfqpoint{2.675181in}{1.858500in}}%
\pgfpathlineto{\pgfqpoint{2.675181in}{1.386000in}}%
\pgfpathclose%
\pgfusepath{fill}%
\end{pgfscope}%
\begin{pgfscope}%
\pgfpathrectangle{\pgfqpoint{0.595000in}{1.386000in}}{\pgfqpoint{2.835000in}{0.630000in}}%
\pgfusepath{clip}%
\pgfsetbuttcap%
\pgfsetmiterjoin%
\definecolor{currentfill}{rgb}{0.101961,0.101961,0.101961}%
\pgfsetfillcolor{currentfill}%
\pgfsetlinewidth{0.000000pt}%
\definecolor{currentstroke}{rgb}{0.000000,0.000000,0.000000}%
\pgfsetstrokecolor{currentstroke}%
\pgfsetstrokeopacity{0.000000}%
\pgfsetdash{}{0pt}%
\pgfpathmoveto{\pgfqpoint{2.698806in}{1.386000in}}%
\pgfpathlineto{\pgfqpoint{2.720069in}{1.386000in}}%
\pgfpathlineto{\pgfqpoint{2.720069in}{1.386000in}}%
\pgfpathlineto{\pgfqpoint{2.698806in}{1.386000in}}%
\pgfpathlineto{\pgfqpoint{2.698806in}{1.386000in}}%
\pgfpathclose%
\pgfusepath{fill}%
\end{pgfscope}%
\begin{pgfscope}%
\pgfpathrectangle{\pgfqpoint{0.595000in}{1.386000in}}{\pgfqpoint{2.835000in}{0.630000in}}%
\pgfusepath{clip}%
\pgfsetbuttcap%
\pgfsetmiterjoin%
\definecolor{currentfill}{rgb}{0.101961,0.101961,0.101961}%
\pgfsetfillcolor{currentfill}%
\pgfsetlinewidth{0.000000pt}%
\definecolor{currentstroke}{rgb}{0.000000,0.000000,0.000000}%
\pgfsetstrokecolor{currentstroke}%
\pgfsetstrokeopacity{0.000000}%
\pgfsetdash{}{0pt}%
\pgfpathmoveto{\pgfqpoint{2.722431in}{1.386000in}}%
\pgfpathlineto{\pgfqpoint{2.743694in}{1.386000in}}%
\pgfpathlineto{\pgfqpoint{2.743694in}{1.543500in}}%
\pgfpathlineto{\pgfqpoint{2.722431in}{1.543500in}}%
\pgfpathlineto{\pgfqpoint{2.722431in}{1.386000in}}%
\pgfpathclose%
\pgfusepath{fill}%
\end{pgfscope}%
\begin{pgfscope}%
\pgfpathrectangle{\pgfqpoint{0.595000in}{1.386000in}}{\pgfqpoint{2.835000in}{0.630000in}}%
\pgfusepath{clip}%
\pgfsetbuttcap%
\pgfsetmiterjoin%
\definecolor{currentfill}{rgb}{0.101961,0.101961,0.101961}%
\pgfsetfillcolor{currentfill}%
\pgfsetlinewidth{0.000000pt}%
\definecolor{currentstroke}{rgb}{0.000000,0.000000,0.000000}%
\pgfsetstrokecolor{currentstroke}%
\pgfsetstrokeopacity{0.000000}%
\pgfsetdash{}{0pt}%
\pgfpathmoveto{\pgfqpoint{2.746056in}{1.386000in}}%
\pgfpathlineto{\pgfqpoint{2.767319in}{1.386000in}}%
\pgfpathlineto{\pgfqpoint{2.767319in}{1.543500in}}%
\pgfpathlineto{\pgfqpoint{2.746056in}{1.543500in}}%
\pgfpathlineto{\pgfqpoint{2.746056in}{1.386000in}}%
\pgfpathclose%
\pgfusepath{fill}%
\end{pgfscope}%
\begin{pgfscope}%
\pgfpathrectangle{\pgfqpoint{0.595000in}{1.386000in}}{\pgfqpoint{2.835000in}{0.630000in}}%
\pgfusepath{clip}%
\pgfsetbuttcap%
\pgfsetmiterjoin%
\definecolor{currentfill}{rgb}{0.101961,0.101961,0.101961}%
\pgfsetfillcolor{currentfill}%
\pgfsetlinewidth{0.000000pt}%
\definecolor{currentstroke}{rgb}{0.000000,0.000000,0.000000}%
\pgfsetstrokecolor{currentstroke}%
\pgfsetstrokeopacity{0.000000}%
\pgfsetdash{}{0pt}%
\pgfpathmoveto{\pgfqpoint{2.769681in}{1.386000in}}%
\pgfpathlineto{\pgfqpoint{2.790944in}{1.386000in}}%
\pgfpathlineto{\pgfqpoint{2.790944in}{1.386000in}}%
\pgfpathlineto{\pgfqpoint{2.769681in}{1.386000in}}%
\pgfpathlineto{\pgfqpoint{2.769681in}{1.386000in}}%
\pgfpathclose%
\pgfusepath{fill}%
\end{pgfscope}%
\begin{pgfscope}%
\pgfpathrectangle{\pgfqpoint{0.595000in}{1.386000in}}{\pgfqpoint{2.835000in}{0.630000in}}%
\pgfusepath{clip}%
\pgfsetbuttcap%
\pgfsetmiterjoin%
\definecolor{currentfill}{rgb}{0.101961,0.101961,0.101961}%
\pgfsetfillcolor{currentfill}%
\pgfsetlinewidth{0.000000pt}%
\definecolor{currentstroke}{rgb}{0.000000,0.000000,0.000000}%
\pgfsetstrokecolor{currentstroke}%
\pgfsetstrokeopacity{0.000000}%
\pgfsetdash{}{0pt}%
\pgfpathmoveto{\pgfqpoint{2.793306in}{1.386000in}}%
\pgfpathlineto{\pgfqpoint{2.814569in}{1.386000in}}%
\pgfpathlineto{\pgfqpoint{2.814569in}{1.386000in}}%
\pgfpathlineto{\pgfqpoint{2.793306in}{1.386000in}}%
\pgfpathlineto{\pgfqpoint{2.793306in}{1.386000in}}%
\pgfpathclose%
\pgfusepath{fill}%
\end{pgfscope}%
\begin{pgfscope}%
\pgfpathrectangle{\pgfqpoint{0.595000in}{1.386000in}}{\pgfqpoint{2.835000in}{0.630000in}}%
\pgfusepath{clip}%
\pgfsetbuttcap%
\pgfsetmiterjoin%
\definecolor{currentfill}{rgb}{0.101961,0.101961,0.101961}%
\pgfsetfillcolor{currentfill}%
\pgfsetlinewidth{0.000000pt}%
\definecolor{currentstroke}{rgb}{0.000000,0.000000,0.000000}%
\pgfsetstrokecolor{currentstroke}%
\pgfsetstrokeopacity{0.000000}%
\pgfsetdash{}{0pt}%
\pgfpathmoveto{\pgfqpoint{2.816931in}{1.386000in}}%
\pgfpathlineto{\pgfqpoint{2.838194in}{1.386000in}}%
\pgfpathlineto{\pgfqpoint{2.838194in}{1.386000in}}%
\pgfpathlineto{\pgfqpoint{2.816931in}{1.386000in}}%
\pgfpathlineto{\pgfqpoint{2.816931in}{1.386000in}}%
\pgfpathclose%
\pgfusepath{fill}%
\end{pgfscope}%
\begin{pgfscope}%
\pgfpathrectangle{\pgfqpoint{0.595000in}{1.386000in}}{\pgfqpoint{2.835000in}{0.630000in}}%
\pgfusepath{clip}%
\pgfsetbuttcap%
\pgfsetmiterjoin%
\definecolor{currentfill}{rgb}{0.101961,0.101961,0.101961}%
\pgfsetfillcolor{currentfill}%
\pgfsetlinewidth{0.000000pt}%
\definecolor{currentstroke}{rgb}{0.000000,0.000000,0.000000}%
\pgfsetstrokecolor{currentstroke}%
\pgfsetstrokeopacity{0.000000}%
\pgfsetdash{}{0pt}%
\pgfpathmoveto{\pgfqpoint{2.840556in}{1.386000in}}%
\pgfpathlineto{\pgfqpoint{2.861819in}{1.386000in}}%
\pgfpathlineto{\pgfqpoint{2.861819in}{1.386000in}}%
\pgfpathlineto{\pgfqpoint{2.840556in}{1.386000in}}%
\pgfpathlineto{\pgfqpoint{2.840556in}{1.386000in}}%
\pgfpathclose%
\pgfusepath{fill}%
\end{pgfscope}%
\begin{pgfscope}%
\pgfpathrectangle{\pgfqpoint{0.595000in}{1.386000in}}{\pgfqpoint{2.835000in}{0.630000in}}%
\pgfusepath{clip}%
\pgfsetbuttcap%
\pgfsetmiterjoin%
\definecolor{currentfill}{rgb}{0.101961,0.101961,0.101961}%
\pgfsetfillcolor{currentfill}%
\pgfsetlinewidth{0.000000pt}%
\definecolor{currentstroke}{rgb}{0.000000,0.000000,0.000000}%
\pgfsetstrokecolor{currentstroke}%
\pgfsetstrokeopacity{0.000000}%
\pgfsetdash{}{0pt}%
\pgfpathmoveto{\pgfqpoint{2.864181in}{1.386000in}}%
\pgfpathlineto{\pgfqpoint{2.885444in}{1.386000in}}%
\pgfpathlineto{\pgfqpoint{2.885444in}{1.701000in}}%
\pgfpathlineto{\pgfqpoint{2.864181in}{1.701000in}}%
\pgfpathlineto{\pgfqpoint{2.864181in}{1.386000in}}%
\pgfpathclose%
\pgfusepath{fill}%
\end{pgfscope}%
\begin{pgfscope}%
\pgfpathrectangle{\pgfqpoint{0.595000in}{1.386000in}}{\pgfqpoint{2.835000in}{0.630000in}}%
\pgfusepath{clip}%
\pgfsetbuttcap%
\pgfsetmiterjoin%
\definecolor{currentfill}{rgb}{0.101961,0.101961,0.101961}%
\pgfsetfillcolor{currentfill}%
\pgfsetlinewidth{0.000000pt}%
\definecolor{currentstroke}{rgb}{0.000000,0.000000,0.000000}%
\pgfsetstrokecolor{currentstroke}%
\pgfsetstrokeopacity{0.000000}%
\pgfsetdash{}{0pt}%
\pgfpathmoveto{\pgfqpoint{2.887806in}{1.386000in}}%
\pgfpathlineto{\pgfqpoint{2.909069in}{1.386000in}}%
\pgfpathlineto{\pgfqpoint{2.909069in}{1.386000in}}%
\pgfpathlineto{\pgfqpoint{2.887806in}{1.386000in}}%
\pgfpathlineto{\pgfqpoint{2.887806in}{1.386000in}}%
\pgfpathclose%
\pgfusepath{fill}%
\end{pgfscope}%
\begin{pgfscope}%
\pgfpathrectangle{\pgfqpoint{0.595000in}{1.386000in}}{\pgfqpoint{2.835000in}{0.630000in}}%
\pgfusepath{clip}%
\pgfsetbuttcap%
\pgfsetmiterjoin%
\definecolor{currentfill}{rgb}{0.101961,0.101961,0.101961}%
\pgfsetfillcolor{currentfill}%
\pgfsetlinewidth{0.000000pt}%
\definecolor{currentstroke}{rgb}{0.000000,0.000000,0.000000}%
\pgfsetstrokecolor{currentstroke}%
\pgfsetstrokeopacity{0.000000}%
\pgfsetdash{}{0pt}%
\pgfpathmoveto{\pgfqpoint{2.911431in}{1.386000in}}%
\pgfpathlineto{\pgfqpoint{2.932694in}{1.386000in}}%
\pgfpathlineto{\pgfqpoint{2.932694in}{1.386000in}}%
\pgfpathlineto{\pgfqpoint{2.911431in}{1.386000in}}%
\pgfpathlineto{\pgfqpoint{2.911431in}{1.386000in}}%
\pgfpathclose%
\pgfusepath{fill}%
\end{pgfscope}%
\begin{pgfscope}%
\pgfpathrectangle{\pgfqpoint{0.595000in}{1.386000in}}{\pgfqpoint{2.835000in}{0.630000in}}%
\pgfusepath{clip}%
\pgfsetbuttcap%
\pgfsetmiterjoin%
\definecolor{currentfill}{rgb}{0.101961,0.101961,0.101961}%
\pgfsetfillcolor{currentfill}%
\pgfsetlinewidth{0.000000pt}%
\definecolor{currentstroke}{rgb}{0.000000,0.000000,0.000000}%
\pgfsetstrokecolor{currentstroke}%
\pgfsetstrokeopacity{0.000000}%
\pgfsetdash{}{0pt}%
\pgfpathmoveto{\pgfqpoint{2.935056in}{1.386000in}}%
\pgfpathlineto{\pgfqpoint{2.956319in}{1.386000in}}%
\pgfpathlineto{\pgfqpoint{2.956319in}{1.386000in}}%
\pgfpathlineto{\pgfqpoint{2.935056in}{1.386000in}}%
\pgfpathlineto{\pgfqpoint{2.935056in}{1.386000in}}%
\pgfpathclose%
\pgfusepath{fill}%
\end{pgfscope}%
\begin{pgfscope}%
\pgfpathrectangle{\pgfqpoint{0.595000in}{1.386000in}}{\pgfqpoint{2.835000in}{0.630000in}}%
\pgfusepath{clip}%
\pgfsetbuttcap%
\pgfsetmiterjoin%
\definecolor{currentfill}{rgb}{0.101961,0.101961,0.101961}%
\pgfsetfillcolor{currentfill}%
\pgfsetlinewidth{0.000000pt}%
\definecolor{currentstroke}{rgb}{0.000000,0.000000,0.000000}%
\pgfsetstrokecolor{currentstroke}%
\pgfsetstrokeopacity{0.000000}%
\pgfsetdash{}{0pt}%
\pgfpathmoveto{\pgfqpoint{2.958681in}{1.386000in}}%
\pgfpathlineto{\pgfqpoint{2.979944in}{1.386000in}}%
\pgfpathlineto{\pgfqpoint{2.979944in}{1.543500in}}%
\pgfpathlineto{\pgfqpoint{2.958681in}{1.543500in}}%
\pgfpathlineto{\pgfqpoint{2.958681in}{1.386000in}}%
\pgfpathclose%
\pgfusepath{fill}%
\end{pgfscope}%
\begin{pgfscope}%
\pgfpathrectangle{\pgfqpoint{0.595000in}{1.386000in}}{\pgfqpoint{2.835000in}{0.630000in}}%
\pgfusepath{clip}%
\pgfsetbuttcap%
\pgfsetmiterjoin%
\definecolor{currentfill}{rgb}{0.101961,0.101961,0.101961}%
\pgfsetfillcolor{currentfill}%
\pgfsetlinewidth{0.000000pt}%
\definecolor{currentstroke}{rgb}{0.000000,0.000000,0.000000}%
\pgfsetstrokecolor{currentstroke}%
\pgfsetstrokeopacity{0.000000}%
\pgfsetdash{}{0pt}%
\pgfpathmoveto{\pgfqpoint{2.982306in}{1.386000in}}%
\pgfpathlineto{\pgfqpoint{3.003569in}{1.386000in}}%
\pgfpathlineto{\pgfqpoint{3.003569in}{1.386000in}}%
\pgfpathlineto{\pgfqpoint{2.982306in}{1.386000in}}%
\pgfpathlineto{\pgfqpoint{2.982306in}{1.386000in}}%
\pgfpathclose%
\pgfusepath{fill}%
\end{pgfscope}%
\begin{pgfscope}%
\pgfpathrectangle{\pgfqpoint{0.595000in}{1.386000in}}{\pgfqpoint{2.835000in}{0.630000in}}%
\pgfusepath{clip}%
\pgfsetbuttcap%
\pgfsetmiterjoin%
\definecolor{currentfill}{rgb}{0.101961,0.101961,0.101961}%
\pgfsetfillcolor{currentfill}%
\pgfsetlinewidth{0.000000pt}%
\definecolor{currentstroke}{rgb}{0.000000,0.000000,0.000000}%
\pgfsetstrokecolor{currentstroke}%
\pgfsetstrokeopacity{0.000000}%
\pgfsetdash{}{0pt}%
\pgfpathmoveto{\pgfqpoint{3.005931in}{1.386000in}}%
\pgfpathlineto{\pgfqpoint{3.027194in}{1.386000in}}%
\pgfpathlineto{\pgfqpoint{3.027194in}{1.386000in}}%
\pgfpathlineto{\pgfqpoint{3.005931in}{1.386000in}}%
\pgfpathlineto{\pgfqpoint{3.005931in}{1.386000in}}%
\pgfpathclose%
\pgfusepath{fill}%
\end{pgfscope}%
\begin{pgfscope}%
\pgfpathrectangle{\pgfqpoint{0.595000in}{1.386000in}}{\pgfqpoint{2.835000in}{0.630000in}}%
\pgfusepath{clip}%
\pgfsetbuttcap%
\pgfsetmiterjoin%
\definecolor{currentfill}{rgb}{0.101961,0.101961,0.101961}%
\pgfsetfillcolor{currentfill}%
\pgfsetlinewidth{0.000000pt}%
\definecolor{currentstroke}{rgb}{0.000000,0.000000,0.000000}%
\pgfsetstrokecolor{currentstroke}%
\pgfsetstrokeopacity{0.000000}%
\pgfsetdash{}{0pt}%
\pgfpathmoveto{\pgfqpoint{3.029556in}{1.386000in}}%
\pgfpathlineto{\pgfqpoint{3.050819in}{1.386000in}}%
\pgfpathlineto{\pgfqpoint{3.050819in}{1.386000in}}%
\pgfpathlineto{\pgfqpoint{3.029556in}{1.386000in}}%
\pgfpathlineto{\pgfqpoint{3.029556in}{1.386000in}}%
\pgfpathclose%
\pgfusepath{fill}%
\end{pgfscope}%
\begin{pgfscope}%
\pgfpathrectangle{\pgfqpoint{0.595000in}{1.386000in}}{\pgfqpoint{2.835000in}{0.630000in}}%
\pgfusepath{clip}%
\pgfsetbuttcap%
\pgfsetmiterjoin%
\definecolor{currentfill}{rgb}{0.101961,0.101961,0.101961}%
\pgfsetfillcolor{currentfill}%
\pgfsetlinewidth{0.000000pt}%
\definecolor{currentstroke}{rgb}{0.000000,0.000000,0.000000}%
\pgfsetstrokecolor{currentstroke}%
\pgfsetstrokeopacity{0.000000}%
\pgfsetdash{}{0pt}%
\pgfpathmoveto{\pgfqpoint{3.053181in}{1.386000in}}%
\pgfpathlineto{\pgfqpoint{3.074444in}{1.386000in}}%
\pgfpathlineto{\pgfqpoint{3.074444in}{1.386000in}}%
\pgfpathlineto{\pgfqpoint{3.053181in}{1.386000in}}%
\pgfpathlineto{\pgfqpoint{3.053181in}{1.386000in}}%
\pgfpathclose%
\pgfusepath{fill}%
\end{pgfscope}%
\begin{pgfscope}%
\pgfpathrectangle{\pgfqpoint{0.595000in}{1.386000in}}{\pgfqpoint{2.835000in}{0.630000in}}%
\pgfusepath{clip}%
\pgfsetbuttcap%
\pgfsetmiterjoin%
\definecolor{currentfill}{rgb}{0.101961,0.101961,0.101961}%
\pgfsetfillcolor{currentfill}%
\pgfsetlinewidth{0.000000pt}%
\definecolor{currentstroke}{rgb}{0.000000,0.000000,0.000000}%
\pgfsetstrokecolor{currentstroke}%
\pgfsetstrokeopacity{0.000000}%
\pgfsetdash{}{0pt}%
\pgfpathmoveto{\pgfqpoint{3.076806in}{1.386000in}}%
\pgfpathlineto{\pgfqpoint{3.098069in}{1.386000in}}%
\pgfpathlineto{\pgfqpoint{3.098069in}{1.386000in}}%
\pgfpathlineto{\pgfqpoint{3.076806in}{1.386000in}}%
\pgfpathlineto{\pgfqpoint{3.076806in}{1.386000in}}%
\pgfpathclose%
\pgfusepath{fill}%
\end{pgfscope}%
\begin{pgfscope}%
\pgfpathrectangle{\pgfqpoint{0.595000in}{1.386000in}}{\pgfqpoint{2.835000in}{0.630000in}}%
\pgfusepath{clip}%
\pgfsetbuttcap%
\pgfsetmiterjoin%
\definecolor{currentfill}{rgb}{0.101961,0.101961,0.101961}%
\pgfsetfillcolor{currentfill}%
\pgfsetlinewidth{0.000000pt}%
\definecolor{currentstroke}{rgb}{0.000000,0.000000,0.000000}%
\pgfsetstrokecolor{currentstroke}%
\pgfsetstrokeopacity{0.000000}%
\pgfsetdash{}{0pt}%
\pgfpathmoveto{\pgfqpoint{3.100431in}{1.386000in}}%
\pgfpathlineto{\pgfqpoint{3.121694in}{1.386000in}}%
\pgfpathlineto{\pgfqpoint{3.121694in}{1.386000in}}%
\pgfpathlineto{\pgfqpoint{3.100431in}{1.386000in}}%
\pgfpathlineto{\pgfqpoint{3.100431in}{1.386000in}}%
\pgfpathclose%
\pgfusepath{fill}%
\end{pgfscope}%
\begin{pgfscope}%
\pgfpathrectangle{\pgfqpoint{0.595000in}{1.386000in}}{\pgfqpoint{2.835000in}{0.630000in}}%
\pgfusepath{clip}%
\pgfsetbuttcap%
\pgfsetmiterjoin%
\definecolor{currentfill}{rgb}{0.101961,0.101961,0.101961}%
\pgfsetfillcolor{currentfill}%
\pgfsetlinewidth{0.000000pt}%
\definecolor{currentstroke}{rgb}{0.000000,0.000000,0.000000}%
\pgfsetstrokecolor{currentstroke}%
\pgfsetstrokeopacity{0.000000}%
\pgfsetdash{}{0pt}%
\pgfpathmoveto{\pgfqpoint{3.124056in}{1.386000in}}%
\pgfpathlineto{\pgfqpoint{3.145319in}{1.386000in}}%
\pgfpathlineto{\pgfqpoint{3.145319in}{1.386000in}}%
\pgfpathlineto{\pgfqpoint{3.124056in}{1.386000in}}%
\pgfpathlineto{\pgfqpoint{3.124056in}{1.386000in}}%
\pgfpathclose%
\pgfusepath{fill}%
\end{pgfscope}%
\begin{pgfscope}%
\pgfpathrectangle{\pgfqpoint{0.595000in}{1.386000in}}{\pgfqpoint{2.835000in}{0.630000in}}%
\pgfusepath{clip}%
\pgfsetbuttcap%
\pgfsetmiterjoin%
\definecolor{currentfill}{rgb}{0.101961,0.101961,0.101961}%
\pgfsetfillcolor{currentfill}%
\pgfsetlinewidth{0.000000pt}%
\definecolor{currentstroke}{rgb}{0.000000,0.000000,0.000000}%
\pgfsetstrokecolor{currentstroke}%
\pgfsetstrokeopacity{0.000000}%
\pgfsetdash{}{0pt}%
\pgfpathmoveto{\pgfqpoint{3.147681in}{1.386000in}}%
\pgfpathlineto{\pgfqpoint{3.168944in}{1.386000in}}%
\pgfpathlineto{\pgfqpoint{3.168944in}{1.386000in}}%
\pgfpathlineto{\pgfqpoint{3.147681in}{1.386000in}}%
\pgfpathlineto{\pgfqpoint{3.147681in}{1.386000in}}%
\pgfpathclose%
\pgfusepath{fill}%
\end{pgfscope}%
\begin{pgfscope}%
\pgfpathrectangle{\pgfqpoint{0.595000in}{1.386000in}}{\pgfqpoint{2.835000in}{0.630000in}}%
\pgfusepath{clip}%
\pgfsetbuttcap%
\pgfsetmiterjoin%
\definecolor{currentfill}{rgb}{0.101961,0.101961,0.101961}%
\pgfsetfillcolor{currentfill}%
\pgfsetlinewidth{0.000000pt}%
\definecolor{currentstroke}{rgb}{0.000000,0.000000,0.000000}%
\pgfsetstrokecolor{currentstroke}%
\pgfsetstrokeopacity{0.000000}%
\pgfsetdash{}{0pt}%
\pgfpathmoveto{\pgfqpoint{3.171306in}{1.386000in}}%
\pgfpathlineto{\pgfqpoint{3.192569in}{1.386000in}}%
\pgfpathlineto{\pgfqpoint{3.192569in}{1.386000in}}%
\pgfpathlineto{\pgfqpoint{3.171306in}{1.386000in}}%
\pgfpathlineto{\pgfqpoint{3.171306in}{1.386000in}}%
\pgfpathclose%
\pgfusepath{fill}%
\end{pgfscope}%
\begin{pgfscope}%
\pgfpathrectangle{\pgfqpoint{0.595000in}{1.386000in}}{\pgfqpoint{2.835000in}{0.630000in}}%
\pgfusepath{clip}%
\pgfsetbuttcap%
\pgfsetmiterjoin%
\definecolor{currentfill}{rgb}{0.101961,0.101961,0.101961}%
\pgfsetfillcolor{currentfill}%
\pgfsetlinewidth{0.000000pt}%
\definecolor{currentstroke}{rgb}{0.000000,0.000000,0.000000}%
\pgfsetstrokecolor{currentstroke}%
\pgfsetstrokeopacity{0.000000}%
\pgfsetdash{}{0pt}%
\pgfpathmoveto{\pgfqpoint{3.194931in}{1.386000in}}%
\pgfpathlineto{\pgfqpoint{3.216194in}{1.386000in}}%
\pgfpathlineto{\pgfqpoint{3.216194in}{1.386000in}}%
\pgfpathlineto{\pgfqpoint{3.194931in}{1.386000in}}%
\pgfpathlineto{\pgfqpoint{3.194931in}{1.386000in}}%
\pgfpathclose%
\pgfusepath{fill}%
\end{pgfscope}%
\begin{pgfscope}%
\pgfpathrectangle{\pgfqpoint{0.595000in}{1.386000in}}{\pgfqpoint{2.835000in}{0.630000in}}%
\pgfusepath{clip}%
\pgfsetbuttcap%
\pgfsetmiterjoin%
\definecolor{currentfill}{rgb}{0.101961,0.101961,0.101961}%
\pgfsetfillcolor{currentfill}%
\pgfsetlinewidth{0.000000pt}%
\definecolor{currentstroke}{rgb}{0.000000,0.000000,0.000000}%
\pgfsetstrokecolor{currentstroke}%
\pgfsetstrokeopacity{0.000000}%
\pgfsetdash{}{0pt}%
\pgfpathmoveto{\pgfqpoint{3.218556in}{1.386000in}}%
\pgfpathlineto{\pgfqpoint{3.239819in}{1.386000in}}%
\pgfpathlineto{\pgfqpoint{3.239819in}{1.386000in}}%
\pgfpathlineto{\pgfqpoint{3.218556in}{1.386000in}}%
\pgfpathlineto{\pgfqpoint{3.218556in}{1.386000in}}%
\pgfpathclose%
\pgfusepath{fill}%
\end{pgfscope}%
\begin{pgfscope}%
\pgfpathrectangle{\pgfqpoint{0.595000in}{1.386000in}}{\pgfqpoint{2.835000in}{0.630000in}}%
\pgfusepath{clip}%
\pgfsetbuttcap%
\pgfsetmiterjoin%
\definecolor{currentfill}{rgb}{0.101961,0.101961,0.101961}%
\pgfsetfillcolor{currentfill}%
\pgfsetlinewidth{0.000000pt}%
\definecolor{currentstroke}{rgb}{0.000000,0.000000,0.000000}%
\pgfsetstrokecolor{currentstroke}%
\pgfsetstrokeopacity{0.000000}%
\pgfsetdash{}{0pt}%
\pgfpathmoveto{\pgfqpoint{3.242181in}{1.386000in}}%
\pgfpathlineto{\pgfqpoint{3.263444in}{1.386000in}}%
\pgfpathlineto{\pgfqpoint{3.263444in}{1.386000in}}%
\pgfpathlineto{\pgfqpoint{3.242181in}{1.386000in}}%
\pgfpathlineto{\pgfqpoint{3.242181in}{1.386000in}}%
\pgfpathclose%
\pgfusepath{fill}%
\end{pgfscope}%
\begin{pgfscope}%
\pgfpathrectangle{\pgfqpoint{0.595000in}{1.386000in}}{\pgfqpoint{2.835000in}{0.630000in}}%
\pgfusepath{clip}%
\pgfsetbuttcap%
\pgfsetmiterjoin%
\definecolor{currentfill}{rgb}{0.101961,0.101961,0.101961}%
\pgfsetfillcolor{currentfill}%
\pgfsetlinewidth{0.000000pt}%
\definecolor{currentstroke}{rgb}{0.000000,0.000000,0.000000}%
\pgfsetstrokecolor{currentstroke}%
\pgfsetstrokeopacity{0.000000}%
\pgfsetdash{}{0pt}%
\pgfpathmoveto{\pgfqpoint{3.265806in}{1.386000in}}%
\pgfpathlineto{\pgfqpoint{3.287069in}{1.386000in}}%
\pgfpathlineto{\pgfqpoint{3.287069in}{1.386000in}}%
\pgfpathlineto{\pgfqpoint{3.265806in}{1.386000in}}%
\pgfpathlineto{\pgfqpoint{3.265806in}{1.386000in}}%
\pgfpathclose%
\pgfusepath{fill}%
\end{pgfscope}%
\begin{pgfscope}%
\pgfpathrectangle{\pgfqpoint{0.595000in}{1.386000in}}{\pgfqpoint{2.835000in}{0.630000in}}%
\pgfusepath{clip}%
\pgfsetbuttcap%
\pgfsetmiterjoin%
\definecolor{currentfill}{rgb}{0.101961,0.101961,0.101961}%
\pgfsetfillcolor{currentfill}%
\pgfsetlinewidth{0.000000pt}%
\definecolor{currentstroke}{rgb}{0.000000,0.000000,0.000000}%
\pgfsetstrokecolor{currentstroke}%
\pgfsetstrokeopacity{0.000000}%
\pgfsetdash{}{0pt}%
\pgfpathmoveto{\pgfqpoint{3.289431in}{1.386000in}}%
\pgfpathlineto{\pgfqpoint{3.310694in}{1.386000in}}%
\pgfpathlineto{\pgfqpoint{3.310694in}{1.386000in}}%
\pgfpathlineto{\pgfqpoint{3.289431in}{1.386000in}}%
\pgfpathlineto{\pgfqpoint{3.289431in}{1.386000in}}%
\pgfpathclose%
\pgfusepath{fill}%
\end{pgfscope}%
\begin{pgfscope}%
\pgfpathrectangle{\pgfqpoint{0.595000in}{1.386000in}}{\pgfqpoint{2.835000in}{0.630000in}}%
\pgfusepath{clip}%
\pgfsetbuttcap%
\pgfsetmiterjoin%
\definecolor{currentfill}{rgb}{0.101961,0.101961,0.101961}%
\pgfsetfillcolor{currentfill}%
\pgfsetlinewidth{0.000000pt}%
\definecolor{currentstroke}{rgb}{0.000000,0.000000,0.000000}%
\pgfsetstrokecolor{currentstroke}%
\pgfsetstrokeopacity{0.000000}%
\pgfsetdash{}{0pt}%
\pgfpathmoveto{\pgfqpoint{3.313056in}{1.386000in}}%
\pgfpathlineto{\pgfqpoint{3.334319in}{1.386000in}}%
\pgfpathlineto{\pgfqpoint{3.334319in}{1.386000in}}%
\pgfpathlineto{\pgfqpoint{3.313056in}{1.386000in}}%
\pgfpathlineto{\pgfqpoint{3.313056in}{1.386000in}}%
\pgfpathclose%
\pgfusepath{fill}%
\end{pgfscope}%
\begin{pgfscope}%
\pgfpathrectangle{\pgfqpoint{0.595000in}{1.386000in}}{\pgfqpoint{2.835000in}{0.630000in}}%
\pgfusepath{clip}%
\pgfsetbuttcap%
\pgfsetmiterjoin%
\definecolor{currentfill}{rgb}{0.101961,0.101961,0.101961}%
\pgfsetfillcolor{currentfill}%
\pgfsetlinewidth{0.000000pt}%
\definecolor{currentstroke}{rgb}{0.000000,0.000000,0.000000}%
\pgfsetstrokecolor{currentstroke}%
\pgfsetstrokeopacity{0.000000}%
\pgfsetdash{}{0pt}%
\pgfpathmoveto{\pgfqpoint{3.336681in}{1.386000in}}%
\pgfpathlineto{\pgfqpoint{3.357944in}{1.386000in}}%
\pgfpathlineto{\pgfqpoint{3.357944in}{1.386000in}}%
\pgfpathlineto{\pgfqpoint{3.336681in}{1.386000in}}%
\pgfpathlineto{\pgfqpoint{3.336681in}{1.386000in}}%
\pgfpathclose%
\pgfusepath{fill}%
\end{pgfscope}%
\begin{pgfscope}%
\pgfpathrectangle{\pgfqpoint{0.595000in}{1.386000in}}{\pgfqpoint{2.835000in}{0.630000in}}%
\pgfusepath{clip}%
\pgfsetbuttcap%
\pgfsetmiterjoin%
\definecolor{currentfill}{rgb}{0.101961,0.101961,0.101961}%
\pgfsetfillcolor{currentfill}%
\pgfsetlinewidth{0.000000pt}%
\definecolor{currentstroke}{rgb}{0.000000,0.000000,0.000000}%
\pgfsetstrokecolor{currentstroke}%
\pgfsetstrokeopacity{0.000000}%
\pgfsetdash{}{0pt}%
\pgfpathmoveto{\pgfqpoint{3.360306in}{1.386000in}}%
\pgfpathlineto{\pgfqpoint{3.381569in}{1.386000in}}%
\pgfpathlineto{\pgfqpoint{3.381569in}{1.386000in}}%
\pgfpathlineto{\pgfqpoint{3.360306in}{1.386000in}}%
\pgfpathlineto{\pgfqpoint{3.360306in}{1.386000in}}%
\pgfpathclose%
\pgfusepath{fill}%
\end{pgfscope}%
\begin{pgfscope}%
\pgfpathrectangle{\pgfqpoint{0.595000in}{1.386000in}}{\pgfqpoint{2.835000in}{0.630000in}}%
\pgfusepath{clip}%
\pgfsetbuttcap%
\pgfsetmiterjoin%
\definecolor{currentfill}{rgb}{0.101961,0.101961,0.101961}%
\pgfsetfillcolor{currentfill}%
\pgfsetlinewidth{0.000000pt}%
\definecolor{currentstroke}{rgb}{0.000000,0.000000,0.000000}%
\pgfsetstrokecolor{currentstroke}%
\pgfsetstrokeopacity{0.000000}%
\pgfsetdash{}{0pt}%
\pgfpathmoveto{\pgfqpoint{3.383931in}{1.386000in}}%
\pgfpathlineto{\pgfqpoint{3.405194in}{1.386000in}}%
\pgfpathlineto{\pgfqpoint{3.405194in}{1.386000in}}%
\pgfpathlineto{\pgfqpoint{3.383931in}{1.386000in}}%
\pgfpathlineto{\pgfqpoint{3.383931in}{1.386000in}}%
\pgfpathclose%
\pgfusepath{fill}%
\end{pgfscope}%
\begin{pgfscope}%
\pgfpathrectangle{\pgfqpoint{0.595000in}{1.386000in}}{\pgfqpoint{2.835000in}{0.630000in}}%
\pgfusepath{clip}%
\pgfsetbuttcap%
\pgfsetmiterjoin%
\definecolor{currentfill}{rgb}{0.101961,0.101961,0.101961}%
\pgfsetfillcolor{currentfill}%
\pgfsetlinewidth{0.000000pt}%
\definecolor{currentstroke}{rgb}{0.000000,0.000000,0.000000}%
\pgfsetstrokecolor{currentstroke}%
\pgfsetstrokeopacity{0.000000}%
\pgfsetdash{}{0pt}%
\pgfpathmoveto{\pgfqpoint{3.407556in}{1.386000in}}%
\pgfpathlineto{\pgfqpoint{3.428819in}{1.386000in}}%
\pgfpathlineto{\pgfqpoint{3.428819in}{1.543500in}}%
\pgfpathlineto{\pgfqpoint{3.407556in}{1.543500in}}%
\pgfpathlineto{\pgfqpoint{3.407556in}{1.386000in}}%
\pgfpathclose%
\pgfusepath{fill}%
\end{pgfscope}%
\begin{pgfscope}%
\pgfsetrectcap%
\pgfsetmiterjoin%
\pgfsetlinewidth{0.803000pt}%
\definecolor{currentstroke}{rgb}{0.101961,0.101961,0.101961}%
\pgfsetstrokecolor{currentstroke}%
\pgfsetdash{}{0pt}%
\pgfpathmoveto{\pgfqpoint{0.595000in}{1.386000in}}%
\pgfpathlineto{\pgfqpoint{0.595000in}{2.016000in}}%
\pgfusepath{stroke}%
\end{pgfscope}%
\begin{pgfscope}%
\pgfsetrectcap%
\pgfsetmiterjoin%
\pgfsetlinewidth{0.803000pt}%
\definecolor{currentstroke}{rgb}{0.101961,0.101961,0.101961}%
\pgfsetstrokecolor{currentstroke}%
\pgfsetdash{}{0pt}%
\pgfpathmoveto{\pgfqpoint{0.595000in}{1.386000in}}%
\pgfpathlineto{\pgfqpoint{3.430000in}{1.386000in}}%
\pgfusepath{stroke}%
\end{pgfscope}%
\begin{pgfscope}%
\pgfsetbuttcap%
\pgfsetmiterjoin%
\definecolor{currentfill}{rgb}{1.000000,1.000000,1.000000}%
\pgfsetfillcolor{currentfill}%
\pgfsetlinewidth{0.000000pt}%
\definecolor{currentstroke}{rgb}{0.000000,0.000000,0.000000}%
\pgfsetstrokecolor{currentstroke}%
\pgfsetstrokeopacity{0.000000}%
\pgfsetdash{}{0pt}%
\pgfpathmoveto{\pgfqpoint{0.595000in}{0.483000in}}%
\pgfpathlineto{\pgfqpoint{3.430000in}{0.483000in}}%
\pgfpathlineto{\pgfqpoint{3.430000in}{1.113000in}}%
\pgfpathlineto{\pgfqpoint{0.595000in}{1.113000in}}%
\pgfpathlineto{\pgfqpoint{0.595000in}{0.483000in}}%
\pgfpathclose%
\pgfusepath{fill}%
\end{pgfscope}%
\begin{pgfscope}%
\pgfpathrectangle{\pgfqpoint{0.595000in}{0.483000in}}{\pgfqpoint{2.835000in}{0.630000in}}%
\pgfusepath{clip}%
\pgfsetbuttcap%
\pgfsetmiterjoin%
\pgfsetlinewidth{0.000000pt}%
\definecolor{currentstroke}{rgb}{0.000000,0.000000,0.000000}%
\pgfsetstrokecolor{currentstroke}%
\pgfsetstrokeopacity{0.000000}%
\pgfsetdash{}{0pt}%
\pgfpathmoveto{\pgfqpoint{1.557719in}{0.483000in}}%
\pgfpathlineto{\pgfqpoint{1.557719in}{1.113000in}}%
\pgfpathlineto{\pgfqpoint{2.001869in}{1.113000in}}%
\pgfpathlineto{\pgfqpoint{2.001869in}{0.483000in}}%
\pgfpathlineto{\pgfqpoint{1.557719in}{0.483000in}}%
\pgfpathclose%
\pgfusepath{}%
\end{pgfscope}%
\begin{pgfscope}%
\definecolor{textcolor}{rgb}{0.101961,0.101961,0.101961}%
\pgfsetstrokecolor{textcolor}%
\pgfsetfillcolor{textcolor}%
\pgftext[x=0.595000in,y=0.434389in,,top]{\color{textcolor}\sffamily\fontsize{7.000000}{8.400000}\selectfont 0}%
\end{pgfscope}%
\begin{pgfscope}%
\definecolor{textcolor}{rgb}{0.101961,0.101961,0.101961}%
\pgfsetstrokecolor{textcolor}%
\pgfsetfillcolor{textcolor}%
\pgftext[x=1.067500in,y=0.434389in,,top]{\color{textcolor}\sffamily\fontsize{7.000000}{8.400000}\selectfont 200}%
\end{pgfscope}%
\begin{pgfscope}%
\definecolor{textcolor}{rgb}{0.101961,0.101961,0.101961}%
\pgfsetstrokecolor{textcolor}%
\pgfsetfillcolor{textcolor}%
\pgftext[x=1.540000in,y=0.434389in,,top]{\color{textcolor}\sffamily\fontsize{7.000000}{8.400000}\selectfont 400}%
\end{pgfscope}%
\begin{pgfscope}%
\definecolor{textcolor}{rgb}{0.101961,0.101961,0.101961}%
\pgfsetstrokecolor{textcolor}%
\pgfsetfillcolor{textcolor}%
\pgftext[x=2.012500in,y=0.434389in,,top]{\color{textcolor}\sffamily\fontsize{7.000000}{8.400000}\selectfont 600}%
\end{pgfscope}%
\begin{pgfscope}%
\definecolor{textcolor}{rgb}{0.101961,0.101961,0.101961}%
\pgfsetstrokecolor{textcolor}%
\pgfsetfillcolor{textcolor}%
\pgftext[x=2.485000in,y=0.434389in,,top]{\color{textcolor}\sffamily\fontsize{7.000000}{8.400000}\selectfont 800}%
\end{pgfscope}%
\begin{pgfscope}%
\definecolor{textcolor}{rgb}{0.101961,0.101961,0.101961}%
\pgfsetstrokecolor{textcolor}%
\pgfsetfillcolor{textcolor}%
\pgftext[x=2.957500in,y=0.434389in,,top]{\color{textcolor}\sffamily\fontsize{7.000000}{8.400000}\selectfont 1000}%
\end{pgfscope}%
\begin{pgfscope}%
\definecolor{textcolor}{rgb}{0.101961,0.101961,0.101961}%
\pgfsetstrokecolor{textcolor}%
\pgfsetfillcolor{textcolor}%
\pgftext[x=3.430000in,y=0.434389in,,top]{\color{textcolor}\sffamily\fontsize{7.000000}{8.400000}\selectfont 1200}%
\end{pgfscope}%
\begin{pgfscope}%
\definecolor{textcolor}{rgb}{0.000000,0.000000,0.000000}%
\pgfsetstrokecolor{textcolor}%
\pgfsetfillcolor{textcolor}%
\pgftext[x=2.012500in,y=0.320192in,,top]{\color{textcolor}\sffamily\fontsize{7.500000}{9.000000}\selectfont Time (s)}%
\end{pgfscope}%
\begin{pgfscope}%
\pgfpathrectangle{\pgfqpoint{0.595000in}{0.483000in}}{\pgfqpoint{2.835000in}{0.630000in}}%
\pgfusepath{clip}%
\pgfsetrectcap%
\pgfsetroundjoin%
\pgfsetlinewidth{0.401500pt}%
\definecolor{currentstroke}{rgb}{0.901961,0.901961,0.901961}%
\pgfsetstrokecolor{currentstroke}%
\pgfsetdash{}{0pt}%
\pgfpathmoveto{\pgfqpoint{0.595000in}{0.483000in}}%
\pgfpathlineto{\pgfqpoint{3.430000in}{0.483000in}}%
\pgfusepath{stroke}%
\end{pgfscope}%
\begin{pgfscope}%
\definecolor{textcolor}{rgb}{0.101961,0.101961,0.101961}%
\pgfsetstrokecolor{textcolor}%
\pgfsetfillcolor{textcolor}%
\pgftext[x=0.494739in, y=0.449242in, left, base]{\color{textcolor}\sffamily\fontsize{7.000000}{8.400000}\selectfont 0}%
\end{pgfscope}%
\begin{pgfscope}%
\pgfpathrectangle{\pgfqpoint{0.595000in}{0.483000in}}{\pgfqpoint{2.835000in}{0.630000in}}%
\pgfusepath{clip}%
\pgfsetrectcap%
\pgfsetroundjoin%
\pgfsetlinewidth{0.401500pt}%
\definecolor{currentstroke}{rgb}{0.901961,0.901961,0.901961}%
\pgfsetstrokecolor{currentstroke}%
\pgfsetdash{}{0pt}%
\pgfpathmoveto{\pgfqpoint{0.595000in}{0.708000in}}%
\pgfpathlineto{\pgfqpoint{3.430000in}{0.708000in}}%
\pgfusepath{stroke}%
\end{pgfscope}%
\begin{pgfscope}%
\definecolor{textcolor}{rgb}{0.101961,0.101961,0.101961}%
\pgfsetstrokecolor{textcolor}%
\pgfsetfillcolor{textcolor}%
\pgftext[x=0.494739in, y=0.674242in, left, base]{\color{textcolor}\sffamily\fontsize{7.000000}{8.400000}\selectfont 5}%
\end{pgfscope}%
\begin{pgfscope}%
\pgfpathrectangle{\pgfqpoint{0.595000in}{0.483000in}}{\pgfqpoint{2.835000in}{0.630000in}}%
\pgfusepath{clip}%
\pgfsetrectcap%
\pgfsetroundjoin%
\pgfsetlinewidth{0.401500pt}%
\definecolor{currentstroke}{rgb}{0.901961,0.901961,0.901961}%
\pgfsetstrokecolor{currentstroke}%
\pgfsetdash{}{0pt}%
\pgfpathmoveto{\pgfqpoint{0.595000in}{0.933000in}}%
\pgfpathlineto{\pgfqpoint{3.430000in}{0.933000in}}%
\pgfusepath{stroke}%
\end{pgfscope}%
\begin{pgfscope}%
\definecolor{textcolor}{rgb}{0.101961,0.101961,0.101961}%
\pgfsetstrokecolor{textcolor}%
\pgfsetfillcolor{textcolor}%
\pgftext[x=0.443089in, y=0.899242in, left, base]{\color{textcolor}\sffamily\fontsize{7.000000}{8.400000}\selectfont 10}%
\end{pgfscope}%
\begin{pgfscope}%
\definecolor{textcolor}{rgb}{0.000000,0.000000,0.000000}%
\pgfsetstrokecolor{textcolor}%
\pgfsetfillcolor{textcolor}%
\pgftext[x=0.268707in, y=0.570082in, left, base,rotate=90.000000]{\color{textcolor}\sffamily\fontsize{7.500000}{9.000000}\selectfont Late alerts}%
\end{pgfscope}%
\begin{pgfscope}%
\definecolor{textcolor}{rgb}{0.000000,0.000000,0.000000}%
\pgfsetstrokecolor{textcolor}%
\pgfsetfillcolor{textcolor}%
\pgftext[x=0.368629in, y=0.623420in, left, base,rotate=90.000000]{\color{textcolor}\sffamily\fontsize{7.500000}{9.000000}\selectfont per 10 s}%
\end{pgfscope}%
\begin{pgfscope}%
\pgfpathrectangle{\pgfqpoint{0.595000in}{0.483000in}}{\pgfqpoint{2.835000in}{0.630000in}}%
\pgfusepath{clip}%
\pgfsetbuttcap%
\pgfsetmiterjoin%
\definecolor{currentfill}{rgb}{0.101961,0.101961,0.101961}%
\pgfsetfillcolor{currentfill}%
\pgfsetlinewidth{0.000000pt}%
\definecolor{currentstroke}{rgb}{0.000000,0.000000,0.000000}%
\pgfsetstrokecolor{currentstroke}%
\pgfsetstrokeopacity{0.000000}%
\pgfsetdash{}{0pt}%
\pgfpathmoveto{\pgfqpoint{0.596181in}{0.483000in}}%
\pgfpathlineto{\pgfqpoint{0.617444in}{0.483000in}}%
\pgfpathlineto{\pgfqpoint{0.617444in}{0.483000in}}%
\pgfpathlineto{\pgfqpoint{0.596181in}{0.483000in}}%
\pgfpathlineto{\pgfqpoint{0.596181in}{0.483000in}}%
\pgfpathclose%
\pgfusepath{fill}%
\end{pgfscope}%
\begin{pgfscope}%
\pgfpathrectangle{\pgfqpoint{0.595000in}{0.483000in}}{\pgfqpoint{2.835000in}{0.630000in}}%
\pgfusepath{clip}%
\pgfsetbuttcap%
\pgfsetmiterjoin%
\definecolor{currentfill}{rgb}{0.101961,0.101961,0.101961}%
\pgfsetfillcolor{currentfill}%
\pgfsetlinewidth{0.000000pt}%
\definecolor{currentstroke}{rgb}{0.000000,0.000000,0.000000}%
\pgfsetstrokecolor{currentstroke}%
\pgfsetstrokeopacity{0.000000}%
\pgfsetdash{}{0pt}%
\pgfpathmoveto{\pgfqpoint{0.619806in}{0.483000in}}%
\pgfpathlineto{\pgfqpoint{0.641069in}{0.483000in}}%
\pgfpathlineto{\pgfqpoint{0.641069in}{0.483000in}}%
\pgfpathlineto{\pgfqpoint{0.619806in}{0.483000in}}%
\pgfpathlineto{\pgfqpoint{0.619806in}{0.483000in}}%
\pgfpathclose%
\pgfusepath{fill}%
\end{pgfscope}%
\begin{pgfscope}%
\pgfpathrectangle{\pgfqpoint{0.595000in}{0.483000in}}{\pgfqpoint{2.835000in}{0.630000in}}%
\pgfusepath{clip}%
\pgfsetbuttcap%
\pgfsetmiterjoin%
\definecolor{currentfill}{rgb}{0.101961,0.101961,0.101961}%
\pgfsetfillcolor{currentfill}%
\pgfsetlinewidth{0.000000pt}%
\definecolor{currentstroke}{rgb}{0.000000,0.000000,0.000000}%
\pgfsetstrokecolor{currentstroke}%
\pgfsetstrokeopacity{0.000000}%
\pgfsetdash{}{0pt}%
\pgfpathmoveto{\pgfqpoint{0.643431in}{0.483000in}}%
\pgfpathlineto{\pgfqpoint{0.664694in}{0.483000in}}%
\pgfpathlineto{\pgfqpoint{0.664694in}{0.528000in}}%
\pgfpathlineto{\pgfqpoint{0.643431in}{0.528000in}}%
\pgfpathlineto{\pgfqpoint{0.643431in}{0.483000in}}%
\pgfpathclose%
\pgfusepath{fill}%
\end{pgfscope}%
\begin{pgfscope}%
\pgfpathrectangle{\pgfqpoint{0.595000in}{0.483000in}}{\pgfqpoint{2.835000in}{0.630000in}}%
\pgfusepath{clip}%
\pgfsetbuttcap%
\pgfsetmiterjoin%
\definecolor{currentfill}{rgb}{0.101961,0.101961,0.101961}%
\pgfsetfillcolor{currentfill}%
\pgfsetlinewidth{0.000000pt}%
\definecolor{currentstroke}{rgb}{0.000000,0.000000,0.000000}%
\pgfsetstrokecolor{currentstroke}%
\pgfsetstrokeopacity{0.000000}%
\pgfsetdash{}{0pt}%
\pgfpathmoveto{\pgfqpoint{0.667056in}{0.483000in}}%
\pgfpathlineto{\pgfqpoint{0.688319in}{0.483000in}}%
\pgfpathlineto{\pgfqpoint{0.688319in}{0.483000in}}%
\pgfpathlineto{\pgfqpoint{0.667056in}{0.483000in}}%
\pgfpathlineto{\pgfqpoint{0.667056in}{0.483000in}}%
\pgfpathclose%
\pgfusepath{fill}%
\end{pgfscope}%
\begin{pgfscope}%
\pgfpathrectangle{\pgfqpoint{0.595000in}{0.483000in}}{\pgfqpoint{2.835000in}{0.630000in}}%
\pgfusepath{clip}%
\pgfsetbuttcap%
\pgfsetmiterjoin%
\definecolor{currentfill}{rgb}{0.101961,0.101961,0.101961}%
\pgfsetfillcolor{currentfill}%
\pgfsetlinewidth{0.000000pt}%
\definecolor{currentstroke}{rgb}{0.000000,0.000000,0.000000}%
\pgfsetstrokecolor{currentstroke}%
\pgfsetstrokeopacity{0.000000}%
\pgfsetdash{}{0pt}%
\pgfpathmoveto{\pgfqpoint{0.690681in}{0.483000in}}%
\pgfpathlineto{\pgfqpoint{0.711944in}{0.483000in}}%
\pgfpathlineto{\pgfqpoint{0.711944in}{0.573000in}}%
\pgfpathlineto{\pgfqpoint{0.690681in}{0.573000in}}%
\pgfpathlineto{\pgfqpoint{0.690681in}{0.483000in}}%
\pgfpathclose%
\pgfusepath{fill}%
\end{pgfscope}%
\begin{pgfscope}%
\pgfpathrectangle{\pgfqpoint{0.595000in}{0.483000in}}{\pgfqpoint{2.835000in}{0.630000in}}%
\pgfusepath{clip}%
\pgfsetbuttcap%
\pgfsetmiterjoin%
\definecolor{currentfill}{rgb}{0.101961,0.101961,0.101961}%
\pgfsetfillcolor{currentfill}%
\pgfsetlinewidth{0.000000pt}%
\definecolor{currentstroke}{rgb}{0.000000,0.000000,0.000000}%
\pgfsetstrokecolor{currentstroke}%
\pgfsetstrokeopacity{0.000000}%
\pgfsetdash{}{0pt}%
\pgfpathmoveto{\pgfqpoint{0.714306in}{0.483000in}}%
\pgfpathlineto{\pgfqpoint{0.735569in}{0.483000in}}%
\pgfpathlineto{\pgfqpoint{0.735569in}{0.483000in}}%
\pgfpathlineto{\pgfqpoint{0.714306in}{0.483000in}}%
\pgfpathlineto{\pgfqpoint{0.714306in}{0.483000in}}%
\pgfpathclose%
\pgfusepath{fill}%
\end{pgfscope}%
\begin{pgfscope}%
\pgfpathrectangle{\pgfqpoint{0.595000in}{0.483000in}}{\pgfqpoint{2.835000in}{0.630000in}}%
\pgfusepath{clip}%
\pgfsetbuttcap%
\pgfsetmiterjoin%
\definecolor{currentfill}{rgb}{0.101961,0.101961,0.101961}%
\pgfsetfillcolor{currentfill}%
\pgfsetlinewidth{0.000000pt}%
\definecolor{currentstroke}{rgb}{0.000000,0.000000,0.000000}%
\pgfsetstrokecolor{currentstroke}%
\pgfsetstrokeopacity{0.000000}%
\pgfsetdash{}{0pt}%
\pgfpathmoveto{\pgfqpoint{0.737931in}{0.483000in}}%
\pgfpathlineto{\pgfqpoint{0.759194in}{0.483000in}}%
\pgfpathlineto{\pgfqpoint{0.759194in}{0.483000in}}%
\pgfpathlineto{\pgfqpoint{0.737931in}{0.483000in}}%
\pgfpathlineto{\pgfqpoint{0.737931in}{0.483000in}}%
\pgfpathclose%
\pgfusepath{fill}%
\end{pgfscope}%
\begin{pgfscope}%
\pgfpathrectangle{\pgfqpoint{0.595000in}{0.483000in}}{\pgfqpoint{2.835000in}{0.630000in}}%
\pgfusepath{clip}%
\pgfsetbuttcap%
\pgfsetmiterjoin%
\definecolor{currentfill}{rgb}{0.101961,0.101961,0.101961}%
\pgfsetfillcolor{currentfill}%
\pgfsetlinewidth{0.000000pt}%
\definecolor{currentstroke}{rgb}{0.000000,0.000000,0.000000}%
\pgfsetstrokecolor{currentstroke}%
\pgfsetstrokeopacity{0.000000}%
\pgfsetdash{}{0pt}%
\pgfpathmoveto{\pgfqpoint{0.761556in}{0.483000in}}%
\pgfpathlineto{\pgfqpoint{0.782819in}{0.483000in}}%
\pgfpathlineto{\pgfqpoint{0.782819in}{0.483000in}}%
\pgfpathlineto{\pgfqpoint{0.761556in}{0.483000in}}%
\pgfpathlineto{\pgfqpoint{0.761556in}{0.483000in}}%
\pgfpathclose%
\pgfusepath{fill}%
\end{pgfscope}%
\begin{pgfscope}%
\pgfpathrectangle{\pgfqpoint{0.595000in}{0.483000in}}{\pgfqpoint{2.835000in}{0.630000in}}%
\pgfusepath{clip}%
\pgfsetbuttcap%
\pgfsetmiterjoin%
\definecolor{currentfill}{rgb}{0.101961,0.101961,0.101961}%
\pgfsetfillcolor{currentfill}%
\pgfsetlinewidth{0.000000pt}%
\definecolor{currentstroke}{rgb}{0.000000,0.000000,0.000000}%
\pgfsetstrokecolor{currentstroke}%
\pgfsetstrokeopacity{0.000000}%
\pgfsetdash{}{0pt}%
\pgfpathmoveto{\pgfqpoint{0.785181in}{0.483000in}}%
\pgfpathlineto{\pgfqpoint{0.806444in}{0.483000in}}%
\pgfpathlineto{\pgfqpoint{0.806444in}{0.483000in}}%
\pgfpathlineto{\pgfqpoint{0.785181in}{0.483000in}}%
\pgfpathlineto{\pgfqpoint{0.785181in}{0.483000in}}%
\pgfpathclose%
\pgfusepath{fill}%
\end{pgfscope}%
\begin{pgfscope}%
\pgfpathrectangle{\pgfqpoint{0.595000in}{0.483000in}}{\pgfqpoint{2.835000in}{0.630000in}}%
\pgfusepath{clip}%
\pgfsetbuttcap%
\pgfsetmiterjoin%
\definecolor{currentfill}{rgb}{0.101961,0.101961,0.101961}%
\pgfsetfillcolor{currentfill}%
\pgfsetlinewidth{0.000000pt}%
\definecolor{currentstroke}{rgb}{0.000000,0.000000,0.000000}%
\pgfsetstrokecolor{currentstroke}%
\pgfsetstrokeopacity{0.000000}%
\pgfsetdash{}{0pt}%
\pgfpathmoveto{\pgfqpoint{0.808806in}{0.483000in}}%
\pgfpathlineto{\pgfqpoint{0.830069in}{0.483000in}}%
\pgfpathlineto{\pgfqpoint{0.830069in}{0.483000in}}%
\pgfpathlineto{\pgfqpoint{0.808806in}{0.483000in}}%
\pgfpathlineto{\pgfqpoint{0.808806in}{0.483000in}}%
\pgfpathclose%
\pgfusepath{fill}%
\end{pgfscope}%
\begin{pgfscope}%
\pgfpathrectangle{\pgfqpoint{0.595000in}{0.483000in}}{\pgfqpoint{2.835000in}{0.630000in}}%
\pgfusepath{clip}%
\pgfsetbuttcap%
\pgfsetmiterjoin%
\definecolor{currentfill}{rgb}{0.101961,0.101961,0.101961}%
\pgfsetfillcolor{currentfill}%
\pgfsetlinewidth{0.000000pt}%
\definecolor{currentstroke}{rgb}{0.000000,0.000000,0.000000}%
\pgfsetstrokecolor{currentstroke}%
\pgfsetstrokeopacity{0.000000}%
\pgfsetdash{}{0pt}%
\pgfpathmoveto{\pgfqpoint{0.832431in}{0.483000in}}%
\pgfpathlineto{\pgfqpoint{0.853694in}{0.483000in}}%
\pgfpathlineto{\pgfqpoint{0.853694in}{0.483000in}}%
\pgfpathlineto{\pgfqpoint{0.832431in}{0.483000in}}%
\pgfpathlineto{\pgfqpoint{0.832431in}{0.483000in}}%
\pgfpathclose%
\pgfusepath{fill}%
\end{pgfscope}%
\begin{pgfscope}%
\pgfpathrectangle{\pgfqpoint{0.595000in}{0.483000in}}{\pgfqpoint{2.835000in}{0.630000in}}%
\pgfusepath{clip}%
\pgfsetbuttcap%
\pgfsetmiterjoin%
\definecolor{currentfill}{rgb}{0.101961,0.101961,0.101961}%
\pgfsetfillcolor{currentfill}%
\pgfsetlinewidth{0.000000pt}%
\definecolor{currentstroke}{rgb}{0.000000,0.000000,0.000000}%
\pgfsetstrokecolor{currentstroke}%
\pgfsetstrokeopacity{0.000000}%
\pgfsetdash{}{0pt}%
\pgfpathmoveto{\pgfqpoint{0.856056in}{0.483000in}}%
\pgfpathlineto{\pgfqpoint{0.877319in}{0.483000in}}%
\pgfpathlineto{\pgfqpoint{0.877319in}{0.483000in}}%
\pgfpathlineto{\pgfqpoint{0.856056in}{0.483000in}}%
\pgfpathlineto{\pgfqpoint{0.856056in}{0.483000in}}%
\pgfpathclose%
\pgfusepath{fill}%
\end{pgfscope}%
\begin{pgfscope}%
\pgfpathrectangle{\pgfqpoint{0.595000in}{0.483000in}}{\pgfqpoint{2.835000in}{0.630000in}}%
\pgfusepath{clip}%
\pgfsetbuttcap%
\pgfsetmiterjoin%
\definecolor{currentfill}{rgb}{0.101961,0.101961,0.101961}%
\pgfsetfillcolor{currentfill}%
\pgfsetlinewidth{0.000000pt}%
\definecolor{currentstroke}{rgb}{0.000000,0.000000,0.000000}%
\pgfsetstrokecolor{currentstroke}%
\pgfsetstrokeopacity{0.000000}%
\pgfsetdash{}{0pt}%
\pgfpathmoveto{\pgfqpoint{0.879681in}{0.483000in}}%
\pgfpathlineto{\pgfqpoint{0.900944in}{0.483000in}}%
\pgfpathlineto{\pgfqpoint{0.900944in}{0.618000in}}%
\pgfpathlineto{\pgfqpoint{0.879681in}{0.618000in}}%
\pgfpathlineto{\pgfqpoint{0.879681in}{0.483000in}}%
\pgfpathclose%
\pgfusepath{fill}%
\end{pgfscope}%
\begin{pgfscope}%
\pgfpathrectangle{\pgfqpoint{0.595000in}{0.483000in}}{\pgfqpoint{2.835000in}{0.630000in}}%
\pgfusepath{clip}%
\pgfsetbuttcap%
\pgfsetmiterjoin%
\definecolor{currentfill}{rgb}{0.101961,0.101961,0.101961}%
\pgfsetfillcolor{currentfill}%
\pgfsetlinewidth{0.000000pt}%
\definecolor{currentstroke}{rgb}{0.000000,0.000000,0.000000}%
\pgfsetstrokecolor{currentstroke}%
\pgfsetstrokeopacity{0.000000}%
\pgfsetdash{}{0pt}%
\pgfpathmoveto{\pgfqpoint{0.903306in}{0.483000in}}%
\pgfpathlineto{\pgfqpoint{0.924569in}{0.483000in}}%
\pgfpathlineto{\pgfqpoint{0.924569in}{0.573000in}}%
\pgfpathlineto{\pgfqpoint{0.903306in}{0.573000in}}%
\pgfpathlineto{\pgfqpoint{0.903306in}{0.483000in}}%
\pgfpathclose%
\pgfusepath{fill}%
\end{pgfscope}%
\begin{pgfscope}%
\pgfpathrectangle{\pgfqpoint{0.595000in}{0.483000in}}{\pgfqpoint{2.835000in}{0.630000in}}%
\pgfusepath{clip}%
\pgfsetbuttcap%
\pgfsetmiterjoin%
\definecolor{currentfill}{rgb}{0.101961,0.101961,0.101961}%
\pgfsetfillcolor{currentfill}%
\pgfsetlinewidth{0.000000pt}%
\definecolor{currentstroke}{rgb}{0.000000,0.000000,0.000000}%
\pgfsetstrokecolor{currentstroke}%
\pgfsetstrokeopacity{0.000000}%
\pgfsetdash{}{0pt}%
\pgfpathmoveto{\pgfqpoint{0.926931in}{0.483000in}}%
\pgfpathlineto{\pgfqpoint{0.948194in}{0.483000in}}%
\pgfpathlineto{\pgfqpoint{0.948194in}{0.483000in}}%
\pgfpathlineto{\pgfqpoint{0.926931in}{0.483000in}}%
\pgfpathlineto{\pgfqpoint{0.926931in}{0.483000in}}%
\pgfpathclose%
\pgfusepath{fill}%
\end{pgfscope}%
\begin{pgfscope}%
\pgfpathrectangle{\pgfqpoint{0.595000in}{0.483000in}}{\pgfqpoint{2.835000in}{0.630000in}}%
\pgfusepath{clip}%
\pgfsetbuttcap%
\pgfsetmiterjoin%
\definecolor{currentfill}{rgb}{0.101961,0.101961,0.101961}%
\pgfsetfillcolor{currentfill}%
\pgfsetlinewidth{0.000000pt}%
\definecolor{currentstroke}{rgb}{0.000000,0.000000,0.000000}%
\pgfsetstrokecolor{currentstroke}%
\pgfsetstrokeopacity{0.000000}%
\pgfsetdash{}{0pt}%
\pgfpathmoveto{\pgfqpoint{0.950556in}{0.483000in}}%
\pgfpathlineto{\pgfqpoint{0.971819in}{0.483000in}}%
\pgfpathlineto{\pgfqpoint{0.971819in}{0.573000in}}%
\pgfpathlineto{\pgfqpoint{0.950556in}{0.573000in}}%
\pgfpathlineto{\pgfqpoint{0.950556in}{0.483000in}}%
\pgfpathclose%
\pgfusepath{fill}%
\end{pgfscope}%
\begin{pgfscope}%
\pgfpathrectangle{\pgfqpoint{0.595000in}{0.483000in}}{\pgfqpoint{2.835000in}{0.630000in}}%
\pgfusepath{clip}%
\pgfsetbuttcap%
\pgfsetmiterjoin%
\definecolor{currentfill}{rgb}{0.101961,0.101961,0.101961}%
\pgfsetfillcolor{currentfill}%
\pgfsetlinewidth{0.000000pt}%
\definecolor{currentstroke}{rgb}{0.000000,0.000000,0.000000}%
\pgfsetstrokecolor{currentstroke}%
\pgfsetstrokeopacity{0.000000}%
\pgfsetdash{}{0pt}%
\pgfpathmoveto{\pgfqpoint{0.974181in}{0.483000in}}%
\pgfpathlineto{\pgfqpoint{0.995444in}{0.483000in}}%
\pgfpathlineto{\pgfqpoint{0.995444in}{0.663000in}}%
\pgfpathlineto{\pgfqpoint{0.974181in}{0.663000in}}%
\pgfpathlineto{\pgfqpoint{0.974181in}{0.483000in}}%
\pgfpathclose%
\pgfusepath{fill}%
\end{pgfscope}%
\begin{pgfscope}%
\pgfpathrectangle{\pgfqpoint{0.595000in}{0.483000in}}{\pgfqpoint{2.835000in}{0.630000in}}%
\pgfusepath{clip}%
\pgfsetbuttcap%
\pgfsetmiterjoin%
\definecolor{currentfill}{rgb}{0.101961,0.101961,0.101961}%
\pgfsetfillcolor{currentfill}%
\pgfsetlinewidth{0.000000pt}%
\definecolor{currentstroke}{rgb}{0.000000,0.000000,0.000000}%
\pgfsetstrokecolor{currentstroke}%
\pgfsetstrokeopacity{0.000000}%
\pgfsetdash{}{0pt}%
\pgfpathmoveto{\pgfqpoint{0.997806in}{0.483000in}}%
\pgfpathlineto{\pgfqpoint{1.019069in}{0.483000in}}%
\pgfpathlineto{\pgfqpoint{1.019069in}{0.483000in}}%
\pgfpathlineto{\pgfqpoint{0.997806in}{0.483000in}}%
\pgfpathlineto{\pgfqpoint{0.997806in}{0.483000in}}%
\pgfpathclose%
\pgfusepath{fill}%
\end{pgfscope}%
\begin{pgfscope}%
\pgfpathrectangle{\pgfqpoint{0.595000in}{0.483000in}}{\pgfqpoint{2.835000in}{0.630000in}}%
\pgfusepath{clip}%
\pgfsetbuttcap%
\pgfsetmiterjoin%
\definecolor{currentfill}{rgb}{0.101961,0.101961,0.101961}%
\pgfsetfillcolor{currentfill}%
\pgfsetlinewidth{0.000000pt}%
\definecolor{currentstroke}{rgb}{0.000000,0.000000,0.000000}%
\pgfsetstrokecolor{currentstroke}%
\pgfsetstrokeopacity{0.000000}%
\pgfsetdash{}{0pt}%
\pgfpathmoveto{\pgfqpoint{1.021431in}{0.483000in}}%
\pgfpathlineto{\pgfqpoint{1.042694in}{0.483000in}}%
\pgfpathlineto{\pgfqpoint{1.042694in}{0.573000in}}%
\pgfpathlineto{\pgfqpoint{1.021431in}{0.573000in}}%
\pgfpathlineto{\pgfqpoint{1.021431in}{0.483000in}}%
\pgfpathclose%
\pgfusepath{fill}%
\end{pgfscope}%
\begin{pgfscope}%
\pgfpathrectangle{\pgfqpoint{0.595000in}{0.483000in}}{\pgfqpoint{2.835000in}{0.630000in}}%
\pgfusepath{clip}%
\pgfsetbuttcap%
\pgfsetmiterjoin%
\definecolor{currentfill}{rgb}{0.101961,0.101961,0.101961}%
\pgfsetfillcolor{currentfill}%
\pgfsetlinewidth{0.000000pt}%
\definecolor{currentstroke}{rgb}{0.000000,0.000000,0.000000}%
\pgfsetstrokecolor{currentstroke}%
\pgfsetstrokeopacity{0.000000}%
\pgfsetdash{}{0pt}%
\pgfpathmoveto{\pgfqpoint{1.045056in}{0.483000in}}%
\pgfpathlineto{\pgfqpoint{1.066319in}{0.483000in}}%
\pgfpathlineto{\pgfqpoint{1.066319in}{0.483000in}}%
\pgfpathlineto{\pgfqpoint{1.045056in}{0.483000in}}%
\pgfpathlineto{\pgfqpoint{1.045056in}{0.483000in}}%
\pgfpathclose%
\pgfusepath{fill}%
\end{pgfscope}%
\begin{pgfscope}%
\pgfpathrectangle{\pgfqpoint{0.595000in}{0.483000in}}{\pgfqpoint{2.835000in}{0.630000in}}%
\pgfusepath{clip}%
\pgfsetbuttcap%
\pgfsetmiterjoin%
\definecolor{currentfill}{rgb}{0.101961,0.101961,0.101961}%
\pgfsetfillcolor{currentfill}%
\pgfsetlinewidth{0.000000pt}%
\definecolor{currentstroke}{rgb}{0.000000,0.000000,0.000000}%
\pgfsetstrokecolor{currentstroke}%
\pgfsetstrokeopacity{0.000000}%
\pgfsetdash{}{0pt}%
\pgfpathmoveto{\pgfqpoint{1.068681in}{0.483000in}}%
\pgfpathlineto{\pgfqpoint{1.089944in}{0.483000in}}%
\pgfpathlineto{\pgfqpoint{1.089944in}{0.528000in}}%
\pgfpathlineto{\pgfqpoint{1.068681in}{0.528000in}}%
\pgfpathlineto{\pgfqpoint{1.068681in}{0.483000in}}%
\pgfpathclose%
\pgfusepath{fill}%
\end{pgfscope}%
\begin{pgfscope}%
\pgfpathrectangle{\pgfqpoint{0.595000in}{0.483000in}}{\pgfqpoint{2.835000in}{0.630000in}}%
\pgfusepath{clip}%
\pgfsetbuttcap%
\pgfsetmiterjoin%
\definecolor{currentfill}{rgb}{0.101961,0.101961,0.101961}%
\pgfsetfillcolor{currentfill}%
\pgfsetlinewidth{0.000000pt}%
\definecolor{currentstroke}{rgb}{0.000000,0.000000,0.000000}%
\pgfsetstrokecolor{currentstroke}%
\pgfsetstrokeopacity{0.000000}%
\pgfsetdash{}{0pt}%
\pgfpathmoveto{\pgfqpoint{1.092306in}{0.483000in}}%
\pgfpathlineto{\pgfqpoint{1.113569in}{0.483000in}}%
\pgfpathlineto{\pgfqpoint{1.113569in}{0.663000in}}%
\pgfpathlineto{\pgfqpoint{1.092306in}{0.663000in}}%
\pgfpathlineto{\pgfqpoint{1.092306in}{0.483000in}}%
\pgfpathclose%
\pgfusepath{fill}%
\end{pgfscope}%
\begin{pgfscope}%
\pgfpathrectangle{\pgfqpoint{0.595000in}{0.483000in}}{\pgfqpoint{2.835000in}{0.630000in}}%
\pgfusepath{clip}%
\pgfsetbuttcap%
\pgfsetmiterjoin%
\definecolor{currentfill}{rgb}{0.101961,0.101961,0.101961}%
\pgfsetfillcolor{currentfill}%
\pgfsetlinewidth{0.000000pt}%
\definecolor{currentstroke}{rgb}{0.000000,0.000000,0.000000}%
\pgfsetstrokecolor{currentstroke}%
\pgfsetstrokeopacity{0.000000}%
\pgfsetdash{}{0pt}%
\pgfpathmoveto{\pgfqpoint{1.115931in}{0.483000in}}%
\pgfpathlineto{\pgfqpoint{1.137194in}{0.483000in}}%
\pgfpathlineto{\pgfqpoint{1.137194in}{0.483000in}}%
\pgfpathlineto{\pgfqpoint{1.115931in}{0.483000in}}%
\pgfpathlineto{\pgfqpoint{1.115931in}{0.483000in}}%
\pgfpathclose%
\pgfusepath{fill}%
\end{pgfscope}%
\begin{pgfscope}%
\pgfpathrectangle{\pgfqpoint{0.595000in}{0.483000in}}{\pgfqpoint{2.835000in}{0.630000in}}%
\pgfusepath{clip}%
\pgfsetbuttcap%
\pgfsetmiterjoin%
\definecolor{currentfill}{rgb}{0.101961,0.101961,0.101961}%
\pgfsetfillcolor{currentfill}%
\pgfsetlinewidth{0.000000pt}%
\definecolor{currentstroke}{rgb}{0.000000,0.000000,0.000000}%
\pgfsetstrokecolor{currentstroke}%
\pgfsetstrokeopacity{0.000000}%
\pgfsetdash{}{0pt}%
\pgfpathmoveto{\pgfqpoint{1.139556in}{0.483000in}}%
\pgfpathlineto{\pgfqpoint{1.160819in}{0.483000in}}%
\pgfpathlineto{\pgfqpoint{1.160819in}{0.528000in}}%
\pgfpathlineto{\pgfqpoint{1.139556in}{0.528000in}}%
\pgfpathlineto{\pgfqpoint{1.139556in}{0.483000in}}%
\pgfpathclose%
\pgfusepath{fill}%
\end{pgfscope}%
\begin{pgfscope}%
\pgfpathrectangle{\pgfqpoint{0.595000in}{0.483000in}}{\pgfqpoint{2.835000in}{0.630000in}}%
\pgfusepath{clip}%
\pgfsetbuttcap%
\pgfsetmiterjoin%
\definecolor{currentfill}{rgb}{0.101961,0.101961,0.101961}%
\pgfsetfillcolor{currentfill}%
\pgfsetlinewidth{0.000000pt}%
\definecolor{currentstroke}{rgb}{0.000000,0.000000,0.000000}%
\pgfsetstrokecolor{currentstroke}%
\pgfsetstrokeopacity{0.000000}%
\pgfsetdash{}{0pt}%
\pgfpathmoveto{\pgfqpoint{1.163181in}{0.483000in}}%
\pgfpathlineto{\pgfqpoint{1.184444in}{0.483000in}}%
\pgfpathlineto{\pgfqpoint{1.184444in}{0.483000in}}%
\pgfpathlineto{\pgfqpoint{1.163181in}{0.483000in}}%
\pgfpathlineto{\pgfqpoint{1.163181in}{0.483000in}}%
\pgfpathclose%
\pgfusepath{fill}%
\end{pgfscope}%
\begin{pgfscope}%
\pgfpathrectangle{\pgfqpoint{0.595000in}{0.483000in}}{\pgfqpoint{2.835000in}{0.630000in}}%
\pgfusepath{clip}%
\pgfsetbuttcap%
\pgfsetmiterjoin%
\definecolor{currentfill}{rgb}{0.101961,0.101961,0.101961}%
\pgfsetfillcolor{currentfill}%
\pgfsetlinewidth{0.000000pt}%
\definecolor{currentstroke}{rgb}{0.000000,0.000000,0.000000}%
\pgfsetstrokecolor{currentstroke}%
\pgfsetstrokeopacity{0.000000}%
\pgfsetdash{}{0pt}%
\pgfpathmoveto{\pgfqpoint{1.186806in}{0.483000in}}%
\pgfpathlineto{\pgfqpoint{1.208069in}{0.483000in}}%
\pgfpathlineto{\pgfqpoint{1.208069in}{0.528000in}}%
\pgfpathlineto{\pgfqpoint{1.186806in}{0.528000in}}%
\pgfpathlineto{\pgfqpoint{1.186806in}{0.483000in}}%
\pgfpathclose%
\pgfusepath{fill}%
\end{pgfscope}%
\begin{pgfscope}%
\pgfpathrectangle{\pgfqpoint{0.595000in}{0.483000in}}{\pgfqpoint{2.835000in}{0.630000in}}%
\pgfusepath{clip}%
\pgfsetbuttcap%
\pgfsetmiterjoin%
\definecolor{currentfill}{rgb}{0.101961,0.101961,0.101961}%
\pgfsetfillcolor{currentfill}%
\pgfsetlinewidth{0.000000pt}%
\definecolor{currentstroke}{rgb}{0.000000,0.000000,0.000000}%
\pgfsetstrokecolor{currentstroke}%
\pgfsetstrokeopacity{0.000000}%
\pgfsetdash{}{0pt}%
\pgfpathmoveto{\pgfqpoint{1.210431in}{0.483000in}}%
\pgfpathlineto{\pgfqpoint{1.231694in}{0.483000in}}%
\pgfpathlineto{\pgfqpoint{1.231694in}{0.483000in}}%
\pgfpathlineto{\pgfqpoint{1.210431in}{0.483000in}}%
\pgfpathlineto{\pgfqpoint{1.210431in}{0.483000in}}%
\pgfpathclose%
\pgfusepath{fill}%
\end{pgfscope}%
\begin{pgfscope}%
\pgfpathrectangle{\pgfqpoint{0.595000in}{0.483000in}}{\pgfqpoint{2.835000in}{0.630000in}}%
\pgfusepath{clip}%
\pgfsetbuttcap%
\pgfsetmiterjoin%
\definecolor{currentfill}{rgb}{0.101961,0.101961,0.101961}%
\pgfsetfillcolor{currentfill}%
\pgfsetlinewidth{0.000000pt}%
\definecolor{currentstroke}{rgb}{0.000000,0.000000,0.000000}%
\pgfsetstrokecolor{currentstroke}%
\pgfsetstrokeopacity{0.000000}%
\pgfsetdash{}{0pt}%
\pgfpathmoveto{\pgfqpoint{1.234056in}{0.483000in}}%
\pgfpathlineto{\pgfqpoint{1.255319in}{0.483000in}}%
\pgfpathlineto{\pgfqpoint{1.255319in}{0.483000in}}%
\pgfpathlineto{\pgfqpoint{1.234056in}{0.483000in}}%
\pgfpathlineto{\pgfqpoint{1.234056in}{0.483000in}}%
\pgfpathclose%
\pgfusepath{fill}%
\end{pgfscope}%
\begin{pgfscope}%
\pgfpathrectangle{\pgfqpoint{0.595000in}{0.483000in}}{\pgfqpoint{2.835000in}{0.630000in}}%
\pgfusepath{clip}%
\pgfsetbuttcap%
\pgfsetmiterjoin%
\definecolor{currentfill}{rgb}{0.101961,0.101961,0.101961}%
\pgfsetfillcolor{currentfill}%
\pgfsetlinewidth{0.000000pt}%
\definecolor{currentstroke}{rgb}{0.000000,0.000000,0.000000}%
\pgfsetstrokecolor{currentstroke}%
\pgfsetstrokeopacity{0.000000}%
\pgfsetdash{}{0pt}%
\pgfpathmoveto{\pgfqpoint{1.257681in}{0.483000in}}%
\pgfpathlineto{\pgfqpoint{1.278944in}{0.483000in}}%
\pgfpathlineto{\pgfqpoint{1.278944in}{0.573000in}}%
\pgfpathlineto{\pgfqpoint{1.257681in}{0.573000in}}%
\pgfpathlineto{\pgfqpoint{1.257681in}{0.483000in}}%
\pgfpathclose%
\pgfusepath{fill}%
\end{pgfscope}%
\begin{pgfscope}%
\pgfpathrectangle{\pgfqpoint{0.595000in}{0.483000in}}{\pgfqpoint{2.835000in}{0.630000in}}%
\pgfusepath{clip}%
\pgfsetbuttcap%
\pgfsetmiterjoin%
\definecolor{currentfill}{rgb}{0.101961,0.101961,0.101961}%
\pgfsetfillcolor{currentfill}%
\pgfsetlinewidth{0.000000pt}%
\definecolor{currentstroke}{rgb}{0.000000,0.000000,0.000000}%
\pgfsetstrokecolor{currentstroke}%
\pgfsetstrokeopacity{0.000000}%
\pgfsetdash{}{0pt}%
\pgfpathmoveto{\pgfqpoint{1.281306in}{0.483000in}}%
\pgfpathlineto{\pgfqpoint{1.302569in}{0.483000in}}%
\pgfpathlineto{\pgfqpoint{1.302569in}{0.483000in}}%
\pgfpathlineto{\pgfqpoint{1.281306in}{0.483000in}}%
\pgfpathlineto{\pgfqpoint{1.281306in}{0.483000in}}%
\pgfpathclose%
\pgfusepath{fill}%
\end{pgfscope}%
\begin{pgfscope}%
\pgfpathrectangle{\pgfqpoint{0.595000in}{0.483000in}}{\pgfqpoint{2.835000in}{0.630000in}}%
\pgfusepath{clip}%
\pgfsetbuttcap%
\pgfsetmiterjoin%
\definecolor{currentfill}{rgb}{0.101961,0.101961,0.101961}%
\pgfsetfillcolor{currentfill}%
\pgfsetlinewidth{0.000000pt}%
\definecolor{currentstroke}{rgb}{0.000000,0.000000,0.000000}%
\pgfsetstrokecolor{currentstroke}%
\pgfsetstrokeopacity{0.000000}%
\pgfsetdash{}{0pt}%
\pgfpathmoveto{\pgfqpoint{1.304931in}{0.483000in}}%
\pgfpathlineto{\pgfqpoint{1.326194in}{0.483000in}}%
\pgfpathlineto{\pgfqpoint{1.326194in}{0.483000in}}%
\pgfpathlineto{\pgfqpoint{1.304931in}{0.483000in}}%
\pgfpathlineto{\pgfqpoint{1.304931in}{0.483000in}}%
\pgfpathclose%
\pgfusepath{fill}%
\end{pgfscope}%
\begin{pgfscope}%
\pgfpathrectangle{\pgfqpoint{0.595000in}{0.483000in}}{\pgfqpoint{2.835000in}{0.630000in}}%
\pgfusepath{clip}%
\pgfsetbuttcap%
\pgfsetmiterjoin%
\definecolor{currentfill}{rgb}{0.101961,0.101961,0.101961}%
\pgfsetfillcolor{currentfill}%
\pgfsetlinewidth{0.000000pt}%
\definecolor{currentstroke}{rgb}{0.000000,0.000000,0.000000}%
\pgfsetstrokecolor{currentstroke}%
\pgfsetstrokeopacity{0.000000}%
\pgfsetdash{}{0pt}%
\pgfpathmoveto{\pgfqpoint{1.328556in}{0.483000in}}%
\pgfpathlineto{\pgfqpoint{1.349819in}{0.483000in}}%
\pgfpathlineto{\pgfqpoint{1.349819in}{0.618000in}}%
\pgfpathlineto{\pgfqpoint{1.328556in}{0.618000in}}%
\pgfpathlineto{\pgfqpoint{1.328556in}{0.483000in}}%
\pgfpathclose%
\pgfusepath{fill}%
\end{pgfscope}%
\begin{pgfscope}%
\pgfpathrectangle{\pgfqpoint{0.595000in}{0.483000in}}{\pgfqpoint{2.835000in}{0.630000in}}%
\pgfusepath{clip}%
\pgfsetbuttcap%
\pgfsetmiterjoin%
\definecolor{currentfill}{rgb}{0.101961,0.101961,0.101961}%
\pgfsetfillcolor{currentfill}%
\pgfsetlinewidth{0.000000pt}%
\definecolor{currentstroke}{rgb}{0.000000,0.000000,0.000000}%
\pgfsetstrokecolor{currentstroke}%
\pgfsetstrokeopacity{0.000000}%
\pgfsetdash{}{0pt}%
\pgfpathmoveto{\pgfqpoint{1.352181in}{0.483000in}}%
\pgfpathlineto{\pgfqpoint{1.373444in}{0.483000in}}%
\pgfpathlineto{\pgfqpoint{1.373444in}{0.573000in}}%
\pgfpathlineto{\pgfqpoint{1.352181in}{0.573000in}}%
\pgfpathlineto{\pgfqpoint{1.352181in}{0.483000in}}%
\pgfpathclose%
\pgfusepath{fill}%
\end{pgfscope}%
\begin{pgfscope}%
\pgfpathrectangle{\pgfqpoint{0.595000in}{0.483000in}}{\pgfqpoint{2.835000in}{0.630000in}}%
\pgfusepath{clip}%
\pgfsetbuttcap%
\pgfsetmiterjoin%
\definecolor{currentfill}{rgb}{0.101961,0.101961,0.101961}%
\pgfsetfillcolor{currentfill}%
\pgfsetlinewidth{0.000000pt}%
\definecolor{currentstroke}{rgb}{0.000000,0.000000,0.000000}%
\pgfsetstrokecolor{currentstroke}%
\pgfsetstrokeopacity{0.000000}%
\pgfsetdash{}{0pt}%
\pgfpathmoveto{\pgfqpoint{1.375806in}{0.483000in}}%
\pgfpathlineto{\pgfqpoint{1.397069in}{0.483000in}}%
\pgfpathlineto{\pgfqpoint{1.397069in}{0.483000in}}%
\pgfpathlineto{\pgfqpoint{1.375806in}{0.483000in}}%
\pgfpathlineto{\pgfqpoint{1.375806in}{0.483000in}}%
\pgfpathclose%
\pgfusepath{fill}%
\end{pgfscope}%
\begin{pgfscope}%
\pgfpathrectangle{\pgfqpoint{0.595000in}{0.483000in}}{\pgfqpoint{2.835000in}{0.630000in}}%
\pgfusepath{clip}%
\pgfsetbuttcap%
\pgfsetmiterjoin%
\definecolor{currentfill}{rgb}{0.101961,0.101961,0.101961}%
\pgfsetfillcolor{currentfill}%
\pgfsetlinewidth{0.000000pt}%
\definecolor{currentstroke}{rgb}{0.000000,0.000000,0.000000}%
\pgfsetstrokecolor{currentstroke}%
\pgfsetstrokeopacity{0.000000}%
\pgfsetdash{}{0pt}%
\pgfpathmoveto{\pgfqpoint{1.399431in}{0.483000in}}%
\pgfpathlineto{\pgfqpoint{1.420694in}{0.483000in}}%
\pgfpathlineto{\pgfqpoint{1.420694in}{0.528000in}}%
\pgfpathlineto{\pgfqpoint{1.399431in}{0.528000in}}%
\pgfpathlineto{\pgfqpoint{1.399431in}{0.483000in}}%
\pgfpathclose%
\pgfusepath{fill}%
\end{pgfscope}%
\begin{pgfscope}%
\pgfpathrectangle{\pgfqpoint{0.595000in}{0.483000in}}{\pgfqpoint{2.835000in}{0.630000in}}%
\pgfusepath{clip}%
\pgfsetbuttcap%
\pgfsetmiterjoin%
\definecolor{currentfill}{rgb}{0.101961,0.101961,0.101961}%
\pgfsetfillcolor{currentfill}%
\pgfsetlinewidth{0.000000pt}%
\definecolor{currentstroke}{rgb}{0.000000,0.000000,0.000000}%
\pgfsetstrokecolor{currentstroke}%
\pgfsetstrokeopacity{0.000000}%
\pgfsetdash{}{0pt}%
\pgfpathmoveto{\pgfqpoint{1.423056in}{0.483000in}}%
\pgfpathlineto{\pgfqpoint{1.444319in}{0.483000in}}%
\pgfpathlineto{\pgfqpoint{1.444319in}{0.528000in}}%
\pgfpathlineto{\pgfqpoint{1.423056in}{0.528000in}}%
\pgfpathlineto{\pgfqpoint{1.423056in}{0.483000in}}%
\pgfpathclose%
\pgfusepath{fill}%
\end{pgfscope}%
\begin{pgfscope}%
\pgfpathrectangle{\pgfqpoint{0.595000in}{0.483000in}}{\pgfqpoint{2.835000in}{0.630000in}}%
\pgfusepath{clip}%
\pgfsetbuttcap%
\pgfsetmiterjoin%
\definecolor{currentfill}{rgb}{0.101961,0.101961,0.101961}%
\pgfsetfillcolor{currentfill}%
\pgfsetlinewidth{0.000000pt}%
\definecolor{currentstroke}{rgb}{0.000000,0.000000,0.000000}%
\pgfsetstrokecolor{currentstroke}%
\pgfsetstrokeopacity{0.000000}%
\pgfsetdash{}{0pt}%
\pgfpathmoveto{\pgfqpoint{1.446681in}{0.483000in}}%
\pgfpathlineto{\pgfqpoint{1.467944in}{0.483000in}}%
\pgfpathlineto{\pgfqpoint{1.467944in}{0.528000in}}%
\pgfpathlineto{\pgfqpoint{1.446681in}{0.528000in}}%
\pgfpathlineto{\pgfqpoint{1.446681in}{0.483000in}}%
\pgfpathclose%
\pgfusepath{fill}%
\end{pgfscope}%
\begin{pgfscope}%
\pgfpathrectangle{\pgfqpoint{0.595000in}{0.483000in}}{\pgfqpoint{2.835000in}{0.630000in}}%
\pgfusepath{clip}%
\pgfsetbuttcap%
\pgfsetmiterjoin%
\definecolor{currentfill}{rgb}{0.101961,0.101961,0.101961}%
\pgfsetfillcolor{currentfill}%
\pgfsetlinewidth{0.000000pt}%
\definecolor{currentstroke}{rgb}{0.000000,0.000000,0.000000}%
\pgfsetstrokecolor{currentstroke}%
\pgfsetstrokeopacity{0.000000}%
\pgfsetdash{}{0pt}%
\pgfpathmoveto{\pgfqpoint{1.470306in}{0.483000in}}%
\pgfpathlineto{\pgfqpoint{1.491569in}{0.483000in}}%
\pgfpathlineto{\pgfqpoint{1.491569in}{0.573000in}}%
\pgfpathlineto{\pgfqpoint{1.470306in}{0.573000in}}%
\pgfpathlineto{\pgfqpoint{1.470306in}{0.483000in}}%
\pgfpathclose%
\pgfusepath{fill}%
\end{pgfscope}%
\begin{pgfscope}%
\pgfpathrectangle{\pgfqpoint{0.595000in}{0.483000in}}{\pgfqpoint{2.835000in}{0.630000in}}%
\pgfusepath{clip}%
\pgfsetbuttcap%
\pgfsetmiterjoin%
\definecolor{currentfill}{rgb}{0.101961,0.101961,0.101961}%
\pgfsetfillcolor{currentfill}%
\pgfsetlinewidth{0.000000pt}%
\definecolor{currentstroke}{rgb}{0.000000,0.000000,0.000000}%
\pgfsetstrokecolor{currentstroke}%
\pgfsetstrokeopacity{0.000000}%
\pgfsetdash{}{0pt}%
\pgfpathmoveto{\pgfqpoint{1.493931in}{0.483000in}}%
\pgfpathlineto{\pgfqpoint{1.515194in}{0.483000in}}%
\pgfpathlineto{\pgfqpoint{1.515194in}{0.528000in}}%
\pgfpathlineto{\pgfqpoint{1.493931in}{0.528000in}}%
\pgfpathlineto{\pgfqpoint{1.493931in}{0.483000in}}%
\pgfpathclose%
\pgfusepath{fill}%
\end{pgfscope}%
\begin{pgfscope}%
\pgfpathrectangle{\pgfqpoint{0.595000in}{0.483000in}}{\pgfqpoint{2.835000in}{0.630000in}}%
\pgfusepath{clip}%
\pgfsetbuttcap%
\pgfsetmiterjoin%
\definecolor{currentfill}{rgb}{0.101961,0.101961,0.101961}%
\pgfsetfillcolor{currentfill}%
\pgfsetlinewidth{0.000000pt}%
\definecolor{currentstroke}{rgb}{0.000000,0.000000,0.000000}%
\pgfsetstrokecolor{currentstroke}%
\pgfsetstrokeopacity{0.000000}%
\pgfsetdash{}{0pt}%
\pgfpathmoveto{\pgfqpoint{1.517556in}{0.483000in}}%
\pgfpathlineto{\pgfqpoint{1.538819in}{0.483000in}}%
\pgfpathlineto{\pgfqpoint{1.538819in}{0.528000in}}%
\pgfpathlineto{\pgfqpoint{1.517556in}{0.528000in}}%
\pgfpathlineto{\pgfqpoint{1.517556in}{0.483000in}}%
\pgfpathclose%
\pgfusepath{fill}%
\end{pgfscope}%
\begin{pgfscope}%
\pgfpathrectangle{\pgfqpoint{0.595000in}{0.483000in}}{\pgfqpoint{2.835000in}{0.630000in}}%
\pgfusepath{clip}%
\pgfsetbuttcap%
\pgfsetmiterjoin%
\definecolor{currentfill}{rgb}{0.101961,0.101961,0.101961}%
\pgfsetfillcolor{currentfill}%
\pgfsetlinewidth{0.000000pt}%
\definecolor{currentstroke}{rgb}{0.000000,0.000000,0.000000}%
\pgfsetstrokecolor{currentstroke}%
\pgfsetstrokeopacity{0.000000}%
\pgfsetdash{}{0pt}%
\pgfpathmoveto{\pgfqpoint{1.541181in}{0.483000in}}%
\pgfpathlineto{\pgfqpoint{1.562444in}{0.483000in}}%
\pgfpathlineto{\pgfqpoint{1.562444in}{0.483000in}}%
\pgfpathlineto{\pgfqpoint{1.541181in}{0.483000in}}%
\pgfpathlineto{\pgfqpoint{1.541181in}{0.483000in}}%
\pgfpathclose%
\pgfusepath{fill}%
\end{pgfscope}%
\begin{pgfscope}%
\pgfpathrectangle{\pgfqpoint{0.595000in}{0.483000in}}{\pgfqpoint{2.835000in}{0.630000in}}%
\pgfusepath{clip}%
\pgfsetbuttcap%
\pgfsetmiterjoin%
\definecolor{currentfill}{rgb}{0.101961,0.101961,0.101961}%
\pgfsetfillcolor{currentfill}%
\pgfsetlinewidth{0.000000pt}%
\definecolor{currentstroke}{rgb}{0.000000,0.000000,0.000000}%
\pgfsetstrokecolor{currentstroke}%
\pgfsetstrokeopacity{0.000000}%
\pgfsetdash{}{0pt}%
\pgfpathmoveto{\pgfqpoint{1.564806in}{0.483000in}}%
\pgfpathlineto{\pgfqpoint{1.586069in}{0.483000in}}%
\pgfpathlineto{\pgfqpoint{1.586069in}{0.483000in}}%
\pgfpathlineto{\pgfqpoint{1.564806in}{0.483000in}}%
\pgfpathlineto{\pgfqpoint{1.564806in}{0.483000in}}%
\pgfpathclose%
\pgfusepath{fill}%
\end{pgfscope}%
\begin{pgfscope}%
\pgfpathrectangle{\pgfqpoint{0.595000in}{0.483000in}}{\pgfqpoint{2.835000in}{0.630000in}}%
\pgfusepath{clip}%
\pgfsetbuttcap%
\pgfsetmiterjoin%
\definecolor{currentfill}{rgb}{0.101961,0.101961,0.101961}%
\pgfsetfillcolor{currentfill}%
\pgfsetlinewidth{0.000000pt}%
\definecolor{currentstroke}{rgb}{0.000000,0.000000,0.000000}%
\pgfsetstrokecolor{currentstroke}%
\pgfsetstrokeopacity{0.000000}%
\pgfsetdash{}{0pt}%
\pgfpathmoveto{\pgfqpoint{1.588431in}{0.483000in}}%
\pgfpathlineto{\pgfqpoint{1.609694in}{0.483000in}}%
\pgfpathlineto{\pgfqpoint{1.609694in}{0.483000in}}%
\pgfpathlineto{\pgfqpoint{1.588431in}{0.483000in}}%
\pgfpathlineto{\pgfqpoint{1.588431in}{0.483000in}}%
\pgfpathclose%
\pgfusepath{fill}%
\end{pgfscope}%
\begin{pgfscope}%
\pgfpathrectangle{\pgfqpoint{0.595000in}{0.483000in}}{\pgfqpoint{2.835000in}{0.630000in}}%
\pgfusepath{clip}%
\pgfsetbuttcap%
\pgfsetmiterjoin%
\definecolor{currentfill}{rgb}{0.101961,0.101961,0.101961}%
\pgfsetfillcolor{currentfill}%
\pgfsetlinewidth{0.000000pt}%
\definecolor{currentstroke}{rgb}{0.000000,0.000000,0.000000}%
\pgfsetstrokecolor{currentstroke}%
\pgfsetstrokeopacity{0.000000}%
\pgfsetdash{}{0pt}%
\pgfpathmoveto{\pgfqpoint{1.612056in}{0.483000in}}%
\pgfpathlineto{\pgfqpoint{1.633319in}{0.483000in}}%
\pgfpathlineto{\pgfqpoint{1.633319in}{0.528000in}}%
\pgfpathlineto{\pgfqpoint{1.612056in}{0.528000in}}%
\pgfpathlineto{\pgfqpoint{1.612056in}{0.483000in}}%
\pgfpathclose%
\pgfusepath{fill}%
\end{pgfscope}%
\begin{pgfscope}%
\pgfpathrectangle{\pgfqpoint{0.595000in}{0.483000in}}{\pgfqpoint{2.835000in}{0.630000in}}%
\pgfusepath{clip}%
\pgfsetbuttcap%
\pgfsetmiterjoin%
\definecolor{currentfill}{rgb}{0.101961,0.101961,0.101961}%
\pgfsetfillcolor{currentfill}%
\pgfsetlinewidth{0.000000pt}%
\definecolor{currentstroke}{rgb}{0.000000,0.000000,0.000000}%
\pgfsetstrokecolor{currentstroke}%
\pgfsetstrokeopacity{0.000000}%
\pgfsetdash{}{0pt}%
\pgfpathmoveto{\pgfqpoint{1.635681in}{0.483000in}}%
\pgfpathlineto{\pgfqpoint{1.656944in}{0.483000in}}%
\pgfpathlineto{\pgfqpoint{1.656944in}{0.483000in}}%
\pgfpathlineto{\pgfqpoint{1.635681in}{0.483000in}}%
\pgfpathlineto{\pgfqpoint{1.635681in}{0.483000in}}%
\pgfpathclose%
\pgfusepath{fill}%
\end{pgfscope}%
\begin{pgfscope}%
\pgfpathrectangle{\pgfqpoint{0.595000in}{0.483000in}}{\pgfqpoint{2.835000in}{0.630000in}}%
\pgfusepath{clip}%
\pgfsetbuttcap%
\pgfsetmiterjoin%
\definecolor{currentfill}{rgb}{0.101961,0.101961,0.101961}%
\pgfsetfillcolor{currentfill}%
\pgfsetlinewidth{0.000000pt}%
\definecolor{currentstroke}{rgb}{0.000000,0.000000,0.000000}%
\pgfsetstrokecolor{currentstroke}%
\pgfsetstrokeopacity{0.000000}%
\pgfsetdash{}{0pt}%
\pgfpathmoveto{\pgfqpoint{1.659306in}{0.483000in}}%
\pgfpathlineto{\pgfqpoint{1.680569in}{0.483000in}}%
\pgfpathlineto{\pgfqpoint{1.680569in}{0.618000in}}%
\pgfpathlineto{\pgfqpoint{1.659306in}{0.618000in}}%
\pgfpathlineto{\pgfqpoint{1.659306in}{0.483000in}}%
\pgfpathclose%
\pgfusepath{fill}%
\end{pgfscope}%
\begin{pgfscope}%
\pgfpathrectangle{\pgfqpoint{0.595000in}{0.483000in}}{\pgfqpoint{2.835000in}{0.630000in}}%
\pgfusepath{clip}%
\pgfsetbuttcap%
\pgfsetmiterjoin%
\definecolor{currentfill}{rgb}{0.101961,0.101961,0.101961}%
\pgfsetfillcolor{currentfill}%
\pgfsetlinewidth{0.000000pt}%
\definecolor{currentstroke}{rgb}{0.000000,0.000000,0.000000}%
\pgfsetstrokecolor{currentstroke}%
\pgfsetstrokeopacity{0.000000}%
\pgfsetdash{}{0pt}%
\pgfpathmoveto{\pgfqpoint{1.682931in}{0.483000in}}%
\pgfpathlineto{\pgfqpoint{1.704194in}{0.483000in}}%
\pgfpathlineto{\pgfqpoint{1.704194in}{0.483000in}}%
\pgfpathlineto{\pgfqpoint{1.682931in}{0.483000in}}%
\pgfpathlineto{\pgfqpoint{1.682931in}{0.483000in}}%
\pgfpathclose%
\pgfusepath{fill}%
\end{pgfscope}%
\begin{pgfscope}%
\pgfpathrectangle{\pgfqpoint{0.595000in}{0.483000in}}{\pgfqpoint{2.835000in}{0.630000in}}%
\pgfusepath{clip}%
\pgfsetbuttcap%
\pgfsetmiterjoin%
\definecolor{currentfill}{rgb}{0.101961,0.101961,0.101961}%
\pgfsetfillcolor{currentfill}%
\pgfsetlinewidth{0.000000pt}%
\definecolor{currentstroke}{rgb}{0.000000,0.000000,0.000000}%
\pgfsetstrokecolor{currentstroke}%
\pgfsetstrokeopacity{0.000000}%
\pgfsetdash{}{0pt}%
\pgfpathmoveto{\pgfqpoint{1.706556in}{0.483000in}}%
\pgfpathlineto{\pgfqpoint{1.727819in}{0.483000in}}%
\pgfpathlineto{\pgfqpoint{1.727819in}{0.843000in}}%
\pgfpathlineto{\pgfqpoint{1.706556in}{0.843000in}}%
\pgfpathlineto{\pgfqpoint{1.706556in}{0.483000in}}%
\pgfpathclose%
\pgfusepath{fill}%
\end{pgfscope}%
\begin{pgfscope}%
\pgfpathrectangle{\pgfqpoint{0.595000in}{0.483000in}}{\pgfqpoint{2.835000in}{0.630000in}}%
\pgfusepath{clip}%
\pgfsetbuttcap%
\pgfsetmiterjoin%
\definecolor{currentfill}{rgb}{0.101961,0.101961,0.101961}%
\pgfsetfillcolor{currentfill}%
\pgfsetlinewidth{0.000000pt}%
\definecolor{currentstroke}{rgb}{0.000000,0.000000,0.000000}%
\pgfsetstrokecolor{currentstroke}%
\pgfsetstrokeopacity{0.000000}%
\pgfsetdash{}{0pt}%
\pgfpathmoveto{\pgfqpoint{1.730181in}{0.483000in}}%
\pgfpathlineto{\pgfqpoint{1.751444in}{0.483000in}}%
\pgfpathlineto{\pgfqpoint{1.751444in}{0.528000in}}%
\pgfpathlineto{\pgfqpoint{1.730181in}{0.528000in}}%
\pgfpathlineto{\pgfqpoint{1.730181in}{0.483000in}}%
\pgfpathclose%
\pgfusepath{fill}%
\end{pgfscope}%
\begin{pgfscope}%
\pgfpathrectangle{\pgfqpoint{0.595000in}{0.483000in}}{\pgfqpoint{2.835000in}{0.630000in}}%
\pgfusepath{clip}%
\pgfsetbuttcap%
\pgfsetmiterjoin%
\definecolor{currentfill}{rgb}{0.101961,0.101961,0.101961}%
\pgfsetfillcolor{currentfill}%
\pgfsetlinewidth{0.000000pt}%
\definecolor{currentstroke}{rgb}{0.000000,0.000000,0.000000}%
\pgfsetstrokecolor{currentstroke}%
\pgfsetstrokeopacity{0.000000}%
\pgfsetdash{}{0pt}%
\pgfpathmoveto{\pgfqpoint{1.753806in}{0.483000in}}%
\pgfpathlineto{\pgfqpoint{1.775069in}{0.483000in}}%
\pgfpathlineto{\pgfqpoint{1.775069in}{0.483000in}}%
\pgfpathlineto{\pgfqpoint{1.753806in}{0.483000in}}%
\pgfpathlineto{\pgfqpoint{1.753806in}{0.483000in}}%
\pgfpathclose%
\pgfusepath{fill}%
\end{pgfscope}%
\begin{pgfscope}%
\pgfpathrectangle{\pgfqpoint{0.595000in}{0.483000in}}{\pgfqpoint{2.835000in}{0.630000in}}%
\pgfusepath{clip}%
\pgfsetbuttcap%
\pgfsetmiterjoin%
\definecolor{currentfill}{rgb}{0.101961,0.101961,0.101961}%
\pgfsetfillcolor{currentfill}%
\pgfsetlinewidth{0.000000pt}%
\definecolor{currentstroke}{rgb}{0.000000,0.000000,0.000000}%
\pgfsetstrokecolor{currentstroke}%
\pgfsetstrokeopacity{0.000000}%
\pgfsetdash{}{0pt}%
\pgfpathmoveto{\pgfqpoint{1.777431in}{0.483000in}}%
\pgfpathlineto{\pgfqpoint{1.798694in}{0.483000in}}%
\pgfpathlineto{\pgfqpoint{1.798694in}{0.483000in}}%
\pgfpathlineto{\pgfqpoint{1.777431in}{0.483000in}}%
\pgfpathlineto{\pgfqpoint{1.777431in}{0.483000in}}%
\pgfpathclose%
\pgfusepath{fill}%
\end{pgfscope}%
\begin{pgfscope}%
\pgfpathrectangle{\pgfqpoint{0.595000in}{0.483000in}}{\pgfqpoint{2.835000in}{0.630000in}}%
\pgfusepath{clip}%
\pgfsetbuttcap%
\pgfsetmiterjoin%
\definecolor{currentfill}{rgb}{0.101961,0.101961,0.101961}%
\pgfsetfillcolor{currentfill}%
\pgfsetlinewidth{0.000000pt}%
\definecolor{currentstroke}{rgb}{0.000000,0.000000,0.000000}%
\pgfsetstrokecolor{currentstroke}%
\pgfsetstrokeopacity{0.000000}%
\pgfsetdash{}{0pt}%
\pgfpathmoveto{\pgfqpoint{1.801056in}{0.483000in}}%
\pgfpathlineto{\pgfqpoint{1.822319in}{0.483000in}}%
\pgfpathlineto{\pgfqpoint{1.822319in}{0.528000in}}%
\pgfpathlineto{\pgfqpoint{1.801056in}{0.528000in}}%
\pgfpathlineto{\pgfqpoint{1.801056in}{0.483000in}}%
\pgfpathclose%
\pgfusepath{fill}%
\end{pgfscope}%
\begin{pgfscope}%
\pgfpathrectangle{\pgfqpoint{0.595000in}{0.483000in}}{\pgfqpoint{2.835000in}{0.630000in}}%
\pgfusepath{clip}%
\pgfsetbuttcap%
\pgfsetmiterjoin%
\definecolor{currentfill}{rgb}{0.101961,0.101961,0.101961}%
\pgfsetfillcolor{currentfill}%
\pgfsetlinewidth{0.000000pt}%
\definecolor{currentstroke}{rgb}{0.000000,0.000000,0.000000}%
\pgfsetstrokecolor{currentstroke}%
\pgfsetstrokeopacity{0.000000}%
\pgfsetdash{}{0pt}%
\pgfpathmoveto{\pgfqpoint{1.824681in}{0.483000in}}%
\pgfpathlineto{\pgfqpoint{1.845944in}{0.483000in}}%
\pgfpathlineto{\pgfqpoint{1.845944in}{0.528000in}}%
\pgfpathlineto{\pgfqpoint{1.824681in}{0.528000in}}%
\pgfpathlineto{\pgfqpoint{1.824681in}{0.483000in}}%
\pgfpathclose%
\pgfusepath{fill}%
\end{pgfscope}%
\begin{pgfscope}%
\pgfpathrectangle{\pgfqpoint{0.595000in}{0.483000in}}{\pgfqpoint{2.835000in}{0.630000in}}%
\pgfusepath{clip}%
\pgfsetbuttcap%
\pgfsetmiterjoin%
\definecolor{currentfill}{rgb}{0.101961,0.101961,0.101961}%
\pgfsetfillcolor{currentfill}%
\pgfsetlinewidth{0.000000pt}%
\definecolor{currentstroke}{rgb}{0.000000,0.000000,0.000000}%
\pgfsetstrokecolor{currentstroke}%
\pgfsetstrokeopacity{0.000000}%
\pgfsetdash{}{0pt}%
\pgfpathmoveto{\pgfqpoint{1.848306in}{0.483000in}}%
\pgfpathlineto{\pgfqpoint{1.869569in}{0.483000in}}%
\pgfpathlineto{\pgfqpoint{1.869569in}{0.483000in}}%
\pgfpathlineto{\pgfqpoint{1.848306in}{0.483000in}}%
\pgfpathlineto{\pgfqpoint{1.848306in}{0.483000in}}%
\pgfpathclose%
\pgfusepath{fill}%
\end{pgfscope}%
\begin{pgfscope}%
\pgfpathrectangle{\pgfqpoint{0.595000in}{0.483000in}}{\pgfqpoint{2.835000in}{0.630000in}}%
\pgfusepath{clip}%
\pgfsetbuttcap%
\pgfsetmiterjoin%
\definecolor{currentfill}{rgb}{0.101961,0.101961,0.101961}%
\pgfsetfillcolor{currentfill}%
\pgfsetlinewidth{0.000000pt}%
\definecolor{currentstroke}{rgb}{0.000000,0.000000,0.000000}%
\pgfsetstrokecolor{currentstroke}%
\pgfsetstrokeopacity{0.000000}%
\pgfsetdash{}{0pt}%
\pgfpathmoveto{\pgfqpoint{1.871931in}{0.483000in}}%
\pgfpathlineto{\pgfqpoint{1.893194in}{0.483000in}}%
\pgfpathlineto{\pgfqpoint{1.893194in}{0.528000in}}%
\pgfpathlineto{\pgfqpoint{1.871931in}{0.528000in}}%
\pgfpathlineto{\pgfqpoint{1.871931in}{0.483000in}}%
\pgfpathclose%
\pgfusepath{fill}%
\end{pgfscope}%
\begin{pgfscope}%
\pgfpathrectangle{\pgfqpoint{0.595000in}{0.483000in}}{\pgfqpoint{2.835000in}{0.630000in}}%
\pgfusepath{clip}%
\pgfsetbuttcap%
\pgfsetmiterjoin%
\definecolor{currentfill}{rgb}{0.101961,0.101961,0.101961}%
\pgfsetfillcolor{currentfill}%
\pgfsetlinewidth{0.000000pt}%
\definecolor{currentstroke}{rgb}{0.000000,0.000000,0.000000}%
\pgfsetstrokecolor{currentstroke}%
\pgfsetstrokeopacity{0.000000}%
\pgfsetdash{}{0pt}%
\pgfpathmoveto{\pgfqpoint{1.895556in}{0.483000in}}%
\pgfpathlineto{\pgfqpoint{1.916819in}{0.483000in}}%
\pgfpathlineto{\pgfqpoint{1.916819in}{0.483000in}}%
\pgfpathlineto{\pgfqpoint{1.895556in}{0.483000in}}%
\pgfpathlineto{\pgfqpoint{1.895556in}{0.483000in}}%
\pgfpathclose%
\pgfusepath{fill}%
\end{pgfscope}%
\begin{pgfscope}%
\pgfpathrectangle{\pgfqpoint{0.595000in}{0.483000in}}{\pgfqpoint{2.835000in}{0.630000in}}%
\pgfusepath{clip}%
\pgfsetbuttcap%
\pgfsetmiterjoin%
\definecolor{currentfill}{rgb}{0.101961,0.101961,0.101961}%
\pgfsetfillcolor{currentfill}%
\pgfsetlinewidth{0.000000pt}%
\definecolor{currentstroke}{rgb}{0.000000,0.000000,0.000000}%
\pgfsetstrokecolor{currentstroke}%
\pgfsetstrokeopacity{0.000000}%
\pgfsetdash{}{0pt}%
\pgfpathmoveto{\pgfqpoint{1.919181in}{0.483000in}}%
\pgfpathlineto{\pgfqpoint{1.940444in}{0.483000in}}%
\pgfpathlineto{\pgfqpoint{1.940444in}{0.483000in}}%
\pgfpathlineto{\pgfqpoint{1.919181in}{0.483000in}}%
\pgfpathlineto{\pgfqpoint{1.919181in}{0.483000in}}%
\pgfpathclose%
\pgfusepath{fill}%
\end{pgfscope}%
\begin{pgfscope}%
\pgfpathrectangle{\pgfqpoint{0.595000in}{0.483000in}}{\pgfqpoint{2.835000in}{0.630000in}}%
\pgfusepath{clip}%
\pgfsetbuttcap%
\pgfsetmiterjoin%
\definecolor{currentfill}{rgb}{0.101961,0.101961,0.101961}%
\pgfsetfillcolor{currentfill}%
\pgfsetlinewidth{0.000000pt}%
\definecolor{currentstroke}{rgb}{0.000000,0.000000,0.000000}%
\pgfsetstrokecolor{currentstroke}%
\pgfsetstrokeopacity{0.000000}%
\pgfsetdash{}{0pt}%
\pgfpathmoveto{\pgfqpoint{1.942806in}{0.483000in}}%
\pgfpathlineto{\pgfqpoint{1.964069in}{0.483000in}}%
\pgfpathlineto{\pgfqpoint{1.964069in}{0.528000in}}%
\pgfpathlineto{\pgfqpoint{1.942806in}{0.528000in}}%
\pgfpathlineto{\pgfqpoint{1.942806in}{0.483000in}}%
\pgfpathclose%
\pgfusepath{fill}%
\end{pgfscope}%
\begin{pgfscope}%
\pgfpathrectangle{\pgfqpoint{0.595000in}{0.483000in}}{\pgfqpoint{2.835000in}{0.630000in}}%
\pgfusepath{clip}%
\pgfsetbuttcap%
\pgfsetmiterjoin%
\definecolor{currentfill}{rgb}{0.101961,0.101961,0.101961}%
\pgfsetfillcolor{currentfill}%
\pgfsetlinewidth{0.000000pt}%
\definecolor{currentstroke}{rgb}{0.000000,0.000000,0.000000}%
\pgfsetstrokecolor{currentstroke}%
\pgfsetstrokeopacity{0.000000}%
\pgfsetdash{}{0pt}%
\pgfpathmoveto{\pgfqpoint{1.966431in}{0.483000in}}%
\pgfpathlineto{\pgfqpoint{1.987694in}{0.483000in}}%
\pgfpathlineto{\pgfqpoint{1.987694in}{0.483000in}}%
\pgfpathlineto{\pgfqpoint{1.966431in}{0.483000in}}%
\pgfpathlineto{\pgfqpoint{1.966431in}{0.483000in}}%
\pgfpathclose%
\pgfusepath{fill}%
\end{pgfscope}%
\begin{pgfscope}%
\pgfpathrectangle{\pgfqpoint{0.595000in}{0.483000in}}{\pgfqpoint{2.835000in}{0.630000in}}%
\pgfusepath{clip}%
\pgfsetbuttcap%
\pgfsetmiterjoin%
\definecolor{currentfill}{rgb}{0.101961,0.101961,0.101961}%
\pgfsetfillcolor{currentfill}%
\pgfsetlinewidth{0.000000pt}%
\definecolor{currentstroke}{rgb}{0.000000,0.000000,0.000000}%
\pgfsetstrokecolor{currentstroke}%
\pgfsetstrokeopacity{0.000000}%
\pgfsetdash{}{0pt}%
\pgfpathmoveto{\pgfqpoint{1.990056in}{0.483000in}}%
\pgfpathlineto{\pgfqpoint{2.011319in}{0.483000in}}%
\pgfpathlineto{\pgfqpoint{2.011319in}{0.483000in}}%
\pgfpathlineto{\pgfqpoint{1.990056in}{0.483000in}}%
\pgfpathlineto{\pgfqpoint{1.990056in}{0.483000in}}%
\pgfpathclose%
\pgfusepath{fill}%
\end{pgfscope}%
\begin{pgfscope}%
\pgfpathrectangle{\pgfqpoint{0.595000in}{0.483000in}}{\pgfqpoint{2.835000in}{0.630000in}}%
\pgfusepath{clip}%
\pgfsetbuttcap%
\pgfsetmiterjoin%
\definecolor{currentfill}{rgb}{0.101961,0.101961,0.101961}%
\pgfsetfillcolor{currentfill}%
\pgfsetlinewidth{0.000000pt}%
\definecolor{currentstroke}{rgb}{0.000000,0.000000,0.000000}%
\pgfsetstrokecolor{currentstroke}%
\pgfsetstrokeopacity{0.000000}%
\pgfsetdash{}{0pt}%
\pgfpathmoveto{\pgfqpoint{2.013681in}{0.483000in}}%
\pgfpathlineto{\pgfqpoint{2.034944in}{0.483000in}}%
\pgfpathlineto{\pgfqpoint{2.034944in}{0.483000in}}%
\pgfpathlineto{\pgfqpoint{2.013681in}{0.483000in}}%
\pgfpathlineto{\pgfqpoint{2.013681in}{0.483000in}}%
\pgfpathclose%
\pgfusepath{fill}%
\end{pgfscope}%
\begin{pgfscope}%
\pgfpathrectangle{\pgfqpoint{0.595000in}{0.483000in}}{\pgfqpoint{2.835000in}{0.630000in}}%
\pgfusepath{clip}%
\pgfsetbuttcap%
\pgfsetmiterjoin%
\definecolor{currentfill}{rgb}{0.101961,0.101961,0.101961}%
\pgfsetfillcolor{currentfill}%
\pgfsetlinewidth{0.000000pt}%
\definecolor{currentstroke}{rgb}{0.000000,0.000000,0.000000}%
\pgfsetstrokecolor{currentstroke}%
\pgfsetstrokeopacity{0.000000}%
\pgfsetdash{}{0pt}%
\pgfpathmoveto{\pgfqpoint{2.037306in}{0.483000in}}%
\pgfpathlineto{\pgfqpoint{2.058569in}{0.483000in}}%
\pgfpathlineto{\pgfqpoint{2.058569in}{0.573000in}}%
\pgfpathlineto{\pgfqpoint{2.037306in}{0.573000in}}%
\pgfpathlineto{\pgfqpoint{2.037306in}{0.483000in}}%
\pgfpathclose%
\pgfusepath{fill}%
\end{pgfscope}%
\begin{pgfscope}%
\pgfpathrectangle{\pgfqpoint{0.595000in}{0.483000in}}{\pgfqpoint{2.835000in}{0.630000in}}%
\pgfusepath{clip}%
\pgfsetbuttcap%
\pgfsetmiterjoin%
\definecolor{currentfill}{rgb}{0.101961,0.101961,0.101961}%
\pgfsetfillcolor{currentfill}%
\pgfsetlinewidth{0.000000pt}%
\definecolor{currentstroke}{rgb}{0.000000,0.000000,0.000000}%
\pgfsetstrokecolor{currentstroke}%
\pgfsetstrokeopacity{0.000000}%
\pgfsetdash{}{0pt}%
\pgfpathmoveto{\pgfqpoint{2.060931in}{0.483000in}}%
\pgfpathlineto{\pgfqpoint{2.082194in}{0.483000in}}%
\pgfpathlineto{\pgfqpoint{2.082194in}{0.663000in}}%
\pgfpathlineto{\pgfqpoint{2.060931in}{0.663000in}}%
\pgfpathlineto{\pgfqpoint{2.060931in}{0.483000in}}%
\pgfpathclose%
\pgfusepath{fill}%
\end{pgfscope}%
\begin{pgfscope}%
\pgfpathrectangle{\pgfqpoint{0.595000in}{0.483000in}}{\pgfqpoint{2.835000in}{0.630000in}}%
\pgfusepath{clip}%
\pgfsetbuttcap%
\pgfsetmiterjoin%
\definecolor{currentfill}{rgb}{0.101961,0.101961,0.101961}%
\pgfsetfillcolor{currentfill}%
\pgfsetlinewidth{0.000000pt}%
\definecolor{currentstroke}{rgb}{0.000000,0.000000,0.000000}%
\pgfsetstrokecolor{currentstroke}%
\pgfsetstrokeopacity{0.000000}%
\pgfsetdash{}{0pt}%
\pgfpathmoveto{\pgfqpoint{2.084556in}{0.483000in}}%
\pgfpathlineto{\pgfqpoint{2.105819in}{0.483000in}}%
\pgfpathlineto{\pgfqpoint{2.105819in}{0.483000in}}%
\pgfpathlineto{\pgfqpoint{2.084556in}{0.483000in}}%
\pgfpathlineto{\pgfqpoint{2.084556in}{0.483000in}}%
\pgfpathclose%
\pgfusepath{fill}%
\end{pgfscope}%
\begin{pgfscope}%
\pgfpathrectangle{\pgfqpoint{0.595000in}{0.483000in}}{\pgfqpoint{2.835000in}{0.630000in}}%
\pgfusepath{clip}%
\pgfsetbuttcap%
\pgfsetmiterjoin%
\definecolor{currentfill}{rgb}{0.101961,0.101961,0.101961}%
\pgfsetfillcolor{currentfill}%
\pgfsetlinewidth{0.000000pt}%
\definecolor{currentstroke}{rgb}{0.000000,0.000000,0.000000}%
\pgfsetstrokecolor{currentstroke}%
\pgfsetstrokeopacity{0.000000}%
\pgfsetdash{}{0pt}%
\pgfpathmoveto{\pgfqpoint{2.108181in}{0.483000in}}%
\pgfpathlineto{\pgfqpoint{2.129444in}{0.483000in}}%
\pgfpathlineto{\pgfqpoint{2.129444in}{0.483000in}}%
\pgfpathlineto{\pgfqpoint{2.108181in}{0.483000in}}%
\pgfpathlineto{\pgfqpoint{2.108181in}{0.483000in}}%
\pgfpathclose%
\pgfusepath{fill}%
\end{pgfscope}%
\begin{pgfscope}%
\pgfpathrectangle{\pgfqpoint{0.595000in}{0.483000in}}{\pgfqpoint{2.835000in}{0.630000in}}%
\pgfusepath{clip}%
\pgfsetbuttcap%
\pgfsetmiterjoin%
\definecolor{currentfill}{rgb}{0.101961,0.101961,0.101961}%
\pgfsetfillcolor{currentfill}%
\pgfsetlinewidth{0.000000pt}%
\definecolor{currentstroke}{rgb}{0.000000,0.000000,0.000000}%
\pgfsetstrokecolor{currentstroke}%
\pgfsetstrokeopacity{0.000000}%
\pgfsetdash{}{0pt}%
\pgfpathmoveto{\pgfqpoint{2.131806in}{0.483000in}}%
\pgfpathlineto{\pgfqpoint{2.153069in}{0.483000in}}%
\pgfpathlineto{\pgfqpoint{2.153069in}{0.483000in}}%
\pgfpathlineto{\pgfqpoint{2.131806in}{0.483000in}}%
\pgfpathlineto{\pgfqpoint{2.131806in}{0.483000in}}%
\pgfpathclose%
\pgfusepath{fill}%
\end{pgfscope}%
\begin{pgfscope}%
\pgfpathrectangle{\pgfqpoint{0.595000in}{0.483000in}}{\pgfqpoint{2.835000in}{0.630000in}}%
\pgfusepath{clip}%
\pgfsetbuttcap%
\pgfsetmiterjoin%
\definecolor{currentfill}{rgb}{0.101961,0.101961,0.101961}%
\pgfsetfillcolor{currentfill}%
\pgfsetlinewidth{0.000000pt}%
\definecolor{currentstroke}{rgb}{0.000000,0.000000,0.000000}%
\pgfsetstrokecolor{currentstroke}%
\pgfsetstrokeopacity{0.000000}%
\pgfsetdash{}{0pt}%
\pgfpathmoveto{\pgfqpoint{2.155431in}{0.483000in}}%
\pgfpathlineto{\pgfqpoint{2.176694in}{0.483000in}}%
\pgfpathlineto{\pgfqpoint{2.176694in}{0.483000in}}%
\pgfpathlineto{\pgfqpoint{2.155431in}{0.483000in}}%
\pgfpathlineto{\pgfqpoint{2.155431in}{0.483000in}}%
\pgfpathclose%
\pgfusepath{fill}%
\end{pgfscope}%
\begin{pgfscope}%
\pgfpathrectangle{\pgfqpoint{0.595000in}{0.483000in}}{\pgfqpoint{2.835000in}{0.630000in}}%
\pgfusepath{clip}%
\pgfsetbuttcap%
\pgfsetmiterjoin%
\definecolor{currentfill}{rgb}{0.101961,0.101961,0.101961}%
\pgfsetfillcolor{currentfill}%
\pgfsetlinewidth{0.000000pt}%
\definecolor{currentstroke}{rgb}{0.000000,0.000000,0.000000}%
\pgfsetstrokecolor{currentstroke}%
\pgfsetstrokeopacity{0.000000}%
\pgfsetdash{}{0pt}%
\pgfpathmoveto{\pgfqpoint{2.179056in}{0.483000in}}%
\pgfpathlineto{\pgfqpoint{2.200319in}{0.483000in}}%
\pgfpathlineto{\pgfqpoint{2.200319in}{0.483000in}}%
\pgfpathlineto{\pgfqpoint{2.179056in}{0.483000in}}%
\pgfpathlineto{\pgfqpoint{2.179056in}{0.483000in}}%
\pgfpathclose%
\pgfusepath{fill}%
\end{pgfscope}%
\begin{pgfscope}%
\pgfpathrectangle{\pgfqpoint{0.595000in}{0.483000in}}{\pgfqpoint{2.835000in}{0.630000in}}%
\pgfusepath{clip}%
\pgfsetbuttcap%
\pgfsetmiterjoin%
\definecolor{currentfill}{rgb}{0.101961,0.101961,0.101961}%
\pgfsetfillcolor{currentfill}%
\pgfsetlinewidth{0.000000pt}%
\definecolor{currentstroke}{rgb}{0.000000,0.000000,0.000000}%
\pgfsetstrokecolor{currentstroke}%
\pgfsetstrokeopacity{0.000000}%
\pgfsetdash{}{0pt}%
\pgfpathmoveto{\pgfqpoint{2.202681in}{0.483000in}}%
\pgfpathlineto{\pgfqpoint{2.223944in}{0.483000in}}%
\pgfpathlineto{\pgfqpoint{2.223944in}{0.483000in}}%
\pgfpathlineto{\pgfqpoint{2.202681in}{0.483000in}}%
\pgfpathlineto{\pgfqpoint{2.202681in}{0.483000in}}%
\pgfpathclose%
\pgfusepath{fill}%
\end{pgfscope}%
\begin{pgfscope}%
\pgfpathrectangle{\pgfqpoint{0.595000in}{0.483000in}}{\pgfqpoint{2.835000in}{0.630000in}}%
\pgfusepath{clip}%
\pgfsetbuttcap%
\pgfsetmiterjoin%
\definecolor{currentfill}{rgb}{0.101961,0.101961,0.101961}%
\pgfsetfillcolor{currentfill}%
\pgfsetlinewidth{0.000000pt}%
\definecolor{currentstroke}{rgb}{0.000000,0.000000,0.000000}%
\pgfsetstrokecolor{currentstroke}%
\pgfsetstrokeopacity{0.000000}%
\pgfsetdash{}{0pt}%
\pgfpathmoveto{\pgfqpoint{2.226306in}{0.483000in}}%
\pgfpathlineto{\pgfqpoint{2.247569in}{0.483000in}}%
\pgfpathlineto{\pgfqpoint{2.247569in}{0.483000in}}%
\pgfpathlineto{\pgfqpoint{2.226306in}{0.483000in}}%
\pgfpathlineto{\pgfqpoint{2.226306in}{0.483000in}}%
\pgfpathclose%
\pgfusepath{fill}%
\end{pgfscope}%
\begin{pgfscope}%
\pgfpathrectangle{\pgfqpoint{0.595000in}{0.483000in}}{\pgfqpoint{2.835000in}{0.630000in}}%
\pgfusepath{clip}%
\pgfsetbuttcap%
\pgfsetmiterjoin%
\definecolor{currentfill}{rgb}{0.101961,0.101961,0.101961}%
\pgfsetfillcolor{currentfill}%
\pgfsetlinewidth{0.000000pt}%
\definecolor{currentstroke}{rgb}{0.000000,0.000000,0.000000}%
\pgfsetstrokecolor{currentstroke}%
\pgfsetstrokeopacity{0.000000}%
\pgfsetdash{}{0pt}%
\pgfpathmoveto{\pgfqpoint{2.249931in}{0.483000in}}%
\pgfpathlineto{\pgfqpoint{2.271194in}{0.483000in}}%
\pgfpathlineto{\pgfqpoint{2.271194in}{0.483000in}}%
\pgfpathlineto{\pgfqpoint{2.249931in}{0.483000in}}%
\pgfpathlineto{\pgfqpoint{2.249931in}{0.483000in}}%
\pgfpathclose%
\pgfusepath{fill}%
\end{pgfscope}%
\begin{pgfscope}%
\pgfpathrectangle{\pgfqpoint{0.595000in}{0.483000in}}{\pgfqpoint{2.835000in}{0.630000in}}%
\pgfusepath{clip}%
\pgfsetbuttcap%
\pgfsetmiterjoin%
\definecolor{currentfill}{rgb}{0.101961,0.101961,0.101961}%
\pgfsetfillcolor{currentfill}%
\pgfsetlinewidth{0.000000pt}%
\definecolor{currentstroke}{rgb}{0.000000,0.000000,0.000000}%
\pgfsetstrokecolor{currentstroke}%
\pgfsetstrokeopacity{0.000000}%
\pgfsetdash{}{0pt}%
\pgfpathmoveto{\pgfqpoint{2.273556in}{0.483000in}}%
\pgfpathlineto{\pgfqpoint{2.294819in}{0.483000in}}%
\pgfpathlineto{\pgfqpoint{2.294819in}{0.483000in}}%
\pgfpathlineto{\pgfqpoint{2.273556in}{0.483000in}}%
\pgfpathlineto{\pgfqpoint{2.273556in}{0.483000in}}%
\pgfpathclose%
\pgfusepath{fill}%
\end{pgfscope}%
\begin{pgfscope}%
\pgfpathrectangle{\pgfqpoint{0.595000in}{0.483000in}}{\pgfqpoint{2.835000in}{0.630000in}}%
\pgfusepath{clip}%
\pgfsetbuttcap%
\pgfsetmiterjoin%
\definecolor{currentfill}{rgb}{0.101961,0.101961,0.101961}%
\pgfsetfillcolor{currentfill}%
\pgfsetlinewidth{0.000000pt}%
\definecolor{currentstroke}{rgb}{0.000000,0.000000,0.000000}%
\pgfsetstrokecolor{currentstroke}%
\pgfsetstrokeopacity{0.000000}%
\pgfsetdash{}{0pt}%
\pgfpathmoveto{\pgfqpoint{2.297181in}{0.483000in}}%
\pgfpathlineto{\pgfqpoint{2.318444in}{0.483000in}}%
\pgfpathlineto{\pgfqpoint{2.318444in}{0.483000in}}%
\pgfpathlineto{\pgfqpoint{2.297181in}{0.483000in}}%
\pgfpathlineto{\pgfqpoint{2.297181in}{0.483000in}}%
\pgfpathclose%
\pgfusepath{fill}%
\end{pgfscope}%
\begin{pgfscope}%
\pgfpathrectangle{\pgfqpoint{0.595000in}{0.483000in}}{\pgfqpoint{2.835000in}{0.630000in}}%
\pgfusepath{clip}%
\pgfsetbuttcap%
\pgfsetmiterjoin%
\definecolor{currentfill}{rgb}{0.101961,0.101961,0.101961}%
\pgfsetfillcolor{currentfill}%
\pgfsetlinewidth{0.000000pt}%
\definecolor{currentstroke}{rgb}{0.000000,0.000000,0.000000}%
\pgfsetstrokecolor{currentstroke}%
\pgfsetstrokeopacity{0.000000}%
\pgfsetdash{}{0pt}%
\pgfpathmoveto{\pgfqpoint{2.320806in}{0.483000in}}%
\pgfpathlineto{\pgfqpoint{2.342069in}{0.483000in}}%
\pgfpathlineto{\pgfqpoint{2.342069in}{0.663000in}}%
\pgfpathlineto{\pgfqpoint{2.320806in}{0.663000in}}%
\pgfpathlineto{\pgfqpoint{2.320806in}{0.483000in}}%
\pgfpathclose%
\pgfusepath{fill}%
\end{pgfscope}%
\begin{pgfscope}%
\pgfpathrectangle{\pgfqpoint{0.595000in}{0.483000in}}{\pgfqpoint{2.835000in}{0.630000in}}%
\pgfusepath{clip}%
\pgfsetbuttcap%
\pgfsetmiterjoin%
\definecolor{currentfill}{rgb}{0.101961,0.101961,0.101961}%
\pgfsetfillcolor{currentfill}%
\pgfsetlinewidth{0.000000pt}%
\definecolor{currentstroke}{rgb}{0.000000,0.000000,0.000000}%
\pgfsetstrokecolor{currentstroke}%
\pgfsetstrokeopacity{0.000000}%
\pgfsetdash{}{0pt}%
\pgfpathmoveto{\pgfqpoint{2.344431in}{0.483000in}}%
\pgfpathlineto{\pgfqpoint{2.365694in}{0.483000in}}%
\pgfpathlineto{\pgfqpoint{2.365694in}{0.573000in}}%
\pgfpathlineto{\pgfqpoint{2.344431in}{0.573000in}}%
\pgfpathlineto{\pgfqpoint{2.344431in}{0.483000in}}%
\pgfpathclose%
\pgfusepath{fill}%
\end{pgfscope}%
\begin{pgfscope}%
\pgfpathrectangle{\pgfqpoint{0.595000in}{0.483000in}}{\pgfqpoint{2.835000in}{0.630000in}}%
\pgfusepath{clip}%
\pgfsetbuttcap%
\pgfsetmiterjoin%
\definecolor{currentfill}{rgb}{0.101961,0.101961,0.101961}%
\pgfsetfillcolor{currentfill}%
\pgfsetlinewidth{0.000000pt}%
\definecolor{currentstroke}{rgb}{0.000000,0.000000,0.000000}%
\pgfsetstrokecolor{currentstroke}%
\pgfsetstrokeopacity{0.000000}%
\pgfsetdash{}{0pt}%
\pgfpathmoveto{\pgfqpoint{2.368056in}{0.483000in}}%
\pgfpathlineto{\pgfqpoint{2.389319in}{0.483000in}}%
\pgfpathlineto{\pgfqpoint{2.389319in}{0.573000in}}%
\pgfpathlineto{\pgfqpoint{2.368056in}{0.573000in}}%
\pgfpathlineto{\pgfqpoint{2.368056in}{0.483000in}}%
\pgfpathclose%
\pgfusepath{fill}%
\end{pgfscope}%
\begin{pgfscope}%
\pgfpathrectangle{\pgfqpoint{0.595000in}{0.483000in}}{\pgfqpoint{2.835000in}{0.630000in}}%
\pgfusepath{clip}%
\pgfsetbuttcap%
\pgfsetmiterjoin%
\definecolor{currentfill}{rgb}{0.101961,0.101961,0.101961}%
\pgfsetfillcolor{currentfill}%
\pgfsetlinewidth{0.000000pt}%
\definecolor{currentstroke}{rgb}{0.000000,0.000000,0.000000}%
\pgfsetstrokecolor{currentstroke}%
\pgfsetstrokeopacity{0.000000}%
\pgfsetdash{}{0pt}%
\pgfpathmoveto{\pgfqpoint{2.391681in}{0.483000in}}%
\pgfpathlineto{\pgfqpoint{2.412944in}{0.483000in}}%
\pgfpathlineto{\pgfqpoint{2.412944in}{0.528000in}}%
\pgfpathlineto{\pgfqpoint{2.391681in}{0.528000in}}%
\pgfpathlineto{\pgfqpoint{2.391681in}{0.483000in}}%
\pgfpathclose%
\pgfusepath{fill}%
\end{pgfscope}%
\begin{pgfscope}%
\pgfpathrectangle{\pgfqpoint{0.595000in}{0.483000in}}{\pgfqpoint{2.835000in}{0.630000in}}%
\pgfusepath{clip}%
\pgfsetbuttcap%
\pgfsetmiterjoin%
\definecolor{currentfill}{rgb}{0.101961,0.101961,0.101961}%
\pgfsetfillcolor{currentfill}%
\pgfsetlinewidth{0.000000pt}%
\definecolor{currentstroke}{rgb}{0.000000,0.000000,0.000000}%
\pgfsetstrokecolor{currentstroke}%
\pgfsetstrokeopacity{0.000000}%
\pgfsetdash{}{0pt}%
\pgfpathmoveto{\pgfqpoint{2.415306in}{0.483000in}}%
\pgfpathlineto{\pgfqpoint{2.436569in}{0.483000in}}%
\pgfpathlineto{\pgfqpoint{2.436569in}{0.528000in}}%
\pgfpathlineto{\pgfqpoint{2.415306in}{0.528000in}}%
\pgfpathlineto{\pgfqpoint{2.415306in}{0.483000in}}%
\pgfpathclose%
\pgfusepath{fill}%
\end{pgfscope}%
\begin{pgfscope}%
\pgfpathrectangle{\pgfqpoint{0.595000in}{0.483000in}}{\pgfqpoint{2.835000in}{0.630000in}}%
\pgfusepath{clip}%
\pgfsetbuttcap%
\pgfsetmiterjoin%
\definecolor{currentfill}{rgb}{0.101961,0.101961,0.101961}%
\pgfsetfillcolor{currentfill}%
\pgfsetlinewidth{0.000000pt}%
\definecolor{currentstroke}{rgb}{0.000000,0.000000,0.000000}%
\pgfsetstrokecolor{currentstroke}%
\pgfsetstrokeopacity{0.000000}%
\pgfsetdash{}{0pt}%
\pgfpathmoveto{\pgfqpoint{2.438931in}{0.483000in}}%
\pgfpathlineto{\pgfqpoint{2.460194in}{0.483000in}}%
\pgfpathlineto{\pgfqpoint{2.460194in}{1.068000in}}%
\pgfpathlineto{\pgfqpoint{2.438931in}{1.068000in}}%
\pgfpathlineto{\pgfqpoint{2.438931in}{0.483000in}}%
\pgfpathclose%
\pgfusepath{fill}%
\end{pgfscope}%
\begin{pgfscope}%
\pgfpathrectangle{\pgfqpoint{0.595000in}{0.483000in}}{\pgfqpoint{2.835000in}{0.630000in}}%
\pgfusepath{clip}%
\pgfsetbuttcap%
\pgfsetmiterjoin%
\definecolor{currentfill}{rgb}{0.101961,0.101961,0.101961}%
\pgfsetfillcolor{currentfill}%
\pgfsetlinewidth{0.000000pt}%
\definecolor{currentstroke}{rgb}{0.000000,0.000000,0.000000}%
\pgfsetstrokecolor{currentstroke}%
\pgfsetstrokeopacity{0.000000}%
\pgfsetdash{}{0pt}%
\pgfpathmoveto{\pgfqpoint{2.462556in}{0.483000in}}%
\pgfpathlineto{\pgfqpoint{2.483819in}{0.483000in}}%
\pgfpathlineto{\pgfqpoint{2.483819in}{0.528000in}}%
\pgfpathlineto{\pgfqpoint{2.462556in}{0.528000in}}%
\pgfpathlineto{\pgfqpoint{2.462556in}{0.483000in}}%
\pgfpathclose%
\pgfusepath{fill}%
\end{pgfscope}%
\begin{pgfscope}%
\pgfpathrectangle{\pgfqpoint{0.595000in}{0.483000in}}{\pgfqpoint{2.835000in}{0.630000in}}%
\pgfusepath{clip}%
\pgfsetbuttcap%
\pgfsetmiterjoin%
\definecolor{currentfill}{rgb}{0.101961,0.101961,0.101961}%
\pgfsetfillcolor{currentfill}%
\pgfsetlinewidth{0.000000pt}%
\definecolor{currentstroke}{rgb}{0.000000,0.000000,0.000000}%
\pgfsetstrokecolor{currentstroke}%
\pgfsetstrokeopacity{0.000000}%
\pgfsetdash{}{0pt}%
\pgfpathmoveto{\pgfqpoint{2.486181in}{0.483000in}}%
\pgfpathlineto{\pgfqpoint{2.507444in}{0.483000in}}%
\pgfpathlineto{\pgfqpoint{2.507444in}{0.528000in}}%
\pgfpathlineto{\pgfqpoint{2.486181in}{0.528000in}}%
\pgfpathlineto{\pgfqpoint{2.486181in}{0.483000in}}%
\pgfpathclose%
\pgfusepath{fill}%
\end{pgfscope}%
\begin{pgfscope}%
\pgfpathrectangle{\pgfqpoint{0.595000in}{0.483000in}}{\pgfqpoint{2.835000in}{0.630000in}}%
\pgfusepath{clip}%
\pgfsetbuttcap%
\pgfsetmiterjoin%
\definecolor{currentfill}{rgb}{0.101961,0.101961,0.101961}%
\pgfsetfillcolor{currentfill}%
\pgfsetlinewidth{0.000000pt}%
\definecolor{currentstroke}{rgb}{0.000000,0.000000,0.000000}%
\pgfsetstrokecolor{currentstroke}%
\pgfsetstrokeopacity{0.000000}%
\pgfsetdash{}{0pt}%
\pgfpathmoveto{\pgfqpoint{2.509806in}{0.483000in}}%
\pgfpathlineto{\pgfqpoint{2.531069in}{0.483000in}}%
\pgfpathlineto{\pgfqpoint{2.531069in}{0.573000in}}%
\pgfpathlineto{\pgfqpoint{2.509806in}{0.573000in}}%
\pgfpathlineto{\pgfqpoint{2.509806in}{0.483000in}}%
\pgfpathclose%
\pgfusepath{fill}%
\end{pgfscope}%
\begin{pgfscope}%
\pgfpathrectangle{\pgfqpoint{0.595000in}{0.483000in}}{\pgfqpoint{2.835000in}{0.630000in}}%
\pgfusepath{clip}%
\pgfsetbuttcap%
\pgfsetmiterjoin%
\definecolor{currentfill}{rgb}{0.101961,0.101961,0.101961}%
\pgfsetfillcolor{currentfill}%
\pgfsetlinewidth{0.000000pt}%
\definecolor{currentstroke}{rgb}{0.000000,0.000000,0.000000}%
\pgfsetstrokecolor{currentstroke}%
\pgfsetstrokeopacity{0.000000}%
\pgfsetdash{}{0pt}%
\pgfpathmoveto{\pgfqpoint{2.533431in}{0.483000in}}%
\pgfpathlineto{\pgfqpoint{2.554694in}{0.483000in}}%
\pgfpathlineto{\pgfqpoint{2.554694in}{0.528000in}}%
\pgfpathlineto{\pgfqpoint{2.533431in}{0.528000in}}%
\pgfpathlineto{\pgfqpoint{2.533431in}{0.483000in}}%
\pgfpathclose%
\pgfusepath{fill}%
\end{pgfscope}%
\begin{pgfscope}%
\pgfpathrectangle{\pgfqpoint{0.595000in}{0.483000in}}{\pgfqpoint{2.835000in}{0.630000in}}%
\pgfusepath{clip}%
\pgfsetbuttcap%
\pgfsetmiterjoin%
\definecolor{currentfill}{rgb}{0.101961,0.101961,0.101961}%
\pgfsetfillcolor{currentfill}%
\pgfsetlinewidth{0.000000pt}%
\definecolor{currentstroke}{rgb}{0.000000,0.000000,0.000000}%
\pgfsetstrokecolor{currentstroke}%
\pgfsetstrokeopacity{0.000000}%
\pgfsetdash{}{0pt}%
\pgfpathmoveto{\pgfqpoint{2.557056in}{0.483000in}}%
\pgfpathlineto{\pgfqpoint{2.578319in}{0.483000in}}%
\pgfpathlineto{\pgfqpoint{2.578319in}{0.483000in}}%
\pgfpathlineto{\pgfqpoint{2.557056in}{0.483000in}}%
\pgfpathlineto{\pgfqpoint{2.557056in}{0.483000in}}%
\pgfpathclose%
\pgfusepath{fill}%
\end{pgfscope}%
\begin{pgfscope}%
\pgfpathrectangle{\pgfqpoint{0.595000in}{0.483000in}}{\pgfqpoint{2.835000in}{0.630000in}}%
\pgfusepath{clip}%
\pgfsetbuttcap%
\pgfsetmiterjoin%
\definecolor{currentfill}{rgb}{0.101961,0.101961,0.101961}%
\pgfsetfillcolor{currentfill}%
\pgfsetlinewidth{0.000000pt}%
\definecolor{currentstroke}{rgb}{0.000000,0.000000,0.000000}%
\pgfsetstrokecolor{currentstroke}%
\pgfsetstrokeopacity{0.000000}%
\pgfsetdash{}{0pt}%
\pgfpathmoveto{\pgfqpoint{2.580681in}{0.483000in}}%
\pgfpathlineto{\pgfqpoint{2.601944in}{0.483000in}}%
\pgfpathlineto{\pgfqpoint{2.601944in}{0.483000in}}%
\pgfpathlineto{\pgfqpoint{2.580681in}{0.483000in}}%
\pgfpathlineto{\pgfqpoint{2.580681in}{0.483000in}}%
\pgfpathclose%
\pgfusepath{fill}%
\end{pgfscope}%
\begin{pgfscope}%
\pgfpathrectangle{\pgfqpoint{0.595000in}{0.483000in}}{\pgfqpoint{2.835000in}{0.630000in}}%
\pgfusepath{clip}%
\pgfsetbuttcap%
\pgfsetmiterjoin%
\definecolor{currentfill}{rgb}{0.101961,0.101961,0.101961}%
\pgfsetfillcolor{currentfill}%
\pgfsetlinewidth{0.000000pt}%
\definecolor{currentstroke}{rgb}{0.000000,0.000000,0.000000}%
\pgfsetstrokecolor{currentstroke}%
\pgfsetstrokeopacity{0.000000}%
\pgfsetdash{}{0pt}%
\pgfpathmoveto{\pgfqpoint{2.604306in}{0.483000in}}%
\pgfpathlineto{\pgfqpoint{2.625569in}{0.483000in}}%
\pgfpathlineto{\pgfqpoint{2.625569in}{0.573000in}}%
\pgfpathlineto{\pgfqpoint{2.604306in}{0.573000in}}%
\pgfpathlineto{\pgfqpoint{2.604306in}{0.483000in}}%
\pgfpathclose%
\pgfusepath{fill}%
\end{pgfscope}%
\begin{pgfscope}%
\pgfpathrectangle{\pgfqpoint{0.595000in}{0.483000in}}{\pgfqpoint{2.835000in}{0.630000in}}%
\pgfusepath{clip}%
\pgfsetbuttcap%
\pgfsetmiterjoin%
\definecolor{currentfill}{rgb}{0.101961,0.101961,0.101961}%
\pgfsetfillcolor{currentfill}%
\pgfsetlinewidth{0.000000pt}%
\definecolor{currentstroke}{rgb}{0.000000,0.000000,0.000000}%
\pgfsetstrokecolor{currentstroke}%
\pgfsetstrokeopacity{0.000000}%
\pgfsetdash{}{0pt}%
\pgfpathmoveto{\pgfqpoint{2.627931in}{0.483000in}}%
\pgfpathlineto{\pgfqpoint{2.649194in}{0.483000in}}%
\pgfpathlineto{\pgfqpoint{2.649194in}{0.483000in}}%
\pgfpathlineto{\pgfqpoint{2.627931in}{0.483000in}}%
\pgfpathlineto{\pgfqpoint{2.627931in}{0.483000in}}%
\pgfpathclose%
\pgfusepath{fill}%
\end{pgfscope}%
\begin{pgfscope}%
\pgfpathrectangle{\pgfqpoint{0.595000in}{0.483000in}}{\pgfqpoint{2.835000in}{0.630000in}}%
\pgfusepath{clip}%
\pgfsetbuttcap%
\pgfsetmiterjoin%
\definecolor{currentfill}{rgb}{0.101961,0.101961,0.101961}%
\pgfsetfillcolor{currentfill}%
\pgfsetlinewidth{0.000000pt}%
\definecolor{currentstroke}{rgb}{0.000000,0.000000,0.000000}%
\pgfsetstrokecolor{currentstroke}%
\pgfsetstrokeopacity{0.000000}%
\pgfsetdash{}{0pt}%
\pgfpathmoveto{\pgfqpoint{2.651556in}{0.483000in}}%
\pgfpathlineto{\pgfqpoint{2.672819in}{0.483000in}}%
\pgfpathlineto{\pgfqpoint{2.672819in}{1.068000in}}%
\pgfpathlineto{\pgfqpoint{2.651556in}{1.068000in}}%
\pgfpathlineto{\pgfqpoint{2.651556in}{0.483000in}}%
\pgfpathclose%
\pgfusepath{fill}%
\end{pgfscope}%
\begin{pgfscope}%
\pgfpathrectangle{\pgfqpoint{0.595000in}{0.483000in}}{\pgfqpoint{2.835000in}{0.630000in}}%
\pgfusepath{clip}%
\pgfsetbuttcap%
\pgfsetmiterjoin%
\definecolor{currentfill}{rgb}{0.101961,0.101961,0.101961}%
\pgfsetfillcolor{currentfill}%
\pgfsetlinewidth{0.000000pt}%
\definecolor{currentstroke}{rgb}{0.000000,0.000000,0.000000}%
\pgfsetstrokecolor{currentstroke}%
\pgfsetstrokeopacity{0.000000}%
\pgfsetdash{}{0pt}%
\pgfpathmoveto{\pgfqpoint{2.675181in}{0.483000in}}%
\pgfpathlineto{\pgfqpoint{2.696444in}{0.483000in}}%
\pgfpathlineto{\pgfqpoint{2.696444in}{0.573000in}}%
\pgfpathlineto{\pgfqpoint{2.675181in}{0.573000in}}%
\pgfpathlineto{\pgfqpoint{2.675181in}{0.483000in}}%
\pgfpathclose%
\pgfusepath{fill}%
\end{pgfscope}%
\begin{pgfscope}%
\pgfpathrectangle{\pgfqpoint{0.595000in}{0.483000in}}{\pgfqpoint{2.835000in}{0.630000in}}%
\pgfusepath{clip}%
\pgfsetbuttcap%
\pgfsetmiterjoin%
\definecolor{currentfill}{rgb}{0.101961,0.101961,0.101961}%
\pgfsetfillcolor{currentfill}%
\pgfsetlinewidth{0.000000pt}%
\definecolor{currentstroke}{rgb}{0.000000,0.000000,0.000000}%
\pgfsetstrokecolor{currentstroke}%
\pgfsetstrokeopacity{0.000000}%
\pgfsetdash{}{0pt}%
\pgfpathmoveto{\pgfqpoint{2.698806in}{0.483000in}}%
\pgfpathlineto{\pgfqpoint{2.720069in}{0.483000in}}%
\pgfpathlineto{\pgfqpoint{2.720069in}{0.573000in}}%
\pgfpathlineto{\pgfqpoint{2.698806in}{0.573000in}}%
\pgfpathlineto{\pgfqpoint{2.698806in}{0.483000in}}%
\pgfpathclose%
\pgfusepath{fill}%
\end{pgfscope}%
\begin{pgfscope}%
\pgfpathrectangle{\pgfqpoint{0.595000in}{0.483000in}}{\pgfqpoint{2.835000in}{0.630000in}}%
\pgfusepath{clip}%
\pgfsetbuttcap%
\pgfsetmiterjoin%
\definecolor{currentfill}{rgb}{0.101961,0.101961,0.101961}%
\pgfsetfillcolor{currentfill}%
\pgfsetlinewidth{0.000000pt}%
\definecolor{currentstroke}{rgb}{0.000000,0.000000,0.000000}%
\pgfsetstrokecolor{currentstroke}%
\pgfsetstrokeopacity{0.000000}%
\pgfsetdash{}{0pt}%
\pgfpathmoveto{\pgfqpoint{2.722431in}{0.483000in}}%
\pgfpathlineto{\pgfqpoint{2.743694in}{0.483000in}}%
\pgfpathlineto{\pgfqpoint{2.743694in}{0.483000in}}%
\pgfpathlineto{\pgfqpoint{2.722431in}{0.483000in}}%
\pgfpathlineto{\pgfqpoint{2.722431in}{0.483000in}}%
\pgfpathclose%
\pgfusepath{fill}%
\end{pgfscope}%
\begin{pgfscope}%
\pgfpathrectangle{\pgfqpoint{0.595000in}{0.483000in}}{\pgfqpoint{2.835000in}{0.630000in}}%
\pgfusepath{clip}%
\pgfsetbuttcap%
\pgfsetmiterjoin%
\definecolor{currentfill}{rgb}{0.101961,0.101961,0.101961}%
\pgfsetfillcolor{currentfill}%
\pgfsetlinewidth{0.000000pt}%
\definecolor{currentstroke}{rgb}{0.000000,0.000000,0.000000}%
\pgfsetstrokecolor{currentstroke}%
\pgfsetstrokeopacity{0.000000}%
\pgfsetdash{}{0pt}%
\pgfpathmoveto{\pgfqpoint{2.746056in}{0.483000in}}%
\pgfpathlineto{\pgfqpoint{2.767319in}{0.483000in}}%
\pgfpathlineto{\pgfqpoint{2.767319in}{0.528000in}}%
\pgfpathlineto{\pgfqpoint{2.746056in}{0.528000in}}%
\pgfpathlineto{\pgfqpoint{2.746056in}{0.483000in}}%
\pgfpathclose%
\pgfusepath{fill}%
\end{pgfscope}%
\begin{pgfscope}%
\pgfpathrectangle{\pgfqpoint{0.595000in}{0.483000in}}{\pgfqpoint{2.835000in}{0.630000in}}%
\pgfusepath{clip}%
\pgfsetbuttcap%
\pgfsetmiterjoin%
\definecolor{currentfill}{rgb}{0.101961,0.101961,0.101961}%
\pgfsetfillcolor{currentfill}%
\pgfsetlinewidth{0.000000pt}%
\definecolor{currentstroke}{rgb}{0.000000,0.000000,0.000000}%
\pgfsetstrokecolor{currentstroke}%
\pgfsetstrokeopacity{0.000000}%
\pgfsetdash{}{0pt}%
\pgfpathmoveto{\pgfqpoint{2.769681in}{0.483000in}}%
\pgfpathlineto{\pgfqpoint{2.790944in}{0.483000in}}%
\pgfpathlineto{\pgfqpoint{2.790944in}{0.618000in}}%
\pgfpathlineto{\pgfqpoint{2.769681in}{0.618000in}}%
\pgfpathlineto{\pgfqpoint{2.769681in}{0.483000in}}%
\pgfpathclose%
\pgfusepath{fill}%
\end{pgfscope}%
\begin{pgfscope}%
\pgfpathrectangle{\pgfqpoint{0.595000in}{0.483000in}}{\pgfqpoint{2.835000in}{0.630000in}}%
\pgfusepath{clip}%
\pgfsetbuttcap%
\pgfsetmiterjoin%
\definecolor{currentfill}{rgb}{0.101961,0.101961,0.101961}%
\pgfsetfillcolor{currentfill}%
\pgfsetlinewidth{0.000000pt}%
\definecolor{currentstroke}{rgb}{0.000000,0.000000,0.000000}%
\pgfsetstrokecolor{currentstroke}%
\pgfsetstrokeopacity{0.000000}%
\pgfsetdash{}{0pt}%
\pgfpathmoveto{\pgfqpoint{2.793306in}{0.483000in}}%
\pgfpathlineto{\pgfqpoint{2.814569in}{0.483000in}}%
\pgfpathlineto{\pgfqpoint{2.814569in}{0.663000in}}%
\pgfpathlineto{\pgfqpoint{2.793306in}{0.663000in}}%
\pgfpathlineto{\pgfqpoint{2.793306in}{0.483000in}}%
\pgfpathclose%
\pgfusepath{fill}%
\end{pgfscope}%
\begin{pgfscope}%
\pgfpathrectangle{\pgfqpoint{0.595000in}{0.483000in}}{\pgfqpoint{2.835000in}{0.630000in}}%
\pgfusepath{clip}%
\pgfsetbuttcap%
\pgfsetmiterjoin%
\definecolor{currentfill}{rgb}{0.101961,0.101961,0.101961}%
\pgfsetfillcolor{currentfill}%
\pgfsetlinewidth{0.000000pt}%
\definecolor{currentstroke}{rgb}{0.000000,0.000000,0.000000}%
\pgfsetstrokecolor{currentstroke}%
\pgfsetstrokeopacity{0.000000}%
\pgfsetdash{}{0pt}%
\pgfpathmoveto{\pgfqpoint{2.816931in}{0.483000in}}%
\pgfpathlineto{\pgfqpoint{2.838194in}{0.483000in}}%
\pgfpathlineto{\pgfqpoint{2.838194in}{0.618000in}}%
\pgfpathlineto{\pgfqpoint{2.816931in}{0.618000in}}%
\pgfpathlineto{\pgfqpoint{2.816931in}{0.483000in}}%
\pgfpathclose%
\pgfusepath{fill}%
\end{pgfscope}%
\begin{pgfscope}%
\pgfpathrectangle{\pgfqpoint{0.595000in}{0.483000in}}{\pgfqpoint{2.835000in}{0.630000in}}%
\pgfusepath{clip}%
\pgfsetbuttcap%
\pgfsetmiterjoin%
\definecolor{currentfill}{rgb}{0.101961,0.101961,0.101961}%
\pgfsetfillcolor{currentfill}%
\pgfsetlinewidth{0.000000pt}%
\definecolor{currentstroke}{rgb}{0.000000,0.000000,0.000000}%
\pgfsetstrokecolor{currentstroke}%
\pgfsetstrokeopacity{0.000000}%
\pgfsetdash{}{0pt}%
\pgfpathmoveto{\pgfqpoint{2.840556in}{0.483000in}}%
\pgfpathlineto{\pgfqpoint{2.861819in}{0.483000in}}%
\pgfpathlineto{\pgfqpoint{2.861819in}{0.483000in}}%
\pgfpathlineto{\pgfqpoint{2.840556in}{0.483000in}}%
\pgfpathlineto{\pgfqpoint{2.840556in}{0.483000in}}%
\pgfpathclose%
\pgfusepath{fill}%
\end{pgfscope}%
\begin{pgfscope}%
\pgfpathrectangle{\pgfqpoint{0.595000in}{0.483000in}}{\pgfqpoint{2.835000in}{0.630000in}}%
\pgfusepath{clip}%
\pgfsetbuttcap%
\pgfsetmiterjoin%
\definecolor{currentfill}{rgb}{0.101961,0.101961,0.101961}%
\pgfsetfillcolor{currentfill}%
\pgfsetlinewidth{0.000000pt}%
\definecolor{currentstroke}{rgb}{0.000000,0.000000,0.000000}%
\pgfsetstrokecolor{currentstroke}%
\pgfsetstrokeopacity{0.000000}%
\pgfsetdash{}{0pt}%
\pgfpathmoveto{\pgfqpoint{2.864181in}{0.483000in}}%
\pgfpathlineto{\pgfqpoint{2.885444in}{0.483000in}}%
\pgfpathlineto{\pgfqpoint{2.885444in}{0.483000in}}%
\pgfpathlineto{\pgfqpoint{2.864181in}{0.483000in}}%
\pgfpathlineto{\pgfqpoint{2.864181in}{0.483000in}}%
\pgfpathclose%
\pgfusepath{fill}%
\end{pgfscope}%
\begin{pgfscope}%
\pgfpathrectangle{\pgfqpoint{0.595000in}{0.483000in}}{\pgfqpoint{2.835000in}{0.630000in}}%
\pgfusepath{clip}%
\pgfsetbuttcap%
\pgfsetmiterjoin%
\definecolor{currentfill}{rgb}{0.101961,0.101961,0.101961}%
\pgfsetfillcolor{currentfill}%
\pgfsetlinewidth{0.000000pt}%
\definecolor{currentstroke}{rgb}{0.000000,0.000000,0.000000}%
\pgfsetstrokecolor{currentstroke}%
\pgfsetstrokeopacity{0.000000}%
\pgfsetdash{}{0pt}%
\pgfpathmoveto{\pgfqpoint{2.887806in}{0.483000in}}%
\pgfpathlineto{\pgfqpoint{2.909069in}{0.483000in}}%
\pgfpathlineto{\pgfqpoint{2.909069in}{0.618000in}}%
\pgfpathlineto{\pgfqpoint{2.887806in}{0.618000in}}%
\pgfpathlineto{\pgfqpoint{2.887806in}{0.483000in}}%
\pgfpathclose%
\pgfusepath{fill}%
\end{pgfscope}%
\begin{pgfscope}%
\pgfpathrectangle{\pgfqpoint{0.595000in}{0.483000in}}{\pgfqpoint{2.835000in}{0.630000in}}%
\pgfusepath{clip}%
\pgfsetbuttcap%
\pgfsetmiterjoin%
\definecolor{currentfill}{rgb}{0.101961,0.101961,0.101961}%
\pgfsetfillcolor{currentfill}%
\pgfsetlinewidth{0.000000pt}%
\definecolor{currentstroke}{rgb}{0.000000,0.000000,0.000000}%
\pgfsetstrokecolor{currentstroke}%
\pgfsetstrokeopacity{0.000000}%
\pgfsetdash{}{0pt}%
\pgfpathmoveto{\pgfqpoint{2.911431in}{0.483000in}}%
\pgfpathlineto{\pgfqpoint{2.932694in}{0.483000in}}%
\pgfpathlineto{\pgfqpoint{2.932694in}{0.483000in}}%
\pgfpathlineto{\pgfqpoint{2.911431in}{0.483000in}}%
\pgfpathlineto{\pgfqpoint{2.911431in}{0.483000in}}%
\pgfpathclose%
\pgfusepath{fill}%
\end{pgfscope}%
\begin{pgfscope}%
\pgfpathrectangle{\pgfqpoint{0.595000in}{0.483000in}}{\pgfqpoint{2.835000in}{0.630000in}}%
\pgfusepath{clip}%
\pgfsetbuttcap%
\pgfsetmiterjoin%
\definecolor{currentfill}{rgb}{0.101961,0.101961,0.101961}%
\pgfsetfillcolor{currentfill}%
\pgfsetlinewidth{0.000000pt}%
\definecolor{currentstroke}{rgb}{0.000000,0.000000,0.000000}%
\pgfsetstrokecolor{currentstroke}%
\pgfsetstrokeopacity{0.000000}%
\pgfsetdash{}{0pt}%
\pgfpathmoveto{\pgfqpoint{2.935056in}{0.483000in}}%
\pgfpathlineto{\pgfqpoint{2.956319in}{0.483000in}}%
\pgfpathlineto{\pgfqpoint{2.956319in}{0.528000in}}%
\pgfpathlineto{\pgfqpoint{2.935056in}{0.528000in}}%
\pgfpathlineto{\pgfqpoint{2.935056in}{0.483000in}}%
\pgfpathclose%
\pgfusepath{fill}%
\end{pgfscope}%
\begin{pgfscope}%
\pgfpathrectangle{\pgfqpoint{0.595000in}{0.483000in}}{\pgfqpoint{2.835000in}{0.630000in}}%
\pgfusepath{clip}%
\pgfsetbuttcap%
\pgfsetmiterjoin%
\definecolor{currentfill}{rgb}{0.101961,0.101961,0.101961}%
\pgfsetfillcolor{currentfill}%
\pgfsetlinewidth{0.000000pt}%
\definecolor{currentstroke}{rgb}{0.000000,0.000000,0.000000}%
\pgfsetstrokecolor{currentstroke}%
\pgfsetstrokeopacity{0.000000}%
\pgfsetdash{}{0pt}%
\pgfpathmoveto{\pgfqpoint{2.958681in}{0.483000in}}%
\pgfpathlineto{\pgfqpoint{2.979944in}{0.483000in}}%
\pgfpathlineto{\pgfqpoint{2.979944in}{0.573000in}}%
\pgfpathlineto{\pgfqpoint{2.958681in}{0.573000in}}%
\pgfpathlineto{\pgfqpoint{2.958681in}{0.483000in}}%
\pgfpathclose%
\pgfusepath{fill}%
\end{pgfscope}%
\begin{pgfscope}%
\pgfpathrectangle{\pgfqpoint{0.595000in}{0.483000in}}{\pgfqpoint{2.835000in}{0.630000in}}%
\pgfusepath{clip}%
\pgfsetbuttcap%
\pgfsetmiterjoin%
\definecolor{currentfill}{rgb}{0.101961,0.101961,0.101961}%
\pgfsetfillcolor{currentfill}%
\pgfsetlinewidth{0.000000pt}%
\definecolor{currentstroke}{rgb}{0.000000,0.000000,0.000000}%
\pgfsetstrokecolor{currentstroke}%
\pgfsetstrokeopacity{0.000000}%
\pgfsetdash{}{0pt}%
\pgfpathmoveto{\pgfqpoint{2.982306in}{0.483000in}}%
\pgfpathlineto{\pgfqpoint{3.003569in}{0.483000in}}%
\pgfpathlineto{\pgfqpoint{3.003569in}{0.573000in}}%
\pgfpathlineto{\pgfqpoint{2.982306in}{0.573000in}}%
\pgfpathlineto{\pgfqpoint{2.982306in}{0.483000in}}%
\pgfpathclose%
\pgfusepath{fill}%
\end{pgfscope}%
\begin{pgfscope}%
\pgfpathrectangle{\pgfqpoint{0.595000in}{0.483000in}}{\pgfqpoint{2.835000in}{0.630000in}}%
\pgfusepath{clip}%
\pgfsetbuttcap%
\pgfsetmiterjoin%
\definecolor{currentfill}{rgb}{0.101961,0.101961,0.101961}%
\pgfsetfillcolor{currentfill}%
\pgfsetlinewidth{0.000000pt}%
\definecolor{currentstroke}{rgb}{0.000000,0.000000,0.000000}%
\pgfsetstrokecolor{currentstroke}%
\pgfsetstrokeopacity{0.000000}%
\pgfsetdash{}{0pt}%
\pgfpathmoveto{\pgfqpoint{3.005931in}{0.483000in}}%
\pgfpathlineto{\pgfqpoint{3.027194in}{0.483000in}}%
\pgfpathlineto{\pgfqpoint{3.027194in}{0.573000in}}%
\pgfpathlineto{\pgfqpoint{3.005931in}{0.573000in}}%
\pgfpathlineto{\pgfqpoint{3.005931in}{0.483000in}}%
\pgfpathclose%
\pgfusepath{fill}%
\end{pgfscope}%
\begin{pgfscope}%
\pgfpathrectangle{\pgfqpoint{0.595000in}{0.483000in}}{\pgfqpoint{2.835000in}{0.630000in}}%
\pgfusepath{clip}%
\pgfsetbuttcap%
\pgfsetmiterjoin%
\definecolor{currentfill}{rgb}{0.101961,0.101961,0.101961}%
\pgfsetfillcolor{currentfill}%
\pgfsetlinewidth{0.000000pt}%
\definecolor{currentstroke}{rgb}{0.000000,0.000000,0.000000}%
\pgfsetstrokecolor{currentstroke}%
\pgfsetstrokeopacity{0.000000}%
\pgfsetdash{}{0pt}%
\pgfpathmoveto{\pgfqpoint{3.029556in}{0.483000in}}%
\pgfpathlineto{\pgfqpoint{3.050819in}{0.483000in}}%
\pgfpathlineto{\pgfqpoint{3.050819in}{0.753000in}}%
\pgfpathlineto{\pgfqpoint{3.029556in}{0.753000in}}%
\pgfpathlineto{\pgfqpoint{3.029556in}{0.483000in}}%
\pgfpathclose%
\pgfusepath{fill}%
\end{pgfscope}%
\begin{pgfscope}%
\pgfpathrectangle{\pgfqpoint{0.595000in}{0.483000in}}{\pgfqpoint{2.835000in}{0.630000in}}%
\pgfusepath{clip}%
\pgfsetbuttcap%
\pgfsetmiterjoin%
\definecolor{currentfill}{rgb}{0.101961,0.101961,0.101961}%
\pgfsetfillcolor{currentfill}%
\pgfsetlinewidth{0.000000pt}%
\definecolor{currentstroke}{rgb}{0.000000,0.000000,0.000000}%
\pgfsetstrokecolor{currentstroke}%
\pgfsetstrokeopacity{0.000000}%
\pgfsetdash{}{0pt}%
\pgfpathmoveto{\pgfqpoint{3.053181in}{0.483000in}}%
\pgfpathlineto{\pgfqpoint{3.074444in}{0.483000in}}%
\pgfpathlineto{\pgfqpoint{3.074444in}{0.528000in}}%
\pgfpathlineto{\pgfqpoint{3.053181in}{0.528000in}}%
\pgfpathlineto{\pgfqpoint{3.053181in}{0.483000in}}%
\pgfpathclose%
\pgfusepath{fill}%
\end{pgfscope}%
\begin{pgfscope}%
\pgfpathrectangle{\pgfqpoint{0.595000in}{0.483000in}}{\pgfqpoint{2.835000in}{0.630000in}}%
\pgfusepath{clip}%
\pgfsetbuttcap%
\pgfsetmiterjoin%
\definecolor{currentfill}{rgb}{0.101961,0.101961,0.101961}%
\pgfsetfillcolor{currentfill}%
\pgfsetlinewidth{0.000000pt}%
\definecolor{currentstroke}{rgb}{0.000000,0.000000,0.000000}%
\pgfsetstrokecolor{currentstroke}%
\pgfsetstrokeopacity{0.000000}%
\pgfsetdash{}{0pt}%
\pgfpathmoveto{\pgfqpoint{3.076806in}{0.483000in}}%
\pgfpathlineto{\pgfqpoint{3.098069in}{0.483000in}}%
\pgfpathlineto{\pgfqpoint{3.098069in}{0.483000in}}%
\pgfpathlineto{\pgfqpoint{3.076806in}{0.483000in}}%
\pgfpathlineto{\pgfqpoint{3.076806in}{0.483000in}}%
\pgfpathclose%
\pgfusepath{fill}%
\end{pgfscope}%
\begin{pgfscope}%
\pgfpathrectangle{\pgfqpoint{0.595000in}{0.483000in}}{\pgfqpoint{2.835000in}{0.630000in}}%
\pgfusepath{clip}%
\pgfsetbuttcap%
\pgfsetmiterjoin%
\definecolor{currentfill}{rgb}{0.101961,0.101961,0.101961}%
\pgfsetfillcolor{currentfill}%
\pgfsetlinewidth{0.000000pt}%
\definecolor{currentstroke}{rgb}{0.000000,0.000000,0.000000}%
\pgfsetstrokecolor{currentstroke}%
\pgfsetstrokeopacity{0.000000}%
\pgfsetdash{}{0pt}%
\pgfpathmoveto{\pgfqpoint{3.100431in}{0.483000in}}%
\pgfpathlineto{\pgfqpoint{3.121694in}{0.483000in}}%
\pgfpathlineto{\pgfqpoint{3.121694in}{0.573000in}}%
\pgfpathlineto{\pgfqpoint{3.100431in}{0.573000in}}%
\pgfpathlineto{\pgfqpoint{3.100431in}{0.483000in}}%
\pgfpathclose%
\pgfusepath{fill}%
\end{pgfscope}%
\begin{pgfscope}%
\pgfpathrectangle{\pgfqpoint{0.595000in}{0.483000in}}{\pgfqpoint{2.835000in}{0.630000in}}%
\pgfusepath{clip}%
\pgfsetbuttcap%
\pgfsetmiterjoin%
\definecolor{currentfill}{rgb}{0.101961,0.101961,0.101961}%
\pgfsetfillcolor{currentfill}%
\pgfsetlinewidth{0.000000pt}%
\definecolor{currentstroke}{rgb}{0.000000,0.000000,0.000000}%
\pgfsetstrokecolor{currentstroke}%
\pgfsetstrokeopacity{0.000000}%
\pgfsetdash{}{0pt}%
\pgfpathmoveto{\pgfqpoint{3.124056in}{0.483000in}}%
\pgfpathlineto{\pgfqpoint{3.145319in}{0.483000in}}%
\pgfpathlineto{\pgfqpoint{3.145319in}{0.483000in}}%
\pgfpathlineto{\pgfqpoint{3.124056in}{0.483000in}}%
\pgfpathlineto{\pgfqpoint{3.124056in}{0.483000in}}%
\pgfpathclose%
\pgfusepath{fill}%
\end{pgfscope}%
\begin{pgfscope}%
\pgfpathrectangle{\pgfqpoint{0.595000in}{0.483000in}}{\pgfqpoint{2.835000in}{0.630000in}}%
\pgfusepath{clip}%
\pgfsetbuttcap%
\pgfsetmiterjoin%
\definecolor{currentfill}{rgb}{0.101961,0.101961,0.101961}%
\pgfsetfillcolor{currentfill}%
\pgfsetlinewidth{0.000000pt}%
\definecolor{currentstroke}{rgb}{0.000000,0.000000,0.000000}%
\pgfsetstrokecolor{currentstroke}%
\pgfsetstrokeopacity{0.000000}%
\pgfsetdash{}{0pt}%
\pgfpathmoveto{\pgfqpoint{3.147681in}{0.483000in}}%
\pgfpathlineto{\pgfqpoint{3.168944in}{0.483000in}}%
\pgfpathlineto{\pgfqpoint{3.168944in}{0.528000in}}%
\pgfpathlineto{\pgfqpoint{3.147681in}{0.528000in}}%
\pgfpathlineto{\pgfqpoint{3.147681in}{0.483000in}}%
\pgfpathclose%
\pgfusepath{fill}%
\end{pgfscope}%
\begin{pgfscope}%
\pgfpathrectangle{\pgfqpoint{0.595000in}{0.483000in}}{\pgfqpoint{2.835000in}{0.630000in}}%
\pgfusepath{clip}%
\pgfsetbuttcap%
\pgfsetmiterjoin%
\definecolor{currentfill}{rgb}{0.101961,0.101961,0.101961}%
\pgfsetfillcolor{currentfill}%
\pgfsetlinewidth{0.000000pt}%
\definecolor{currentstroke}{rgb}{0.000000,0.000000,0.000000}%
\pgfsetstrokecolor{currentstroke}%
\pgfsetstrokeopacity{0.000000}%
\pgfsetdash{}{0pt}%
\pgfpathmoveto{\pgfqpoint{3.171306in}{0.483000in}}%
\pgfpathlineto{\pgfqpoint{3.192569in}{0.483000in}}%
\pgfpathlineto{\pgfqpoint{3.192569in}{0.753000in}}%
\pgfpathlineto{\pgfqpoint{3.171306in}{0.753000in}}%
\pgfpathlineto{\pgfqpoint{3.171306in}{0.483000in}}%
\pgfpathclose%
\pgfusepath{fill}%
\end{pgfscope}%
\begin{pgfscope}%
\pgfpathrectangle{\pgfqpoint{0.595000in}{0.483000in}}{\pgfqpoint{2.835000in}{0.630000in}}%
\pgfusepath{clip}%
\pgfsetbuttcap%
\pgfsetmiterjoin%
\definecolor{currentfill}{rgb}{0.101961,0.101961,0.101961}%
\pgfsetfillcolor{currentfill}%
\pgfsetlinewidth{0.000000pt}%
\definecolor{currentstroke}{rgb}{0.000000,0.000000,0.000000}%
\pgfsetstrokecolor{currentstroke}%
\pgfsetstrokeopacity{0.000000}%
\pgfsetdash{}{0pt}%
\pgfpathmoveto{\pgfqpoint{3.194931in}{0.483000in}}%
\pgfpathlineto{\pgfqpoint{3.216194in}{0.483000in}}%
\pgfpathlineto{\pgfqpoint{3.216194in}{0.618000in}}%
\pgfpathlineto{\pgfqpoint{3.194931in}{0.618000in}}%
\pgfpathlineto{\pgfqpoint{3.194931in}{0.483000in}}%
\pgfpathclose%
\pgfusepath{fill}%
\end{pgfscope}%
\begin{pgfscope}%
\pgfpathrectangle{\pgfqpoint{0.595000in}{0.483000in}}{\pgfqpoint{2.835000in}{0.630000in}}%
\pgfusepath{clip}%
\pgfsetbuttcap%
\pgfsetmiterjoin%
\definecolor{currentfill}{rgb}{0.101961,0.101961,0.101961}%
\pgfsetfillcolor{currentfill}%
\pgfsetlinewidth{0.000000pt}%
\definecolor{currentstroke}{rgb}{0.000000,0.000000,0.000000}%
\pgfsetstrokecolor{currentstroke}%
\pgfsetstrokeopacity{0.000000}%
\pgfsetdash{}{0pt}%
\pgfpathmoveto{\pgfqpoint{3.218556in}{0.483000in}}%
\pgfpathlineto{\pgfqpoint{3.239819in}{0.483000in}}%
\pgfpathlineto{\pgfqpoint{3.239819in}{0.483000in}}%
\pgfpathlineto{\pgfqpoint{3.218556in}{0.483000in}}%
\pgfpathlineto{\pgfqpoint{3.218556in}{0.483000in}}%
\pgfpathclose%
\pgfusepath{fill}%
\end{pgfscope}%
\begin{pgfscope}%
\pgfpathrectangle{\pgfqpoint{0.595000in}{0.483000in}}{\pgfqpoint{2.835000in}{0.630000in}}%
\pgfusepath{clip}%
\pgfsetbuttcap%
\pgfsetmiterjoin%
\definecolor{currentfill}{rgb}{0.101961,0.101961,0.101961}%
\pgfsetfillcolor{currentfill}%
\pgfsetlinewidth{0.000000pt}%
\definecolor{currentstroke}{rgb}{0.000000,0.000000,0.000000}%
\pgfsetstrokecolor{currentstroke}%
\pgfsetstrokeopacity{0.000000}%
\pgfsetdash{}{0pt}%
\pgfpathmoveto{\pgfqpoint{3.242181in}{0.483000in}}%
\pgfpathlineto{\pgfqpoint{3.263444in}{0.483000in}}%
\pgfpathlineto{\pgfqpoint{3.263444in}{0.528000in}}%
\pgfpathlineto{\pgfqpoint{3.242181in}{0.528000in}}%
\pgfpathlineto{\pgfqpoint{3.242181in}{0.483000in}}%
\pgfpathclose%
\pgfusepath{fill}%
\end{pgfscope}%
\begin{pgfscope}%
\pgfpathrectangle{\pgfqpoint{0.595000in}{0.483000in}}{\pgfqpoint{2.835000in}{0.630000in}}%
\pgfusepath{clip}%
\pgfsetbuttcap%
\pgfsetmiterjoin%
\definecolor{currentfill}{rgb}{0.101961,0.101961,0.101961}%
\pgfsetfillcolor{currentfill}%
\pgfsetlinewidth{0.000000pt}%
\definecolor{currentstroke}{rgb}{0.000000,0.000000,0.000000}%
\pgfsetstrokecolor{currentstroke}%
\pgfsetstrokeopacity{0.000000}%
\pgfsetdash{}{0pt}%
\pgfpathmoveto{\pgfqpoint{3.265806in}{0.483000in}}%
\pgfpathlineto{\pgfqpoint{3.287069in}{0.483000in}}%
\pgfpathlineto{\pgfqpoint{3.287069in}{0.663000in}}%
\pgfpathlineto{\pgfqpoint{3.265806in}{0.663000in}}%
\pgfpathlineto{\pgfqpoint{3.265806in}{0.483000in}}%
\pgfpathclose%
\pgfusepath{fill}%
\end{pgfscope}%
\begin{pgfscope}%
\pgfpathrectangle{\pgfqpoint{0.595000in}{0.483000in}}{\pgfqpoint{2.835000in}{0.630000in}}%
\pgfusepath{clip}%
\pgfsetbuttcap%
\pgfsetmiterjoin%
\definecolor{currentfill}{rgb}{0.101961,0.101961,0.101961}%
\pgfsetfillcolor{currentfill}%
\pgfsetlinewidth{0.000000pt}%
\definecolor{currentstroke}{rgb}{0.000000,0.000000,0.000000}%
\pgfsetstrokecolor{currentstroke}%
\pgfsetstrokeopacity{0.000000}%
\pgfsetdash{}{0pt}%
\pgfpathmoveto{\pgfqpoint{3.289431in}{0.483000in}}%
\pgfpathlineto{\pgfqpoint{3.310694in}{0.483000in}}%
\pgfpathlineto{\pgfqpoint{3.310694in}{0.843000in}}%
\pgfpathlineto{\pgfqpoint{3.289431in}{0.843000in}}%
\pgfpathlineto{\pgfqpoint{3.289431in}{0.483000in}}%
\pgfpathclose%
\pgfusepath{fill}%
\end{pgfscope}%
\begin{pgfscope}%
\pgfpathrectangle{\pgfqpoint{0.595000in}{0.483000in}}{\pgfqpoint{2.835000in}{0.630000in}}%
\pgfusepath{clip}%
\pgfsetbuttcap%
\pgfsetmiterjoin%
\definecolor{currentfill}{rgb}{0.101961,0.101961,0.101961}%
\pgfsetfillcolor{currentfill}%
\pgfsetlinewidth{0.000000pt}%
\definecolor{currentstroke}{rgb}{0.000000,0.000000,0.000000}%
\pgfsetstrokecolor{currentstroke}%
\pgfsetstrokeopacity{0.000000}%
\pgfsetdash{}{0pt}%
\pgfpathmoveto{\pgfqpoint{3.313056in}{0.483000in}}%
\pgfpathlineto{\pgfqpoint{3.334319in}{0.483000in}}%
\pgfpathlineto{\pgfqpoint{3.334319in}{0.618000in}}%
\pgfpathlineto{\pgfqpoint{3.313056in}{0.618000in}}%
\pgfpathlineto{\pgfqpoint{3.313056in}{0.483000in}}%
\pgfpathclose%
\pgfusepath{fill}%
\end{pgfscope}%
\begin{pgfscope}%
\pgfpathrectangle{\pgfqpoint{0.595000in}{0.483000in}}{\pgfqpoint{2.835000in}{0.630000in}}%
\pgfusepath{clip}%
\pgfsetbuttcap%
\pgfsetmiterjoin%
\definecolor{currentfill}{rgb}{0.101961,0.101961,0.101961}%
\pgfsetfillcolor{currentfill}%
\pgfsetlinewidth{0.000000pt}%
\definecolor{currentstroke}{rgb}{0.000000,0.000000,0.000000}%
\pgfsetstrokecolor{currentstroke}%
\pgfsetstrokeopacity{0.000000}%
\pgfsetdash{}{0pt}%
\pgfpathmoveto{\pgfqpoint{3.336681in}{0.483000in}}%
\pgfpathlineto{\pgfqpoint{3.357944in}{0.483000in}}%
\pgfpathlineto{\pgfqpoint{3.357944in}{0.483000in}}%
\pgfpathlineto{\pgfqpoint{3.336681in}{0.483000in}}%
\pgfpathlineto{\pgfqpoint{3.336681in}{0.483000in}}%
\pgfpathclose%
\pgfusepath{fill}%
\end{pgfscope}%
\begin{pgfscope}%
\pgfpathrectangle{\pgfqpoint{0.595000in}{0.483000in}}{\pgfqpoint{2.835000in}{0.630000in}}%
\pgfusepath{clip}%
\pgfsetbuttcap%
\pgfsetmiterjoin%
\definecolor{currentfill}{rgb}{0.101961,0.101961,0.101961}%
\pgfsetfillcolor{currentfill}%
\pgfsetlinewidth{0.000000pt}%
\definecolor{currentstroke}{rgb}{0.000000,0.000000,0.000000}%
\pgfsetstrokecolor{currentstroke}%
\pgfsetstrokeopacity{0.000000}%
\pgfsetdash{}{0pt}%
\pgfpathmoveto{\pgfqpoint{3.360306in}{0.483000in}}%
\pgfpathlineto{\pgfqpoint{3.381569in}{0.483000in}}%
\pgfpathlineto{\pgfqpoint{3.381569in}{0.483000in}}%
\pgfpathlineto{\pgfqpoint{3.360306in}{0.483000in}}%
\pgfpathlineto{\pgfqpoint{3.360306in}{0.483000in}}%
\pgfpathclose%
\pgfusepath{fill}%
\end{pgfscope}%
\begin{pgfscope}%
\pgfpathrectangle{\pgfqpoint{0.595000in}{0.483000in}}{\pgfqpoint{2.835000in}{0.630000in}}%
\pgfusepath{clip}%
\pgfsetbuttcap%
\pgfsetmiterjoin%
\definecolor{currentfill}{rgb}{0.101961,0.101961,0.101961}%
\pgfsetfillcolor{currentfill}%
\pgfsetlinewidth{0.000000pt}%
\definecolor{currentstroke}{rgb}{0.000000,0.000000,0.000000}%
\pgfsetstrokecolor{currentstroke}%
\pgfsetstrokeopacity{0.000000}%
\pgfsetdash{}{0pt}%
\pgfpathmoveto{\pgfqpoint{3.383931in}{0.483000in}}%
\pgfpathlineto{\pgfqpoint{3.405194in}{0.483000in}}%
\pgfpathlineto{\pgfqpoint{3.405194in}{0.528000in}}%
\pgfpathlineto{\pgfqpoint{3.383931in}{0.528000in}}%
\pgfpathlineto{\pgfqpoint{3.383931in}{0.483000in}}%
\pgfpathclose%
\pgfusepath{fill}%
\end{pgfscope}%
\begin{pgfscope}%
\pgfpathrectangle{\pgfqpoint{0.595000in}{0.483000in}}{\pgfqpoint{2.835000in}{0.630000in}}%
\pgfusepath{clip}%
\pgfsetbuttcap%
\pgfsetmiterjoin%
\definecolor{currentfill}{rgb}{0.101961,0.101961,0.101961}%
\pgfsetfillcolor{currentfill}%
\pgfsetlinewidth{0.000000pt}%
\definecolor{currentstroke}{rgb}{0.000000,0.000000,0.000000}%
\pgfsetstrokecolor{currentstroke}%
\pgfsetstrokeopacity{0.000000}%
\pgfsetdash{}{0pt}%
\pgfpathmoveto{\pgfqpoint{3.407556in}{0.483000in}}%
\pgfpathlineto{\pgfqpoint{3.428819in}{0.483000in}}%
\pgfpathlineto{\pgfqpoint{3.428819in}{0.573000in}}%
\pgfpathlineto{\pgfqpoint{3.407556in}{0.573000in}}%
\pgfpathlineto{\pgfqpoint{3.407556in}{0.483000in}}%
\pgfpathclose%
\pgfusepath{fill}%
\end{pgfscope}%
\begin{pgfscope}%
\pgfsetrectcap%
\pgfsetmiterjoin%
\pgfsetlinewidth{0.803000pt}%
\definecolor{currentstroke}{rgb}{0.101961,0.101961,0.101961}%
\pgfsetstrokecolor{currentstroke}%
\pgfsetdash{}{0pt}%
\pgfpathmoveto{\pgfqpoint{0.595000in}{0.483000in}}%
\pgfpathlineto{\pgfqpoint{0.595000in}{1.113000in}}%
\pgfusepath{stroke}%
\end{pgfscope}%
\begin{pgfscope}%
\pgfsetrectcap%
\pgfsetmiterjoin%
\pgfsetlinewidth{0.803000pt}%
\definecolor{currentstroke}{rgb}{0.101961,0.101961,0.101961}%
\pgfsetstrokecolor{currentstroke}%
\pgfsetdash{}{0pt}%
\pgfpathmoveto{\pgfqpoint{0.595000in}{0.483000in}}%
\pgfpathlineto{\pgfqpoint{3.430000in}{0.483000in}}%
\pgfusepath{stroke}%
\end{pgfscope}%
\begin{pgfscope}%
\pgfsetbuttcap%
\pgfsetmiterjoin%
\definecolor{currentfill}{rgb}{1.000000,1.000000,1.000000}%
\pgfsetfillcolor{currentfill}%
\pgfsetlinewidth{0.000000pt}%
\definecolor{currentstroke}{rgb}{0.000000,0.000000,0.000000}%
\pgfsetstrokecolor{currentstroke}%
\pgfsetstrokeopacity{0.000000}%
\pgfsetdash{}{0pt}%
\pgfpathmoveto{\pgfqpoint{3.815000in}{3.192000in}}%
\pgfpathlineto{\pgfqpoint{6.650000in}{3.192000in}}%
\pgfpathlineto{\pgfqpoint{6.650000in}{3.822000in}}%
\pgfpathlineto{\pgfqpoint{3.815000in}{3.822000in}}%
\pgfpathlineto{\pgfqpoint{3.815000in}{3.192000in}}%
\pgfpathclose%
\pgfusepath{fill}%
\end{pgfscope}%
\begin{pgfscope}%
\pgfpathrectangle{\pgfqpoint{3.815000in}{3.192000in}}{\pgfqpoint{2.835000in}{0.630000in}}%
\pgfusepath{clip}%
\pgfsetbuttcap%
\pgfsetmiterjoin%
\definecolor{currentfill}{rgb}{0.941176,0.894118,0.894118}%
\pgfsetfillcolor{currentfill}%
\pgfsetlinewidth{0.000000pt}%
\definecolor{currentstroke}{rgb}{0.941176,0.894118,0.894118}%
\pgfsetstrokecolor{currentstroke}%
\pgfsetdash{}{0pt}%
\pgfpathmoveto{\pgfqpoint{4.777719in}{3.192000in}}%
\pgfpathlineto{\pgfqpoint{4.777719in}{3.822000in}}%
\pgfpathlineto{\pgfqpoint{5.221869in}{3.822000in}}%
\pgfpathlineto{\pgfqpoint{5.221869in}{3.192000in}}%
\pgfpathlineto{\pgfqpoint{4.777719in}{3.192000in}}%
\pgfpathclose%
\pgfusepath{fill}%
\end{pgfscope}%
\begin{pgfscope}%
\pgfpathrectangle{\pgfqpoint{3.815000in}{3.192000in}}{\pgfqpoint{2.835000in}{0.630000in}}%
\pgfusepath{clip}%
\pgfsetrectcap%
\pgfsetroundjoin%
\pgfsetlinewidth{0.401500pt}%
\definecolor{currentstroke}{rgb}{0.901961,0.901961,0.901961}%
\pgfsetstrokecolor{currentstroke}%
\pgfsetdash{}{0pt}%
\pgfpathmoveto{\pgfqpoint{3.815000in}{3.192000in}}%
\pgfpathlineto{\pgfqpoint{6.650000in}{3.192000in}}%
\pgfusepath{stroke}%
\end{pgfscope}%
\begin{pgfscope}%
\pgfpathrectangle{\pgfqpoint{3.815000in}{3.192000in}}{\pgfqpoint{2.835000in}{0.630000in}}%
\pgfusepath{clip}%
\pgfsetrectcap%
\pgfsetroundjoin%
\pgfsetlinewidth{0.401500pt}%
\definecolor{currentstroke}{rgb}{0.901961,0.901961,0.901961}%
\pgfsetstrokecolor{currentstroke}%
\pgfsetdash{}{0pt}%
\pgfpathmoveto{\pgfqpoint{3.815000in}{3.402000in}}%
\pgfpathlineto{\pgfqpoint{6.650000in}{3.402000in}}%
\pgfusepath{stroke}%
\end{pgfscope}%
\begin{pgfscope}%
\pgfpathrectangle{\pgfqpoint{3.815000in}{3.192000in}}{\pgfqpoint{2.835000in}{0.630000in}}%
\pgfusepath{clip}%
\pgfsetrectcap%
\pgfsetroundjoin%
\pgfsetlinewidth{0.401500pt}%
\definecolor{currentstroke}{rgb}{0.901961,0.901961,0.901961}%
\pgfsetstrokecolor{currentstroke}%
\pgfsetdash{}{0pt}%
\pgfpathmoveto{\pgfqpoint{3.815000in}{3.612000in}}%
\pgfpathlineto{\pgfqpoint{6.650000in}{3.612000in}}%
\pgfusepath{stroke}%
\end{pgfscope}%
\begin{pgfscope}%
\pgfpathrectangle{\pgfqpoint{3.815000in}{3.192000in}}{\pgfqpoint{2.835000in}{0.630000in}}%
\pgfusepath{clip}%
\pgfsetrectcap%
\pgfsetroundjoin%
\pgfsetlinewidth{0.401500pt}%
\definecolor{currentstroke}{rgb}{0.901961,0.901961,0.901961}%
\pgfsetstrokecolor{currentstroke}%
\pgfsetdash{}{0pt}%
\pgfpathmoveto{\pgfqpoint{3.815000in}{3.822000in}}%
\pgfpathlineto{\pgfqpoint{6.650000in}{3.822000in}}%
\pgfusepath{stroke}%
\end{pgfscope}%
\begin{pgfscope}%
\pgfpathrectangle{\pgfqpoint{3.815000in}{3.192000in}}{\pgfqpoint{2.835000in}{0.630000in}}%
\pgfusepath{clip}%
\pgfsetbuttcap%
\pgfsetroundjoin%
\pgfsetlinewidth{0.702625pt}%
\definecolor{currentstroke}{rgb}{0.690196,0.250980,0.250980}%
\pgfsetstrokecolor{currentstroke}%
\pgfsetdash{{2.590000pt}{1.120000pt}}{0.000000pt}%
\pgfpathmoveto{\pgfqpoint{5.221869in}{3.192000in}}%
\pgfpathlineto{\pgfqpoint{5.221869in}{3.822000in}}%
\pgfusepath{stroke}%
\end{pgfscope}%
\begin{pgfscope}%
\pgfpathrectangle{\pgfqpoint{3.815000in}{3.192000in}}{\pgfqpoint{2.835000in}{0.630000in}}%
\pgfusepath{clip}%
\pgfsetrectcap%
\pgfsetroundjoin%
\pgfsetlinewidth{0.903375pt}%
\definecolor{currentstroke}{rgb}{0.156863,0.356863,0.568627}%
\pgfsetstrokecolor{currentstroke}%
\pgfsetdash{}{0pt}%
\pgfpathmoveto{\pgfqpoint{3.817200in}{3.192000in}}%
\pgfpathlineto{\pgfqpoint{3.841363in}{3.192000in}}%
\pgfpathlineto{\pgfqpoint{3.865473in}{3.192000in}}%
\pgfpathlineto{\pgfqpoint{3.889596in}{3.319719in}}%
\pgfpathlineto{\pgfqpoint{3.913963in}{3.192000in}}%
\pgfpathlineto{\pgfqpoint{3.938209in}{3.192000in}}%
\pgfpathlineto{\pgfqpoint{3.962359in}{3.192000in}}%
\pgfpathlineto{\pgfqpoint{3.986463in}{3.192000in}}%
\pgfpathlineto{\pgfqpoint{4.010620in}{3.192000in}}%
\pgfpathlineto{\pgfqpoint{4.034845in}{3.264189in}}%
\pgfpathlineto{\pgfqpoint{4.059053in}{3.192000in}}%
\pgfpathlineto{\pgfqpoint{4.083562in}{3.246889in}}%
\pgfpathlineto{\pgfqpoint{4.107797in}{3.256500in}}%
\pgfpathlineto{\pgfqpoint{4.131977in}{3.242404in}}%
\pgfpathlineto{\pgfqpoint{4.156169in}{3.342999in}}%
\pgfpathlineto{\pgfqpoint{4.180351in}{3.228735in}}%
\pgfpathlineto{\pgfqpoint{4.204456in}{3.192000in}}%
\pgfpathlineto{\pgfqpoint{4.228671in}{3.347484in}}%
\pgfpathlineto{\pgfqpoint{4.253088in}{3.192000in}}%
\pgfpathlineto{\pgfqpoint{4.277230in}{3.356455in}}%
\pgfpathlineto{\pgfqpoint{4.301399in}{3.420528in}}%
\pgfpathlineto{\pgfqpoint{4.325572in}{3.326554in}}%
\pgfpathlineto{\pgfqpoint{4.349747in}{3.192000in}}%
\pgfpathlineto{\pgfqpoint{4.373944in}{3.242191in}}%
\pgfpathlineto{\pgfqpoint{4.398129in}{3.433129in}}%
\pgfpathlineto{\pgfqpoint{4.422600in}{3.192000in}}%
\pgfpathlineto{\pgfqpoint{4.446740in}{3.192000in}}%
\pgfpathlineto{\pgfqpoint{4.470903in}{3.454700in}}%
\pgfpathlineto{\pgfqpoint{4.495058in}{3.412198in}}%
\pgfpathlineto{\pgfqpoint{4.519188in}{3.376104in}}%
\pgfpathlineto{\pgfqpoint{4.543430in}{3.410703in}}%
\pgfpathlineto{\pgfqpoint{4.567643in}{3.377172in}}%
\pgfpathlineto{\pgfqpoint{4.591816in}{3.192000in}}%
\pgfpathlineto{\pgfqpoint{4.615937in}{3.273800in}}%
\pgfpathlineto{\pgfqpoint{4.640082in}{3.422877in}}%
\pgfpathlineto{\pgfqpoint{4.664279in}{3.269315in}}%
\pgfpathlineto{\pgfqpoint{4.688513in}{3.273800in}}%
\pgfpathlineto{\pgfqpoint{4.712769in}{3.384006in}}%
\pgfpathlineto{\pgfqpoint{4.736909in}{3.370337in}}%
\pgfpathlineto{\pgfqpoint{4.761437in}{3.347271in}}%
\pgfpathlineto{\pgfqpoint{4.785643in}{3.282984in}}%
\pgfpathlineto{\pgfqpoint{4.809813in}{3.256500in}}%
\pgfpathlineto{\pgfqpoint{4.833968in}{3.192000in}}%
\pgfpathlineto{\pgfqpoint{4.858212in}{3.287469in}}%
\pgfpathlineto{\pgfqpoint{4.882444in}{3.192000in}}%
\pgfpathlineto{\pgfqpoint{4.906707in}{3.282984in}}%
\pgfpathlineto{\pgfqpoint{4.931099in}{3.210795in}}%
\pgfpathlineto{\pgfqpoint{4.955351in}{3.192000in}}%
\pgfpathlineto{\pgfqpoint{4.979568in}{3.192000in}}%
\pgfpathlineto{\pgfqpoint{5.003855in}{3.306264in}}%
\pgfpathlineto{\pgfqpoint{5.028052in}{3.192000in}}%
\pgfpathlineto{\pgfqpoint{5.052161in}{3.192000in}}%
\pgfpathlineto{\pgfqpoint{5.076372in}{3.192000in}}%
\pgfpathlineto{\pgfqpoint{5.100673in}{3.192000in}}%
\pgfpathlineto{\pgfqpoint{5.124907in}{3.237706in}}%
\pgfpathlineto{\pgfqpoint{5.149090in}{3.279140in}}%
\pgfpathlineto{\pgfqpoint{5.173400in}{3.247530in}}%
\pgfpathlineto{\pgfqpoint{5.197521in}{3.192000in}}%
\pgfpathlineto{\pgfqpoint{5.221652in}{3.192000in}}%
\pgfpathlineto{\pgfqpoint{5.245763in}{3.466447in}}%
\pgfpathlineto{\pgfqpoint{5.269946in}{3.192000in}}%
\pgfpathlineto{\pgfqpoint{5.294086in}{3.355387in}}%
\pgfpathlineto{\pgfqpoint{5.318342in}{3.707789in}}%
\pgfpathlineto{\pgfqpoint{5.342607in}{3.256500in}}%
\pgfpathlineto{\pgfqpoint{5.366900in}{3.352824in}}%
\pgfpathlineto{\pgfqpoint{5.391147in}{3.349193in}}%
\pgfpathlineto{\pgfqpoint{5.415450in}{3.339368in}}%
\pgfpathlineto{\pgfqpoint{5.439668in}{3.242191in}}%
\pgfpathlineto{\pgfqpoint{5.463909in}{3.192000in}}%
\pgfpathlineto{\pgfqpoint{5.488042in}{3.192000in}}%
\pgfpathlineto{\pgfqpoint{5.512218in}{3.192000in}}%
\pgfpathlineto{\pgfqpoint{5.536370in}{3.270810in}}%
\pgfpathlineto{\pgfqpoint{5.560550in}{3.434838in}}%
\pgfpathlineto{\pgfqpoint{5.584664in}{3.301352in}}%
\pgfpathlineto{\pgfqpoint{5.608858in}{3.448506in}}%
\pgfpathlineto{\pgfqpoint{5.632975in}{3.192000in}}%
\pgfpathlineto{\pgfqpoint{5.657110in}{3.301779in}}%
\pgfpathlineto{\pgfqpoint{5.681285in}{3.260345in}}%
\pgfpathlineto{\pgfqpoint{5.705492in}{3.322282in}}%
\pgfpathlineto{\pgfqpoint{5.729627in}{3.192000in}}%
\pgfpathlineto{\pgfqpoint{5.753760in}{3.329330in}}%
\pgfpathlineto{\pgfqpoint{5.778025in}{3.410703in}}%
\pgfpathlineto{\pgfqpoint{5.802148in}{3.192000in}}%
\pgfpathlineto{\pgfqpoint{5.826317in}{3.192000in}}%
\pgfpathlineto{\pgfqpoint{5.850452in}{3.351969in}}%
\pgfpathlineto{\pgfqpoint{5.874630in}{3.192000in}}%
\pgfpathlineto{\pgfqpoint{5.898855in}{3.388064in}}%
\pgfpathlineto{\pgfqpoint{5.922988in}{3.270810in}}%
\pgfpathlineto{\pgfqpoint{5.947378in}{3.271878in}}%
\pgfpathlineto{\pgfqpoint{5.971608in}{3.452137in}}%
\pgfpathlineto{\pgfqpoint{5.995781in}{3.324632in}}%
\pgfpathlineto{\pgfqpoint{6.019912in}{3.333602in}}%
\pgfpathlineto{\pgfqpoint{6.044059in}{3.320146in}}%
\pgfpathlineto{\pgfqpoint{6.068275in}{3.448293in}}%
\pgfpathlineto{\pgfqpoint{6.092417in}{3.736622in}}%
\pgfpathlineto{\pgfqpoint{6.116581in}{3.233861in}}%
\pgfpathlineto{\pgfqpoint{6.140737in}{3.616165in}}%
\pgfpathlineto{\pgfqpoint{6.164873in}{3.192000in}}%
\pgfpathlineto{\pgfqpoint{6.189362in}{3.661016in}}%
\pgfpathlineto{\pgfqpoint{6.213516in}{3.413053in}}%
\pgfpathlineto{\pgfqpoint{6.237746in}{3.442953in}}%
\pgfpathlineto{\pgfqpoint{6.261893in}{3.352397in}}%
\pgfpathlineto{\pgfqpoint{6.286048in}{3.761611in}}%
\pgfpathlineto{\pgfqpoint{6.310206in}{3.745806in}}%
\pgfpathlineto{\pgfqpoint{6.334434in}{3.192000in}}%
\pgfpathlineto{\pgfqpoint{6.358607in}{3.410703in}}%
\pgfpathlineto{\pgfqpoint{6.382804in}{3.372900in}}%
\pgfpathlineto{\pgfqpoint{6.406925in}{3.192000in}}%
\pgfpathlineto{\pgfqpoint{6.431105in}{3.385288in}}%
\pgfpathlineto{\pgfqpoint{6.455295in}{3.192000in}}%
\pgfpathlineto{\pgfqpoint{6.479444in}{3.192000in}}%
\pgfpathlineto{\pgfqpoint{6.503669in}{3.397675in}}%
\pgfpathlineto{\pgfqpoint{6.527835in}{3.410063in}}%
\pgfpathlineto{\pgfqpoint{6.551954in}{3.246889in}}%
\pgfpathlineto{\pgfqpoint{6.576049in}{3.596302in}}%
\pgfpathlineto{\pgfqpoint{6.600326in}{3.438468in}}%
\pgfpathlineto{\pgfqpoint{6.624766in}{3.346844in}}%
\pgfpathlineto{\pgfqpoint{6.649005in}{3.579002in}}%
\pgfusepath{stroke}%
\end{pgfscope}%
\begin{pgfscope}%
\pgfpathrectangle{\pgfqpoint{3.815000in}{3.192000in}}{\pgfqpoint{2.835000in}{0.630000in}}%
\pgfusepath{clip}%
\pgfsetrectcap%
\pgfsetroundjoin%
\pgfsetlinewidth{0.903375pt}%
\definecolor{currentstroke}{rgb}{0.478431,0.478431,0.478431}%
\pgfsetstrokecolor{currentstroke}%
\pgfsetdash{}{0pt}%
\pgfpathmoveto{\pgfqpoint{3.841363in}{3.192000in}}%
\pgfpathlineto{\pgfqpoint{3.865473in}{3.192000in}}%
\pgfpathlineto{\pgfqpoint{3.889596in}{3.192000in}}%
\pgfpathlineto{\pgfqpoint{3.913963in}{3.192000in}}%
\pgfpathlineto{\pgfqpoint{3.938209in}{3.192000in}}%
\pgfpathlineto{\pgfqpoint{3.962359in}{3.192000in}}%
\pgfpathlineto{\pgfqpoint{3.986463in}{3.192000in}}%
\pgfpathlineto{\pgfqpoint{4.010620in}{3.192000in}}%
\pgfpathlineto{\pgfqpoint{4.034845in}{3.222139in}}%
\pgfpathlineto{\pgfqpoint{4.059053in}{3.192000in}}%
\pgfpathlineto{\pgfqpoint{4.083562in}{3.192000in}}%
\pgfpathlineto{\pgfqpoint{4.107797in}{3.192000in}}%
\pgfpathlineto{\pgfqpoint{4.131977in}{3.192000in}}%
\pgfpathlineto{\pgfqpoint{4.156169in}{3.192000in}}%
\pgfpathlineto{\pgfqpoint{4.180351in}{3.192000in}}%
\pgfpathlineto{\pgfqpoint{4.204456in}{3.192000in}}%
\pgfpathlineto{\pgfqpoint{4.228671in}{3.192000in}}%
\pgfpathlineto{\pgfqpoint{4.253088in}{3.223840in}}%
\pgfpathlineto{\pgfqpoint{4.277230in}{3.192000in}}%
\pgfpathlineto{\pgfqpoint{4.301399in}{3.192000in}}%
\pgfpathlineto{\pgfqpoint{4.325572in}{3.221288in}}%
\pgfpathlineto{\pgfqpoint{4.349747in}{3.192000in}}%
\pgfpathlineto{\pgfqpoint{4.373944in}{3.192000in}}%
\pgfpathlineto{\pgfqpoint{4.398129in}{3.192000in}}%
\pgfpathlineto{\pgfqpoint{4.422600in}{3.192000in}}%
\pgfpathlineto{\pgfqpoint{4.446740in}{3.192000in}}%
\pgfpathlineto{\pgfqpoint{4.470903in}{3.281323in}}%
\pgfpathlineto{\pgfqpoint{4.495058in}{3.310125in}}%
\pgfpathlineto{\pgfqpoint{4.519188in}{3.251549in}}%
\pgfpathlineto{\pgfqpoint{4.543430in}{3.192000in}}%
\pgfpathlineto{\pgfqpoint{4.567643in}{3.192000in}}%
\pgfpathlineto{\pgfqpoint{4.591816in}{3.192000in}}%
\pgfpathlineto{\pgfqpoint{4.615937in}{3.192000in}}%
\pgfpathlineto{\pgfqpoint{4.640082in}{3.297243in}}%
\pgfpathlineto{\pgfqpoint{4.664279in}{3.192000in}}%
\pgfpathlineto{\pgfqpoint{4.688513in}{3.192000in}}%
\pgfpathlineto{\pgfqpoint{4.712769in}{3.192000in}}%
\pgfpathlineto{\pgfqpoint{4.736909in}{3.192000in}}%
\pgfpathlineto{\pgfqpoint{4.761437in}{3.192000in}}%
\pgfpathlineto{\pgfqpoint{4.785643in}{3.192000in}}%
\pgfpathlineto{\pgfqpoint{4.809813in}{3.192000in}}%
\pgfpathlineto{\pgfqpoint{4.833968in}{3.192000in}}%
\pgfpathlineto{\pgfqpoint{4.858212in}{3.192000in}}%
\pgfpathlineto{\pgfqpoint{4.882444in}{3.192000in}}%
\pgfpathlineto{\pgfqpoint{4.906707in}{3.192000in}}%
\pgfpathlineto{\pgfqpoint{4.931099in}{3.192000in}}%
\pgfpathlineto{\pgfqpoint{4.955351in}{3.192000in}}%
\pgfpathlineto{\pgfqpoint{4.979568in}{3.192000in}}%
\pgfpathlineto{\pgfqpoint{5.003855in}{3.192000in}}%
\pgfpathlineto{\pgfqpoint{5.028052in}{3.192000in}}%
\pgfpathlineto{\pgfqpoint{5.052161in}{3.192000in}}%
\pgfpathlineto{\pgfqpoint{5.076372in}{3.192000in}}%
\pgfpathlineto{\pgfqpoint{5.100673in}{3.192000in}}%
\pgfpathlineto{\pgfqpoint{5.124907in}{3.192000in}}%
\pgfpathlineto{\pgfqpoint{5.149090in}{3.192000in}}%
\pgfpathlineto{\pgfqpoint{5.173400in}{3.192000in}}%
\pgfpathlineto{\pgfqpoint{5.197521in}{3.192000in}}%
\pgfpathlineto{\pgfqpoint{5.221652in}{3.192000in}}%
\pgfpathlineto{\pgfqpoint{5.245763in}{3.254101in}}%
\pgfpathlineto{\pgfqpoint{5.269946in}{3.192000in}}%
\pgfpathlineto{\pgfqpoint{5.294086in}{3.192000in}}%
\pgfpathlineto{\pgfqpoint{5.318342in}{3.215333in}}%
\pgfpathlineto{\pgfqpoint{5.342607in}{3.192000in}}%
\pgfpathlineto{\pgfqpoint{5.366900in}{3.192000in}}%
\pgfpathlineto{\pgfqpoint{5.391147in}{3.192000in}}%
\pgfpathlineto{\pgfqpoint{5.415450in}{3.244986in}}%
\pgfpathlineto{\pgfqpoint{5.439668in}{3.192000in}}%
\pgfpathlineto{\pgfqpoint{5.463909in}{3.192000in}}%
\pgfpathlineto{\pgfqpoint{5.488042in}{3.192000in}}%
\pgfpathlineto{\pgfqpoint{5.512218in}{3.192000in}}%
\pgfpathlineto{\pgfqpoint{5.536370in}{3.207069in}}%
\pgfpathlineto{\pgfqpoint{5.560550in}{3.314622in}}%
\pgfpathlineto{\pgfqpoint{5.584664in}{3.255559in}}%
\pgfpathlineto{\pgfqpoint{5.608858in}{3.316201in}}%
\pgfpathlineto{\pgfqpoint{5.632975in}{3.192000in}}%
\pgfpathlineto{\pgfqpoint{5.657110in}{3.207799in}}%
\pgfpathlineto{\pgfqpoint{5.681285in}{3.192000in}}%
\pgfpathlineto{\pgfqpoint{5.705492in}{3.234535in}}%
\pgfpathlineto{\pgfqpoint{5.729627in}{3.192000in}}%
\pgfpathlineto{\pgfqpoint{5.753760in}{3.192000in}}%
\pgfpathlineto{\pgfqpoint{5.778025in}{3.316444in}}%
\pgfpathlineto{\pgfqpoint{5.802148in}{3.192000in}}%
\pgfpathlineto{\pgfqpoint{5.826317in}{3.192000in}}%
\pgfpathlineto{\pgfqpoint{5.850452in}{3.241948in}}%
\pgfpathlineto{\pgfqpoint{5.874630in}{3.192000in}}%
\pgfpathlineto{\pgfqpoint{5.898855in}{3.254222in}}%
\pgfpathlineto{\pgfqpoint{5.922988in}{3.192000in}}%
\pgfpathlineto{\pgfqpoint{5.947378in}{3.192000in}}%
\pgfpathlineto{\pgfqpoint{5.971608in}{3.201115in}}%
\pgfpathlineto{\pgfqpoint{5.995781in}{3.254101in}}%
\pgfpathlineto{\pgfqpoint{6.019912in}{3.225785in}}%
\pgfpathlineto{\pgfqpoint{6.044059in}{3.254101in}}%
\pgfpathlineto{\pgfqpoint{6.068275in}{3.197469in}}%
\pgfpathlineto{\pgfqpoint{6.092417in}{3.274031in}}%
\pgfpathlineto{\pgfqpoint{6.116581in}{3.192000in}}%
\pgfpathlineto{\pgfqpoint{6.140737in}{3.316323in}}%
\pgfpathlineto{\pgfqpoint{6.164873in}{3.192000in}}%
\pgfpathlineto{\pgfqpoint{6.189362in}{3.352781in}}%
\pgfpathlineto{\pgfqpoint{6.213516in}{3.348042in}}%
\pgfpathlineto{\pgfqpoint{6.237746in}{3.274517in}}%
\pgfpathlineto{\pgfqpoint{6.261893in}{3.283146in}}%
\pgfpathlineto{\pgfqpoint{6.286048in}{3.316201in}}%
\pgfpathlineto{\pgfqpoint{6.310206in}{3.316444in}}%
\pgfpathlineto{\pgfqpoint{6.334434in}{3.192000in}}%
\pgfpathlineto{\pgfqpoint{6.358607in}{3.229552in}}%
\pgfpathlineto{\pgfqpoint{6.382804in}{3.256531in}}%
\pgfpathlineto{\pgfqpoint{6.406925in}{3.192000in}}%
\pgfpathlineto{\pgfqpoint{6.431105in}{3.378545in}}%
\pgfpathlineto{\pgfqpoint{6.455295in}{3.192000in}}%
\pgfpathlineto{\pgfqpoint{6.479444in}{3.192000in}}%
\pgfpathlineto{\pgfqpoint{6.503669in}{3.192000in}}%
\pgfpathlineto{\pgfqpoint{6.527835in}{3.270142in}}%
\pgfpathlineto{\pgfqpoint{6.551954in}{3.192000in}}%
\pgfpathlineto{\pgfqpoint{6.576049in}{3.223719in}}%
\pgfpathlineto{\pgfqpoint{6.600326in}{3.192000in}}%
\pgfpathlineto{\pgfqpoint{6.624766in}{3.192000in}}%
\pgfpathlineto{\pgfqpoint{6.649005in}{3.378545in}}%
\pgfusepath{stroke}%
\end{pgfscope}%
\begin{pgfscope}%
\pgfsetrectcap%
\pgfsetmiterjoin%
\pgfsetlinewidth{0.803000pt}%
\definecolor{currentstroke}{rgb}{0.101961,0.101961,0.101961}%
\pgfsetstrokecolor{currentstroke}%
\pgfsetdash{}{0pt}%
\pgfpathmoveto{\pgfqpoint{3.815000in}{3.192000in}}%
\pgfpathlineto{\pgfqpoint{3.815000in}{3.822000in}}%
\pgfusepath{stroke}%
\end{pgfscope}%
\begin{pgfscope}%
\pgfsetrectcap%
\pgfsetmiterjoin%
\pgfsetlinewidth{0.803000pt}%
\definecolor{currentstroke}{rgb}{0.101961,0.101961,0.101961}%
\pgfsetstrokecolor{currentstroke}%
\pgfsetdash{}{0pt}%
\pgfpathmoveto{\pgfqpoint{3.815000in}{3.192000in}}%
\pgfpathlineto{\pgfqpoint{6.650000in}{3.192000in}}%
\pgfusepath{stroke}%
\end{pgfscope}%
\begin{pgfscope}%
\definecolor{textcolor}{rgb}{0.101961,0.101961,0.101961}%
\pgfsetstrokecolor{textcolor}%
\pgfsetfillcolor{textcolor}%
\pgftext[x=3.843350in,y=3.840900in,left,bottom]{\color{textcolor}\sffamily\fontsize{7.000000}{8.400000}\selectfont (b) 14-robot outage run, reconnect at 595.5 s}%
\end{pgfscope}%
\begin{pgfscope}%
\pgfsetbuttcap%
\pgfsetmiterjoin%
\definecolor{currentfill}{rgb}{1.000000,1.000000,1.000000}%
\pgfsetfillcolor{currentfill}%
\pgfsetlinewidth{0.000000pt}%
\definecolor{currentstroke}{rgb}{0.000000,0.000000,0.000000}%
\pgfsetstrokecolor{currentstroke}%
\pgfsetstrokeopacity{0.000000}%
\pgfsetdash{}{0pt}%
\pgfpathmoveto{\pgfqpoint{3.815000in}{2.289000in}}%
\pgfpathlineto{\pgfqpoint{6.650000in}{2.289000in}}%
\pgfpathlineto{\pgfqpoint{6.650000in}{2.919000in}}%
\pgfpathlineto{\pgfqpoint{3.815000in}{2.919000in}}%
\pgfpathlineto{\pgfqpoint{3.815000in}{2.289000in}}%
\pgfpathclose%
\pgfusepath{fill}%
\end{pgfscope}%
\begin{pgfscope}%
\pgfpathrectangle{\pgfqpoint{3.815000in}{2.289000in}}{\pgfqpoint{2.835000in}{0.630000in}}%
\pgfusepath{clip}%
\pgfsetbuttcap%
\pgfsetmiterjoin%
\definecolor{currentfill}{rgb}{0.941176,0.894118,0.894118}%
\pgfsetfillcolor{currentfill}%
\pgfsetlinewidth{0.000000pt}%
\definecolor{currentstroke}{rgb}{0.941176,0.894118,0.894118}%
\pgfsetstrokecolor{currentstroke}%
\pgfsetdash{}{0pt}%
\pgfpathmoveto{\pgfqpoint{4.777719in}{2.289000in}}%
\pgfpathlineto{\pgfqpoint{4.777719in}{2.919000in}}%
\pgfpathlineto{\pgfqpoint{5.221869in}{2.919000in}}%
\pgfpathlineto{\pgfqpoint{5.221869in}{2.289000in}}%
\pgfpathlineto{\pgfqpoint{4.777719in}{2.289000in}}%
\pgfpathclose%
\pgfusepath{fill}%
\end{pgfscope}%
\begin{pgfscope}%
\pgfpathrectangle{\pgfqpoint{3.815000in}{2.289000in}}{\pgfqpoint{2.835000in}{0.630000in}}%
\pgfusepath{clip}%
\pgfsetrectcap%
\pgfsetroundjoin%
\pgfsetlinewidth{0.401500pt}%
\definecolor{currentstroke}{rgb}{0.901961,0.901961,0.901961}%
\pgfsetstrokecolor{currentstroke}%
\pgfsetdash{}{0pt}%
\pgfpathmoveto{\pgfqpoint{3.815000in}{2.289000in}}%
\pgfpathlineto{\pgfqpoint{6.650000in}{2.289000in}}%
\pgfusepath{stroke}%
\end{pgfscope}%
\begin{pgfscope}%
\pgfpathrectangle{\pgfqpoint{3.815000in}{2.289000in}}{\pgfqpoint{2.835000in}{0.630000in}}%
\pgfusepath{clip}%
\pgfsetrectcap%
\pgfsetroundjoin%
\pgfsetlinewidth{0.401500pt}%
\definecolor{currentstroke}{rgb}{0.901961,0.901961,0.901961}%
\pgfsetstrokecolor{currentstroke}%
\pgfsetdash{}{0pt}%
\pgfpathmoveto{\pgfqpoint{3.815000in}{2.604000in}}%
\pgfpathlineto{\pgfqpoint{6.650000in}{2.604000in}}%
\pgfusepath{stroke}%
\end{pgfscope}%
\begin{pgfscope}%
\pgfpathrectangle{\pgfqpoint{3.815000in}{2.289000in}}{\pgfqpoint{2.835000in}{0.630000in}}%
\pgfusepath{clip}%
\pgfsetrectcap%
\pgfsetroundjoin%
\pgfsetlinewidth{0.401500pt}%
\definecolor{currentstroke}{rgb}{0.901961,0.901961,0.901961}%
\pgfsetstrokecolor{currentstroke}%
\pgfsetdash{}{0pt}%
\pgfpathmoveto{\pgfqpoint{3.815000in}{2.919000in}}%
\pgfpathlineto{\pgfqpoint{6.650000in}{2.919000in}}%
\pgfusepath{stroke}%
\end{pgfscope}%
\begin{pgfscope}%
\pgfpathrectangle{\pgfqpoint{3.815000in}{2.289000in}}{\pgfqpoint{2.835000in}{0.630000in}}%
\pgfusepath{clip}%
\pgfsetbuttcap%
\pgfsetroundjoin%
\pgfsetlinewidth{0.702625pt}%
\definecolor{currentstroke}{rgb}{0.690196,0.250980,0.250980}%
\pgfsetstrokecolor{currentstroke}%
\pgfsetdash{{2.590000pt}{1.120000pt}}{0.000000pt}%
\pgfpathmoveto{\pgfqpoint{5.221869in}{2.289000in}}%
\pgfpathlineto{\pgfqpoint{5.221869in}{2.919000in}}%
\pgfusepath{stroke}%
\end{pgfscope}%
\begin{pgfscope}%
\pgfpathrectangle{\pgfqpoint{3.815000in}{2.289000in}}{\pgfqpoint{2.835000in}{0.630000in}}%
\pgfusepath{clip}%
\pgfsetrectcap%
\pgfsetroundjoin%
\pgfsetlinewidth{0.903375pt}%
\definecolor{currentstroke}{rgb}{0.156863,0.356863,0.568627}%
\pgfsetstrokecolor{currentstroke}%
\pgfsetdash{}{0pt}%
\pgfpathmoveto{\pgfqpoint{3.826812in}{2.289000in}}%
\pgfpathlineto{\pgfqpoint{5.279750in}{2.289000in}}%
\pgfpathlineto{\pgfqpoint{5.279750in}{2.566200in}}%
\pgfpathlineto{\pgfqpoint{5.303375in}{2.566200in}}%
\pgfpathlineto{\pgfqpoint{5.303375in}{2.452800in}}%
\pgfpathlineto{\pgfqpoint{5.327000in}{2.452800in}}%
\pgfpathlineto{\pgfqpoint{5.327000in}{2.604000in}}%
\pgfpathlineto{\pgfqpoint{5.350625in}{2.604000in}}%
\pgfpathlineto{\pgfqpoint{5.350625in}{2.594550in}}%
\pgfpathlineto{\pgfqpoint{5.374250in}{2.594550in}}%
\pgfpathlineto{\pgfqpoint{5.374250in}{2.604000in}}%
\pgfpathlineto{\pgfqpoint{5.421500in}{2.604000in}}%
\pgfpathlineto{\pgfqpoint{5.421500in}{2.556750in}}%
\pgfpathlineto{\pgfqpoint{5.445125in}{2.556750in}}%
\pgfpathlineto{\pgfqpoint{5.445125in}{2.604000in}}%
\pgfpathlineto{\pgfqpoint{5.492375in}{2.604000in}}%
\pgfpathlineto{\pgfqpoint{5.492375in}{2.392950in}}%
\pgfpathlineto{\pgfqpoint{5.516000in}{2.392950in}}%
\pgfpathlineto{\pgfqpoint{5.516000in}{2.342550in}}%
\pgfpathlineto{\pgfqpoint{5.539625in}{2.342550in}}%
\pgfpathlineto{\pgfqpoint{5.539625in}{2.361450in}}%
\pgfpathlineto{\pgfqpoint{5.563250in}{2.361450in}}%
\pgfpathlineto{\pgfqpoint{5.563250in}{2.289000in}}%
\pgfpathlineto{\pgfqpoint{5.917625in}{2.289000in}}%
\pgfpathlineto{\pgfqpoint{5.917625in}{2.383500in}}%
\pgfpathlineto{\pgfqpoint{5.941250in}{2.383500in}}%
\pgfpathlineto{\pgfqpoint{5.941250in}{2.307900in}}%
\pgfpathlineto{\pgfqpoint{5.964875in}{2.307900in}}%
\pgfpathlineto{\pgfqpoint{5.964875in}{2.377200in}}%
\pgfpathlineto{\pgfqpoint{5.988500in}{2.377200in}}%
\pgfpathlineto{\pgfqpoint{5.988500in}{2.588250in}}%
\pgfpathlineto{\pgfqpoint{6.012125in}{2.588250in}}%
\pgfpathlineto{\pgfqpoint{6.012125in}{2.604000in}}%
\pgfpathlineto{\pgfqpoint{6.035750in}{2.604000in}}%
\pgfpathlineto{\pgfqpoint{6.035750in}{2.493750in}}%
\pgfpathlineto{\pgfqpoint{6.059375in}{2.493750in}}%
\pgfpathlineto{\pgfqpoint{6.059375in}{2.289000in}}%
\pgfpathlineto{\pgfqpoint{6.083000in}{2.289000in}}%
\pgfpathlineto{\pgfqpoint{6.083000in}{2.399250in}}%
\pgfpathlineto{\pgfqpoint{6.106625in}{2.399250in}}%
\pgfpathlineto{\pgfqpoint{6.106625in}{2.311050in}}%
\pgfpathlineto{\pgfqpoint{6.130250in}{2.311050in}}%
\pgfpathlineto{\pgfqpoint{6.130250in}{2.289000in}}%
\pgfpathlineto{\pgfqpoint{6.390125in}{2.289000in}}%
\pgfpathlineto{\pgfqpoint{6.390125in}{2.383500in}}%
\pgfpathlineto{\pgfqpoint{6.413750in}{2.383500in}}%
\pgfpathlineto{\pgfqpoint{6.413750in}{2.289000in}}%
\pgfpathlineto{\pgfqpoint{6.555500in}{2.289000in}}%
\pgfpathlineto{\pgfqpoint{6.555500in}{2.389800in}}%
\pgfpathlineto{\pgfqpoint{6.579125in}{2.389800in}}%
\pgfpathlineto{\pgfqpoint{6.579125in}{2.392950in}}%
\pgfpathlineto{\pgfqpoint{6.602750in}{2.392950in}}%
\pgfpathlineto{\pgfqpoint{6.602750in}{2.311050in}}%
\pgfpathlineto{\pgfqpoint{6.626375in}{2.311050in}}%
\pgfpathlineto{\pgfqpoint{6.626375in}{2.289000in}}%
\pgfpathlineto{\pgfqpoint{6.638188in}{2.289000in}}%
\pgfpathlineto{\pgfqpoint{6.638188in}{2.289000in}}%
\pgfusepath{stroke}%
\end{pgfscope}%
\begin{pgfscope}%
\pgfsetrectcap%
\pgfsetmiterjoin%
\pgfsetlinewidth{0.803000pt}%
\definecolor{currentstroke}{rgb}{0.101961,0.101961,0.101961}%
\pgfsetstrokecolor{currentstroke}%
\pgfsetdash{}{0pt}%
\pgfpathmoveto{\pgfqpoint{3.815000in}{2.289000in}}%
\pgfpathlineto{\pgfqpoint{3.815000in}{2.919000in}}%
\pgfusepath{stroke}%
\end{pgfscope}%
\begin{pgfscope}%
\pgfsetrectcap%
\pgfsetmiterjoin%
\pgfsetlinewidth{0.803000pt}%
\definecolor{currentstroke}{rgb}{0.101961,0.101961,0.101961}%
\pgfsetstrokecolor{currentstroke}%
\pgfsetdash{}{0pt}%
\pgfpathmoveto{\pgfqpoint{3.815000in}{2.289000in}}%
\pgfpathlineto{\pgfqpoint{6.650000in}{2.289000in}}%
\pgfusepath{stroke}%
\end{pgfscope}%
\begin{pgfscope}%
\definecolor{textcolor}{rgb}{0.690196,0.250980,0.250980}%
\pgfsetstrokecolor{textcolor}%
\pgfsetfillcolor{textcolor}%
\pgftext[x=4.999794in,y=2.887500in,,top]{\color{textcolor}\sffamily\fontsize{7.000000}{8.400000}\selectfont outage}%
\end{pgfscope}%
\begin{pgfscope}%
\pgfsetbuttcap%
\pgfsetmiterjoin%
\definecolor{currentfill}{rgb}{1.000000,1.000000,1.000000}%
\pgfsetfillcolor{currentfill}%
\pgfsetlinewidth{0.000000pt}%
\definecolor{currentstroke}{rgb}{0.000000,0.000000,0.000000}%
\pgfsetstrokecolor{currentstroke}%
\pgfsetstrokeopacity{0.000000}%
\pgfsetdash{}{0pt}%
\pgfpathmoveto{\pgfqpoint{3.815000in}{1.386000in}}%
\pgfpathlineto{\pgfqpoint{6.650000in}{1.386000in}}%
\pgfpathlineto{\pgfqpoint{6.650000in}{2.016000in}}%
\pgfpathlineto{\pgfqpoint{3.815000in}{2.016000in}}%
\pgfpathlineto{\pgfqpoint{3.815000in}{1.386000in}}%
\pgfpathclose%
\pgfusepath{fill}%
\end{pgfscope}%
\begin{pgfscope}%
\pgfpathrectangle{\pgfqpoint{3.815000in}{1.386000in}}{\pgfqpoint{2.835000in}{0.630000in}}%
\pgfusepath{clip}%
\pgfsetbuttcap%
\pgfsetmiterjoin%
\definecolor{currentfill}{rgb}{0.941176,0.894118,0.894118}%
\pgfsetfillcolor{currentfill}%
\pgfsetlinewidth{0.000000pt}%
\definecolor{currentstroke}{rgb}{0.941176,0.894118,0.894118}%
\pgfsetstrokecolor{currentstroke}%
\pgfsetdash{}{0pt}%
\pgfpathmoveto{\pgfqpoint{4.777719in}{1.386000in}}%
\pgfpathlineto{\pgfqpoint{4.777719in}{2.016000in}}%
\pgfpathlineto{\pgfqpoint{5.221869in}{2.016000in}}%
\pgfpathlineto{\pgfqpoint{5.221869in}{1.386000in}}%
\pgfpathlineto{\pgfqpoint{4.777719in}{1.386000in}}%
\pgfpathclose%
\pgfusepath{fill}%
\end{pgfscope}%
\begin{pgfscope}%
\pgfpathrectangle{\pgfqpoint{3.815000in}{1.386000in}}{\pgfqpoint{2.835000in}{0.630000in}}%
\pgfusepath{clip}%
\pgfsetrectcap%
\pgfsetroundjoin%
\pgfsetlinewidth{0.401500pt}%
\definecolor{currentstroke}{rgb}{0.901961,0.901961,0.901961}%
\pgfsetstrokecolor{currentstroke}%
\pgfsetdash{}{0pt}%
\pgfpathmoveto{\pgfqpoint{3.815000in}{1.386000in}}%
\pgfpathlineto{\pgfqpoint{6.650000in}{1.386000in}}%
\pgfusepath{stroke}%
\end{pgfscope}%
\begin{pgfscope}%
\pgfpathrectangle{\pgfqpoint{3.815000in}{1.386000in}}{\pgfqpoint{2.835000in}{0.630000in}}%
\pgfusepath{clip}%
\pgfsetrectcap%
\pgfsetroundjoin%
\pgfsetlinewidth{0.401500pt}%
\definecolor{currentstroke}{rgb}{0.901961,0.901961,0.901961}%
\pgfsetstrokecolor{currentstroke}%
\pgfsetdash{}{0pt}%
\pgfpathmoveto{\pgfqpoint{3.815000in}{1.701000in}}%
\pgfpathlineto{\pgfqpoint{6.650000in}{1.701000in}}%
\pgfusepath{stroke}%
\end{pgfscope}%
\begin{pgfscope}%
\pgfpathrectangle{\pgfqpoint{3.815000in}{1.386000in}}{\pgfqpoint{2.835000in}{0.630000in}}%
\pgfusepath{clip}%
\pgfsetrectcap%
\pgfsetroundjoin%
\pgfsetlinewidth{0.401500pt}%
\definecolor{currentstroke}{rgb}{0.901961,0.901961,0.901961}%
\pgfsetstrokecolor{currentstroke}%
\pgfsetdash{}{0pt}%
\pgfpathmoveto{\pgfqpoint{3.815000in}{2.016000in}}%
\pgfpathlineto{\pgfqpoint{6.650000in}{2.016000in}}%
\pgfusepath{stroke}%
\end{pgfscope}%
\begin{pgfscope}%
\pgfpathrectangle{\pgfqpoint{3.815000in}{1.386000in}}{\pgfqpoint{2.835000in}{0.630000in}}%
\pgfusepath{clip}%
\pgfsetbuttcap%
\pgfsetmiterjoin%
\definecolor{currentfill}{rgb}{0.101961,0.101961,0.101961}%
\pgfsetfillcolor{currentfill}%
\pgfsetlinewidth{0.000000pt}%
\definecolor{currentstroke}{rgb}{0.000000,0.000000,0.000000}%
\pgfsetstrokecolor{currentstroke}%
\pgfsetstrokeopacity{0.000000}%
\pgfsetdash{}{0pt}%
\pgfpathmoveto{\pgfqpoint{3.816181in}{1.386000in}}%
\pgfpathlineto{\pgfqpoint{3.837444in}{1.386000in}}%
\pgfpathlineto{\pgfqpoint{3.837444in}{1.386000in}}%
\pgfpathlineto{\pgfqpoint{3.816181in}{1.386000in}}%
\pgfpathlineto{\pgfqpoint{3.816181in}{1.386000in}}%
\pgfpathclose%
\pgfusepath{fill}%
\end{pgfscope}%
\begin{pgfscope}%
\pgfpathrectangle{\pgfqpoint{3.815000in}{1.386000in}}{\pgfqpoint{2.835000in}{0.630000in}}%
\pgfusepath{clip}%
\pgfsetbuttcap%
\pgfsetmiterjoin%
\definecolor{currentfill}{rgb}{0.101961,0.101961,0.101961}%
\pgfsetfillcolor{currentfill}%
\pgfsetlinewidth{0.000000pt}%
\definecolor{currentstroke}{rgb}{0.000000,0.000000,0.000000}%
\pgfsetstrokecolor{currentstroke}%
\pgfsetstrokeopacity{0.000000}%
\pgfsetdash{}{0pt}%
\pgfpathmoveto{\pgfqpoint{3.839806in}{1.386000in}}%
\pgfpathlineto{\pgfqpoint{3.861069in}{1.386000in}}%
\pgfpathlineto{\pgfqpoint{3.861069in}{1.386000in}}%
\pgfpathlineto{\pgfqpoint{3.839806in}{1.386000in}}%
\pgfpathlineto{\pgfqpoint{3.839806in}{1.386000in}}%
\pgfpathclose%
\pgfusepath{fill}%
\end{pgfscope}%
\begin{pgfscope}%
\pgfpathrectangle{\pgfqpoint{3.815000in}{1.386000in}}{\pgfqpoint{2.835000in}{0.630000in}}%
\pgfusepath{clip}%
\pgfsetbuttcap%
\pgfsetmiterjoin%
\definecolor{currentfill}{rgb}{0.101961,0.101961,0.101961}%
\pgfsetfillcolor{currentfill}%
\pgfsetlinewidth{0.000000pt}%
\definecolor{currentstroke}{rgb}{0.000000,0.000000,0.000000}%
\pgfsetstrokecolor{currentstroke}%
\pgfsetstrokeopacity{0.000000}%
\pgfsetdash{}{0pt}%
\pgfpathmoveto{\pgfqpoint{3.863431in}{1.386000in}}%
\pgfpathlineto{\pgfqpoint{3.884694in}{1.386000in}}%
\pgfpathlineto{\pgfqpoint{3.884694in}{1.386000in}}%
\pgfpathlineto{\pgfqpoint{3.863431in}{1.386000in}}%
\pgfpathlineto{\pgfqpoint{3.863431in}{1.386000in}}%
\pgfpathclose%
\pgfusepath{fill}%
\end{pgfscope}%
\begin{pgfscope}%
\pgfpathrectangle{\pgfqpoint{3.815000in}{1.386000in}}{\pgfqpoint{2.835000in}{0.630000in}}%
\pgfusepath{clip}%
\pgfsetbuttcap%
\pgfsetmiterjoin%
\definecolor{currentfill}{rgb}{0.101961,0.101961,0.101961}%
\pgfsetfillcolor{currentfill}%
\pgfsetlinewidth{0.000000pt}%
\definecolor{currentstroke}{rgb}{0.000000,0.000000,0.000000}%
\pgfsetstrokecolor{currentstroke}%
\pgfsetstrokeopacity{0.000000}%
\pgfsetdash{}{0pt}%
\pgfpathmoveto{\pgfqpoint{3.887056in}{1.386000in}}%
\pgfpathlineto{\pgfqpoint{3.908319in}{1.386000in}}%
\pgfpathlineto{\pgfqpoint{3.908319in}{1.386000in}}%
\pgfpathlineto{\pgfqpoint{3.887056in}{1.386000in}}%
\pgfpathlineto{\pgfqpoint{3.887056in}{1.386000in}}%
\pgfpathclose%
\pgfusepath{fill}%
\end{pgfscope}%
\begin{pgfscope}%
\pgfpathrectangle{\pgfqpoint{3.815000in}{1.386000in}}{\pgfqpoint{2.835000in}{0.630000in}}%
\pgfusepath{clip}%
\pgfsetbuttcap%
\pgfsetmiterjoin%
\definecolor{currentfill}{rgb}{0.101961,0.101961,0.101961}%
\pgfsetfillcolor{currentfill}%
\pgfsetlinewidth{0.000000pt}%
\definecolor{currentstroke}{rgb}{0.000000,0.000000,0.000000}%
\pgfsetstrokecolor{currentstroke}%
\pgfsetstrokeopacity{0.000000}%
\pgfsetdash{}{0pt}%
\pgfpathmoveto{\pgfqpoint{3.910681in}{1.386000in}}%
\pgfpathlineto{\pgfqpoint{3.931944in}{1.386000in}}%
\pgfpathlineto{\pgfqpoint{3.931944in}{1.386000in}}%
\pgfpathlineto{\pgfqpoint{3.910681in}{1.386000in}}%
\pgfpathlineto{\pgfqpoint{3.910681in}{1.386000in}}%
\pgfpathclose%
\pgfusepath{fill}%
\end{pgfscope}%
\begin{pgfscope}%
\pgfpathrectangle{\pgfqpoint{3.815000in}{1.386000in}}{\pgfqpoint{2.835000in}{0.630000in}}%
\pgfusepath{clip}%
\pgfsetbuttcap%
\pgfsetmiterjoin%
\definecolor{currentfill}{rgb}{0.101961,0.101961,0.101961}%
\pgfsetfillcolor{currentfill}%
\pgfsetlinewidth{0.000000pt}%
\definecolor{currentstroke}{rgb}{0.000000,0.000000,0.000000}%
\pgfsetstrokecolor{currentstroke}%
\pgfsetstrokeopacity{0.000000}%
\pgfsetdash{}{0pt}%
\pgfpathmoveto{\pgfqpoint{3.934306in}{1.386000in}}%
\pgfpathlineto{\pgfqpoint{3.955569in}{1.386000in}}%
\pgfpathlineto{\pgfqpoint{3.955569in}{1.386000in}}%
\pgfpathlineto{\pgfqpoint{3.934306in}{1.386000in}}%
\pgfpathlineto{\pgfqpoint{3.934306in}{1.386000in}}%
\pgfpathclose%
\pgfusepath{fill}%
\end{pgfscope}%
\begin{pgfscope}%
\pgfpathrectangle{\pgfqpoint{3.815000in}{1.386000in}}{\pgfqpoint{2.835000in}{0.630000in}}%
\pgfusepath{clip}%
\pgfsetbuttcap%
\pgfsetmiterjoin%
\definecolor{currentfill}{rgb}{0.101961,0.101961,0.101961}%
\pgfsetfillcolor{currentfill}%
\pgfsetlinewidth{0.000000pt}%
\definecolor{currentstroke}{rgb}{0.000000,0.000000,0.000000}%
\pgfsetstrokecolor{currentstroke}%
\pgfsetstrokeopacity{0.000000}%
\pgfsetdash{}{0pt}%
\pgfpathmoveto{\pgfqpoint{3.957931in}{1.386000in}}%
\pgfpathlineto{\pgfqpoint{3.979194in}{1.386000in}}%
\pgfpathlineto{\pgfqpoint{3.979194in}{1.386000in}}%
\pgfpathlineto{\pgfqpoint{3.957931in}{1.386000in}}%
\pgfpathlineto{\pgfqpoint{3.957931in}{1.386000in}}%
\pgfpathclose%
\pgfusepath{fill}%
\end{pgfscope}%
\begin{pgfscope}%
\pgfpathrectangle{\pgfqpoint{3.815000in}{1.386000in}}{\pgfqpoint{2.835000in}{0.630000in}}%
\pgfusepath{clip}%
\pgfsetbuttcap%
\pgfsetmiterjoin%
\definecolor{currentfill}{rgb}{0.101961,0.101961,0.101961}%
\pgfsetfillcolor{currentfill}%
\pgfsetlinewidth{0.000000pt}%
\definecolor{currentstroke}{rgb}{0.000000,0.000000,0.000000}%
\pgfsetstrokecolor{currentstroke}%
\pgfsetstrokeopacity{0.000000}%
\pgfsetdash{}{0pt}%
\pgfpathmoveto{\pgfqpoint{3.981556in}{1.386000in}}%
\pgfpathlineto{\pgfqpoint{4.002819in}{1.386000in}}%
\pgfpathlineto{\pgfqpoint{4.002819in}{1.386000in}}%
\pgfpathlineto{\pgfqpoint{3.981556in}{1.386000in}}%
\pgfpathlineto{\pgfqpoint{3.981556in}{1.386000in}}%
\pgfpathclose%
\pgfusepath{fill}%
\end{pgfscope}%
\begin{pgfscope}%
\pgfpathrectangle{\pgfqpoint{3.815000in}{1.386000in}}{\pgfqpoint{2.835000in}{0.630000in}}%
\pgfusepath{clip}%
\pgfsetbuttcap%
\pgfsetmiterjoin%
\definecolor{currentfill}{rgb}{0.101961,0.101961,0.101961}%
\pgfsetfillcolor{currentfill}%
\pgfsetlinewidth{0.000000pt}%
\definecolor{currentstroke}{rgb}{0.000000,0.000000,0.000000}%
\pgfsetstrokecolor{currentstroke}%
\pgfsetstrokeopacity{0.000000}%
\pgfsetdash{}{0pt}%
\pgfpathmoveto{\pgfqpoint{4.005181in}{1.386000in}}%
\pgfpathlineto{\pgfqpoint{4.026444in}{1.386000in}}%
\pgfpathlineto{\pgfqpoint{4.026444in}{1.386000in}}%
\pgfpathlineto{\pgfqpoint{4.005181in}{1.386000in}}%
\pgfpathlineto{\pgfqpoint{4.005181in}{1.386000in}}%
\pgfpathclose%
\pgfusepath{fill}%
\end{pgfscope}%
\begin{pgfscope}%
\pgfpathrectangle{\pgfqpoint{3.815000in}{1.386000in}}{\pgfqpoint{2.835000in}{0.630000in}}%
\pgfusepath{clip}%
\pgfsetbuttcap%
\pgfsetmiterjoin%
\definecolor{currentfill}{rgb}{0.101961,0.101961,0.101961}%
\pgfsetfillcolor{currentfill}%
\pgfsetlinewidth{0.000000pt}%
\definecolor{currentstroke}{rgb}{0.000000,0.000000,0.000000}%
\pgfsetstrokecolor{currentstroke}%
\pgfsetstrokeopacity{0.000000}%
\pgfsetdash{}{0pt}%
\pgfpathmoveto{\pgfqpoint{4.028806in}{1.386000in}}%
\pgfpathlineto{\pgfqpoint{4.050069in}{1.386000in}}%
\pgfpathlineto{\pgfqpoint{4.050069in}{1.386000in}}%
\pgfpathlineto{\pgfqpoint{4.028806in}{1.386000in}}%
\pgfpathlineto{\pgfqpoint{4.028806in}{1.386000in}}%
\pgfpathclose%
\pgfusepath{fill}%
\end{pgfscope}%
\begin{pgfscope}%
\pgfpathrectangle{\pgfqpoint{3.815000in}{1.386000in}}{\pgfqpoint{2.835000in}{0.630000in}}%
\pgfusepath{clip}%
\pgfsetbuttcap%
\pgfsetmiterjoin%
\definecolor{currentfill}{rgb}{0.101961,0.101961,0.101961}%
\pgfsetfillcolor{currentfill}%
\pgfsetlinewidth{0.000000pt}%
\definecolor{currentstroke}{rgb}{0.000000,0.000000,0.000000}%
\pgfsetstrokecolor{currentstroke}%
\pgfsetstrokeopacity{0.000000}%
\pgfsetdash{}{0pt}%
\pgfpathmoveto{\pgfqpoint{4.052431in}{1.386000in}}%
\pgfpathlineto{\pgfqpoint{4.073694in}{1.386000in}}%
\pgfpathlineto{\pgfqpoint{4.073694in}{1.386000in}}%
\pgfpathlineto{\pgfqpoint{4.052431in}{1.386000in}}%
\pgfpathlineto{\pgfqpoint{4.052431in}{1.386000in}}%
\pgfpathclose%
\pgfusepath{fill}%
\end{pgfscope}%
\begin{pgfscope}%
\pgfpathrectangle{\pgfqpoint{3.815000in}{1.386000in}}{\pgfqpoint{2.835000in}{0.630000in}}%
\pgfusepath{clip}%
\pgfsetbuttcap%
\pgfsetmiterjoin%
\definecolor{currentfill}{rgb}{0.101961,0.101961,0.101961}%
\pgfsetfillcolor{currentfill}%
\pgfsetlinewidth{0.000000pt}%
\definecolor{currentstroke}{rgb}{0.000000,0.000000,0.000000}%
\pgfsetstrokecolor{currentstroke}%
\pgfsetstrokeopacity{0.000000}%
\pgfsetdash{}{0pt}%
\pgfpathmoveto{\pgfqpoint{4.076056in}{1.386000in}}%
\pgfpathlineto{\pgfqpoint{4.097319in}{1.386000in}}%
\pgfpathlineto{\pgfqpoint{4.097319in}{1.386000in}}%
\pgfpathlineto{\pgfqpoint{4.076056in}{1.386000in}}%
\pgfpathlineto{\pgfqpoint{4.076056in}{1.386000in}}%
\pgfpathclose%
\pgfusepath{fill}%
\end{pgfscope}%
\begin{pgfscope}%
\pgfpathrectangle{\pgfqpoint{3.815000in}{1.386000in}}{\pgfqpoint{2.835000in}{0.630000in}}%
\pgfusepath{clip}%
\pgfsetbuttcap%
\pgfsetmiterjoin%
\definecolor{currentfill}{rgb}{0.101961,0.101961,0.101961}%
\pgfsetfillcolor{currentfill}%
\pgfsetlinewidth{0.000000pt}%
\definecolor{currentstroke}{rgb}{0.000000,0.000000,0.000000}%
\pgfsetstrokecolor{currentstroke}%
\pgfsetstrokeopacity{0.000000}%
\pgfsetdash{}{0pt}%
\pgfpathmoveto{\pgfqpoint{4.099681in}{1.386000in}}%
\pgfpathlineto{\pgfqpoint{4.120944in}{1.386000in}}%
\pgfpathlineto{\pgfqpoint{4.120944in}{1.386000in}}%
\pgfpathlineto{\pgfqpoint{4.099681in}{1.386000in}}%
\pgfpathlineto{\pgfqpoint{4.099681in}{1.386000in}}%
\pgfpathclose%
\pgfusepath{fill}%
\end{pgfscope}%
\begin{pgfscope}%
\pgfpathrectangle{\pgfqpoint{3.815000in}{1.386000in}}{\pgfqpoint{2.835000in}{0.630000in}}%
\pgfusepath{clip}%
\pgfsetbuttcap%
\pgfsetmiterjoin%
\definecolor{currentfill}{rgb}{0.101961,0.101961,0.101961}%
\pgfsetfillcolor{currentfill}%
\pgfsetlinewidth{0.000000pt}%
\definecolor{currentstroke}{rgb}{0.000000,0.000000,0.000000}%
\pgfsetstrokecolor{currentstroke}%
\pgfsetstrokeopacity{0.000000}%
\pgfsetdash{}{0pt}%
\pgfpathmoveto{\pgfqpoint{4.123306in}{1.386000in}}%
\pgfpathlineto{\pgfqpoint{4.144569in}{1.386000in}}%
\pgfpathlineto{\pgfqpoint{4.144569in}{1.386000in}}%
\pgfpathlineto{\pgfqpoint{4.123306in}{1.386000in}}%
\pgfpathlineto{\pgfqpoint{4.123306in}{1.386000in}}%
\pgfpathclose%
\pgfusepath{fill}%
\end{pgfscope}%
\begin{pgfscope}%
\pgfpathrectangle{\pgfqpoint{3.815000in}{1.386000in}}{\pgfqpoint{2.835000in}{0.630000in}}%
\pgfusepath{clip}%
\pgfsetbuttcap%
\pgfsetmiterjoin%
\definecolor{currentfill}{rgb}{0.101961,0.101961,0.101961}%
\pgfsetfillcolor{currentfill}%
\pgfsetlinewidth{0.000000pt}%
\definecolor{currentstroke}{rgb}{0.000000,0.000000,0.000000}%
\pgfsetstrokecolor{currentstroke}%
\pgfsetstrokeopacity{0.000000}%
\pgfsetdash{}{0pt}%
\pgfpathmoveto{\pgfqpoint{4.146931in}{1.386000in}}%
\pgfpathlineto{\pgfqpoint{4.168194in}{1.386000in}}%
\pgfpathlineto{\pgfqpoint{4.168194in}{1.386000in}}%
\pgfpathlineto{\pgfqpoint{4.146931in}{1.386000in}}%
\pgfpathlineto{\pgfqpoint{4.146931in}{1.386000in}}%
\pgfpathclose%
\pgfusepath{fill}%
\end{pgfscope}%
\begin{pgfscope}%
\pgfpathrectangle{\pgfqpoint{3.815000in}{1.386000in}}{\pgfqpoint{2.835000in}{0.630000in}}%
\pgfusepath{clip}%
\pgfsetbuttcap%
\pgfsetmiterjoin%
\definecolor{currentfill}{rgb}{0.101961,0.101961,0.101961}%
\pgfsetfillcolor{currentfill}%
\pgfsetlinewidth{0.000000pt}%
\definecolor{currentstroke}{rgb}{0.000000,0.000000,0.000000}%
\pgfsetstrokecolor{currentstroke}%
\pgfsetstrokeopacity{0.000000}%
\pgfsetdash{}{0pt}%
\pgfpathmoveto{\pgfqpoint{4.170556in}{1.386000in}}%
\pgfpathlineto{\pgfqpoint{4.191819in}{1.386000in}}%
\pgfpathlineto{\pgfqpoint{4.191819in}{1.386000in}}%
\pgfpathlineto{\pgfqpoint{4.170556in}{1.386000in}}%
\pgfpathlineto{\pgfqpoint{4.170556in}{1.386000in}}%
\pgfpathclose%
\pgfusepath{fill}%
\end{pgfscope}%
\begin{pgfscope}%
\pgfpathrectangle{\pgfqpoint{3.815000in}{1.386000in}}{\pgfqpoint{2.835000in}{0.630000in}}%
\pgfusepath{clip}%
\pgfsetbuttcap%
\pgfsetmiterjoin%
\definecolor{currentfill}{rgb}{0.101961,0.101961,0.101961}%
\pgfsetfillcolor{currentfill}%
\pgfsetlinewidth{0.000000pt}%
\definecolor{currentstroke}{rgb}{0.000000,0.000000,0.000000}%
\pgfsetstrokecolor{currentstroke}%
\pgfsetstrokeopacity{0.000000}%
\pgfsetdash{}{0pt}%
\pgfpathmoveto{\pgfqpoint{4.194181in}{1.386000in}}%
\pgfpathlineto{\pgfqpoint{4.215444in}{1.386000in}}%
\pgfpathlineto{\pgfqpoint{4.215444in}{1.386000in}}%
\pgfpathlineto{\pgfqpoint{4.194181in}{1.386000in}}%
\pgfpathlineto{\pgfqpoint{4.194181in}{1.386000in}}%
\pgfpathclose%
\pgfusepath{fill}%
\end{pgfscope}%
\begin{pgfscope}%
\pgfpathrectangle{\pgfqpoint{3.815000in}{1.386000in}}{\pgfqpoint{2.835000in}{0.630000in}}%
\pgfusepath{clip}%
\pgfsetbuttcap%
\pgfsetmiterjoin%
\definecolor{currentfill}{rgb}{0.101961,0.101961,0.101961}%
\pgfsetfillcolor{currentfill}%
\pgfsetlinewidth{0.000000pt}%
\definecolor{currentstroke}{rgb}{0.000000,0.000000,0.000000}%
\pgfsetstrokecolor{currentstroke}%
\pgfsetstrokeopacity{0.000000}%
\pgfsetdash{}{0pt}%
\pgfpathmoveto{\pgfqpoint{4.217806in}{1.386000in}}%
\pgfpathlineto{\pgfqpoint{4.239069in}{1.386000in}}%
\pgfpathlineto{\pgfqpoint{4.239069in}{1.386000in}}%
\pgfpathlineto{\pgfqpoint{4.217806in}{1.386000in}}%
\pgfpathlineto{\pgfqpoint{4.217806in}{1.386000in}}%
\pgfpathclose%
\pgfusepath{fill}%
\end{pgfscope}%
\begin{pgfscope}%
\pgfpathrectangle{\pgfqpoint{3.815000in}{1.386000in}}{\pgfqpoint{2.835000in}{0.630000in}}%
\pgfusepath{clip}%
\pgfsetbuttcap%
\pgfsetmiterjoin%
\definecolor{currentfill}{rgb}{0.101961,0.101961,0.101961}%
\pgfsetfillcolor{currentfill}%
\pgfsetlinewidth{0.000000pt}%
\definecolor{currentstroke}{rgb}{0.000000,0.000000,0.000000}%
\pgfsetstrokecolor{currentstroke}%
\pgfsetstrokeopacity{0.000000}%
\pgfsetdash{}{0pt}%
\pgfpathmoveto{\pgfqpoint{4.241431in}{1.386000in}}%
\pgfpathlineto{\pgfqpoint{4.262694in}{1.386000in}}%
\pgfpathlineto{\pgfqpoint{4.262694in}{1.386000in}}%
\pgfpathlineto{\pgfqpoint{4.241431in}{1.386000in}}%
\pgfpathlineto{\pgfqpoint{4.241431in}{1.386000in}}%
\pgfpathclose%
\pgfusepath{fill}%
\end{pgfscope}%
\begin{pgfscope}%
\pgfpathrectangle{\pgfqpoint{3.815000in}{1.386000in}}{\pgfqpoint{2.835000in}{0.630000in}}%
\pgfusepath{clip}%
\pgfsetbuttcap%
\pgfsetmiterjoin%
\definecolor{currentfill}{rgb}{0.101961,0.101961,0.101961}%
\pgfsetfillcolor{currentfill}%
\pgfsetlinewidth{0.000000pt}%
\definecolor{currentstroke}{rgb}{0.000000,0.000000,0.000000}%
\pgfsetstrokecolor{currentstroke}%
\pgfsetstrokeopacity{0.000000}%
\pgfsetdash{}{0pt}%
\pgfpathmoveto{\pgfqpoint{4.265056in}{1.386000in}}%
\pgfpathlineto{\pgfqpoint{4.286319in}{1.386000in}}%
\pgfpathlineto{\pgfqpoint{4.286319in}{1.386000in}}%
\pgfpathlineto{\pgfqpoint{4.265056in}{1.386000in}}%
\pgfpathlineto{\pgfqpoint{4.265056in}{1.386000in}}%
\pgfpathclose%
\pgfusepath{fill}%
\end{pgfscope}%
\begin{pgfscope}%
\pgfpathrectangle{\pgfqpoint{3.815000in}{1.386000in}}{\pgfqpoint{2.835000in}{0.630000in}}%
\pgfusepath{clip}%
\pgfsetbuttcap%
\pgfsetmiterjoin%
\definecolor{currentfill}{rgb}{0.101961,0.101961,0.101961}%
\pgfsetfillcolor{currentfill}%
\pgfsetlinewidth{0.000000pt}%
\definecolor{currentstroke}{rgb}{0.000000,0.000000,0.000000}%
\pgfsetstrokecolor{currentstroke}%
\pgfsetstrokeopacity{0.000000}%
\pgfsetdash{}{0pt}%
\pgfpathmoveto{\pgfqpoint{4.288681in}{1.386000in}}%
\pgfpathlineto{\pgfqpoint{4.309944in}{1.386000in}}%
\pgfpathlineto{\pgfqpoint{4.309944in}{1.386000in}}%
\pgfpathlineto{\pgfqpoint{4.288681in}{1.386000in}}%
\pgfpathlineto{\pgfqpoint{4.288681in}{1.386000in}}%
\pgfpathclose%
\pgfusepath{fill}%
\end{pgfscope}%
\begin{pgfscope}%
\pgfpathrectangle{\pgfqpoint{3.815000in}{1.386000in}}{\pgfqpoint{2.835000in}{0.630000in}}%
\pgfusepath{clip}%
\pgfsetbuttcap%
\pgfsetmiterjoin%
\definecolor{currentfill}{rgb}{0.101961,0.101961,0.101961}%
\pgfsetfillcolor{currentfill}%
\pgfsetlinewidth{0.000000pt}%
\definecolor{currentstroke}{rgb}{0.000000,0.000000,0.000000}%
\pgfsetstrokecolor{currentstroke}%
\pgfsetstrokeopacity{0.000000}%
\pgfsetdash{}{0pt}%
\pgfpathmoveto{\pgfqpoint{4.312306in}{1.386000in}}%
\pgfpathlineto{\pgfqpoint{4.333569in}{1.386000in}}%
\pgfpathlineto{\pgfqpoint{4.333569in}{1.386000in}}%
\pgfpathlineto{\pgfqpoint{4.312306in}{1.386000in}}%
\pgfpathlineto{\pgfqpoint{4.312306in}{1.386000in}}%
\pgfpathclose%
\pgfusepath{fill}%
\end{pgfscope}%
\begin{pgfscope}%
\pgfpathrectangle{\pgfqpoint{3.815000in}{1.386000in}}{\pgfqpoint{2.835000in}{0.630000in}}%
\pgfusepath{clip}%
\pgfsetbuttcap%
\pgfsetmiterjoin%
\definecolor{currentfill}{rgb}{0.101961,0.101961,0.101961}%
\pgfsetfillcolor{currentfill}%
\pgfsetlinewidth{0.000000pt}%
\definecolor{currentstroke}{rgb}{0.000000,0.000000,0.000000}%
\pgfsetstrokecolor{currentstroke}%
\pgfsetstrokeopacity{0.000000}%
\pgfsetdash{}{0pt}%
\pgfpathmoveto{\pgfqpoint{4.335931in}{1.386000in}}%
\pgfpathlineto{\pgfqpoint{4.357194in}{1.386000in}}%
\pgfpathlineto{\pgfqpoint{4.357194in}{1.386000in}}%
\pgfpathlineto{\pgfqpoint{4.335931in}{1.386000in}}%
\pgfpathlineto{\pgfqpoint{4.335931in}{1.386000in}}%
\pgfpathclose%
\pgfusepath{fill}%
\end{pgfscope}%
\begin{pgfscope}%
\pgfpathrectangle{\pgfqpoint{3.815000in}{1.386000in}}{\pgfqpoint{2.835000in}{0.630000in}}%
\pgfusepath{clip}%
\pgfsetbuttcap%
\pgfsetmiterjoin%
\definecolor{currentfill}{rgb}{0.101961,0.101961,0.101961}%
\pgfsetfillcolor{currentfill}%
\pgfsetlinewidth{0.000000pt}%
\definecolor{currentstroke}{rgb}{0.000000,0.000000,0.000000}%
\pgfsetstrokecolor{currentstroke}%
\pgfsetstrokeopacity{0.000000}%
\pgfsetdash{}{0pt}%
\pgfpathmoveto{\pgfqpoint{4.359556in}{1.386000in}}%
\pgfpathlineto{\pgfqpoint{4.380819in}{1.386000in}}%
\pgfpathlineto{\pgfqpoint{4.380819in}{1.386000in}}%
\pgfpathlineto{\pgfqpoint{4.359556in}{1.386000in}}%
\pgfpathlineto{\pgfqpoint{4.359556in}{1.386000in}}%
\pgfpathclose%
\pgfusepath{fill}%
\end{pgfscope}%
\begin{pgfscope}%
\pgfpathrectangle{\pgfqpoint{3.815000in}{1.386000in}}{\pgfqpoint{2.835000in}{0.630000in}}%
\pgfusepath{clip}%
\pgfsetbuttcap%
\pgfsetmiterjoin%
\definecolor{currentfill}{rgb}{0.101961,0.101961,0.101961}%
\pgfsetfillcolor{currentfill}%
\pgfsetlinewidth{0.000000pt}%
\definecolor{currentstroke}{rgb}{0.000000,0.000000,0.000000}%
\pgfsetstrokecolor{currentstroke}%
\pgfsetstrokeopacity{0.000000}%
\pgfsetdash{}{0pt}%
\pgfpathmoveto{\pgfqpoint{4.383181in}{1.386000in}}%
\pgfpathlineto{\pgfqpoint{4.404444in}{1.386000in}}%
\pgfpathlineto{\pgfqpoint{4.404444in}{1.386000in}}%
\pgfpathlineto{\pgfqpoint{4.383181in}{1.386000in}}%
\pgfpathlineto{\pgfqpoint{4.383181in}{1.386000in}}%
\pgfpathclose%
\pgfusepath{fill}%
\end{pgfscope}%
\begin{pgfscope}%
\pgfpathrectangle{\pgfqpoint{3.815000in}{1.386000in}}{\pgfqpoint{2.835000in}{0.630000in}}%
\pgfusepath{clip}%
\pgfsetbuttcap%
\pgfsetmiterjoin%
\definecolor{currentfill}{rgb}{0.101961,0.101961,0.101961}%
\pgfsetfillcolor{currentfill}%
\pgfsetlinewidth{0.000000pt}%
\definecolor{currentstroke}{rgb}{0.000000,0.000000,0.000000}%
\pgfsetstrokecolor{currentstroke}%
\pgfsetstrokeopacity{0.000000}%
\pgfsetdash{}{0pt}%
\pgfpathmoveto{\pgfqpoint{4.406806in}{1.386000in}}%
\pgfpathlineto{\pgfqpoint{4.428069in}{1.386000in}}%
\pgfpathlineto{\pgfqpoint{4.428069in}{1.386000in}}%
\pgfpathlineto{\pgfqpoint{4.406806in}{1.386000in}}%
\pgfpathlineto{\pgfqpoint{4.406806in}{1.386000in}}%
\pgfpathclose%
\pgfusepath{fill}%
\end{pgfscope}%
\begin{pgfscope}%
\pgfpathrectangle{\pgfqpoint{3.815000in}{1.386000in}}{\pgfqpoint{2.835000in}{0.630000in}}%
\pgfusepath{clip}%
\pgfsetbuttcap%
\pgfsetmiterjoin%
\definecolor{currentfill}{rgb}{0.101961,0.101961,0.101961}%
\pgfsetfillcolor{currentfill}%
\pgfsetlinewidth{0.000000pt}%
\definecolor{currentstroke}{rgb}{0.000000,0.000000,0.000000}%
\pgfsetstrokecolor{currentstroke}%
\pgfsetstrokeopacity{0.000000}%
\pgfsetdash{}{0pt}%
\pgfpathmoveto{\pgfqpoint{4.430431in}{1.386000in}}%
\pgfpathlineto{\pgfqpoint{4.451694in}{1.386000in}}%
\pgfpathlineto{\pgfqpoint{4.451694in}{1.386000in}}%
\pgfpathlineto{\pgfqpoint{4.430431in}{1.386000in}}%
\pgfpathlineto{\pgfqpoint{4.430431in}{1.386000in}}%
\pgfpathclose%
\pgfusepath{fill}%
\end{pgfscope}%
\begin{pgfscope}%
\pgfpathrectangle{\pgfqpoint{3.815000in}{1.386000in}}{\pgfqpoint{2.835000in}{0.630000in}}%
\pgfusepath{clip}%
\pgfsetbuttcap%
\pgfsetmiterjoin%
\definecolor{currentfill}{rgb}{0.101961,0.101961,0.101961}%
\pgfsetfillcolor{currentfill}%
\pgfsetlinewidth{0.000000pt}%
\definecolor{currentstroke}{rgb}{0.000000,0.000000,0.000000}%
\pgfsetstrokecolor{currentstroke}%
\pgfsetstrokeopacity{0.000000}%
\pgfsetdash{}{0pt}%
\pgfpathmoveto{\pgfqpoint{4.454056in}{1.386000in}}%
\pgfpathlineto{\pgfqpoint{4.475319in}{1.386000in}}%
\pgfpathlineto{\pgfqpoint{4.475319in}{1.386000in}}%
\pgfpathlineto{\pgfqpoint{4.454056in}{1.386000in}}%
\pgfpathlineto{\pgfqpoint{4.454056in}{1.386000in}}%
\pgfpathclose%
\pgfusepath{fill}%
\end{pgfscope}%
\begin{pgfscope}%
\pgfpathrectangle{\pgfqpoint{3.815000in}{1.386000in}}{\pgfqpoint{2.835000in}{0.630000in}}%
\pgfusepath{clip}%
\pgfsetbuttcap%
\pgfsetmiterjoin%
\definecolor{currentfill}{rgb}{0.101961,0.101961,0.101961}%
\pgfsetfillcolor{currentfill}%
\pgfsetlinewidth{0.000000pt}%
\definecolor{currentstroke}{rgb}{0.000000,0.000000,0.000000}%
\pgfsetstrokecolor{currentstroke}%
\pgfsetstrokeopacity{0.000000}%
\pgfsetdash{}{0pt}%
\pgfpathmoveto{\pgfqpoint{4.477681in}{1.386000in}}%
\pgfpathlineto{\pgfqpoint{4.498944in}{1.386000in}}%
\pgfpathlineto{\pgfqpoint{4.498944in}{1.543500in}}%
\pgfpathlineto{\pgfqpoint{4.477681in}{1.543500in}}%
\pgfpathlineto{\pgfqpoint{4.477681in}{1.386000in}}%
\pgfpathclose%
\pgfusepath{fill}%
\end{pgfscope}%
\begin{pgfscope}%
\pgfpathrectangle{\pgfqpoint{3.815000in}{1.386000in}}{\pgfqpoint{2.835000in}{0.630000in}}%
\pgfusepath{clip}%
\pgfsetbuttcap%
\pgfsetmiterjoin%
\definecolor{currentfill}{rgb}{0.101961,0.101961,0.101961}%
\pgfsetfillcolor{currentfill}%
\pgfsetlinewidth{0.000000pt}%
\definecolor{currentstroke}{rgb}{0.000000,0.000000,0.000000}%
\pgfsetstrokecolor{currentstroke}%
\pgfsetstrokeopacity{0.000000}%
\pgfsetdash{}{0pt}%
\pgfpathmoveto{\pgfqpoint{4.501306in}{1.386000in}}%
\pgfpathlineto{\pgfqpoint{4.522569in}{1.386000in}}%
\pgfpathlineto{\pgfqpoint{4.522569in}{1.543500in}}%
\pgfpathlineto{\pgfqpoint{4.501306in}{1.543500in}}%
\pgfpathlineto{\pgfqpoint{4.501306in}{1.386000in}}%
\pgfpathclose%
\pgfusepath{fill}%
\end{pgfscope}%
\begin{pgfscope}%
\pgfpathrectangle{\pgfqpoint{3.815000in}{1.386000in}}{\pgfqpoint{2.835000in}{0.630000in}}%
\pgfusepath{clip}%
\pgfsetbuttcap%
\pgfsetmiterjoin%
\definecolor{currentfill}{rgb}{0.101961,0.101961,0.101961}%
\pgfsetfillcolor{currentfill}%
\pgfsetlinewidth{0.000000pt}%
\definecolor{currentstroke}{rgb}{0.000000,0.000000,0.000000}%
\pgfsetstrokecolor{currentstroke}%
\pgfsetstrokeopacity{0.000000}%
\pgfsetdash{}{0pt}%
\pgfpathmoveto{\pgfqpoint{4.524931in}{1.386000in}}%
\pgfpathlineto{\pgfqpoint{4.546194in}{1.386000in}}%
\pgfpathlineto{\pgfqpoint{4.546194in}{1.386000in}}%
\pgfpathlineto{\pgfqpoint{4.524931in}{1.386000in}}%
\pgfpathlineto{\pgfqpoint{4.524931in}{1.386000in}}%
\pgfpathclose%
\pgfusepath{fill}%
\end{pgfscope}%
\begin{pgfscope}%
\pgfpathrectangle{\pgfqpoint{3.815000in}{1.386000in}}{\pgfqpoint{2.835000in}{0.630000in}}%
\pgfusepath{clip}%
\pgfsetbuttcap%
\pgfsetmiterjoin%
\definecolor{currentfill}{rgb}{0.101961,0.101961,0.101961}%
\pgfsetfillcolor{currentfill}%
\pgfsetlinewidth{0.000000pt}%
\definecolor{currentstroke}{rgb}{0.000000,0.000000,0.000000}%
\pgfsetstrokecolor{currentstroke}%
\pgfsetstrokeopacity{0.000000}%
\pgfsetdash{}{0pt}%
\pgfpathmoveto{\pgfqpoint{4.548556in}{1.386000in}}%
\pgfpathlineto{\pgfqpoint{4.569819in}{1.386000in}}%
\pgfpathlineto{\pgfqpoint{4.569819in}{1.701000in}}%
\pgfpathlineto{\pgfqpoint{4.548556in}{1.701000in}}%
\pgfpathlineto{\pgfqpoint{4.548556in}{1.386000in}}%
\pgfpathclose%
\pgfusepath{fill}%
\end{pgfscope}%
\begin{pgfscope}%
\pgfpathrectangle{\pgfqpoint{3.815000in}{1.386000in}}{\pgfqpoint{2.835000in}{0.630000in}}%
\pgfusepath{clip}%
\pgfsetbuttcap%
\pgfsetmiterjoin%
\definecolor{currentfill}{rgb}{0.101961,0.101961,0.101961}%
\pgfsetfillcolor{currentfill}%
\pgfsetlinewidth{0.000000pt}%
\definecolor{currentstroke}{rgb}{0.000000,0.000000,0.000000}%
\pgfsetstrokecolor{currentstroke}%
\pgfsetstrokeopacity{0.000000}%
\pgfsetdash{}{0pt}%
\pgfpathmoveto{\pgfqpoint{4.572181in}{1.386000in}}%
\pgfpathlineto{\pgfqpoint{4.593444in}{1.386000in}}%
\pgfpathlineto{\pgfqpoint{4.593444in}{1.386000in}}%
\pgfpathlineto{\pgfqpoint{4.572181in}{1.386000in}}%
\pgfpathlineto{\pgfqpoint{4.572181in}{1.386000in}}%
\pgfpathclose%
\pgfusepath{fill}%
\end{pgfscope}%
\begin{pgfscope}%
\pgfpathrectangle{\pgfqpoint{3.815000in}{1.386000in}}{\pgfqpoint{2.835000in}{0.630000in}}%
\pgfusepath{clip}%
\pgfsetbuttcap%
\pgfsetmiterjoin%
\definecolor{currentfill}{rgb}{0.101961,0.101961,0.101961}%
\pgfsetfillcolor{currentfill}%
\pgfsetlinewidth{0.000000pt}%
\definecolor{currentstroke}{rgb}{0.000000,0.000000,0.000000}%
\pgfsetstrokecolor{currentstroke}%
\pgfsetstrokeopacity{0.000000}%
\pgfsetdash{}{0pt}%
\pgfpathmoveto{\pgfqpoint{4.595806in}{1.386000in}}%
\pgfpathlineto{\pgfqpoint{4.617069in}{1.386000in}}%
\pgfpathlineto{\pgfqpoint{4.617069in}{1.386000in}}%
\pgfpathlineto{\pgfqpoint{4.595806in}{1.386000in}}%
\pgfpathlineto{\pgfqpoint{4.595806in}{1.386000in}}%
\pgfpathclose%
\pgfusepath{fill}%
\end{pgfscope}%
\begin{pgfscope}%
\pgfpathrectangle{\pgfqpoint{3.815000in}{1.386000in}}{\pgfqpoint{2.835000in}{0.630000in}}%
\pgfusepath{clip}%
\pgfsetbuttcap%
\pgfsetmiterjoin%
\definecolor{currentfill}{rgb}{0.101961,0.101961,0.101961}%
\pgfsetfillcolor{currentfill}%
\pgfsetlinewidth{0.000000pt}%
\definecolor{currentstroke}{rgb}{0.000000,0.000000,0.000000}%
\pgfsetstrokecolor{currentstroke}%
\pgfsetstrokeopacity{0.000000}%
\pgfsetdash{}{0pt}%
\pgfpathmoveto{\pgfqpoint{4.619431in}{1.386000in}}%
\pgfpathlineto{\pgfqpoint{4.640694in}{1.386000in}}%
\pgfpathlineto{\pgfqpoint{4.640694in}{1.543500in}}%
\pgfpathlineto{\pgfqpoint{4.619431in}{1.543500in}}%
\pgfpathlineto{\pgfqpoint{4.619431in}{1.386000in}}%
\pgfpathclose%
\pgfusepath{fill}%
\end{pgfscope}%
\begin{pgfscope}%
\pgfpathrectangle{\pgfqpoint{3.815000in}{1.386000in}}{\pgfqpoint{2.835000in}{0.630000in}}%
\pgfusepath{clip}%
\pgfsetbuttcap%
\pgfsetmiterjoin%
\definecolor{currentfill}{rgb}{0.101961,0.101961,0.101961}%
\pgfsetfillcolor{currentfill}%
\pgfsetlinewidth{0.000000pt}%
\definecolor{currentstroke}{rgb}{0.000000,0.000000,0.000000}%
\pgfsetstrokecolor{currentstroke}%
\pgfsetstrokeopacity{0.000000}%
\pgfsetdash{}{0pt}%
\pgfpathmoveto{\pgfqpoint{4.643056in}{1.386000in}}%
\pgfpathlineto{\pgfqpoint{4.664319in}{1.386000in}}%
\pgfpathlineto{\pgfqpoint{4.664319in}{1.386000in}}%
\pgfpathlineto{\pgfqpoint{4.643056in}{1.386000in}}%
\pgfpathlineto{\pgfqpoint{4.643056in}{1.386000in}}%
\pgfpathclose%
\pgfusepath{fill}%
\end{pgfscope}%
\begin{pgfscope}%
\pgfpathrectangle{\pgfqpoint{3.815000in}{1.386000in}}{\pgfqpoint{2.835000in}{0.630000in}}%
\pgfusepath{clip}%
\pgfsetbuttcap%
\pgfsetmiterjoin%
\definecolor{currentfill}{rgb}{0.101961,0.101961,0.101961}%
\pgfsetfillcolor{currentfill}%
\pgfsetlinewidth{0.000000pt}%
\definecolor{currentstroke}{rgb}{0.000000,0.000000,0.000000}%
\pgfsetstrokecolor{currentstroke}%
\pgfsetstrokeopacity{0.000000}%
\pgfsetdash{}{0pt}%
\pgfpathmoveto{\pgfqpoint{4.666681in}{1.386000in}}%
\pgfpathlineto{\pgfqpoint{4.687944in}{1.386000in}}%
\pgfpathlineto{\pgfqpoint{4.687944in}{1.543500in}}%
\pgfpathlineto{\pgfqpoint{4.666681in}{1.543500in}}%
\pgfpathlineto{\pgfqpoint{4.666681in}{1.386000in}}%
\pgfpathclose%
\pgfusepath{fill}%
\end{pgfscope}%
\begin{pgfscope}%
\pgfpathrectangle{\pgfqpoint{3.815000in}{1.386000in}}{\pgfqpoint{2.835000in}{0.630000in}}%
\pgfusepath{clip}%
\pgfsetbuttcap%
\pgfsetmiterjoin%
\definecolor{currentfill}{rgb}{0.101961,0.101961,0.101961}%
\pgfsetfillcolor{currentfill}%
\pgfsetlinewidth{0.000000pt}%
\definecolor{currentstroke}{rgb}{0.000000,0.000000,0.000000}%
\pgfsetstrokecolor{currentstroke}%
\pgfsetstrokeopacity{0.000000}%
\pgfsetdash{}{0pt}%
\pgfpathmoveto{\pgfqpoint{4.690306in}{1.386000in}}%
\pgfpathlineto{\pgfqpoint{4.711569in}{1.386000in}}%
\pgfpathlineto{\pgfqpoint{4.711569in}{1.386000in}}%
\pgfpathlineto{\pgfqpoint{4.690306in}{1.386000in}}%
\pgfpathlineto{\pgfqpoint{4.690306in}{1.386000in}}%
\pgfpathclose%
\pgfusepath{fill}%
\end{pgfscope}%
\begin{pgfscope}%
\pgfpathrectangle{\pgfqpoint{3.815000in}{1.386000in}}{\pgfqpoint{2.835000in}{0.630000in}}%
\pgfusepath{clip}%
\pgfsetbuttcap%
\pgfsetmiterjoin%
\definecolor{currentfill}{rgb}{0.101961,0.101961,0.101961}%
\pgfsetfillcolor{currentfill}%
\pgfsetlinewidth{0.000000pt}%
\definecolor{currentstroke}{rgb}{0.000000,0.000000,0.000000}%
\pgfsetstrokecolor{currentstroke}%
\pgfsetstrokeopacity{0.000000}%
\pgfsetdash{}{0pt}%
\pgfpathmoveto{\pgfqpoint{4.713931in}{1.386000in}}%
\pgfpathlineto{\pgfqpoint{4.735194in}{1.386000in}}%
\pgfpathlineto{\pgfqpoint{4.735194in}{1.386000in}}%
\pgfpathlineto{\pgfqpoint{4.713931in}{1.386000in}}%
\pgfpathlineto{\pgfqpoint{4.713931in}{1.386000in}}%
\pgfpathclose%
\pgfusepath{fill}%
\end{pgfscope}%
\begin{pgfscope}%
\pgfpathrectangle{\pgfqpoint{3.815000in}{1.386000in}}{\pgfqpoint{2.835000in}{0.630000in}}%
\pgfusepath{clip}%
\pgfsetbuttcap%
\pgfsetmiterjoin%
\definecolor{currentfill}{rgb}{0.101961,0.101961,0.101961}%
\pgfsetfillcolor{currentfill}%
\pgfsetlinewidth{0.000000pt}%
\definecolor{currentstroke}{rgb}{0.000000,0.000000,0.000000}%
\pgfsetstrokecolor{currentstroke}%
\pgfsetstrokeopacity{0.000000}%
\pgfsetdash{}{0pt}%
\pgfpathmoveto{\pgfqpoint{4.737556in}{1.386000in}}%
\pgfpathlineto{\pgfqpoint{4.758819in}{1.386000in}}%
\pgfpathlineto{\pgfqpoint{4.758819in}{1.386000in}}%
\pgfpathlineto{\pgfqpoint{4.737556in}{1.386000in}}%
\pgfpathlineto{\pgfqpoint{4.737556in}{1.386000in}}%
\pgfpathclose%
\pgfusepath{fill}%
\end{pgfscope}%
\begin{pgfscope}%
\pgfpathrectangle{\pgfqpoint{3.815000in}{1.386000in}}{\pgfqpoint{2.835000in}{0.630000in}}%
\pgfusepath{clip}%
\pgfsetbuttcap%
\pgfsetmiterjoin%
\definecolor{currentfill}{rgb}{0.101961,0.101961,0.101961}%
\pgfsetfillcolor{currentfill}%
\pgfsetlinewidth{0.000000pt}%
\definecolor{currentstroke}{rgb}{0.000000,0.000000,0.000000}%
\pgfsetstrokecolor{currentstroke}%
\pgfsetstrokeopacity{0.000000}%
\pgfsetdash{}{0pt}%
\pgfpathmoveto{\pgfqpoint{4.761181in}{1.386000in}}%
\pgfpathlineto{\pgfqpoint{4.782444in}{1.386000in}}%
\pgfpathlineto{\pgfqpoint{4.782444in}{1.386000in}}%
\pgfpathlineto{\pgfqpoint{4.761181in}{1.386000in}}%
\pgfpathlineto{\pgfqpoint{4.761181in}{1.386000in}}%
\pgfpathclose%
\pgfusepath{fill}%
\end{pgfscope}%
\begin{pgfscope}%
\pgfpathrectangle{\pgfqpoint{3.815000in}{1.386000in}}{\pgfqpoint{2.835000in}{0.630000in}}%
\pgfusepath{clip}%
\pgfsetbuttcap%
\pgfsetmiterjoin%
\definecolor{currentfill}{rgb}{0.101961,0.101961,0.101961}%
\pgfsetfillcolor{currentfill}%
\pgfsetlinewidth{0.000000pt}%
\definecolor{currentstroke}{rgb}{0.000000,0.000000,0.000000}%
\pgfsetstrokecolor{currentstroke}%
\pgfsetstrokeopacity{0.000000}%
\pgfsetdash{}{0pt}%
\pgfpathmoveto{\pgfqpoint{4.784806in}{1.386000in}}%
\pgfpathlineto{\pgfqpoint{4.806069in}{1.386000in}}%
\pgfpathlineto{\pgfqpoint{4.806069in}{1.386000in}}%
\pgfpathlineto{\pgfqpoint{4.784806in}{1.386000in}}%
\pgfpathlineto{\pgfqpoint{4.784806in}{1.386000in}}%
\pgfpathclose%
\pgfusepath{fill}%
\end{pgfscope}%
\begin{pgfscope}%
\pgfpathrectangle{\pgfqpoint{3.815000in}{1.386000in}}{\pgfqpoint{2.835000in}{0.630000in}}%
\pgfusepath{clip}%
\pgfsetbuttcap%
\pgfsetmiterjoin%
\definecolor{currentfill}{rgb}{0.101961,0.101961,0.101961}%
\pgfsetfillcolor{currentfill}%
\pgfsetlinewidth{0.000000pt}%
\definecolor{currentstroke}{rgb}{0.000000,0.000000,0.000000}%
\pgfsetstrokecolor{currentstroke}%
\pgfsetstrokeopacity{0.000000}%
\pgfsetdash{}{0pt}%
\pgfpathmoveto{\pgfqpoint{4.808431in}{1.386000in}}%
\pgfpathlineto{\pgfqpoint{4.829694in}{1.386000in}}%
\pgfpathlineto{\pgfqpoint{4.829694in}{1.386000in}}%
\pgfpathlineto{\pgfqpoint{4.808431in}{1.386000in}}%
\pgfpathlineto{\pgfqpoint{4.808431in}{1.386000in}}%
\pgfpathclose%
\pgfusepath{fill}%
\end{pgfscope}%
\begin{pgfscope}%
\pgfpathrectangle{\pgfqpoint{3.815000in}{1.386000in}}{\pgfqpoint{2.835000in}{0.630000in}}%
\pgfusepath{clip}%
\pgfsetbuttcap%
\pgfsetmiterjoin%
\definecolor{currentfill}{rgb}{0.101961,0.101961,0.101961}%
\pgfsetfillcolor{currentfill}%
\pgfsetlinewidth{0.000000pt}%
\definecolor{currentstroke}{rgb}{0.000000,0.000000,0.000000}%
\pgfsetstrokecolor{currentstroke}%
\pgfsetstrokeopacity{0.000000}%
\pgfsetdash{}{0pt}%
\pgfpathmoveto{\pgfqpoint{4.832056in}{1.386000in}}%
\pgfpathlineto{\pgfqpoint{4.853319in}{1.386000in}}%
\pgfpathlineto{\pgfqpoint{4.853319in}{1.386000in}}%
\pgfpathlineto{\pgfqpoint{4.832056in}{1.386000in}}%
\pgfpathlineto{\pgfqpoint{4.832056in}{1.386000in}}%
\pgfpathclose%
\pgfusepath{fill}%
\end{pgfscope}%
\begin{pgfscope}%
\pgfpathrectangle{\pgfqpoint{3.815000in}{1.386000in}}{\pgfqpoint{2.835000in}{0.630000in}}%
\pgfusepath{clip}%
\pgfsetbuttcap%
\pgfsetmiterjoin%
\definecolor{currentfill}{rgb}{0.101961,0.101961,0.101961}%
\pgfsetfillcolor{currentfill}%
\pgfsetlinewidth{0.000000pt}%
\definecolor{currentstroke}{rgb}{0.000000,0.000000,0.000000}%
\pgfsetstrokecolor{currentstroke}%
\pgfsetstrokeopacity{0.000000}%
\pgfsetdash{}{0pt}%
\pgfpathmoveto{\pgfqpoint{4.855681in}{1.386000in}}%
\pgfpathlineto{\pgfqpoint{4.876944in}{1.386000in}}%
\pgfpathlineto{\pgfqpoint{4.876944in}{1.386000in}}%
\pgfpathlineto{\pgfqpoint{4.855681in}{1.386000in}}%
\pgfpathlineto{\pgfqpoint{4.855681in}{1.386000in}}%
\pgfpathclose%
\pgfusepath{fill}%
\end{pgfscope}%
\begin{pgfscope}%
\pgfpathrectangle{\pgfqpoint{3.815000in}{1.386000in}}{\pgfqpoint{2.835000in}{0.630000in}}%
\pgfusepath{clip}%
\pgfsetbuttcap%
\pgfsetmiterjoin%
\definecolor{currentfill}{rgb}{0.101961,0.101961,0.101961}%
\pgfsetfillcolor{currentfill}%
\pgfsetlinewidth{0.000000pt}%
\definecolor{currentstroke}{rgb}{0.000000,0.000000,0.000000}%
\pgfsetstrokecolor{currentstroke}%
\pgfsetstrokeopacity{0.000000}%
\pgfsetdash{}{0pt}%
\pgfpathmoveto{\pgfqpoint{4.879306in}{1.386000in}}%
\pgfpathlineto{\pgfqpoint{4.900569in}{1.386000in}}%
\pgfpathlineto{\pgfqpoint{4.900569in}{1.386000in}}%
\pgfpathlineto{\pgfqpoint{4.879306in}{1.386000in}}%
\pgfpathlineto{\pgfqpoint{4.879306in}{1.386000in}}%
\pgfpathclose%
\pgfusepath{fill}%
\end{pgfscope}%
\begin{pgfscope}%
\pgfpathrectangle{\pgfqpoint{3.815000in}{1.386000in}}{\pgfqpoint{2.835000in}{0.630000in}}%
\pgfusepath{clip}%
\pgfsetbuttcap%
\pgfsetmiterjoin%
\definecolor{currentfill}{rgb}{0.101961,0.101961,0.101961}%
\pgfsetfillcolor{currentfill}%
\pgfsetlinewidth{0.000000pt}%
\definecolor{currentstroke}{rgb}{0.000000,0.000000,0.000000}%
\pgfsetstrokecolor{currentstroke}%
\pgfsetstrokeopacity{0.000000}%
\pgfsetdash{}{0pt}%
\pgfpathmoveto{\pgfqpoint{4.902931in}{1.386000in}}%
\pgfpathlineto{\pgfqpoint{4.924194in}{1.386000in}}%
\pgfpathlineto{\pgfqpoint{4.924194in}{1.386000in}}%
\pgfpathlineto{\pgfqpoint{4.902931in}{1.386000in}}%
\pgfpathlineto{\pgfqpoint{4.902931in}{1.386000in}}%
\pgfpathclose%
\pgfusepath{fill}%
\end{pgfscope}%
\begin{pgfscope}%
\pgfpathrectangle{\pgfqpoint{3.815000in}{1.386000in}}{\pgfqpoint{2.835000in}{0.630000in}}%
\pgfusepath{clip}%
\pgfsetbuttcap%
\pgfsetmiterjoin%
\definecolor{currentfill}{rgb}{0.101961,0.101961,0.101961}%
\pgfsetfillcolor{currentfill}%
\pgfsetlinewidth{0.000000pt}%
\definecolor{currentstroke}{rgb}{0.000000,0.000000,0.000000}%
\pgfsetstrokecolor{currentstroke}%
\pgfsetstrokeopacity{0.000000}%
\pgfsetdash{}{0pt}%
\pgfpathmoveto{\pgfqpoint{4.926556in}{1.386000in}}%
\pgfpathlineto{\pgfqpoint{4.947819in}{1.386000in}}%
\pgfpathlineto{\pgfqpoint{4.947819in}{1.386000in}}%
\pgfpathlineto{\pgfqpoint{4.926556in}{1.386000in}}%
\pgfpathlineto{\pgfqpoint{4.926556in}{1.386000in}}%
\pgfpathclose%
\pgfusepath{fill}%
\end{pgfscope}%
\begin{pgfscope}%
\pgfpathrectangle{\pgfqpoint{3.815000in}{1.386000in}}{\pgfqpoint{2.835000in}{0.630000in}}%
\pgfusepath{clip}%
\pgfsetbuttcap%
\pgfsetmiterjoin%
\definecolor{currentfill}{rgb}{0.101961,0.101961,0.101961}%
\pgfsetfillcolor{currentfill}%
\pgfsetlinewidth{0.000000pt}%
\definecolor{currentstroke}{rgb}{0.000000,0.000000,0.000000}%
\pgfsetstrokecolor{currentstroke}%
\pgfsetstrokeopacity{0.000000}%
\pgfsetdash{}{0pt}%
\pgfpathmoveto{\pgfqpoint{4.950181in}{1.386000in}}%
\pgfpathlineto{\pgfqpoint{4.971444in}{1.386000in}}%
\pgfpathlineto{\pgfqpoint{4.971444in}{1.386000in}}%
\pgfpathlineto{\pgfqpoint{4.950181in}{1.386000in}}%
\pgfpathlineto{\pgfqpoint{4.950181in}{1.386000in}}%
\pgfpathclose%
\pgfusepath{fill}%
\end{pgfscope}%
\begin{pgfscope}%
\pgfpathrectangle{\pgfqpoint{3.815000in}{1.386000in}}{\pgfqpoint{2.835000in}{0.630000in}}%
\pgfusepath{clip}%
\pgfsetbuttcap%
\pgfsetmiterjoin%
\definecolor{currentfill}{rgb}{0.101961,0.101961,0.101961}%
\pgfsetfillcolor{currentfill}%
\pgfsetlinewidth{0.000000pt}%
\definecolor{currentstroke}{rgb}{0.000000,0.000000,0.000000}%
\pgfsetstrokecolor{currentstroke}%
\pgfsetstrokeopacity{0.000000}%
\pgfsetdash{}{0pt}%
\pgfpathmoveto{\pgfqpoint{4.973806in}{1.386000in}}%
\pgfpathlineto{\pgfqpoint{4.995069in}{1.386000in}}%
\pgfpathlineto{\pgfqpoint{4.995069in}{1.386000in}}%
\pgfpathlineto{\pgfqpoint{4.973806in}{1.386000in}}%
\pgfpathlineto{\pgfqpoint{4.973806in}{1.386000in}}%
\pgfpathclose%
\pgfusepath{fill}%
\end{pgfscope}%
\begin{pgfscope}%
\pgfpathrectangle{\pgfqpoint{3.815000in}{1.386000in}}{\pgfqpoint{2.835000in}{0.630000in}}%
\pgfusepath{clip}%
\pgfsetbuttcap%
\pgfsetmiterjoin%
\definecolor{currentfill}{rgb}{0.101961,0.101961,0.101961}%
\pgfsetfillcolor{currentfill}%
\pgfsetlinewidth{0.000000pt}%
\definecolor{currentstroke}{rgb}{0.000000,0.000000,0.000000}%
\pgfsetstrokecolor{currentstroke}%
\pgfsetstrokeopacity{0.000000}%
\pgfsetdash{}{0pt}%
\pgfpathmoveto{\pgfqpoint{4.997431in}{1.386000in}}%
\pgfpathlineto{\pgfqpoint{5.018694in}{1.386000in}}%
\pgfpathlineto{\pgfqpoint{5.018694in}{1.386000in}}%
\pgfpathlineto{\pgfqpoint{4.997431in}{1.386000in}}%
\pgfpathlineto{\pgfqpoint{4.997431in}{1.386000in}}%
\pgfpathclose%
\pgfusepath{fill}%
\end{pgfscope}%
\begin{pgfscope}%
\pgfpathrectangle{\pgfqpoint{3.815000in}{1.386000in}}{\pgfqpoint{2.835000in}{0.630000in}}%
\pgfusepath{clip}%
\pgfsetbuttcap%
\pgfsetmiterjoin%
\definecolor{currentfill}{rgb}{0.101961,0.101961,0.101961}%
\pgfsetfillcolor{currentfill}%
\pgfsetlinewidth{0.000000pt}%
\definecolor{currentstroke}{rgb}{0.000000,0.000000,0.000000}%
\pgfsetstrokecolor{currentstroke}%
\pgfsetstrokeopacity{0.000000}%
\pgfsetdash{}{0pt}%
\pgfpathmoveto{\pgfqpoint{5.021056in}{1.386000in}}%
\pgfpathlineto{\pgfqpoint{5.042319in}{1.386000in}}%
\pgfpathlineto{\pgfqpoint{5.042319in}{1.386000in}}%
\pgfpathlineto{\pgfqpoint{5.021056in}{1.386000in}}%
\pgfpathlineto{\pgfqpoint{5.021056in}{1.386000in}}%
\pgfpathclose%
\pgfusepath{fill}%
\end{pgfscope}%
\begin{pgfscope}%
\pgfpathrectangle{\pgfqpoint{3.815000in}{1.386000in}}{\pgfqpoint{2.835000in}{0.630000in}}%
\pgfusepath{clip}%
\pgfsetbuttcap%
\pgfsetmiterjoin%
\definecolor{currentfill}{rgb}{0.101961,0.101961,0.101961}%
\pgfsetfillcolor{currentfill}%
\pgfsetlinewidth{0.000000pt}%
\definecolor{currentstroke}{rgb}{0.000000,0.000000,0.000000}%
\pgfsetstrokecolor{currentstroke}%
\pgfsetstrokeopacity{0.000000}%
\pgfsetdash{}{0pt}%
\pgfpathmoveto{\pgfqpoint{5.044681in}{1.386000in}}%
\pgfpathlineto{\pgfqpoint{5.065944in}{1.386000in}}%
\pgfpathlineto{\pgfqpoint{5.065944in}{1.386000in}}%
\pgfpathlineto{\pgfqpoint{5.044681in}{1.386000in}}%
\pgfpathlineto{\pgfqpoint{5.044681in}{1.386000in}}%
\pgfpathclose%
\pgfusepath{fill}%
\end{pgfscope}%
\begin{pgfscope}%
\pgfpathrectangle{\pgfqpoint{3.815000in}{1.386000in}}{\pgfqpoint{2.835000in}{0.630000in}}%
\pgfusepath{clip}%
\pgfsetbuttcap%
\pgfsetmiterjoin%
\definecolor{currentfill}{rgb}{0.101961,0.101961,0.101961}%
\pgfsetfillcolor{currentfill}%
\pgfsetlinewidth{0.000000pt}%
\definecolor{currentstroke}{rgb}{0.000000,0.000000,0.000000}%
\pgfsetstrokecolor{currentstroke}%
\pgfsetstrokeopacity{0.000000}%
\pgfsetdash{}{0pt}%
\pgfpathmoveto{\pgfqpoint{5.068306in}{1.386000in}}%
\pgfpathlineto{\pgfqpoint{5.089569in}{1.386000in}}%
\pgfpathlineto{\pgfqpoint{5.089569in}{1.386000in}}%
\pgfpathlineto{\pgfqpoint{5.068306in}{1.386000in}}%
\pgfpathlineto{\pgfqpoint{5.068306in}{1.386000in}}%
\pgfpathclose%
\pgfusepath{fill}%
\end{pgfscope}%
\begin{pgfscope}%
\pgfpathrectangle{\pgfqpoint{3.815000in}{1.386000in}}{\pgfqpoint{2.835000in}{0.630000in}}%
\pgfusepath{clip}%
\pgfsetbuttcap%
\pgfsetmiterjoin%
\definecolor{currentfill}{rgb}{0.101961,0.101961,0.101961}%
\pgfsetfillcolor{currentfill}%
\pgfsetlinewidth{0.000000pt}%
\definecolor{currentstroke}{rgb}{0.000000,0.000000,0.000000}%
\pgfsetstrokecolor{currentstroke}%
\pgfsetstrokeopacity{0.000000}%
\pgfsetdash{}{0pt}%
\pgfpathmoveto{\pgfqpoint{5.091931in}{1.386000in}}%
\pgfpathlineto{\pgfqpoint{5.113194in}{1.386000in}}%
\pgfpathlineto{\pgfqpoint{5.113194in}{1.386000in}}%
\pgfpathlineto{\pgfqpoint{5.091931in}{1.386000in}}%
\pgfpathlineto{\pgfqpoint{5.091931in}{1.386000in}}%
\pgfpathclose%
\pgfusepath{fill}%
\end{pgfscope}%
\begin{pgfscope}%
\pgfpathrectangle{\pgfqpoint{3.815000in}{1.386000in}}{\pgfqpoint{2.835000in}{0.630000in}}%
\pgfusepath{clip}%
\pgfsetbuttcap%
\pgfsetmiterjoin%
\definecolor{currentfill}{rgb}{0.101961,0.101961,0.101961}%
\pgfsetfillcolor{currentfill}%
\pgfsetlinewidth{0.000000pt}%
\definecolor{currentstroke}{rgb}{0.000000,0.000000,0.000000}%
\pgfsetstrokecolor{currentstroke}%
\pgfsetstrokeopacity{0.000000}%
\pgfsetdash{}{0pt}%
\pgfpathmoveto{\pgfqpoint{5.115556in}{1.386000in}}%
\pgfpathlineto{\pgfqpoint{5.136819in}{1.386000in}}%
\pgfpathlineto{\pgfqpoint{5.136819in}{1.386000in}}%
\pgfpathlineto{\pgfqpoint{5.115556in}{1.386000in}}%
\pgfpathlineto{\pgfqpoint{5.115556in}{1.386000in}}%
\pgfpathclose%
\pgfusepath{fill}%
\end{pgfscope}%
\begin{pgfscope}%
\pgfpathrectangle{\pgfqpoint{3.815000in}{1.386000in}}{\pgfqpoint{2.835000in}{0.630000in}}%
\pgfusepath{clip}%
\pgfsetbuttcap%
\pgfsetmiterjoin%
\definecolor{currentfill}{rgb}{0.101961,0.101961,0.101961}%
\pgfsetfillcolor{currentfill}%
\pgfsetlinewidth{0.000000pt}%
\definecolor{currentstroke}{rgb}{0.000000,0.000000,0.000000}%
\pgfsetstrokecolor{currentstroke}%
\pgfsetstrokeopacity{0.000000}%
\pgfsetdash{}{0pt}%
\pgfpathmoveto{\pgfqpoint{5.139181in}{1.386000in}}%
\pgfpathlineto{\pgfqpoint{5.160444in}{1.386000in}}%
\pgfpathlineto{\pgfqpoint{5.160444in}{1.386000in}}%
\pgfpathlineto{\pgfqpoint{5.139181in}{1.386000in}}%
\pgfpathlineto{\pgfqpoint{5.139181in}{1.386000in}}%
\pgfpathclose%
\pgfusepath{fill}%
\end{pgfscope}%
\begin{pgfscope}%
\pgfpathrectangle{\pgfqpoint{3.815000in}{1.386000in}}{\pgfqpoint{2.835000in}{0.630000in}}%
\pgfusepath{clip}%
\pgfsetbuttcap%
\pgfsetmiterjoin%
\definecolor{currentfill}{rgb}{0.101961,0.101961,0.101961}%
\pgfsetfillcolor{currentfill}%
\pgfsetlinewidth{0.000000pt}%
\definecolor{currentstroke}{rgb}{0.000000,0.000000,0.000000}%
\pgfsetstrokecolor{currentstroke}%
\pgfsetstrokeopacity{0.000000}%
\pgfsetdash{}{0pt}%
\pgfpathmoveto{\pgfqpoint{5.162806in}{1.386000in}}%
\pgfpathlineto{\pgfqpoint{5.184069in}{1.386000in}}%
\pgfpathlineto{\pgfqpoint{5.184069in}{1.386000in}}%
\pgfpathlineto{\pgfqpoint{5.162806in}{1.386000in}}%
\pgfpathlineto{\pgfqpoint{5.162806in}{1.386000in}}%
\pgfpathclose%
\pgfusepath{fill}%
\end{pgfscope}%
\begin{pgfscope}%
\pgfpathrectangle{\pgfqpoint{3.815000in}{1.386000in}}{\pgfqpoint{2.835000in}{0.630000in}}%
\pgfusepath{clip}%
\pgfsetbuttcap%
\pgfsetmiterjoin%
\definecolor{currentfill}{rgb}{0.101961,0.101961,0.101961}%
\pgfsetfillcolor{currentfill}%
\pgfsetlinewidth{0.000000pt}%
\definecolor{currentstroke}{rgb}{0.000000,0.000000,0.000000}%
\pgfsetstrokecolor{currentstroke}%
\pgfsetstrokeopacity{0.000000}%
\pgfsetdash{}{0pt}%
\pgfpathmoveto{\pgfqpoint{5.186431in}{1.386000in}}%
\pgfpathlineto{\pgfqpoint{5.207694in}{1.386000in}}%
\pgfpathlineto{\pgfqpoint{5.207694in}{1.386000in}}%
\pgfpathlineto{\pgfqpoint{5.186431in}{1.386000in}}%
\pgfpathlineto{\pgfqpoint{5.186431in}{1.386000in}}%
\pgfpathclose%
\pgfusepath{fill}%
\end{pgfscope}%
\begin{pgfscope}%
\pgfpathrectangle{\pgfqpoint{3.815000in}{1.386000in}}{\pgfqpoint{2.835000in}{0.630000in}}%
\pgfusepath{clip}%
\pgfsetbuttcap%
\pgfsetmiterjoin%
\definecolor{currentfill}{rgb}{0.101961,0.101961,0.101961}%
\pgfsetfillcolor{currentfill}%
\pgfsetlinewidth{0.000000pt}%
\definecolor{currentstroke}{rgb}{0.000000,0.000000,0.000000}%
\pgfsetstrokecolor{currentstroke}%
\pgfsetstrokeopacity{0.000000}%
\pgfsetdash{}{0pt}%
\pgfpathmoveto{\pgfqpoint{5.210056in}{1.386000in}}%
\pgfpathlineto{\pgfqpoint{5.231319in}{1.386000in}}%
\pgfpathlineto{\pgfqpoint{5.231319in}{1.386000in}}%
\pgfpathlineto{\pgfqpoint{5.210056in}{1.386000in}}%
\pgfpathlineto{\pgfqpoint{5.210056in}{1.386000in}}%
\pgfpathclose%
\pgfusepath{fill}%
\end{pgfscope}%
\begin{pgfscope}%
\pgfpathrectangle{\pgfqpoint{3.815000in}{1.386000in}}{\pgfqpoint{2.835000in}{0.630000in}}%
\pgfusepath{clip}%
\pgfsetbuttcap%
\pgfsetmiterjoin%
\definecolor{currentfill}{rgb}{0.101961,0.101961,0.101961}%
\pgfsetfillcolor{currentfill}%
\pgfsetlinewidth{0.000000pt}%
\definecolor{currentstroke}{rgb}{0.000000,0.000000,0.000000}%
\pgfsetstrokecolor{currentstroke}%
\pgfsetstrokeopacity{0.000000}%
\pgfsetdash{}{0pt}%
\pgfpathmoveto{\pgfqpoint{5.233681in}{1.386000in}}%
\pgfpathlineto{\pgfqpoint{5.254944in}{1.386000in}}%
\pgfpathlineto{\pgfqpoint{5.254944in}{1.543500in}}%
\pgfpathlineto{\pgfqpoint{5.233681in}{1.543500in}}%
\pgfpathlineto{\pgfqpoint{5.233681in}{1.386000in}}%
\pgfpathclose%
\pgfusepath{fill}%
\end{pgfscope}%
\begin{pgfscope}%
\pgfpathrectangle{\pgfqpoint{3.815000in}{1.386000in}}{\pgfqpoint{2.835000in}{0.630000in}}%
\pgfusepath{clip}%
\pgfsetbuttcap%
\pgfsetmiterjoin%
\definecolor{currentfill}{rgb}{0.101961,0.101961,0.101961}%
\pgfsetfillcolor{currentfill}%
\pgfsetlinewidth{0.000000pt}%
\definecolor{currentstroke}{rgb}{0.000000,0.000000,0.000000}%
\pgfsetstrokecolor{currentstroke}%
\pgfsetstrokeopacity{0.000000}%
\pgfsetdash{}{0pt}%
\pgfpathmoveto{\pgfqpoint{5.257306in}{1.386000in}}%
\pgfpathlineto{\pgfqpoint{5.278569in}{1.386000in}}%
\pgfpathlineto{\pgfqpoint{5.278569in}{1.386000in}}%
\pgfpathlineto{\pgfqpoint{5.257306in}{1.386000in}}%
\pgfpathlineto{\pgfqpoint{5.257306in}{1.386000in}}%
\pgfpathclose%
\pgfusepath{fill}%
\end{pgfscope}%
\begin{pgfscope}%
\pgfpathrectangle{\pgfqpoint{3.815000in}{1.386000in}}{\pgfqpoint{2.835000in}{0.630000in}}%
\pgfusepath{clip}%
\pgfsetbuttcap%
\pgfsetmiterjoin%
\definecolor{currentfill}{rgb}{0.101961,0.101961,0.101961}%
\pgfsetfillcolor{currentfill}%
\pgfsetlinewidth{0.000000pt}%
\definecolor{currentstroke}{rgb}{0.000000,0.000000,0.000000}%
\pgfsetstrokecolor{currentstroke}%
\pgfsetstrokeopacity{0.000000}%
\pgfsetdash{}{0pt}%
\pgfpathmoveto{\pgfqpoint{5.280931in}{1.386000in}}%
\pgfpathlineto{\pgfqpoint{5.302194in}{1.386000in}}%
\pgfpathlineto{\pgfqpoint{5.302194in}{1.386000in}}%
\pgfpathlineto{\pgfqpoint{5.280931in}{1.386000in}}%
\pgfpathlineto{\pgfqpoint{5.280931in}{1.386000in}}%
\pgfpathclose%
\pgfusepath{fill}%
\end{pgfscope}%
\begin{pgfscope}%
\pgfpathrectangle{\pgfqpoint{3.815000in}{1.386000in}}{\pgfqpoint{2.835000in}{0.630000in}}%
\pgfusepath{clip}%
\pgfsetbuttcap%
\pgfsetmiterjoin%
\definecolor{currentfill}{rgb}{0.101961,0.101961,0.101961}%
\pgfsetfillcolor{currentfill}%
\pgfsetlinewidth{0.000000pt}%
\definecolor{currentstroke}{rgb}{0.000000,0.000000,0.000000}%
\pgfsetstrokecolor{currentstroke}%
\pgfsetstrokeopacity{0.000000}%
\pgfsetdash{}{0pt}%
\pgfpathmoveto{\pgfqpoint{5.304556in}{1.386000in}}%
\pgfpathlineto{\pgfqpoint{5.325819in}{1.386000in}}%
\pgfpathlineto{\pgfqpoint{5.325819in}{1.386000in}}%
\pgfpathlineto{\pgfqpoint{5.304556in}{1.386000in}}%
\pgfpathlineto{\pgfqpoint{5.304556in}{1.386000in}}%
\pgfpathclose%
\pgfusepath{fill}%
\end{pgfscope}%
\begin{pgfscope}%
\pgfpathrectangle{\pgfqpoint{3.815000in}{1.386000in}}{\pgfqpoint{2.835000in}{0.630000in}}%
\pgfusepath{clip}%
\pgfsetbuttcap%
\pgfsetmiterjoin%
\definecolor{currentfill}{rgb}{0.101961,0.101961,0.101961}%
\pgfsetfillcolor{currentfill}%
\pgfsetlinewidth{0.000000pt}%
\definecolor{currentstroke}{rgb}{0.000000,0.000000,0.000000}%
\pgfsetstrokecolor{currentstroke}%
\pgfsetstrokeopacity{0.000000}%
\pgfsetdash{}{0pt}%
\pgfpathmoveto{\pgfqpoint{5.328181in}{1.386000in}}%
\pgfpathlineto{\pgfqpoint{5.349444in}{1.386000in}}%
\pgfpathlineto{\pgfqpoint{5.349444in}{1.386000in}}%
\pgfpathlineto{\pgfqpoint{5.328181in}{1.386000in}}%
\pgfpathlineto{\pgfqpoint{5.328181in}{1.386000in}}%
\pgfpathclose%
\pgfusepath{fill}%
\end{pgfscope}%
\begin{pgfscope}%
\pgfpathrectangle{\pgfqpoint{3.815000in}{1.386000in}}{\pgfqpoint{2.835000in}{0.630000in}}%
\pgfusepath{clip}%
\pgfsetbuttcap%
\pgfsetmiterjoin%
\definecolor{currentfill}{rgb}{0.101961,0.101961,0.101961}%
\pgfsetfillcolor{currentfill}%
\pgfsetlinewidth{0.000000pt}%
\definecolor{currentstroke}{rgb}{0.000000,0.000000,0.000000}%
\pgfsetstrokecolor{currentstroke}%
\pgfsetstrokeopacity{0.000000}%
\pgfsetdash{}{0pt}%
\pgfpathmoveto{\pgfqpoint{5.351806in}{1.386000in}}%
\pgfpathlineto{\pgfqpoint{5.373069in}{1.386000in}}%
\pgfpathlineto{\pgfqpoint{5.373069in}{1.386000in}}%
\pgfpathlineto{\pgfqpoint{5.351806in}{1.386000in}}%
\pgfpathlineto{\pgfqpoint{5.351806in}{1.386000in}}%
\pgfpathclose%
\pgfusepath{fill}%
\end{pgfscope}%
\begin{pgfscope}%
\pgfpathrectangle{\pgfqpoint{3.815000in}{1.386000in}}{\pgfqpoint{2.835000in}{0.630000in}}%
\pgfusepath{clip}%
\pgfsetbuttcap%
\pgfsetmiterjoin%
\definecolor{currentfill}{rgb}{0.101961,0.101961,0.101961}%
\pgfsetfillcolor{currentfill}%
\pgfsetlinewidth{0.000000pt}%
\definecolor{currentstroke}{rgb}{0.000000,0.000000,0.000000}%
\pgfsetstrokecolor{currentstroke}%
\pgfsetstrokeopacity{0.000000}%
\pgfsetdash{}{0pt}%
\pgfpathmoveto{\pgfqpoint{5.375431in}{1.386000in}}%
\pgfpathlineto{\pgfqpoint{5.396694in}{1.386000in}}%
\pgfpathlineto{\pgfqpoint{5.396694in}{1.386000in}}%
\pgfpathlineto{\pgfqpoint{5.375431in}{1.386000in}}%
\pgfpathlineto{\pgfqpoint{5.375431in}{1.386000in}}%
\pgfpathclose%
\pgfusepath{fill}%
\end{pgfscope}%
\begin{pgfscope}%
\pgfpathrectangle{\pgfqpoint{3.815000in}{1.386000in}}{\pgfqpoint{2.835000in}{0.630000in}}%
\pgfusepath{clip}%
\pgfsetbuttcap%
\pgfsetmiterjoin%
\definecolor{currentfill}{rgb}{0.101961,0.101961,0.101961}%
\pgfsetfillcolor{currentfill}%
\pgfsetlinewidth{0.000000pt}%
\definecolor{currentstroke}{rgb}{0.000000,0.000000,0.000000}%
\pgfsetstrokecolor{currentstroke}%
\pgfsetstrokeopacity{0.000000}%
\pgfsetdash{}{0pt}%
\pgfpathmoveto{\pgfqpoint{5.399056in}{1.386000in}}%
\pgfpathlineto{\pgfqpoint{5.420319in}{1.386000in}}%
\pgfpathlineto{\pgfqpoint{5.420319in}{1.386000in}}%
\pgfpathlineto{\pgfqpoint{5.399056in}{1.386000in}}%
\pgfpathlineto{\pgfqpoint{5.399056in}{1.386000in}}%
\pgfpathclose%
\pgfusepath{fill}%
\end{pgfscope}%
\begin{pgfscope}%
\pgfpathrectangle{\pgfqpoint{3.815000in}{1.386000in}}{\pgfqpoint{2.835000in}{0.630000in}}%
\pgfusepath{clip}%
\pgfsetbuttcap%
\pgfsetmiterjoin%
\definecolor{currentfill}{rgb}{0.101961,0.101961,0.101961}%
\pgfsetfillcolor{currentfill}%
\pgfsetlinewidth{0.000000pt}%
\definecolor{currentstroke}{rgb}{0.000000,0.000000,0.000000}%
\pgfsetstrokecolor{currentstroke}%
\pgfsetstrokeopacity{0.000000}%
\pgfsetdash{}{0pt}%
\pgfpathmoveto{\pgfqpoint{5.422681in}{1.386000in}}%
\pgfpathlineto{\pgfqpoint{5.443944in}{1.386000in}}%
\pgfpathlineto{\pgfqpoint{5.443944in}{1.386000in}}%
\pgfpathlineto{\pgfqpoint{5.422681in}{1.386000in}}%
\pgfpathlineto{\pgfqpoint{5.422681in}{1.386000in}}%
\pgfpathclose%
\pgfusepath{fill}%
\end{pgfscope}%
\begin{pgfscope}%
\pgfpathrectangle{\pgfqpoint{3.815000in}{1.386000in}}{\pgfqpoint{2.835000in}{0.630000in}}%
\pgfusepath{clip}%
\pgfsetbuttcap%
\pgfsetmiterjoin%
\definecolor{currentfill}{rgb}{0.101961,0.101961,0.101961}%
\pgfsetfillcolor{currentfill}%
\pgfsetlinewidth{0.000000pt}%
\definecolor{currentstroke}{rgb}{0.000000,0.000000,0.000000}%
\pgfsetstrokecolor{currentstroke}%
\pgfsetstrokeopacity{0.000000}%
\pgfsetdash{}{0pt}%
\pgfpathmoveto{\pgfqpoint{5.446306in}{1.386000in}}%
\pgfpathlineto{\pgfqpoint{5.467569in}{1.386000in}}%
\pgfpathlineto{\pgfqpoint{5.467569in}{1.386000in}}%
\pgfpathlineto{\pgfqpoint{5.446306in}{1.386000in}}%
\pgfpathlineto{\pgfqpoint{5.446306in}{1.386000in}}%
\pgfpathclose%
\pgfusepath{fill}%
\end{pgfscope}%
\begin{pgfscope}%
\pgfpathrectangle{\pgfqpoint{3.815000in}{1.386000in}}{\pgfqpoint{2.835000in}{0.630000in}}%
\pgfusepath{clip}%
\pgfsetbuttcap%
\pgfsetmiterjoin%
\definecolor{currentfill}{rgb}{0.101961,0.101961,0.101961}%
\pgfsetfillcolor{currentfill}%
\pgfsetlinewidth{0.000000pt}%
\definecolor{currentstroke}{rgb}{0.000000,0.000000,0.000000}%
\pgfsetstrokecolor{currentstroke}%
\pgfsetstrokeopacity{0.000000}%
\pgfsetdash{}{0pt}%
\pgfpathmoveto{\pgfqpoint{5.469931in}{1.386000in}}%
\pgfpathlineto{\pgfqpoint{5.491194in}{1.386000in}}%
\pgfpathlineto{\pgfqpoint{5.491194in}{1.386000in}}%
\pgfpathlineto{\pgfqpoint{5.469931in}{1.386000in}}%
\pgfpathlineto{\pgfqpoint{5.469931in}{1.386000in}}%
\pgfpathclose%
\pgfusepath{fill}%
\end{pgfscope}%
\begin{pgfscope}%
\pgfpathrectangle{\pgfqpoint{3.815000in}{1.386000in}}{\pgfqpoint{2.835000in}{0.630000in}}%
\pgfusepath{clip}%
\pgfsetbuttcap%
\pgfsetmiterjoin%
\definecolor{currentfill}{rgb}{0.101961,0.101961,0.101961}%
\pgfsetfillcolor{currentfill}%
\pgfsetlinewidth{0.000000pt}%
\definecolor{currentstroke}{rgb}{0.000000,0.000000,0.000000}%
\pgfsetstrokecolor{currentstroke}%
\pgfsetstrokeopacity{0.000000}%
\pgfsetdash{}{0pt}%
\pgfpathmoveto{\pgfqpoint{5.493556in}{1.386000in}}%
\pgfpathlineto{\pgfqpoint{5.514819in}{1.386000in}}%
\pgfpathlineto{\pgfqpoint{5.514819in}{1.386000in}}%
\pgfpathlineto{\pgfqpoint{5.493556in}{1.386000in}}%
\pgfpathlineto{\pgfqpoint{5.493556in}{1.386000in}}%
\pgfpathclose%
\pgfusepath{fill}%
\end{pgfscope}%
\begin{pgfscope}%
\pgfpathrectangle{\pgfqpoint{3.815000in}{1.386000in}}{\pgfqpoint{2.835000in}{0.630000in}}%
\pgfusepath{clip}%
\pgfsetbuttcap%
\pgfsetmiterjoin%
\definecolor{currentfill}{rgb}{0.101961,0.101961,0.101961}%
\pgfsetfillcolor{currentfill}%
\pgfsetlinewidth{0.000000pt}%
\definecolor{currentstroke}{rgb}{0.000000,0.000000,0.000000}%
\pgfsetstrokecolor{currentstroke}%
\pgfsetstrokeopacity{0.000000}%
\pgfsetdash{}{0pt}%
\pgfpathmoveto{\pgfqpoint{5.517181in}{1.386000in}}%
\pgfpathlineto{\pgfqpoint{5.538444in}{1.386000in}}%
\pgfpathlineto{\pgfqpoint{5.538444in}{1.386000in}}%
\pgfpathlineto{\pgfqpoint{5.517181in}{1.386000in}}%
\pgfpathlineto{\pgfqpoint{5.517181in}{1.386000in}}%
\pgfpathclose%
\pgfusepath{fill}%
\end{pgfscope}%
\begin{pgfscope}%
\pgfpathrectangle{\pgfqpoint{3.815000in}{1.386000in}}{\pgfqpoint{2.835000in}{0.630000in}}%
\pgfusepath{clip}%
\pgfsetbuttcap%
\pgfsetmiterjoin%
\definecolor{currentfill}{rgb}{0.101961,0.101961,0.101961}%
\pgfsetfillcolor{currentfill}%
\pgfsetlinewidth{0.000000pt}%
\definecolor{currentstroke}{rgb}{0.000000,0.000000,0.000000}%
\pgfsetstrokecolor{currentstroke}%
\pgfsetstrokeopacity{0.000000}%
\pgfsetdash{}{0pt}%
\pgfpathmoveto{\pgfqpoint{5.540806in}{1.386000in}}%
\pgfpathlineto{\pgfqpoint{5.562069in}{1.386000in}}%
\pgfpathlineto{\pgfqpoint{5.562069in}{1.386000in}}%
\pgfpathlineto{\pgfqpoint{5.540806in}{1.386000in}}%
\pgfpathlineto{\pgfqpoint{5.540806in}{1.386000in}}%
\pgfpathclose%
\pgfusepath{fill}%
\end{pgfscope}%
\begin{pgfscope}%
\pgfpathrectangle{\pgfqpoint{3.815000in}{1.386000in}}{\pgfqpoint{2.835000in}{0.630000in}}%
\pgfusepath{clip}%
\pgfsetbuttcap%
\pgfsetmiterjoin%
\definecolor{currentfill}{rgb}{0.101961,0.101961,0.101961}%
\pgfsetfillcolor{currentfill}%
\pgfsetlinewidth{0.000000pt}%
\definecolor{currentstroke}{rgb}{0.000000,0.000000,0.000000}%
\pgfsetstrokecolor{currentstroke}%
\pgfsetstrokeopacity{0.000000}%
\pgfsetdash{}{0pt}%
\pgfpathmoveto{\pgfqpoint{5.564431in}{1.386000in}}%
\pgfpathlineto{\pgfqpoint{5.585694in}{1.386000in}}%
\pgfpathlineto{\pgfqpoint{5.585694in}{1.386000in}}%
\pgfpathlineto{\pgfqpoint{5.564431in}{1.386000in}}%
\pgfpathlineto{\pgfqpoint{5.564431in}{1.386000in}}%
\pgfpathclose%
\pgfusepath{fill}%
\end{pgfscope}%
\begin{pgfscope}%
\pgfpathrectangle{\pgfqpoint{3.815000in}{1.386000in}}{\pgfqpoint{2.835000in}{0.630000in}}%
\pgfusepath{clip}%
\pgfsetbuttcap%
\pgfsetmiterjoin%
\definecolor{currentfill}{rgb}{0.101961,0.101961,0.101961}%
\pgfsetfillcolor{currentfill}%
\pgfsetlinewidth{0.000000pt}%
\definecolor{currentstroke}{rgb}{0.000000,0.000000,0.000000}%
\pgfsetstrokecolor{currentstroke}%
\pgfsetstrokeopacity{0.000000}%
\pgfsetdash{}{0pt}%
\pgfpathmoveto{\pgfqpoint{5.588056in}{1.386000in}}%
\pgfpathlineto{\pgfqpoint{5.609319in}{1.386000in}}%
\pgfpathlineto{\pgfqpoint{5.609319in}{1.386000in}}%
\pgfpathlineto{\pgfqpoint{5.588056in}{1.386000in}}%
\pgfpathlineto{\pgfqpoint{5.588056in}{1.386000in}}%
\pgfpathclose%
\pgfusepath{fill}%
\end{pgfscope}%
\begin{pgfscope}%
\pgfpathrectangle{\pgfqpoint{3.815000in}{1.386000in}}{\pgfqpoint{2.835000in}{0.630000in}}%
\pgfusepath{clip}%
\pgfsetbuttcap%
\pgfsetmiterjoin%
\definecolor{currentfill}{rgb}{0.101961,0.101961,0.101961}%
\pgfsetfillcolor{currentfill}%
\pgfsetlinewidth{0.000000pt}%
\definecolor{currentstroke}{rgb}{0.000000,0.000000,0.000000}%
\pgfsetstrokecolor{currentstroke}%
\pgfsetstrokeopacity{0.000000}%
\pgfsetdash{}{0pt}%
\pgfpathmoveto{\pgfqpoint{5.611681in}{1.386000in}}%
\pgfpathlineto{\pgfqpoint{5.632944in}{1.386000in}}%
\pgfpathlineto{\pgfqpoint{5.632944in}{1.386000in}}%
\pgfpathlineto{\pgfqpoint{5.611681in}{1.386000in}}%
\pgfpathlineto{\pgfqpoint{5.611681in}{1.386000in}}%
\pgfpathclose%
\pgfusepath{fill}%
\end{pgfscope}%
\begin{pgfscope}%
\pgfpathrectangle{\pgfqpoint{3.815000in}{1.386000in}}{\pgfqpoint{2.835000in}{0.630000in}}%
\pgfusepath{clip}%
\pgfsetbuttcap%
\pgfsetmiterjoin%
\definecolor{currentfill}{rgb}{0.101961,0.101961,0.101961}%
\pgfsetfillcolor{currentfill}%
\pgfsetlinewidth{0.000000pt}%
\definecolor{currentstroke}{rgb}{0.000000,0.000000,0.000000}%
\pgfsetstrokecolor{currentstroke}%
\pgfsetstrokeopacity{0.000000}%
\pgfsetdash{}{0pt}%
\pgfpathmoveto{\pgfqpoint{5.635306in}{1.386000in}}%
\pgfpathlineto{\pgfqpoint{5.656569in}{1.386000in}}%
\pgfpathlineto{\pgfqpoint{5.656569in}{1.386000in}}%
\pgfpathlineto{\pgfqpoint{5.635306in}{1.386000in}}%
\pgfpathlineto{\pgfqpoint{5.635306in}{1.386000in}}%
\pgfpathclose%
\pgfusepath{fill}%
\end{pgfscope}%
\begin{pgfscope}%
\pgfpathrectangle{\pgfqpoint{3.815000in}{1.386000in}}{\pgfqpoint{2.835000in}{0.630000in}}%
\pgfusepath{clip}%
\pgfsetbuttcap%
\pgfsetmiterjoin%
\definecolor{currentfill}{rgb}{0.101961,0.101961,0.101961}%
\pgfsetfillcolor{currentfill}%
\pgfsetlinewidth{0.000000pt}%
\definecolor{currentstroke}{rgb}{0.000000,0.000000,0.000000}%
\pgfsetstrokecolor{currentstroke}%
\pgfsetstrokeopacity{0.000000}%
\pgfsetdash{}{0pt}%
\pgfpathmoveto{\pgfqpoint{5.658931in}{1.386000in}}%
\pgfpathlineto{\pgfqpoint{5.680194in}{1.386000in}}%
\pgfpathlineto{\pgfqpoint{5.680194in}{1.386000in}}%
\pgfpathlineto{\pgfqpoint{5.658931in}{1.386000in}}%
\pgfpathlineto{\pgfqpoint{5.658931in}{1.386000in}}%
\pgfpathclose%
\pgfusepath{fill}%
\end{pgfscope}%
\begin{pgfscope}%
\pgfpathrectangle{\pgfqpoint{3.815000in}{1.386000in}}{\pgfqpoint{2.835000in}{0.630000in}}%
\pgfusepath{clip}%
\pgfsetbuttcap%
\pgfsetmiterjoin%
\definecolor{currentfill}{rgb}{0.101961,0.101961,0.101961}%
\pgfsetfillcolor{currentfill}%
\pgfsetlinewidth{0.000000pt}%
\definecolor{currentstroke}{rgb}{0.000000,0.000000,0.000000}%
\pgfsetstrokecolor{currentstroke}%
\pgfsetstrokeopacity{0.000000}%
\pgfsetdash{}{0pt}%
\pgfpathmoveto{\pgfqpoint{5.682556in}{1.386000in}}%
\pgfpathlineto{\pgfqpoint{5.703819in}{1.386000in}}%
\pgfpathlineto{\pgfqpoint{5.703819in}{1.386000in}}%
\pgfpathlineto{\pgfqpoint{5.682556in}{1.386000in}}%
\pgfpathlineto{\pgfqpoint{5.682556in}{1.386000in}}%
\pgfpathclose%
\pgfusepath{fill}%
\end{pgfscope}%
\begin{pgfscope}%
\pgfpathrectangle{\pgfqpoint{3.815000in}{1.386000in}}{\pgfqpoint{2.835000in}{0.630000in}}%
\pgfusepath{clip}%
\pgfsetbuttcap%
\pgfsetmiterjoin%
\definecolor{currentfill}{rgb}{0.101961,0.101961,0.101961}%
\pgfsetfillcolor{currentfill}%
\pgfsetlinewidth{0.000000pt}%
\definecolor{currentstroke}{rgb}{0.000000,0.000000,0.000000}%
\pgfsetstrokecolor{currentstroke}%
\pgfsetstrokeopacity{0.000000}%
\pgfsetdash{}{0pt}%
\pgfpathmoveto{\pgfqpoint{5.706181in}{1.386000in}}%
\pgfpathlineto{\pgfqpoint{5.727444in}{1.386000in}}%
\pgfpathlineto{\pgfqpoint{5.727444in}{1.386000in}}%
\pgfpathlineto{\pgfqpoint{5.706181in}{1.386000in}}%
\pgfpathlineto{\pgfqpoint{5.706181in}{1.386000in}}%
\pgfpathclose%
\pgfusepath{fill}%
\end{pgfscope}%
\begin{pgfscope}%
\pgfpathrectangle{\pgfqpoint{3.815000in}{1.386000in}}{\pgfqpoint{2.835000in}{0.630000in}}%
\pgfusepath{clip}%
\pgfsetbuttcap%
\pgfsetmiterjoin%
\definecolor{currentfill}{rgb}{0.101961,0.101961,0.101961}%
\pgfsetfillcolor{currentfill}%
\pgfsetlinewidth{0.000000pt}%
\definecolor{currentstroke}{rgb}{0.000000,0.000000,0.000000}%
\pgfsetstrokecolor{currentstroke}%
\pgfsetstrokeopacity{0.000000}%
\pgfsetdash{}{0pt}%
\pgfpathmoveto{\pgfqpoint{5.729806in}{1.386000in}}%
\pgfpathlineto{\pgfqpoint{5.751069in}{1.386000in}}%
\pgfpathlineto{\pgfqpoint{5.751069in}{1.386000in}}%
\pgfpathlineto{\pgfqpoint{5.729806in}{1.386000in}}%
\pgfpathlineto{\pgfqpoint{5.729806in}{1.386000in}}%
\pgfpathclose%
\pgfusepath{fill}%
\end{pgfscope}%
\begin{pgfscope}%
\pgfpathrectangle{\pgfqpoint{3.815000in}{1.386000in}}{\pgfqpoint{2.835000in}{0.630000in}}%
\pgfusepath{clip}%
\pgfsetbuttcap%
\pgfsetmiterjoin%
\definecolor{currentfill}{rgb}{0.101961,0.101961,0.101961}%
\pgfsetfillcolor{currentfill}%
\pgfsetlinewidth{0.000000pt}%
\definecolor{currentstroke}{rgb}{0.000000,0.000000,0.000000}%
\pgfsetstrokecolor{currentstroke}%
\pgfsetstrokeopacity{0.000000}%
\pgfsetdash{}{0pt}%
\pgfpathmoveto{\pgfqpoint{5.753431in}{1.386000in}}%
\pgfpathlineto{\pgfqpoint{5.774694in}{1.386000in}}%
\pgfpathlineto{\pgfqpoint{5.774694in}{1.386000in}}%
\pgfpathlineto{\pgfqpoint{5.753431in}{1.386000in}}%
\pgfpathlineto{\pgfqpoint{5.753431in}{1.386000in}}%
\pgfpathclose%
\pgfusepath{fill}%
\end{pgfscope}%
\begin{pgfscope}%
\pgfpathrectangle{\pgfqpoint{3.815000in}{1.386000in}}{\pgfqpoint{2.835000in}{0.630000in}}%
\pgfusepath{clip}%
\pgfsetbuttcap%
\pgfsetmiterjoin%
\definecolor{currentfill}{rgb}{0.101961,0.101961,0.101961}%
\pgfsetfillcolor{currentfill}%
\pgfsetlinewidth{0.000000pt}%
\definecolor{currentstroke}{rgb}{0.000000,0.000000,0.000000}%
\pgfsetstrokecolor{currentstroke}%
\pgfsetstrokeopacity{0.000000}%
\pgfsetdash{}{0pt}%
\pgfpathmoveto{\pgfqpoint{5.777056in}{1.386000in}}%
\pgfpathlineto{\pgfqpoint{5.798319in}{1.386000in}}%
\pgfpathlineto{\pgfqpoint{5.798319in}{1.386000in}}%
\pgfpathlineto{\pgfqpoint{5.777056in}{1.386000in}}%
\pgfpathlineto{\pgfqpoint{5.777056in}{1.386000in}}%
\pgfpathclose%
\pgfusepath{fill}%
\end{pgfscope}%
\begin{pgfscope}%
\pgfpathrectangle{\pgfqpoint{3.815000in}{1.386000in}}{\pgfqpoint{2.835000in}{0.630000in}}%
\pgfusepath{clip}%
\pgfsetbuttcap%
\pgfsetmiterjoin%
\definecolor{currentfill}{rgb}{0.101961,0.101961,0.101961}%
\pgfsetfillcolor{currentfill}%
\pgfsetlinewidth{0.000000pt}%
\definecolor{currentstroke}{rgb}{0.000000,0.000000,0.000000}%
\pgfsetstrokecolor{currentstroke}%
\pgfsetstrokeopacity{0.000000}%
\pgfsetdash{}{0pt}%
\pgfpathmoveto{\pgfqpoint{5.800681in}{1.386000in}}%
\pgfpathlineto{\pgfqpoint{5.821944in}{1.386000in}}%
\pgfpathlineto{\pgfqpoint{5.821944in}{1.386000in}}%
\pgfpathlineto{\pgfqpoint{5.800681in}{1.386000in}}%
\pgfpathlineto{\pgfqpoint{5.800681in}{1.386000in}}%
\pgfpathclose%
\pgfusepath{fill}%
\end{pgfscope}%
\begin{pgfscope}%
\pgfpathrectangle{\pgfqpoint{3.815000in}{1.386000in}}{\pgfqpoint{2.835000in}{0.630000in}}%
\pgfusepath{clip}%
\pgfsetbuttcap%
\pgfsetmiterjoin%
\definecolor{currentfill}{rgb}{0.101961,0.101961,0.101961}%
\pgfsetfillcolor{currentfill}%
\pgfsetlinewidth{0.000000pt}%
\definecolor{currentstroke}{rgb}{0.000000,0.000000,0.000000}%
\pgfsetstrokecolor{currentstroke}%
\pgfsetstrokeopacity{0.000000}%
\pgfsetdash{}{0pt}%
\pgfpathmoveto{\pgfqpoint{5.824306in}{1.386000in}}%
\pgfpathlineto{\pgfqpoint{5.845569in}{1.386000in}}%
\pgfpathlineto{\pgfqpoint{5.845569in}{1.543500in}}%
\pgfpathlineto{\pgfqpoint{5.824306in}{1.543500in}}%
\pgfpathlineto{\pgfqpoint{5.824306in}{1.386000in}}%
\pgfpathclose%
\pgfusepath{fill}%
\end{pgfscope}%
\begin{pgfscope}%
\pgfpathrectangle{\pgfqpoint{3.815000in}{1.386000in}}{\pgfqpoint{2.835000in}{0.630000in}}%
\pgfusepath{clip}%
\pgfsetbuttcap%
\pgfsetmiterjoin%
\definecolor{currentfill}{rgb}{0.101961,0.101961,0.101961}%
\pgfsetfillcolor{currentfill}%
\pgfsetlinewidth{0.000000pt}%
\definecolor{currentstroke}{rgb}{0.000000,0.000000,0.000000}%
\pgfsetstrokecolor{currentstroke}%
\pgfsetstrokeopacity{0.000000}%
\pgfsetdash{}{0pt}%
\pgfpathmoveto{\pgfqpoint{5.847931in}{1.386000in}}%
\pgfpathlineto{\pgfqpoint{5.869194in}{1.386000in}}%
\pgfpathlineto{\pgfqpoint{5.869194in}{1.386000in}}%
\pgfpathlineto{\pgfqpoint{5.847931in}{1.386000in}}%
\pgfpathlineto{\pgfqpoint{5.847931in}{1.386000in}}%
\pgfpathclose%
\pgfusepath{fill}%
\end{pgfscope}%
\begin{pgfscope}%
\pgfpathrectangle{\pgfqpoint{3.815000in}{1.386000in}}{\pgfqpoint{2.835000in}{0.630000in}}%
\pgfusepath{clip}%
\pgfsetbuttcap%
\pgfsetmiterjoin%
\definecolor{currentfill}{rgb}{0.101961,0.101961,0.101961}%
\pgfsetfillcolor{currentfill}%
\pgfsetlinewidth{0.000000pt}%
\definecolor{currentstroke}{rgb}{0.000000,0.000000,0.000000}%
\pgfsetstrokecolor{currentstroke}%
\pgfsetstrokeopacity{0.000000}%
\pgfsetdash{}{0pt}%
\pgfpathmoveto{\pgfqpoint{5.871556in}{1.386000in}}%
\pgfpathlineto{\pgfqpoint{5.892819in}{1.386000in}}%
\pgfpathlineto{\pgfqpoint{5.892819in}{1.386000in}}%
\pgfpathlineto{\pgfqpoint{5.871556in}{1.386000in}}%
\pgfpathlineto{\pgfqpoint{5.871556in}{1.386000in}}%
\pgfpathclose%
\pgfusepath{fill}%
\end{pgfscope}%
\begin{pgfscope}%
\pgfpathrectangle{\pgfqpoint{3.815000in}{1.386000in}}{\pgfqpoint{2.835000in}{0.630000in}}%
\pgfusepath{clip}%
\pgfsetbuttcap%
\pgfsetmiterjoin%
\definecolor{currentfill}{rgb}{0.101961,0.101961,0.101961}%
\pgfsetfillcolor{currentfill}%
\pgfsetlinewidth{0.000000pt}%
\definecolor{currentstroke}{rgb}{0.000000,0.000000,0.000000}%
\pgfsetstrokecolor{currentstroke}%
\pgfsetstrokeopacity{0.000000}%
\pgfsetdash{}{0pt}%
\pgfpathmoveto{\pgfqpoint{5.895181in}{1.386000in}}%
\pgfpathlineto{\pgfqpoint{5.916444in}{1.386000in}}%
\pgfpathlineto{\pgfqpoint{5.916444in}{1.543500in}}%
\pgfpathlineto{\pgfqpoint{5.895181in}{1.543500in}}%
\pgfpathlineto{\pgfqpoint{5.895181in}{1.386000in}}%
\pgfpathclose%
\pgfusepath{fill}%
\end{pgfscope}%
\begin{pgfscope}%
\pgfpathrectangle{\pgfqpoint{3.815000in}{1.386000in}}{\pgfqpoint{2.835000in}{0.630000in}}%
\pgfusepath{clip}%
\pgfsetbuttcap%
\pgfsetmiterjoin%
\definecolor{currentfill}{rgb}{0.101961,0.101961,0.101961}%
\pgfsetfillcolor{currentfill}%
\pgfsetlinewidth{0.000000pt}%
\definecolor{currentstroke}{rgb}{0.000000,0.000000,0.000000}%
\pgfsetstrokecolor{currentstroke}%
\pgfsetstrokeopacity{0.000000}%
\pgfsetdash{}{0pt}%
\pgfpathmoveto{\pgfqpoint{5.918806in}{1.386000in}}%
\pgfpathlineto{\pgfqpoint{5.940069in}{1.386000in}}%
\pgfpathlineto{\pgfqpoint{5.940069in}{1.386000in}}%
\pgfpathlineto{\pgfqpoint{5.918806in}{1.386000in}}%
\pgfpathlineto{\pgfqpoint{5.918806in}{1.386000in}}%
\pgfpathclose%
\pgfusepath{fill}%
\end{pgfscope}%
\begin{pgfscope}%
\pgfpathrectangle{\pgfqpoint{3.815000in}{1.386000in}}{\pgfqpoint{2.835000in}{0.630000in}}%
\pgfusepath{clip}%
\pgfsetbuttcap%
\pgfsetmiterjoin%
\definecolor{currentfill}{rgb}{0.101961,0.101961,0.101961}%
\pgfsetfillcolor{currentfill}%
\pgfsetlinewidth{0.000000pt}%
\definecolor{currentstroke}{rgb}{0.000000,0.000000,0.000000}%
\pgfsetstrokecolor{currentstroke}%
\pgfsetstrokeopacity{0.000000}%
\pgfsetdash{}{0pt}%
\pgfpathmoveto{\pgfqpoint{5.942431in}{1.386000in}}%
\pgfpathlineto{\pgfqpoint{5.963694in}{1.386000in}}%
\pgfpathlineto{\pgfqpoint{5.963694in}{1.386000in}}%
\pgfpathlineto{\pgfqpoint{5.942431in}{1.386000in}}%
\pgfpathlineto{\pgfqpoint{5.942431in}{1.386000in}}%
\pgfpathclose%
\pgfusepath{fill}%
\end{pgfscope}%
\begin{pgfscope}%
\pgfpathrectangle{\pgfqpoint{3.815000in}{1.386000in}}{\pgfqpoint{2.835000in}{0.630000in}}%
\pgfusepath{clip}%
\pgfsetbuttcap%
\pgfsetmiterjoin%
\definecolor{currentfill}{rgb}{0.101961,0.101961,0.101961}%
\pgfsetfillcolor{currentfill}%
\pgfsetlinewidth{0.000000pt}%
\definecolor{currentstroke}{rgb}{0.000000,0.000000,0.000000}%
\pgfsetstrokecolor{currentstroke}%
\pgfsetstrokeopacity{0.000000}%
\pgfsetdash{}{0pt}%
\pgfpathmoveto{\pgfqpoint{5.966056in}{1.386000in}}%
\pgfpathlineto{\pgfqpoint{5.987319in}{1.386000in}}%
\pgfpathlineto{\pgfqpoint{5.987319in}{1.386000in}}%
\pgfpathlineto{\pgfqpoint{5.966056in}{1.386000in}}%
\pgfpathlineto{\pgfqpoint{5.966056in}{1.386000in}}%
\pgfpathclose%
\pgfusepath{fill}%
\end{pgfscope}%
\begin{pgfscope}%
\pgfpathrectangle{\pgfqpoint{3.815000in}{1.386000in}}{\pgfqpoint{2.835000in}{0.630000in}}%
\pgfusepath{clip}%
\pgfsetbuttcap%
\pgfsetmiterjoin%
\definecolor{currentfill}{rgb}{0.101961,0.101961,0.101961}%
\pgfsetfillcolor{currentfill}%
\pgfsetlinewidth{0.000000pt}%
\definecolor{currentstroke}{rgb}{0.000000,0.000000,0.000000}%
\pgfsetstrokecolor{currentstroke}%
\pgfsetstrokeopacity{0.000000}%
\pgfsetdash{}{0pt}%
\pgfpathmoveto{\pgfqpoint{5.989681in}{1.386000in}}%
\pgfpathlineto{\pgfqpoint{6.010944in}{1.386000in}}%
\pgfpathlineto{\pgfqpoint{6.010944in}{1.386000in}}%
\pgfpathlineto{\pgfqpoint{5.989681in}{1.386000in}}%
\pgfpathlineto{\pgfqpoint{5.989681in}{1.386000in}}%
\pgfpathclose%
\pgfusepath{fill}%
\end{pgfscope}%
\begin{pgfscope}%
\pgfpathrectangle{\pgfqpoint{3.815000in}{1.386000in}}{\pgfqpoint{2.835000in}{0.630000in}}%
\pgfusepath{clip}%
\pgfsetbuttcap%
\pgfsetmiterjoin%
\definecolor{currentfill}{rgb}{0.101961,0.101961,0.101961}%
\pgfsetfillcolor{currentfill}%
\pgfsetlinewidth{0.000000pt}%
\definecolor{currentstroke}{rgb}{0.000000,0.000000,0.000000}%
\pgfsetstrokecolor{currentstroke}%
\pgfsetstrokeopacity{0.000000}%
\pgfsetdash{}{0pt}%
\pgfpathmoveto{\pgfqpoint{6.013306in}{1.386000in}}%
\pgfpathlineto{\pgfqpoint{6.034569in}{1.386000in}}%
\pgfpathlineto{\pgfqpoint{6.034569in}{1.386000in}}%
\pgfpathlineto{\pgfqpoint{6.013306in}{1.386000in}}%
\pgfpathlineto{\pgfqpoint{6.013306in}{1.386000in}}%
\pgfpathclose%
\pgfusepath{fill}%
\end{pgfscope}%
\begin{pgfscope}%
\pgfpathrectangle{\pgfqpoint{3.815000in}{1.386000in}}{\pgfqpoint{2.835000in}{0.630000in}}%
\pgfusepath{clip}%
\pgfsetbuttcap%
\pgfsetmiterjoin%
\definecolor{currentfill}{rgb}{0.101961,0.101961,0.101961}%
\pgfsetfillcolor{currentfill}%
\pgfsetlinewidth{0.000000pt}%
\definecolor{currentstroke}{rgb}{0.000000,0.000000,0.000000}%
\pgfsetstrokecolor{currentstroke}%
\pgfsetstrokeopacity{0.000000}%
\pgfsetdash{}{0pt}%
\pgfpathmoveto{\pgfqpoint{6.036931in}{1.386000in}}%
\pgfpathlineto{\pgfqpoint{6.058194in}{1.386000in}}%
\pgfpathlineto{\pgfqpoint{6.058194in}{1.386000in}}%
\pgfpathlineto{\pgfqpoint{6.036931in}{1.386000in}}%
\pgfpathlineto{\pgfqpoint{6.036931in}{1.386000in}}%
\pgfpathclose%
\pgfusepath{fill}%
\end{pgfscope}%
\begin{pgfscope}%
\pgfpathrectangle{\pgfqpoint{3.815000in}{1.386000in}}{\pgfqpoint{2.835000in}{0.630000in}}%
\pgfusepath{clip}%
\pgfsetbuttcap%
\pgfsetmiterjoin%
\definecolor{currentfill}{rgb}{0.101961,0.101961,0.101961}%
\pgfsetfillcolor{currentfill}%
\pgfsetlinewidth{0.000000pt}%
\definecolor{currentstroke}{rgb}{0.000000,0.000000,0.000000}%
\pgfsetstrokecolor{currentstroke}%
\pgfsetstrokeopacity{0.000000}%
\pgfsetdash{}{0pt}%
\pgfpathmoveto{\pgfqpoint{6.060556in}{1.386000in}}%
\pgfpathlineto{\pgfqpoint{6.081819in}{1.386000in}}%
\pgfpathlineto{\pgfqpoint{6.081819in}{1.386000in}}%
\pgfpathlineto{\pgfqpoint{6.060556in}{1.386000in}}%
\pgfpathlineto{\pgfqpoint{6.060556in}{1.386000in}}%
\pgfpathclose%
\pgfusepath{fill}%
\end{pgfscope}%
\begin{pgfscope}%
\pgfpathrectangle{\pgfqpoint{3.815000in}{1.386000in}}{\pgfqpoint{2.835000in}{0.630000in}}%
\pgfusepath{clip}%
\pgfsetbuttcap%
\pgfsetmiterjoin%
\definecolor{currentfill}{rgb}{0.101961,0.101961,0.101961}%
\pgfsetfillcolor{currentfill}%
\pgfsetlinewidth{0.000000pt}%
\definecolor{currentstroke}{rgb}{0.000000,0.000000,0.000000}%
\pgfsetstrokecolor{currentstroke}%
\pgfsetstrokeopacity{0.000000}%
\pgfsetdash{}{0pt}%
\pgfpathmoveto{\pgfqpoint{6.084181in}{1.386000in}}%
\pgfpathlineto{\pgfqpoint{6.105444in}{1.386000in}}%
\pgfpathlineto{\pgfqpoint{6.105444in}{1.701000in}}%
\pgfpathlineto{\pgfqpoint{6.084181in}{1.701000in}}%
\pgfpathlineto{\pgfqpoint{6.084181in}{1.386000in}}%
\pgfpathclose%
\pgfusepath{fill}%
\end{pgfscope}%
\begin{pgfscope}%
\pgfpathrectangle{\pgfqpoint{3.815000in}{1.386000in}}{\pgfqpoint{2.835000in}{0.630000in}}%
\pgfusepath{clip}%
\pgfsetbuttcap%
\pgfsetmiterjoin%
\definecolor{currentfill}{rgb}{0.101961,0.101961,0.101961}%
\pgfsetfillcolor{currentfill}%
\pgfsetlinewidth{0.000000pt}%
\definecolor{currentstroke}{rgb}{0.000000,0.000000,0.000000}%
\pgfsetstrokecolor{currentstroke}%
\pgfsetstrokeopacity{0.000000}%
\pgfsetdash{}{0pt}%
\pgfpathmoveto{\pgfqpoint{6.107806in}{1.386000in}}%
\pgfpathlineto{\pgfqpoint{6.129069in}{1.386000in}}%
\pgfpathlineto{\pgfqpoint{6.129069in}{1.543500in}}%
\pgfpathlineto{\pgfqpoint{6.107806in}{1.543500in}}%
\pgfpathlineto{\pgfqpoint{6.107806in}{1.386000in}}%
\pgfpathclose%
\pgfusepath{fill}%
\end{pgfscope}%
\begin{pgfscope}%
\pgfpathrectangle{\pgfqpoint{3.815000in}{1.386000in}}{\pgfqpoint{2.835000in}{0.630000in}}%
\pgfusepath{clip}%
\pgfsetbuttcap%
\pgfsetmiterjoin%
\definecolor{currentfill}{rgb}{0.101961,0.101961,0.101961}%
\pgfsetfillcolor{currentfill}%
\pgfsetlinewidth{0.000000pt}%
\definecolor{currentstroke}{rgb}{0.000000,0.000000,0.000000}%
\pgfsetstrokecolor{currentstroke}%
\pgfsetstrokeopacity{0.000000}%
\pgfsetdash{}{0pt}%
\pgfpathmoveto{\pgfqpoint{6.131431in}{1.386000in}}%
\pgfpathlineto{\pgfqpoint{6.152694in}{1.386000in}}%
\pgfpathlineto{\pgfqpoint{6.152694in}{1.701000in}}%
\pgfpathlineto{\pgfqpoint{6.131431in}{1.701000in}}%
\pgfpathlineto{\pgfqpoint{6.131431in}{1.386000in}}%
\pgfpathclose%
\pgfusepath{fill}%
\end{pgfscope}%
\begin{pgfscope}%
\pgfpathrectangle{\pgfqpoint{3.815000in}{1.386000in}}{\pgfqpoint{2.835000in}{0.630000in}}%
\pgfusepath{clip}%
\pgfsetbuttcap%
\pgfsetmiterjoin%
\definecolor{currentfill}{rgb}{0.101961,0.101961,0.101961}%
\pgfsetfillcolor{currentfill}%
\pgfsetlinewidth{0.000000pt}%
\definecolor{currentstroke}{rgb}{0.000000,0.000000,0.000000}%
\pgfsetstrokecolor{currentstroke}%
\pgfsetstrokeopacity{0.000000}%
\pgfsetdash{}{0pt}%
\pgfpathmoveto{\pgfqpoint{6.155056in}{1.386000in}}%
\pgfpathlineto{\pgfqpoint{6.176319in}{1.386000in}}%
\pgfpathlineto{\pgfqpoint{6.176319in}{1.386000in}}%
\pgfpathlineto{\pgfqpoint{6.155056in}{1.386000in}}%
\pgfpathlineto{\pgfqpoint{6.155056in}{1.386000in}}%
\pgfpathclose%
\pgfusepath{fill}%
\end{pgfscope}%
\begin{pgfscope}%
\pgfpathrectangle{\pgfqpoint{3.815000in}{1.386000in}}{\pgfqpoint{2.835000in}{0.630000in}}%
\pgfusepath{clip}%
\pgfsetbuttcap%
\pgfsetmiterjoin%
\definecolor{currentfill}{rgb}{0.101961,0.101961,0.101961}%
\pgfsetfillcolor{currentfill}%
\pgfsetlinewidth{0.000000pt}%
\definecolor{currentstroke}{rgb}{0.000000,0.000000,0.000000}%
\pgfsetstrokecolor{currentstroke}%
\pgfsetstrokeopacity{0.000000}%
\pgfsetdash{}{0pt}%
\pgfpathmoveto{\pgfqpoint{6.178681in}{1.386000in}}%
\pgfpathlineto{\pgfqpoint{6.199944in}{1.386000in}}%
\pgfpathlineto{\pgfqpoint{6.199944in}{1.386000in}}%
\pgfpathlineto{\pgfqpoint{6.178681in}{1.386000in}}%
\pgfpathlineto{\pgfqpoint{6.178681in}{1.386000in}}%
\pgfpathclose%
\pgfusepath{fill}%
\end{pgfscope}%
\begin{pgfscope}%
\pgfpathrectangle{\pgfqpoint{3.815000in}{1.386000in}}{\pgfqpoint{2.835000in}{0.630000in}}%
\pgfusepath{clip}%
\pgfsetbuttcap%
\pgfsetmiterjoin%
\definecolor{currentfill}{rgb}{0.101961,0.101961,0.101961}%
\pgfsetfillcolor{currentfill}%
\pgfsetlinewidth{0.000000pt}%
\definecolor{currentstroke}{rgb}{0.000000,0.000000,0.000000}%
\pgfsetstrokecolor{currentstroke}%
\pgfsetstrokeopacity{0.000000}%
\pgfsetdash{}{0pt}%
\pgfpathmoveto{\pgfqpoint{6.202306in}{1.386000in}}%
\pgfpathlineto{\pgfqpoint{6.223569in}{1.386000in}}%
\pgfpathlineto{\pgfqpoint{6.223569in}{1.386000in}}%
\pgfpathlineto{\pgfqpoint{6.202306in}{1.386000in}}%
\pgfpathlineto{\pgfqpoint{6.202306in}{1.386000in}}%
\pgfpathclose%
\pgfusepath{fill}%
\end{pgfscope}%
\begin{pgfscope}%
\pgfpathrectangle{\pgfqpoint{3.815000in}{1.386000in}}{\pgfqpoint{2.835000in}{0.630000in}}%
\pgfusepath{clip}%
\pgfsetbuttcap%
\pgfsetmiterjoin%
\definecolor{currentfill}{rgb}{0.101961,0.101961,0.101961}%
\pgfsetfillcolor{currentfill}%
\pgfsetlinewidth{0.000000pt}%
\definecolor{currentstroke}{rgb}{0.000000,0.000000,0.000000}%
\pgfsetstrokecolor{currentstroke}%
\pgfsetstrokeopacity{0.000000}%
\pgfsetdash{}{0pt}%
\pgfpathmoveto{\pgfqpoint{6.225931in}{1.386000in}}%
\pgfpathlineto{\pgfqpoint{6.247194in}{1.386000in}}%
\pgfpathlineto{\pgfqpoint{6.247194in}{1.701000in}}%
\pgfpathlineto{\pgfqpoint{6.225931in}{1.701000in}}%
\pgfpathlineto{\pgfqpoint{6.225931in}{1.386000in}}%
\pgfpathclose%
\pgfusepath{fill}%
\end{pgfscope}%
\begin{pgfscope}%
\pgfpathrectangle{\pgfqpoint{3.815000in}{1.386000in}}{\pgfqpoint{2.835000in}{0.630000in}}%
\pgfusepath{clip}%
\pgfsetbuttcap%
\pgfsetmiterjoin%
\definecolor{currentfill}{rgb}{0.101961,0.101961,0.101961}%
\pgfsetfillcolor{currentfill}%
\pgfsetlinewidth{0.000000pt}%
\definecolor{currentstroke}{rgb}{0.000000,0.000000,0.000000}%
\pgfsetstrokecolor{currentstroke}%
\pgfsetstrokeopacity{0.000000}%
\pgfsetdash{}{0pt}%
\pgfpathmoveto{\pgfqpoint{6.249556in}{1.386000in}}%
\pgfpathlineto{\pgfqpoint{6.270819in}{1.386000in}}%
\pgfpathlineto{\pgfqpoint{6.270819in}{1.386000in}}%
\pgfpathlineto{\pgfqpoint{6.249556in}{1.386000in}}%
\pgfpathlineto{\pgfqpoint{6.249556in}{1.386000in}}%
\pgfpathclose%
\pgfusepath{fill}%
\end{pgfscope}%
\begin{pgfscope}%
\pgfpathrectangle{\pgfqpoint{3.815000in}{1.386000in}}{\pgfqpoint{2.835000in}{0.630000in}}%
\pgfusepath{clip}%
\pgfsetbuttcap%
\pgfsetmiterjoin%
\definecolor{currentfill}{rgb}{0.101961,0.101961,0.101961}%
\pgfsetfillcolor{currentfill}%
\pgfsetlinewidth{0.000000pt}%
\definecolor{currentstroke}{rgb}{0.000000,0.000000,0.000000}%
\pgfsetstrokecolor{currentstroke}%
\pgfsetstrokeopacity{0.000000}%
\pgfsetdash{}{0pt}%
\pgfpathmoveto{\pgfqpoint{6.273181in}{1.386000in}}%
\pgfpathlineto{\pgfqpoint{6.294444in}{1.386000in}}%
\pgfpathlineto{\pgfqpoint{6.294444in}{1.386000in}}%
\pgfpathlineto{\pgfqpoint{6.273181in}{1.386000in}}%
\pgfpathlineto{\pgfqpoint{6.273181in}{1.386000in}}%
\pgfpathclose%
\pgfusepath{fill}%
\end{pgfscope}%
\begin{pgfscope}%
\pgfpathrectangle{\pgfqpoint{3.815000in}{1.386000in}}{\pgfqpoint{2.835000in}{0.630000in}}%
\pgfusepath{clip}%
\pgfsetbuttcap%
\pgfsetmiterjoin%
\definecolor{currentfill}{rgb}{0.101961,0.101961,0.101961}%
\pgfsetfillcolor{currentfill}%
\pgfsetlinewidth{0.000000pt}%
\definecolor{currentstroke}{rgb}{0.000000,0.000000,0.000000}%
\pgfsetstrokecolor{currentstroke}%
\pgfsetstrokeopacity{0.000000}%
\pgfsetdash{}{0pt}%
\pgfpathmoveto{\pgfqpoint{6.296806in}{1.386000in}}%
\pgfpathlineto{\pgfqpoint{6.318069in}{1.386000in}}%
\pgfpathlineto{\pgfqpoint{6.318069in}{1.386000in}}%
\pgfpathlineto{\pgfqpoint{6.296806in}{1.386000in}}%
\pgfpathlineto{\pgfqpoint{6.296806in}{1.386000in}}%
\pgfpathclose%
\pgfusepath{fill}%
\end{pgfscope}%
\begin{pgfscope}%
\pgfpathrectangle{\pgfqpoint{3.815000in}{1.386000in}}{\pgfqpoint{2.835000in}{0.630000in}}%
\pgfusepath{clip}%
\pgfsetbuttcap%
\pgfsetmiterjoin%
\definecolor{currentfill}{rgb}{0.101961,0.101961,0.101961}%
\pgfsetfillcolor{currentfill}%
\pgfsetlinewidth{0.000000pt}%
\definecolor{currentstroke}{rgb}{0.000000,0.000000,0.000000}%
\pgfsetstrokecolor{currentstroke}%
\pgfsetstrokeopacity{0.000000}%
\pgfsetdash{}{0pt}%
\pgfpathmoveto{\pgfqpoint{6.320431in}{1.386000in}}%
\pgfpathlineto{\pgfqpoint{6.341694in}{1.386000in}}%
\pgfpathlineto{\pgfqpoint{6.341694in}{1.386000in}}%
\pgfpathlineto{\pgfqpoint{6.320431in}{1.386000in}}%
\pgfpathlineto{\pgfqpoint{6.320431in}{1.386000in}}%
\pgfpathclose%
\pgfusepath{fill}%
\end{pgfscope}%
\begin{pgfscope}%
\pgfpathrectangle{\pgfqpoint{3.815000in}{1.386000in}}{\pgfqpoint{2.835000in}{0.630000in}}%
\pgfusepath{clip}%
\pgfsetbuttcap%
\pgfsetmiterjoin%
\definecolor{currentfill}{rgb}{0.101961,0.101961,0.101961}%
\pgfsetfillcolor{currentfill}%
\pgfsetlinewidth{0.000000pt}%
\definecolor{currentstroke}{rgb}{0.000000,0.000000,0.000000}%
\pgfsetstrokecolor{currentstroke}%
\pgfsetstrokeopacity{0.000000}%
\pgfsetdash{}{0pt}%
\pgfpathmoveto{\pgfqpoint{6.344056in}{1.386000in}}%
\pgfpathlineto{\pgfqpoint{6.365319in}{1.386000in}}%
\pgfpathlineto{\pgfqpoint{6.365319in}{1.386000in}}%
\pgfpathlineto{\pgfqpoint{6.344056in}{1.386000in}}%
\pgfpathlineto{\pgfqpoint{6.344056in}{1.386000in}}%
\pgfpathclose%
\pgfusepath{fill}%
\end{pgfscope}%
\begin{pgfscope}%
\pgfpathrectangle{\pgfqpoint{3.815000in}{1.386000in}}{\pgfqpoint{2.835000in}{0.630000in}}%
\pgfusepath{clip}%
\pgfsetbuttcap%
\pgfsetmiterjoin%
\definecolor{currentfill}{rgb}{0.101961,0.101961,0.101961}%
\pgfsetfillcolor{currentfill}%
\pgfsetlinewidth{0.000000pt}%
\definecolor{currentstroke}{rgb}{0.000000,0.000000,0.000000}%
\pgfsetstrokecolor{currentstroke}%
\pgfsetstrokeopacity{0.000000}%
\pgfsetdash{}{0pt}%
\pgfpathmoveto{\pgfqpoint{6.367681in}{1.386000in}}%
\pgfpathlineto{\pgfqpoint{6.388944in}{1.386000in}}%
\pgfpathlineto{\pgfqpoint{6.388944in}{1.386000in}}%
\pgfpathlineto{\pgfqpoint{6.367681in}{1.386000in}}%
\pgfpathlineto{\pgfqpoint{6.367681in}{1.386000in}}%
\pgfpathclose%
\pgfusepath{fill}%
\end{pgfscope}%
\begin{pgfscope}%
\pgfpathrectangle{\pgfqpoint{3.815000in}{1.386000in}}{\pgfqpoint{2.835000in}{0.630000in}}%
\pgfusepath{clip}%
\pgfsetbuttcap%
\pgfsetmiterjoin%
\definecolor{currentfill}{rgb}{0.101961,0.101961,0.101961}%
\pgfsetfillcolor{currentfill}%
\pgfsetlinewidth{0.000000pt}%
\definecolor{currentstroke}{rgb}{0.000000,0.000000,0.000000}%
\pgfsetstrokecolor{currentstroke}%
\pgfsetstrokeopacity{0.000000}%
\pgfsetdash{}{0pt}%
\pgfpathmoveto{\pgfqpoint{6.391306in}{1.386000in}}%
\pgfpathlineto{\pgfqpoint{6.412569in}{1.386000in}}%
\pgfpathlineto{\pgfqpoint{6.412569in}{1.386000in}}%
\pgfpathlineto{\pgfqpoint{6.391306in}{1.386000in}}%
\pgfpathlineto{\pgfqpoint{6.391306in}{1.386000in}}%
\pgfpathclose%
\pgfusepath{fill}%
\end{pgfscope}%
\begin{pgfscope}%
\pgfpathrectangle{\pgfqpoint{3.815000in}{1.386000in}}{\pgfqpoint{2.835000in}{0.630000in}}%
\pgfusepath{clip}%
\pgfsetbuttcap%
\pgfsetmiterjoin%
\definecolor{currentfill}{rgb}{0.101961,0.101961,0.101961}%
\pgfsetfillcolor{currentfill}%
\pgfsetlinewidth{0.000000pt}%
\definecolor{currentstroke}{rgb}{0.000000,0.000000,0.000000}%
\pgfsetstrokecolor{currentstroke}%
\pgfsetstrokeopacity{0.000000}%
\pgfsetdash{}{0pt}%
\pgfpathmoveto{\pgfqpoint{6.414931in}{1.386000in}}%
\pgfpathlineto{\pgfqpoint{6.436194in}{1.386000in}}%
\pgfpathlineto{\pgfqpoint{6.436194in}{1.386000in}}%
\pgfpathlineto{\pgfqpoint{6.414931in}{1.386000in}}%
\pgfpathlineto{\pgfqpoint{6.414931in}{1.386000in}}%
\pgfpathclose%
\pgfusepath{fill}%
\end{pgfscope}%
\begin{pgfscope}%
\pgfpathrectangle{\pgfqpoint{3.815000in}{1.386000in}}{\pgfqpoint{2.835000in}{0.630000in}}%
\pgfusepath{clip}%
\pgfsetbuttcap%
\pgfsetmiterjoin%
\definecolor{currentfill}{rgb}{0.101961,0.101961,0.101961}%
\pgfsetfillcolor{currentfill}%
\pgfsetlinewidth{0.000000pt}%
\definecolor{currentstroke}{rgb}{0.000000,0.000000,0.000000}%
\pgfsetstrokecolor{currentstroke}%
\pgfsetstrokeopacity{0.000000}%
\pgfsetdash{}{0pt}%
\pgfpathmoveto{\pgfqpoint{6.438556in}{1.386000in}}%
\pgfpathlineto{\pgfqpoint{6.459819in}{1.386000in}}%
\pgfpathlineto{\pgfqpoint{6.459819in}{1.386000in}}%
\pgfpathlineto{\pgfqpoint{6.438556in}{1.386000in}}%
\pgfpathlineto{\pgfqpoint{6.438556in}{1.386000in}}%
\pgfpathclose%
\pgfusepath{fill}%
\end{pgfscope}%
\begin{pgfscope}%
\pgfpathrectangle{\pgfqpoint{3.815000in}{1.386000in}}{\pgfqpoint{2.835000in}{0.630000in}}%
\pgfusepath{clip}%
\pgfsetbuttcap%
\pgfsetmiterjoin%
\definecolor{currentfill}{rgb}{0.101961,0.101961,0.101961}%
\pgfsetfillcolor{currentfill}%
\pgfsetlinewidth{0.000000pt}%
\definecolor{currentstroke}{rgb}{0.000000,0.000000,0.000000}%
\pgfsetstrokecolor{currentstroke}%
\pgfsetstrokeopacity{0.000000}%
\pgfsetdash{}{0pt}%
\pgfpathmoveto{\pgfqpoint{6.462181in}{1.386000in}}%
\pgfpathlineto{\pgfqpoint{6.483444in}{1.386000in}}%
\pgfpathlineto{\pgfqpoint{6.483444in}{1.386000in}}%
\pgfpathlineto{\pgfqpoint{6.462181in}{1.386000in}}%
\pgfpathlineto{\pgfqpoint{6.462181in}{1.386000in}}%
\pgfpathclose%
\pgfusepath{fill}%
\end{pgfscope}%
\begin{pgfscope}%
\pgfpathrectangle{\pgfqpoint{3.815000in}{1.386000in}}{\pgfqpoint{2.835000in}{0.630000in}}%
\pgfusepath{clip}%
\pgfsetbuttcap%
\pgfsetmiterjoin%
\definecolor{currentfill}{rgb}{0.101961,0.101961,0.101961}%
\pgfsetfillcolor{currentfill}%
\pgfsetlinewidth{0.000000pt}%
\definecolor{currentstroke}{rgb}{0.000000,0.000000,0.000000}%
\pgfsetstrokecolor{currentstroke}%
\pgfsetstrokeopacity{0.000000}%
\pgfsetdash{}{0pt}%
\pgfpathmoveto{\pgfqpoint{6.485806in}{1.386000in}}%
\pgfpathlineto{\pgfqpoint{6.507069in}{1.386000in}}%
\pgfpathlineto{\pgfqpoint{6.507069in}{1.386000in}}%
\pgfpathlineto{\pgfqpoint{6.485806in}{1.386000in}}%
\pgfpathlineto{\pgfqpoint{6.485806in}{1.386000in}}%
\pgfpathclose%
\pgfusepath{fill}%
\end{pgfscope}%
\begin{pgfscope}%
\pgfpathrectangle{\pgfqpoint{3.815000in}{1.386000in}}{\pgfqpoint{2.835000in}{0.630000in}}%
\pgfusepath{clip}%
\pgfsetbuttcap%
\pgfsetmiterjoin%
\definecolor{currentfill}{rgb}{0.101961,0.101961,0.101961}%
\pgfsetfillcolor{currentfill}%
\pgfsetlinewidth{0.000000pt}%
\definecolor{currentstroke}{rgb}{0.000000,0.000000,0.000000}%
\pgfsetstrokecolor{currentstroke}%
\pgfsetstrokeopacity{0.000000}%
\pgfsetdash{}{0pt}%
\pgfpathmoveto{\pgfqpoint{6.509431in}{1.386000in}}%
\pgfpathlineto{\pgfqpoint{6.530694in}{1.386000in}}%
\pgfpathlineto{\pgfqpoint{6.530694in}{1.543500in}}%
\pgfpathlineto{\pgfqpoint{6.509431in}{1.543500in}}%
\pgfpathlineto{\pgfqpoint{6.509431in}{1.386000in}}%
\pgfpathclose%
\pgfusepath{fill}%
\end{pgfscope}%
\begin{pgfscope}%
\pgfpathrectangle{\pgfqpoint{3.815000in}{1.386000in}}{\pgfqpoint{2.835000in}{0.630000in}}%
\pgfusepath{clip}%
\pgfsetbuttcap%
\pgfsetmiterjoin%
\definecolor{currentfill}{rgb}{0.101961,0.101961,0.101961}%
\pgfsetfillcolor{currentfill}%
\pgfsetlinewidth{0.000000pt}%
\definecolor{currentstroke}{rgb}{0.000000,0.000000,0.000000}%
\pgfsetstrokecolor{currentstroke}%
\pgfsetstrokeopacity{0.000000}%
\pgfsetdash{}{0pt}%
\pgfpathmoveto{\pgfqpoint{6.533056in}{1.386000in}}%
\pgfpathlineto{\pgfqpoint{6.554319in}{1.386000in}}%
\pgfpathlineto{\pgfqpoint{6.554319in}{1.543500in}}%
\pgfpathlineto{\pgfqpoint{6.533056in}{1.543500in}}%
\pgfpathlineto{\pgfqpoint{6.533056in}{1.386000in}}%
\pgfpathclose%
\pgfusepath{fill}%
\end{pgfscope}%
\begin{pgfscope}%
\pgfpathrectangle{\pgfqpoint{3.815000in}{1.386000in}}{\pgfqpoint{2.835000in}{0.630000in}}%
\pgfusepath{clip}%
\pgfsetbuttcap%
\pgfsetmiterjoin%
\definecolor{currentfill}{rgb}{0.101961,0.101961,0.101961}%
\pgfsetfillcolor{currentfill}%
\pgfsetlinewidth{0.000000pt}%
\definecolor{currentstroke}{rgb}{0.000000,0.000000,0.000000}%
\pgfsetstrokecolor{currentstroke}%
\pgfsetstrokeopacity{0.000000}%
\pgfsetdash{}{0pt}%
\pgfpathmoveto{\pgfqpoint{6.556681in}{1.386000in}}%
\pgfpathlineto{\pgfqpoint{6.577944in}{1.386000in}}%
\pgfpathlineto{\pgfqpoint{6.577944in}{1.386000in}}%
\pgfpathlineto{\pgfqpoint{6.556681in}{1.386000in}}%
\pgfpathlineto{\pgfqpoint{6.556681in}{1.386000in}}%
\pgfpathclose%
\pgfusepath{fill}%
\end{pgfscope}%
\begin{pgfscope}%
\pgfpathrectangle{\pgfqpoint{3.815000in}{1.386000in}}{\pgfqpoint{2.835000in}{0.630000in}}%
\pgfusepath{clip}%
\pgfsetbuttcap%
\pgfsetmiterjoin%
\definecolor{currentfill}{rgb}{0.101961,0.101961,0.101961}%
\pgfsetfillcolor{currentfill}%
\pgfsetlinewidth{0.000000pt}%
\definecolor{currentstroke}{rgb}{0.000000,0.000000,0.000000}%
\pgfsetstrokecolor{currentstroke}%
\pgfsetstrokeopacity{0.000000}%
\pgfsetdash{}{0pt}%
\pgfpathmoveto{\pgfqpoint{6.580306in}{1.386000in}}%
\pgfpathlineto{\pgfqpoint{6.601569in}{1.386000in}}%
\pgfpathlineto{\pgfqpoint{6.601569in}{1.386000in}}%
\pgfpathlineto{\pgfqpoint{6.580306in}{1.386000in}}%
\pgfpathlineto{\pgfqpoint{6.580306in}{1.386000in}}%
\pgfpathclose%
\pgfusepath{fill}%
\end{pgfscope}%
\begin{pgfscope}%
\pgfpathrectangle{\pgfqpoint{3.815000in}{1.386000in}}{\pgfqpoint{2.835000in}{0.630000in}}%
\pgfusepath{clip}%
\pgfsetbuttcap%
\pgfsetmiterjoin%
\definecolor{currentfill}{rgb}{0.101961,0.101961,0.101961}%
\pgfsetfillcolor{currentfill}%
\pgfsetlinewidth{0.000000pt}%
\definecolor{currentstroke}{rgb}{0.000000,0.000000,0.000000}%
\pgfsetstrokecolor{currentstroke}%
\pgfsetstrokeopacity{0.000000}%
\pgfsetdash{}{0pt}%
\pgfpathmoveto{\pgfqpoint{6.603931in}{1.386000in}}%
\pgfpathlineto{\pgfqpoint{6.625194in}{1.386000in}}%
\pgfpathlineto{\pgfqpoint{6.625194in}{1.386000in}}%
\pgfpathlineto{\pgfqpoint{6.603931in}{1.386000in}}%
\pgfpathlineto{\pgfqpoint{6.603931in}{1.386000in}}%
\pgfpathclose%
\pgfusepath{fill}%
\end{pgfscope}%
\begin{pgfscope}%
\pgfpathrectangle{\pgfqpoint{3.815000in}{1.386000in}}{\pgfqpoint{2.835000in}{0.630000in}}%
\pgfusepath{clip}%
\pgfsetbuttcap%
\pgfsetmiterjoin%
\definecolor{currentfill}{rgb}{0.101961,0.101961,0.101961}%
\pgfsetfillcolor{currentfill}%
\pgfsetlinewidth{0.000000pt}%
\definecolor{currentstroke}{rgb}{0.000000,0.000000,0.000000}%
\pgfsetstrokecolor{currentstroke}%
\pgfsetstrokeopacity{0.000000}%
\pgfsetdash{}{0pt}%
\pgfpathmoveto{\pgfqpoint{6.627556in}{1.386000in}}%
\pgfpathlineto{\pgfqpoint{6.648819in}{1.386000in}}%
\pgfpathlineto{\pgfqpoint{6.648819in}{1.386000in}}%
\pgfpathlineto{\pgfqpoint{6.627556in}{1.386000in}}%
\pgfpathlineto{\pgfqpoint{6.627556in}{1.386000in}}%
\pgfpathclose%
\pgfusepath{fill}%
\end{pgfscope}%
\begin{pgfscope}%
\pgfpathrectangle{\pgfqpoint{3.815000in}{1.386000in}}{\pgfqpoint{2.835000in}{0.630000in}}%
\pgfusepath{clip}%
\pgfsetbuttcap%
\pgfsetroundjoin%
\pgfsetlinewidth{0.702625pt}%
\definecolor{currentstroke}{rgb}{0.690196,0.250980,0.250980}%
\pgfsetstrokecolor{currentstroke}%
\pgfsetdash{{2.590000pt}{1.120000pt}}{0.000000pt}%
\pgfpathmoveto{\pgfqpoint{5.221869in}{1.386000in}}%
\pgfpathlineto{\pgfqpoint{5.221869in}{2.016000in}}%
\pgfusepath{stroke}%
\end{pgfscope}%
\begin{pgfscope}%
\pgfsetrectcap%
\pgfsetmiterjoin%
\pgfsetlinewidth{0.803000pt}%
\definecolor{currentstroke}{rgb}{0.101961,0.101961,0.101961}%
\pgfsetstrokecolor{currentstroke}%
\pgfsetdash{}{0pt}%
\pgfpathmoveto{\pgfqpoint{3.815000in}{1.386000in}}%
\pgfpathlineto{\pgfqpoint{3.815000in}{2.016000in}}%
\pgfusepath{stroke}%
\end{pgfscope}%
\begin{pgfscope}%
\pgfsetrectcap%
\pgfsetmiterjoin%
\pgfsetlinewidth{0.803000pt}%
\definecolor{currentstroke}{rgb}{0.101961,0.101961,0.101961}%
\pgfsetstrokecolor{currentstroke}%
\pgfsetdash{}{0pt}%
\pgfpathmoveto{\pgfqpoint{3.815000in}{1.386000in}}%
\pgfpathlineto{\pgfqpoint{6.650000in}{1.386000in}}%
\pgfusepath{stroke}%
\end{pgfscope}%
\begin{pgfscope}%
\pgfsetbuttcap%
\pgfsetmiterjoin%
\definecolor{currentfill}{rgb}{1.000000,1.000000,1.000000}%
\pgfsetfillcolor{currentfill}%
\pgfsetlinewidth{0.000000pt}%
\definecolor{currentstroke}{rgb}{0.000000,0.000000,0.000000}%
\pgfsetstrokecolor{currentstroke}%
\pgfsetstrokeopacity{0.000000}%
\pgfsetdash{}{0pt}%
\pgfpathmoveto{\pgfqpoint{3.815000in}{0.483000in}}%
\pgfpathlineto{\pgfqpoint{6.650000in}{0.483000in}}%
\pgfpathlineto{\pgfqpoint{6.650000in}{1.113000in}}%
\pgfpathlineto{\pgfqpoint{3.815000in}{1.113000in}}%
\pgfpathlineto{\pgfqpoint{3.815000in}{0.483000in}}%
\pgfpathclose%
\pgfusepath{fill}%
\end{pgfscope}%
\begin{pgfscope}%
\pgfpathrectangle{\pgfqpoint{3.815000in}{0.483000in}}{\pgfqpoint{2.835000in}{0.630000in}}%
\pgfusepath{clip}%
\pgfsetbuttcap%
\pgfsetmiterjoin%
\definecolor{currentfill}{rgb}{0.941176,0.894118,0.894118}%
\pgfsetfillcolor{currentfill}%
\pgfsetlinewidth{0.000000pt}%
\definecolor{currentstroke}{rgb}{0.941176,0.894118,0.894118}%
\pgfsetstrokecolor{currentstroke}%
\pgfsetdash{}{0pt}%
\pgfpathmoveto{\pgfqpoint{4.777719in}{0.483000in}}%
\pgfpathlineto{\pgfqpoint{4.777719in}{1.113000in}}%
\pgfpathlineto{\pgfqpoint{5.221869in}{1.113000in}}%
\pgfpathlineto{\pgfqpoint{5.221869in}{0.483000in}}%
\pgfpathlineto{\pgfqpoint{4.777719in}{0.483000in}}%
\pgfpathclose%
\pgfusepath{fill}%
\end{pgfscope}%
\begin{pgfscope}%
\definecolor{textcolor}{rgb}{0.101961,0.101961,0.101961}%
\pgfsetstrokecolor{textcolor}%
\pgfsetfillcolor{textcolor}%
\pgftext[x=3.815000in,y=0.434389in,,top]{\color{textcolor}\sffamily\fontsize{7.000000}{8.400000}\selectfont 0}%
\end{pgfscope}%
\begin{pgfscope}%
\definecolor{textcolor}{rgb}{0.101961,0.101961,0.101961}%
\pgfsetstrokecolor{textcolor}%
\pgfsetfillcolor{textcolor}%
\pgftext[x=4.287500in,y=0.434389in,,top]{\color{textcolor}\sffamily\fontsize{7.000000}{8.400000}\selectfont 200}%
\end{pgfscope}%
\begin{pgfscope}%
\definecolor{textcolor}{rgb}{0.101961,0.101961,0.101961}%
\pgfsetstrokecolor{textcolor}%
\pgfsetfillcolor{textcolor}%
\pgftext[x=4.760000in,y=0.434389in,,top]{\color{textcolor}\sffamily\fontsize{7.000000}{8.400000}\selectfont 400}%
\end{pgfscope}%
\begin{pgfscope}%
\definecolor{textcolor}{rgb}{0.101961,0.101961,0.101961}%
\pgfsetstrokecolor{textcolor}%
\pgfsetfillcolor{textcolor}%
\pgftext[x=5.232500in,y=0.434389in,,top]{\color{textcolor}\sffamily\fontsize{7.000000}{8.400000}\selectfont 600}%
\end{pgfscope}%
\begin{pgfscope}%
\definecolor{textcolor}{rgb}{0.101961,0.101961,0.101961}%
\pgfsetstrokecolor{textcolor}%
\pgfsetfillcolor{textcolor}%
\pgftext[x=5.705000in,y=0.434389in,,top]{\color{textcolor}\sffamily\fontsize{7.000000}{8.400000}\selectfont 800}%
\end{pgfscope}%
\begin{pgfscope}%
\definecolor{textcolor}{rgb}{0.101961,0.101961,0.101961}%
\pgfsetstrokecolor{textcolor}%
\pgfsetfillcolor{textcolor}%
\pgftext[x=6.177500in,y=0.434389in,,top]{\color{textcolor}\sffamily\fontsize{7.000000}{8.400000}\selectfont 1000}%
\end{pgfscope}%
\begin{pgfscope}%
\definecolor{textcolor}{rgb}{0.101961,0.101961,0.101961}%
\pgfsetstrokecolor{textcolor}%
\pgfsetfillcolor{textcolor}%
\pgftext[x=6.650000in,y=0.434389in,,top]{\color{textcolor}\sffamily\fontsize{7.000000}{8.400000}\selectfont 1200}%
\end{pgfscope}%
\begin{pgfscope}%
\definecolor{textcolor}{rgb}{0.000000,0.000000,0.000000}%
\pgfsetstrokecolor{textcolor}%
\pgfsetfillcolor{textcolor}%
\pgftext[x=5.232500in,y=0.320192in,,top]{\color{textcolor}\sffamily\fontsize{7.500000}{9.000000}\selectfont Time (s)}%
\end{pgfscope}%
\begin{pgfscope}%
\pgfpathrectangle{\pgfqpoint{3.815000in}{0.483000in}}{\pgfqpoint{2.835000in}{0.630000in}}%
\pgfusepath{clip}%
\pgfsetrectcap%
\pgfsetroundjoin%
\pgfsetlinewidth{0.401500pt}%
\definecolor{currentstroke}{rgb}{0.901961,0.901961,0.901961}%
\pgfsetstrokecolor{currentstroke}%
\pgfsetdash{}{0pt}%
\pgfpathmoveto{\pgfqpoint{3.815000in}{0.483000in}}%
\pgfpathlineto{\pgfqpoint{6.650000in}{0.483000in}}%
\pgfusepath{stroke}%
\end{pgfscope}%
\begin{pgfscope}%
\pgfpathrectangle{\pgfqpoint{3.815000in}{0.483000in}}{\pgfqpoint{2.835000in}{0.630000in}}%
\pgfusepath{clip}%
\pgfsetrectcap%
\pgfsetroundjoin%
\pgfsetlinewidth{0.401500pt}%
\definecolor{currentstroke}{rgb}{0.901961,0.901961,0.901961}%
\pgfsetstrokecolor{currentstroke}%
\pgfsetdash{}{0pt}%
\pgfpathmoveto{\pgfqpoint{3.815000in}{0.708000in}}%
\pgfpathlineto{\pgfqpoint{6.650000in}{0.708000in}}%
\pgfusepath{stroke}%
\end{pgfscope}%
\begin{pgfscope}%
\pgfpathrectangle{\pgfqpoint{3.815000in}{0.483000in}}{\pgfqpoint{2.835000in}{0.630000in}}%
\pgfusepath{clip}%
\pgfsetrectcap%
\pgfsetroundjoin%
\pgfsetlinewidth{0.401500pt}%
\definecolor{currentstroke}{rgb}{0.901961,0.901961,0.901961}%
\pgfsetstrokecolor{currentstroke}%
\pgfsetdash{}{0pt}%
\pgfpathmoveto{\pgfqpoint{3.815000in}{0.933000in}}%
\pgfpathlineto{\pgfqpoint{6.650000in}{0.933000in}}%
\pgfusepath{stroke}%
\end{pgfscope}%
\begin{pgfscope}%
\pgfpathrectangle{\pgfqpoint{3.815000in}{0.483000in}}{\pgfqpoint{2.835000in}{0.630000in}}%
\pgfusepath{clip}%
\pgfsetbuttcap%
\pgfsetmiterjoin%
\definecolor{currentfill}{rgb}{0.101961,0.101961,0.101961}%
\pgfsetfillcolor{currentfill}%
\pgfsetlinewidth{0.000000pt}%
\definecolor{currentstroke}{rgb}{0.000000,0.000000,0.000000}%
\pgfsetstrokecolor{currentstroke}%
\pgfsetstrokeopacity{0.000000}%
\pgfsetdash{}{0pt}%
\pgfpathmoveto{\pgfqpoint{3.816181in}{0.483000in}}%
\pgfpathlineto{\pgfqpoint{3.837444in}{0.483000in}}%
\pgfpathlineto{\pgfqpoint{3.837444in}{0.483000in}}%
\pgfpathlineto{\pgfqpoint{3.816181in}{0.483000in}}%
\pgfpathlineto{\pgfqpoint{3.816181in}{0.483000in}}%
\pgfpathclose%
\pgfusepath{fill}%
\end{pgfscope}%
\begin{pgfscope}%
\pgfpathrectangle{\pgfqpoint{3.815000in}{0.483000in}}{\pgfqpoint{2.835000in}{0.630000in}}%
\pgfusepath{clip}%
\pgfsetbuttcap%
\pgfsetmiterjoin%
\definecolor{currentfill}{rgb}{0.101961,0.101961,0.101961}%
\pgfsetfillcolor{currentfill}%
\pgfsetlinewidth{0.000000pt}%
\definecolor{currentstroke}{rgb}{0.000000,0.000000,0.000000}%
\pgfsetstrokecolor{currentstroke}%
\pgfsetstrokeopacity{0.000000}%
\pgfsetdash{}{0pt}%
\pgfpathmoveto{\pgfqpoint{3.839806in}{0.483000in}}%
\pgfpathlineto{\pgfqpoint{3.861069in}{0.483000in}}%
\pgfpathlineto{\pgfqpoint{3.861069in}{0.483000in}}%
\pgfpathlineto{\pgfqpoint{3.839806in}{0.483000in}}%
\pgfpathlineto{\pgfqpoint{3.839806in}{0.483000in}}%
\pgfpathclose%
\pgfusepath{fill}%
\end{pgfscope}%
\begin{pgfscope}%
\pgfpathrectangle{\pgfqpoint{3.815000in}{0.483000in}}{\pgfqpoint{2.835000in}{0.630000in}}%
\pgfusepath{clip}%
\pgfsetbuttcap%
\pgfsetmiterjoin%
\definecolor{currentfill}{rgb}{0.101961,0.101961,0.101961}%
\pgfsetfillcolor{currentfill}%
\pgfsetlinewidth{0.000000pt}%
\definecolor{currentstroke}{rgb}{0.000000,0.000000,0.000000}%
\pgfsetstrokecolor{currentstroke}%
\pgfsetstrokeopacity{0.000000}%
\pgfsetdash{}{0pt}%
\pgfpathmoveto{\pgfqpoint{3.863431in}{0.483000in}}%
\pgfpathlineto{\pgfqpoint{3.884694in}{0.483000in}}%
\pgfpathlineto{\pgfqpoint{3.884694in}{0.528000in}}%
\pgfpathlineto{\pgfqpoint{3.863431in}{0.528000in}}%
\pgfpathlineto{\pgfqpoint{3.863431in}{0.483000in}}%
\pgfpathclose%
\pgfusepath{fill}%
\end{pgfscope}%
\begin{pgfscope}%
\pgfpathrectangle{\pgfqpoint{3.815000in}{0.483000in}}{\pgfqpoint{2.835000in}{0.630000in}}%
\pgfusepath{clip}%
\pgfsetbuttcap%
\pgfsetmiterjoin%
\definecolor{currentfill}{rgb}{0.101961,0.101961,0.101961}%
\pgfsetfillcolor{currentfill}%
\pgfsetlinewidth{0.000000pt}%
\definecolor{currentstroke}{rgb}{0.000000,0.000000,0.000000}%
\pgfsetstrokecolor{currentstroke}%
\pgfsetstrokeopacity{0.000000}%
\pgfsetdash{}{0pt}%
\pgfpathmoveto{\pgfqpoint{3.887056in}{0.483000in}}%
\pgfpathlineto{\pgfqpoint{3.908319in}{0.483000in}}%
\pgfpathlineto{\pgfqpoint{3.908319in}{0.483000in}}%
\pgfpathlineto{\pgfqpoint{3.887056in}{0.483000in}}%
\pgfpathlineto{\pgfqpoint{3.887056in}{0.483000in}}%
\pgfpathclose%
\pgfusepath{fill}%
\end{pgfscope}%
\begin{pgfscope}%
\pgfpathrectangle{\pgfqpoint{3.815000in}{0.483000in}}{\pgfqpoint{2.835000in}{0.630000in}}%
\pgfusepath{clip}%
\pgfsetbuttcap%
\pgfsetmiterjoin%
\definecolor{currentfill}{rgb}{0.101961,0.101961,0.101961}%
\pgfsetfillcolor{currentfill}%
\pgfsetlinewidth{0.000000pt}%
\definecolor{currentstroke}{rgb}{0.000000,0.000000,0.000000}%
\pgfsetstrokecolor{currentstroke}%
\pgfsetstrokeopacity{0.000000}%
\pgfsetdash{}{0pt}%
\pgfpathmoveto{\pgfqpoint{3.910681in}{0.483000in}}%
\pgfpathlineto{\pgfqpoint{3.931944in}{0.483000in}}%
\pgfpathlineto{\pgfqpoint{3.931944in}{0.573000in}}%
\pgfpathlineto{\pgfqpoint{3.910681in}{0.573000in}}%
\pgfpathlineto{\pgfqpoint{3.910681in}{0.483000in}}%
\pgfpathclose%
\pgfusepath{fill}%
\end{pgfscope}%
\begin{pgfscope}%
\pgfpathrectangle{\pgfqpoint{3.815000in}{0.483000in}}{\pgfqpoint{2.835000in}{0.630000in}}%
\pgfusepath{clip}%
\pgfsetbuttcap%
\pgfsetmiterjoin%
\definecolor{currentfill}{rgb}{0.101961,0.101961,0.101961}%
\pgfsetfillcolor{currentfill}%
\pgfsetlinewidth{0.000000pt}%
\definecolor{currentstroke}{rgb}{0.000000,0.000000,0.000000}%
\pgfsetstrokecolor{currentstroke}%
\pgfsetstrokeopacity{0.000000}%
\pgfsetdash{}{0pt}%
\pgfpathmoveto{\pgfqpoint{3.934306in}{0.483000in}}%
\pgfpathlineto{\pgfqpoint{3.955569in}{0.483000in}}%
\pgfpathlineto{\pgfqpoint{3.955569in}{0.483000in}}%
\pgfpathlineto{\pgfqpoint{3.934306in}{0.483000in}}%
\pgfpathlineto{\pgfqpoint{3.934306in}{0.483000in}}%
\pgfpathclose%
\pgfusepath{fill}%
\end{pgfscope}%
\begin{pgfscope}%
\pgfpathrectangle{\pgfqpoint{3.815000in}{0.483000in}}{\pgfqpoint{2.835000in}{0.630000in}}%
\pgfusepath{clip}%
\pgfsetbuttcap%
\pgfsetmiterjoin%
\definecolor{currentfill}{rgb}{0.101961,0.101961,0.101961}%
\pgfsetfillcolor{currentfill}%
\pgfsetlinewidth{0.000000pt}%
\definecolor{currentstroke}{rgb}{0.000000,0.000000,0.000000}%
\pgfsetstrokecolor{currentstroke}%
\pgfsetstrokeopacity{0.000000}%
\pgfsetdash{}{0pt}%
\pgfpathmoveto{\pgfqpoint{3.957931in}{0.483000in}}%
\pgfpathlineto{\pgfqpoint{3.979194in}{0.483000in}}%
\pgfpathlineto{\pgfqpoint{3.979194in}{0.483000in}}%
\pgfpathlineto{\pgfqpoint{3.957931in}{0.483000in}}%
\pgfpathlineto{\pgfqpoint{3.957931in}{0.483000in}}%
\pgfpathclose%
\pgfusepath{fill}%
\end{pgfscope}%
\begin{pgfscope}%
\pgfpathrectangle{\pgfqpoint{3.815000in}{0.483000in}}{\pgfqpoint{2.835000in}{0.630000in}}%
\pgfusepath{clip}%
\pgfsetbuttcap%
\pgfsetmiterjoin%
\definecolor{currentfill}{rgb}{0.101961,0.101961,0.101961}%
\pgfsetfillcolor{currentfill}%
\pgfsetlinewidth{0.000000pt}%
\definecolor{currentstroke}{rgb}{0.000000,0.000000,0.000000}%
\pgfsetstrokecolor{currentstroke}%
\pgfsetstrokeopacity{0.000000}%
\pgfsetdash{}{0pt}%
\pgfpathmoveto{\pgfqpoint{3.981556in}{0.483000in}}%
\pgfpathlineto{\pgfqpoint{4.002819in}{0.483000in}}%
\pgfpathlineto{\pgfqpoint{4.002819in}{0.483000in}}%
\pgfpathlineto{\pgfqpoint{3.981556in}{0.483000in}}%
\pgfpathlineto{\pgfqpoint{3.981556in}{0.483000in}}%
\pgfpathclose%
\pgfusepath{fill}%
\end{pgfscope}%
\begin{pgfscope}%
\pgfpathrectangle{\pgfqpoint{3.815000in}{0.483000in}}{\pgfqpoint{2.835000in}{0.630000in}}%
\pgfusepath{clip}%
\pgfsetbuttcap%
\pgfsetmiterjoin%
\definecolor{currentfill}{rgb}{0.101961,0.101961,0.101961}%
\pgfsetfillcolor{currentfill}%
\pgfsetlinewidth{0.000000pt}%
\definecolor{currentstroke}{rgb}{0.000000,0.000000,0.000000}%
\pgfsetstrokecolor{currentstroke}%
\pgfsetstrokeopacity{0.000000}%
\pgfsetdash{}{0pt}%
\pgfpathmoveto{\pgfqpoint{4.005181in}{0.483000in}}%
\pgfpathlineto{\pgfqpoint{4.026444in}{0.483000in}}%
\pgfpathlineto{\pgfqpoint{4.026444in}{0.483000in}}%
\pgfpathlineto{\pgfqpoint{4.005181in}{0.483000in}}%
\pgfpathlineto{\pgfqpoint{4.005181in}{0.483000in}}%
\pgfpathclose%
\pgfusepath{fill}%
\end{pgfscope}%
\begin{pgfscope}%
\pgfpathrectangle{\pgfqpoint{3.815000in}{0.483000in}}{\pgfqpoint{2.835000in}{0.630000in}}%
\pgfusepath{clip}%
\pgfsetbuttcap%
\pgfsetmiterjoin%
\definecolor{currentfill}{rgb}{0.101961,0.101961,0.101961}%
\pgfsetfillcolor{currentfill}%
\pgfsetlinewidth{0.000000pt}%
\definecolor{currentstroke}{rgb}{0.000000,0.000000,0.000000}%
\pgfsetstrokecolor{currentstroke}%
\pgfsetstrokeopacity{0.000000}%
\pgfsetdash{}{0pt}%
\pgfpathmoveto{\pgfqpoint{4.028806in}{0.483000in}}%
\pgfpathlineto{\pgfqpoint{4.050069in}{0.483000in}}%
\pgfpathlineto{\pgfqpoint{4.050069in}{0.528000in}}%
\pgfpathlineto{\pgfqpoint{4.028806in}{0.528000in}}%
\pgfpathlineto{\pgfqpoint{4.028806in}{0.483000in}}%
\pgfpathclose%
\pgfusepath{fill}%
\end{pgfscope}%
\begin{pgfscope}%
\pgfpathrectangle{\pgfqpoint{3.815000in}{0.483000in}}{\pgfqpoint{2.835000in}{0.630000in}}%
\pgfusepath{clip}%
\pgfsetbuttcap%
\pgfsetmiterjoin%
\definecolor{currentfill}{rgb}{0.101961,0.101961,0.101961}%
\pgfsetfillcolor{currentfill}%
\pgfsetlinewidth{0.000000pt}%
\definecolor{currentstroke}{rgb}{0.000000,0.000000,0.000000}%
\pgfsetstrokecolor{currentstroke}%
\pgfsetstrokeopacity{0.000000}%
\pgfsetdash{}{0pt}%
\pgfpathmoveto{\pgfqpoint{4.052431in}{0.483000in}}%
\pgfpathlineto{\pgfqpoint{4.073694in}{0.483000in}}%
\pgfpathlineto{\pgfqpoint{4.073694in}{0.483000in}}%
\pgfpathlineto{\pgfqpoint{4.052431in}{0.483000in}}%
\pgfpathlineto{\pgfqpoint{4.052431in}{0.483000in}}%
\pgfpathclose%
\pgfusepath{fill}%
\end{pgfscope}%
\begin{pgfscope}%
\pgfpathrectangle{\pgfqpoint{3.815000in}{0.483000in}}{\pgfqpoint{2.835000in}{0.630000in}}%
\pgfusepath{clip}%
\pgfsetbuttcap%
\pgfsetmiterjoin%
\definecolor{currentfill}{rgb}{0.101961,0.101961,0.101961}%
\pgfsetfillcolor{currentfill}%
\pgfsetlinewidth{0.000000pt}%
\definecolor{currentstroke}{rgb}{0.000000,0.000000,0.000000}%
\pgfsetstrokecolor{currentstroke}%
\pgfsetstrokeopacity{0.000000}%
\pgfsetdash{}{0pt}%
\pgfpathmoveto{\pgfqpoint{4.076056in}{0.483000in}}%
\pgfpathlineto{\pgfqpoint{4.097319in}{0.483000in}}%
\pgfpathlineto{\pgfqpoint{4.097319in}{0.483000in}}%
\pgfpathlineto{\pgfqpoint{4.076056in}{0.483000in}}%
\pgfpathlineto{\pgfqpoint{4.076056in}{0.483000in}}%
\pgfpathclose%
\pgfusepath{fill}%
\end{pgfscope}%
\begin{pgfscope}%
\pgfpathrectangle{\pgfqpoint{3.815000in}{0.483000in}}{\pgfqpoint{2.835000in}{0.630000in}}%
\pgfusepath{clip}%
\pgfsetbuttcap%
\pgfsetmiterjoin%
\definecolor{currentfill}{rgb}{0.101961,0.101961,0.101961}%
\pgfsetfillcolor{currentfill}%
\pgfsetlinewidth{0.000000pt}%
\definecolor{currentstroke}{rgb}{0.000000,0.000000,0.000000}%
\pgfsetstrokecolor{currentstroke}%
\pgfsetstrokeopacity{0.000000}%
\pgfsetdash{}{0pt}%
\pgfpathmoveto{\pgfqpoint{4.099681in}{0.483000in}}%
\pgfpathlineto{\pgfqpoint{4.120944in}{0.483000in}}%
\pgfpathlineto{\pgfqpoint{4.120944in}{0.618000in}}%
\pgfpathlineto{\pgfqpoint{4.099681in}{0.618000in}}%
\pgfpathlineto{\pgfqpoint{4.099681in}{0.483000in}}%
\pgfpathclose%
\pgfusepath{fill}%
\end{pgfscope}%
\begin{pgfscope}%
\pgfpathrectangle{\pgfqpoint{3.815000in}{0.483000in}}{\pgfqpoint{2.835000in}{0.630000in}}%
\pgfusepath{clip}%
\pgfsetbuttcap%
\pgfsetmiterjoin%
\definecolor{currentfill}{rgb}{0.101961,0.101961,0.101961}%
\pgfsetfillcolor{currentfill}%
\pgfsetlinewidth{0.000000pt}%
\definecolor{currentstroke}{rgb}{0.000000,0.000000,0.000000}%
\pgfsetstrokecolor{currentstroke}%
\pgfsetstrokeopacity{0.000000}%
\pgfsetdash{}{0pt}%
\pgfpathmoveto{\pgfqpoint{4.123306in}{0.483000in}}%
\pgfpathlineto{\pgfqpoint{4.144569in}{0.483000in}}%
\pgfpathlineto{\pgfqpoint{4.144569in}{0.573000in}}%
\pgfpathlineto{\pgfqpoint{4.123306in}{0.573000in}}%
\pgfpathlineto{\pgfqpoint{4.123306in}{0.483000in}}%
\pgfpathclose%
\pgfusepath{fill}%
\end{pgfscope}%
\begin{pgfscope}%
\pgfpathrectangle{\pgfqpoint{3.815000in}{0.483000in}}{\pgfqpoint{2.835000in}{0.630000in}}%
\pgfusepath{clip}%
\pgfsetbuttcap%
\pgfsetmiterjoin%
\definecolor{currentfill}{rgb}{0.101961,0.101961,0.101961}%
\pgfsetfillcolor{currentfill}%
\pgfsetlinewidth{0.000000pt}%
\definecolor{currentstroke}{rgb}{0.000000,0.000000,0.000000}%
\pgfsetstrokecolor{currentstroke}%
\pgfsetstrokeopacity{0.000000}%
\pgfsetdash{}{0pt}%
\pgfpathmoveto{\pgfqpoint{4.146931in}{0.483000in}}%
\pgfpathlineto{\pgfqpoint{4.168194in}{0.483000in}}%
\pgfpathlineto{\pgfqpoint{4.168194in}{0.618000in}}%
\pgfpathlineto{\pgfqpoint{4.146931in}{0.618000in}}%
\pgfpathlineto{\pgfqpoint{4.146931in}{0.483000in}}%
\pgfpathclose%
\pgfusepath{fill}%
\end{pgfscope}%
\begin{pgfscope}%
\pgfpathrectangle{\pgfqpoint{3.815000in}{0.483000in}}{\pgfqpoint{2.835000in}{0.630000in}}%
\pgfusepath{clip}%
\pgfsetbuttcap%
\pgfsetmiterjoin%
\definecolor{currentfill}{rgb}{0.101961,0.101961,0.101961}%
\pgfsetfillcolor{currentfill}%
\pgfsetlinewidth{0.000000pt}%
\definecolor{currentstroke}{rgb}{0.000000,0.000000,0.000000}%
\pgfsetstrokecolor{currentstroke}%
\pgfsetstrokeopacity{0.000000}%
\pgfsetdash{}{0pt}%
\pgfpathmoveto{\pgfqpoint{4.170556in}{0.483000in}}%
\pgfpathlineto{\pgfqpoint{4.191819in}{0.483000in}}%
\pgfpathlineto{\pgfqpoint{4.191819in}{0.663000in}}%
\pgfpathlineto{\pgfqpoint{4.170556in}{0.663000in}}%
\pgfpathlineto{\pgfqpoint{4.170556in}{0.483000in}}%
\pgfpathclose%
\pgfusepath{fill}%
\end{pgfscope}%
\begin{pgfscope}%
\pgfpathrectangle{\pgfqpoint{3.815000in}{0.483000in}}{\pgfqpoint{2.835000in}{0.630000in}}%
\pgfusepath{clip}%
\pgfsetbuttcap%
\pgfsetmiterjoin%
\definecolor{currentfill}{rgb}{0.101961,0.101961,0.101961}%
\pgfsetfillcolor{currentfill}%
\pgfsetlinewidth{0.000000pt}%
\definecolor{currentstroke}{rgb}{0.000000,0.000000,0.000000}%
\pgfsetstrokecolor{currentstroke}%
\pgfsetstrokeopacity{0.000000}%
\pgfsetdash{}{0pt}%
\pgfpathmoveto{\pgfqpoint{4.194181in}{0.483000in}}%
\pgfpathlineto{\pgfqpoint{4.215444in}{0.483000in}}%
\pgfpathlineto{\pgfqpoint{4.215444in}{0.573000in}}%
\pgfpathlineto{\pgfqpoint{4.194181in}{0.573000in}}%
\pgfpathlineto{\pgfqpoint{4.194181in}{0.483000in}}%
\pgfpathclose%
\pgfusepath{fill}%
\end{pgfscope}%
\begin{pgfscope}%
\pgfpathrectangle{\pgfqpoint{3.815000in}{0.483000in}}{\pgfqpoint{2.835000in}{0.630000in}}%
\pgfusepath{clip}%
\pgfsetbuttcap%
\pgfsetmiterjoin%
\definecolor{currentfill}{rgb}{0.101961,0.101961,0.101961}%
\pgfsetfillcolor{currentfill}%
\pgfsetlinewidth{0.000000pt}%
\definecolor{currentstroke}{rgb}{0.000000,0.000000,0.000000}%
\pgfsetstrokecolor{currentstroke}%
\pgfsetstrokeopacity{0.000000}%
\pgfsetdash{}{0pt}%
\pgfpathmoveto{\pgfqpoint{4.217806in}{0.483000in}}%
\pgfpathlineto{\pgfqpoint{4.239069in}{0.483000in}}%
\pgfpathlineto{\pgfqpoint{4.239069in}{0.483000in}}%
\pgfpathlineto{\pgfqpoint{4.217806in}{0.483000in}}%
\pgfpathlineto{\pgfqpoint{4.217806in}{0.483000in}}%
\pgfpathclose%
\pgfusepath{fill}%
\end{pgfscope}%
\begin{pgfscope}%
\pgfpathrectangle{\pgfqpoint{3.815000in}{0.483000in}}{\pgfqpoint{2.835000in}{0.630000in}}%
\pgfusepath{clip}%
\pgfsetbuttcap%
\pgfsetmiterjoin%
\definecolor{currentfill}{rgb}{0.101961,0.101961,0.101961}%
\pgfsetfillcolor{currentfill}%
\pgfsetlinewidth{0.000000pt}%
\definecolor{currentstroke}{rgb}{0.000000,0.000000,0.000000}%
\pgfsetstrokecolor{currentstroke}%
\pgfsetstrokeopacity{0.000000}%
\pgfsetdash{}{0pt}%
\pgfpathmoveto{\pgfqpoint{4.241431in}{0.483000in}}%
\pgfpathlineto{\pgfqpoint{4.262694in}{0.483000in}}%
\pgfpathlineto{\pgfqpoint{4.262694in}{0.663000in}}%
\pgfpathlineto{\pgfqpoint{4.241431in}{0.663000in}}%
\pgfpathlineto{\pgfqpoint{4.241431in}{0.483000in}}%
\pgfpathclose%
\pgfusepath{fill}%
\end{pgfscope}%
\begin{pgfscope}%
\pgfpathrectangle{\pgfqpoint{3.815000in}{0.483000in}}{\pgfqpoint{2.835000in}{0.630000in}}%
\pgfusepath{clip}%
\pgfsetbuttcap%
\pgfsetmiterjoin%
\definecolor{currentfill}{rgb}{0.101961,0.101961,0.101961}%
\pgfsetfillcolor{currentfill}%
\pgfsetlinewidth{0.000000pt}%
\definecolor{currentstroke}{rgb}{0.000000,0.000000,0.000000}%
\pgfsetstrokecolor{currentstroke}%
\pgfsetstrokeopacity{0.000000}%
\pgfsetdash{}{0pt}%
\pgfpathmoveto{\pgfqpoint{4.265056in}{0.483000in}}%
\pgfpathlineto{\pgfqpoint{4.286319in}{0.483000in}}%
\pgfpathlineto{\pgfqpoint{4.286319in}{0.483000in}}%
\pgfpathlineto{\pgfqpoint{4.265056in}{0.483000in}}%
\pgfpathlineto{\pgfqpoint{4.265056in}{0.483000in}}%
\pgfpathclose%
\pgfusepath{fill}%
\end{pgfscope}%
\begin{pgfscope}%
\pgfpathrectangle{\pgfqpoint{3.815000in}{0.483000in}}{\pgfqpoint{2.835000in}{0.630000in}}%
\pgfusepath{clip}%
\pgfsetbuttcap%
\pgfsetmiterjoin%
\definecolor{currentfill}{rgb}{0.101961,0.101961,0.101961}%
\pgfsetfillcolor{currentfill}%
\pgfsetlinewidth{0.000000pt}%
\definecolor{currentstroke}{rgb}{0.000000,0.000000,0.000000}%
\pgfsetstrokecolor{currentstroke}%
\pgfsetstrokeopacity{0.000000}%
\pgfsetdash{}{0pt}%
\pgfpathmoveto{\pgfqpoint{4.288681in}{0.483000in}}%
\pgfpathlineto{\pgfqpoint{4.309944in}{0.483000in}}%
\pgfpathlineto{\pgfqpoint{4.309944in}{0.528000in}}%
\pgfpathlineto{\pgfqpoint{4.288681in}{0.528000in}}%
\pgfpathlineto{\pgfqpoint{4.288681in}{0.483000in}}%
\pgfpathclose%
\pgfusepath{fill}%
\end{pgfscope}%
\begin{pgfscope}%
\pgfpathrectangle{\pgfqpoint{3.815000in}{0.483000in}}{\pgfqpoint{2.835000in}{0.630000in}}%
\pgfusepath{clip}%
\pgfsetbuttcap%
\pgfsetmiterjoin%
\definecolor{currentfill}{rgb}{0.101961,0.101961,0.101961}%
\pgfsetfillcolor{currentfill}%
\pgfsetlinewidth{0.000000pt}%
\definecolor{currentstroke}{rgb}{0.000000,0.000000,0.000000}%
\pgfsetstrokecolor{currentstroke}%
\pgfsetstrokeopacity{0.000000}%
\pgfsetdash{}{0pt}%
\pgfpathmoveto{\pgfqpoint{4.312306in}{0.483000in}}%
\pgfpathlineto{\pgfqpoint{4.333569in}{0.483000in}}%
\pgfpathlineto{\pgfqpoint{4.333569in}{0.663000in}}%
\pgfpathlineto{\pgfqpoint{4.312306in}{0.663000in}}%
\pgfpathlineto{\pgfqpoint{4.312306in}{0.483000in}}%
\pgfpathclose%
\pgfusepath{fill}%
\end{pgfscope}%
\begin{pgfscope}%
\pgfpathrectangle{\pgfqpoint{3.815000in}{0.483000in}}{\pgfqpoint{2.835000in}{0.630000in}}%
\pgfusepath{clip}%
\pgfsetbuttcap%
\pgfsetmiterjoin%
\definecolor{currentfill}{rgb}{0.101961,0.101961,0.101961}%
\pgfsetfillcolor{currentfill}%
\pgfsetlinewidth{0.000000pt}%
\definecolor{currentstroke}{rgb}{0.000000,0.000000,0.000000}%
\pgfsetstrokecolor{currentstroke}%
\pgfsetstrokeopacity{0.000000}%
\pgfsetdash{}{0pt}%
\pgfpathmoveto{\pgfqpoint{4.335931in}{0.483000in}}%
\pgfpathlineto{\pgfqpoint{4.357194in}{0.483000in}}%
\pgfpathlineto{\pgfqpoint{4.357194in}{0.483000in}}%
\pgfpathlineto{\pgfqpoint{4.335931in}{0.483000in}}%
\pgfpathlineto{\pgfqpoint{4.335931in}{0.483000in}}%
\pgfpathclose%
\pgfusepath{fill}%
\end{pgfscope}%
\begin{pgfscope}%
\pgfpathrectangle{\pgfqpoint{3.815000in}{0.483000in}}{\pgfqpoint{2.835000in}{0.630000in}}%
\pgfusepath{clip}%
\pgfsetbuttcap%
\pgfsetmiterjoin%
\definecolor{currentfill}{rgb}{0.101961,0.101961,0.101961}%
\pgfsetfillcolor{currentfill}%
\pgfsetlinewidth{0.000000pt}%
\definecolor{currentstroke}{rgb}{0.000000,0.000000,0.000000}%
\pgfsetstrokecolor{currentstroke}%
\pgfsetstrokeopacity{0.000000}%
\pgfsetdash{}{0pt}%
\pgfpathmoveto{\pgfqpoint{4.359556in}{0.483000in}}%
\pgfpathlineto{\pgfqpoint{4.380819in}{0.483000in}}%
\pgfpathlineto{\pgfqpoint{4.380819in}{0.573000in}}%
\pgfpathlineto{\pgfqpoint{4.359556in}{0.573000in}}%
\pgfpathlineto{\pgfqpoint{4.359556in}{0.483000in}}%
\pgfpathclose%
\pgfusepath{fill}%
\end{pgfscope}%
\begin{pgfscope}%
\pgfpathrectangle{\pgfqpoint{3.815000in}{0.483000in}}{\pgfqpoint{2.835000in}{0.630000in}}%
\pgfusepath{clip}%
\pgfsetbuttcap%
\pgfsetmiterjoin%
\definecolor{currentfill}{rgb}{0.101961,0.101961,0.101961}%
\pgfsetfillcolor{currentfill}%
\pgfsetlinewidth{0.000000pt}%
\definecolor{currentstroke}{rgb}{0.000000,0.000000,0.000000}%
\pgfsetstrokecolor{currentstroke}%
\pgfsetstrokeopacity{0.000000}%
\pgfsetdash{}{0pt}%
\pgfpathmoveto{\pgfqpoint{4.383181in}{0.483000in}}%
\pgfpathlineto{\pgfqpoint{4.404444in}{0.483000in}}%
\pgfpathlineto{\pgfqpoint{4.404444in}{0.618000in}}%
\pgfpathlineto{\pgfqpoint{4.383181in}{0.618000in}}%
\pgfpathlineto{\pgfqpoint{4.383181in}{0.483000in}}%
\pgfpathclose%
\pgfusepath{fill}%
\end{pgfscope}%
\begin{pgfscope}%
\pgfpathrectangle{\pgfqpoint{3.815000in}{0.483000in}}{\pgfqpoint{2.835000in}{0.630000in}}%
\pgfusepath{clip}%
\pgfsetbuttcap%
\pgfsetmiterjoin%
\definecolor{currentfill}{rgb}{0.101961,0.101961,0.101961}%
\pgfsetfillcolor{currentfill}%
\pgfsetlinewidth{0.000000pt}%
\definecolor{currentstroke}{rgb}{0.000000,0.000000,0.000000}%
\pgfsetstrokecolor{currentstroke}%
\pgfsetstrokeopacity{0.000000}%
\pgfsetdash{}{0pt}%
\pgfpathmoveto{\pgfqpoint{4.406806in}{0.483000in}}%
\pgfpathlineto{\pgfqpoint{4.428069in}{0.483000in}}%
\pgfpathlineto{\pgfqpoint{4.428069in}{0.483000in}}%
\pgfpathlineto{\pgfqpoint{4.406806in}{0.483000in}}%
\pgfpathlineto{\pgfqpoint{4.406806in}{0.483000in}}%
\pgfpathclose%
\pgfusepath{fill}%
\end{pgfscope}%
\begin{pgfscope}%
\pgfpathrectangle{\pgfqpoint{3.815000in}{0.483000in}}{\pgfqpoint{2.835000in}{0.630000in}}%
\pgfusepath{clip}%
\pgfsetbuttcap%
\pgfsetmiterjoin%
\definecolor{currentfill}{rgb}{0.101961,0.101961,0.101961}%
\pgfsetfillcolor{currentfill}%
\pgfsetlinewidth{0.000000pt}%
\definecolor{currentstroke}{rgb}{0.000000,0.000000,0.000000}%
\pgfsetstrokecolor{currentstroke}%
\pgfsetstrokeopacity{0.000000}%
\pgfsetdash{}{0pt}%
\pgfpathmoveto{\pgfqpoint{4.430431in}{0.483000in}}%
\pgfpathlineto{\pgfqpoint{4.451694in}{0.483000in}}%
\pgfpathlineto{\pgfqpoint{4.451694in}{0.528000in}}%
\pgfpathlineto{\pgfqpoint{4.430431in}{0.528000in}}%
\pgfpathlineto{\pgfqpoint{4.430431in}{0.483000in}}%
\pgfpathclose%
\pgfusepath{fill}%
\end{pgfscope}%
\begin{pgfscope}%
\pgfpathrectangle{\pgfqpoint{3.815000in}{0.483000in}}{\pgfqpoint{2.835000in}{0.630000in}}%
\pgfusepath{clip}%
\pgfsetbuttcap%
\pgfsetmiterjoin%
\definecolor{currentfill}{rgb}{0.101961,0.101961,0.101961}%
\pgfsetfillcolor{currentfill}%
\pgfsetlinewidth{0.000000pt}%
\definecolor{currentstroke}{rgb}{0.000000,0.000000,0.000000}%
\pgfsetstrokecolor{currentstroke}%
\pgfsetstrokeopacity{0.000000}%
\pgfsetdash{}{0pt}%
\pgfpathmoveto{\pgfqpoint{4.454056in}{0.483000in}}%
\pgfpathlineto{\pgfqpoint{4.475319in}{0.483000in}}%
\pgfpathlineto{\pgfqpoint{4.475319in}{0.573000in}}%
\pgfpathlineto{\pgfqpoint{4.454056in}{0.573000in}}%
\pgfpathlineto{\pgfqpoint{4.454056in}{0.483000in}}%
\pgfpathclose%
\pgfusepath{fill}%
\end{pgfscope}%
\begin{pgfscope}%
\pgfpathrectangle{\pgfqpoint{3.815000in}{0.483000in}}{\pgfqpoint{2.835000in}{0.630000in}}%
\pgfusepath{clip}%
\pgfsetbuttcap%
\pgfsetmiterjoin%
\definecolor{currentfill}{rgb}{0.101961,0.101961,0.101961}%
\pgfsetfillcolor{currentfill}%
\pgfsetlinewidth{0.000000pt}%
\definecolor{currentstroke}{rgb}{0.000000,0.000000,0.000000}%
\pgfsetstrokecolor{currentstroke}%
\pgfsetstrokeopacity{0.000000}%
\pgfsetdash{}{0pt}%
\pgfpathmoveto{\pgfqpoint{4.477681in}{0.483000in}}%
\pgfpathlineto{\pgfqpoint{4.498944in}{0.483000in}}%
\pgfpathlineto{\pgfqpoint{4.498944in}{0.483000in}}%
\pgfpathlineto{\pgfqpoint{4.477681in}{0.483000in}}%
\pgfpathlineto{\pgfqpoint{4.477681in}{0.483000in}}%
\pgfpathclose%
\pgfusepath{fill}%
\end{pgfscope}%
\begin{pgfscope}%
\pgfpathrectangle{\pgfqpoint{3.815000in}{0.483000in}}{\pgfqpoint{2.835000in}{0.630000in}}%
\pgfusepath{clip}%
\pgfsetbuttcap%
\pgfsetmiterjoin%
\definecolor{currentfill}{rgb}{0.101961,0.101961,0.101961}%
\pgfsetfillcolor{currentfill}%
\pgfsetlinewidth{0.000000pt}%
\definecolor{currentstroke}{rgb}{0.000000,0.000000,0.000000}%
\pgfsetstrokecolor{currentstroke}%
\pgfsetstrokeopacity{0.000000}%
\pgfsetdash{}{0pt}%
\pgfpathmoveto{\pgfqpoint{4.501306in}{0.483000in}}%
\pgfpathlineto{\pgfqpoint{4.522569in}{0.483000in}}%
\pgfpathlineto{\pgfqpoint{4.522569in}{0.483000in}}%
\pgfpathlineto{\pgfqpoint{4.501306in}{0.483000in}}%
\pgfpathlineto{\pgfqpoint{4.501306in}{0.483000in}}%
\pgfpathclose%
\pgfusepath{fill}%
\end{pgfscope}%
\begin{pgfscope}%
\pgfpathrectangle{\pgfqpoint{3.815000in}{0.483000in}}{\pgfqpoint{2.835000in}{0.630000in}}%
\pgfusepath{clip}%
\pgfsetbuttcap%
\pgfsetmiterjoin%
\definecolor{currentfill}{rgb}{0.101961,0.101961,0.101961}%
\pgfsetfillcolor{currentfill}%
\pgfsetlinewidth{0.000000pt}%
\definecolor{currentstroke}{rgb}{0.000000,0.000000,0.000000}%
\pgfsetstrokecolor{currentstroke}%
\pgfsetstrokeopacity{0.000000}%
\pgfsetdash{}{0pt}%
\pgfpathmoveto{\pgfqpoint{4.524931in}{0.483000in}}%
\pgfpathlineto{\pgfqpoint{4.546194in}{0.483000in}}%
\pgfpathlineto{\pgfqpoint{4.546194in}{0.528000in}}%
\pgfpathlineto{\pgfqpoint{4.524931in}{0.528000in}}%
\pgfpathlineto{\pgfqpoint{4.524931in}{0.483000in}}%
\pgfpathclose%
\pgfusepath{fill}%
\end{pgfscope}%
\begin{pgfscope}%
\pgfpathrectangle{\pgfqpoint{3.815000in}{0.483000in}}{\pgfqpoint{2.835000in}{0.630000in}}%
\pgfusepath{clip}%
\pgfsetbuttcap%
\pgfsetmiterjoin%
\definecolor{currentfill}{rgb}{0.101961,0.101961,0.101961}%
\pgfsetfillcolor{currentfill}%
\pgfsetlinewidth{0.000000pt}%
\definecolor{currentstroke}{rgb}{0.000000,0.000000,0.000000}%
\pgfsetstrokecolor{currentstroke}%
\pgfsetstrokeopacity{0.000000}%
\pgfsetdash{}{0pt}%
\pgfpathmoveto{\pgfqpoint{4.548556in}{0.483000in}}%
\pgfpathlineto{\pgfqpoint{4.569819in}{0.483000in}}%
\pgfpathlineto{\pgfqpoint{4.569819in}{0.708000in}}%
\pgfpathlineto{\pgfqpoint{4.548556in}{0.708000in}}%
\pgfpathlineto{\pgfqpoint{4.548556in}{0.483000in}}%
\pgfpathclose%
\pgfusepath{fill}%
\end{pgfscope}%
\begin{pgfscope}%
\pgfpathrectangle{\pgfqpoint{3.815000in}{0.483000in}}{\pgfqpoint{2.835000in}{0.630000in}}%
\pgfusepath{clip}%
\pgfsetbuttcap%
\pgfsetmiterjoin%
\definecolor{currentfill}{rgb}{0.101961,0.101961,0.101961}%
\pgfsetfillcolor{currentfill}%
\pgfsetlinewidth{0.000000pt}%
\definecolor{currentstroke}{rgb}{0.000000,0.000000,0.000000}%
\pgfsetstrokecolor{currentstroke}%
\pgfsetstrokeopacity{0.000000}%
\pgfsetdash{}{0pt}%
\pgfpathmoveto{\pgfqpoint{4.572181in}{0.483000in}}%
\pgfpathlineto{\pgfqpoint{4.593444in}{0.483000in}}%
\pgfpathlineto{\pgfqpoint{4.593444in}{0.843000in}}%
\pgfpathlineto{\pgfqpoint{4.572181in}{0.843000in}}%
\pgfpathlineto{\pgfqpoint{4.572181in}{0.483000in}}%
\pgfpathclose%
\pgfusepath{fill}%
\end{pgfscope}%
\begin{pgfscope}%
\pgfpathrectangle{\pgfqpoint{3.815000in}{0.483000in}}{\pgfqpoint{2.835000in}{0.630000in}}%
\pgfusepath{clip}%
\pgfsetbuttcap%
\pgfsetmiterjoin%
\definecolor{currentfill}{rgb}{0.101961,0.101961,0.101961}%
\pgfsetfillcolor{currentfill}%
\pgfsetlinewidth{0.000000pt}%
\definecolor{currentstroke}{rgb}{0.000000,0.000000,0.000000}%
\pgfsetstrokecolor{currentstroke}%
\pgfsetstrokeopacity{0.000000}%
\pgfsetdash{}{0pt}%
\pgfpathmoveto{\pgfqpoint{4.595806in}{0.483000in}}%
\pgfpathlineto{\pgfqpoint{4.617069in}{0.483000in}}%
\pgfpathlineto{\pgfqpoint{4.617069in}{0.483000in}}%
\pgfpathlineto{\pgfqpoint{4.595806in}{0.483000in}}%
\pgfpathlineto{\pgfqpoint{4.595806in}{0.483000in}}%
\pgfpathclose%
\pgfusepath{fill}%
\end{pgfscope}%
\begin{pgfscope}%
\pgfpathrectangle{\pgfqpoint{3.815000in}{0.483000in}}{\pgfqpoint{2.835000in}{0.630000in}}%
\pgfusepath{clip}%
\pgfsetbuttcap%
\pgfsetmiterjoin%
\definecolor{currentfill}{rgb}{0.101961,0.101961,0.101961}%
\pgfsetfillcolor{currentfill}%
\pgfsetlinewidth{0.000000pt}%
\definecolor{currentstroke}{rgb}{0.000000,0.000000,0.000000}%
\pgfsetstrokecolor{currentstroke}%
\pgfsetstrokeopacity{0.000000}%
\pgfsetdash{}{0pt}%
\pgfpathmoveto{\pgfqpoint{4.619431in}{0.483000in}}%
\pgfpathlineto{\pgfqpoint{4.640694in}{0.483000in}}%
\pgfpathlineto{\pgfqpoint{4.640694in}{0.573000in}}%
\pgfpathlineto{\pgfqpoint{4.619431in}{0.573000in}}%
\pgfpathlineto{\pgfqpoint{4.619431in}{0.483000in}}%
\pgfpathclose%
\pgfusepath{fill}%
\end{pgfscope}%
\begin{pgfscope}%
\pgfpathrectangle{\pgfqpoint{3.815000in}{0.483000in}}{\pgfqpoint{2.835000in}{0.630000in}}%
\pgfusepath{clip}%
\pgfsetbuttcap%
\pgfsetmiterjoin%
\definecolor{currentfill}{rgb}{0.101961,0.101961,0.101961}%
\pgfsetfillcolor{currentfill}%
\pgfsetlinewidth{0.000000pt}%
\definecolor{currentstroke}{rgb}{0.000000,0.000000,0.000000}%
\pgfsetstrokecolor{currentstroke}%
\pgfsetstrokeopacity{0.000000}%
\pgfsetdash{}{0pt}%
\pgfpathmoveto{\pgfqpoint{4.643056in}{0.483000in}}%
\pgfpathlineto{\pgfqpoint{4.664319in}{0.483000in}}%
\pgfpathlineto{\pgfqpoint{4.664319in}{0.483000in}}%
\pgfpathlineto{\pgfqpoint{4.643056in}{0.483000in}}%
\pgfpathlineto{\pgfqpoint{4.643056in}{0.483000in}}%
\pgfpathclose%
\pgfusepath{fill}%
\end{pgfscope}%
\begin{pgfscope}%
\pgfpathrectangle{\pgfqpoint{3.815000in}{0.483000in}}{\pgfqpoint{2.835000in}{0.630000in}}%
\pgfusepath{clip}%
\pgfsetbuttcap%
\pgfsetmiterjoin%
\definecolor{currentfill}{rgb}{0.101961,0.101961,0.101961}%
\pgfsetfillcolor{currentfill}%
\pgfsetlinewidth{0.000000pt}%
\definecolor{currentstroke}{rgb}{0.000000,0.000000,0.000000}%
\pgfsetstrokecolor{currentstroke}%
\pgfsetstrokeopacity{0.000000}%
\pgfsetdash{}{0pt}%
\pgfpathmoveto{\pgfqpoint{4.666681in}{0.483000in}}%
\pgfpathlineto{\pgfqpoint{4.687944in}{0.483000in}}%
\pgfpathlineto{\pgfqpoint{4.687944in}{0.528000in}}%
\pgfpathlineto{\pgfqpoint{4.666681in}{0.528000in}}%
\pgfpathlineto{\pgfqpoint{4.666681in}{0.483000in}}%
\pgfpathclose%
\pgfusepath{fill}%
\end{pgfscope}%
\begin{pgfscope}%
\pgfpathrectangle{\pgfqpoint{3.815000in}{0.483000in}}{\pgfqpoint{2.835000in}{0.630000in}}%
\pgfusepath{clip}%
\pgfsetbuttcap%
\pgfsetmiterjoin%
\definecolor{currentfill}{rgb}{0.101961,0.101961,0.101961}%
\pgfsetfillcolor{currentfill}%
\pgfsetlinewidth{0.000000pt}%
\definecolor{currentstroke}{rgb}{0.000000,0.000000,0.000000}%
\pgfsetstrokecolor{currentstroke}%
\pgfsetstrokeopacity{0.000000}%
\pgfsetdash{}{0pt}%
\pgfpathmoveto{\pgfqpoint{4.690306in}{0.483000in}}%
\pgfpathlineto{\pgfqpoint{4.711569in}{0.483000in}}%
\pgfpathlineto{\pgfqpoint{4.711569in}{0.573000in}}%
\pgfpathlineto{\pgfqpoint{4.690306in}{0.573000in}}%
\pgfpathlineto{\pgfqpoint{4.690306in}{0.483000in}}%
\pgfpathclose%
\pgfusepath{fill}%
\end{pgfscope}%
\begin{pgfscope}%
\pgfpathrectangle{\pgfqpoint{3.815000in}{0.483000in}}{\pgfqpoint{2.835000in}{0.630000in}}%
\pgfusepath{clip}%
\pgfsetbuttcap%
\pgfsetmiterjoin%
\definecolor{currentfill}{rgb}{0.101961,0.101961,0.101961}%
\pgfsetfillcolor{currentfill}%
\pgfsetlinewidth{0.000000pt}%
\definecolor{currentstroke}{rgb}{0.000000,0.000000,0.000000}%
\pgfsetstrokecolor{currentstroke}%
\pgfsetstrokeopacity{0.000000}%
\pgfsetdash{}{0pt}%
\pgfpathmoveto{\pgfqpoint{4.713931in}{0.483000in}}%
\pgfpathlineto{\pgfqpoint{4.735194in}{0.483000in}}%
\pgfpathlineto{\pgfqpoint{4.735194in}{0.528000in}}%
\pgfpathlineto{\pgfqpoint{4.713931in}{0.528000in}}%
\pgfpathlineto{\pgfqpoint{4.713931in}{0.483000in}}%
\pgfpathclose%
\pgfusepath{fill}%
\end{pgfscope}%
\begin{pgfscope}%
\pgfpathrectangle{\pgfqpoint{3.815000in}{0.483000in}}{\pgfqpoint{2.835000in}{0.630000in}}%
\pgfusepath{clip}%
\pgfsetbuttcap%
\pgfsetmiterjoin%
\definecolor{currentfill}{rgb}{0.101961,0.101961,0.101961}%
\pgfsetfillcolor{currentfill}%
\pgfsetlinewidth{0.000000pt}%
\definecolor{currentstroke}{rgb}{0.000000,0.000000,0.000000}%
\pgfsetstrokecolor{currentstroke}%
\pgfsetstrokeopacity{0.000000}%
\pgfsetdash{}{0pt}%
\pgfpathmoveto{\pgfqpoint{4.737556in}{0.483000in}}%
\pgfpathlineto{\pgfqpoint{4.758819in}{0.483000in}}%
\pgfpathlineto{\pgfqpoint{4.758819in}{0.483000in}}%
\pgfpathlineto{\pgfqpoint{4.737556in}{0.483000in}}%
\pgfpathlineto{\pgfqpoint{4.737556in}{0.483000in}}%
\pgfpathclose%
\pgfusepath{fill}%
\end{pgfscope}%
\begin{pgfscope}%
\pgfpathrectangle{\pgfqpoint{3.815000in}{0.483000in}}{\pgfqpoint{2.835000in}{0.630000in}}%
\pgfusepath{clip}%
\pgfsetbuttcap%
\pgfsetmiterjoin%
\definecolor{currentfill}{rgb}{0.101961,0.101961,0.101961}%
\pgfsetfillcolor{currentfill}%
\pgfsetlinewidth{0.000000pt}%
\definecolor{currentstroke}{rgb}{0.000000,0.000000,0.000000}%
\pgfsetstrokecolor{currentstroke}%
\pgfsetstrokeopacity{0.000000}%
\pgfsetdash{}{0pt}%
\pgfpathmoveto{\pgfqpoint{4.761181in}{0.483000in}}%
\pgfpathlineto{\pgfqpoint{4.782444in}{0.483000in}}%
\pgfpathlineto{\pgfqpoint{4.782444in}{0.528000in}}%
\pgfpathlineto{\pgfqpoint{4.761181in}{0.528000in}}%
\pgfpathlineto{\pgfqpoint{4.761181in}{0.483000in}}%
\pgfpathclose%
\pgfusepath{fill}%
\end{pgfscope}%
\begin{pgfscope}%
\pgfpathrectangle{\pgfqpoint{3.815000in}{0.483000in}}{\pgfqpoint{2.835000in}{0.630000in}}%
\pgfusepath{clip}%
\pgfsetbuttcap%
\pgfsetmiterjoin%
\definecolor{currentfill}{rgb}{0.101961,0.101961,0.101961}%
\pgfsetfillcolor{currentfill}%
\pgfsetlinewidth{0.000000pt}%
\definecolor{currentstroke}{rgb}{0.000000,0.000000,0.000000}%
\pgfsetstrokecolor{currentstroke}%
\pgfsetstrokeopacity{0.000000}%
\pgfsetdash{}{0pt}%
\pgfpathmoveto{\pgfqpoint{4.784806in}{0.483000in}}%
\pgfpathlineto{\pgfqpoint{4.806069in}{0.483000in}}%
\pgfpathlineto{\pgfqpoint{4.806069in}{0.483000in}}%
\pgfpathlineto{\pgfqpoint{4.784806in}{0.483000in}}%
\pgfpathlineto{\pgfqpoint{4.784806in}{0.483000in}}%
\pgfpathclose%
\pgfusepath{fill}%
\end{pgfscope}%
\begin{pgfscope}%
\pgfpathrectangle{\pgfqpoint{3.815000in}{0.483000in}}{\pgfqpoint{2.835000in}{0.630000in}}%
\pgfusepath{clip}%
\pgfsetbuttcap%
\pgfsetmiterjoin%
\definecolor{currentfill}{rgb}{0.101961,0.101961,0.101961}%
\pgfsetfillcolor{currentfill}%
\pgfsetlinewidth{0.000000pt}%
\definecolor{currentstroke}{rgb}{0.000000,0.000000,0.000000}%
\pgfsetstrokecolor{currentstroke}%
\pgfsetstrokeopacity{0.000000}%
\pgfsetdash{}{0pt}%
\pgfpathmoveto{\pgfqpoint{4.808431in}{0.483000in}}%
\pgfpathlineto{\pgfqpoint{4.829694in}{0.483000in}}%
\pgfpathlineto{\pgfqpoint{4.829694in}{0.483000in}}%
\pgfpathlineto{\pgfqpoint{4.808431in}{0.483000in}}%
\pgfpathlineto{\pgfqpoint{4.808431in}{0.483000in}}%
\pgfpathclose%
\pgfusepath{fill}%
\end{pgfscope}%
\begin{pgfscope}%
\pgfpathrectangle{\pgfqpoint{3.815000in}{0.483000in}}{\pgfqpoint{2.835000in}{0.630000in}}%
\pgfusepath{clip}%
\pgfsetbuttcap%
\pgfsetmiterjoin%
\definecolor{currentfill}{rgb}{0.101961,0.101961,0.101961}%
\pgfsetfillcolor{currentfill}%
\pgfsetlinewidth{0.000000pt}%
\definecolor{currentstroke}{rgb}{0.000000,0.000000,0.000000}%
\pgfsetstrokecolor{currentstroke}%
\pgfsetstrokeopacity{0.000000}%
\pgfsetdash{}{0pt}%
\pgfpathmoveto{\pgfqpoint{4.832056in}{0.483000in}}%
\pgfpathlineto{\pgfqpoint{4.853319in}{0.483000in}}%
\pgfpathlineto{\pgfqpoint{4.853319in}{0.483000in}}%
\pgfpathlineto{\pgfqpoint{4.832056in}{0.483000in}}%
\pgfpathlineto{\pgfqpoint{4.832056in}{0.483000in}}%
\pgfpathclose%
\pgfusepath{fill}%
\end{pgfscope}%
\begin{pgfscope}%
\pgfpathrectangle{\pgfqpoint{3.815000in}{0.483000in}}{\pgfqpoint{2.835000in}{0.630000in}}%
\pgfusepath{clip}%
\pgfsetbuttcap%
\pgfsetmiterjoin%
\definecolor{currentfill}{rgb}{0.101961,0.101961,0.101961}%
\pgfsetfillcolor{currentfill}%
\pgfsetlinewidth{0.000000pt}%
\definecolor{currentstroke}{rgb}{0.000000,0.000000,0.000000}%
\pgfsetstrokecolor{currentstroke}%
\pgfsetstrokeopacity{0.000000}%
\pgfsetdash{}{0pt}%
\pgfpathmoveto{\pgfqpoint{4.855681in}{0.483000in}}%
\pgfpathlineto{\pgfqpoint{4.876944in}{0.483000in}}%
\pgfpathlineto{\pgfqpoint{4.876944in}{0.483000in}}%
\pgfpathlineto{\pgfqpoint{4.855681in}{0.483000in}}%
\pgfpathlineto{\pgfqpoint{4.855681in}{0.483000in}}%
\pgfpathclose%
\pgfusepath{fill}%
\end{pgfscope}%
\begin{pgfscope}%
\pgfpathrectangle{\pgfqpoint{3.815000in}{0.483000in}}{\pgfqpoint{2.835000in}{0.630000in}}%
\pgfusepath{clip}%
\pgfsetbuttcap%
\pgfsetmiterjoin%
\definecolor{currentfill}{rgb}{0.101961,0.101961,0.101961}%
\pgfsetfillcolor{currentfill}%
\pgfsetlinewidth{0.000000pt}%
\definecolor{currentstroke}{rgb}{0.000000,0.000000,0.000000}%
\pgfsetstrokecolor{currentstroke}%
\pgfsetstrokeopacity{0.000000}%
\pgfsetdash{}{0pt}%
\pgfpathmoveto{\pgfqpoint{4.879306in}{0.483000in}}%
\pgfpathlineto{\pgfqpoint{4.900569in}{0.483000in}}%
\pgfpathlineto{\pgfqpoint{4.900569in}{0.483000in}}%
\pgfpathlineto{\pgfqpoint{4.879306in}{0.483000in}}%
\pgfpathlineto{\pgfqpoint{4.879306in}{0.483000in}}%
\pgfpathclose%
\pgfusepath{fill}%
\end{pgfscope}%
\begin{pgfscope}%
\pgfpathrectangle{\pgfqpoint{3.815000in}{0.483000in}}{\pgfqpoint{2.835000in}{0.630000in}}%
\pgfusepath{clip}%
\pgfsetbuttcap%
\pgfsetmiterjoin%
\definecolor{currentfill}{rgb}{0.101961,0.101961,0.101961}%
\pgfsetfillcolor{currentfill}%
\pgfsetlinewidth{0.000000pt}%
\definecolor{currentstroke}{rgb}{0.000000,0.000000,0.000000}%
\pgfsetstrokecolor{currentstroke}%
\pgfsetstrokeopacity{0.000000}%
\pgfsetdash{}{0pt}%
\pgfpathmoveto{\pgfqpoint{4.902931in}{0.483000in}}%
\pgfpathlineto{\pgfqpoint{4.924194in}{0.483000in}}%
\pgfpathlineto{\pgfqpoint{4.924194in}{0.483000in}}%
\pgfpathlineto{\pgfqpoint{4.902931in}{0.483000in}}%
\pgfpathlineto{\pgfqpoint{4.902931in}{0.483000in}}%
\pgfpathclose%
\pgfusepath{fill}%
\end{pgfscope}%
\begin{pgfscope}%
\pgfpathrectangle{\pgfqpoint{3.815000in}{0.483000in}}{\pgfqpoint{2.835000in}{0.630000in}}%
\pgfusepath{clip}%
\pgfsetbuttcap%
\pgfsetmiterjoin%
\definecolor{currentfill}{rgb}{0.101961,0.101961,0.101961}%
\pgfsetfillcolor{currentfill}%
\pgfsetlinewidth{0.000000pt}%
\definecolor{currentstroke}{rgb}{0.000000,0.000000,0.000000}%
\pgfsetstrokecolor{currentstroke}%
\pgfsetstrokeopacity{0.000000}%
\pgfsetdash{}{0pt}%
\pgfpathmoveto{\pgfqpoint{4.926556in}{0.483000in}}%
\pgfpathlineto{\pgfqpoint{4.947819in}{0.483000in}}%
\pgfpathlineto{\pgfqpoint{4.947819in}{0.483000in}}%
\pgfpathlineto{\pgfqpoint{4.926556in}{0.483000in}}%
\pgfpathlineto{\pgfqpoint{4.926556in}{0.483000in}}%
\pgfpathclose%
\pgfusepath{fill}%
\end{pgfscope}%
\begin{pgfscope}%
\pgfpathrectangle{\pgfqpoint{3.815000in}{0.483000in}}{\pgfqpoint{2.835000in}{0.630000in}}%
\pgfusepath{clip}%
\pgfsetbuttcap%
\pgfsetmiterjoin%
\definecolor{currentfill}{rgb}{0.101961,0.101961,0.101961}%
\pgfsetfillcolor{currentfill}%
\pgfsetlinewidth{0.000000pt}%
\definecolor{currentstroke}{rgb}{0.000000,0.000000,0.000000}%
\pgfsetstrokecolor{currentstroke}%
\pgfsetstrokeopacity{0.000000}%
\pgfsetdash{}{0pt}%
\pgfpathmoveto{\pgfqpoint{4.950181in}{0.483000in}}%
\pgfpathlineto{\pgfqpoint{4.971444in}{0.483000in}}%
\pgfpathlineto{\pgfqpoint{4.971444in}{0.483000in}}%
\pgfpathlineto{\pgfqpoint{4.950181in}{0.483000in}}%
\pgfpathlineto{\pgfqpoint{4.950181in}{0.483000in}}%
\pgfpathclose%
\pgfusepath{fill}%
\end{pgfscope}%
\begin{pgfscope}%
\pgfpathrectangle{\pgfqpoint{3.815000in}{0.483000in}}{\pgfqpoint{2.835000in}{0.630000in}}%
\pgfusepath{clip}%
\pgfsetbuttcap%
\pgfsetmiterjoin%
\definecolor{currentfill}{rgb}{0.101961,0.101961,0.101961}%
\pgfsetfillcolor{currentfill}%
\pgfsetlinewidth{0.000000pt}%
\definecolor{currentstroke}{rgb}{0.000000,0.000000,0.000000}%
\pgfsetstrokecolor{currentstroke}%
\pgfsetstrokeopacity{0.000000}%
\pgfsetdash{}{0pt}%
\pgfpathmoveto{\pgfqpoint{4.973806in}{0.483000in}}%
\pgfpathlineto{\pgfqpoint{4.995069in}{0.483000in}}%
\pgfpathlineto{\pgfqpoint{4.995069in}{0.483000in}}%
\pgfpathlineto{\pgfqpoint{4.973806in}{0.483000in}}%
\pgfpathlineto{\pgfqpoint{4.973806in}{0.483000in}}%
\pgfpathclose%
\pgfusepath{fill}%
\end{pgfscope}%
\begin{pgfscope}%
\pgfpathrectangle{\pgfqpoint{3.815000in}{0.483000in}}{\pgfqpoint{2.835000in}{0.630000in}}%
\pgfusepath{clip}%
\pgfsetbuttcap%
\pgfsetmiterjoin%
\definecolor{currentfill}{rgb}{0.101961,0.101961,0.101961}%
\pgfsetfillcolor{currentfill}%
\pgfsetlinewidth{0.000000pt}%
\definecolor{currentstroke}{rgb}{0.000000,0.000000,0.000000}%
\pgfsetstrokecolor{currentstroke}%
\pgfsetstrokeopacity{0.000000}%
\pgfsetdash{}{0pt}%
\pgfpathmoveto{\pgfqpoint{4.997431in}{0.483000in}}%
\pgfpathlineto{\pgfqpoint{5.018694in}{0.483000in}}%
\pgfpathlineto{\pgfqpoint{5.018694in}{0.483000in}}%
\pgfpathlineto{\pgfqpoint{4.997431in}{0.483000in}}%
\pgfpathlineto{\pgfqpoint{4.997431in}{0.483000in}}%
\pgfpathclose%
\pgfusepath{fill}%
\end{pgfscope}%
\begin{pgfscope}%
\pgfpathrectangle{\pgfqpoint{3.815000in}{0.483000in}}{\pgfqpoint{2.835000in}{0.630000in}}%
\pgfusepath{clip}%
\pgfsetbuttcap%
\pgfsetmiterjoin%
\definecolor{currentfill}{rgb}{0.101961,0.101961,0.101961}%
\pgfsetfillcolor{currentfill}%
\pgfsetlinewidth{0.000000pt}%
\definecolor{currentstroke}{rgb}{0.000000,0.000000,0.000000}%
\pgfsetstrokecolor{currentstroke}%
\pgfsetstrokeopacity{0.000000}%
\pgfsetdash{}{0pt}%
\pgfpathmoveto{\pgfqpoint{5.021056in}{0.483000in}}%
\pgfpathlineto{\pgfqpoint{5.042319in}{0.483000in}}%
\pgfpathlineto{\pgfqpoint{5.042319in}{0.483000in}}%
\pgfpathlineto{\pgfqpoint{5.021056in}{0.483000in}}%
\pgfpathlineto{\pgfqpoint{5.021056in}{0.483000in}}%
\pgfpathclose%
\pgfusepath{fill}%
\end{pgfscope}%
\begin{pgfscope}%
\pgfpathrectangle{\pgfqpoint{3.815000in}{0.483000in}}{\pgfqpoint{2.835000in}{0.630000in}}%
\pgfusepath{clip}%
\pgfsetbuttcap%
\pgfsetmiterjoin%
\definecolor{currentfill}{rgb}{0.101961,0.101961,0.101961}%
\pgfsetfillcolor{currentfill}%
\pgfsetlinewidth{0.000000pt}%
\definecolor{currentstroke}{rgb}{0.000000,0.000000,0.000000}%
\pgfsetstrokecolor{currentstroke}%
\pgfsetstrokeopacity{0.000000}%
\pgfsetdash{}{0pt}%
\pgfpathmoveto{\pgfqpoint{5.044681in}{0.483000in}}%
\pgfpathlineto{\pgfqpoint{5.065944in}{0.483000in}}%
\pgfpathlineto{\pgfqpoint{5.065944in}{0.483000in}}%
\pgfpathlineto{\pgfqpoint{5.044681in}{0.483000in}}%
\pgfpathlineto{\pgfqpoint{5.044681in}{0.483000in}}%
\pgfpathclose%
\pgfusepath{fill}%
\end{pgfscope}%
\begin{pgfscope}%
\pgfpathrectangle{\pgfqpoint{3.815000in}{0.483000in}}{\pgfqpoint{2.835000in}{0.630000in}}%
\pgfusepath{clip}%
\pgfsetbuttcap%
\pgfsetmiterjoin%
\definecolor{currentfill}{rgb}{0.101961,0.101961,0.101961}%
\pgfsetfillcolor{currentfill}%
\pgfsetlinewidth{0.000000pt}%
\definecolor{currentstroke}{rgb}{0.000000,0.000000,0.000000}%
\pgfsetstrokecolor{currentstroke}%
\pgfsetstrokeopacity{0.000000}%
\pgfsetdash{}{0pt}%
\pgfpathmoveto{\pgfqpoint{5.068306in}{0.483000in}}%
\pgfpathlineto{\pgfqpoint{5.089569in}{0.483000in}}%
\pgfpathlineto{\pgfqpoint{5.089569in}{0.483000in}}%
\pgfpathlineto{\pgfqpoint{5.068306in}{0.483000in}}%
\pgfpathlineto{\pgfqpoint{5.068306in}{0.483000in}}%
\pgfpathclose%
\pgfusepath{fill}%
\end{pgfscope}%
\begin{pgfscope}%
\pgfpathrectangle{\pgfqpoint{3.815000in}{0.483000in}}{\pgfqpoint{2.835000in}{0.630000in}}%
\pgfusepath{clip}%
\pgfsetbuttcap%
\pgfsetmiterjoin%
\definecolor{currentfill}{rgb}{0.101961,0.101961,0.101961}%
\pgfsetfillcolor{currentfill}%
\pgfsetlinewidth{0.000000pt}%
\definecolor{currentstroke}{rgb}{0.000000,0.000000,0.000000}%
\pgfsetstrokecolor{currentstroke}%
\pgfsetstrokeopacity{0.000000}%
\pgfsetdash{}{0pt}%
\pgfpathmoveto{\pgfqpoint{5.091931in}{0.483000in}}%
\pgfpathlineto{\pgfqpoint{5.113194in}{0.483000in}}%
\pgfpathlineto{\pgfqpoint{5.113194in}{0.483000in}}%
\pgfpathlineto{\pgfqpoint{5.091931in}{0.483000in}}%
\pgfpathlineto{\pgfqpoint{5.091931in}{0.483000in}}%
\pgfpathclose%
\pgfusepath{fill}%
\end{pgfscope}%
\begin{pgfscope}%
\pgfpathrectangle{\pgfqpoint{3.815000in}{0.483000in}}{\pgfqpoint{2.835000in}{0.630000in}}%
\pgfusepath{clip}%
\pgfsetbuttcap%
\pgfsetmiterjoin%
\definecolor{currentfill}{rgb}{0.101961,0.101961,0.101961}%
\pgfsetfillcolor{currentfill}%
\pgfsetlinewidth{0.000000pt}%
\definecolor{currentstroke}{rgb}{0.000000,0.000000,0.000000}%
\pgfsetstrokecolor{currentstroke}%
\pgfsetstrokeopacity{0.000000}%
\pgfsetdash{}{0pt}%
\pgfpathmoveto{\pgfqpoint{5.115556in}{0.483000in}}%
\pgfpathlineto{\pgfqpoint{5.136819in}{0.483000in}}%
\pgfpathlineto{\pgfqpoint{5.136819in}{0.483000in}}%
\pgfpathlineto{\pgfqpoint{5.115556in}{0.483000in}}%
\pgfpathlineto{\pgfqpoint{5.115556in}{0.483000in}}%
\pgfpathclose%
\pgfusepath{fill}%
\end{pgfscope}%
\begin{pgfscope}%
\pgfpathrectangle{\pgfqpoint{3.815000in}{0.483000in}}{\pgfqpoint{2.835000in}{0.630000in}}%
\pgfusepath{clip}%
\pgfsetbuttcap%
\pgfsetmiterjoin%
\definecolor{currentfill}{rgb}{0.101961,0.101961,0.101961}%
\pgfsetfillcolor{currentfill}%
\pgfsetlinewidth{0.000000pt}%
\definecolor{currentstroke}{rgb}{0.000000,0.000000,0.000000}%
\pgfsetstrokecolor{currentstroke}%
\pgfsetstrokeopacity{0.000000}%
\pgfsetdash{}{0pt}%
\pgfpathmoveto{\pgfqpoint{5.139181in}{0.483000in}}%
\pgfpathlineto{\pgfqpoint{5.160444in}{0.483000in}}%
\pgfpathlineto{\pgfqpoint{5.160444in}{0.483000in}}%
\pgfpathlineto{\pgfqpoint{5.139181in}{0.483000in}}%
\pgfpathlineto{\pgfqpoint{5.139181in}{0.483000in}}%
\pgfpathclose%
\pgfusepath{fill}%
\end{pgfscope}%
\begin{pgfscope}%
\pgfpathrectangle{\pgfqpoint{3.815000in}{0.483000in}}{\pgfqpoint{2.835000in}{0.630000in}}%
\pgfusepath{clip}%
\pgfsetbuttcap%
\pgfsetmiterjoin%
\definecolor{currentfill}{rgb}{0.101961,0.101961,0.101961}%
\pgfsetfillcolor{currentfill}%
\pgfsetlinewidth{0.000000pt}%
\definecolor{currentstroke}{rgb}{0.000000,0.000000,0.000000}%
\pgfsetstrokecolor{currentstroke}%
\pgfsetstrokeopacity{0.000000}%
\pgfsetdash{}{0pt}%
\pgfpathmoveto{\pgfqpoint{5.162806in}{0.483000in}}%
\pgfpathlineto{\pgfqpoint{5.184069in}{0.483000in}}%
\pgfpathlineto{\pgfqpoint{5.184069in}{0.483000in}}%
\pgfpathlineto{\pgfqpoint{5.162806in}{0.483000in}}%
\pgfpathlineto{\pgfqpoint{5.162806in}{0.483000in}}%
\pgfpathclose%
\pgfusepath{fill}%
\end{pgfscope}%
\begin{pgfscope}%
\pgfpathrectangle{\pgfqpoint{3.815000in}{0.483000in}}{\pgfqpoint{2.835000in}{0.630000in}}%
\pgfusepath{clip}%
\pgfsetbuttcap%
\pgfsetmiterjoin%
\definecolor{currentfill}{rgb}{0.101961,0.101961,0.101961}%
\pgfsetfillcolor{currentfill}%
\pgfsetlinewidth{0.000000pt}%
\definecolor{currentstroke}{rgb}{0.000000,0.000000,0.000000}%
\pgfsetstrokecolor{currentstroke}%
\pgfsetstrokeopacity{0.000000}%
\pgfsetdash{}{0pt}%
\pgfpathmoveto{\pgfqpoint{5.186431in}{0.483000in}}%
\pgfpathlineto{\pgfqpoint{5.207694in}{0.483000in}}%
\pgfpathlineto{\pgfqpoint{5.207694in}{0.483000in}}%
\pgfpathlineto{\pgfqpoint{5.186431in}{0.483000in}}%
\pgfpathlineto{\pgfqpoint{5.186431in}{0.483000in}}%
\pgfpathclose%
\pgfusepath{fill}%
\end{pgfscope}%
\begin{pgfscope}%
\pgfpathrectangle{\pgfqpoint{3.815000in}{0.483000in}}{\pgfqpoint{2.835000in}{0.630000in}}%
\pgfusepath{clip}%
\pgfsetbuttcap%
\pgfsetmiterjoin%
\definecolor{currentfill}{rgb}{0.101961,0.101961,0.101961}%
\pgfsetfillcolor{currentfill}%
\pgfsetlinewidth{0.000000pt}%
\definecolor{currentstroke}{rgb}{0.000000,0.000000,0.000000}%
\pgfsetstrokecolor{currentstroke}%
\pgfsetstrokeopacity{0.000000}%
\pgfsetdash{}{0pt}%
\pgfpathmoveto{\pgfqpoint{5.210056in}{0.483000in}}%
\pgfpathlineto{\pgfqpoint{5.231319in}{0.483000in}}%
\pgfpathlineto{\pgfqpoint{5.231319in}{0.483000in}}%
\pgfpathlineto{\pgfqpoint{5.210056in}{0.483000in}}%
\pgfpathlineto{\pgfqpoint{5.210056in}{0.483000in}}%
\pgfpathclose%
\pgfusepath{fill}%
\end{pgfscope}%
\begin{pgfscope}%
\pgfpathrectangle{\pgfqpoint{3.815000in}{0.483000in}}{\pgfqpoint{2.835000in}{0.630000in}}%
\pgfusepath{clip}%
\pgfsetbuttcap%
\pgfsetmiterjoin%
\definecolor{currentfill}{rgb}{0.101961,0.101961,0.101961}%
\pgfsetfillcolor{currentfill}%
\pgfsetlinewidth{0.000000pt}%
\definecolor{currentstroke}{rgb}{0.000000,0.000000,0.000000}%
\pgfsetstrokecolor{currentstroke}%
\pgfsetstrokeopacity{0.000000}%
\pgfsetdash{}{0pt}%
\pgfpathmoveto{\pgfqpoint{5.233681in}{0.483000in}}%
\pgfpathlineto{\pgfqpoint{5.254944in}{0.483000in}}%
\pgfpathlineto{\pgfqpoint{5.254944in}{0.528000in}}%
\pgfpathlineto{\pgfqpoint{5.233681in}{0.528000in}}%
\pgfpathlineto{\pgfqpoint{5.233681in}{0.483000in}}%
\pgfpathclose%
\pgfusepath{fill}%
\end{pgfscope}%
\begin{pgfscope}%
\pgfpathrectangle{\pgfqpoint{3.815000in}{0.483000in}}{\pgfqpoint{2.835000in}{0.630000in}}%
\pgfusepath{clip}%
\pgfsetbuttcap%
\pgfsetmiterjoin%
\definecolor{currentfill}{rgb}{0.101961,0.101961,0.101961}%
\pgfsetfillcolor{currentfill}%
\pgfsetlinewidth{0.000000pt}%
\definecolor{currentstroke}{rgb}{0.000000,0.000000,0.000000}%
\pgfsetstrokecolor{currentstroke}%
\pgfsetstrokeopacity{0.000000}%
\pgfsetdash{}{0pt}%
\pgfpathmoveto{\pgfqpoint{5.257306in}{0.483000in}}%
\pgfpathlineto{\pgfqpoint{5.278569in}{0.483000in}}%
\pgfpathlineto{\pgfqpoint{5.278569in}{0.528000in}}%
\pgfpathlineto{\pgfqpoint{5.257306in}{0.528000in}}%
\pgfpathlineto{\pgfqpoint{5.257306in}{0.483000in}}%
\pgfpathclose%
\pgfusepath{fill}%
\end{pgfscope}%
\begin{pgfscope}%
\pgfpathrectangle{\pgfqpoint{3.815000in}{0.483000in}}{\pgfqpoint{2.835000in}{0.630000in}}%
\pgfusepath{clip}%
\pgfsetbuttcap%
\pgfsetmiterjoin%
\definecolor{currentfill}{rgb}{0.101961,0.101961,0.101961}%
\pgfsetfillcolor{currentfill}%
\pgfsetlinewidth{0.000000pt}%
\definecolor{currentstroke}{rgb}{0.000000,0.000000,0.000000}%
\pgfsetstrokecolor{currentstroke}%
\pgfsetstrokeopacity{0.000000}%
\pgfsetdash{}{0pt}%
\pgfpathmoveto{\pgfqpoint{5.280931in}{0.483000in}}%
\pgfpathlineto{\pgfqpoint{5.302194in}{0.483000in}}%
\pgfpathlineto{\pgfqpoint{5.302194in}{0.618000in}}%
\pgfpathlineto{\pgfqpoint{5.280931in}{0.618000in}}%
\pgfpathlineto{\pgfqpoint{5.280931in}{0.483000in}}%
\pgfpathclose%
\pgfusepath{fill}%
\end{pgfscope}%
\begin{pgfscope}%
\pgfpathrectangle{\pgfqpoint{3.815000in}{0.483000in}}{\pgfqpoint{2.835000in}{0.630000in}}%
\pgfusepath{clip}%
\pgfsetbuttcap%
\pgfsetmiterjoin%
\definecolor{currentfill}{rgb}{0.101961,0.101961,0.101961}%
\pgfsetfillcolor{currentfill}%
\pgfsetlinewidth{0.000000pt}%
\definecolor{currentstroke}{rgb}{0.000000,0.000000,0.000000}%
\pgfsetstrokecolor{currentstroke}%
\pgfsetstrokeopacity{0.000000}%
\pgfsetdash{}{0pt}%
\pgfpathmoveto{\pgfqpoint{5.304556in}{0.483000in}}%
\pgfpathlineto{\pgfqpoint{5.325819in}{0.483000in}}%
\pgfpathlineto{\pgfqpoint{5.325819in}{0.573000in}}%
\pgfpathlineto{\pgfqpoint{5.304556in}{0.573000in}}%
\pgfpathlineto{\pgfqpoint{5.304556in}{0.483000in}}%
\pgfpathclose%
\pgfusepath{fill}%
\end{pgfscope}%
\begin{pgfscope}%
\pgfpathrectangle{\pgfqpoint{3.815000in}{0.483000in}}{\pgfqpoint{2.835000in}{0.630000in}}%
\pgfusepath{clip}%
\pgfsetbuttcap%
\pgfsetmiterjoin%
\definecolor{currentfill}{rgb}{0.101961,0.101961,0.101961}%
\pgfsetfillcolor{currentfill}%
\pgfsetlinewidth{0.000000pt}%
\definecolor{currentstroke}{rgb}{0.000000,0.000000,0.000000}%
\pgfsetstrokecolor{currentstroke}%
\pgfsetstrokeopacity{0.000000}%
\pgfsetdash{}{0pt}%
\pgfpathmoveto{\pgfqpoint{5.328181in}{0.483000in}}%
\pgfpathlineto{\pgfqpoint{5.349444in}{0.483000in}}%
\pgfpathlineto{\pgfqpoint{5.349444in}{0.483000in}}%
\pgfpathlineto{\pgfqpoint{5.328181in}{0.483000in}}%
\pgfpathlineto{\pgfqpoint{5.328181in}{0.483000in}}%
\pgfpathclose%
\pgfusepath{fill}%
\end{pgfscope}%
\begin{pgfscope}%
\pgfpathrectangle{\pgfqpoint{3.815000in}{0.483000in}}{\pgfqpoint{2.835000in}{0.630000in}}%
\pgfusepath{clip}%
\pgfsetbuttcap%
\pgfsetmiterjoin%
\definecolor{currentfill}{rgb}{0.101961,0.101961,0.101961}%
\pgfsetfillcolor{currentfill}%
\pgfsetlinewidth{0.000000pt}%
\definecolor{currentstroke}{rgb}{0.000000,0.000000,0.000000}%
\pgfsetstrokecolor{currentstroke}%
\pgfsetstrokeopacity{0.000000}%
\pgfsetdash{}{0pt}%
\pgfpathmoveto{\pgfqpoint{5.351806in}{0.483000in}}%
\pgfpathlineto{\pgfqpoint{5.373069in}{0.483000in}}%
\pgfpathlineto{\pgfqpoint{5.373069in}{0.483000in}}%
\pgfpathlineto{\pgfqpoint{5.351806in}{0.483000in}}%
\pgfpathlineto{\pgfqpoint{5.351806in}{0.483000in}}%
\pgfpathclose%
\pgfusepath{fill}%
\end{pgfscope}%
\begin{pgfscope}%
\pgfpathrectangle{\pgfqpoint{3.815000in}{0.483000in}}{\pgfqpoint{2.835000in}{0.630000in}}%
\pgfusepath{clip}%
\pgfsetbuttcap%
\pgfsetmiterjoin%
\definecolor{currentfill}{rgb}{0.101961,0.101961,0.101961}%
\pgfsetfillcolor{currentfill}%
\pgfsetlinewidth{0.000000pt}%
\definecolor{currentstroke}{rgb}{0.000000,0.000000,0.000000}%
\pgfsetstrokecolor{currentstroke}%
\pgfsetstrokeopacity{0.000000}%
\pgfsetdash{}{0pt}%
\pgfpathmoveto{\pgfqpoint{5.375431in}{0.483000in}}%
\pgfpathlineto{\pgfqpoint{5.396694in}{0.483000in}}%
\pgfpathlineto{\pgfqpoint{5.396694in}{0.573000in}}%
\pgfpathlineto{\pgfqpoint{5.375431in}{0.573000in}}%
\pgfpathlineto{\pgfqpoint{5.375431in}{0.483000in}}%
\pgfpathclose%
\pgfusepath{fill}%
\end{pgfscope}%
\begin{pgfscope}%
\pgfpathrectangle{\pgfqpoint{3.815000in}{0.483000in}}{\pgfqpoint{2.835000in}{0.630000in}}%
\pgfusepath{clip}%
\pgfsetbuttcap%
\pgfsetmiterjoin%
\definecolor{currentfill}{rgb}{0.101961,0.101961,0.101961}%
\pgfsetfillcolor{currentfill}%
\pgfsetlinewidth{0.000000pt}%
\definecolor{currentstroke}{rgb}{0.000000,0.000000,0.000000}%
\pgfsetstrokecolor{currentstroke}%
\pgfsetstrokeopacity{0.000000}%
\pgfsetdash{}{0pt}%
\pgfpathmoveto{\pgfqpoint{5.399056in}{0.483000in}}%
\pgfpathlineto{\pgfqpoint{5.420319in}{0.483000in}}%
\pgfpathlineto{\pgfqpoint{5.420319in}{0.483000in}}%
\pgfpathlineto{\pgfqpoint{5.399056in}{0.483000in}}%
\pgfpathlineto{\pgfqpoint{5.399056in}{0.483000in}}%
\pgfpathclose%
\pgfusepath{fill}%
\end{pgfscope}%
\begin{pgfscope}%
\pgfpathrectangle{\pgfqpoint{3.815000in}{0.483000in}}{\pgfqpoint{2.835000in}{0.630000in}}%
\pgfusepath{clip}%
\pgfsetbuttcap%
\pgfsetmiterjoin%
\definecolor{currentfill}{rgb}{0.101961,0.101961,0.101961}%
\pgfsetfillcolor{currentfill}%
\pgfsetlinewidth{0.000000pt}%
\definecolor{currentstroke}{rgb}{0.000000,0.000000,0.000000}%
\pgfsetstrokecolor{currentstroke}%
\pgfsetstrokeopacity{0.000000}%
\pgfsetdash{}{0pt}%
\pgfpathmoveto{\pgfqpoint{5.422681in}{0.483000in}}%
\pgfpathlineto{\pgfqpoint{5.443944in}{0.483000in}}%
\pgfpathlineto{\pgfqpoint{5.443944in}{0.573000in}}%
\pgfpathlineto{\pgfqpoint{5.422681in}{0.573000in}}%
\pgfpathlineto{\pgfqpoint{5.422681in}{0.483000in}}%
\pgfpathclose%
\pgfusepath{fill}%
\end{pgfscope}%
\begin{pgfscope}%
\pgfpathrectangle{\pgfqpoint{3.815000in}{0.483000in}}{\pgfqpoint{2.835000in}{0.630000in}}%
\pgfusepath{clip}%
\pgfsetbuttcap%
\pgfsetmiterjoin%
\definecolor{currentfill}{rgb}{0.101961,0.101961,0.101961}%
\pgfsetfillcolor{currentfill}%
\pgfsetlinewidth{0.000000pt}%
\definecolor{currentstroke}{rgb}{0.000000,0.000000,0.000000}%
\pgfsetstrokecolor{currentstroke}%
\pgfsetstrokeopacity{0.000000}%
\pgfsetdash{}{0pt}%
\pgfpathmoveto{\pgfqpoint{5.446306in}{0.483000in}}%
\pgfpathlineto{\pgfqpoint{5.467569in}{0.483000in}}%
\pgfpathlineto{\pgfqpoint{5.467569in}{0.618000in}}%
\pgfpathlineto{\pgfqpoint{5.446306in}{0.618000in}}%
\pgfpathlineto{\pgfqpoint{5.446306in}{0.483000in}}%
\pgfpathclose%
\pgfusepath{fill}%
\end{pgfscope}%
\begin{pgfscope}%
\pgfpathrectangle{\pgfqpoint{3.815000in}{0.483000in}}{\pgfqpoint{2.835000in}{0.630000in}}%
\pgfusepath{clip}%
\pgfsetbuttcap%
\pgfsetmiterjoin%
\definecolor{currentfill}{rgb}{0.101961,0.101961,0.101961}%
\pgfsetfillcolor{currentfill}%
\pgfsetlinewidth{0.000000pt}%
\definecolor{currentstroke}{rgb}{0.000000,0.000000,0.000000}%
\pgfsetstrokecolor{currentstroke}%
\pgfsetstrokeopacity{0.000000}%
\pgfsetdash{}{0pt}%
\pgfpathmoveto{\pgfqpoint{5.469931in}{0.483000in}}%
\pgfpathlineto{\pgfqpoint{5.491194in}{0.483000in}}%
\pgfpathlineto{\pgfqpoint{5.491194in}{0.573000in}}%
\pgfpathlineto{\pgfqpoint{5.469931in}{0.573000in}}%
\pgfpathlineto{\pgfqpoint{5.469931in}{0.483000in}}%
\pgfpathclose%
\pgfusepath{fill}%
\end{pgfscope}%
\begin{pgfscope}%
\pgfpathrectangle{\pgfqpoint{3.815000in}{0.483000in}}{\pgfqpoint{2.835000in}{0.630000in}}%
\pgfusepath{clip}%
\pgfsetbuttcap%
\pgfsetmiterjoin%
\definecolor{currentfill}{rgb}{0.101961,0.101961,0.101961}%
\pgfsetfillcolor{currentfill}%
\pgfsetlinewidth{0.000000pt}%
\definecolor{currentstroke}{rgb}{0.000000,0.000000,0.000000}%
\pgfsetstrokecolor{currentstroke}%
\pgfsetstrokeopacity{0.000000}%
\pgfsetdash{}{0pt}%
\pgfpathmoveto{\pgfqpoint{5.493556in}{0.483000in}}%
\pgfpathlineto{\pgfqpoint{5.514819in}{0.483000in}}%
\pgfpathlineto{\pgfqpoint{5.514819in}{0.618000in}}%
\pgfpathlineto{\pgfqpoint{5.493556in}{0.618000in}}%
\pgfpathlineto{\pgfqpoint{5.493556in}{0.483000in}}%
\pgfpathclose%
\pgfusepath{fill}%
\end{pgfscope}%
\begin{pgfscope}%
\pgfpathrectangle{\pgfqpoint{3.815000in}{0.483000in}}{\pgfqpoint{2.835000in}{0.630000in}}%
\pgfusepath{clip}%
\pgfsetbuttcap%
\pgfsetmiterjoin%
\definecolor{currentfill}{rgb}{0.101961,0.101961,0.101961}%
\pgfsetfillcolor{currentfill}%
\pgfsetlinewidth{0.000000pt}%
\definecolor{currentstroke}{rgb}{0.000000,0.000000,0.000000}%
\pgfsetstrokecolor{currentstroke}%
\pgfsetstrokeopacity{0.000000}%
\pgfsetdash{}{0pt}%
\pgfpathmoveto{\pgfqpoint{5.517181in}{0.483000in}}%
\pgfpathlineto{\pgfqpoint{5.538444in}{0.483000in}}%
\pgfpathlineto{\pgfqpoint{5.538444in}{0.708000in}}%
\pgfpathlineto{\pgfqpoint{5.517181in}{0.708000in}}%
\pgfpathlineto{\pgfqpoint{5.517181in}{0.483000in}}%
\pgfpathclose%
\pgfusepath{fill}%
\end{pgfscope}%
\begin{pgfscope}%
\pgfpathrectangle{\pgfqpoint{3.815000in}{0.483000in}}{\pgfqpoint{2.835000in}{0.630000in}}%
\pgfusepath{clip}%
\pgfsetbuttcap%
\pgfsetmiterjoin%
\definecolor{currentfill}{rgb}{0.101961,0.101961,0.101961}%
\pgfsetfillcolor{currentfill}%
\pgfsetlinewidth{0.000000pt}%
\definecolor{currentstroke}{rgb}{0.000000,0.000000,0.000000}%
\pgfsetstrokecolor{currentstroke}%
\pgfsetstrokeopacity{0.000000}%
\pgfsetdash{}{0pt}%
\pgfpathmoveto{\pgfqpoint{5.540806in}{0.483000in}}%
\pgfpathlineto{\pgfqpoint{5.562069in}{0.483000in}}%
\pgfpathlineto{\pgfqpoint{5.562069in}{0.528000in}}%
\pgfpathlineto{\pgfqpoint{5.540806in}{0.528000in}}%
\pgfpathlineto{\pgfqpoint{5.540806in}{0.483000in}}%
\pgfpathclose%
\pgfusepath{fill}%
\end{pgfscope}%
\begin{pgfscope}%
\pgfpathrectangle{\pgfqpoint{3.815000in}{0.483000in}}{\pgfqpoint{2.835000in}{0.630000in}}%
\pgfusepath{clip}%
\pgfsetbuttcap%
\pgfsetmiterjoin%
\definecolor{currentfill}{rgb}{0.101961,0.101961,0.101961}%
\pgfsetfillcolor{currentfill}%
\pgfsetlinewidth{0.000000pt}%
\definecolor{currentstroke}{rgb}{0.000000,0.000000,0.000000}%
\pgfsetstrokecolor{currentstroke}%
\pgfsetstrokeopacity{0.000000}%
\pgfsetdash{}{0pt}%
\pgfpathmoveto{\pgfqpoint{5.564431in}{0.483000in}}%
\pgfpathlineto{\pgfqpoint{5.585694in}{0.483000in}}%
\pgfpathlineto{\pgfqpoint{5.585694in}{0.483000in}}%
\pgfpathlineto{\pgfqpoint{5.564431in}{0.483000in}}%
\pgfpathlineto{\pgfqpoint{5.564431in}{0.483000in}}%
\pgfpathclose%
\pgfusepath{fill}%
\end{pgfscope}%
\begin{pgfscope}%
\pgfpathrectangle{\pgfqpoint{3.815000in}{0.483000in}}{\pgfqpoint{2.835000in}{0.630000in}}%
\pgfusepath{clip}%
\pgfsetbuttcap%
\pgfsetmiterjoin%
\definecolor{currentfill}{rgb}{0.101961,0.101961,0.101961}%
\pgfsetfillcolor{currentfill}%
\pgfsetlinewidth{0.000000pt}%
\definecolor{currentstroke}{rgb}{0.000000,0.000000,0.000000}%
\pgfsetstrokecolor{currentstroke}%
\pgfsetstrokeopacity{0.000000}%
\pgfsetdash{}{0pt}%
\pgfpathmoveto{\pgfqpoint{5.588056in}{0.483000in}}%
\pgfpathlineto{\pgfqpoint{5.609319in}{0.483000in}}%
\pgfpathlineto{\pgfqpoint{5.609319in}{0.483000in}}%
\pgfpathlineto{\pgfqpoint{5.588056in}{0.483000in}}%
\pgfpathlineto{\pgfqpoint{5.588056in}{0.483000in}}%
\pgfpathclose%
\pgfusepath{fill}%
\end{pgfscope}%
\begin{pgfscope}%
\pgfpathrectangle{\pgfqpoint{3.815000in}{0.483000in}}{\pgfqpoint{2.835000in}{0.630000in}}%
\pgfusepath{clip}%
\pgfsetbuttcap%
\pgfsetmiterjoin%
\definecolor{currentfill}{rgb}{0.101961,0.101961,0.101961}%
\pgfsetfillcolor{currentfill}%
\pgfsetlinewidth{0.000000pt}%
\definecolor{currentstroke}{rgb}{0.000000,0.000000,0.000000}%
\pgfsetstrokecolor{currentstroke}%
\pgfsetstrokeopacity{0.000000}%
\pgfsetdash{}{0pt}%
\pgfpathmoveto{\pgfqpoint{5.611681in}{0.483000in}}%
\pgfpathlineto{\pgfqpoint{5.632944in}{0.483000in}}%
\pgfpathlineto{\pgfqpoint{5.632944in}{0.483000in}}%
\pgfpathlineto{\pgfqpoint{5.611681in}{0.483000in}}%
\pgfpathlineto{\pgfqpoint{5.611681in}{0.483000in}}%
\pgfpathclose%
\pgfusepath{fill}%
\end{pgfscope}%
\begin{pgfscope}%
\pgfpathrectangle{\pgfqpoint{3.815000in}{0.483000in}}{\pgfqpoint{2.835000in}{0.630000in}}%
\pgfusepath{clip}%
\pgfsetbuttcap%
\pgfsetmiterjoin%
\definecolor{currentfill}{rgb}{0.101961,0.101961,0.101961}%
\pgfsetfillcolor{currentfill}%
\pgfsetlinewidth{0.000000pt}%
\definecolor{currentstroke}{rgb}{0.000000,0.000000,0.000000}%
\pgfsetstrokecolor{currentstroke}%
\pgfsetstrokeopacity{0.000000}%
\pgfsetdash{}{0pt}%
\pgfpathmoveto{\pgfqpoint{5.635306in}{0.483000in}}%
\pgfpathlineto{\pgfqpoint{5.656569in}{0.483000in}}%
\pgfpathlineto{\pgfqpoint{5.656569in}{0.573000in}}%
\pgfpathlineto{\pgfqpoint{5.635306in}{0.573000in}}%
\pgfpathlineto{\pgfqpoint{5.635306in}{0.483000in}}%
\pgfpathclose%
\pgfusepath{fill}%
\end{pgfscope}%
\begin{pgfscope}%
\pgfpathrectangle{\pgfqpoint{3.815000in}{0.483000in}}{\pgfqpoint{2.835000in}{0.630000in}}%
\pgfusepath{clip}%
\pgfsetbuttcap%
\pgfsetmiterjoin%
\definecolor{currentfill}{rgb}{0.101961,0.101961,0.101961}%
\pgfsetfillcolor{currentfill}%
\pgfsetlinewidth{0.000000pt}%
\definecolor{currentstroke}{rgb}{0.000000,0.000000,0.000000}%
\pgfsetstrokecolor{currentstroke}%
\pgfsetstrokeopacity{0.000000}%
\pgfsetdash{}{0pt}%
\pgfpathmoveto{\pgfqpoint{5.658931in}{0.483000in}}%
\pgfpathlineto{\pgfqpoint{5.680194in}{0.483000in}}%
\pgfpathlineto{\pgfqpoint{5.680194in}{0.483000in}}%
\pgfpathlineto{\pgfqpoint{5.658931in}{0.483000in}}%
\pgfpathlineto{\pgfqpoint{5.658931in}{0.483000in}}%
\pgfpathclose%
\pgfusepath{fill}%
\end{pgfscope}%
\begin{pgfscope}%
\pgfpathrectangle{\pgfqpoint{3.815000in}{0.483000in}}{\pgfqpoint{2.835000in}{0.630000in}}%
\pgfusepath{clip}%
\pgfsetbuttcap%
\pgfsetmiterjoin%
\definecolor{currentfill}{rgb}{0.101961,0.101961,0.101961}%
\pgfsetfillcolor{currentfill}%
\pgfsetlinewidth{0.000000pt}%
\definecolor{currentstroke}{rgb}{0.000000,0.000000,0.000000}%
\pgfsetstrokecolor{currentstroke}%
\pgfsetstrokeopacity{0.000000}%
\pgfsetdash{}{0pt}%
\pgfpathmoveto{\pgfqpoint{5.682556in}{0.483000in}}%
\pgfpathlineto{\pgfqpoint{5.703819in}{0.483000in}}%
\pgfpathlineto{\pgfqpoint{5.703819in}{0.483000in}}%
\pgfpathlineto{\pgfqpoint{5.682556in}{0.483000in}}%
\pgfpathlineto{\pgfqpoint{5.682556in}{0.483000in}}%
\pgfpathclose%
\pgfusepath{fill}%
\end{pgfscope}%
\begin{pgfscope}%
\pgfpathrectangle{\pgfqpoint{3.815000in}{0.483000in}}{\pgfqpoint{2.835000in}{0.630000in}}%
\pgfusepath{clip}%
\pgfsetbuttcap%
\pgfsetmiterjoin%
\definecolor{currentfill}{rgb}{0.101961,0.101961,0.101961}%
\pgfsetfillcolor{currentfill}%
\pgfsetlinewidth{0.000000pt}%
\definecolor{currentstroke}{rgb}{0.000000,0.000000,0.000000}%
\pgfsetstrokecolor{currentstroke}%
\pgfsetstrokeopacity{0.000000}%
\pgfsetdash{}{0pt}%
\pgfpathmoveto{\pgfqpoint{5.706181in}{0.483000in}}%
\pgfpathlineto{\pgfqpoint{5.727444in}{0.483000in}}%
\pgfpathlineto{\pgfqpoint{5.727444in}{0.573000in}}%
\pgfpathlineto{\pgfqpoint{5.706181in}{0.573000in}}%
\pgfpathlineto{\pgfqpoint{5.706181in}{0.483000in}}%
\pgfpathclose%
\pgfusepath{fill}%
\end{pgfscope}%
\begin{pgfscope}%
\pgfpathrectangle{\pgfqpoint{3.815000in}{0.483000in}}{\pgfqpoint{2.835000in}{0.630000in}}%
\pgfusepath{clip}%
\pgfsetbuttcap%
\pgfsetmiterjoin%
\definecolor{currentfill}{rgb}{0.101961,0.101961,0.101961}%
\pgfsetfillcolor{currentfill}%
\pgfsetlinewidth{0.000000pt}%
\definecolor{currentstroke}{rgb}{0.000000,0.000000,0.000000}%
\pgfsetstrokecolor{currentstroke}%
\pgfsetstrokeopacity{0.000000}%
\pgfsetdash{}{0pt}%
\pgfpathmoveto{\pgfqpoint{5.729806in}{0.483000in}}%
\pgfpathlineto{\pgfqpoint{5.751069in}{0.483000in}}%
\pgfpathlineto{\pgfqpoint{5.751069in}{0.618000in}}%
\pgfpathlineto{\pgfqpoint{5.729806in}{0.618000in}}%
\pgfpathlineto{\pgfqpoint{5.729806in}{0.483000in}}%
\pgfpathclose%
\pgfusepath{fill}%
\end{pgfscope}%
\begin{pgfscope}%
\pgfpathrectangle{\pgfqpoint{3.815000in}{0.483000in}}{\pgfqpoint{2.835000in}{0.630000in}}%
\pgfusepath{clip}%
\pgfsetbuttcap%
\pgfsetmiterjoin%
\definecolor{currentfill}{rgb}{0.101961,0.101961,0.101961}%
\pgfsetfillcolor{currentfill}%
\pgfsetlinewidth{0.000000pt}%
\definecolor{currentstroke}{rgb}{0.000000,0.000000,0.000000}%
\pgfsetstrokecolor{currentstroke}%
\pgfsetstrokeopacity{0.000000}%
\pgfsetdash{}{0pt}%
\pgfpathmoveto{\pgfqpoint{5.753431in}{0.483000in}}%
\pgfpathlineto{\pgfqpoint{5.774694in}{0.483000in}}%
\pgfpathlineto{\pgfqpoint{5.774694in}{0.528000in}}%
\pgfpathlineto{\pgfqpoint{5.753431in}{0.528000in}}%
\pgfpathlineto{\pgfqpoint{5.753431in}{0.483000in}}%
\pgfpathclose%
\pgfusepath{fill}%
\end{pgfscope}%
\begin{pgfscope}%
\pgfpathrectangle{\pgfqpoint{3.815000in}{0.483000in}}{\pgfqpoint{2.835000in}{0.630000in}}%
\pgfusepath{clip}%
\pgfsetbuttcap%
\pgfsetmiterjoin%
\definecolor{currentfill}{rgb}{0.101961,0.101961,0.101961}%
\pgfsetfillcolor{currentfill}%
\pgfsetlinewidth{0.000000pt}%
\definecolor{currentstroke}{rgb}{0.000000,0.000000,0.000000}%
\pgfsetstrokecolor{currentstroke}%
\pgfsetstrokeopacity{0.000000}%
\pgfsetdash{}{0pt}%
\pgfpathmoveto{\pgfqpoint{5.777056in}{0.483000in}}%
\pgfpathlineto{\pgfqpoint{5.798319in}{0.483000in}}%
\pgfpathlineto{\pgfqpoint{5.798319in}{0.528000in}}%
\pgfpathlineto{\pgfqpoint{5.777056in}{0.528000in}}%
\pgfpathlineto{\pgfqpoint{5.777056in}{0.483000in}}%
\pgfpathclose%
\pgfusepath{fill}%
\end{pgfscope}%
\begin{pgfscope}%
\pgfpathrectangle{\pgfqpoint{3.815000in}{0.483000in}}{\pgfqpoint{2.835000in}{0.630000in}}%
\pgfusepath{clip}%
\pgfsetbuttcap%
\pgfsetmiterjoin%
\definecolor{currentfill}{rgb}{0.101961,0.101961,0.101961}%
\pgfsetfillcolor{currentfill}%
\pgfsetlinewidth{0.000000pt}%
\definecolor{currentstroke}{rgb}{0.000000,0.000000,0.000000}%
\pgfsetstrokecolor{currentstroke}%
\pgfsetstrokeopacity{0.000000}%
\pgfsetdash{}{0pt}%
\pgfpathmoveto{\pgfqpoint{5.800681in}{0.483000in}}%
\pgfpathlineto{\pgfqpoint{5.821944in}{0.483000in}}%
\pgfpathlineto{\pgfqpoint{5.821944in}{0.483000in}}%
\pgfpathlineto{\pgfqpoint{5.800681in}{0.483000in}}%
\pgfpathlineto{\pgfqpoint{5.800681in}{0.483000in}}%
\pgfpathclose%
\pgfusepath{fill}%
\end{pgfscope}%
\begin{pgfscope}%
\pgfpathrectangle{\pgfqpoint{3.815000in}{0.483000in}}{\pgfqpoint{2.835000in}{0.630000in}}%
\pgfusepath{clip}%
\pgfsetbuttcap%
\pgfsetmiterjoin%
\definecolor{currentfill}{rgb}{0.101961,0.101961,0.101961}%
\pgfsetfillcolor{currentfill}%
\pgfsetlinewidth{0.000000pt}%
\definecolor{currentstroke}{rgb}{0.000000,0.000000,0.000000}%
\pgfsetstrokecolor{currentstroke}%
\pgfsetstrokeopacity{0.000000}%
\pgfsetdash{}{0pt}%
\pgfpathmoveto{\pgfqpoint{5.824306in}{0.483000in}}%
\pgfpathlineto{\pgfqpoint{5.845569in}{0.483000in}}%
\pgfpathlineto{\pgfqpoint{5.845569in}{0.618000in}}%
\pgfpathlineto{\pgfqpoint{5.824306in}{0.618000in}}%
\pgfpathlineto{\pgfqpoint{5.824306in}{0.483000in}}%
\pgfpathclose%
\pgfusepath{fill}%
\end{pgfscope}%
\begin{pgfscope}%
\pgfpathrectangle{\pgfqpoint{3.815000in}{0.483000in}}{\pgfqpoint{2.835000in}{0.630000in}}%
\pgfusepath{clip}%
\pgfsetbuttcap%
\pgfsetmiterjoin%
\definecolor{currentfill}{rgb}{0.101961,0.101961,0.101961}%
\pgfsetfillcolor{currentfill}%
\pgfsetlinewidth{0.000000pt}%
\definecolor{currentstroke}{rgb}{0.000000,0.000000,0.000000}%
\pgfsetstrokecolor{currentstroke}%
\pgfsetstrokeopacity{0.000000}%
\pgfsetdash{}{0pt}%
\pgfpathmoveto{\pgfqpoint{5.847931in}{0.483000in}}%
\pgfpathlineto{\pgfqpoint{5.869194in}{0.483000in}}%
\pgfpathlineto{\pgfqpoint{5.869194in}{0.483000in}}%
\pgfpathlineto{\pgfqpoint{5.847931in}{0.483000in}}%
\pgfpathlineto{\pgfqpoint{5.847931in}{0.483000in}}%
\pgfpathclose%
\pgfusepath{fill}%
\end{pgfscope}%
\begin{pgfscope}%
\pgfpathrectangle{\pgfqpoint{3.815000in}{0.483000in}}{\pgfqpoint{2.835000in}{0.630000in}}%
\pgfusepath{clip}%
\pgfsetbuttcap%
\pgfsetmiterjoin%
\definecolor{currentfill}{rgb}{0.101961,0.101961,0.101961}%
\pgfsetfillcolor{currentfill}%
\pgfsetlinewidth{0.000000pt}%
\definecolor{currentstroke}{rgb}{0.000000,0.000000,0.000000}%
\pgfsetstrokecolor{currentstroke}%
\pgfsetstrokeopacity{0.000000}%
\pgfsetdash{}{0pt}%
\pgfpathmoveto{\pgfqpoint{5.871556in}{0.483000in}}%
\pgfpathlineto{\pgfqpoint{5.892819in}{0.483000in}}%
\pgfpathlineto{\pgfqpoint{5.892819in}{0.483000in}}%
\pgfpathlineto{\pgfqpoint{5.871556in}{0.483000in}}%
\pgfpathlineto{\pgfqpoint{5.871556in}{0.483000in}}%
\pgfpathclose%
\pgfusepath{fill}%
\end{pgfscope}%
\begin{pgfscope}%
\pgfpathrectangle{\pgfqpoint{3.815000in}{0.483000in}}{\pgfqpoint{2.835000in}{0.630000in}}%
\pgfusepath{clip}%
\pgfsetbuttcap%
\pgfsetmiterjoin%
\definecolor{currentfill}{rgb}{0.101961,0.101961,0.101961}%
\pgfsetfillcolor{currentfill}%
\pgfsetlinewidth{0.000000pt}%
\definecolor{currentstroke}{rgb}{0.000000,0.000000,0.000000}%
\pgfsetstrokecolor{currentstroke}%
\pgfsetstrokeopacity{0.000000}%
\pgfsetdash{}{0pt}%
\pgfpathmoveto{\pgfqpoint{5.895181in}{0.483000in}}%
\pgfpathlineto{\pgfqpoint{5.916444in}{0.483000in}}%
\pgfpathlineto{\pgfqpoint{5.916444in}{0.843000in}}%
\pgfpathlineto{\pgfqpoint{5.895181in}{0.843000in}}%
\pgfpathlineto{\pgfqpoint{5.895181in}{0.483000in}}%
\pgfpathclose%
\pgfusepath{fill}%
\end{pgfscope}%
\begin{pgfscope}%
\pgfpathrectangle{\pgfqpoint{3.815000in}{0.483000in}}{\pgfqpoint{2.835000in}{0.630000in}}%
\pgfusepath{clip}%
\pgfsetbuttcap%
\pgfsetmiterjoin%
\definecolor{currentfill}{rgb}{0.101961,0.101961,0.101961}%
\pgfsetfillcolor{currentfill}%
\pgfsetlinewidth{0.000000pt}%
\definecolor{currentstroke}{rgb}{0.000000,0.000000,0.000000}%
\pgfsetstrokecolor{currentstroke}%
\pgfsetstrokeopacity{0.000000}%
\pgfsetdash{}{0pt}%
\pgfpathmoveto{\pgfqpoint{5.918806in}{0.483000in}}%
\pgfpathlineto{\pgfqpoint{5.940069in}{0.483000in}}%
\pgfpathlineto{\pgfqpoint{5.940069in}{0.663000in}}%
\pgfpathlineto{\pgfqpoint{5.918806in}{0.663000in}}%
\pgfpathlineto{\pgfqpoint{5.918806in}{0.483000in}}%
\pgfpathclose%
\pgfusepath{fill}%
\end{pgfscope}%
\begin{pgfscope}%
\pgfpathrectangle{\pgfqpoint{3.815000in}{0.483000in}}{\pgfqpoint{2.835000in}{0.630000in}}%
\pgfusepath{clip}%
\pgfsetbuttcap%
\pgfsetmiterjoin%
\definecolor{currentfill}{rgb}{0.101961,0.101961,0.101961}%
\pgfsetfillcolor{currentfill}%
\pgfsetlinewidth{0.000000pt}%
\definecolor{currentstroke}{rgb}{0.000000,0.000000,0.000000}%
\pgfsetstrokecolor{currentstroke}%
\pgfsetstrokeopacity{0.000000}%
\pgfsetdash{}{0pt}%
\pgfpathmoveto{\pgfqpoint{5.942431in}{0.483000in}}%
\pgfpathlineto{\pgfqpoint{5.963694in}{0.483000in}}%
\pgfpathlineto{\pgfqpoint{5.963694in}{0.618000in}}%
\pgfpathlineto{\pgfqpoint{5.942431in}{0.618000in}}%
\pgfpathlineto{\pgfqpoint{5.942431in}{0.483000in}}%
\pgfpathclose%
\pgfusepath{fill}%
\end{pgfscope}%
\begin{pgfscope}%
\pgfpathrectangle{\pgfqpoint{3.815000in}{0.483000in}}{\pgfqpoint{2.835000in}{0.630000in}}%
\pgfusepath{clip}%
\pgfsetbuttcap%
\pgfsetmiterjoin%
\definecolor{currentfill}{rgb}{0.101961,0.101961,0.101961}%
\pgfsetfillcolor{currentfill}%
\pgfsetlinewidth{0.000000pt}%
\definecolor{currentstroke}{rgb}{0.000000,0.000000,0.000000}%
\pgfsetstrokecolor{currentstroke}%
\pgfsetstrokeopacity{0.000000}%
\pgfsetdash{}{0pt}%
\pgfpathmoveto{\pgfqpoint{5.966056in}{0.483000in}}%
\pgfpathlineto{\pgfqpoint{5.987319in}{0.483000in}}%
\pgfpathlineto{\pgfqpoint{5.987319in}{0.663000in}}%
\pgfpathlineto{\pgfqpoint{5.966056in}{0.663000in}}%
\pgfpathlineto{\pgfqpoint{5.966056in}{0.483000in}}%
\pgfpathclose%
\pgfusepath{fill}%
\end{pgfscope}%
\begin{pgfscope}%
\pgfpathrectangle{\pgfqpoint{3.815000in}{0.483000in}}{\pgfqpoint{2.835000in}{0.630000in}}%
\pgfusepath{clip}%
\pgfsetbuttcap%
\pgfsetmiterjoin%
\definecolor{currentfill}{rgb}{0.101961,0.101961,0.101961}%
\pgfsetfillcolor{currentfill}%
\pgfsetlinewidth{0.000000pt}%
\definecolor{currentstroke}{rgb}{0.000000,0.000000,0.000000}%
\pgfsetstrokecolor{currentstroke}%
\pgfsetstrokeopacity{0.000000}%
\pgfsetdash{}{0pt}%
\pgfpathmoveto{\pgfqpoint{5.989681in}{0.483000in}}%
\pgfpathlineto{\pgfqpoint{6.010944in}{0.483000in}}%
\pgfpathlineto{\pgfqpoint{6.010944in}{0.483000in}}%
\pgfpathlineto{\pgfqpoint{5.989681in}{0.483000in}}%
\pgfpathlineto{\pgfqpoint{5.989681in}{0.483000in}}%
\pgfpathclose%
\pgfusepath{fill}%
\end{pgfscope}%
\begin{pgfscope}%
\pgfpathrectangle{\pgfqpoint{3.815000in}{0.483000in}}{\pgfqpoint{2.835000in}{0.630000in}}%
\pgfusepath{clip}%
\pgfsetbuttcap%
\pgfsetmiterjoin%
\definecolor{currentfill}{rgb}{0.101961,0.101961,0.101961}%
\pgfsetfillcolor{currentfill}%
\pgfsetlinewidth{0.000000pt}%
\definecolor{currentstroke}{rgb}{0.000000,0.000000,0.000000}%
\pgfsetstrokecolor{currentstroke}%
\pgfsetstrokeopacity{0.000000}%
\pgfsetdash{}{0pt}%
\pgfpathmoveto{\pgfqpoint{6.013306in}{0.483000in}}%
\pgfpathlineto{\pgfqpoint{6.034569in}{0.483000in}}%
\pgfpathlineto{\pgfqpoint{6.034569in}{0.933000in}}%
\pgfpathlineto{\pgfqpoint{6.013306in}{0.933000in}}%
\pgfpathlineto{\pgfqpoint{6.013306in}{0.483000in}}%
\pgfpathclose%
\pgfusepath{fill}%
\end{pgfscope}%
\begin{pgfscope}%
\pgfpathrectangle{\pgfqpoint{3.815000in}{0.483000in}}{\pgfqpoint{2.835000in}{0.630000in}}%
\pgfusepath{clip}%
\pgfsetbuttcap%
\pgfsetmiterjoin%
\definecolor{currentfill}{rgb}{0.101961,0.101961,0.101961}%
\pgfsetfillcolor{currentfill}%
\pgfsetlinewidth{0.000000pt}%
\definecolor{currentstroke}{rgb}{0.000000,0.000000,0.000000}%
\pgfsetstrokecolor{currentstroke}%
\pgfsetstrokeopacity{0.000000}%
\pgfsetdash{}{0pt}%
\pgfpathmoveto{\pgfqpoint{6.036931in}{0.483000in}}%
\pgfpathlineto{\pgfqpoint{6.058194in}{0.483000in}}%
\pgfpathlineto{\pgfqpoint{6.058194in}{0.528000in}}%
\pgfpathlineto{\pgfqpoint{6.036931in}{0.528000in}}%
\pgfpathlineto{\pgfqpoint{6.036931in}{0.483000in}}%
\pgfpathclose%
\pgfusepath{fill}%
\end{pgfscope}%
\begin{pgfscope}%
\pgfpathrectangle{\pgfqpoint{3.815000in}{0.483000in}}{\pgfqpoint{2.835000in}{0.630000in}}%
\pgfusepath{clip}%
\pgfsetbuttcap%
\pgfsetmiterjoin%
\definecolor{currentfill}{rgb}{0.101961,0.101961,0.101961}%
\pgfsetfillcolor{currentfill}%
\pgfsetlinewidth{0.000000pt}%
\definecolor{currentstroke}{rgb}{0.000000,0.000000,0.000000}%
\pgfsetstrokecolor{currentstroke}%
\pgfsetstrokeopacity{0.000000}%
\pgfsetdash{}{0pt}%
\pgfpathmoveto{\pgfqpoint{6.060556in}{0.483000in}}%
\pgfpathlineto{\pgfqpoint{6.081819in}{0.483000in}}%
\pgfpathlineto{\pgfqpoint{6.081819in}{0.483000in}}%
\pgfpathlineto{\pgfqpoint{6.060556in}{0.483000in}}%
\pgfpathlineto{\pgfqpoint{6.060556in}{0.483000in}}%
\pgfpathclose%
\pgfusepath{fill}%
\end{pgfscope}%
\begin{pgfscope}%
\pgfpathrectangle{\pgfqpoint{3.815000in}{0.483000in}}{\pgfqpoint{2.835000in}{0.630000in}}%
\pgfusepath{clip}%
\pgfsetbuttcap%
\pgfsetmiterjoin%
\definecolor{currentfill}{rgb}{0.101961,0.101961,0.101961}%
\pgfsetfillcolor{currentfill}%
\pgfsetlinewidth{0.000000pt}%
\definecolor{currentstroke}{rgb}{0.000000,0.000000,0.000000}%
\pgfsetstrokecolor{currentstroke}%
\pgfsetstrokeopacity{0.000000}%
\pgfsetdash{}{0pt}%
\pgfpathmoveto{\pgfqpoint{6.084181in}{0.483000in}}%
\pgfpathlineto{\pgfqpoint{6.105444in}{0.483000in}}%
\pgfpathlineto{\pgfqpoint{6.105444in}{0.483000in}}%
\pgfpathlineto{\pgfqpoint{6.084181in}{0.483000in}}%
\pgfpathlineto{\pgfqpoint{6.084181in}{0.483000in}}%
\pgfpathclose%
\pgfusepath{fill}%
\end{pgfscope}%
\begin{pgfscope}%
\pgfpathrectangle{\pgfqpoint{3.815000in}{0.483000in}}{\pgfqpoint{2.835000in}{0.630000in}}%
\pgfusepath{clip}%
\pgfsetbuttcap%
\pgfsetmiterjoin%
\definecolor{currentfill}{rgb}{0.101961,0.101961,0.101961}%
\pgfsetfillcolor{currentfill}%
\pgfsetlinewidth{0.000000pt}%
\definecolor{currentstroke}{rgb}{0.000000,0.000000,0.000000}%
\pgfsetstrokecolor{currentstroke}%
\pgfsetstrokeopacity{0.000000}%
\pgfsetdash{}{0pt}%
\pgfpathmoveto{\pgfqpoint{6.107806in}{0.483000in}}%
\pgfpathlineto{\pgfqpoint{6.129069in}{0.483000in}}%
\pgfpathlineto{\pgfqpoint{6.129069in}{0.798000in}}%
\pgfpathlineto{\pgfqpoint{6.107806in}{0.798000in}}%
\pgfpathlineto{\pgfqpoint{6.107806in}{0.483000in}}%
\pgfpathclose%
\pgfusepath{fill}%
\end{pgfscope}%
\begin{pgfscope}%
\pgfpathrectangle{\pgfqpoint{3.815000in}{0.483000in}}{\pgfqpoint{2.835000in}{0.630000in}}%
\pgfusepath{clip}%
\pgfsetbuttcap%
\pgfsetmiterjoin%
\definecolor{currentfill}{rgb}{0.101961,0.101961,0.101961}%
\pgfsetfillcolor{currentfill}%
\pgfsetlinewidth{0.000000pt}%
\definecolor{currentstroke}{rgb}{0.000000,0.000000,0.000000}%
\pgfsetstrokecolor{currentstroke}%
\pgfsetstrokeopacity{0.000000}%
\pgfsetdash{}{0pt}%
\pgfpathmoveto{\pgfqpoint{6.131431in}{0.483000in}}%
\pgfpathlineto{\pgfqpoint{6.152694in}{0.483000in}}%
\pgfpathlineto{\pgfqpoint{6.152694in}{0.483000in}}%
\pgfpathlineto{\pgfqpoint{6.131431in}{0.483000in}}%
\pgfpathlineto{\pgfqpoint{6.131431in}{0.483000in}}%
\pgfpathclose%
\pgfusepath{fill}%
\end{pgfscope}%
\begin{pgfscope}%
\pgfpathrectangle{\pgfqpoint{3.815000in}{0.483000in}}{\pgfqpoint{2.835000in}{0.630000in}}%
\pgfusepath{clip}%
\pgfsetbuttcap%
\pgfsetmiterjoin%
\definecolor{currentfill}{rgb}{0.101961,0.101961,0.101961}%
\pgfsetfillcolor{currentfill}%
\pgfsetlinewidth{0.000000pt}%
\definecolor{currentstroke}{rgb}{0.000000,0.000000,0.000000}%
\pgfsetstrokecolor{currentstroke}%
\pgfsetstrokeopacity{0.000000}%
\pgfsetdash{}{0pt}%
\pgfpathmoveto{\pgfqpoint{6.155056in}{0.483000in}}%
\pgfpathlineto{\pgfqpoint{6.176319in}{0.483000in}}%
\pgfpathlineto{\pgfqpoint{6.176319in}{0.528000in}}%
\pgfpathlineto{\pgfqpoint{6.155056in}{0.528000in}}%
\pgfpathlineto{\pgfqpoint{6.155056in}{0.483000in}}%
\pgfpathclose%
\pgfusepath{fill}%
\end{pgfscope}%
\begin{pgfscope}%
\pgfpathrectangle{\pgfqpoint{3.815000in}{0.483000in}}{\pgfqpoint{2.835000in}{0.630000in}}%
\pgfusepath{clip}%
\pgfsetbuttcap%
\pgfsetmiterjoin%
\definecolor{currentfill}{rgb}{0.101961,0.101961,0.101961}%
\pgfsetfillcolor{currentfill}%
\pgfsetlinewidth{0.000000pt}%
\definecolor{currentstroke}{rgb}{0.000000,0.000000,0.000000}%
\pgfsetstrokecolor{currentstroke}%
\pgfsetstrokeopacity{0.000000}%
\pgfsetdash{}{0pt}%
\pgfpathmoveto{\pgfqpoint{6.178681in}{0.483000in}}%
\pgfpathlineto{\pgfqpoint{6.199944in}{0.483000in}}%
\pgfpathlineto{\pgfqpoint{6.199944in}{0.528000in}}%
\pgfpathlineto{\pgfqpoint{6.178681in}{0.528000in}}%
\pgfpathlineto{\pgfqpoint{6.178681in}{0.483000in}}%
\pgfpathclose%
\pgfusepath{fill}%
\end{pgfscope}%
\begin{pgfscope}%
\pgfpathrectangle{\pgfqpoint{3.815000in}{0.483000in}}{\pgfqpoint{2.835000in}{0.630000in}}%
\pgfusepath{clip}%
\pgfsetbuttcap%
\pgfsetmiterjoin%
\definecolor{currentfill}{rgb}{0.101961,0.101961,0.101961}%
\pgfsetfillcolor{currentfill}%
\pgfsetlinewidth{0.000000pt}%
\definecolor{currentstroke}{rgb}{0.000000,0.000000,0.000000}%
\pgfsetstrokecolor{currentstroke}%
\pgfsetstrokeopacity{0.000000}%
\pgfsetdash{}{0pt}%
\pgfpathmoveto{\pgfqpoint{6.202306in}{0.483000in}}%
\pgfpathlineto{\pgfqpoint{6.223569in}{0.483000in}}%
\pgfpathlineto{\pgfqpoint{6.223569in}{0.483000in}}%
\pgfpathlineto{\pgfqpoint{6.202306in}{0.483000in}}%
\pgfpathlineto{\pgfqpoint{6.202306in}{0.483000in}}%
\pgfpathclose%
\pgfusepath{fill}%
\end{pgfscope}%
\begin{pgfscope}%
\pgfpathrectangle{\pgfqpoint{3.815000in}{0.483000in}}{\pgfqpoint{2.835000in}{0.630000in}}%
\pgfusepath{clip}%
\pgfsetbuttcap%
\pgfsetmiterjoin%
\definecolor{currentfill}{rgb}{0.101961,0.101961,0.101961}%
\pgfsetfillcolor{currentfill}%
\pgfsetlinewidth{0.000000pt}%
\definecolor{currentstroke}{rgb}{0.000000,0.000000,0.000000}%
\pgfsetstrokecolor{currentstroke}%
\pgfsetstrokeopacity{0.000000}%
\pgfsetdash{}{0pt}%
\pgfpathmoveto{\pgfqpoint{6.225931in}{0.483000in}}%
\pgfpathlineto{\pgfqpoint{6.247194in}{0.483000in}}%
\pgfpathlineto{\pgfqpoint{6.247194in}{0.483000in}}%
\pgfpathlineto{\pgfqpoint{6.225931in}{0.483000in}}%
\pgfpathlineto{\pgfqpoint{6.225931in}{0.483000in}}%
\pgfpathclose%
\pgfusepath{fill}%
\end{pgfscope}%
\begin{pgfscope}%
\pgfpathrectangle{\pgfqpoint{3.815000in}{0.483000in}}{\pgfqpoint{2.835000in}{0.630000in}}%
\pgfusepath{clip}%
\pgfsetbuttcap%
\pgfsetmiterjoin%
\definecolor{currentfill}{rgb}{0.101961,0.101961,0.101961}%
\pgfsetfillcolor{currentfill}%
\pgfsetlinewidth{0.000000pt}%
\definecolor{currentstroke}{rgb}{0.000000,0.000000,0.000000}%
\pgfsetstrokecolor{currentstroke}%
\pgfsetstrokeopacity{0.000000}%
\pgfsetdash{}{0pt}%
\pgfpathmoveto{\pgfqpoint{6.249556in}{0.483000in}}%
\pgfpathlineto{\pgfqpoint{6.270819in}{0.483000in}}%
\pgfpathlineto{\pgfqpoint{6.270819in}{0.663000in}}%
\pgfpathlineto{\pgfqpoint{6.249556in}{0.663000in}}%
\pgfpathlineto{\pgfqpoint{6.249556in}{0.483000in}}%
\pgfpathclose%
\pgfusepath{fill}%
\end{pgfscope}%
\begin{pgfscope}%
\pgfpathrectangle{\pgfqpoint{3.815000in}{0.483000in}}{\pgfqpoint{2.835000in}{0.630000in}}%
\pgfusepath{clip}%
\pgfsetbuttcap%
\pgfsetmiterjoin%
\definecolor{currentfill}{rgb}{0.101961,0.101961,0.101961}%
\pgfsetfillcolor{currentfill}%
\pgfsetlinewidth{0.000000pt}%
\definecolor{currentstroke}{rgb}{0.000000,0.000000,0.000000}%
\pgfsetstrokecolor{currentstroke}%
\pgfsetstrokeopacity{0.000000}%
\pgfsetdash{}{0pt}%
\pgfpathmoveto{\pgfqpoint{6.273181in}{0.483000in}}%
\pgfpathlineto{\pgfqpoint{6.294444in}{0.483000in}}%
\pgfpathlineto{\pgfqpoint{6.294444in}{0.483000in}}%
\pgfpathlineto{\pgfqpoint{6.273181in}{0.483000in}}%
\pgfpathlineto{\pgfqpoint{6.273181in}{0.483000in}}%
\pgfpathclose%
\pgfusepath{fill}%
\end{pgfscope}%
\begin{pgfscope}%
\pgfpathrectangle{\pgfqpoint{3.815000in}{0.483000in}}{\pgfqpoint{2.835000in}{0.630000in}}%
\pgfusepath{clip}%
\pgfsetbuttcap%
\pgfsetmiterjoin%
\definecolor{currentfill}{rgb}{0.101961,0.101961,0.101961}%
\pgfsetfillcolor{currentfill}%
\pgfsetlinewidth{0.000000pt}%
\definecolor{currentstroke}{rgb}{0.000000,0.000000,0.000000}%
\pgfsetstrokecolor{currentstroke}%
\pgfsetstrokeopacity{0.000000}%
\pgfsetdash{}{0pt}%
\pgfpathmoveto{\pgfqpoint{6.296806in}{0.483000in}}%
\pgfpathlineto{\pgfqpoint{6.318069in}{0.483000in}}%
\pgfpathlineto{\pgfqpoint{6.318069in}{0.528000in}}%
\pgfpathlineto{\pgfqpoint{6.296806in}{0.528000in}}%
\pgfpathlineto{\pgfqpoint{6.296806in}{0.483000in}}%
\pgfpathclose%
\pgfusepath{fill}%
\end{pgfscope}%
\begin{pgfscope}%
\pgfpathrectangle{\pgfqpoint{3.815000in}{0.483000in}}{\pgfqpoint{2.835000in}{0.630000in}}%
\pgfusepath{clip}%
\pgfsetbuttcap%
\pgfsetmiterjoin%
\definecolor{currentfill}{rgb}{0.101961,0.101961,0.101961}%
\pgfsetfillcolor{currentfill}%
\pgfsetlinewidth{0.000000pt}%
\definecolor{currentstroke}{rgb}{0.000000,0.000000,0.000000}%
\pgfsetstrokecolor{currentstroke}%
\pgfsetstrokeopacity{0.000000}%
\pgfsetdash{}{0pt}%
\pgfpathmoveto{\pgfqpoint{6.320431in}{0.483000in}}%
\pgfpathlineto{\pgfqpoint{6.341694in}{0.483000in}}%
\pgfpathlineto{\pgfqpoint{6.341694in}{0.528000in}}%
\pgfpathlineto{\pgfqpoint{6.320431in}{0.528000in}}%
\pgfpathlineto{\pgfqpoint{6.320431in}{0.483000in}}%
\pgfpathclose%
\pgfusepath{fill}%
\end{pgfscope}%
\begin{pgfscope}%
\pgfpathrectangle{\pgfqpoint{3.815000in}{0.483000in}}{\pgfqpoint{2.835000in}{0.630000in}}%
\pgfusepath{clip}%
\pgfsetbuttcap%
\pgfsetmiterjoin%
\definecolor{currentfill}{rgb}{0.101961,0.101961,0.101961}%
\pgfsetfillcolor{currentfill}%
\pgfsetlinewidth{0.000000pt}%
\definecolor{currentstroke}{rgb}{0.000000,0.000000,0.000000}%
\pgfsetstrokecolor{currentstroke}%
\pgfsetstrokeopacity{0.000000}%
\pgfsetdash{}{0pt}%
\pgfpathmoveto{\pgfqpoint{6.344056in}{0.483000in}}%
\pgfpathlineto{\pgfqpoint{6.365319in}{0.483000in}}%
\pgfpathlineto{\pgfqpoint{6.365319in}{0.483000in}}%
\pgfpathlineto{\pgfqpoint{6.344056in}{0.483000in}}%
\pgfpathlineto{\pgfqpoint{6.344056in}{0.483000in}}%
\pgfpathclose%
\pgfusepath{fill}%
\end{pgfscope}%
\begin{pgfscope}%
\pgfpathrectangle{\pgfqpoint{3.815000in}{0.483000in}}{\pgfqpoint{2.835000in}{0.630000in}}%
\pgfusepath{clip}%
\pgfsetbuttcap%
\pgfsetmiterjoin%
\definecolor{currentfill}{rgb}{0.101961,0.101961,0.101961}%
\pgfsetfillcolor{currentfill}%
\pgfsetlinewidth{0.000000pt}%
\definecolor{currentstroke}{rgb}{0.000000,0.000000,0.000000}%
\pgfsetstrokecolor{currentstroke}%
\pgfsetstrokeopacity{0.000000}%
\pgfsetdash{}{0pt}%
\pgfpathmoveto{\pgfqpoint{6.367681in}{0.483000in}}%
\pgfpathlineto{\pgfqpoint{6.388944in}{0.483000in}}%
\pgfpathlineto{\pgfqpoint{6.388944in}{0.483000in}}%
\pgfpathlineto{\pgfqpoint{6.367681in}{0.483000in}}%
\pgfpathlineto{\pgfqpoint{6.367681in}{0.483000in}}%
\pgfpathclose%
\pgfusepath{fill}%
\end{pgfscope}%
\begin{pgfscope}%
\pgfpathrectangle{\pgfqpoint{3.815000in}{0.483000in}}{\pgfqpoint{2.835000in}{0.630000in}}%
\pgfusepath{clip}%
\pgfsetbuttcap%
\pgfsetmiterjoin%
\definecolor{currentfill}{rgb}{0.101961,0.101961,0.101961}%
\pgfsetfillcolor{currentfill}%
\pgfsetlinewidth{0.000000pt}%
\definecolor{currentstroke}{rgb}{0.000000,0.000000,0.000000}%
\pgfsetstrokecolor{currentstroke}%
\pgfsetstrokeopacity{0.000000}%
\pgfsetdash{}{0pt}%
\pgfpathmoveto{\pgfqpoint{6.391306in}{0.483000in}}%
\pgfpathlineto{\pgfqpoint{6.412569in}{0.483000in}}%
\pgfpathlineto{\pgfqpoint{6.412569in}{0.483000in}}%
\pgfpathlineto{\pgfqpoint{6.391306in}{0.483000in}}%
\pgfpathlineto{\pgfqpoint{6.391306in}{0.483000in}}%
\pgfpathclose%
\pgfusepath{fill}%
\end{pgfscope}%
\begin{pgfscope}%
\pgfpathrectangle{\pgfqpoint{3.815000in}{0.483000in}}{\pgfqpoint{2.835000in}{0.630000in}}%
\pgfusepath{clip}%
\pgfsetbuttcap%
\pgfsetmiterjoin%
\definecolor{currentfill}{rgb}{0.101961,0.101961,0.101961}%
\pgfsetfillcolor{currentfill}%
\pgfsetlinewidth{0.000000pt}%
\definecolor{currentstroke}{rgb}{0.000000,0.000000,0.000000}%
\pgfsetstrokecolor{currentstroke}%
\pgfsetstrokeopacity{0.000000}%
\pgfsetdash{}{0pt}%
\pgfpathmoveto{\pgfqpoint{6.414931in}{0.483000in}}%
\pgfpathlineto{\pgfqpoint{6.436194in}{0.483000in}}%
\pgfpathlineto{\pgfqpoint{6.436194in}{0.483000in}}%
\pgfpathlineto{\pgfqpoint{6.414931in}{0.483000in}}%
\pgfpathlineto{\pgfqpoint{6.414931in}{0.483000in}}%
\pgfpathclose%
\pgfusepath{fill}%
\end{pgfscope}%
\begin{pgfscope}%
\pgfpathrectangle{\pgfqpoint{3.815000in}{0.483000in}}{\pgfqpoint{2.835000in}{0.630000in}}%
\pgfusepath{clip}%
\pgfsetbuttcap%
\pgfsetmiterjoin%
\definecolor{currentfill}{rgb}{0.101961,0.101961,0.101961}%
\pgfsetfillcolor{currentfill}%
\pgfsetlinewidth{0.000000pt}%
\definecolor{currentstroke}{rgb}{0.000000,0.000000,0.000000}%
\pgfsetstrokecolor{currentstroke}%
\pgfsetstrokeopacity{0.000000}%
\pgfsetdash{}{0pt}%
\pgfpathmoveto{\pgfqpoint{6.438556in}{0.483000in}}%
\pgfpathlineto{\pgfqpoint{6.459819in}{0.483000in}}%
\pgfpathlineto{\pgfqpoint{6.459819in}{0.708000in}}%
\pgfpathlineto{\pgfqpoint{6.438556in}{0.708000in}}%
\pgfpathlineto{\pgfqpoint{6.438556in}{0.483000in}}%
\pgfpathclose%
\pgfusepath{fill}%
\end{pgfscope}%
\begin{pgfscope}%
\pgfpathrectangle{\pgfqpoint{3.815000in}{0.483000in}}{\pgfqpoint{2.835000in}{0.630000in}}%
\pgfusepath{clip}%
\pgfsetbuttcap%
\pgfsetmiterjoin%
\definecolor{currentfill}{rgb}{0.101961,0.101961,0.101961}%
\pgfsetfillcolor{currentfill}%
\pgfsetlinewidth{0.000000pt}%
\definecolor{currentstroke}{rgb}{0.000000,0.000000,0.000000}%
\pgfsetstrokecolor{currentstroke}%
\pgfsetstrokeopacity{0.000000}%
\pgfsetdash{}{0pt}%
\pgfpathmoveto{\pgfqpoint{6.462181in}{0.483000in}}%
\pgfpathlineto{\pgfqpoint{6.483444in}{0.483000in}}%
\pgfpathlineto{\pgfqpoint{6.483444in}{0.483000in}}%
\pgfpathlineto{\pgfqpoint{6.462181in}{0.483000in}}%
\pgfpathlineto{\pgfqpoint{6.462181in}{0.483000in}}%
\pgfpathclose%
\pgfusepath{fill}%
\end{pgfscope}%
\begin{pgfscope}%
\pgfpathrectangle{\pgfqpoint{3.815000in}{0.483000in}}{\pgfqpoint{2.835000in}{0.630000in}}%
\pgfusepath{clip}%
\pgfsetbuttcap%
\pgfsetmiterjoin%
\definecolor{currentfill}{rgb}{0.101961,0.101961,0.101961}%
\pgfsetfillcolor{currentfill}%
\pgfsetlinewidth{0.000000pt}%
\definecolor{currentstroke}{rgb}{0.000000,0.000000,0.000000}%
\pgfsetstrokecolor{currentstroke}%
\pgfsetstrokeopacity{0.000000}%
\pgfsetdash{}{0pt}%
\pgfpathmoveto{\pgfqpoint{6.485806in}{0.483000in}}%
\pgfpathlineto{\pgfqpoint{6.507069in}{0.483000in}}%
\pgfpathlineto{\pgfqpoint{6.507069in}{0.483000in}}%
\pgfpathlineto{\pgfqpoint{6.485806in}{0.483000in}}%
\pgfpathlineto{\pgfqpoint{6.485806in}{0.483000in}}%
\pgfpathclose%
\pgfusepath{fill}%
\end{pgfscope}%
\begin{pgfscope}%
\pgfpathrectangle{\pgfqpoint{3.815000in}{0.483000in}}{\pgfqpoint{2.835000in}{0.630000in}}%
\pgfusepath{clip}%
\pgfsetbuttcap%
\pgfsetmiterjoin%
\definecolor{currentfill}{rgb}{0.101961,0.101961,0.101961}%
\pgfsetfillcolor{currentfill}%
\pgfsetlinewidth{0.000000pt}%
\definecolor{currentstroke}{rgb}{0.000000,0.000000,0.000000}%
\pgfsetstrokecolor{currentstroke}%
\pgfsetstrokeopacity{0.000000}%
\pgfsetdash{}{0pt}%
\pgfpathmoveto{\pgfqpoint{6.509431in}{0.483000in}}%
\pgfpathlineto{\pgfqpoint{6.530694in}{0.483000in}}%
\pgfpathlineto{\pgfqpoint{6.530694in}{0.528000in}}%
\pgfpathlineto{\pgfqpoint{6.509431in}{0.528000in}}%
\pgfpathlineto{\pgfqpoint{6.509431in}{0.483000in}}%
\pgfpathclose%
\pgfusepath{fill}%
\end{pgfscope}%
\begin{pgfscope}%
\pgfpathrectangle{\pgfqpoint{3.815000in}{0.483000in}}{\pgfqpoint{2.835000in}{0.630000in}}%
\pgfusepath{clip}%
\pgfsetbuttcap%
\pgfsetmiterjoin%
\definecolor{currentfill}{rgb}{0.101961,0.101961,0.101961}%
\pgfsetfillcolor{currentfill}%
\pgfsetlinewidth{0.000000pt}%
\definecolor{currentstroke}{rgb}{0.000000,0.000000,0.000000}%
\pgfsetstrokecolor{currentstroke}%
\pgfsetstrokeopacity{0.000000}%
\pgfsetdash{}{0pt}%
\pgfpathmoveto{\pgfqpoint{6.533056in}{0.483000in}}%
\pgfpathlineto{\pgfqpoint{6.554319in}{0.483000in}}%
\pgfpathlineto{\pgfqpoint{6.554319in}{0.573000in}}%
\pgfpathlineto{\pgfqpoint{6.533056in}{0.573000in}}%
\pgfpathlineto{\pgfqpoint{6.533056in}{0.483000in}}%
\pgfpathclose%
\pgfusepath{fill}%
\end{pgfscope}%
\begin{pgfscope}%
\pgfpathrectangle{\pgfqpoint{3.815000in}{0.483000in}}{\pgfqpoint{2.835000in}{0.630000in}}%
\pgfusepath{clip}%
\pgfsetbuttcap%
\pgfsetmiterjoin%
\definecolor{currentfill}{rgb}{0.101961,0.101961,0.101961}%
\pgfsetfillcolor{currentfill}%
\pgfsetlinewidth{0.000000pt}%
\definecolor{currentstroke}{rgb}{0.000000,0.000000,0.000000}%
\pgfsetstrokecolor{currentstroke}%
\pgfsetstrokeopacity{0.000000}%
\pgfsetdash{}{0pt}%
\pgfpathmoveto{\pgfqpoint{6.556681in}{0.483000in}}%
\pgfpathlineto{\pgfqpoint{6.577944in}{0.483000in}}%
\pgfpathlineto{\pgfqpoint{6.577944in}{0.663000in}}%
\pgfpathlineto{\pgfqpoint{6.556681in}{0.663000in}}%
\pgfpathlineto{\pgfqpoint{6.556681in}{0.483000in}}%
\pgfpathclose%
\pgfusepath{fill}%
\end{pgfscope}%
\begin{pgfscope}%
\pgfpathrectangle{\pgfqpoint{3.815000in}{0.483000in}}{\pgfqpoint{2.835000in}{0.630000in}}%
\pgfusepath{clip}%
\pgfsetbuttcap%
\pgfsetmiterjoin%
\definecolor{currentfill}{rgb}{0.101961,0.101961,0.101961}%
\pgfsetfillcolor{currentfill}%
\pgfsetlinewidth{0.000000pt}%
\definecolor{currentstroke}{rgb}{0.000000,0.000000,0.000000}%
\pgfsetstrokecolor{currentstroke}%
\pgfsetstrokeopacity{0.000000}%
\pgfsetdash{}{0pt}%
\pgfpathmoveto{\pgfqpoint{6.580306in}{0.483000in}}%
\pgfpathlineto{\pgfqpoint{6.601569in}{0.483000in}}%
\pgfpathlineto{\pgfqpoint{6.601569in}{0.528000in}}%
\pgfpathlineto{\pgfqpoint{6.580306in}{0.528000in}}%
\pgfpathlineto{\pgfqpoint{6.580306in}{0.483000in}}%
\pgfpathclose%
\pgfusepath{fill}%
\end{pgfscope}%
\begin{pgfscope}%
\pgfpathrectangle{\pgfqpoint{3.815000in}{0.483000in}}{\pgfqpoint{2.835000in}{0.630000in}}%
\pgfusepath{clip}%
\pgfsetbuttcap%
\pgfsetmiterjoin%
\definecolor{currentfill}{rgb}{0.101961,0.101961,0.101961}%
\pgfsetfillcolor{currentfill}%
\pgfsetlinewidth{0.000000pt}%
\definecolor{currentstroke}{rgb}{0.000000,0.000000,0.000000}%
\pgfsetstrokecolor{currentstroke}%
\pgfsetstrokeopacity{0.000000}%
\pgfsetdash{}{0pt}%
\pgfpathmoveto{\pgfqpoint{6.603931in}{0.483000in}}%
\pgfpathlineto{\pgfqpoint{6.625194in}{0.483000in}}%
\pgfpathlineto{\pgfqpoint{6.625194in}{0.618000in}}%
\pgfpathlineto{\pgfqpoint{6.603931in}{0.618000in}}%
\pgfpathlineto{\pgfqpoint{6.603931in}{0.483000in}}%
\pgfpathclose%
\pgfusepath{fill}%
\end{pgfscope}%
\begin{pgfscope}%
\pgfpathrectangle{\pgfqpoint{3.815000in}{0.483000in}}{\pgfqpoint{2.835000in}{0.630000in}}%
\pgfusepath{clip}%
\pgfsetbuttcap%
\pgfsetmiterjoin%
\definecolor{currentfill}{rgb}{0.101961,0.101961,0.101961}%
\pgfsetfillcolor{currentfill}%
\pgfsetlinewidth{0.000000pt}%
\definecolor{currentstroke}{rgb}{0.000000,0.000000,0.000000}%
\pgfsetstrokecolor{currentstroke}%
\pgfsetstrokeopacity{0.000000}%
\pgfsetdash{}{0pt}%
\pgfpathmoveto{\pgfqpoint{6.627556in}{0.483000in}}%
\pgfpathlineto{\pgfqpoint{6.648819in}{0.483000in}}%
\pgfpathlineto{\pgfqpoint{6.648819in}{0.483000in}}%
\pgfpathlineto{\pgfqpoint{6.627556in}{0.483000in}}%
\pgfpathlineto{\pgfqpoint{6.627556in}{0.483000in}}%
\pgfpathclose%
\pgfusepath{fill}%
\end{pgfscope}%
\begin{pgfscope}%
\pgfpathrectangle{\pgfqpoint{3.815000in}{0.483000in}}{\pgfqpoint{2.835000in}{0.630000in}}%
\pgfusepath{clip}%
\pgfsetbuttcap%
\pgfsetroundjoin%
\pgfsetlinewidth{0.702625pt}%
\definecolor{currentstroke}{rgb}{0.690196,0.250980,0.250980}%
\pgfsetstrokecolor{currentstroke}%
\pgfsetdash{{2.590000pt}{1.120000pt}}{0.000000pt}%
\pgfpathmoveto{\pgfqpoint{5.221869in}{0.483000in}}%
\pgfpathlineto{\pgfqpoint{5.221869in}{1.113000in}}%
\pgfusepath{stroke}%
\end{pgfscope}%
\begin{pgfscope}%
\pgfsetrectcap%
\pgfsetmiterjoin%
\pgfsetlinewidth{0.803000pt}%
\definecolor{currentstroke}{rgb}{0.101961,0.101961,0.101961}%
\pgfsetstrokecolor{currentstroke}%
\pgfsetdash{}{0pt}%
\pgfpathmoveto{\pgfqpoint{3.815000in}{0.483000in}}%
\pgfpathlineto{\pgfqpoint{3.815000in}{1.113000in}}%
\pgfusepath{stroke}%
\end{pgfscope}%
\begin{pgfscope}%
\pgfsetrectcap%
\pgfsetmiterjoin%
\pgfsetlinewidth{0.803000pt}%
\definecolor{currentstroke}{rgb}{0.101961,0.101961,0.101961}%
\pgfsetstrokecolor{currentstroke}%
\pgfsetdash{}{0pt}%
\pgfpathmoveto{\pgfqpoint{3.815000in}{0.483000in}}%
\pgfpathlineto{\pgfqpoint{6.650000in}{0.483000in}}%
\pgfusepath{stroke}%
\end{pgfscope}%
\begin{pgfscope}%
\definecolor{textcolor}{rgb}{0.549020,0.549020,0.549020}%
\pgfsetstrokecolor{textcolor}%
\pgfsetfillcolor{textcolor}%
\pgftext[x=6.930000in,y=4.137000in,right,top]{\color{textcolor}\sffamily\fontsize{5.500000}{6.600000}\selectfont DEV DATA}%
\end{pgfscope}%
\end{pgfpicture}%
\makeatother%
\endgroup%
}
\caption{[PLACEHOLDER: DEV DATA, COVEX dev (cand-tl), n20 seed 4] One run on the engine, time on the horizontal axis. (a) A normal run at 1x. (b) The same fleet with 14 of 20 robots offline from 407.5 s to 595.5 s (shaded; dashed line: reconnect). Top: engine-measured KV pool and host-tier occupancy. Second and third rows: warm actions (each paired with a touch, so one line) and deferrals per 10 s. Bottom row: alerts delivered after the 1 s deadline per 10 s (undelivered alerts are not counted).}
\label{fig:F06_timeline}
\end{figure*}
\clearpage
\begin{figure}[t]
\centering
{\fontfamily{LinuxBiolinumT}\selectfont
%% Creator: Matplotlib, PGF backend
%%
%% To include the figure in your LaTeX document, write
%%   \input{<filename>.pgf}
%%
%% Make sure the required packages are loaded in your preamble
%%   \usepackage{pgf}
%%
%% Also ensure that all the required font packages are loaded; for instance,
%% the lmodern package is sometimes necessary when using math font.
%%   \usepackage{lmodern}
%%
%% Figures using additional raster images can only be included by \input if
%% they are in the same directory as the main LaTeX file. For loading figures
%% from other directories you can use the `import` package
%%   \usepackage{import}
%%
%% and then include the figures with
%%   \import{<path to file>}{<filename>.pgf}
%%
%% Matplotlib used the following preamble
%%
\begingroup%
\makeatletter%
\begin{pgfpicture}%
\pgfpathrectangle{\pgfpointorigin}{\pgfqpoint{3.330000in}{2.200000in}}%
\pgfusepath{use as bounding box, clip}%
\begin{pgfscope}%
\pgfsetbuttcap%
\pgfsetmiterjoin%
\definecolor{currentfill}{rgb}{1.000000,1.000000,1.000000}%
\pgfsetfillcolor{currentfill}%
\pgfsetlinewidth{0.000000pt}%
\definecolor{currentstroke}{rgb}{1.000000,1.000000,1.000000}%
\pgfsetstrokecolor{currentstroke}%
\pgfsetdash{}{0pt}%
\pgfpathmoveto{\pgfqpoint{0.000000in}{0.000000in}}%
\pgfpathlineto{\pgfqpoint{3.330000in}{0.000000in}}%
\pgfpathlineto{\pgfqpoint{3.330000in}{2.200000in}}%
\pgfpathlineto{\pgfqpoint{0.000000in}{2.200000in}}%
\pgfpathlineto{\pgfqpoint{0.000000in}{0.000000in}}%
\pgfpathclose%
\pgfusepath{fill}%
\end{pgfscope}%
\begin{pgfscope}%
\pgfsetbuttcap%
\pgfsetmiterjoin%
\definecolor{currentfill}{rgb}{1.000000,1.000000,1.000000}%
\pgfsetfillcolor{currentfill}%
\pgfsetlinewidth{0.000000pt}%
\definecolor{currentstroke}{rgb}{0.000000,0.000000,0.000000}%
\pgfsetstrokecolor{currentstroke}%
\pgfsetstrokeopacity{0.000000}%
\pgfsetdash{}{0pt}%
\pgfpathmoveto{\pgfqpoint{0.566100in}{0.594000in}}%
\pgfpathlineto{\pgfqpoint{3.230100in}{0.594000in}}%
\pgfpathlineto{\pgfqpoint{3.230100in}{1.958000in}}%
\pgfpathlineto{\pgfqpoint{0.566100in}{1.958000in}}%
\pgfpathlineto{\pgfqpoint{0.566100in}{0.594000in}}%
\pgfpathclose%
\pgfusepath{fill}%
\end{pgfscope}%
\begin{pgfscope}%
\definecolor{textcolor}{rgb}{0.101961,0.101961,0.101961}%
\pgfsetstrokecolor{textcolor}%
\pgfsetfillcolor{textcolor}%
\pgftext[x=0.873485in,y=0.545389in,,top]{\color{textcolor}\sffamily\fontsize{7.000000}{8.400000}\selectfont KEEP}%
\end{pgfscope}%
\begin{pgfscope}%
\definecolor{textcolor}{rgb}{0.101961,0.101961,0.101961}%
\pgfsetstrokecolor{textcolor}%
\pgfsetfillcolor{textcolor}%
\pgftext[x=1.385792in,y=0.545389in,,top]{\color{textcolor}\sffamily\fontsize{7.000000}{8.400000}\selectfont TOUCH}%
\end{pgfscope}%
\begin{pgfscope}%
\definecolor{textcolor}{rgb}{0.101961,0.101961,0.101961}%
\pgfsetstrokecolor{textcolor}%
\pgfsetfillcolor{textcolor}%
\pgftext[x=1.898100in,y=0.545389in,,top]{\color{textcolor}\sffamily\fontsize{7.000000}{8.400000}\selectfont RELOAD}%
\end{pgfscope}%
\begin{pgfscope}%
\definecolor{textcolor}{rgb}{0.101961,0.101961,0.101961}%
\pgfsetstrokecolor{textcolor}%
\pgfsetfillcolor{textcolor}%
\pgftext[x=2.270396in, y=0.477874in, left, base]{\color{textcolor}\sffamily\fontsize{7.000000}{8.400000}\selectfont HOLD}%
\end{pgfscope}%
\begin{pgfscope}%
\definecolor{textcolor}{rgb}{0.101961,0.101961,0.101961}%
\pgfsetstrokecolor{textcolor}%
\pgfsetfillcolor{textcolor}%
\pgftext[x=2.274109in, y=0.377951in, left, base]{\color{textcolor}\sffamily\fontsize{7.000000}{8.400000}\selectfont BACK}%
\end{pgfscope}%
\begin{pgfscope}%
\definecolor{textcolor}{rgb}{0.101961,0.101961,0.101961}%
\pgfsetstrokecolor{textcolor}%
\pgfsetfillcolor{textcolor}%
\pgftext[x=2.922715in,y=0.545389in,,top]{\color{textcolor}\sffamily\fontsize{7.000000}{8.400000}\selectfont SLICE}%
\end{pgfscope}%
\begin{pgfscope}%
\definecolor{textcolor}{rgb}{0.000000,0.000000,0.000000}%
\pgfsetstrokecolor{textcolor}%
\pgfsetfillcolor{textcolor}%
\pgftext[x=1.898100in,y=0.331269in,,top]{\color{textcolor}\sffamily\fontsize{7.500000}{9.000000}\selectfont Decision kind}%
\end{pgfscope}%
\begin{pgfscope}%
\pgfpathrectangle{\pgfqpoint{0.566100in}{0.594000in}}{\pgfqpoint{2.664000in}{1.364000in}}%
\pgfusepath{clip}%
\pgfsetrectcap%
\pgfsetroundjoin%
\pgfsetlinewidth{0.401500pt}%
\definecolor{currentstroke}{rgb}{0.901961,0.901961,0.901961}%
\pgfsetstrokecolor{currentstroke}%
\pgfsetdash{}{0pt}%
\pgfpathmoveto{\pgfqpoint{0.566100in}{0.594000in}}%
\pgfpathlineto{\pgfqpoint{3.230100in}{0.594000in}}%
\pgfusepath{stroke}%
\end{pgfscope}%
\begin{pgfscope}%
\definecolor{textcolor}{rgb}{0.101961,0.101961,0.101961}%
\pgfsetstrokecolor{textcolor}%
\pgfsetfillcolor{textcolor}%
\pgftext[x=0.465839in, y=0.560242in, left, base]{\color{textcolor}\sffamily\fontsize{7.000000}{8.400000}\selectfont 0}%
\end{pgfscope}%
\begin{pgfscope}%
\pgfpathrectangle{\pgfqpoint{0.566100in}{0.594000in}}{\pgfqpoint{2.664000in}{1.364000in}}%
\pgfusepath{clip}%
\pgfsetrectcap%
\pgfsetroundjoin%
\pgfsetlinewidth{0.401500pt}%
\definecolor{currentstroke}{rgb}{0.901961,0.901961,0.901961}%
\pgfsetstrokecolor{currentstroke}%
\pgfsetdash{}{0pt}%
\pgfpathmoveto{\pgfqpoint{0.566100in}{1.276000in}}%
\pgfpathlineto{\pgfqpoint{3.230100in}{1.276000in}}%
\pgfusepath{stroke}%
\end{pgfscope}%
\begin{pgfscope}%
\pgfpathrectangle{\pgfqpoint{0.566100in}{0.594000in}}{\pgfqpoint{2.664000in}{1.364000in}}%
\pgfusepath{clip}%
\pgfsetrectcap%
\pgfsetroundjoin%
\pgfsetlinewidth{0.401500pt}%
\definecolor{currentstroke}{rgb}{0.901961,0.901961,0.901961}%
\pgfsetstrokecolor{currentstroke}%
\pgfsetdash{}{0pt}%
\pgfpathmoveto{\pgfqpoint{0.566100in}{1.958000in}}%
\pgfpathlineto{\pgfqpoint{3.230100in}{1.958000in}}%
\pgfusepath{stroke}%
\end{pgfscope}%
\begin{pgfscope}%
\definecolor{textcolor}{rgb}{0.000000,0.000000,0.000000}%
\pgfsetstrokecolor{textcolor}%
\pgfsetfillcolor{textcolor}%
\pgftext[x=0.410283in,y=1.276000in,,bottom,rotate=90.000000]{\color{textcolor}\sffamily\fontsize{7.500000}{9.000000}\selectfont Decisions per run}%
\end{pgfscope}%
\begin{pgfscope}%
\pgfsetrectcap%
\pgfsetmiterjoin%
\pgfsetlinewidth{0.803000pt}%
\definecolor{currentstroke}{rgb}{0.101961,0.101961,0.101961}%
\pgfsetstrokecolor{currentstroke}%
\pgfsetdash{}{0pt}%
\pgfpathmoveto{\pgfqpoint{0.566100in}{0.594000in}}%
\pgfpathlineto{\pgfqpoint{0.566100in}{1.958000in}}%
\pgfusepath{stroke}%
\end{pgfscope}%
\begin{pgfscope}%
\pgfsetrectcap%
\pgfsetmiterjoin%
\pgfsetlinewidth{0.803000pt}%
\definecolor{currentstroke}{rgb}{0.101961,0.101961,0.101961}%
\pgfsetstrokecolor{currentstroke}%
\pgfsetdash{}{0pt}%
\pgfpathmoveto{\pgfqpoint{0.566100in}{0.594000in}}%
\pgfpathlineto{\pgfqpoint{3.230100in}{0.594000in}}%
\pgfusepath{stroke}%
\end{pgfscope}%
\begin{pgfscope}%
\pgfsetbuttcap%
\pgfsetmiterjoin%
\definecolor{currentfill}{rgb}{0.156863,0.356863,0.568627}%
\pgfsetfillcolor{currentfill}%
\pgfsetlinewidth{0.000000pt}%
\definecolor{currentstroke}{rgb}{0.000000,0.000000,0.000000}%
\pgfsetstrokecolor{currentstroke}%
\pgfsetstrokeopacity{0.000000}%
\pgfsetdash{}{0pt}%
\pgfpathmoveto{\pgfqpoint{2.392497in}{1.802444in}}%
\pgfpathlineto{\pgfqpoint{2.509163in}{1.802444in}}%
\pgfpathlineto{\pgfqpoint{2.509163in}{1.870500in}}%
\pgfpathlineto{\pgfqpoint{2.392497in}{1.870500in}}%
\pgfpathlineto{\pgfqpoint{2.392497in}{1.802444in}}%
\pgfpathclose%
\pgfusepath{fill}%
\end{pgfscope}%
\begin{pgfscope}%
\definecolor{textcolor}{rgb}{0.101961,0.101961,0.101961}%
\pgfsetstrokecolor{textcolor}%
\pgfsetfillcolor{textcolor}%
\pgftext[x=2.586941in,y=1.802444in,left,base]{\color{textcolor}\sffamily\fontsize{7.000000}{8.400000}\selectfont 1x memory}%
\end{pgfscope}%
\begin{pgfscope}%
\pgfsetbuttcap%
\pgfsetmiterjoin%
\definecolor{currentfill}{rgb}{1.000000,1.000000,1.000000}%
\pgfsetfillcolor{currentfill}%
\pgfsetlinewidth{1.003750pt}%
\definecolor{currentstroke}{rgb}{0.156863,0.356863,0.568627}%
\pgfsetstrokecolor{currentstroke}%
\pgfsetdash{}{0pt}%
\pgfpathmoveto{\pgfqpoint{2.392497in}{1.666874in}}%
\pgfpathlineto{\pgfqpoint{2.509163in}{1.666874in}}%
\pgfpathlineto{\pgfqpoint{2.509163in}{1.734929in}}%
\pgfpathlineto{\pgfqpoint{2.392497in}{1.734929in}}%
\pgfpathlineto{\pgfqpoint{2.392497in}{1.666874in}}%
\pgfpathclose%
\pgfusepath{stroke,fill}%
\end{pgfscope}%
\begin{pgfscope}%
\pgfsetbuttcap%
\pgfsetmiterjoin%
\definecolor{currentfill}{rgb}{1.000000,1.000000,1.000000}%
\pgfsetfillcolor{currentfill}%
\pgfsetlinewidth{1.003750pt}%
\definecolor{currentstroke}{rgb}{0.156863,0.356863,0.568627}%
\pgfsetstrokecolor{currentstroke}%
\pgfsetdash{}{0pt}%
\pgfpathmoveto{\pgfqpoint{2.392497in}{1.666874in}}%
\pgfpathlineto{\pgfqpoint{2.509163in}{1.666874in}}%
\pgfpathlineto{\pgfqpoint{2.509163in}{1.734929in}}%
\pgfpathlineto{\pgfqpoint{2.392497in}{1.734929in}}%
\pgfpathlineto{\pgfqpoint{2.392497in}{1.666874in}}%
\pgfpathclose%
\pgfusepath{clip}%
\pgfsys@defobject{currentpattern}{\pgfqpoint{0in}{0in}}{\pgfqpoint{1in}{1in}}{%
\begin{pgfscope}%
\pgfpathrectangle{\pgfqpoint{0in}{0in}}{\pgfqpoint{1in}{1in}}%
\pgfusepath{clip}%
\pgfpathmoveto{\pgfqpoint{-0.500000in}{0.500000in}}%
\pgfpathlineto{\pgfqpoint{0.500000in}{1.500000in}}%
\pgfpathmoveto{\pgfqpoint{-0.458333in}{0.458333in}}%
\pgfpathlineto{\pgfqpoint{0.541667in}{1.458333in}}%
\pgfpathmoveto{\pgfqpoint{-0.416667in}{0.416667in}}%
\pgfpathlineto{\pgfqpoint{0.583333in}{1.416667in}}%
\pgfpathmoveto{\pgfqpoint{-0.375000in}{0.375000in}}%
\pgfpathlineto{\pgfqpoint{0.625000in}{1.375000in}}%
\pgfpathmoveto{\pgfqpoint{-0.333333in}{0.333333in}}%
\pgfpathlineto{\pgfqpoint{0.666667in}{1.333333in}}%
\pgfpathmoveto{\pgfqpoint{-0.291667in}{0.291667in}}%
\pgfpathlineto{\pgfqpoint{0.708333in}{1.291667in}}%
\pgfpathmoveto{\pgfqpoint{-0.250000in}{0.250000in}}%
\pgfpathlineto{\pgfqpoint{0.750000in}{1.250000in}}%
\pgfpathmoveto{\pgfqpoint{-0.208333in}{0.208333in}}%
\pgfpathlineto{\pgfqpoint{0.791667in}{1.208333in}}%
\pgfpathmoveto{\pgfqpoint{-0.166667in}{0.166667in}}%
\pgfpathlineto{\pgfqpoint{0.833333in}{1.166667in}}%
\pgfpathmoveto{\pgfqpoint{-0.125000in}{0.125000in}}%
\pgfpathlineto{\pgfqpoint{0.875000in}{1.125000in}}%
\pgfpathmoveto{\pgfqpoint{-0.083333in}{0.083333in}}%
\pgfpathlineto{\pgfqpoint{0.916667in}{1.083333in}}%
\pgfpathmoveto{\pgfqpoint{-0.041667in}{0.041667in}}%
\pgfpathlineto{\pgfqpoint{0.958333in}{1.041667in}}%
\pgfpathmoveto{\pgfqpoint{0.000000in}{0.000000in}}%
\pgfpathlineto{\pgfqpoint{1.000000in}{1.000000in}}%
\pgfpathmoveto{\pgfqpoint{0.041667in}{-0.041667in}}%
\pgfpathlineto{\pgfqpoint{1.041667in}{0.958333in}}%
\pgfpathmoveto{\pgfqpoint{0.083333in}{-0.083333in}}%
\pgfpathlineto{\pgfqpoint{1.083333in}{0.916667in}}%
\pgfpathmoveto{\pgfqpoint{0.125000in}{-0.125000in}}%
\pgfpathlineto{\pgfqpoint{1.125000in}{0.875000in}}%
\pgfpathmoveto{\pgfqpoint{0.166667in}{-0.166667in}}%
\pgfpathlineto{\pgfqpoint{1.166667in}{0.833333in}}%
\pgfpathmoveto{\pgfqpoint{0.208333in}{-0.208333in}}%
\pgfpathlineto{\pgfqpoint{1.208333in}{0.791667in}}%
\pgfpathmoveto{\pgfqpoint{0.250000in}{-0.250000in}}%
\pgfpathlineto{\pgfqpoint{1.250000in}{0.750000in}}%
\pgfpathmoveto{\pgfqpoint{0.291667in}{-0.291667in}}%
\pgfpathlineto{\pgfqpoint{1.291667in}{0.708333in}}%
\pgfpathmoveto{\pgfqpoint{0.333333in}{-0.333333in}}%
\pgfpathlineto{\pgfqpoint{1.333333in}{0.666667in}}%
\pgfpathmoveto{\pgfqpoint{0.375000in}{-0.375000in}}%
\pgfpathlineto{\pgfqpoint{1.375000in}{0.625000in}}%
\pgfpathmoveto{\pgfqpoint{0.416667in}{-0.416667in}}%
\pgfpathlineto{\pgfqpoint{1.416667in}{0.583333in}}%
\pgfpathmoveto{\pgfqpoint{0.458333in}{-0.458333in}}%
\pgfpathlineto{\pgfqpoint{1.458333in}{0.541667in}}%
\pgfpathmoveto{\pgfqpoint{0.500000in}{-0.500000in}}%
\pgfpathlineto{\pgfqpoint{1.500000in}{0.500000in}}%
\pgfusepath{stroke}%
\end{pgfscope}%
}%
\pgfsys@transformshift{2.392497in}{1.666874in}%
\pgfsys@useobject{currentpattern}{}%
\pgfsys@transformshift{1in}{0in}%
\pgfsys@transformshift{-1in}{0in}%
\pgfsys@transformshift{0in}{1in}%
\end{pgfscope}%
\begin{pgfscope}%
\definecolor{textcolor}{rgb}{0.101961,0.101961,0.101961}%
\pgfsetstrokecolor{textcolor}%
\pgfsetfillcolor{textcolor}%
\pgftext[x=2.586941in,y=1.666874in,left,base]{\color{textcolor}\sffamily\fontsize{7.000000}{8.400000}\selectfont 0.5x memory}%
\end{pgfscope}%
\begin{pgfscope}%
\pgfsetbuttcap%
\pgfsetmiterjoin%
\definecolor{currentfill}{rgb}{1.000000,1.000000,1.000000}%
\pgfsetfillcolor{currentfill}%
\pgfsetlinewidth{0.000000pt}%
\definecolor{currentstroke}{rgb}{0.000000,0.000000,0.000000}%
\pgfsetstrokecolor{currentstroke}%
\pgfsetstrokeopacity{0.000000}%
\pgfsetdash{}{0pt}%
\pgfpathmoveto{\pgfqpoint{0.593314in}{1.098732in}}%
\pgfpathlineto{\pgfqpoint{3.202886in}{1.098732in}}%
\pgfpathlineto{\pgfqpoint{3.202886in}{1.453268in}}%
\pgfpathlineto{\pgfqpoint{0.593314in}{1.453268in}}%
\pgfpathlineto{\pgfqpoint{0.593314in}{1.098732in}}%
\pgfpathclose%
\pgfusepath{fill}%
\end{pgfscope}%
\begin{pgfscope}%
\definecolor{textcolor}{rgb}{0.549020,0.549020,0.549020}%
\pgfsetstrokecolor{textcolor}%
\pgfsetfillcolor{textcolor}%
\pgftext[x=0.744011in, y=1.357558in, left, base]{\color{textcolor}\sffamily\fontsize{7.000000}{8.400000}\selectfont NO DATA YET: rung 3 state-planner decision counts}%
\end{pgfscope}%
\begin{pgfscope}%
\definecolor{textcolor}{rgb}{0.549020,0.549020,0.549020}%
\pgfsetstrokecolor{textcolor}%
\pgfsetfillcolor{textcolor}%
\pgftext[x=0.621508in, y=1.251154in, left, base]{\color{textcolor}\sffamily\fontsize{7.000000}{8.400000}\selectfont by kind (KEEP, TOUCH, RELOAD, HOLD\_BACK, SLICE)}%
\end{pgfscope}%
\begin{pgfscope}%
\definecolor{textcolor}{rgb}{0.549020,0.549020,0.549020}%
\pgfsetstrokecolor{textcolor}%
\pgfsetfillcolor{textcolor}%
\pgftext[x=0.983694in, y=1.145830in, left, base]{\color{textcolor}\sffamily\fontsize{7.000000}{8.400000}\selectfont at 1x and 0.5x memory; F7.csv has 0 rows}%
\end{pgfscope}%
\begin{pgfscope}%
\definecolor{textcolor}{rgb}{0.549020,0.549020,0.549020}%
\pgfsetstrokecolor{textcolor}%
\pgfsetfillcolor{textcolor}%
\pgftext[x=3.296700in,y=2.167000in,right,top]{\color{textcolor}\sffamily\fontsize{5.500000}{6.600000}\selectfont OLD DATA}%
\end{pgfscope}%
\end{pgfpicture}%
\makeatother%
\endgroup%
}
\caption{[PLACEHOLDER: OLD DATA, NO DATA, rung 3 stand-in for COVEX] Rung 3 state-planner decisions by kind at 1x and 0.5x memory. Empty: no per-kind decision log exists yet (F7.csv has no rows).}
\label{fig:F07_planner_decisions}
\end{figure}
\clearpage
\begin{figure}[t]
\centering
{\fontfamily{LinuxBiolinumT}\selectfont
%% Creator: Matplotlib, PGF backend
%%
%% To include the figure in your LaTeX document, write
%%   \input{<filename>.pgf}
%%
%% Make sure the required packages are loaded in your preamble
%%   \usepackage{pgf}
%%
%% Also ensure that all the required font packages are loaded; for instance,
%% the lmodern package is sometimes necessary when using math font.
%%   \usepackage{lmodern}
%%
%% Figures using additional raster images can only be included by \input if
%% they are in the same directory as the main LaTeX file. For loading figures
%% from other directories you can use the `import` package
%%   \usepackage{import}
%%
%% and then include the figures with
%%   \import{<path to file>}{<filename>.pgf}
%%
%% Matplotlib used the following preamble
%%
\begingroup%
\makeatletter%
\begin{pgfpicture}%
\pgfpathrectangle{\pgfpointorigin}{\pgfqpoint{3.330000in}{3.500000in}}%
\pgfusepath{use as bounding box, clip}%
\begin{pgfscope}%
\pgfsetbuttcap%
\pgfsetmiterjoin%
\definecolor{currentfill}{rgb}{1.000000,1.000000,1.000000}%
\pgfsetfillcolor{currentfill}%
\pgfsetlinewidth{0.000000pt}%
\definecolor{currentstroke}{rgb}{1.000000,1.000000,1.000000}%
\pgfsetstrokecolor{currentstroke}%
\pgfsetdash{}{0pt}%
\pgfpathmoveto{\pgfqpoint{0.000000in}{0.000000in}}%
\pgfpathlineto{\pgfqpoint{3.330000in}{0.000000in}}%
\pgfpathlineto{\pgfqpoint{3.330000in}{3.500000in}}%
\pgfpathlineto{\pgfqpoint{0.000000in}{3.500000in}}%
\pgfpathlineto{\pgfqpoint{0.000000in}{0.000000in}}%
\pgfpathclose%
\pgfusepath{fill}%
\end{pgfscope}%
\begin{pgfscope}%
\pgfsetbuttcap%
\pgfsetmiterjoin%
\definecolor{currentfill}{rgb}{1.000000,1.000000,1.000000}%
\pgfsetfillcolor{currentfill}%
\pgfsetlinewidth{0.000000pt}%
\definecolor{currentstroke}{rgb}{0.000000,0.000000,0.000000}%
\pgfsetstrokecolor{currentstroke}%
\pgfsetstrokeopacity{0.000000}%
\pgfsetdash{}{0pt}%
\pgfpathmoveto{\pgfqpoint{1.465200in}{0.822500in}}%
\pgfpathlineto{\pgfqpoint{2.364300in}{0.822500in}}%
\pgfpathlineto{\pgfqpoint{2.364300in}{3.360000in}}%
\pgfpathlineto{\pgfqpoint{1.465200in}{3.360000in}}%
\pgfpathlineto{\pgfqpoint{1.465200in}{0.822500in}}%
\pgfpathclose%
\pgfusepath{fill}%
\end{pgfscope}%
\begin{pgfscope}%
\pgfpathrectangle{\pgfqpoint{1.465200in}{0.822500in}}{\pgfqpoint{0.899100in}{2.537500in}}%
\pgfusepath{clip}%
\pgfsetrectcap%
\pgfsetroundjoin%
\pgfsetlinewidth{0.401500pt}%
\definecolor{currentstroke}{rgb}{0.901961,0.901961,0.901961}%
\pgfsetstrokecolor{currentstroke}%
\pgfsetdash{}{0pt}%
\pgfpathmoveto{\pgfqpoint{1.577588in}{0.822500in}}%
\pgfpathlineto{\pgfqpoint{1.577588in}{3.360000in}}%
\pgfusepath{stroke}%
\end{pgfscope}%
\begin{pgfscope}%
\definecolor{textcolor}{rgb}{0.101961,0.101961,0.101961}%
\pgfsetstrokecolor{textcolor}%
\pgfsetfillcolor{textcolor}%
\pgftext[x=1.577588in,y=0.773889in,,top]{\color{textcolor}\sffamily\fontsize{7.000000}{8.400000}\selectfont 0.4}%
\end{pgfscope}%
\begin{pgfscope}%
\pgfpathrectangle{\pgfqpoint{1.465200in}{0.822500in}}{\pgfqpoint{0.899100in}{2.537500in}}%
\pgfusepath{clip}%
\pgfsetrectcap%
\pgfsetroundjoin%
\pgfsetlinewidth{0.401500pt}%
\definecolor{currentstroke}{rgb}{0.901961,0.901961,0.901961}%
\pgfsetstrokecolor{currentstroke}%
\pgfsetdash{}{0pt}%
\pgfpathmoveto{\pgfqpoint{1.802362in}{0.822500in}}%
\pgfpathlineto{\pgfqpoint{1.802362in}{3.360000in}}%
\pgfusepath{stroke}%
\end{pgfscope}%
\begin{pgfscope}%
\definecolor{textcolor}{rgb}{0.101961,0.101961,0.101961}%
\pgfsetstrokecolor{textcolor}%
\pgfsetfillcolor{textcolor}%
\pgftext[x=1.802362in,y=0.773889in,,top]{\color{textcolor}\sffamily\fontsize{7.000000}{8.400000}\selectfont 0.5}%
\end{pgfscope}%
\begin{pgfscope}%
\pgfpathrectangle{\pgfqpoint{1.465200in}{0.822500in}}{\pgfqpoint{0.899100in}{2.537500in}}%
\pgfusepath{clip}%
\pgfsetrectcap%
\pgfsetroundjoin%
\pgfsetlinewidth{0.401500pt}%
\definecolor{currentstroke}{rgb}{0.901961,0.901961,0.901961}%
\pgfsetstrokecolor{currentstroke}%
\pgfsetdash{}{0pt}%
\pgfpathmoveto{\pgfqpoint{2.027138in}{0.822500in}}%
\pgfpathlineto{\pgfqpoint{2.027138in}{3.360000in}}%
\pgfusepath{stroke}%
\end{pgfscope}%
\begin{pgfscope}%
\definecolor{textcolor}{rgb}{0.101961,0.101961,0.101961}%
\pgfsetstrokecolor{textcolor}%
\pgfsetfillcolor{textcolor}%
\pgftext[x=2.027138in,y=0.773889in,,top]{\color{textcolor}\sffamily\fontsize{7.000000}{8.400000}\selectfont 0.6}%
\end{pgfscope}%
\begin{pgfscope}%
\pgfpathrectangle{\pgfqpoint{1.465200in}{0.822500in}}{\pgfqpoint{0.899100in}{2.537500in}}%
\pgfusepath{clip}%
\pgfsetrectcap%
\pgfsetroundjoin%
\pgfsetlinewidth{0.401500pt}%
\definecolor{currentstroke}{rgb}{0.901961,0.901961,0.901961}%
\pgfsetstrokecolor{currentstroke}%
\pgfsetdash{}{0pt}%
\pgfpathmoveto{\pgfqpoint{2.251912in}{0.822500in}}%
\pgfpathlineto{\pgfqpoint{2.251912in}{3.360000in}}%
\pgfusepath{stroke}%
\end{pgfscope}%
\begin{pgfscope}%
\definecolor{textcolor}{rgb}{0.101961,0.101961,0.101961}%
\pgfsetstrokecolor{textcolor}%
\pgfsetfillcolor{textcolor}%
\pgftext[x=2.251912in,y=0.773889in,,top]{\color{textcolor}\sffamily\fontsize{7.000000}{8.400000}\selectfont 0.7}%
\end{pgfscope}%
\begin{pgfscope}%
\definecolor{textcolor}{rgb}{0.000000,0.000000,0.000000}%
\pgfsetstrokecolor{textcolor}%
\pgfsetfillcolor{textcolor}%
\pgftext[x=1.768155in, y=0.592176in, left, base]{\color{textcolor}\sffamily\fontsize{7.500000}{9.000000}\selectfont Pooled}%
\end{pgfscope}%
\begin{pgfscope}%
\definecolor{textcolor}{rgb}{0.000000,0.000000,0.000000}%
\pgfsetstrokecolor{textcolor}%
\pgfsetfillcolor{textcolor}%
\pgftext[x=1.599196in, y=0.492254in, left, base]{\color{textcolor}\sffamily\fontsize{7.500000}{9.000000}\selectfont correct in time}%
\end{pgfscope}%
\begin{pgfscope}%
\pgfpathrectangle{\pgfqpoint{1.465200in}{0.822500in}}{\pgfqpoint{0.899100in}{2.537500in}}%
\pgfusepath{clip}%
\pgfsetbuttcap%
\pgfsetroundjoin%
\pgfsetlinewidth{1.003750pt}%
\definecolor{currentstroke}{rgb}{0.478431,0.305882,0.611765}%
\pgfsetstrokecolor{currentstroke}%
\pgfsetdash{}{0pt}%
\pgfpathmoveto{\pgfqpoint{1.727073in}{0.955056in}}%
\pgfpathlineto{\pgfqpoint{1.727073in}{0.955056in}}%
\pgfusepath{stroke}%
\end{pgfscope}%
\begin{pgfscope}%
\pgfpathrectangle{\pgfqpoint{1.465200in}{0.822500in}}{\pgfqpoint{0.899100in}{2.537500in}}%
\pgfusepath{clip}%
\pgfsetbuttcap%
\pgfsetroundjoin%
\pgfsetlinewidth{1.003750pt}%
\definecolor{currentstroke}{rgb}{0.478431,0.478431,0.478431}%
\pgfsetstrokecolor{currentstroke}%
\pgfsetdash{}{0pt}%
\pgfpathmoveto{\pgfqpoint{1.755954in}{1.144422in}}%
\pgfpathlineto{\pgfqpoint{1.762906in}{1.144422in}}%
\pgfusepath{stroke}%
\end{pgfscope}%
\begin{pgfscope}%
\pgfpathrectangle{\pgfqpoint{1.465200in}{0.822500in}}{\pgfqpoint{0.899100in}{2.537500in}}%
\pgfusepath{clip}%
\pgfsetbuttcap%
\pgfsetroundjoin%
\pgfsetlinewidth{1.003750pt}%
\definecolor{currentstroke}{rgb}{0.784314,0.784314,0.784314}%
\pgfsetstrokecolor{currentstroke}%
\pgfsetdash{}{0pt}%
\pgfpathmoveto{\pgfqpoint{1.758281in}{1.333787in}}%
\pgfpathlineto{\pgfqpoint{1.758281in}{1.333787in}}%
\pgfusepath{stroke}%
\end{pgfscope}%
\begin{pgfscope}%
\pgfpathrectangle{\pgfqpoint{1.465200in}{0.822500in}}{\pgfqpoint{0.899100in}{2.537500in}}%
\pgfusepath{clip}%
\pgfsetbuttcap%
\pgfsetroundjoin%
\pgfsetlinewidth{1.003750pt}%
\definecolor{currentstroke}{rgb}{0.478431,0.478431,0.478431}%
\pgfsetstrokecolor{currentstroke}%
\pgfsetdash{}{0pt}%
\pgfpathmoveto{\pgfqpoint{1.760636in}{1.523153in}}%
\pgfpathlineto{\pgfqpoint{1.760636in}{1.523153in}}%
\pgfusepath{stroke}%
\end{pgfscope}%
\begin{pgfscope}%
\pgfpathrectangle{\pgfqpoint{1.465200in}{0.822500in}}{\pgfqpoint{0.899100in}{2.537500in}}%
\pgfusepath{clip}%
\pgfsetbuttcap%
\pgfsetroundjoin%
\pgfsetlinewidth{1.003750pt}%
\definecolor{currentstroke}{rgb}{0.784314,0.784314,0.784314}%
\pgfsetstrokecolor{currentstroke}%
\pgfsetdash{}{0pt}%
\pgfpathmoveto{\pgfqpoint{1.762906in}{1.712519in}}%
\pgfpathlineto{\pgfqpoint{1.769886in}{1.712519in}}%
\pgfusepath{stroke}%
\end{pgfscope}%
\begin{pgfscope}%
\pgfpathrectangle{\pgfqpoint{1.465200in}{0.822500in}}{\pgfqpoint{0.899100in}{2.537500in}}%
\pgfusepath{clip}%
\pgfsetbuttcap%
\pgfsetroundjoin%
\pgfsetlinewidth{1.003750pt}%
\definecolor{currentstroke}{rgb}{0.784314,0.784314,0.784314}%
\pgfsetstrokecolor{currentstroke}%
\pgfsetdash{}{0pt}%
\pgfpathmoveto{\pgfqpoint{1.762934in}{1.901884in}}%
\pgfpathlineto{\pgfqpoint{1.765233in}{1.901884in}}%
\pgfusepath{stroke}%
\end{pgfscope}%
\begin{pgfscope}%
\pgfpathrectangle{\pgfqpoint{1.465200in}{0.822500in}}{\pgfqpoint{0.899100in}{2.537500in}}%
\pgfusepath{clip}%
\pgfsetbuttcap%
\pgfsetroundjoin%
\pgfsetlinewidth{1.003750pt}%
\definecolor{currentstroke}{rgb}{0.784314,0.784314,0.784314}%
\pgfsetstrokecolor{currentstroke}%
\pgfsetdash{}{0pt}%
\pgfpathmoveto{\pgfqpoint{1.774483in}{2.091250in}}%
\pgfpathlineto{\pgfqpoint{1.769858in}{2.091250in}}%
\pgfusepath{stroke}%
\end{pgfscope}%
\begin{pgfscope}%
\pgfpathrectangle{\pgfqpoint{1.465200in}{0.822500in}}{\pgfqpoint{0.899100in}{2.537500in}}%
\pgfusepath{clip}%
\pgfsetbuttcap%
\pgfsetroundjoin%
\pgfsetlinewidth{1.003750pt}%
\definecolor{currentstroke}{rgb}{0.784314,0.784314,0.784314}%
\pgfsetstrokecolor{currentstroke}%
\pgfsetdash{}{0pt}%
\pgfpathmoveto{\pgfqpoint{1.786893in}{2.280616in}}%
\pgfpathlineto{\pgfqpoint{1.775946in}{2.280616in}}%
\pgfusepath{stroke}%
\end{pgfscope}%
\begin{pgfscope}%
\pgfpathrectangle{\pgfqpoint{1.465200in}{0.822500in}}{\pgfqpoint{0.899100in}{2.537500in}}%
\pgfusepath{clip}%
\pgfsetbuttcap%
\pgfsetroundjoin%
\pgfsetlinewidth{1.003750pt}%
\definecolor{currentstroke}{rgb}{0.784314,0.784314,0.784314}%
\pgfsetstrokecolor{currentstroke}%
\pgfsetdash{}{0pt}%
\pgfpathmoveto{\pgfqpoint{1.787063in}{2.469981in}}%
\pgfpathlineto{\pgfqpoint{1.777316in}{2.469981in}}%
\pgfusepath{stroke}%
\end{pgfscope}%
\begin{pgfscope}%
\pgfpathrectangle{\pgfqpoint{1.465200in}{0.822500in}}{\pgfqpoint{0.899100in}{2.537500in}}%
\pgfusepath{clip}%
\pgfsetbuttcap%
\pgfsetroundjoin%
\pgfsetlinewidth{1.003750pt}%
\definecolor{currentstroke}{rgb}{0.478431,0.478431,0.478431}%
\pgfsetstrokecolor{currentstroke}%
\pgfsetdash{}{0pt}%
\pgfpathmoveto{\pgfqpoint{1.799754in}{2.659347in}}%
\pgfpathlineto{\pgfqpoint{1.773725in}{2.659347in}}%
\pgfusepath{stroke}%
\end{pgfscope}%
\begin{pgfscope}%
\pgfpathrectangle{\pgfqpoint{1.465200in}{0.822500in}}{\pgfqpoint{0.899100in}{2.537500in}}%
\pgfusepath{clip}%
\pgfsetbuttcap%
\pgfsetroundjoin%
\pgfsetlinewidth{1.003750pt}%
\definecolor{currentstroke}{rgb}{0.784314,0.784314,0.784314}%
\pgfsetstrokecolor{currentstroke}%
\pgfsetdash{}{0pt}%
\pgfpathmoveto{\pgfqpoint{1.805082in}{2.848713in}}%
\pgfpathlineto{\pgfqpoint{1.805068in}{2.848713in}}%
\pgfusepath{stroke}%
\end{pgfscope}%
\begin{pgfscope}%
\pgfpathrectangle{\pgfqpoint{1.465200in}{0.822500in}}{\pgfqpoint{0.899100in}{2.537500in}}%
\pgfusepath{clip}%
\pgfsetbuttcap%
\pgfsetroundjoin%
\pgfsetlinewidth{1.003750pt}%
\definecolor{currentstroke}{rgb}{0.156863,0.356863,0.568627}%
\pgfsetstrokecolor{currentstroke}%
\pgfsetdash{}{0pt}%
\pgfpathmoveto{\pgfqpoint{1.918147in}{3.038078in}}%
\pgfpathlineto{\pgfqpoint{1.977049in}{3.038078in}}%
\pgfusepath{stroke}%
\end{pgfscope}%
\begin{pgfscope}%
\pgfpathrectangle{\pgfqpoint{1.465200in}{0.822500in}}{\pgfqpoint{0.899100in}{2.537500in}}%
\pgfusepath{clip}%
\pgfsetbuttcap%
\pgfsetroundjoin%
\pgfsetlinewidth{1.003750pt}%
\definecolor{currentstroke}{rgb}{0.156863,0.356863,0.568627}%
\pgfsetstrokecolor{currentstroke}%
\pgfsetdash{}{0pt}%
\pgfpathmoveto{\pgfqpoint{1.960946in}{3.227444in}}%
\pgfpathlineto{\pgfqpoint{2.048282in}{3.227444in}}%
\pgfusepath{stroke}%
\end{pgfscope}%
\begin{pgfscope}%
\pgfsetrectcap%
\pgfsetmiterjoin%
\pgfsetlinewidth{0.803000pt}%
\definecolor{currentstroke}{rgb}{0.101961,0.101961,0.101961}%
\pgfsetstrokecolor{currentstroke}%
\pgfsetdash{}{0pt}%
\pgfpathmoveto{\pgfqpoint{1.465200in}{0.822500in}}%
\pgfpathlineto{\pgfqpoint{2.364300in}{0.822500in}}%
\pgfusepath{stroke}%
\end{pgfscope}%
\begin{pgfscope}%
\pgfpathrectangle{\pgfqpoint{1.465200in}{0.822500in}}{\pgfqpoint{0.899100in}{2.537500in}}%
\pgfusepath{clip}%
\pgfsetbuttcap%
\pgfsetroundjoin%
\definecolor{currentfill}{rgb}{0.478431,0.305882,0.611765}%
\pgfsetfillcolor{currentfill}%
\pgfsetlinewidth{0.000000pt}%
\definecolor{currentstroke}{rgb}{0.478431,0.305882,0.611765}%
\pgfsetstrokecolor{currentstroke}%
\pgfsetdash{}{0pt}%
\pgfsys@defobject{currentmarker}{\pgfqpoint{-0.025984in}{-0.025984in}}{\pgfqpoint{0.025984in}{0.025984in}}{%
\pgfpathmoveto{\pgfqpoint{0.000000in}{-0.025984in}}%
\pgfpathcurveto{\pgfqpoint{0.006891in}{-0.025984in}}{\pgfqpoint{0.013501in}{-0.023246in}}{\pgfqpoint{0.018373in}{-0.018373in}}%
\pgfpathcurveto{\pgfqpoint{0.023246in}{-0.013501in}}{\pgfqpoint{0.025984in}{-0.006891in}}{\pgfqpoint{0.025984in}{0.000000in}}%
\pgfpathcurveto{\pgfqpoint{0.025984in}{0.006891in}}{\pgfqpoint{0.023246in}{0.013501in}}{\pgfqpoint{0.018373in}{0.018373in}}%
\pgfpathcurveto{\pgfqpoint{0.013501in}{0.023246in}}{\pgfqpoint{0.006891in}{0.025984in}}{\pgfqpoint{0.000000in}{0.025984in}}%
\pgfpathcurveto{\pgfqpoint{-0.006891in}{0.025984in}}{\pgfqpoint{-0.013501in}{0.023246in}}{\pgfqpoint{-0.018373in}{0.018373in}}%
\pgfpathcurveto{\pgfqpoint{-0.023246in}{0.013501in}}{\pgfqpoint{-0.025984in}{0.006891in}}{\pgfqpoint{-0.025984in}{0.000000in}}%
\pgfpathcurveto{\pgfqpoint{-0.025984in}{-0.006891in}}{\pgfqpoint{-0.023246in}{-0.013501in}}{\pgfqpoint{-0.018373in}{-0.018373in}}%
\pgfpathcurveto{\pgfqpoint{-0.013501in}{-0.023246in}}{\pgfqpoint{-0.006891in}{-0.025984in}}{\pgfqpoint{0.000000in}{-0.025984in}}%
\pgfpathlineto{\pgfqpoint{0.000000in}{-0.025984in}}%
\pgfpathclose%
\pgfusepath{fill}%
}%
\begin{pgfscope}%
\pgfsys@transformshift{1.727073in}{0.955056in}%
\pgfsys@useobject{currentmarker}{}%
\end{pgfscope}%
\end{pgfscope}%
\begin{pgfscope}%
\pgfpathrectangle{\pgfqpoint{1.465200in}{0.822500in}}{\pgfqpoint{0.899100in}{2.537500in}}%
\pgfusepath{clip}%
\pgfsetbuttcap%
\pgfsetroundjoin%
\definecolor{currentfill}{rgb}{1.000000,1.000000,1.000000}%
\pgfsetfillcolor{currentfill}%
\pgfsetlinewidth{0.903375pt}%
\definecolor{currentstroke}{rgb}{0.478431,0.305882,0.611765}%
\pgfsetstrokecolor{currentstroke}%
\pgfsetdash{}{0pt}%
\pgfsys@defobject{currentmarker}{\pgfqpoint{-0.025984in}{-0.025984in}}{\pgfqpoint{0.025984in}{0.025984in}}{%
\pgfpathmoveto{\pgfqpoint{0.000000in}{-0.025984in}}%
\pgfpathcurveto{\pgfqpoint{0.006891in}{-0.025984in}}{\pgfqpoint{0.013501in}{-0.023246in}}{\pgfqpoint{0.018373in}{-0.018373in}}%
\pgfpathcurveto{\pgfqpoint{0.023246in}{-0.013501in}}{\pgfqpoint{0.025984in}{-0.006891in}}{\pgfqpoint{0.025984in}{0.000000in}}%
\pgfpathcurveto{\pgfqpoint{0.025984in}{0.006891in}}{\pgfqpoint{0.023246in}{0.013501in}}{\pgfqpoint{0.018373in}{0.018373in}}%
\pgfpathcurveto{\pgfqpoint{0.013501in}{0.023246in}}{\pgfqpoint{0.006891in}{0.025984in}}{\pgfqpoint{0.000000in}{0.025984in}}%
\pgfpathcurveto{\pgfqpoint{-0.006891in}{0.025984in}}{\pgfqpoint{-0.013501in}{0.023246in}}{\pgfqpoint{-0.018373in}{0.018373in}}%
\pgfpathcurveto{\pgfqpoint{-0.023246in}{0.013501in}}{\pgfqpoint{-0.025984in}{0.006891in}}{\pgfqpoint{-0.025984in}{0.000000in}}%
\pgfpathcurveto{\pgfqpoint{-0.025984in}{-0.006891in}}{\pgfqpoint{-0.023246in}{-0.013501in}}{\pgfqpoint{-0.018373in}{-0.018373in}}%
\pgfpathcurveto{\pgfqpoint{-0.013501in}{-0.023246in}}{\pgfqpoint{-0.006891in}{-0.025984in}}{\pgfqpoint{0.000000in}{-0.025984in}}%
\pgfpathlineto{\pgfqpoint{0.000000in}{-0.025984in}}%
\pgfpathclose%
\pgfusepath{stroke,fill}%
}%
\begin{pgfscope}%
\pgfsys@transformshift{1.727073in}{0.955056in}%
\pgfsys@useobject{currentmarker}{}%
\end{pgfscope}%
\end{pgfscope}%
\begin{pgfscope}%
\pgfpathrectangle{\pgfqpoint{1.465200in}{0.822500in}}{\pgfqpoint{0.899100in}{2.537500in}}%
\pgfusepath{clip}%
\pgfsetbuttcap%
\pgfsetroundjoin%
\definecolor{currentfill}{rgb}{0.478431,0.478431,0.478431}%
\pgfsetfillcolor{currentfill}%
\pgfsetlinewidth{0.000000pt}%
\definecolor{currentstroke}{rgb}{0.478431,0.478431,0.478431}%
\pgfsetstrokecolor{currentstroke}%
\pgfsetdash{}{0pt}%
\pgfsys@defobject{currentmarker}{\pgfqpoint{-0.025984in}{-0.025984in}}{\pgfqpoint{0.025984in}{0.025984in}}{%
\pgfpathmoveto{\pgfqpoint{0.000000in}{-0.025984in}}%
\pgfpathcurveto{\pgfqpoint{0.006891in}{-0.025984in}}{\pgfqpoint{0.013501in}{-0.023246in}}{\pgfqpoint{0.018373in}{-0.018373in}}%
\pgfpathcurveto{\pgfqpoint{0.023246in}{-0.013501in}}{\pgfqpoint{0.025984in}{-0.006891in}}{\pgfqpoint{0.025984in}{0.000000in}}%
\pgfpathcurveto{\pgfqpoint{0.025984in}{0.006891in}}{\pgfqpoint{0.023246in}{0.013501in}}{\pgfqpoint{0.018373in}{0.018373in}}%
\pgfpathcurveto{\pgfqpoint{0.013501in}{0.023246in}}{\pgfqpoint{0.006891in}{0.025984in}}{\pgfqpoint{0.000000in}{0.025984in}}%
\pgfpathcurveto{\pgfqpoint{-0.006891in}{0.025984in}}{\pgfqpoint{-0.013501in}{0.023246in}}{\pgfqpoint{-0.018373in}{0.018373in}}%
\pgfpathcurveto{\pgfqpoint{-0.023246in}{0.013501in}}{\pgfqpoint{-0.025984in}{0.006891in}}{\pgfqpoint{-0.025984in}{0.000000in}}%
\pgfpathcurveto{\pgfqpoint{-0.025984in}{-0.006891in}}{\pgfqpoint{-0.023246in}{-0.013501in}}{\pgfqpoint{-0.018373in}{-0.018373in}}%
\pgfpathcurveto{\pgfqpoint{-0.013501in}{-0.023246in}}{\pgfqpoint{-0.006891in}{-0.025984in}}{\pgfqpoint{0.000000in}{-0.025984in}}%
\pgfpathlineto{\pgfqpoint{0.000000in}{-0.025984in}}%
\pgfpathclose%
\pgfusepath{fill}%
}%
\begin{pgfscope}%
\pgfsys@transformshift{1.762906in}{1.144422in}%
\pgfsys@useobject{currentmarker}{}%
\end{pgfscope}%
\end{pgfscope}%
\begin{pgfscope}%
\pgfpathrectangle{\pgfqpoint{1.465200in}{0.822500in}}{\pgfqpoint{0.899100in}{2.537500in}}%
\pgfusepath{clip}%
\pgfsetbuttcap%
\pgfsetroundjoin%
\definecolor{currentfill}{rgb}{1.000000,1.000000,1.000000}%
\pgfsetfillcolor{currentfill}%
\pgfsetlinewidth{0.903375pt}%
\definecolor{currentstroke}{rgb}{0.478431,0.478431,0.478431}%
\pgfsetstrokecolor{currentstroke}%
\pgfsetdash{}{0pt}%
\pgfsys@defobject{currentmarker}{\pgfqpoint{-0.025984in}{-0.025984in}}{\pgfqpoint{0.025984in}{0.025984in}}{%
\pgfpathmoveto{\pgfqpoint{0.000000in}{-0.025984in}}%
\pgfpathcurveto{\pgfqpoint{0.006891in}{-0.025984in}}{\pgfqpoint{0.013501in}{-0.023246in}}{\pgfqpoint{0.018373in}{-0.018373in}}%
\pgfpathcurveto{\pgfqpoint{0.023246in}{-0.013501in}}{\pgfqpoint{0.025984in}{-0.006891in}}{\pgfqpoint{0.025984in}{0.000000in}}%
\pgfpathcurveto{\pgfqpoint{0.025984in}{0.006891in}}{\pgfqpoint{0.023246in}{0.013501in}}{\pgfqpoint{0.018373in}{0.018373in}}%
\pgfpathcurveto{\pgfqpoint{0.013501in}{0.023246in}}{\pgfqpoint{0.006891in}{0.025984in}}{\pgfqpoint{0.000000in}{0.025984in}}%
\pgfpathcurveto{\pgfqpoint{-0.006891in}{0.025984in}}{\pgfqpoint{-0.013501in}{0.023246in}}{\pgfqpoint{-0.018373in}{0.018373in}}%
\pgfpathcurveto{\pgfqpoint{-0.023246in}{0.013501in}}{\pgfqpoint{-0.025984in}{0.006891in}}{\pgfqpoint{-0.025984in}{0.000000in}}%
\pgfpathcurveto{\pgfqpoint{-0.025984in}{-0.006891in}}{\pgfqpoint{-0.023246in}{-0.013501in}}{\pgfqpoint{-0.018373in}{-0.018373in}}%
\pgfpathcurveto{\pgfqpoint{-0.013501in}{-0.023246in}}{\pgfqpoint{-0.006891in}{-0.025984in}}{\pgfqpoint{0.000000in}{-0.025984in}}%
\pgfpathlineto{\pgfqpoint{0.000000in}{-0.025984in}}%
\pgfpathclose%
\pgfusepath{stroke,fill}%
}%
\begin{pgfscope}%
\pgfsys@transformshift{1.755954in}{1.144422in}%
\pgfsys@useobject{currentmarker}{}%
\end{pgfscope}%
\end{pgfscope}%
\begin{pgfscope}%
\pgfpathrectangle{\pgfqpoint{1.465200in}{0.822500in}}{\pgfqpoint{0.899100in}{2.537500in}}%
\pgfusepath{clip}%
\pgfsetbuttcap%
\pgfsetroundjoin%
\definecolor{currentfill}{rgb}{0.784314,0.784314,0.784314}%
\pgfsetfillcolor{currentfill}%
\pgfsetlinewidth{0.000000pt}%
\definecolor{currentstroke}{rgb}{0.784314,0.784314,0.784314}%
\pgfsetstrokecolor{currentstroke}%
\pgfsetdash{}{0pt}%
\pgfsys@defobject{currentmarker}{\pgfqpoint{-0.025984in}{-0.025984in}}{\pgfqpoint{0.025984in}{0.025984in}}{%
\pgfpathmoveto{\pgfqpoint{0.000000in}{-0.025984in}}%
\pgfpathcurveto{\pgfqpoint{0.006891in}{-0.025984in}}{\pgfqpoint{0.013501in}{-0.023246in}}{\pgfqpoint{0.018373in}{-0.018373in}}%
\pgfpathcurveto{\pgfqpoint{0.023246in}{-0.013501in}}{\pgfqpoint{0.025984in}{-0.006891in}}{\pgfqpoint{0.025984in}{0.000000in}}%
\pgfpathcurveto{\pgfqpoint{0.025984in}{0.006891in}}{\pgfqpoint{0.023246in}{0.013501in}}{\pgfqpoint{0.018373in}{0.018373in}}%
\pgfpathcurveto{\pgfqpoint{0.013501in}{0.023246in}}{\pgfqpoint{0.006891in}{0.025984in}}{\pgfqpoint{0.000000in}{0.025984in}}%
\pgfpathcurveto{\pgfqpoint{-0.006891in}{0.025984in}}{\pgfqpoint{-0.013501in}{0.023246in}}{\pgfqpoint{-0.018373in}{0.018373in}}%
\pgfpathcurveto{\pgfqpoint{-0.023246in}{0.013501in}}{\pgfqpoint{-0.025984in}{0.006891in}}{\pgfqpoint{-0.025984in}{0.000000in}}%
\pgfpathcurveto{\pgfqpoint{-0.025984in}{-0.006891in}}{\pgfqpoint{-0.023246in}{-0.013501in}}{\pgfqpoint{-0.018373in}{-0.018373in}}%
\pgfpathcurveto{\pgfqpoint{-0.013501in}{-0.023246in}}{\pgfqpoint{-0.006891in}{-0.025984in}}{\pgfqpoint{0.000000in}{-0.025984in}}%
\pgfpathlineto{\pgfqpoint{0.000000in}{-0.025984in}}%
\pgfpathclose%
\pgfusepath{fill}%
}%
\begin{pgfscope}%
\pgfsys@transformshift{1.758281in}{1.333787in}%
\pgfsys@useobject{currentmarker}{}%
\end{pgfscope}%
\end{pgfscope}%
\begin{pgfscope}%
\pgfpathrectangle{\pgfqpoint{1.465200in}{0.822500in}}{\pgfqpoint{0.899100in}{2.537500in}}%
\pgfusepath{clip}%
\pgfsetbuttcap%
\pgfsetroundjoin%
\definecolor{currentfill}{rgb}{1.000000,1.000000,1.000000}%
\pgfsetfillcolor{currentfill}%
\pgfsetlinewidth{0.903375pt}%
\definecolor{currentstroke}{rgb}{0.784314,0.784314,0.784314}%
\pgfsetstrokecolor{currentstroke}%
\pgfsetdash{}{0pt}%
\pgfsys@defobject{currentmarker}{\pgfqpoint{-0.025984in}{-0.025984in}}{\pgfqpoint{0.025984in}{0.025984in}}{%
\pgfpathmoveto{\pgfqpoint{0.000000in}{-0.025984in}}%
\pgfpathcurveto{\pgfqpoint{0.006891in}{-0.025984in}}{\pgfqpoint{0.013501in}{-0.023246in}}{\pgfqpoint{0.018373in}{-0.018373in}}%
\pgfpathcurveto{\pgfqpoint{0.023246in}{-0.013501in}}{\pgfqpoint{0.025984in}{-0.006891in}}{\pgfqpoint{0.025984in}{0.000000in}}%
\pgfpathcurveto{\pgfqpoint{0.025984in}{0.006891in}}{\pgfqpoint{0.023246in}{0.013501in}}{\pgfqpoint{0.018373in}{0.018373in}}%
\pgfpathcurveto{\pgfqpoint{0.013501in}{0.023246in}}{\pgfqpoint{0.006891in}{0.025984in}}{\pgfqpoint{0.000000in}{0.025984in}}%
\pgfpathcurveto{\pgfqpoint{-0.006891in}{0.025984in}}{\pgfqpoint{-0.013501in}{0.023246in}}{\pgfqpoint{-0.018373in}{0.018373in}}%
\pgfpathcurveto{\pgfqpoint{-0.023246in}{0.013501in}}{\pgfqpoint{-0.025984in}{0.006891in}}{\pgfqpoint{-0.025984in}{0.000000in}}%
\pgfpathcurveto{\pgfqpoint{-0.025984in}{-0.006891in}}{\pgfqpoint{-0.023246in}{-0.013501in}}{\pgfqpoint{-0.018373in}{-0.018373in}}%
\pgfpathcurveto{\pgfqpoint{-0.013501in}{-0.023246in}}{\pgfqpoint{-0.006891in}{-0.025984in}}{\pgfqpoint{0.000000in}{-0.025984in}}%
\pgfpathlineto{\pgfqpoint{0.000000in}{-0.025984in}}%
\pgfpathclose%
\pgfusepath{stroke,fill}%
}%
\begin{pgfscope}%
\pgfsys@transformshift{1.758281in}{1.333787in}%
\pgfsys@useobject{currentmarker}{}%
\end{pgfscope}%
\end{pgfscope}%
\begin{pgfscope}%
\pgfpathrectangle{\pgfqpoint{1.465200in}{0.822500in}}{\pgfqpoint{0.899100in}{2.537500in}}%
\pgfusepath{clip}%
\pgfsetbuttcap%
\pgfsetroundjoin%
\definecolor{currentfill}{rgb}{0.478431,0.478431,0.478431}%
\pgfsetfillcolor{currentfill}%
\pgfsetlinewidth{0.000000pt}%
\definecolor{currentstroke}{rgb}{0.478431,0.478431,0.478431}%
\pgfsetstrokecolor{currentstroke}%
\pgfsetdash{}{0pt}%
\pgfsys@defobject{currentmarker}{\pgfqpoint{-0.025984in}{-0.025984in}}{\pgfqpoint{0.025984in}{0.025984in}}{%
\pgfpathmoveto{\pgfqpoint{0.000000in}{-0.025984in}}%
\pgfpathcurveto{\pgfqpoint{0.006891in}{-0.025984in}}{\pgfqpoint{0.013501in}{-0.023246in}}{\pgfqpoint{0.018373in}{-0.018373in}}%
\pgfpathcurveto{\pgfqpoint{0.023246in}{-0.013501in}}{\pgfqpoint{0.025984in}{-0.006891in}}{\pgfqpoint{0.025984in}{0.000000in}}%
\pgfpathcurveto{\pgfqpoint{0.025984in}{0.006891in}}{\pgfqpoint{0.023246in}{0.013501in}}{\pgfqpoint{0.018373in}{0.018373in}}%
\pgfpathcurveto{\pgfqpoint{0.013501in}{0.023246in}}{\pgfqpoint{0.006891in}{0.025984in}}{\pgfqpoint{0.000000in}{0.025984in}}%
\pgfpathcurveto{\pgfqpoint{-0.006891in}{0.025984in}}{\pgfqpoint{-0.013501in}{0.023246in}}{\pgfqpoint{-0.018373in}{0.018373in}}%
\pgfpathcurveto{\pgfqpoint{-0.023246in}{0.013501in}}{\pgfqpoint{-0.025984in}{0.006891in}}{\pgfqpoint{-0.025984in}{0.000000in}}%
\pgfpathcurveto{\pgfqpoint{-0.025984in}{-0.006891in}}{\pgfqpoint{-0.023246in}{-0.013501in}}{\pgfqpoint{-0.018373in}{-0.018373in}}%
\pgfpathcurveto{\pgfqpoint{-0.013501in}{-0.023246in}}{\pgfqpoint{-0.006891in}{-0.025984in}}{\pgfqpoint{0.000000in}{-0.025984in}}%
\pgfpathlineto{\pgfqpoint{0.000000in}{-0.025984in}}%
\pgfpathclose%
\pgfusepath{fill}%
}%
\begin{pgfscope}%
\pgfsys@transformshift{1.760636in}{1.523153in}%
\pgfsys@useobject{currentmarker}{}%
\end{pgfscope}%
\end{pgfscope}%
\begin{pgfscope}%
\pgfpathrectangle{\pgfqpoint{1.465200in}{0.822500in}}{\pgfqpoint{0.899100in}{2.537500in}}%
\pgfusepath{clip}%
\pgfsetbuttcap%
\pgfsetroundjoin%
\definecolor{currentfill}{rgb}{1.000000,1.000000,1.000000}%
\pgfsetfillcolor{currentfill}%
\pgfsetlinewidth{0.903375pt}%
\definecolor{currentstroke}{rgb}{0.478431,0.478431,0.478431}%
\pgfsetstrokecolor{currentstroke}%
\pgfsetdash{}{0pt}%
\pgfsys@defobject{currentmarker}{\pgfqpoint{-0.025984in}{-0.025984in}}{\pgfqpoint{0.025984in}{0.025984in}}{%
\pgfpathmoveto{\pgfqpoint{0.000000in}{-0.025984in}}%
\pgfpathcurveto{\pgfqpoint{0.006891in}{-0.025984in}}{\pgfqpoint{0.013501in}{-0.023246in}}{\pgfqpoint{0.018373in}{-0.018373in}}%
\pgfpathcurveto{\pgfqpoint{0.023246in}{-0.013501in}}{\pgfqpoint{0.025984in}{-0.006891in}}{\pgfqpoint{0.025984in}{0.000000in}}%
\pgfpathcurveto{\pgfqpoint{0.025984in}{0.006891in}}{\pgfqpoint{0.023246in}{0.013501in}}{\pgfqpoint{0.018373in}{0.018373in}}%
\pgfpathcurveto{\pgfqpoint{0.013501in}{0.023246in}}{\pgfqpoint{0.006891in}{0.025984in}}{\pgfqpoint{0.000000in}{0.025984in}}%
\pgfpathcurveto{\pgfqpoint{-0.006891in}{0.025984in}}{\pgfqpoint{-0.013501in}{0.023246in}}{\pgfqpoint{-0.018373in}{0.018373in}}%
\pgfpathcurveto{\pgfqpoint{-0.023246in}{0.013501in}}{\pgfqpoint{-0.025984in}{0.006891in}}{\pgfqpoint{-0.025984in}{0.000000in}}%
\pgfpathcurveto{\pgfqpoint{-0.025984in}{-0.006891in}}{\pgfqpoint{-0.023246in}{-0.013501in}}{\pgfqpoint{-0.018373in}{-0.018373in}}%
\pgfpathcurveto{\pgfqpoint{-0.013501in}{-0.023246in}}{\pgfqpoint{-0.006891in}{-0.025984in}}{\pgfqpoint{0.000000in}{-0.025984in}}%
\pgfpathlineto{\pgfqpoint{0.000000in}{-0.025984in}}%
\pgfpathclose%
\pgfusepath{stroke,fill}%
}%
\begin{pgfscope}%
\pgfsys@transformshift{1.760636in}{1.523153in}%
\pgfsys@useobject{currentmarker}{}%
\end{pgfscope}%
\end{pgfscope}%
\begin{pgfscope}%
\pgfpathrectangle{\pgfqpoint{1.465200in}{0.822500in}}{\pgfqpoint{0.899100in}{2.537500in}}%
\pgfusepath{clip}%
\pgfsetbuttcap%
\pgfsetroundjoin%
\definecolor{currentfill}{rgb}{0.784314,0.784314,0.784314}%
\pgfsetfillcolor{currentfill}%
\pgfsetlinewidth{0.000000pt}%
\definecolor{currentstroke}{rgb}{0.784314,0.784314,0.784314}%
\pgfsetstrokecolor{currentstroke}%
\pgfsetdash{}{0pt}%
\pgfsys@defobject{currentmarker}{\pgfqpoint{-0.025984in}{-0.025984in}}{\pgfqpoint{0.025984in}{0.025984in}}{%
\pgfpathmoveto{\pgfqpoint{0.000000in}{-0.025984in}}%
\pgfpathcurveto{\pgfqpoint{0.006891in}{-0.025984in}}{\pgfqpoint{0.013501in}{-0.023246in}}{\pgfqpoint{0.018373in}{-0.018373in}}%
\pgfpathcurveto{\pgfqpoint{0.023246in}{-0.013501in}}{\pgfqpoint{0.025984in}{-0.006891in}}{\pgfqpoint{0.025984in}{0.000000in}}%
\pgfpathcurveto{\pgfqpoint{0.025984in}{0.006891in}}{\pgfqpoint{0.023246in}{0.013501in}}{\pgfqpoint{0.018373in}{0.018373in}}%
\pgfpathcurveto{\pgfqpoint{0.013501in}{0.023246in}}{\pgfqpoint{0.006891in}{0.025984in}}{\pgfqpoint{0.000000in}{0.025984in}}%
\pgfpathcurveto{\pgfqpoint{-0.006891in}{0.025984in}}{\pgfqpoint{-0.013501in}{0.023246in}}{\pgfqpoint{-0.018373in}{0.018373in}}%
\pgfpathcurveto{\pgfqpoint{-0.023246in}{0.013501in}}{\pgfqpoint{-0.025984in}{0.006891in}}{\pgfqpoint{-0.025984in}{0.000000in}}%
\pgfpathcurveto{\pgfqpoint{-0.025984in}{-0.006891in}}{\pgfqpoint{-0.023246in}{-0.013501in}}{\pgfqpoint{-0.018373in}{-0.018373in}}%
\pgfpathcurveto{\pgfqpoint{-0.013501in}{-0.023246in}}{\pgfqpoint{-0.006891in}{-0.025984in}}{\pgfqpoint{0.000000in}{-0.025984in}}%
\pgfpathlineto{\pgfqpoint{0.000000in}{-0.025984in}}%
\pgfpathclose%
\pgfusepath{fill}%
}%
\begin{pgfscope}%
\pgfsys@transformshift{1.769886in}{1.712519in}%
\pgfsys@useobject{currentmarker}{}%
\end{pgfscope}%
\end{pgfscope}%
\begin{pgfscope}%
\pgfpathrectangle{\pgfqpoint{1.465200in}{0.822500in}}{\pgfqpoint{0.899100in}{2.537500in}}%
\pgfusepath{clip}%
\pgfsetbuttcap%
\pgfsetroundjoin%
\definecolor{currentfill}{rgb}{1.000000,1.000000,1.000000}%
\pgfsetfillcolor{currentfill}%
\pgfsetlinewidth{0.903375pt}%
\definecolor{currentstroke}{rgb}{0.784314,0.784314,0.784314}%
\pgfsetstrokecolor{currentstroke}%
\pgfsetdash{}{0pt}%
\pgfsys@defobject{currentmarker}{\pgfqpoint{-0.025984in}{-0.025984in}}{\pgfqpoint{0.025984in}{0.025984in}}{%
\pgfpathmoveto{\pgfqpoint{0.000000in}{-0.025984in}}%
\pgfpathcurveto{\pgfqpoint{0.006891in}{-0.025984in}}{\pgfqpoint{0.013501in}{-0.023246in}}{\pgfqpoint{0.018373in}{-0.018373in}}%
\pgfpathcurveto{\pgfqpoint{0.023246in}{-0.013501in}}{\pgfqpoint{0.025984in}{-0.006891in}}{\pgfqpoint{0.025984in}{0.000000in}}%
\pgfpathcurveto{\pgfqpoint{0.025984in}{0.006891in}}{\pgfqpoint{0.023246in}{0.013501in}}{\pgfqpoint{0.018373in}{0.018373in}}%
\pgfpathcurveto{\pgfqpoint{0.013501in}{0.023246in}}{\pgfqpoint{0.006891in}{0.025984in}}{\pgfqpoint{0.000000in}{0.025984in}}%
\pgfpathcurveto{\pgfqpoint{-0.006891in}{0.025984in}}{\pgfqpoint{-0.013501in}{0.023246in}}{\pgfqpoint{-0.018373in}{0.018373in}}%
\pgfpathcurveto{\pgfqpoint{-0.023246in}{0.013501in}}{\pgfqpoint{-0.025984in}{0.006891in}}{\pgfqpoint{-0.025984in}{0.000000in}}%
\pgfpathcurveto{\pgfqpoint{-0.025984in}{-0.006891in}}{\pgfqpoint{-0.023246in}{-0.013501in}}{\pgfqpoint{-0.018373in}{-0.018373in}}%
\pgfpathcurveto{\pgfqpoint{-0.013501in}{-0.023246in}}{\pgfqpoint{-0.006891in}{-0.025984in}}{\pgfqpoint{0.000000in}{-0.025984in}}%
\pgfpathlineto{\pgfqpoint{0.000000in}{-0.025984in}}%
\pgfpathclose%
\pgfusepath{stroke,fill}%
}%
\begin{pgfscope}%
\pgfsys@transformshift{1.762906in}{1.712519in}%
\pgfsys@useobject{currentmarker}{}%
\end{pgfscope}%
\end{pgfscope}%
\begin{pgfscope}%
\pgfpathrectangle{\pgfqpoint{1.465200in}{0.822500in}}{\pgfqpoint{0.899100in}{2.537500in}}%
\pgfusepath{clip}%
\pgfsetbuttcap%
\pgfsetroundjoin%
\definecolor{currentfill}{rgb}{0.784314,0.784314,0.784314}%
\pgfsetfillcolor{currentfill}%
\pgfsetlinewidth{0.000000pt}%
\definecolor{currentstroke}{rgb}{0.784314,0.784314,0.784314}%
\pgfsetstrokecolor{currentstroke}%
\pgfsetdash{}{0pt}%
\pgfsys@defobject{currentmarker}{\pgfqpoint{-0.025984in}{-0.025984in}}{\pgfqpoint{0.025984in}{0.025984in}}{%
\pgfpathmoveto{\pgfqpoint{0.000000in}{-0.025984in}}%
\pgfpathcurveto{\pgfqpoint{0.006891in}{-0.025984in}}{\pgfqpoint{0.013501in}{-0.023246in}}{\pgfqpoint{0.018373in}{-0.018373in}}%
\pgfpathcurveto{\pgfqpoint{0.023246in}{-0.013501in}}{\pgfqpoint{0.025984in}{-0.006891in}}{\pgfqpoint{0.025984in}{0.000000in}}%
\pgfpathcurveto{\pgfqpoint{0.025984in}{0.006891in}}{\pgfqpoint{0.023246in}{0.013501in}}{\pgfqpoint{0.018373in}{0.018373in}}%
\pgfpathcurveto{\pgfqpoint{0.013501in}{0.023246in}}{\pgfqpoint{0.006891in}{0.025984in}}{\pgfqpoint{0.000000in}{0.025984in}}%
\pgfpathcurveto{\pgfqpoint{-0.006891in}{0.025984in}}{\pgfqpoint{-0.013501in}{0.023246in}}{\pgfqpoint{-0.018373in}{0.018373in}}%
\pgfpathcurveto{\pgfqpoint{-0.023246in}{0.013501in}}{\pgfqpoint{-0.025984in}{0.006891in}}{\pgfqpoint{-0.025984in}{0.000000in}}%
\pgfpathcurveto{\pgfqpoint{-0.025984in}{-0.006891in}}{\pgfqpoint{-0.023246in}{-0.013501in}}{\pgfqpoint{-0.018373in}{-0.018373in}}%
\pgfpathcurveto{\pgfqpoint{-0.013501in}{-0.023246in}}{\pgfqpoint{-0.006891in}{-0.025984in}}{\pgfqpoint{0.000000in}{-0.025984in}}%
\pgfpathlineto{\pgfqpoint{0.000000in}{-0.025984in}}%
\pgfpathclose%
\pgfusepath{fill}%
}%
\begin{pgfscope}%
\pgfsys@transformshift{1.765233in}{1.901884in}%
\pgfsys@useobject{currentmarker}{}%
\end{pgfscope}%
\end{pgfscope}%
\begin{pgfscope}%
\pgfpathrectangle{\pgfqpoint{1.465200in}{0.822500in}}{\pgfqpoint{0.899100in}{2.537500in}}%
\pgfusepath{clip}%
\pgfsetbuttcap%
\pgfsetroundjoin%
\definecolor{currentfill}{rgb}{1.000000,1.000000,1.000000}%
\pgfsetfillcolor{currentfill}%
\pgfsetlinewidth{0.903375pt}%
\definecolor{currentstroke}{rgb}{0.784314,0.784314,0.784314}%
\pgfsetstrokecolor{currentstroke}%
\pgfsetdash{}{0pt}%
\pgfsys@defobject{currentmarker}{\pgfqpoint{-0.025984in}{-0.025984in}}{\pgfqpoint{0.025984in}{0.025984in}}{%
\pgfpathmoveto{\pgfqpoint{0.000000in}{-0.025984in}}%
\pgfpathcurveto{\pgfqpoint{0.006891in}{-0.025984in}}{\pgfqpoint{0.013501in}{-0.023246in}}{\pgfqpoint{0.018373in}{-0.018373in}}%
\pgfpathcurveto{\pgfqpoint{0.023246in}{-0.013501in}}{\pgfqpoint{0.025984in}{-0.006891in}}{\pgfqpoint{0.025984in}{0.000000in}}%
\pgfpathcurveto{\pgfqpoint{0.025984in}{0.006891in}}{\pgfqpoint{0.023246in}{0.013501in}}{\pgfqpoint{0.018373in}{0.018373in}}%
\pgfpathcurveto{\pgfqpoint{0.013501in}{0.023246in}}{\pgfqpoint{0.006891in}{0.025984in}}{\pgfqpoint{0.000000in}{0.025984in}}%
\pgfpathcurveto{\pgfqpoint{-0.006891in}{0.025984in}}{\pgfqpoint{-0.013501in}{0.023246in}}{\pgfqpoint{-0.018373in}{0.018373in}}%
\pgfpathcurveto{\pgfqpoint{-0.023246in}{0.013501in}}{\pgfqpoint{-0.025984in}{0.006891in}}{\pgfqpoint{-0.025984in}{0.000000in}}%
\pgfpathcurveto{\pgfqpoint{-0.025984in}{-0.006891in}}{\pgfqpoint{-0.023246in}{-0.013501in}}{\pgfqpoint{-0.018373in}{-0.018373in}}%
\pgfpathcurveto{\pgfqpoint{-0.013501in}{-0.023246in}}{\pgfqpoint{-0.006891in}{-0.025984in}}{\pgfqpoint{0.000000in}{-0.025984in}}%
\pgfpathlineto{\pgfqpoint{0.000000in}{-0.025984in}}%
\pgfpathclose%
\pgfusepath{stroke,fill}%
}%
\begin{pgfscope}%
\pgfsys@transformshift{1.762934in}{1.901884in}%
\pgfsys@useobject{currentmarker}{}%
\end{pgfscope}%
\end{pgfscope}%
\begin{pgfscope}%
\pgfpathrectangle{\pgfqpoint{1.465200in}{0.822500in}}{\pgfqpoint{0.899100in}{2.537500in}}%
\pgfusepath{clip}%
\pgfsetbuttcap%
\pgfsetroundjoin%
\definecolor{currentfill}{rgb}{0.784314,0.784314,0.784314}%
\pgfsetfillcolor{currentfill}%
\pgfsetlinewidth{0.000000pt}%
\definecolor{currentstroke}{rgb}{0.784314,0.784314,0.784314}%
\pgfsetstrokecolor{currentstroke}%
\pgfsetdash{}{0pt}%
\pgfsys@defobject{currentmarker}{\pgfqpoint{-0.025984in}{-0.025984in}}{\pgfqpoint{0.025984in}{0.025984in}}{%
\pgfpathmoveto{\pgfqpoint{0.000000in}{-0.025984in}}%
\pgfpathcurveto{\pgfqpoint{0.006891in}{-0.025984in}}{\pgfqpoint{0.013501in}{-0.023246in}}{\pgfqpoint{0.018373in}{-0.018373in}}%
\pgfpathcurveto{\pgfqpoint{0.023246in}{-0.013501in}}{\pgfqpoint{0.025984in}{-0.006891in}}{\pgfqpoint{0.025984in}{0.000000in}}%
\pgfpathcurveto{\pgfqpoint{0.025984in}{0.006891in}}{\pgfqpoint{0.023246in}{0.013501in}}{\pgfqpoint{0.018373in}{0.018373in}}%
\pgfpathcurveto{\pgfqpoint{0.013501in}{0.023246in}}{\pgfqpoint{0.006891in}{0.025984in}}{\pgfqpoint{0.000000in}{0.025984in}}%
\pgfpathcurveto{\pgfqpoint{-0.006891in}{0.025984in}}{\pgfqpoint{-0.013501in}{0.023246in}}{\pgfqpoint{-0.018373in}{0.018373in}}%
\pgfpathcurveto{\pgfqpoint{-0.023246in}{0.013501in}}{\pgfqpoint{-0.025984in}{0.006891in}}{\pgfqpoint{-0.025984in}{0.000000in}}%
\pgfpathcurveto{\pgfqpoint{-0.025984in}{-0.006891in}}{\pgfqpoint{-0.023246in}{-0.013501in}}{\pgfqpoint{-0.018373in}{-0.018373in}}%
\pgfpathcurveto{\pgfqpoint{-0.013501in}{-0.023246in}}{\pgfqpoint{-0.006891in}{-0.025984in}}{\pgfqpoint{0.000000in}{-0.025984in}}%
\pgfpathlineto{\pgfqpoint{0.000000in}{-0.025984in}}%
\pgfpathclose%
\pgfusepath{fill}%
}%
\begin{pgfscope}%
\pgfsys@transformshift{1.769858in}{2.091250in}%
\pgfsys@useobject{currentmarker}{}%
\end{pgfscope}%
\end{pgfscope}%
\begin{pgfscope}%
\pgfpathrectangle{\pgfqpoint{1.465200in}{0.822500in}}{\pgfqpoint{0.899100in}{2.537500in}}%
\pgfusepath{clip}%
\pgfsetbuttcap%
\pgfsetroundjoin%
\definecolor{currentfill}{rgb}{1.000000,1.000000,1.000000}%
\pgfsetfillcolor{currentfill}%
\pgfsetlinewidth{0.903375pt}%
\definecolor{currentstroke}{rgb}{0.784314,0.784314,0.784314}%
\pgfsetstrokecolor{currentstroke}%
\pgfsetdash{}{0pt}%
\pgfsys@defobject{currentmarker}{\pgfqpoint{-0.025984in}{-0.025984in}}{\pgfqpoint{0.025984in}{0.025984in}}{%
\pgfpathmoveto{\pgfqpoint{0.000000in}{-0.025984in}}%
\pgfpathcurveto{\pgfqpoint{0.006891in}{-0.025984in}}{\pgfqpoint{0.013501in}{-0.023246in}}{\pgfqpoint{0.018373in}{-0.018373in}}%
\pgfpathcurveto{\pgfqpoint{0.023246in}{-0.013501in}}{\pgfqpoint{0.025984in}{-0.006891in}}{\pgfqpoint{0.025984in}{0.000000in}}%
\pgfpathcurveto{\pgfqpoint{0.025984in}{0.006891in}}{\pgfqpoint{0.023246in}{0.013501in}}{\pgfqpoint{0.018373in}{0.018373in}}%
\pgfpathcurveto{\pgfqpoint{0.013501in}{0.023246in}}{\pgfqpoint{0.006891in}{0.025984in}}{\pgfqpoint{0.000000in}{0.025984in}}%
\pgfpathcurveto{\pgfqpoint{-0.006891in}{0.025984in}}{\pgfqpoint{-0.013501in}{0.023246in}}{\pgfqpoint{-0.018373in}{0.018373in}}%
\pgfpathcurveto{\pgfqpoint{-0.023246in}{0.013501in}}{\pgfqpoint{-0.025984in}{0.006891in}}{\pgfqpoint{-0.025984in}{0.000000in}}%
\pgfpathcurveto{\pgfqpoint{-0.025984in}{-0.006891in}}{\pgfqpoint{-0.023246in}{-0.013501in}}{\pgfqpoint{-0.018373in}{-0.018373in}}%
\pgfpathcurveto{\pgfqpoint{-0.013501in}{-0.023246in}}{\pgfqpoint{-0.006891in}{-0.025984in}}{\pgfqpoint{0.000000in}{-0.025984in}}%
\pgfpathlineto{\pgfqpoint{0.000000in}{-0.025984in}}%
\pgfpathclose%
\pgfusepath{stroke,fill}%
}%
\begin{pgfscope}%
\pgfsys@transformshift{1.774483in}{2.091250in}%
\pgfsys@useobject{currentmarker}{}%
\end{pgfscope}%
\end{pgfscope}%
\begin{pgfscope}%
\pgfpathrectangle{\pgfqpoint{1.465200in}{0.822500in}}{\pgfqpoint{0.899100in}{2.537500in}}%
\pgfusepath{clip}%
\pgfsetbuttcap%
\pgfsetroundjoin%
\definecolor{currentfill}{rgb}{0.784314,0.784314,0.784314}%
\pgfsetfillcolor{currentfill}%
\pgfsetlinewidth{0.000000pt}%
\definecolor{currentstroke}{rgb}{0.784314,0.784314,0.784314}%
\pgfsetstrokecolor{currentstroke}%
\pgfsetdash{}{0pt}%
\pgfsys@defobject{currentmarker}{\pgfqpoint{-0.025984in}{-0.025984in}}{\pgfqpoint{0.025984in}{0.025984in}}{%
\pgfpathmoveto{\pgfqpoint{0.000000in}{-0.025984in}}%
\pgfpathcurveto{\pgfqpoint{0.006891in}{-0.025984in}}{\pgfqpoint{0.013501in}{-0.023246in}}{\pgfqpoint{0.018373in}{-0.018373in}}%
\pgfpathcurveto{\pgfqpoint{0.023246in}{-0.013501in}}{\pgfqpoint{0.025984in}{-0.006891in}}{\pgfqpoint{0.025984in}{0.000000in}}%
\pgfpathcurveto{\pgfqpoint{0.025984in}{0.006891in}}{\pgfqpoint{0.023246in}{0.013501in}}{\pgfqpoint{0.018373in}{0.018373in}}%
\pgfpathcurveto{\pgfqpoint{0.013501in}{0.023246in}}{\pgfqpoint{0.006891in}{0.025984in}}{\pgfqpoint{0.000000in}{0.025984in}}%
\pgfpathcurveto{\pgfqpoint{-0.006891in}{0.025984in}}{\pgfqpoint{-0.013501in}{0.023246in}}{\pgfqpoint{-0.018373in}{0.018373in}}%
\pgfpathcurveto{\pgfqpoint{-0.023246in}{0.013501in}}{\pgfqpoint{-0.025984in}{0.006891in}}{\pgfqpoint{-0.025984in}{0.000000in}}%
\pgfpathcurveto{\pgfqpoint{-0.025984in}{-0.006891in}}{\pgfqpoint{-0.023246in}{-0.013501in}}{\pgfqpoint{-0.018373in}{-0.018373in}}%
\pgfpathcurveto{\pgfqpoint{-0.013501in}{-0.023246in}}{\pgfqpoint{-0.006891in}{-0.025984in}}{\pgfqpoint{0.000000in}{-0.025984in}}%
\pgfpathlineto{\pgfqpoint{0.000000in}{-0.025984in}}%
\pgfpathclose%
\pgfusepath{fill}%
}%
\begin{pgfscope}%
\pgfsys@transformshift{1.775946in}{2.280616in}%
\pgfsys@useobject{currentmarker}{}%
\end{pgfscope}%
\end{pgfscope}%
\begin{pgfscope}%
\pgfpathrectangle{\pgfqpoint{1.465200in}{0.822500in}}{\pgfqpoint{0.899100in}{2.537500in}}%
\pgfusepath{clip}%
\pgfsetbuttcap%
\pgfsetroundjoin%
\definecolor{currentfill}{rgb}{1.000000,1.000000,1.000000}%
\pgfsetfillcolor{currentfill}%
\pgfsetlinewidth{0.903375pt}%
\definecolor{currentstroke}{rgb}{0.784314,0.784314,0.784314}%
\pgfsetstrokecolor{currentstroke}%
\pgfsetdash{}{0pt}%
\pgfsys@defobject{currentmarker}{\pgfqpoint{-0.025984in}{-0.025984in}}{\pgfqpoint{0.025984in}{0.025984in}}{%
\pgfpathmoveto{\pgfqpoint{0.000000in}{-0.025984in}}%
\pgfpathcurveto{\pgfqpoint{0.006891in}{-0.025984in}}{\pgfqpoint{0.013501in}{-0.023246in}}{\pgfqpoint{0.018373in}{-0.018373in}}%
\pgfpathcurveto{\pgfqpoint{0.023246in}{-0.013501in}}{\pgfqpoint{0.025984in}{-0.006891in}}{\pgfqpoint{0.025984in}{0.000000in}}%
\pgfpathcurveto{\pgfqpoint{0.025984in}{0.006891in}}{\pgfqpoint{0.023246in}{0.013501in}}{\pgfqpoint{0.018373in}{0.018373in}}%
\pgfpathcurveto{\pgfqpoint{0.013501in}{0.023246in}}{\pgfqpoint{0.006891in}{0.025984in}}{\pgfqpoint{0.000000in}{0.025984in}}%
\pgfpathcurveto{\pgfqpoint{-0.006891in}{0.025984in}}{\pgfqpoint{-0.013501in}{0.023246in}}{\pgfqpoint{-0.018373in}{0.018373in}}%
\pgfpathcurveto{\pgfqpoint{-0.023246in}{0.013501in}}{\pgfqpoint{-0.025984in}{0.006891in}}{\pgfqpoint{-0.025984in}{0.000000in}}%
\pgfpathcurveto{\pgfqpoint{-0.025984in}{-0.006891in}}{\pgfqpoint{-0.023246in}{-0.013501in}}{\pgfqpoint{-0.018373in}{-0.018373in}}%
\pgfpathcurveto{\pgfqpoint{-0.013501in}{-0.023246in}}{\pgfqpoint{-0.006891in}{-0.025984in}}{\pgfqpoint{0.000000in}{-0.025984in}}%
\pgfpathlineto{\pgfqpoint{0.000000in}{-0.025984in}}%
\pgfpathclose%
\pgfusepath{stroke,fill}%
}%
\begin{pgfscope}%
\pgfsys@transformshift{1.786893in}{2.280616in}%
\pgfsys@useobject{currentmarker}{}%
\end{pgfscope}%
\end{pgfscope}%
\begin{pgfscope}%
\pgfpathrectangle{\pgfqpoint{1.465200in}{0.822500in}}{\pgfqpoint{0.899100in}{2.537500in}}%
\pgfusepath{clip}%
\pgfsetbuttcap%
\pgfsetroundjoin%
\definecolor{currentfill}{rgb}{0.784314,0.784314,0.784314}%
\pgfsetfillcolor{currentfill}%
\pgfsetlinewidth{0.000000pt}%
\definecolor{currentstroke}{rgb}{0.784314,0.784314,0.784314}%
\pgfsetstrokecolor{currentstroke}%
\pgfsetdash{}{0pt}%
\pgfsys@defobject{currentmarker}{\pgfqpoint{-0.025984in}{-0.025984in}}{\pgfqpoint{0.025984in}{0.025984in}}{%
\pgfpathmoveto{\pgfqpoint{0.000000in}{-0.025984in}}%
\pgfpathcurveto{\pgfqpoint{0.006891in}{-0.025984in}}{\pgfqpoint{0.013501in}{-0.023246in}}{\pgfqpoint{0.018373in}{-0.018373in}}%
\pgfpathcurveto{\pgfqpoint{0.023246in}{-0.013501in}}{\pgfqpoint{0.025984in}{-0.006891in}}{\pgfqpoint{0.025984in}{0.000000in}}%
\pgfpathcurveto{\pgfqpoint{0.025984in}{0.006891in}}{\pgfqpoint{0.023246in}{0.013501in}}{\pgfqpoint{0.018373in}{0.018373in}}%
\pgfpathcurveto{\pgfqpoint{0.013501in}{0.023246in}}{\pgfqpoint{0.006891in}{0.025984in}}{\pgfqpoint{0.000000in}{0.025984in}}%
\pgfpathcurveto{\pgfqpoint{-0.006891in}{0.025984in}}{\pgfqpoint{-0.013501in}{0.023246in}}{\pgfqpoint{-0.018373in}{0.018373in}}%
\pgfpathcurveto{\pgfqpoint{-0.023246in}{0.013501in}}{\pgfqpoint{-0.025984in}{0.006891in}}{\pgfqpoint{-0.025984in}{0.000000in}}%
\pgfpathcurveto{\pgfqpoint{-0.025984in}{-0.006891in}}{\pgfqpoint{-0.023246in}{-0.013501in}}{\pgfqpoint{-0.018373in}{-0.018373in}}%
\pgfpathcurveto{\pgfqpoint{-0.013501in}{-0.023246in}}{\pgfqpoint{-0.006891in}{-0.025984in}}{\pgfqpoint{0.000000in}{-0.025984in}}%
\pgfpathlineto{\pgfqpoint{0.000000in}{-0.025984in}}%
\pgfpathclose%
\pgfusepath{fill}%
}%
\begin{pgfscope}%
\pgfsys@transformshift{1.777316in}{2.469981in}%
\pgfsys@useobject{currentmarker}{}%
\end{pgfscope}%
\end{pgfscope}%
\begin{pgfscope}%
\pgfpathrectangle{\pgfqpoint{1.465200in}{0.822500in}}{\pgfqpoint{0.899100in}{2.537500in}}%
\pgfusepath{clip}%
\pgfsetbuttcap%
\pgfsetroundjoin%
\definecolor{currentfill}{rgb}{1.000000,1.000000,1.000000}%
\pgfsetfillcolor{currentfill}%
\pgfsetlinewidth{0.903375pt}%
\definecolor{currentstroke}{rgb}{0.784314,0.784314,0.784314}%
\pgfsetstrokecolor{currentstroke}%
\pgfsetdash{}{0pt}%
\pgfsys@defobject{currentmarker}{\pgfqpoint{-0.025984in}{-0.025984in}}{\pgfqpoint{0.025984in}{0.025984in}}{%
\pgfpathmoveto{\pgfqpoint{0.000000in}{-0.025984in}}%
\pgfpathcurveto{\pgfqpoint{0.006891in}{-0.025984in}}{\pgfqpoint{0.013501in}{-0.023246in}}{\pgfqpoint{0.018373in}{-0.018373in}}%
\pgfpathcurveto{\pgfqpoint{0.023246in}{-0.013501in}}{\pgfqpoint{0.025984in}{-0.006891in}}{\pgfqpoint{0.025984in}{0.000000in}}%
\pgfpathcurveto{\pgfqpoint{0.025984in}{0.006891in}}{\pgfqpoint{0.023246in}{0.013501in}}{\pgfqpoint{0.018373in}{0.018373in}}%
\pgfpathcurveto{\pgfqpoint{0.013501in}{0.023246in}}{\pgfqpoint{0.006891in}{0.025984in}}{\pgfqpoint{0.000000in}{0.025984in}}%
\pgfpathcurveto{\pgfqpoint{-0.006891in}{0.025984in}}{\pgfqpoint{-0.013501in}{0.023246in}}{\pgfqpoint{-0.018373in}{0.018373in}}%
\pgfpathcurveto{\pgfqpoint{-0.023246in}{0.013501in}}{\pgfqpoint{-0.025984in}{0.006891in}}{\pgfqpoint{-0.025984in}{0.000000in}}%
\pgfpathcurveto{\pgfqpoint{-0.025984in}{-0.006891in}}{\pgfqpoint{-0.023246in}{-0.013501in}}{\pgfqpoint{-0.018373in}{-0.018373in}}%
\pgfpathcurveto{\pgfqpoint{-0.013501in}{-0.023246in}}{\pgfqpoint{-0.006891in}{-0.025984in}}{\pgfqpoint{0.000000in}{-0.025984in}}%
\pgfpathlineto{\pgfqpoint{0.000000in}{-0.025984in}}%
\pgfpathclose%
\pgfusepath{stroke,fill}%
}%
\begin{pgfscope}%
\pgfsys@transformshift{1.787063in}{2.469981in}%
\pgfsys@useobject{currentmarker}{}%
\end{pgfscope}%
\end{pgfscope}%
\begin{pgfscope}%
\pgfpathrectangle{\pgfqpoint{1.465200in}{0.822500in}}{\pgfqpoint{0.899100in}{2.537500in}}%
\pgfusepath{clip}%
\pgfsetbuttcap%
\pgfsetroundjoin%
\definecolor{currentfill}{rgb}{0.478431,0.478431,0.478431}%
\pgfsetfillcolor{currentfill}%
\pgfsetlinewidth{0.000000pt}%
\definecolor{currentstroke}{rgb}{0.478431,0.478431,0.478431}%
\pgfsetstrokecolor{currentstroke}%
\pgfsetdash{}{0pt}%
\pgfsys@defobject{currentmarker}{\pgfqpoint{-0.025984in}{-0.025984in}}{\pgfqpoint{0.025984in}{0.025984in}}{%
\pgfpathmoveto{\pgfqpoint{0.000000in}{-0.025984in}}%
\pgfpathcurveto{\pgfqpoint{0.006891in}{-0.025984in}}{\pgfqpoint{0.013501in}{-0.023246in}}{\pgfqpoint{0.018373in}{-0.018373in}}%
\pgfpathcurveto{\pgfqpoint{0.023246in}{-0.013501in}}{\pgfqpoint{0.025984in}{-0.006891in}}{\pgfqpoint{0.025984in}{0.000000in}}%
\pgfpathcurveto{\pgfqpoint{0.025984in}{0.006891in}}{\pgfqpoint{0.023246in}{0.013501in}}{\pgfqpoint{0.018373in}{0.018373in}}%
\pgfpathcurveto{\pgfqpoint{0.013501in}{0.023246in}}{\pgfqpoint{0.006891in}{0.025984in}}{\pgfqpoint{0.000000in}{0.025984in}}%
\pgfpathcurveto{\pgfqpoint{-0.006891in}{0.025984in}}{\pgfqpoint{-0.013501in}{0.023246in}}{\pgfqpoint{-0.018373in}{0.018373in}}%
\pgfpathcurveto{\pgfqpoint{-0.023246in}{0.013501in}}{\pgfqpoint{-0.025984in}{0.006891in}}{\pgfqpoint{-0.025984in}{0.000000in}}%
\pgfpathcurveto{\pgfqpoint{-0.025984in}{-0.006891in}}{\pgfqpoint{-0.023246in}{-0.013501in}}{\pgfqpoint{-0.018373in}{-0.018373in}}%
\pgfpathcurveto{\pgfqpoint{-0.013501in}{-0.023246in}}{\pgfqpoint{-0.006891in}{-0.025984in}}{\pgfqpoint{0.000000in}{-0.025984in}}%
\pgfpathlineto{\pgfqpoint{0.000000in}{-0.025984in}}%
\pgfpathclose%
\pgfusepath{fill}%
}%
\begin{pgfscope}%
\pgfsys@transformshift{1.773725in}{2.659347in}%
\pgfsys@useobject{currentmarker}{}%
\end{pgfscope}%
\end{pgfscope}%
\begin{pgfscope}%
\pgfpathrectangle{\pgfqpoint{1.465200in}{0.822500in}}{\pgfqpoint{0.899100in}{2.537500in}}%
\pgfusepath{clip}%
\pgfsetbuttcap%
\pgfsetroundjoin%
\definecolor{currentfill}{rgb}{1.000000,1.000000,1.000000}%
\pgfsetfillcolor{currentfill}%
\pgfsetlinewidth{0.903375pt}%
\definecolor{currentstroke}{rgb}{0.478431,0.478431,0.478431}%
\pgfsetstrokecolor{currentstroke}%
\pgfsetdash{}{0pt}%
\pgfsys@defobject{currentmarker}{\pgfqpoint{-0.025984in}{-0.025984in}}{\pgfqpoint{0.025984in}{0.025984in}}{%
\pgfpathmoveto{\pgfqpoint{0.000000in}{-0.025984in}}%
\pgfpathcurveto{\pgfqpoint{0.006891in}{-0.025984in}}{\pgfqpoint{0.013501in}{-0.023246in}}{\pgfqpoint{0.018373in}{-0.018373in}}%
\pgfpathcurveto{\pgfqpoint{0.023246in}{-0.013501in}}{\pgfqpoint{0.025984in}{-0.006891in}}{\pgfqpoint{0.025984in}{0.000000in}}%
\pgfpathcurveto{\pgfqpoint{0.025984in}{0.006891in}}{\pgfqpoint{0.023246in}{0.013501in}}{\pgfqpoint{0.018373in}{0.018373in}}%
\pgfpathcurveto{\pgfqpoint{0.013501in}{0.023246in}}{\pgfqpoint{0.006891in}{0.025984in}}{\pgfqpoint{0.000000in}{0.025984in}}%
\pgfpathcurveto{\pgfqpoint{-0.006891in}{0.025984in}}{\pgfqpoint{-0.013501in}{0.023246in}}{\pgfqpoint{-0.018373in}{0.018373in}}%
\pgfpathcurveto{\pgfqpoint{-0.023246in}{0.013501in}}{\pgfqpoint{-0.025984in}{0.006891in}}{\pgfqpoint{-0.025984in}{0.000000in}}%
\pgfpathcurveto{\pgfqpoint{-0.025984in}{-0.006891in}}{\pgfqpoint{-0.023246in}{-0.013501in}}{\pgfqpoint{-0.018373in}{-0.018373in}}%
\pgfpathcurveto{\pgfqpoint{-0.013501in}{-0.023246in}}{\pgfqpoint{-0.006891in}{-0.025984in}}{\pgfqpoint{0.000000in}{-0.025984in}}%
\pgfpathlineto{\pgfqpoint{0.000000in}{-0.025984in}}%
\pgfpathclose%
\pgfusepath{stroke,fill}%
}%
\begin{pgfscope}%
\pgfsys@transformshift{1.799754in}{2.659347in}%
\pgfsys@useobject{currentmarker}{}%
\end{pgfscope}%
\end{pgfscope}%
\begin{pgfscope}%
\pgfpathrectangle{\pgfqpoint{1.465200in}{0.822500in}}{\pgfqpoint{0.899100in}{2.537500in}}%
\pgfusepath{clip}%
\pgfsetbuttcap%
\pgfsetroundjoin%
\definecolor{currentfill}{rgb}{0.784314,0.784314,0.784314}%
\pgfsetfillcolor{currentfill}%
\pgfsetlinewidth{0.000000pt}%
\definecolor{currentstroke}{rgb}{0.784314,0.784314,0.784314}%
\pgfsetstrokecolor{currentstroke}%
\pgfsetdash{}{0pt}%
\pgfsys@defobject{currentmarker}{\pgfqpoint{-0.025984in}{-0.025984in}}{\pgfqpoint{0.025984in}{0.025984in}}{%
\pgfpathmoveto{\pgfqpoint{0.000000in}{-0.025984in}}%
\pgfpathcurveto{\pgfqpoint{0.006891in}{-0.025984in}}{\pgfqpoint{0.013501in}{-0.023246in}}{\pgfqpoint{0.018373in}{-0.018373in}}%
\pgfpathcurveto{\pgfqpoint{0.023246in}{-0.013501in}}{\pgfqpoint{0.025984in}{-0.006891in}}{\pgfqpoint{0.025984in}{0.000000in}}%
\pgfpathcurveto{\pgfqpoint{0.025984in}{0.006891in}}{\pgfqpoint{0.023246in}{0.013501in}}{\pgfqpoint{0.018373in}{0.018373in}}%
\pgfpathcurveto{\pgfqpoint{0.013501in}{0.023246in}}{\pgfqpoint{0.006891in}{0.025984in}}{\pgfqpoint{0.000000in}{0.025984in}}%
\pgfpathcurveto{\pgfqpoint{-0.006891in}{0.025984in}}{\pgfqpoint{-0.013501in}{0.023246in}}{\pgfqpoint{-0.018373in}{0.018373in}}%
\pgfpathcurveto{\pgfqpoint{-0.023246in}{0.013501in}}{\pgfqpoint{-0.025984in}{0.006891in}}{\pgfqpoint{-0.025984in}{0.000000in}}%
\pgfpathcurveto{\pgfqpoint{-0.025984in}{-0.006891in}}{\pgfqpoint{-0.023246in}{-0.013501in}}{\pgfqpoint{-0.018373in}{-0.018373in}}%
\pgfpathcurveto{\pgfqpoint{-0.013501in}{-0.023246in}}{\pgfqpoint{-0.006891in}{-0.025984in}}{\pgfqpoint{0.000000in}{-0.025984in}}%
\pgfpathlineto{\pgfqpoint{0.000000in}{-0.025984in}}%
\pgfpathclose%
\pgfusepath{fill}%
}%
\begin{pgfscope}%
\pgfsys@transformshift{1.805068in}{2.848713in}%
\pgfsys@useobject{currentmarker}{}%
\end{pgfscope}%
\end{pgfscope}%
\begin{pgfscope}%
\pgfpathrectangle{\pgfqpoint{1.465200in}{0.822500in}}{\pgfqpoint{0.899100in}{2.537500in}}%
\pgfusepath{clip}%
\pgfsetbuttcap%
\pgfsetroundjoin%
\definecolor{currentfill}{rgb}{1.000000,1.000000,1.000000}%
\pgfsetfillcolor{currentfill}%
\pgfsetlinewidth{0.903375pt}%
\definecolor{currentstroke}{rgb}{0.784314,0.784314,0.784314}%
\pgfsetstrokecolor{currentstroke}%
\pgfsetdash{}{0pt}%
\pgfsys@defobject{currentmarker}{\pgfqpoint{-0.025984in}{-0.025984in}}{\pgfqpoint{0.025984in}{0.025984in}}{%
\pgfpathmoveto{\pgfqpoint{0.000000in}{-0.025984in}}%
\pgfpathcurveto{\pgfqpoint{0.006891in}{-0.025984in}}{\pgfqpoint{0.013501in}{-0.023246in}}{\pgfqpoint{0.018373in}{-0.018373in}}%
\pgfpathcurveto{\pgfqpoint{0.023246in}{-0.013501in}}{\pgfqpoint{0.025984in}{-0.006891in}}{\pgfqpoint{0.025984in}{0.000000in}}%
\pgfpathcurveto{\pgfqpoint{0.025984in}{0.006891in}}{\pgfqpoint{0.023246in}{0.013501in}}{\pgfqpoint{0.018373in}{0.018373in}}%
\pgfpathcurveto{\pgfqpoint{0.013501in}{0.023246in}}{\pgfqpoint{0.006891in}{0.025984in}}{\pgfqpoint{0.000000in}{0.025984in}}%
\pgfpathcurveto{\pgfqpoint{-0.006891in}{0.025984in}}{\pgfqpoint{-0.013501in}{0.023246in}}{\pgfqpoint{-0.018373in}{0.018373in}}%
\pgfpathcurveto{\pgfqpoint{-0.023246in}{0.013501in}}{\pgfqpoint{-0.025984in}{0.006891in}}{\pgfqpoint{-0.025984in}{0.000000in}}%
\pgfpathcurveto{\pgfqpoint{-0.025984in}{-0.006891in}}{\pgfqpoint{-0.023246in}{-0.013501in}}{\pgfqpoint{-0.018373in}{-0.018373in}}%
\pgfpathcurveto{\pgfqpoint{-0.013501in}{-0.023246in}}{\pgfqpoint{-0.006891in}{-0.025984in}}{\pgfqpoint{0.000000in}{-0.025984in}}%
\pgfpathlineto{\pgfqpoint{0.000000in}{-0.025984in}}%
\pgfpathclose%
\pgfusepath{stroke,fill}%
}%
\begin{pgfscope}%
\pgfsys@transformshift{1.805082in}{2.848713in}%
\pgfsys@useobject{currentmarker}{}%
\end{pgfscope}%
\end{pgfscope}%
\begin{pgfscope}%
\pgfpathrectangle{\pgfqpoint{1.465200in}{0.822500in}}{\pgfqpoint{0.899100in}{2.537500in}}%
\pgfusepath{clip}%
\pgfsetbuttcap%
\pgfsetroundjoin%
\definecolor{currentfill}{rgb}{0.156863,0.356863,0.568627}%
\pgfsetfillcolor{currentfill}%
\pgfsetlinewidth{0.000000pt}%
\definecolor{currentstroke}{rgb}{0.156863,0.356863,0.568627}%
\pgfsetstrokecolor{currentstroke}%
\pgfsetdash{}{0pt}%
\pgfsys@defobject{currentmarker}{\pgfqpoint{-0.025984in}{-0.025984in}}{\pgfqpoint{0.025984in}{0.025984in}}{%
\pgfpathmoveto{\pgfqpoint{0.000000in}{-0.025984in}}%
\pgfpathcurveto{\pgfqpoint{0.006891in}{-0.025984in}}{\pgfqpoint{0.013501in}{-0.023246in}}{\pgfqpoint{0.018373in}{-0.018373in}}%
\pgfpathcurveto{\pgfqpoint{0.023246in}{-0.013501in}}{\pgfqpoint{0.025984in}{-0.006891in}}{\pgfqpoint{0.025984in}{0.000000in}}%
\pgfpathcurveto{\pgfqpoint{0.025984in}{0.006891in}}{\pgfqpoint{0.023246in}{0.013501in}}{\pgfqpoint{0.018373in}{0.018373in}}%
\pgfpathcurveto{\pgfqpoint{0.013501in}{0.023246in}}{\pgfqpoint{0.006891in}{0.025984in}}{\pgfqpoint{0.000000in}{0.025984in}}%
\pgfpathcurveto{\pgfqpoint{-0.006891in}{0.025984in}}{\pgfqpoint{-0.013501in}{0.023246in}}{\pgfqpoint{-0.018373in}{0.018373in}}%
\pgfpathcurveto{\pgfqpoint{-0.023246in}{0.013501in}}{\pgfqpoint{-0.025984in}{0.006891in}}{\pgfqpoint{-0.025984in}{0.000000in}}%
\pgfpathcurveto{\pgfqpoint{-0.025984in}{-0.006891in}}{\pgfqpoint{-0.023246in}{-0.013501in}}{\pgfqpoint{-0.018373in}{-0.018373in}}%
\pgfpathcurveto{\pgfqpoint{-0.013501in}{-0.023246in}}{\pgfqpoint{-0.006891in}{-0.025984in}}{\pgfqpoint{0.000000in}{-0.025984in}}%
\pgfpathlineto{\pgfqpoint{0.000000in}{-0.025984in}}%
\pgfpathclose%
\pgfusepath{fill}%
}%
\begin{pgfscope}%
\pgfsys@transformshift{1.977049in}{3.038078in}%
\pgfsys@useobject{currentmarker}{}%
\end{pgfscope}%
\end{pgfscope}%
\begin{pgfscope}%
\pgfpathrectangle{\pgfqpoint{1.465200in}{0.822500in}}{\pgfqpoint{0.899100in}{2.537500in}}%
\pgfusepath{clip}%
\pgfsetbuttcap%
\pgfsetroundjoin%
\definecolor{currentfill}{rgb}{1.000000,1.000000,1.000000}%
\pgfsetfillcolor{currentfill}%
\pgfsetlinewidth{0.903375pt}%
\definecolor{currentstroke}{rgb}{0.156863,0.356863,0.568627}%
\pgfsetstrokecolor{currentstroke}%
\pgfsetdash{}{0pt}%
\pgfsys@defobject{currentmarker}{\pgfqpoint{-0.025984in}{-0.025984in}}{\pgfqpoint{0.025984in}{0.025984in}}{%
\pgfpathmoveto{\pgfqpoint{0.000000in}{-0.025984in}}%
\pgfpathcurveto{\pgfqpoint{0.006891in}{-0.025984in}}{\pgfqpoint{0.013501in}{-0.023246in}}{\pgfqpoint{0.018373in}{-0.018373in}}%
\pgfpathcurveto{\pgfqpoint{0.023246in}{-0.013501in}}{\pgfqpoint{0.025984in}{-0.006891in}}{\pgfqpoint{0.025984in}{0.000000in}}%
\pgfpathcurveto{\pgfqpoint{0.025984in}{0.006891in}}{\pgfqpoint{0.023246in}{0.013501in}}{\pgfqpoint{0.018373in}{0.018373in}}%
\pgfpathcurveto{\pgfqpoint{0.013501in}{0.023246in}}{\pgfqpoint{0.006891in}{0.025984in}}{\pgfqpoint{0.000000in}{0.025984in}}%
\pgfpathcurveto{\pgfqpoint{-0.006891in}{0.025984in}}{\pgfqpoint{-0.013501in}{0.023246in}}{\pgfqpoint{-0.018373in}{0.018373in}}%
\pgfpathcurveto{\pgfqpoint{-0.023246in}{0.013501in}}{\pgfqpoint{-0.025984in}{0.006891in}}{\pgfqpoint{-0.025984in}{0.000000in}}%
\pgfpathcurveto{\pgfqpoint{-0.025984in}{-0.006891in}}{\pgfqpoint{-0.023246in}{-0.013501in}}{\pgfqpoint{-0.018373in}{-0.018373in}}%
\pgfpathcurveto{\pgfqpoint{-0.013501in}{-0.023246in}}{\pgfqpoint{-0.006891in}{-0.025984in}}{\pgfqpoint{0.000000in}{-0.025984in}}%
\pgfpathlineto{\pgfqpoint{0.000000in}{-0.025984in}}%
\pgfpathclose%
\pgfusepath{stroke,fill}%
}%
\begin{pgfscope}%
\pgfsys@transformshift{1.918147in}{3.038078in}%
\pgfsys@useobject{currentmarker}{}%
\end{pgfscope}%
\end{pgfscope}%
\begin{pgfscope}%
\pgfpathrectangle{\pgfqpoint{1.465200in}{0.822500in}}{\pgfqpoint{0.899100in}{2.537500in}}%
\pgfusepath{clip}%
\pgfsetbuttcap%
\pgfsetroundjoin%
\definecolor{currentfill}{rgb}{0.156863,0.356863,0.568627}%
\pgfsetfillcolor{currentfill}%
\pgfsetlinewidth{0.000000pt}%
\definecolor{currentstroke}{rgb}{0.156863,0.356863,0.568627}%
\pgfsetstrokecolor{currentstroke}%
\pgfsetdash{}{0pt}%
\pgfsys@defobject{currentmarker}{\pgfqpoint{-0.025984in}{-0.025984in}}{\pgfqpoint{0.025984in}{0.025984in}}{%
\pgfpathmoveto{\pgfqpoint{0.000000in}{-0.025984in}}%
\pgfpathcurveto{\pgfqpoint{0.006891in}{-0.025984in}}{\pgfqpoint{0.013501in}{-0.023246in}}{\pgfqpoint{0.018373in}{-0.018373in}}%
\pgfpathcurveto{\pgfqpoint{0.023246in}{-0.013501in}}{\pgfqpoint{0.025984in}{-0.006891in}}{\pgfqpoint{0.025984in}{0.000000in}}%
\pgfpathcurveto{\pgfqpoint{0.025984in}{0.006891in}}{\pgfqpoint{0.023246in}{0.013501in}}{\pgfqpoint{0.018373in}{0.018373in}}%
\pgfpathcurveto{\pgfqpoint{0.013501in}{0.023246in}}{\pgfqpoint{0.006891in}{0.025984in}}{\pgfqpoint{0.000000in}{0.025984in}}%
\pgfpathcurveto{\pgfqpoint{-0.006891in}{0.025984in}}{\pgfqpoint{-0.013501in}{0.023246in}}{\pgfqpoint{-0.018373in}{0.018373in}}%
\pgfpathcurveto{\pgfqpoint{-0.023246in}{0.013501in}}{\pgfqpoint{-0.025984in}{0.006891in}}{\pgfqpoint{-0.025984in}{0.000000in}}%
\pgfpathcurveto{\pgfqpoint{-0.025984in}{-0.006891in}}{\pgfqpoint{-0.023246in}{-0.013501in}}{\pgfqpoint{-0.018373in}{-0.018373in}}%
\pgfpathcurveto{\pgfqpoint{-0.013501in}{-0.023246in}}{\pgfqpoint{-0.006891in}{-0.025984in}}{\pgfqpoint{0.000000in}{-0.025984in}}%
\pgfpathlineto{\pgfqpoint{0.000000in}{-0.025984in}}%
\pgfpathclose%
\pgfusepath{fill}%
}%
\begin{pgfscope}%
\pgfsys@transformshift{2.048282in}{3.227444in}%
\pgfsys@useobject{currentmarker}{}%
\end{pgfscope}%
\end{pgfscope}%
\begin{pgfscope}%
\pgfpathrectangle{\pgfqpoint{1.465200in}{0.822500in}}{\pgfqpoint{0.899100in}{2.537500in}}%
\pgfusepath{clip}%
\pgfsetbuttcap%
\pgfsetroundjoin%
\definecolor{currentfill}{rgb}{1.000000,1.000000,1.000000}%
\pgfsetfillcolor{currentfill}%
\pgfsetlinewidth{0.903375pt}%
\definecolor{currentstroke}{rgb}{0.156863,0.356863,0.568627}%
\pgfsetstrokecolor{currentstroke}%
\pgfsetdash{}{0pt}%
\pgfsys@defobject{currentmarker}{\pgfqpoint{-0.025984in}{-0.025984in}}{\pgfqpoint{0.025984in}{0.025984in}}{%
\pgfpathmoveto{\pgfqpoint{0.000000in}{-0.025984in}}%
\pgfpathcurveto{\pgfqpoint{0.006891in}{-0.025984in}}{\pgfqpoint{0.013501in}{-0.023246in}}{\pgfqpoint{0.018373in}{-0.018373in}}%
\pgfpathcurveto{\pgfqpoint{0.023246in}{-0.013501in}}{\pgfqpoint{0.025984in}{-0.006891in}}{\pgfqpoint{0.025984in}{0.000000in}}%
\pgfpathcurveto{\pgfqpoint{0.025984in}{0.006891in}}{\pgfqpoint{0.023246in}{0.013501in}}{\pgfqpoint{0.018373in}{0.018373in}}%
\pgfpathcurveto{\pgfqpoint{0.013501in}{0.023246in}}{\pgfqpoint{0.006891in}{0.025984in}}{\pgfqpoint{0.000000in}{0.025984in}}%
\pgfpathcurveto{\pgfqpoint{-0.006891in}{0.025984in}}{\pgfqpoint{-0.013501in}{0.023246in}}{\pgfqpoint{-0.018373in}{0.018373in}}%
\pgfpathcurveto{\pgfqpoint{-0.023246in}{0.013501in}}{\pgfqpoint{-0.025984in}{0.006891in}}{\pgfqpoint{-0.025984in}{0.000000in}}%
\pgfpathcurveto{\pgfqpoint{-0.025984in}{-0.006891in}}{\pgfqpoint{-0.023246in}{-0.013501in}}{\pgfqpoint{-0.018373in}{-0.018373in}}%
\pgfpathcurveto{\pgfqpoint{-0.013501in}{-0.023246in}}{\pgfqpoint{-0.006891in}{-0.025984in}}{\pgfqpoint{0.000000in}{-0.025984in}}%
\pgfpathlineto{\pgfqpoint{0.000000in}{-0.025984in}}%
\pgfpathclose%
\pgfusepath{stroke,fill}%
}%
\begin{pgfscope}%
\pgfsys@transformshift{1.960946in}{3.227444in}%
\pgfsys@useobject{currentmarker}{}%
\end{pgfscope}%
\end{pgfscope}%
\begin{pgfscope}%
\definecolor{textcolor}{rgb}{0.101961,0.101961,0.101961}%
\pgfsetstrokecolor{textcolor}%
\pgfsetfillcolor{textcolor}%
\pgftext[x=1.429236in,y=0.955056in,left,]{\makebox[0pt][r]{\color{textcolor}\sffamily\fontsize{7.000000}{8.400000}\selectfont On-device only}}%
\end{pgfscope}%
\begin{pgfscope}%
\definecolor{textcolor}{rgb}{0.101961,0.101961,0.101961}%
\pgfsetstrokecolor{textcolor}%
\pgfsetfillcolor{textcolor}%
\pgftext[x=1.429236in,y=1.144422in,left,]{\makebox[0pt][r]{\color{textcolor}\sffamily\fontsize{7.000000}{8.400000}\selectfont Warm-all LRU\(\displaystyle ^\dagger\)}}%
\end{pgfscope}%
\begin{pgfscope}%
\definecolor{textcolor}{rgb}{0.101961,0.101961,0.101961}%
\pgfsetstrokecolor{textcolor}%
\pgfsetfillcolor{textcolor}%
\pgftext[x=1.429236in,y=1.333787in,left,]{\makebox[0pt][r]{\color{textcolor}\sffamily\fontsize{7.000000}{8.400000}\selectfont QueryStream~\cite{querystream2026}\(\displaystyle ^\dagger\)}}%
\end{pgfscope}%
\begin{pgfscope}%
\definecolor{textcolor}{rgb}{0.101961,0.101961,0.101961}%
\pgfsetstrokecolor{textcolor}%
\pgfsetfillcolor{textcolor}%
\pgftext[x=1.429236in,y=1.523153in,left,]{\makebox[0pt][r]{\color{textcolor}\sffamily\fontsize{7.000000}{8.400000}\selectfont AttentionStore~\cite{cachedattention2024}\(\displaystyle ^\dagger\)}}%
\end{pgfscope}%
\begin{pgfscope}%
\definecolor{textcolor}{rgb}{0.101961,0.101961,0.101961}%
\pgfsetstrokecolor{textcolor}%
\pgfsetfillcolor{textcolor}%
\pgftext[x=1.429236in,y=1.712519in,left,]{\makebox[0pt][r]{\color{textcolor}\sffamily\fontsize{7.000000}{8.400000}\selectfont Always-cloud FCFS\(\displaystyle ^\dagger\)}}%
\end{pgfscope}%
\begin{pgfscope}%
\definecolor{textcolor}{rgb}{0.101961,0.101961,0.101961}%
\pgfsetstrokecolor{textcolor}%
\pgfsetfillcolor{textcolor}%
\pgftext[x=1.429236in,y=1.901884in,left,]{\makebox[0pt][r]{\color{textcolor}\sffamily\fontsize{7.000000}{8.400000}\selectfont EDF + feasibility~\cite{liu1973edf}\(\displaystyle ^\dagger\)}}%
\end{pgfscope}%
\begin{pgfscope}%
\definecolor{textcolor}{rgb}{0.101961,0.101961,0.101961}%
\pgfsetstrokecolor{textcolor}%
\pgfsetfillcolor{textcolor}%
\pgftext[x=1.429236in,y=2.091250in,left,]{\makebox[0pt][r]{\color{textcolor}\sffamily\fontsize{7.000000}{8.400000}\selectfont EDF~\cite{liu1973edf}\(\displaystyle ^\dagger\)}}%
\end{pgfscope}%
\begin{pgfscope}%
\definecolor{textcolor}{rgb}{0.101961,0.101961,0.101961}%
\pgfsetstrokecolor{textcolor}%
\pgfsetfillcolor{textcolor}%
\pgftext[x=1.429236in,y=2.280616in,left,]{\makebox[0pt][r]{\color{textcolor}\sffamily\fontsize{7.000000}{8.400000}\selectfont ECO-LLM~\cite{ecollm2025}}}%
\end{pgfscope}%
\begin{pgfscope}%
\definecolor{textcolor}{rgb}{0.101961,0.101961,0.101961}%
\pgfsetstrokecolor{textcolor}%
\pgfsetfillcolor{textcolor}%
\pgftext[x=1.429236in,y=2.469981in,left,]{\makebox[0pt][r]{\color{textcolor}\sffamily\fontsize{7.000000}{8.400000}\selectfont LLF~\cite{mok1983llf}}}%
\end{pgfscope}%
\begin{pgfscope}%
\definecolor{textcolor}{rgb}{0.101961,0.101961,0.101961}%
\pgfsetstrokecolor{textcolor}%
\pgfsetfillcolor{textcolor}%
\pgftext[x=1.429236in,y=2.659347in,left,]{\makebox[0pt][r]{\color{textcolor}\sffamily\fontsize{7.000000}{8.400000}\selectfont Mooncake~\cite{mooncake2025}}}%
\end{pgfscope}%
\begin{pgfscope}%
\definecolor{textcolor}{rgb}{0.101961,0.101961,0.101961}%
\pgfsetstrokecolor{textcolor}%
\pgfsetfillcolor{textcolor}%
\pgftext[x=1.429236in,y=2.848713in,left,]{\makebox[0pt][r]{\color{textcolor}\sffamily\fontsize{7.000000}{8.400000}\selectfont SimpleStream~\cite{simplestream2026}}}%
\end{pgfscope}%
\begin{pgfscope}%
\definecolor{textcolor}{rgb}{0.101961,0.101961,0.101961}%
\pgfsetstrokecolor{textcolor}%
\pgfsetfillcolor{textcolor}%
\pgftext[x=1.429236in,y=3.038078in,left,]{\makebox[0pt][r]{\color{textcolor}\sffamily\fontsize{7.000000}{8.400000}\selectfont \textbf{COVEX (old: rung 3)}}}%
\end{pgfscope}%
\begin{pgfscope}%
\definecolor{textcolor}{rgb}{0.101961,0.101961,0.101961}%
\pgfsetstrokecolor{textcolor}%
\pgfsetfillcolor{textcolor}%
\pgftext[x=1.429236in,y=3.227444in,left,]{\makebox[0pt][r]{\color{textcolor}\sffamily\fontsize{7.000000}{8.400000}\selectfont \textbf{COVEX (dev)} (s4)}}%
\end{pgfscope}%
\begin{pgfscope}%
\pgfsetbuttcap%
\pgfsetmiterjoin%
\definecolor{currentfill}{rgb}{1.000000,1.000000,1.000000}%
\pgfsetfillcolor{currentfill}%
\pgfsetlinewidth{0.000000pt}%
\definecolor{currentstroke}{rgb}{0.000000,0.000000,0.000000}%
\pgfsetstrokecolor{currentstroke}%
\pgfsetstrokeopacity{0.000000}%
\pgfsetdash{}{0pt}%
\pgfpathmoveto{\pgfqpoint{2.564100in}{0.822500in}}%
\pgfpathlineto{\pgfqpoint{3.263400in}{0.822500in}}%
\pgfpathlineto{\pgfqpoint{3.263400in}{3.360000in}}%
\pgfpathlineto{\pgfqpoint{2.564100in}{3.360000in}}%
\pgfpathlineto{\pgfqpoint{2.564100in}{0.822500in}}%
\pgfpathclose%
\pgfusepath{fill}%
\end{pgfscope}%
\begin{pgfscope}%
\pgfpathrectangle{\pgfqpoint{2.564100in}{0.822500in}}{\pgfqpoint{0.699300in}{2.537500in}}%
\pgfusepath{clip}%
\pgfsetrectcap%
\pgfsetroundjoin%
\pgfsetlinewidth{0.401500pt}%
\definecolor{currentstroke}{rgb}{0.901961,0.901961,0.901961}%
\pgfsetstrokecolor{currentstroke}%
\pgfsetdash{}{0pt}%
\pgfpathmoveto{\pgfqpoint{2.564100in}{0.822500in}}%
\pgfpathlineto{\pgfqpoint{2.564100in}{3.360000in}}%
\pgfusepath{stroke}%
\end{pgfscope}%
\begin{pgfscope}%
\definecolor{textcolor}{rgb}{0.101961,0.101961,0.101961}%
\pgfsetstrokecolor{textcolor}%
\pgfsetfillcolor{textcolor}%
\pgftext[x=2.564100in,y=0.773889in,,top]{\color{textcolor}\sffamily\fontsize{7.000000}{8.400000}\selectfont 0.0}%
\end{pgfscope}%
\begin{pgfscope}%
\pgfpathrectangle{\pgfqpoint{2.564100in}{0.822500in}}{\pgfqpoint{0.699300in}{2.537500in}}%
\pgfusepath{clip}%
\pgfsetrectcap%
\pgfsetroundjoin%
\pgfsetlinewidth{0.401500pt}%
\definecolor{currentstroke}{rgb}{0.901961,0.901961,0.901961}%
\pgfsetstrokecolor{currentstroke}%
\pgfsetdash{}{0pt}%
\pgfpathmoveto{\pgfqpoint{2.893182in}{0.822500in}}%
\pgfpathlineto{\pgfqpoint{2.893182in}{3.360000in}}%
\pgfusepath{stroke}%
\end{pgfscope}%
\begin{pgfscope}%
\definecolor{textcolor}{rgb}{0.101961,0.101961,0.101961}%
\pgfsetstrokecolor{textcolor}%
\pgfsetfillcolor{textcolor}%
\pgftext[x=2.893182in,y=0.773889in,,top]{\color{textcolor}\sffamily\fontsize{7.000000}{8.400000}\selectfont 0.4}%
\end{pgfscope}%
\begin{pgfscope}%
\pgfpathrectangle{\pgfqpoint{2.564100in}{0.822500in}}{\pgfqpoint{0.699300in}{2.537500in}}%
\pgfusepath{clip}%
\pgfsetrectcap%
\pgfsetroundjoin%
\pgfsetlinewidth{0.401500pt}%
\definecolor{currentstroke}{rgb}{0.901961,0.901961,0.901961}%
\pgfsetstrokecolor{currentstroke}%
\pgfsetdash{}{0pt}%
\pgfpathmoveto{\pgfqpoint{3.222265in}{0.822500in}}%
\pgfpathlineto{\pgfqpoint{3.222265in}{3.360000in}}%
\pgfusepath{stroke}%
\end{pgfscope}%
\begin{pgfscope}%
\definecolor{textcolor}{rgb}{0.101961,0.101961,0.101961}%
\pgfsetstrokecolor{textcolor}%
\pgfsetfillcolor{textcolor}%
\pgftext[x=3.222265in,y=0.773889in,,top]{\color{textcolor}\sffamily\fontsize{7.000000}{8.400000}\selectfont 0.8}%
\end{pgfscope}%
\begin{pgfscope}%
\definecolor{textcolor}{rgb}{0.000000,0.000000,0.000000}%
\pgfsetstrokecolor{textcolor}%
\pgfsetfillcolor{textcolor}%
\pgftext[x=2.716383in, y=0.592176in, left, base]{\color{textcolor}\sffamily\fontsize{7.500000}{9.000000}\selectfont Dashcam}%
\end{pgfscope}%
\begin{pgfscope}%
\definecolor{textcolor}{rgb}{0.000000,0.000000,0.000000}%
\pgfsetstrokecolor{textcolor}%
\pgfsetfillcolor{textcolor}%
\pgftext[x=2.678541in, y=0.492254in, left, base]{\color{textcolor}\sffamily\fontsize{7.500000}{9.000000}\selectfont attainment}%
\end{pgfscope}%
\begin{pgfscope}%
\pgfpathrectangle{\pgfqpoint{2.564100in}{0.822500in}}{\pgfqpoint{0.699300in}{2.537500in}}%
\pgfusepath{clip}%
\pgfsetbuttcap%
\pgfsetroundjoin%
\pgfsetlinewidth{1.003750pt}%
\definecolor{currentstroke}{rgb}{0.478431,0.305882,0.611765}%
\pgfsetstrokecolor{currentstroke}%
\pgfsetdash{}{0pt}%
\pgfpathmoveto{\pgfqpoint{2.564100in}{0.955056in}}%
\pgfpathlineto{\pgfqpoint{2.564100in}{0.955056in}}%
\pgfusepath{stroke}%
\end{pgfscope}%
\begin{pgfscope}%
\pgfpathrectangle{\pgfqpoint{2.564100in}{0.822500in}}{\pgfqpoint{0.699300in}{2.537500in}}%
\pgfusepath{clip}%
\pgfsetbuttcap%
\pgfsetroundjoin%
\pgfsetlinewidth{1.003750pt}%
\definecolor{currentstroke}{rgb}{0.478431,0.478431,0.478431}%
\pgfsetstrokecolor{currentstroke}%
\pgfsetdash{}{0pt}%
\pgfpathmoveto{\pgfqpoint{2.603004in}{1.144422in}}%
\pgfpathlineto{\pgfqpoint{2.601133in}{1.144422in}}%
\pgfusepath{stroke}%
\end{pgfscope}%
\begin{pgfscope}%
\pgfpathrectangle{\pgfqpoint{2.564100in}{0.822500in}}{\pgfqpoint{0.699300in}{2.537500in}}%
\pgfusepath{clip}%
\pgfsetbuttcap%
\pgfsetroundjoin%
\pgfsetlinewidth{1.003750pt}%
\definecolor{currentstroke}{rgb}{0.784314,0.784314,0.784314}%
\pgfsetstrokecolor{currentstroke}%
\pgfsetdash{}{0pt}%
\pgfpathmoveto{\pgfqpoint{2.609780in}{1.333787in}}%
\pgfpathlineto{\pgfqpoint{2.609401in}{1.333787in}}%
\pgfusepath{stroke}%
\end{pgfscope}%
\begin{pgfscope}%
\pgfpathrectangle{\pgfqpoint{2.564100in}{0.822500in}}{\pgfqpoint{0.699300in}{2.537500in}}%
\pgfusepath{clip}%
\pgfsetbuttcap%
\pgfsetroundjoin%
\pgfsetlinewidth{1.003750pt}%
\definecolor{currentstroke}{rgb}{0.478431,0.478431,0.478431}%
\pgfsetstrokecolor{currentstroke}%
\pgfsetdash{}{0pt}%
\pgfpathmoveto{\pgfqpoint{2.612432in}{1.523153in}}%
\pgfpathlineto{\pgfqpoint{2.612344in}{1.523153in}}%
\pgfusepath{stroke}%
\end{pgfscope}%
\begin{pgfscope}%
\pgfpathrectangle{\pgfqpoint{2.564100in}{0.822500in}}{\pgfqpoint{0.699300in}{2.537500in}}%
\pgfusepath{clip}%
\pgfsetbuttcap%
\pgfsetroundjoin%
\pgfsetlinewidth{1.003750pt}%
\definecolor{currentstroke}{rgb}{0.784314,0.784314,0.784314}%
\pgfsetstrokecolor{currentstroke}%
\pgfsetdash{}{0pt}%
\pgfpathmoveto{\pgfqpoint{2.610503in}{1.712519in}}%
\pgfpathlineto{\pgfqpoint{2.614253in}{1.712519in}}%
\pgfusepath{stroke}%
\end{pgfscope}%
\begin{pgfscope}%
\pgfpathrectangle{\pgfqpoint{2.564100in}{0.822500in}}{\pgfqpoint{0.699300in}{2.537500in}}%
\pgfusepath{clip}%
\pgfsetbuttcap%
\pgfsetroundjoin%
\pgfsetlinewidth{1.003750pt}%
\definecolor{currentstroke}{rgb}{0.784314,0.784314,0.784314}%
\pgfsetstrokecolor{currentstroke}%
\pgfsetdash{}{0pt}%
\pgfpathmoveto{\pgfqpoint{2.613973in}{1.901884in}}%
\pgfpathlineto{\pgfqpoint{2.615002in}{1.901884in}}%
\pgfusepath{stroke}%
\end{pgfscope}%
\begin{pgfscope}%
\pgfpathrectangle{\pgfqpoint{2.564100in}{0.822500in}}{\pgfqpoint{0.699300in}{2.537500in}}%
\pgfusepath{clip}%
\pgfsetbuttcap%
\pgfsetroundjoin%
\pgfsetlinewidth{1.003750pt}%
\definecolor{currentstroke}{rgb}{0.784314,0.784314,0.784314}%
\pgfsetstrokecolor{currentstroke}%
\pgfsetdash{}{0pt}%
\pgfpathmoveto{\pgfqpoint{2.777681in}{2.091250in}}%
\pgfpathlineto{\pgfqpoint{2.773400in}{2.091250in}}%
\pgfusepath{stroke}%
\end{pgfscope}%
\begin{pgfscope}%
\pgfpathrectangle{\pgfqpoint{2.564100in}{0.822500in}}{\pgfqpoint{0.699300in}{2.537500in}}%
\pgfusepath{clip}%
\pgfsetbuttcap%
\pgfsetroundjoin%
\pgfsetlinewidth{1.003750pt}%
\definecolor{currentstroke}{rgb}{0.784314,0.784314,0.784314}%
\pgfsetstrokecolor{currentstroke}%
\pgfsetdash{}{0pt}%
\pgfpathmoveto{\pgfqpoint{2.937406in}{2.280616in}}%
\pgfpathlineto{\pgfqpoint{2.817585in}{2.280616in}}%
\pgfusepath{stroke}%
\end{pgfscope}%
\begin{pgfscope}%
\pgfpathrectangle{\pgfqpoint{2.564100in}{0.822500in}}{\pgfqpoint{0.699300in}{2.537500in}}%
\pgfusepath{clip}%
\pgfsetbuttcap%
\pgfsetroundjoin%
\pgfsetlinewidth{1.003750pt}%
\definecolor{currentstroke}{rgb}{0.784314,0.784314,0.784314}%
\pgfsetstrokecolor{currentstroke}%
\pgfsetdash{}{0pt}%
\pgfpathmoveto{\pgfqpoint{3.106121in}{2.469981in}}%
\pgfpathlineto{\pgfqpoint{3.075286in}{2.469981in}}%
\pgfusepath{stroke}%
\end{pgfscope}%
\begin{pgfscope}%
\pgfpathrectangle{\pgfqpoint{2.564100in}{0.822500in}}{\pgfqpoint{0.699300in}{2.537500in}}%
\pgfusepath{clip}%
\pgfsetbuttcap%
\pgfsetroundjoin%
\pgfsetlinewidth{1.003750pt}%
\definecolor{currentstroke}{rgb}{0.478431,0.478431,0.478431}%
\pgfsetstrokecolor{currentstroke}%
\pgfsetdash{}{0pt}%
\pgfpathmoveto{\pgfqpoint{2.963773in}{2.659347in}}%
\pgfpathlineto{\pgfqpoint{2.832619in}{2.659347in}}%
\pgfusepath{stroke}%
\end{pgfscope}%
\begin{pgfscope}%
\pgfpathrectangle{\pgfqpoint{2.564100in}{0.822500in}}{\pgfqpoint{0.699300in}{2.537500in}}%
\pgfusepath{clip}%
\pgfsetbuttcap%
\pgfsetroundjoin%
\pgfsetlinewidth{1.003750pt}%
\definecolor{currentstroke}{rgb}{0.784314,0.784314,0.784314}%
\pgfsetstrokecolor{currentstroke}%
\pgfsetdash{}{0pt}%
\pgfusepath{stroke}%
\end{pgfscope}%
\begin{pgfscope}%
\pgfpathrectangle{\pgfqpoint{2.564100in}{0.822500in}}{\pgfqpoint{0.699300in}{2.537500in}}%
\pgfusepath{clip}%
\pgfsetbuttcap%
\pgfsetroundjoin%
\pgfsetlinewidth{1.003750pt}%
\definecolor{currentstroke}{rgb}{0.156863,0.356863,0.568627}%
\pgfsetstrokecolor{currentstroke}%
\pgfsetdash{}{0pt}%
\pgfpathmoveto{\pgfqpoint{2.968293in}{3.038078in}}%
\pgfpathlineto{\pgfqpoint{3.154835in}{3.038078in}}%
\pgfusepath{stroke}%
\end{pgfscope}%
\begin{pgfscope}%
\pgfpathrectangle{\pgfqpoint{2.564100in}{0.822500in}}{\pgfqpoint{0.699300in}{2.537500in}}%
\pgfusepath{clip}%
\pgfsetbuttcap%
\pgfsetroundjoin%
\pgfsetlinewidth{1.003750pt}%
\definecolor{currentstroke}{rgb}{0.156863,0.356863,0.568627}%
\pgfsetstrokecolor{currentstroke}%
\pgfsetdash{}{0pt}%
\pgfpathmoveto{\pgfqpoint{2.986479in}{3.227444in}}%
\pgfpathlineto{\pgfqpoint{3.160923in}{3.227444in}}%
\pgfusepath{stroke}%
\end{pgfscope}%
\begin{pgfscope}%
\pgfsetrectcap%
\pgfsetmiterjoin%
\pgfsetlinewidth{0.803000pt}%
\definecolor{currentstroke}{rgb}{0.101961,0.101961,0.101961}%
\pgfsetstrokecolor{currentstroke}%
\pgfsetdash{}{0pt}%
\pgfpathmoveto{\pgfqpoint{2.564100in}{0.822500in}}%
\pgfpathlineto{\pgfqpoint{3.263400in}{0.822500in}}%
\pgfusepath{stroke}%
\end{pgfscope}%
\begin{pgfscope}%
\pgfpathrectangle{\pgfqpoint{2.564100in}{0.822500in}}{\pgfqpoint{0.699300in}{2.537500in}}%
\pgfusepath{clip}%
\pgfsetbuttcap%
\pgfsetroundjoin%
\definecolor{currentfill}{rgb}{0.478431,0.305882,0.611765}%
\pgfsetfillcolor{currentfill}%
\pgfsetlinewidth{0.000000pt}%
\definecolor{currentstroke}{rgb}{0.478431,0.305882,0.611765}%
\pgfsetstrokecolor{currentstroke}%
\pgfsetdash{}{0pt}%
\pgfsys@defobject{currentmarker}{\pgfqpoint{-0.025984in}{-0.025984in}}{\pgfqpoint{0.025984in}{0.025984in}}{%
\pgfpathmoveto{\pgfqpoint{0.000000in}{-0.025984in}}%
\pgfpathcurveto{\pgfqpoint{0.006891in}{-0.025984in}}{\pgfqpoint{0.013501in}{-0.023246in}}{\pgfqpoint{0.018373in}{-0.018373in}}%
\pgfpathcurveto{\pgfqpoint{0.023246in}{-0.013501in}}{\pgfqpoint{0.025984in}{-0.006891in}}{\pgfqpoint{0.025984in}{0.000000in}}%
\pgfpathcurveto{\pgfqpoint{0.025984in}{0.006891in}}{\pgfqpoint{0.023246in}{0.013501in}}{\pgfqpoint{0.018373in}{0.018373in}}%
\pgfpathcurveto{\pgfqpoint{0.013501in}{0.023246in}}{\pgfqpoint{0.006891in}{0.025984in}}{\pgfqpoint{0.000000in}{0.025984in}}%
\pgfpathcurveto{\pgfqpoint{-0.006891in}{0.025984in}}{\pgfqpoint{-0.013501in}{0.023246in}}{\pgfqpoint{-0.018373in}{0.018373in}}%
\pgfpathcurveto{\pgfqpoint{-0.023246in}{0.013501in}}{\pgfqpoint{-0.025984in}{0.006891in}}{\pgfqpoint{-0.025984in}{0.000000in}}%
\pgfpathcurveto{\pgfqpoint{-0.025984in}{-0.006891in}}{\pgfqpoint{-0.023246in}{-0.013501in}}{\pgfqpoint{-0.018373in}{-0.018373in}}%
\pgfpathcurveto{\pgfqpoint{-0.013501in}{-0.023246in}}{\pgfqpoint{-0.006891in}{-0.025984in}}{\pgfqpoint{0.000000in}{-0.025984in}}%
\pgfpathlineto{\pgfqpoint{0.000000in}{-0.025984in}}%
\pgfpathclose%
\pgfusepath{fill}%
}%
\begin{pgfscope}%
\pgfsys@transformshift{2.564100in}{0.955056in}%
\pgfsys@useobject{currentmarker}{}%
\end{pgfscope}%
\end{pgfscope}%
\begin{pgfscope}%
\pgfpathrectangle{\pgfqpoint{2.564100in}{0.822500in}}{\pgfqpoint{0.699300in}{2.537500in}}%
\pgfusepath{clip}%
\pgfsetbuttcap%
\pgfsetroundjoin%
\definecolor{currentfill}{rgb}{1.000000,1.000000,1.000000}%
\pgfsetfillcolor{currentfill}%
\pgfsetlinewidth{0.903375pt}%
\definecolor{currentstroke}{rgb}{0.478431,0.305882,0.611765}%
\pgfsetstrokecolor{currentstroke}%
\pgfsetdash{}{0pt}%
\pgfsys@defobject{currentmarker}{\pgfqpoint{-0.025984in}{-0.025984in}}{\pgfqpoint{0.025984in}{0.025984in}}{%
\pgfpathmoveto{\pgfqpoint{0.000000in}{-0.025984in}}%
\pgfpathcurveto{\pgfqpoint{0.006891in}{-0.025984in}}{\pgfqpoint{0.013501in}{-0.023246in}}{\pgfqpoint{0.018373in}{-0.018373in}}%
\pgfpathcurveto{\pgfqpoint{0.023246in}{-0.013501in}}{\pgfqpoint{0.025984in}{-0.006891in}}{\pgfqpoint{0.025984in}{0.000000in}}%
\pgfpathcurveto{\pgfqpoint{0.025984in}{0.006891in}}{\pgfqpoint{0.023246in}{0.013501in}}{\pgfqpoint{0.018373in}{0.018373in}}%
\pgfpathcurveto{\pgfqpoint{0.013501in}{0.023246in}}{\pgfqpoint{0.006891in}{0.025984in}}{\pgfqpoint{0.000000in}{0.025984in}}%
\pgfpathcurveto{\pgfqpoint{-0.006891in}{0.025984in}}{\pgfqpoint{-0.013501in}{0.023246in}}{\pgfqpoint{-0.018373in}{0.018373in}}%
\pgfpathcurveto{\pgfqpoint{-0.023246in}{0.013501in}}{\pgfqpoint{-0.025984in}{0.006891in}}{\pgfqpoint{-0.025984in}{0.000000in}}%
\pgfpathcurveto{\pgfqpoint{-0.025984in}{-0.006891in}}{\pgfqpoint{-0.023246in}{-0.013501in}}{\pgfqpoint{-0.018373in}{-0.018373in}}%
\pgfpathcurveto{\pgfqpoint{-0.013501in}{-0.023246in}}{\pgfqpoint{-0.006891in}{-0.025984in}}{\pgfqpoint{0.000000in}{-0.025984in}}%
\pgfpathlineto{\pgfqpoint{0.000000in}{-0.025984in}}%
\pgfpathclose%
\pgfusepath{stroke,fill}%
}%
\begin{pgfscope}%
\pgfsys@transformshift{2.564100in}{0.955056in}%
\pgfsys@useobject{currentmarker}{}%
\end{pgfscope}%
\end{pgfscope}%
\begin{pgfscope}%
\pgfpathrectangle{\pgfqpoint{2.564100in}{0.822500in}}{\pgfqpoint{0.699300in}{2.537500in}}%
\pgfusepath{clip}%
\pgfsetbuttcap%
\pgfsetroundjoin%
\definecolor{currentfill}{rgb}{0.478431,0.478431,0.478431}%
\pgfsetfillcolor{currentfill}%
\pgfsetlinewidth{0.000000pt}%
\definecolor{currentstroke}{rgb}{0.478431,0.478431,0.478431}%
\pgfsetstrokecolor{currentstroke}%
\pgfsetdash{}{0pt}%
\pgfsys@defobject{currentmarker}{\pgfqpoint{-0.025984in}{-0.025984in}}{\pgfqpoint{0.025984in}{0.025984in}}{%
\pgfpathmoveto{\pgfqpoint{0.000000in}{-0.025984in}}%
\pgfpathcurveto{\pgfqpoint{0.006891in}{-0.025984in}}{\pgfqpoint{0.013501in}{-0.023246in}}{\pgfqpoint{0.018373in}{-0.018373in}}%
\pgfpathcurveto{\pgfqpoint{0.023246in}{-0.013501in}}{\pgfqpoint{0.025984in}{-0.006891in}}{\pgfqpoint{0.025984in}{0.000000in}}%
\pgfpathcurveto{\pgfqpoint{0.025984in}{0.006891in}}{\pgfqpoint{0.023246in}{0.013501in}}{\pgfqpoint{0.018373in}{0.018373in}}%
\pgfpathcurveto{\pgfqpoint{0.013501in}{0.023246in}}{\pgfqpoint{0.006891in}{0.025984in}}{\pgfqpoint{0.000000in}{0.025984in}}%
\pgfpathcurveto{\pgfqpoint{-0.006891in}{0.025984in}}{\pgfqpoint{-0.013501in}{0.023246in}}{\pgfqpoint{-0.018373in}{0.018373in}}%
\pgfpathcurveto{\pgfqpoint{-0.023246in}{0.013501in}}{\pgfqpoint{-0.025984in}{0.006891in}}{\pgfqpoint{-0.025984in}{0.000000in}}%
\pgfpathcurveto{\pgfqpoint{-0.025984in}{-0.006891in}}{\pgfqpoint{-0.023246in}{-0.013501in}}{\pgfqpoint{-0.018373in}{-0.018373in}}%
\pgfpathcurveto{\pgfqpoint{-0.013501in}{-0.023246in}}{\pgfqpoint{-0.006891in}{-0.025984in}}{\pgfqpoint{0.000000in}{-0.025984in}}%
\pgfpathlineto{\pgfqpoint{0.000000in}{-0.025984in}}%
\pgfpathclose%
\pgfusepath{fill}%
}%
\begin{pgfscope}%
\pgfsys@transformshift{2.601133in}{1.144422in}%
\pgfsys@useobject{currentmarker}{}%
\end{pgfscope}%
\end{pgfscope}%
\begin{pgfscope}%
\pgfpathrectangle{\pgfqpoint{2.564100in}{0.822500in}}{\pgfqpoint{0.699300in}{2.537500in}}%
\pgfusepath{clip}%
\pgfsetbuttcap%
\pgfsetroundjoin%
\definecolor{currentfill}{rgb}{1.000000,1.000000,1.000000}%
\pgfsetfillcolor{currentfill}%
\pgfsetlinewidth{0.903375pt}%
\definecolor{currentstroke}{rgb}{0.478431,0.478431,0.478431}%
\pgfsetstrokecolor{currentstroke}%
\pgfsetdash{}{0pt}%
\pgfsys@defobject{currentmarker}{\pgfqpoint{-0.025984in}{-0.025984in}}{\pgfqpoint{0.025984in}{0.025984in}}{%
\pgfpathmoveto{\pgfqpoint{0.000000in}{-0.025984in}}%
\pgfpathcurveto{\pgfqpoint{0.006891in}{-0.025984in}}{\pgfqpoint{0.013501in}{-0.023246in}}{\pgfqpoint{0.018373in}{-0.018373in}}%
\pgfpathcurveto{\pgfqpoint{0.023246in}{-0.013501in}}{\pgfqpoint{0.025984in}{-0.006891in}}{\pgfqpoint{0.025984in}{0.000000in}}%
\pgfpathcurveto{\pgfqpoint{0.025984in}{0.006891in}}{\pgfqpoint{0.023246in}{0.013501in}}{\pgfqpoint{0.018373in}{0.018373in}}%
\pgfpathcurveto{\pgfqpoint{0.013501in}{0.023246in}}{\pgfqpoint{0.006891in}{0.025984in}}{\pgfqpoint{0.000000in}{0.025984in}}%
\pgfpathcurveto{\pgfqpoint{-0.006891in}{0.025984in}}{\pgfqpoint{-0.013501in}{0.023246in}}{\pgfqpoint{-0.018373in}{0.018373in}}%
\pgfpathcurveto{\pgfqpoint{-0.023246in}{0.013501in}}{\pgfqpoint{-0.025984in}{0.006891in}}{\pgfqpoint{-0.025984in}{0.000000in}}%
\pgfpathcurveto{\pgfqpoint{-0.025984in}{-0.006891in}}{\pgfqpoint{-0.023246in}{-0.013501in}}{\pgfqpoint{-0.018373in}{-0.018373in}}%
\pgfpathcurveto{\pgfqpoint{-0.013501in}{-0.023246in}}{\pgfqpoint{-0.006891in}{-0.025984in}}{\pgfqpoint{0.000000in}{-0.025984in}}%
\pgfpathlineto{\pgfqpoint{0.000000in}{-0.025984in}}%
\pgfpathclose%
\pgfusepath{stroke,fill}%
}%
\begin{pgfscope}%
\pgfsys@transformshift{2.603004in}{1.144422in}%
\pgfsys@useobject{currentmarker}{}%
\end{pgfscope}%
\end{pgfscope}%
\begin{pgfscope}%
\pgfpathrectangle{\pgfqpoint{2.564100in}{0.822500in}}{\pgfqpoint{0.699300in}{2.537500in}}%
\pgfusepath{clip}%
\pgfsetbuttcap%
\pgfsetroundjoin%
\definecolor{currentfill}{rgb}{0.784314,0.784314,0.784314}%
\pgfsetfillcolor{currentfill}%
\pgfsetlinewidth{0.000000pt}%
\definecolor{currentstroke}{rgb}{0.784314,0.784314,0.784314}%
\pgfsetstrokecolor{currentstroke}%
\pgfsetdash{}{0pt}%
\pgfsys@defobject{currentmarker}{\pgfqpoint{-0.025984in}{-0.025984in}}{\pgfqpoint{0.025984in}{0.025984in}}{%
\pgfpathmoveto{\pgfqpoint{0.000000in}{-0.025984in}}%
\pgfpathcurveto{\pgfqpoint{0.006891in}{-0.025984in}}{\pgfqpoint{0.013501in}{-0.023246in}}{\pgfqpoint{0.018373in}{-0.018373in}}%
\pgfpathcurveto{\pgfqpoint{0.023246in}{-0.013501in}}{\pgfqpoint{0.025984in}{-0.006891in}}{\pgfqpoint{0.025984in}{0.000000in}}%
\pgfpathcurveto{\pgfqpoint{0.025984in}{0.006891in}}{\pgfqpoint{0.023246in}{0.013501in}}{\pgfqpoint{0.018373in}{0.018373in}}%
\pgfpathcurveto{\pgfqpoint{0.013501in}{0.023246in}}{\pgfqpoint{0.006891in}{0.025984in}}{\pgfqpoint{0.000000in}{0.025984in}}%
\pgfpathcurveto{\pgfqpoint{-0.006891in}{0.025984in}}{\pgfqpoint{-0.013501in}{0.023246in}}{\pgfqpoint{-0.018373in}{0.018373in}}%
\pgfpathcurveto{\pgfqpoint{-0.023246in}{0.013501in}}{\pgfqpoint{-0.025984in}{0.006891in}}{\pgfqpoint{-0.025984in}{0.000000in}}%
\pgfpathcurveto{\pgfqpoint{-0.025984in}{-0.006891in}}{\pgfqpoint{-0.023246in}{-0.013501in}}{\pgfqpoint{-0.018373in}{-0.018373in}}%
\pgfpathcurveto{\pgfqpoint{-0.013501in}{-0.023246in}}{\pgfqpoint{-0.006891in}{-0.025984in}}{\pgfqpoint{0.000000in}{-0.025984in}}%
\pgfpathlineto{\pgfqpoint{0.000000in}{-0.025984in}}%
\pgfpathclose%
\pgfusepath{fill}%
}%
\begin{pgfscope}%
\pgfsys@transformshift{2.609401in}{1.333787in}%
\pgfsys@useobject{currentmarker}{}%
\end{pgfscope}%
\end{pgfscope}%
\begin{pgfscope}%
\pgfpathrectangle{\pgfqpoint{2.564100in}{0.822500in}}{\pgfqpoint{0.699300in}{2.537500in}}%
\pgfusepath{clip}%
\pgfsetbuttcap%
\pgfsetroundjoin%
\definecolor{currentfill}{rgb}{1.000000,1.000000,1.000000}%
\pgfsetfillcolor{currentfill}%
\pgfsetlinewidth{0.903375pt}%
\definecolor{currentstroke}{rgb}{0.784314,0.784314,0.784314}%
\pgfsetstrokecolor{currentstroke}%
\pgfsetdash{}{0pt}%
\pgfsys@defobject{currentmarker}{\pgfqpoint{-0.025984in}{-0.025984in}}{\pgfqpoint{0.025984in}{0.025984in}}{%
\pgfpathmoveto{\pgfqpoint{0.000000in}{-0.025984in}}%
\pgfpathcurveto{\pgfqpoint{0.006891in}{-0.025984in}}{\pgfqpoint{0.013501in}{-0.023246in}}{\pgfqpoint{0.018373in}{-0.018373in}}%
\pgfpathcurveto{\pgfqpoint{0.023246in}{-0.013501in}}{\pgfqpoint{0.025984in}{-0.006891in}}{\pgfqpoint{0.025984in}{0.000000in}}%
\pgfpathcurveto{\pgfqpoint{0.025984in}{0.006891in}}{\pgfqpoint{0.023246in}{0.013501in}}{\pgfqpoint{0.018373in}{0.018373in}}%
\pgfpathcurveto{\pgfqpoint{0.013501in}{0.023246in}}{\pgfqpoint{0.006891in}{0.025984in}}{\pgfqpoint{0.000000in}{0.025984in}}%
\pgfpathcurveto{\pgfqpoint{-0.006891in}{0.025984in}}{\pgfqpoint{-0.013501in}{0.023246in}}{\pgfqpoint{-0.018373in}{0.018373in}}%
\pgfpathcurveto{\pgfqpoint{-0.023246in}{0.013501in}}{\pgfqpoint{-0.025984in}{0.006891in}}{\pgfqpoint{-0.025984in}{0.000000in}}%
\pgfpathcurveto{\pgfqpoint{-0.025984in}{-0.006891in}}{\pgfqpoint{-0.023246in}{-0.013501in}}{\pgfqpoint{-0.018373in}{-0.018373in}}%
\pgfpathcurveto{\pgfqpoint{-0.013501in}{-0.023246in}}{\pgfqpoint{-0.006891in}{-0.025984in}}{\pgfqpoint{0.000000in}{-0.025984in}}%
\pgfpathlineto{\pgfqpoint{0.000000in}{-0.025984in}}%
\pgfpathclose%
\pgfusepath{stroke,fill}%
}%
\begin{pgfscope}%
\pgfsys@transformshift{2.609780in}{1.333787in}%
\pgfsys@useobject{currentmarker}{}%
\end{pgfscope}%
\end{pgfscope}%
\begin{pgfscope}%
\pgfpathrectangle{\pgfqpoint{2.564100in}{0.822500in}}{\pgfqpoint{0.699300in}{2.537500in}}%
\pgfusepath{clip}%
\pgfsetbuttcap%
\pgfsetroundjoin%
\definecolor{currentfill}{rgb}{0.478431,0.478431,0.478431}%
\pgfsetfillcolor{currentfill}%
\pgfsetlinewidth{0.000000pt}%
\definecolor{currentstroke}{rgb}{0.478431,0.478431,0.478431}%
\pgfsetstrokecolor{currentstroke}%
\pgfsetdash{}{0pt}%
\pgfsys@defobject{currentmarker}{\pgfqpoint{-0.025984in}{-0.025984in}}{\pgfqpoint{0.025984in}{0.025984in}}{%
\pgfpathmoveto{\pgfqpoint{0.000000in}{-0.025984in}}%
\pgfpathcurveto{\pgfqpoint{0.006891in}{-0.025984in}}{\pgfqpoint{0.013501in}{-0.023246in}}{\pgfqpoint{0.018373in}{-0.018373in}}%
\pgfpathcurveto{\pgfqpoint{0.023246in}{-0.013501in}}{\pgfqpoint{0.025984in}{-0.006891in}}{\pgfqpoint{0.025984in}{0.000000in}}%
\pgfpathcurveto{\pgfqpoint{0.025984in}{0.006891in}}{\pgfqpoint{0.023246in}{0.013501in}}{\pgfqpoint{0.018373in}{0.018373in}}%
\pgfpathcurveto{\pgfqpoint{0.013501in}{0.023246in}}{\pgfqpoint{0.006891in}{0.025984in}}{\pgfqpoint{0.000000in}{0.025984in}}%
\pgfpathcurveto{\pgfqpoint{-0.006891in}{0.025984in}}{\pgfqpoint{-0.013501in}{0.023246in}}{\pgfqpoint{-0.018373in}{0.018373in}}%
\pgfpathcurveto{\pgfqpoint{-0.023246in}{0.013501in}}{\pgfqpoint{-0.025984in}{0.006891in}}{\pgfqpoint{-0.025984in}{0.000000in}}%
\pgfpathcurveto{\pgfqpoint{-0.025984in}{-0.006891in}}{\pgfqpoint{-0.023246in}{-0.013501in}}{\pgfqpoint{-0.018373in}{-0.018373in}}%
\pgfpathcurveto{\pgfqpoint{-0.013501in}{-0.023246in}}{\pgfqpoint{-0.006891in}{-0.025984in}}{\pgfqpoint{0.000000in}{-0.025984in}}%
\pgfpathlineto{\pgfqpoint{0.000000in}{-0.025984in}}%
\pgfpathclose%
\pgfusepath{fill}%
}%
\begin{pgfscope}%
\pgfsys@transformshift{2.612344in}{1.523153in}%
\pgfsys@useobject{currentmarker}{}%
\end{pgfscope}%
\end{pgfscope}%
\begin{pgfscope}%
\pgfpathrectangle{\pgfqpoint{2.564100in}{0.822500in}}{\pgfqpoint{0.699300in}{2.537500in}}%
\pgfusepath{clip}%
\pgfsetbuttcap%
\pgfsetroundjoin%
\definecolor{currentfill}{rgb}{1.000000,1.000000,1.000000}%
\pgfsetfillcolor{currentfill}%
\pgfsetlinewidth{0.903375pt}%
\definecolor{currentstroke}{rgb}{0.478431,0.478431,0.478431}%
\pgfsetstrokecolor{currentstroke}%
\pgfsetdash{}{0pt}%
\pgfsys@defobject{currentmarker}{\pgfqpoint{-0.025984in}{-0.025984in}}{\pgfqpoint{0.025984in}{0.025984in}}{%
\pgfpathmoveto{\pgfqpoint{0.000000in}{-0.025984in}}%
\pgfpathcurveto{\pgfqpoint{0.006891in}{-0.025984in}}{\pgfqpoint{0.013501in}{-0.023246in}}{\pgfqpoint{0.018373in}{-0.018373in}}%
\pgfpathcurveto{\pgfqpoint{0.023246in}{-0.013501in}}{\pgfqpoint{0.025984in}{-0.006891in}}{\pgfqpoint{0.025984in}{0.000000in}}%
\pgfpathcurveto{\pgfqpoint{0.025984in}{0.006891in}}{\pgfqpoint{0.023246in}{0.013501in}}{\pgfqpoint{0.018373in}{0.018373in}}%
\pgfpathcurveto{\pgfqpoint{0.013501in}{0.023246in}}{\pgfqpoint{0.006891in}{0.025984in}}{\pgfqpoint{0.000000in}{0.025984in}}%
\pgfpathcurveto{\pgfqpoint{-0.006891in}{0.025984in}}{\pgfqpoint{-0.013501in}{0.023246in}}{\pgfqpoint{-0.018373in}{0.018373in}}%
\pgfpathcurveto{\pgfqpoint{-0.023246in}{0.013501in}}{\pgfqpoint{-0.025984in}{0.006891in}}{\pgfqpoint{-0.025984in}{0.000000in}}%
\pgfpathcurveto{\pgfqpoint{-0.025984in}{-0.006891in}}{\pgfqpoint{-0.023246in}{-0.013501in}}{\pgfqpoint{-0.018373in}{-0.018373in}}%
\pgfpathcurveto{\pgfqpoint{-0.013501in}{-0.023246in}}{\pgfqpoint{-0.006891in}{-0.025984in}}{\pgfqpoint{0.000000in}{-0.025984in}}%
\pgfpathlineto{\pgfqpoint{0.000000in}{-0.025984in}}%
\pgfpathclose%
\pgfusepath{stroke,fill}%
}%
\begin{pgfscope}%
\pgfsys@transformshift{2.612432in}{1.523153in}%
\pgfsys@useobject{currentmarker}{}%
\end{pgfscope}%
\end{pgfscope}%
\begin{pgfscope}%
\pgfpathrectangle{\pgfqpoint{2.564100in}{0.822500in}}{\pgfqpoint{0.699300in}{2.537500in}}%
\pgfusepath{clip}%
\pgfsetbuttcap%
\pgfsetroundjoin%
\definecolor{currentfill}{rgb}{0.784314,0.784314,0.784314}%
\pgfsetfillcolor{currentfill}%
\pgfsetlinewidth{0.000000pt}%
\definecolor{currentstroke}{rgb}{0.784314,0.784314,0.784314}%
\pgfsetstrokecolor{currentstroke}%
\pgfsetdash{}{0pt}%
\pgfsys@defobject{currentmarker}{\pgfqpoint{-0.025984in}{-0.025984in}}{\pgfqpoint{0.025984in}{0.025984in}}{%
\pgfpathmoveto{\pgfqpoint{0.000000in}{-0.025984in}}%
\pgfpathcurveto{\pgfqpoint{0.006891in}{-0.025984in}}{\pgfqpoint{0.013501in}{-0.023246in}}{\pgfqpoint{0.018373in}{-0.018373in}}%
\pgfpathcurveto{\pgfqpoint{0.023246in}{-0.013501in}}{\pgfqpoint{0.025984in}{-0.006891in}}{\pgfqpoint{0.025984in}{0.000000in}}%
\pgfpathcurveto{\pgfqpoint{0.025984in}{0.006891in}}{\pgfqpoint{0.023246in}{0.013501in}}{\pgfqpoint{0.018373in}{0.018373in}}%
\pgfpathcurveto{\pgfqpoint{0.013501in}{0.023246in}}{\pgfqpoint{0.006891in}{0.025984in}}{\pgfqpoint{0.000000in}{0.025984in}}%
\pgfpathcurveto{\pgfqpoint{-0.006891in}{0.025984in}}{\pgfqpoint{-0.013501in}{0.023246in}}{\pgfqpoint{-0.018373in}{0.018373in}}%
\pgfpathcurveto{\pgfqpoint{-0.023246in}{0.013501in}}{\pgfqpoint{-0.025984in}{0.006891in}}{\pgfqpoint{-0.025984in}{0.000000in}}%
\pgfpathcurveto{\pgfqpoint{-0.025984in}{-0.006891in}}{\pgfqpoint{-0.023246in}{-0.013501in}}{\pgfqpoint{-0.018373in}{-0.018373in}}%
\pgfpathcurveto{\pgfqpoint{-0.013501in}{-0.023246in}}{\pgfqpoint{-0.006891in}{-0.025984in}}{\pgfqpoint{0.000000in}{-0.025984in}}%
\pgfpathlineto{\pgfqpoint{0.000000in}{-0.025984in}}%
\pgfpathclose%
\pgfusepath{fill}%
}%
\begin{pgfscope}%
\pgfsys@transformshift{2.614253in}{1.712519in}%
\pgfsys@useobject{currentmarker}{}%
\end{pgfscope}%
\end{pgfscope}%
\begin{pgfscope}%
\pgfpathrectangle{\pgfqpoint{2.564100in}{0.822500in}}{\pgfqpoint{0.699300in}{2.537500in}}%
\pgfusepath{clip}%
\pgfsetbuttcap%
\pgfsetroundjoin%
\definecolor{currentfill}{rgb}{1.000000,1.000000,1.000000}%
\pgfsetfillcolor{currentfill}%
\pgfsetlinewidth{0.903375pt}%
\definecolor{currentstroke}{rgb}{0.784314,0.784314,0.784314}%
\pgfsetstrokecolor{currentstroke}%
\pgfsetdash{}{0pt}%
\pgfsys@defobject{currentmarker}{\pgfqpoint{-0.025984in}{-0.025984in}}{\pgfqpoint{0.025984in}{0.025984in}}{%
\pgfpathmoveto{\pgfqpoint{0.000000in}{-0.025984in}}%
\pgfpathcurveto{\pgfqpoint{0.006891in}{-0.025984in}}{\pgfqpoint{0.013501in}{-0.023246in}}{\pgfqpoint{0.018373in}{-0.018373in}}%
\pgfpathcurveto{\pgfqpoint{0.023246in}{-0.013501in}}{\pgfqpoint{0.025984in}{-0.006891in}}{\pgfqpoint{0.025984in}{0.000000in}}%
\pgfpathcurveto{\pgfqpoint{0.025984in}{0.006891in}}{\pgfqpoint{0.023246in}{0.013501in}}{\pgfqpoint{0.018373in}{0.018373in}}%
\pgfpathcurveto{\pgfqpoint{0.013501in}{0.023246in}}{\pgfqpoint{0.006891in}{0.025984in}}{\pgfqpoint{0.000000in}{0.025984in}}%
\pgfpathcurveto{\pgfqpoint{-0.006891in}{0.025984in}}{\pgfqpoint{-0.013501in}{0.023246in}}{\pgfqpoint{-0.018373in}{0.018373in}}%
\pgfpathcurveto{\pgfqpoint{-0.023246in}{0.013501in}}{\pgfqpoint{-0.025984in}{0.006891in}}{\pgfqpoint{-0.025984in}{0.000000in}}%
\pgfpathcurveto{\pgfqpoint{-0.025984in}{-0.006891in}}{\pgfqpoint{-0.023246in}{-0.013501in}}{\pgfqpoint{-0.018373in}{-0.018373in}}%
\pgfpathcurveto{\pgfqpoint{-0.013501in}{-0.023246in}}{\pgfqpoint{-0.006891in}{-0.025984in}}{\pgfqpoint{0.000000in}{-0.025984in}}%
\pgfpathlineto{\pgfqpoint{0.000000in}{-0.025984in}}%
\pgfpathclose%
\pgfusepath{stroke,fill}%
}%
\begin{pgfscope}%
\pgfsys@transformshift{2.610503in}{1.712519in}%
\pgfsys@useobject{currentmarker}{}%
\end{pgfscope}%
\end{pgfscope}%
\begin{pgfscope}%
\pgfpathrectangle{\pgfqpoint{2.564100in}{0.822500in}}{\pgfqpoint{0.699300in}{2.537500in}}%
\pgfusepath{clip}%
\pgfsetbuttcap%
\pgfsetroundjoin%
\definecolor{currentfill}{rgb}{0.784314,0.784314,0.784314}%
\pgfsetfillcolor{currentfill}%
\pgfsetlinewidth{0.000000pt}%
\definecolor{currentstroke}{rgb}{0.784314,0.784314,0.784314}%
\pgfsetstrokecolor{currentstroke}%
\pgfsetdash{}{0pt}%
\pgfsys@defobject{currentmarker}{\pgfqpoint{-0.025984in}{-0.025984in}}{\pgfqpoint{0.025984in}{0.025984in}}{%
\pgfpathmoveto{\pgfqpoint{0.000000in}{-0.025984in}}%
\pgfpathcurveto{\pgfqpoint{0.006891in}{-0.025984in}}{\pgfqpoint{0.013501in}{-0.023246in}}{\pgfqpoint{0.018373in}{-0.018373in}}%
\pgfpathcurveto{\pgfqpoint{0.023246in}{-0.013501in}}{\pgfqpoint{0.025984in}{-0.006891in}}{\pgfqpoint{0.025984in}{0.000000in}}%
\pgfpathcurveto{\pgfqpoint{0.025984in}{0.006891in}}{\pgfqpoint{0.023246in}{0.013501in}}{\pgfqpoint{0.018373in}{0.018373in}}%
\pgfpathcurveto{\pgfqpoint{0.013501in}{0.023246in}}{\pgfqpoint{0.006891in}{0.025984in}}{\pgfqpoint{0.000000in}{0.025984in}}%
\pgfpathcurveto{\pgfqpoint{-0.006891in}{0.025984in}}{\pgfqpoint{-0.013501in}{0.023246in}}{\pgfqpoint{-0.018373in}{0.018373in}}%
\pgfpathcurveto{\pgfqpoint{-0.023246in}{0.013501in}}{\pgfqpoint{-0.025984in}{0.006891in}}{\pgfqpoint{-0.025984in}{0.000000in}}%
\pgfpathcurveto{\pgfqpoint{-0.025984in}{-0.006891in}}{\pgfqpoint{-0.023246in}{-0.013501in}}{\pgfqpoint{-0.018373in}{-0.018373in}}%
\pgfpathcurveto{\pgfqpoint{-0.013501in}{-0.023246in}}{\pgfqpoint{-0.006891in}{-0.025984in}}{\pgfqpoint{0.000000in}{-0.025984in}}%
\pgfpathlineto{\pgfqpoint{0.000000in}{-0.025984in}}%
\pgfpathclose%
\pgfusepath{fill}%
}%
\begin{pgfscope}%
\pgfsys@transformshift{2.615002in}{1.901884in}%
\pgfsys@useobject{currentmarker}{}%
\end{pgfscope}%
\end{pgfscope}%
\begin{pgfscope}%
\pgfpathrectangle{\pgfqpoint{2.564100in}{0.822500in}}{\pgfqpoint{0.699300in}{2.537500in}}%
\pgfusepath{clip}%
\pgfsetbuttcap%
\pgfsetroundjoin%
\definecolor{currentfill}{rgb}{1.000000,1.000000,1.000000}%
\pgfsetfillcolor{currentfill}%
\pgfsetlinewidth{0.903375pt}%
\definecolor{currentstroke}{rgb}{0.784314,0.784314,0.784314}%
\pgfsetstrokecolor{currentstroke}%
\pgfsetdash{}{0pt}%
\pgfsys@defobject{currentmarker}{\pgfqpoint{-0.025984in}{-0.025984in}}{\pgfqpoint{0.025984in}{0.025984in}}{%
\pgfpathmoveto{\pgfqpoint{0.000000in}{-0.025984in}}%
\pgfpathcurveto{\pgfqpoint{0.006891in}{-0.025984in}}{\pgfqpoint{0.013501in}{-0.023246in}}{\pgfqpoint{0.018373in}{-0.018373in}}%
\pgfpathcurveto{\pgfqpoint{0.023246in}{-0.013501in}}{\pgfqpoint{0.025984in}{-0.006891in}}{\pgfqpoint{0.025984in}{0.000000in}}%
\pgfpathcurveto{\pgfqpoint{0.025984in}{0.006891in}}{\pgfqpoint{0.023246in}{0.013501in}}{\pgfqpoint{0.018373in}{0.018373in}}%
\pgfpathcurveto{\pgfqpoint{0.013501in}{0.023246in}}{\pgfqpoint{0.006891in}{0.025984in}}{\pgfqpoint{0.000000in}{0.025984in}}%
\pgfpathcurveto{\pgfqpoint{-0.006891in}{0.025984in}}{\pgfqpoint{-0.013501in}{0.023246in}}{\pgfqpoint{-0.018373in}{0.018373in}}%
\pgfpathcurveto{\pgfqpoint{-0.023246in}{0.013501in}}{\pgfqpoint{-0.025984in}{0.006891in}}{\pgfqpoint{-0.025984in}{0.000000in}}%
\pgfpathcurveto{\pgfqpoint{-0.025984in}{-0.006891in}}{\pgfqpoint{-0.023246in}{-0.013501in}}{\pgfqpoint{-0.018373in}{-0.018373in}}%
\pgfpathcurveto{\pgfqpoint{-0.013501in}{-0.023246in}}{\pgfqpoint{-0.006891in}{-0.025984in}}{\pgfqpoint{0.000000in}{-0.025984in}}%
\pgfpathlineto{\pgfqpoint{0.000000in}{-0.025984in}}%
\pgfpathclose%
\pgfusepath{stroke,fill}%
}%
\begin{pgfscope}%
\pgfsys@transformshift{2.613973in}{1.901884in}%
\pgfsys@useobject{currentmarker}{}%
\end{pgfscope}%
\end{pgfscope}%
\begin{pgfscope}%
\pgfpathrectangle{\pgfqpoint{2.564100in}{0.822500in}}{\pgfqpoint{0.699300in}{2.537500in}}%
\pgfusepath{clip}%
\pgfsetbuttcap%
\pgfsetroundjoin%
\definecolor{currentfill}{rgb}{0.784314,0.784314,0.784314}%
\pgfsetfillcolor{currentfill}%
\pgfsetlinewidth{0.000000pt}%
\definecolor{currentstroke}{rgb}{0.784314,0.784314,0.784314}%
\pgfsetstrokecolor{currentstroke}%
\pgfsetdash{}{0pt}%
\pgfsys@defobject{currentmarker}{\pgfqpoint{-0.025984in}{-0.025984in}}{\pgfqpoint{0.025984in}{0.025984in}}{%
\pgfpathmoveto{\pgfqpoint{0.000000in}{-0.025984in}}%
\pgfpathcurveto{\pgfqpoint{0.006891in}{-0.025984in}}{\pgfqpoint{0.013501in}{-0.023246in}}{\pgfqpoint{0.018373in}{-0.018373in}}%
\pgfpathcurveto{\pgfqpoint{0.023246in}{-0.013501in}}{\pgfqpoint{0.025984in}{-0.006891in}}{\pgfqpoint{0.025984in}{0.000000in}}%
\pgfpathcurveto{\pgfqpoint{0.025984in}{0.006891in}}{\pgfqpoint{0.023246in}{0.013501in}}{\pgfqpoint{0.018373in}{0.018373in}}%
\pgfpathcurveto{\pgfqpoint{0.013501in}{0.023246in}}{\pgfqpoint{0.006891in}{0.025984in}}{\pgfqpoint{0.000000in}{0.025984in}}%
\pgfpathcurveto{\pgfqpoint{-0.006891in}{0.025984in}}{\pgfqpoint{-0.013501in}{0.023246in}}{\pgfqpoint{-0.018373in}{0.018373in}}%
\pgfpathcurveto{\pgfqpoint{-0.023246in}{0.013501in}}{\pgfqpoint{-0.025984in}{0.006891in}}{\pgfqpoint{-0.025984in}{0.000000in}}%
\pgfpathcurveto{\pgfqpoint{-0.025984in}{-0.006891in}}{\pgfqpoint{-0.023246in}{-0.013501in}}{\pgfqpoint{-0.018373in}{-0.018373in}}%
\pgfpathcurveto{\pgfqpoint{-0.013501in}{-0.023246in}}{\pgfqpoint{-0.006891in}{-0.025984in}}{\pgfqpoint{0.000000in}{-0.025984in}}%
\pgfpathlineto{\pgfqpoint{0.000000in}{-0.025984in}}%
\pgfpathclose%
\pgfusepath{fill}%
}%
\begin{pgfscope}%
\pgfsys@transformshift{2.773400in}{2.091250in}%
\pgfsys@useobject{currentmarker}{}%
\end{pgfscope}%
\end{pgfscope}%
\begin{pgfscope}%
\pgfpathrectangle{\pgfqpoint{2.564100in}{0.822500in}}{\pgfqpoint{0.699300in}{2.537500in}}%
\pgfusepath{clip}%
\pgfsetbuttcap%
\pgfsetroundjoin%
\definecolor{currentfill}{rgb}{1.000000,1.000000,1.000000}%
\pgfsetfillcolor{currentfill}%
\pgfsetlinewidth{0.903375pt}%
\definecolor{currentstroke}{rgb}{0.784314,0.784314,0.784314}%
\pgfsetstrokecolor{currentstroke}%
\pgfsetdash{}{0pt}%
\pgfsys@defobject{currentmarker}{\pgfqpoint{-0.025984in}{-0.025984in}}{\pgfqpoint{0.025984in}{0.025984in}}{%
\pgfpathmoveto{\pgfqpoint{0.000000in}{-0.025984in}}%
\pgfpathcurveto{\pgfqpoint{0.006891in}{-0.025984in}}{\pgfqpoint{0.013501in}{-0.023246in}}{\pgfqpoint{0.018373in}{-0.018373in}}%
\pgfpathcurveto{\pgfqpoint{0.023246in}{-0.013501in}}{\pgfqpoint{0.025984in}{-0.006891in}}{\pgfqpoint{0.025984in}{0.000000in}}%
\pgfpathcurveto{\pgfqpoint{0.025984in}{0.006891in}}{\pgfqpoint{0.023246in}{0.013501in}}{\pgfqpoint{0.018373in}{0.018373in}}%
\pgfpathcurveto{\pgfqpoint{0.013501in}{0.023246in}}{\pgfqpoint{0.006891in}{0.025984in}}{\pgfqpoint{0.000000in}{0.025984in}}%
\pgfpathcurveto{\pgfqpoint{-0.006891in}{0.025984in}}{\pgfqpoint{-0.013501in}{0.023246in}}{\pgfqpoint{-0.018373in}{0.018373in}}%
\pgfpathcurveto{\pgfqpoint{-0.023246in}{0.013501in}}{\pgfqpoint{-0.025984in}{0.006891in}}{\pgfqpoint{-0.025984in}{0.000000in}}%
\pgfpathcurveto{\pgfqpoint{-0.025984in}{-0.006891in}}{\pgfqpoint{-0.023246in}{-0.013501in}}{\pgfqpoint{-0.018373in}{-0.018373in}}%
\pgfpathcurveto{\pgfqpoint{-0.013501in}{-0.023246in}}{\pgfqpoint{-0.006891in}{-0.025984in}}{\pgfqpoint{0.000000in}{-0.025984in}}%
\pgfpathlineto{\pgfqpoint{0.000000in}{-0.025984in}}%
\pgfpathclose%
\pgfusepath{stroke,fill}%
}%
\begin{pgfscope}%
\pgfsys@transformshift{2.777681in}{2.091250in}%
\pgfsys@useobject{currentmarker}{}%
\end{pgfscope}%
\end{pgfscope}%
\begin{pgfscope}%
\pgfpathrectangle{\pgfqpoint{2.564100in}{0.822500in}}{\pgfqpoint{0.699300in}{2.537500in}}%
\pgfusepath{clip}%
\pgfsetbuttcap%
\pgfsetroundjoin%
\definecolor{currentfill}{rgb}{0.784314,0.784314,0.784314}%
\pgfsetfillcolor{currentfill}%
\pgfsetlinewidth{0.000000pt}%
\definecolor{currentstroke}{rgb}{0.784314,0.784314,0.784314}%
\pgfsetstrokecolor{currentstroke}%
\pgfsetdash{}{0pt}%
\pgfsys@defobject{currentmarker}{\pgfqpoint{-0.025984in}{-0.025984in}}{\pgfqpoint{0.025984in}{0.025984in}}{%
\pgfpathmoveto{\pgfqpoint{0.000000in}{-0.025984in}}%
\pgfpathcurveto{\pgfqpoint{0.006891in}{-0.025984in}}{\pgfqpoint{0.013501in}{-0.023246in}}{\pgfqpoint{0.018373in}{-0.018373in}}%
\pgfpathcurveto{\pgfqpoint{0.023246in}{-0.013501in}}{\pgfqpoint{0.025984in}{-0.006891in}}{\pgfqpoint{0.025984in}{0.000000in}}%
\pgfpathcurveto{\pgfqpoint{0.025984in}{0.006891in}}{\pgfqpoint{0.023246in}{0.013501in}}{\pgfqpoint{0.018373in}{0.018373in}}%
\pgfpathcurveto{\pgfqpoint{0.013501in}{0.023246in}}{\pgfqpoint{0.006891in}{0.025984in}}{\pgfqpoint{0.000000in}{0.025984in}}%
\pgfpathcurveto{\pgfqpoint{-0.006891in}{0.025984in}}{\pgfqpoint{-0.013501in}{0.023246in}}{\pgfqpoint{-0.018373in}{0.018373in}}%
\pgfpathcurveto{\pgfqpoint{-0.023246in}{0.013501in}}{\pgfqpoint{-0.025984in}{0.006891in}}{\pgfqpoint{-0.025984in}{0.000000in}}%
\pgfpathcurveto{\pgfqpoint{-0.025984in}{-0.006891in}}{\pgfqpoint{-0.023246in}{-0.013501in}}{\pgfqpoint{-0.018373in}{-0.018373in}}%
\pgfpathcurveto{\pgfqpoint{-0.013501in}{-0.023246in}}{\pgfqpoint{-0.006891in}{-0.025984in}}{\pgfqpoint{0.000000in}{-0.025984in}}%
\pgfpathlineto{\pgfqpoint{0.000000in}{-0.025984in}}%
\pgfpathclose%
\pgfusepath{fill}%
}%
\begin{pgfscope}%
\pgfsys@transformshift{2.817585in}{2.280616in}%
\pgfsys@useobject{currentmarker}{}%
\end{pgfscope}%
\end{pgfscope}%
\begin{pgfscope}%
\pgfpathrectangle{\pgfqpoint{2.564100in}{0.822500in}}{\pgfqpoint{0.699300in}{2.537500in}}%
\pgfusepath{clip}%
\pgfsetbuttcap%
\pgfsetroundjoin%
\definecolor{currentfill}{rgb}{1.000000,1.000000,1.000000}%
\pgfsetfillcolor{currentfill}%
\pgfsetlinewidth{0.903375pt}%
\definecolor{currentstroke}{rgb}{0.784314,0.784314,0.784314}%
\pgfsetstrokecolor{currentstroke}%
\pgfsetdash{}{0pt}%
\pgfsys@defobject{currentmarker}{\pgfqpoint{-0.025984in}{-0.025984in}}{\pgfqpoint{0.025984in}{0.025984in}}{%
\pgfpathmoveto{\pgfqpoint{0.000000in}{-0.025984in}}%
\pgfpathcurveto{\pgfqpoint{0.006891in}{-0.025984in}}{\pgfqpoint{0.013501in}{-0.023246in}}{\pgfqpoint{0.018373in}{-0.018373in}}%
\pgfpathcurveto{\pgfqpoint{0.023246in}{-0.013501in}}{\pgfqpoint{0.025984in}{-0.006891in}}{\pgfqpoint{0.025984in}{0.000000in}}%
\pgfpathcurveto{\pgfqpoint{0.025984in}{0.006891in}}{\pgfqpoint{0.023246in}{0.013501in}}{\pgfqpoint{0.018373in}{0.018373in}}%
\pgfpathcurveto{\pgfqpoint{0.013501in}{0.023246in}}{\pgfqpoint{0.006891in}{0.025984in}}{\pgfqpoint{0.000000in}{0.025984in}}%
\pgfpathcurveto{\pgfqpoint{-0.006891in}{0.025984in}}{\pgfqpoint{-0.013501in}{0.023246in}}{\pgfqpoint{-0.018373in}{0.018373in}}%
\pgfpathcurveto{\pgfqpoint{-0.023246in}{0.013501in}}{\pgfqpoint{-0.025984in}{0.006891in}}{\pgfqpoint{-0.025984in}{0.000000in}}%
\pgfpathcurveto{\pgfqpoint{-0.025984in}{-0.006891in}}{\pgfqpoint{-0.023246in}{-0.013501in}}{\pgfqpoint{-0.018373in}{-0.018373in}}%
\pgfpathcurveto{\pgfqpoint{-0.013501in}{-0.023246in}}{\pgfqpoint{-0.006891in}{-0.025984in}}{\pgfqpoint{0.000000in}{-0.025984in}}%
\pgfpathlineto{\pgfqpoint{0.000000in}{-0.025984in}}%
\pgfpathclose%
\pgfusepath{stroke,fill}%
}%
\begin{pgfscope}%
\pgfsys@transformshift{2.937406in}{2.280616in}%
\pgfsys@useobject{currentmarker}{}%
\end{pgfscope}%
\end{pgfscope}%
\begin{pgfscope}%
\pgfpathrectangle{\pgfqpoint{2.564100in}{0.822500in}}{\pgfqpoint{0.699300in}{2.537500in}}%
\pgfusepath{clip}%
\pgfsetbuttcap%
\pgfsetroundjoin%
\definecolor{currentfill}{rgb}{0.784314,0.784314,0.784314}%
\pgfsetfillcolor{currentfill}%
\pgfsetlinewidth{0.000000pt}%
\definecolor{currentstroke}{rgb}{0.784314,0.784314,0.784314}%
\pgfsetstrokecolor{currentstroke}%
\pgfsetdash{}{0pt}%
\pgfsys@defobject{currentmarker}{\pgfqpoint{-0.025984in}{-0.025984in}}{\pgfqpoint{0.025984in}{0.025984in}}{%
\pgfpathmoveto{\pgfqpoint{0.000000in}{-0.025984in}}%
\pgfpathcurveto{\pgfqpoint{0.006891in}{-0.025984in}}{\pgfqpoint{0.013501in}{-0.023246in}}{\pgfqpoint{0.018373in}{-0.018373in}}%
\pgfpathcurveto{\pgfqpoint{0.023246in}{-0.013501in}}{\pgfqpoint{0.025984in}{-0.006891in}}{\pgfqpoint{0.025984in}{0.000000in}}%
\pgfpathcurveto{\pgfqpoint{0.025984in}{0.006891in}}{\pgfqpoint{0.023246in}{0.013501in}}{\pgfqpoint{0.018373in}{0.018373in}}%
\pgfpathcurveto{\pgfqpoint{0.013501in}{0.023246in}}{\pgfqpoint{0.006891in}{0.025984in}}{\pgfqpoint{0.000000in}{0.025984in}}%
\pgfpathcurveto{\pgfqpoint{-0.006891in}{0.025984in}}{\pgfqpoint{-0.013501in}{0.023246in}}{\pgfqpoint{-0.018373in}{0.018373in}}%
\pgfpathcurveto{\pgfqpoint{-0.023246in}{0.013501in}}{\pgfqpoint{-0.025984in}{0.006891in}}{\pgfqpoint{-0.025984in}{0.000000in}}%
\pgfpathcurveto{\pgfqpoint{-0.025984in}{-0.006891in}}{\pgfqpoint{-0.023246in}{-0.013501in}}{\pgfqpoint{-0.018373in}{-0.018373in}}%
\pgfpathcurveto{\pgfqpoint{-0.013501in}{-0.023246in}}{\pgfqpoint{-0.006891in}{-0.025984in}}{\pgfqpoint{0.000000in}{-0.025984in}}%
\pgfpathlineto{\pgfqpoint{0.000000in}{-0.025984in}}%
\pgfpathclose%
\pgfusepath{fill}%
}%
\begin{pgfscope}%
\pgfsys@transformshift{3.075286in}{2.469981in}%
\pgfsys@useobject{currentmarker}{}%
\end{pgfscope}%
\end{pgfscope}%
\begin{pgfscope}%
\pgfpathrectangle{\pgfqpoint{2.564100in}{0.822500in}}{\pgfqpoint{0.699300in}{2.537500in}}%
\pgfusepath{clip}%
\pgfsetbuttcap%
\pgfsetroundjoin%
\definecolor{currentfill}{rgb}{1.000000,1.000000,1.000000}%
\pgfsetfillcolor{currentfill}%
\pgfsetlinewidth{0.903375pt}%
\definecolor{currentstroke}{rgb}{0.784314,0.784314,0.784314}%
\pgfsetstrokecolor{currentstroke}%
\pgfsetdash{}{0pt}%
\pgfsys@defobject{currentmarker}{\pgfqpoint{-0.025984in}{-0.025984in}}{\pgfqpoint{0.025984in}{0.025984in}}{%
\pgfpathmoveto{\pgfqpoint{0.000000in}{-0.025984in}}%
\pgfpathcurveto{\pgfqpoint{0.006891in}{-0.025984in}}{\pgfqpoint{0.013501in}{-0.023246in}}{\pgfqpoint{0.018373in}{-0.018373in}}%
\pgfpathcurveto{\pgfqpoint{0.023246in}{-0.013501in}}{\pgfqpoint{0.025984in}{-0.006891in}}{\pgfqpoint{0.025984in}{0.000000in}}%
\pgfpathcurveto{\pgfqpoint{0.025984in}{0.006891in}}{\pgfqpoint{0.023246in}{0.013501in}}{\pgfqpoint{0.018373in}{0.018373in}}%
\pgfpathcurveto{\pgfqpoint{0.013501in}{0.023246in}}{\pgfqpoint{0.006891in}{0.025984in}}{\pgfqpoint{0.000000in}{0.025984in}}%
\pgfpathcurveto{\pgfqpoint{-0.006891in}{0.025984in}}{\pgfqpoint{-0.013501in}{0.023246in}}{\pgfqpoint{-0.018373in}{0.018373in}}%
\pgfpathcurveto{\pgfqpoint{-0.023246in}{0.013501in}}{\pgfqpoint{-0.025984in}{0.006891in}}{\pgfqpoint{-0.025984in}{0.000000in}}%
\pgfpathcurveto{\pgfqpoint{-0.025984in}{-0.006891in}}{\pgfqpoint{-0.023246in}{-0.013501in}}{\pgfqpoint{-0.018373in}{-0.018373in}}%
\pgfpathcurveto{\pgfqpoint{-0.013501in}{-0.023246in}}{\pgfqpoint{-0.006891in}{-0.025984in}}{\pgfqpoint{0.000000in}{-0.025984in}}%
\pgfpathlineto{\pgfqpoint{0.000000in}{-0.025984in}}%
\pgfpathclose%
\pgfusepath{stroke,fill}%
}%
\begin{pgfscope}%
\pgfsys@transformshift{3.106121in}{2.469981in}%
\pgfsys@useobject{currentmarker}{}%
\end{pgfscope}%
\end{pgfscope}%
\begin{pgfscope}%
\pgfpathrectangle{\pgfqpoint{2.564100in}{0.822500in}}{\pgfqpoint{0.699300in}{2.537500in}}%
\pgfusepath{clip}%
\pgfsetbuttcap%
\pgfsetroundjoin%
\definecolor{currentfill}{rgb}{0.478431,0.478431,0.478431}%
\pgfsetfillcolor{currentfill}%
\pgfsetlinewidth{0.000000pt}%
\definecolor{currentstroke}{rgb}{0.478431,0.478431,0.478431}%
\pgfsetstrokecolor{currentstroke}%
\pgfsetdash{}{0pt}%
\pgfsys@defobject{currentmarker}{\pgfqpoint{-0.025984in}{-0.025984in}}{\pgfqpoint{0.025984in}{0.025984in}}{%
\pgfpathmoveto{\pgfqpoint{0.000000in}{-0.025984in}}%
\pgfpathcurveto{\pgfqpoint{0.006891in}{-0.025984in}}{\pgfqpoint{0.013501in}{-0.023246in}}{\pgfqpoint{0.018373in}{-0.018373in}}%
\pgfpathcurveto{\pgfqpoint{0.023246in}{-0.013501in}}{\pgfqpoint{0.025984in}{-0.006891in}}{\pgfqpoint{0.025984in}{0.000000in}}%
\pgfpathcurveto{\pgfqpoint{0.025984in}{0.006891in}}{\pgfqpoint{0.023246in}{0.013501in}}{\pgfqpoint{0.018373in}{0.018373in}}%
\pgfpathcurveto{\pgfqpoint{0.013501in}{0.023246in}}{\pgfqpoint{0.006891in}{0.025984in}}{\pgfqpoint{0.000000in}{0.025984in}}%
\pgfpathcurveto{\pgfqpoint{-0.006891in}{0.025984in}}{\pgfqpoint{-0.013501in}{0.023246in}}{\pgfqpoint{-0.018373in}{0.018373in}}%
\pgfpathcurveto{\pgfqpoint{-0.023246in}{0.013501in}}{\pgfqpoint{-0.025984in}{0.006891in}}{\pgfqpoint{-0.025984in}{0.000000in}}%
\pgfpathcurveto{\pgfqpoint{-0.025984in}{-0.006891in}}{\pgfqpoint{-0.023246in}{-0.013501in}}{\pgfqpoint{-0.018373in}{-0.018373in}}%
\pgfpathcurveto{\pgfqpoint{-0.013501in}{-0.023246in}}{\pgfqpoint{-0.006891in}{-0.025984in}}{\pgfqpoint{0.000000in}{-0.025984in}}%
\pgfpathlineto{\pgfqpoint{0.000000in}{-0.025984in}}%
\pgfpathclose%
\pgfusepath{fill}%
}%
\begin{pgfscope}%
\pgfsys@transformshift{2.832619in}{2.659347in}%
\pgfsys@useobject{currentmarker}{}%
\end{pgfscope}%
\end{pgfscope}%
\begin{pgfscope}%
\pgfpathrectangle{\pgfqpoint{2.564100in}{0.822500in}}{\pgfqpoint{0.699300in}{2.537500in}}%
\pgfusepath{clip}%
\pgfsetbuttcap%
\pgfsetroundjoin%
\definecolor{currentfill}{rgb}{1.000000,1.000000,1.000000}%
\pgfsetfillcolor{currentfill}%
\pgfsetlinewidth{0.903375pt}%
\definecolor{currentstroke}{rgb}{0.478431,0.478431,0.478431}%
\pgfsetstrokecolor{currentstroke}%
\pgfsetdash{}{0pt}%
\pgfsys@defobject{currentmarker}{\pgfqpoint{-0.025984in}{-0.025984in}}{\pgfqpoint{0.025984in}{0.025984in}}{%
\pgfpathmoveto{\pgfqpoint{0.000000in}{-0.025984in}}%
\pgfpathcurveto{\pgfqpoint{0.006891in}{-0.025984in}}{\pgfqpoint{0.013501in}{-0.023246in}}{\pgfqpoint{0.018373in}{-0.018373in}}%
\pgfpathcurveto{\pgfqpoint{0.023246in}{-0.013501in}}{\pgfqpoint{0.025984in}{-0.006891in}}{\pgfqpoint{0.025984in}{0.000000in}}%
\pgfpathcurveto{\pgfqpoint{0.025984in}{0.006891in}}{\pgfqpoint{0.023246in}{0.013501in}}{\pgfqpoint{0.018373in}{0.018373in}}%
\pgfpathcurveto{\pgfqpoint{0.013501in}{0.023246in}}{\pgfqpoint{0.006891in}{0.025984in}}{\pgfqpoint{0.000000in}{0.025984in}}%
\pgfpathcurveto{\pgfqpoint{-0.006891in}{0.025984in}}{\pgfqpoint{-0.013501in}{0.023246in}}{\pgfqpoint{-0.018373in}{0.018373in}}%
\pgfpathcurveto{\pgfqpoint{-0.023246in}{0.013501in}}{\pgfqpoint{-0.025984in}{0.006891in}}{\pgfqpoint{-0.025984in}{0.000000in}}%
\pgfpathcurveto{\pgfqpoint{-0.025984in}{-0.006891in}}{\pgfqpoint{-0.023246in}{-0.013501in}}{\pgfqpoint{-0.018373in}{-0.018373in}}%
\pgfpathcurveto{\pgfqpoint{-0.013501in}{-0.023246in}}{\pgfqpoint{-0.006891in}{-0.025984in}}{\pgfqpoint{0.000000in}{-0.025984in}}%
\pgfpathlineto{\pgfqpoint{0.000000in}{-0.025984in}}%
\pgfpathclose%
\pgfusepath{stroke,fill}%
}%
\begin{pgfscope}%
\pgfsys@transformshift{2.963773in}{2.659347in}%
\pgfsys@useobject{currentmarker}{}%
\end{pgfscope}%
\end{pgfscope}%
\begin{pgfscope}%
\pgfpathrectangle{\pgfqpoint{2.564100in}{0.822500in}}{\pgfqpoint{0.699300in}{2.537500in}}%
\pgfusepath{clip}%
\pgfsetbuttcap%
\pgfsetroundjoin%
\definecolor{currentfill}{rgb}{0.784314,0.784314,0.784314}%
\pgfsetfillcolor{currentfill}%
\pgfsetlinewidth{0.000000pt}%
\definecolor{currentstroke}{rgb}{0.784314,0.784314,0.784314}%
\pgfsetstrokecolor{currentstroke}%
\pgfsetdash{}{0pt}%
\pgfsys@defobject{currentmarker}{\pgfqpoint{-0.025984in}{-0.025984in}}{\pgfqpoint{0.025984in}{0.025984in}}{%
\pgfpathmoveto{\pgfqpoint{0.000000in}{-0.025984in}}%
\pgfpathcurveto{\pgfqpoint{0.006891in}{-0.025984in}}{\pgfqpoint{0.013501in}{-0.023246in}}{\pgfqpoint{0.018373in}{-0.018373in}}%
\pgfpathcurveto{\pgfqpoint{0.023246in}{-0.013501in}}{\pgfqpoint{0.025984in}{-0.006891in}}{\pgfqpoint{0.025984in}{0.000000in}}%
\pgfpathcurveto{\pgfqpoint{0.025984in}{0.006891in}}{\pgfqpoint{0.023246in}{0.013501in}}{\pgfqpoint{0.018373in}{0.018373in}}%
\pgfpathcurveto{\pgfqpoint{0.013501in}{0.023246in}}{\pgfqpoint{0.006891in}{0.025984in}}{\pgfqpoint{0.000000in}{0.025984in}}%
\pgfpathcurveto{\pgfqpoint{-0.006891in}{0.025984in}}{\pgfqpoint{-0.013501in}{0.023246in}}{\pgfqpoint{-0.018373in}{0.018373in}}%
\pgfpathcurveto{\pgfqpoint{-0.023246in}{0.013501in}}{\pgfqpoint{-0.025984in}{0.006891in}}{\pgfqpoint{-0.025984in}{0.000000in}}%
\pgfpathcurveto{\pgfqpoint{-0.025984in}{-0.006891in}}{\pgfqpoint{-0.023246in}{-0.013501in}}{\pgfqpoint{-0.018373in}{-0.018373in}}%
\pgfpathcurveto{\pgfqpoint{-0.013501in}{-0.023246in}}{\pgfqpoint{-0.006891in}{-0.025984in}}{\pgfqpoint{0.000000in}{-0.025984in}}%
\pgfpathlineto{\pgfqpoint{0.000000in}{-0.025984in}}%
\pgfpathclose%
\pgfusepath{fill}%
}%
\begin{pgfscope}%
\pgfsys@transformshift{3.344950in}{2.848713in}%
\pgfsys@useobject{currentmarker}{}%
\end{pgfscope}%
\end{pgfscope}%
\begin{pgfscope}%
\pgfpathrectangle{\pgfqpoint{2.564100in}{0.822500in}}{\pgfqpoint{0.699300in}{2.537500in}}%
\pgfusepath{clip}%
\pgfsetbuttcap%
\pgfsetroundjoin%
\definecolor{currentfill}{rgb}{1.000000,1.000000,1.000000}%
\pgfsetfillcolor{currentfill}%
\pgfsetlinewidth{0.903375pt}%
\definecolor{currentstroke}{rgb}{0.784314,0.784314,0.784314}%
\pgfsetstrokecolor{currentstroke}%
\pgfsetdash{}{0pt}%
\pgfsys@defobject{currentmarker}{\pgfqpoint{-0.025984in}{-0.025984in}}{\pgfqpoint{0.025984in}{0.025984in}}{%
\pgfpathmoveto{\pgfqpoint{0.000000in}{-0.025984in}}%
\pgfpathcurveto{\pgfqpoint{0.006891in}{-0.025984in}}{\pgfqpoint{0.013501in}{-0.023246in}}{\pgfqpoint{0.018373in}{-0.018373in}}%
\pgfpathcurveto{\pgfqpoint{0.023246in}{-0.013501in}}{\pgfqpoint{0.025984in}{-0.006891in}}{\pgfqpoint{0.025984in}{0.000000in}}%
\pgfpathcurveto{\pgfqpoint{0.025984in}{0.006891in}}{\pgfqpoint{0.023246in}{0.013501in}}{\pgfqpoint{0.018373in}{0.018373in}}%
\pgfpathcurveto{\pgfqpoint{0.013501in}{0.023246in}}{\pgfqpoint{0.006891in}{0.025984in}}{\pgfqpoint{0.000000in}{0.025984in}}%
\pgfpathcurveto{\pgfqpoint{-0.006891in}{0.025984in}}{\pgfqpoint{-0.013501in}{0.023246in}}{\pgfqpoint{-0.018373in}{0.018373in}}%
\pgfpathcurveto{\pgfqpoint{-0.023246in}{0.013501in}}{\pgfqpoint{-0.025984in}{0.006891in}}{\pgfqpoint{-0.025984in}{0.000000in}}%
\pgfpathcurveto{\pgfqpoint{-0.025984in}{-0.006891in}}{\pgfqpoint{-0.023246in}{-0.013501in}}{\pgfqpoint{-0.018373in}{-0.018373in}}%
\pgfpathcurveto{\pgfqpoint{-0.013501in}{-0.023246in}}{\pgfqpoint{-0.006891in}{-0.025984in}}{\pgfqpoint{0.000000in}{-0.025984in}}%
\pgfpathlineto{\pgfqpoint{0.000000in}{-0.025984in}}%
\pgfpathclose%
\pgfusepath{stroke,fill}%
}%
\begin{pgfscope}%
\pgfsys@transformshift{3.332739in}{2.848713in}%
\pgfsys@useobject{currentmarker}{}%
\end{pgfscope}%
\end{pgfscope}%
\begin{pgfscope}%
\pgfpathrectangle{\pgfqpoint{2.564100in}{0.822500in}}{\pgfqpoint{0.699300in}{2.537500in}}%
\pgfusepath{clip}%
\pgfsetbuttcap%
\pgfsetroundjoin%
\definecolor{currentfill}{rgb}{0.156863,0.356863,0.568627}%
\pgfsetfillcolor{currentfill}%
\pgfsetlinewidth{0.000000pt}%
\definecolor{currentstroke}{rgb}{0.156863,0.356863,0.568627}%
\pgfsetstrokecolor{currentstroke}%
\pgfsetdash{}{0pt}%
\pgfsys@defobject{currentmarker}{\pgfqpoint{-0.025984in}{-0.025984in}}{\pgfqpoint{0.025984in}{0.025984in}}{%
\pgfpathmoveto{\pgfqpoint{0.000000in}{-0.025984in}}%
\pgfpathcurveto{\pgfqpoint{0.006891in}{-0.025984in}}{\pgfqpoint{0.013501in}{-0.023246in}}{\pgfqpoint{0.018373in}{-0.018373in}}%
\pgfpathcurveto{\pgfqpoint{0.023246in}{-0.013501in}}{\pgfqpoint{0.025984in}{-0.006891in}}{\pgfqpoint{0.025984in}{0.000000in}}%
\pgfpathcurveto{\pgfqpoint{0.025984in}{0.006891in}}{\pgfqpoint{0.023246in}{0.013501in}}{\pgfqpoint{0.018373in}{0.018373in}}%
\pgfpathcurveto{\pgfqpoint{0.013501in}{0.023246in}}{\pgfqpoint{0.006891in}{0.025984in}}{\pgfqpoint{0.000000in}{0.025984in}}%
\pgfpathcurveto{\pgfqpoint{-0.006891in}{0.025984in}}{\pgfqpoint{-0.013501in}{0.023246in}}{\pgfqpoint{-0.018373in}{0.018373in}}%
\pgfpathcurveto{\pgfqpoint{-0.023246in}{0.013501in}}{\pgfqpoint{-0.025984in}{0.006891in}}{\pgfqpoint{-0.025984in}{0.000000in}}%
\pgfpathcurveto{\pgfqpoint{-0.025984in}{-0.006891in}}{\pgfqpoint{-0.023246in}{-0.013501in}}{\pgfqpoint{-0.018373in}{-0.018373in}}%
\pgfpathcurveto{\pgfqpoint{-0.013501in}{-0.023246in}}{\pgfqpoint{-0.006891in}{-0.025984in}}{\pgfqpoint{0.000000in}{-0.025984in}}%
\pgfpathlineto{\pgfqpoint{0.000000in}{-0.025984in}}%
\pgfpathclose%
\pgfusepath{fill}%
}%
\begin{pgfscope}%
\pgfsys@transformshift{3.154835in}{3.038078in}%
\pgfsys@useobject{currentmarker}{}%
\end{pgfscope}%
\end{pgfscope}%
\begin{pgfscope}%
\pgfpathrectangle{\pgfqpoint{2.564100in}{0.822500in}}{\pgfqpoint{0.699300in}{2.537500in}}%
\pgfusepath{clip}%
\pgfsetbuttcap%
\pgfsetroundjoin%
\definecolor{currentfill}{rgb}{1.000000,1.000000,1.000000}%
\pgfsetfillcolor{currentfill}%
\pgfsetlinewidth{0.903375pt}%
\definecolor{currentstroke}{rgb}{0.156863,0.356863,0.568627}%
\pgfsetstrokecolor{currentstroke}%
\pgfsetdash{}{0pt}%
\pgfsys@defobject{currentmarker}{\pgfqpoint{-0.025984in}{-0.025984in}}{\pgfqpoint{0.025984in}{0.025984in}}{%
\pgfpathmoveto{\pgfqpoint{0.000000in}{-0.025984in}}%
\pgfpathcurveto{\pgfqpoint{0.006891in}{-0.025984in}}{\pgfqpoint{0.013501in}{-0.023246in}}{\pgfqpoint{0.018373in}{-0.018373in}}%
\pgfpathcurveto{\pgfqpoint{0.023246in}{-0.013501in}}{\pgfqpoint{0.025984in}{-0.006891in}}{\pgfqpoint{0.025984in}{0.000000in}}%
\pgfpathcurveto{\pgfqpoint{0.025984in}{0.006891in}}{\pgfqpoint{0.023246in}{0.013501in}}{\pgfqpoint{0.018373in}{0.018373in}}%
\pgfpathcurveto{\pgfqpoint{0.013501in}{0.023246in}}{\pgfqpoint{0.006891in}{0.025984in}}{\pgfqpoint{0.000000in}{0.025984in}}%
\pgfpathcurveto{\pgfqpoint{-0.006891in}{0.025984in}}{\pgfqpoint{-0.013501in}{0.023246in}}{\pgfqpoint{-0.018373in}{0.018373in}}%
\pgfpathcurveto{\pgfqpoint{-0.023246in}{0.013501in}}{\pgfqpoint{-0.025984in}{0.006891in}}{\pgfqpoint{-0.025984in}{0.000000in}}%
\pgfpathcurveto{\pgfqpoint{-0.025984in}{-0.006891in}}{\pgfqpoint{-0.023246in}{-0.013501in}}{\pgfqpoint{-0.018373in}{-0.018373in}}%
\pgfpathcurveto{\pgfqpoint{-0.013501in}{-0.023246in}}{\pgfqpoint{-0.006891in}{-0.025984in}}{\pgfqpoint{0.000000in}{-0.025984in}}%
\pgfpathlineto{\pgfqpoint{0.000000in}{-0.025984in}}%
\pgfpathclose%
\pgfusepath{stroke,fill}%
}%
\begin{pgfscope}%
\pgfsys@transformshift{2.968293in}{3.038078in}%
\pgfsys@useobject{currentmarker}{}%
\end{pgfscope}%
\end{pgfscope}%
\begin{pgfscope}%
\pgfpathrectangle{\pgfqpoint{2.564100in}{0.822500in}}{\pgfqpoint{0.699300in}{2.537500in}}%
\pgfusepath{clip}%
\pgfsetbuttcap%
\pgfsetroundjoin%
\definecolor{currentfill}{rgb}{0.156863,0.356863,0.568627}%
\pgfsetfillcolor{currentfill}%
\pgfsetlinewidth{0.000000pt}%
\definecolor{currentstroke}{rgb}{0.156863,0.356863,0.568627}%
\pgfsetstrokecolor{currentstroke}%
\pgfsetdash{}{0pt}%
\pgfsys@defobject{currentmarker}{\pgfqpoint{-0.025984in}{-0.025984in}}{\pgfqpoint{0.025984in}{0.025984in}}{%
\pgfpathmoveto{\pgfqpoint{0.000000in}{-0.025984in}}%
\pgfpathcurveto{\pgfqpoint{0.006891in}{-0.025984in}}{\pgfqpoint{0.013501in}{-0.023246in}}{\pgfqpoint{0.018373in}{-0.018373in}}%
\pgfpathcurveto{\pgfqpoint{0.023246in}{-0.013501in}}{\pgfqpoint{0.025984in}{-0.006891in}}{\pgfqpoint{0.025984in}{0.000000in}}%
\pgfpathcurveto{\pgfqpoint{0.025984in}{0.006891in}}{\pgfqpoint{0.023246in}{0.013501in}}{\pgfqpoint{0.018373in}{0.018373in}}%
\pgfpathcurveto{\pgfqpoint{0.013501in}{0.023246in}}{\pgfqpoint{0.006891in}{0.025984in}}{\pgfqpoint{0.000000in}{0.025984in}}%
\pgfpathcurveto{\pgfqpoint{-0.006891in}{0.025984in}}{\pgfqpoint{-0.013501in}{0.023246in}}{\pgfqpoint{-0.018373in}{0.018373in}}%
\pgfpathcurveto{\pgfqpoint{-0.023246in}{0.013501in}}{\pgfqpoint{-0.025984in}{0.006891in}}{\pgfqpoint{-0.025984in}{0.000000in}}%
\pgfpathcurveto{\pgfqpoint{-0.025984in}{-0.006891in}}{\pgfqpoint{-0.023246in}{-0.013501in}}{\pgfqpoint{-0.018373in}{-0.018373in}}%
\pgfpathcurveto{\pgfqpoint{-0.013501in}{-0.023246in}}{\pgfqpoint{-0.006891in}{-0.025984in}}{\pgfqpoint{0.000000in}{-0.025984in}}%
\pgfpathlineto{\pgfqpoint{0.000000in}{-0.025984in}}%
\pgfpathclose%
\pgfusepath{fill}%
}%
\begin{pgfscope}%
\pgfsys@transformshift{3.160923in}{3.227444in}%
\pgfsys@useobject{currentmarker}{}%
\end{pgfscope}%
\end{pgfscope}%
\begin{pgfscope}%
\pgfpathrectangle{\pgfqpoint{2.564100in}{0.822500in}}{\pgfqpoint{0.699300in}{2.537500in}}%
\pgfusepath{clip}%
\pgfsetbuttcap%
\pgfsetroundjoin%
\definecolor{currentfill}{rgb}{1.000000,1.000000,1.000000}%
\pgfsetfillcolor{currentfill}%
\pgfsetlinewidth{0.903375pt}%
\definecolor{currentstroke}{rgb}{0.156863,0.356863,0.568627}%
\pgfsetstrokecolor{currentstroke}%
\pgfsetdash{}{0pt}%
\pgfsys@defobject{currentmarker}{\pgfqpoint{-0.025984in}{-0.025984in}}{\pgfqpoint{0.025984in}{0.025984in}}{%
\pgfpathmoveto{\pgfqpoint{0.000000in}{-0.025984in}}%
\pgfpathcurveto{\pgfqpoint{0.006891in}{-0.025984in}}{\pgfqpoint{0.013501in}{-0.023246in}}{\pgfqpoint{0.018373in}{-0.018373in}}%
\pgfpathcurveto{\pgfqpoint{0.023246in}{-0.013501in}}{\pgfqpoint{0.025984in}{-0.006891in}}{\pgfqpoint{0.025984in}{0.000000in}}%
\pgfpathcurveto{\pgfqpoint{0.025984in}{0.006891in}}{\pgfqpoint{0.023246in}{0.013501in}}{\pgfqpoint{0.018373in}{0.018373in}}%
\pgfpathcurveto{\pgfqpoint{0.013501in}{0.023246in}}{\pgfqpoint{0.006891in}{0.025984in}}{\pgfqpoint{0.000000in}{0.025984in}}%
\pgfpathcurveto{\pgfqpoint{-0.006891in}{0.025984in}}{\pgfqpoint{-0.013501in}{0.023246in}}{\pgfqpoint{-0.018373in}{0.018373in}}%
\pgfpathcurveto{\pgfqpoint{-0.023246in}{0.013501in}}{\pgfqpoint{-0.025984in}{0.006891in}}{\pgfqpoint{-0.025984in}{0.000000in}}%
\pgfpathcurveto{\pgfqpoint{-0.025984in}{-0.006891in}}{\pgfqpoint{-0.023246in}{-0.013501in}}{\pgfqpoint{-0.018373in}{-0.018373in}}%
\pgfpathcurveto{\pgfqpoint{-0.013501in}{-0.023246in}}{\pgfqpoint{-0.006891in}{-0.025984in}}{\pgfqpoint{0.000000in}{-0.025984in}}%
\pgfpathlineto{\pgfqpoint{0.000000in}{-0.025984in}}%
\pgfpathclose%
\pgfusepath{stroke,fill}%
}%
\begin{pgfscope}%
\pgfsys@transformshift{2.986479in}{3.227444in}%
\pgfsys@useobject{currentmarker}{}%
\end{pgfscope}%
\end{pgfscope}%
\begin{pgfscope}%
\definecolor{textcolor}{rgb}{0.549020,0.549020,0.549020}%
\pgfsetstrokecolor{textcolor}%
\pgfsetfillcolor{textcolor}%
\pgftext[x=3.296700in,y=3.447500in,right,top]{\color{textcolor}\sffamily\fontsize{5.500000}{6.600000}\selectfont OLD DATA}%
\end{pgfscope}%
\begin{pgfscope}%
\pgfsetrectcap%
\pgfsetroundjoin%
\pgfsetlinewidth{0.000000pt}%
\definecolor{currentstroke}{rgb}{0.101961,0.101961,0.101961}%
\pgfsetstrokecolor{currentstroke}%
\pgfsetdash{}{0pt}%
\pgfpathmoveto{\pgfqpoint{1.205430in}{0.122932in}}%
\pgfpathlineto{\pgfqpoint{1.302652in}{0.122932in}}%
\pgfpathlineto{\pgfqpoint{1.399874in}{0.122932in}}%
\pgfusepath{}%
\end{pgfscope}%
\begin{pgfscope}%
\pgfsetbuttcap%
\pgfsetroundjoin%
\definecolor{currentfill}{rgb}{0.101961,0.101961,0.101961}%
\pgfsetfillcolor{currentfill}%
\pgfsetlinewidth{1.003750pt}%
\definecolor{currentstroke}{rgb}{0.101961,0.101961,0.101961}%
\pgfsetstrokecolor{currentstroke}%
\pgfsetdash{}{0pt}%
\pgfsys@defobject{currentmarker}{\pgfqpoint{-0.025000in}{-0.025000in}}{\pgfqpoint{0.025000in}{0.025000in}}{%
\pgfpathmoveto{\pgfqpoint{0.000000in}{-0.025000in}}%
\pgfpathcurveto{\pgfqpoint{0.006630in}{-0.025000in}}{\pgfqpoint{0.012989in}{-0.022366in}}{\pgfqpoint{0.017678in}{-0.017678in}}%
\pgfpathcurveto{\pgfqpoint{0.022366in}{-0.012989in}}{\pgfqpoint{0.025000in}{-0.006630in}}{\pgfqpoint{0.025000in}{0.000000in}}%
\pgfpathcurveto{\pgfqpoint{0.025000in}{0.006630in}}{\pgfqpoint{0.022366in}{0.012989in}}{\pgfqpoint{0.017678in}{0.017678in}}%
\pgfpathcurveto{\pgfqpoint{0.012989in}{0.022366in}}{\pgfqpoint{0.006630in}{0.025000in}}{\pgfqpoint{0.000000in}{0.025000in}}%
\pgfpathcurveto{\pgfqpoint{-0.006630in}{0.025000in}}{\pgfqpoint{-0.012989in}{0.022366in}}{\pgfqpoint{-0.017678in}{0.017678in}}%
\pgfpathcurveto{\pgfqpoint{-0.022366in}{0.012989in}}{\pgfqpoint{-0.025000in}{0.006630in}}{\pgfqpoint{-0.025000in}{0.000000in}}%
\pgfpathcurveto{\pgfqpoint{-0.025000in}{-0.006630in}}{\pgfqpoint{-0.022366in}{-0.012989in}}{\pgfqpoint{-0.017678in}{-0.017678in}}%
\pgfpathcurveto{\pgfqpoint{-0.012989in}{-0.022366in}}{\pgfqpoint{-0.006630in}{-0.025000in}}{\pgfqpoint{0.000000in}{-0.025000in}}%
\pgfpathlineto{\pgfqpoint{0.000000in}{-0.025000in}}%
\pgfpathclose%
\pgfusepath{stroke,fill}%
}%
\begin{pgfscope}%
\pgfsys@transformshift{1.302652in}{0.122932in}%
\pgfsys@useobject{currentmarker}{}%
\end{pgfscope}%
\end{pgfscope}%
\begin{pgfscope}%
\definecolor{textcolor}{rgb}{0.101961,0.101961,0.101961}%
\pgfsetstrokecolor{textcolor}%
\pgfsetfillcolor{textcolor}%
\pgftext[x=1.409596in,y=0.088904in,left,base]{\color{textcolor}\sffamily\fontsize{7.000000}{8.400000}\selectfont 1x memory}%
\end{pgfscope}%
\begin{pgfscope}%
\pgfsetrectcap%
\pgfsetroundjoin%
\pgfsetlinewidth{0.000000pt}%
\definecolor{currentstroke}{rgb}{0.101961,0.101961,0.101961}%
\pgfsetstrokecolor{currentstroke}%
\pgfsetdash{}{0pt}%
\pgfpathmoveto{\pgfqpoint{2.030745in}{0.122932in}}%
\pgfpathlineto{\pgfqpoint{2.127967in}{0.122932in}}%
\pgfpathlineto{\pgfqpoint{2.225189in}{0.122932in}}%
\pgfusepath{}%
\end{pgfscope}%
\begin{pgfscope}%
\pgfsetbuttcap%
\pgfsetroundjoin%
\definecolor{currentfill}{rgb}{1.000000,1.000000,1.000000}%
\pgfsetfillcolor{currentfill}%
\pgfsetlinewidth{1.003750pt}%
\definecolor{currentstroke}{rgb}{0.101961,0.101961,0.101961}%
\pgfsetstrokecolor{currentstroke}%
\pgfsetdash{}{0pt}%
\pgfsys@defobject{currentmarker}{\pgfqpoint{-0.025000in}{-0.025000in}}{\pgfqpoint{0.025000in}{0.025000in}}{%
\pgfpathmoveto{\pgfqpoint{0.000000in}{-0.025000in}}%
\pgfpathcurveto{\pgfqpoint{0.006630in}{-0.025000in}}{\pgfqpoint{0.012989in}{-0.022366in}}{\pgfqpoint{0.017678in}{-0.017678in}}%
\pgfpathcurveto{\pgfqpoint{0.022366in}{-0.012989in}}{\pgfqpoint{0.025000in}{-0.006630in}}{\pgfqpoint{0.025000in}{0.000000in}}%
\pgfpathcurveto{\pgfqpoint{0.025000in}{0.006630in}}{\pgfqpoint{0.022366in}{0.012989in}}{\pgfqpoint{0.017678in}{0.017678in}}%
\pgfpathcurveto{\pgfqpoint{0.012989in}{0.022366in}}{\pgfqpoint{0.006630in}{0.025000in}}{\pgfqpoint{0.000000in}{0.025000in}}%
\pgfpathcurveto{\pgfqpoint{-0.006630in}{0.025000in}}{\pgfqpoint{-0.012989in}{0.022366in}}{\pgfqpoint{-0.017678in}{0.017678in}}%
\pgfpathcurveto{\pgfqpoint{-0.022366in}{0.012989in}}{\pgfqpoint{-0.025000in}{0.006630in}}{\pgfqpoint{-0.025000in}{0.000000in}}%
\pgfpathcurveto{\pgfqpoint{-0.025000in}{-0.006630in}}{\pgfqpoint{-0.022366in}{-0.012989in}}{\pgfqpoint{-0.017678in}{-0.017678in}}%
\pgfpathcurveto{\pgfqpoint{-0.012989in}{-0.022366in}}{\pgfqpoint{-0.006630in}{-0.025000in}}{\pgfqpoint{0.000000in}{-0.025000in}}%
\pgfpathlineto{\pgfqpoint{0.000000in}{-0.025000in}}%
\pgfpathclose%
\pgfusepath{stroke,fill}%
}%
\begin{pgfscope}%
\pgfsys@transformshift{2.127967in}{0.122932in}%
\pgfsys@useobject{currentmarker}{}%
\end{pgfscope}%
\end{pgfscope}%
\begin{pgfscope}%
\definecolor{textcolor}{rgb}{0.101961,0.101961,0.101961}%
\pgfsetstrokecolor{textcolor}%
\pgfsetfillcolor{textcolor}%
\pgftext[x=2.234911in,y=0.088904in,left,base]{\color{textcolor}\sffamily\fontsize{7.000000}{8.400000}\selectfont 0.5x memory}%
\end{pgfscope}%
\end{pgfpicture}%
\makeatother%
\endgroup%
}
\caption{[PLACEHOLDER: OLD DATA, rung 3 stand-in and COVEX dev seed 4, n16, 3.0 s, seeds 4-7 baselines] Effect of halving memory (1x filled, 0.5x open), left: pooled correct-in-time fraction over robot and single-shot units, right: dashcam per-episode alert attainment. Means over seeds 4-7 where available (COVEX dev: seed 4 only). $\dagger$: 0.5x cell uses the older tag (kv0.5 or kv0.5r), not kv0.5r\_h0.5, because no kv0.5r\_h0.5 cell exists for that arm. Arms without a 0.5x cell are omitted.}
\label{fig:F08_memory}
\end{figure}
\clearpage
\begin{figure}[t]
\centering
{\fontfamily{LinuxBiolinumT}\selectfont
%% Creator: Matplotlib, PGF backend
%%
%% To include the figure in your LaTeX document, write
%%   \input{<filename>.pgf}
%%
%% Make sure the required packages are loaded in your preamble
%%   \usepackage{pgf}
%%
%% Also ensure that all the required font packages are loaded; for instance,
%% the lmodern package is sometimes necessary when using math font.
%%   \usepackage{lmodern}
%%
%% Figures using additional raster images can only be included by \input if
%% they are in the same directory as the main LaTeX file. For loading figures
%% from other directories you can use the `import` package
%%   \usepackage{import}
%%
%% and then include the figures with
%%   \import{<path to file>}{<filename>.pgf}
%%
%% Matplotlib used the following preamble
%%
\begingroup%
\makeatletter%
\begin{pgfpicture}%
\pgfpathrectangle{\pgfpointorigin}{\pgfqpoint{3.330000in}{3.900000in}}%
\pgfusepath{use as bounding box, clip}%
\begin{pgfscope}%
\pgfsetbuttcap%
\pgfsetmiterjoin%
\definecolor{currentfill}{rgb}{1.000000,1.000000,1.000000}%
\pgfsetfillcolor{currentfill}%
\pgfsetlinewidth{0.000000pt}%
\definecolor{currentstroke}{rgb}{1.000000,1.000000,1.000000}%
\pgfsetstrokecolor{currentstroke}%
\pgfsetdash{}{0pt}%
\pgfpathmoveto{\pgfqpoint{0.000000in}{0.000000in}}%
\pgfpathlineto{\pgfqpoint{3.330000in}{0.000000in}}%
\pgfpathlineto{\pgfqpoint{3.330000in}{3.900000in}}%
\pgfpathlineto{\pgfqpoint{0.000000in}{3.900000in}}%
\pgfpathlineto{\pgfqpoint{0.000000in}{0.000000in}}%
\pgfpathclose%
\pgfusepath{fill}%
\end{pgfscope}%
\begin{pgfscope}%
\pgfsetbuttcap%
\pgfsetmiterjoin%
\definecolor{currentfill}{rgb}{1.000000,1.000000,1.000000}%
\pgfsetfillcolor{currentfill}%
\pgfsetlinewidth{0.000000pt}%
\definecolor{currentstroke}{rgb}{0.000000,0.000000,0.000000}%
\pgfsetstrokecolor{currentstroke}%
\pgfsetstrokeopacity{0.000000}%
\pgfsetdash{}{0pt}%
\pgfpathmoveto{\pgfqpoint{0.616050in}{2.359500in}}%
\pgfpathlineto{\pgfqpoint{3.263400in}{2.359500in}}%
\pgfpathlineto{\pgfqpoint{3.263400in}{3.744000in}}%
\pgfpathlineto{\pgfqpoint{0.616050in}{3.744000in}}%
\pgfpathlineto{\pgfqpoint{0.616050in}{2.359500in}}%
\pgfpathclose%
\pgfusepath{fill}%
\end{pgfscope}%
\begin{pgfscope}%
\pgfpathrectangle{\pgfqpoint{0.616050in}{2.359500in}}{\pgfqpoint{2.647350in}{1.384500in}}%
\pgfusepath{clip}%
\pgfsetrectcap%
\pgfsetroundjoin%
\pgfsetlinewidth{0.401500pt}%
\definecolor{currentstroke}{rgb}{0.901961,0.901961,0.901961}%
\pgfsetstrokecolor{currentstroke}%
\pgfsetdash{}{0pt}%
\pgfpathmoveto{\pgfqpoint{0.760451in}{2.359500in}}%
\pgfpathlineto{\pgfqpoint{0.760451in}{3.744000in}}%
\pgfusepath{stroke}%
\end{pgfscope}%
\begin{pgfscope}%
\pgfpathrectangle{\pgfqpoint{0.616050in}{2.359500in}}{\pgfqpoint{2.647350in}{1.384500in}}%
\pgfusepath{clip}%
\pgfsetrectcap%
\pgfsetroundjoin%
\pgfsetlinewidth{0.401500pt}%
\definecolor{currentstroke}{rgb}{0.901961,0.901961,0.901961}%
\pgfsetstrokecolor{currentstroke}%
\pgfsetdash{}{0pt}%
\pgfpathmoveto{\pgfqpoint{1.145520in}{2.359500in}}%
\pgfpathlineto{\pgfqpoint{1.145520in}{3.744000in}}%
\pgfusepath{stroke}%
\end{pgfscope}%
\begin{pgfscope}%
\pgfpathrectangle{\pgfqpoint{0.616050in}{2.359500in}}{\pgfqpoint{2.647350in}{1.384500in}}%
\pgfusepath{clip}%
\pgfsetrectcap%
\pgfsetroundjoin%
\pgfsetlinewidth{0.401500pt}%
\definecolor{currentstroke}{rgb}{0.901961,0.901961,0.901961}%
\pgfsetstrokecolor{currentstroke}%
\pgfsetdash{}{0pt}%
\pgfpathmoveto{\pgfqpoint{1.530589in}{2.359500in}}%
\pgfpathlineto{\pgfqpoint{1.530589in}{3.744000in}}%
\pgfusepath{stroke}%
\end{pgfscope}%
\begin{pgfscope}%
\pgfpathrectangle{\pgfqpoint{0.616050in}{2.359500in}}{\pgfqpoint{2.647350in}{1.384500in}}%
\pgfusepath{clip}%
\pgfsetrectcap%
\pgfsetroundjoin%
\pgfsetlinewidth{0.401500pt}%
\definecolor{currentstroke}{rgb}{0.901961,0.901961,0.901961}%
\pgfsetstrokecolor{currentstroke}%
\pgfsetdash{}{0pt}%
\pgfpathmoveto{\pgfqpoint{1.915658in}{2.359500in}}%
\pgfpathlineto{\pgfqpoint{1.915658in}{3.744000in}}%
\pgfusepath{stroke}%
\end{pgfscope}%
\begin{pgfscope}%
\pgfpathrectangle{\pgfqpoint{0.616050in}{2.359500in}}{\pgfqpoint{2.647350in}{1.384500in}}%
\pgfusepath{clip}%
\pgfsetrectcap%
\pgfsetroundjoin%
\pgfsetlinewidth{0.401500pt}%
\definecolor{currentstroke}{rgb}{0.901961,0.901961,0.901961}%
\pgfsetstrokecolor{currentstroke}%
\pgfsetdash{}{0pt}%
\pgfpathmoveto{\pgfqpoint{2.300727in}{2.359500in}}%
\pgfpathlineto{\pgfqpoint{2.300727in}{3.744000in}}%
\pgfusepath{stroke}%
\end{pgfscope}%
\begin{pgfscope}%
\pgfpathrectangle{\pgfqpoint{0.616050in}{2.359500in}}{\pgfqpoint{2.647350in}{1.384500in}}%
\pgfusepath{clip}%
\pgfsetrectcap%
\pgfsetroundjoin%
\pgfsetlinewidth{0.401500pt}%
\definecolor{currentstroke}{rgb}{0.901961,0.901961,0.901961}%
\pgfsetstrokecolor{currentstroke}%
\pgfsetdash{}{0pt}%
\pgfpathmoveto{\pgfqpoint{3.070865in}{2.359500in}}%
\pgfpathlineto{\pgfqpoint{3.070865in}{3.744000in}}%
\pgfusepath{stroke}%
\end{pgfscope}%
\begin{pgfscope}%
\pgfpathrectangle{\pgfqpoint{0.616050in}{2.359500in}}{\pgfqpoint{2.647350in}{1.384500in}}%
\pgfusepath{clip}%
\pgfsetrectcap%
\pgfsetroundjoin%
\pgfsetlinewidth{0.401500pt}%
\definecolor{currentstroke}{rgb}{0.901961,0.901961,0.901961}%
\pgfsetstrokecolor{currentstroke}%
\pgfsetdash{}{0pt}%
\pgfpathmoveto{\pgfqpoint{0.616050in}{2.485364in}}%
\pgfpathlineto{\pgfqpoint{3.263400in}{2.485364in}}%
\pgfusepath{stroke}%
\end{pgfscope}%
\begin{pgfscope}%
\definecolor{textcolor}{rgb}{0.101961,0.101961,0.101961}%
\pgfsetstrokecolor{textcolor}%
\pgfsetfillcolor{textcolor}%
\pgftext[x=0.435445in, y=2.451606in, left, base]{\color{textcolor}\sffamily\fontsize{7.000000}{8.400000}\selectfont 0.2}%
\end{pgfscope}%
\begin{pgfscope}%
\pgfpathrectangle{\pgfqpoint{0.616050in}{2.359500in}}{\pgfqpoint{2.647350in}{1.384500in}}%
\pgfusepath{clip}%
\pgfsetrectcap%
\pgfsetroundjoin%
\pgfsetlinewidth{0.401500pt}%
\definecolor{currentstroke}{rgb}{0.901961,0.901961,0.901961}%
\pgfsetstrokecolor{currentstroke}%
\pgfsetdash{}{0pt}%
\pgfpathmoveto{\pgfqpoint{0.616050in}{2.737091in}}%
\pgfpathlineto{\pgfqpoint{3.263400in}{2.737091in}}%
\pgfusepath{stroke}%
\end{pgfscope}%
\begin{pgfscope}%
\definecolor{textcolor}{rgb}{0.101961,0.101961,0.101961}%
\pgfsetstrokecolor{textcolor}%
\pgfsetfillcolor{textcolor}%
\pgftext[x=0.435445in, y=2.703333in, left, base]{\color{textcolor}\sffamily\fontsize{7.000000}{8.400000}\selectfont 0.3}%
\end{pgfscope}%
\begin{pgfscope}%
\pgfpathrectangle{\pgfqpoint{0.616050in}{2.359500in}}{\pgfqpoint{2.647350in}{1.384500in}}%
\pgfusepath{clip}%
\pgfsetrectcap%
\pgfsetroundjoin%
\pgfsetlinewidth{0.401500pt}%
\definecolor{currentstroke}{rgb}{0.901961,0.901961,0.901961}%
\pgfsetstrokecolor{currentstroke}%
\pgfsetdash{}{0pt}%
\pgfpathmoveto{\pgfqpoint{0.616050in}{2.988818in}}%
\pgfpathlineto{\pgfqpoint{3.263400in}{2.988818in}}%
\pgfusepath{stroke}%
\end{pgfscope}%
\begin{pgfscope}%
\definecolor{textcolor}{rgb}{0.101961,0.101961,0.101961}%
\pgfsetstrokecolor{textcolor}%
\pgfsetfillcolor{textcolor}%
\pgftext[x=0.435445in, y=2.955061in, left, base]{\color{textcolor}\sffamily\fontsize{7.000000}{8.400000}\selectfont 0.4}%
\end{pgfscope}%
\begin{pgfscope}%
\pgfpathrectangle{\pgfqpoint{0.616050in}{2.359500in}}{\pgfqpoint{2.647350in}{1.384500in}}%
\pgfusepath{clip}%
\pgfsetrectcap%
\pgfsetroundjoin%
\pgfsetlinewidth{0.401500pt}%
\definecolor{currentstroke}{rgb}{0.901961,0.901961,0.901961}%
\pgfsetstrokecolor{currentstroke}%
\pgfsetdash{}{0pt}%
\pgfpathmoveto{\pgfqpoint{0.616050in}{3.240545in}}%
\pgfpathlineto{\pgfqpoint{3.263400in}{3.240545in}}%
\pgfusepath{stroke}%
\end{pgfscope}%
\begin{pgfscope}%
\definecolor{textcolor}{rgb}{0.101961,0.101961,0.101961}%
\pgfsetstrokecolor{textcolor}%
\pgfsetfillcolor{textcolor}%
\pgftext[x=0.435445in, y=3.206788in, left, base]{\color{textcolor}\sffamily\fontsize{7.000000}{8.400000}\selectfont 0.5}%
\end{pgfscope}%
\begin{pgfscope}%
\pgfpathrectangle{\pgfqpoint{0.616050in}{2.359500in}}{\pgfqpoint{2.647350in}{1.384500in}}%
\pgfusepath{clip}%
\pgfsetrectcap%
\pgfsetroundjoin%
\pgfsetlinewidth{0.401500pt}%
\definecolor{currentstroke}{rgb}{0.901961,0.901961,0.901961}%
\pgfsetstrokecolor{currentstroke}%
\pgfsetdash{}{0pt}%
\pgfpathmoveto{\pgfqpoint{0.616050in}{3.492273in}}%
\pgfpathlineto{\pgfqpoint{3.263400in}{3.492273in}}%
\pgfusepath{stroke}%
\end{pgfscope}%
\begin{pgfscope}%
\definecolor{textcolor}{rgb}{0.101961,0.101961,0.101961}%
\pgfsetstrokecolor{textcolor}%
\pgfsetfillcolor{textcolor}%
\pgftext[x=0.435445in, y=3.458515in, left, base]{\color{textcolor}\sffamily\fontsize{7.000000}{8.400000}\selectfont 0.6}%
\end{pgfscope}%
\begin{pgfscope}%
\pgfpathrectangle{\pgfqpoint{0.616050in}{2.359500in}}{\pgfqpoint{2.647350in}{1.384500in}}%
\pgfusepath{clip}%
\pgfsetrectcap%
\pgfsetroundjoin%
\pgfsetlinewidth{0.401500pt}%
\definecolor{currentstroke}{rgb}{0.901961,0.901961,0.901961}%
\pgfsetstrokecolor{currentstroke}%
\pgfsetdash{}{0pt}%
\pgfpathmoveto{\pgfqpoint{0.616050in}{3.744000in}}%
\pgfpathlineto{\pgfqpoint{3.263400in}{3.744000in}}%
\pgfusepath{stroke}%
\end{pgfscope}%
\begin{pgfscope}%
\definecolor{textcolor}{rgb}{0.101961,0.101961,0.101961}%
\pgfsetstrokecolor{textcolor}%
\pgfsetfillcolor{textcolor}%
\pgftext[x=0.435445in, y=3.710242in, left, base]{\color{textcolor}\sffamily\fontsize{7.000000}{8.400000}\selectfont 0.7}%
\end{pgfscope}%
\begin{pgfscope}%
\definecolor{textcolor}{rgb}{0.000000,0.000000,0.000000}%
\pgfsetstrokecolor{textcolor}%
\pgfsetfillcolor{textcolor}%
\pgftext[x=0.261062in, y=2.740753in, left, base,rotate=90.000000]{\color{textcolor}\sffamily\fontsize{7.500000}{9.000000}\selectfont Pooled correct}%
\end{pgfscope}%
\begin{pgfscope}%
\definecolor{textcolor}{rgb}{0.000000,0.000000,0.000000}%
\pgfsetstrokecolor{textcolor}%
\pgfsetfillcolor{textcolor}%
\pgftext[x=0.360985in, y=2.900598in, left, base,rotate=90.000000]{\color{textcolor}\sffamily\fontsize{7.500000}{9.000000}\selectfont in time}%
\end{pgfscope}%
\begin{pgfscope}%
\pgfpathrectangle{\pgfqpoint{0.616050in}{2.359500in}}{\pgfqpoint{2.647350in}{1.384500in}}%
\pgfusepath{clip}%
\pgfsetrectcap%
\pgfsetroundjoin%
\pgfsetlinewidth{0.903375pt}%
\definecolor{currentstroke}{rgb}{0.478431,0.305882,0.611765}%
\pgfsetstrokecolor{currentstroke}%
\pgfsetdash{}{0pt}%
\pgfpathmoveto{\pgfqpoint{0.760451in}{3.240545in}}%
\pgfpathlineto{\pgfqpoint{1.145520in}{3.034840in}}%
\pgfpathlineto{\pgfqpoint{1.530589in}{3.156228in}}%
\pgfpathlineto{\pgfqpoint{1.915658in}{3.075377in}}%
\pgfusepath{stroke}%
\end{pgfscope}%
\begin{pgfscope}%
\pgfpathrectangle{\pgfqpoint{0.616050in}{2.359500in}}{\pgfqpoint{2.647350in}{1.384500in}}%
\pgfusepath{clip}%
\pgfsetbuttcap%
\pgfsetmiterjoin%
\definecolor{currentfill}{rgb}{0.478431,0.305882,0.611765}%
\pgfsetfillcolor{currentfill}%
\pgfsetlinewidth{0.702625pt}%
\definecolor{currentstroke}{rgb}{0.478431,0.305882,0.611765}%
\pgfsetstrokecolor{currentstroke}%
\pgfsetdash{}{0pt}%
\pgfsys@defobject{currentmarker}{\pgfqpoint{-0.022222in}{-0.022222in}}{\pgfqpoint{0.022222in}{0.022222in}}{%
\pgfpathmoveto{\pgfqpoint{-0.011111in}{-0.022222in}}%
\pgfpathlineto{\pgfqpoint{0.000000in}{-0.011111in}}%
\pgfpathlineto{\pgfqpoint{0.011111in}{-0.022222in}}%
\pgfpathlineto{\pgfqpoint{0.022222in}{-0.011111in}}%
\pgfpathlineto{\pgfqpoint{0.011111in}{0.000000in}}%
\pgfpathlineto{\pgfqpoint{0.022222in}{0.011111in}}%
\pgfpathlineto{\pgfqpoint{0.011111in}{0.022222in}}%
\pgfpathlineto{\pgfqpoint{0.000000in}{0.011111in}}%
\pgfpathlineto{\pgfqpoint{-0.011111in}{0.022222in}}%
\pgfpathlineto{\pgfqpoint{-0.022222in}{0.011111in}}%
\pgfpathlineto{\pgfqpoint{-0.011111in}{0.000000in}}%
\pgfpathlineto{\pgfqpoint{-0.022222in}{-0.011111in}}%
\pgfpathlineto{\pgfqpoint{-0.011111in}{-0.022222in}}%
\pgfpathclose%
\pgfusepath{stroke,fill}%
}%
\begin{pgfscope}%
\pgfsys@transformshift{0.760451in}{3.240545in}%
\pgfsys@useobject{currentmarker}{}%
\end{pgfscope}%
\begin{pgfscope}%
\pgfsys@transformshift{1.145520in}{3.034840in}%
\pgfsys@useobject{currentmarker}{}%
\end{pgfscope}%
\begin{pgfscope}%
\pgfsys@transformshift{1.530589in}{3.156228in}%
\pgfsys@useobject{currentmarker}{}%
\end{pgfscope}%
\begin{pgfscope}%
\pgfsys@transformshift{1.915658in}{3.075377in}%
\pgfsys@useobject{currentmarker}{}%
\end{pgfscope}%
\end{pgfscope}%
\begin{pgfscope}%
\pgfpathrectangle{\pgfqpoint{0.616050in}{2.359500in}}{\pgfqpoint{2.647350in}{1.384500in}}%
\pgfusepath{clip}%
\pgfsetrectcap%
\pgfsetroundjoin%
\pgfsetlinewidth{0.903375pt}%
\definecolor{currentstroke}{rgb}{0.784314,0.784314,0.784314}%
\pgfsetstrokecolor{currentstroke}%
\pgfsetdash{}{0pt}%
\pgfpathmoveto{\pgfqpoint{0.760451in}{3.125376in}}%
\pgfpathlineto{\pgfqpoint{1.145520in}{3.076679in}}%
\pgfpathlineto{\pgfqpoint{1.530589in}{3.243575in}}%
\pgfusepath{stroke}%
\end{pgfscope}%
\begin{pgfscope}%
\pgfpathrectangle{\pgfqpoint{0.616050in}{2.359500in}}{\pgfqpoint{2.647350in}{1.384500in}}%
\pgfusepath{clip}%
\pgfsetbuttcap%
\pgfsetmiterjoin%
\definecolor{currentfill}{rgb}{0.784314,0.784314,0.784314}%
\pgfsetfillcolor{currentfill}%
\pgfsetlinewidth{0.702625pt}%
\definecolor{currentstroke}{rgb}{0.784314,0.784314,0.784314}%
\pgfsetstrokecolor{currentstroke}%
\pgfsetdash{}{0pt}%
\pgfsys@defobject{currentmarker}{\pgfqpoint{-0.022222in}{-0.022222in}}{\pgfqpoint{0.022222in}{0.022222in}}{%
\pgfpathmoveto{\pgfqpoint{0.000000in}{0.022222in}}%
\pgfpathlineto{\pgfqpoint{-0.022222in}{-0.022222in}}%
\pgfpathlineto{\pgfqpoint{0.022222in}{-0.022222in}}%
\pgfpathlineto{\pgfqpoint{0.000000in}{0.022222in}}%
\pgfpathclose%
\pgfusepath{stroke,fill}%
}%
\begin{pgfscope}%
\pgfsys@transformshift{0.760451in}{3.125376in}%
\pgfsys@useobject{currentmarker}{}%
\end{pgfscope}%
\begin{pgfscope}%
\pgfsys@transformshift{1.145520in}{3.076679in}%
\pgfsys@useobject{currentmarker}{}%
\end{pgfscope}%
\begin{pgfscope}%
\pgfsys@transformshift{1.530589in}{3.243575in}%
\pgfsys@useobject{currentmarker}{}%
\end{pgfscope}%
\end{pgfscope}%
\begin{pgfscope}%
\pgfpathrectangle{\pgfqpoint{0.616050in}{2.359500in}}{\pgfqpoint{2.647350in}{1.384500in}}%
\pgfusepath{clip}%
\pgfsetrectcap%
\pgfsetroundjoin%
\pgfsetlinewidth{0.903375pt}%
\definecolor{currentstroke}{rgb}{0.784314,0.784314,0.784314}%
\pgfsetstrokecolor{currentstroke}%
\pgfsetdash{}{0pt}%
\pgfpathmoveto{\pgfqpoint{0.760451in}{3.232319in}}%
\pgfpathlineto{\pgfqpoint{1.530589in}{3.210962in}}%
\pgfusepath{stroke}%
\end{pgfscope}%
\begin{pgfscope}%
\pgfpathrectangle{\pgfqpoint{0.616050in}{2.359500in}}{\pgfqpoint{2.647350in}{1.384500in}}%
\pgfusepath{clip}%
\pgfsetbuttcap%
\pgfsetmiterjoin%
\definecolor{currentfill}{rgb}{0.784314,0.784314,0.784314}%
\pgfsetfillcolor{currentfill}%
\pgfsetlinewidth{0.702625pt}%
\definecolor{currentstroke}{rgb}{0.784314,0.784314,0.784314}%
\pgfsetstrokecolor{currentstroke}%
\pgfsetdash{}{0pt}%
\pgfsys@defobject{currentmarker}{\pgfqpoint{-0.022222in}{-0.022222in}}{\pgfqpoint{0.022222in}{0.022222in}}{%
\pgfpathmoveto{\pgfqpoint{-0.000000in}{-0.022222in}}%
\pgfpathlineto{\pgfqpoint{0.022222in}{0.022222in}}%
\pgfpathlineto{\pgfqpoint{-0.022222in}{0.022222in}}%
\pgfpathlineto{\pgfqpoint{-0.000000in}{-0.022222in}}%
\pgfpathclose%
\pgfusepath{stroke,fill}%
}%
\begin{pgfscope}%
\pgfsys@transformshift{0.760451in}{3.232319in}%
\pgfsys@useobject{currentmarker}{}%
\end{pgfscope}%
\begin{pgfscope}%
\pgfsys@transformshift{1.530589in}{3.210962in}%
\pgfsys@useobject{currentmarker}{}%
\end{pgfscope}%
\end{pgfscope}%
\begin{pgfscope}%
\pgfpathrectangle{\pgfqpoint{0.616050in}{2.359500in}}{\pgfqpoint{2.647350in}{1.384500in}}%
\pgfusepath{clip}%
\pgfsetrectcap%
\pgfsetroundjoin%
\pgfsetlinewidth{0.903375pt}%
\definecolor{currentstroke}{rgb}{0.784314,0.784314,0.784314}%
\pgfsetstrokecolor{currentstroke}%
\pgfsetdash{}{0pt}%
\pgfpathmoveto{\pgfqpoint{0.760451in}{3.355715in}}%
\pgfpathlineto{\pgfqpoint{1.145520in}{3.209167in}}%
\pgfpathlineto{\pgfqpoint{1.530589in}{3.212496in}}%
\pgfusepath{stroke}%
\end{pgfscope}%
\begin{pgfscope}%
\pgfpathrectangle{\pgfqpoint{0.616050in}{2.359500in}}{\pgfqpoint{2.647350in}{1.384500in}}%
\pgfusepath{clip}%
\pgfsetbuttcap%
\pgfsetmiterjoin%
\definecolor{currentfill}{rgb}{0.784314,0.784314,0.784314}%
\pgfsetfillcolor{currentfill}%
\pgfsetlinewidth{0.702625pt}%
\definecolor{currentstroke}{rgb}{0.784314,0.784314,0.784314}%
\pgfsetstrokecolor{currentstroke}%
\pgfsetdash{}{0pt}%
\pgfsys@defobject{currentmarker}{\pgfqpoint{-0.022222in}{-0.022222in}}{\pgfqpoint{0.022222in}{0.022222in}}{%
\pgfpathmoveto{\pgfqpoint{-0.022222in}{-0.022222in}}%
\pgfpathlineto{\pgfqpoint{0.022222in}{-0.022222in}}%
\pgfpathlineto{\pgfqpoint{0.022222in}{0.022222in}}%
\pgfpathlineto{\pgfqpoint{-0.022222in}{0.022222in}}%
\pgfpathlineto{\pgfqpoint{-0.022222in}{-0.022222in}}%
\pgfpathclose%
\pgfusepath{stroke,fill}%
}%
\begin{pgfscope}%
\pgfsys@transformshift{0.760451in}{3.355715in}%
\pgfsys@useobject{currentmarker}{}%
\end{pgfscope}%
\begin{pgfscope}%
\pgfsys@transformshift{1.145520in}{3.209167in}%
\pgfsys@useobject{currentmarker}{}%
\end{pgfscope}%
\begin{pgfscope}%
\pgfsys@transformshift{1.530589in}{3.212496in}%
\pgfsys@useobject{currentmarker}{}%
\end{pgfscope}%
\end{pgfscope}%
\begin{pgfscope}%
\pgfpathrectangle{\pgfqpoint{0.616050in}{2.359500in}}{\pgfqpoint{2.647350in}{1.384500in}}%
\pgfusepath{clip}%
\pgfsetrectcap%
\pgfsetroundjoin%
\pgfsetlinewidth{0.903375pt}%
\definecolor{currentstroke}{rgb}{0.478431,0.478431,0.478431}%
\pgfsetstrokecolor{currentstroke}%
\pgfsetdash{}{0pt}%
\pgfpathmoveto{\pgfqpoint{0.760451in}{3.429752in}}%
\pgfpathlineto{\pgfqpoint{1.530589in}{3.208474in}}%
\pgfusepath{stroke}%
\end{pgfscope}%
\begin{pgfscope}%
\pgfpathrectangle{\pgfqpoint{0.616050in}{2.359500in}}{\pgfqpoint{2.647350in}{1.384500in}}%
\pgfusepath{clip}%
\pgfsetbuttcap%
\pgfsetmiterjoin%
\definecolor{currentfill}{rgb}{0.478431,0.478431,0.478431}%
\pgfsetfillcolor{currentfill}%
\pgfsetlinewidth{0.702625pt}%
\definecolor{currentstroke}{rgb}{0.478431,0.478431,0.478431}%
\pgfsetstrokecolor{currentstroke}%
\pgfsetdash{}{0pt}%
\pgfsys@defobject{currentmarker}{\pgfqpoint{-0.022222in}{-0.022222in}}{\pgfqpoint{0.022222in}{0.022222in}}{%
\pgfpathmoveto{\pgfqpoint{-0.007407in}{-0.022222in}}%
\pgfpathlineto{\pgfqpoint{0.007407in}{-0.022222in}}%
\pgfpathlineto{\pgfqpoint{0.007407in}{-0.007407in}}%
\pgfpathlineto{\pgfqpoint{0.022222in}{-0.007407in}}%
\pgfpathlineto{\pgfqpoint{0.022222in}{0.007407in}}%
\pgfpathlineto{\pgfqpoint{0.007407in}{0.007407in}}%
\pgfpathlineto{\pgfqpoint{0.007407in}{0.022222in}}%
\pgfpathlineto{\pgfqpoint{-0.007407in}{0.022222in}}%
\pgfpathlineto{\pgfqpoint{-0.007407in}{0.007407in}}%
\pgfpathlineto{\pgfqpoint{-0.022222in}{0.007407in}}%
\pgfpathlineto{\pgfqpoint{-0.022222in}{-0.007407in}}%
\pgfpathlineto{\pgfqpoint{-0.007407in}{-0.007407in}}%
\pgfpathlineto{\pgfqpoint{-0.007407in}{-0.022222in}}%
\pgfpathclose%
\pgfusepath{stroke,fill}%
}%
\begin{pgfscope}%
\pgfsys@transformshift{0.760451in}{3.429752in}%
\pgfsys@useobject{currentmarker}{}%
\end{pgfscope}%
\begin{pgfscope}%
\pgfsys@transformshift{1.530589in}{3.208474in}%
\pgfsys@useobject{currentmarker}{}%
\end{pgfscope}%
\end{pgfscope}%
\begin{pgfscope}%
\pgfsetrectcap%
\pgfsetmiterjoin%
\pgfsetlinewidth{0.803000pt}%
\definecolor{currentstroke}{rgb}{0.101961,0.101961,0.101961}%
\pgfsetstrokecolor{currentstroke}%
\pgfsetdash{}{0pt}%
\pgfpathmoveto{\pgfqpoint{0.616050in}{2.359500in}}%
\pgfpathlineto{\pgfqpoint{0.616050in}{3.744000in}}%
\pgfusepath{stroke}%
\end{pgfscope}%
\begin{pgfscope}%
\pgfsetrectcap%
\pgfsetmiterjoin%
\pgfsetlinewidth{0.803000pt}%
\definecolor{currentstroke}{rgb}{0.101961,0.101961,0.101961}%
\pgfsetstrokecolor{currentstroke}%
\pgfsetdash{}{0pt}%
\pgfpathmoveto{\pgfqpoint{0.616050in}{2.359500in}}%
\pgfpathlineto{\pgfqpoint{3.263400in}{2.359500in}}%
\pgfusepath{stroke}%
\end{pgfscope}%
\begin{pgfscope}%
\pgfpathrectangle{\pgfqpoint{0.616050in}{2.359500in}}{\pgfqpoint{2.647350in}{1.384500in}}%
\pgfusepath{clip}%
\pgfsetbuttcap%
\pgfsetroundjoin%
\pgfsetlinewidth{0.903375pt}%
\definecolor{currentstroke}{rgb}{0.156863,0.356863,0.568627}%
\pgfsetstrokecolor{currentstroke}%
\pgfsetdash{{3.330000pt}{1.440000pt}}{0.000000pt}%
\pgfpathmoveto{\pgfqpoint{1.915658in}{3.192136in}}%
\pgfpathlineto{\pgfqpoint{2.300727in}{3.219726in}}%
\pgfpathlineto{\pgfqpoint{3.070865in}{2.975272in}}%
\pgfusepath{stroke}%
\end{pgfscope}%
\begin{pgfscope}%
\pgfpathrectangle{\pgfqpoint{0.616050in}{2.359500in}}{\pgfqpoint{2.647350in}{1.384500in}}%
\pgfusepath{clip}%
\pgfsetbuttcap%
\pgfsetroundjoin%
\definecolor{currentfill}{rgb}{1.000000,1.000000,1.000000}%
\pgfsetfillcolor{currentfill}%
\pgfsetlinewidth{0.702625pt}%
\definecolor{currentstroke}{rgb}{0.156863,0.356863,0.568627}%
\pgfsetstrokecolor{currentstroke}%
\pgfsetdash{}{0pt}%
\pgfsys@defobject{currentmarker}{\pgfqpoint{-0.022222in}{-0.022222in}}{\pgfqpoint{0.022222in}{0.022222in}}{%
\pgfpathmoveto{\pgfqpoint{0.000000in}{-0.022222in}}%
\pgfpathcurveto{\pgfqpoint{0.005893in}{-0.022222in}}{\pgfqpoint{0.011546in}{-0.019881in}}{\pgfqpoint{0.015713in}{-0.015713in}}%
\pgfpathcurveto{\pgfqpoint{0.019881in}{-0.011546in}}{\pgfqpoint{0.022222in}{-0.005893in}}{\pgfqpoint{0.022222in}{0.000000in}}%
\pgfpathcurveto{\pgfqpoint{0.022222in}{0.005893in}}{\pgfqpoint{0.019881in}{0.011546in}}{\pgfqpoint{0.015713in}{0.015713in}}%
\pgfpathcurveto{\pgfqpoint{0.011546in}{0.019881in}}{\pgfqpoint{0.005893in}{0.022222in}}{\pgfqpoint{0.000000in}{0.022222in}}%
\pgfpathcurveto{\pgfqpoint{-0.005893in}{0.022222in}}{\pgfqpoint{-0.011546in}{0.019881in}}{\pgfqpoint{-0.015713in}{0.015713in}}%
\pgfpathcurveto{\pgfqpoint{-0.019881in}{0.011546in}}{\pgfqpoint{-0.022222in}{0.005893in}}{\pgfqpoint{-0.022222in}{0.000000in}}%
\pgfpathcurveto{\pgfqpoint{-0.022222in}{-0.005893in}}{\pgfqpoint{-0.019881in}{-0.011546in}}{\pgfqpoint{-0.015713in}{-0.015713in}}%
\pgfpathcurveto{\pgfqpoint{-0.011546in}{-0.019881in}}{\pgfqpoint{-0.005893in}{-0.022222in}}{\pgfqpoint{0.000000in}{-0.022222in}}%
\pgfpathlineto{\pgfqpoint{0.000000in}{-0.022222in}}%
\pgfpathclose%
\pgfusepath{stroke,fill}%
}%
\begin{pgfscope}%
\pgfsys@transformshift{1.915658in}{3.192136in}%
\pgfsys@useobject{currentmarker}{}%
\end{pgfscope}%
\begin{pgfscope}%
\pgfsys@transformshift{2.300727in}{3.219726in}%
\pgfsys@useobject{currentmarker}{}%
\end{pgfscope}%
\begin{pgfscope}%
\pgfsys@transformshift{3.070865in}{2.975272in}%
\pgfsys@useobject{currentmarker}{}%
\end{pgfscope}%
\end{pgfscope}%
\begin{pgfscope}%
\pgfpathrectangle{\pgfqpoint{0.616050in}{2.359500in}}{\pgfqpoint{2.647350in}{1.384500in}}%
\pgfusepath{clip}%
\pgfsetrectcap%
\pgfsetroundjoin%
\pgfsetlinewidth{0.903375pt}%
\definecolor{currentstroke}{rgb}{0.156863,0.356863,0.568627}%
\pgfsetstrokecolor{currentstroke}%
\pgfsetdash{}{0pt}%
\pgfpathmoveto{\pgfqpoint{1.530589in}{3.436178in}}%
\pgfpathlineto{\pgfqpoint{1.915658in}{3.308043in}}%
\pgfpathlineto{\pgfqpoint{2.300727in}{3.390806in}}%
\pgfusepath{stroke}%
\end{pgfscope}%
\begin{pgfscope}%
\pgfpathrectangle{\pgfqpoint{0.616050in}{2.359500in}}{\pgfqpoint{2.647350in}{1.384500in}}%
\pgfusepath{clip}%
\pgfsetbuttcap%
\pgfsetroundjoin%
\definecolor{currentfill}{rgb}{0.156863,0.356863,0.568627}%
\pgfsetfillcolor{currentfill}%
\pgfsetlinewidth{0.702625pt}%
\definecolor{currentstroke}{rgb}{0.156863,0.356863,0.568627}%
\pgfsetstrokecolor{currentstroke}%
\pgfsetdash{}{0pt}%
\pgfsys@defobject{currentmarker}{\pgfqpoint{-0.022222in}{-0.022222in}}{\pgfqpoint{0.022222in}{0.022222in}}{%
\pgfpathmoveto{\pgfqpoint{0.000000in}{-0.022222in}}%
\pgfpathcurveto{\pgfqpoint{0.005893in}{-0.022222in}}{\pgfqpoint{0.011546in}{-0.019881in}}{\pgfqpoint{0.015713in}{-0.015713in}}%
\pgfpathcurveto{\pgfqpoint{0.019881in}{-0.011546in}}{\pgfqpoint{0.022222in}{-0.005893in}}{\pgfqpoint{0.022222in}{0.000000in}}%
\pgfpathcurveto{\pgfqpoint{0.022222in}{0.005893in}}{\pgfqpoint{0.019881in}{0.011546in}}{\pgfqpoint{0.015713in}{0.015713in}}%
\pgfpathcurveto{\pgfqpoint{0.011546in}{0.019881in}}{\pgfqpoint{0.005893in}{0.022222in}}{\pgfqpoint{0.000000in}{0.022222in}}%
\pgfpathcurveto{\pgfqpoint{-0.005893in}{0.022222in}}{\pgfqpoint{-0.011546in}{0.019881in}}{\pgfqpoint{-0.015713in}{0.015713in}}%
\pgfpathcurveto{\pgfqpoint{-0.019881in}{0.011546in}}{\pgfqpoint{-0.022222in}{0.005893in}}{\pgfqpoint{-0.022222in}{0.000000in}}%
\pgfpathcurveto{\pgfqpoint{-0.022222in}{-0.005893in}}{\pgfqpoint{-0.019881in}{-0.011546in}}{\pgfqpoint{-0.015713in}{-0.015713in}}%
\pgfpathcurveto{\pgfqpoint{-0.011546in}{-0.019881in}}{\pgfqpoint{-0.005893in}{-0.022222in}}{\pgfqpoint{0.000000in}{-0.022222in}}%
\pgfpathlineto{\pgfqpoint{0.000000in}{-0.022222in}}%
\pgfpathclose%
\pgfusepath{stroke,fill}%
}%
\begin{pgfscope}%
\pgfsys@transformshift{1.530589in}{3.436178in}%
\pgfsys@useobject{currentmarker}{}%
\end{pgfscope}%
\begin{pgfscope}%
\pgfsys@transformshift{1.915658in}{3.308043in}%
\pgfsys@useobject{currentmarker}{}%
\end{pgfscope}%
\begin{pgfscope}%
\pgfsys@transformshift{2.300727in}{3.390806in}%
\pgfsys@useobject{currentmarker}{}%
\end{pgfscope}%
\end{pgfscope}%
\begin{pgfscope}%
\pgfpathrectangle{\pgfqpoint{0.616050in}{2.359500in}}{\pgfqpoint{2.647350in}{1.384500in}}%
\pgfusepath{clip}%
\pgfsetbuttcap%
\pgfsetroundjoin%
\pgfsetlinewidth{0.903375pt}%
\definecolor{currentstroke}{rgb}{0.156863,0.356863,0.568627}%
\pgfsetstrokecolor{currentstroke}%
\pgfsetdash{{3.330000pt}{1.440000pt}}{0.000000pt}%
\pgfpathmoveto{\pgfqpoint{1.915658in}{3.347045in}}%
\pgfusepath{stroke}%
\end{pgfscope}%
\begin{pgfscope}%
\pgfpathrectangle{\pgfqpoint{0.616050in}{2.359500in}}{\pgfqpoint{2.647350in}{1.384500in}}%
\pgfusepath{clip}%
\pgfsetbuttcap%
\pgfsetmiterjoin%
\definecolor{currentfill}{rgb}{1.000000,1.000000,1.000000}%
\pgfsetfillcolor{currentfill}%
\pgfsetlinewidth{0.702625pt}%
\definecolor{currentstroke}{rgb}{0.156863,0.356863,0.568627}%
\pgfsetstrokecolor{currentstroke}%
\pgfsetdash{}{0pt}%
\pgfsys@defobject{currentmarker}{\pgfqpoint{-0.031427in}{-0.031427in}}{\pgfqpoint{0.031427in}{0.031427in}}{%
\pgfpathmoveto{\pgfqpoint{-0.000000in}{-0.031427in}}%
\pgfpathlineto{\pgfqpoint{0.031427in}{0.000000in}}%
\pgfpathlineto{\pgfqpoint{0.000000in}{0.031427in}}%
\pgfpathlineto{\pgfqpoint{-0.031427in}{0.000000in}}%
\pgfpathlineto{\pgfqpoint{-0.000000in}{-0.031427in}}%
\pgfpathclose%
\pgfusepath{stroke,fill}%
}%
\begin{pgfscope}%
\pgfsys@transformshift{1.915658in}{3.347045in}%
\pgfsys@useobject{currentmarker}{}%
\end{pgfscope}%
\end{pgfscope}%
\begin{pgfscope}%
\pgfpathrectangle{\pgfqpoint{0.616050in}{2.359500in}}{\pgfqpoint{2.647350in}{1.384500in}}%
\pgfusepath{clip}%
\pgfsetrectcap%
\pgfsetroundjoin%
\pgfsetlinewidth{0.903375pt}%
\definecolor{currentstroke}{rgb}{0.156863,0.356863,0.568627}%
\pgfsetstrokecolor{currentstroke}%
\pgfsetdash{}{0pt}%
\pgfpathmoveto{\pgfqpoint{1.530589in}{3.515953in}}%
\pgfusepath{stroke}%
\end{pgfscope}%
\begin{pgfscope}%
\pgfpathrectangle{\pgfqpoint{0.616050in}{2.359500in}}{\pgfqpoint{2.647350in}{1.384500in}}%
\pgfusepath{clip}%
\pgfsetbuttcap%
\pgfsetmiterjoin%
\definecolor{currentfill}{rgb}{0.156863,0.356863,0.568627}%
\pgfsetfillcolor{currentfill}%
\pgfsetlinewidth{0.702625pt}%
\definecolor{currentstroke}{rgb}{0.156863,0.356863,0.568627}%
\pgfsetstrokecolor{currentstroke}%
\pgfsetdash{}{0pt}%
\pgfsys@defobject{currentmarker}{\pgfqpoint{-0.031427in}{-0.031427in}}{\pgfqpoint{0.031427in}{0.031427in}}{%
\pgfpathmoveto{\pgfqpoint{-0.000000in}{-0.031427in}}%
\pgfpathlineto{\pgfqpoint{0.031427in}{0.000000in}}%
\pgfpathlineto{\pgfqpoint{0.000000in}{0.031427in}}%
\pgfpathlineto{\pgfqpoint{-0.031427in}{0.000000in}}%
\pgfpathlineto{\pgfqpoint{-0.000000in}{-0.031427in}}%
\pgfpathclose%
\pgfusepath{stroke,fill}%
}%
\begin{pgfscope}%
\pgfsys@transformshift{1.530589in}{3.515953in}%
\pgfsys@useobject{currentmarker}{}%
\end{pgfscope}%
\end{pgfscope}%
\begin{pgfscope}%
\definecolor{textcolor}{rgb}{0.549020,0.549020,0.549020}%
\pgfsetstrokecolor{textcolor}%
\pgfsetfillcolor{textcolor}%
\pgftext[x=3.236926in,y=2.414880in,right,bottom]{\color{textcolor}\sffamily\fontsize{7.000000}{8.400000}\selectfont capacity line: not measured (fleet sweep pending)}%
\end{pgfscope}%
\begin{pgfscope}%
\pgfsetbuttcap%
\pgfsetmiterjoin%
\definecolor{currentfill}{rgb}{1.000000,1.000000,1.000000}%
\pgfsetfillcolor{currentfill}%
\pgfsetlinewidth{0.000000pt}%
\definecolor{currentstroke}{rgb}{0.000000,0.000000,0.000000}%
\pgfsetstrokecolor{currentstroke}%
\pgfsetstrokeopacity{0.000000}%
\pgfsetdash{}{0pt}%
\pgfpathmoveto{\pgfqpoint{0.616050in}{0.916500in}}%
\pgfpathlineto{\pgfqpoint{3.263400in}{0.916500in}}%
\pgfpathlineto{\pgfqpoint{3.263400in}{2.223000in}}%
\pgfpathlineto{\pgfqpoint{0.616050in}{2.223000in}}%
\pgfpathlineto{\pgfqpoint{0.616050in}{0.916500in}}%
\pgfpathclose%
\pgfusepath{fill}%
\end{pgfscope}%
\begin{pgfscope}%
\pgfpathrectangle{\pgfqpoint{0.616050in}{0.916500in}}{\pgfqpoint{2.647350in}{1.306500in}}%
\pgfusepath{clip}%
\pgfsetrectcap%
\pgfsetroundjoin%
\pgfsetlinewidth{0.401500pt}%
\definecolor{currentstroke}{rgb}{0.901961,0.901961,0.901961}%
\pgfsetstrokecolor{currentstroke}%
\pgfsetdash{}{0pt}%
\pgfpathmoveto{\pgfqpoint{0.760451in}{0.916500in}}%
\pgfpathlineto{\pgfqpoint{0.760451in}{2.223000in}}%
\pgfusepath{stroke}%
\end{pgfscope}%
\begin{pgfscope}%
\definecolor{textcolor}{rgb}{0.101961,0.101961,0.101961}%
\pgfsetstrokecolor{textcolor}%
\pgfsetfillcolor{textcolor}%
\pgftext[x=0.760451in,y=0.867889in,,top]{\color{textcolor}\sffamily\fontsize{7.000000}{8.400000}\selectfont 8}%
\end{pgfscope}%
\begin{pgfscope}%
\pgfpathrectangle{\pgfqpoint{0.616050in}{0.916500in}}{\pgfqpoint{2.647350in}{1.306500in}}%
\pgfusepath{clip}%
\pgfsetrectcap%
\pgfsetroundjoin%
\pgfsetlinewidth{0.401500pt}%
\definecolor{currentstroke}{rgb}{0.901961,0.901961,0.901961}%
\pgfsetstrokecolor{currentstroke}%
\pgfsetdash{}{0pt}%
\pgfpathmoveto{\pgfqpoint{1.145520in}{0.916500in}}%
\pgfpathlineto{\pgfqpoint{1.145520in}{2.223000in}}%
\pgfusepath{stroke}%
\end{pgfscope}%
\begin{pgfscope}%
\definecolor{textcolor}{rgb}{0.101961,0.101961,0.101961}%
\pgfsetstrokecolor{textcolor}%
\pgfsetfillcolor{textcolor}%
\pgftext[x=1.145520in,y=0.867889in,,top]{\color{textcolor}\sffamily\fontsize{7.000000}{8.400000}\selectfont 12}%
\end{pgfscope}%
\begin{pgfscope}%
\pgfpathrectangle{\pgfqpoint{0.616050in}{0.916500in}}{\pgfqpoint{2.647350in}{1.306500in}}%
\pgfusepath{clip}%
\pgfsetrectcap%
\pgfsetroundjoin%
\pgfsetlinewidth{0.401500pt}%
\definecolor{currentstroke}{rgb}{0.901961,0.901961,0.901961}%
\pgfsetstrokecolor{currentstroke}%
\pgfsetdash{}{0pt}%
\pgfpathmoveto{\pgfqpoint{1.530589in}{0.916500in}}%
\pgfpathlineto{\pgfqpoint{1.530589in}{2.223000in}}%
\pgfusepath{stroke}%
\end{pgfscope}%
\begin{pgfscope}%
\definecolor{textcolor}{rgb}{0.101961,0.101961,0.101961}%
\pgfsetstrokecolor{textcolor}%
\pgfsetfillcolor{textcolor}%
\pgftext[x=1.530589in,y=0.867889in,,top]{\color{textcolor}\sffamily\fontsize{7.000000}{8.400000}\selectfont 16}%
\end{pgfscope}%
\begin{pgfscope}%
\pgfpathrectangle{\pgfqpoint{0.616050in}{0.916500in}}{\pgfqpoint{2.647350in}{1.306500in}}%
\pgfusepath{clip}%
\pgfsetrectcap%
\pgfsetroundjoin%
\pgfsetlinewidth{0.401500pt}%
\definecolor{currentstroke}{rgb}{0.901961,0.901961,0.901961}%
\pgfsetstrokecolor{currentstroke}%
\pgfsetdash{}{0pt}%
\pgfpathmoveto{\pgfqpoint{1.915658in}{0.916500in}}%
\pgfpathlineto{\pgfqpoint{1.915658in}{2.223000in}}%
\pgfusepath{stroke}%
\end{pgfscope}%
\begin{pgfscope}%
\definecolor{textcolor}{rgb}{0.101961,0.101961,0.101961}%
\pgfsetstrokecolor{textcolor}%
\pgfsetfillcolor{textcolor}%
\pgftext[x=1.915658in,y=0.867889in,,top]{\color{textcolor}\sffamily\fontsize{7.000000}{8.400000}\selectfont 20}%
\end{pgfscope}%
\begin{pgfscope}%
\pgfpathrectangle{\pgfqpoint{0.616050in}{0.916500in}}{\pgfqpoint{2.647350in}{1.306500in}}%
\pgfusepath{clip}%
\pgfsetrectcap%
\pgfsetroundjoin%
\pgfsetlinewidth{0.401500pt}%
\definecolor{currentstroke}{rgb}{0.901961,0.901961,0.901961}%
\pgfsetstrokecolor{currentstroke}%
\pgfsetdash{}{0pt}%
\pgfpathmoveto{\pgfqpoint{2.300727in}{0.916500in}}%
\pgfpathlineto{\pgfqpoint{2.300727in}{2.223000in}}%
\pgfusepath{stroke}%
\end{pgfscope}%
\begin{pgfscope}%
\definecolor{textcolor}{rgb}{0.101961,0.101961,0.101961}%
\pgfsetstrokecolor{textcolor}%
\pgfsetfillcolor{textcolor}%
\pgftext[x=2.300727in,y=0.867889in,,top]{\color{textcolor}\sffamily\fontsize{7.000000}{8.400000}\selectfont 24}%
\end{pgfscope}%
\begin{pgfscope}%
\pgfpathrectangle{\pgfqpoint{0.616050in}{0.916500in}}{\pgfqpoint{2.647350in}{1.306500in}}%
\pgfusepath{clip}%
\pgfsetrectcap%
\pgfsetroundjoin%
\pgfsetlinewidth{0.401500pt}%
\definecolor{currentstroke}{rgb}{0.901961,0.901961,0.901961}%
\pgfsetstrokecolor{currentstroke}%
\pgfsetdash{}{0pt}%
\pgfpathmoveto{\pgfqpoint{3.070865in}{0.916500in}}%
\pgfpathlineto{\pgfqpoint{3.070865in}{2.223000in}}%
\pgfusepath{stroke}%
\end{pgfscope}%
\begin{pgfscope}%
\definecolor{textcolor}{rgb}{0.101961,0.101961,0.101961}%
\pgfsetstrokecolor{textcolor}%
\pgfsetfillcolor{textcolor}%
\pgftext[x=3.070865in,y=0.867889in,,top]{\color{textcolor}\sffamily\fontsize{7.000000}{8.400000}\selectfont 32}%
\end{pgfscope}%
\begin{pgfscope}%
\definecolor{textcolor}{rgb}{0.000000,0.000000,0.000000}%
\pgfsetstrokecolor{textcolor}%
\pgfsetfillcolor{textcolor}%
\pgftext[x=1.939725in,y=0.753692in,,top]{\color{textcolor}\sffamily\fontsize{7.500000}{9.000000}\selectfont Fleet size (devices)}%
\end{pgfscope}%
\begin{pgfscope}%
\pgfpathrectangle{\pgfqpoint{0.616050in}{0.916500in}}{\pgfqpoint{2.647350in}{1.306500in}}%
\pgfusepath{clip}%
\pgfsetrectcap%
\pgfsetroundjoin%
\pgfsetlinewidth{0.401500pt}%
\definecolor{currentstroke}{rgb}{0.901961,0.901961,0.901961}%
\pgfsetstrokecolor{currentstroke}%
\pgfsetdash{}{0pt}%
\pgfpathmoveto{\pgfqpoint{0.616050in}{0.954553in}}%
\pgfpathlineto{\pgfqpoint{3.263400in}{0.954553in}}%
\pgfusepath{stroke}%
\end{pgfscope}%
\begin{pgfscope}%
\definecolor{textcolor}{rgb}{0.101961,0.101961,0.101961}%
\pgfsetstrokecolor{textcolor}%
\pgfsetfillcolor{textcolor}%
\pgftext[x=0.435445in, y=0.920796in, left, base]{\color{textcolor}\sffamily\fontsize{7.000000}{8.400000}\selectfont 0.0}%
\end{pgfscope}%
\begin{pgfscope}%
\pgfpathrectangle{\pgfqpoint{0.616050in}{0.916500in}}{\pgfqpoint{2.647350in}{1.306500in}}%
\pgfusepath{clip}%
\pgfsetrectcap%
\pgfsetroundjoin%
\pgfsetlinewidth{0.401500pt}%
\definecolor{currentstroke}{rgb}{0.901961,0.901961,0.901961}%
\pgfsetstrokecolor{currentstroke}%
\pgfsetdash{}{0pt}%
\pgfpathmoveto{\pgfqpoint{0.616050in}{1.208243in}}%
\pgfpathlineto{\pgfqpoint{3.263400in}{1.208243in}}%
\pgfusepath{stroke}%
\end{pgfscope}%
\begin{pgfscope}%
\definecolor{textcolor}{rgb}{0.101961,0.101961,0.101961}%
\pgfsetstrokecolor{textcolor}%
\pgfsetfillcolor{textcolor}%
\pgftext[x=0.435445in, y=1.174485in, left, base]{\color{textcolor}\sffamily\fontsize{7.000000}{8.400000}\selectfont 0.2}%
\end{pgfscope}%
\begin{pgfscope}%
\pgfpathrectangle{\pgfqpoint{0.616050in}{0.916500in}}{\pgfqpoint{2.647350in}{1.306500in}}%
\pgfusepath{clip}%
\pgfsetrectcap%
\pgfsetroundjoin%
\pgfsetlinewidth{0.401500pt}%
\definecolor{currentstroke}{rgb}{0.901961,0.901961,0.901961}%
\pgfsetstrokecolor{currentstroke}%
\pgfsetdash{}{0pt}%
\pgfpathmoveto{\pgfqpoint{0.616050in}{1.461932in}}%
\pgfpathlineto{\pgfqpoint{3.263400in}{1.461932in}}%
\pgfusepath{stroke}%
\end{pgfscope}%
\begin{pgfscope}%
\definecolor{textcolor}{rgb}{0.101961,0.101961,0.101961}%
\pgfsetstrokecolor{textcolor}%
\pgfsetfillcolor{textcolor}%
\pgftext[x=0.435445in, y=1.428174in, left, base]{\color{textcolor}\sffamily\fontsize{7.000000}{8.400000}\selectfont 0.4}%
\end{pgfscope}%
\begin{pgfscope}%
\pgfpathrectangle{\pgfqpoint{0.616050in}{0.916500in}}{\pgfqpoint{2.647350in}{1.306500in}}%
\pgfusepath{clip}%
\pgfsetrectcap%
\pgfsetroundjoin%
\pgfsetlinewidth{0.401500pt}%
\definecolor{currentstroke}{rgb}{0.901961,0.901961,0.901961}%
\pgfsetstrokecolor{currentstroke}%
\pgfsetdash{}{0pt}%
\pgfpathmoveto{\pgfqpoint{0.616050in}{1.715621in}}%
\pgfpathlineto{\pgfqpoint{3.263400in}{1.715621in}}%
\pgfusepath{stroke}%
\end{pgfscope}%
\begin{pgfscope}%
\definecolor{textcolor}{rgb}{0.101961,0.101961,0.101961}%
\pgfsetstrokecolor{textcolor}%
\pgfsetfillcolor{textcolor}%
\pgftext[x=0.435445in, y=1.681864in, left, base]{\color{textcolor}\sffamily\fontsize{7.000000}{8.400000}\selectfont 0.6}%
\end{pgfscope}%
\begin{pgfscope}%
\pgfpathrectangle{\pgfqpoint{0.616050in}{0.916500in}}{\pgfqpoint{2.647350in}{1.306500in}}%
\pgfusepath{clip}%
\pgfsetrectcap%
\pgfsetroundjoin%
\pgfsetlinewidth{0.401500pt}%
\definecolor{currentstroke}{rgb}{0.901961,0.901961,0.901961}%
\pgfsetstrokecolor{currentstroke}%
\pgfsetdash{}{0pt}%
\pgfpathmoveto{\pgfqpoint{0.616050in}{1.969311in}}%
\pgfpathlineto{\pgfqpoint{3.263400in}{1.969311in}}%
\pgfusepath{stroke}%
\end{pgfscope}%
\begin{pgfscope}%
\definecolor{textcolor}{rgb}{0.101961,0.101961,0.101961}%
\pgfsetstrokecolor{textcolor}%
\pgfsetfillcolor{textcolor}%
\pgftext[x=0.435445in, y=1.935553in, left, base]{\color{textcolor}\sffamily\fontsize{7.000000}{8.400000}\selectfont 0.8}%
\end{pgfscope}%
\begin{pgfscope}%
\pgfpathrectangle{\pgfqpoint{0.616050in}{0.916500in}}{\pgfqpoint{2.647350in}{1.306500in}}%
\pgfusepath{clip}%
\pgfsetrectcap%
\pgfsetroundjoin%
\pgfsetlinewidth{0.401500pt}%
\definecolor{currentstroke}{rgb}{0.901961,0.901961,0.901961}%
\pgfsetstrokecolor{currentstroke}%
\pgfsetdash{}{0pt}%
\pgfpathmoveto{\pgfqpoint{0.616050in}{2.223000in}}%
\pgfpathlineto{\pgfqpoint{3.263400in}{2.223000in}}%
\pgfusepath{stroke}%
\end{pgfscope}%
\begin{pgfscope}%
\definecolor{textcolor}{rgb}{0.101961,0.101961,0.101961}%
\pgfsetstrokecolor{textcolor}%
\pgfsetfillcolor{textcolor}%
\pgftext[x=0.435445in, y=2.189242in, left, base]{\color{textcolor}\sffamily\fontsize{7.000000}{8.400000}\selectfont 1.0}%
\end{pgfscope}%
\begin{pgfscope}%
\definecolor{textcolor}{rgb}{0.000000,0.000000,0.000000}%
\pgfsetstrokecolor{textcolor}%
\pgfsetfillcolor{textcolor}%
\pgftext[x=0.261062in, y=1.372383in, left, base,rotate=90.000000]{\color{textcolor}\sffamily\fontsize{7.500000}{9.000000}\selectfont Dashcam}%
\end{pgfscope}%
\begin{pgfscope}%
\definecolor{textcolor}{rgb}{0.000000,0.000000,0.000000}%
\pgfsetstrokecolor{textcolor}%
\pgfsetfillcolor{textcolor}%
\pgftext[x=0.360985in, y=1.334541in, left, base,rotate=90.000000]{\color{textcolor}\sffamily\fontsize{7.500000}{9.000000}\selectfont attainment}%
\end{pgfscope}%
\begin{pgfscope}%
\pgfpathrectangle{\pgfqpoint{0.616050in}{0.916500in}}{\pgfqpoint{2.647350in}{1.306500in}}%
\pgfusepath{clip}%
\pgfsetrectcap%
\pgfsetroundjoin%
\pgfsetlinewidth{0.903375pt}%
\definecolor{currentstroke}{rgb}{0.478431,0.305882,0.611765}%
\pgfsetstrokecolor{currentstroke}%
\pgfsetdash{}{0pt}%
\pgfpathmoveto{\pgfqpoint{0.760451in}{0.954553in}}%
\pgfpathlineto{\pgfqpoint{1.145520in}{0.954553in}}%
\pgfpathlineto{\pgfqpoint{1.530589in}{0.954553in}}%
\pgfpathlineto{\pgfqpoint{1.915658in}{0.954553in}}%
\pgfusepath{stroke}%
\end{pgfscope}%
\begin{pgfscope}%
\pgfpathrectangle{\pgfqpoint{0.616050in}{0.916500in}}{\pgfqpoint{2.647350in}{1.306500in}}%
\pgfusepath{clip}%
\pgfsetbuttcap%
\pgfsetmiterjoin%
\definecolor{currentfill}{rgb}{0.478431,0.305882,0.611765}%
\pgfsetfillcolor{currentfill}%
\pgfsetlinewidth{0.702625pt}%
\definecolor{currentstroke}{rgb}{0.478431,0.305882,0.611765}%
\pgfsetstrokecolor{currentstroke}%
\pgfsetdash{}{0pt}%
\pgfsys@defobject{currentmarker}{\pgfqpoint{-0.022222in}{-0.022222in}}{\pgfqpoint{0.022222in}{0.022222in}}{%
\pgfpathmoveto{\pgfqpoint{-0.011111in}{-0.022222in}}%
\pgfpathlineto{\pgfqpoint{0.000000in}{-0.011111in}}%
\pgfpathlineto{\pgfqpoint{0.011111in}{-0.022222in}}%
\pgfpathlineto{\pgfqpoint{0.022222in}{-0.011111in}}%
\pgfpathlineto{\pgfqpoint{0.011111in}{0.000000in}}%
\pgfpathlineto{\pgfqpoint{0.022222in}{0.011111in}}%
\pgfpathlineto{\pgfqpoint{0.011111in}{0.022222in}}%
\pgfpathlineto{\pgfqpoint{0.000000in}{0.011111in}}%
\pgfpathlineto{\pgfqpoint{-0.011111in}{0.022222in}}%
\pgfpathlineto{\pgfqpoint{-0.022222in}{0.011111in}}%
\pgfpathlineto{\pgfqpoint{-0.011111in}{0.000000in}}%
\pgfpathlineto{\pgfqpoint{-0.022222in}{-0.011111in}}%
\pgfpathlineto{\pgfqpoint{-0.011111in}{-0.022222in}}%
\pgfpathclose%
\pgfusepath{stroke,fill}%
}%
\begin{pgfscope}%
\pgfsys@transformshift{0.760451in}{0.954553in}%
\pgfsys@useobject{currentmarker}{}%
\end{pgfscope}%
\begin{pgfscope}%
\pgfsys@transformshift{1.145520in}{0.954553in}%
\pgfsys@useobject{currentmarker}{}%
\end{pgfscope}%
\begin{pgfscope}%
\pgfsys@transformshift{1.530589in}{0.954553in}%
\pgfsys@useobject{currentmarker}{}%
\end{pgfscope}%
\begin{pgfscope}%
\pgfsys@transformshift{1.915658in}{0.954553in}%
\pgfsys@useobject{currentmarker}{}%
\end{pgfscope}%
\end{pgfscope}%
\begin{pgfscope}%
\pgfpathrectangle{\pgfqpoint{0.616050in}{0.916500in}}{\pgfqpoint{2.647350in}{1.306500in}}%
\pgfusepath{clip}%
\pgfsetrectcap%
\pgfsetroundjoin%
\pgfsetlinewidth{0.903375pt}%
\definecolor{currentstroke}{rgb}{0.784314,0.784314,0.784314}%
\pgfsetstrokecolor{currentstroke}%
\pgfsetdash{}{0pt}%
\pgfpathmoveto{\pgfqpoint{0.760451in}{2.223000in}}%
\pgfpathlineto{\pgfqpoint{1.145520in}{2.218449in}}%
\pgfpathlineto{\pgfqpoint{1.530589in}{2.158467in}}%
\pgfusepath{stroke}%
\end{pgfscope}%
\begin{pgfscope}%
\pgfpathrectangle{\pgfqpoint{0.616050in}{0.916500in}}{\pgfqpoint{2.647350in}{1.306500in}}%
\pgfusepath{clip}%
\pgfsetbuttcap%
\pgfsetmiterjoin%
\definecolor{currentfill}{rgb}{0.784314,0.784314,0.784314}%
\pgfsetfillcolor{currentfill}%
\pgfsetlinewidth{0.702625pt}%
\definecolor{currentstroke}{rgb}{0.784314,0.784314,0.784314}%
\pgfsetstrokecolor{currentstroke}%
\pgfsetdash{}{0pt}%
\pgfsys@defobject{currentmarker}{\pgfqpoint{-0.022222in}{-0.022222in}}{\pgfqpoint{0.022222in}{0.022222in}}{%
\pgfpathmoveto{\pgfqpoint{0.000000in}{0.022222in}}%
\pgfpathlineto{\pgfqpoint{-0.022222in}{-0.022222in}}%
\pgfpathlineto{\pgfqpoint{0.022222in}{-0.022222in}}%
\pgfpathlineto{\pgfqpoint{0.000000in}{0.022222in}}%
\pgfpathclose%
\pgfusepath{stroke,fill}%
}%
\begin{pgfscope}%
\pgfsys@transformshift{0.760451in}{2.223000in}%
\pgfsys@useobject{currentmarker}{}%
\end{pgfscope}%
\begin{pgfscope}%
\pgfsys@transformshift{1.145520in}{2.218449in}%
\pgfsys@useobject{currentmarker}{}%
\end{pgfscope}%
\begin{pgfscope}%
\pgfsys@transformshift{1.530589in}{2.158467in}%
\pgfsys@useobject{currentmarker}{}%
\end{pgfscope}%
\end{pgfscope}%
\begin{pgfscope}%
\pgfpathrectangle{\pgfqpoint{0.616050in}{0.916500in}}{\pgfqpoint{2.647350in}{1.306500in}}%
\pgfusepath{clip}%
\pgfsetrectcap%
\pgfsetroundjoin%
\pgfsetlinewidth{0.903375pt}%
\definecolor{currentstroke}{rgb}{0.784314,0.784314,0.784314}%
\pgfsetstrokecolor{currentstroke}%
\pgfsetdash{}{0pt}%
\pgfpathmoveto{\pgfqpoint{0.760451in}{2.195935in}}%
\pgfpathlineto{\pgfqpoint{1.530589in}{1.345376in}}%
\pgfusepath{stroke}%
\end{pgfscope}%
\begin{pgfscope}%
\pgfpathrectangle{\pgfqpoint{0.616050in}{0.916500in}}{\pgfqpoint{2.647350in}{1.306500in}}%
\pgfusepath{clip}%
\pgfsetbuttcap%
\pgfsetmiterjoin%
\definecolor{currentfill}{rgb}{0.784314,0.784314,0.784314}%
\pgfsetfillcolor{currentfill}%
\pgfsetlinewidth{0.702625pt}%
\definecolor{currentstroke}{rgb}{0.784314,0.784314,0.784314}%
\pgfsetstrokecolor{currentstroke}%
\pgfsetdash{}{0pt}%
\pgfsys@defobject{currentmarker}{\pgfqpoint{-0.022222in}{-0.022222in}}{\pgfqpoint{0.022222in}{0.022222in}}{%
\pgfpathmoveto{\pgfqpoint{-0.000000in}{-0.022222in}}%
\pgfpathlineto{\pgfqpoint{0.022222in}{0.022222in}}%
\pgfpathlineto{\pgfqpoint{-0.022222in}{0.022222in}}%
\pgfpathlineto{\pgfqpoint{-0.000000in}{-0.022222in}}%
\pgfpathclose%
\pgfusepath{stroke,fill}%
}%
\begin{pgfscope}%
\pgfsys@transformshift{0.760451in}{2.195935in}%
\pgfsys@useobject{currentmarker}{}%
\end{pgfscope}%
\begin{pgfscope}%
\pgfsys@transformshift{1.530589in}{1.345376in}%
\pgfsys@useobject{currentmarker}{}%
\end{pgfscope}%
\end{pgfscope}%
\begin{pgfscope}%
\pgfpathrectangle{\pgfqpoint{0.616050in}{0.916500in}}{\pgfqpoint{2.647350in}{1.306500in}}%
\pgfusepath{clip}%
\pgfsetrectcap%
\pgfsetroundjoin%
\pgfsetlinewidth{0.903375pt}%
\definecolor{currentstroke}{rgb}{0.784314,0.784314,0.784314}%
\pgfsetstrokecolor{currentstroke}%
\pgfsetdash{}{0pt}%
\pgfpathmoveto{\pgfqpoint{0.760451in}{2.154618in}}%
\pgfpathlineto{\pgfqpoint{1.145520in}{2.017121in}}%
\pgfpathlineto{\pgfqpoint{1.530589in}{1.742699in}}%
\pgfusepath{stroke}%
\end{pgfscope}%
\begin{pgfscope}%
\pgfpathrectangle{\pgfqpoint{0.616050in}{0.916500in}}{\pgfqpoint{2.647350in}{1.306500in}}%
\pgfusepath{clip}%
\pgfsetbuttcap%
\pgfsetmiterjoin%
\definecolor{currentfill}{rgb}{0.784314,0.784314,0.784314}%
\pgfsetfillcolor{currentfill}%
\pgfsetlinewidth{0.702625pt}%
\definecolor{currentstroke}{rgb}{0.784314,0.784314,0.784314}%
\pgfsetstrokecolor{currentstroke}%
\pgfsetdash{}{0pt}%
\pgfsys@defobject{currentmarker}{\pgfqpoint{-0.022222in}{-0.022222in}}{\pgfqpoint{0.022222in}{0.022222in}}{%
\pgfpathmoveto{\pgfqpoint{-0.022222in}{-0.022222in}}%
\pgfpathlineto{\pgfqpoint{0.022222in}{-0.022222in}}%
\pgfpathlineto{\pgfqpoint{0.022222in}{0.022222in}}%
\pgfpathlineto{\pgfqpoint{-0.022222in}{0.022222in}}%
\pgfpathlineto{\pgfqpoint{-0.022222in}{-0.022222in}}%
\pgfpathclose%
\pgfusepath{stroke,fill}%
}%
\begin{pgfscope}%
\pgfsys@transformshift{0.760451in}{2.154618in}%
\pgfsys@useobject{currentmarker}{}%
\end{pgfscope}%
\begin{pgfscope}%
\pgfsys@transformshift{1.145520in}{2.017121in}%
\pgfsys@useobject{currentmarker}{}%
\end{pgfscope}%
\begin{pgfscope}%
\pgfsys@transformshift{1.530589in}{1.742699in}%
\pgfsys@useobject{currentmarker}{}%
\end{pgfscope}%
\end{pgfscope}%
\begin{pgfscope}%
\pgfpathrectangle{\pgfqpoint{0.616050in}{0.916500in}}{\pgfqpoint{2.647350in}{1.306500in}}%
\pgfusepath{clip}%
\pgfsetrectcap%
\pgfsetroundjoin%
\pgfsetlinewidth{0.903375pt}%
\definecolor{currentstroke}{rgb}{0.478431,0.478431,0.478431}%
\pgfsetstrokecolor{currentstroke}%
\pgfsetdash{}{0pt}%
\pgfpathmoveto{\pgfqpoint{0.760451in}{2.196124in}}%
\pgfpathlineto{\pgfqpoint{1.530589in}{1.368556in}}%
\pgfusepath{stroke}%
\end{pgfscope}%
\begin{pgfscope}%
\pgfpathrectangle{\pgfqpoint{0.616050in}{0.916500in}}{\pgfqpoint{2.647350in}{1.306500in}}%
\pgfusepath{clip}%
\pgfsetbuttcap%
\pgfsetmiterjoin%
\definecolor{currentfill}{rgb}{0.478431,0.478431,0.478431}%
\pgfsetfillcolor{currentfill}%
\pgfsetlinewidth{0.702625pt}%
\definecolor{currentstroke}{rgb}{0.478431,0.478431,0.478431}%
\pgfsetstrokecolor{currentstroke}%
\pgfsetdash{}{0pt}%
\pgfsys@defobject{currentmarker}{\pgfqpoint{-0.022222in}{-0.022222in}}{\pgfqpoint{0.022222in}{0.022222in}}{%
\pgfpathmoveto{\pgfqpoint{-0.007407in}{-0.022222in}}%
\pgfpathlineto{\pgfqpoint{0.007407in}{-0.022222in}}%
\pgfpathlineto{\pgfqpoint{0.007407in}{-0.007407in}}%
\pgfpathlineto{\pgfqpoint{0.022222in}{-0.007407in}}%
\pgfpathlineto{\pgfqpoint{0.022222in}{0.007407in}}%
\pgfpathlineto{\pgfqpoint{0.007407in}{0.007407in}}%
\pgfpathlineto{\pgfqpoint{0.007407in}{0.022222in}}%
\pgfpathlineto{\pgfqpoint{-0.007407in}{0.022222in}}%
\pgfpathlineto{\pgfqpoint{-0.007407in}{0.007407in}}%
\pgfpathlineto{\pgfqpoint{-0.022222in}{0.007407in}}%
\pgfpathlineto{\pgfqpoint{-0.022222in}{-0.007407in}}%
\pgfpathlineto{\pgfqpoint{-0.007407in}{-0.007407in}}%
\pgfpathlineto{\pgfqpoint{-0.007407in}{-0.022222in}}%
\pgfpathclose%
\pgfusepath{stroke,fill}%
}%
\begin{pgfscope}%
\pgfsys@transformshift{0.760451in}{2.196124in}%
\pgfsys@useobject{currentmarker}{}%
\end{pgfscope}%
\begin{pgfscope}%
\pgfsys@transformshift{1.530589in}{1.368556in}%
\pgfsys@useobject{currentmarker}{}%
\end{pgfscope}%
\end{pgfscope}%
\begin{pgfscope}%
\pgfsetrectcap%
\pgfsetmiterjoin%
\pgfsetlinewidth{0.803000pt}%
\definecolor{currentstroke}{rgb}{0.101961,0.101961,0.101961}%
\pgfsetstrokecolor{currentstroke}%
\pgfsetdash{}{0pt}%
\pgfpathmoveto{\pgfqpoint{0.616050in}{0.916500in}}%
\pgfpathlineto{\pgfqpoint{0.616050in}{2.223000in}}%
\pgfusepath{stroke}%
\end{pgfscope}%
\begin{pgfscope}%
\pgfsetrectcap%
\pgfsetmiterjoin%
\pgfsetlinewidth{0.803000pt}%
\definecolor{currentstroke}{rgb}{0.101961,0.101961,0.101961}%
\pgfsetstrokecolor{currentstroke}%
\pgfsetdash{}{0pt}%
\pgfpathmoveto{\pgfqpoint{0.616050in}{0.916500in}}%
\pgfpathlineto{\pgfqpoint{3.263400in}{0.916500in}}%
\pgfusepath{stroke}%
\end{pgfscope}%
\begin{pgfscope}%
\pgfpathrectangle{\pgfqpoint{0.616050in}{0.916500in}}{\pgfqpoint{2.647350in}{1.306500in}}%
\pgfusepath{clip}%
\pgfsetbuttcap%
\pgfsetroundjoin%
\pgfsetlinewidth{0.903375pt}%
\definecolor{currentstroke}{rgb}{0.156863,0.356863,0.568627}%
\pgfsetstrokecolor{currentstroke}%
\pgfsetdash{{3.330000pt}{1.440000pt}}{0.000000pt}%
\pgfpathmoveto{\pgfqpoint{1.915658in}{1.573930in}}%
\pgfpathlineto{\pgfqpoint{2.300727in}{1.457562in}}%
\pgfpathlineto{\pgfqpoint{3.070865in}{1.125687in}}%
\pgfusepath{stroke}%
\end{pgfscope}%
\begin{pgfscope}%
\pgfpathrectangle{\pgfqpoint{0.616050in}{0.916500in}}{\pgfqpoint{2.647350in}{1.306500in}}%
\pgfusepath{clip}%
\pgfsetbuttcap%
\pgfsetroundjoin%
\definecolor{currentfill}{rgb}{1.000000,1.000000,1.000000}%
\pgfsetfillcolor{currentfill}%
\pgfsetlinewidth{0.702625pt}%
\definecolor{currentstroke}{rgb}{0.156863,0.356863,0.568627}%
\pgfsetstrokecolor{currentstroke}%
\pgfsetdash{}{0pt}%
\pgfsys@defobject{currentmarker}{\pgfqpoint{-0.022222in}{-0.022222in}}{\pgfqpoint{0.022222in}{0.022222in}}{%
\pgfpathmoveto{\pgfqpoint{0.000000in}{-0.022222in}}%
\pgfpathcurveto{\pgfqpoint{0.005893in}{-0.022222in}}{\pgfqpoint{0.011546in}{-0.019881in}}{\pgfqpoint{0.015713in}{-0.015713in}}%
\pgfpathcurveto{\pgfqpoint{0.019881in}{-0.011546in}}{\pgfqpoint{0.022222in}{-0.005893in}}{\pgfqpoint{0.022222in}{0.000000in}}%
\pgfpathcurveto{\pgfqpoint{0.022222in}{0.005893in}}{\pgfqpoint{0.019881in}{0.011546in}}{\pgfqpoint{0.015713in}{0.015713in}}%
\pgfpathcurveto{\pgfqpoint{0.011546in}{0.019881in}}{\pgfqpoint{0.005893in}{0.022222in}}{\pgfqpoint{0.000000in}{0.022222in}}%
\pgfpathcurveto{\pgfqpoint{-0.005893in}{0.022222in}}{\pgfqpoint{-0.011546in}{0.019881in}}{\pgfqpoint{-0.015713in}{0.015713in}}%
\pgfpathcurveto{\pgfqpoint{-0.019881in}{0.011546in}}{\pgfqpoint{-0.022222in}{0.005893in}}{\pgfqpoint{-0.022222in}{0.000000in}}%
\pgfpathcurveto{\pgfqpoint{-0.022222in}{-0.005893in}}{\pgfqpoint{-0.019881in}{-0.011546in}}{\pgfqpoint{-0.015713in}{-0.015713in}}%
\pgfpathcurveto{\pgfqpoint{-0.011546in}{-0.019881in}}{\pgfqpoint{-0.005893in}{-0.022222in}}{\pgfqpoint{0.000000in}{-0.022222in}}%
\pgfpathlineto{\pgfqpoint{0.000000in}{-0.022222in}}%
\pgfpathclose%
\pgfusepath{stroke,fill}%
}%
\begin{pgfscope}%
\pgfsys@transformshift{1.915658in}{1.573930in}%
\pgfsys@useobject{currentmarker}{}%
\end{pgfscope}%
\begin{pgfscope}%
\pgfsys@transformshift{2.300727in}{1.457562in}%
\pgfsys@useobject{currentmarker}{}%
\end{pgfscope}%
\begin{pgfscope}%
\pgfsys@transformshift{3.070865in}{1.125687in}%
\pgfsys@useobject{currentmarker}{}%
\end{pgfscope}%
\end{pgfscope}%
\begin{pgfscope}%
\pgfpathrectangle{\pgfqpoint{0.616050in}{0.916500in}}{\pgfqpoint{2.647350in}{1.306500in}}%
\pgfusepath{clip}%
\pgfsetrectcap%
\pgfsetroundjoin%
\pgfsetlinewidth{0.903375pt}%
\definecolor{currentstroke}{rgb}{0.156863,0.356863,0.568627}%
\pgfsetstrokecolor{currentstroke}%
\pgfsetdash{}{0pt}%
\pgfpathmoveto{\pgfqpoint{1.530589in}{1.865348in}}%
\pgfpathlineto{\pgfqpoint{1.915658in}{1.609717in}}%
\pgfpathlineto{\pgfqpoint{2.300727in}{1.415364in}}%
\pgfusepath{stroke}%
\end{pgfscope}%
\begin{pgfscope}%
\pgfpathrectangle{\pgfqpoint{0.616050in}{0.916500in}}{\pgfqpoint{2.647350in}{1.306500in}}%
\pgfusepath{clip}%
\pgfsetbuttcap%
\pgfsetroundjoin%
\definecolor{currentfill}{rgb}{0.156863,0.356863,0.568627}%
\pgfsetfillcolor{currentfill}%
\pgfsetlinewidth{0.702625pt}%
\definecolor{currentstroke}{rgb}{0.156863,0.356863,0.568627}%
\pgfsetstrokecolor{currentstroke}%
\pgfsetdash{}{0pt}%
\pgfsys@defobject{currentmarker}{\pgfqpoint{-0.022222in}{-0.022222in}}{\pgfqpoint{0.022222in}{0.022222in}}{%
\pgfpathmoveto{\pgfqpoint{0.000000in}{-0.022222in}}%
\pgfpathcurveto{\pgfqpoint{0.005893in}{-0.022222in}}{\pgfqpoint{0.011546in}{-0.019881in}}{\pgfqpoint{0.015713in}{-0.015713in}}%
\pgfpathcurveto{\pgfqpoint{0.019881in}{-0.011546in}}{\pgfqpoint{0.022222in}{-0.005893in}}{\pgfqpoint{0.022222in}{0.000000in}}%
\pgfpathcurveto{\pgfqpoint{0.022222in}{0.005893in}}{\pgfqpoint{0.019881in}{0.011546in}}{\pgfqpoint{0.015713in}{0.015713in}}%
\pgfpathcurveto{\pgfqpoint{0.011546in}{0.019881in}}{\pgfqpoint{0.005893in}{0.022222in}}{\pgfqpoint{0.000000in}{0.022222in}}%
\pgfpathcurveto{\pgfqpoint{-0.005893in}{0.022222in}}{\pgfqpoint{-0.011546in}{0.019881in}}{\pgfqpoint{-0.015713in}{0.015713in}}%
\pgfpathcurveto{\pgfqpoint{-0.019881in}{0.011546in}}{\pgfqpoint{-0.022222in}{0.005893in}}{\pgfqpoint{-0.022222in}{0.000000in}}%
\pgfpathcurveto{\pgfqpoint{-0.022222in}{-0.005893in}}{\pgfqpoint{-0.019881in}{-0.011546in}}{\pgfqpoint{-0.015713in}{-0.015713in}}%
\pgfpathcurveto{\pgfqpoint{-0.011546in}{-0.019881in}}{\pgfqpoint{-0.005893in}{-0.022222in}}{\pgfqpoint{0.000000in}{-0.022222in}}%
\pgfpathlineto{\pgfqpoint{0.000000in}{-0.022222in}}%
\pgfpathclose%
\pgfusepath{stroke,fill}%
}%
\begin{pgfscope}%
\pgfsys@transformshift{1.530589in}{1.865348in}%
\pgfsys@useobject{currentmarker}{}%
\end{pgfscope}%
\begin{pgfscope}%
\pgfsys@transformshift{1.915658in}{1.609717in}%
\pgfsys@useobject{currentmarker}{}%
\end{pgfscope}%
\begin{pgfscope}%
\pgfsys@transformshift{2.300727in}{1.415364in}%
\pgfsys@useobject{currentmarker}{}%
\end{pgfscope}%
\end{pgfscope}%
\begin{pgfscope}%
\pgfpathrectangle{\pgfqpoint{0.616050in}{0.916500in}}{\pgfqpoint{2.647350in}{1.306500in}}%
\pgfusepath{clip}%
\pgfsetbuttcap%
\pgfsetroundjoin%
\pgfsetlinewidth{0.903375pt}%
\definecolor{currentstroke}{rgb}{0.156863,0.356863,0.568627}%
\pgfsetstrokecolor{currentstroke}%
\pgfsetdash{{3.330000pt}{1.440000pt}}{0.000000pt}%
\pgfpathmoveto{\pgfqpoint{1.915658in}{1.823050in}}%
\pgfusepath{stroke}%
\end{pgfscope}%
\begin{pgfscope}%
\pgfpathrectangle{\pgfqpoint{0.616050in}{0.916500in}}{\pgfqpoint{2.647350in}{1.306500in}}%
\pgfusepath{clip}%
\pgfsetbuttcap%
\pgfsetmiterjoin%
\definecolor{currentfill}{rgb}{1.000000,1.000000,1.000000}%
\pgfsetfillcolor{currentfill}%
\pgfsetlinewidth{0.702625pt}%
\definecolor{currentstroke}{rgb}{0.156863,0.356863,0.568627}%
\pgfsetstrokecolor{currentstroke}%
\pgfsetdash{}{0pt}%
\pgfsys@defobject{currentmarker}{\pgfqpoint{-0.031427in}{-0.031427in}}{\pgfqpoint{0.031427in}{0.031427in}}{%
\pgfpathmoveto{\pgfqpoint{-0.000000in}{-0.031427in}}%
\pgfpathlineto{\pgfqpoint{0.031427in}{0.000000in}}%
\pgfpathlineto{\pgfqpoint{0.000000in}{0.031427in}}%
\pgfpathlineto{\pgfqpoint{-0.031427in}{0.000000in}}%
\pgfpathlineto{\pgfqpoint{-0.000000in}{-0.031427in}}%
\pgfpathclose%
\pgfusepath{stroke,fill}%
}%
\begin{pgfscope}%
\pgfsys@transformshift{1.915658in}{1.823050in}%
\pgfsys@useobject{currentmarker}{}%
\end{pgfscope}%
\end{pgfscope}%
\begin{pgfscope}%
\pgfpathrectangle{\pgfqpoint{0.616050in}{0.916500in}}{\pgfqpoint{2.647350in}{1.306500in}}%
\pgfusepath{clip}%
\pgfsetrectcap%
\pgfsetroundjoin%
\pgfsetlinewidth{0.903375pt}%
\definecolor{currentstroke}{rgb}{0.156863,0.356863,0.568627}%
\pgfsetstrokecolor{currentstroke}%
\pgfsetdash{}{0pt}%
\pgfpathmoveto{\pgfqpoint{1.530589in}{1.874734in}}%
\pgfusepath{stroke}%
\end{pgfscope}%
\begin{pgfscope}%
\pgfpathrectangle{\pgfqpoint{0.616050in}{0.916500in}}{\pgfqpoint{2.647350in}{1.306500in}}%
\pgfusepath{clip}%
\pgfsetbuttcap%
\pgfsetmiterjoin%
\definecolor{currentfill}{rgb}{0.156863,0.356863,0.568627}%
\pgfsetfillcolor{currentfill}%
\pgfsetlinewidth{0.702625pt}%
\definecolor{currentstroke}{rgb}{0.156863,0.356863,0.568627}%
\pgfsetstrokecolor{currentstroke}%
\pgfsetdash{}{0pt}%
\pgfsys@defobject{currentmarker}{\pgfqpoint{-0.031427in}{-0.031427in}}{\pgfqpoint{0.031427in}{0.031427in}}{%
\pgfpathmoveto{\pgfqpoint{-0.000000in}{-0.031427in}}%
\pgfpathlineto{\pgfqpoint{0.031427in}{0.000000in}}%
\pgfpathlineto{\pgfqpoint{0.000000in}{0.031427in}}%
\pgfpathlineto{\pgfqpoint{-0.031427in}{0.000000in}}%
\pgfpathlineto{\pgfqpoint{-0.000000in}{-0.031427in}}%
\pgfpathclose%
\pgfusepath{stroke,fill}%
}%
\begin{pgfscope}%
\pgfsys@transformshift{1.530589in}{1.874734in}%
\pgfsys@useobject{currentmarker}{}%
\end{pgfscope}%
\end{pgfscope}%
\begin{pgfscope}%
\definecolor{textcolor}{rgb}{0.549020,0.549020,0.549020}%
\pgfsetstrokecolor{textcolor}%
\pgfsetfillcolor{textcolor}%
\pgftext[x=3.296700in,y=3.841500in,right,top]{\color{textcolor}\sffamily\fontsize{5.500000}{6.600000}\selectfont OLD DATA}%
\end{pgfscope}%
\begin{pgfscope}%
\pgfsetrectcap%
\pgfsetroundjoin%
\pgfsetlinewidth{0.903375pt}%
\definecolor{currentstroke}{rgb}{0.478431,0.305882,0.611765}%
\pgfsetstrokecolor{currentstroke}%
\pgfsetdash{}{0pt}%
\pgfpathmoveto{\pgfqpoint{0.111443in}{0.398500in}}%
\pgfpathlineto{\pgfqpoint{0.174637in}{0.398500in}}%
\pgfpathlineto{\pgfqpoint{0.237832in}{0.398500in}}%
\pgfusepath{stroke}%
\end{pgfscope}%
\begin{pgfscope}%
\pgfsetbuttcap%
\pgfsetmiterjoin%
\definecolor{currentfill}{rgb}{0.478431,0.305882,0.611765}%
\pgfsetfillcolor{currentfill}%
\pgfsetlinewidth{1.003750pt}%
\definecolor{currentstroke}{rgb}{0.478431,0.305882,0.611765}%
\pgfsetstrokecolor{currentstroke}%
\pgfsetdash{}{0pt}%
\pgfsys@defobject{currentmarker}{\pgfqpoint{-0.022222in}{-0.022222in}}{\pgfqpoint{0.022222in}{0.022222in}}{%
\pgfpathmoveto{\pgfqpoint{-0.011111in}{-0.022222in}}%
\pgfpathlineto{\pgfqpoint{0.000000in}{-0.011111in}}%
\pgfpathlineto{\pgfqpoint{0.011111in}{-0.022222in}}%
\pgfpathlineto{\pgfqpoint{0.022222in}{-0.011111in}}%
\pgfpathlineto{\pgfqpoint{0.011111in}{0.000000in}}%
\pgfpathlineto{\pgfqpoint{0.022222in}{0.011111in}}%
\pgfpathlineto{\pgfqpoint{0.011111in}{0.022222in}}%
\pgfpathlineto{\pgfqpoint{0.000000in}{0.011111in}}%
\pgfpathlineto{\pgfqpoint{-0.011111in}{0.022222in}}%
\pgfpathlineto{\pgfqpoint{-0.022222in}{0.011111in}}%
\pgfpathlineto{\pgfqpoint{-0.011111in}{0.000000in}}%
\pgfpathlineto{\pgfqpoint{-0.022222in}{-0.011111in}}%
\pgfpathlineto{\pgfqpoint{-0.011111in}{-0.022222in}}%
\pgfpathclose%
\pgfusepath{stroke,fill}%
}%
\begin{pgfscope}%
\pgfsys@transformshift{0.174637in}{0.398500in}%
\pgfsys@useobject{currentmarker}{}%
\end{pgfscope}%
\end{pgfscope}%
\begin{pgfscope}%
\definecolor{textcolor}{rgb}{0.101961,0.101961,0.101961}%
\pgfsetstrokecolor{textcolor}%
\pgfsetfillcolor{textcolor}%
\pgftext[x=0.266999in,y=0.364472in,left,base]{\color{textcolor}\sffamily\fontsize{7.000000}{8.400000}\selectfont On-device only}%
\end{pgfscope}%
\begin{pgfscope}%
\pgfsetrectcap%
\pgfsetroundjoin%
\pgfsetlinewidth{0.903375pt}%
\definecolor{currentstroke}{rgb}{0.784314,0.784314,0.784314}%
\pgfsetstrokecolor{currentstroke}%
\pgfsetdash{}{0pt}%
\pgfpathmoveto{\pgfqpoint{0.111443in}{0.258068in}}%
\pgfpathlineto{\pgfqpoint{0.174637in}{0.258068in}}%
\pgfpathlineto{\pgfqpoint{0.237832in}{0.258068in}}%
\pgfusepath{stroke}%
\end{pgfscope}%
\begin{pgfscope}%
\pgfsetbuttcap%
\pgfsetmiterjoin%
\definecolor{currentfill}{rgb}{0.784314,0.784314,0.784314}%
\pgfsetfillcolor{currentfill}%
\pgfsetlinewidth{1.003750pt}%
\definecolor{currentstroke}{rgb}{0.784314,0.784314,0.784314}%
\pgfsetstrokecolor{currentstroke}%
\pgfsetdash{}{0pt}%
\pgfsys@defobject{currentmarker}{\pgfqpoint{-0.022222in}{-0.022222in}}{\pgfqpoint{0.022222in}{0.022222in}}{%
\pgfpathmoveto{\pgfqpoint{0.000000in}{0.022222in}}%
\pgfpathlineto{\pgfqpoint{-0.022222in}{-0.022222in}}%
\pgfpathlineto{\pgfqpoint{0.022222in}{-0.022222in}}%
\pgfpathlineto{\pgfqpoint{0.000000in}{0.022222in}}%
\pgfpathclose%
\pgfusepath{stroke,fill}%
}%
\begin{pgfscope}%
\pgfsys@transformshift{0.174637in}{0.258068in}%
\pgfsys@useobject{currentmarker}{}%
\end{pgfscope}%
\end{pgfscope}%
\begin{pgfscope}%
\definecolor{textcolor}{rgb}{0.101961,0.101961,0.101961}%
\pgfsetstrokecolor{textcolor}%
\pgfsetfillcolor{textcolor}%
\pgftext[x=0.266999in,y=0.224040in,left,base]{\color{textcolor}\sffamily\fontsize{7.000000}{8.400000}\selectfont SimpleStream~\cite{simplestream2026}}%
\end{pgfscope}%
\begin{pgfscope}%
\pgfsetrectcap%
\pgfsetroundjoin%
\pgfsetlinewidth{0.903375pt}%
\definecolor{currentstroke}{rgb}{0.784314,0.784314,0.784314}%
\pgfsetstrokecolor{currentstroke}%
\pgfsetdash{}{0pt}%
\pgfpathmoveto{\pgfqpoint{0.111443in}{0.112235in}}%
\pgfpathlineto{\pgfqpoint{0.174637in}{0.112235in}}%
\pgfpathlineto{\pgfqpoint{0.237832in}{0.112235in}}%
\pgfusepath{stroke}%
\end{pgfscope}%
\begin{pgfscope}%
\pgfsetbuttcap%
\pgfsetmiterjoin%
\definecolor{currentfill}{rgb}{0.784314,0.784314,0.784314}%
\pgfsetfillcolor{currentfill}%
\pgfsetlinewidth{1.003750pt}%
\definecolor{currentstroke}{rgb}{0.784314,0.784314,0.784314}%
\pgfsetstrokecolor{currentstroke}%
\pgfsetdash{}{0pt}%
\pgfsys@defobject{currentmarker}{\pgfqpoint{-0.022222in}{-0.022222in}}{\pgfqpoint{0.022222in}{0.022222in}}{%
\pgfpathmoveto{\pgfqpoint{-0.000000in}{-0.022222in}}%
\pgfpathlineto{\pgfqpoint{0.022222in}{0.022222in}}%
\pgfpathlineto{\pgfqpoint{-0.022222in}{0.022222in}}%
\pgfpathlineto{\pgfqpoint{-0.000000in}{-0.022222in}}%
\pgfpathclose%
\pgfusepath{stroke,fill}%
}%
\begin{pgfscope}%
\pgfsys@transformshift{0.174637in}{0.112235in}%
\pgfsys@useobject{currentmarker}{}%
\end{pgfscope}%
\end{pgfscope}%
\begin{pgfscope}%
\definecolor{textcolor}{rgb}{0.101961,0.101961,0.101961}%
\pgfsetstrokecolor{textcolor}%
\pgfsetfillcolor{textcolor}%
\pgftext[x=0.266999in,y=0.078207in,left,base]{\color{textcolor}\sffamily\fontsize{7.000000}{8.400000}\selectfont ECO-LLM~\cite{ecollm2025}}%
\end{pgfscope}%
\begin{pgfscope}%
\pgfsetrectcap%
\pgfsetroundjoin%
\pgfsetlinewidth{0.903375pt}%
\definecolor{currentstroke}{rgb}{0.784314,0.784314,0.784314}%
\pgfsetstrokecolor{currentstroke}%
\pgfsetdash{}{0pt}%
\pgfpathmoveto{\pgfqpoint{1.084542in}{0.398500in}}%
\pgfpathlineto{\pgfqpoint{1.147736in}{0.398500in}}%
\pgfpathlineto{\pgfqpoint{1.210931in}{0.398500in}}%
\pgfusepath{stroke}%
\end{pgfscope}%
\begin{pgfscope}%
\pgfsetbuttcap%
\pgfsetmiterjoin%
\definecolor{currentfill}{rgb}{0.784314,0.784314,0.784314}%
\pgfsetfillcolor{currentfill}%
\pgfsetlinewidth{1.003750pt}%
\definecolor{currentstroke}{rgb}{0.784314,0.784314,0.784314}%
\pgfsetstrokecolor{currentstroke}%
\pgfsetdash{}{0pt}%
\pgfsys@defobject{currentmarker}{\pgfqpoint{-0.022222in}{-0.022222in}}{\pgfqpoint{0.022222in}{0.022222in}}{%
\pgfpathmoveto{\pgfqpoint{-0.022222in}{-0.022222in}}%
\pgfpathlineto{\pgfqpoint{0.022222in}{-0.022222in}}%
\pgfpathlineto{\pgfqpoint{0.022222in}{0.022222in}}%
\pgfpathlineto{\pgfqpoint{-0.022222in}{0.022222in}}%
\pgfpathlineto{\pgfqpoint{-0.022222in}{-0.022222in}}%
\pgfpathclose%
\pgfusepath{stroke,fill}%
}%
\begin{pgfscope}%
\pgfsys@transformshift{1.147736in}{0.398500in}%
\pgfsys@useobject{currentmarker}{}%
\end{pgfscope}%
\end{pgfscope}%
\begin{pgfscope}%
\definecolor{textcolor}{rgb}{0.101961,0.101961,0.101961}%
\pgfsetstrokecolor{textcolor}%
\pgfsetfillcolor{textcolor}%
\pgftext[x=1.240097in,y=0.364472in,left,base]{\color{textcolor}\sffamily\fontsize{7.000000}{8.400000}\selectfont LLF~\cite{mok1983llf}}%
\end{pgfscope}%
\begin{pgfscope}%
\pgfsetrectcap%
\pgfsetroundjoin%
\pgfsetlinewidth{0.903375pt}%
\definecolor{currentstroke}{rgb}{0.478431,0.478431,0.478431}%
\pgfsetstrokecolor{currentstroke}%
\pgfsetdash{}{0pt}%
\pgfpathmoveto{\pgfqpoint{1.084542in}{0.252667in}}%
\pgfpathlineto{\pgfqpoint{1.147736in}{0.252667in}}%
\pgfpathlineto{\pgfqpoint{1.210931in}{0.252667in}}%
\pgfusepath{stroke}%
\end{pgfscope}%
\begin{pgfscope}%
\pgfsetbuttcap%
\pgfsetmiterjoin%
\definecolor{currentfill}{rgb}{0.478431,0.478431,0.478431}%
\pgfsetfillcolor{currentfill}%
\pgfsetlinewidth{1.003750pt}%
\definecolor{currentstroke}{rgb}{0.478431,0.478431,0.478431}%
\pgfsetstrokecolor{currentstroke}%
\pgfsetdash{}{0pt}%
\pgfsys@defobject{currentmarker}{\pgfqpoint{-0.022222in}{-0.022222in}}{\pgfqpoint{0.022222in}{0.022222in}}{%
\pgfpathmoveto{\pgfqpoint{-0.007407in}{-0.022222in}}%
\pgfpathlineto{\pgfqpoint{0.007407in}{-0.022222in}}%
\pgfpathlineto{\pgfqpoint{0.007407in}{-0.007407in}}%
\pgfpathlineto{\pgfqpoint{0.022222in}{-0.007407in}}%
\pgfpathlineto{\pgfqpoint{0.022222in}{0.007407in}}%
\pgfpathlineto{\pgfqpoint{0.007407in}{0.007407in}}%
\pgfpathlineto{\pgfqpoint{0.007407in}{0.022222in}}%
\pgfpathlineto{\pgfqpoint{-0.007407in}{0.022222in}}%
\pgfpathlineto{\pgfqpoint{-0.007407in}{0.007407in}}%
\pgfpathlineto{\pgfqpoint{-0.022222in}{0.007407in}}%
\pgfpathlineto{\pgfqpoint{-0.022222in}{-0.007407in}}%
\pgfpathlineto{\pgfqpoint{-0.007407in}{-0.007407in}}%
\pgfpathlineto{\pgfqpoint{-0.007407in}{-0.022222in}}%
\pgfpathclose%
\pgfusepath{stroke,fill}%
}%
\begin{pgfscope}%
\pgfsys@transformshift{1.147736in}{0.252667in}%
\pgfsys@useobject{currentmarker}{}%
\end{pgfscope}%
\end{pgfscope}%
\begin{pgfscope}%
\definecolor{textcolor}{rgb}{0.101961,0.101961,0.101961}%
\pgfsetstrokecolor{textcolor}%
\pgfsetfillcolor{textcolor}%
\pgftext[x=1.240097in,y=0.218639in,left,base]{\color{textcolor}\sffamily\fontsize{7.000000}{8.400000}\selectfont Mooncake~\cite{mooncake2025}}%
\end{pgfscope}%
\begin{pgfscope}%
\pgfsetrectcap%
\pgfsetroundjoin%
\pgfsetlinewidth{0.903375pt}%
\definecolor{currentstroke}{rgb}{0.156863,0.356863,0.568627}%
\pgfsetstrokecolor{currentstroke}%
\pgfsetdash{}{0pt}%
\pgfpathmoveto{\pgfqpoint{1.084542in}{0.106833in}}%
\pgfpathlineto{\pgfqpoint{1.147736in}{0.106833in}}%
\pgfpathlineto{\pgfqpoint{1.210931in}{0.106833in}}%
\pgfusepath{stroke}%
\end{pgfscope}%
\begin{pgfscope}%
\pgfsetbuttcap%
\pgfsetroundjoin%
\definecolor{currentfill}{rgb}{0.156863,0.356863,0.568627}%
\pgfsetfillcolor{currentfill}%
\pgfsetlinewidth{1.003750pt}%
\definecolor{currentstroke}{rgb}{0.156863,0.356863,0.568627}%
\pgfsetstrokecolor{currentstroke}%
\pgfsetdash{}{0pt}%
\pgfsys@defobject{currentmarker}{\pgfqpoint{-0.022222in}{-0.022222in}}{\pgfqpoint{0.022222in}{0.022222in}}{%
\pgfpathmoveto{\pgfqpoint{0.000000in}{-0.022222in}}%
\pgfpathcurveto{\pgfqpoint{0.005893in}{-0.022222in}}{\pgfqpoint{0.011546in}{-0.019881in}}{\pgfqpoint{0.015713in}{-0.015713in}}%
\pgfpathcurveto{\pgfqpoint{0.019881in}{-0.011546in}}{\pgfqpoint{0.022222in}{-0.005893in}}{\pgfqpoint{0.022222in}{0.000000in}}%
\pgfpathcurveto{\pgfqpoint{0.022222in}{0.005893in}}{\pgfqpoint{0.019881in}{0.011546in}}{\pgfqpoint{0.015713in}{0.015713in}}%
\pgfpathcurveto{\pgfqpoint{0.011546in}{0.019881in}}{\pgfqpoint{0.005893in}{0.022222in}}{\pgfqpoint{0.000000in}{0.022222in}}%
\pgfpathcurveto{\pgfqpoint{-0.005893in}{0.022222in}}{\pgfqpoint{-0.011546in}{0.019881in}}{\pgfqpoint{-0.015713in}{0.015713in}}%
\pgfpathcurveto{\pgfqpoint{-0.019881in}{0.011546in}}{\pgfqpoint{-0.022222in}{0.005893in}}{\pgfqpoint{-0.022222in}{0.000000in}}%
\pgfpathcurveto{\pgfqpoint{-0.022222in}{-0.005893in}}{\pgfqpoint{-0.019881in}{-0.011546in}}{\pgfqpoint{-0.015713in}{-0.015713in}}%
\pgfpathcurveto{\pgfqpoint{-0.011546in}{-0.019881in}}{\pgfqpoint{-0.005893in}{-0.022222in}}{\pgfqpoint{0.000000in}{-0.022222in}}%
\pgfpathlineto{\pgfqpoint{0.000000in}{-0.022222in}}%
\pgfpathclose%
\pgfusepath{stroke,fill}%
}%
\begin{pgfscope}%
\pgfsys@transformshift{1.147736in}{0.106833in}%
\pgfsys@useobject{currentmarker}{}%
\end{pgfscope}%
\end{pgfscope}%
\begin{pgfscope}%
\definecolor{textcolor}{rgb}{0.101961,0.101961,0.101961}%
\pgfsetstrokecolor{textcolor}%
\pgfsetfillcolor{textcolor}%
\pgftext[x=1.240097in,y=0.072806in,left,base]{\color{textcolor}\sffamily\fontsize{7.000000}{8.400000}\selectfont COVEX (old: rung 3)}%
\end{pgfscope}%
\begin{pgfscope}%
\pgfsetrectcap%
\pgfsetroundjoin%
\pgfsetlinewidth{0.903375pt}%
\definecolor{currentstroke}{rgb}{0.156863,0.356863,0.568627}%
\pgfsetstrokecolor{currentstroke}%
\pgfsetdash{}{0pt}%
\pgfpathmoveto{\pgfqpoint{2.245954in}{0.398500in}}%
\pgfpathlineto{\pgfqpoint{2.309148in}{0.398500in}}%
\pgfpathlineto{\pgfqpoint{2.372343in}{0.398500in}}%
\pgfusepath{stroke}%
\end{pgfscope}%
\begin{pgfscope}%
\pgfsetbuttcap%
\pgfsetmiterjoin%
\definecolor{currentfill}{rgb}{0.156863,0.356863,0.568627}%
\pgfsetfillcolor{currentfill}%
\pgfsetlinewidth{1.003750pt}%
\definecolor{currentstroke}{rgb}{0.156863,0.356863,0.568627}%
\pgfsetstrokecolor{currentstroke}%
\pgfsetdash{}{0pt}%
\pgfsys@defobject{currentmarker}{\pgfqpoint{-0.031427in}{-0.031427in}}{\pgfqpoint{0.031427in}{0.031427in}}{%
\pgfpathmoveto{\pgfqpoint{-0.000000in}{-0.031427in}}%
\pgfpathlineto{\pgfqpoint{0.031427in}{0.000000in}}%
\pgfpathlineto{\pgfqpoint{0.000000in}{0.031427in}}%
\pgfpathlineto{\pgfqpoint{-0.031427in}{0.000000in}}%
\pgfpathlineto{\pgfqpoint{-0.000000in}{-0.031427in}}%
\pgfpathclose%
\pgfusepath{stroke,fill}%
}%
\begin{pgfscope}%
\pgfsys@transformshift{2.309148in}{0.398500in}%
\pgfsys@useobject{currentmarker}{}%
\end{pgfscope}%
\end{pgfscope}%
\begin{pgfscope}%
\definecolor{textcolor}{rgb}{0.101961,0.101961,0.101961}%
\pgfsetstrokecolor{textcolor}%
\pgfsetfillcolor{textcolor}%
\pgftext[x=2.401509in,y=0.364472in,left,base]{\color{textcolor}\sffamily\fontsize{7.000000}{8.400000}\selectfont COVEX (dev)}%
\end{pgfscope}%
\begin{pgfscope}%
\pgfsetrectcap%
\pgfsetroundjoin%
\pgfsetlinewidth{0.903375pt}%
\definecolor{currentstroke}{rgb}{0.101961,0.101961,0.101961}%
\pgfsetstrokecolor{currentstroke}%
\pgfsetdash{}{0pt}%
\pgfpathmoveto{\pgfqpoint{2.245954in}{0.257528in}}%
\pgfpathlineto{\pgfqpoint{2.309148in}{0.257528in}}%
\pgfpathlineto{\pgfqpoint{2.372343in}{0.257528in}}%
\pgfusepath{stroke}%
\end{pgfscope}%
\begin{pgfscope}%
\definecolor{textcolor}{rgb}{0.101961,0.101961,0.101961}%
\pgfsetstrokecolor{textcolor}%
\pgfsetfillcolor{textcolor}%
\pgftext[x=2.401509in,y=0.223500in,left,base]{\color{textcolor}\sffamily\fontsize{7.000000}{8.400000}\selectfont std. composition}%
\end{pgfscope}%
\begin{pgfscope}%
\pgfsetbuttcap%
\pgfsetroundjoin%
\pgfsetlinewidth{0.903375pt}%
\definecolor{currentstroke}{rgb}{0.101961,0.101961,0.101961}%
\pgfsetstrokecolor{currentstroke}%
\pgfsetdash{{3.330000pt}{1.440000pt}}{0.000000pt}%
\pgfpathmoveto{\pgfqpoint{2.245954in}{0.121957in}}%
\pgfpathlineto{\pgfqpoint{2.309148in}{0.121957in}}%
\pgfpathlineto{\pgfqpoint{2.372343in}{0.121957in}}%
\pgfusepath{stroke}%
\end{pgfscope}%
\begin{pgfscope}%
\definecolor{textcolor}{rgb}{0.101961,0.101961,0.101961}%
\pgfsetstrokecolor{textcolor}%
\pgfsetfillcolor{textcolor}%
\pgftext[x=2.401509in,y=0.087929in,left,base]{\color{textcolor}\sffamily\fontsize{7.000000}{8.400000}\selectfont 3:11:2 composition}%
\end{pgfscope}%
\end{pgfpicture}%
\makeatother%
\endgroup%
}
\caption{[PLACEHOLDER: OLD DATA, rung 3 stand-in and COVEX dev seed 4, no capacity values] Pooled correct-in-time fraction (top) and dashcam per-episode alert attainment (bottom) against fleet size. Solid: standard composition, dashed with open markers: 3:11:2 composition (separate series, not comparable in level). Means over the seeds present at each size (mostly seed 4 only). The capacity line (largest fleet meeting all three targets) is absent: no fleet sweep has produced capacity values yet.}
\label{fig:F09_fleet_load}
\end{figure}
\clearpage
\begin{figure}[t]
\centering
{\fontfamily{LinuxBiolinumT}\selectfont
%% Creator: Matplotlib, PGF backend
%%
%% To include the figure in your LaTeX document, write
%%   \input{<filename>.pgf}
%%
%% Make sure the required packages are loaded in your preamble
%%   \usepackage{pgf}
%%
%% Also ensure that all the required font packages are loaded; for instance,
%% the lmodern package is sometimes necessary when using math font.
%%   \usepackage{lmodern}
%%
%% Figures using additional raster images can only be included by \input if
%% they are in the same directory as the main LaTeX file. For loading figures
%% from other directories you can use the `import` package
%%   \usepackage{import}
%%
%% and then include the figures with
%%   \import{<path to file>}{<filename>.pgf}
%%
%% Matplotlib used the following preamble
%%
\begingroup%
\makeatletter%
\begin{pgfpicture}%
\pgfpathrectangle{\pgfpointorigin}{\pgfqpoint{3.330000in}{2.900000in}}%
\pgfusepath{use as bounding box, clip}%
\begin{pgfscope}%
\pgfsetbuttcap%
\pgfsetmiterjoin%
\definecolor{currentfill}{rgb}{1.000000,1.000000,1.000000}%
\pgfsetfillcolor{currentfill}%
\pgfsetlinewidth{0.000000pt}%
\definecolor{currentstroke}{rgb}{1.000000,1.000000,1.000000}%
\pgfsetstrokecolor{currentstroke}%
\pgfsetdash{}{0pt}%
\pgfpathmoveto{\pgfqpoint{0.000000in}{0.000000in}}%
\pgfpathlineto{\pgfqpoint{3.330000in}{0.000000in}}%
\pgfpathlineto{\pgfqpoint{3.330000in}{2.900000in}}%
\pgfpathlineto{\pgfqpoint{0.000000in}{2.900000in}}%
\pgfpathlineto{\pgfqpoint{0.000000in}{0.000000in}}%
\pgfpathclose%
\pgfusepath{fill}%
\end{pgfscope}%
\begin{pgfscope}%
\pgfsetbuttcap%
\pgfsetmiterjoin%
\definecolor{currentfill}{rgb}{1.000000,1.000000,1.000000}%
\pgfsetfillcolor{currentfill}%
\pgfsetlinewidth{0.000000pt}%
\definecolor{currentstroke}{rgb}{0.000000,0.000000,0.000000}%
\pgfsetstrokecolor{currentstroke}%
\pgfsetstrokeopacity{0.000000}%
\pgfsetdash{}{0pt}%
\pgfpathmoveto{\pgfqpoint{0.499500in}{0.957000in}}%
\pgfpathlineto{\pgfqpoint{3.230100in}{0.957000in}}%
\pgfpathlineto{\pgfqpoint{3.230100in}{2.755000in}}%
\pgfpathlineto{\pgfqpoint{0.499500in}{2.755000in}}%
\pgfpathlineto{\pgfqpoint{0.499500in}{0.957000in}}%
\pgfpathclose%
\pgfusepath{fill}%
\end{pgfscope}%
\begin{pgfscope}%
\pgfpathrectangle{\pgfqpoint{0.499500in}{0.957000in}}{\pgfqpoint{2.730600in}{1.798000in}}%
\pgfusepath{clip}%
\pgfsetbuttcap%
\pgfsetmiterjoin%
\definecolor{currentfill}{rgb}{0.949020,0.949020,0.949020}%
\pgfsetfillcolor{currentfill}%
\pgfsetlinewidth{0.000000pt}%
\definecolor{currentstroke}{rgb}{0.949020,0.949020,0.949020}%
\pgfsetstrokecolor{currentstroke}%
\pgfsetdash{}{0pt}%
\pgfpathmoveto{\pgfqpoint{2.037424in}{0.957000in}}%
\pgfpathlineto{\pgfqpoint{2.037424in}{2.755000in}}%
\pgfpathlineto{\pgfqpoint{3.230100in}{2.755000in}}%
\pgfpathlineto{\pgfqpoint{3.230100in}{0.957000in}}%
\pgfpathlineto{\pgfqpoint{2.037424in}{0.957000in}}%
\pgfpathclose%
\pgfusepath{fill}%
\end{pgfscope}%
\begin{pgfscope}%
\pgfpathrectangle{\pgfqpoint{0.499500in}{0.957000in}}{\pgfqpoint{2.730600in}{1.798000in}}%
\pgfusepath{clip}%
\pgfsetrectcap%
\pgfsetroundjoin%
\pgfsetlinewidth{0.401500pt}%
\definecolor{currentstroke}{rgb}{0.901961,0.901961,0.901961}%
\pgfsetstrokecolor{currentstroke}%
\pgfsetdash{}{0pt}%
\pgfpathmoveto{\pgfqpoint{0.750590in}{0.957000in}}%
\pgfpathlineto{\pgfqpoint{0.750590in}{2.755000in}}%
\pgfusepath{stroke}%
\end{pgfscope}%
\begin{pgfscope}%
\definecolor{textcolor}{rgb}{0.101961,0.101961,0.101961}%
\pgfsetstrokecolor{textcolor}%
\pgfsetfillcolor{textcolor}%
\pgftext[x=0.750590in,y=0.908389in,,top]{\color{textcolor}\sffamily\fontsize{7.000000}{8.400000}\selectfont 1.5}%
\end{pgfscope}%
\begin{pgfscope}%
\pgfpathrectangle{\pgfqpoint{0.499500in}{0.957000in}}{\pgfqpoint{2.730600in}{1.798000in}}%
\pgfusepath{clip}%
\pgfsetrectcap%
\pgfsetroundjoin%
\pgfsetlinewidth{0.401500pt}%
\definecolor{currentstroke}{rgb}{0.901961,0.901961,0.901961}%
\pgfsetstrokecolor{currentstroke}%
\pgfsetdash{}{0pt}%
\pgfpathmoveto{\pgfqpoint{1.378314in}{0.957000in}}%
\pgfpathlineto{\pgfqpoint{1.378314in}{2.755000in}}%
\pgfusepath{stroke}%
\end{pgfscope}%
\begin{pgfscope}%
\definecolor{textcolor}{rgb}{0.101961,0.101961,0.101961}%
\pgfsetstrokecolor{textcolor}%
\pgfsetfillcolor{textcolor}%
\pgftext[x=1.378314in,y=0.908389in,,top]{\color{textcolor}\sffamily\fontsize{7.000000}{8.400000}\selectfont 2.5}%
\end{pgfscope}%
\begin{pgfscope}%
\pgfpathrectangle{\pgfqpoint{0.499500in}{0.957000in}}{\pgfqpoint{2.730600in}{1.798000in}}%
\pgfusepath{clip}%
\pgfsetrectcap%
\pgfsetroundjoin%
\pgfsetlinewidth{0.401500pt}%
\definecolor{currentstroke}{rgb}{0.901961,0.901961,0.901961}%
\pgfsetstrokecolor{currentstroke}%
\pgfsetdash{}{0pt}%
\pgfpathmoveto{\pgfqpoint{1.692176in}{0.957000in}}%
\pgfpathlineto{\pgfqpoint{1.692176in}{2.755000in}}%
\pgfusepath{stroke}%
\end{pgfscope}%
\begin{pgfscope}%
\definecolor{textcolor}{rgb}{0.101961,0.101961,0.101961}%
\pgfsetstrokecolor{textcolor}%
\pgfsetfillcolor{textcolor}%
\pgftext[x=1.692176in,y=0.908389in,,top]{\color{textcolor}\sffamily\fontsize{7.000000}{8.400000}\selectfont 3}%
\end{pgfscope}%
\begin{pgfscope}%
\pgfpathrectangle{\pgfqpoint{0.499500in}{0.957000in}}{\pgfqpoint{2.730600in}{1.798000in}}%
\pgfusepath{clip}%
\pgfsetrectcap%
\pgfsetroundjoin%
\pgfsetlinewidth{0.401500pt}%
\definecolor{currentstroke}{rgb}{0.901961,0.901961,0.901961}%
\pgfsetstrokecolor{currentstroke}%
\pgfsetdash{}{0pt}%
\pgfpathmoveto{\pgfqpoint{2.319900in}{0.957000in}}%
\pgfpathlineto{\pgfqpoint{2.319900in}{2.755000in}}%
\pgfusepath{stroke}%
\end{pgfscope}%
\begin{pgfscope}%
\definecolor{textcolor}{rgb}{0.101961,0.101961,0.101961}%
\pgfsetstrokecolor{textcolor}%
\pgfsetfillcolor{textcolor}%
\pgftext[x=2.319900in,y=0.908389in,,top]{\color{textcolor}\sffamily\fontsize{7.000000}{8.400000}\selectfont 4}%
\end{pgfscope}%
\begin{pgfscope}%
\pgfpathrectangle{\pgfqpoint{0.499500in}{0.957000in}}{\pgfqpoint{2.730600in}{1.798000in}}%
\pgfusepath{clip}%
\pgfsetrectcap%
\pgfsetroundjoin%
\pgfsetlinewidth{0.401500pt}%
\definecolor{currentstroke}{rgb}{0.901961,0.901961,0.901961}%
\pgfsetstrokecolor{currentstroke}%
\pgfsetdash{}{0pt}%
\pgfpathmoveto{\pgfqpoint{2.947624in}{0.957000in}}%
\pgfpathlineto{\pgfqpoint{2.947624in}{2.755000in}}%
\pgfusepath{stroke}%
\end{pgfscope}%
\begin{pgfscope}%
\definecolor{textcolor}{rgb}{0.101961,0.101961,0.101961}%
\pgfsetstrokecolor{textcolor}%
\pgfsetfillcolor{textcolor}%
\pgftext[x=2.947624in,y=0.908389in,,top]{\color{textcolor}\sffamily\fontsize{7.000000}{8.400000}\selectfont 5}%
\end{pgfscope}%
\begin{pgfscope}%
\definecolor{textcolor}{rgb}{0.000000,0.000000,0.000000}%
\pgfsetstrokecolor{textcolor}%
\pgfsetfillcolor{textcolor}%
\pgftext[x=1.864800in,y=0.794192in,,top]{\color{textcolor}\sffamily\fontsize{7.500000}{9.000000}\selectfont Robot deadline (s)}%
\end{pgfscope}%
\begin{pgfscope}%
\pgfpathrectangle{\pgfqpoint{0.499500in}{0.957000in}}{\pgfqpoint{2.730600in}{1.798000in}}%
\pgfusepath{clip}%
\pgfsetrectcap%
\pgfsetroundjoin%
\pgfsetlinewidth{0.401500pt}%
\definecolor{currentstroke}{rgb}{0.901961,0.901961,0.901961}%
\pgfsetstrokecolor{currentstroke}%
\pgfsetdash{}{0pt}%
\pgfpathmoveto{\pgfqpoint{0.499500in}{0.957000in}}%
\pgfpathlineto{\pgfqpoint{3.230100in}{0.957000in}}%
\pgfusepath{stroke}%
\end{pgfscope}%
\begin{pgfscope}%
\definecolor{textcolor}{rgb}{0.101961,0.101961,0.101961}%
\pgfsetstrokecolor{textcolor}%
\pgfsetfillcolor{textcolor}%
\pgftext[x=0.318895in, y=0.923242in, left, base]{\color{textcolor}\sffamily\fontsize{7.000000}{8.400000}\selectfont 0.1}%
\end{pgfscope}%
\begin{pgfscope}%
\pgfpathrectangle{\pgfqpoint{0.499500in}{0.957000in}}{\pgfqpoint{2.730600in}{1.798000in}}%
\pgfusepath{clip}%
\pgfsetrectcap%
\pgfsetroundjoin%
\pgfsetlinewidth{0.401500pt}%
\definecolor{currentstroke}{rgb}{0.901961,0.901961,0.901961}%
\pgfsetstrokecolor{currentstroke}%
\pgfsetdash{}{0pt}%
\pgfpathmoveto{\pgfqpoint{0.499500in}{1.283909in}}%
\pgfpathlineto{\pgfqpoint{3.230100in}{1.283909in}}%
\pgfusepath{stroke}%
\end{pgfscope}%
\begin{pgfscope}%
\definecolor{textcolor}{rgb}{0.101961,0.101961,0.101961}%
\pgfsetstrokecolor{textcolor}%
\pgfsetfillcolor{textcolor}%
\pgftext[x=0.318895in, y=1.250151in, left, base]{\color{textcolor}\sffamily\fontsize{7.000000}{8.400000}\selectfont 0.2}%
\end{pgfscope}%
\begin{pgfscope}%
\pgfpathrectangle{\pgfqpoint{0.499500in}{0.957000in}}{\pgfqpoint{2.730600in}{1.798000in}}%
\pgfusepath{clip}%
\pgfsetrectcap%
\pgfsetroundjoin%
\pgfsetlinewidth{0.401500pt}%
\definecolor{currentstroke}{rgb}{0.901961,0.901961,0.901961}%
\pgfsetstrokecolor{currentstroke}%
\pgfsetdash{}{0pt}%
\pgfpathmoveto{\pgfqpoint{0.499500in}{1.610818in}}%
\pgfpathlineto{\pgfqpoint{3.230100in}{1.610818in}}%
\pgfusepath{stroke}%
\end{pgfscope}%
\begin{pgfscope}%
\definecolor{textcolor}{rgb}{0.101961,0.101961,0.101961}%
\pgfsetstrokecolor{textcolor}%
\pgfsetfillcolor{textcolor}%
\pgftext[x=0.318895in, y=1.577061in, left, base]{\color{textcolor}\sffamily\fontsize{7.000000}{8.400000}\selectfont 0.3}%
\end{pgfscope}%
\begin{pgfscope}%
\pgfpathrectangle{\pgfqpoint{0.499500in}{0.957000in}}{\pgfqpoint{2.730600in}{1.798000in}}%
\pgfusepath{clip}%
\pgfsetrectcap%
\pgfsetroundjoin%
\pgfsetlinewidth{0.401500pt}%
\definecolor{currentstroke}{rgb}{0.901961,0.901961,0.901961}%
\pgfsetstrokecolor{currentstroke}%
\pgfsetdash{}{0pt}%
\pgfpathmoveto{\pgfqpoint{0.499500in}{1.937727in}}%
\pgfpathlineto{\pgfqpoint{3.230100in}{1.937727in}}%
\pgfusepath{stroke}%
\end{pgfscope}%
\begin{pgfscope}%
\definecolor{textcolor}{rgb}{0.101961,0.101961,0.101961}%
\pgfsetstrokecolor{textcolor}%
\pgfsetfillcolor{textcolor}%
\pgftext[x=0.318895in, y=1.903970in, left, base]{\color{textcolor}\sffamily\fontsize{7.000000}{8.400000}\selectfont 0.4}%
\end{pgfscope}%
\begin{pgfscope}%
\pgfpathrectangle{\pgfqpoint{0.499500in}{0.957000in}}{\pgfqpoint{2.730600in}{1.798000in}}%
\pgfusepath{clip}%
\pgfsetrectcap%
\pgfsetroundjoin%
\pgfsetlinewidth{0.401500pt}%
\definecolor{currentstroke}{rgb}{0.901961,0.901961,0.901961}%
\pgfsetstrokecolor{currentstroke}%
\pgfsetdash{}{0pt}%
\pgfpathmoveto{\pgfqpoint{0.499500in}{2.264636in}}%
\pgfpathlineto{\pgfqpoint{3.230100in}{2.264636in}}%
\pgfusepath{stroke}%
\end{pgfscope}%
\begin{pgfscope}%
\definecolor{textcolor}{rgb}{0.101961,0.101961,0.101961}%
\pgfsetstrokecolor{textcolor}%
\pgfsetfillcolor{textcolor}%
\pgftext[x=0.318895in, y=2.230879in, left, base]{\color{textcolor}\sffamily\fontsize{7.000000}{8.400000}\selectfont 0.5}%
\end{pgfscope}%
\begin{pgfscope}%
\pgfpathrectangle{\pgfqpoint{0.499500in}{0.957000in}}{\pgfqpoint{2.730600in}{1.798000in}}%
\pgfusepath{clip}%
\pgfsetrectcap%
\pgfsetroundjoin%
\pgfsetlinewidth{0.401500pt}%
\definecolor{currentstroke}{rgb}{0.901961,0.901961,0.901961}%
\pgfsetstrokecolor{currentstroke}%
\pgfsetdash{}{0pt}%
\pgfpathmoveto{\pgfqpoint{0.499500in}{2.591545in}}%
\pgfpathlineto{\pgfqpoint{3.230100in}{2.591545in}}%
\pgfusepath{stroke}%
\end{pgfscope}%
\begin{pgfscope}%
\definecolor{textcolor}{rgb}{0.101961,0.101961,0.101961}%
\pgfsetstrokecolor{textcolor}%
\pgfsetfillcolor{textcolor}%
\pgftext[x=0.318895in, y=2.557788in, left, base]{\color{textcolor}\sffamily\fontsize{7.000000}{8.400000}\selectfont 0.6}%
\end{pgfscope}%
\begin{pgfscope}%
\definecolor{textcolor}{rgb}{0.000000,0.000000,0.000000}%
\pgfsetstrokecolor{textcolor}%
\pgfsetfillcolor{textcolor}%
\pgftext[x=0.263339in,y=1.856000in,,bottom,rotate=90.000000]{\color{textcolor}\sffamily\fontsize{7.500000}{9.000000}\selectfont Robot correct by deadline}%
\end{pgfscope}%
\begin{pgfscope}%
\pgfsetrectcap%
\pgfsetmiterjoin%
\pgfsetlinewidth{0.803000pt}%
\definecolor{currentstroke}{rgb}{0.101961,0.101961,0.101961}%
\pgfsetstrokecolor{currentstroke}%
\pgfsetdash{}{0pt}%
\pgfpathmoveto{\pgfqpoint{0.499500in}{0.957000in}}%
\pgfpathlineto{\pgfqpoint{0.499500in}{2.755000in}}%
\pgfusepath{stroke}%
\end{pgfscope}%
\begin{pgfscope}%
\pgfsetrectcap%
\pgfsetmiterjoin%
\pgfsetlinewidth{0.803000pt}%
\definecolor{currentstroke}{rgb}{0.101961,0.101961,0.101961}%
\pgfsetstrokecolor{currentstroke}%
\pgfsetdash{}{0pt}%
\pgfpathmoveto{\pgfqpoint{0.499500in}{0.957000in}}%
\pgfpathlineto{\pgfqpoint{3.230100in}{0.957000in}}%
\pgfusepath{stroke}%
\end{pgfscope}%
\begin{pgfscope}%
\definecolor{textcolor}{rgb}{0.549020,0.549020,0.549020}%
\pgfsetstrokecolor{textcolor}%
\pgfsetfillcolor{textcolor}%
\pgftext[x=2.473698in, y=2.589412in, left, base]{\color{textcolor}\sffamily\fontsize{7.000000}{8.400000}\selectfont no cells}%
\end{pgfscope}%
\begin{pgfscope}%
\definecolor{textcolor}{rgb}{0.549020,0.549020,0.549020}%
\pgfsetstrokecolor{textcolor}%
\pgfsetfillcolor{textcolor}%
\pgftext[x=2.334817in, y=2.489489in, left, base]{\color{textcolor}\sffamily\fontsize{7.000000}{8.400000}\selectfont at 4 s and 5 s}%
\end{pgfscope}%
\begin{pgfscope}%
\pgfpathrectangle{\pgfqpoint{0.499500in}{0.957000in}}{\pgfqpoint{2.730600in}{1.798000in}}%
\pgfusepath{clip}%
\pgfsetrectcap%
\pgfsetroundjoin%
\pgfsetlinewidth{0.903375pt}%
\definecolor{currentstroke}{rgb}{0.156863,0.356863,0.568627}%
\pgfsetstrokecolor{currentstroke}%
\pgfsetdash{}{0pt}%
\pgfpathmoveto{\pgfqpoint{0.750590in}{2.462763in}}%
\pgfpathlineto{\pgfqpoint{1.692176in}{2.685656in}}%
\pgfusepath{stroke}%
\end{pgfscope}%
\begin{pgfscope}%
\pgfpathrectangle{\pgfqpoint{0.499500in}{0.957000in}}{\pgfqpoint{2.730600in}{1.798000in}}%
\pgfusepath{clip}%
\pgfsetbuttcap%
\pgfsetroundjoin%
\definecolor{currentfill}{rgb}{0.156863,0.356863,0.568627}%
\pgfsetfillcolor{currentfill}%
\pgfsetlinewidth{1.003750pt}%
\definecolor{currentstroke}{rgb}{0.156863,0.356863,0.568627}%
\pgfsetstrokecolor{currentstroke}%
\pgfsetdash{}{0pt}%
\pgfsys@defobject{currentmarker}{\pgfqpoint{-0.023611in}{-0.023611in}}{\pgfqpoint{0.023611in}{0.023611in}}{%
\pgfpathmoveto{\pgfqpoint{0.000000in}{-0.023611in}}%
\pgfpathcurveto{\pgfqpoint{0.006262in}{-0.023611in}}{\pgfqpoint{0.012268in}{-0.021123in}}{\pgfqpoint{0.016696in}{-0.016696in}}%
\pgfpathcurveto{\pgfqpoint{0.021123in}{-0.012268in}}{\pgfqpoint{0.023611in}{-0.006262in}}{\pgfqpoint{0.023611in}{0.000000in}}%
\pgfpathcurveto{\pgfqpoint{0.023611in}{0.006262in}}{\pgfqpoint{0.021123in}{0.012268in}}{\pgfqpoint{0.016696in}{0.016696in}}%
\pgfpathcurveto{\pgfqpoint{0.012268in}{0.021123in}}{\pgfqpoint{0.006262in}{0.023611in}}{\pgfqpoint{0.000000in}{0.023611in}}%
\pgfpathcurveto{\pgfqpoint{-0.006262in}{0.023611in}}{\pgfqpoint{-0.012268in}{0.021123in}}{\pgfqpoint{-0.016696in}{0.016696in}}%
\pgfpathcurveto{\pgfqpoint{-0.021123in}{0.012268in}}{\pgfqpoint{-0.023611in}{0.006262in}}{\pgfqpoint{-0.023611in}{0.000000in}}%
\pgfpathcurveto{\pgfqpoint{-0.023611in}{-0.006262in}}{\pgfqpoint{-0.021123in}{-0.012268in}}{\pgfqpoint{-0.016696in}{-0.016696in}}%
\pgfpathcurveto{\pgfqpoint{-0.012268in}{-0.021123in}}{\pgfqpoint{-0.006262in}{-0.023611in}}{\pgfqpoint{0.000000in}{-0.023611in}}%
\pgfpathlineto{\pgfqpoint{0.000000in}{-0.023611in}}%
\pgfpathclose%
\pgfusepath{stroke,fill}%
}%
\begin{pgfscope}%
\pgfsys@transformshift{0.750590in}{2.462763in}%
\pgfsys@useobject{currentmarker}{}%
\end{pgfscope}%
\begin{pgfscope}%
\pgfsys@transformshift{1.692176in}{2.685656in}%
\pgfsys@useobject{currentmarker}{}%
\end{pgfscope}%
\end{pgfscope}%
\begin{pgfscope}%
\pgfpathrectangle{\pgfqpoint{0.499500in}{0.957000in}}{\pgfqpoint{2.730600in}{1.798000in}}%
\pgfusepath{clip}%
\pgfsetrectcap%
\pgfsetroundjoin%
\pgfsetlinewidth{0.903375pt}%
\definecolor{currentstroke}{rgb}{0.784314,0.784314,0.784314}%
\pgfsetstrokecolor{currentstroke}%
\pgfsetdash{}{0pt}%
\pgfpathmoveto{\pgfqpoint{0.750590in}{1.199704in}}%
\pgfpathlineto{\pgfqpoint{1.692176in}{1.199705in}}%
\pgfusepath{stroke}%
\end{pgfscope}%
\begin{pgfscope}%
\pgfpathrectangle{\pgfqpoint{0.499500in}{0.957000in}}{\pgfqpoint{2.730600in}{1.798000in}}%
\pgfusepath{clip}%
\pgfsetbuttcap%
\pgfsetmiterjoin%
\definecolor{currentfill}{rgb}{0.784314,0.784314,0.784314}%
\pgfsetfillcolor{currentfill}%
\pgfsetlinewidth{1.003750pt}%
\definecolor{currentstroke}{rgb}{0.784314,0.784314,0.784314}%
\pgfsetstrokecolor{currentstroke}%
\pgfsetdash{}{0pt}%
\pgfsys@defobject{currentmarker}{\pgfqpoint{-0.023611in}{-0.023611in}}{\pgfqpoint{0.023611in}{0.023611in}}{%
\pgfpathmoveto{\pgfqpoint{0.000000in}{0.023611in}}%
\pgfpathlineto{\pgfqpoint{-0.023611in}{-0.023611in}}%
\pgfpathlineto{\pgfqpoint{0.023611in}{-0.023611in}}%
\pgfpathlineto{\pgfqpoint{0.000000in}{0.023611in}}%
\pgfpathclose%
\pgfusepath{stroke,fill}%
}%
\begin{pgfscope}%
\pgfsys@transformshift{0.750590in}{1.199704in}%
\pgfsys@useobject{currentmarker}{}%
\end{pgfscope}%
\begin{pgfscope}%
\pgfsys@transformshift{1.692176in}{1.199705in}%
\pgfsys@useobject{currentmarker}{}%
\end{pgfscope}%
\end{pgfscope}%
\begin{pgfscope}%
\pgfpathrectangle{\pgfqpoint{0.499500in}{0.957000in}}{\pgfqpoint{2.730600in}{1.798000in}}%
\pgfusepath{clip}%
\pgfsetrectcap%
\pgfsetroundjoin%
\pgfsetlinewidth{0.903375pt}%
\definecolor{currentstroke}{rgb}{0.784314,0.784314,0.784314}%
\pgfsetstrokecolor{currentstroke}%
\pgfsetdash{}{0pt}%
\pgfpathmoveto{\pgfqpoint{1.378314in}{1.818850in}}%
\pgfpathlineto{\pgfqpoint{1.692176in}{1.917915in}}%
\pgfusepath{stroke}%
\end{pgfscope}%
\begin{pgfscope}%
\pgfpathrectangle{\pgfqpoint{0.499500in}{0.957000in}}{\pgfqpoint{2.730600in}{1.798000in}}%
\pgfusepath{clip}%
\pgfsetbuttcap%
\pgfsetmiterjoin%
\definecolor{currentfill}{rgb}{0.784314,0.784314,0.784314}%
\pgfsetfillcolor{currentfill}%
\pgfsetlinewidth{1.003750pt}%
\definecolor{currentstroke}{rgb}{0.784314,0.784314,0.784314}%
\pgfsetstrokecolor{currentstroke}%
\pgfsetdash{}{0pt}%
\pgfsys@defobject{currentmarker}{\pgfqpoint{-0.023611in}{-0.023611in}}{\pgfqpoint{0.023611in}{0.023611in}}{%
\pgfpathmoveto{\pgfqpoint{-0.023611in}{-0.023611in}}%
\pgfpathlineto{\pgfqpoint{0.023611in}{-0.023611in}}%
\pgfpathlineto{\pgfqpoint{0.023611in}{0.023611in}}%
\pgfpathlineto{\pgfqpoint{-0.023611in}{0.023611in}}%
\pgfpathlineto{\pgfqpoint{-0.023611in}{-0.023611in}}%
\pgfpathclose%
\pgfusepath{stroke,fill}%
}%
\begin{pgfscope}%
\pgfsys@transformshift{1.378314in}{1.818850in}%
\pgfsys@useobject{currentmarker}{}%
\end{pgfscope}%
\begin{pgfscope}%
\pgfsys@transformshift{1.692176in}{1.917915in}%
\pgfsys@useobject{currentmarker}{}%
\end{pgfscope}%
\end{pgfscope}%
\begin{pgfscope}%
\pgfpathrectangle{\pgfqpoint{0.499500in}{0.957000in}}{\pgfqpoint{2.730600in}{1.798000in}}%
\pgfusepath{clip}%
\pgfsetrectcap%
\pgfsetroundjoin%
\pgfsetlinewidth{0.903375pt}%
\definecolor{currentstroke}{rgb}{0.478431,0.478431,0.478431}%
\pgfsetstrokecolor{currentstroke}%
\pgfsetdash{}{0pt}%
\pgfpathmoveto{\pgfqpoint{1.378314in}{1.868383in}}%
\pgfpathlineto{\pgfqpoint{1.692176in}{1.843617in}}%
\pgfusepath{stroke}%
\end{pgfscope}%
\begin{pgfscope}%
\pgfpathrectangle{\pgfqpoint{0.499500in}{0.957000in}}{\pgfqpoint{2.730600in}{1.798000in}}%
\pgfusepath{clip}%
\pgfsetbuttcap%
\pgfsetmiterjoin%
\definecolor{currentfill}{rgb}{0.478431,0.478431,0.478431}%
\pgfsetfillcolor{currentfill}%
\pgfsetlinewidth{1.003750pt}%
\definecolor{currentstroke}{rgb}{0.478431,0.478431,0.478431}%
\pgfsetstrokecolor{currentstroke}%
\pgfsetdash{}{0pt}%
\pgfsys@defobject{currentmarker}{\pgfqpoint{-0.023611in}{-0.023611in}}{\pgfqpoint{0.023611in}{0.023611in}}{%
\pgfpathmoveto{\pgfqpoint{-0.007870in}{-0.023611in}}%
\pgfpathlineto{\pgfqpoint{0.007870in}{-0.023611in}}%
\pgfpathlineto{\pgfqpoint{0.007870in}{-0.007870in}}%
\pgfpathlineto{\pgfqpoint{0.023611in}{-0.007870in}}%
\pgfpathlineto{\pgfqpoint{0.023611in}{0.007870in}}%
\pgfpathlineto{\pgfqpoint{0.007870in}{0.007870in}}%
\pgfpathlineto{\pgfqpoint{0.007870in}{0.023611in}}%
\pgfpathlineto{\pgfqpoint{-0.007870in}{0.023611in}}%
\pgfpathlineto{\pgfqpoint{-0.007870in}{0.007870in}}%
\pgfpathlineto{\pgfqpoint{-0.023611in}{0.007870in}}%
\pgfpathlineto{\pgfqpoint{-0.023611in}{-0.007870in}}%
\pgfpathlineto{\pgfqpoint{-0.007870in}{-0.007870in}}%
\pgfpathlineto{\pgfqpoint{-0.007870in}{-0.023611in}}%
\pgfpathclose%
\pgfusepath{stroke,fill}%
}%
\begin{pgfscope}%
\pgfsys@transformshift{1.378314in}{1.868383in}%
\pgfsys@useobject{currentmarker}{}%
\end{pgfscope}%
\begin{pgfscope}%
\pgfsys@transformshift{1.692176in}{1.843617in}%
\pgfsys@useobject{currentmarker}{}%
\end{pgfscope}%
\end{pgfscope}%
\begin{pgfscope}%
\pgfpathrectangle{\pgfqpoint{0.499500in}{0.957000in}}{\pgfqpoint{2.730600in}{1.798000in}}%
\pgfusepath{clip}%
\pgfsetrectcap%
\pgfsetroundjoin%
\pgfsetlinewidth{0.903375pt}%
\definecolor{currentstroke}{rgb}{0.478431,0.305882,0.611765}%
\pgfsetstrokecolor{currentstroke}%
\pgfsetdash{}{0pt}%
\pgfpathmoveto{\pgfqpoint{1.378314in}{1.725931in}}%
\pgfpathlineto{\pgfqpoint{1.692176in}{1.725931in}}%
\pgfusepath{stroke}%
\end{pgfscope}%
\begin{pgfscope}%
\pgfpathrectangle{\pgfqpoint{0.499500in}{0.957000in}}{\pgfqpoint{2.730600in}{1.798000in}}%
\pgfusepath{clip}%
\pgfsetbuttcap%
\pgfsetmiterjoin%
\definecolor{currentfill}{rgb}{0.478431,0.305882,0.611765}%
\pgfsetfillcolor{currentfill}%
\pgfsetlinewidth{1.003750pt}%
\definecolor{currentstroke}{rgb}{0.478431,0.305882,0.611765}%
\pgfsetstrokecolor{currentstroke}%
\pgfsetdash{}{0pt}%
\pgfsys@defobject{currentmarker}{\pgfqpoint{-0.023611in}{-0.023611in}}{\pgfqpoint{0.023611in}{0.023611in}}{%
\pgfpathmoveto{\pgfqpoint{-0.011806in}{-0.023611in}}%
\pgfpathlineto{\pgfqpoint{0.000000in}{-0.011806in}}%
\pgfpathlineto{\pgfqpoint{0.011806in}{-0.023611in}}%
\pgfpathlineto{\pgfqpoint{0.023611in}{-0.011806in}}%
\pgfpathlineto{\pgfqpoint{0.011806in}{0.000000in}}%
\pgfpathlineto{\pgfqpoint{0.023611in}{0.011806in}}%
\pgfpathlineto{\pgfqpoint{0.011806in}{0.023611in}}%
\pgfpathlineto{\pgfqpoint{0.000000in}{0.011806in}}%
\pgfpathlineto{\pgfqpoint{-0.011806in}{0.023611in}}%
\pgfpathlineto{\pgfqpoint{-0.023611in}{0.011806in}}%
\pgfpathlineto{\pgfqpoint{-0.011806in}{0.000000in}}%
\pgfpathlineto{\pgfqpoint{-0.023611in}{-0.011806in}}%
\pgfpathlineto{\pgfqpoint{-0.011806in}{-0.023611in}}%
\pgfpathclose%
\pgfusepath{stroke,fill}%
}%
\begin{pgfscope}%
\pgfsys@transformshift{1.378314in}{1.725931in}%
\pgfsys@useobject{currentmarker}{}%
\end{pgfscope}%
\begin{pgfscope}%
\pgfsys@transformshift{1.692176in}{1.725931in}%
\pgfsys@useobject{currentmarker}{}%
\end{pgfscope}%
\end{pgfscope}%
\begin{pgfscope}%
\pgfpathrectangle{\pgfqpoint{0.499500in}{0.957000in}}{\pgfqpoint{2.730600in}{1.798000in}}%
\pgfusepath{clip}%
\pgfsetrectcap%
\pgfsetroundjoin%
\pgfsetlinewidth{0.903375pt}%
\definecolor{currentstroke}{rgb}{0.435294,0.560784,0.709804}%
\pgfsetstrokecolor{currentstroke}%
\pgfsetdash{}{0pt}%
\pgfpathmoveto{\pgfqpoint{1.378314in}{1.850248in}}%
\pgfpathlineto{\pgfqpoint{1.692176in}{1.963156in}}%
\pgfusepath{stroke}%
\end{pgfscope}%
\begin{pgfscope}%
\pgfpathrectangle{\pgfqpoint{0.499500in}{0.957000in}}{\pgfqpoint{2.730600in}{1.798000in}}%
\pgfusepath{clip}%
\pgfsetbuttcap%
\pgfsetmiterjoin%
\definecolor{currentfill}{rgb}{0.435294,0.560784,0.709804}%
\pgfsetfillcolor{currentfill}%
\pgfsetlinewidth{1.003750pt}%
\definecolor{currentstroke}{rgb}{0.435294,0.560784,0.709804}%
\pgfsetstrokecolor{currentstroke}%
\pgfsetdash{}{0pt}%
\pgfsys@defobject{currentmarker}{\pgfqpoint{-0.020448in}{-0.023611in}}{\pgfqpoint{0.020448in}{0.023611in}}{%
\pgfpathmoveto{\pgfqpoint{0.000000in}{0.023611in}}%
\pgfpathlineto{\pgfqpoint{-0.020448in}{0.011806in}}%
\pgfpathlineto{\pgfqpoint{-0.020448in}{-0.011806in}}%
\pgfpathlineto{\pgfqpoint{-0.000000in}{-0.023611in}}%
\pgfpathlineto{\pgfqpoint{0.020448in}{-0.011806in}}%
\pgfpathlineto{\pgfqpoint{0.020448in}{0.011806in}}%
\pgfpathlineto{\pgfqpoint{0.000000in}{0.023611in}}%
\pgfpathclose%
\pgfusepath{stroke,fill}%
}%
\begin{pgfscope}%
\pgfsys@transformshift{1.378314in}{1.850248in}%
\pgfsys@useobject{currentmarker}{}%
\end{pgfscope}%
\begin{pgfscope}%
\pgfsys@transformshift{1.692176in}{1.963156in}%
\pgfsys@useobject{currentmarker}{}%
\end{pgfscope}%
\end{pgfscope}%
\begin{pgfscope}%
\pgfpathrectangle{\pgfqpoint{0.499500in}{0.957000in}}{\pgfqpoint{2.730600in}{1.798000in}}%
\pgfusepath{clip}%
\pgfsetrectcap%
\pgfsetroundjoin%
\pgfsetlinewidth{0.903375pt}%
\definecolor{currentstroke}{rgb}{0.435294,0.560784,0.709804}%
\pgfsetstrokecolor{currentstroke}%
\pgfsetdash{}{0pt}%
\pgfpathmoveto{\pgfqpoint{1.378314in}{1.838841in}}%
\pgfpathlineto{\pgfqpoint{1.692176in}{1.951261in}}%
\pgfusepath{stroke}%
\end{pgfscope}%
\begin{pgfscope}%
\pgfpathrectangle{\pgfqpoint{0.499500in}{0.957000in}}{\pgfqpoint{2.730600in}{1.798000in}}%
\pgfusepath{clip}%
\pgfsetbuttcap%
\pgfsetmiterjoin%
\definecolor{currentfill}{rgb}{0.435294,0.560784,0.709804}%
\pgfsetfillcolor{currentfill}%
\pgfsetlinewidth{1.003750pt}%
\definecolor{currentstroke}{rgb}{0.435294,0.560784,0.709804}%
\pgfsetstrokecolor{currentstroke}%
\pgfsetdash{}{0pt}%
\pgfsys@defobject{currentmarker}{\pgfqpoint{-0.020035in}{-0.033391in}}{\pgfqpoint{0.020035in}{0.033391in}}{%
\pgfpathmoveto{\pgfqpoint{-0.000000in}{-0.033391in}}%
\pgfpathlineto{\pgfqpoint{0.020035in}{0.000000in}}%
\pgfpathlineto{\pgfqpoint{-0.000000in}{0.033391in}}%
\pgfpathlineto{\pgfqpoint{-0.020035in}{0.000000in}}%
\pgfpathlineto{\pgfqpoint{-0.000000in}{-0.033391in}}%
\pgfpathclose%
\pgfusepath{stroke,fill}%
}%
\begin{pgfscope}%
\pgfsys@transformshift{1.378314in}{1.838841in}%
\pgfsys@useobject{currentmarker}{}%
\end{pgfscope}%
\begin{pgfscope}%
\pgfsys@transformshift{1.692176in}{1.951261in}%
\pgfsys@useobject{currentmarker}{}%
\end{pgfscope}%
\end{pgfscope}%
\begin{pgfscope}%
\definecolor{textcolor}{rgb}{0.549020,0.549020,0.549020}%
\pgfsetstrokecolor{textcolor}%
\pgfsetfillcolor{textcolor}%
\pgftext[x=3.296700in,y=2.856500in,right,top]{\color{textcolor}\sffamily\fontsize{5.500000}{6.600000}\selectfont OLD DATA}%
\end{pgfscope}%
\begin{pgfscope}%
\pgfsetrectcap%
\pgfsetroundjoin%
\pgfsetlinewidth{0.903375pt}%
\definecolor{currentstroke}{rgb}{0.156863,0.356863,0.568627}%
\pgfsetstrokecolor{currentstroke}%
\pgfsetdash{}{0pt}%
\pgfpathmoveto{\pgfqpoint{0.475524in}{0.408500in}}%
\pgfpathlineto{\pgfqpoint{0.538718in}{0.408500in}}%
\pgfpathlineto{\pgfqpoint{0.601913in}{0.408500in}}%
\pgfusepath{stroke}%
\end{pgfscope}%
\begin{pgfscope}%
\pgfsetbuttcap%
\pgfsetroundjoin%
\definecolor{currentfill}{rgb}{0.156863,0.356863,0.568627}%
\pgfsetfillcolor{currentfill}%
\pgfsetlinewidth{1.003750pt}%
\definecolor{currentstroke}{rgb}{0.156863,0.356863,0.568627}%
\pgfsetstrokecolor{currentstroke}%
\pgfsetdash{}{0pt}%
\pgfsys@defobject{currentmarker}{\pgfqpoint{-0.023611in}{-0.023611in}}{\pgfqpoint{0.023611in}{0.023611in}}{%
\pgfpathmoveto{\pgfqpoint{0.000000in}{-0.023611in}}%
\pgfpathcurveto{\pgfqpoint{0.006262in}{-0.023611in}}{\pgfqpoint{0.012268in}{-0.021123in}}{\pgfqpoint{0.016696in}{-0.016696in}}%
\pgfpathcurveto{\pgfqpoint{0.021123in}{-0.012268in}}{\pgfqpoint{0.023611in}{-0.006262in}}{\pgfqpoint{0.023611in}{0.000000in}}%
\pgfpathcurveto{\pgfqpoint{0.023611in}{0.006262in}}{\pgfqpoint{0.021123in}{0.012268in}}{\pgfqpoint{0.016696in}{0.016696in}}%
\pgfpathcurveto{\pgfqpoint{0.012268in}{0.021123in}}{\pgfqpoint{0.006262in}{0.023611in}}{\pgfqpoint{0.000000in}{0.023611in}}%
\pgfpathcurveto{\pgfqpoint{-0.006262in}{0.023611in}}{\pgfqpoint{-0.012268in}{0.021123in}}{\pgfqpoint{-0.016696in}{0.016696in}}%
\pgfpathcurveto{\pgfqpoint{-0.021123in}{0.012268in}}{\pgfqpoint{-0.023611in}{0.006262in}}{\pgfqpoint{-0.023611in}{0.000000in}}%
\pgfpathcurveto{\pgfqpoint{-0.023611in}{-0.006262in}}{\pgfqpoint{-0.021123in}{-0.012268in}}{\pgfqpoint{-0.016696in}{-0.016696in}}%
\pgfpathcurveto{\pgfqpoint{-0.012268in}{-0.021123in}}{\pgfqpoint{-0.006262in}{-0.023611in}}{\pgfqpoint{0.000000in}{-0.023611in}}%
\pgfpathlineto{\pgfqpoint{0.000000in}{-0.023611in}}%
\pgfpathclose%
\pgfusepath{stroke,fill}%
}%
\begin{pgfscope}%
\pgfsys@transformshift{0.538718in}{0.408500in}%
\pgfsys@useobject{currentmarker}{}%
\end{pgfscope}%
\end{pgfscope}%
\begin{pgfscope}%
\definecolor{textcolor}{rgb}{0.101961,0.101961,0.101961}%
\pgfsetstrokecolor{textcolor}%
\pgfsetfillcolor{textcolor}%
\pgftext[x=0.631080in,y=0.374472in,left,base]{\color{textcolor}\sffamily\fontsize{7.000000}{8.400000}\selectfont COVEX (old: rung 3)}%
\end{pgfscope}%
\begin{pgfscope}%
\pgfsetrectcap%
\pgfsetroundjoin%
\pgfsetlinewidth{0.903375pt}%
\definecolor{currentstroke}{rgb}{0.784314,0.784314,0.784314}%
\pgfsetstrokecolor{currentstroke}%
\pgfsetdash{}{0pt}%
\pgfpathmoveto{\pgfqpoint{0.475524in}{0.262667in}}%
\pgfpathlineto{\pgfqpoint{0.538718in}{0.262667in}}%
\pgfpathlineto{\pgfqpoint{0.601913in}{0.262667in}}%
\pgfusepath{stroke}%
\end{pgfscope}%
\begin{pgfscope}%
\pgfsetbuttcap%
\pgfsetmiterjoin%
\definecolor{currentfill}{rgb}{0.784314,0.784314,0.784314}%
\pgfsetfillcolor{currentfill}%
\pgfsetlinewidth{1.003750pt}%
\definecolor{currentstroke}{rgb}{0.784314,0.784314,0.784314}%
\pgfsetstrokecolor{currentstroke}%
\pgfsetdash{}{0pt}%
\pgfsys@defobject{currentmarker}{\pgfqpoint{-0.023611in}{-0.023611in}}{\pgfqpoint{0.023611in}{0.023611in}}{%
\pgfpathmoveto{\pgfqpoint{0.000000in}{0.023611in}}%
\pgfpathlineto{\pgfqpoint{-0.023611in}{-0.023611in}}%
\pgfpathlineto{\pgfqpoint{0.023611in}{-0.023611in}}%
\pgfpathlineto{\pgfqpoint{0.000000in}{0.023611in}}%
\pgfpathclose%
\pgfusepath{stroke,fill}%
}%
\begin{pgfscope}%
\pgfsys@transformshift{0.538718in}{0.262667in}%
\pgfsys@useobject{currentmarker}{}%
\end{pgfscope}%
\end{pgfscope}%
\begin{pgfscope}%
\definecolor{textcolor}{rgb}{0.101961,0.101961,0.101961}%
\pgfsetstrokecolor{textcolor}%
\pgfsetfillcolor{textcolor}%
\pgftext[x=0.631080in,y=0.228639in,left,base]{\color{textcolor}\sffamily\fontsize{7.000000}{8.400000}\selectfont SimpleStream~\cite{simplestream2026}}%
\end{pgfscope}%
\begin{pgfscope}%
\pgfsetrectcap%
\pgfsetroundjoin%
\pgfsetlinewidth{0.903375pt}%
\definecolor{currentstroke}{rgb}{0.784314,0.784314,0.784314}%
\pgfsetstrokecolor{currentstroke}%
\pgfsetdash{}{0pt}%
\pgfpathmoveto{\pgfqpoint{0.475524in}{0.116833in}}%
\pgfpathlineto{\pgfqpoint{0.538718in}{0.116833in}}%
\pgfpathlineto{\pgfqpoint{0.601913in}{0.116833in}}%
\pgfusepath{stroke}%
\end{pgfscope}%
\begin{pgfscope}%
\pgfsetbuttcap%
\pgfsetmiterjoin%
\definecolor{currentfill}{rgb}{0.784314,0.784314,0.784314}%
\pgfsetfillcolor{currentfill}%
\pgfsetlinewidth{1.003750pt}%
\definecolor{currentstroke}{rgb}{0.784314,0.784314,0.784314}%
\pgfsetstrokecolor{currentstroke}%
\pgfsetdash{}{0pt}%
\pgfsys@defobject{currentmarker}{\pgfqpoint{-0.023611in}{-0.023611in}}{\pgfqpoint{0.023611in}{0.023611in}}{%
\pgfpathmoveto{\pgfqpoint{-0.023611in}{-0.023611in}}%
\pgfpathlineto{\pgfqpoint{0.023611in}{-0.023611in}}%
\pgfpathlineto{\pgfqpoint{0.023611in}{0.023611in}}%
\pgfpathlineto{\pgfqpoint{-0.023611in}{0.023611in}}%
\pgfpathlineto{\pgfqpoint{-0.023611in}{-0.023611in}}%
\pgfpathclose%
\pgfusepath{stroke,fill}%
}%
\begin{pgfscope}%
\pgfsys@transformshift{0.538718in}{0.116833in}%
\pgfsys@useobject{currentmarker}{}%
\end{pgfscope}%
\end{pgfscope}%
\begin{pgfscope}%
\definecolor{textcolor}{rgb}{0.101961,0.101961,0.101961}%
\pgfsetstrokecolor{textcolor}%
\pgfsetfillcolor{textcolor}%
\pgftext[x=0.631080in,y=0.082806in,left,base]{\color{textcolor}\sffamily\fontsize{7.000000}{8.400000}\selectfont LLF~\cite{mok1983llf}}%
\end{pgfscope}%
\begin{pgfscope}%
\pgfsetrectcap%
\pgfsetroundjoin%
\pgfsetlinewidth{0.903375pt}%
\definecolor{currentstroke}{rgb}{0.478431,0.478431,0.478431}%
\pgfsetstrokecolor{currentstroke}%
\pgfsetdash{}{0pt}%
\pgfpathmoveto{\pgfqpoint{1.646658in}{0.408500in}}%
\pgfpathlineto{\pgfqpoint{1.709853in}{0.408500in}}%
\pgfpathlineto{\pgfqpoint{1.773047in}{0.408500in}}%
\pgfusepath{stroke}%
\end{pgfscope}%
\begin{pgfscope}%
\pgfsetbuttcap%
\pgfsetmiterjoin%
\definecolor{currentfill}{rgb}{0.478431,0.478431,0.478431}%
\pgfsetfillcolor{currentfill}%
\pgfsetlinewidth{1.003750pt}%
\definecolor{currentstroke}{rgb}{0.478431,0.478431,0.478431}%
\pgfsetstrokecolor{currentstroke}%
\pgfsetdash{}{0pt}%
\pgfsys@defobject{currentmarker}{\pgfqpoint{-0.023611in}{-0.023611in}}{\pgfqpoint{0.023611in}{0.023611in}}{%
\pgfpathmoveto{\pgfqpoint{-0.007870in}{-0.023611in}}%
\pgfpathlineto{\pgfqpoint{0.007870in}{-0.023611in}}%
\pgfpathlineto{\pgfqpoint{0.007870in}{-0.007870in}}%
\pgfpathlineto{\pgfqpoint{0.023611in}{-0.007870in}}%
\pgfpathlineto{\pgfqpoint{0.023611in}{0.007870in}}%
\pgfpathlineto{\pgfqpoint{0.007870in}{0.007870in}}%
\pgfpathlineto{\pgfqpoint{0.007870in}{0.023611in}}%
\pgfpathlineto{\pgfqpoint{-0.007870in}{0.023611in}}%
\pgfpathlineto{\pgfqpoint{-0.007870in}{0.007870in}}%
\pgfpathlineto{\pgfqpoint{-0.023611in}{0.007870in}}%
\pgfpathlineto{\pgfqpoint{-0.023611in}{-0.007870in}}%
\pgfpathlineto{\pgfqpoint{-0.007870in}{-0.007870in}}%
\pgfpathlineto{\pgfqpoint{-0.007870in}{-0.023611in}}%
\pgfpathclose%
\pgfusepath{stroke,fill}%
}%
\begin{pgfscope}%
\pgfsys@transformshift{1.709853in}{0.408500in}%
\pgfsys@useobject{currentmarker}{}%
\end{pgfscope}%
\end{pgfscope}%
\begin{pgfscope}%
\definecolor{textcolor}{rgb}{0.101961,0.101961,0.101961}%
\pgfsetstrokecolor{textcolor}%
\pgfsetfillcolor{textcolor}%
\pgftext[x=1.802214in,y=0.374472in,left,base]{\color{textcolor}\sffamily\fontsize{7.000000}{8.400000}\selectfont Mooncake~\cite{mooncake2025}}%
\end{pgfscope}%
\begin{pgfscope}%
\pgfsetrectcap%
\pgfsetroundjoin%
\pgfsetlinewidth{0.903375pt}%
\definecolor{currentstroke}{rgb}{0.478431,0.305882,0.611765}%
\pgfsetstrokecolor{currentstroke}%
\pgfsetdash{}{0pt}%
\pgfpathmoveto{\pgfqpoint{1.646658in}{0.267528in}}%
\pgfpathlineto{\pgfqpoint{1.709853in}{0.267528in}}%
\pgfpathlineto{\pgfqpoint{1.773047in}{0.267528in}}%
\pgfusepath{stroke}%
\end{pgfscope}%
\begin{pgfscope}%
\pgfsetbuttcap%
\pgfsetmiterjoin%
\definecolor{currentfill}{rgb}{0.478431,0.305882,0.611765}%
\pgfsetfillcolor{currentfill}%
\pgfsetlinewidth{1.003750pt}%
\definecolor{currentstroke}{rgb}{0.478431,0.305882,0.611765}%
\pgfsetstrokecolor{currentstroke}%
\pgfsetdash{}{0pt}%
\pgfsys@defobject{currentmarker}{\pgfqpoint{-0.023611in}{-0.023611in}}{\pgfqpoint{0.023611in}{0.023611in}}{%
\pgfpathmoveto{\pgfqpoint{-0.011806in}{-0.023611in}}%
\pgfpathlineto{\pgfqpoint{0.000000in}{-0.011806in}}%
\pgfpathlineto{\pgfqpoint{0.011806in}{-0.023611in}}%
\pgfpathlineto{\pgfqpoint{0.023611in}{-0.011806in}}%
\pgfpathlineto{\pgfqpoint{0.011806in}{0.000000in}}%
\pgfpathlineto{\pgfqpoint{0.023611in}{0.011806in}}%
\pgfpathlineto{\pgfqpoint{0.011806in}{0.023611in}}%
\pgfpathlineto{\pgfqpoint{0.000000in}{0.011806in}}%
\pgfpathlineto{\pgfqpoint{-0.011806in}{0.023611in}}%
\pgfpathlineto{\pgfqpoint{-0.023611in}{0.011806in}}%
\pgfpathlineto{\pgfqpoint{-0.011806in}{0.000000in}}%
\pgfpathlineto{\pgfqpoint{-0.023611in}{-0.011806in}}%
\pgfpathlineto{\pgfqpoint{-0.011806in}{-0.023611in}}%
\pgfpathclose%
\pgfusepath{stroke,fill}%
}%
\begin{pgfscope}%
\pgfsys@transformshift{1.709853in}{0.267528in}%
\pgfsys@useobject{currentmarker}{}%
\end{pgfscope}%
\end{pgfscope}%
\begin{pgfscope}%
\definecolor{textcolor}{rgb}{0.101961,0.101961,0.101961}%
\pgfsetstrokecolor{textcolor}%
\pgfsetfillcolor{textcolor}%
\pgftext[x=1.802214in,y=0.233500in,left,base]{\color{textcolor}\sffamily\fontsize{7.000000}{8.400000}\selectfont On-device only}%
\end{pgfscope}%
\begin{pgfscope}%
\pgfsetrectcap%
\pgfsetroundjoin%
\pgfsetlinewidth{0.903375pt}%
\definecolor{currentstroke}{rgb}{0.435294,0.560784,0.709804}%
\pgfsetstrokecolor{currentstroke}%
\pgfsetdash{}{0pt}%
\pgfpathmoveto{\pgfqpoint{2.519033in}{0.408500in}}%
\pgfpathlineto{\pgfqpoint{2.582228in}{0.408500in}}%
\pgfpathlineto{\pgfqpoint{2.645422in}{0.408500in}}%
\pgfusepath{stroke}%
\end{pgfscope}%
\begin{pgfscope}%
\pgfsetbuttcap%
\pgfsetmiterjoin%
\definecolor{currentfill}{rgb}{0.435294,0.560784,0.709804}%
\pgfsetfillcolor{currentfill}%
\pgfsetlinewidth{1.003750pt}%
\definecolor{currentstroke}{rgb}{0.435294,0.560784,0.709804}%
\pgfsetstrokecolor{currentstroke}%
\pgfsetdash{}{0pt}%
\pgfsys@defobject{currentmarker}{\pgfqpoint{-0.020448in}{-0.023611in}}{\pgfqpoint{0.020448in}{0.023611in}}{%
\pgfpathmoveto{\pgfqpoint{0.000000in}{0.023611in}}%
\pgfpathlineto{\pgfqpoint{-0.020448in}{0.011806in}}%
\pgfpathlineto{\pgfqpoint{-0.020448in}{-0.011806in}}%
\pgfpathlineto{\pgfqpoint{-0.000000in}{-0.023611in}}%
\pgfpathlineto{\pgfqpoint{0.020448in}{-0.011806in}}%
\pgfpathlineto{\pgfqpoint{0.020448in}{0.011806in}}%
\pgfpathlineto{\pgfqpoint{0.000000in}{0.023611in}}%
\pgfpathclose%
\pgfusepath{stroke,fill}%
}%
\begin{pgfscope}%
\pgfsys@transformshift{2.582228in}{0.408500in}%
\pgfsys@useobject{currentmarker}{}%
\end{pgfscope}%
\end{pgfscope}%
\begin{pgfscope}%
\definecolor{textcolor}{rgb}{0.101961,0.101961,0.101961}%
\pgfsetstrokecolor{textcolor}%
\pgfsetfillcolor{textcolor}%
\pgftext[x=2.674589in,y=0.374472in,left,base]{\color{textcolor}\sffamily\fontsize{7.000000}{8.400000}\selectfont rung 1 rules}%
\end{pgfscope}%
\begin{pgfscope}%
\pgfsetrectcap%
\pgfsetroundjoin%
\pgfsetlinewidth{0.903375pt}%
\definecolor{currentstroke}{rgb}{0.435294,0.560784,0.709804}%
\pgfsetstrokecolor{currentstroke}%
\pgfsetdash{}{0pt}%
\pgfpathmoveto{\pgfqpoint{2.519033in}{0.272929in}}%
\pgfpathlineto{\pgfqpoint{2.582228in}{0.272929in}}%
\pgfpathlineto{\pgfqpoint{2.645422in}{0.272929in}}%
\pgfusepath{stroke}%
\end{pgfscope}%
\begin{pgfscope}%
\pgfsetbuttcap%
\pgfsetmiterjoin%
\definecolor{currentfill}{rgb}{0.435294,0.560784,0.709804}%
\pgfsetfillcolor{currentfill}%
\pgfsetlinewidth{1.003750pt}%
\definecolor{currentstroke}{rgb}{0.435294,0.560784,0.709804}%
\pgfsetstrokecolor{currentstroke}%
\pgfsetdash{}{0pt}%
\pgfsys@defobject{currentmarker}{\pgfqpoint{-0.020035in}{-0.033391in}}{\pgfqpoint{0.020035in}{0.033391in}}{%
\pgfpathmoveto{\pgfqpoint{-0.000000in}{-0.033391in}}%
\pgfpathlineto{\pgfqpoint{0.020035in}{0.000000in}}%
\pgfpathlineto{\pgfqpoint{-0.000000in}{0.033391in}}%
\pgfpathlineto{\pgfqpoint{-0.020035in}{0.000000in}}%
\pgfpathlineto{\pgfqpoint{-0.000000in}{-0.033391in}}%
\pgfpathclose%
\pgfusepath{stroke,fill}%
}%
\begin{pgfscope}%
\pgfsys@transformshift{2.582228in}{0.272929in}%
\pgfsys@useobject{currentmarker}{}%
\end{pgfscope}%
\end{pgfscope}%
\begin{pgfscope}%
\definecolor{textcolor}{rgb}{0.101961,0.101961,0.101961}%
\pgfsetstrokecolor{textcolor}%
\pgfsetfillcolor{textcolor}%
\pgftext[x=2.674589in,y=0.238901in,left,base]{\color{textcolor}\sffamily\fontsize{7.000000}{8.400000}\selectfont rung 0 base}%
\end{pgfscope}%
\end{pgfpicture}%
\makeatother%
\endgroup%
}
\caption{[PLACEHOLDER: OLD DATA, rung 3 stand-in for COVEX, 4 s and 5 s missing] Robot correct-by-deadline fraction against the robot deadline. Points exist at 1.5 s (COVEX rung 3 stand-in, SimpleStream), 2.5 s (LLF, Mooncake, on-device, development rungs) and 3 s (all arms, standard cells); 4 s and 5 s have no cells yet. Each arm's 3 s point uses only the seeds that arm has at the other deadlines (seeds 4-5 at 2.5 s, one seed at 1.5 s), because robot correctness varies strongly across seeds.}
\label{fig:F10_robot_deadline}
\end{figure}
\clearpage
\begin{figure*}[t]
\centering
{\fontfamily{LinuxBiolinumT}\selectfont
%% Creator: Matplotlib, PGF backend
%%
%% To include the figure in your LaTeX document, write
%%   \input{<filename>.pgf}
%%
%% Make sure the required packages are loaded in your preamble
%%   \usepackage{pgf}
%%
%% Also ensure that all the required font packages are loaded; for instance,
%% the lmodern package is sometimes necessary when using math font.
%%   \usepackage{lmodern}
%%
%% Figures using additional raster images can only be included by \input if
%% they are in the same directory as the main LaTeX file. For loading figures
%% from other directories you can use the `import` package
%%   \usepackage{import}
%%
%% and then include the figures with
%%   \import{<path to file>}{<filename>.pgf}
%%
%% Matplotlib used the following preamble
%%
\begingroup%
\makeatletter%
\begin{pgfpicture}%
\pgfpathrectangle{\pgfpointorigin}{\pgfqpoint{7.000000in}{2.200000in}}%
\pgfusepath{use as bounding box, clip}%
\begin{pgfscope}%
\pgfsetbuttcap%
\pgfsetmiterjoin%
\definecolor{currentfill}{rgb}{1.000000,1.000000,1.000000}%
\pgfsetfillcolor{currentfill}%
\pgfsetlinewidth{0.000000pt}%
\definecolor{currentstroke}{rgb}{1.000000,1.000000,1.000000}%
\pgfsetstrokecolor{currentstroke}%
\pgfsetdash{}{0pt}%
\pgfpathmoveto{\pgfqpoint{0.000000in}{0.000000in}}%
\pgfpathlineto{\pgfqpoint{7.000000in}{0.000000in}}%
\pgfpathlineto{\pgfqpoint{7.000000in}{2.200000in}}%
\pgfpathlineto{\pgfqpoint{0.000000in}{2.200000in}}%
\pgfpathlineto{\pgfqpoint{0.000000in}{0.000000in}}%
\pgfpathclose%
\pgfusepath{fill}%
\end{pgfscope}%
\begin{pgfscope}%
\pgfsetbuttcap%
\pgfsetmiterjoin%
\definecolor{currentfill}{rgb}{1.000000,1.000000,1.000000}%
\pgfsetfillcolor{currentfill}%
\pgfsetlinewidth{0.000000pt}%
\definecolor{currentstroke}{rgb}{0.000000,0.000000,0.000000}%
\pgfsetstrokecolor{currentstroke}%
\pgfsetstrokeopacity{0.000000}%
\pgfsetdash{}{0pt}%
\pgfpathmoveto{\pgfqpoint{0.525000in}{0.550000in}}%
\pgfpathlineto{\pgfqpoint{6.860000in}{0.550000in}}%
\pgfpathlineto{\pgfqpoint{6.860000in}{1.958000in}}%
\pgfpathlineto{\pgfqpoint{0.525000in}{1.958000in}}%
\pgfpathlineto{\pgfqpoint{0.525000in}{0.550000in}}%
\pgfpathclose%
\pgfusepath{fill}%
\end{pgfscope}%
\begin{pgfscope}%
\definecolor{textcolor}{rgb}{0.101961,0.101961,0.101961}%
\pgfsetstrokecolor{textcolor}%
\pgfsetfillcolor{textcolor}%
\pgftext[x=1.138065in,y=0.501389in,,top]{\color{textcolor}\sffamily\fontsize{7.000000}{8.400000}\selectfont dashcam-heavy}%
\end{pgfscope}%
\begin{pgfscope}%
\definecolor{textcolor}{rgb}{0.101961,0.101961,0.101961}%
\pgfsetstrokecolor{textcolor}%
\pgfsetfillcolor{textcolor}%
\pgftext[x=1.904543in, y=0.433874in, left, base]{\color{textcolor}\sffamily\fontsize{7.000000}{8.400000}\selectfont robot-heavy}%
\end{pgfscope}%
\begin{pgfscope}%
\definecolor{textcolor}{rgb}{0.101961,0.101961,0.101961}%
\pgfsetstrokecolor{textcolor}%
\pgfsetfillcolor{textcolor}%
\pgftext[x=2.027844in, y=0.333951in, left, base]{\color{textcolor}\sffamily\fontsize{7.000000}{8.400000}\selectfont 3:11:2}%
\end{pgfscope}%
\begin{pgfscope}%
\definecolor{textcolor}{rgb}{0.101961,0.101961,0.101961}%
\pgfsetstrokecolor{textcolor}%
\pgfsetfillcolor{textcolor}%
\pgftext[x=3.053957in, y=0.433874in, left, base]{\color{textcolor}\sffamily\fontsize{7.000000}{8.400000}\selectfont re-ask}%
\end{pgfscope}%
\begin{pgfscope}%
\definecolor{textcolor}{rgb}{0.101961,0.101961,0.101961}%
\pgfsetstrokecolor{textcolor}%
\pgfsetfillcolor{textcolor}%
\pgftext[x=3.085655in, y=0.333951in, left, base]{\color{textcolor}\sffamily\fontsize{7.000000}{8.400000}\selectfont rq10}%
\end{pgfscope}%
\begin{pgfscope}%
\definecolor{textcolor}{rgb}{0.101961,0.101961,0.101961}%
\pgfsetstrokecolor{textcolor}%
\pgfsetfillcolor{textcolor}%
\pgftext[x=4.075731in, y=0.433874in, left, base]{\color{textcolor}\sffamily\fontsize{7.000000}{8.400000}\selectfont re-ask}%
\end{pgfscope}%
\begin{pgfscope}%
\definecolor{textcolor}{rgb}{0.101961,0.101961,0.101961}%
\pgfsetstrokecolor{textcolor}%
\pgfsetfillcolor{textcolor}%
\pgftext[x=4.107430in, y=0.333951in, left, base]{\color{textcolor}\sffamily\fontsize{7.000000}{8.400000}\selectfont rq30}%
\end{pgfscope}%
\begin{pgfscope}%
\definecolor{textcolor}{rgb}{0.101961,0.101961,0.101961}%
\pgfsetstrokecolor{textcolor}%
\pgfsetfillcolor{textcolor}%
\pgftext[x=5.097505in, y=0.433874in, left, base]{\color{textcolor}\sffamily\fontsize{7.000000}{8.400000}\selectfont re-ask}%
\end{pgfscope}%
\begin{pgfscope}%
\definecolor{textcolor}{rgb}{0.101961,0.101961,0.101961}%
\pgfsetstrokecolor{textcolor}%
\pgfsetfillcolor{textcolor}%
\pgftext[x=5.129204in, y=0.333951in, left, base]{\color{textcolor}\sffamily\fontsize{7.000000}{8.400000}\selectfont rq60}%
\end{pgfscope}%
\begin{pgfscope}%
\definecolor{textcolor}{rgb}{0.101961,0.101961,0.101961}%
\pgfsetstrokecolor{textcolor}%
\pgfsetfillcolor{textcolor}%
\pgftext[x=5.986914in, y=0.433874in, left, base]{\color{textcolor}\sffamily\fontsize{7.000000}{8.400000}\selectfont mixed robot}%
\end{pgfscope}%
\begin{pgfscope}%
\definecolor{textcolor}{rgb}{0.101961,0.101961,0.101961}%
\pgfsetstrokecolor{textcolor}%
\pgfsetfillcolor{textcolor}%
\pgftext[x=6.056202in, y=0.333951in, left, base]{\color{textcolor}\sffamily\fontsize{7.000000}{8.400000}\selectfont workload}%
\end{pgfscope}%
\begin{pgfscope}%
\definecolor{textcolor}{rgb}{0.000000,0.000000,0.000000}%
\pgfsetstrokecolor{textcolor}%
\pgfsetfillcolor{textcolor}%
\pgftext[x=3.692500in,y=0.287269in,,top]{\color{textcolor}\sffamily\fontsize{7.500000}{9.000000}\selectfont Workload composition}%
\end{pgfscope}%
\begin{pgfscope}%
\pgfpathrectangle{\pgfqpoint{0.525000in}{0.550000in}}{\pgfqpoint{6.335000in}{1.408000in}}%
\pgfusepath{clip}%
\pgfsetrectcap%
\pgfsetroundjoin%
\pgfsetlinewidth{0.401500pt}%
\definecolor{currentstroke}{rgb}{0.901961,0.901961,0.901961}%
\pgfsetstrokecolor{currentstroke}%
\pgfsetdash{}{0pt}%
\pgfpathmoveto{\pgfqpoint{0.525000in}{0.550000in}}%
\pgfpathlineto{\pgfqpoint{6.860000in}{0.550000in}}%
\pgfusepath{stroke}%
\end{pgfscope}%
\begin{pgfscope}%
\definecolor{textcolor}{rgb}{0.101961,0.101961,0.101961}%
\pgfsetstrokecolor{textcolor}%
\pgfsetfillcolor{textcolor}%
\pgftext[x=0.424739in, y=0.516242in, left, base]{\color{textcolor}\sffamily\fontsize{7.000000}{8.400000}\selectfont 0}%
\end{pgfscope}%
\begin{pgfscope}%
\pgfpathrectangle{\pgfqpoint{0.525000in}{0.550000in}}{\pgfqpoint{6.335000in}{1.408000in}}%
\pgfusepath{clip}%
\pgfsetrectcap%
\pgfsetroundjoin%
\pgfsetlinewidth{0.401500pt}%
\definecolor{currentstroke}{rgb}{0.901961,0.901961,0.901961}%
\pgfsetstrokecolor{currentstroke}%
\pgfsetdash{}{0pt}%
\pgfpathmoveto{\pgfqpoint{0.525000in}{1.254000in}}%
\pgfpathlineto{\pgfqpoint{6.860000in}{1.254000in}}%
\pgfusepath{stroke}%
\end{pgfscope}%
\begin{pgfscope}%
\definecolor{textcolor}{rgb}{0.101961,0.101961,0.101961}%
\pgfsetstrokecolor{textcolor}%
\pgfsetfillcolor{textcolor}%
\pgftext[x=0.344395in, y=1.220242in, left, base]{\color{textcolor}\sffamily\fontsize{7.000000}{8.400000}\selectfont 0.5}%
\end{pgfscope}%
\begin{pgfscope}%
\pgfpathrectangle{\pgfqpoint{0.525000in}{0.550000in}}{\pgfqpoint{6.335000in}{1.408000in}}%
\pgfusepath{clip}%
\pgfsetrectcap%
\pgfsetroundjoin%
\pgfsetlinewidth{0.401500pt}%
\definecolor{currentstroke}{rgb}{0.901961,0.901961,0.901961}%
\pgfsetstrokecolor{currentstroke}%
\pgfsetdash{}{0pt}%
\pgfpathmoveto{\pgfqpoint{0.525000in}{1.958000in}}%
\pgfpathlineto{\pgfqpoint{6.860000in}{1.958000in}}%
\pgfusepath{stroke}%
\end{pgfscope}%
\begin{pgfscope}%
\definecolor{textcolor}{rgb}{0.101961,0.101961,0.101961}%
\pgfsetstrokecolor{textcolor}%
\pgfsetfillcolor{textcolor}%
\pgftext[x=0.424739in, y=1.924242in, left, base]{\color{textcolor}\sffamily\fontsize{7.000000}{8.400000}\selectfont 1}%
\end{pgfscope}%
\begin{pgfscope}%
\definecolor{textcolor}{rgb}{0.000000,0.000000,0.000000}%
\pgfsetstrokecolor{textcolor}%
\pgfsetfillcolor{textcolor}%
\pgftext[x=0.288839in,y=1.254000in,,bottom,rotate=90.000000]{\color{textcolor}\sffamily\fontsize{7.500000}{9.000000}\selectfont Correct by deadline}%
\end{pgfscope}%
\begin{pgfscope}%
\pgfsetrectcap%
\pgfsetmiterjoin%
\pgfsetlinewidth{0.803000pt}%
\definecolor{currentstroke}{rgb}{0.101961,0.101961,0.101961}%
\pgfsetstrokecolor{currentstroke}%
\pgfsetdash{}{0pt}%
\pgfpathmoveto{\pgfqpoint{0.525000in}{0.550000in}}%
\pgfpathlineto{\pgfqpoint{0.525000in}{1.958000in}}%
\pgfusepath{stroke}%
\end{pgfscope}%
\begin{pgfscope}%
\pgfsetrectcap%
\pgfsetmiterjoin%
\pgfsetlinewidth{0.803000pt}%
\definecolor{currentstroke}{rgb}{0.101961,0.101961,0.101961}%
\pgfsetstrokecolor{currentstroke}%
\pgfsetdash{}{0pt}%
\pgfpathmoveto{\pgfqpoint{0.525000in}{0.550000in}}%
\pgfpathlineto{\pgfqpoint{6.860000in}{0.550000in}}%
\pgfusepath{stroke}%
\end{pgfscope}%
\begin{pgfscope}%
\pgfsetrectcap%
\pgfsetroundjoin%
\pgfsetlinewidth{0.000000pt}%
\definecolor{currentstroke}{rgb}{0.156863,0.356863,0.568627}%
\pgfsetstrokecolor{currentstroke}%
\pgfsetdash{}{0pt}%
\pgfpathmoveto{\pgfqpoint{3.462155in}{1.831611in}}%
\pgfpathlineto{\pgfqpoint{3.559377in}{1.831611in}}%
\pgfpathlineto{\pgfqpoint{3.656599in}{1.831611in}}%
\pgfusepath{}%
\end{pgfscope}%
\begin{pgfscope}%
\pgfsetbuttcap%
\pgfsetroundjoin%
\definecolor{currentfill}{rgb}{0.156863,0.356863,0.568627}%
\pgfsetfillcolor{currentfill}%
\pgfsetlinewidth{1.003750pt}%
\definecolor{currentstroke}{rgb}{0.156863,0.356863,0.568627}%
\pgfsetstrokecolor{currentstroke}%
\pgfsetdash{}{0pt}%
\pgfsys@defobject{currentmarker}{\pgfqpoint{-0.027778in}{-0.027778in}}{\pgfqpoint{0.027778in}{0.027778in}}{%
\pgfpathmoveto{\pgfqpoint{0.000000in}{-0.027778in}}%
\pgfpathcurveto{\pgfqpoint{0.007367in}{-0.027778in}}{\pgfqpoint{0.014433in}{-0.024851in}}{\pgfqpoint{0.019642in}{-0.019642in}}%
\pgfpathcurveto{\pgfqpoint{0.024851in}{-0.014433in}}{\pgfqpoint{0.027778in}{-0.007367in}}{\pgfqpoint{0.027778in}{0.000000in}}%
\pgfpathcurveto{\pgfqpoint{0.027778in}{0.007367in}}{\pgfqpoint{0.024851in}{0.014433in}}{\pgfqpoint{0.019642in}{0.019642in}}%
\pgfpathcurveto{\pgfqpoint{0.014433in}{0.024851in}}{\pgfqpoint{0.007367in}{0.027778in}}{\pgfqpoint{0.000000in}{0.027778in}}%
\pgfpathcurveto{\pgfqpoint{-0.007367in}{0.027778in}}{\pgfqpoint{-0.014433in}{0.024851in}}{\pgfqpoint{-0.019642in}{0.019642in}}%
\pgfpathcurveto{\pgfqpoint{-0.024851in}{0.014433in}}{\pgfqpoint{-0.027778in}{0.007367in}}{\pgfqpoint{-0.027778in}{0.000000in}}%
\pgfpathcurveto{\pgfqpoint{-0.027778in}{-0.007367in}}{\pgfqpoint{-0.024851in}{-0.014433in}}{\pgfqpoint{-0.019642in}{-0.019642in}}%
\pgfpathcurveto{\pgfqpoint{-0.014433in}{-0.024851in}}{\pgfqpoint{-0.007367in}{-0.027778in}}{\pgfqpoint{0.000000in}{-0.027778in}}%
\pgfpathlineto{\pgfqpoint{0.000000in}{-0.027778in}}%
\pgfpathclose%
\pgfusepath{stroke,fill}%
}%
\begin{pgfscope}%
\pgfsys@transformshift{3.559377in}{1.831611in}%
\pgfsys@useobject{currentmarker}{}%
\end{pgfscope}%
\end{pgfscope}%
\begin{pgfscope}%
\definecolor{textcolor}{rgb}{0.101961,0.101961,0.101961}%
\pgfsetstrokecolor{textcolor}%
\pgfsetfillcolor{textcolor}%
\pgftext[x=3.676044in,y=1.797583in,left,base]{\color{textcolor}\sffamily\fontsize{7.000000}{8.400000}\selectfont COVEX (old: rung 3)}%
\end{pgfscope}%
\begin{pgfscope}%
\pgfsetrectcap%
\pgfsetroundjoin%
\pgfsetlinewidth{0.000000pt}%
\definecolor{currentstroke}{rgb}{0.478431,0.478431,0.478431}%
\pgfsetstrokecolor{currentstroke}%
\pgfsetdash{}{0pt}%
\pgfpathmoveto{\pgfqpoint{4.711067in}{1.831611in}}%
\pgfpathlineto{\pgfqpoint{4.808289in}{1.831611in}}%
\pgfpathlineto{\pgfqpoint{4.905512in}{1.831611in}}%
\pgfusepath{}%
\end{pgfscope}%
\begin{pgfscope}%
\pgfsetbuttcap%
\pgfsetroundjoin%
\definecolor{currentfill}{rgb}{0.478431,0.478431,0.478431}%
\pgfsetfillcolor{currentfill}%
\pgfsetlinewidth{1.003750pt}%
\definecolor{currentstroke}{rgb}{0.478431,0.478431,0.478431}%
\pgfsetstrokecolor{currentstroke}%
\pgfsetdash{}{0pt}%
\pgfsys@defobject{currentmarker}{\pgfqpoint{-0.027778in}{-0.027778in}}{\pgfqpoint{0.027778in}{0.027778in}}{%
\pgfpathmoveto{\pgfqpoint{0.000000in}{-0.027778in}}%
\pgfpathcurveto{\pgfqpoint{0.007367in}{-0.027778in}}{\pgfqpoint{0.014433in}{-0.024851in}}{\pgfqpoint{0.019642in}{-0.019642in}}%
\pgfpathcurveto{\pgfqpoint{0.024851in}{-0.014433in}}{\pgfqpoint{0.027778in}{-0.007367in}}{\pgfqpoint{0.027778in}{0.000000in}}%
\pgfpathcurveto{\pgfqpoint{0.027778in}{0.007367in}}{\pgfqpoint{0.024851in}{0.014433in}}{\pgfqpoint{0.019642in}{0.019642in}}%
\pgfpathcurveto{\pgfqpoint{0.014433in}{0.024851in}}{\pgfqpoint{0.007367in}{0.027778in}}{\pgfqpoint{0.000000in}{0.027778in}}%
\pgfpathcurveto{\pgfqpoint{-0.007367in}{0.027778in}}{\pgfqpoint{-0.014433in}{0.024851in}}{\pgfqpoint{-0.019642in}{0.019642in}}%
\pgfpathcurveto{\pgfqpoint{-0.024851in}{0.014433in}}{\pgfqpoint{-0.027778in}{0.007367in}}{\pgfqpoint{-0.027778in}{0.000000in}}%
\pgfpathcurveto{\pgfqpoint{-0.027778in}{-0.007367in}}{\pgfqpoint{-0.024851in}{-0.014433in}}{\pgfqpoint{-0.019642in}{-0.019642in}}%
\pgfpathcurveto{\pgfqpoint{-0.014433in}{-0.024851in}}{\pgfqpoint{-0.007367in}{-0.027778in}}{\pgfqpoint{0.000000in}{-0.027778in}}%
\pgfpathlineto{\pgfqpoint{0.000000in}{-0.027778in}}%
\pgfpathclose%
\pgfusepath{stroke,fill}%
}%
\begin{pgfscope}%
\pgfsys@transformshift{4.808289in}{1.831611in}%
\pgfsys@useobject{currentmarker}{}%
\end{pgfscope}%
\end{pgfscope}%
\begin{pgfscope}%
\definecolor{textcolor}{rgb}{0.101961,0.101961,0.101961}%
\pgfsetstrokecolor{textcolor}%
\pgfsetfillcolor{textcolor}%
\pgftext[x=4.924956in,y=1.797583in,left,base]{\color{textcolor}\sffamily\fontsize{7.000000}{8.400000}\selectfont Stateful baselines}%
\end{pgfscope}%
\begin{pgfscope}%
\pgfsetrectcap%
\pgfsetroundjoin%
\pgfsetlinewidth{0.000000pt}%
\definecolor{currentstroke}{rgb}{0.784314,0.784314,0.784314}%
\pgfsetstrokecolor{currentstroke}%
\pgfsetdash{}{0pt}%
\pgfpathmoveto{\pgfqpoint{5.770292in}{1.831611in}}%
\pgfpathlineto{\pgfqpoint{5.867514in}{1.831611in}}%
\pgfpathlineto{\pgfqpoint{5.964736in}{1.831611in}}%
\pgfusepath{}%
\end{pgfscope}%
\begin{pgfscope}%
\pgfsetbuttcap%
\pgfsetroundjoin%
\definecolor{currentfill}{rgb}{0.784314,0.784314,0.784314}%
\pgfsetfillcolor{currentfill}%
\pgfsetlinewidth{1.003750pt}%
\definecolor{currentstroke}{rgb}{0.784314,0.784314,0.784314}%
\pgfsetstrokecolor{currentstroke}%
\pgfsetdash{}{0pt}%
\pgfsys@defobject{currentmarker}{\pgfqpoint{-0.027778in}{-0.027778in}}{\pgfqpoint{0.027778in}{0.027778in}}{%
\pgfpathmoveto{\pgfqpoint{0.000000in}{-0.027778in}}%
\pgfpathcurveto{\pgfqpoint{0.007367in}{-0.027778in}}{\pgfqpoint{0.014433in}{-0.024851in}}{\pgfqpoint{0.019642in}{-0.019642in}}%
\pgfpathcurveto{\pgfqpoint{0.024851in}{-0.014433in}}{\pgfqpoint{0.027778in}{-0.007367in}}{\pgfqpoint{0.027778in}{0.000000in}}%
\pgfpathcurveto{\pgfqpoint{0.027778in}{0.007367in}}{\pgfqpoint{0.024851in}{0.014433in}}{\pgfqpoint{0.019642in}{0.019642in}}%
\pgfpathcurveto{\pgfqpoint{0.014433in}{0.024851in}}{\pgfqpoint{0.007367in}{0.027778in}}{\pgfqpoint{0.000000in}{0.027778in}}%
\pgfpathcurveto{\pgfqpoint{-0.007367in}{0.027778in}}{\pgfqpoint{-0.014433in}{0.024851in}}{\pgfqpoint{-0.019642in}{0.019642in}}%
\pgfpathcurveto{\pgfqpoint{-0.024851in}{0.014433in}}{\pgfqpoint{-0.027778in}{0.007367in}}{\pgfqpoint{-0.027778in}{0.000000in}}%
\pgfpathcurveto{\pgfqpoint{-0.027778in}{-0.007367in}}{\pgfqpoint{-0.024851in}{-0.014433in}}{\pgfqpoint{-0.019642in}{-0.019642in}}%
\pgfpathcurveto{\pgfqpoint{-0.014433in}{-0.024851in}}{\pgfqpoint{-0.007367in}{-0.027778in}}{\pgfqpoint{0.000000in}{-0.027778in}}%
\pgfpathlineto{\pgfqpoint{0.000000in}{-0.027778in}}%
\pgfpathclose%
\pgfusepath{stroke,fill}%
}%
\begin{pgfscope}%
\pgfsys@transformshift{5.867514in}{1.831611in}%
\pgfsys@useobject{currentmarker}{}%
\end{pgfscope}%
\end{pgfscope}%
\begin{pgfscope}%
\definecolor{textcolor}{rgb}{0.101961,0.101961,0.101961}%
\pgfsetstrokecolor{textcolor}%
\pgfsetfillcolor{textcolor}%
\pgftext[x=5.984181in,y=1.797583in,left,base]{\color{textcolor}\sffamily\fontsize{7.000000}{8.400000}\selectfont Stateless baselines}%
\end{pgfscope}%
\begin{pgfscope}%
\pgfsetbuttcap%
\pgfsetmiterjoin%
\definecolor{currentfill}{rgb}{1.000000,1.000000,1.000000}%
\pgfsetfillcolor{currentfill}%
\pgfsetlinewidth{0.000000pt}%
\definecolor{currentstroke}{rgb}{0.000000,0.000000,0.000000}%
\pgfsetstrokecolor{currentstroke}%
\pgfsetstrokeopacity{0.000000}%
\pgfsetdash{}{0pt}%
\pgfpathmoveto{\pgfqpoint{1.099596in}{1.077292in}}%
\pgfpathlineto{\pgfqpoint{6.320404in}{1.077292in}}%
\pgfpathlineto{\pgfqpoint{6.320404in}{1.342708in}}%
\pgfpathlineto{\pgfqpoint{1.099596in}{1.342708in}}%
\pgfpathlineto{\pgfqpoint{1.099596in}{1.077292in}}%
\pgfpathclose%
\pgfusepath{fill}%
\end{pgfscope}%
\begin{pgfscope}%
\definecolor{textcolor}{rgb}{0.549020,0.549020,0.549020}%
\pgfsetstrokecolor{textcolor}%
\pgfsetfillcolor{textcolor}%
\pgftext[x=1.127790in, y=1.241597in, left, base]{\color{textcolor}\sffamily\fontsize{7.000000}{8.400000}\selectfont NO DATA YET: composition cells: dashcam-heavy mix, robot-heavy 3:11:2 mix, re-ask rq10 and rq60 (rq30 exists only}%
\end{pgfscope}%
\begin{pgfscope}%
\definecolor{textcolor}{rgb}{0.549020,0.549020,0.549020}%
\pgfsetstrokecolor{textcolor}%
\pgfsetfillcolor{textcolor}%
\pgftext[x=2.124348in, y=1.129792in, left, base]{\color{textcolor}\sffamily\fontsize{7.000000}{8.400000}\selectfont in the standard 40:40:20 grid), mixed robot workload; F11.csv has 0 rows}%
\end{pgfscope}%
\begin{pgfscope}%
\definecolor{textcolor}{rgb}{0.549020,0.549020,0.549020}%
\pgfsetstrokecolor{textcolor}%
\pgfsetfillcolor{textcolor}%
\pgftext[x=6.930000in,y=2.167000in,right,top]{\color{textcolor}\sffamily\fontsize{5.500000}{6.600000}\selectfont OLD DATA}%
\end{pgfscope}%
\end{pgfpicture}%
\makeatother%
\endgroup%
}
\caption{[PLACEHOLDER: OLD DATA, NO DATA, rung 3 stand-in for COVEX] Correct-by-deadline under workload composition: dashcam-heavy, robot-heavy 3:11:2, re-ask rq10, rq30 and rq60, and a mixed robot workload. Empty: none of the composition cells exist yet.}
\label{fig:F11_composition}
\end{figure*}
\clearpage
\begin{figure*}[t]
\centering
{\fontfamily{LinuxBiolinumT}\selectfont
%% Creator: Matplotlib, PGF backend
%%
%% To include the figure in your LaTeX document, write
%%   \input{<filename>.pgf}
%%
%% Make sure the required packages are loaded in your preamble
%%   \usepackage{pgf}
%%
%% Also ensure that all the required font packages are loaded; for instance,
%% the lmodern package is sometimes necessary when using math font.
%%   \usepackage{lmodern}
%%
%% Figures using additional raster images can only be included by \input if
%% they are in the same directory as the main LaTeX file. For loading figures
%% from other directories you can use the `import` package
%%   \usepackage{import}
%%
%% and then include the figures with
%%   \import{<path to file>}{<filename>.pgf}
%%
%% Matplotlib used the following preamble
%%
\begingroup%
\makeatletter%
\begin{pgfpicture}%
\pgfpathrectangle{\pgfpointorigin}{\pgfqpoint{7.000000in}{3.900000in}}%
\pgfusepath{use as bounding box, clip}%
\begin{pgfscope}%
\pgfsetbuttcap%
\pgfsetmiterjoin%
\definecolor{currentfill}{rgb}{1.000000,1.000000,1.000000}%
\pgfsetfillcolor{currentfill}%
\pgfsetlinewidth{0.000000pt}%
\definecolor{currentstroke}{rgb}{1.000000,1.000000,1.000000}%
\pgfsetstrokecolor{currentstroke}%
\pgfsetdash{}{0pt}%
\pgfpathmoveto{\pgfqpoint{0.000000in}{0.000000in}}%
\pgfpathlineto{\pgfqpoint{7.000000in}{0.000000in}}%
\pgfpathlineto{\pgfqpoint{7.000000in}{3.900000in}}%
\pgfpathlineto{\pgfqpoint{0.000000in}{3.900000in}}%
\pgfpathlineto{\pgfqpoint{0.000000in}{0.000000in}}%
\pgfpathclose%
\pgfusepath{fill}%
\end{pgfscope}%
\begin{pgfscope}%
\pgfsetbuttcap%
\pgfsetmiterjoin%
\definecolor{currentfill}{rgb}{1.000000,1.000000,1.000000}%
\pgfsetfillcolor{currentfill}%
\pgfsetlinewidth{0.000000pt}%
\definecolor{currentstroke}{rgb}{0.000000,0.000000,0.000000}%
\pgfsetstrokecolor{currentstroke}%
\pgfsetstrokeopacity{0.000000}%
\pgfsetdash{}{0pt}%
\pgfpathmoveto{\pgfqpoint{0.595000in}{2.488200in}}%
\pgfpathlineto{\pgfqpoint{2.515732in}{2.488200in}}%
\pgfpathlineto{\pgfqpoint{2.515732in}{3.705000in}}%
\pgfpathlineto{\pgfqpoint{0.595000in}{3.705000in}}%
\pgfpathlineto{\pgfqpoint{0.595000in}{2.488200in}}%
\pgfpathclose%
\pgfusepath{fill}%
\end{pgfscope}%
\begin{pgfscope}%
\pgfpathrectangle{\pgfqpoint{0.595000in}{2.488200in}}{\pgfqpoint{1.920732in}{1.216800in}}%
\pgfusepath{clip}%
\pgfsetrectcap%
\pgfsetroundjoin%
\pgfsetlinewidth{0.401500pt}%
\definecolor{currentstroke}{rgb}{0.901961,0.901961,0.901961}%
\pgfsetstrokecolor{currentstroke}%
\pgfsetdash{}{0pt}%
\pgfpathmoveto{\pgfqpoint{0.595000in}{2.488200in}}%
\pgfpathlineto{\pgfqpoint{0.595000in}{3.705000in}}%
\pgfusepath{stroke}%
\end{pgfscope}%
\begin{pgfscope}%
\definecolor{textcolor}{rgb}{0.101961,0.101961,0.101961}%
\pgfsetstrokecolor{textcolor}%
\pgfsetfillcolor{textcolor}%
\pgftext[x=0.595000in,y=2.439589in,,top]{\color{textcolor}\sffamily\fontsize{7.000000}{8.400000}\selectfont \ensuremath{-}60}%
\end{pgfscope}%
\begin{pgfscope}%
\pgfpathrectangle{\pgfqpoint{0.595000in}{2.488200in}}{\pgfqpoint{1.920732in}{1.216800in}}%
\pgfusepath{clip}%
\pgfsetrectcap%
\pgfsetroundjoin%
\pgfsetlinewidth{0.401500pt}%
\definecolor{currentstroke}{rgb}{0.901961,0.901961,0.901961}%
\pgfsetstrokecolor{currentstroke}%
\pgfsetdash{}{0pt}%
\pgfpathmoveto{\pgfqpoint{0.915122in}{2.488200in}}%
\pgfpathlineto{\pgfqpoint{0.915122in}{3.705000in}}%
\pgfusepath{stroke}%
\end{pgfscope}%
\begin{pgfscope}%
\definecolor{textcolor}{rgb}{0.101961,0.101961,0.101961}%
\pgfsetstrokecolor{textcolor}%
\pgfsetfillcolor{textcolor}%
\pgftext[x=0.915122in,y=2.439589in,,top]{\color{textcolor}\sffamily\fontsize{7.000000}{8.400000}\selectfont 0}%
\end{pgfscope}%
\begin{pgfscope}%
\pgfpathrectangle{\pgfqpoint{0.595000in}{2.488200in}}{\pgfqpoint{1.920732in}{1.216800in}}%
\pgfusepath{clip}%
\pgfsetrectcap%
\pgfsetroundjoin%
\pgfsetlinewidth{0.401500pt}%
\definecolor{currentstroke}{rgb}{0.901961,0.901961,0.901961}%
\pgfsetstrokecolor{currentstroke}%
\pgfsetdash{}{0pt}%
\pgfpathmoveto{\pgfqpoint{1.448659in}{2.488200in}}%
\pgfpathlineto{\pgfqpoint{1.448659in}{3.705000in}}%
\pgfusepath{stroke}%
\end{pgfscope}%
\begin{pgfscope}%
\definecolor{textcolor}{rgb}{0.101961,0.101961,0.101961}%
\pgfsetstrokecolor{textcolor}%
\pgfsetfillcolor{textcolor}%
\pgftext[x=1.448659in,y=2.439589in,,top]{\color{textcolor}\sffamily\fontsize{7.000000}{8.400000}\selectfont 100}%
\end{pgfscope}%
\begin{pgfscope}%
\pgfpathrectangle{\pgfqpoint{0.595000in}{2.488200in}}{\pgfqpoint{1.920732in}{1.216800in}}%
\pgfusepath{clip}%
\pgfsetrectcap%
\pgfsetroundjoin%
\pgfsetlinewidth{0.401500pt}%
\definecolor{currentstroke}{rgb}{0.901961,0.901961,0.901961}%
\pgfsetstrokecolor{currentstroke}%
\pgfsetdash{}{0pt}%
\pgfpathmoveto{\pgfqpoint{1.982195in}{2.488200in}}%
\pgfpathlineto{\pgfqpoint{1.982195in}{3.705000in}}%
\pgfusepath{stroke}%
\end{pgfscope}%
\begin{pgfscope}%
\definecolor{textcolor}{rgb}{0.101961,0.101961,0.101961}%
\pgfsetstrokecolor{textcolor}%
\pgfsetfillcolor{textcolor}%
\pgftext[x=1.982195in,y=2.439589in,,top]{\color{textcolor}\sffamily\fontsize{7.000000}{8.400000}\selectfont 200}%
\end{pgfscope}%
\begin{pgfscope}%
\pgfpathrectangle{\pgfqpoint{0.595000in}{2.488200in}}{\pgfqpoint{1.920732in}{1.216800in}}%
\pgfusepath{clip}%
\pgfsetrectcap%
\pgfsetroundjoin%
\pgfsetlinewidth{0.401500pt}%
\definecolor{currentstroke}{rgb}{0.901961,0.901961,0.901961}%
\pgfsetstrokecolor{currentstroke}%
\pgfsetdash{}{0pt}%
\pgfpathmoveto{\pgfqpoint{2.515732in}{2.488200in}}%
\pgfpathlineto{\pgfqpoint{2.515732in}{3.705000in}}%
\pgfusepath{stroke}%
\end{pgfscope}%
\begin{pgfscope}%
\definecolor{textcolor}{rgb}{0.101961,0.101961,0.101961}%
\pgfsetstrokecolor{textcolor}%
\pgfsetfillcolor{textcolor}%
\pgftext[x=2.515732in,y=2.439589in,,top]{\color{textcolor}\sffamily\fontsize{7.000000}{8.400000}\selectfont 300}%
\end{pgfscope}%
\begin{pgfscope}%
\definecolor{textcolor}{rgb}{0.000000,0.000000,0.000000}%
\pgfsetstrokecolor{textcolor}%
\pgfsetfillcolor{textcolor}%
\pgftext[x=1.555366in,y=2.297614in,,top]{\color{textcolor}\sffamily\fontsize{7.500000}{9.000000}\selectfont s from reconnect R}%
\end{pgfscope}%
\begin{pgfscope}%
\pgfpathrectangle{\pgfqpoint{0.595000in}{2.488200in}}{\pgfqpoint{1.920732in}{1.216800in}}%
\pgfusepath{clip}%
\pgfsetrectcap%
\pgfsetroundjoin%
\pgfsetlinewidth{0.401500pt}%
\definecolor{currentstroke}{rgb}{0.901961,0.901961,0.901961}%
\pgfsetstrokecolor{currentstroke}%
\pgfsetdash{}{0pt}%
\pgfpathmoveto{\pgfqpoint{0.595000in}{2.522000in}}%
\pgfpathlineto{\pgfqpoint{2.515732in}{2.522000in}}%
\pgfusepath{stroke}%
\end{pgfscope}%
\begin{pgfscope}%
\definecolor{textcolor}{rgb}{0.101961,0.101961,0.101961}%
\pgfsetstrokecolor{textcolor}%
\pgfsetfillcolor{textcolor}%
\pgftext[x=0.414395in, y=2.488242in, left, base]{\color{textcolor}\sffamily\fontsize{7.000000}{8.400000}\selectfont 0.0}%
\end{pgfscope}%
\begin{pgfscope}%
\pgfpathrectangle{\pgfqpoint{0.595000in}{2.488200in}}{\pgfqpoint{1.920732in}{1.216800in}}%
\pgfusepath{clip}%
\pgfsetrectcap%
\pgfsetroundjoin%
\pgfsetlinewidth{0.401500pt}%
\definecolor{currentstroke}{rgb}{0.901961,0.901961,0.901961}%
\pgfsetstrokecolor{currentstroke}%
\pgfsetdash{}{0pt}%
\pgfpathmoveto{\pgfqpoint{0.595000in}{2.747333in}}%
\pgfpathlineto{\pgfqpoint{2.515732in}{2.747333in}}%
\pgfusepath{stroke}%
\end{pgfscope}%
\begin{pgfscope}%
\definecolor{textcolor}{rgb}{0.101961,0.101961,0.101961}%
\pgfsetstrokecolor{textcolor}%
\pgfsetfillcolor{textcolor}%
\pgftext[x=0.414395in, y=2.713576in, left, base]{\color{textcolor}\sffamily\fontsize{7.000000}{8.400000}\selectfont 0.2}%
\end{pgfscope}%
\begin{pgfscope}%
\pgfpathrectangle{\pgfqpoint{0.595000in}{2.488200in}}{\pgfqpoint{1.920732in}{1.216800in}}%
\pgfusepath{clip}%
\pgfsetrectcap%
\pgfsetroundjoin%
\pgfsetlinewidth{0.401500pt}%
\definecolor{currentstroke}{rgb}{0.901961,0.901961,0.901961}%
\pgfsetstrokecolor{currentstroke}%
\pgfsetdash{}{0pt}%
\pgfpathmoveto{\pgfqpoint{0.595000in}{2.972667in}}%
\pgfpathlineto{\pgfqpoint{2.515732in}{2.972667in}}%
\pgfusepath{stroke}%
\end{pgfscope}%
\begin{pgfscope}%
\definecolor{textcolor}{rgb}{0.101961,0.101961,0.101961}%
\pgfsetstrokecolor{textcolor}%
\pgfsetfillcolor{textcolor}%
\pgftext[x=0.414395in, y=2.938909in, left, base]{\color{textcolor}\sffamily\fontsize{7.000000}{8.400000}\selectfont 0.4}%
\end{pgfscope}%
\begin{pgfscope}%
\pgfpathrectangle{\pgfqpoint{0.595000in}{2.488200in}}{\pgfqpoint{1.920732in}{1.216800in}}%
\pgfusepath{clip}%
\pgfsetrectcap%
\pgfsetroundjoin%
\pgfsetlinewidth{0.401500pt}%
\definecolor{currentstroke}{rgb}{0.901961,0.901961,0.901961}%
\pgfsetstrokecolor{currentstroke}%
\pgfsetdash{}{0pt}%
\pgfpathmoveto{\pgfqpoint{0.595000in}{3.198000in}}%
\pgfpathlineto{\pgfqpoint{2.515732in}{3.198000in}}%
\pgfusepath{stroke}%
\end{pgfscope}%
\begin{pgfscope}%
\definecolor{textcolor}{rgb}{0.101961,0.101961,0.101961}%
\pgfsetstrokecolor{textcolor}%
\pgfsetfillcolor{textcolor}%
\pgftext[x=0.414395in, y=3.164242in, left, base]{\color{textcolor}\sffamily\fontsize{7.000000}{8.400000}\selectfont 0.6}%
\end{pgfscope}%
\begin{pgfscope}%
\pgfpathrectangle{\pgfqpoint{0.595000in}{2.488200in}}{\pgfqpoint{1.920732in}{1.216800in}}%
\pgfusepath{clip}%
\pgfsetrectcap%
\pgfsetroundjoin%
\pgfsetlinewidth{0.401500pt}%
\definecolor{currentstroke}{rgb}{0.901961,0.901961,0.901961}%
\pgfsetstrokecolor{currentstroke}%
\pgfsetdash{}{0pt}%
\pgfpathmoveto{\pgfqpoint{0.595000in}{3.423333in}}%
\pgfpathlineto{\pgfqpoint{2.515732in}{3.423333in}}%
\pgfusepath{stroke}%
\end{pgfscope}%
\begin{pgfscope}%
\definecolor{textcolor}{rgb}{0.101961,0.101961,0.101961}%
\pgfsetstrokecolor{textcolor}%
\pgfsetfillcolor{textcolor}%
\pgftext[x=0.414395in, y=3.389576in, left, base]{\color{textcolor}\sffamily\fontsize{7.000000}{8.400000}\selectfont 0.8}%
\end{pgfscope}%
\begin{pgfscope}%
\pgfpathrectangle{\pgfqpoint{0.595000in}{2.488200in}}{\pgfqpoint{1.920732in}{1.216800in}}%
\pgfusepath{clip}%
\pgfsetrectcap%
\pgfsetroundjoin%
\pgfsetlinewidth{0.401500pt}%
\definecolor{currentstroke}{rgb}{0.901961,0.901961,0.901961}%
\pgfsetstrokecolor{currentstroke}%
\pgfsetdash{}{0pt}%
\pgfpathmoveto{\pgfqpoint{0.595000in}{3.648667in}}%
\pgfpathlineto{\pgfqpoint{2.515732in}{3.648667in}}%
\pgfusepath{stroke}%
\end{pgfscope}%
\begin{pgfscope}%
\definecolor{textcolor}{rgb}{0.101961,0.101961,0.101961}%
\pgfsetstrokecolor{textcolor}%
\pgfsetfillcolor{textcolor}%
\pgftext[x=0.414395in, y=3.614909in, left, base]{\color{textcolor}\sffamily\fontsize{7.000000}{8.400000}\selectfont 1.0}%
\end{pgfscope}%
\begin{pgfscope}%
\definecolor{textcolor}{rgb}{0.000000,0.000000,0.000000}%
\pgfsetstrokecolor{textcolor}%
\pgfsetfillcolor{textcolor}%
\pgftext[x=0.358839in,y=3.096600in,,bottom,rotate=90.000000]{\color{textcolor}\sffamily\fontsize{7.500000}{9.000000}\selectfont on-time rate}%
\end{pgfscope}%
\begin{pgfscope}%
\pgfpathrectangle{\pgfqpoint{0.595000in}{2.488200in}}{\pgfqpoint{1.920732in}{1.216800in}}%
\pgfusepath{clip}%
\pgfsetbuttcap%
\pgfsetmiterjoin%
\definecolor{currentfill}{rgb}{0.929412,0.929412,0.929412}%
\pgfsetfillcolor{currentfill}%
\pgfsetlinewidth{0.000000pt}%
\definecolor{currentstroke}{rgb}{0.929412,0.929412,0.929412}%
\pgfsetstrokecolor{currentstroke}%
\pgfsetdash{}{0pt}%
\pgfpathmoveto{\pgfqpoint{0.595000in}{2.488200in}}%
\pgfpathlineto{\pgfqpoint{0.595000in}{3.705000in}}%
\pgfpathlineto{\pgfqpoint{0.915122in}{3.705000in}}%
\pgfpathlineto{\pgfqpoint{0.915122in}{2.488200in}}%
\pgfpathlineto{\pgfqpoint{0.595000in}{2.488200in}}%
\pgfpathclose%
\pgfusepath{fill}%
\end{pgfscope}%
\begin{pgfscope}%
\pgfpathrectangle{\pgfqpoint{0.595000in}{2.488200in}}{\pgfqpoint{1.920732in}{1.216800in}}%
\pgfusepath{clip}%
\pgfsetrectcap%
\pgfsetroundjoin%
\pgfsetlinewidth{0.702625pt}%
\definecolor{currentstroke}{rgb}{0.615686,0.705882,0.807843}%
\pgfsetstrokecolor{currentstroke}%
\pgfsetdash{}{0pt}%
\pgfpathmoveto{\pgfqpoint{0.621677in}{3.648667in}}%
\pgfpathlineto{\pgfqpoint{0.675030in}{3.648667in}}%
\pgfpathlineto{\pgfqpoint{0.728384in}{3.648667in}}%
\pgfpathlineto{\pgfqpoint{0.781738in}{3.648667in}}%
\pgfpathlineto{\pgfqpoint{0.835091in}{3.648667in}}%
\pgfpathlineto{\pgfqpoint{0.888445in}{3.648667in}}%
\pgfpathlineto{\pgfqpoint{0.941799in}{2.522000in}}%
\pgfpathlineto{\pgfqpoint{0.995152in}{3.648667in}}%
\pgfpathlineto{\pgfqpoint{1.048506in}{3.648667in}}%
\pgfpathlineto{\pgfqpoint{1.101860in}{3.648667in}}%
\pgfpathlineto{\pgfqpoint{1.155213in}{3.648667in}}%
\pgfpathlineto{\pgfqpoint{1.208567in}{3.648667in}}%
\pgfpathlineto{\pgfqpoint{1.261921in}{3.648667in}}%
\pgfpathlineto{\pgfqpoint{1.315274in}{3.648667in}}%
\pgfpathlineto{\pgfqpoint{1.368628in}{3.648667in}}%
\pgfpathlineto{\pgfqpoint{1.421982in}{3.648667in}}%
\pgfpathlineto{\pgfqpoint{1.475335in}{3.648667in}}%
\pgfpathlineto{\pgfqpoint{1.528689in}{3.648667in}}%
\pgfpathlineto{\pgfqpoint{1.582043in}{3.648667in}}%
\pgfpathlineto{\pgfqpoint{1.635396in}{3.648667in}}%
\pgfpathlineto{\pgfqpoint{1.688750in}{2.522000in}}%
\pgfpathlineto{\pgfqpoint{1.742104in}{3.085333in}}%
\pgfpathlineto{\pgfqpoint{1.795457in}{2.522000in}}%
\pgfpathlineto{\pgfqpoint{1.848811in}{2.522000in}}%
\pgfpathlineto{\pgfqpoint{1.902165in}{2.522000in}}%
\pgfpathlineto{\pgfqpoint{1.955518in}{2.522000in}}%
\pgfpathlineto{\pgfqpoint{2.008872in}{3.648667in}}%
\pgfpathlineto{\pgfqpoint{2.062226in}{2.522000in}}%
\pgfpathlineto{\pgfqpoint{2.115579in}{3.648667in}}%
\pgfpathlineto{\pgfqpoint{2.168933in}{3.648667in}}%
\pgfpathlineto{\pgfqpoint{2.222287in}{3.648667in}}%
\pgfpathlineto{\pgfqpoint{2.275640in}{3.085333in}}%
\pgfpathlineto{\pgfqpoint{2.328994in}{3.648667in}}%
\pgfpathlineto{\pgfqpoint{2.382348in}{3.648667in}}%
\pgfpathlineto{\pgfqpoint{2.435701in}{3.648667in}}%
\pgfpathlineto{\pgfqpoint{2.489055in}{3.085333in}}%
\pgfusepath{stroke}%
\end{pgfscope}%
\begin{pgfscope}%
\pgfpathrectangle{\pgfqpoint{0.595000in}{2.488200in}}{\pgfqpoint{1.920732in}{1.216800in}}%
\pgfusepath{clip}%
\pgfsetbuttcap%
\pgfsetroundjoin%
\definecolor{currentfill}{rgb}{0.615686,0.705882,0.807843}%
\pgfsetfillcolor{currentfill}%
\pgfsetlinewidth{1.003750pt}%
\definecolor{currentstroke}{rgb}{0.615686,0.705882,0.807843}%
\pgfsetstrokecolor{currentstroke}%
\pgfsetdash{}{0pt}%
\pgfsys@defobject{currentmarker}{\pgfqpoint{-0.010417in}{-0.010417in}}{\pgfqpoint{0.010417in}{0.010417in}}{%
\pgfpathmoveto{\pgfqpoint{0.000000in}{-0.010417in}}%
\pgfpathcurveto{\pgfqpoint{0.002763in}{-0.010417in}}{\pgfqpoint{0.005412in}{-0.009319in}}{\pgfqpoint{0.007366in}{-0.007366in}}%
\pgfpathcurveto{\pgfqpoint{0.009319in}{-0.005412in}}{\pgfqpoint{0.010417in}{-0.002763in}}{\pgfqpoint{0.010417in}{0.000000in}}%
\pgfpathcurveto{\pgfqpoint{0.010417in}{0.002763in}}{\pgfqpoint{0.009319in}{0.005412in}}{\pgfqpoint{0.007366in}{0.007366in}}%
\pgfpathcurveto{\pgfqpoint{0.005412in}{0.009319in}}{\pgfqpoint{0.002763in}{0.010417in}}{\pgfqpoint{0.000000in}{0.010417in}}%
\pgfpathcurveto{\pgfqpoint{-0.002763in}{0.010417in}}{\pgfqpoint{-0.005412in}{0.009319in}}{\pgfqpoint{-0.007366in}{0.007366in}}%
\pgfpathcurveto{\pgfqpoint{-0.009319in}{0.005412in}}{\pgfqpoint{-0.010417in}{0.002763in}}{\pgfqpoint{-0.010417in}{0.000000in}}%
\pgfpathcurveto{\pgfqpoint{-0.010417in}{-0.002763in}}{\pgfqpoint{-0.009319in}{-0.005412in}}{\pgfqpoint{-0.007366in}{-0.007366in}}%
\pgfpathcurveto{\pgfqpoint{-0.005412in}{-0.009319in}}{\pgfqpoint{-0.002763in}{-0.010417in}}{\pgfqpoint{0.000000in}{-0.010417in}}%
\pgfpathlineto{\pgfqpoint{0.000000in}{-0.010417in}}%
\pgfpathclose%
\pgfusepath{stroke,fill}%
}%
\begin{pgfscope}%
\pgfsys@transformshift{0.621677in}{3.648667in}%
\pgfsys@useobject{currentmarker}{}%
\end{pgfscope}%
\begin{pgfscope}%
\pgfsys@transformshift{0.675030in}{3.648667in}%
\pgfsys@useobject{currentmarker}{}%
\end{pgfscope}%
\begin{pgfscope}%
\pgfsys@transformshift{0.728384in}{3.648667in}%
\pgfsys@useobject{currentmarker}{}%
\end{pgfscope}%
\begin{pgfscope}%
\pgfsys@transformshift{0.781738in}{3.648667in}%
\pgfsys@useobject{currentmarker}{}%
\end{pgfscope}%
\begin{pgfscope}%
\pgfsys@transformshift{0.835091in}{3.648667in}%
\pgfsys@useobject{currentmarker}{}%
\end{pgfscope}%
\begin{pgfscope}%
\pgfsys@transformshift{0.888445in}{3.648667in}%
\pgfsys@useobject{currentmarker}{}%
\end{pgfscope}%
\begin{pgfscope}%
\pgfsys@transformshift{0.941799in}{2.522000in}%
\pgfsys@useobject{currentmarker}{}%
\end{pgfscope}%
\begin{pgfscope}%
\pgfsys@transformshift{0.995152in}{3.648667in}%
\pgfsys@useobject{currentmarker}{}%
\end{pgfscope}%
\begin{pgfscope}%
\pgfsys@transformshift{1.048506in}{3.648667in}%
\pgfsys@useobject{currentmarker}{}%
\end{pgfscope}%
\begin{pgfscope}%
\pgfsys@transformshift{1.101860in}{3.648667in}%
\pgfsys@useobject{currentmarker}{}%
\end{pgfscope}%
\begin{pgfscope}%
\pgfsys@transformshift{1.155213in}{3.648667in}%
\pgfsys@useobject{currentmarker}{}%
\end{pgfscope}%
\begin{pgfscope}%
\pgfsys@transformshift{1.208567in}{3.648667in}%
\pgfsys@useobject{currentmarker}{}%
\end{pgfscope}%
\begin{pgfscope}%
\pgfsys@transformshift{1.261921in}{3.648667in}%
\pgfsys@useobject{currentmarker}{}%
\end{pgfscope}%
\begin{pgfscope}%
\pgfsys@transformshift{1.315274in}{3.648667in}%
\pgfsys@useobject{currentmarker}{}%
\end{pgfscope}%
\begin{pgfscope}%
\pgfsys@transformshift{1.368628in}{3.648667in}%
\pgfsys@useobject{currentmarker}{}%
\end{pgfscope}%
\begin{pgfscope}%
\pgfsys@transformshift{1.421982in}{3.648667in}%
\pgfsys@useobject{currentmarker}{}%
\end{pgfscope}%
\begin{pgfscope}%
\pgfsys@transformshift{1.475335in}{3.648667in}%
\pgfsys@useobject{currentmarker}{}%
\end{pgfscope}%
\begin{pgfscope}%
\pgfsys@transformshift{1.528689in}{3.648667in}%
\pgfsys@useobject{currentmarker}{}%
\end{pgfscope}%
\begin{pgfscope}%
\pgfsys@transformshift{1.582043in}{3.648667in}%
\pgfsys@useobject{currentmarker}{}%
\end{pgfscope}%
\begin{pgfscope}%
\pgfsys@transformshift{1.635396in}{3.648667in}%
\pgfsys@useobject{currentmarker}{}%
\end{pgfscope}%
\begin{pgfscope}%
\pgfsys@transformshift{1.688750in}{2.522000in}%
\pgfsys@useobject{currentmarker}{}%
\end{pgfscope}%
\begin{pgfscope}%
\pgfsys@transformshift{1.742104in}{3.085333in}%
\pgfsys@useobject{currentmarker}{}%
\end{pgfscope}%
\begin{pgfscope}%
\pgfsys@transformshift{1.795457in}{2.522000in}%
\pgfsys@useobject{currentmarker}{}%
\end{pgfscope}%
\begin{pgfscope}%
\pgfsys@transformshift{1.848811in}{2.522000in}%
\pgfsys@useobject{currentmarker}{}%
\end{pgfscope}%
\begin{pgfscope}%
\pgfsys@transformshift{1.902165in}{2.522000in}%
\pgfsys@useobject{currentmarker}{}%
\end{pgfscope}%
\begin{pgfscope}%
\pgfsys@transformshift{1.955518in}{2.522000in}%
\pgfsys@useobject{currentmarker}{}%
\end{pgfscope}%
\begin{pgfscope}%
\pgfsys@transformshift{2.008872in}{3.648667in}%
\pgfsys@useobject{currentmarker}{}%
\end{pgfscope}%
\begin{pgfscope}%
\pgfsys@transformshift{2.062226in}{2.522000in}%
\pgfsys@useobject{currentmarker}{}%
\end{pgfscope}%
\begin{pgfscope}%
\pgfsys@transformshift{2.115579in}{3.648667in}%
\pgfsys@useobject{currentmarker}{}%
\end{pgfscope}%
\begin{pgfscope}%
\pgfsys@transformshift{2.168933in}{3.648667in}%
\pgfsys@useobject{currentmarker}{}%
\end{pgfscope}%
\begin{pgfscope}%
\pgfsys@transformshift{2.222287in}{3.648667in}%
\pgfsys@useobject{currentmarker}{}%
\end{pgfscope}%
\begin{pgfscope}%
\pgfsys@transformshift{2.275640in}{3.085333in}%
\pgfsys@useobject{currentmarker}{}%
\end{pgfscope}%
\begin{pgfscope}%
\pgfsys@transformshift{2.328994in}{3.648667in}%
\pgfsys@useobject{currentmarker}{}%
\end{pgfscope}%
\begin{pgfscope}%
\pgfsys@transformshift{2.382348in}{3.648667in}%
\pgfsys@useobject{currentmarker}{}%
\end{pgfscope}%
\begin{pgfscope}%
\pgfsys@transformshift{2.435701in}{3.648667in}%
\pgfsys@useobject{currentmarker}{}%
\end{pgfscope}%
\begin{pgfscope}%
\pgfsys@transformshift{2.489055in}{3.085333in}%
\pgfsys@useobject{currentmarker}{}%
\end{pgfscope}%
\end{pgfscope}%
\begin{pgfscope}%
\pgfpathrectangle{\pgfqpoint{0.595000in}{2.488200in}}{\pgfqpoint{1.920732in}{1.216800in}}%
\pgfusepath{clip}%
\pgfsetrectcap%
\pgfsetroundjoin%
\pgfsetlinewidth{0.702625pt}%
\definecolor{currentstroke}{rgb}{0.435294,0.576471,0.729412}%
\pgfsetstrokecolor{currentstroke}%
\pgfsetdash{}{0pt}%
\pgfpathmoveto{\pgfqpoint{0.621677in}{3.648667in}}%
\pgfpathlineto{\pgfqpoint{0.675030in}{3.648667in}}%
\pgfpathlineto{\pgfqpoint{0.728384in}{3.648667in}}%
\pgfpathlineto{\pgfqpoint{0.781738in}{3.648667in}}%
\pgfpathlineto{\pgfqpoint{0.835091in}{3.648667in}}%
\pgfpathlineto{\pgfqpoint{0.888445in}{3.648667in}}%
\pgfpathlineto{\pgfqpoint{0.941799in}{2.522000in}}%
\pgfpathlineto{\pgfqpoint{0.995152in}{3.648667in}}%
\pgfpathlineto{\pgfqpoint{1.048506in}{3.648667in}}%
\pgfpathlineto{\pgfqpoint{1.101860in}{2.522000in}}%
\pgfpathlineto{\pgfqpoint{1.155213in}{3.648667in}}%
\pgfpathlineto{\pgfqpoint{1.208567in}{3.085333in}}%
\pgfpathlineto{\pgfqpoint{1.261921in}{3.648667in}}%
\pgfpathlineto{\pgfqpoint{1.315274in}{3.648667in}}%
\pgfpathlineto{\pgfqpoint{1.368628in}{3.648667in}}%
\pgfpathlineto{\pgfqpoint{1.421982in}{3.085333in}}%
\pgfpathlineto{\pgfqpoint{1.475335in}{3.085333in}}%
\pgfpathlineto{\pgfqpoint{1.528689in}{3.648667in}}%
\pgfpathlineto{\pgfqpoint{1.582043in}{2.522000in}}%
\pgfpathlineto{\pgfqpoint{1.635396in}{3.648667in}}%
\pgfpathlineto{\pgfqpoint{1.688750in}{3.085333in}}%
\pgfpathlineto{\pgfqpoint{1.742104in}{2.522000in}}%
\pgfpathlineto{\pgfqpoint{1.795457in}{3.648667in}}%
\pgfpathlineto{\pgfqpoint{1.848811in}{2.522000in}}%
\pgfpathlineto{\pgfqpoint{1.902165in}{3.648667in}}%
\pgfpathlineto{\pgfqpoint{1.955518in}{2.522000in}}%
\pgfpathlineto{\pgfqpoint{2.008872in}{2.522000in}}%
\pgfpathlineto{\pgfqpoint{2.062226in}{2.522000in}}%
\pgfpathlineto{\pgfqpoint{2.115579in}{2.522000in}}%
\pgfpathlineto{\pgfqpoint{2.168933in}{3.648667in}}%
\pgfpathlineto{\pgfqpoint{2.222287in}{3.648667in}}%
\pgfpathlineto{\pgfqpoint{2.275640in}{3.648667in}}%
\pgfpathlineto{\pgfqpoint{2.328994in}{3.648667in}}%
\pgfpathlineto{\pgfqpoint{2.382348in}{3.648667in}}%
\pgfpathlineto{\pgfqpoint{2.435701in}{3.648667in}}%
\pgfpathlineto{\pgfqpoint{2.489055in}{3.085333in}}%
\pgfusepath{stroke}%
\end{pgfscope}%
\begin{pgfscope}%
\pgfpathrectangle{\pgfqpoint{0.595000in}{2.488200in}}{\pgfqpoint{1.920732in}{1.216800in}}%
\pgfusepath{clip}%
\pgfsetbuttcap%
\pgfsetroundjoin%
\definecolor{currentfill}{rgb}{0.435294,0.576471,0.729412}%
\pgfsetfillcolor{currentfill}%
\pgfsetlinewidth{1.003750pt}%
\definecolor{currentstroke}{rgb}{0.435294,0.576471,0.729412}%
\pgfsetstrokecolor{currentstroke}%
\pgfsetdash{}{0pt}%
\pgfsys@defobject{currentmarker}{\pgfqpoint{-0.010417in}{-0.010417in}}{\pgfqpoint{0.010417in}{0.010417in}}{%
\pgfpathmoveto{\pgfqpoint{0.000000in}{-0.010417in}}%
\pgfpathcurveto{\pgfqpoint{0.002763in}{-0.010417in}}{\pgfqpoint{0.005412in}{-0.009319in}}{\pgfqpoint{0.007366in}{-0.007366in}}%
\pgfpathcurveto{\pgfqpoint{0.009319in}{-0.005412in}}{\pgfqpoint{0.010417in}{-0.002763in}}{\pgfqpoint{0.010417in}{0.000000in}}%
\pgfpathcurveto{\pgfqpoint{0.010417in}{0.002763in}}{\pgfqpoint{0.009319in}{0.005412in}}{\pgfqpoint{0.007366in}{0.007366in}}%
\pgfpathcurveto{\pgfqpoint{0.005412in}{0.009319in}}{\pgfqpoint{0.002763in}{0.010417in}}{\pgfqpoint{0.000000in}{0.010417in}}%
\pgfpathcurveto{\pgfqpoint{-0.002763in}{0.010417in}}{\pgfqpoint{-0.005412in}{0.009319in}}{\pgfqpoint{-0.007366in}{0.007366in}}%
\pgfpathcurveto{\pgfqpoint{-0.009319in}{0.005412in}}{\pgfqpoint{-0.010417in}{0.002763in}}{\pgfqpoint{-0.010417in}{0.000000in}}%
\pgfpathcurveto{\pgfqpoint{-0.010417in}{-0.002763in}}{\pgfqpoint{-0.009319in}{-0.005412in}}{\pgfqpoint{-0.007366in}{-0.007366in}}%
\pgfpathcurveto{\pgfqpoint{-0.005412in}{-0.009319in}}{\pgfqpoint{-0.002763in}{-0.010417in}}{\pgfqpoint{0.000000in}{-0.010417in}}%
\pgfpathlineto{\pgfqpoint{0.000000in}{-0.010417in}}%
\pgfpathclose%
\pgfusepath{stroke,fill}%
}%
\begin{pgfscope}%
\pgfsys@transformshift{0.621677in}{3.648667in}%
\pgfsys@useobject{currentmarker}{}%
\end{pgfscope}%
\begin{pgfscope}%
\pgfsys@transformshift{0.675030in}{3.648667in}%
\pgfsys@useobject{currentmarker}{}%
\end{pgfscope}%
\begin{pgfscope}%
\pgfsys@transformshift{0.728384in}{3.648667in}%
\pgfsys@useobject{currentmarker}{}%
\end{pgfscope}%
\begin{pgfscope}%
\pgfsys@transformshift{0.781738in}{3.648667in}%
\pgfsys@useobject{currentmarker}{}%
\end{pgfscope}%
\begin{pgfscope}%
\pgfsys@transformshift{0.835091in}{3.648667in}%
\pgfsys@useobject{currentmarker}{}%
\end{pgfscope}%
\begin{pgfscope}%
\pgfsys@transformshift{0.888445in}{3.648667in}%
\pgfsys@useobject{currentmarker}{}%
\end{pgfscope}%
\begin{pgfscope}%
\pgfsys@transformshift{0.941799in}{2.522000in}%
\pgfsys@useobject{currentmarker}{}%
\end{pgfscope}%
\begin{pgfscope}%
\pgfsys@transformshift{0.995152in}{3.648667in}%
\pgfsys@useobject{currentmarker}{}%
\end{pgfscope}%
\begin{pgfscope}%
\pgfsys@transformshift{1.048506in}{3.648667in}%
\pgfsys@useobject{currentmarker}{}%
\end{pgfscope}%
\begin{pgfscope}%
\pgfsys@transformshift{1.101860in}{2.522000in}%
\pgfsys@useobject{currentmarker}{}%
\end{pgfscope}%
\begin{pgfscope}%
\pgfsys@transformshift{1.155213in}{3.648667in}%
\pgfsys@useobject{currentmarker}{}%
\end{pgfscope}%
\begin{pgfscope}%
\pgfsys@transformshift{1.208567in}{3.085333in}%
\pgfsys@useobject{currentmarker}{}%
\end{pgfscope}%
\begin{pgfscope}%
\pgfsys@transformshift{1.261921in}{3.648667in}%
\pgfsys@useobject{currentmarker}{}%
\end{pgfscope}%
\begin{pgfscope}%
\pgfsys@transformshift{1.315274in}{3.648667in}%
\pgfsys@useobject{currentmarker}{}%
\end{pgfscope}%
\begin{pgfscope}%
\pgfsys@transformshift{1.368628in}{3.648667in}%
\pgfsys@useobject{currentmarker}{}%
\end{pgfscope}%
\begin{pgfscope}%
\pgfsys@transformshift{1.421982in}{3.085333in}%
\pgfsys@useobject{currentmarker}{}%
\end{pgfscope}%
\begin{pgfscope}%
\pgfsys@transformshift{1.475335in}{3.085333in}%
\pgfsys@useobject{currentmarker}{}%
\end{pgfscope}%
\begin{pgfscope}%
\pgfsys@transformshift{1.528689in}{3.648667in}%
\pgfsys@useobject{currentmarker}{}%
\end{pgfscope}%
\begin{pgfscope}%
\pgfsys@transformshift{1.582043in}{2.522000in}%
\pgfsys@useobject{currentmarker}{}%
\end{pgfscope}%
\begin{pgfscope}%
\pgfsys@transformshift{1.635396in}{3.648667in}%
\pgfsys@useobject{currentmarker}{}%
\end{pgfscope}%
\begin{pgfscope}%
\pgfsys@transformshift{1.688750in}{3.085333in}%
\pgfsys@useobject{currentmarker}{}%
\end{pgfscope}%
\begin{pgfscope}%
\pgfsys@transformshift{1.742104in}{2.522000in}%
\pgfsys@useobject{currentmarker}{}%
\end{pgfscope}%
\begin{pgfscope}%
\pgfsys@transformshift{1.795457in}{3.648667in}%
\pgfsys@useobject{currentmarker}{}%
\end{pgfscope}%
\begin{pgfscope}%
\pgfsys@transformshift{1.848811in}{2.522000in}%
\pgfsys@useobject{currentmarker}{}%
\end{pgfscope}%
\begin{pgfscope}%
\pgfsys@transformshift{1.902165in}{3.648667in}%
\pgfsys@useobject{currentmarker}{}%
\end{pgfscope}%
\begin{pgfscope}%
\pgfsys@transformshift{1.955518in}{2.522000in}%
\pgfsys@useobject{currentmarker}{}%
\end{pgfscope}%
\begin{pgfscope}%
\pgfsys@transformshift{2.008872in}{2.522000in}%
\pgfsys@useobject{currentmarker}{}%
\end{pgfscope}%
\begin{pgfscope}%
\pgfsys@transformshift{2.062226in}{2.522000in}%
\pgfsys@useobject{currentmarker}{}%
\end{pgfscope}%
\begin{pgfscope}%
\pgfsys@transformshift{2.115579in}{2.522000in}%
\pgfsys@useobject{currentmarker}{}%
\end{pgfscope}%
\begin{pgfscope}%
\pgfsys@transformshift{2.168933in}{3.648667in}%
\pgfsys@useobject{currentmarker}{}%
\end{pgfscope}%
\begin{pgfscope}%
\pgfsys@transformshift{2.222287in}{3.648667in}%
\pgfsys@useobject{currentmarker}{}%
\end{pgfscope}%
\begin{pgfscope}%
\pgfsys@transformshift{2.275640in}{3.648667in}%
\pgfsys@useobject{currentmarker}{}%
\end{pgfscope}%
\begin{pgfscope}%
\pgfsys@transformshift{2.328994in}{3.648667in}%
\pgfsys@useobject{currentmarker}{}%
\end{pgfscope}%
\begin{pgfscope}%
\pgfsys@transformshift{2.382348in}{3.648667in}%
\pgfsys@useobject{currentmarker}{}%
\end{pgfscope}%
\begin{pgfscope}%
\pgfsys@transformshift{2.435701in}{3.648667in}%
\pgfsys@useobject{currentmarker}{}%
\end{pgfscope}%
\begin{pgfscope}%
\pgfsys@transformshift{2.489055in}{3.085333in}%
\pgfsys@useobject{currentmarker}{}%
\end{pgfscope}%
\end{pgfscope}%
\begin{pgfscope}%
\pgfpathrectangle{\pgfqpoint{0.595000in}{2.488200in}}{\pgfqpoint{1.920732in}{1.216800in}}%
\pgfusepath{clip}%
\pgfsetrectcap%
\pgfsetroundjoin%
\pgfsetlinewidth{0.702625pt}%
\definecolor{currentstroke}{rgb}{0.156863,0.356863,0.568627}%
\pgfsetstrokecolor{currentstroke}%
\pgfsetdash{}{0pt}%
\pgfpathmoveto{\pgfqpoint{0.621677in}{3.648667in}}%
\pgfpathlineto{\pgfqpoint{0.675030in}{3.648667in}}%
\pgfpathlineto{\pgfqpoint{0.728384in}{3.648667in}}%
\pgfpathlineto{\pgfqpoint{0.781738in}{3.648667in}}%
\pgfpathlineto{\pgfqpoint{0.835091in}{3.648667in}}%
\pgfpathlineto{\pgfqpoint{0.888445in}{3.648667in}}%
\pgfpathlineto{\pgfqpoint{0.941799in}{2.522000in}}%
\pgfpathlineto{\pgfqpoint{0.995152in}{3.085333in}}%
\pgfpathlineto{\pgfqpoint{1.048506in}{3.648667in}}%
\pgfpathlineto{\pgfqpoint{1.101860in}{3.648667in}}%
\pgfpathlineto{\pgfqpoint{1.155213in}{3.648667in}}%
\pgfpathlineto{\pgfqpoint{1.208567in}{3.648667in}}%
\pgfpathlineto{\pgfqpoint{1.261921in}{3.648667in}}%
\pgfpathlineto{\pgfqpoint{1.315274in}{3.648667in}}%
\pgfpathlineto{\pgfqpoint{1.368628in}{3.648667in}}%
\pgfpathlineto{\pgfqpoint{1.421982in}{3.648667in}}%
\pgfpathlineto{\pgfqpoint{1.475335in}{3.648667in}}%
\pgfpathlineto{\pgfqpoint{1.528689in}{3.648667in}}%
\pgfpathlineto{\pgfqpoint{1.582043in}{3.648667in}}%
\pgfpathlineto{\pgfqpoint{1.635396in}{3.648667in}}%
\pgfpathlineto{\pgfqpoint{1.688750in}{2.522000in}}%
\pgfpathlineto{\pgfqpoint{1.742104in}{3.648667in}}%
\pgfpathlineto{\pgfqpoint{1.795457in}{2.522000in}}%
\pgfpathlineto{\pgfqpoint{1.848811in}{3.648667in}}%
\pgfpathlineto{\pgfqpoint{1.902165in}{3.648667in}}%
\pgfpathlineto{\pgfqpoint{1.955518in}{3.648667in}}%
\pgfpathlineto{\pgfqpoint{2.008872in}{3.648667in}}%
\pgfpathlineto{\pgfqpoint{2.062226in}{3.648667in}}%
\pgfpathlineto{\pgfqpoint{2.115579in}{3.648667in}}%
\pgfpathlineto{\pgfqpoint{2.168933in}{3.648667in}}%
\pgfpathlineto{\pgfqpoint{2.222287in}{3.648667in}}%
\pgfpathlineto{\pgfqpoint{2.275640in}{3.648667in}}%
\pgfpathlineto{\pgfqpoint{2.328994in}{2.522000in}}%
\pgfpathlineto{\pgfqpoint{2.382348in}{3.648667in}}%
\pgfpathlineto{\pgfqpoint{2.435701in}{3.085333in}}%
\pgfpathlineto{\pgfqpoint{2.489055in}{3.085333in}}%
\pgfusepath{stroke}%
\end{pgfscope}%
\begin{pgfscope}%
\pgfpathrectangle{\pgfqpoint{0.595000in}{2.488200in}}{\pgfqpoint{1.920732in}{1.216800in}}%
\pgfusepath{clip}%
\pgfsetbuttcap%
\pgfsetroundjoin%
\definecolor{currentfill}{rgb}{0.156863,0.356863,0.568627}%
\pgfsetfillcolor{currentfill}%
\pgfsetlinewidth{1.003750pt}%
\definecolor{currentstroke}{rgb}{0.156863,0.356863,0.568627}%
\pgfsetstrokecolor{currentstroke}%
\pgfsetdash{}{0pt}%
\pgfsys@defobject{currentmarker}{\pgfqpoint{-0.010417in}{-0.010417in}}{\pgfqpoint{0.010417in}{0.010417in}}{%
\pgfpathmoveto{\pgfqpoint{0.000000in}{-0.010417in}}%
\pgfpathcurveto{\pgfqpoint{0.002763in}{-0.010417in}}{\pgfqpoint{0.005412in}{-0.009319in}}{\pgfqpoint{0.007366in}{-0.007366in}}%
\pgfpathcurveto{\pgfqpoint{0.009319in}{-0.005412in}}{\pgfqpoint{0.010417in}{-0.002763in}}{\pgfqpoint{0.010417in}{0.000000in}}%
\pgfpathcurveto{\pgfqpoint{0.010417in}{0.002763in}}{\pgfqpoint{0.009319in}{0.005412in}}{\pgfqpoint{0.007366in}{0.007366in}}%
\pgfpathcurveto{\pgfqpoint{0.005412in}{0.009319in}}{\pgfqpoint{0.002763in}{0.010417in}}{\pgfqpoint{0.000000in}{0.010417in}}%
\pgfpathcurveto{\pgfqpoint{-0.002763in}{0.010417in}}{\pgfqpoint{-0.005412in}{0.009319in}}{\pgfqpoint{-0.007366in}{0.007366in}}%
\pgfpathcurveto{\pgfqpoint{-0.009319in}{0.005412in}}{\pgfqpoint{-0.010417in}{0.002763in}}{\pgfqpoint{-0.010417in}{0.000000in}}%
\pgfpathcurveto{\pgfqpoint{-0.010417in}{-0.002763in}}{\pgfqpoint{-0.009319in}{-0.005412in}}{\pgfqpoint{-0.007366in}{-0.007366in}}%
\pgfpathcurveto{\pgfqpoint{-0.005412in}{-0.009319in}}{\pgfqpoint{-0.002763in}{-0.010417in}}{\pgfqpoint{0.000000in}{-0.010417in}}%
\pgfpathlineto{\pgfqpoint{0.000000in}{-0.010417in}}%
\pgfpathclose%
\pgfusepath{stroke,fill}%
}%
\begin{pgfscope}%
\pgfsys@transformshift{0.621677in}{3.648667in}%
\pgfsys@useobject{currentmarker}{}%
\end{pgfscope}%
\begin{pgfscope}%
\pgfsys@transformshift{0.675030in}{3.648667in}%
\pgfsys@useobject{currentmarker}{}%
\end{pgfscope}%
\begin{pgfscope}%
\pgfsys@transformshift{0.728384in}{3.648667in}%
\pgfsys@useobject{currentmarker}{}%
\end{pgfscope}%
\begin{pgfscope}%
\pgfsys@transformshift{0.781738in}{3.648667in}%
\pgfsys@useobject{currentmarker}{}%
\end{pgfscope}%
\begin{pgfscope}%
\pgfsys@transformshift{0.835091in}{3.648667in}%
\pgfsys@useobject{currentmarker}{}%
\end{pgfscope}%
\begin{pgfscope}%
\pgfsys@transformshift{0.888445in}{3.648667in}%
\pgfsys@useobject{currentmarker}{}%
\end{pgfscope}%
\begin{pgfscope}%
\pgfsys@transformshift{0.941799in}{2.522000in}%
\pgfsys@useobject{currentmarker}{}%
\end{pgfscope}%
\begin{pgfscope}%
\pgfsys@transformshift{0.995152in}{3.085333in}%
\pgfsys@useobject{currentmarker}{}%
\end{pgfscope}%
\begin{pgfscope}%
\pgfsys@transformshift{1.048506in}{3.648667in}%
\pgfsys@useobject{currentmarker}{}%
\end{pgfscope}%
\begin{pgfscope}%
\pgfsys@transformshift{1.101860in}{3.648667in}%
\pgfsys@useobject{currentmarker}{}%
\end{pgfscope}%
\begin{pgfscope}%
\pgfsys@transformshift{1.155213in}{3.648667in}%
\pgfsys@useobject{currentmarker}{}%
\end{pgfscope}%
\begin{pgfscope}%
\pgfsys@transformshift{1.208567in}{3.648667in}%
\pgfsys@useobject{currentmarker}{}%
\end{pgfscope}%
\begin{pgfscope}%
\pgfsys@transformshift{1.261921in}{3.648667in}%
\pgfsys@useobject{currentmarker}{}%
\end{pgfscope}%
\begin{pgfscope}%
\pgfsys@transformshift{1.315274in}{3.648667in}%
\pgfsys@useobject{currentmarker}{}%
\end{pgfscope}%
\begin{pgfscope}%
\pgfsys@transformshift{1.368628in}{3.648667in}%
\pgfsys@useobject{currentmarker}{}%
\end{pgfscope}%
\begin{pgfscope}%
\pgfsys@transformshift{1.421982in}{3.648667in}%
\pgfsys@useobject{currentmarker}{}%
\end{pgfscope}%
\begin{pgfscope}%
\pgfsys@transformshift{1.475335in}{3.648667in}%
\pgfsys@useobject{currentmarker}{}%
\end{pgfscope}%
\begin{pgfscope}%
\pgfsys@transformshift{1.528689in}{3.648667in}%
\pgfsys@useobject{currentmarker}{}%
\end{pgfscope}%
\begin{pgfscope}%
\pgfsys@transformshift{1.582043in}{3.648667in}%
\pgfsys@useobject{currentmarker}{}%
\end{pgfscope}%
\begin{pgfscope}%
\pgfsys@transformshift{1.635396in}{3.648667in}%
\pgfsys@useobject{currentmarker}{}%
\end{pgfscope}%
\begin{pgfscope}%
\pgfsys@transformshift{1.688750in}{2.522000in}%
\pgfsys@useobject{currentmarker}{}%
\end{pgfscope}%
\begin{pgfscope}%
\pgfsys@transformshift{1.742104in}{3.648667in}%
\pgfsys@useobject{currentmarker}{}%
\end{pgfscope}%
\begin{pgfscope}%
\pgfsys@transformshift{1.795457in}{2.522000in}%
\pgfsys@useobject{currentmarker}{}%
\end{pgfscope}%
\begin{pgfscope}%
\pgfsys@transformshift{1.848811in}{3.648667in}%
\pgfsys@useobject{currentmarker}{}%
\end{pgfscope}%
\begin{pgfscope}%
\pgfsys@transformshift{1.902165in}{3.648667in}%
\pgfsys@useobject{currentmarker}{}%
\end{pgfscope}%
\begin{pgfscope}%
\pgfsys@transformshift{1.955518in}{3.648667in}%
\pgfsys@useobject{currentmarker}{}%
\end{pgfscope}%
\begin{pgfscope}%
\pgfsys@transformshift{2.008872in}{3.648667in}%
\pgfsys@useobject{currentmarker}{}%
\end{pgfscope}%
\begin{pgfscope}%
\pgfsys@transformshift{2.062226in}{3.648667in}%
\pgfsys@useobject{currentmarker}{}%
\end{pgfscope}%
\begin{pgfscope}%
\pgfsys@transformshift{2.115579in}{3.648667in}%
\pgfsys@useobject{currentmarker}{}%
\end{pgfscope}%
\begin{pgfscope}%
\pgfsys@transformshift{2.168933in}{3.648667in}%
\pgfsys@useobject{currentmarker}{}%
\end{pgfscope}%
\begin{pgfscope}%
\pgfsys@transformshift{2.222287in}{3.648667in}%
\pgfsys@useobject{currentmarker}{}%
\end{pgfscope}%
\begin{pgfscope}%
\pgfsys@transformshift{2.275640in}{3.648667in}%
\pgfsys@useobject{currentmarker}{}%
\end{pgfscope}%
\begin{pgfscope}%
\pgfsys@transformshift{2.328994in}{2.522000in}%
\pgfsys@useobject{currentmarker}{}%
\end{pgfscope}%
\begin{pgfscope}%
\pgfsys@transformshift{2.382348in}{3.648667in}%
\pgfsys@useobject{currentmarker}{}%
\end{pgfscope}%
\begin{pgfscope}%
\pgfsys@transformshift{2.435701in}{3.085333in}%
\pgfsys@useobject{currentmarker}{}%
\end{pgfscope}%
\begin{pgfscope}%
\pgfsys@transformshift{2.489055in}{3.085333in}%
\pgfsys@useobject{currentmarker}{}%
\end{pgfscope}%
\end{pgfscope}%
\begin{pgfscope}%
\pgfpathrectangle{\pgfqpoint{0.595000in}{2.488200in}}{\pgfqpoint{1.920732in}{1.216800in}}%
\pgfusepath{clip}%
\pgfsetrectcap%
\pgfsetroundjoin%
\pgfsetlinewidth{0.702625pt}%
\definecolor{currentstroke}{rgb}{0.062745,0.160784,0.290196}%
\pgfsetstrokecolor{currentstroke}%
\pgfsetdash{}{0pt}%
\pgfpathmoveto{\pgfqpoint{0.621677in}{3.648667in}}%
\pgfpathlineto{\pgfqpoint{0.675030in}{3.648667in}}%
\pgfpathlineto{\pgfqpoint{0.728384in}{3.648667in}}%
\pgfpathlineto{\pgfqpoint{0.781738in}{3.648667in}}%
\pgfpathlineto{\pgfqpoint{0.835091in}{3.648667in}}%
\pgfpathlineto{\pgfqpoint{0.888445in}{3.648667in}}%
\pgfpathlineto{\pgfqpoint{0.941799in}{2.522000in}}%
\pgfpathlineto{\pgfqpoint{0.995152in}{3.648667in}}%
\pgfpathlineto{\pgfqpoint{1.048506in}{3.648667in}}%
\pgfpathlineto{\pgfqpoint{1.101860in}{3.085333in}}%
\pgfpathlineto{\pgfqpoint{1.155213in}{3.648667in}}%
\pgfpathlineto{\pgfqpoint{1.208567in}{3.648667in}}%
\pgfpathlineto{\pgfqpoint{1.261921in}{3.648667in}}%
\pgfpathlineto{\pgfqpoint{1.315274in}{3.648667in}}%
\pgfpathlineto{\pgfqpoint{1.368628in}{3.648667in}}%
\pgfpathlineto{\pgfqpoint{1.421982in}{3.648667in}}%
\pgfpathlineto{\pgfqpoint{1.475335in}{3.648667in}}%
\pgfpathlineto{\pgfqpoint{1.528689in}{3.648667in}}%
\pgfpathlineto{\pgfqpoint{1.582043in}{3.085333in}}%
\pgfpathlineto{\pgfqpoint{1.635396in}{3.648667in}}%
\pgfpathlineto{\pgfqpoint{1.688750in}{2.522000in}}%
\pgfpathlineto{\pgfqpoint{1.742104in}{3.648667in}}%
\pgfpathlineto{\pgfqpoint{1.795457in}{2.522000in}}%
\pgfpathlineto{\pgfqpoint{1.848811in}{3.085333in}}%
\pgfpathlineto{\pgfqpoint{1.902165in}{3.085333in}}%
\pgfpathlineto{\pgfqpoint{1.955518in}{3.648667in}}%
\pgfpathlineto{\pgfqpoint{2.008872in}{3.648667in}}%
\pgfpathlineto{\pgfqpoint{2.062226in}{3.648667in}}%
\pgfpathlineto{\pgfqpoint{2.115579in}{3.648667in}}%
\pgfpathlineto{\pgfqpoint{2.168933in}{2.522000in}}%
\pgfpathlineto{\pgfqpoint{2.222287in}{3.648667in}}%
\pgfpathlineto{\pgfqpoint{2.275640in}{3.648667in}}%
\pgfpathlineto{\pgfqpoint{2.328994in}{3.648667in}}%
\pgfpathlineto{\pgfqpoint{2.382348in}{3.648667in}}%
\pgfpathlineto{\pgfqpoint{2.435701in}{3.648667in}}%
\pgfpathlineto{\pgfqpoint{2.489055in}{2.522000in}}%
\pgfusepath{stroke}%
\end{pgfscope}%
\begin{pgfscope}%
\pgfpathrectangle{\pgfqpoint{0.595000in}{2.488200in}}{\pgfqpoint{1.920732in}{1.216800in}}%
\pgfusepath{clip}%
\pgfsetbuttcap%
\pgfsetroundjoin%
\definecolor{currentfill}{rgb}{0.062745,0.160784,0.290196}%
\pgfsetfillcolor{currentfill}%
\pgfsetlinewidth{1.003750pt}%
\definecolor{currentstroke}{rgb}{0.062745,0.160784,0.290196}%
\pgfsetstrokecolor{currentstroke}%
\pgfsetdash{}{0pt}%
\pgfsys@defobject{currentmarker}{\pgfqpoint{-0.010417in}{-0.010417in}}{\pgfqpoint{0.010417in}{0.010417in}}{%
\pgfpathmoveto{\pgfqpoint{0.000000in}{-0.010417in}}%
\pgfpathcurveto{\pgfqpoint{0.002763in}{-0.010417in}}{\pgfqpoint{0.005412in}{-0.009319in}}{\pgfqpoint{0.007366in}{-0.007366in}}%
\pgfpathcurveto{\pgfqpoint{0.009319in}{-0.005412in}}{\pgfqpoint{0.010417in}{-0.002763in}}{\pgfqpoint{0.010417in}{0.000000in}}%
\pgfpathcurveto{\pgfqpoint{0.010417in}{0.002763in}}{\pgfqpoint{0.009319in}{0.005412in}}{\pgfqpoint{0.007366in}{0.007366in}}%
\pgfpathcurveto{\pgfqpoint{0.005412in}{0.009319in}}{\pgfqpoint{0.002763in}{0.010417in}}{\pgfqpoint{0.000000in}{0.010417in}}%
\pgfpathcurveto{\pgfqpoint{-0.002763in}{0.010417in}}{\pgfqpoint{-0.005412in}{0.009319in}}{\pgfqpoint{-0.007366in}{0.007366in}}%
\pgfpathcurveto{\pgfqpoint{-0.009319in}{0.005412in}}{\pgfqpoint{-0.010417in}{0.002763in}}{\pgfqpoint{-0.010417in}{0.000000in}}%
\pgfpathcurveto{\pgfqpoint{-0.010417in}{-0.002763in}}{\pgfqpoint{-0.009319in}{-0.005412in}}{\pgfqpoint{-0.007366in}{-0.007366in}}%
\pgfpathcurveto{\pgfqpoint{-0.005412in}{-0.009319in}}{\pgfqpoint{-0.002763in}{-0.010417in}}{\pgfqpoint{0.000000in}{-0.010417in}}%
\pgfpathlineto{\pgfqpoint{0.000000in}{-0.010417in}}%
\pgfpathclose%
\pgfusepath{stroke,fill}%
}%
\begin{pgfscope}%
\pgfsys@transformshift{0.621677in}{3.648667in}%
\pgfsys@useobject{currentmarker}{}%
\end{pgfscope}%
\begin{pgfscope}%
\pgfsys@transformshift{0.675030in}{3.648667in}%
\pgfsys@useobject{currentmarker}{}%
\end{pgfscope}%
\begin{pgfscope}%
\pgfsys@transformshift{0.728384in}{3.648667in}%
\pgfsys@useobject{currentmarker}{}%
\end{pgfscope}%
\begin{pgfscope}%
\pgfsys@transformshift{0.781738in}{3.648667in}%
\pgfsys@useobject{currentmarker}{}%
\end{pgfscope}%
\begin{pgfscope}%
\pgfsys@transformshift{0.835091in}{3.648667in}%
\pgfsys@useobject{currentmarker}{}%
\end{pgfscope}%
\begin{pgfscope}%
\pgfsys@transformshift{0.888445in}{3.648667in}%
\pgfsys@useobject{currentmarker}{}%
\end{pgfscope}%
\begin{pgfscope}%
\pgfsys@transformshift{0.941799in}{2.522000in}%
\pgfsys@useobject{currentmarker}{}%
\end{pgfscope}%
\begin{pgfscope}%
\pgfsys@transformshift{0.995152in}{3.648667in}%
\pgfsys@useobject{currentmarker}{}%
\end{pgfscope}%
\begin{pgfscope}%
\pgfsys@transformshift{1.048506in}{3.648667in}%
\pgfsys@useobject{currentmarker}{}%
\end{pgfscope}%
\begin{pgfscope}%
\pgfsys@transformshift{1.101860in}{3.085333in}%
\pgfsys@useobject{currentmarker}{}%
\end{pgfscope}%
\begin{pgfscope}%
\pgfsys@transformshift{1.155213in}{3.648667in}%
\pgfsys@useobject{currentmarker}{}%
\end{pgfscope}%
\begin{pgfscope}%
\pgfsys@transformshift{1.208567in}{3.648667in}%
\pgfsys@useobject{currentmarker}{}%
\end{pgfscope}%
\begin{pgfscope}%
\pgfsys@transformshift{1.261921in}{3.648667in}%
\pgfsys@useobject{currentmarker}{}%
\end{pgfscope}%
\begin{pgfscope}%
\pgfsys@transformshift{1.315274in}{3.648667in}%
\pgfsys@useobject{currentmarker}{}%
\end{pgfscope}%
\begin{pgfscope}%
\pgfsys@transformshift{1.368628in}{3.648667in}%
\pgfsys@useobject{currentmarker}{}%
\end{pgfscope}%
\begin{pgfscope}%
\pgfsys@transformshift{1.421982in}{3.648667in}%
\pgfsys@useobject{currentmarker}{}%
\end{pgfscope}%
\begin{pgfscope}%
\pgfsys@transformshift{1.475335in}{3.648667in}%
\pgfsys@useobject{currentmarker}{}%
\end{pgfscope}%
\begin{pgfscope}%
\pgfsys@transformshift{1.528689in}{3.648667in}%
\pgfsys@useobject{currentmarker}{}%
\end{pgfscope}%
\begin{pgfscope}%
\pgfsys@transformshift{1.582043in}{3.085333in}%
\pgfsys@useobject{currentmarker}{}%
\end{pgfscope}%
\begin{pgfscope}%
\pgfsys@transformshift{1.635396in}{3.648667in}%
\pgfsys@useobject{currentmarker}{}%
\end{pgfscope}%
\begin{pgfscope}%
\pgfsys@transformshift{1.688750in}{2.522000in}%
\pgfsys@useobject{currentmarker}{}%
\end{pgfscope}%
\begin{pgfscope}%
\pgfsys@transformshift{1.742104in}{3.648667in}%
\pgfsys@useobject{currentmarker}{}%
\end{pgfscope}%
\begin{pgfscope}%
\pgfsys@transformshift{1.795457in}{2.522000in}%
\pgfsys@useobject{currentmarker}{}%
\end{pgfscope}%
\begin{pgfscope}%
\pgfsys@transformshift{1.848811in}{3.085333in}%
\pgfsys@useobject{currentmarker}{}%
\end{pgfscope}%
\begin{pgfscope}%
\pgfsys@transformshift{1.902165in}{3.085333in}%
\pgfsys@useobject{currentmarker}{}%
\end{pgfscope}%
\begin{pgfscope}%
\pgfsys@transformshift{1.955518in}{3.648667in}%
\pgfsys@useobject{currentmarker}{}%
\end{pgfscope}%
\begin{pgfscope}%
\pgfsys@transformshift{2.008872in}{3.648667in}%
\pgfsys@useobject{currentmarker}{}%
\end{pgfscope}%
\begin{pgfscope}%
\pgfsys@transformshift{2.062226in}{3.648667in}%
\pgfsys@useobject{currentmarker}{}%
\end{pgfscope}%
\begin{pgfscope}%
\pgfsys@transformshift{2.115579in}{3.648667in}%
\pgfsys@useobject{currentmarker}{}%
\end{pgfscope}%
\begin{pgfscope}%
\pgfsys@transformshift{2.168933in}{2.522000in}%
\pgfsys@useobject{currentmarker}{}%
\end{pgfscope}%
\begin{pgfscope}%
\pgfsys@transformshift{2.222287in}{3.648667in}%
\pgfsys@useobject{currentmarker}{}%
\end{pgfscope}%
\begin{pgfscope}%
\pgfsys@transformshift{2.275640in}{3.648667in}%
\pgfsys@useobject{currentmarker}{}%
\end{pgfscope}%
\begin{pgfscope}%
\pgfsys@transformshift{2.328994in}{3.648667in}%
\pgfsys@useobject{currentmarker}{}%
\end{pgfscope}%
\begin{pgfscope}%
\pgfsys@transformshift{2.382348in}{3.648667in}%
\pgfsys@useobject{currentmarker}{}%
\end{pgfscope}%
\begin{pgfscope}%
\pgfsys@transformshift{2.435701in}{3.648667in}%
\pgfsys@useobject{currentmarker}{}%
\end{pgfscope}%
\begin{pgfscope}%
\pgfsys@transformshift{2.489055in}{2.522000in}%
\pgfsys@useobject{currentmarker}{}%
\end{pgfscope}%
\end{pgfscope}%
\begin{pgfscope}%
\pgfpathrectangle{\pgfqpoint{0.595000in}{2.488200in}}{\pgfqpoint{1.920732in}{1.216800in}}%
\pgfusepath{clip}%
\pgfsetrectcap%
\pgfsetroundjoin%
\pgfsetlinewidth{0.501875pt}%
\definecolor{currentstroke}{rgb}{0.101961,0.101961,0.101961}%
\pgfsetstrokecolor{currentstroke}%
\pgfsetdash{}{0pt}%
\pgfpathmoveto{\pgfqpoint{0.915122in}{2.488200in}}%
\pgfpathlineto{\pgfqpoint{0.915122in}{3.705000in}}%
\pgfusepath{stroke}%
\end{pgfscope}%
\begin{pgfscope}%
\pgfsetrectcap%
\pgfsetmiterjoin%
\pgfsetlinewidth{0.803000pt}%
\definecolor{currentstroke}{rgb}{0.101961,0.101961,0.101961}%
\pgfsetstrokecolor{currentstroke}%
\pgfsetdash{}{0pt}%
\pgfpathmoveto{\pgfqpoint{0.595000in}{2.488200in}}%
\pgfpathlineto{\pgfqpoint{0.595000in}{3.705000in}}%
\pgfusepath{stroke}%
\end{pgfscope}%
\begin{pgfscope}%
\pgfsetrectcap%
\pgfsetmiterjoin%
\pgfsetlinewidth{0.803000pt}%
\definecolor{currentstroke}{rgb}{0.101961,0.101961,0.101961}%
\pgfsetstrokecolor{currentstroke}%
\pgfsetdash{}{0pt}%
\pgfpathmoveto{\pgfqpoint{0.595000in}{2.488200in}}%
\pgfpathlineto{\pgfqpoint{2.515732in}{2.488200in}}%
\pgfusepath{stroke}%
\end{pgfscope}%
\begin{pgfscope}%
\pgfsetbuttcap%
\pgfsetmiterjoin%
\definecolor{currentfill}{rgb}{1.000000,1.000000,1.000000}%
\pgfsetfillcolor{currentfill}%
\pgfsetfillopacity{0.850000}%
\pgfsetlinewidth{0.000000pt}%
\definecolor{currentstroke}{rgb}{0.000000,0.000000,0.000000}%
\pgfsetstrokecolor{currentstroke}%
\pgfsetstrokeopacity{0.850000}%
\pgfsetdash{}{0pt}%
\pgfpathmoveto{\pgfqpoint{1.989810in}{2.523261in}}%
\pgfpathlineto{\pgfqpoint{2.490928in}{2.523261in}}%
\pgfpathlineto{\pgfqpoint{2.490928in}{2.636903in}}%
\pgfpathlineto{\pgfqpoint{1.989810in}{2.636903in}}%
\pgfpathlineto{\pgfqpoint{1.989810in}{2.523261in}}%
\pgfpathclose%
\pgfusepath{fill}%
\end{pgfscope}%
\begin{pgfscope}%
\definecolor{textcolor}{rgb}{0.101961,0.101961,0.101961}%
\pgfsetstrokecolor{textcolor}%
\pgfsetfillcolor{textcolor}%
\pgftext[x=2.477317in,y=2.536872in,right,bottom]{\color{textcolor}\sffamily\fontsize{7.000000}{8.400000}\selectfont Single-shot}%
\end{pgfscope}%
\begin{pgfscope}%
\definecolor{textcolor}{rgb}{0.466667,0.466667,0.466667}%
\pgfsetstrokecolor{textcolor}%
\pgfsetfillcolor{textcolor}%
\pgftext[x=0.755061in,y=3.682467in,,bottom]{\color{textcolor}\sffamily\fontsize{5.500000}{6.600000}\selectfont outage}%
\end{pgfscope}%
\begin{pgfscope}%
\pgfsetbuttcap%
\pgfsetmiterjoin%
\definecolor{currentfill}{rgb}{1.000000,1.000000,1.000000}%
\pgfsetfillcolor{currentfill}%
\pgfsetlinewidth{0.000000pt}%
\definecolor{currentstroke}{rgb}{0.000000,0.000000,0.000000}%
\pgfsetstrokecolor{currentstroke}%
\pgfsetstrokeopacity{0.000000}%
\pgfsetdash{}{0pt}%
\pgfpathmoveto{\pgfqpoint{2.784634in}{2.488200in}}%
\pgfpathlineto{\pgfqpoint{4.705366in}{2.488200in}}%
\pgfpathlineto{\pgfqpoint{4.705366in}{3.705000in}}%
\pgfpathlineto{\pgfqpoint{2.784634in}{3.705000in}}%
\pgfpathlineto{\pgfqpoint{2.784634in}{2.488200in}}%
\pgfpathclose%
\pgfusepath{fill}%
\end{pgfscope}%
\begin{pgfscope}%
\pgfpathrectangle{\pgfqpoint{2.784634in}{2.488200in}}{\pgfqpoint{1.920732in}{1.216800in}}%
\pgfusepath{clip}%
\pgfsetrectcap%
\pgfsetroundjoin%
\pgfsetlinewidth{0.401500pt}%
\definecolor{currentstroke}{rgb}{0.901961,0.901961,0.901961}%
\pgfsetstrokecolor{currentstroke}%
\pgfsetdash{}{0pt}%
\pgfpathmoveto{\pgfqpoint{2.784634in}{2.488200in}}%
\pgfpathlineto{\pgfqpoint{2.784634in}{3.705000in}}%
\pgfusepath{stroke}%
\end{pgfscope}%
\begin{pgfscope}%
\definecolor{textcolor}{rgb}{0.101961,0.101961,0.101961}%
\pgfsetstrokecolor{textcolor}%
\pgfsetfillcolor{textcolor}%
\pgftext[x=2.784634in,y=2.439589in,,top]{\color{textcolor}\sffamily\fontsize{7.000000}{8.400000}\selectfont \ensuremath{-}60}%
\end{pgfscope}%
\begin{pgfscope}%
\pgfpathrectangle{\pgfqpoint{2.784634in}{2.488200in}}{\pgfqpoint{1.920732in}{1.216800in}}%
\pgfusepath{clip}%
\pgfsetrectcap%
\pgfsetroundjoin%
\pgfsetlinewidth{0.401500pt}%
\definecolor{currentstroke}{rgb}{0.901961,0.901961,0.901961}%
\pgfsetstrokecolor{currentstroke}%
\pgfsetdash{}{0pt}%
\pgfpathmoveto{\pgfqpoint{3.104756in}{2.488200in}}%
\pgfpathlineto{\pgfqpoint{3.104756in}{3.705000in}}%
\pgfusepath{stroke}%
\end{pgfscope}%
\begin{pgfscope}%
\definecolor{textcolor}{rgb}{0.101961,0.101961,0.101961}%
\pgfsetstrokecolor{textcolor}%
\pgfsetfillcolor{textcolor}%
\pgftext[x=3.104756in,y=2.439589in,,top]{\color{textcolor}\sffamily\fontsize{7.000000}{8.400000}\selectfont 0}%
\end{pgfscope}%
\begin{pgfscope}%
\pgfpathrectangle{\pgfqpoint{2.784634in}{2.488200in}}{\pgfqpoint{1.920732in}{1.216800in}}%
\pgfusepath{clip}%
\pgfsetrectcap%
\pgfsetroundjoin%
\pgfsetlinewidth{0.401500pt}%
\definecolor{currentstroke}{rgb}{0.901961,0.901961,0.901961}%
\pgfsetstrokecolor{currentstroke}%
\pgfsetdash{}{0pt}%
\pgfpathmoveto{\pgfqpoint{3.638293in}{2.488200in}}%
\pgfpathlineto{\pgfqpoint{3.638293in}{3.705000in}}%
\pgfusepath{stroke}%
\end{pgfscope}%
\begin{pgfscope}%
\definecolor{textcolor}{rgb}{0.101961,0.101961,0.101961}%
\pgfsetstrokecolor{textcolor}%
\pgfsetfillcolor{textcolor}%
\pgftext[x=3.638293in,y=2.439589in,,top]{\color{textcolor}\sffamily\fontsize{7.000000}{8.400000}\selectfont 100}%
\end{pgfscope}%
\begin{pgfscope}%
\pgfpathrectangle{\pgfqpoint{2.784634in}{2.488200in}}{\pgfqpoint{1.920732in}{1.216800in}}%
\pgfusepath{clip}%
\pgfsetrectcap%
\pgfsetroundjoin%
\pgfsetlinewidth{0.401500pt}%
\definecolor{currentstroke}{rgb}{0.901961,0.901961,0.901961}%
\pgfsetstrokecolor{currentstroke}%
\pgfsetdash{}{0pt}%
\pgfpathmoveto{\pgfqpoint{4.171829in}{2.488200in}}%
\pgfpathlineto{\pgfqpoint{4.171829in}{3.705000in}}%
\pgfusepath{stroke}%
\end{pgfscope}%
\begin{pgfscope}%
\definecolor{textcolor}{rgb}{0.101961,0.101961,0.101961}%
\pgfsetstrokecolor{textcolor}%
\pgfsetfillcolor{textcolor}%
\pgftext[x=4.171829in,y=2.439589in,,top]{\color{textcolor}\sffamily\fontsize{7.000000}{8.400000}\selectfont 200}%
\end{pgfscope}%
\begin{pgfscope}%
\pgfpathrectangle{\pgfqpoint{2.784634in}{2.488200in}}{\pgfqpoint{1.920732in}{1.216800in}}%
\pgfusepath{clip}%
\pgfsetrectcap%
\pgfsetroundjoin%
\pgfsetlinewidth{0.401500pt}%
\definecolor{currentstroke}{rgb}{0.901961,0.901961,0.901961}%
\pgfsetstrokecolor{currentstroke}%
\pgfsetdash{}{0pt}%
\pgfpathmoveto{\pgfqpoint{4.705366in}{2.488200in}}%
\pgfpathlineto{\pgfqpoint{4.705366in}{3.705000in}}%
\pgfusepath{stroke}%
\end{pgfscope}%
\begin{pgfscope}%
\definecolor{textcolor}{rgb}{0.101961,0.101961,0.101961}%
\pgfsetstrokecolor{textcolor}%
\pgfsetfillcolor{textcolor}%
\pgftext[x=4.705366in,y=2.439589in,,top]{\color{textcolor}\sffamily\fontsize{7.000000}{8.400000}\selectfont 300}%
\end{pgfscope}%
\begin{pgfscope}%
\definecolor{textcolor}{rgb}{0.000000,0.000000,0.000000}%
\pgfsetstrokecolor{textcolor}%
\pgfsetfillcolor{textcolor}%
\pgftext[x=3.745000in,y=2.297614in,,top]{\color{textcolor}\sffamily\fontsize{7.500000}{9.000000}\selectfont s from reconnect R}%
\end{pgfscope}%
\begin{pgfscope}%
\pgfpathrectangle{\pgfqpoint{2.784634in}{2.488200in}}{\pgfqpoint{1.920732in}{1.216800in}}%
\pgfusepath{clip}%
\pgfsetrectcap%
\pgfsetroundjoin%
\pgfsetlinewidth{0.401500pt}%
\definecolor{currentstroke}{rgb}{0.901961,0.901961,0.901961}%
\pgfsetstrokecolor{currentstroke}%
\pgfsetdash{}{0pt}%
\pgfpathmoveto{\pgfqpoint{2.784634in}{2.522000in}}%
\pgfpathlineto{\pgfqpoint{4.705366in}{2.522000in}}%
\pgfusepath{stroke}%
\end{pgfscope}%
\begin{pgfscope}%
\pgfpathrectangle{\pgfqpoint{2.784634in}{2.488200in}}{\pgfqpoint{1.920732in}{1.216800in}}%
\pgfusepath{clip}%
\pgfsetrectcap%
\pgfsetroundjoin%
\pgfsetlinewidth{0.401500pt}%
\definecolor{currentstroke}{rgb}{0.901961,0.901961,0.901961}%
\pgfsetstrokecolor{currentstroke}%
\pgfsetdash{}{0pt}%
\pgfpathmoveto{\pgfqpoint{2.784634in}{2.747333in}}%
\pgfpathlineto{\pgfqpoint{4.705366in}{2.747333in}}%
\pgfusepath{stroke}%
\end{pgfscope}%
\begin{pgfscope}%
\pgfpathrectangle{\pgfqpoint{2.784634in}{2.488200in}}{\pgfqpoint{1.920732in}{1.216800in}}%
\pgfusepath{clip}%
\pgfsetrectcap%
\pgfsetroundjoin%
\pgfsetlinewidth{0.401500pt}%
\definecolor{currentstroke}{rgb}{0.901961,0.901961,0.901961}%
\pgfsetstrokecolor{currentstroke}%
\pgfsetdash{}{0pt}%
\pgfpathmoveto{\pgfqpoint{2.784634in}{2.972667in}}%
\pgfpathlineto{\pgfqpoint{4.705366in}{2.972667in}}%
\pgfusepath{stroke}%
\end{pgfscope}%
\begin{pgfscope}%
\pgfpathrectangle{\pgfqpoint{2.784634in}{2.488200in}}{\pgfqpoint{1.920732in}{1.216800in}}%
\pgfusepath{clip}%
\pgfsetrectcap%
\pgfsetroundjoin%
\pgfsetlinewidth{0.401500pt}%
\definecolor{currentstroke}{rgb}{0.901961,0.901961,0.901961}%
\pgfsetstrokecolor{currentstroke}%
\pgfsetdash{}{0pt}%
\pgfpathmoveto{\pgfqpoint{2.784634in}{3.198000in}}%
\pgfpathlineto{\pgfqpoint{4.705366in}{3.198000in}}%
\pgfusepath{stroke}%
\end{pgfscope}%
\begin{pgfscope}%
\pgfpathrectangle{\pgfqpoint{2.784634in}{2.488200in}}{\pgfqpoint{1.920732in}{1.216800in}}%
\pgfusepath{clip}%
\pgfsetrectcap%
\pgfsetroundjoin%
\pgfsetlinewidth{0.401500pt}%
\definecolor{currentstroke}{rgb}{0.901961,0.901961,0.901961}%
\pgfsetstrokecolor{currentstroke}%
\pgfsetdash{}{0pt}%
\pgfpathmoveto{\pgfqpoint{2.784634in}{3.423333in}}%
\pgfpathlineto{\pgfqpoint{4.705366in}{3.423333in}}%
\pgfusepath{stroke}%
\end{pgfscope}%
\begin{pgfscope}%
\pgfpathrectangle{\pgfqpoint{2.784634in}{2.488200in}}{\pgfqpoint{1.920732in}{1.216800in}}%
\pgfusepath{clip}%
\pgfsetrectcap%
\pgfsetroundjoin%
\pgfsetlinewidth{0.401500pt}%
\definecolor{currentstroke}{rgb}{0.901961,0.901961,0.901961}%
\pgfsetstrokecolor{currentstroke}%
\pgfsetdash{}{0pt}%
\pgfpathmoveto{\pgfqpoint{2.784634in}{3.648667in}}%
\pgfpathlineto{\pgfqpoint{4.705366in}{3.648667in}}%
\pgfusepath{stroke}%
\end{pgfscope}%
\begin{pgfscope}%
\pgfpathrectangle{\pgfqpoint{2.784634in}{2.488200in}}{\pgfqpoint{1.920732in}{1.216800in}}%
\pgfusepath{clip}%
\pgfsetbuttcap%
\pgfsetmiterjoin%
\definecolor{currentfill}{rgb}{0.929412,0.929412,0.929412}%
\pgfsetfillcolor{currentfill}%
\pgfsetlinewidth{0.000000pt}%
\definecolor{currentstroke}{rgb}{0.929412,0.929412,0.929412}%
\pgfsetstrokecolor{currentstroke}%
\pgfsetdash{}{0pt}%
\pgfpathmoveto{\pgfqpoint{2.784634in}{2.488200in}}%
\pgfpathlineto{\pgfqpoint{2.784634in}{3.705000in}}%
\pgfpathlineto{\pgfqpoint{3.104756in}{3.705000in}}%
\pgfpathlineto{\pgfqpoint{3.104756in}{2.488200in}}%
\pgfpathlineto{\pgfqpoint{2.784634in}{2.488200in}}%
\pgfpathclose%
\pgfusepath{fill}%
\end{pgfscope}%
\begin{pgfscope}%
\pgfpathrectangle{\pgfqpoint{2.784634in}{2.488200in}}{\pgfqpoint{1.920732in}{1.216800in}}%
\pgfusepath{clip}%
\pgfsetrectcap%
\pgfsetroundjoin%
\pgfsetlinewidth{0.702625pt}%
\definecolor{currentstroke}{rgb}{0.615686,0.705882,0.807843}%
\pgfsetstrokecolor{currentstroke}%
\pgfsetdash{}{0pt}%
\pgfpathmoveto{\pgfqpoint{2.811311in}{2.522000in}}%
\pgfpathlineto{\pgfqpoint{2.864665in}{2.522000in}}%
\pgfpathmoveto{\pgfqpoint{3.024726in}{2.522000in}}%
\pgfpathmoveto{\pgfqpoint{3.131433in}{2.522000in}}%
\pgfpathlineto{\pgfqpoint{3.184787in}{3.648667in}}%
\pgfpathlineto{\pgfqpoint{3.238140in}{3.648667in}}%
\pgfpathlineto{\pgfqpoint{3.291494in}{3.648667in}}%
\pgfpathlineto{\pgfqpoint{3.344848in}{3.648667in}}%
\pgfpathmoveto{\pgfqpoint{3.558262in}{3.648667in}}%
\pgfpathlineto{\pgfqpoint{3.611616in}{3.648667in}}%
\pgfpathlineto{\pgfqpoint{3.664970in}{3.648667in}}%
\pgfpathlineto{\pgfqpoint{3.718323in}{3.648667in}}%
\pgfpathlineto{\pgfqpoint{3.771677in}{3.648667in}}%
\pgfpathlineto{\pgfqpoint{3.825030in}{2.522000in}}%
\pgfpathlineto{\pgfqpoint{3.878384in}{3.085333in}}%
\pgfpathlineto{\pgfqpoint{3.931738in}{3.423333in}}%
\pgfpathlineto{\pgfqpoint{3.985091in}{3.273111in}}%
\pgfpathlineto{\pgfqpoint{4.038445in}{3.085333in}}%
\pgfpathlineto{\pgfqpoint{4.091799in}{2.897556in}}%
\pgfpathlineto{\pgfqpoint{4.145152in}{2.522000in}}%
\pgfpathlineto{\pgfqpoint{4.198506in}{2.747333in}}%
\pgfpathlineto{\pgfqpoint{4.251860in}{2.972667in}}%
\pgfpathlineto{\pgfqpoint{4.305213in}{3.273111in}}%
\pgfpathmoveto{\pgfqpoint{4.411921in}{3.648667in}}%
\pgfpathmoveto{\pgfqpoint{4.518628in}{3.648667in}}%
\pgfpathmoveto{\pgfqpoint{4.625335in}{3.648667in}}%
\pgfpathlineto{\pgfqpoint{4.678689in}{2.522000in}}%
\pgfusepath{stroke}%
\end{pgfscope}%
\begin{pgfscope}%
\pgfpathrectangle{\pgfqpoint{2.784634in}{2.488200in}}{\pgfqpoint{1.920732in}{1.216800in}}%
\pgfusepath{clip}%
\pgfsetbuttcap%
\pgfsetroundjoin%
\definecolor{currentfill}{rgb}{0.615686,0.705882,0.807843}%
\pgfsetfillcolor{currentfill}%
\pgfsetlinewidth{1.003750pt}%
\definecolor{currentstroke}{rgb}{0.615686,0.705882,0.807843}%
\pgfsetstrokecolor{currentstroke}%
\pgfsetdash{}{0pt}%
\pgfsys@defobject{currentmarker}{\pgfqpoint{-0.010417in}{-0.010417in}}{\pgfqpoint{0.010417in}{0.010417in}}{%
\pgfpathmoveto{\pgfqpoint{0.000000in}{-0.010417in}}%
\pgfpathcurveto{\pgfqpoint{0.002763in}{-0.010417in}}{\pgfqpoint{0.005412in}{-0.009319in}}{\pgfqpoint{0.007366in}{-0.007366in}}%
\pgfpathcurveto{\pgfqpoint{0.009319in}{-0.005412in}}{\pgfqpoint{0.010417in}{-0.002763in}}{\pgfqpoint{0.010417in}{0.000000in}}%
\pgfpathcurveto{\pgfqpoint{0.010417in}{0.002763in}}{\pgfqpoint{0.009319in}{0.005412in}}{\pgfqpoint{0.007366in}{0.007366in}}%
\pgfpathcurveto{\pgfqpoint{0.005412in}{0.009319in}}{\pgfqpoint{0.002763in}{0.010417in}}{\pgfqpoint{0.000000in}{0.010417in}}%
\pgfpathcurveto{\pgfqpoint{-0.002763in}{0.010417in}}{\pgfqpoint{-0.005412in}{0.009319in}}{\pgfqpoint{-0.007366in}{0.007366in}}%
\pgfpathcurveto{\pgfqpoint{-0.009319in}{0.005412in}}{\pgfqpoint{-0.010417in}{0.002763in}}{\pgfqpoint{-0.010417in}{0.000000in}}%
\pgfpathcurveto{\pgfqpoint{-0.010417in}{-0.002763in}}{\pgfqpoint{-0.009319in}{-0.005412in}}{\pgfqpoint{-0.007366in}{-0.007366in}}%
\pgfpathcurveto{\pgfqpoint{-0.005412in}{-0.009319in}}{\pgfqpoint{-0.002763in}{-0.010417in}}{\pgfqpoint{0.000000in}{-0.010417in}}%
\pgfpathlineto{\pgfqpoint{0.000000in}{-0.010417in}}%
\pgfpathclose%
\pgfusepath{stroke,fill}%
}%
\begin{pgfscope}%
\pgfsys@transformshift{2.811311in}{2.522000in}%
\pgfsys@useobject{currentmarker}{}%
\end{pgfscope}%
\begin{pgfscope}%
\pgfsys@transformshift{2.864665in}{2.522000in}%
\pgfsys@useobject{currentmarker}{}%
\end{pgfscope}%
\begin{pgfscope}%
\pgfsys@transformshift{3.024726in}{2.522000in}%
\pgfsys@useobject{currentmarker}{}%
\end{pgfscope}%
\begin{pgfscope}%
\pgfsys@transformshift{3.131433in}{2.522000in}%
\pgfsys@useobject{currentmarker}{}%
\end{pgfscope}%
\begin{pgfscope}%
\pgfsys@transformshift{3.184787in}{3.648667in}%
\pgfsys@useobject{currentmarker}{}%
\end{pgfscope}%
\begin{pgfscope}%
\pgfsys@transformshift{3.238140in}{3.648667in}%
\pgfsys@useobject{currentmarker}{}%
\end{pgfscope}%
\begin{pgfscope}%
\pgfsys@transformshift{3.291494in}{3.648667in}%
\pgfsys@useobject{currentmarker}{}%
\end{pgfscope}%
\begin{pgfscope}%
\pgfsys@transformshift{3.344848in}{3.648667in}%
\pgfsys@useobject{currentmarker}{}%
\end{pgfscope}%
\begin{pgfscope}%
\pgfsys@transformshift{3.558262in}{3.648667in}%
\pgfsys@useobject{currentmarker}{}%
\end{pgfscope}%
\begin{pgfscope}%
\pgfsys@transformshift{3.611616in}{3.648667in}%
\pgfsys@useobject{currentmarker}{}%
\end{pgfscope}%
\begin{pgfscope}%
\pgfsys@transformshift{3.664970in}{3.648667in}%
\pgfsys@useobject{currentmarker}{}%
\end{pgfscope}%
\begin{pgfscope}%
\pgfsys@transformshift{3.718323in}{3.648667in}%
\pgfsys@useobject{currentmarker}{}%
\end{pgfscope}%
\begin{pgfscope}%
\pgfsys@transformshift{3.771677in}{3.648667in}%
\pgfsys@useobject{currentmarker}{}%
\end{pgfscope}%
\begin{pgfscope}%
\pgfsys@transformshift{3.825030in}{2.522000in}%
\pgfsys@useobject{currentmarker}{}%
\end{pgfscope}%
\begin{pgfscope}%
\pgfsys@transformshift{3.878384in}{3.085333in}%
\pgfsys@useobject{currentmarker}{}%
\end{pgfscope}%
\begin{pgfscope}%
\pgfsys@transformshift{3.931738in}{3.423333in}%
\pgfsys@useobject{currentmarker}{}%
\end{pgfscope}%
\begin{pgfscope}%
\pgfsys@transformshift{3.985091in}{3.273111in}%
\pgfsys@useobject{currentmarker}{}%
\end{pgfscope}%
\begin{pgfscope}%
\pgfsys@transformshift{4.038445in}{3.085333in}%
\pgfsys@useobject{currentmarker}{}%
\end{pgfscope}%
\begin{pgfscope}%
\pgfsys@transformshift{4.091799in}{2.897556in}%
\pgfsys@useobject{currentmarker}{}%
\end{pgfscope}%
\begin{pgfscope}%
\pgfsys@transformshift{4.145152in}{2.522000in}%
\pgfsys@useobject{currentmarker}{}%
\end{pgfscope}%
\begin{pgfscope}%
\pgfsys@transformshift{4.198506in}{2.747333in}%
\pgfsys@useobject{currentmarker}{}%
\end{pgfscope}%
\begin{pgfscope}%
\pgfsys@transformshift{4.251860in}{2.972667in}%
\pgfsys@useobject{currentmarker}{}%
\end{pgfscope}%
\begin{pgfscope}%
\pgfsys@transformshift{4.305213in}{3.273111in}%
\pgfsys@useobject{currentmarker}{}%
\end{pgfscope}%
\begin{pgfscope}%
\pgfsys@transformshift{4.411921in}{3.648667in}%
\pgfsys@useobject{currentmarker}{}%
\end{pgfscope}%
\begin{pgfscope}%
\pgfsys@transformshift{4.518628in}{3.648667in}%
\pgfsys@useobject{currentmarker}{}%
\end{pgfscope}%
\begin{pgfscope}%
\pgfsys@transformshift{4.625335in}{3.648667in}%
\pgfsys@useobject{currentmarker}{}%
\end{pgfscope}%
\begin{pgfscope}%
\pgfsys@transformshift{4.678689in}{2.522000in}%
\pgfsys@useobject{currentmarker}{}%
\end{pgfscope}%
\end{pgfscope}%
\begin{pgfscope}%
\pgfpathrectangle{\pgfqpoint{2.784634in}{2.488200in}}{\pgfqpoint{1.920732in}{1.216800in}}%
\pgfusepath{clip}%
\pgfsetrectcap%
\pgfsetroundjoin%
\pgfsetlinewidth{0.702625pt}%
\definecolor{currentstroke}{rgb}{0.435294,0.576471,0.729412}%
\pgfsetstrokecolor{currentstroke}%
\pgfsetdash{}{0pt}%
\pgfpathmoveto{\pgfqpoint{2.811311in}{2.522000in}}%
\pgfpathlineto{\pgfqpoint{2.864665in}{2.522000in}}%
\pgfpathmoveto{\pgfqpoint{3.024726in}{2.522000in}}%
\pgfpathmoveto{\pgfqpoint{3.131433in}{2.972667in}}%
\pgfpathlineto{\pgfqpoint{3.184787in}{3.273111in}}%
\pgfpathlineto{\pgfqpoint{3.238140in}{3.648667in}}%
\pgfpathlineto{\pgfqpoint{3.291494in}{3.085333in}}%
\pgfpathlineto{\pgfqpoint{3.344848in}{3.648667in}}%
\pgfpathmoveto{\pgfqpoint{3.558262in}{2.972667in}}%
\pgfpathlineto{\pgfqpoint{3.611616in}{3.460889in}}%
\pgfpathlineto{\pgfqpoint{3.664970in}{3.648667in}}%
\pgfpathlineto{\pgfqpoint{3.718323in}{3.648667in}}%
\pgfpathlineto{\pgfqpoint{3.771677in}{3.648667in}}%
\pgfpathlineto{\pgfqpoint{3.825030in}{3.648667in}}%
\pgfpathlineto{\pgfqpoint{3.878384in}{2.522000in}}%
\pgfpathlineto{\pgfqpoint{3.931738in}{3.198000in}}%
\pgfpathlineto{\pgfqpoint{3.985091in}{2.897556in}}%
\pgfpathlineto{\pgfqpoint{4.038445in}{2.522000in}}%
\pgfpathlineto{\pgfqpoint{4.091799in}{3.460889in}}%
\pgfpathlineto{\pgfqpoint{4.145152in}{2.522000in}}%
\pgfpathlineto{\pgfqpoint{4.198506in}{2.522000in}}%
\pgfpathlineto{\pgfqpoint{4.251860in}{3.648667in}}%
\pgfpathlineto{\pgfqpoint{4.305213in}{2.897556in}}%
\pgfpathmoveto{\pgfqpoint{4.411921in}{3.085333in}}%
\pgfpathmoveto{\pgfqpoint{4.518628in}{3.648667in}}%
\pgfpathmoveto{\pgfqpoint{4.625335in}{2.522000in}}%
\pgfpathlineto{\pgfqpoint{4.678689in}{3.085333in}}%
\pgfusepath{stroke}%
\end{pgfscope}%
\begin{pgfscope}%
\pgfpathrectangle{\pgfqpoint{2.784634in}{2.488200in}}{\pgfqpoint{1.920732in}{1.216800in}}%
\pgfusepath{clip}%
\pgfsetbuttcap%
\pgfsetroundjoin%
\definecolor{currentfill}{rgb}{0.435294,0.576471,0.729412}%
\pgfsetfillcolor{currentfill}%
\pgfsetlinewidth{1.003750pt}%
\definecolor{currentstroke}{rgb}{0.435294,0.576471,0.729412}%
\pgfsetstrokecolor{currentstroke}%
\pgfsetdash{}{0pt}%
\pgfsys@defobject{currentmarker}{\pgfqpoint{-0.010417in}{-0.010417in}}{\pgfqpoint{0.010417in}{0.010417in}}{%
\pgfpathmoveto{\pgfqpoint{0.000000in}{-0.010417in}}%
\pgfpathcurveto{\pgfqpoint{0.002763in}{-0.010417in}}{\pgfqpoint{0.005412in}{-0.009319in}}{\pgfqpoint{0.007366in}{-0.007366in}}%
\pgfpathcurveto{\pgfqpoint{0.009319in}{-0.005412in}}{\pgfqpoint{0.010417in}{-0.002763in}}{\pgfqpoint{0.010417in}{0.000000in}}%
\pgfpathcurveto{\pgfqpoint{0.010417in}{0.002763in}}{\pgfqpoint{0.009319in}{0.005412in}}{\pgfqpoint{0.007366in}{0.007366in}}%
\pgfpathcurveto{\pgfqpoint{0.005412in}{0.009319in}}{\pgfqpoint{0.002763in}{0.010417in}}{\pgfqpoint{0.000000in}{0.010417in}}%
\pgfpathcurveto{\pgfqpoint{-0.002763in}{0.010417in}}{\pgfqpoint{-0.005412in}{0.009319in}}{\pgfqpoint{-0.007366in}{0.007366in}}%
\pgfpathcurveto{\pgfqpoint{-0.009319in}{0.005412in}}{\pgfqpoint{-0.010417in}{0.002763in}}{\pgfqpoint{-0.010417in}{0.000000in}}%
\pgfpathcurveto{\pgfqpoint{-0.010417in}{-0.002763in}}{\pgfqpoint{-0.009319in}{-0.005412in}}{\pgfqpoint{-0.007366in}{-0.007366in}}%
\pgfpathcurveto{\pgfqpoint{-0.005412in}{-0.009319in}}{\pgfqpoint{-0.002763in}{-0.010417in}}{\pgfqpoint{0.000000in}{-0.010417in}}%
\pgfpathlineto{\pgfqpoint{0.000000in}{-0.010417in}}%
\pgfpathclose%
\pgfusepath{stroke,fill}%
}%
\begin{pgfscope}%
\pgfsys@transformshift{2.811311in}{2.522000in}%
\pgfsys@useobject{currentmarker}{}%
\end{pgfscope}%
\begin{pgfscope}%
\pgfsys@transformshift{2.864665in}{2.522000in}%
\pgfsys@useobject{currentmarker}{}%
\end{pgfscope}%
\begin{pgfscope}%
\pgfsys@transformshift{3.024726in}{2.522000in}%
\pgfsys@useobject{currentmarker}{}%
\end{pgfscope}%
\begin{pgfscope}%
\pgfsys@transformshift{3.131433in}{2.972667in}%
\pgfsys@useobject{currentmarker}{}%
\end{pgfscope}%
\begin{pgfscope}%
\pgfsys@transformshift{3.184787in}{3.273111in}%
\pgfsys@useobject{currentmarker}{}%
\end{pgfscope}%
\begin{pgfscope}%
\pgfsys@transformshift{3.238140in}{3.648667in}%
\pgfsys@useobject{currentmarker}{}%
\end{pgfscope}%
\begin{pgfscope}%
\pgfsys@transformshift{3.291494in}{3.085333in}%
\pgfsys@useobject{currentmarker}{}%
\end{pgfscope}%
\begin{pgfscope}%
\pgfsys@transformshift{3.344848in}{3.648667in}%
\pgfsys@useobject{currentmarker}{}%
\end{pgfscope}%
\begin{pgfscope}%
\pgfsys@transformshift{3.558262in}{2.972667in}%
\pgfsys@useobject{currentmarker}{}%
\end{pgfscope}%
\begin{pgfscope}%
\pgfsys@transformshift{3.611616in}{3.460889in}%
\pgfsys@useobject{currentmarker}{}%
\end{pgfscope}%
\begin{pgfscope}%
\pgfsys@transformshift{3.664970in}{3.648667in}%
\pgfsys@useobject{currentmarker}{}%
\end{pgfscope}%
\begin{pgfscope}%
\pgfsys@transformshift{3.718323in}{3.648667in}%
\pgfsys@useobject{currentmarker}{}%
\end{pgfscope}%
\begin{pgfscope}%
\pgfsys@transformshift{3.771677in}{3.648667in}%
\pgfsys@useobject{currentmarker}{}%
\end{pgfscope}%
\begin{pgfscope}%
\pgfsys@transformshift{3.825030in}{3.648667in}%
\pgfsys@useobject{currentmarker}{}%
\end{pgfscope}%
\begin{pgfscope}%
\pgfsys@transformshift{3.878384in}{2.522000in}%
\pgfsys@useobject{currentmarker}{}%
\end{pgfscope}%
\begin{pgfscope}%
\pgfsys@transformshift{3.931738in}{3.198000in}%
\pgfsys@useobject{currentmarker}{}%
\end{pgfscope}%
\begin{pgfscope}%
\pgfsys@transformshift{3.985091in}{2.897556in}%
\pgfsys@useobject{currentmarker}{}%
\end{pgfscope}%
\begin{pgfscope}%
\pgfsys@transformshift{4.038445in}{2.522000in}%
\pgfsys@useobject{currentmarker}{}%
\end{pgfscope}%
\begin{pgfscope}%
\pgfsys@transformshift{4.091799in}{3.460889in}%
\pgfsys@useobject{currentmarker}{}%
\end{pgfscope}%
\begin{pgfscope}%
\pgfsys@transformshift{4.145152in}{2.522000in}%
\pgfsys@useobject{currentmarker}{}%
\end{pgfscope}%
\begin{pgfscope}%
\pgfsys@transformshift{4.198506in}{2.522000in}%
\pgfsys@useobject{currentmarker}{}%
\end{pgfscope}%
\begin{pgfscope}%
\pgfsys@transformshift{4.251860in}{3.648667in}%
\pgfsys@useobject{currentmarker}{}%
\end{pgfscope}%
\begin{pgfscope}%
\pgfsys@transformshift{4.305213in}{2.897556in}%
\pgfsys@useobject{currentmarker}{}%
\end{pgfscope}%
\begin{pgfscope}%
\pgfsys@transformshift{4.411921in}{3.085333in}%
\pgfsys@useobject{currentmarker}{}%
\end{pgfscope}%
\begin{pgfscope}%
\pgfsys@transformshift{4.518628in}{3.648667in}%
\pgfsys@useobject{currentmarker}{}%
\end{pgfscope}%
\begin{pgfscope}%
\pgfsys@transformshift{4.625335in}{2.522000in}%
\pgfsys@useobject{currentmarker}{}%
\end{pgfscope}%
\begin{pgfscope}%
\pgfsys@transformshift{4.678689in}{3.085333in}%
\pgfsys@useobject{currentmarker}{}%
\end{pgfscope}%
\end{pgfscope}%
\begin{pgfscope}%
\pgfpathrectangle{\pgfqpoint{2.784634in}{2.488200in}}{\pgfqpoint{1.920732in}{1.216800in}}%
\pgfusepath{clip}%
\pgfsetrectcap%
\pgfsetroundjoin%
\pgfsetlinewidth{0.702625pt}%
\definecolor{currentstroke}{rgb}{0.156863,0.356863,0.568627}%
\pgfsetstrokecolor{currentstroke}%
\pgfsetdash{}{0pt}%
\pgfpathmoveto{\pgfqpoint{2.811311in}{2.522000in}}%
\pgfpathlineto{\pgfqpoint{2.864665in}{2.522000in}}%
\pgfpathmoveto{\pgfqpoint{3.024726in}{2.522000in}}%
\pgfpathmoveto{\pgfqpoint{3.131433in}{2.522000in}}%
\pgfpathlineto{\pgfqpoint{3.184787in}{3.648667in}}%
\pgfpathlineto{\pgfqpoint{3.238140in}{3.648667in}}%
\pgfpathlineto{\pgfqpoint{3.291494in}{3.648667in}}%
\pgfpathlineto{\pgfqpoint{3.344848in}{3.648667in}}%
\pgfpathmoveto{\pgfqpoint{3.558262in}{3.648667in}}%
\pgfpathlineto{\pgfqpoint{3.611616in}{3.648667in}}%
\pgfpathlineto{\pgfqpoint{3.664970in}{3.648667in}}%
\pgfpathlineto{\pgfqpoint{3.718323in}{3.648667in}}%
\pgfpathlineto{\pgfqpoint{3.771677in}{3.648667in}}%
\pgfpathlineto{\pgfqpoint{3.825030in}{2.522000in}}%
\pgfpathlineto{\pgfqpoint{3.878384in}{3.085333in}}%
\pgfpathlineto{\pgfqpoint{3.931738in}{3.423333in}}%
\pgfpathlineto{\pgfqpoint{3.985091in}{3.648667in}}%
\pgfpathlineto{\pgfqpoint{4.038445in}{3.085333in}}%
\pgfpathlineto{\pgfqpoint{4.091799in}{3.273111in}}%
\pgfpathlineto{\pgfqpoint{4.145152in}{3.648667in}}%
\pgfpathlineto{\pgfqpoint{4.198506in}{3.648667in}}%
\pgfpathlineto{\pgfqpoint{4.251860in}{3.648667in}}%
\pgfpathlineto{\pgfqpoint{4.305213in}{3.648667in}}%
\pgfpathmoveto{\pgfqpoint{4.411921in}{3.648667in}}%
\pgfpathmoveto{\pgfqpoint{4.518628in}{2.522000in}}%
\pgfpathmoveto{\pgfqpoint{4.625335in}{3.085333in}}%
\pgfpathlineto{\pgfqpoint{4.678689in}{2.522000in}}%
\pgfusepath{stroke}%
\end{pgfscope}%
\begin{pgfscope}%
\pgfpathrectangle{\pgfqpoint{2.784634in}{2.488200in}}{\pgfqpoint{1.920732in}{1.216800in}}%
\pgfusepath{clip}%
\pgfsetbuttcap%
\pgfsetroundjoin%
\definecolor{currentfill}{rgb}{0.156863,0.356863,0.568627}%
\pgfsetfillcolor{currentfill}%
\pgfsetlinewidth{1.003750pt}%
\definecolor{currentstroke}{rgb}{0.156863,0.356863,0.568627}%
\pgfsetstrokecolor{currentstroke}%
\pgfsetdash{}{0pt}%
\pgfsys@defobject{currentmarker}{\pgfqpoint{-0.010417in}{-0.010417in}}{\pgfqpoint{0.010417in}{0.010417in}}{%
\pgfpathmoveto{\pgfqpoint{0.000000in}{-0.010417in}}%
\pgfpathcurveto{\pgfqpoint{0.002763in}{-0.010417in}}{\pgfqpoint{0.005412in}{-0.009319in}}{\pgfqpoint{0.007366in}{-0.007366in}}%
\pgfpathcurveto{\pgfqpoint{0.009319in}{-0.005412in}}{\pgfqpoint{0.010417in}{-0.002763in}}{\pgfqpoint{0.010417in}{0.000000in}}%
\pgfpathcurveto{\pgfqpoint{0.010417in}{0.002763in}}{\pgfqpoint{0.009319in}{0.005412in}}{\pgfqpoint{0.007366in}{0.007366in}}%
\pgfpathcurveto{\pgfqpoint{0.005412in}{0.009319in}}{\pgfqpoint{0.002763in}{0.010417in}}{\pgfqpoint{0.000000in}{0.010417in}}%
\pgfpathcurveto{\pgfqpoint{-0.002763in}{0.010417in}}{\pgfqpoint{-0.005412in}{0.009319in}}{\pgfqpoint{-0.007366in}{0.007366in}}%
\pgfpathcurveto{\pgfqpoint{-0.009319in}{0.005412in}}{\pgfqpoint{-0.010417in}{0.002763in}}{\pgfqpoint{-0.010417in}{0.000000in}}%
\pgfpathcurveto{\pgfqpoint{-0.010417in}{-0.002763in}}{\pgfqpoint{-0.009319in}{-0.005412in}}{\pgfqpoint{-0.007366in}{-0.007366in}}%
\pgfpathcurveto{\pgfqpoint{-0.005412in}{-0.009319in}}{\pgfqpoint{-0.002763in}{-0.010417in}}{\pgfqpoint{0.000000in}{-0.010417in}}%
\pgfpathlineto{\pgfqpoint{0.000000in}{-0.010417in}}%
\pgfpathclose%
\pgfusepath{stroke,fill}%
}%
\begin{pgfscope}%
\pgfsys@transformshift{2.811311in}{2.522000in}%
\pgfsys@useobject{currentmarker}{}%
\end{pgfscope}%
\begin{pgfscope}%
\pgfsys@transformshift{2.864665in}{2.522000in}%
\pgfsys@useobject{currentmarker}{}%
\end{pgfscope}%
\begin{pgfscope}%
\pgfsys@transformshift{3.024726in}{2.522000in}%
\pgfsys@useobject{currentmarker}{}%
\end{pgfscope}%
\begin{pgfscope}%
\pgfsys@transformshift{3.131433in}{2.522000in}%
\pgfsys@useobject{currentmarker}{}%
\end{pgfscope}%
\begin{pgfscope}%
\pgfsys@transformshift{3.184787in}{3.648667in}%
\pgfsys@useobject{currentmarker}{}%
\end{pgfscope}%
\begin{pgfscope}%
\pgfsys@transformshift{3.238140in}{3.648667in}%
\pgfsys@useobject{currentmarker}{}%
\end{pgfscope}%
\begin{pgfscope}%
\pgfsys@transformshift{3.291494in}{3.648667in}%
\pgfsys@useobject{currentmarker}{}%
\end{pgfscope}%
\begin{pgfscope}%
\pgfsys@transformshift{3.344848in}{3.648667in}%
\pgfsys@useobject{currentmarker}{}%
\end{pgfscope}%
\begin{pgfscope}%
\pgfsys@transformshift{3.558262in}{3.648667in}%
\pgfsys@useobject{currentmarker}{}%
\end{pgfscope}%
\begin{pgfscope}%
\pgfsys@transformshift{3.611616in}{3.648667in}%
\pgfsys@useobject{currentmarker}{}%
\end{pgfscope}%
\begin{pgfscope}%
\pgfsys@transformshift{3.664970in}{3.648667in}%
\pgfsys@useobject{currentmarker}{}%
\end{pgfscope}%
\begin{pgfscope}%
\pgfsys@transformshift{3.718323in}{3.648667in}%
\pgfsys@useobject{currentmarker}{}%
\end{pgfscope}%
\begin{pgfscope}%
\pgfsys@transformshift{3.771677in}{3.648667in}%
\pgfsys@useobject{currentmarker}{}%
\end{pgfscope}%
\begin{pgfscope}%
\pgfsys@transformshift{3.825030in}{2.522000in}%
\pgfsys@useobject{currentmarker}{}%
\end{pgfscope}%
\begin{pgfscope}%
\pgfsys@transformshift{3.878384in}{3.085333in}%
\pgfsys@useobject{currentmarker}{}%
\end{pgfscope}%
\begin{pgfscope}%
\pgfsys@transformshift{3.931738in}{3.423333in}%
\pgfsys@useobject{currentmarker}{}%
\end{pgfscope}%
\begin{pgfscope}%
\pgfsys@transformshift{3.985091in}{3.648667in}%
\pgfsys@useobject{currentmarker}{}%
\end{pgfscope}%
\begin{pgfscope}%
\pgfsys@transformshift{4.038445in}{3.085333in}%
\pgfsys@useobject{currentmarker}{}%
\end{pgfscope}%
\begin{pgfscope}%
\pgfsys@transformshift{4.091799in}{3.273111in}%
\pgfsys@useobject{currentmarker}{}%
\end{pgfscope}%
\begin{pgfscope}%
\pgfsys@transformshift{4.145152in}{3.648667in}%
\pgfsys@useobject{currentmarker}{}%
\end{pgfscope}%
\begin{pgfscope}%
\pgfsys@transformshift{4.198506in}{3.648667in}%
\pgfsys@useobject{currentmarker}{}%
\end{pgfscope}%
\begin{pgfscope}%
\pgfsys@transformshift{4.251860in}{3.648667in}%
\pgfsys@useobject{currentmarker}{}%
\end{pgfscope}%
\begin{pgfscope}%
\pgfsys@transformshift{4.305213in}{3.648667in}%
\pgfsys@useobject{currentmarker}{}%
\end{pgfscope}%
\begin{pgfscope}%
\pgfsys@transformshift{4.411921in}{3.648667in}%
\pgfsys@useobject{currentmarker}{}%
\end{pgfscope}%
\begin{pgfscope}%
\pgfsys@transformshift{4.518628in}{2.522000in}%
\pgfsys@useobject{currentmarker}{}%
\end{pgfscope}%
\begin{pgfscope}%
\pgfsys@transformshift{4.625335in}{3.085333in}%
\pgfsys@useobject{currentmarker}{}%
\end{pgfscope}%
\begin{pgfscope}%
\pgfsys@transformshift{4.678689in}{2.522000in}%
\pgfsys@useobject{currentmarker}{}%
\end{pgfscope}%
\end{pgfscope}%
\begin{pgfscope}%
\pgfpathrectangle{\pgfqpoint{2.784634in}{2.488200in}}{\pgfqpoint{1.920732in}{1.216800in}}%
\pgfusepath{clip}%
\pgfsetrectcap%
\pgfsetroundjoin%
\pgfsetlinewidth{0.702625pt}%
\definecolor{currentstroke}{rgb}{0.062745,0.160784,0.290196}%
\pgfsetstrokecolor{currentstroke}%
\pgfsetdash{}{0pt}%
\pgfpathmoveto{\pgfqpoint{2.811311in}{2.522000in}}%
\pgfpathlineto{\pgfqpoint{2.864665in}{2.522000in}}%
\pgfpathmoveto{\pgfqpoint{3.024726in}{2.522000in}}%
\pgfpathmoveto{\pgfqpoint{3.131433in}{2.972667in}}%
\pgfpathlineto{\pgfqpoint{3.184787in}{3.273111in}}%
\pgfpathlineto{\pgfqpoint{3.238140in}{3.648667in}}%
\pgfpathlineto{\pgfqpoint{3.291494in}{3.648667in}}%
\pgfpathlineto{\pgfqpoint{3.344848in}{3.648667in}}%
\pgfpathmoveto{\pgfqpoint{3.558262in}{3.648667in}}%
\pgfpathlineto{\pgfqpoint{3.611616in}{3.648667in}}%
\pgfpathlineto{\pgfqpoint{3.664970in}{3.648667in}}%
\pgfpathlineto{\pgfqpoint{3.718323in}{3.648667in}}%
\pgfpathlineto{\pgfqpoint{3.771677in}{3.648667in}}%
\pgfpathlineto{\pgfqpoint{3.825030in}{2.522000in}}%
\pgfpathlineto{\pgfqpoint{3.878384in}{3.085333in}}%
\pgfpathlineto{\pgfqpoint{3.931738in}{2.972667in}}%
\pgfpathlineto{\pgfqpoint{3.985091in}{2.897556in}}%
\pgfpathlineto{\pgfqpoint{4.038445in}{3.085333in}}%
\pgfpathlineto{\pgfqpoint{4.091799in}{3.085333in}}%
\pgfpathlineto{\pgfqpoint{4.145152in}{3.648667in}}%
\pgfpathlineto{\pgfqpoint{4.198506in}{3.198000in}}%
\pgfpathlineto{\pgfqpoint{4.251860in}{3.198000in}}%
\pgfpathlineto{\pgfqpoint{4.305213in}{2.897556in}}%
\pgfpathmoveto{\pgfqpoint{4.411921in}{3.085333in}}%
\pgfpathmoveto{\pgfqpoint{4.518628in}{3.648667in}}%
\pgfpathmoveto{\pgfqpoint{4.625335in}{3.648667in}}%
\pgfpathlineto{\pgfqpoint{4.678689in}{2.522000in}}%
\pgfusepath{stroke}%
\end{pgfscope}%
\begin{pgfscope}%
\pgfpathrectangle{\pgfqpoint{2.784634in}{2.488200in}}{\pgfqpoint{1.920732in}{1.216800in}}%
\pgfusepath{clip}%
\pgfsetbuttcap%
\pgfsetroundjoin%
\definecolor{currentfill}{rgb}{0.062745,0.160784,0.290196}%
\pgfsetfillcolor{currentfill}%
\pgfsetlinewidth{1.003750pt}%
\definecolor{currentstroke}{rgb}{0.062745,0.160784,0.290196}%
\pgfsetstrokecolor{currentstroke}%
\pgfsetdash{}{0pt}%
\pgfsys@defobject{currentmarker}{\pgfqpoint{-0.010417in}{-0.010417in}}{\pgfqpoint{0.010417in}{0.010417in}}{%
\pgfpathmoveto{\pgfqpoint{0.000000in}{-0.010417in}}%
\pgfpathcurveto{\pgfqpoint{0.002763in}{-0.010417in}}{\pgfqpoint{0.005412in}{-0.009319in}}{\pgfqpoint{0.007366in}{-0.007366in}}%
\pgfpathcurveto{\pgfqpoint{0.009319in}{-0.005412in}}{\pgfqpoint{0.010417in}{-0.002763in}}{\pgfqpoint{0.010417in}{0.000000in}}%
\pgfpathcurveto{\pgfqpoint{0.010417in}{0.002763in}}{\pgfqpoint{0.009319in}{0.005412in}}{\pgfqpoint{0.007366in}{0.007366in}}%
\pgfpathcurveto{\pgfqpoint{0.005412in}{0.009319in}}{\pgfqpoint{0.002763in}{0.010417in}}{\pgfqpoint{0.000000in}{0.010417in}}%
\pgfpathcurveto{\pgfqpoint{-0.002763in}{0.010417in}}{\pgfqpoint{-0.005412in}{0.009319in}}{\pgfqpoint{-0.007366in}{0.007366in}}%
\pgfpathcurveto{\pgfqpoint{-0.009319in}{0.005412in}}{\pgfqpoint{-0.010417in}{0.002763in}}{\pgfqpoint{-0.010417in}{0.000000in}}%
\pgfpathcurveto{\pgfqpoint{-0.010417in}{-0.002763in}}{\pgfqpoint{-0.009319in}{-0.005412in}}{\pgfqpoint{-0.007366in}{-0.007366in}}%
\pgfpathcurveto{\pgfqpoint{-0.005412in}{-0.009319in}}{\pgfqpoint{-0.002763in}{-0.010417in}}{\pgfqpoint{0.000000in}{-0.010417in}}%
\pgfpathlineto{\pgfqpoint{0.000000in}{-0.010417in}}%
\pgfpathclose%
\pgfusepath{stroke,fill}%
}%
\begin{pgfscope}%
\pgfsys@transformshift{2.811311in}{2.522000in}%
\pgfsys@useobject{currentmarker}{}%
\end{pgfscope}%
\begin{pgfscope}%
\pgfsys@transformshift{2.864665in}{2.522000in}%
\pgfsys@useobject{currentmarker}{}%
\end{pgfscope}%
\begin{pgfscope}%
\pgfsys@transformshift{3.024726in}{2.522000in}%
\pgfsys@useobject{currentmarker}{}%
\end{pgfscope}%
\begin{pgfscope}%
\pgfsys@transformshift{3.131433in}{2.972667in}%
\pgfsys@useobject{currentmarker}{}%
\end{pgfscope}%
\begin{pgfscope}%
\pgfsys@transformshift{3.184787in}{3.273111in}%
\pgfsys@useobject{currentmarker}{}%
\end{pgfscope}%
\begin{pgfscope}%
\pgfsys@transformshift{3.238140in}{3.648667in}%
\pgfsys@useobject{currentmarker}{}%
\end{pgfscope}%
\begin{pgfscope}%
\pgfsys@transformshift{3.291494in}{3.648667in}%
\pgfsys@useobject{currentmarker}{}%
\end{pgfscope}%
\begin{pgfscope}%
\pgfsys@transformshift{3.344848in}{3.648667in}%
\pgfsys@useobject{currentmarker}{}%
\end{pgfscope}%
\begin{pgfscope}%
\pgfsys@transformshift{3.558262in}{3.648667in}%
\pgfsys@useobject{currentmarker}{}%
\end{pgfscope}%
\begin{pgfscope}%
\pgfsys@transformshift{3.611616in}{3.648667in}%
\pgfsys@useobject{currentmarker}{}%
\end{pgfscope}%
\begin{pgfscope}%
\pgfsys@transformshift{3.664970in}{3.648667in}%
\pgfsys@useobject{currentmarker}{}%
\end{pgfscope}%
\begin{pgfscope}%
\pgfsys@transformshift{3.718323in}{3.648667in}%
\pgfsys@useobject{currentmarker}{}%
\end{pgfscope}%
\begin{pgfscope}%
\pgfsys@transformshift{3.771677in}{3.648667in}%
\pgfsys@useobject{currentmarker}{}%
\end{pgfscope}%
\begin{pgfscope}%
\pgfsys@transformshift{3.825030in}{2.522000in}%
\pgfsys@useobject{currentmarker}{}%
\end{pgfscope}%
\begin{pgfscope}%
\pgfsys@transformshift{3.878384in}{3.085333in}%
\pgfsys@useobject{currentmarker}{}%
\end{pgfscope}%
\begin{pgfscope}%
\pgfsys@transformshift{3.931738in}{2.972667in}%
\pgfsys@useobject{currentmarker}{}%
\end{pgfscope}%
\begin{pgfscope}%
\pgfsys@transformshift{3.985091in}{2.897556in}%
\pgfsys@useobject{currentmarker}{}%
\end{pgfscope}%
\begin{pgfscope}%
\pgfsys@transformshift{4.038445in}{3.085333in}%
\pgfsys@useobject{currentmarker}{}%
\end{pgfscope}%
\begin{pgfscope}%
\pgfsys@transformshift{4.091799in}{3.085333in}%
\pgfsys@useobject{currentmarker}{}%
\end{pgfscope}%
\begin{pgfscope}%
\pgfsys@transformshift{4.145152in}{3.648667in}%
\pgfsys@useobject{currentmarker}{}%
\end{pgfscope}%
\begin{pgfscope}%
\pgfsys@transformshift{4.198506in}{3.198000in}%
\pgfsys@useobject{currentmarker}{}%
\end{pgfscope}%
\begin{pgfscope}%
\pgfsys@transformshift{4.251860in}{3.198000in}%
\pgfsys@useobject{currentmarker}{}%
\end{pgfscope}%
\begin{pgfscope}%
\pgfsys@transformshift{4.305213in}{2.897556in}%
\pgfsys@useobject{currentmarker}{}%
\end{pgfscope}%
\begin{pgfscope}%
\pgfsys@transformshift{4.411921in}{3.085333in}%
\pgfsys@useobject{currentmarker}{}%
\end{pgfscope}%
\begin{pgfscope}%
\pgfsys@transformshift{4.518628in}{3.648667in}%
\pgfsys@useobject{currentmarker}{}%
\end{pgfscope}%
\begin{pgfscope}%
\pgfsys@transformshift{4.625335in}{3.648667in}%
\pgfsys@useobject{currentmarker}{}%
\end{pgfscope}%
\begin{pgfscope}%
\pgfsys@transformshift{4.678689in}{2.522000in}%
\pgfsys@useobject{currentmarker}{}%
\end{pgfscope}%
\end{pgfscope}%
\begin{pgfscope}%
\pgfpathrectangle{\pgfqpoint{2.784634in}{2.488200in}}{\pgfqpoint{1.920732in}{1.216800in}}%
\pgfusepath{clip}%
\pgfsetrectcap%
\pgfsetroundjoin%
\pgfsetlinewidth{0.501875pt}%
\definecolor{currentstroke}{rgb}{0.101961,0.101961,0.101961}%
\pgfsetstrokecolor{currentstroke}%
\pgfsetdash{}{0pt}%
\pgfpathmoveto{\pgfqpoint{3.104756in}{2.488200in}}%
\pgfpathlineto{\pgfqpoint{3.104756in}{3.705000in}}%
\pgfusepath{stroke}%
\end{pgfscope}%
\begin{pgfscope}%
\pgfsetrectcap%
\pgfsetmiterjoin%
\pgfsetlinewidth{0.803000pt}%
\definecolor{currentstroke}{rgb}{0.101961,0.101961,0.101961}%
\pgfsetstrokecolor{currentstroke}%
\pgfsetdash{}{0pt}%
\pgfpathmoveto{\pgfqpoint{2.784634in}{2.488200in}}%
\pgfpathlineto{\pgfqpoint{2.784634in}{3.705000in}}%
\pgfusepath{stroke}%
\end{pgfscope}%
\begin{pgfscope}%
\pgfsetrectcap%
\pgfsetmiterjoin%
\pgfsetlinewidth{0.803000pt}%
\definecolor{currentstroke}{rgb}{0.101961,0.101961,0.101961}%
\pgfsetstrokecolor{currentstroke}%
\pgfsetdash{}{0pt}%
\pgfpathmoveto{\pgfqpoint{2.784634in}{2.488200in}}%
\pgfpathlineto{\pgfqpoint{4.705366in}{2.488200in}}%
\pgfusepath{stroke}%
\end{pgfscope}%
\begin{pgfscope}%
\pgfsetbuttcap%
\pgfsetmiterjoin%
\definecolor{currentfill}{rgb}{1.000000,1.000000,1.000000}%
\pgfsetfillcolor{currentfill}%
\pgfsetfillopacity{0.850000}%
\pgfsetlinewidth{0.000000pt}%
\definecolor{currentstroke}{rgb}{0.000000,0.000000,0.000000}%
\pgfsetstrokecolor{currentstroke}%
\pgfsetstrokeopacity{0.850000}%
\pgfsetdash{}{0pt}%
\pgfpathmoveto{\pgfqpoint{3.962312in}{2.523261in}}%
\pgfpathlineto{\pgfqpoint{4.680562in}{2.523261in}}%
\pgfpathlineto{\pgfqpoint{4.680562in}{2.636903in}}%
\pgfpathlineto{\pgfqpoint{3.962312in}{2.636903in}}%
\pgfpathlineto{\pgfqpoint{3.962312in}{2.523261in}}%
\pgfpathclose%
\pgfusepath{fill}%
\end{pgfscope}%
\begin{pgfscope}%
\definecolor{textcolor}{rgb}{0.101961,0.101961,0.101961}%
\pgfsetstrokecolor{textcolor}%
\pgfsetfillcolor{textcolor}%
\pgftext[x=4.666951in,y=2.536872in,right,bottom]{\color{textcolor}\sffamily\fontsize{7.000000}{8.400000}\selectfont Robot follow-up}%
\end{pgfscope}%
\begin{pgfscope}%
\definecolor{textcolor}{rgb}{0.466667,0.466667,0.466667}%
\pgfsetstrokecolor{textcolor}%
\pgfsetfillcolor{textcolor}%
\pgftext[x=2.944695in,y=3.682467in,,bottom]{\color{textcolor}\sffamily\fontsize{5.500000}{6.600000}\selectfont outage}%
\end{pgfscope}%
\begin{pgfscope}%
\pgfsetbuttcap%
\pgfsetmiterjoin%
\definecolor{currentfill}{rgb}{1.000000,1.000000,1.000000}%
\pgfsetfillcolor{currentfill}%
\pgfsetlinewidth{0.000000pt}%
\definecolor{currentstroke}{rgb}{0.000000,0.000000,0.000000}%
\pgfsetstrokecolor{currentstroke}%
\pgfsetstrokeopacity{0.000000}%
\pgfsetdash{}{0pt}%
\pgfpathmoveto{\pgfqpoint{4.974268in}{2.488200in}}%
\pgfpathlineto{\pgfqpoint{6.895000in}{2.488200in}}%
\pgfpathlineto{\pgfqpoint{6.895000in}{3.705000in}}%
\pgfpathlineto{\pgfqpoint{4.974268in}{3.705000in}}%
\pgfpathlineto{\pgfqpoint{4.974268in}{2.488200in}}%
\pgfpathclose%
\pgfusepath{fill}%
\end{pgfscope}%
\begin{pgfscope}%
\pgfpathrectangle{\pgfqpoint{4.974268in}{2.488200in}}{\pgfqpoint{1.920732in}{1.216800in}}%
\pgfusepath{clip}%
\pgfsetrectcap%
\pgfsetroundjoin%
\pgfsetlinewidth{0.401500pt}%
\definecolor{currentstroke}{rgb}{0.901961,0.901961,0.901961}%
\pgfsetstrokecolor{currentstroke}%
\pgfsetdash{}{0pt}%
\pgfpathmoveto{\pgfqpoint{4.974268in}{2.488200in}}%
\pgfpathlineto{\pgfqpoint{4.974268in}{3.705000in}}%
\pgfusepath{stroke}%
\end{pgfscope}%
\begin{pgfscope}%
\definecolor{textcolor}{rgb}{0.101961,0.101961,0.101961}%
\pgfsetstrokecolor{textcolor}%
\pgfsetfillcolor{textcolor}%
\pgftext[x=4.974268in,y=2.439589in,,top]{\color{textcolor}\sffamily\fontsize{7.000000}{8.400000}\selectfont \ensuremath{-}60}%
\end{pgfscope}%
\begin{pgfscope}%
\pgfpathrectangle{\pgfqpoint{4.974268in}{2.488200in}}{\pgfqpoint{1.920732in}{1.216800in}}%
\pgfusepath{clip}%
\pgfsetrectcap%
\pgfsetroundjoin%
\pgfsetlinewidth{0.401500pt}%
\definecolor{currentstroke}{rgb}{0.901961,0.901961,0.901961}%
\pgfsetstrokecolor{currentstroke}%
\pgfsetdash{}{0pt}%
\pgfpathmoveto{\pgfqpoint{5.294390in}{2.488200in}}%
\pgfpathlineto{\pgfqpoint{5.294390in}{3.705000in}}%
\pgfusepath{stroke}%
\end{pgfscope}%
\begin{pgfscope}%
\definecolor{textcolor}{rgb}{0.101961,0.101961,0.101961}%
\pgfsetstrokecolor{textcolor}%
\pgfsetfillcolor{textcolor}%
\pgftext[x=5.294390in,y=2.439589in,,top]{\color{textcolor}\sffamily\fontsize{7.000000}{8.400000}\selectfont 0}%
\end{pgfscope}%
\begin{pgfscope}%
\pgfpathrectangle{\pgfqpoint{4.974268in}{2.488200in}}{\pgfqpoint{1.920732in}{1.216800in}}%
\pgfusepath{clip}%
\pgfsetrectcap%
\pgfsetroundjoin%
\pgfsetlinewidth{0.401500pt}%
\definecolor{currentstroke}{rgb}{0.901961,0.901961,0.901961}%
\pgfsetstrokecolor{currentstroke}%
\pgfsetdash{}{0pt}%
\pgfpathmoveto{\pgfqpoint{5.827927in}{2.488200in}}%
\pgfpathlineto{\pgfqpoint{5.827927in}{3.705000in}}%
\pgfusepath{stroke}%
\end{pgfscope}%
\begin{pgfscope}%
\definecolor{textcolor}{rgb}{0.101961,0.101961,0.101961}%
\pgfsetstrokecolor{textcolor}%
\pgfsetfillcolor{textcolor}%
\pgftext[x=5.827927in,y=2.439589in,,top]{\color{textcolor}\sffamily\fontsize{7.000000}{8.400000}\selectfont 100}%
\end{pgfscope}%
\begin{pgfscope}%
\pgfpathrectangle{\pgfqpoint{4.974268in}{2.488200in}}{\pgfqpoint{1.920732in}{1.216800in}}%
\pgfusepath{clip}%
\pgfsetrectcap%
\pgfsetroundjoin%
\pgfsetlinewidth{0.401500pt}%
\definecolor{currentstroke}{rgb}{0.901961,0.901961,0.901961}%
\pgfsetstrokecolor{currentstroke}%
\pgfsetdash{}{0pt}%
\pgfpathmoveto{\pgfqpoint{6.361463in}{2.488200in}}%
\pgfpathlineto{\pgfqpoint{6.361463in}{3.705000in}}%
\pgfusepath{stroke}%
\end{pgfscope}%
\begin{pgfscope}%
\definecolor{textcolor}{rgb}{0.101961,0.101961,0.101961}%
\pgfsetstrokecolor{textcolor}%
\pgfsetfillcolor{textcolor}%
\pgftext[x=6.361463in,y=2.439589in,,top]{\color{textcolor}\sffamily\fontsize{7.000000}{8.400000}\selectfont 200}%
\end{pgfscope}%
\begin{pgfscope}%
\pgfpathrectangle{\pgfqpoint{4.974268in}{2.488200in}}{\pgfqpoint{1.920732in}{1.216800in}}%
\pgfusepath{clip}%
\pgfsetrectcap%
\pgfsetroundjoin%
\pgfsetlinewidth{0.401500pt}%
\definecolor{currentstroke}{rgb}{0.901961,0.901961,0.901961}%
\pgfsetstrokecolor{currentstroke}%
\pgfsetdash{}{0pt}%
\pgfpathmoveto{\pgfqpoint{6.895000in}{2.488200in}}%
\pgfpathlineto{\pgfqpoint{6.895000in}{3.705000in}}%
\pgfusepath{stroke}%
\end{pgfscope}%
\begin{pgfscope}%
\definecolor{textcolor}{rgb}{0.101961,0.101961,0.101961}%
\pgfsetstrokecolor{textcolor}%
\pgfsetfillcolor{textcolor}%
\pgftext[x=6.895000in,y=2.439589in,,top]{\color{textcolor}\sffamily\fontsize{7.000000}{8.400000}\selectfont 300}%
\end{pgfscope}%
\begin{pgfscope}%
\definecolor{textcolor}{rgb}{0.000000,0.000000,0.000000}%
\pgfsetstrokecolor{textcolor}%
\pgfsetfillcolor{textcolor}%
\pgftext[x=5.934634in,y=2.297614in,,top]{\color{textcolor}\sffamily\fontsize{7.500000}{9.000000}\selectfont s from reconnect R}%
\end{pgfscope}%
\begin{pgfscope}%
\pgfpathrectangle{\pgfqpoint{4.974268in}{2.488200in}}{\pgfqpoint{1.920732in}{1.216800in}}%
\pgfusepath{clip}%
\pgfsetrectcap%
\pgfsetroundjoin%
\pgfsetlinewidth{0.401500pt}%
\definecolor{currentstroke}{rgb}{0.901961,0.901961,0.901961}%
\pgfsetstrokecolor{currentstroke}%
\pgfsetdash{}{0pt}%
\pgfpathmoveto{\pgfqpoint{4.974268in}{2.522000in}}%
\pgfpathlineto{\pgfqpoint{6.895000in}{2.522000in}}%
\pgfusepath{stroke}%
\end{pgfscope}%
\begin{pgfscope}%
\pgfpathrectangle{\pgfqpoint{4.974268in}{2.488200in}}{\pgfqpoint{1.920732in}{1.216800in}}%
\pgfusepath{clip}%
\pgfsetrectcap%
\pgfsetroundjoin%
\pgfsetlinewidth{0.401500pt}%
\definecolor{currentstroke}{rgb}{0.901961,0.901961,0.901961}%
\pgfsetstrokecolor{currentstroke}%
\pgfsetdash{}{0pt}%
\pgfpathmoveto{\pgfqpoint{4.974268in}{2.747333in}}%
\pgfpathlineto{\pgfqpoint{6.895000in}{2.747333in}}%
\pgfusepath{stroke}%
\end{pgfscope}%
\begin{pgfscope}%
\pgfpathrectangle{\pgfqpoint{4.974268in}{2.488200in}}{\pgfqpoint{1.920732in}{1.216800in}}%
\pgfusepath{clip}%
\pgfsetrectcap%
\pgfsetroundjoin%
\pgfsetlinewidth{0.401500pt}%
\definecolor{currentstroke}{rgb}{0.901961,0.901961,0.901961}%
\pgfsetstrokecolor{currentstroke}%
\pgfsetdash{}{0pt}%
\pgfpathmoveto{\pgfqpoint{4.974268in}{2.972667in}}%
\pgfpathlineto{\pgfqpoint{6.895000in}{2.972667in}}%
\pgfusepath{stroke}%
\end{pgfscope}%
\begin{pgfscope}%
\pgfpathrectangle{\pgfqpoint{4.974268in}{2.488200in}}{\pgfqpoint{1.920732in}{1.216800in}}%
\pgfusepath{clip}%
\pgfsetrectcap%
\pgfsetroundjoin%
\pgfsetlinewidth{0.401500pt}%
\definecolor{currentstroke}{rgb}{0.901961,0.901961,0.901961}%
\pgfsetstrokecolor{currentstroke}%
\pgfsetdash{}{0pt}%
\pgfpathmoveto{\pgfqpoint{4.974268in}{3.198000in}}%
\pgfpathlineto{\pgfqpoint{6.895000in}{3.198000in}}%
\pgfusepath{stroke}%
\end{pgfscope}%
\begin{pgfscope}%
\pgfpathrectangle{\pgfqpoint{4.974268in}{2.488200in}}{\pgfqpoint{1.920732in}{1.216800in}}%
\pgfusepath{clip}%
\pgfsetrectcap%
\pgfsetroundjoin%
\pgfsetlinewidth{0.401500pt}%
\definecolor{currentstroke}{rgb}{0.901961,0.901961,0.901961}%
\pgfsetstrokecolor{currentstroke}%
\pgfsetdash{}{0pt}%
\pgfpathmoveto{\pgfqpoint{4.974268in}{3.423333in}}%
\pgfpathlineto{\pgfqpoint{6.895000in}{3.423333in}}%
\pgfusepath{stroke}%
\end{pgfscope}%
\begin{pgfscope}%
\pgfpathrectangle{\pgfqpoint{4.974268in}{2.488200in}}{\pgfqpoint{1.920732in}{1.216800in}}%
\pgfusepath{clip}%
\pgfsetrectcap%
\pgfsetroundjoin%
\pgfsetlinewidth{0.401500pt}%
\definecolor{currentstroke}{rgb}{0.901961,0.901961,0.901961}%
\pgfsetstrokecolor{currentstroke}%
\pgfsetdash{}{0pt}%
\pgfpathmoveto{\pgfqpoint{4.974268in}{3.648667in}}%
\pgfpathlineto{\pgfqpoint{6.895000in}{3.648667in}}%
\pgfusepath{stroke}%
\end{pgfscope}%
\begin{pgfscope}%
\pgfpathrectangle{\pgfqpoint{4.974268in}{2.488200in}}{\pgfqpoint{1.920732in}{1.216800in}}%
\pgfusepath{clip}%
\pgfsetbuttcap%
\pgfsetmiterjoin%
\definecolor{currentfill}{rgb}{0.929412,0.929412,0.929412}%
\pgfsetfillcolor{currentfill}%
\pgfsetlinewidth{0.000000pt}%
\definecolor{currentstroke}{rgb}{0.929412,0.929412,0.929412}%
\pgfsetstrokecolor{currentstroke}%
\pgfsetdash{}{0pt}%
\pgfpathmoveto{\pgfqpoint{4.974268in}{2.488200in}}%
\pgfpathlineto{\pgfqpoint{4.974268in}{3.705000in}}%
\pgfpathlineto{\pgfqpoint{5.294390in}{3.705000in}}%
\pgfpathlineto{\pgfqpoint{5.294390in}{2.488200in}}%
\pgfpathlineto{\pgfqpoint{4.974268in}{2.488200in}}%
\pgfpathclose%
\pgfusepath{fill}%
\end{pgfscope}%
\begin{pgfscope}%
\pgfpathrectangle{\pgfqpoint{4.974268in}{2.488200in}}{\pgfqpoint{1.920732in}{1.216800in}}%
\pgfusepath{clip}%
\pgfsetrectcap%
\pgfsetroundjoin%
\pgfsetlinewidth{0.702625pt}%
\definecolor{currentstroke}{rgb}{0.615686,0.705882,0.807843}%
\pgfsetstrokecolor{currentstroke}%
\pgfsetdash{}{0pt}%
\pgfpathmoveto{\pgfqpoint{5.000945in}{3.648667in}}%
\pgfpathlineto{\pgfqpoint{5.054299in}{3.648667in}}%
\pgfpathlineto{\pgfqpoint{5.107652in}{3.648667in}}%
\pgfpathlineto{\pgfqpoint{5.161006in}{3.648667in}}%
\pgfpathlineto{\pgfqpoint{5.214360in}{3.648667in}}%
\pgfpathlineto{\pgfqpoint{5.267713in}{3.648667in}}%
\pgfpathlineto{\pgfqpoint{5.321067in}{2.522000in}}%
\pgfpathlineto{\pgfqpoint{5.374421in}{2.788303in}}%
\pgfpathlineto{\pgfqpoint{5.427774in}{3.492185in}}%
\pgfpathlineto{\pgfqpoint{5.481128in}{3.007093in}}%
\pgfpathlineto{\pgfqpoint{5.534482in}{3.581403in}}%
\pgfpathlineto{\pgfqpoint{5.587835in}{3.648667in}}%
\pgfpathlineto{\pgfqpoint{5.641189in}{3.648667in}}%
\pgfpathlineto{\pgfqpoint{5.694543in}{3.648667in}}%
\pgfpathlineto{\pgfqpoint{5.747896in}{3.093741in}}%
\pgfpathlineto{\pgfqpoint{5.801250in}{2.980305in}}%
\pgfpathlineto{\pgfqpoint{5.854604in}{3.584286in}}%
\pgfpathlineto{\pgfqpoint{5.907957in}{3.648667in}}%
\pgfpathlineto{\pgfqpoint{5.961311in}{3.563313in}}%
\pgfpathlineto{\pgfqpoint{6.014665in}{3.318437in}}%
\pgfpathlineto{\pgfqpoint{6.068018in}{2.636300in}}%
\pgfpathlineto{\pgfqpoint{6.121372in}{2.572448in}}%
\pgfpathlineto{\pgfqpoint{6.174726in}{2.572448in}}%
\pgfpathlineto{\pgfqpoint{6.228079in}{2.991444in}}%
\pgfpathlineto{\pgfqpoint{6.281433in}{2.654549in}}%
\pgfpathlineto{\pgfqpoint{6.334787in}{2.819315in}}%
\pgfpathlineto{\pgfqpoint{6.388140in}{2.925582in}}%
\pgfpathlineto{\pgfqpoint{6.441494in}{2.860000in}}%
\pgfpathlineto{\pgfqpoint{6.494848in}{3.436396in}}%
\pgfpathlineto{\pgfqpoint{6.548201in}{3.648667in}}%
\pgfpathlineto{\pgfqpoint{6.601555in}{3.367000in}}%
\pgfpathlineto{\pgfqpoint{6.654909in}{3.063667in}}%
\pgfpathlineto{\pgfqpoint{6.708262in}{3.300725in}}%
\pgfpathlineto{\pgfqpoint{6.761616in}{3.648667in}}%
\pgfpathlineto{\pgfqpoint{6.814970in}{3.257463in}}%
\pgfpathlineto{\pgfqpoint{6.868323in}{2.602476in}}%
\pgfusepath{stroke}%
\end{pgfscope}%
\begin{pgfscope}%
\pgfpathrectangle{\pgfqpoint{4.974268in}{2.488200in}}{\pgfqpoint{1.920732in}{1.216800in}}%
\pgfusepath{clip}%
\pgfsetbuttcap%
\pgfsetroundjoin%
\definecolor{currentfill}{rgb}{0.615686,0.705882,0.807843}%
\pgfsetfillcolor{currentfill}%
\pgfsetlinewidth{1.003750pt}%
\definecolor{currentstroke}{rgb}{0.615686,0.705882,0.807843}%
\pgfsetstrokecolor{currentstroke}%
\pgfsetdash{}{0pt}%
\pgfsys@defobject{currentmarker}{\pgfqpoint{-0.010417in}{-0.010417in}}{\pgfqpoint{0.010417in}{0.010417in}}{%
\pgfpathmoveto{\pgfqpoint{0.000000in}{-0.010417in}}%
\pgfpathcurveto{\pgfqpoint{0.002763in}{-0.010417in}}{\pgfqpoint{0.005412in}{-0.009319in}}{\pgfqpoint{0.007366in}{-0.007366in}}%
\pgfpathcurveto{\pgfqpoint{0.009319in}{-0.005412in}}{\pgfqpoint{0.010417in}{-0.002763in}}{\pgfqpoint{0.010417in}{0.000000in}}%
\pgfpathcurveto{\pgfqpoint{0.010417in}{0.002763in}}{\pgfqpoint{0.009319in}{0.005412in}}{\pgfqpoint{0.007366in}{0.007366in}}%
\pgfpathcurveto{\pgfqpoint{0.005412in}{0.009319in}}{\pgfqpoint{0.002763in}{0.010417in}}{\pgfqpoint{0.000000in}{0.010417in}}%
\pgfpathcurveto{\pgfqpoint{-0.002763in}{0.010417in}}{\pgfqpoint{-0.005412in}{0.009319in}}{\pgfqpoint{-0.007366in}{0.007366in}}%
\pgfpathcurveto{\pgfqpoint{-0.009319in}{0.005412in}}{\pgfqpoint{-0.010417in}{0.002763in}}{\pgfqpoint{-0.010417in}{0.000000in}}%
\pgfpathcurveto{\pgfqpoint{-0.010417in}{-0.002763in}}{\pgfqpoint{-0.009319in}{-0.005412in}}{\pgfqpoint{-0.007366in}{-0.007366in}}%
\pgfpathcurveto{\pgfqpoint{-0.005412in}{-0.009319in}}{\pgfqpoint{-0.002763in}{-0.010417in}}{\pgfqpoint{0.000000in}{-0.010417in}}%
\pgfpathlineto{\pgfqpoint{0.000000in}{-0.010417in}}%
\pgfpathclose%
\pgfusepath{stroke,fill}%
}%
\begin{pgfscope}%
\pgfsys@transformshift{5.000945in}{3.648667in}%
\pgfsys@useobject{currentmarker}{}%
\end{pgfscope}%
\begin{pgfscope}%
\pgfsys@transformshift{5.054299in}{3.648667in}%
\pgfsys@useobject{currentmarker}{}%
\end{pgfscope}%
\begin{pgfscope}%
\pgfsys@transformshift{5.107652in}{3.648667in}%
\pgfsys@useobject{currentmarker}{}%
\end{pgfscope}%
\begin{pgfscope}%
\pgfsys@transformshift{5.161006in}{3.648667in}%
\pgfsys@useobject{currentmarker}{}%
\end{pgfscope}%
\begin{pgfscope}%
\pgfsys@transformshift{5.214360in}{3.648667in}%
\pgfsys@useobject{currentmarker}{}%
\end{pgfscope}%
\begin{pgfscope}%
\pgfsys@transformshift{5.267713in}{3.648667in}%
\pgfsys@useobject{currentmarker}{}%
\end{pgfscope}%
\begin{pgfscope}%
\pgfsys@transformshift{5.321067in}{2.522000in}%
\pgfsys@useobject{currentmarker}{}%
\end{pgfscope}%
\begin{pgfscope}%
\pgfsys@transformshift{5.374421in}{2.788303in}%
\pgfsys@useobject{currentmarker}{}%
\end{pgfscope}%
\begin{pgfscope}%
\pgfsys@transformshift{5.427774in}{3.492185in}%
\pgfsys@useobject{currentmarker}{}%
\end{pgfscope}%
\begin{pgfscope}%
\pgfsys@transformshift{5.481128in}{3.007093in}%
\pgfsys@useobject{currentmarker}{}%
\end{pgfscope}%
\begin{pgfscope}%
\pgfsys@transformshift{5.534482in}{3.581403in}%
\pgfsys@useobject{currentmarker}{}%
\end{pgfscope}%
\begin{pgfscope}%
\pgfsys@transformshift{5.587835in}{3.648667in}%
\pgfsys@useobject{currentmarker}{}%
\end{pgfscope}%
\begin{pgfscope}%
\pgfsys@transformshift{5.641189in}{3.648667in}%
\pgfsys@useobject{currentmarker}{}%
\end{pgfscope}%
\begin{pgfscope}%
\pgfsys@transformshift{5.694543in}{3.648667in}%
\pgfsys@useobject{currentmarker}{}%
\end{pgfscope}%
\begin{pgfscope}%
\pgfsys@transformshift{5.747896in}{3.093741in}%
\pgfsys@useobject{currentmarker}{}%
\end{pgfscope}%
\begin{pgfscope}%
\pgfsys@transformshift{5.801250in}{2.980305in}%
\pgfsys@useobject{currentmarker}{}%
\end{pgfscope}%
\begin{pgfscope}%
\pgfsys@transformshift{5.854604in}{3.584286in}%
\pgfsys@useobject{currentmarker}{}%
\end{pgfscope}%
\begin{pgfscope}%
\pgfsys@transformshift{5.907957in}{3.648667in}%
\pgfsys@useobject{currentmarker}{}%
\end{pgfscope}%
\begin{pgfscope}%
\pgfsys@transformshift{5.961311in}{3.563313in}%
\pgfsys@useobject{currentmarker}{}%
\end{pgfscope}%
\begin{pgfscope}%
\pgfsys@transformshift{6.014665in}{3.318437in}%
\pgfsys@useobject{currentmarker}{}%
\end{pgfscope}%
\begin{pgfscope}%
\pgfsys@transformshift{6.068018in}{2.636300in}%
\pgfsys@useobject{currentmarker}{}%
\end{pgfscope}%
\begin{pgfscope}%
\pgfsys@transformshift{6.121372in}{2.572448in}%
\pgfsys@useobject{currentmarker}{}%
\end{pgfscope}%
\begin{pgfscope}%
\pgfsys@transformshift{6.174726in}{2.572448in}%
\pgfsys@useobject{currentmarker}{}%
\end{pgfscope}%
\begin{pgfscope}%
\pgfsys@transformshift{6.228079in}{2.991444in}%
\pgfsys@useobject{currentmarker}{}%
\end{pgfscope}%
\begin{pgfscope}%
\pgfsys@transformshift{6.281433in}{2.654549in}%
\pgfsys@useobject{currentmarker}{}%
\end{pgfscope}%
\begin{pgfscope}%
\pgfsys@transformshift{6.334787in}{2.819315in}%
\pgfsys@useobject{currentmarker}{}%
\end{pgfscope}%
\begin{pgfscope}%
\pgfsys@transformshift{6.388140in}{2.925582in}%
\pgfsys@useobject{currentmarker}{}%
\end{pgfscope}%
\begin{pgfscope}%
\pgfsys@transformshift{6.441494in}{2.860000in}%
\pgfsys@useobject{currentmarker}{}%
\end{pgfscope}%
\begin{pgfscope}%
\pgfsys@transformshift{6.494848in}{3.436396in}%
\pgfsys@useobject{currentmarker}{}%
\end{pgfscope}%
\begin{pgfscope}%
\pgfsys@transformshift{6.548201in}{3.648667in}%
\pgfsys@useobject{currentmarker}{}%
\end{pgfscope}%
\begin{pgfscope}%
\pgfsys@transformshift{6.601555in}{3.367000in}%
\pgfsys@useobject{currentmarker}{}%
\end{pgfscope}%
\begin{pgfscope}%
\pgfsys@transformshift{6.654909in}{3.063667in}%
\pgfsys@useobject{currentmarker}{}%
\end{pgfscope}%
\begin{pgfscope}%
\pgfsys@transformshift{6.708262in}{3.300725in}%
\pgfsys@useobject{currentmarker}{}%
\end{pgfscope}%
\begin{pgfscope}%
\pgfsys@transformshift{6.761616in}{3.648667in}%
\pgfsys@useobject{currentmarker}{}%
\end{pgfscope}%
\begin{pgfscope}%
\pgfsys@transformshift{6.814970in}{3.257463in}%
\pgfsys@useobject{currentmarker}{}%
\end{pgfscope}%
\begin{pgfscope}%
\pgfsys@transformshift{6.868323in}{2.602476in}%
\pgfsys@useobject{currentmarker}{}%
\end{pgfscope}%
\end{pgfscope}%
\begin{pgfscope}%
\pgfpathrectangle{\pgfqpoint{4.974268in}{2.488200in}}{\pgfqpoint{1.920732in}{1.216800in}}%
\pgfusepath{clip}%
\pgfsetbuttcap%
\pgfsetroundjoin%
\pgfsetlinewidth{0.602250pt}%
\definecolor{currentstroke}{rgb}{0.615686,0.705882,0.807843}%
\pgfsetstrokecolor{currentstroke}%
\pgfsetstrokeopacity{0.900000}%
\pgfsetdash{{1.200000pt}{1.200000pt}}{0.000000pt}%
\pgfpathmoveto{\pgfqpoint{5.000945in}{3.523481in}}%
\pgfpathlineto{\pgfqpoint{5.054299in}{3.267822in}}%
\pgfpathlineto{\pgfqpoint{5.107652in}{3.243771in}}%
\pgfpathlineto{\pgfqpoint{5.161006in}{3.141667in}}%
\pgfpathlineto{\pgfqpoint{5.214360in}{3.179222in}}%
\pgfpathlineto{\pgfqpoint{5.267713in}{3.289440in}}%
\pgfpathlineto{\pgfqpoint{5.321067in}{3.215333in}}%
\pgfpathlineto{\pgfqpoint{5.374421in}{3.013636in}}%
\pgfpathlineto{\pgfqpoint{5.427774in}{2.803667in}}%
\pgfpathlineto{\pgfqpoint{5.481128in}{3.460889in}}%
\pgfpathlineto{\pgfqpoint{5.534482in}{3.631851in}}%
\pgfpathlineto{\pgfqpoint{5.587835in}{3.648667in}}%
\pgfpathlineto{\pgfqpoint{5.641189in}{3.648667in}}%
\pgfpathlineto{\pgfqpoint{5.694543in}{3.648667in}}%
\pgfpathlineto{\pgfqpoint{5.747896in}{3.026478in}}%
\pgfpathlineto{\pgfqpoint{5.801250in}{3.209458in}}%
\pgfpathlineto{\pgfqpoint{5.854604in}{3.262381in}}%
\pgfpathlineto{\pgfqpoint{5.907957in}{3.617370in}}%
\pgfpathlineto{\pgfqpoint{5.961311in}{3.580384in}}%
\pgfpathlineto{\pgfqpoint{6.014665in}{3.318437in}}%
\pgfpathlineto{\pgfqpoint{6.068018in}{2.603643in}}%
\pgfpathlineto{\pgfqpoint{6.121372in}{2.522000in}}%
\pgfpathlineto{\pgfqpoint{6.174726in}{2.589264in}}%
\pgfpathlineto{\pgfqpoint{6.228079in}{3.198000in}}%
\pgfpathlineto{\pgfqpoint{6.281433in}{2.538569in}}%
\pgfpathlineto{\pgfqpoint{6.334787in}{2.772370in}}%
\pgfpathlineto{\pgfqpoint{6.388140in}{2.908766in}}%
\pgfpathlineto{\pgfqpoint{6.441494in}{2.860000in}}%
\pgfpathlineto{\pgfqpoint{6.494848in}{2.636300in}}%
\pgfpathlineto{\pgfqpoint{6.548201in}{3.100981in}}%
\pgfpathlineto{\pgfqpoint{6.601555in}{3.445241in}}%
\pgfpathlineto{\pgfqpoint{6.654909in}{3.323667in}}%
\pgfpathlineto{\pgfqpoint{6.708262in}{2.903078in}}%
\pgfpathlineto{\pgfqpoint{6.761616in}{3.069685in}}%
\pgfpathlineto{\pgfqpoint{6.814970in}{3.648667in}}%
\pgfpathlineto{\pgfqpoint{6.868323in}{2.843905in}}%
\pgfusepath{stroke}%
\end{pgfscope}%
\begin{pgfscope}%
\pgfpathrectangle{\pgfqpoint{4.974268in}{2.488200in}}{\pgfqpoint{1.920732in}{1.216800in}}%
\pgfusepath{clip}%
\pgfsetrectcap%
\pgfsetroundjoin%
\pgfsetlinewidth{0.702625pt}%
\definecolor{currentstroke}{rgb}{0.435294,0.576471,0.729412}%
\pgfsetstrokecolor{currentstroke}%
\pgfsetdash{}{0pt}%
\pgfpathmoveto{\pgfqpoint{5.000945in}{3.648667in}}%
\pgfpathlineto{\pgfqpoint{5.054299in}{3.648667in}}%
\pgfpathlineto{\pgfqpoint{5.107652in}{3.648667in}}%
\pgfpathlineto{\pgfqpoint{5.161006in}{3.648667in}}%
\pgfpathlineto{\pgfqpoint{5.214360in}{3.648667in}}%
\pgfpathlineto{\pgfqpoint{5.267713in}{3.648667in}}%
\pgfpathlineto{\pgfqpoint{5.321067in}{2.556667in}}%
\pgfpathlineto{\pgfqpoint{5.374421in}{2.808788in}}%
\pgfpathlineto{\pgfqpoint{5.427774in}{2.709778in}}%
\pgfpathlineto{\pgfqpoint{5.481128in}{2.725426in}}%
\pgfpathlineto{\pgfqpoint{5.534482in}{2.589264in}}%
\pgfpathlineto{\pgfqpoint{5.587835in}{2.685879in}}%
\pgfpathlineto{\pgfqpoint{5.641189in}{3.331296in}}%
\pgfpathlineto{\pgfqpoint{5.694543in}{3.413944in}}%
\pgfpathlineto{\pgfqpoint{5.747896in}{3.295532in}}%
\pgfpathlineto{\pgfqpoint{5.801250in}{2.770249in}}%
\pgfpathlineto{\pgfqpoint{5.854604in}{2.795619in}}%
\pgfpathlineto{\pgfqpoint{5.907957in}{3.007093in}}%
\pgfpathlineto{\pgfqpoint{5.961311in}{2.573212in}}%
\pgfpathlineto{\pgfqpoint{6.014665in}{3.415563in}}%
\pgfpathlineto{\pgfqpoint{6.068018in}{2.668957in}}%
\pgfpathlineto{\pgfqpoint{6.121372in}{2.774239in}}%
\pgfpathlineto{\pgfqpoint{6.174726in}{3.110557in}}%
\pgfpathlineto{\pgfqpoint{6.228079in}{2.522000in}}%
\pgfpathlineto{\pgfqpoint{6.281433in}{3.300725in}}%
\pgfpathlineto{\pgfqpoint{6.334787in}{2.522000in}}%
\pgfpathlineto{\pgfqpoint{6.388140in}{2.757423in}}%
\pgfpathlineto{\pgfqpoint{6.441494in}{2.766111in}}%
\pgfpathlineto{\pgfqpoint{6.494848in}{2.701614in}}%
\pgfpathlineto{\pgfqpoint{6.548201in}{2.913204in}}%
\pgfpathlineto{\pgfqpoint{6.601555in}{2.834963in}}%
\pgfpathlineto{\pgfqpoint{6.654909in}{3.475333in}}%
\pgfpathlineto{\pgfqpoint{6.708262in}{2.820235in}}%
\pgfpathlineto{\pgfqpoint{6.761616in}{3.194870in}}%
\pgfpathlineto{\pgfqpoint{6.814970in}{2.850611in}}%
\pgfpathlineto{\pgfqpoint{6.868323in}{3.246286in}}%
\pgfusepath{stroke}%
\end{pgfscope}%
\begin{pgfscope}%
\pgfpathrectangle{\pgfqpoint{4.974268in}{2.488200in}}{\pgfqpoint{1.920732in}{1.216800in}}%
\pgfusepath{clip}%
\pgfsetbuttcap%
\pgfsetroundjoin%
\definecolor{currentfill}{rgb}{0.435294,0.576471,0.729412}%
\pgfsetfillcolor{currentfill}%
\pgfsetlinewidth{1.003750pt}%
\definecolor{currentstroke}{rgb}{0.435294,0.576471,0.729412}%
\pgfsetstrokecolor{currentstroke}%
\pgfsetdash{}{0pt}%
\pgfsys@defobject{currentmarker}{\pgfqpoint{-0.010417in}{-0.010417in}}{\pgfqpoint{0.010417in}{0.010417in}}{%
\pgfpathmoveto{\pgfqpoint{0.000000in}{-0.010417in}}%
\pgfpathcurveto{\pgfqpoint{0.002763in}{-0.010417in}}{\pgfqpoint{0.005412in}{-0.009319in}}{\pgfqpoint{0.007366in}{-0.007366in}}%
\pgfpathcurveto{\pgfqpoint{0.009319in}{-0.005412in}}{\pgfqpoint{0.010417in}{-0.002763in}}{\pgfqpoint{0.010417in}{0.000000in}}%
\pgfpathcurveto{\pgfqpoint{0.010417in}{0.002763in}}{\pgfqpoint{0.009319in}{0.005412in}}{\pgfqpoint{0.007366in}{0.007366in}}%
\pgfpathcurveto{\pgfqpoint{0.005412in}{0.009319in}}{\pgfqpoint{0.002763in}{0.010417in}}{\pgfqpoint{0.000000in}{0.010417in}}%
\pgfpathcurveto{\pgfqpoint{-0.002763in}{0.010417in}}{\pgfqpoint{-0.005412in}{0.009319in}}{\pgfqpoint{-0.007366in}{0.007366in}}%
\pgfpathcurveto{\pgfqpoint{-0.009319in}{0.005412in}}{\pgfqpoint{-0.010417in}{0.002763in}}{\pgfqpoint{-0.010417in}{0.000000in}}%
\pgfpathcurveto{\pgfqpoint{-0.010417in}{-0.002763in}}{\pgfqpoint{-0.009319in}{-0.005412in}}{\pgfqpoint{-0.007366in}{-0.007366in}}%
\pgfpathcurveto{\pgfqpoint{-0.005412in}{-0.009319in}}{\pgfqpoint{-0.002763in}{-0.010417in}}{\pgfqpoint{0.000000in}{-0.010417in}}%
\pgfpathlineto{\pgfqpoint{0.000000in}{-0.010417in}}%
\pgfpathclose%
\pgfusepath{stroke,fill}%
}%
\begin{pgfscope}%
\pgfsys@transformshift{5.000945in}{3.648667in}%
\pgfsys@useobject{currentmarker}{}%
\end{pgfscope}%
\begin{pgfscope}%
\pgfsys@transformshift{5.054299in}{3.648667in}%
\pgfsys@useobject{currentmarker}{}%
\end{pgfscope}%
\begin{pgfscope}%
\pgfsys@transformshift{5.107652in}{3.648667in}%
\pgfsys@useobject{currentmarker}{}%
\end{pgfscope}%
\begin{pgfscope}%
\pgfsys@transformshift{5.161006in}{3.648667in}%
\pgfsys@useobject{currentmarker}{}%
\end{pgfscope}%
\begin{pgfscope}%
\pgfsys@transformshift{5.214360in}{3.648667in}%
\pgfsys@useobject{currentmarker}{}%
\end{pgfscope}%
\begin{pgfscope}%
\pgfsys@transformshift{5.267713in}{3.648667in}%
\pgfsys@useobject{currentmarker}{}%
\end{pgfscope}%
\begin{pgfscope}%
\pgfsys@transformshift{5.321067in}{2.556667in}%
\pgfsys@useobject{currentmarker}{}%
\end{pgfscope}%
\begin{pgfscope}%
\pgfsys@transformshift{5.374421in}{2.808788in}%
\pgfsys@useobject{currentmarker}{}%
\end{pgfscope}%
\begin{pgfscope}%
\pgfsys@transformshift{5.427774in}{2.709778in}%
\pgfsys@useobject{currentmarker}{}%
\end{pgfscope}%
\begin{pgfscope}%
\pgfsys@transformshift{5.481128in}{2.725426in}%
\pgfsys@useobject{currentmarker}{}%
\end{pgfscope}%
\begin{pgfscope}%
\pgfsys@transformshift{5.534482in}{2.589264in}%
\pgfsys@useobject{currentmarker}{}%
\end{pgfscope}%
\begin{pgfscope}%
\pgfsys@transformshift{5.587835in}{2.685879in}%
\pgfsys@useobject{currentmarker}{}%
\end{pgfscope}%
\begin{pgfscope}%
\pgfsys@transformshift{5.641189in}{3.331296in}%
\pgfsys@useobject{currentmarker}{}%
\end{pgfscope}%
\begin{pgfscope}%
\pgfsys@transformshift{5.694543in}{3.413944in}%
\pgfsys@useobject{currentmarker}{}%
\end{pgfscope}%
\begin{pgfscope}%
\pgfsys@transformshift{5.747896in}{3.295532in}%
\pgfsys@useobject{currentmarker}{}%
\end{pgfscope}%
\begin{pgfscope}%
\pgfsys@transformshift{5.801250in}{2.770249in}%
\pgfsys@useobject{currentmarker}{}%
\end{pgfscope}%
\begin{pgfscope}%
\pgfsys@transformshift{5.854604in}{2.795619in}%
\pgfsys@useobject{currentmarker}{}%
\end{pgfscope}%
\begin{pgfscope}%
\pgfsys@transformshift{5.907957in}{3.007093in}%
\pgfsys@useobject{currentmarker}{}%
\end{pgfscope}%
\begin{pgfscope}%
\pgfsys@transformshift{5.961311in}{2.573212in}%
\pgfsys@useobject{currentmarker}{}%
\end{pgfscope}%
\begin{pgfscope}%
\pgfsys@transformshift{6.014665in}{3.415563in}%
\pgfsys@useobject{currentmarker}{}%
\end{pgfscope}%
\begin{pgfscope}%
\pgfsys@transformshift{6.068018in}{2.668957in}%
\pgfsys@useobject{currentmarker}{}%
\end{pgfscope}%
\begin{pgfscope}%
\pgfsys@transformshift{6.121372in}{2.774239in}%
\pgfsys@useobject{currentmarker}{}%
\end{pgfscope}%
\begin{pgfscope}%
\pgfsys@transformshift{6.174726in}{3.110557in}%
\pgfsys@useobject{currentmarker}{}%
\end{pgfscope}%
\begin{pgfscope}%
\pgfsys@transformshift{6.228079in}{2.522000in}%
\pgfsys@useobject{currentmarker}{}%
\end{pgfscope}%
\begin{pgfscope}%
\pgfsys@transformshift{6.281433in}{3.300725in}%
\pgfsys@useobject{currentmarker}{}%
\end{pgfscope}%
\begin{pgfscope}%
\pgfsys@transformshift{6.334787in}{2.522000in}%
\pgfsys@useobject{currentmarker}{}%
\end{pgfscope}%
\begin{pgfscope}%
\pgfsys@transformshift{6.388140in}{2.757423in}%
\pgfsys@useobject{currentmarker}{}%
\end{pgfscope}%
\begin{pgfscope}%
\pgfsys@transformshift{6.441494in}{2.766111in}%
\pgfsys@useobject{currentmarker}{}%
\end{pgfscope}%
\begin{pgfscope}%
\pgfsys@transformshift{6.494848in}{2.701614in}%
\pgfsys@useobject{currentmarker}{}%
\end{pgfscope}%
\begin{pgfscope}%
\pgfsys@transformshift{6.548201in}{2.913204in}%
\pgfsys@useobject{currentmarker}{}%
\end{pgfscope}%
\begin{pgfscope}%
\pgfsys@transformshift{6.601555in}{2.834963in}%
\pgfsys@useobject{currentmarker}{}%
\end{pgfscope}%
\begin{pgfscope}%
\pgfsys@transformshift{6.654909in}{3.475333in}%
\pgfsys@useobject{currentmarker}{}%
\end{pgfscope}%
\begin{pgfscope}%
\pgfsys@transformshift{6.708262in}{2.820235in}%
\pgfsys@useobject{currentmarker}{}%
\end{pgfscope}%
\begin{pgfscope}%
\pgfsys@transformshift{6.761616in}{3.194870in}%
\pgfsys@useobject{currentmarker}{}%
\end{pgfscope}%
\begin{pgfscope}%
\pgfsys@transformshift{6.814970in}{2.850611in}%
\pgfsys@useobject{currentmarker}{}%
\end{pgfscope}%
\begin{pgfscope}%
\pgfsys@transformshift{6.868323in}{3.246286in}%
\pgfsys@useobject{currentmarker}{}%
\end{pgfscope}%
\end{pgfscope}%
\begin{pgfscope}%
\pgfpathrectangle{\pgfqpoint{4.974268in}{2.488200in}}{\pgfqpoint{1.920732in}{1.216800in}}%
\pgfusepath{clip}%
\pgfsetrectcap%
\pgfsetroundjoin%
\pgfsetlinewidth{0.702625pt}%
\definecolor{currentstroke}{rgb}{0.156863,0.356863,0.568627}%
\pgfsetstrokecolor{currentstroke}%
\pgfsetdash{}{0pt}%
\pgfpathmoveto{\pgfqpoint{5.000945in}{3.648667in}}%
\pgfpathlineto{\pgfqpoint{5.054299in}{3.648667in}}%
\pgfpathlineto{\pgfqpoint{5.107652in}{3.648667in}}%
\pgfpathlineto{\pgfqpoint{5.161006in}{3.648667in}}%
\pgfpathlineto{\pgfqpoint{5.214360in}{3.648667in}}%
\pgfpathlineto{\pgfqpoint{5.267713in}{3.648667in}}%
\pgfpathlineto{\pgfqpoint{5.321067in}{2.556667in}}%
\pgfpathlineto{\pgfqpoint{5.374421in}{2.870242in}}%
\pgfpathlineto{\pgfqpoint{5.427774in}{3.398296in}}%
\pgfpathlineto{\pgfqpoint{5.481128in}{3.100981in}}%
\pgfpathlineto{\pgfqpoint{5.534482in}{3.497323in}}%
\pgfpathlineto{\pgfqpoint{5.587835in}{3.648667in}}%
\pgfpathlineto{\pgfqpoint{5.641189in}{3.648667in}}%
\pgfpathlineto{\pgfqpoint{5.694543in}{3.570426in}}%
\pgfpathlineto{\pgfqpoint{5.747896in}{3.043294in}}%
\pgfpathlineto{\pgfqpoint{5.801250in}{3.209458in}}%
\pgfpathlineto{\pgfqpoint{5.854604in}{3.648667in}}%
\pgfpathlineto{\pgfqpoint{5.907957in}{3.398296in}}%
\pgfpathlineto{\pgfqpoint{5.961311in}{3.597455in}}%
\pgfpathlineto{\pgfqpoint{6.014665in}{3.182460in}}%
\pgfpathlineto{\pgfqpoint{6.068018in}{2.668957in}}%
\pgfpathlineto{\pgfqpoint{6.121372in}{2.824687in}}%
\pgfpathlineto{\pgfqpoint{6.174726in}{2.706975in}}%
\pgfpathlineto{\pgfqpoint{6.228079in}{3.160444in}}%
\pgfpathlineto{\pgfqpoint{6.281433in}{3.184745in}}%
\pgfpathlineto{\pgfqpoint{6.334787in}{3.648667in}}%
\pgfpathlineto{\pgfqpoint{6.388140in}{3.076925in}}%
\pgfpathlineto{\pgfqpoint{6.441494in}{3.216778in}}%
\pgfpathlineto{\pgfqpoint{6.494848in}{2.962870in}}%
\pgfpathlineto{\pgfqpoint{6.548201in}{2.944500in}}%
\pgfpathlineto{\pgfqpoint{6.601555in}{3.523481in}}%
\pgfpathlineto{\pgfqpoint{6.654909in}{2.912000in}}%
\pgfpathlineto{\pgfqpoint{6.708262in}{2.737392in}}%
\pgfpathlineto{\pgfqpoint{6.761616in}{3.570426in}}%
\pgfpathlineto{\pgfqpoint{6.814970in}{2.834963in}}%
\pgfpathlineto{\pgfqpoint{6.868323in}{3.407238in}}%
\pgfusepath{stroke}%
\end{pgfscope}%
\begin{pgfscope}%
\pgfpathrectangle{\pgfqpoint{4.974268in}{2.488200in}}{\pgfqpoint{1.920732in}{1.216800in}}%
\pgfusepath{clip}%
\pgfsetbuttcap%
\pgfsetroundjoin%
\definecolor{currentfill}{rgb}{0.156863,0.356863,0.568627}%
\pgfsetfillcolor{currentfill}%
\pgfsetlinewidth{1.003750pt}%
\definecolor{currentstroke}{rgb}{0.156863,0.356863,0.568627}%
\pgfsetstrokecolor{currentstroke}%
\pgfsetdash{}{0pt}%
\pgfsys@defobject{currentmarker}{\pgfqpoint{-0.010417in}{-0.010417in}}{\pgfqpoint{0.010417in}{0.010417in}}{%
\pgfpathmoveto{\pgfqpoint{0.000000in}{-0.010417in}}%
\pgfpathcurveto{\pgfqpoint{0.002763in}{-0.010417in}}{\pgfqpoint{0.005412in}{-0.009319in}}{\pgfqpoint{0.007366in}{-0.007366in}}%
\pgfpathcurveto{\pgfqpoint{0.009319in}{-0.005412in}}{\pgfqpoint{0.010417in}{-0.002763in}}{\pgfqpoint{0.010417in}{0.000000in}}%
\pgfpathcurveto{\pgfqpoint{0.010417in}{0.002763in}}{\pgfqpoint{0.009319in}{0.005412in}}{\pgfqpoint{0.007366in}{0.007366in}}%
\pgfpathcurveto{\pgfqpoint{0.005412in}{0.009319in}}{\pgfqpoint{0.002763in}{0.010417in}}{\pgfqpoint{0.000000in}{0.010417in}}%
\pgfpathcurveto{\pgfqpoint{-0.002763in}{0.010417in}}{\pgfqpoint{-0.005412in}{0.009319in}}{\pgfqpoint{-0.007366in}{0.007366in}}%
\pgfpathcurveto{\pgfqpoint{-0.009319in}{0.005412in}}{\pgfqpoint{-0.010417in}{0.002763in}}{\pgfqpoint{-0.010417in}{0.000000in}}%
\pgfpathcurveto{\pgfqpoint{-0.010417in}{-0.002763in}}{\pgfqpoint{-0.009319in}{-0.005412in}}{\pgfqpoint{-0.007366in}{-0.007366in}}%
\pgfpathcurveto{\pgfqpoint{-0.005412in}{-0.009319in}}{\pgfqpoint{-0.002763in}{-0.010417in}}{\pgfqpoint{0.000000in}{-0.010417in}}%
\pgfpathlineto{\pgfqpoint{0.000000in}{-0.010417in}}%
\pgfpathclose%
\pgfusepath{stroke,fill}%
}%
\begin{pgfscope}%
\pgfsys@transformshift{5.000945in}{3.648667in}%
\pgfsys@useobject{currentmarker}{}%
\end{pgfscope}%
\begin{pgfscope}%
\pgfsys@transformshift{5.054299in}{3.648667in}%
\pgfsys@useobject{currentmarker}{}%
\end{pgfscope}%
\begin{pgfscope}%
\pgfsys@transformshift{5.107652in}{3.648667in}%
\pgfsys@useobject{currentmarker}{}%
\end{pgfscope}%
\begin{pgfscope}%
\pgfsys@transformshift{5.161006in}{3.648667in}%
\pgfsys@useobject{currentmarker}{}%
\end{pgfscope}%
\begin{pgfscope}%
\pgfsys@transformshift{5.214360in}{3.648667in}%
\pgfsys@useobject{currentmarker}{}%
\end{pgfscope}%
\begin{pgfscope}%
\pgfsys@transformshift{5.267713in}{3.648667in}%
\pgfsys@useobject{currentmarker}{}%
\end{pgfscope}%
\begin{pgfscope}%
\pgfsys@transformshift{5.321067in}{2.556667in}%
\pgfsys@useobject{currentmarker}{}%
\end{pgfscope}%
\begin{pgfscope}%
\pgfsys@transformshift{5.374421in}{2.870242in}%
\pgfsys@useobject{currentmarker}{}%
\end{pgfscope}%
\begin{pgfscope}%
\pgfsys@transformshift{5.427774in}{3.398296in}%
\pgfsys@useobject{currentmarker}{}%
\end{pgfscope}%
\begin{pgfscope}%
\pgfsys@transformshift{5.481128in}{3.100981in}%
\pgfsys@useobject{currentmarker}{}%
\end{pgfscope}%
\begin{pgfscope}%
\pgfsys@transformshift{5.534482in}{3.497323in}%
\pgfsys@useobject{currentmarker}{}%
\end{pgfscope}%
\begin{pgfscope}%
\pgfsys@transformshift{5.587835in}{3.648667in}%
\pgfsys@useobject{currentmarker}{}%
\end{pgfscope}%
\begin{pgfscope}%
\pgfsys@transformshift{5.641189in}{3.648667in}%
\pgfsys@useobject{currentmarker}{}%
\end{pgfscope}%
\begin{pgfscope}%
\pgfsys@transformshift{5.694543in}{3.570426in}%
\pgfsys@useobject{currentmarker}{}%
\end{pgfscope}%
\begin{pgfscope}%
\pgfsys@transformshift{5.747896in}{3.043294in}%
\pgfsys@useobject{currentmarker}{}%
\end{pgfscope}%
\begin{pgfscope}%
\pgfsys@transformshift{5.801250in}{3.209458in}%
\pgfsys@useobject{currentmarker}{}%
\end{pgfscope}%
\begin{pgfscope}%
\pgfsys@transformshift{5.854604in}{3.648667in}%
\pgfsys@useobject{currentmarker}{}%
\end{pgfscope}%
\begin{pgfscope}%
\pgfsys@transformshift{5.907957in}{3.398296in}%
\pgfsys@useobject{currentmarker}{}%
\end{pgfscope}%
\begin{pgfscope}%
\pgfsys@transformshift{5.961311in}{3.597455in}%
\pgfsys@useobject{currentmarker}{}%
\end{pgfscope}%
\begin{pgfscope}%
\pgfsys@transformshift{6.014665in}{3.182460in}%
\pgfsys@useobject{currentmarker}{}%
\end{pgfscope}%
\begin{pgfscope}%
\pgfsys@transformshift{6.068018in}{2.668957in}%
\pgfsys@useobject{currentmarker}{}%
\end{pgfscope}%
\begin{pgfscope}%
\pgfsys@transformshift{6.121372in}{2.824687in}%
\pgfsys@useobject{currentmarker}{}%
\end{pgfscope}%
\begin{pgfscope}%
\pgfsys@transformshift{6.174726in}{2.706975in}%
\pgfsys@useobject{currentmarker}{}%
\end{pgfscope}%
\begin{pgfscope}%
\pgfsys@transformshift{6.228079in}{3.160444in}%
\pgfsys@useobject{currentmarker}{}%
\end{pgfscope}%
\begin{pgfscope}%
\pgfsys@transformshift{6.281433in}{3.184745in}%
\pgfsys@useobject{currentmarker}{}%
\end{pgfscope}%
\begin{pgfscope}%
\pgfsys@transformshift{6.334787in}{3.648667in}%
\pgfsys@useobject{currentmarker}{}%
\end{pgfscope}%
\begin{pgfscope}%
\pgfsys@transformshift{6.388140in}{3.076925in}%
\pgfsys@useobject{currentmarker}{}%
\end{pgfscope}%
\begin{pgfscope}%
\pgfsys@transformshift{6.441494in}{3.216778in}%
\pgfsys@useobject{currentmarker}{}%
\end{pgfscope}%
\begin{pgfscope}%
\pgfsys@transformshift{6.494848in}{2.962870in}%
\pgfsys@useobject{currentmarker}{}%
\end{pgfscope}%
\begin{pgfscope}%
\pgfsys@transformshift{6.548201in}{2.944500in}%
\pgfsys@useobject{currentmarker}{}%
\end{pgfscope}%
\begin{pgfscope}%
\pgfsys@transformshift{6.601555in}{3.523481in}%
\pgfsys@useobject{currentmarker}{}%
\end{pgfscope}%
\begin{pgfscope}%
\pgfsys@transformshift{6.654909in}{2.912000in}%
\pgfsys@useobject{currentmarker}{}%
\end{pgfscope}%
\begin{pgfscope}%
\pgfsys@transformshift{6.708262in}{2.737392in}%
\pgfsys@useobject{currentmarker}{}%
\end{pgfscope}%
\begin{pgfscope}%
\pgfsys@transformshift{6.761616in}{3.570426in}%
\pgfsys@useobject{currentmarker}{}%
\end{pgfscope}%
\begin{pgfscope}%
\pgfsys@transformshift{6.814970in}{2.834963in}%
\pgfsys@useobject{currentmarker}{}%
\end{pgfscope}%
\begin{pgfscope}%
\pgfsys@transformshift{6.868323in}{3.407238in}%
\pgfsys@useobject{currentmarker}{}%
\end{pgfscope}%
\end{pgfscope}%
\begin{pgfscope}%
\pgfpathrectangle{\pgfqpoint{4.974268in}{2.488200in}}{\pgfqpoint{1.920732in}{1.216800in}}%
\pgfusepath{clip}%
\pgfsetbuttcap%
\pgfsetroundjoin%
\pgfsetlinewidth{0.602250pt}%
\definecolor{currentstroke}{rgb}{0.156863,0.356863,0.568627}%
\pgfsetstrokecolor{currentstroke}%
\pgfsetstrokeopacity{0.900000}%
\pgfsetdash{{1.200000pt}{1.200000pt}}{0.000000pt}%
\pgfpathmoveto{\pgfqpoint{5.000945in}{3.633019in}}%
\pgfpathlineto{\pgfqpoint{5.054299in}{3.458244in}}%
\pgfpathlineto{\pgfqpoint{5.107652in}{3.648667in}}%
\pgfpathlineto{\pgfqpoint{5.161006in}{3.554778in}}%
\pgfpathlineto{\pgfqpoint{5.214360in}{2.803667in}}%
\pgfpathlineto{\pgfqpoint{5.267713in}{3.648667in}}%
\pgfpathlineto{\pgfqpoint{5.321067in}{3.267333in}}%
\pgfpathlineto{\pgfqpoint{5.374421in}{3.361879in}}%
\pgfpathlineto{\pgfqpoint{5.427774in}{3.460889in}}%
\pgfpathlineto{\pgfqpoint{5.481128in}{3.492185in}}%
\pgfpathlineto{\pgfqpoint{5.534482in}{3.648667in}}%
\pgfpathlineto{\pgfqpoint{5.587835in}{3.443818in}}%
\pgfpathlineto{\pgfqpoint{5.641189in}{3.648667in}}%
\pgfpathlineto{\pgfqpoint{5.694543in}{3.648667in}}%
\pgfpathlineto{\pgfqpoint{5.747896in}{3.295532in}}%
\pgfpathlineto{\pgfqpoint{5.801250in}{3.343130in}}%
\pgfpathlineto{\pgfqpoint{5.854604in}{3.648667in}}%
\pgfpathlineto{\pgfqpoint{5.907957in}{3.648667in}}%
\pgfpathlineto{\pgfqpoint{5.961311in}{3.597455in}}%
\pgfpathlineto{\pgfqpoint{6.014665in}{3.648667in}}%
\pgfpathlineto{\pgfqpoint{6.068018in}{3.305768in}}%
\pgfpathlineto{\pgfqpoint{6.121372in}{2.959214in}}%
\pgfpathlineto{\pgfqpoint{6.174726in}{3.060109in}}%
\pgfpathlineto{\pgfqpoint{6.228079in}{3.329444in}}%
\pgfpathlineto{\pgfqpoint{6.281433in}{2.969353in}}%
\pgfpathlineto{\pgfqpoint{6.334787in}{3.054037in}}%
\pgfpathlineto{\pgfqpoint{6.388140in}{3.161005in}}%
\pgfpathlineto{\pgfqpoint{6.441494in}{2.972667in}}%
\pgfpathlineto{\pgfqpoint{6.494848in}{2.946541in}}%
\pgfpathlineto{\pgfqpoint{6.548201in}{3.507833in}}%
\pgfpathlineto{\pgfqpoint{6.601555in}{3.617370in}}%
\pgfpathlineto{\pgfqpoint{6.654909in}{3.432000in}}%
\pgfpathlineto{\pgfqpoint{6.708262in}{3.582392in}}%
\pgfpathlineto{\pgfqpoint{6.761616in}{3.476537in}}%
\pgfpathlineto{\pgfqpoint{6.814970in}{3.007093in}}%
\pgfpathlineto{\pgfqpoint{6.868323in}{2.984738in}}%
\pgfusepath{stroke}%
\end{pgfscope}%
\begin{pgfscope}%
\pgfpathrectangle{\pgfqpoint{4.974268in}{2.488200in}}{\pgfqpoint{1.920732in}{1.216800in}}%
\pgfusepath{clip}%
\pgfsetrectcap%
\pgfsetroundjoin%
\pgfsetlinewidth{0.702625pt}%
\definecolor{currentstroke}{rgb}{0.062745,0.160784,0.290196}%
\pgfsetstrokecolor{currentstroke}%
\pgfsetdash{}{0pt}%
\pgfpathmoveto{\pgfqpoint{5.000945in}{3.648667in}}%
\pgfpathlineto{\pgfqpoint{5.054299in}{3.648667in}}%
\pgfpathlineto{\pgfqpoint{5.107652in}{3.648667in}}%
\pgfpathlineto{\pgfqpoint{5.161006in}{3.648667in}}%
\pgfpathlineto{\pgfqpoint{5.214360in}{3.648667in}}%
\pgfpathlineto{\pgfqpoint{5.267713in}{3.648667in}}%
\pgfpathlineto{\pgfqpoint{5.321067in}{2.556667in}}%
\pgfpathlineto{\pgfqpoint{5.374421in}{2.849758in}}%
\pgfpathlineto{\pgfqpoint{5.427774in}{3.476537in}}%
\pgfpathlineto{\pgfqpoint{5.481128in}{2.913204in}}%
\pgfpathlineto{\pgfqpoint{5.534482in}{3.547771in}}%
\pgfpathlineto{\pgfqpoint{5.587835in}{3.648667in}}%
\pgfpathlineto{\pgfqpoint{5.641189in}{3.648667in}}%
\pgfpathlineto{\pgfqpoint{5.694543in}{3.648667in}}%
\pgfpathlineto{\pgfqpoint{5.747896in}{3.009662in}}%
\pgfpathlineto{\pgfqpoint{5.801250in}{3.133073in}}%
\pgfpathlineto{\pgfqpoint{5.854604in}{3.648667in}}%
\pgfpathlineto{\pgfqpoint{5.907957in}{3.648667in}}%
\pgfpathlineto{\pgfqpoint{5.961311in}{3.034121in}}%
\pgfpathlineto{\pgfqpoint{6.014665in}{3.065908in}}%
\pgfpathlineto{\pgfqpoint{6.068018in}{2.636300in}}%
\pgfpathlineto{\pgfqpoint{6.121372in}{2.942398in}}%
\pgfpathlineto{\pgfqpoint{6.174726in}{2.522000in}}%
\pgfpathlineto{\pgfqpoint{6.228079in}{2.916333in}}%
\pgfpathlineto{\pgfqpoint{6.281433in}{2.654549in}}%
\pgfpathlineto{\pgfqpoint{6.334787in}{3.554778in}}%
\pgfpathlineto{\pgfqpoint{6.388140in}{2.976030in}}%
\pgfpathlineto{\pgfqpoint{6.441494in}{3.085333in}}%
\pgfpathlineto{\pgfqpoint{6.494848in}{3.256783in}}%
\pgfpathlineto{\pgfqpoint{6.548201in}{3.085333in}}%
\pgfpathlineto{\pgfqpoint{6.601555in}{3.116630in}}%
\pgfpathlineto{\pgfqpoint{6.654909in}{3.345333in}}%
\pgfpathlineto{\pgfqpoint{6.708262in}{2.836804in}}%
\pgfpathlineto{\pgfqpoint{6.761616in}{2.913204in}}%
\pgfpathlineto{\pgfqpoint{6.814970in}{2.819315in}}%
\pgfpathlineto{\pgfqpoint{6.868323in}{3.045095in}}%
\pgfusepath{stroke}%
\end{pgfscope}%
\begin{pgfscope}%
\pgfpathrectangle{\pgfqpoint{4.974268in}{2.488200in}}{\pgfqpoint{1.920732in}{1.216800in}}%
\pgfusepath{clip}%
\pgfsetbuttcap%
\pgfsetroundjoin%
\definecolor{currentfill}{rgb}{0.062745,0.160784,0.290196}%
\pgfsetfillcolor{currentfill}%
\pgfsetlinewidth{1.003750pt}%
\definecolor{currentstroke}{rgb}{0.062745,0.160784,0.290196}%
\pgfsetstrokecolor{currentstroke}%
\pgfsetdash{}{0pt}%
\pgfsys@defobject{currentmarker}{\pgfqpoint{-0.010417in}{-0.010417in}}{\pgfqpoint{0.010417in}{0.010417in}}{%
\pgfpathmoveto{\pgfqpoint{0.000000in}{-0.010417in}}%
\pgfpathcurveto{\pgfqpoint{0.002763in}{-0.010417in}}{\pgfqpoint{0.005412in}{-0.009319in}}{\pgfqpoint{0.007366in}{-0.007366in}}%
\pgfpathcurveto{\pgfqpoint{0.009319in}{-0.005412in}}{\pgfqpoint{0.010417in}{-0.002763in}}{\pgfqpoint{0.010417in}{0.000000in}}%
\pgfpathcurveto{\pgfqpoint{0.010417in}{0.002763in}}{\pgfqpoint{0.009319in}{0.005412in}}{\pgfqpoint{0.007366in}{0.007366in}}%
\pgfpathcurveto{\pgfqpoint{0.005412in}{0.009319in}}{\pgfqpoint{0.002763in}{0.010417in}}{\pgfqpoint{0.000000in}{0.010417in}}%
\pgfpathcurveto{\pgfqpoint{-0.002763in}{0.010417in}}{\pgfqpoint{-0.005412in}{0.009319in}}{\pgfqpoint{-0.007366in}{0.007366in}}%
\pgfpathcurveto{\pgfqpoint{-0.009319in}{0.005412in}}{\pgfqpoint{-0.010417in}{0.002763in}}{\pgfqpoint{-0.010417in}{0.000000in}}%
\pgfpathcurveto{\pgfqpoint{-0.010417in}{-0.002763in}}{\pgfqpoint{-0.009319in}{-0.005412in}}{\pgfqpoint{-0.007366in}{-0.007366in}}%
\pgfpathcurveto{\pgfqpoint{-0.005412in}{-0.009319in}}{\pgfqpoint{-0.002763in}{-0.010417in}}{\pgfqpoint{0.000000in}{-0.010417in}}%
\pgfpathlineto{\pgfqpoint{0.000000in}{-0.010417in}}%
\pgfpathclose%
\pgfusepath{stroke,fill}%
}%
\begin{pgfscope}%
\pgfsys@transformshift{5.000945in}{3.648667in}%
\pgfsys@useobject{currentmarker}{}%
\end{pgfscope}%
\begin{pgfscope}%
\pgfsys@transformshift{5.054299in}{3.648667in}%
\pgfsys@useobject{currentmarker}{}%
\end{pgfscope}%
\begin{pgfscope}%
\pgfsys@transformshift{5.107652in}{3.648667in}%
\pgfsys@useobject{currentmarker}{}%
\end{pgfscope}%
\begin{pgfscope}%
\pgfsys@transformshift{5.161006in}{3.648667in}%
\pgfsys@useobject{currentmarker}{}%
\end{pgfscope}%
\begin{pgfscope}%
\pgfsys@transformshift{5.214360in}{3.648667in}%
\pgfsys@useobject{currentmarker}{}%
\end{pgfscope}%
\begin{pgfscope}%
\pgfsys@transformshift{5.267713in}{3.648667in}%
\pgfsys@useobject{currentmarker}{}%
\end{pgfscope}%
\begin{pgfscope}%
\pgfsys@transformshift{5.321067in}{2.556667in}%
\pgfsys@useobject{currentmarker}{}%
\end{pgfscope}%
\begin{pgfscope}%
\pgfsys@transformshift{5.374421in}{2.849758in}%
\pgfsys@useobject{currentmarker}{}%
\end{pgfscope}%
\begin{pgfscope}%
\pgfsys@transformshift{5.427774in}{3.476537in}%
\pgfsys@useobject{currentmarker}{}%
\end{pgfscope}%
\begin{pgfscope}%
\pgfsys@transformshift{5.481128in}{2.913204in}%
\pgfsys@useobject{currentmarker}{}%
\end{pgfscope}%
\begin{pgfscope}%
\pgfsys@transformshift{5.534482in}{3.547771in}%
\pgfsys@useobject{currentmarker}{}%
\end{pgfscope}%
\begin{pgfscope}%
\pgfsys@transformshift{5.587835in}{3.648667in}%
\pgfsys@useobject{currentmarker}{}%
\end{pgfscope}%
\begin{pgfscope}%
\pgfsys@transformshift{5.641189in}{3.648667in}%
\pgfsys@useobject{currentmarker}{}%
\end{pgfscope}%
\begin{pgfscope}%
\pgfsys@transformshift{5.694543in}{3.648667in}%
\pgfsys@useobject{currentmarker}{}%
\end{pgfscope}%
\begin{pgfscope}%
\pgfsys@transformshift{5.747896in}{3.009662in}%
\pgfsys@useobject{currentmarker}{}%
\end{pgfscope}%
\begin{pgfscope}%
\pgfsys@transformshift{5.801250in}{3.133073in}%
\pgfsys@useobject{currentmarker}{}%
\end{pgfscope}%
\begin{pgfscope}%
\pgfsys@transformshift{5.854604in}{3.648667in}%
\pgfsys@useobject{currentmarker}{}%
\end{pgfscope}%
\begin{pgfscope}%
\pgfsys@transformshift{5.907957in}{3.648667in}%
\pgfsys@useobject{currentmarker}{}%
\end{pgfscope}%
\begin{pgfscope}%
\pgfsys@transformshift{5.961311in}{3.034121in}%
\pgfsys@useobject{currentmarker}{}%
\end{pgfscope}%
\begin{pgfscope}%
\pgfsys@transformshift{6.014665in}{3.065908in}%
\pgfsys@useobject{currentmarker}{}%
\end{pgfscope}%
\begin{pgfscope}%
\pgfsys@transformshift{6.068018in}{2.636300in}%
\pgfsys@useobject{currentmarker}{}%
\end{pgfscope}%
\begin{pgfscope}%
\pgfsys@transformshift{6.121372in}{2.942398in}%
\pgfsys@useobject{currentmarker}{}%
\end{pgfscope}%
\begin{pgfscope}%
\pgfsys@transformshift{6.174726in}{2.522000in}%
\pgfsys@useobject{currentmarker}{}%
\end{pgfscope}%
\begin{pgfscope}%
\pgfsys@transformshift{6.228079in}{2.916333in}%
\pgfsys@useobject{currentmarker}{}%
\end{pgfscope}%
\begin{pgfscope}%
\pgfsys@transformshift{6.281433in}{2.654549in}%
\pgfsys@useobject{currentmarker}{}%
\end{pgfscope}%
\begin{pgfscope}%
\pgfsys@transformshift{6.334787in}{3.554778in}%
\pgfsys@useobject{currentmarker}{}%
\end{pgfscope}%
\begin{pgfscope}%
\pgfsys@transformshift{6.388140in}{2.976030in}%
\pgfsys@useobject{currentmarker}{}%
\end{pgfscope}%
\begin{pgfscope}%
\pgfsys@transformshift{6.441494in}{3.085333in}%
\pgfsys@useobject{currentmarker}{}%
\end{pgfscope}%
\begin{pgfscope}%
\pgfsys@transformshift{6.494848in}{3.256783in}%
\pgfsys@useobject{currentmarker}{}%
\end{pgfscope}%
\begin{pgfscope}%
\pgfsys@transformshift{6.548201in}{3.085333in}%
\pgfsys@useobject{currentmarker}{}%
\end{pgfscope}%
\begin{pgfscope}%
\pgfsys@transformshift{6.601555in}{3.116630in}%
\pgfsys@useobject{currentmarker}{}%
\end{pgfscope}%
\begin{pgfscope}%
\pgfsys@transformshift{6.654909in}{3.345333in}%
\pgfsys@useobject{currentmarker}{}%
\end{pgfscope}%
\begin{pgfscope}%
\pgfsys@transformshift{6.708262in}{2.836804in}%
\pgfsys@useobject{currentmarker}{}%
\end{pgfscope}%
\begin{pgfscope}%
\pgfsys@transformshift{6.761616in}{2.913204in}%
\pgfsys@useobject{currentmarker}{}%
\end{pgfscope}%
\begin{pgfscope}%
\pgfsys@transformshift{6.814970in}{2.819315in}%
\pgfsys@useobject{currentmarker}{}%
\end{pgfscope}%
\begin{pgfscope}%
\pgfsys@transformshift{6.868323in}{3.045095in}%
\pgfsys@useobject{currentmarker}{}%
\end{pgfscope}%
\end{pgfscope}%
\begin{pgfscope}%
\pgfpathrectangle{\pgfqpoint{4.974268in}{2.488200in}}{\pgfqpoint{1.920732in}{1.216800in}}%
\pgfusepath{clip}%
\pgfsetrectcap%
\pgfsetroundjoin%
\pgfsetlinewidth{0.501875pt}%
\definecolor{currentstroke}{rgb}{0.101961,0.101961,0.101961}%
\pgfsetstrokecolor{currentstroke}%
\pgfsetdash{}{0pt}%
\pgfpathmoveto{\pgfqpoint{5.294390in}{2.488200in}}%
\pgfpathlineto{\pgfqpoint{5.294390in}{3.705000in}}%
\pgfusepath{stroke}%
\end{pgfscope}%
\begin{pgfscope}%
\pgfsetrectcap%
\pgfsetmiterjoin%
\pgfsetlinewidth{0.803000pt}%
\definecolor{currentstroke}{rgb}{0.101961,0.101961,0.101961}%
\pgfsetstrokecolor{currentstroke}%
\pgfsetdash{}{0pt}%
\pgfpathmoveto{\pgfqpoint{4.974268in}{2.488200in}}%
\pgfpathlineto{\pgfqpoint{4.974268in}{3.705000in}}%
\pgfusepath{stroke}%
\end{pgfscope}%
\begin{pgfscope}%
\pgfsetrectcap%
\pgfsetmiterjoin%
\pgfsetlinewidth{0.803000pt}%
\definecolor{currentstroke}{rgb}{0.101961,0.101961,0.101961}%
\pgfsetstrokecolor{currentstroke}%
\pgfsetdash{}{0pt}%
\pgfpathmoveto{\pgfqpoint{4.974268in}{2.488200in}}%
\pgfpathlineto{\pgfqpoint{6.895000in}{2.488200in}}%
\pgfusepath{stroke}%
\end{pgfscope}%
\begin{pgfscope}%
\pgfsetbuttcap%
\pgfsetmiterjoin%
\definecolor{currentfill}{rgb}{1.000000,1.000000,1.000000}%
\pgfsetfillcolor{currentfill}%
\pgfsetfillopacity{0.850000}%
\pgfsetlinewidth{0.000000pt}%
\definecolor{currentstroke}{rgb}{0.000000,0.000000,0.000000}%
\pgfsetstrokecolor{currentstroke}%
\pgfsetstrokeopacity{0.850000}%
\pgfsetdash{}{0pt}%
\pgfpathmoveto{\pgfqpoint{6.448241in}{2.523261in}}%
\pgfpathlineto{\pgfqpoint{6.870196in}{2.523261in}}%
\pgfpathlineto{\pgfqpoint{6.870196in}{2.636903in}}%
\pgfpathlineto{\pgfqpoint{6.448241in}{2.636903in}}%
\pgfpathlineto{\pgfqpoint{6.448241in}{2.523261in}}%
\pgfpathclose%
\pgfusepath{fill}%
\end{pgfscope}%
\begin{pgfscope}%
\definecolor{textcolor}{rgb}{0.101961,0.101961,0.101961}%
\pgfsetstrokecolor{textcolor}%
\pgfsetfillcolor{textcolor}%
\pgftext[x=6.856585in,y=2.536872in,right,bottom]{\color{textcolor}\sffamily\fontsize{7.000000}{8.400000}\selectfont Dashcam}%
\end{pgfscope}%
\begin{pgfscope}%
\definecolor{textcolor}{rgb}{0.466667,0.466667,0.466667}%
\pgfsetstrokecolor{textcolor}%
\pgfsetfillcolor{textcolor}%
\pgftext[x=5.134329in,y=3.682467in,,bottom]{\color{textcolor}\sffamily\fontsize{5.500000}{6.600000}\selectfont outage}%
\end{pgfscope}%
\begin{pgfscope}%
\pgfsetbuttcap%
\pgfsetmiterjoin%
\definecolor{currentfill}{rgb}{1.000000,1.000000,1.000000}%
\pgfsetfillcolor{currentfill}%
\pgfsetlinewidth{0.000000pt}%
\definecolor{currentstroke}{rgb}{0.000000,0.000000,0.000000}%
\pgfsetstrokecolor{currentstroke}%
\pgfsetstrokeopacity{0.000000}%
\pgfsetdash{}{0pt}%
\pgfpathmoveto{\pgfqpoint{0.595000in}{0.663000in}}%
\pgfpathlineto{\pgfqpoint{2.515732in}{0.663000in}}%
\pgfpathlineto{\pgfqpoint{2.515732in}{1.879800in}}%
\pgfpathlineto{\pgfqpoint{0.595000in}{1.879800in}}%
\pgfpathlineto{\pgfqpoint{0.595000in}{0.663000in}}%
\pgfpathclose%
\pgfusepath{fill}%
\end{pgfscope}%
\begin{pgfscope}%
\definecolor{textcolor}{rgb}{0.101961,0.101961,0.101961}%
\pgfsetstrokecolor{textcolor}%
\pgfsetfillcolor{textcolor}%
\pgftext[x=0.869390in,y=0.614389in,,top]{\color{textcolor}\sffamily\fontsize{7.000000}{8.400000}\selectfont rung 3}%
\end{pgfscope}%
\begin{pgfscope}%
\definecolor{textcolor}{rgb}{0.101961,0.101961,0.101961}%
\pgfsetstrokecolor{textcolor}%
\pgfsetfillcolor{textcolor}%
\pgftext[x=1.326707in,y=0.614389in,,top]{\color{textcolor}\sffamily\fontsize{7.000000}{8.400000}\selectfont eager}%
\end{pgfscope}%
\begin{pgfscope}%
\definecolor{textcolor}{rgb}{0.101961,0.101961,0.101961}%
\pgfsetstrokecolor{textcolor}%
\pgfsetfillcolor{textcolor}%
\pgftext[x=1.784024in,y=0.614389in,,top]{\color{textcolor}\sffamily\fontsize{7.000000}{8.400000}\selectfont no S2}%
\end{pgfscope}%
\begin{pgfscope}%
\definecolor{textcolor}{rgb}{0.101961,0.101961,0.101961}%
\pgfsetstrokecolor{textcolor}%
\pgfsetfillcolor{textcolor}%
\pgftext[x=2.241341in,y=0.614389in,,top]{\color{textcolor}\sffamily\fontsize{7.000000}{8.400000}\selectfont + S2}%
\end{pgfscope}%
\begin{pgfscope}%
\pgfpathrectangle{\pgfqpoint{0.595000in}{0.663000in}}{\pgfqpoint{1.920732in}{1.216800in}}%
\pgfusepath{clip}%
\pgfsetrectcap%
\pgfsetroundjoin%
\pgfsetlinewidth{0.401500pt}%
\definecolor{currentstroke}{rgb}{0.901961,0.901961,0.901961}%
\pgfsetstrokecolor{currentstroke}%
\pgfsetdash{}{0pt}%
\pgfpathmoveto{\pgfqpoint{0.595000in}{0.663000in}}%
\pgfpathlineto{\pgfqpoint{2.515732in}{0.663000in}}%
\pgfusepath{stroke}%
\end{pgfscope}%
\begin{pgfscope}%
\definecolor{textcolor}{rgb}{0.101961,0.101961,0.101961}%
\pgfsetstrokecolor{textcolor}%
\pgfsetfillcolor{textcolor}%
\pgftext[x=0.414395in, y=0.629242in, left, base]{\color{textcolor}\sffamily\fontsize{7.000000}{8.400000}\selectfont 0.0}%
\end{pgfscope}%
\begin{pgfscope}%
\pgfpathrectangle{\pgfqpoint{0.595000in}{0.663000in}}{\pgfqpoint{1.920732in}{1.216800in}}%
\pgfusepath{clip}%
\pgfsetrectcap%
\pgfsetroundjoin%
\pgfsetlinewidth{0.401500pt}%
\definecolor{currentstroke}{rgb}{0.901961,0.901961,0.901961}%
\pgfsetstrokecolor{currentstroke}%
\pgfsetdash{}{0pt}%
\pgfpathmoveto{\pgfqpoint{0.595000in}{0.894771in}}%
\pgfpathlineto{\pgfqpoint{2.515732in}{0.894771in}}%
\pgfusepath{stroke}%
\end{pgfscope}%
\begin{pgfscope}%
\definecolor{textcolor}{rgb}{0.101961,0.101961,0.101961}%
\pgfsetstrokecolor{textcolor}%
\pgfsetfillcolor{textcolor}%
\pgftext[x=0.414395in, y=0.861014in, left, base]{\color{textcolor}\sffamily\fontsize{7.000000}{8.400000}\selectfont 0.2}%
\end{pgfscope}%
\begin{pgfscope}%
\pgfpathrectangle{\pgfqpoint{0.595000in}{0.663000in}}{\pgfqpoint{1.920732in}{1.216800in}}%
\pgfusepath{clip}%
\pgfsetrectcap%
\pgfsetroundjoin%
\pgfsetlinewidth{0.401500pt}%
\definecolor{currentstroke}{rgb}{0.901961,0.901961,0.901961}%
\pgfsetstrokecolor{currentstroke}%
\pgfsetdash{}{0pt}%
\pgfpathmoveto{\pgfqpoint{0.595000in}{1.126543in}}%
\pgfpathlineto{\pgfqpoint{2.515732in}{1.126543in}}%
\pgfusepath{stroke}%
\end{pgfscope}%
\begin{pgfscope}%
\definecolor{textcolor}{rgb}{0.101961,0.101961,0.101961}%
\pgfsetstrokecolor{textcolor}%
\pgfsetfillcolor{textcolor}%
\pgftext[x=0.414395in, y=1.092785in, left, base]{\color{textcolor}\sffamily\fontsize{7.000000}{8.400000}\selectfont 0.4}%
\end{pgfscope}%
\begin{pgfscope}%
\pgfpathrectangle{\pgfqpoint{0.595000in}{0.663000in}}{\pgfqpoint{1.920732in}{1.216800in}}%
\pgfusepath{clip}%
\pgfsetrectcap%
\pgfsetroundjoin%
\pgfsetlinewidth{0.401500pt}%
\definecolor{currentstroke}{rgb}{0.901961,0.901961,0.901961}%
\pgfsetstrokecolor{currentstroke}%
\pgfsetdash{}{0pt}%
\pgfpathmoveto{\pgfqpoint{0.595000in}{1.358314in}}%
\pgfpathlineto{\pgfqpoint{2.515732in}{1.358314in}}%
\pgfusepath{stroke}%
\end{pgfscope}%
\begin{pgfscope}%
\definecolor{textcolor}{rgb}{0.101961,0.101961,0.101961}%
\pgfsetstrokecolor{textcolor}%
\pgfsetfillcolor{textcolor}%
\pgftext[x=0.414395in, y=1.324557in, left, base]{\color{textcolor}\sffamily\fontsize{7.000000}{8.400000}\selectfont 0.6}%
\end{pgfscope}%
\begin{pgfscope}%
\pgfpathrectangle{\pgfqpoint{0.595000in}{0.663000in}}{\pgfqpoint{1.920732in}{1.216800in}}%
\pgfusepath{clip}%
\pgfsetrectcap%
\pgfsetroundjoin%
\pgfsetlinewidth{0.401500pt}%
\definecolor{currentstroke}{rgb}{0.901961,0.901961,0.901961}%
\pgfsetstrokecolor{currentstroke}%
\pgfsetdash{}{0pt}%
\pgfpathmoveto{\pgfqpoint{0.595000in}{1.590086in}}%
\pgfpathlineto{\pgfqpoint{2.515732in}{1.590086in}}%
\pgfusepath{stroke}%
\end{pgfscope}%
\begin{pgfscope}%
\definecolor{textcolor}{rgb}{0.101961,0.101961,0.101961}%
\pgfsetstrokecolor{textcolor}%
\pgfsetfillcolor{textcolor}%
\pgftext[x=0.414395in, y=1.556328in, left, base]{\color{textcolor}\sffamily\fontsize{7.000000}{8.400000}\selectfont 0.8}%
\end{pgfscope}%
\begin{pgfscope}%
\pgfpathrectangle{\pgfqpoint{0.595000in}{0.663000in}}{\pgfqpoint{1.920732in}{1.216800in}}%
\pgfusepath{clip}%
\pgfsetrectcap%
\pgfsetroundjoin%
\pgfsetlinewidth{0.401500pt}%
\definecolor{currentstroke}{rgb}{0.901961,0.901961,0.901961}%
\pgfsetstrokecolor{currentstroke}%
\pgfsetdash{}{0pt}%
\pgfpathmoveto{\pgfqpoint{0.595000in}{1.821857in}}%
\pgfpathlineto{\pgfqpoint{2.515732in}{1.821857in}}%
\pgfusepath{stroke}%
\end{pgfscope}%
\begin{pgfscope}%
\definecolor{textcolor}{rgb}{0.101961,0.101961,0.101961}%
\pgfsetstrokecolor{textcolor}%
\pgfsetfillcolor{textcolor}%
\pgftext[x=0.414395in, y=1.788100in, left, base]{\color{textcolor}\sffamily\fontsize{7.000000}{8.400000}\selectfont 1.0}%
\end{pgfscope}%
\begin{pgfscope}%
\definecolor{textcolor}{rgb}{0.000000,0.000000,0.000000}%
\pgfsetstrokecolor{textcolor}%
\pgfsetfillcolor{textcolor}%
\pgftext[x=0.358839in,y=1.271400in,,bottom,rotate=90.000000]{\color{textcolor}\sffamily\fontsize{7.500000}{9.000000}\selectfont on-time rate}%
\end{pgfscope}%
\begin{pgfscope}%
\pgfpathrectangle{\pgfqpoint{0.595000in}{0.663000in}}{\pgfqpoint{1.920732in}{1.216800in}}%
\pgfusepath{clip}%
\pgfsetbuttcap%
\pgfsetmiterjoin%
\definecolor{currentfill}{rgb}{0.615686,0.705882,0.807843}%
\pgfsetfillcolor{currentfill}%
\pgfsetfillopacity{0.450000}%
\pgfsetlinewidth{0.000000pt}%
\definecolor{currentstroke}{rgb}{0.000000,0.000000,0.000000}%
\pgfsetstrokecolor{currentstroke}%
\pgfsetstrokeopacity{0.450000}%
\pgfsetdash{}{0pt}%
\pgfpathmoveto{\pgfqpoint{0.695610in}{0.663000in}}%
\pgfpathlineto{\pgfqpoint{0.860244in}{0.663000in}}%
\pgfpathlineto{\pgfqpoint{0.860244in}{1.508917in}}%
\pgfpathlineto{\pgfqpoint{0.695610in}{1.508917in}}%
\pgfpathlineto{\pgfqpoint{0.695610in}{0.663000in}}%
\pgfpathclose%
\pgfusepath{fill}%
\end{pgfscope}%
\begin{pgfscope}%
\pgfpathrectangle{\pgfqpoint{0.595000in}{0.663000in}}{\pgfqpoint{1.920732in}{1.216800in}}%
\pgfusepath{clip}%
\pgfsetbuttcap%
\pgfsetmiterjoin%
\definecolor{currentfill}{rgb}{0.615686,0.705882,0.807843}%
\pgfsetfillcolor{currentfill}%
\pgfsetlinewidth{0.000000pt}%
\definecolor{currentstroke}{rgb}{0.000000,0.000000,0.000000}%
\pgfsetstrokecolor{currentstroke}%
\pgfsetstrokeopacity{0.000000}%
\pgfsetdash{}{0pt}%
\pgfpathmoveto{\pgfqpoint{0.878537in}{0.663000in}}%
\pgfpathlineto{\pgfqpoint{1.043171in}{0.663000in}}%
\pgfpathlineto{\pgfqpoint{1.043171in}{1.725286in}}%
\pgfpathlineto{\pgfqpoint{0.878537in}{1.725286in}}%
\pgfpathlineto{\pgfqpoint{0.878537in}{0.663000in}}%
\pgfpathclose%
\pgfusepath{fill}%
\end{pgfscope}%
\begin{pgfscope}%
\pgfpathrectangle{\pgfqpoint{0.595000in}{0.663000in}}{\pgfqpoint{1.920732in}{1.216800in}}%
\pgfusepath{clip}%
\pgfsetbuttcap%
\pgfsetmiterjoin%
\definecolor{currentfill}{rgb}{0.435294,0.576471,0.729412}%
\pgfsetfillcolor{currentfill}%
\pgfsetfillopacity{0.450000}%
\pgfsetlinewidth{0.000000pt}%
\definecolor{currentstroke}{rgb}{0.000000,0.000000,0.000000}%
\pgfsetstrokecolor{currentstroke}%
\pgfsetstrokeopacity{0.450000}%
\pgfsetdash{}{0pt}%
\pgfpathmoveto{\pgfqpoint{1.152927in}{0.663000in}}%
\pgfpathlineto{\pgfqpoint{1.317561in}{0.663000in}}%
\pgfpathlineto{\pgfqpoint{1.317561in}{1.494248in}}%
\pgfpathlineto{\pgfqpoint{1.152927in}{1.494248in}}%
\pgfpathlineto{\pgfqpoint{1.152927in}{0.663000in}}%
\pgfpathclose%
\pgfusepath{fill}%
\end{pgfscope}%
\begin{pgfscope}%
\pgfpathrectangle{\pgfqpoint{0.595000in}{0.663000in}}{\pgfqpoint{1.920732in}{1.216800in}}%
\pgfusepath{clip}%
\pgfsetbuttcap%
\pgfsetmiterjoin%
\definecolor{currentfill}{rgb}{0.435294,0.576471,0.729412}%
\pgfsetfillcolor{currentfill}%
\pgfsetlinewidth{0.000000pt}%
\definecolor{currentstroke}{rgb}{0.000000,0.000000,0.000000}%
\pgfsetstrokecolor{currentstroke}%
\pgfsetstrokeopacity{0.000000}%
\pgfsetdash{}{0pt}%
\pgfpathmoveto{\pgfqpoint{1.335854in}{0.663000in}}%
\pgfpathlineto{\pgfqpoint{1.500488in}{0.663000in}}%
\pgfpathlineto{\pgfqpoint{1.500488in}{1.483857in}}%
\pgfpathlineto{\pgfqpoint{1.335854in}{1.483857in}}%
\pgfpathlineto{\pgfqpoint{1.335854in}{0.663000in}}%
\pgfpathclose%
\pgfusepath{fill}%
\end{pgfscope}%
\begin{pgfscope}%
\pgfpathrectangle{\pgfqpoint{0.595000in}{0.663000in}}{\pgfqpoint{1.920732in}{1.216800in}}%
\pgfusepath{clip}%
\pgfsetbuttcap%
\pgfsetmiterjoin%
\definecolor{currentfill}{rgb}{0.156863,0.356863,0.568627}%
\pgfsetfillcolor{currentfill}%
\pgfsetfillopacity{0.450000}%
\pgfsetlinewidth{0.000000pt}%
\definecolor{currentstroke}{rgb}{0.000000,0.000000,0.000000}%
\pgfsetstrokecolor{currentstroke}%
\pgfsetstrokeopacity{0.450000}%
\pgfsetdash{}{0pt}%
\pgfpathmoveto{\pgfqpoint{1.610244in}{0.663000in}}%
\pgfpathlineto{\pgfqpoint{1.774878in}{0.663000in}}%
\pgfpathlineto{\pgfqpoint{1.774878in}{1.596931in}}%
\pgfpathlineto{\pgfqpoint{1.610244in}{1.596931in}}%
\pgfpathlineto{\pgfqpoint{1.610244in}{0.663000in}}%
\pgfpathclose%
\pgfusepath{fill}%
\end{pgfscope}%
\begin{pgfscope}%
\pgfpathrectangle{\pgfqpoint{0.595000in}{0.663000in}}{\pgfqpoint{1.920732in}{1.216800in}}%
\pgfusepath{clip}%
\pgfsetbuttcap%
\pgfsetmiterjoin%
\definecolor{currentfill}{rgb}{0.156863,0.356863,0.568627}%
\pgfsetfillcolor{currentfill}%
\pgfsetlinewidth{0.000000pt}%
\definecolor{currentstroke}{rgb}{0.000000,0.000000,0.000000}%
\pgfsetstrokecolor{currentstroke}%
\pgfsetstrokeopacity{0.000000}%
\pgfsetdash{}{0pt}%
\pgfpathmoveto{\pgfqpoint{1.793171in}{0.663000in}}%
\pgfpathlineto{\pgfqpoint{1.957805in}{0.663000in}}%
\pgfpathlineto{\pgfqpoint{1.957805in}{1.677000in}}%
\pgfpathlineto{\pgfqpoint{1.793171in}{1.677000in}}%
\pgfpathlineto{\pgfqpoint{1.793171in}{0.663000in}}%
\pgfpathclose%
\pgfusepath{fill}%
\end{pgfscope}%
\begin{pgfscope}%
\pgfpathrectangle{\pgfqpoint{0.595000in}{0.663000in}}{\pgfqpoint{1.920732in}{1.216800in}}%
\pgfusepath{clip}%
\pgfsetbuttcap%
\pgfsetmiterjoin%
\definecolor{currentfill}{rgb}{0.062745,0.160784,0.290196}%
\pgfsetfillcolor{currentfill}%
\pgfsetfillopacity{0.450000}%
\pgfsetlinewidth{0.000000pt}%
\definecolor{currentstroke}{rgb}{0.000000,0.000000,0.000000}%
\pgfsetstrokecolor{currentstroke}%
\pgfsetstrokeopacity{0.450000}%
\pgfsetdash{}{0pt}%
\pgfpathmoveto{\pgfqpoint{2.067561in}{0.663000in}}%
\pgfpathlineto{\pgfqpoint{2.232195in}{0.663000in}}%
\pgfpathlineto{\pgfqpoint{2.232195in}{1.528476in}}%
\pgfpathlineto{\pgfqpoint{2.067561in}{1.528476in}}%
\pgfpathlineto{\pgfqpoint{2.067561in}{0.663000in}}%
\pgfpathclose%
\pgfusepath{fill}%
\end{pgfscope}%
\begin{pgfscope}%
\pgfpathrectangle{\pgfqpoint{0.595000in}{0.663000in}}{\pgfqpoint{1.920732in}{1.216800in}}%
\pgfusepath{clip}%
\pgfsetbuttcap%
\pgfsetmiterjoin%
\definecolor{currentfill}{rgb}{0.062745,0.160784,0.290196}%
\pgfsetfillcolor{currentfill}%
\pgfsetlinewidth{0.000000pt}%
\definecolor{currentstroke}{rgb}{0.000000,0.000000,0.000000}%
\pgfsetstrokecolor{currentstroke}%
\pgfsetstrokeopacity{0.000000}%
\pgfsetdash{}{0pt}%
\pgfpathmoveto{\pgfqpoint{2.250488in}{0.663000in}}%
\pgfpathlineto{\pgfqpoint{2.415122in}{0.663000in}}%
\pgfpathlineto{\pgfqpoint{2.415122in}{1.677000in}}%
\pgfpathlineto{\pgfqpoint{2.250488in}{1.677000in}}%
\pgfpathlineto{\pgfqpoint{2.250488in}{0.663000in}}%
\pgfpathclose%
\pgfusepath{fill}%
\end{pgfscope}%
\begin{pgfscope}%
\pgfpathrectangle{\pgfqpoint{0.595000in}{0.663000in}}{\pgfqpoint{1.920732in}{1.216800in}}%
\pgfusepath{clip}%
\pgfsetrectcap%
\pgfsetroundjoin%
\pgfsetlinewidth{0.903375pt}%
\definecolor{currentstroke}{rgb}{0.101961,0.101961,0.101961}%
\pgfsetstrokecolor{currentstroke}%
\pgfsetdash{}{0pt}%
\pgfpathmoveto{\pgfqpoint{0.878537in}{1.725286in}}%
\pgfpathlineto{\pgfqpoint{1.043171in}{1.725286in}}%
\pgfusepath{stroke}%
\end{pgfscope}%
\begin{pgfscope}%
\pgfpathrectangle{\pgfqpoint{0.595000in}{0.663000in}}{\pgfqpoint{1.920732in}{1.216800in}}%
\pgfusepath{clip}%
\pgfsetrectcap%
\pgfsetroundjoin%
\pgfsetlinewidth{0.903375pt}%
\definecolor{currentstroke}{rgb}{0.101961,0.101961,0.101961}%
\pgfsetstrokecolor{currentstroke}%
\pgfsetdash{}{0pt}%
\pgfpathmoveto{\pgfqpoint{0.695610in}{1.469799in}}%
\pgfpathlineto{\pgfqpoint{0.860244in}{1.469799in}}%
\pgfusepath{stroke}%
\end{pgfscope}%
\begin{pgfscope}%
\pgfpathrectangle{\pgfqpoint{0.595000in}{0.663000in}}{\pgfqpoint{1.920732in}{1.216800in}}%
\pgfusepath{clip}%
\pgfsetrectcap%
\pgfsetroundjoin%
\pgfsetlinewidth{0.903375pt}%
\definecolor{currentstroke}{rgb}{0.101961,0.101961,0.101961}%
\pgfsetstrokecolor{currentstroke}%
\pgfsetdash{}{0pt}%
\pgfpathmoveto{\pgfqpoint{1.793171in}{1.821857in}}%
\pgfpathlineto{\pgfqpoint{1.957805in}{1.821857in}}%
\pgfusepath{stroke}%
\end{pgfscope}%
\begin{pgfscope}%
\pgfpathrectangle{\pgfqpoint{0.595000in}{0.663000in}}{\pgfqpoint{1.920732in}{1.216800in}}%
\pgfusepath{clip}%
\pgfsetrectcap%
\pgfsetroundjoin%
\pgfsetlinewidth{0.903375pt}%
\definecolor{currentstroke}{rgb}{0.101961,0.101961,0.101961}%
\pgfsetstrokecolor{currentstroke}%
\pgfsetdash{}{0pt}%
\pgfpathmoveto{\pgfqpoint{1.610244in}{1.640939in}}%
\pgfpathlineto{\pgfqpoint{1.774878in}{1.640939in}}%
\pgfusepath{stroke}%
\end{pgfscope}%
\begin{pgfscope}%
\pgfsetrectcap%
\pgfsetmiterjoin%
\pgfsetlinewidth{0.803000pt}%
\definecolor{currentstroke}{rgb}{0.101961,0.101961,0.101961}%
\pgfsetstrokecolor{currentstroke}%
\pgfsetdash{}{0pt}%
\pgfpathmoveto{\pgfqpoint{0.595000in}{0.663000in}}%
\pgfpathlineto{\pgfqpoint{0.595000in}{1.879800in}}%
\pgfusepath{stroke}%
\end{pgfscope}%
\begin{pgfscope}%
\pgfsetrectcap%
\pgfsetmiterjoin%
\pgfsetlinewidth{0.803000pt}%
\definecolor{currentstroke}{rgb}{0.101961,0.101961,0.101961}%
\pgfsetstrokecolor{currentstroke}%
\pgfsetdash{}{0pt}%
\pgfpathmoveto{\pgfqpoint{0.595000in}{0.663000in}}%
\pgfpathlineto{\pgfqpoint{2.515732in}{0.663000in}}%
\pgfusepath{stroke}%
\end{pgfscope}%
\begin{pgfscope}%
\pgfsetbuttcap%
\pgfsetmiterjoin%
\definecolor{currentfill}{rgb}{1.000000,1.000000,1.000000}%
\pgfsetfillcolor{currentfill}%
\pgfsetlinewidth{0.000000pt}%
\definecolor{currentstroke}{rgb}{0.000000,0.000000,0.000000}%
\pgfsetstrokecolor{currentstroke}%
\pgfsetstrokeopacity{0.000000}%
\pgfsetdash{}{0pt}%
\pgfpathmoveto{\pgfqpoint{2.784634in}{0.663000in}}%
\pgfpathlineto{\pgfqpoint{4.705366in}{0.663000in}}%
\pgfpathlineto{\pgfqpoint{4.705366in}{1.879800in}}%
\pgfpathlineto{\pgfqpoint{2.784634in}{1.879800in}}%
\pgfpathlineto{\pgfqpoint{2.784634in}{0.663000in}}%
\pgfpathclose%
\pgfusepath{fill}%
\end{pgfscope}%
\begin{pgfscope}%
\definecolor{textcolor}{rgb}{0.101961,0.101961,0.101961}%
\pgfsetstrokecolor{textcolor}%
\pgfsetfillcolor{textcolor}%
\pgftext[x=3.059024in,y=0.614389in,,top]{\color{textcolor}\sffamily\fontsize{7.000000}{8.400000}\selectfont rung 3}%
\end{pgfscope}%
\begin{pgfscope}%
\definecolor{textcolor}{rgb}{0.101961,0.101961,0.101961}%
\pgfsetstrokecolor{textcolor}%
\pgfsetfillcolor{textcolor}%
\pgftext[x=3.516341in,y=0.614389in,,top]{\color{textcolor}\sffamily\fontsize{7.000000}{8.400000}\selectfont eager}%
\end{pgfscope}%
\begin{pgfscope}%
\definecolor{textcolor}{rgb}{0.101961,0.101961,0.101961}%
\pgfsetstrokecolor{textcolor}%
\pgfsetfillcolor{textcolor}%
\pgftext[x=3.973659in,y=0.614389in,,top]{\color{textcolor}\sffamily\fontsize{7.000000}{8.400000}\selectfont no S2}%
\end{pgfscope}%
\begin{pgfscope}%
\definecolor{textcolor}{rgb}{0.101961,0.101961,0.101961}%
\pgfsetstrokecolor{textcolor}%
\pgfsetfillcolor{textcolor}%
\pgftext[x=4.430976in,y=0.614389in,,top]{\color{textcolor}\sffamily\fontsize{7.000000}{8.400000}\selectfont + S2}%
\end{pgfscope}%
\begin{pgfscope}%
\pgfpathrectangle{\pgfqpoint{2.784634in}{0.663000in}}{\pgfqpoint{1.920732in}{1.216800in}}%
\pgfusepath{clip}%
\pgfsetrectcap%
\pgfsetroundjoin%
\pgfsetlinewidth{0.401500pt}%
\definecolor{currentstroke}{rgb}{0.901961,0.901961,0.901961}%
\pgfsetstrokecolor{currentstroke}%
\pgfsetdash{}{0pt}%
\pgfpathmoveto{\pgfqpoint{2.784634in}{0.663000in}}%
\pgfpathlineto{\pgfqpoint{4.705366in}{0.663000in}}%
\pgfusepath{stroke}%
\end{pgfscope}%
\begin{pgfscope}%
\pgfpathrectangle{\pgfqpoint{2.784634in}{0.663000in}}{\pgfqpoint{1.920732in}{1.216800in}}%
\pgfusepath{clip}%
\pgfsetrectcap%
\pgfsetroundjoin%
\pgfsetlinewidth{0.401500pt}%
\definecolor{currentstroke}{rgb}{0.901961,0.901961,0.901961}%
\pgfsetstrokecolor{currentstroke}%
\pgfsetdash{}{0pt}%
\pgfpathmoveto{\pgfqpoint{2.784634in}{0.894771in}}%
\pgfpathlineto{\pgfqpoint{4.705366in}{0.894771in}}%
\pgfusepath{stroke}%
\end{pgfscope}%
\begin{pgfscope}%
\pgfpathrectangle{\pgfqpoint{2.784634in}{0.663000in}}{\pgfqpoint{1.920732in}{1.216800in}}%
\pgfusepath{clip}%
\pgfsetrectcap%
\pgfsetroundjoin%
\pgfsetlinewidth{0.401500pt}%
\definecolor{currentstroke}{rgb}{0.901961,0.901961,0.901961}%
\pgfsetstrokecolor{currentstroke}%
\pgfsetdash{}{0pt}%
\pgfpathmoveto{\pgfqpoint{2.784634in}{1.126543in}}%
\pgfpathlineto{\pgfqpoint{4.705366in}{1.126543in}}%
\pgfusepath{stroke}%
\end{pgfscope}%
\begin{pgfscope}%
\pgfpathrectangle{\pgfqpoint{2.784634in}{0.663000in}}{\pgfqpoint{1.920732in}{1.216800in}}%
\pgfusepath{clip}%
\pgfsetrectcap%
\pgfsetroundjoin%
\pgfsetlinewidth{0.401500pt}%
\definecolor{currentstroke}{rgb}{0.901961,0.901961,0.901961}%
\pgfsetstrokecolor{currentstroke}%
\pgfsetdash{}{0pt}%
\pgfpathmoveto{\pgfqpoint{2.784634in}{1.358314in}}%
\pgfpathlineto{\pgfqpoint{4.705366in}{1.358314in}}%
\pgfusepath{stroke}%
\end{pgfscope}%
\begin{pgfscope}%
\pgfpathrectangle{\pgfqpoint{2.784634in}{0.663000in}}{\pgfqpoint{1.920732in}{1.216800in}}%
\pgfusepath{clip}%
\pgfsetrectcap%
\pgfsetroundjoin%
\pgfsetlinewidth{0.401500pt}%
\definecolor{currentstroke}{rgb}{0.901961,0.901961,0.901961}%
\pgfsetstrokecolor{currentstroke}%
\pgfsetdash{}{0pt}%
\pgfpathmoveto{\pgfqpoint{2.784634in}{1.590086in}}%
\pgfpathlineto{\pgfqpoint{4.705366in}{1.590086in}}%
\pgfusepath{stroke}%
\end{pgfscope}%
\begin{pgfscope}%
\pgfpathrectangle{\pgfqpoint{2.784634in}{0.663000in}}{\pgfqpoint{1.920732in}{1.216800in}}%
\pgfusepath{clip}%
\pgfsetrectcap%
\pgfsetroundjoin%
\pgfsetlinewidth{0.401500pt}%
\definecolor{currentstroke}{rgb}{0.901961,0.901961,0.901961}%
\pgfsetstrokecolor{currentstroke}%
\pgfsetdash{}{0pt}%
\pgfpathmoveto{\pgfqpoint{2.784634in}{1.821857in}}%
\pgfpathlineto{\pgfqpoint{4.705366in}{1.821857in}}%
\pgfusepath{stroke}%
\end{pgfscope}%
\begin{pgfscope}%
\pgfpathrectangle{\pgfqpoint{2.784634in}{0.663000in}}{\pgfqpoint{1.920732in}{1.216800in}}%
\pgfusepath{clip}%
\pgfsetbuttcap%
\pgfsetmiterjoin%
\definecolor{currentfill}{rgb}{0.615686,0.705882,0.807843}%
\pgfsetfillcolor{currentfill}%
\pgfsetfillopacity{0.450000}%
\pgfsetlinewidth{0.000000pt}%
\definecolor{currentstroke}{rgb}{0.000000,0.000000,0.000000}%
\pgfsetstrokecolor{currentstroke}%
\pgfsetstrokeopacity{0.450000}%
\pgfsetdash{}{0pt}%
\pgfpathmoveto{\pgfqpoint{2.885244in}{0.663000in}}%
\pgfpathlineto{\pgfqpoint{3.049878in}{0.663000in}}%
\pgfpathlineto{\pgfqpoint{3.049878in}{1.121629in}}%
\pgfpathlineto{\pgfqpoint{2.885244in}{1.121629in}}%
\pgfpathlineto{\pgfqpoint{2.885244in}{0.663000in}}%
\pgfpathclose%
\pgfusepath{fill}%
\end{pgfscope}%
\begin{pgfscope}%
\pgfpathrectangle{\pgfqpoint{2.784634in}{0.663000in}}{\pgfqpoint{1.920732in}{1.216800in}}%
\pgfusepath{clip}%
\pgfsetbuttcap%
\pgfsetmiterjoin%
\definecolor{currentfill}{rgb}{0.615686,0.705882,0.807843}%
\pgfsetfillcolor{currentfill}%
\pgfsetlinewidth{0.000000pt}%
\definecolor{currentstroke}{rgb}{0.000000,0.000000,0.000000}%
\pgfsetstrokecolor{currentstroke}%
\pgfsetstrokeopacity{0.000000}%
\pgfsetdash{}{0pt}%
\pgfpathmoveto{\pgfqpoint{3.068171in}{0.663000in}}%
\pgfpathlineto{\pgfqpoint{3.232805in}{0.663000in}}%
\pgfpathlineto{\pgfqpoint{3.232805in}{1.599000in}}%
\pgfpathlineto{\pgfqpoint{3.068171in}{1.599000in}}%
\pgfpathlineto{\pgfqpoint{3.068171in}{0.663000in}}%
\pgfpathclose%
\pgfusepath{fill}%
\end{pgfscope}%
\begin{pgfscope}%
\pgfpathrectangle{\pgfqpoint{2.784634in}{0.663000in}}{\pgfqpoint{1.920732in}{1.216800in}}%
\pgfusepath{clip}%
\pgfsetbuttcap%
\pgfsetmiterjoin%
\definecolor{currentfill}{rgb}{0.435294,0.576471,0.729412}%
\pgfsetfillcolor{currentfill}%
\pgfsetfillopacity{0.450000}%
\pgfsetlinewidth{0.000000pt}%
\definecolor{currentstroke}{rgb}{0.000000,0.000000,0.000000}%
\pgfsetstrokecolor{currentstroke}%
\pgfsetstrokeopacity{0.450000}%
\pgfsetdash{}{0pt}%
\pgfpathmoveto{\pgfqpoint{3.342561in}{0.663000in}}%
\pgfpathlineto{\pgfqpoint{3.507195in}{0.663000in}}%
\pgfpathlineto{\pgfqpoint{3.507195in}{1.174863in}}%
\pgfpathlineto{\pgfqpoint{3.342561in}{1.174863in}}%
\pgfpathlineto{\pgfqpoint{3.342561in}{0.663000in}}%
\pgfpathclose%
\pgfusepath{fill}%
\end{pgfscope}%
\begin{pgfscope}%
\pgfpathrectangle{\pgfqpoint{2.784634in}{0.663000in}}{\pgfqpoint{1.920732in}{1.216800in}}%
\pgfusepath{clip}%
\pgfsetbuttcap%
\pgfsetmiterjoin%
\definecolor{currentfill}{rgb}{0.435294,0.576471,0.729412}%
\pgfsetfillcolor{currentfill}%
\pgfsetlinewidth{0.000000pt}%
\definecolor{currentstroke}{rgb}{0.000000,0.000000,0.000000}%
\pgfsetstrokecolor{currentstroke}%
\pgfsetstrokeopacity{0.000000}%
\pgfsetdash{}{0pt}%
\pgfpathmoveto{\pgfqpoint{3.525488in}{0.663000in}}%
\pgfpathlineto{\pgfqpoint{3.690122in}{0.663000in}}%
\pgfpathlineto{\pgfqpoint{3.690122in}{1.420714in}}%
\pgfpathlineto{\pgfqpoint{3.525488in}{1.420714in}}%
\pgfpathlineto{\pgfqpoint{3.525488in}{0.663000in}}%
\pgfpathclose%
\pgfusepath{fill}%
\end{pgfscope}%
\begin{pgfscope}%
\pgfpathrectangle{\pgfqpoint{2.784634in}{0.663000in}}{\pgfqpoint{1.920732in}{1.216800in}}%
\pgfusepath{clip}%
\pgfsetbuttcap%
\pgfsetmiterjoin%
\definecolor{currentfill}{rgb}{0.156863,0.356863,0.568627}%
\pgfsetfillcolor{currentfill}%
\pgfsetfillopacity{0.450000}%
\pgfsetlinewidth{0.000000pt}%
\definecolor{currentstroke}{rgb}{0.000000,0.000000,0.000000}%
\pgfsetstrokecolor{currentstroke}%
\pgfsetstrokeopacity{0.450000}%
\pgfsetdash{}{0pt}%
\pgfpathmoveto{\pgfqpoint{3.799878in}{0.663000in}}%
\pgfpathlineto{\pgfqpoint{3.964512in}{0.663000in}}%
\pgfpathlineto{\pgfqpoint{3.964512in}{1.281330in}}%
\pgfpathlineto{\pgfqpoint{3.799878in}{1.281330in}}%
\pgfpathlineto{\pgfqpoint{3.799878in}{0.663000in}}%
\pgfpathclose%
\pgfusepath{fill}%
\end{pgfscope}%
\begin{pgfscope}%
\pgfpathrectangle{\pgfqpoint{2.784634in}{0.663000in}}{\pgfqpoint{1.920732in}{1.216800in}}%
\pgfusepath{clip}%
\pgfsetbuttcap%
\pgfsetmiterjoin%
\definecolor{currentfill}{rgb}{0.156863,0.356863,0.568627}%
\pgfsetfillcolor{currentfill}%
\pgfsetlinewidth{0.000000pt}%
\definecolor{currentstroke}{rgb}{0.000000,0.000000,0.000000}%
\pgfsetstrokecolor{currentstroke}%
\pgfsetstrokeopacity{0.000000}%
\pgfsetdash{}{0pt}%
\pgfpathmoveto{\pgfqpoint{3.982805in}{0.663000in}}%
\pgfpathlineto{\pgfqpoint{4.147439in}{0.663000in}}%
\pgfpathlineto{\pgfqpoint{4.147439in}{1.599000in}}%
\pgfpathlineto{\pgfqpoint{3.982805in}{1.599000in}}%
\pgfpathlineto{\pgfqpoint{3.982805in}{0.663000in}}%
\pgfpathclose%
\pgfusepath{fill}%
\end{pgfscope}%
\begin{pgfscope}%
\pgfpathrectangle{\pgfqpoint{2.784634in}{0.663000in}}{\pgfqpoint{1.920732in}{1.216800in}}%
\pgfusepath{clip}%
\pgfsetbuttcap%
\pgfsetmiterjoin%
\definecolor{currentfill}{rgb}{0.062745,0.160784,0.290196}%
\pgfsetfillcolor{currentfill}%
\pgfsetfillopacity{0.450000}%
\pgfsetlinewidth{0.000000pt}%
\definecolor{currentstroke}{rgb}{0.000000,0.000000,0.000000}%
\pgfsetstrokecolor{currentstroke}%
\pgfsetstrokeopacity{0.450000}%
\pgfsetdash{}{0pt}%
\pgfpathmoveto{\pgfqpoint{4.257195in}{0.663000in}}%
\pgfpathlineto{\pgfqpoint{4.421829in}{0.663000in}}%
\pgfpathlineto{\pgfqpoint{4.421829in}{1.224002in}}%
\pgfpathlineto{\pgfqpoint{4.257195in}{1.224002in}}%
\pgfpathlineto{\pgfqpoint{4.257195in}{0.663000in}}%
\pgfpathclose%
\pgfusepath{fill}%
\end{pgfscope}%
\begin{pgfscope}%
\pgfpathrectangle{\pgfqpoint{2.784634in}{0.663000in}}{\pgfqpoint{1.920732in}{1.216800in}}%
\pgfusepath{clip}%
\pgfsetbuttcap%
\pgfsetmiterjoin%
\definecolor{currentfill}{rgb}{0.062745,0.160784,0.290196}%
\pgfsetfillcolor{currentfill}%
\pgfsetlinewidth{0.000000pt}%
\definecolor{currentstroke}{rgb}{0.000000,0.000000,0.000000}%
\pgfsetstrokecolor{currentstroke}%
\pgfsetstrokeopacity{0.000000}%
\pgfsetdash{}{0pt}%
\pgfpathmoveto{\pgfqpoint{4.440122in}{0.663000in}}%
\pgfpathlineto{\pgfqpoint{4.604756in}{0.663000in}}%
\pgfpathlineto{\pgfqpoint{4.604756in}{1.643571in}}%
\pgfpathlineto{\pgfqpoint{4.440122in}{1.643571in}}%
\pgfpathlineto{\pgfqpoint{4.440122in}{0.663000in}}%
\pgfpathclose%
\pgfusepath{fill}%
\end{pgfscope}%
\begin{pgfscope}%
\pgfpathrectangle{\pgfqpoint{2.784634in}{0.663000in}}{\pgfqpoint{1.920732in}{1.216800in}}%
\pgfusepath{clip}%
\pgfsetrectcap%
\pgfsetroundjoin%
\pgfsetlinewidth{0.903375pt}%
\definecolor{currentstroke}{rgb}{0.101961,0.101961,0.101961}%
\pgfsetstrokecolor{currentstroke}%
\pgfsetdash{}{0pt}%
\pgfpathmoveto{\pgfqpoint{3.068171in}{1.643571in}}%
\pgfpathlineto{\pgfqpoint{3.232805in}{1.643571in}}%
\pgfusepath{stroke}%
\end{pgfscope}%
\begin{pgfscope}%
\pgfpathrectangle{\pgfqpoint{2.784634in}{0.663000in}}{\pgfqpoint{1.920732in}{1.216800in}}%
\pgfusepath{clip}%
\pgfsetrectcap%
\pgfsetroundjoin%
\pgfsetlinewidth{0.903375pt}%
\definecolor{currentstroke}{rgb}{0.101961,0.101961,0.101961}%
\pgfsetstrokecolor{currentstroke}%
\pgfsetdash{}{0pt}%
\pgfpathmoveto{\pgfqpoint{2.885244in}{1.174863in}}%
\pgfpathlineto{\pgfqpoint{3.049878in}{1.174863in}}%
\pgfusepath{stroke}%
\end{pgfscope}%
\begin{pgfscope}%
\pgfpathrectangle{\pgfqpoint{2.784634in}{0.663000in}}{\pgfqpoint{1.920732in}{1.216800in}}%
\pgfusepath{clip}%
\pgfsetrectcap%
\pgfsetroundjoin%
\pgfsetlinewidth{0.903375pt}%
\definecolor{currentstroke}{rgb}{0.101961,0.101961,0.101961}%
\pgfsetstrokecolor{currentstroke}%
\pgfsetdash{}{0pt}%
\pgfpathmoveto{\pgfqpoint{3.982805in}{1.777286in}}%
\pgfpathlineto{\pgfqpoint{4.147439in}{1.777286in}}%
\pgfusepath{stroke}%
\end{pgfscope}%
\begin{pgfscope}%
\pgfpathrectangle{\pgfqpoint{2.784634in}{0.663000in}}{\pgfqpoint{1.920732in}{1.216800in}}%
\pgfusepath{clip}%
\pgfsetrectcap%
\pgfsetroundjoin%
\pgfsetlinewidth{0.903375pt}%
\definecolor{currentstroke}{rgb}{0.101961,0.101961,0.101961}%
\pgfsetstrokecolor{currentstroke}%
\pgfsetdash{}{0pt}%
\pgfpathmoveto{\pgfqpoint{3.799878in}{1.416462in}}%
\pgfpathlineto{\pgfqpoint{3.964512in}{1.416462in}}%
\pgfusepath{stroke}%
\end{pgfscope}%
\begin{pgfscope}%
\pgfsetrectcap%
\pgfsetmiterjoin%
\pgfsetlinewidth{0.803000pt}%
\definecolor{currentstroke}{rgb}{0.101961,0.101961,0.101961}%
\pgfsetstrokecolor{currentstroke}%
\pgfsetdash{}{0pt}%
\pgfpathmoveto{\pgfqpoint{2.784634in}{0.663000in}}%
\pgfpathlineto{\pgfqpoint{2.784634in}{1.879800in}}%
\pgfusepath{stroke}%
\end{pgfscope}%
\begin{pgfscope}%
\pgfsetrectcap%
\pgfsetmiterjoin%
\pgfsetlinewidth{0.803000pt}%
\definecolor{currentstroke}{rgb}{0.101961,0.101961,0.101961}%
\pgfsetstrokecolor{currentstroke}%
\pgfsetdash{}{0pt}%
\pgfpathmoveto{\pgfqpoint{2.784634in}{0.663000in}}%
\pgfpathlineto{\pgfqpoint{4.705366in}{0.663000in}}%
\pgfusepath{stroke}%
\end{pgfscope}%
\begin{pgfscope}%
\pgfsetbuttcap%
\pgfsetmiterjoin%
\definecolor{currentfill}{rgb}{1.000000,1.000000,1.000000}%
\pgfsetfillcolor{currentfill}%
\pgfsetlinewidth{0.000000pt}%
\definecolor{currentstroke}{rgb}{0.000000,0.000000,0.000000}%
\pgfsetstrokecolor{currentstroke}%
\pgfsetstrokeopacity{0.000000}%
\pgfsetdash{}{0pt}%
\pgfpathmoveto{\pgfqpoint{4.974268in}{0.663000in}}%
\pgfpathlineto{\pgfqpoint{6.895000in}{0.663000in}}%
\pgfpathlineto{\pgfqpoint{6.895000in}{1.879800in}}%
\pgfpathlineto{\pgfqpoint{4.974268in}{1.879800in}}%
\pgfpathlineto{\pgfqpoint{4.974268in}{0.663000in}}%
\pgfpathclose%
\pgfusepath{fill}%
\end{pgfscope}%
\begin{pgfscope}%
\definecolor{textcolor}{rgb}{0.101961,0.101961,0.101961}%
\pgfsetstrokecolor{textcolor}%
\pgfsetfillcolor{textcolor}%
\pgftext[x=5.248659in,y=0.614389in,,top]{\color{textcolor}\sffamily\fontsize{7.000000}{8.400000}\selectfont rung 3}%
\end{pgfscope}%
\begin{pgfscope}%
\definecolor{textcolor}{rgb}{0.101961,0.101961,0.101961}%
\pgfsetstrokecolor{textcolor}%
\pgfsetfillcolor{textcolor}%
\pgftext[x=5.705976in,y=0.614389in,,top]{\color{textcolor}\sffamily\fontsize{7.000000}{8.400000}\selectfont eager}%
\end{pgfscope}%
\begin{pgfscope}%
\definecolor{textcolor}{rgb}{0.101961,0.101961,0.101961}%
\pgfsetstrokecolor{textcolor}%
\pgfsetfillcolor{textcolor}%
\pgftext[x=6.163293in,y=0.614389in,,top]{\color{textcolor}\sffamily\fontsize{7.000000}{8.400000}\selectfont no S2}%
\end{pgfscope}%
\begin{pgfscope}%
\definecolor{textcolor}{rgb}{0.101961,0.101961,0.101961}%
\pgfsetstrokecolor{textcolor}%
\pgfsetfillcolor{textcolor}%
\pgftext[x=6.620610in,y=0.614389in,,top]{\color{textcolor}\sffamily\fontsize{7.000000}{8.400000}\selectfont + S2}%
\end{pgfscope}%
\begin{pgfscope}%
\pgfpathrectangle{\pgfqpoint{4.974268in}{0.663000in}}{\pgfqpoint{1.920732in}{1.216800in}}%
\pgfusepath{clip}%
\pgfsetrectcap%
\pgfsetroundjoin%
\pgfsetlinewidth{0.401500pt}%
\definecolor{currentstroke}{rgb}{0.901961,0.901961,0.901961}%
\pgfsetstrokecolor{currentstroke}%
\pgfsetdash{}{0pt}%
\pgfpathmoveto{\pgfqpoint{4.974268in}{0.663000in}}%
\pgfpathlineto{\pgfqpoint{6.895000in}{0.663000in}}%
\pgfusepath{stroke}%
\end{pgfscope}%
\begin{pgfscope}%
\pgfpathrectangle{\pgfqpoint{4.974268in}{0.663000in}}{\pgfqpoint{1.920732in}{1.216800in}}%
\pgfusepath{clip}%
\pgfsetrectcap%
\pgfsetroundjoin%
\pgfsetlinewidth{0.401500pt}%
\definecolor{currentstroke}{rgb}{0.901961,0.901961,0.901961}%
\pgfsetstrokecolor{currentstroke}%
\pgfsetdash{}{0pt}%
\pgfpathmoveto{\pgfqpoint{4.974268in}{0.894771in}}%
\pgfpathlineto{\pgfqpoint{6.895000in}{0.894771in}}%
\pgfusepath{stroke}%
\end{pgfscope}%
\begin{pgfscope}%
\pgfpathrectangle{\pgfqpoint{4.974268in}{0.663000in}}{\pgfqpoint{1.920732in}{1.216800in}}%
\pgfusepath{clip}%
\pgfsetrectcap%
\pgfsetroundjoin%
\pgfsetlinewidth{0.401500pt}%
\definecolor{currentstroke}{rgb}{0.901961,0.901961,0.901961}%
\pgfsetstrokecolor{currentstroke}%
\pgfsetdash{}{0pt}%
\pgfpathmoveto{\pgfqpoint{4.974268in}{1.126543in}}%
\pgfpathlineto{\pgfqpoint{6.895000in}{1.126543in}}%
\pgfusepath{stroke}%
\end{pgfscope}%
\begin{pgfscope}%
\pgfpathrectangle{\pgfqpoint{4.974268in}{0.663000in}}{\pgfqpoint{1.920732in}{1.216800in}}%
\pgfusepath{clip}%
\pgfsetrectcap%
\pgfsetroundjoin%
\pgfsetlinewidth{0.401500pt}%
\definecolor{currentstroke}{rgb}{0.901961,0.901961,0.901961}%
\pgfsetstrokecolor{currentstroke}%
\pgfsetdash{}{0pt}%
\pgfpathmoveto{\pgfqpoint{4.974268in}{1.358314in}}%
\pgfpathlineto{\pgfqpoint{6.895000in}{1.358314in}}%
\pgfusepath{stroke}%
\end{pgfscope}%
\begin{pgfscope}%
\pgfpathrectangle{\pgfqpoint{4.974268in}{0.663000in}}{\pgfqpoint{1.920732in}{1.216800in}}%
\pgfusepath{clip}%
\pgfsetrectcap%
\pgfsetroundjoin%
\pgfsetlinewidth{0.401500pt}%
\definecolor{currentstroke}{rgb}{0.901961,0.901961,0.901961}%
\pgfsetstrokecolor{currentstroke}%
\pgfsetdash{}{0pt}%
\pgfpathmoveto{\pgfqpoint{4.974268in}{1.590086in}}%
\pgfpathlineto{\pgfqpoint{6.895000in}{1.590086in}}%
\pgfusepath{stroke}%
\end{pgfscope}%
\begin{pgfscope}%
\pgfpathrectangle{\pgfqpoint{4.974268in}{0.663000in}}{\pgfqpoint{1.920732in}{1.216800in}}%
\pgfusepath{clip}%
\pgfsetrectcap%
\pgfsetroundjoin%
\pgfsetlinewidth{0.401500pt}%
\definecolor{currentstroke}{rgb}{0.901961,0.901961,0.901961}%
\pgfsetstrokecolor{currentstroke}%
\pgfsetdash{}{0pt}%
\pgfpathmoveto{\pgfqpoint{4.974268in}{1.821857in}}%
\pgfpathlineto{\pgfqpoint{6.895000in}{1.821857in}}%
\pgfusepath{stroke}%
\end{pgfscope}%
\begin{pgfscope}%
\pgfpathrectangle{\pgfqpoint{4.974268in}{0.663000in}}{\pgfqpoint{1.920732in}{1.216800in}}%
\pgfusepath{clip}%
\pgfsetbuttcap%
\pgfsetmiterjoin%
\definecolor{currentfill}{rgb}{0.615686,0.705882,0.807843}%
\pgfsetfillcolor{currentfill}%
\pgfsetfillopacity{0.450000}%
\pgfsetlinewidth{0.000000pt}%
\definecolor{currentstroke}{rgb}{0.000000,0.000000,0.000000}%
\pgfsetstrokecolor{currentstroke}%
\pgfsetstrokeopacity{0.450000}%
\pgfsetdash{}{0pt}%
\pgfpathmoveto{\pgfqpoint{5.074878in}{0.663000in}}%
\pgfpathlineto{\pgfqpoint{5.239512in}{0.663000in}}%
\pgfpathlineto{\pgfqpoint{5.239512in}{1.329970in}}%
\pgfpathlineto{\pgfqpoint{5.074878in}{1.329970in}}%
\pgfpathlineto{\pgfqpoint{5.074878in}{0.663000in}}%
\pgfpathclose%
\pgfusepath{fill}%
\end{pgfscope}%
\begin{pgfscope}%
\pgfpathrectangle{\pgfqpoint{4.974268in}{0.663000in}}{\pgfqpoint{1.920732in}{1.216800in}}%
\pgfusepath{clip}%
\pgfsetbuttcap%
\pgfsetmiterjoin%
\definecolor{currentfill}{rgb}{0.615686,0.705882,0.807843}%
\pgfsetfillcolor{currentfill}%
\pgfsetlinewidth{0.000000pt}%
\definecolor{currentstroke}{rgb}{0.000000,0.000000,0.000000}%
\pgfsetstrokecolor{currentstroke}%
\pgfsetstrokeopacity{0.000000}%
\pgfsetdash{}{0pt}%
\pgfpathmoveto{\pgfqpoint{5.257805in}{0.663000in}}%
\pgfpathlineto{\pgfqpoint{5.422439in}{0.663000in}}%
\pgfpathlineto{\pgfqpoint{5.422439in}{1.481616in}}%
\pgfpathlineto{\pgfqpoint{5.257805in}{1.481616in}}%
\pgfpathlineto{\pgfqpoint{5.257805in}{0.663000in}}%
\pgfpathclose%
\pgfusepath{fill}%
\end{pgfscope}%
\begin{pgfscope}%
\pgfpathrectangle{\pgfqpoint{4.974268in}{0.663000in}}{\pgfqpoint{1.920732in}{1.216800in}}%
\pgfusepath{clip}%
\pgfsetbuttcap%
\pgfsetmiterjoin%
\definecolor{currentfill}{rgb}{0.435294,0.576471,0.729412}%
\pgfsetfillcolor{currentfill}%
\pgfsetfillopacity{0.450000}%
\pgfsetlinewidth{0.000000pt}%
\definecolor{currentstroke}{rgb}{0.000000,0.000000,0.000000}%
\pgfsetstrokecolor{currentstroke}%
\pgfsetstrokeopacity{0.450000}%
\pgfsetdash{}{0pt}%
\pgfpathmoveto{\pgfqpoint{5.532195in}{0.663000in}}%
\pgfpathlineto{\pgfqpoint{5.696829in}{0.663000in}}%
\pgfpathlineto{\pgfqpoint{5.696829in}{1.304283in}}%
\pgfpathlineto{\pgfqpoint{5.532195in}{1.304283in}}%
\pgfpathlineto{\pgfqpoint{5.532195in}{0.663000in}}%
\pgfpathclose%
\pgfusepath{fill}%
\end{pgfscope}%
\begin{pgfscope}%
\pgfpathrectangle{\pgfqpoint{4.974268in}{0.663000in}}{\pgfqpoint{1.920732in}{1.216800in}}%
\pgfusepath{clip}%
\pgfsetbuttcap%
\pgfsetmiterjoin%
\definecolor{currentfill}{rgb}{0.435294,0.576471,0.729412}%
\pgfsetfillcolor{currentfill}%
\pgfsetlinewidth{0.000000pt}%
\definecolor{currentstroke}{rgb}{0.000000,0.000000,0.000000}%
\pgfsetstrokecolor{currentstroke}%
\pgfsetstrokeopacity{0.000000}%
\pgfsetdash{}{0pt}%
\pgfpathmoveto{\pgfqpoint{5.715122in}{0.663000in}}%
\pgfpathlineto{\pgfqpoint{5.879756in}{0.663000in}}%
\pgfpathlineto{\pgfqpoint{5.879756in}{1.052678in}}%
\pgfpathlineto{\pgfqpoint{5.715122in}{1.052678in}}%
\pgfpathlineto{\pgfqpoint{5.715122in}{0.663000in}}%
\pgfpathclose%
\pgfusepath{fill}%
\end{pgfscope}%
\begin{pgfscope}%
\pgfpathrectangle{\pgfqpoint{4.974268in}{0.663000in}}{\pgfqpoint{1.920732in}{1.216800in}}%
\pgfusepath{clip}%
\pgfsetbuttcap%
\pgfsetmiterjoin%
\definecolor{currentfill}{rgb}{0.156863,0.356863,0.568627}%
\pgfsetfillcolor{currentfill}%
\pgfsetfillopacity{0.450000}%
\pgfsetlinewidth{0.000000pt}%
\definecolor{currentstroke}{rgb}{0.000000,0.000000,0.000000}%
\pgfsetstrokecolor{currentstroke}%
\pgfsetstrokeopacity{0.450000}%
\pgfsetdash{}{0pt}%
\pgfpathmoveto{\pgfqpoint{5.989512in}{0.663000in}}%
\pgfpathlineto{\pgfqpoint{6.154146in}{0.663000in}}%
\pgfpathlineto{\pgfqpoint{6.154146in}{1.410869in}}%
\pgfpathlineto{\pgfqpoint{5.989512in}{1.410869in}}%
\pgfpathlineto{\pgfqpoint{5.989512in}{0.663000in}}%
\pgfpathclose%
\pgfusepath{fill}%
\end{pgfscope}%
\begin{pgfscope}%
\pgfpathrectangle{\pgfqpoint{4.974268in}{0.663000in}}{\pgfqpoint{1.920732in}{1.216800in}}%
\pgfusepath{clip}%
\pgfsetbuttcap%
\pgfsetmiterjoin%
\definecolor{currentfill}{rgb}{0.156863,0.356863,0.568627}%
\pgfsetfillcolor{currentfill}%
\pgfsetlinewidth{0.000000pt}%
\definecolor{currentstroke}{rgb}{0.000000,0.000000,0.000000}%
\pgfsetstrokecolor{currentstroke}%
\pgfsetstrokeopacity{0.000000}%
\pgfsetdash{}{0pt}%
\pgfpathmoveto{\pgfqpoint{6.172439in}{0.663000in}}%
\pgfpathlineto{\pgfqpoint{6.337073in}{0.663000in}}%
\pgfpathlineto{\pgfqpoint{6.337073in}{1.471437in}}%
\pgfpathlineto{\pgfqpoint{6.172439in}{1.471437in}}%
\pgfpathlineto{\pgfqpoint{6.172439in}{0.663000in}}%
\pgfpathclose%
\pgfusepath{fill}%
\end{pgfscope}%
\begin{pgfscope}%
\pgfpathrectangle{\pgfqpoint{4.974268in}{0.663000in}}{\pgfqpoint{1.920732in}{1.216800in}}%
\pgfusepath{clip}%
\pgfsetbuttcap%
\pgfsetmiterjoin%
\definecolor{currentfill}{rgb}{0.062745,0.160784,0.290196}%
\pgfsetfillcolor{currentfill}%
\pgfsetfillopacity{0.450000}%
\pgfsetlinewidth{0.000000pt}%
\definecolor{currentstroke}{rgb}{0.000000,0.000000,0.000000}%
\pgfsetstrokecolor{currentstroke}%
\pgfsetstrokeopacity{0.450000}%
\pgfsetdash{}{0pt}%
\pgfpathmoveto{\pgfqpoint{6.446829in}{0.663000in}}%
\pgfpathlineto{\pgfqpoint{6.611463in}{0.663000in}}%
\pgfpathlineto{\pgfqpoint{6.611463in}{1.336318in}}%
\pgfpathlineto{\pgfqpoint{6.446829in}{1.336318in}}%
\pgfpathlineto{\pgfqpoint{6.446829in}{0.663000in}}%
\pgfpathclose%
\pgfusepath{fill}%
\end{pgfscope}%
\begin{pgfscope}%
\pgfpathrectangle{\pgfqpoint{4.974268in}{0.663000in}}{\pgfqpoint{1.920732in}{1.216800in}}%
\pgfusepath{clip}%
\pgfsetbuttcap%
\pgfsetmiterjoin%
\definecolor{currentfill}{rgb}{0.062745,0.160784,0.290196}%
\pgfsetfillcolor{currentfill}%
\pgfsetlinewidth{0.000000pt}%
\definecolor{currentstroke}{rgb}{0.000000,0.000000,0.000000}%
\pgfsetstrokecolor{currentstroke}%
\pgfsetstrokeopacity{0.000000}%
\pgfsetdash{}{0pt}%
\pgfpathmoveto{\pgfqpoint{6.629756in}{0.663000in}}%
\pgfpathlineto{\pgfqpoint{6.794390in}{0.663000in}}%
\pgfpathlineto{\pgfqpoint{6.794390in}{1.485978in}}%
\pgfpathlineto{\pgfqpoint{6.629756in}{1.485978in}}%
\pgfpathlineto{\pgfqpoint{6.629756in}{0.663000in}}%
\pgfpathclose%
\pgfusepath{fill}%
\end{pgfscope}%
\begin{pgfscope}%
\pgfpathrectangle{\pgfqpoint{4.974268in}{0.663000in}}{\pgfqpoint{1.920732in}{1.216800in}}%
\pgfusepath{clip}%
\pgfsetrectcap%
\pgfsetroundjoin%
\pgfsetlinewidth{0.903375pt}%
\definecolor{currentstroke}{rgb}{0.101961,0.101961,0.101961}%
\pgfsetstrokecolor{currentstroke}%
\pgfsetdash{}{0pt}%
\pgfpathmoveto{\pgfqpoint{5.257805in}{1.517966in}}%
\pgfpathlineto{\pgfqpoint{5.422439in}{1.517966in}}%
\pgfusepath{stroke}%
\end{pgfscope}%
\begin{pgfscope}%
\pgfpathrectangle{\pgfqpoint{4.974268in}{0.663000in}}{\pgfqpoint{1.920732in}{1.216800in}}%
\pgfusepath{clip}%
\pgfsetrectcap%
\pgfsetroundjoin%
\pgfsetlinewidth{0.903375pt}%
\definecolor{currentstroke}{rgb}{0.101961,0.101961,0.101961}%
\pgfsetstrokecolor{currentstroke}%
\pgfsetdash{}{0pt}%
\pgfpathmoveto{\pgfqpoint{5.074878in}{1.216299in}}%
\pgfpathlineto{\pgfqpoint{5.239512in}{1.216299in}}%
\pgfusepath{stroke}%
\end{pgfscope}%
\begin{pgfscope}%
\pgfpathrectangle{\pgfqpoint{4.974268in}{0.663000in}}{\pgfqpoint{1.920732in}{1.216800in}}%
\pgfusepath{clip}%
\pgfsetrectcap%
\pgfsetroundjoin%
\pgfsetlinewidth{0.903375pt}%
\definecolor{currentstroke}{rgb}{0.101961,0.101961,0.101961}%
\pgfsetstrokecolor{currentstroke}%
\pgfsetdash{}{0pt}%
\pgfpathmoveto{\pgfqpoint{6.172439in}{1.669185in}}%
\pgfpathlineto{\pgfqpoint{6.337073in}{1.669185in}}%
\pgfusepath{stroke}%
\end{pgfscope}%
\begin{pgfscope}%
\pgfpathrectangle{\pgfqpoint{4.974268in}{0.663000in}}{\pgfqpoint{1.920732in}{1.216800in}}%
\pgfusepath{clip}%
\pgfsetrectcap%
\pgfsetroundjoin%
\pgfsetlinewidth{0.903375pt}%
\definecolor{currentstroke}{rgb}{0.101961,0.101961,0.101961}%
\pgfsetstrokecolor{currentstroke}%
\pgfsetdash{}{0pt}%
\pgfpathmoveto{\pgfqpoint{5.989512in}{1.451761in}}%
\pgfpathlineto{\pgfqpoint{6.154146in}{1.451761in}}%
\pgfusepath{stroke}%
\end{pgfscope}%
\begin{pgfscope}%
\pgfsetrectcap%
\pgfsetmiterjoin%
\pgfsetlinewidth{0.803000pt}%
\definecolor{currentstroke}{rgb}{0.101961,0.101961,0.101961}%
\pgfsetstrokecolor{currentstroke}%
\pgfsetdash{}{0pt}%
\pgfpathmoveto{\pgfqpoint{4.974268in}{0.663000in}}%
\pgfpathlineto{\pgfqpoint{4.974268in}{1.879800in}}%
\pgfusepath{stroke}%
\end{pgfscope}%
\begin{pgfscope}%
\pgfsetrectcap%
\pgfsetmiterjoin%
\pgfsetlinewidth{0.803000pt}%
\definecolor{currentstroke}{rgb}{0.101961,0.101961,0.101961}%
\pgfsetstrokecolor{currentstroke}%
\pgfsetdash{}{0pt}%
\pgfpathmoveto{\pgfqpoint{4.974268in}{0.663000in}}%
\pgfpathlineto{\pgfqpoint{6.895000in}{0.663000in}}%
\pgfusepath{stroke}%
\end{pgfscope}%
\begin{pgfscope}%
\definecolor{textcolor}{rgb}{0.101961,0.101961,0.101961}%
\pgfsetstrokecolor{textcolor}%
\pgfsetfillcolor{textcolor}%
\pgftext[x=0.021000in,y=3.841500in,left,top]{\color{textcolor}\sffamily\fontsize{7.500000}{9.000000}\bfseries\selectfont (a)}%
\end{pgfscope}%
\begin{pgfscope}%
\definecolor{textcolor}{rgb}{0.101961,0.101961,0.101961}%
\pgfsetstrokecolor{textcolor}%
\pgfsetfillcolor{textcolor}%
\pgftext[x=0.021000in,y=1.911000in,left,top]{\color{textcolor}\sffamily\fontsize{7.500000}{9.000000}\bfseries\selectfont (b)}%
\end{pgfscope}%
\begin{pgfscope}%
\definecolor{textcolor}{rgb}{0.549020,0.549020,0.549020}%
\pgfsetstrokecolor{textcolor}%
\pgfsetfillcolor{textcolor}%
\pgftext[x=6.930000in,y=3.841500in,right,top]{\color{textcolor}\sffamily\fontsize{5.500000}{6.600000}\selectfont OLD DATA}%
\end{pgfscope}%
\begin{pgfscope}%
\pgfsetrectcap%
\pgfsetroundjoin%
\pgfsetlinewidth{1.204500pt}%
\definecolor{currentstroke}{rgb}{0.615686,0.705882,0.807843}%
\pgfsetstrokecolor{currentstroke}%
\pgfsetdash{}{0pt}%
\pgfpathmoveto{\pgfqpoint{0.979980in}{0.291667in}}%
\pgfpathlineto{\pgfqpoint{1.057758in}{0.291667in}}%
\pgfpathlineto{\pgfqpoint{1.135535in}{0.291667in}}%
\pgfusepath{stroke}%
\end{pgfscope}%
\begin{pgfscope}%
\pgfsetbuttcap%
\pgfsetroundjoin%
\definecolor{currentfill}{rgb}{0.615686,0.705882,0.807843}%
\pgfsetfillcolor{currentfill}%
\pgfsetlinewidth{1.003750pt}%
\definecolor{currentstroke}{rgb}{0.615686,0.705882,0.807843}%
\pgfsetstrokecolor{currentstroke}%
\pgfsetdash{}{0pt}%
\pgfsys@defobject{currentmarker}{\pgfqpoint{-0.013889in}{-0.013889in}}{\pgfqpoint{0.013889in}{0.013889in}}{%
\pgfpathmoveto{\pgfqpoint{0.000000in}{-0.013889in}}%
\pgfpathcurveto{\pgfqpoint{0.003683in}{-0.013889in}}{\pgfqpoint{0.007216in}{-0.012425in}}{\pgfqpoint{0.009821in}{-0.009821in}}%
\pgfpathcurveto{\pgfqpoint{0.012425in}{-0.007216in}}{\pgfqpoint{0.013889in}{-0.003683in}}{\pgfqpoint{0.013889in}{0.000000in}}%
\pgfpathcurveto{\pgfqpoint{0.013889in}{0.003683in}}{\pgfqpoint{0.012425in}{0.007216in}}{\pgfqpoint{0.009821in}{0.009821in}}%
\pgfpathcurveto{\pgfqpoint{0.007216in}{0.012425in}}{\pgfqpoint{0.003683in}{0.013889in}}{\pgfqpoint{0.000000in}{0.013889in}}%
\pgfpathcurveto{\pgfqpoint{-0.003683in}{0.013889in}}{\pgfqpoint{-0.007216in}{0.012425in}}{\pgfqpoint{-0.009821in}{0.009821in}}%
\pgfpathcurveto{\pgfqpoint{-0.012425in}{0.007216in}}{\pgfqpoint{-0.013889in}{0.003683in}}{\pgfqpoint{-0.013889in}{0.000000in}}%
\pgfpathcurveto{\pgfqpoint{-0.013889in}{-0.003683in}}{\pgfqpoint{-0.012425in}{-0.007216in}}{\pgfqpoint{-0.009821in}{-0.009821in}}%
\pgfpathcurveto{\pgfqpoint{-0.007216in}{-0.012425in}}{\pgfqpoint{-0.003683in}{-0.013889in}}{\pgfqpoint{0.000000in}{-0.013889in}}%
\pgfpathlineto{\pgfqpoint{0.000000in}{-0.013889in}}%
\pgfpathclose%
\pgfusepath{stroke,fill}%
}%
\begin{pgfscope}%
\pgfsys@transformshift{1.057758in}{0.291667in}%
\pgfsys@useobject{currentmarker}{}%
\end{pgfscope}%
\end{pgfscope}%
\begin{pgfscope}%
\definecolor{textcolor}{rgb}{0.101961,0.101961,0.101961}%
\pgfsetstrokecolor{textcolor}%
\pgfsetfillcolor{textcolor}%
\pgftext[x=1.213313in,y=0.257639in,left,base]{\color{textcolor}\sffamily\fontsize{7.000000}{8.400000}\selectfont COVEX (old: rung 3)}%
\end{pgfscope}%
\begin{pgfscope}%
\pgfsetrectcap%
\pgfsetroundjoin%
\pgfsetlinewidth{1.204500pt}%
\definecolor{currentstroke}{rgb}{0.435294,0.576471,0.729412}%
\pgfsetstrokecolor{currentstroke}%
\pgfsetdash{}{0pt}%
\pgfpathmoveto{\pgfqpoint{0.979980in}{0.145833in}}%
\pgfpathlineto{\pgfqpoint{1.057758in}{0.145833in}}%
\pgfpathlineto{\pgfqpoint{1.135535in}{0.145833in}}%
\pgfusepath{stroke}%
\end{pgfscope}%
\begin{pgfscope}%
\pgfsetbuttcap%
\pgfsetroundjoin%
\definecolor{currentfill}{rgb}{0.435294,0.576471,0.729412}%
\pgfsetfillcolor{currentfill}%
\pgfsetlinewidth{1.003750pt}%
\definecolor{currentstroke}{rgb}{0.435294,0.576471,0.729412}%
\pgfsetstrokecolor{currentstroke}%
\pgfsetdash{}{0pt}%
\pgfsys@defobject{currentmarker}{\pgfqpoint{-0.013889in}{-0.013889in}}{\pgfqpoint{0.013889in}{0.013889in}}{%
\pgfpathmoveto{\pgfqpoint{0.000000in}{-0.013889in}}%
\pgfpathcurveto{\pgfqpoint{0.003683in}{-0.013889in}}{\pgfqpoint{0.007216in}{-0.012425in}}{\pgfqpoint{0.009821in}{-0.009821in}}%
\pgfpathcurveto{\pgfqpoint{0.012425in}{-0.007216in}}{\pgfqpoint{0.013889in}{-0.003683in}}{\pgfqpoint{0.013889in}{0.000000in}}%
\pgfpathcurveto{\pgfqpoint{0.013889in}{0.003683in}}{\pgfqpoint{0.012425in}{0.007216in}}{\pgfqpoint{0.009821in}{0.009821in}}%
\pgfpathcurveto{\pgfqpoint{0.007216in}{0.012425in}}{\pgfqpoint{0.003683in}{0.013889in}}{\pgfqpoint{0.000000in}{0.013889in}}%
\pgfpathcurveto{\pgfqpoint{-0.003683in}{0.013889in}}{\pgfqpoint{-0.007216in}{0.012425in}}{\pgfqpoint{-0.009821in}{0.009821in}}%
\pgfpathcurveto{\pgfqpoint{-0.012425in}{0.007216in}}{\pgfqpoint{-0.013889in}{0.003683in}}{\pgfqpoint{-0.013889in}{0.000000in}}%
\pgfpathcurveto{\pgfqpoint{-0.013889in}{-0.003683in}}{\pgfqpoint{-0.012425in}{-0.007216in}}{\pgfqpoint{-0.009821in}{-0.009821in}}%
\pgfpathcurveto{\pgfqpoint{-0.007216in}{-0.012425in}}{\pgfqpoint{-0.003683in}{-0.013889in}}{\pgfqpoint{0.000000in}{-0.013889in}}%
\pgfpathlineto{\pgfqpoint{0.000000in}{-0.013889in}}%
\pgfpathclose%
\pgfusepath{stroke,fill}%
}%
\begin{pgfscope}%
\pgfsys@transformshift{1.057758in}{0.145833in}%
\pgfsys@useobject{currentmarker}{}%
\end{pgfscope}%
\end{pgfscope}%
\begin{pgfscope}%
\definecolor{textcolor}{rgb}{0.101961,0.101961,0.101961}%
\pgfsetstrokecolor{textcolor}%
\pgfsetfillcolor{textcolor}%
\pgftext[x=1.213313in,y=0.111806in,left,base]{\color{textcolor}\sffamily\fontsize{7.000000}{8.400000}\selectfont COVEX eager (old)}%
\end{pgfscope}%
\begin{pgfscope}%
\pgfsetrectcap%
\pgfsetroundjoin%
\pgfsetlinewidth{1.204500pt}%
\definecolor{currentstroke}{rgb}{0.156863,0.356863,0.568627}%
\pgfsetstrokecolor{currentstroke}%
\pgfsetdash{}{0pt}%
\pgfpathmoveto{\pgfqpoint{2.248336in}{0.291667in}}%
\pgfpathlineto{\pgfqpoint{2.326114in}{0.291667in}}%
\pgfpathlineto{\pgfqpoint{2.403892in}{0.291667in}}%
\pgfusepath{stroke}%
\end{pgfscope}%
\begin{pgfscope}%
\pgfsetbuttcap%
\pgfsetroundjoin%
\definecolor{currentfill}{rgb}{0.156863,0.356863,0.568627}%
\pgfsetfillcolor{currentfill}%
\pgfsetlinewidth{1.003750pt}%
\definecolor{currentstroke}{rgb}{0.156863,0.356863,0.568627}%
\pgfsetstrokecolor{currentstroke}%
\pgfsetdash{}{0pt}%
\pgfsys@defobject{currentmarker}{\pgfqpoint{-0.013889in}{-0.013889in}}{\pgfqpoint{0.013889in}{0.013889in}}{%
\pgfpathmoveto{\pgfqpoint{0.000000in}{-0.013889in}}%
\pgfpathcurveto{\pgfqpoint{0.003683in}{-0.013889in}}{\pgfqpoint{0.007216in}{-0.012425in}}{\pgfqpoint{0.009821in}{-0.009821in}}%
\pgfpathcurveto{\pgfqpoint{0.012425in}{-0.007216in}}{\pgfqpoint{0.013889in}{-0.003683in}}{\pgfqpoint{0.013889in}{0.000000in}}%
\pgfpathcurveto{\pgfqpoint{0.013889in}{0.003683in}}{\pgfqpoint{0.012425in}{0.007216in}}{\pgfqpoint{0.009821in}{0.009821in}}%
\pgfpathcurveto{\pgfqpoint{0.007216in}{0.012425in}}{\pgfqpoint{0.003683in}{0.013889in}}{\pgfqpoint{0.000000in}{0.013889in}}%
\pgfpathcurveto{\pgfqpoint{-0.003683in}{0.013889in}}{\pgfqpoint{-0.007216in}{0.012425in}}{\pgfqpoint{-0.009821in}{0.009821in}}%
\pgfpathcurveto{\pgfqpoint{-0.012425in}{0.007216in}}{\pgfqpoint{-0.013889in}{0.003683in}}{\pgfqpoint{-0.013889in}{0.000000in}}%
\pgfpathcurveto{\pgfqpoint{-0.013889in}{-0.003683in}}{\pgfqpoint{-0.012425in}{-0.007216in}}{\pgfqpoint{-0.009821in}{-0.009821in}}%
\pgfpathcurveto{\pgfqpoint{-0.007216in}{-0.012425in}}{\pgfqpoint{-0.003683in}{-0.013889in}}{\pgfqpoint{0.000000in}{-0.013889in}}%
\pgfpathlineto{\pgfqpoint{0.000000in}{-0.013889in}}%
\pgfpathclose%
\pgfusepath{stroke,fill}%
}%
\begin{pgfscope}%
\pgfsys@transformshift{2.326114in}{0.291667in}%
\pgfsys@useobject{currentmarker}{}%
\end{pgfscope}%
\end{pgfscope}%
\begin{pgfscope}%
\definecolor{textcolor}{rgb}{0.101961,0.101961,0.101961}%
\pgfsetstrokecolor{textcolor}%
\pgfsetfillcolor{textcolor}%
\pgftext[x=2.481670in,y=0.257639in,left,base]{\color{textcolor}\sffamily\fontsize{7.000000}{8.400000}\selectfont COVEX, no S2 (dev)}%
\end{pgfscope}%
\begin{pgfscope}%
\pgfsetrectcap%
\pgfsetroundjoin%
\pgfsetlinewidth{1.204500pt}%
\definecolor{currentstroke}{rgb}{0.062745,0.160784,0.290196}%
\pgfsetstrokecolor{currentstroke}%
\pgfsetdash{}{0pt}%
\pgfpathmoveto{\pgfqpoint{2.248336in}{0.145833in}}%
\pgfpathlineto{\pgfqpoint{2.326114in}{0.145833in}}%
\pgfpathlineto{\pgfqpoint{2.403892in}{0.145833in}}%
\pgfusepath{stroke}%
\end{pgfscope}%
\begin{pgfscope}%
\pgfsetbuttcap%
\pgfsetroundjoin%
\definecolor{currentfill}{rgb}{0.062745,0.160784,0.290196}%
\pgfsetfillcolor{currentfill}%
\pgfsetlinewidth{1.003750pt}%
\definecolor{currentstroke}{rgb}{0.062745,0.160784,0.290196}%
\pgfsetstrokecolor{currentstroke}%
\pgfsetdash{}{0pt}%
\pgfsys@defobject{currentmarker}{\pgfqpoint{-0.013889in}{-0.013889in}}{\pgfqpoint{0.013889in}{0.013889in}}{%
\pgfpathmoveto{\pgfqpoint{0.000000in}{-0.013889in}}%
\pgfpathcurveto{\pgfqpoint{0.003683in}{-0.013889in}}{\pgfqpoint{0.007216in}{-0.012425in}}{\pgfqpoint{0.009821in}{-0.009821in}}%
\pgfpathcurveto{\pgfqpoint{0.012425in}{-0.007216in}}{\pgfqpoint{0.013889in}{-0.003683in}}{\pgfqpoint{0.013889in}{0.000000in}}%
\pgfpathcurveto{\pgfqpoint{0.013889in}{0.003683in}}{\pgfqpoint{0.012425in}{0.007216in}}{\pgfqpoint{0.009821in}{0.009821in}}%
\pgfpathcurveto{\pgfqpoint{0.007216in}{0.012425in}}{\pgfqpoint{0.003683in}{0.013889in}}{\pgfqpoint{0.000000in}{0.013889in}}%
\pgfpathcurveto{\pgfqpoint{-0.003683in}{0.013889in}}{\pgfqpoint{-0.007216in}{0.012425in}}{\pgfqpoint{-0.009821in}{0.009821in}}%
\pgfpathcurveto{\pgfqpoint{-0.012425in}{0.007216in}}{\pgfqpoint{-0.013889in}{0.003683in}}{\pgfqpoint{-0.013889in}{0.000000in}}%
\pgfpathcurveto{\pgfqpoint{-0.013889in}{-0.003683in}}{\pgfqpoint{-0.012425in}{-0.007216in}}{\pgfqpoint{-0.009821in}{-0.009821in}}%
\pgfpathcurveto{\pgfqpoint{-0.007216in}{-0.012425in}}{\pgfqpoint{-0.003683in}{-0.013889in}}{\pgfqpoint{0.000000in}{-0.013889in}}%
\pgfpathlineto{\pgfqpoint{0.000000in}{-0.013889in}}%
\pgfpathclose%
\pgfusepath{stroke,fill}%
}%
\begin{pgfscope}%
\pgfsys@transformshift{2.326114in}{0.145833in}%
\pgfsys@useobject{currentmarker}{}%
\end{pgfscope}%
\end{pgfscope}%
\begin{pgfscope}%
\definecolor{textcolor}{rgb}{0.101961,0.101961,0.101961}%
\pgfsetstrokecolor{textcolor}%
\pgfsetfillcolor{textcolor}%
\pgftext[x=2.481670in,y=0.111806in,left,base]{\color{textcolor}\sffamily\fontsize{7.000000}{8.400000}\selectfont COVEX + S2 (dev)}%
\end{pgfscope}%
\begin{pgfscope}%
\pgfsetbuttcap%
\pgfsetroundjoin%
\pgfsetlinewidth{0.602250pt}%
\definecolor{currentstroke}{rgb}{0.333333,0.333333,0.333333}%
\pgfsetstrokecolor{currentstroke}%
\pgfsetdash{{1.200000pt}{1.200000pt}}{0.000000pt}%
\pgfpathmoveto{\pgfqpoint{3.491206in}{0.291667in}}%
\pgfpathlineto{\pgfqpoint{3.568984in}{0.291667in}}%
\pgfpathlineto{\pgfqpoint{3.646762in}{0.291667in}}%
\pgfusepath{stroke}%
\end{pgfscope}%
\begin{pgfscope}%
\definecolor{textcolor}{rgb}{0.101961,0.101961,0.101961}%
\pgfsetstrokecolor{textcolor}%
\pgfsetfillcolor{textcolor}%
\pgftext[x=3.724539in,y=0.257639in,left,base]{\color{textcolor}\sffamily\fontsize{7.000000}{8.400000}\selectfont no-outage control, dashcam (a)}%
\end{pgfscope}%
\begin{pgfscope}%
\pgfsetrectcap%
\pgfsetroundjoin%
\pgfsetlinewidth{0.903375pt}%
\definecolor{currentstroke}{rgb}{0.000000,0.000000,0.000000}%
\pgfsetstrokecolor{currentstroke}%
\pgfsetdash{}{0pt}%
\pgfpathmoveto{\pgfqpoint{3.491206in}{0.145833in}}%
\pgfpathlineto{\pgfqpoint{3.568984in}{0.145833in}}%
\pgfpathlineto{\pgfqpoint{3.646762in}{0.145833in}}%
\pgfusepath{stroke}%
\end{pgfscope}%
\begin{pgfscope}%
\definecolor{textcolor}{rgb}{0.101961,0.101961,0.101961}%
\pgfsetstrokecolor{textcolor}%
\pgfsetfillcolor{textcolor}%
\pgftext[x=3.724539in,y=0.111806in,left,base]{\color{textcolor}\sffamily\fontsize{7.000000}{8.400000}\selectfont no-outage control (b)}%
\end{pgfscope}%
\begin{pgfscope}%
\pgfsetbuttcap%
\pgfsetmiterjoin%
\definecolor{currentfill}{rgb}{0.466667,0.466667,0.466667}%
\pgfsetfillcolor{currentfill}%
\pgfsetfillopacity{0.450000}%
\pgfsetlinewidth{0.000000pt}%
\definecolor{currentstroke}{rgb}{0.000000,0.000000,0.000000}%
\pgfsetstrokecolor{currentstroke}%
\pgfsetstrokeopacity{0.450000}%
\pgfsetdash{}{0pt}%
\pgfpathmoveto{\pgfqpoint{5.185859in}{0.257639in}}%
\pgfpathlineto{\pgfqpoint{5.341415in}{0.257639in}}%
\pgfpathlineto{\pgfqpoint{5.341415in}{0.325694in}}%
\pgfpathlineto{\pgfqpoint{5.185859in}{0.325694in}}%
\pgfpathlineto{\pgfqpoint{5.185859in}{0.257639in}}%
\pgfpathclose%
\pgfusepath{fill}%
\end{pgfscope}%
\begin{pgfscope}%
\definecolor{textcolor}{rgb}{0.101961,0.101961,0.101961}%
\pgfsetstrokecolor{textcolor}%
\pgfsetfillcolor{textcolor}%
\pgftext[x=5.419193in,y=0.257639in,left,base]{\color{textcolor}\sffamily\fontsize{7.000000}{8.400000}\selectfont whole run}%
\end{pgfscope}%
\begin{pgfscope}%
\pgfsetbuttcap%
\pgfsetmiterjoin%
\definecolor{currentfill}{rgb}{0.466667,0.466667,0.466667}%
\pgfsetfillcolor{currentfill}%
\pgfsetlinewidth{0.000000pt}%
\definecolor{currentstroke}{rgb}{0.000000,0.000000,0.000000}%
\pgfsetstrokecolor{currentstroke}%
\pgfsetstrokeopacity{0.000000}%
\pgfsetdash{}{0pt}%
\pgfpathmoveto{\pgfqpoint{5.185859in}{0.122068in}}%
\pgfpathlineto{\pgfqpoint{5.341415in}{0.122068in}}%
\pgfpathlineto{\pgfqpoint{5.341415in}{0.190124in}}%
\pgfpathlineto{\pgfqpoint{5.185859in}{0.190124in}}%
\pgfpathlineto{\pgfqpoint{5.185859in}{0.122068in}}%
\pgfpathclose%
\pgfusepath{fill}%
\end{pgfscope}%
\begin{pgfscope}%
\definecolor{textcolor}{rgb}{0.101961,0.101961,0.101961}%
\pgfsetstrokecolor{textcolor}%
\pgfsetfillcolor{textcolor}%
\pgftext[x=5.419193in,y=0.122068in,left,base]{\color{textcolor}\sffamily\fontsize{7.000000}{8.400000}\selectfont R to R+120 s}%
\end{pgfscope}%
\end{pgfpicture}%
\makeatother%
\endgroup%
}
\caption{[PLACEHOLDER: OLD DATA, n20 seed 4 only; seeds 5-7 outage cells not committed; eager and F0 + S2 have no no-outage control; rung 3 and eager are OLD, F0 and F0 + S2 are DEV] Network outage and reconnect, 20 devices, seed 4, 14 robots offline for d = 188 s (outage ends at the reconnect R = 595.5 s). (a) On-time rate in 10 s bins of request time from R-60 s to R+300 s per request type; the dashed lines on the dashcam panel are the same arm without the outage where a cell exists (omitted for the sparse robot and single-shot types). Robot and single-shot bins hold 0 to a few requests each, so those series are noisy and bins with no request are left blank. (b) Whole-run (light) and R to R+120 s (dark) on-time rate; black ticks mark the no-outage control. On time is the committed within-L flag of each request.}
\label{fig:F12_network}
\end{figure*}
\clearpage
\begin{figure}[t]
\centering
{\fontfamily{LinuxBiolinumT}\selectfont
%% Creator: Matplotlib, PGF backend
%%
%% To include the figure in your LaTeX document, write
%%   \input{<filename>.pgf}
%%
%% Make sure the required packages are loaded in your preamble
%%   \usepackage{pgf}
%%
%% Also ensure that all the required font packages are loaded; for instance,
%% the lmodern package is sometimes necessary when using math font.
%%   \usepackage{lmodern}
%%
%% Figures using additional raster images can only be included by \input if
%% they are in the same directory as the main LaTeX file. For loading figures
%% from other directories you can use the `import` package
%%   \usepackage{import}
%%
%% and then include the figures with
%%   \import{<path to file>}{<filename>.pgf}
%%
%% Matplotlib used the following preamble
%%
\begingroup%
\makeatletter%
\begin{pgfpicture}%
\pgfpathrectangle{\pgfpointorigin}{\pgfqpoint{3.330000in}{3.600000in}}%
\pgfusepath{use as bounding box, clip}%
\begin{pgfscope}%
\pgfsetbuttcap%
\pgfsetmiterjoin%
\definecolor{currentfill}{rgb}{1.000000,1.000000,1.000000}%
\pgfsetfillcolor{currentfill}%
\pgfsetlinewidth{0.000000pt}%
\definecolor{currentstroke}{rgb}{1.000000,1.000000,1.000000}%
\pgfsetstrokecolor{currentstroke}%
\pgfsetdash{}{0pt}%
\pgfpathmoveto{\pgfqpoint{0.000000in}{0.000000in}}%
\pgfpathlineto{\pgfqpoint{3.330000in}{0.000000in}}%
\pgfpathlineto{\pgfqpoint{3.330000in}{3.600000in}}%
\pgfpathlineto{\pgfqpoint{0.000000in}{3.600000in}}%
\pgfpathlineto{\pgfqpoint{0.000000in}{0.000000in}}%
\pgfpathclose%
\pgfusepath{fill}%
\end{pgfscope}%
\begin{pgfscope}%
\pgfsetbuttcap%
\pgfsetmiterjoin%
\definecolor{currentfill}{rgb}{1.000000,1.000000,1.000000}%
\pgfsetfillcolor{currentfill}%
\pgfsetlinewidth{0.000000pt}%
\definecolor{currentstroke}{rgb}{0.000000,0.000000,0.000000}%
\pgfsetstrokecolor{currentstroke}%
\pgfsetstrokeopacity{0.000000}%
\pgfsetdash{}{0pt}%
\pgfpathmoveto{\pgfqpoint{1.198800in}{0.720000in}}%
\pgfpathlineto{\pgfqpoint{1.749175in}{0.720000in}}%
\pgfpathlineto{\pgfqpoint{1.749175in}{3.384000in}}%
\pgfpathlineto{\pgfqpoint{1.198800in}{3.384000in}}%
\pgfpathlineto{\pgfqpoint{1.198800in}{0.720000in}}%
\pgfpathclose%
\pgfusepath{fill}%
\end{pgfscope}%
\begin{pgfscope}%
\pgfpathrectangle{\pgfqpoint{1.198800in}{0.720000in}}{\pgfqpoint{0.550375in}{2.664000in}}%
\pgfusepath{clip}%
\pgfsetrectcap%
\pgfsetroundjoin%
\pgfsetlinewidth{0.401500pt}%
\definecolor{currentstroke}{rgb}{0.901961,0.901961,0.901961}%
\pgfsetstrokecolor{currentstroke}%
\pgfsetdash{}{0pt}%
\pgfpathmoveto{\pgfqpoint{1.198800in}{0.720000in}}%
\pgfpathlineto{\pgfqpoint{1.198800in}{3.384000in}}%
\pgfusepath{stroke}%
\end{pgfscope}%
\begin{pgfscope}%
\definecolor{textcolor}{rgb}{0.101961,0.101961,0.101961}%
\pgfsetstrokecolor{textcolor}%
\pgfsetfillcolor{textcolor}%
\pgftext[x=1.198800in,y=0.671389in,,top]{\color{textcolor}\sffamily\fontsize{7.000000}{8.400000}\selectfont 0.00}%
\end{pgfscope}%
\begin{pgfscope}%
\pgfpathrectangle{\pgfqpoint{1.198800in}{0.720000in}}{\pgfqpoint{0.550375in}{2.664000in}}%
\pgfusepath{clip}%
\pgfsetrectcap%
\pgfsetroundjoin%
\pgfsetlinewidth{0.401500pt}%
\definecolor{currentstroke}{rgb}{0.901961,0.901961,0.901961}%
\pgfsetstrokecolor{currentstroke}%
\pgfsetdash{}{0pt}%
\pgfpathmoveto{\pgfqpoint{1.628183in}{0.720000in}}%
\pgfpathlineto{\pgfqpoint{1.628183in}{3.384000in}}%
\pgfusepath{stroke}%
\end{pgfscope}%
\begin{pgfscope}%
\definecolor{textcolor}{rgb}{0.101961,0.101961,0.101961}%
\pgfsetstrokecolor{textcolor}%
\pgfsetfillcolor{textcolor}%
\pgftext[x=1.628183in,y=0.671389in,,top]{\color{textcolor}\sffamily\fontsize{7.000000}{8.400000}\selectfont 0.25}%
\end{pgfscope}%
\begin{pgfscope}%
\definecolor{textcolor}{rgb}{0.000000,0.000000,0.000000}%
\pgfsetstrokecolor{textcolor}%
\pgfsetfillcolor{textcolor}%
\pgftext[x=1.326802in, y=0.461899in, left, base]{\color{textcolor}\sffamily\fontsize{7.500000}{9.000000}\selectfont correct}%
\end{pgfscope}%
\begin{pgfscope}%
\definecolor{textcolor}{rgb}{0.000000,0.000000,0.000000}%
\pgfsetstrokecolor{textcolor}%
\pgfsetfillcolor{textcolor}%
\pgftext[x=1.248179in, y=0.361976in, left, base]{\color{textcolor}\sffamily\fontsize{7.500000}{9.000000}\selectfont per GPU-s}%
\end{pgfscope}%
\begin{pgfscope}%
\definecolor{textcolor}{rgb}{0.101961,0.101961,0.101961}%
\pgfsetstrokecolor{textcolor}%
\pgfsetfillcolor{textcolor}%
\pgftext[x=0.185794in, y=3.216525in, left, base]{\color{textcolor}\sffamily\fontsize{7.000000}{8.400000}\selectfont \textbf{COVEX (old: rung 3)}}%
\end{pgfscope}%
\begin{pgfscope}%
\definecolor{textcolor}{rgb}{0.101961,0.101961,0.101961}%
\pgfsetstrokecolor{textcolor}%
\pgfsetfillcolor{textcolor}%
\pgftext[x=0.400701in, y=2.998165in, left, base]{\color{textcolor}\sffamily\fontsize{7.000000}{8.400000}\selectfont SimpleStream~\cite{simplestream2026}}%
\end{pgfscope}%
\begin{pgfscope}%
\definecolor{textcolor}{rgb}{0.101961,0.101961,0.101961}%
\pgfsetstrokecolor{textcolor}%
\pgfsetfillcolor{textcolor}%
\pgftext[x=0.443068in, y=2.779804in, left, base]{\color{textcolor}\sffamily\fontsize{7.000000}{8.400000}\selectfont StreamScout~\cite{streamscout2026}}%
\end{pgfscope}%
\begin{pgfscope}%
\definecolor{textcolor}{rgb}{0.101961,0.101961,0.101961}%
\pgfsetstrokecolor{textcolor}%
\pgfsetfillcolor{textcolor}%
\pgftext[x=0.557508in, y=2.561443in, left, base]{\color{textcolor}\sffamily\fontsize{7.000000}{8.400000}\selectfont Mooncake~\cite{mooncake2025}}%
\end{pgfscope}%
\begin{pgfscope}%
\definecolor{textcolor}{rgb}{0.101961,0.101961,0.101961}%
\pgfsetstrokecolor{textcolor}%
\pgfsetfillcolor{textcolor}%
\pgftext[x=0.824703in, y=2.343083in, left, base]{\color{textcolor}\sffamily\fontsize{7.000000}{8.400000}\selectfont LLF~\cite{mok1983llf}}%
\end{pgfscope}%
\begin{pgfscope}%
\definecolor{textcolor}{rgb}{0.101961,0.101961,0.101961}%
\pgfsetstrokecolor{textcolor}%
\pgfsetfillcolor{textcolor}%
\pgftext[x=0.800566in, y=2.124722in, left, base]{\color{textcolor}\sffamily\fontsize{7.000000}{8.400000}\selectfont EDF~\cite{liu1973edf}}%
\end{pgfscope}%
\begin{pgfscope}%
\definecolor{textcolor}{rgb}{0.101961,0.101961,0.101961}%
\pgfsetstrokecolor{textcolor}%
\pgfsetfillcolor{textcolor}%
\pgftext[x=0.554976in, y=1.906361in, left, base]{\color{textcolor}\sffamily\fontsize{7.000000}{8.400000}\selectfont ECO-LLM~\cite{ecollm2025}}%
\end{pgfscope}%
\begin{pgfscope}%
\definecolor{textcolor}{rgb}{0.101961,0.101961,0.101961}%
\pgfsetstrokecolor{textcolor}%
\pgfsetfillcolor{textcolor}%
\pgftext[x=0.250714in, y=1.688001in, left, base]{\color{textcolor}\sffamily\fontsize{7.000000}{8.400000}\selectfont EDF + feasibility~\cite{liu1973edf}}%
\end{pgfscope}%
\begin{pgfscope}%
\definecolor{textcolor}{rgb}{0.101961,0.101961,0.101961}%
\pgfsetstrokecolor{textcolor}%
\pgfsetfillcolor{textcolor}%
\pgftext[x=0.652537in, y=1.469640in, left, base]{\color{textcolor}\sffamily\fontsize{7.000000}{8.400000}\selectfont KVFlow~\cite{kvflow2025}}%
\end{pgfscope}%
\begin{pgfscope}%
\definecolor{textcolor}{rgb}{0.101961,0.101961,0.101961}%
\pgfsetstrokecolor{textcolor}%
\pgfsetfillcolor{textcolor}%
\pgftext[x=0.318625in, y=1.253980in, left, base]{\color{textcolor}\sffamily\fontsize{7.000000}{8.400000}\selectfont Always-cloud FCFS}%
\end{pgfscope}%
\begin{pgfscope}%
\definecolor{textcolor}{rgb}{0.101961,0.101961,0.101961}%
\pgfsetstrokecolor{textcolor}%
\pgfsetfillcolor{textcolor}%
\pgftext[x=0.430240in, y=1.032919in, left, base]{\color{textcolor}\sffamily\fontsize{7.000000}{8.400000}\selectfont QueryStream~\cite{querystream2026}}%
\end{pgfscope}%
\begin{pgfscope}%
\definecolor{textcolor}{rgb}{0.101961,0.101961,0.101961}%
\pgfsetstrokecolor{textcolor}%
\pgfsetfillcolor{textcolor}%
\pgftext[x=0.364749in, y=0.814558in, left, base]{\color{textcolor}\sffamily\fontsize{7.000000}{8.400000}\selectfont AttentionStore~\cite{cachedattention2024}}%
\end{pgfscope}%
\begin{pgfscope}%
\pgfpathrectangle{\pgfqpoint{1.198800in}{0.720000in}}{\pgfqpoint{0.550375in}{2.664000in}}%
\pgfusepath{clip}%
\pgfsetbuttcap%
\pgfsetmiterjoin%
\definecolor{currentfill}{rgb}{0.156863,0.356863,0.568627}%
\pgfsetfillcolor{currentfill}%
\pgfsetlinewidth{0.000000pt}%
\definecolor{currentstroke}{rgb}{0.000000,0.000000,0.000000}%
\pgfsetstrokecolor{currentstroke}%
\pgfsetstrokeopacity{0.000000}%
\pgfsetdash{}{0pt}%
\pgfpathmoveto{\pgfqpoint{1.198800in}{3.176557in}}%
\pgfpathlineto{\pgfqpoint{1.694160in}{3.176557in}}%
\pgfpathlineto{\pgfqpoint{1.694160in}{3.329410in}}%
\pgfpathlineto{\pgfqpoint{1.198800in}{3.329410in}}%
\pgfpathlineto{\pgfqpoint{1.198800in}{3.176557in}}%
\pgfpathclose%
\pgfusepath{fill}%
\end{pgfscope}%
\begin{pgfscope}%
\pgfpathrectangle{\pgfqpoint{1.198800in}{0.720000in}}{\pgfqpoint{0.550375in}{2.664000in}}%
\pgfusepath{clip}%
\pgfsetbuttcap%
\pgfsetmiterjoin%
\definecolor{currentfill}{rgb}{0.784314,0.784314,0.784314}%
\pgfsetfillcolor{currentfill}%
\pgfsetlinewidth{0.000000pt}%
\definecolor{currentstroke}{rgb}{0.000000,0.000000,0.000000}%
\pgfsetstrokecolor{currentstroke}%
\pgfsetstrokeopacity{0.000000}%
\pgfsetdash{}{0pt}%
\pgfpathmoveto{\pgfqpoint{1.198800in}{2.958197in}}%
\pgfpathlineto{\pgfqpoint{1.646460in}{2.958197in}}%
\pgfpathlineto{\pgfqpoint{1.646460in}{3.111049in}}%
\pgfpathlineto{\pgfqpoint{1.198800in}{3.111049in}}%
\pgfpathlineto{\pgfqpoint{1.198800in}{2.958197in}}%
\pgfpathclose%
\pgfusepath{fill}%
\end{pgfscope}%
\begin{pgfscope}%
\pgfpathrectangle{\pgfqpoint{1.198800in}{0.720000in}}{\pgfqpoint{0.550375in}{2.664000in}}%
\pgfusepath{clip}%
\pgfsetbuttcap%
\pgfsetmiterjoin%
\definecolor{currentfill}{rgb}{0.784314,0.784314,0.784314}%
\pgfsetfillcolor{currentfill}%
\pgfsetlinewidth{0.000000pt}%
\definecolor{currentstroke}{rgb}{0.000000,0.000000,0.000000}%
\pgfsetstrokecolor{currentstroke}%
\pgfsetstrokeopacity{0.000000}%
\pgfsetdash{}{0pt}%
\pgfpathmoveto{\pgfqpoint{1.198800in}{2.739836in}}%
\pgfpathlineto{\pgfqpoint{1.644179in}{2.739836in}}%
\pgfpathlineto{\pgfqpoint{1.644179in}{2.892689in}}%
\pgfpathlineto{\pgfqpoint{1.198800in}{2.892689in}}%
\pgfpathlineto{\pgfqpoint{1.198800in}{2.739836in}}%
\pgfpathclose%
\pgfusepath{fill}%
\end{pgfscope}%
\begin{pgfscope}%
\pgfpathrectangle{\pgfqpoint{1.198800in}{0.720000in}}{\pgfqpoint{0.550375in}{2.664000in}}%
\pgfusepath{clip}%
\pgfsetbuttcap%
\pgfsetmiterjoin%
\definecolor{currentfill}{rgb}{0.478431,0.478431,0.478431}%
\pgfsetfillcolor{currentfill}%
\pgfsetlinewidth{0.000000pt}%
\definecolor{currentstroke}{rgb}{0.000000,0.000000,0.000000}%
\pgfsetstrokecolor{currentstroke}%
\pgfsetstrokeopacity{0.000000}%
\pgfsetdash{}{0pt}%
\pgfpathmoveto{\pgfqpoint{1.198800in}{2.521475in}}%
\pgfpathlineto{\pgfqpoint{1.614871in}{2.521475in}}%
\pgfpathlineto{\pgfqpoint{1.614871in}{2.674328in}}%
\pgfpathlineto{\pgfqpoint{1.198800in}{2.674328in}}%
\pgfpathlineto{\pgfqpoint{1.198800in}{2.521475in}}%
\pgfpathclose%
\pgfusepath{fill}%
\end{pgfscope}%
\begin{pgfscope}%
\pgfpathrectangle{\pgfqpoint{1.198800in}{0.720000in}}{\pgfqpoint{0.550375in}{2.664000in}}%
\pgfusepath{clip}%
\pgfsetbuttcap%
\pgfsetmiterjoin%
\definecolor{currentfill}{rgb}{0.784314,0.784314,0.784314}%
\pgfsetfillcolor{currentfill}%
\pgfsetlinewidth{0.000000pt}%
\definecolor{currentstroke}{rgb}{0.000000,0.000000,0.000000}%
\pgfsetstrokecolor{currentstroke}%
\pgfsetstrokeopacity{0.000000}%
\pgfsetdash{}{0pt}%
\pgfpathmoveto{\pgfqpoint{1.198800in}{2.303115in}}%
\pgfpathlineto{\pgfqpoint{1.611743in}{2.303115in}}%
\pgfpathlineto{\pgfqpoint{1.611743in}{2.455967in}}%
\pgfpathlineto{\pgfqpoint{1.198800in}{2.455967in}}%
\pgfpathlineto{\pgfqpoint{1.198800in}{2.303115in}}%
\pgfpathclose%
\pgfusepath{fill}%
\end{pgfscope}%
\begin{pgfscope}%
\pgfpathrectangle{\pgfqpoint{1.198800in}{0.720000in}}{\pgfqpoint{0.550375in}{2.664000in}}%
\pgfusepath{clip}%
\pgfsetbuttcap%
\pgfsetmiterjoin%
\definecolor{currentfill}{rgb}{0.784314,0.784314,0.784314}%
\pgfsetfillcolor{currentfill}%
\pgfsetlinewidth{0.000000pt}%
\definecolor{currentstroke}{rgb}{0.000000,0.000000,0.000000}%
\pgfsetstrokecolor{currentstroke}%
\pgfsetstrokeopacity{0.000000}%
\pgfsetdash{}{0pt}%
\pgfpathmoveto{\pgfqpoint{1.198800in}{2.084754in}}%
\pgfpathlineto{\pgfqpoint{1.607880in}{2.084754in}}%
\pgfpathlineto{\pgfqpoint{1.607880in}{2.237607in}}%
\pgfpathlineto{\pgfqpoint{1.198800in}{2.237607in}}%
\pgfpathlineto{\pgfqpoint{1.198800in}{2.084754in}}%
\pgfpathclose%
\pgfusepath{fill}%
\end{pgfscope}%
\begin{pgfscope}%
\pgfpathrectangle{\pgfqpoint{1.198800in}{0.720000in}}{\pgfqpoint{0.550375in}{2.664000in}}%
\pgfusepath{clip}%
\pgfsetbuttcap%
\pgfsetmiterjoin%
\definecolor{currentfill}{rgb}{0.784314,0.784314,0.784314}%
\pgfsetfillcolor{currentfill}%
\pgfsetlinewidth{0.000000pt}%
\definecolor{currentstroke}{rgb}{0.000000,0.000000,0.000000}%
\pgfsetstrokecolor{currentstroke}%
\pgfsetstrokeopacity{0.000000}%
\pgfsetdash{}{0pt}%
\pgfpathmoveto{\pgfqpoint{1.198800in}{1.866393in}}%
\pgfpathlineto{\pgfqpoint{1.575008in}{1.866393in}}%
\pgfpathlineto{\pgfqpoint{1.575008in}{2.019246in}}%
\pgfpathlineto{\pgfqpoint{1.198800in}{2.019246in}}%
\pgfpathlineto{\pgfqpoint{1.198800in}{1.866393in}}%
\pgfpathclose%
\pgfusepath{fill}%
\end{pgfscope}%
\begin{pgfscope}%
\pgfpathrectangle{\pgfqpoint{1.198800in}{0.720000in}}{\pgfqpoint{0.550375in}{2.664000in}}%
\pgfusepath{clip}%
\pgfsetbuttcap%
\pgfsetmiterjoin%
\definecolor{currentfill}{rgb}{0.784314,0.784314,0.784314}%
\pgfsetfillcolor{currentfill}%
\pgfsetlinewidth{0.000000pt}%
\definecolor{currentstroke}{rgb}{0.000000,0.000000,0.000000}%
\pgfsetstrokecolor{currentstroke}%
\pgfsetstrokeopacity{0.000000}%
\pgfsetdash{}{0pt}%
\pgfpathmoveto{\pgfqpoint{1.198800in}{1.648033in}}%
\pgfpathlineto{\pgfqpoint{1.573028in}{1.648033in}}%
\pgfpathlineto{\pgfqpoint{1.573028in}{1.800885in}}%
\pgfpathlineto{\pgfqpoint{1.198800in}{1.800885in}}%
\pgfpathlineto{\pgfqpoint{1.198800in}{1.648033in}}%
\pgfpathclose%
\pgfusepath{fill}%
\end{pgfscope}%
\begin{pgfscope}%
\pgfpathrectangle{\pgfqpoint{1.198800in}{0.720000in}}{\pgfqpoint{0.550375in}{2.664000in}}%
\pgfusepath{clip}%
\pgfsetbuttcap%
\pgfsetmiterjoin%
\definecolor{currentfill}{rgb}{0.478431,0.478431,0.478431}%
\pgfsetfillcolor{currentfill}%
\pgfsetlinewidth{0.000000pt}%
\definecolor{currentstroke}{rgb}{0.000000,0.000000,0.000000}%
\pgfsetstrokecolor{currentstroke}%
\pgfsetstrokeopacity{0.000000}%
\pgfsetdash{}{0pt}%
\pgfpathmoveto{\pgfqpoint{1.198800in}{1.429672in}}%
\pgfpathlineto{\pgfqpoint{1.519412in}{1.429672in}}%
\pgfpathlineto{\pgfqpoint{1.519412in}{1.582525in}}%
\pgfpathlineto{\pgfqpoint{1.198800in}{1.582525in}}%
\pgfpathlineto{\pgfqpoint{1.198800in}{1.429672in}}%
\pgfpathclose%
\pgfusepath{fill}%
\end{pgfscope}%
\begin{pgfscope}%
\pgfpathrectangle{\pgfqpoint{1.198800in}{0.720000in}}{\pgfqpoint{0.550375in}{2.664000in}}%
\pgfusepath{clip}%
\pgfsetbuttcap%
\pgfsetmiterjoin%
\definecolor{currentfill}{rgb}{0.784314,0.784314,0.784314}%
\pgfsetfillcolor{currentfill}%
\pgfsetlinewidth{0.000000pt}%
\definecolor{currentstroke}{rgb}{0.000000,0.000000,0.000000}%
\pgfsetstrokecolor{currentstroke}%
\pgfsetstrokeopacity{0.000000}%
\pgfsetdash{}{0pt}%
\pgfpathmoveto{\pgfqpoint{1.198800in}{1.211311in}}%
\pgfpathlineto{\pgfqpoint{1.519058in}{1.211311in}}%
\pgfpathlineto{\pgfqpoint{1.519058in}{1.364164in}}%
\pgfpathlineto{\pgfqpoint{1.198800in}{1.364164in}}%
\pgfpathlineto{\pgfqpoint{1.198800in}{1.211311in}}%
\pgfpathclose%
\pgfusepath{fill}%
\end{pgfscope}%
\begin{pgfscope}%
\pgfpathrectangle{\pgfqpoint{1.198800in}{0.720000in}}{\pgfqpoint{0.550375in}{2.664000in}}%
\pgfusepath{clip}%
\pgfsetbuttcap%
\pgfsetmiterjoin%
\definecolor{currentfill}{rgb}{0.784314,0.784314,0.784314}%
\pgfsetfillcolor{currentfill}%
\pgfsetlinewidth{0.000000pt}%
\definecolor{currentstroke}{rgb}{0.000000,0.000000,0.000000}%
\pgfsetstrokecolor{currentstroke}%
\pgfsetstrokeopacity{0.000000}%
\pgfsetdash{}{0pt}%
\pgfpathmoveto{\pgfqpoint{1.198800in}{0.992951in}}%
\pgfpathlineto{\pgfqpoint{1.516911in}{0.992951in}}%
\pgfpathlineto{\pgfqpoint{1.516911in}{1.145803in}}%
\pgfpathlineto{\pgfqpoint{1.198800in}{1.145803in}}%
\pgfpathlineto{\pgfqpoint{1.198800in}{0.992951in}}%
\pgfpathclose%
\pgfusepath{fill}%
\end{pgfscope}%
\begin{pgfscope}%
\pgfpathrectangle{\pgfqpoint{1.198800in}{0.720000in}}{\pgfqpoint{0.550375in}{2.664000in}}%
\pgfusepath{clip}%
\pgfsetbuttcap%
\pgfsetmiterjoin%
\definecolor{currentfill}{rgb}{0.478431,0.478431,0.478431}%
\pgfsetfillcolor{currentfill}%
\pgfsetlinewidth{0.000000pt}%
\definecolor{currentstroke}{rgb}{0.000000,0.000000,0.000000}%
\pgfsetstrokecolor{currentstroke}%
\pgfsetstrokeopacity{0.000000}%
\pgfsetdash{}{0pt}%
\pgfpathmoveto{\pgfqpoint{1.198800in}{0.774590in}}%
\pgfpathlineto{\pgfqpoint{1.516903in}{0.774590in}}%
\pgfpathlineto{\pgfqpoint{1.516903in}{0.927443in}}%
\pgfpathlineto{\pgfqpoint{1.198800in}{0.927443in}}%
\pgfpathlineto{\pgfqpoint{1.198800in}{0.774590in}}%
\pgfpathclose%
\pgfusepath{fill}%
\end{pgfscope}%
\begin{pgfscope}%
\pgfpathrectangle{\pgfqpoint{1.198800in}{0.720000in}}{\pgfqpoint{0.550375in}{2.664000in}}%
\pgfusepath{clip}%
\pgfsetrectcap%
\pgfsetroundjoin%
\pgfsetlinewidth{0.702625pt}%
\definecolor{currentstroke}{rgb}{0.101961,0.101961,0.101961}%
\pgfsetstrokecolor{currentstroke}%
\pgfsetdash{}{0pt}%
\pgfpathmoveto{\pgfqpoint{1.678059in}{3.252984in}}%
\pgfpathlineto{\pgfqpoint{1.722967in}{3.252984in}}%
\pgfusepath{stroke}%
\end{pgfscope}%
\begin{pgfscope}%
\pgfpathrectangle{\pgfqpoint{1.198800in}{0.720000in}}{\pgfqpoint{0.550375in}{2.664000in}}%
\pgfusepath{clip}%
\pgfsetbuttcap%
\pgfsetroundjoin%
\definecolor{currentfill}{rgb}{0.101961,0.101961,0.101961}%
\pgfsetfillcolor{currentfill}%
\pgfsetlinewidth{0.501875pt}%
\definecolor{currentstroke}{rgb}{0.101961,0.101961,0.101961}%
\pgfsetstrokecolor{currentstroke}%
\pgfsetdash{}{0pt}%
\pgfsys@defobject{currentmarker}{\pgfqpoint{0.000000in}{-0.020833in}}{\pgfqpoint{0.000000in}{0.020833in}}{%
\pgfpathmoveto{\pgfqpoint{0.000000in}{-0.020833in}}%
\pgfpathlineto{\pgfqpoint{0.000000in}{0.020833in}}%
\pgfusepath{stroke,fill}%
}%
\begin{pgfscope}%
\pgfsys@transformshift{1.722967in}{3.252984in}%
\pgfsys@useobject{currentmarker}{}%
\end{pgfscope}%
\begin{pgfscope}%
\pgfsys@transformshift{1.681865in}{3.252984in}%
\pgfsys@useobject{currentmarker}{}%
\end{pgfscope}%
\begin{pgfscope}%
\pgfsys@transformshift{1.693750in}{3.252984in}%
\pgfsys@useobject{currentmarker}{}%
\end{pgfscope}%
\begin{pgfscope}%
\pgfsys@transformshift{1.678059in}{3.252984in}%
\pgfsys@useobject{currentmarker}{}%
\end{pgfscope}%
\end{pgfscope}%
\begin{pgfscope}%
\pgfpathrectangle{\pgfqpoint{1.198800in}{0.720000in}}{\pgfqpoint{0.550375in}{2.664000in}}%
\pgfusepath{clip}%
\pgfsetrectcap%
\pgfsetroundjoin%
\pgfsetlinewidth{0.702625pt}%
\definecolor{currentstroke}{rgb}{0.101961,0.101961,0.101961}%
\pgfsetstrokecolor{currentstroke}%
\pgfsetdash{}{0pt}%
\pgfpathmoveto{\pgfqpoint{1.627940in}{3.034623in}}%
\pgfpathlineto{\pgfqpoint{1.669495in}{3.034623in}}%
\pgfusepath{stroke}%
\end{pgfscope}%
\begin{pgfscope}%
\pgfpathrectangle{\pgfqpoint{1.198800in}{0.720000in}}{\pgfqpoint{0.550375in}{2.664000in}}%
\pgfusepath{clip}%
\pgfsetbuttcap%
\pgfsetroundjoin%
\definecolor{currentfill}{rgb}{0.101961,0.101961,0.101961}%
\pgfsetfillcolor{currentfill}%
\pgfsetlinewidth{0.501875pt}%
\definecolor{currentstroke}{rgb}{0.101961,0.101961,0.101961}%
\pgfsetstrokecolor{currentstroke}%
\pgfsetdash{}{0pt}%
\pgfsys@defobject{currentmarker}{\pgfqpoint{0.000000in}{-0.020833in}}{\pgfqpoint{0.000000in}{0.020833in}}{%
\pgfpathmoveto{\pgfqpoint{0.000000in}{-0.020833in}}%
\pgfpathlineto{\pgfqpoint{0.000000in}{0.020833in}}%
\pgfusepath{stroke,fill}%
}%
\begin{pgfscope}%
\pgfsys@transformshift{1.641607in}{3.034623in}%
\pgfsys@useobject{currentmarker}{}%
\end{pgfscope}%
\begin{pgfscope}%
\pgfsys@transformshift{1.669495in}{3.034623in}%
\pgfsys@useobject{currentmarker}{}%
\end{pgfscope}%
\begin{pgfscope}%
\pgfsys@transformshift{1.627940in}{3.034623in}%
\pgfsys@useobject{currentmarker}{}%
\end{pgfscope}%
\begin{pgfscope}%
\pgfsys@transformshift{1.646798in}{3.034623in}%
\pgfsys@useobject{currentmarker}{}%
\end{pgfscope}%
\end{pgfscope}%
\begin{pgfscope}%
\pgfpathrectangle{\pgfqpoint{1.198800in}{0.720000in}}{\pgfqpoint{0.550375in}{2.664000in}}%
\pgfusepath{clip}%
\pgfsetrectcap%
\pgfsetroundjoin%
\pgfsetlinewidth{0.702625pt}%
\definecolor{currentstroke}{rgb}{0.101961,0.101961,0.101961}%
\pgfsetstrokecolor{currentstroke}%
\pgfsetdash{}{0pt}%
\pgfpathmoveto{\pgfqpoint{1.627158in}{2.816262in}}%
\pgfpathlineto{\pgfqpoint{1.664161in}{2.816262in}}%
\pgfusepath{stroke}%
\end{pgfscope}%
\begin{pgfscope}%
\pgfpathrectangle{\pgfqpoint{1.198800in}{0.720000in}}{\pgfqpoint{0.550375in}{2.664000in}}%
\pgfusepath{clip}%
\pgfsetbuttcap%
\pgfsetroundjoin%
\definecolor{currentfill}{rgb}{0.101961,0.101961,0.101961}%
\pgfsetfillcolor{currentfill}%
\pgfsetlinewidth{0.501875pt}%
\definecolor{currentstroke}{rgb}{0.101961,0.101961,0.101961}%
\pgfsetstrokecolor{currentstroke}%
\pgfsetdash{}{0pt}%
\pgfsys@defobject{currentmarker}{\pgfqpoint{0.000000in}{-0.020833in}}{\pgfqpoint{0.000000in}{0.020833in}}{%
\pgfpathmoveto{\pgfqpoint{0.000000in}{-0.020833in}}%
\pgfpathlineto{\pgfqpoint{0.000000in}{0.020833in}}%
\pgfusepath{stroke,fill}%
}%
\begin{pgfscope}%
\pgfsys@transformshift{1.627158in}{2.816262in}%
\pgfsys@useobject{currentmarker}{}%
\end{pgfscope}%
\begin{pgfscope}%
\pgfsys@transformshift{1.664161in}{2.816262in}%
\pgfsys@useobject{currentmarker}{}%
\end{pgfscope}%
\begin{pgfscope}%
\pgfsys@transformshift{1.630721in}{2.816262in}%
\pgfsys@useobject{currentmarker}{}%
\end{pgfscope}%
\begin{pgfscope}%
\pgfsys@transformshift{1.654678in}{2.816262in}%
\pgfsys@useobject{currentmarker}{}%
\end{pgfscope}%
\end{pgfscope}%
\begin{pgfscope}%
\pgfpathrectangle{\pgfqpoint{1.198800in}{0.720000in}}{\pgfqpoint{0.550375in}{2.664000in}}%
\pgfusepath{clip}%
\pgfsetrectcap%
\pgfsetroundjoin%
\pgfsetlinewidth{0.702625pt}%
\definecolor{currentstroke}{rgb}{0.101961,0.101961,0.101961}%
\pgfsetstrokecolor{currentstroke}%
\pgfsetdash{}{0pt}%
\pgfpathmoveto{\pgfqpoint{1.614871in}{2.597902in}}%
\pgfpathlineto{\pgfqpoint{1.614871in}{2.597902in}}%
\pgfusepath{stroke}%
\end{pgfscope}%
\begin{pgfscope}%
\pgfpathrectangle{\pgfqpoint{1.198800in}{0.720000in}}{\pgfqpoint{0.550375in}{2.664000in}}%
\pgfusepath{clip}%
\pgfsetbuttcap%
\pgfsetroundjoin%
\definecolor{currentfill}{rgb}{0.101961,0.101961,0.101961}%
\pgfsetfillcolor{currentfill}%
\pgfsetlinewidth{0.501875pt}%
\definecolor{currentstroke}{rgb}{0.101961,0.101961,0.101961}%
\pgfsetstrokecolor{currentstroke}%
\pgfsetdash{}{0pt}%
\pgfsys@defobject{currentmarker}{\pgfqpoint{0.000000in}{-0.020833in}}{\pgfqpoint{0.000000in}{0.020833in}}{%
\pgfpathmoveto{\pgfqpoint{0.000000in}{-0.020833in}}%
\pgfpathlineto{\pgfqpoint{0.000000in}{0.020833in}}%
\pgfusepath{stroke,fill}%
}%
\begin{pgfscope}%
\pgfsys@transformshift{1.614871in}{2.597902in}%
\pgfsys@useobject{currentmarker}{}%
\end{pgfscope}%
\end{pgfscope}%
\begin{pgfscope}%
\pgfpathrectangle{\pgfqpoint{1.198800in}{0.720000in}}{\pgfqpoint{0.550375in}{2.664000in}}%
\pgfusepath{clip}%
\pgfsetrectcap%
\pgfsetroundjoin%
\pgfsetlinewidth{0.702625pt}%
\definecolor{currentstroke}{rgb}{0.101961,0.101961,0.101961}%
\pgfsetstrokecolor{currentstroke}%
\pgfsetdash{}{0pt}%
\pgfpathmoveto{\pgfqpoint{1.573304in}{2.379541in}}%
\pgfpathlineto{\pgfqpoint{1.646842in}{2.379541in}}%
\pgfusepath{stroke}%
\end{pgfscope}%
\begin{pgfscope}%
\pgfpathrectangle{\pgfqpoint{1.198800in}{0.720000in}}{\pgfqpoint{0.550375in}{2.664000in}}%
\pgfusepath{clip}%
\pgfsetbuttcap%
\pgfsetroundjoin%
\definecolor{currentfill}{rgb}{0.101961,0.101961,0.101961}%
\pgfsetfillcolor{currentfill}%
\pgfsetlinewidth{0.501875pt}%
\definecolor{currentstroke}{rgb}{0.101961,0.101961,0.101961}%
\pgfsetstrokecolor{currentstroke}%
\pgfsetdash{}{0pt}%
\pgfsys@defobject{currentmarker}{\pgfqpoint{0.000000in}{-0.020833in}}{\pgfqpoint{0.000000in}{0.020833in}}{%
\pgfpathmoveto{\pgfqpoint{0.000000in}{-0.020833in}}%
\pgfpathlineto{\pgfqpoint{0.000000in}{0.020833in}}%
\pgfusepath{stroke,fill}%
}%
\begin{pgfscope}%
\pgfsys@transformshift{1.646842in}{2.379541in}%
\pgfsys@useobject{currentmarker}{}%
\end{pgfscope}%
\begin{pgfscope}%
\pgfsys@transformshift{1.622275in}{2.379541in}%
\pgfsys@useobject{currentmarker}{}%
\end{pgfscope}%
\begin{pgfscope}%
\pgfsys@transformshift{1.573304in}{2.379541in}%
\pgfsys@useobject{currentmarker}{}%
\end{pgfscope}%
\begin{pgfscope}%
\pgfsys@transformshift{1.604550in}{2.379541in}%
\pgfsys@useobject{currentmarker}{}%
\end{pgfscope}%
\end{pgfscope}%
\begin{pgfscope}%
\pgfpathrectangle{\pgfqpoint{1.198800in}{0.720000in}}{\pgfqpoint{0.550375in}{2.664000in}}%
\pgfusepath{clip}%
\pgfsetrectcap%
\pgfsetroundjoin%
\pgfsetlinewidth{0.702625pt}%
\definecolor{currentstroke}{rgb}{0.101961,0.101961,0.101961}%
\pgfsetstrokecolor{currentstroke}%
\pgfsetdash{}{0pt}%
\pgfpathmoveto{\pgfqpoint{1.572614in}{2.161180in}}%
\pgfpathlineto{\pgfqpoint{1.637380in}{2.161180in}}%
\pgfusepath{stroke}%
\end{pgfscope}%
\begin{pgfscope}%
\pgfpathrectangle{\pgfqpoint{1.198800in}{0.720000in}}{\pgfqpoint{0.550375in}{2.664000in}}%
\pgfusepath{clip}%
\pgfsetbuttcap%
\pgfsetroundjoin%
\definecolor{currentfill}{rgb}{0.101961,0.101961,0.101961}%
\pgfsetfillcolor{currentfill}%
\pgfsetlinewidth{0.501875pt}%
\definecolor{currentstroke}{rgb}{0.101961,0.101961,0.101961}%
\pgfsetstrokecolor{currentstroke}%
\pgfsetdash{}{0pt}%
\pgfsys@defobject{currentmarker}{\pgfqpoint{0.000000in}{-0.020833in}}{\pgfqpoint{0.000000in}{0.020833in}}{%
\pgfpathmoveto{\pgfqpoint{0.000000in}{-0.020833in}}%
\pgfpathlineto{\pgfqpoint{0.000000in}{0.020833in}}%
\pgfusepath{stroke,fill}%
}%
\begin{pgfscope}%
\pgfsys@transformshift{1.637380in}{2.161180in}%
\pgfsys@useobject{currentmarker}{}%
\end{pgfscope}%
\begin{pgfscope}%
\pgfsys@transformshift{1.624257in}{2.161180in}%
\pgfsys@useobject{currentmarker}{}%
\end{pgfscope}%
\begin{pgfscope}%
\pgfsys@transformshift{1.572614in}{2.161180in}%
\pgfsys@useobject{currentmarker}{}%
\end{pgfscope}%
\begin{pgfscope}%
\pgfsys@transformshift{1.597270in}{2.161180in}%
\pgfsys@useobject{currentmarker}{}%
\end{pgfscope}%
\end{pgfscope}%
\begin{pgfscope}%
\pgfpathrectangle{\pgfqpoint{1.198800in}{0.720000in}}{\pgfqpoint{0.550375in}{2.664000in}}%
\pgfusepath{clip}%
\pgfsetrectcap%
\pgfsetroundjoin%
\pgfsetlinewidth{0.702625pt}%
\definecolor{currentstroke}{rgb}{0.101961,0.101961,0.101961}%
\pgfsetstrokecolor{currentstroke}%
\pgfsetdash{}{0pt}%
\pgfpathmoveto{\pgfqpoint{1.544219in}{1.942820in}}%
\pgfpathlineto{\pgfqpoint{1.604931in}{1.942820in}}%
\pgfusepath{stroke}%
\end{pgfscope}%
\begin{pgfscope}%
\pgfpathrectangle{\pgfqpoint{1.198800in}{0.720000in}}{\pgfqpoint{0.550375in}{2.664000in}}%
\pgfusepath{clip}%
\pgfsetbuttcap%
\pgfsetroundjoin%
\definecolor{currentfill}{rgb}{0.101961,0.101961,0.101961}%
\pgfsetfillcolor{currentfill}%
\pgfsetlinewidth{0.501875pt}%
\definecolor{currentstroke}{rgb}{0.101961,0.101961,0.101961}%
\pgfsetstrokecolor{currentstroke}%
\pgfsetdash{}{0pt}%
\pgfsys@defobject{currentmarker}{\pgfqpoint{0.000000in}{-0.020833in}}{\pgfqpoint{0.000000in}{0.020833in}}{%
\pgfpathmoveto{\pgfqpoint{0.000000in}{-0.020833in}}%
\pgfpathlineto{\pgfqpoint{0.000000in}{0.020833in}}%
\pgfusepath{stroke,fill}%
}%
\begin{pgfscope}%
\pgfsys@transformshift{1.575874in}{1.942820in}%
\pgfsys@useobject{currentmarker}{}%
\end{pgfscope}%
\begin{pgfscope}%
\pgfsys@transformshift{1.544219in}{1.942820in}%
\pgfsys@useobject{currentmarker}{}%
\end{pgfscope}%
\begin{pgfscope}%
\pgfsys@transformshift{1.604931in}{1.942820in}%
\pgfsys@useobject{currentmarker}{}%
\end{pgfscope}%
\end{pgfscope}%
\begin{pgfscope}%
\pgfpathrectangle{\pgfqpoint{1.198800in}{0.720000in}}{\pgfqpoint{0.550375in}{2.664000in}}%
\pgfusepath{clip}%
\pgfsetrectcap%
\pgfsetroundjoin%
\pgfsetlinewidth{0.702625pt}%
\definecolor{currentstroke}{rgb}{0.101961,0.101961,0.101961}%
\pgfsetstrokecolor{currentstroke}%
\pgfsetdash{}{0pt}%
\pgfpathmoveto{\pgfqpoint{1.534452in}{1.724459in}}%
\pgfpathlineto{\pgfqpoint{1.610348in}{1.724459in}}%
\pgfusepath{stroke}%
\end{pgfscope}%
\begin{pgfscope}%
\pgfpathrectangle{\pgfqpoint{1.198800in}{0.720000in}}{\pgfqpoint{0.550375in}{2.664000in}}%
\pgfusepath{clip}%
\pgfsetbuttcap%
\pgfsetroundjoin%
\definecolor{currentfill}{rgb}{0.101961,0.101961,0.101961}%
\pgfsetfillcolor{currentfill}%
\pgfsetlinewidth{0.501875pt}%
\definecolor{currentstroke}{rgb}{0.101961,0.101961,0.101961}%
\pgfsetstrokecolor{currentstroke}%
\pgfsetdash{}{0pt}%
\pgfsys@defobject{currentmarker}{\pgfqpoint{0.000000in}{-0.020833in}}{\pgfqpoint{0.000000in}{0.020833in}}{%
\pgfpathmoveto{\pgfqpoint{0.000000in}{-0.020833in}}%
\pgfpathlineto{\pgfqpoint{0.000000in}{0.020833in}}%
\pgfusepath{stroke,fill}%
}%
\begin{pgfscope}%
\pgfsys@transformshift{1.578597in}{1.724459in}%
\pgfsys@useobject{currentmarker}{}%
\end{pgfscope}%
\begin{pgfscope}%
\pgfsys@transformshift{1.568714in}{1.724459in}%
\pgfsys@useobject{currentmarker}{}%
\end{pgfscope}%
\begin{pgfscope}%
\pgfsys@transformshift{1.534452in}{1.724459in}%
\pgfsys@useobject{currentmarker}{}%
\end{pgfscope}%
\begin{pgfscope}%
\pgfsys@transformshift{1.610348in}{1.724459in}%
\pgfsys@useobject{currentmarker}{}%
\end{pgfscope}%
\end{pgfscope}%
\begin{pgfscope}%
\pgfpathrectangle{\pgfqpoint{1.198800in}{0.720000in}}{\pgfqpoint{0.550375in}{2.664000in}}%
\pgfusepath{clip}%
\pgfsetrectcap%
\pgfsetroundjoin%
\pgfsetlinewidth{0.702625pt}%
\definecolor{currentstroke}{rgb}{0.101961,0.101961,0.101961}%
\pgfsetstrokecolor{currentstroke}%
\pgfsetdash{}{0pt}%
\pgfpathmoveto{\pgfqpoint{1.489351in}{1.506098in}}%
\pgfpathlineto{\pgfqpoint{1.543743in}{1.506098in}}%
\pgfusepath{stroke}%
\end{pgfscope}%
\begin{pgfscope}%
\pgfpathrectangle{\pgfqpoint{1.198800in}{0.720000in}}{\pgfqpoint{0.550375in}{2.664000in}}%
\pgfusepath{clip}%
\pgfsetbuttcap%
\pgfsetroundjoin%
\definecolor{currentfill}{rgb}{0.101961,0.101961,0.101961}%
\pgfsetfillcolor{currentfill}%
\pgfsetlinewidth{0.501875pt}%
\definecolor{currentstroke}{rgb}{0.101961,0.101961,0.101961}%
\pgfsetstrokecolor{currentstroke}%
\pgfsetdash{}{0pt}%
\pgfsys@defobject{currentmarker}{\pgfqpoint{0.000000in}{-0.020833in}}{\pgfqpoint{0.000000in}{0.020833in}}{%
\pgfpathmoveto{\pgfqpoint{0.000000in}{-0.020833in}}%
\pgfpathlineto{\pgfqpoint{0.000000in}{0.020833in}}%
\pgfusepath{stroke,fill}%
}%
\begin{pgfscope}%
\pgfsys@transformshift{1.543743in}{1.506098in}%
\pgfsys@useobject{currentmarker}{}%
\end{pgfscope}%
\begin{pgfscope}%
\pgfsys@transformshift{1.528002in}{1.506098in}%
\pgfsys@useobject{currentmarker}{}%
\end{pgfscope}%
\begin{pgfscope}%
\pgfsys@transformshift{1.489351in}{1.506098in}%
\pgfsys@useobject{currentmarker}{}%
\end{pgfscope}%
\begin{pgfscope}%
\pgfsys@transformshift{1.516551in}{1.506098in}%
\pgfsys@useobject{currentmarker}{}%
\end{pgfscope}%
\end{pgfscope}%
\begin{pgfscope}%
\pgfpathrectangle{\pgfqpoint{1.198800in}{0.720000in}}{\pgfqpoint{0.550375in}{2.664000in}}%
\pgfusepath{clip}%
\pgfsetrectcap%
\pgfsetroundjoin%
\pgfsetlinewidth{0.702625pt}%
\definecolor{currentstroke}{rgb}{0.101961,0.101961,0.101961}%
\pgfsetstrokecolor{currentstroke}%
\pgfsetdash{}{0pt}%
\pgfpathmoveto{\pgfqpoint{1.487925in}{1.287738in}}%
\pgfpathlineto{\pgfqpoint{1.543758in}{1.287738in}}%
\pgfusepath{stroke}%
\end{pgfscope}%
\begin{pgfscope}%
\pgfpathrectangle{\pgfqpoint{1.198800in}{0.720000in}}{\pgfqpoint{0.550375in}{2.664000in}}%
\pgfusepath{clip}%
\pgfsetbuttcap%
\pgfsetroundjoin%
\definecolor{currentfill}{rgb}{0.101961,0.101961,0.101961}%
\pgfsetfillcolor{currentfill}%
\pgfsetlinewidth{0.501875pt}%
\definecolor{currentstroke}{rgb}{0.101961,0.101961,0.101961}%
\pgfsetstrokecolor{currentstroke}%
\pgfsetdash{}{0pt}%
\pgfsys@defobject{currentmarker}{\pgfqpoint{0.000000in}{-0.020833in}}{\pgfqpoint{0.000000in}{0.020833in}}{%
\pgfpathmoveto{\pgfqpoint{0.000000in}{-0.020833in}}%
\pgfpathlineto{\pgfqpoint{0.000000in}{0.020833in}}%
\pgfusepath{stroke,fill}%
}%
\begin{pgfscope}%
\pgfsys@transformshift{1.543758in}{1.287738in}%
\pgfsys@useobject{currentmarker}{}%
\end{pgfscope}%
\begin{pgfscope}%
\pgfsys@transformshift{1.529430in}{1.287738in}%
\pgfsys@useobject{currentmarker}{}%
\end{pgfscope}%
\begin{pgfscope}%
\pgfsys@transformshift{1.487925in}{1.287738in}%
\pgfsys@useobject{currentmarker}{}%
\end{pgfscope}%
\begin{pgfscope}%
\pgfsys@transformshift{1.515117in}{1.287738in}%
\pgfsys@useobject{currentmarker}{}%
\end{pgfscope}%
\end{pgfscope}%
\begin{pgfscope}%
\pgfpathrectangle{\pgfqpoint{1.198800in}{0.720000in}}{\pgfqpoint{0.550375in}{2.664000in}}%
\pgfusepath{clip}%
\pgfsetrectcap%
\pgfsetroundjoin%
\pgfsetlinewidth{0.702625pt}%
\definecolor{currentstroke}{rgb}{0.101961,0.101961,0.101961}%
\pgfsetstrokecolor{currentstroke}%
\pgfsetdash{}{0pt}%
\pgfpathmoveto{\pgfqpoint{1.485058in}{1.069377in}}%
\pgfpathlineto{\pgfqpoint{1.542326in}{1.069377in}}%
\pgfusepath{stroke}%
\end{pgfscope}%
\begin{pgfscope}%
\pgfpathrectangle{\pgfqpoint{1.198800in}{0.720000in}}{\pgfqpoint{0.550375in}{2.664000in}}%
\pgfusepath{clip}%
\pgfsetbuttcap%
\pgfsetroundjoin%
\definecolor{currentfill}{rgb}{0.101961,0.101961,0.101961}%
\pgfsetfillcolor{currentfill}%
\pgfsetlinewidth{0.501875pt}%
\definecolor{currentstroke}{rgb}{0.101961,0.101961,0.101961}%
\pgfsetstrokecolor{currentstroke}%
\pgfsetdash{}{0pt}%
\pgfsys@defobject{currentmarker}{\pgfqpoint{0.000000in}{-0.020833in}}{\pgfqpoint{0.000000in}{0.020833in}}{%
\pgfpathmoveto{\pgfqpoint{0.000000in}{-0.020833in}}%
\pgfpathlineto{\pgfqpoint{0.000000in}{0.020833in}}%
\pgfusepath{stroke,fill}%
}%
\begin{pgfscope}%
\pgfsys@transformshift{1.542326in}{1.069377in}%
\pgfsys@useobject{currentmarker}{}%
\end{pgfscope}%
\begin{pgfscope}%
\pgfsys@transformshift{1.523713in}{1.069377in}%
\pgfsys@useobject{currentmarker}{}%
\end{pgfscope}%
\begin{pgfscope}%
\pgfsys@transformshift{1.485058in}{1.069377in}%
\pgfsys@useobject{currentmarker}{}%
\end{pgfscope}%
\begin{pgfscope}%
\pgfsys@transformshift{1.516549in}{1.069377in}%
\pgfsys@useobject{currentmarker}{}%
\end{pgfscope}%
\end{pgfscope}%
\begin{pgfscope}%
\pgfpathrectangle{\pgfqpoint{1.198800in}{0.720000in}}{\pgfqpoint{0.550375in}{2.664000in}}%
\pgfusepath{clip}%
\pgfsetrectcap%
\pgfsetroundjoin%
\pgfsetlinewidth{0.702625pt}%
\definecolor{currentstroke}{rgb}{0.101961,0.101961,0.101961}%
\pgfsetstrokecolor{currentstroke}%
\pgfsetdash{}{0pt}%
\pgfpathmoveto{\pgfqpoint{1.483626in}{0.851016in}}%
\pgfpathlineto{\pgfqpoint{1.540875in}{0.851016in}}%
\pgfusepath{stroke}%
\end{pgfscope}%
\begin{pgfscope}%
\pgfpathrectangle{\pgfqpoint{1.198800in}{0.720000in}}{\pgfqpoint{0.550375in}{2.664000in}}%
\pgfusepath{clip}%
\pgfsetbuttcap%
\pgfsetroundjoin%
\definecolor{currentfill}{rgb}{0.101961,0.101961,0.101961}%
\pgfsetfillcolor{currentfill}%
\pgfsetlinewidth{0.501875pt}%
\definecolor{currentstroke}{rgb}{0.101961,0.101961,0.101961}%
\pgfsetstrokecolor{currentstroke}%
\pgfsetdash{}{0pt}%
\pgfsys@defobject{currentmarker}{\pgfqpoint{0.000000in}{-0.020833in}}{\pgfqpoint{0.000000in}{0.020833in}}{%
\pgfpathmoveto{\pgfqpoint{0.000000in}{-0.020833in}}%
\pgfpathlineto{\pgfqpoint{0.000000in}{0.020833in}}%
\pgfusepath{stroke,fill}%
}%
\begin{pgfscope}%
\pgfsys@transformshift{1.540875in}{0.851016in}%
\pgfsys@useobject{currentmarker}{}%
\end{pgfscope}%
\begin{pgfscope}%
\pgfsys@transformshift{1.526565in}{0.851016in}%
\pgfsys@useobject{currentmarker}{}%
\end{pgfscope}%
\begin{pgfscope}%
\pgfsys@transformshift{1.483626in}{0.851016in}%
\pgfsys@useobject{currentmarker}{}%
\end{pgfscope}%
\begin{pgfscope}%
\pgfsys@transformshift{1.516546in}{0.851016in}%
\pgfsys@useobject{currentmarker}{}%
\end{pgfscope}%
\end{pgfscope}%
\begin{pgfscope}%
\pgfsetrectcap%
\pgfsetmiterjoin%
\pgfsetlinewidth{0.803000pt}%
\definecolor{currentstroke}{rgb}{0.101961,0.101961,0.101961}%
\pgfsetstrokecolor{currentstroke}%
\pgfsetdash{}{0pt}%
\pgfpathmoveto{\pgfqpoint{1.198800in}{0.720000in}}%
\pgfpathlineto{\pgfqpoint{1.198800in}{3.384000in}}%
\pgfusepath{stroke}%
\end{pgfscope}%
\begin{pgfscope}%
\pgfsetrectcap%
\pgfsetmiterjoin%
\pgfsetlinewidth{0.803000pt}%
\definecolor{currentstroke}{rgb}{0.101961,0.101961,0.101961}%
\pgfsetstrokecolor{currentstroke}%
\pgfsetdash{}{0pt}%
\pgfpathmoveto{\pgfqpoint{1.198800in}{0.720000in}}%
\pgfpathlineto{\pgfqpoint{1.749175in}{0.720000in}}%
\pgfusepath{stroke}%
\end{pgfscope}%
\begin{pgfscope}%
\pgfsetbuttcap%
\pgfsetmiterjoin%
\definecolor{currentfill}{rgb}{1.000000,1.000000,1.000000}%
\pgfsetfillcolor{currentfill}%
\pgfsetlinewidth{0.000000pt}%
\definecolor{currentstroke}{rgb}{0.000000,0.000000,0.000000}%
\pgfsetstrokecolor{currentstroke}%
\pgfsetstrokeopacity{0.000000}%
\pgfsetdash{}{0pt}%
\pgfpathmoveto{\pgfqpoint{1.914287in}{0.720000in}}%
\pgfpathlineto{\pgfqpoint{2.464663in}{0.720000in}}%
\pgfpathlineto{\pgfqpoint{2.464663in}{3.384000in}}%
\pgfpathlineto{\pgfqpoint{1.914287in}{3.384000in}}%
\pgfpathlineto{\pgfqpoint{1.914287in}{0.720000in}}%
\pgfpathclose%
\pgfusepath{fill}%
\end{pgfscope}%
\begin{pgfscope}%
\pgfpathrectangle{\pgfqpoint{1.914287in}{0.720000in}}{\pgfqpoint{0.550375in}{2.664000in}}%
\pgfusepath{clip}%
\pgfsetrectcap%
\pgfsetroundjoin%
\pgfsetlinewidth{0.401500pt}%
\definecolor{currentstroke}{rgb}{0.901961,0.901961,0.901961}%
\pgfsetstrokecolor{currentstroke}%
\pgfsetdash{}{0pt}%
\pgfpathmoveto{\pgfqpoint{1.914287in}{0.720000in}}%
\pgfpathlineto{\pgfqpoint{1.914287in}{3.384000in}}%
\pgfusepath{stroke}%
\end{pgfscope}%
\begin{pgfscope}%
\definecolor{textcolor}{rgb}{0.101961,0.101961,0.101961}%
\pgfsetstrokecolor{textcolor}%
\pgfsetfillcolor{textcolor}%
\pgftext[x=1.914287in,y=0.671389in,,top]{\color{textcolor}\sffamily\fontsize{7.000000}{8.400000}\selectfont 0.0}%
\end{pgfscope}%
\begin{pgfscope}%
\pgfpathrectangle{\pgfqpoint{1.914287in}{0.720000in}}{\pgfqpoint{0.550375in}{2.664000in}}%
\pgfusepath{clip}%
\pgfsetrectcap%
\pgfsetroundjoin%
\pgfsetlinewidth{0.401500pt}%
\definecolor{currentstroke}{rgb}{0.901961,0.901961,0.901961}%
\pgfsetstrokecolor{currentstroke}%
\pgfsetdash{}{0pt}%
\pgfpathmoveto{\pgfqpoint{2.189475in}{0.720000in}}%
\pgfpathlineto{\pgfqpoint{2.189475in}{3.384000in}}%
\pgfusepath{stroke}%
\end{pgfscope}%
\begin{pgfscope}%
\definecolor{textcolor}{rgb}{0.101961,0.101961,0.101961}%
\pgfsetstrokecolor{textcolor}%
\pgfsetfillcolor{textcolor}%
\pgftext[x=2.189475in,y=0.671389in,,top]{\color{textcolor}\sffamily\fontsize{7.000000}{8.400000}\selectfont 0.5}%
\end{pgfscope}%
\begin{pgfscope}%
\pgfpathrectangle{\pgfqpoint{1.914287in}{0.720000in}}{\pgfqpoint{0.550375in}{2.664000in}}%
\pgfusepath{clip}%
\pgfsetrectcap%
\pgfsetroundjoin%
\pgfsetlinewidth{0.401500pt}%
\definecolor{currentstroke}{rgb}{0.901961,0.901961,0.901961}%
\pgfsetstrokecolor{currentstroke}%
\pgfsetdash{}{0pt}%
\pgfpathmoveto{\pgfqpoint{2.464663in}{0.720000in}}%
\pgfpathlineto{\pgfqpoint{2.464663in}{3.384000in}}%
\pgfusepath{stroke}%
\end{pgfscope}%
\begin{pgfscope}%
\definecolor{textcolor}{rgb}{0.101961,0.101961,0.101961}%
\pgfsetstrokecolor{textcolor}%
\pgfsetfillcolor{textcolor}%
\pgftext[x=2.464663in,y=0.671389in,,top]{\color{textcolor}\sffamily\fontsize{7.000000}{8.400000}\selectfont 1.0}%
\end{pgfscope}%
\begin{pgfscope}%
\definecolor{textcolor}{rgb}{0.000000,0.000000,0.000000}%
\pgfsetstrokecolor{textcolor}%
\pgfsetfillcolor{textcolor}%
\pgftext[x=1.972444in, y=0.461899in, left, base]{\color{textcolor}\sffamily\fontsize{7.500000}{9.000000}\selectfont GPU busy}%
\end{pgfscope}%
\begin{pgfscope}%
\definecolor{textcolor}{rgb}{0.000000,0.000000,0.000000}%
\pgfsetstrokecolor{textcolor}%
\pgfsetfillcolor{textcolor}%
\pgftext[x=2.024820in, y=0.361976in, left, base]{\color{textcolor}\sffamily\fontsize{7.500000}{9.000000}\selectfont fraction}%
\end{pgfscope}%
\begin{pgfscope}%
\pgfpathrectangle{\pgfqpoint{1.914287in}{0.720000in}}{\pgfqpoint{0.550375in}{2.664000in}}%
\pgfusepath{clip}%
\pgfsetbuttcap%
\pgfsetmiterjoin%
\definecolor{currentfill}{rgb}{0.156863,0.356863,0.568627}%
\pgfsetfillcolor{currentfill}%
\pgfsetlinewidth{0.000000pt}%
\definecolor{currentstroke}{rgb}{0.000000,0.000000,0.000000}%
\pgfsetstrokecolor{currentstroke}%
\pgfsetstrokeopacity{0.000000}%
\pgfsetdash{}{0pt}%
\pgfpathmoveto{\pgfqpoint{1.914287in}{3.176557in}}%
\pgfpathlineto{\pgfqpoint{2.353875in}{3.176557in}}%
\pgfpathlineto{\pgfqpoint{2.353875in}{3.329410in}}%
\pgfpathlineto{\pgfqpoint{1.914287in}{3.329410in}}%
\pgfpathlineto{\pgfqpoint{1.914287in}{3.176557in}}%
\pgfpathclose%
\pgfusepath{fill}%
\end{pgfscope}%
\begin{pgfscope}%
\pgfpathrectangle{\pgfqpoint{1.914287in}{0.720000in}}{\pgfqpoint{0.550375in}{2.664000in}}%
\pgfusepath{clip}%
\pgfsetbuttcap%
\pgfsetmiterjoin%
\definecolor{currentfill}{rgb}{0.784314,0.784314,0.784314}%
\pgfsetfillcolor{currentfill}%
\pgfsetlinewidth{0.000000pt}%
\definecolor{currentstroke}{rgb}{0.000000,0.000000,0.000000}%
\pgfsetstrokecolor{currentstroke}%
\pgfsetstrokeopacity{0.000000}%
\pgfsetdash{}{0pt}%
\pgfpathmoveto{\pgfqpoint{1.914287in}{2.958197in}}%
\pgfpathlineto{\pgfqpoint{2.336175in}{2.958197in}}%
\pgfpathlineto{\pgfqpoint{2.336175in}{3.111049in}}%
\pgfpathlineto{\pgfqpoint{1.914287in}{3.111049in}}%
\pgfpathlineto{\pgfqpoint{1.914287in}{2.958197in}}%
\pgfpathclose%
\pgfusepath{fill}%
\end{pgfscope}%
\begin{pgfscope}%
\pgfpathrectangle{\pgfqpoint{1.914287in}{0.720000in}}{\pgfqpoint{0.550375in}{2.664000in}}%
\pgfusepath{clip}%
\pgfsetbuttcap%
\pgfsetmiterjoin%
\definecolor{currentfill}{rgb}{0.784314,0.784314,0.784314}%
\pgfsetfillcolor{currentfill}%
\pgfsetlinewidth{0.000000pt}%
\definecolor{currentstroke}{rgb}{0.000000,0.000000,0.000000}%
\pgfsetstrokecolor{currentstroke}%
\pgfsetstrokeopacity{0.000000}%
\pgfsetdash{}{0pt}%
\pgfpathmoveto{\pgfqpoint{1.914287in}{2.739836in}}%
\pgfpathlineto{\pgfqpoint{2.340553in}{2.739836in}}%
\pgfpathlineto{\pgfqpoint{2.340553in}{2.892689in}}%
\pgfpathlineto{\pgfqpoint{1.914287in}{2.892689in}}%
\pgfpathlineto{\pgfqpoint{1.914287in}{2.739836in}}%
\pgfpathclose%
\pgfusepath{fill}%
\end{pgfscope}%
\begin{pgfscope}%
\pgfpathrectangle{\pgfqpoint{1.914287in}{0.720000in}}{\pgfqpoint{0.550375in}{2.664000in}}%
\pgfusepath{clip}%
\pgfsetbuttcap%
\pgfsetmiterjoin%
\definecolor{currentfill}{rgb}{0.478431,0.478431,0.478431}%
\pgfsetfillcolor{currentfill}%
\pgfsetlinewidth{0.000000pt}%
\definecolor{currentstroke}{rgb}{0.000000,0.000000,0.000000}%
\pgfsetstrokecolor{currentstroke}%
\pgfsetstrokeopacity{0.000000}%
\pgfsetdash{}{0pt}%
\pgfpathmoveto{\pgfqpoint{1.914287in}{2.521475in}}%
\pgfpathlineto{\pgfqpoint{2.381927in}{2.521475in}}%
\pgfpathlineto{\pgfqpoint{2.381927in}{2.674328in}}%
\pgfpathlineto{\pgfqpoint{1.914287in}{2.674328in}}%
\pgfpathlineto{\pgfqpoint{1.914287in}{2.521475in}}%
\pgfpathclose%
\pgfusepath{fill}%
\end{pgfscope}%
\begin{pgfscope}%
\pgfpathrectangle{\pgfqpoint{1.914287in}{0.720000in}}{\pgfqpoint{0.550375in}{2.664000in}}%
\pgfusepath{clip}%
\pgfsetbuttcap%
\pgfsetmiterjoin%
\definecolor{currentfill}{rgb}{0.784314,0.784314,0.784314}%
\pgfsetfillcolor{currentfill}%
\pgfsetlinewidth{0.000000pt}%
\definecolor{currentstroke}{rgb}{0.000000,0.000000,0.000000}%
\pgfsetstrokecolor{currentstroke}%
\pgfsetstrokeopacity{0.000000}%
\pgfsetdash{}{0pt}%
\pgfpathmoveto{\pgfqpoint{1.914287in}{2.303115in}}%
\pgfpathlineto{\pgfqpoint{2.360728in}{2.303115in}}%
\pgfpathlineto{\pgfqpoint{2.360728in}{2.455967in}}%
\pgfpathlineto{\pgfqpoint{1.914287in}{2.455967in}}%
\pgfpathlineto{\pgfqpoint{1.914287in}{2.303115in}}%
\pgfpathclose%
\pgfusepath{fill}%
\end{pgfscope}%
\begin{pgfscope}%
\pgfpathrectangle{\pgfqpoint{1.914287in}{0.720000in}}{\pgfqpoint{0.550375in}{2.664000in}}%
\pgfusepath{clip}%
\pgfsetbuttcap%
\pgfsetmiterjoin%
\definecolor{currentfill}{rgb}{0.784314,0.784314,0.784314}%
\pgfsetfillcolor{currentfill}%
\pgfsetlinewidth{0.000000pt}%
\definecolor{currentstroke}{rgb}{0.000000,0.000000,0.000000}%
\pgfsetstrokecolor{currentstroke}%
\pgfsetstrokeopacity{0.000000}%
\pgfsetdash{}{0pt}%
\pgfpathmoveto{\pgfqpoint{1.914287in}{2.084754in}}%
\pgfpathlineto{\pgfqpoint{2.349635in}{2.084754in}}%
\pgfpathlineto{\pgfqpoint{2.349635in}{2.237607in}}%
\pgfpathlineto{\pgfqpoint{1.914287in}{2.237607in}}%
\pgfpathlineto{\pgfqpoint{1.914287in}{2.084754in}}%
\pgfpathclose%
\pgfusepath{fill}%
\end{pgfscope}%
\begin{pgfscope}%
\pgfpathrectangle{\pgfqpoint{1.914287in}{0.720000in}}{\pgfqpoint{0.550375in}{2.664000in}}%
\pgfusepath{clip}%
\pgfsetbuttcap%
\pgfsetmiterjoin%
\definecolor{currentfill}{rgb}{0.784314,0.784314,0.784314}%
\pgfsetfillcolor{currentfill}%
\pgfsetlinewidth{0.000000pt}%
\definecolor{currentstroke}{rgb}{0.000000,0.000000,0.000000}%
\pgfsetstrokecolor{currentstroke}%
\pgfsetstrokeopacity{0.000000}%
\pgfsetdash{}{0pt}%
\pgfpathmoveto{\pgfqpoint{1.914287in}{1.866393in}}%
\pgfpathlineto{\pgfqpoint{2.406459in}{1.866393in}}%
\pgfpathlineto{\pgfqpoint{2.406459in}{2.019246in}}%
\pgfpathlineto{\pgfqpoint{1.914287in}{2.019246in}}%
\pgfpathlineto{\pgfqpoint{1.914287in}{1.866393in}}%
\pgfpathclose%
\pgfusepath{fill}%
\end{pgfscope}%
\begin{pgfscope}%
\pgfpathrectangle{\pgfqpoint{1.914287in}{0.720000in}}{\pgfqpoint{0.550375in}{2.664000in}}%
\pgfusepath{clip}%
\pgfsetbuttcap%
\pgfsetmiterjoin%
\definecolor{currentfill}{rgb}{0.784314,0.784314,0.784314}%
\pgfsetfillcolor{currentfill}%
\pgfsetlinewidth{0.000000pt}%
\definecolor{currentstroke}{rgb}{0.000000,0.000000,0.000000}%
\pgfsetstrokecolor{currentstroke}%
\pgfsetstrokeopacity{0.000000}%
\pgfsetdash{}{0pt}%
\pgfpathmoveto{\pgfqpoint{1.914287in}{1.648033in}}%
\pgfpathlineto{\pgfqpoint{2.403377in}{1.648033in}}%
\pgfpathlineto{\pgfqpoint{2.403377in}{1.800885in}}%
\pgfpathlineto{\pgfqpoint{1.914287in}{1.800885in}}%
\pgfpathlineto{\pgfqpoint{1.914287in}{1.648033in}}%
\pgfpathclose%
\pgfusepath{fill}%
\end{pgfscope}%
\begin{pgfscope}%
\pgfpathrectangle{\pgfqpoint{1.914287in}{0.720000in}}{\pgfqpoint{0.550375in}{2.664000in}}%
\pgfusepath{clip}%
\pgfsetbuttcap%
\pgfsetmiterjoin%
\definecolor{currentfill}{rgb}{0.478431,0.478431,0.478431}%
\pgfsetfillcolor{currentfill}%
\pgfsetlinewidth{0.000000pt}%
\definecolor{currentstroke}{rgb}{0.000000,0.000000,0.000000}%
\pgfsetstrokecolor{currentstroke}%
\pgfsetstrokeopacity{0.000000}%
\pgfsetdash{}{0pt}%
\pgfpathmoveto{\pgfqpoint{1.914287in}{1.429672in}}%
\pgfpathlineto{\pgfqpoint{2.464652in}{1.429672in}}%
\pgfpathlineto{\pgfqpoint{2.464652in}{1.582525in}}%
\pgfpathlineto{\pgfqpoint{1.914287in}{1.582525in}}%
\pgfpathlineto{\pgfqpoint{1.914287in}{1.429672in}}%
\pgfpathclose%
\pgfusepath{fill}%
\end{pgfscope}%
\begin{pgfscope}%
\pgfpathrectangle{\pgfqpoint{1.914287in}{0.720000in}}{\pgfqpoint{0.550375in}{2.664000in}}%
\pgfusepath{clip}%
\pgfsetbuttcap%
\pgfsetmiterjoin%
\definecolor{currentfill}{rgb}{0.784314,0.784314,0.784314}%
\pgfsetfillcolor{currentfill}%
\pgfsetlinewidth{0.000000pt}%
\definecolor{currentstroke}{rgb}{0.000000,0.000000,0.000000}%
\pgfsetstrokecolor{currentstroke}%
\pgfsetstrokeopacity{0.000000}%
\pgfsetdash{}{0pt}%
\pgfpathmoveto{\pgfqpoint{1.914287in}{1.211311in}}%
\pgfpathlineto{\pgfqpoint{2.464646in}{1.211311in}}%
\pgfpathlineto{\pgfqpoint{2.464646in}{1.364164in}}%
\pgfpathlineto{\pgfqpoint{1.914287in}{1.364164in}}%
\pgfpathlineto{\pgfqpoint{1.914287in}{1.211311in}}%
\pgfpathclose%
\pgfusepath{fill}%
\end{pgfscope}%
\begin{pgfscope}%
\pgfpathrectangle{\pgfqpoint{1.914287in}{0.720000in}}{\pgfqpoint{0.550375in}{2.664000in}}%
\pgfusepath{clip}%
\pgfsetbuttcap%
\pgfsetmiterjoin%
\definecolor{currentfill}{rgb}{0.784314,0.784314,0.784314}%
\pgfsetfillcolor{currentfill}%
\pgfsetlinewidth{0.000000pt}%
\definecolor{currentstroke}{rgb}{0.000000,0.000000,0.000000}%
\pgfsetstrokecolor{currentstroke}%
\pgfsetstrokeopacity{0.000000}%
\pgfsetdash{}{0pt}%
\pgfpathmoveto{\pgfqpoint{1.914287in}{0.992951in}}%
\pgfpathlineto{\pgfqpoint{2.464645in}{0.992951in}}%
\pgfpathlineto{\pgfqpoint{2.464645in}{1.145803in}}%
\pgfpathlineto{\pgfqpoint{1.914287in}{1.145803in}}%
\pgfpathlineto{\pgfqpoint{1.914287in}{0.992951in}}%
\pgfpathclose%
\pgfusepath{fill}%
\end{pgfscope}%
\begin{pgfscope}%
\pgfpathrectangle{\pgfqpoint{1.914287in}{0.720000in}}{\pgfqpoint{0.550375in}{2.664000in}}%
\pgfusepath{clip}%
\pgfsetbuttcap%
\pgfsetmiterjoin%
\definecolor{currentfill}{rgb}{0.478431,0.478431,0.478431}%
\pgfsetfillcolor{currentfill}%
\pgfsetlinewidth{0.000000pt}%
\definecolor{currentstroke}{rgb}{0.000000,0.000000,0.000000}%
\pgfsetstrokecolor{currentstroke}%
\pgfsetstrokeopacity{0.000000}%
\pgfsetdash{}{0pt}%
\pgfpathmoveto{\pgfqpoint{1.914287in}{0.774590in}}%
\pgfpathlineto{\pgfqpoint{2.464659in}{0.774590in}}%
\pgfpathlineto{\pgfqpoint{2.464659in}{0.927443in}}%
\pgfpathlineto{\pgfqpoint{1.914287in}{0.927443in}}%
\pgfpathlineto{\pgfqpoint{1.914287in}{0.774590in}}%
\pgfpathclose%
\pgfusepath{fill}%
\end{pgfscope}%
\begin{pgfscope}%
\pgfpathrectangle{\pgfqpoint{1.914287in}{0.720000in}}{\pgfqpoint{0.550375in}{2.664000in}}%
\pgfusepath{clip}%
\pgfsetrectcap%
\pgfsetroundjoin%
\pgfsetlinewidth{0.702625pt}%
\definecolor{currentstroke}{rgb}{0.101961,0.101961,0.101961}%
\pgfsetstrokecolor{currentstroke}%
\pgfsetdash{}{0pt}%
\pgfpathmoveto{\pgfqpoint{2.351317in}{3.252984in}}%
\pgfpathlineto{\pgfqpoint{2.359128in}{3.252984in}}%
\pgfusepath{stroke}%
\end{pgfscope}%
\begin{pgfscope}%
\pgfpathrectangle{\pgfqpoint{1.914287in}{0.720000in}}{\pgfqpoint{0.550375in}{2.664000in}}%
\pgfusepath{clip}%
\pgfsetbuttcap%
\pgfsetroundjoin%
\definecolor{currentfill}{rgb}{0.101961,0.101961,0.101961}%
\pgfsetfillcolor{currentfill}%
\pgfsetlinewidth{0.501875pt}%
\definecolor{currentstroke}{rgb}{0.101961,0.101961,0.101961}%
\pgfsetstrokecolor{currentstroke}%
\pgfsetdash{}{0pt}%
\pgfsys@defobject{currentmarker}{\pgfqpoint{0.000000in}{-0.020833in}}{\pgfqpoint{0.000000in}{0.020833in}}{%
\pgfpathmoveto{\pgfqpoint{0.000000in}{-0.020833in}}%
\pgfpathlineto{\pgfqpoint{0.000000in}{0.020833in}}%
\pgfusepath{stroke,fill}%
}%
\begin{pgfscope}%
\pgfsys@transformshift{2.359128in}{3.252984in}%
\pgfsys@useobject{currentmarker}{}%
\end{pgfscope}%
\begin{pgfscope}%
\pgfsys@transformshift{2.351317in}{3.252984in}%
\pgfsys@useobject{currentmarker}{}%
\end{pgfscope}%
\begin{pgfscope}%
\pgfsys@transformshift{2.353556in}{3.252984in}%
\pgfsys@useobject{currentmarker}{}%
\end{pgfscope}%
\begin{pgfscope}%
\pgfsys@transformshift{2.351501in}{3.252984in}%
\pgfsys@useobject{currentmarker}{}%
\end{pgfscope}%
\end{pgfscope}%
\begin{pgfscope}%
\pgfpathrectangle{\pgfqpoint{1.914287in}{0.720000in}}{\pgfqpoint{0.550375in}{2.664000in}}%
\pgfusepath{clip}%
\pgfsetrectcap%
\pgfsetroundjoin%
\pgfsetlinewidth{0.702625pt}%
\definecolor{currentstroke}{rgb}{0.101961,0.101961,0.101961}%
\pgfsetstrokecolor{currentstroke}%
\pgfsetdash{}{0pt}%
\pgfpathmoveto{\pgfqpoint{2.335902in}{3.034623in}}%
\pgfpathlineto{\pgfqpoint{2.336480in}{3.034623in}}%
\pgfusepath{stroke}%
\end{pgfscope}%
\begin{pgfscope}%
\pgfpathrectangle{\pgfqpoint{1.914287in}{0.720000in}}{\pgfqpoint{0.550375in}{2.664000in}}%
\pgfusepath{clip}%
\pgfsetbuttcap%
\pgfsetroundjoin%
\definecolor{currentfill}{rgb}{0.101961,0.101961,0.101961}%
\pgfsetfillcolor{currentfill}%
\pgfsetlinewidth{0.501875pt}%
\definecolor{currentstroke}{rgb}{0.101961,0.101961,0.101961}%
\pgfsetstrokecolor{currentstroke}%
\pgfsetdash{}{0pt}%
\pgfsys@defobject{currentmarker}{\pgfqpoint{0.000000in}{-0.020833in}}{\pgfqpoint{0.000000in}{0.020833in}}{%
\pgfpathmoveto{\pgfqpoint{0.000000in}{-0.020833in}}%
\pgfpathlineto{\pgfqpoint{0.000000in}{0.020833in}}%
\pgfusepath{stroke,fill}%
}%
\begin{pgfscope}%
\pgfsys@transformshift{2.335902in}{3.034623in}%
\pgfsys@useobject{currentmarker}{}%
\end{pgfscope}%
\begin{pgfscope}%
\pgfsys@transformshift{2.336026in}{3.034623in}%
\pgfsys@useobject{currentmarker}{}%
\end{pgfscope}%
\begin{pgfscope}%
\pgfsys@transformshift{2.336480in}{3.034623in}%
\pgfsys@useobject{currentmarker}{}%
\end{pgfscope}%
\begin{pgfscope}%
\pgfsys@transformshift{2.336292in}{3.034623in}%
\pgfsys@useobject{currentmarker}{}%
\end{pgfscope}%
\end{pgfscope}%
\begin{pgfscope}%
\pgfpathrectangle{\pgfqpoint{1.914287in}{0.720000in}}{\pgfqpoint{0.550375in}{2.664000in}}%
\pgfusepath{clip}%
\pgfsetrectcap%
\pgfsetroundjoin%
\pgfsetlinewidth{0.702625pt}%
\definecolor{currentstroke}{rgb}{0.101961,0.101961,0.101961}%
\pgfsetstrokecolor{currentstroke}%
\pgfsetdash{}{0pt}%
\pgfpathmoveto{\pgfqpoint{2.339365in}{2.816262in}}%
\pgfpathlineto{\pgfqpoint{2.341057in}{2.816262in}}%
\pgfusepath{stroke}%
\end{pgfscope}%
\begin{pgfscope}%
\pgfpathrectangle{\pgfqpoint{1.914287in}{0.720000in}}{\pgfqpoint{0.550375in}{2.664000in}}%
\pgfusepath{clip}%
\pgfsetbuttcap%
\pgfsetroundjoin%
\definecolor{currentfill}{rgb}{0.101961,0.101961,0.101961}%
\pgfsetfillcolor{currentfill}%
\pgfsetlinewidth{0.501875pt}%
\definecolor{currentstroke}{rgb}{0.101961,0.101961,0.101961}%
\pgfsetstrokecolor{currentstroke}%
\pgfsetdash{}{0pt}%
\pgfsys@defobject{currentmarker}{\pgfqpoint{0.000000in}{-0.020833in}}{\pgfqpoint{0.000000in}{0.020833in}}{%
\pgfpathmoveto{\pgfqpoint{0.000000in}{-0.020833in}}%
\pgfpathlineto{\pgfqpoint{0.000000in}{0.020833in}}%
\pgfusepath{stroke,fill}%
}%
\begin{pgfscope}%
\pgfsys@transformshift{2.340929in}{2.816262in}%
\pgfsys@useobject{currentmarker}{}%
\end{pgfscope}%
\begin{pgfscope}%
\pgfsys@transformshift{2.340860in}{2.816262in}%
\pgfsys@useobject{currentmarker}{}%
\end{pgfscope}%
\begin{pgfscope}%
\pgfsys@transformshift{2.341057in}{2.816262in}%
\pgfsys@useobject{currentmarker}{}%
\end{pgfscope}%
\begin{pgfscope}%
\pgfsys@transformshift{2.339365in}{2.816262in}%
\pgfsys@useobject{currentmarker}{}%
\end{pgfscope}%
\end{pgfscope}%
\begin{pgfscope}%
\pgfpathrectangle{\pgfqpoint{1.914287in}{0.720000in}}{\pgfqpoint{0.550375in}{2.664000in}}%
\pgfusepath{clip}%
\pgfsetrectcap%
\pgfsetroundjoin%
\pgfsetlinewidth{0.702625pt}%
\definecolor{currentstroke}{rgb}{0.101961,0.101961,0.101961}%
\pgfsetstrokecolor{currentstroke}%
\pgfsetdash{}{0pt}%
\pgfpathmoveto{\pgfqpoint{2.381927in}{2.597902in}}%
\pgfpathlineto{\pgfqpoint{2.381927in}{2.597902in}}%
\pgfusepath{stroke}%
\end{pgfscope}%
\begin{pgfscope}%
\pgfpathrectangle{\pgfqpoint{1.914287in}{0.720000in}}{\pgfqpoint{0.550375in}{2.664000in}}%
\pgfusepath{clip}%
\pgfsetbuttcap%
\pgfsetroundjoin%
\definecolor{currentfill}{rgb}{0.101961,0.101961,0.101961}%
\pgfsetfillcolor{currentfill}%
\pgfsetlinewidth{0.501875pt}%
\definecolor{currentstroke}{rgb}{0.101961,0.101961,0.101961}%
\pgfsetstrokecolor{currentstroke}%
\pgfsetdash{}{0pt}%
\pgfsys@defobject{currentmarker}{\pgfqpoint{0.000000in}{-0.020833in}}{\pgfqpoint{0.000000in}{0.020833in}}{%
\pgfpathmoveto{\pgfqpoint{0.000000in}{-0.020833in}}%
\pgfpathlineto{\pgfqpoint{0.000000in}{0.020833in}}%
\pgfusepath{stroke,fill}%
}%
\begin{pgfscope}%
\pgfsys@transformshift{2.381927in}{2.597902in}%
\pgfsys@useobject{currentmarker}{}%
\end{pgfscope}%
\end{pgfscope}%
\begin{pgfscope}%
\pgfpathrectangle{\pgfqpoint{1.914287in}{0.720000in}}{\pgfqpoint{0.550375in}{2.664000in}}%
\pgfusepath{clip}%
\pgfsetrectcap%
\pgfsetroundjoin%
\pgfsetlinewidth{0.702625pt}%
\definecolor{currentstroke}{rgb}{0.101961,0.101961,0.101961}%
\pgfsetstrokecolor{currentstroke}%
\pgfsetdash{}{0pt}%
\pgfpathmoveto{\pgfqpoint{2.356936in}{2.379541in}}%
\pgfpathlineto{\pgfqpoint{2.366310in}{2.379541in}}%
\pgfusepath{stroke}%
\end{pgfscope}%
\begin{pgfscope}%
\pgfpathrectangle{\pgfqpoint{1.914287in}{0.720000in}}{\pgfqpoint{0.550375in}{2.664000in}}%
\pgfusepath{clip}%
\pgfsetbuttcap%
\pgfsetroundjoin%
\definecolor{currentfill}{rgb}{0.101961,0.101961,0.101961}%
\pgfsetfillcolor{currentfill}%
\pgfsetlinewidth{0.501875pt}%
\definecolor{currentstroke}{rgb}{0.101961,0.101961,0.101961}%
\pgfsetstrokecolor{currentstroke}%
\pgfsetdash{}{0pt}%
\pgfsys@defobject{currentmarker}{\pgfqpoint{0.000000in}{-0.020833in}}{\pgfqpoint{0.000000in}{0.020833in}}{%
\pgfpathmoveto{\pgfqpoint{0.000000in}{-0.020833in}}%
\pgfpathlineto{\pgfqpoint{0.000000in}{0.020833in}}%
\pgfusepath{stroke,fill}%
}%
\begin{pgfscope}%
\pgfsys@transformshift{2.357349in}{2.379541in}%
\pgfsys@useobject{currentmarker}{}%
\end{pgfscope}%
\begin{pgfscope}%
\pgfsys@transformshift{2.366310in}{2.379541in}%
\pgfsys@useobject{currentmarker}{}%
\end{pgfscope}%
\begin{pgfscope}%
\pgfsys@transformshift{2.362316in}{2.379541in}%
\pgfsys@useobject{currentmarker}{}%
\end{pgfscope}%
\begin{pgfscope}%
\pgfsys@transformshift{2.356936in}{2.379541in}%
\pgfsys@useobject{currentmarker}{}%
\end{pgfscope}%
\end{pgfscope}%
\begin{pgfscope}%
\pgfpathrectangle{\pgfqpoint{1.914287in}{0.720000in}}{\pgfqpoint{0.550375in}{2.664000in}}%
\pgfusepath{clip}%
\pgfsetrectcap%
\pgfsetroundjoin%
\pgfsetlinewidth{0.702625pt}%
\definecolor{currentstroke}{rgb}{0.101961,0.101961,0.101961}%
\pgfsetstrokecolor{currentstroke}%
\pgfsetdash{}{0pt}%
\pgfpathmoveto{\pgfqpoint{2.340136in}{2.161180in}}%
\pgfpathlineto{\pgfqpoint{2.361068in}{2.161180in}}%
\pgfusepath{stroke}%
\end{pgfscope}%
\begin{pgfscope}%
\pgfpathrectangle{\pgfqpoint{1.914287in}{0.720000in}}{\pgfqpoint{0.550375in}{2.664000in}}%
\pgfusepath{clip}%
\pgfsetbuttcap%
\pgfsetroundjoin%
\definecolor{currentfill}{rgb}{0.101961,0.101961,0.101961}%
\pgfsetfillcolor{currentfill}%
\pgfsetlinewidth{0.501875pt}%
\definecolor{currentstroke}{rgb}{0.101961,0.101961,0.101961}%
\pgfsetstrokecolor{currentstroke}%
\pgfsetdash{}{0pt}%
\pgfsys@defobject{currentmarker}{\pgfqpoint{0.000000in}{-0.020833in}}{\pgfqpoint{0.000000in}{0.020833in}}{%
\pgfpathmoveto{\pgfqpoint{0.000000in}{-0.020833in}}%
\pgfpathlineto{\pgfqpoint{0.000000in}{0.020833in}}%
\pgfusepath{stroke,fill}%
}%
\begin{pgfscope}%
\pgfsys@transformshift{2.348946in}{2.161180in}%
\pgfsys@useobject{currentmarker}{}%
\end{pgfscope}%
\begin{pgfscope}%
\pgfsys@transformshift{2.340136in}{2.161180in}%
\pgfsys@useobject{currentmarker}{}%
\end{pgfscope}%
\begin{pgfscope}%
\pgfsys@transformshift{2.348391in}{2.161180in}%
\pgfsys@useobject{currentmarker}{}%
\end{pgfscope}%
\begin{pgfscope}%
\pgfsys@transformshift{2.361068in}{2.161180in}%
\pgfsys@useobject{currentmarker}{}%
\end{pgfscope}%
\end{pgfscope}%
\begin{pgfscope}%
\pgfpathrectangle{\pgfqpoint{1.914287in}{0.720000in}}{\pgfqpoint{0.550375in}{2.664000in}}%
\pgfusepath{clip}%
\pgfsetrectcap%
\pgfsetroundjoin%
\pgfsetlinewidth{0.702625pt}%
\definecolor{currentstroke}{rgb}{0.101961,0.101961,0.101961}%
\pgfsetstrokecolor{currentstroke}%
\pgfsetdash{}{0pt}%
\pgfpathmoveto{\pgfqpoint{2.385615in}{1.942820in}}%
\pgfpathlineto{\pgfqpoint{2.417757in}{1.942820in}}%
\pgfusepath{stroke}%
\end{pgfscope}%
\begin{pgfscope}%
\pgfpathrectangle{\pgfqpoint{1.914287in}{0.720000in}}{\pgfqpoint{0.550375in}{2.664000in}}%
\pgfusepath{clip}%
\pgfsetbuttcap%
\pgfsetroundjoin%
\definecolor{currentfill}{rgb}{0.101961,0.101961,0.101961}%
\pgfsetfillcolor{currentfill}%
\pgfsetlinewidth{0.501875pt}%
\definecolor{currentstroke}{rgb}{0.101961,0.101961,0.101961}%
\pgfsetstrokecolor{currentstroke}%
\pgfsetdash{}{0pt}%
\pgfsys@defobject{currentmarker}{\pgfqpoint{0.000000in}{-0.020833in}}{\pgfqpoint{0.000000in}{0.020833in}}{%
\pgfpathmoveto{\pgfqpoint{0.000000in}{-0.020833in}}%
\pgfpathlineto{\pgfqpoint{0.000000in}{0.020833in}}%
\pgfusepath{stroke,fill}%
}%
\begin{pgfscope}%
\pgfsys@transformshift{2.417757in}{1.942820in}%
\pgfsys@useobject{currentmarker}{}%
\end{pgfscope}%
\begin{pgfscope}%
\pgfsys@transformshift{2.416005in}{1.942820in}%
\pgfsys@useobject{currentmarker}{}%
\end{pgfscope}%
\begin{pgfscope}%
\pgfsys@transformshift{2.385615in}{1.942820in}%
\pgfsys@useobject{currentmarker}{}%
\end{pgfscope}%
\end{pgfscope}%
\begin{pgfscope}%
\pgfpathrectangle{\pgfqpoint{1.914287in}{0.720000in}}{\pgfqpoint{0.550375in}{2.664000in}}%
\pgfusepath{clip}%
\pgfsetrectcap%
\pgfsetroundjoin%
\pgfsetlinewidth{0.702625pt}%
\definecolor{currentstroke}{rgb}{0.101961,0.101961,0.101961}%
\pgfsetstrokecolor{currentstroke}%
\pgfsetdash{}{0pt}%
\pgfpathmoveto{\pgfqpoint{2.383239in}{1.724459in}}%
\pgfpathlineto{\pgfqpoint{2.414175in}{1.724459in}}%
\pgfusepath{stroke}%
\end{pgfscope}%
\begin{pgfscope}%
\pgfpathrectangle{\pgfqpoint{1.914287in}{0.720000in}}{\pgfqpoint{0.550375in}{2.664000in}}%
\pgfusepath{clip}%
\pgfsetbuttcap%
\pgfsetroundjoin%
\definecolor{currentfill}{rgb}{0.101961,0.101961,0.101961}%
\pgfsetfillcolor{currentfill}%
\pgfsetlinewidth{0.501875pt}%
\definecolor{currentstroke}{rgb}{0.101961,0.101961,0.101961}%
\pgfsetstrokecolor{currentstroke}%
\pgfsetdash{}{0pt}%
\pgfsys@defobject{currentmarker}{\pgfqpoint{0.000000in}{-0.020833in}}{\pgfqpoint{0.000000in}{0.020833in}}{%
\pgfpathmoveto{\pgfqpoint{0.000000in}{-0.020833in}}%
\pgfpathlineto{\pgfqpoint{0.000000in}{0.020833in}}%
\pgfusepath{stroke,fill}%
}%
\begin{pgfscope}%
\pgfsys@transformshift{2.414147in}{1.724459in}%
\pgfsys@useobject{currentmarker}{}%
\end{pgfscope}%
\begin{pgfscope}%
\pgfsys@transformshift{2.401947in}{1.724459in}%
\pgfsys@useobject{currentmarker}{}%
\end{pgfscope}%
\begin{pgfscope}%
\pgfsys@transformshift{2.414175in}{1.724459in}%
\pgfsys@useobject{currentmarker}{}%
\end{pgfscope}%
\begin{pgfscope}%
\pgfsys@transformshift{2.383239in}{1.724459in}%
\pgfsys@useobject{currentmarker}{}%
\end{pgfscope}%
\end{pgfscope}%
\begin{pgfscope}%
\pgfpathrectangle{\pgfqpoint{1.914287in}{0.720000in}}{\pgfqpoint{0.550375in}{2.664000in}}%
\pgfusepath{clip}%
\pgfsetrectcap%
\pgfsetroundjoin%
\pgfsetlinewidth{0.702625pt}%
\definecolor{currentstroke}{rgb}{0.101961,0.101961,0.101961}%
\pgfsetstrokecolor{currentstroke}%
\pgfsetdash{}{0pt}%
\pgfpathmoveto{\pgfqpoint{2.464649in}{1.506098in}}%
\pgfpathlineto{\pgfqpoint{2.464658in}{1.506098in}}%
\pgfusepath{stroke}%
\end{pgfscope}%
\begin{pgfscope}%
\pgfpathrectangle{\pgfqpoint{1.914287in}{0.720000in}}{\pgfqpoint{0.550375in}{2.664000in}}%
\pgfusepath{clip}%
\pgfsetbuttcap%
\pgfsetroundjoin%
\definecolor{currentfill}{rgb}{0.101961,0.101961,0.101961}%
\pgfsetfillcolor{currentfill}%
\pgfsetlinewidth{0.501875pt}%
\definecolor{currentstroke}{rgb}{0.101961,0.101961,0.101961}%
\pgfsetstrokecolor{currentstroke}%
\pgfsetdash{}{0pt}%
\pgfsys@defobject{currentmarker}{\pgfqpoint{0.000000in}{-0.020833in}}{\pgfqpoint{0.000000in}{0.020833in}}{%
\pgfpathmoveto{\pgfqpoint{0.000000in}{-0.020833in}}%
\pgfpathlineto{\pgfqpoint{0.000000in}{0.020833in}}%
\pgfusepath{stroke,fill}%
}%
\begin{pgfscope}%
\pgfsys@transformshift{2.464653in}{1.506098in}%
\pgfsys@useobject{currentmarker}{}%
\end{pgfscope}%
\begin{pgfscope}%
\pgfsys@transformshift{2.464649in}{1.506098in}%
\pgfsys@useobject{currentmarker}{}%
\end{pgfscope}%
\begin{pgfscope}%
\pgfsys@transformshift{2.464658in}{1.506098in}%
\pgfsys@useobject{currentmarker}{}%
\end{pgfscope}%
\begin{pgfscope}%
\pgfsys@transformshift{2.464649in}{1.506098in}%
\pgfsys@useobject{currentmarker}{}%
\end{pgfscope}%
\end{pgfscope}%
\begin{pgfscope}%
\pgfpathrectangle{\pgfqpoint{1.914287in}{0.720000in}}{\pgfqpoint{0.550375in}{2.664000in}}%
\pgfusepath{clip}%
\pgfsetrectcap%
\pgfsetroundjoin%
\pgfsetlinewidth{0.702625pt}%
\definecolor{currentstroke}{rgb}{0.101961,0.101961,0.101961}%
\pgfsetstrokecolor{currentstroke}%
\pgfsetdash{}{0pt}%
\pgfpathmoveto{\pgfqpoint{2.464630in}{1.287738in}}%
\pgfpathlineto{\pgfqpoint{2.464653in}{1.287738in}}%
\pgfusepath{stroke}%
\end{pgfscope}%
\begin{pgfscope}%
\pgfpathrectangle{\pgfqpoint{1.914287in}{0.720000in}}{\pgfqpoint{0.550375in}{2.664000in}}%
\pgfusepath{clip}%
\pgfsetbuttcap%
\pgfsetroundjoin%
\definecolor{currentfill}{rgb}{0.101961,0.101961,0.101961}%
\pgfsetfillcolor{currentfill}%
\pgfsetlinewidth{0.501875pt}%
\definecolor{currentstroke}{rgb}{0.101961,0.101961,0.101961}%
\pgfsetstrokecolor{currentstroke}%
\pgfsetdash{}{0pt}%
\pgfsys@defobject{currentmarker}{\pgfqpoint{0.000000in}{-0.020833in}}{\pgfqpoint{0.000000in}{0.020833in}}{%
\pgfpathmoveto{\pgfqpoint{0.000000in}{-0.020833in}}%
\pgfpathlineto{\pgfqpoint{0.000000in}{0.020833in}}%
\pgfusepath{stroke,fill}%
}%
\begin{pgfscope}%
\pgfsys@transformshift{2.464630in}{1.287738in}%
\pgfsys@useobject{currentmarker}{}%
\end{pgfscope}%
\begin{pgfscope}%
\pgfsys@transformshift{2.464653in}{1.287738in}%
\pgfsys@useobject{currentmarker}{}%
\end{pgfscope}%
\begin{pgfscope}%
\pgfsys@transformshift{2.464649in}{1.287738in}%
\pgfsys@useobject{currentmarker}{}%
\end{pgfscope}%
\begin{pgfscope}%
\pgfsys@transformshift{2.464653in}{1.287738in}%
\pgfsys@useobject{currentmarker}{}%
\end{pgfscope}%
\end{pgfscope}%
\begin{pgfscope}%
\pgfpathrectangle{\pgfqpoint{1.914287in}{0.720000in}}{\pgfqpoint{0.550375in}{2.664000in}}%
\pgfusepath{clip}%
\pgfsetrectcap%
\pgfsetroundjoin%
\pgfsetlinewidth{0.702625pt}%
\definecolor{currentstroke}{rgb}{0.101961,0.101961,0.101961}%
\pgfsetstrokecolor{currentstroke}%
\pgfsetdash{}{0pt}%
\pgfpathmoveto{\pgfqpoint{2.464630in}{1.069377in}}%
\pgfpathlineto{\pgfqpoint{2.464658in}{1.069377in}}%
\pgfusepath{stroke}%
\end{pgfscope}%
\begin{pgfscope}%
\pgfpathrectangle{\pgfqpoint{1.914287in}{0.720000in}}{\pgfqpoint{0.550375in}{2.664000in}}%
\pgfusepath{clip}%
\pgfsetbuttcap%
\pgfsetroundjoin%
\definecolor{currentfill}{rgb}{0.101961,0.101961,0.101961}%
\pgfsetfillcolor{currentfill}%
\pgfsetlinewidth{0.501875pt}%
\definecolor{currentstroke}{rgb}{0.101961,0.101961,0.101961}%
\pgfsetstrokecolor{currentstroke}%
\pgfsetdash{}{0pt}%
\pgfsys@defobject{currentmarker}{\pgfqpoint{0.000000in}{-0.020833in}}{\pgfqpoint{0.000000in}{0.020833in}}{%
\pgfpathmoveto{\pgfqpoint{0.000000in}{-0.020833in}}%
\pgfpathlineto{\pgfqpoint{0.000000in}{0.020833in}}%
\pgfusepath{stroke,fill}%
}%
\begin{pgfscope}%
\pgfsys@transformshift{2.464630in}{1.069377in}%
\pgfsys@useobject{currentmarker}{}%
\end{pgfscope}%
\begin{pgfscope}%
\pgfsys@transformshift{2.464640in}{1.069377in}%
\pgfsys@useobject{currentmarker}{}%
\end{pgfscope}%
\begin{pgfscope}%
\pgfsys@transformshift{2.464658in}{1.069377in}%
\pgfsys@useobject{currentmarker}{}%
\end{pgfscope}%
\begin{pgfscope}%
\pgfsys@transformshift{2.464653in}{1.069377in}%
\pgfsys@useobject{currentmarker}{}%
\end{pgfscope}%
\end{pgfscope}%
\begin{pgfscope}%
\pgfpathrectangle{\pgfqpoint{1.914287in}{0.720000in}}{\pgfqpoint{0.550375in}{2.664000in}}%
\pgfusepath{clip}%
\pgfsetrectcap%
\pgfsetroundjoin%
\pgfsetlinewidth{0.702625pt}%
\definecolor{currentstroke}{rgb}{0.101961,0.101961,0.101961}%
\pgfsetstrokecolor{currentstroke}%
\pgfsetdash{}{0pt}%
\pgfpathmoveto{\pgfqpoint{2.464658in}{0.851016in}}%
\pgfpathlineto{\pgfqpoint{2.464663in}{0.851016in}}%
\pgfusepath{stroke}%
\end{pgfscope}%
\begin{pgfscope}%
\pgfpathrectangle{\pgfqpoint{1.914287in}{0.720000in}}{\pgfqpoint{0.550375in}{2.664000in}}%
\pgfusepath{clip}%
\pgfsetbuttcap%
\pgfsetroundjoin%
\definecolor{currentfill}{rgb}{0.101961,0.101961,0.101961}%
\pgfsetfillcolor{currentfill}%
\pgfsetlinewidth{0.501875pt}%
\definecolor{currentstroke}{rgb}{0.101961,0.101961,0.101961}%
\pgfsetstrokecolor{currentstroke}%
\pgfsetdash{}{0pt}%
\pgfsys@defobject{currentmarker}{\pgfqpoint{0.000000in}{-0.020833in}}{\pgfqpoint{0.000000in}{0.020833in}}{%
\pgfpathmoveto{\pgfqpoint{0.000000in}{-0.020833in}}%
\pgfpathlineto{\pgfqpoint{0.000000in}{0.020833in}}%
\pgfusepath{stroke,fill}%
}%
\begin{pgfscope}%
\pgfsys@transformshift{2.464663in}{0.851016in}%
\pgfsys@useobject{currentmarker}{}%
\end{pgfscope}%
\begin{pgfscope}%
\pgfsys@transformshift{2.464658in}{0.851016in}%
\pgfsys@useobject{currentmarker}{}%
\end{pgfscope}%
\begin{pgfscope}%
\pgfsys@transformshift{2.464658in}{0.851016in}%
\pgfsys@useobject{currentmarker}{}%
\end{pgfscope}%
\begin{pgfscope}%
\pgfsys@transformshift{2.464658in}{0.851016in}%
\pgfsys@useobject{currentmarker}{}%
\end{pgfscope}%
\end{pgfscope}%
\begin{pgfscope}%
\pgfsetrectcap%
\pgfsetmiterjoin%
\pgfsetlinewidth{0.803000pt}%
\definecolor{currentstroke}{rgb}{0.101961,0.101961,0.101961}%
\pgfsetstrokecolor{currentstroke}%
\pgfsetdash{}{0pt}%
\pgfpathmoveto{\pgfqpoint{1.914287in}{0.720000in}}%
\pgfpathlineto{\pgfqpoint{1.914287in}{3.384000in}}%
\pgfusepath{stroke}%
\end{pgfscope}%
\begin{pgfscope}%
\pgfsetrectcap%
\pgfsetmiterjoin%
\pgfsetlinewidth{0.803000pt}%
\definecolor{currentstroke}{rgb}{0.101961,0.101961,0.101961}%
\pgfsetstrokecolor{currentstroke}%
\pgfsetdash{}{0pt}%
\pgfpathmoveto{\pgfqpoint{1.914287in}{0.720000in}}%
\pgfpathlineto{\pgfqpoint{2.464663in}{0.720000in}}%
\pgfusepath{stroke}%
\end{pgfscope}%
\begin{pgfscope}%
\pgfsetbuttcap%
\pgfsetmiterjoin%
\definecolor{currentfill}{rgb}{1.000000,1.000000,1.000000}%
\pgfsetfillcolor{currentfill}%
\pgfsetlinewidth{0.000000pt}%
\definecolor{currentstroke}{rgb}{0.000000,0.000000,0.000000}%
\pgfsetstrokecolor{currentstroke}%
\pgfsetstrokeopacity{0.000000}%
\pgfsetdash{}{0pt}%
\pgfpathmoveto{\pgfqpoint{2.629775in}{0.720000in}}%
\pgfpathlineto{\pgfqpoint{3.180150in}{0.720000in}}%
\pgfpathlineto{\pgfqpoint{3.180150in}{3.384000in}}%
\pgfpathlineto{\pgfqpoint{2.629775in}{3.384000in}}%
\pgfpathlineto{\pgfqpoint{2.629775in}{0.720000in}}%
\pgfpathclose%
\pgfusepath{fill}%
\end{pgfscope}%
\begin{pgfscope}%
\pgfpathrectangle{\pgfqpoint{2.629775in}{0.720000in}}{\pgfqpoint{0.550375in}{2.664000in}}%
\pgfusepath{clip}%
\pgfsetrectcap%
\pgfsetroundjoin%
\pgfsetlinewidth{0.401500pt}%
\definecolor{currentstroke}{rgb}{0.901961,0.901961,0.901961}%
\pgfsetstrokecolor{currentstroke}%
\pgfsetdash{}{0pt}%
\pgfpathmoveto{\pgfqpoint{2.629775in}{0.720000in}}%
\pgfpathlineto{\pgfqpoint{2.629775in}{3.384000in}}%
\pgfusepath{stroke}%
\end{pgfscope}%
\begin{pgfscope}%
\definecolor{textcolor}{rgb}{0.101961,0.101961,0.101961}%
\pgfsetstrokecolor{textcolor}%
\pgfsetfillcolor{textcolor}%
\pgftext[x=2.629775in,y=0.671389in,,top]{\color{textcolor}\sffamily\fontsize{7.000000}{8.400000}\selectfont 0}%
\end{pgfscope}%
\begin{pgfscope}%
\pgfpathrectangle{\pgfqpoint{2.629775in}{0.720000in}}{\pgfqpoint{0.550375in}{2.664000in}}%
\pgfusepath{clip}%
\pgfsetrectcap%
\pgfsetroundjoin%
\pgfsetlinewidth{0.401500pt}%
\definecolor{currentstroke}{rgb}{0.901961,0.901961,0.901961}%
\pgfsetstrokecolor{currentstroke}%
\pgfsetdash{}{0pt}%
\pgfpathmoveto{\pgfqpoint{2.864808in}{0.720000in}}%
\pgfpathlineto{\pgfqpoint{2.864808in}{3.384000in}}%
\pgfusepath{stroke}%
\end{pgfscope}%
\begin{pgfscope}%
\definecolor{textcolor}{rgb}{0.101961,0.101961,0.101961}%
\pgfsetstrokecolor{textcolor}%
\pgfsetfillcolor{textcolor}%
\pgftext[x=2.864808in,y=0.671389in,,top]{\color{textcolor}\sffamily\fontsize{7.000000}{8.400000}\selectfont 3}%
\end{pgfscope}%
\begin{pgfscope}%
\pgfpathrectangle{\pgfqpoint{2.629775in}{0.720000in}}{\pgfqpoint{0.550375in}{2.664000in}}%
\pgfusepath{clip}%
\pgfsetrectcap%
\pgfsetroundjoin%
\pgfsetlinewidth{0.401500pt}%
\definecolor{currentstroke}{rgb}{0.901961,0.901961,0.901961}%
\pgfsetstrokecolor{currentstroke}%
\pgfsetdash{}{0pt}%
\pgfpathmoveto{\pgfqpoint{3.099842in}{0.720000in}}%
\pgfpathlineto{\pgfqpoint{3.099842in}{3.384000in}}%
\pgfusepath{stroke}%
\end{pgfscope}%
\begin{pgfscope}%
\definecolor{textcolor}{rgb}{0.101961,0.101961,0.101961}%
\pgfsetstrokecolor{textcolor}%
\pgfsetfillcolor{textcolor}%
\pgftext[x=3.099842in,y=0.671389in,,top]{\color{textcolor}\sffamily\fontsize{7.000000}{8.400000}\selectfont 6}%
\end{pgfscope}%
\begin{pgfscope}%
\definecolor{textcolor}{rgb}{0.000000,0.000000,0.000000}%
\pgfsetstrokecolor{textcolor}%
\pgfsetfillcolor{textcolor}%
\pgftext[x=2.603334in, y=0.461899in, left, base]{\color{textcolor}\sffamily\fontsize{7.500000}{9.000000}\selectfont uncached tok.}%
\end{pgfscope}%
\begin{pgfscope}%
\definecolor{textcolor}{rgb}{0.000000,0.000000,0.000000}%
\pgfsetstrokecolor{textcolor}%
\pgfsetfillcolor{textcolor}%
\pgftext[x=2.596531in, y=0.355494in, left, base]{\color{textcolor}\sffamily\fontsize{7.500000}{9.000000}\selectfont per GPU-s (k)}%
\end{pgfscope}%
\begin{pgfscope}%
\pgfpathrectangle{\pgfqpoint{2.629775in}{0.720000in}}{\pgfqpoint{0.550375in}{2.664000in}}%
\pgfusepath{clip}%
\pgfsetbuttcap%
\pgfsetmiterjoin%
\definecolor{currentfill}{rgb}{0.156863,0.356863,0.568627}%
\pgfsetfillcolor{currentfill}%
\pgfsetlinewidth{0.000000pt}%
\definecolor{currentstroke}{rgb}{0.000000,0.000000,0.000000}%
\pgfsetstrokecolor{currentstroke}%
\pgfsetstrokeopacity{0.000000}%
\pgfsetdash{}{0pt}%
\pgfpathmoveto{\pgfqpoint{2.629775in}{3.176557in}}%
\pgfpathlineto{\pgfqpoint{3.044709in}{3.176557in}}%
\pgfpathlineto{\pgfqpoint{3.044709in}{3.329410in}}%
\pgfpathlineto{\pgfqpoint{2.629775in}{3.329410in}}%
\pgfpathlineto{\pgfqpoint{2.629775in}{3.176557in}}%
\pgfpathclose%
\pgfusepath{fill}%
\end{pgfscope}%
\begin{pgfscope}%
\pgfpathrectangle{\pgfqpoint{2.629775in}{0.720000in}}{\pgfqpoint{0.550375in}{2.664000in}}%
\pgfusepath{clip}%
\pgfsetbuttcap%
\pgfsetmiterjoin%
\definecolor{currentfill}{rgb}{0.784314,0.784314,0.784314}%
\pgfsetfillcolor{currentfill}%
\pgfsetlinewidth{0.000000pt}%
\definecolor{currentstroke}{rgb}{0.000000,0.000000,0.000000}%
\pgfsetstrokecolor{currentstroke}%
\pgfsetstrokeopacity{0.000000}%
\pgfsetdash{}{0pt}%
\pgfpathmoveto{\pgfqpoint{2.629775in}{2.958197in}}%
\pgfpathlineto{\pgfqpoint{3.077864in}{2.958197in}}%
\pgfpathlineto{\pgfqpoint{3.077864in}{3.111049in}}%
\pgfpathlineto{\pgfqpoint{2.629775in}{3.111049in}}%
\pgfpathlineto{\pgfqpoint{2.629775in}{2.958197in}}%
\pgfpathclose%
\pgfusepath{fill}%
\end{pgfscope}%
\begin{pgfscope}%
\pgfpathrectangle{\pgfqpoint{2.629775in}{0.720000in}}{\pgfqpoint{0.550375in}{2.664000in}}%
\pgfusepath{clip}%
\pgfsetbuttcap%
\pgfsetmiterjoin%
\definecolor{currentfill}{rgb}{0.784314,0.784314,0.784314}%
\pgfsetfillcolor{currentfill}%
\pgfsetlinewidth{0.000000pt}%
\definecolor{currentstroke}{rgb}{0.000000,0.000000,0.000000}%
\pgfsetstrokecolor{currentstroke}%
\pgfsetstrokeopacity{0.000000}%
\pgfsetdash{}{0pt}%
\pgfpathmoveto{\pgfqpoint{2.629775in}{2.739836in}}%
\pgfpathlineto{\pgfqpoint{3.077362in}{2.739836in}}%
\pgfpathlineto{\pgfqpoint{3.077362in}{2.892689in}}%
\pgfpathlineto{\pgfqpoint{2.629775in}{2.892689in}}%
\pgfpathlineto{\pgfqpoint{2.629775in}{2.739836in}}%
\pgfpathclose%
\pgfusepath{fill}%
\end{pgfscope}%
\begin{pgfscope}%
\pgfpathrectangle{\pgfqpoint{2.629775in}{0.720000in}}{\pgfqpoint{0.550375in}{2.664000in}}%
\pgfusepath{clip}%
\pgfsetbuttcap%
\pgfsetmiterjoin%
\definecolor{currentfill}{rgb}{0.478431,0.478431,0.478431}%
\pgfsetfillcolor{currentfill}%
\pgfsetlinewidth{0.000000pt}%
\definecolor{currentstroke}{rgb}{0.000000,0.000000,0.000000}%
\pgfsetstrokecolor{currentstroke}%
\pgfsetstrokeopacity{0.000000}%
\pgfsetdash{}{0pt}%
\pgfpathmoveto{\pgfqpoint{2.629775in}{2.521475in}}%
\pgfpathlineto{\pgfqpoint{3.046877in}{2.521475in}}%
\pgfpathlineto{\pgfqpoint{3.046877in}{2.674328in}}%
\pgfpathlineto{\pgfqpoint{2.629775in}{2.674328in}}%
\pgfpathlineto{\pgfqpoint{2.629775in}{2.521475in}}%
\pgfpathclose%
\pgfusepath{fill}%
\end{pgfscope}%
\begin{pgfscope}%
\pgfpathrectangle{\pgfqpoint{2.629775in}{0.720000in}}{\pgfqpoint{0.550375in}{2.664000in}}%
\pgfusepath{clip}%
\pgfsetbuttcap%
\pgfsetmiterjoin%
\definecolor{currentfill}{rgb}{0.784314,0.784314,0.784314}%
\pgfsetfillcolor{currentfill}%
\pgfsetlinewidth{0.000000pt}%
\definecolor{currentstroke}{rgb}{0.000000,0.000000,0.000000}%
\pgfsetstrokecolor{currentstroke}%
\pgfsetstrokeopacity{0.000000}%
\pgfsetdash{}{0pt}%
\pgfpathmoveto{\pgfqpoint{2.629775in}{2.303115in}}%
\pgfpathlineto{\pgfqpoint{3.067240in}{2.303115in}}%
\pgfpathlineto{\pgfqpoint{3.067240in}{2.455967in}}%
\pgfpathlineto{\pgfqpoint{2.629775in}{2.455967in}}%
\pgfpathlineto{\pgfqpoint{2.629775in}{2.303115in}}%
\pgfpathclose%
\pgfusepath{fill}%
\end{pgfscope}%
\begin{pgfscope}%
\pgfpathrectangle{\pgfqpoint{2.629775in}{0.720000in}}{\pgfqpoint{0.550375in}{2.664000in}}%
\pgfusepath{clip}%
\pgfsetbuttcap%
\pgfsetmiterjoin%
\definecolor{currentfill}{rgb}{0.784314,0.784314,0.784314}%
\pgfsetfillcolor{currentfill}%
\pgfsetlinewidth{0.000000pt}%
\definecolor{currentstroke}{rgb}{0.000000,0.000000,0.000000}%
\pgfsetstrokecolor{currentstroke}%
\pgfsetstrokeopacity{0.000000}%
\pgfsetdash{}{0pt}%
\pgfpathmoveto{\pgfqpoint{2.629775in}{2.084754in}}%
\pgfpathlineto{\pgfqpoint{3.090012in}{2.084754in}}%
\pgfpathlineto{\pgfqpoint{3.090012in}{2.237607in}}%
\pgfpathlineto{\pgfqpoint{2.629775in}{2.237607in}}%
\pgfpathlineto{\pgfqpoint{2.629775in}{2.084754in}}%
\pgfpathclose%
\pgfusepath{fill}%
\end{pgfscope}%
\begin{pgfscope}%
\pgfpathrectangle{\pgfqpoint{2.629775in}{0.720000in}}{\pgfqpoint{0.550375in}{2.664000in}}%
\pgfusepath{clip}%
\pgfsetbuttcap%
\pgfsetmiterjoin%
\definecolor{currentfill}{rgb}{0.784314,0.784314,0.784314}%
\pgfsetfillcolor{currentfill}%
\pgfsetlinewidth{0.000000pt}%
\definecolor{currentstroke}{rgb}{0.000000,0.000000,0.000000}%
\pgfsetstrokecolor{currentstroke}%
\pgfsetstrokeopacity{0.000000}%
\pgfsetdash{}{0pt}%
\pgfpathmoveto{\pgfqpoint{2.629775in}{1.866393in}}%
\pgfpathlineto{\pgfqpoint{3.064493in}{1.866393in}}%
\pgfpathlineto{\pgfqpoint{3.064493in}{2.019246in}}%
\pgfpathlineto{\pgfqpoint{2.629775in}{2.019246in}}%
\pgfpathlineto{\pgfqpoint{2.629775in}{1.866393in}}%
\pgfpathclose%
\pgfusepath{fill}%
\end{pgfscope}%
\begin{pgfscope}%
\pgfpathrectangle{\pgfqpoint{2.629775in}{0.720000in}}{\pgfqpoint{0.550375in}{2.664000in}}%
\pgfusepath{clip}%
\pgfsetbuttcap%
\pgfsetmiterjoin%
\definecolor{currentfill}{rgb}{0.784314,0.784314,0.784314}%
\pgfsetfillcolor{currentfill}%
\pgfsetlinewidth{0.000000pt}%
\definecolor{currentstroke}{rgb}{0.000000,0.000000,0.000000}%
\pgfsetstrokecolor{currentstroke}%
\pgfsetstrokeopacity{0.000000}%
\pgfsetdash{}{0pt}%
\pgfpathmoveto{\pgfqpoint{2.629775in}{1.648033in}}%
\pgfpathlineto{\pgfqpoint{3.086526in}{1.648033in}}%
\pgfpathlineto{\pgfqpoint{3.086526in}{1.800885in}}%
\pgfpathlineto{\pgfqpoint{2.629775in}{1.800885in}}%
\pgfpathlineto{\pgfqpoint{2.629775in}{1.648033in}}%
\pgfpathclose%
\pgfusepath{fill}%
\end{pgfscope}%
\begin{pgfscope}%
\pgfpathrectangle{\pgfqpoint{2.629775in}{0.720000in}}{\pgfqpoint{0.550375in}{2.664000in}}%
\pgfusepath{clip}%
\pgfsetbuttcap%
\pgfsetmiterjoin%
\definecolor{currentfill}{rgb}{0.478431,0.478431,0.478431}%
\pgfsetfillcolor{currentfill}%
\pgfsetlinewidth{0.000000pt}%
\definecolor{currentstroke}{rgb}{0.000000,0.000000,0.000000}%
\pgfsetstrokecolor{currentstroke}%
\pgfsetstrokeopacity{0.000000}%
\pgfsetdash{}{0pt}%
\pgfpathmoveto{\pgfqpoint{2.629775in}{1.429672in}}%
\pgfpathlineto{\pgfqpoint{3.072431in}{1.429672in}}%
\pgfpathlineto{\pgfqpoint{3.072431in}{1.582525in}}%
\pgfpathlineto{\pgfqpoint{2.629775in}{1.582525in}}%
\pgfpathlineto{\pgfqpoint{2.629775in}{1.429672in}}%
\pgfpathclose%
\pgfusepath{fill}%
\end{pgfscope}%
\begin{pgfscope}%
\pgfpathrectangle{\pgfqpoint{2.629775in}{0.720000in}}{\pgfqpoint{0.550375in}{2.664000in}}%
\pgfusepath{clip}%
\pgfsetbuttcap%
\pgfsetmiterjoin%
\definecolor{currentfill}{rgb}{0.784314,0.784314,0.784314}%
\pgfsetfillcolor{currentfill}%
\pgfsetlinewidth{0.000000pt}%
\definecolor{currentstroke}{rgb}{0.000000,0.000000,0.000000}%
\pgfsetstrokecolor{currentstroke}%
\pgfsetstrokeopacity{0.000000}%
\pgfsetdash{}{0pt}%
\pgfpathmoveto{\pgfqpoint{2.629775in}{1.211311in}}%
\pgfpathlineto{\pgfqpoint{3.069275in}{1.211311in}}%
\pgfpathlineto{\pgfqpoint{3.069275in}{1.364164in}}%
\pgfpathlineto{\pgfqpoint{2.629775in}{1.364164in}}%
\pgfpathlineto{\pgfqpoint{2.629775in}{1.211311in}}%
\pgfpathclose%
\pgfusepath{fill}%
\end{pgfscope}%
\begin{pgfscope}%
\pgfpathrectangle{\pgfqpoint{2.629775in}{0.720000in}}{\pgfqpoint{0.550375in}{2.664000in}}%
\pgfusepath{clip}%
\pgfsetbuttcap%
\pgfsetmiterjoin%
\definecolor{currentfill}{rgb}{0.784314,0.784314,0.784314}%
\pgfsetfillcolor{currentfill}%
\pgfsetlinewidth{0.000000pt}%
\definecolor{currentstroke}{rgb}{0.000000,0.000000,0.000000}%
\pgfsetstrokecolor{currentstroke}%
\pgfsetstrokeopacity{0.000000}%
\pgfsetdash{}{0pt}%
\pgfpathmoveto{\pgfqpoint{2.629775in}{0.992951in}}%
\pgfpathlineto{\pgfqpoint{3.059286in}{0.992951in}}%
\pgfpathlineto{\pgfqpoint{3.059286in}{1.145803in}}%
\pgfpathlineto{\pgfqpoint{2.629775in}{1.145803in}}%
\pgfpathlineto{\pgfqpoint{2.629775in}{0.992951in}}%
\pgfpathclose%
\pgfusepath{fill}%
\end{pgfscope}%
\begin{pgfscope}%
\pgfpathrectangle{\pgfqpoint{2.629775in}{0.720000in}}{\pgfqpoint{0.550375in}{2.664000in}}%
\pgfusepath{clip}%
\pgfsetbuttcap%
\pgfsetmiterjoin%
\definecolor{currentfill}{rgb}{0.478431,0.478431,0.478431}%
\pgfsetfillcolor{currentfill}%
\pgfsetlinewidth{0.000000pt}%
\definecolor{currentstroke}{rgb}{0.000000,0.000000,0.000000}%
\pgfsetstrokecolor{currentstroke}%
\pgfsetstrokeopacity{0.000000}%
\pgfsetdash{}{0pt}%
\pgfpathmoveto{\pgfqpoint{2.629775in}{0.774590in}}%
\pgfpathlineto{\pgfqpoint{3.068840in}{0.774590in}}%
\pgfpathlineto{\pgfqpoint{3.068840in}{0.927443in}}%
\pgfpathlineto{\pgfqpoint{2.629775in}{0.927443in}}%
\pgfpathlineto{\pgfqpoint{2.629775in}{0.774590in}}%
\pgfpathclose%
\pgfusepath{fill}%
\end{pgfscope}%
\begin{pgfscope}%
\pgfpathrectangle{\pgfqpoint{2.629775in}{0.720000in}}{\pgfqpoint{0.550375in}{2.664000in}}%
\pgfusepath{clip}%
\pgfsetrectcap%
\pgfsetroundjoin%
\pgfsetlinewidth{0.702625pt}%
\definecolor{currentstroke}{rgb}{0.101961,0.101961,0.101961}%
\pgfsetstrokecolor{currentstroke}%
\pgfsetdash{}{0pt}%
\pgfpathmoveto{\pgfqpoint{3.040815in}{3.252984in}}%
\pgfpathlineto{\pgfqpoint{3.046989in}{3.252984in}}%
\pgfusepath{stroke}%
\end{pgfscope}%
\begin{pgfscope}%
\pgfpathrectangle{\pgfqpoint{2.629775in}{0.720000in}}{\pgfqpoint{0.550375in}{2.664000in}}%
\pgfusepath{clip}%
\pgfsetbuttcap%
\pgfsetroundjoin%
\definecolor{currentfill}{rgb}{0.101961,0.101961,0.101961}%
\pgfsetfillcolor{currentfill}%
\pgfsetlinewidth{0.501875pt}%
\definecolor{currentstroke}{rgb}{0.101961,0.101961,0.101961}%
\pgfsetstrokecolor{currentstroke}%
\pgfsetdash{}{0pt}%
\pgfsys@defobject{currentmarker}{\pgfqpoint{0.000000in}{-0.020833in}}{\pgfqpoint{0.000000in}{0.020833in}}{%
\pgfpathmoveto{\pgfqpoint{0.000000in}{-0.020833in}}%
\pgfpathlineto{\pgfqpoint{0.000000in}{0.020833in}}%
\pgfusepath{stroke,fill}%
}%
\begin{pgfscope}%
\pgfsys@transformshift{3.040815in}{3.252984in}%
\pgfsys@useobject{currentmarker}{}%
\end{pgfscope}%
\begin{pgfscope}%
\pgfsys@transformshift{3.046989in}{3.252984in}%
\pgfsys@useobject{currentmarker}{}%
\end{pgfscope}%
\begin{pgfscope}%
\pgfsys@transformshift{3.045812in}{3.252984in}%
\pgfsys@useobject{currentmarker}{}%
\end{pgfscope}%
\begin{pgfscope}%
\pgfsys@transformshift{3.045221in}{3.252984in}%
\pgfsys@useobject{currentmarker}{}%
\end{pgfscope}%
\end{pgfscope}%
\begin{pgfscope}%
\pgfpathrectangle{\pgfqpoint{2.629775in}{0.720000in}}{\pgfqpoint{0.550375in}{2.664000in}}%
\pgfusepath{clip}%
\pgfsetrectcap%
\pgfsetroundjoin%
\pgfsetlinewidth{0.702625pt}%
\definecolor{currentstroke}{rgb}{0.101961,0.101961,0.101961}%
\pgfsetstrokecolor{currentstroke}%
\pgfsetdash{}{0pt}%
\pgfpathmoveto{\pgfqpoint{3.077723in}{3.034623in}}%
\pgfpathlineto{\pgfqpoint{3.078062in}{3.034623in}}%
\pgfusepath{stroke}%
\end{pgfscope}%
\begin{pgfscope}%
\pgfpathrectangle{\pgfqpoint{2.629775in}{0.720000in}}{\pgfqpoint{0.550375in}{2.664000in}}%
\pgfusepath{clip}%
\pgfsetbuttcap%
\pgfsetroundjoin%
\definecolor{currentfill}{rgb}{0.101961,0.101961,0.101961}%
\pgfsetfillcolor{currentfill}%
\pgfsetlinewidth{0.501875pt}%
\definecolor{currentstroke}{rgb}{0.101961,0.101961,0.101961}%
\pgfsetstrokecolor{currentstroke}%
\pgfsetdash{}{0pt}%
\pgfsys@defobject{currentmarker}{\pgfqpoint{0.000000in}{-0.020833in}}{\pgfqpoint{0.000000in}{0.020833in}}{%
\pgfpathmoveto{\pgfqpoint{0.000000in}{-0.020833in}}%
\pgfpathlineto{\pgfqpoint{0.000000in}{0.020833in}}%
\pgfusepath{stroke,fill}%
}%
\begin{pgfscope}%
\pgfsys@transformshift{3.077743in}{3.034623in}%
\pgfsys@useobject{currentmarker}{}%
\end{pgfscope}%
\begin{pgfscope}%
\pgfsys@transformshift{3.078062in}{3.034623in}%
\pgfsys@useobject{currentmarker}{}%
\end{pgfscope}%
\begin{pgfscope}%
\pgfsys@transformshift{3.077930in}{3.034623in}%
\pgfsys@useobject{currentmarker}{}%
\end{pgfscope}%
\begin{pgfscope}%
\pgfsys@transformshift{3.077723in}{3.034623in}%
\pgfsys@useobject{currentmarker}{}%
\end{pgfscope}%
\end{pgfscope}%
\begin{pgfscope}%
\pgfpathrectangle{\pgfqpoint{2.629775in}{0.720000in}}{\pgfqpoint{0.550375in}{2.664000in}}%
\pgfusepath{clip}%
\pgfsetrectcap%
\pgfsetroundjoin%
\pgfsetlinewidth{0.702625pt}%
\definecolor{currentstroke}{rgb}{0.101961,0.101961,0.101961}%
\pgfsetstrokecolor{currentstroke}%
\pgfsetdash{}{0pt}%
\pgfpathmoveto{\pgfqpoint{3.077153in}{2.816262in}}%
\pgfpathlineto{\pgfqpoint{3.077575in}{2.816262in}}%
\pgfusepath{stroke}%
\end{pgfscope}%
\begin{pgfscope}%
\pgfpathrectangle{\pgfqpoint{2.629775in}{0.720000in}}{\pgfqpoint{0.550375in}{2.664000in}}%
\pgfusepath{clip}%
\pgfsetbuttcap%
\pgfsetroundjoin%
\definecolor{currentfill}{rgb}{0.101961,0.101961,0.101961}%
\pgfsetfillcolor{currentfill}%
\pgfsetlinewidth{0.501875pt}%
\definecolor{currentstroke}{rgb}{0.101961,0.101961,0.101961}%
\pgfsetstrokecolor{currentstroke}%
\pgfsetdash{}{0pt}%
\pgfsys@defobject{currentmarker}{\pgfqpoint{0.000000in}{-0.020833in}}{\pgfqpoint{0.000000in}{0.020833in}}{%
\pgfpathmoveto{\pgfqpoint{0.000000in}{-0.020833in}}%
\pgfpathlineto{\pgfqpoint{0.000000in}{0.020833in}}%
\pgfusepath{stroke,fill}%
}%
\begin{pgfscope}%
\pgfsys@transformshift{3.077153in}{2.816262in}%
\pgfsys@useobject{currentmarker}{}%
\end{pgfscope}%
\begin{pgfscope}%
\pgfsys@transformshift{3.077575in}{2.816262in}%
\pgfsys@useobject{currentmarker}{}%
\end{pgfscope}%
\begin{pgfscope}%
\pgfsys@transformshift{3.077372in}{2.816262in}%
\pgfsys@useobject{currentmarker}{}%
\end{pgfscope}%
\begin{pgfscope}%
\pgfsys@transformshift{3.077347in}{2.816262in}%
\pgfsys@useobject{currentmarker}{}%
\end{pgfscope}%
\end{pgfscope}%
\begin{pgfscope}%
\pgfpathrectangle{\pgfqpoint{2.629775in}{0.720000in}}{\pgfqpoint{0.550375in}{2.664000in}}%
\pgfusepath{clip}%
\pgfsetrectcap%
\pgfsetroundjoin%
\pgfsetlinewidth{0.702625pt}%
\definecolor{currentstroke}{rgb}{0.101961,0.101961,0.101961}%
\pgfsetstrokecolor{currentstroke}%
\pgfsetdash{}{0pt}%
\pgfpathmoveto{\pgfqpoint{3.046877in}{2.597902in}}%
\pgfpathlineto{\pgfqpoint{3.046877in}{2.597902in}}%
\pgfusepath{stroke}%
\end{pgfscope}%
\begin{pgfscope}%
\pgfpathrectangle{\pgfqpoint{2.629775in}{0.720000in}}{\pgfqpoint{0.550375in}{2.664000in}}%
\pgfusepath{clip}%
\pgfsetbuttcap%
\pgfsetroundjoin%
\definecolor{currentfill}{rgb}{0.101961,0.101961,0.101961}%
\pgfsetfillcolor{currentfill}%
\pgfsetlinewidth{0.501875pt}%
\definecolor{currentstroke}{rgb}{0.101961,0.101961,0.101961}%
\pgfsetstrokecolor{currentstroke}%
\pgfsetdash{}{0pt}%
\pgfsys@defobject{currentmarker}{\pgfqpoint{0.000000in}{-0.020833in}}{\pgfqpoint{0.000000in}{0.020833in}}{%
\pgfpathmoveto{\pgfqpoint{0.000000in}{-0.020833in}}%
\pgfpathlineto{\pgfqpoint{0.000000in}{0.020833in}}%
\pgfusepath{stroke,fill}%
}%
\begin{pgfscope}%
\pgfsys@transformshift{3.046877in}{2.597902in}%
\pgfsys@useobject{currentmarker}{}%
\end{pgfscope}%
\end{pgfscope}%
\begin{pgfscope}%
\pgfpathrectangle{\pgfqpoint{2.629775in}{0.720000in}}{\pgfqpoint{0.550375in}{2.664000in}}%
\pgfusepath{clip}%
\pgfsetrectcap%
\pgfsetroundjoin%
\pgfsetlinewidth{0.702625pt}%
\definecolor{currentstroke}{rgb}{0.101961,0.101961,0.101961}%
\pgfsetstrokecolor{currentstroke}%
\pgfsetdash{}{0pt}%
\pgfpathmoveto{\pgfqpoint{3.058517in}{2.379541in}}%
\pgfpathlineto{\pgfqpoint{3.072755in}{2.379541in}}%
\pgfusepath{stroke}%
\end{pgfscope}%
\begin{pgfscope}%
\pgfpathrectangle{\pgfqpoint{2.629775in}{0.720000in}}{\pgfqpoint{0.550375in}{2.664000in}}%
\pgfusepath{clip}%
\pgfsetbuttcap%
\pgfsetroundjoin%
\definecolor{currentfill}{rgb}{0.101961,0.101961,0.101961}%
\pgfsetfillcolor{currentfill}%
\pgfsetlinewidth{0.501875pt}%
\definecolor{currentstroke}{rgb}{0.101961,0.101961,0.101961}%
\pgfsetstrokecolor{currentstroke}%
\pgfsetdash{}{0pt}%
\pgfsys@defobject{currentmarker}{\pgfqpoint{0.000000in}{-0.020833in}}{\pgfqpoint{0.000000in}{0.020833in}}{%
\pgfpathmoveto{\pgfqpoint{0.000000in}{-0.020833in}}%
\pgfpathlineto{\pgfqpoint{0.000000in}{0.020833in}}%
\pgfusepath{stroke,fill}%
}%
\begin{pgfscope}%
\pgfsys@transformshift{3.072755in}{2.379541in}%
\pgfsys@useobject{currentmarker}{}%
\end{pgfscope}%
\begin{pgfscope}%
\pgfsys@transformshift{3.058517in}{2.379541in}%
\pgfsys@useobject{currentmarker}{}%
\end{pgfscope}%
\begin{pgfscope}%
\pgfsys@transformshift{3.068264in}{2.379541in}%
\pgfsys@useobject{currentmarker}{}%
\end{pgfscope}%
\begin{pgfscope}%
\pgfsys@transformshift{3.069421in}{2.379541in}%
\pgfsys@useobject{currentmarker}{}%
\end{pgfscope}%
\end{pgfscope}%
\begin{pgfscope}%
\pgfpathrectangle{\pgfqpoint{2.629775in}{0.720000in}}{\pgfqpoint{0.550375in}{2.664000in}}%
\pgfusepath{clip}%
\pgfsetrectcap%
\pgfsetroundjoin%
\pgfsetlinewidth{0.702625pt}%
\definecolor{currentstroke}{rgb}{0.101961,0.101961,0.101961}%
\pgfsetstrokecolor{currentstroke}%
\pgfsetdash{}{0pt}%
\pgfpathmoveto{\pgfqpoint{3.073132in}{2.161180in}}%
\pgfpathlineto{\pgfqpoint{3.101366in}{2.161180in}}%
\pgfusepath{stroke}%
\end{pgfscope}%
\begin{pgfscope}%
\pgfpathrectangle{\pgfqpoint{2.629775in}{0.720000in}}{\pgfqpoint{0.550375in}{2.664000in}}%
\pgfusepath{clip}%
\pgfsetbuttcap%
\pgfsetroundjoin%
\definecolor{currentfill}{rgb}{0.101961,0.101961,0.101961}%
\pgfsetfillcolor{currentfill}%
\pgfsetlinewidth{0.501875pt}%
\definecolor{currentstroke}{rgb}{0.101961,0.101961,0.101961}%
\pgfsetstrokecolor{currentstroke}%
\pgfsetdash{}{0pt}%
\pgfsys@defobject{currentmarker}{\pgfqpoint{0.000000in}{-0.020833in}}{\pgfqpoint{0.000000in}{0.020833in}}{%
\pgfpathmoveto{\pgfqpoint{0.000000in}{-0.020833in}}%
\pgfpathlineto{\pgfqpoint{0.000000in}{0.020833in}}%
\pgfusepath{stroke,fill}%
}%
\begin{pgfscope}%
\pgfsys@transformshift{3.092877in}{2.161180in}%
\pgfsys@useobject{currentmarker}{}%
\end{pgfscope}%
\begin{pgfscope}%
\pgfsys@transformshift{3.101366in}{2.161180in}%
\pgfsys@useobject{currentmarker}{}%
\end{pgfscope}%
\begin{pgfscope}%
\pgfsys@transformshift{3.092673in}{2.161180in}%
\pgfsys@useobject{currentmarker}{}%
\end{pgfscope}%
\begin{pgfscope}%
\pgfsys@transformshift{3.073132in}{2.161180in}%
\pgfsys@useobject{currentmarker}{}%
\end{pgfscope}%
\end{pgfscope}%
\begin{pgfscope}%
\pgfpathrectangle{\pgfqpoint{2.629775in}{0.720000in}}{\pgfqpoint{0.550375in}{2.664000in}}%
\pgfusepath{clip}%
\pgfsetrectcap%
\pgfsetroundjoin%
\pgfsetlinewidth{0.702625pt}%
\definecolor{currentstroke}{rgb}{0.101961,0.101961,0.101961}%
\pgfsetstrokecolor{currentstroke}%
\pgfsetdash{}{0pt}%
\pgfpathmoveto{\pgfqpoint{3.046939in}{1.942820in}}%
\pgfpathlineto{\pgfqpoint{3.093403in}{1.942820in}}%
\pgfusepath{stroke}%
\end{pgfscope}%
\begin{pgfscope}%
\pgfpathrectangle{\pgfqpoint{2.629775in}{0.720000in}}{\pgfqpoint{0.550375in}{2.664000in}}%
\pgfusepath{clip}%
\pgfsetbuttcap%
\pgfsetroundjoin%
\definecolor{currentfill}{rgb}{0.101961,0.101961,0.101961}%
\pgfsetfillcolor{currentfill}%
\pgfsetlinewidth{0.501875pt}%
\definecolor{currentstroke}{rgb}{0.101961,0.101961,0.101961}%
\pgfsetstrokecolor{currentstroke}%
\pgfsetdash{}{0pt}%
\pgfsys@defobject{currentmarker}{\pgfqpoint{0.000000in}{-0.020833in}}{\pgfqpoint{0.000000in}{0.020833in}}{%
\pgfpathmoveto{\pgfqpoint{0.000000in}{-0.020833in}}%
\pgfpathlineto{\pgfqpoint{0.000000in}{0.020833in}}%
\pgfusepath{stroke,fill}%
}%
\begin{pgfscope}%
\pgfsys@transformshift{3.093403in}{1.942820in}%
\pgfsys@useobject{currentmarker}{}%
\end{pgfscope}%
\begin{pgfscope}%
\pgfsys@transformshift{3.053138in}{1.942820in}%
\pgfsys@useobject{currentmarker}{}%
\end{pgfscope}%
\begin{pgfscope}%
\pgfsys@transformshift{3.046939in}{1.942820in}%
\pgfsys@useobject{currentmarker}{}%
\end{pgfscope}%
\end{pgfscope}%
\begin{pgfscope}%
\pgfpathrectangle{\pgfqpoint{2.629775in}{0.720000in}}{\pgfqpoint{0.550375in}{2.664000in}}%
\pgfusepath{clip}%
\pgfsetrectcap%
\pgfsetroundjoin%
\pgfsetlinewidth{0.702625pt}%
\definecolor{currentstroke}{rgb}{0.101961,0.101961,0.101961}%
\pgfsetstrokecolor{currentstroke}%
\pgfsetdash{}{0pt}%
\pgfpathmoveto{\pgfqpoint{3.046408in}{1.724459in}}%
\pgfpathlineto{\pgfqpoint{3.153942in}{1.724459in}}%
\pgfusepath{stroke}%
\end{pgfscope}%
\begin{pgfscope}%
\pgfpathrectangle{\pgfqpoint{2.629775in}{0.720000in}}{\pgfqpoint{0.550375in}{2.664000in}}%
\pgfusepath{clip}%
\pgfsetbuttcap%
\pgfsetroundjoin%
\definecolor{currentfill}{rgb}{0.101961,0.101961,0.101961}%
\pgfsetfillcolor{currentfill}%
\pgfsetlinewidth{0.501875pt}%
\definecolor{currentstroke}{rgb}{0.101961,0.101961,0.101961}%
\pgfsetstrokecolor{currentstroke}%
\pgfsetdash{}{0pt}%
\pgfsys@defobject{currentmarker}{\pgfqpoint{0.000000in}{-0.020833in}}{\pgfqpoint{0.000000in}{0.020833in}}{%
\pgfpathmoveto{\pgfqpoint{0.000000in}{-0.020833in}}%
\pgfpathlineto{\pgfqpoint{0.000000in}{0.020833in}}%
\pgfusepath{stroke,fill}%
}%
\begin{pgfscope}%
\pgfsys@transformshift{3.084334in}{1.724459in}%
\pgfsys@useobject{currentmarker}{}%
\end{pgfscope}%
\begin{pgfscope}%
\pgfsys@transformshift{3.153942in}{1.724459in}%
\pgfsys@useobject{currentmarker}{}%
\end{pgfscope}%
\begin{pgfscope}%
\pgfsys@transformshift{3.061421in}{1.724459in}%
\pgfsys@useobject{currentmarker}{}%
\end{pgfscope}%
\begin{pgfscope}%
\pgfsys@transformshift{3.046408in}{1.724459in}%
\pgfsys@useobject{currentmarker}{}%
\end{pgfscope}%
\end{pgfscope}%
\begin{pgfscope}%
\pgfpathrectangle{\pgfqpoint{2.629775in}{0.720000in}}{\pgfqpoint{0.550375in}{2.664000in}}%
\pgfusepath{clip}%
\pgfsetrectcap%
\pgfsetroundjoin%
\pgfsetlinewidth{0.702625pt}%
\definecolor{currentstroke}{rgb}{0.101961,0.101961,0.101961}%
\pgfsetstrokecolor{currentstroke}%
\pgfsetdash{}{0pt}%
\pgfpathmoveto{\pgfqpoint{3.058992in}{1.506098in}}%
\pgfpathlineto{\pgfqpoint{3.091259in}{1.506098in}}%
\pgfusepath{stroke}%
\end{pgfscope}%
\begin{pgfscope}%
\pgfpathrectangle{\pgfqpoint{2.629775in}{0.720000in}}{\pgfqpoint{0.550375in}{2.664000in}}%
\pgfusepath{clip}%
\pgfsetbuttcap%
\pgfsetroundjoin%
\definecolor{currentfill}{rgb}{0.101961,0.101961,0.101961}%
\pgfsetfillcolor{currentfill}%
\pgfsetlinewidth{0.501875pt}%
\definecolor{currentstroke}{rgb}{0.101961,0.101961,0.101961}%
\pgfsetstrokecolor{currentstroke}%
\pgfsetdash{}{0pt}%
\pgfsys@defobject{currentmarker}{\pgfqpoint{0.000000in}{-0.020833in}}{\pgfqpoint{0.000000in}{0.020833in}}{%
\pgfpathmoveto{\pgfqpoint{0.000000in}{-0.020833in}}%
\pgfpathlineto{\pgfqpoint{0.000000in}{0.020833in}}%
\pgfusepath{stroke,fill}%
}%
\begin{pgfscope}%
\pgfsys@transformshift{3.091259in}{1.506098in}%
\pgfsys@useobject{currentmarker}{}%
\end{pgfscope}%
\begin{pgfscope}%
\pgfsys@transformshift{3.073165in}{1.506098in}%
\pgfsys@useobject{currentmarker}{}%
\end{pgfscope}%
\begin{pgfscope}%
\pgfsys@transformshift{3.058992in}{1.506098in}%
\pgfsys@useobject{currentmarker}{}%
\end{pgfscope}%
\begin{pgfscope}%
\pgfsys@transformshift{3.066309in}{1.506098in}%
\pgfsys@useobject{currentmarker}{}%
\end{pgfscope}%
\end{pgfscope}%
\begin{pgfscope}%
\pgfpathrectangle{\pgfqpoint{2.629775in}{0.720000in}}{\pgfqpoint{0.550375in}{2.664000in}}%
\pgfusepath{clip}%
\pgfsetrectcap%
\pgfsetroundjoin%
\pgfsetlinewidth{0.702625pt}%
\definecolor{currentstroke}{rgb}{0.101961,0.101961,0.101961}%
\pgfsetstrokecolor{currentstroke}%
\pgfsetdash{}{0pt}%
\pgfpathmoveto{\pgfqpoint{3.049132in}{1.287738in}}%
\pgfpathlineto{\pgfqpoint{3.094446in}{1.287738in}}%
\pgfusepath{stroke}%
\end{pgfscope}%
\begin{pgfscope}%
\pgfpathrectangle{\pgfqpoint{2.629775in}{0.720000in}}{\pgfqpoint{0.550375in}{2.664000in}}%
\pgfusepath{clip}%
\pgfsetbuttcap%
\pgfsetroundjoin%
\definecolor{currentfill}{rgb}{0.101961,0.101961,0.101961}%
\pgfsetfillcolor{currentfill}%
\pgfsetlinewidth{0.501875pt}%
\definecolor{currentstroke}{rgb}{0.101961,0.101961,0.101961}%
\pgfsetstrokecolor{currentstroke}%
\pgfsetdash{}{0pt}%
\pgfsys@defobject{currentmarker}{\pgfqpoint{0.000000in}{-0.020833in}}{\pgfqpoint{0.000000in}{0.020833in}}{%
\pgfpathmoveto{\pgfqpoint{0.000000in}{-0.020833in}}%
\pgfpathlineto{\pgfqpoint{0.000000in}{0.020833in}}%
\pgfusepath{stroke,fill}%
}%
\begin{pgfscope}%
\pgfsys@transformshift{3.060020in}{1.287738in}%
\pgfsys@useobject{currentmarker}{}%
\end{pgfscope}%
\begin{pgfscope}%
\pgfsys@transformshift{3.073500in}{1.287738in}%
\pgfsys@useobject{currentmarker}{}%
\end{pgfscope}%
\begin{pgfscope}%
\pgfsys@transformshift{3.049132in}{1.287738in}%
\pgfsys@useobject{currentmarker}{}%
\end{pgfscope}%
\begin{pgfscope}%
\pgfsys@transformshift{3.094446in}{1.287738in}%
\pgfsys@useobject{currentmarker}{}%
\end{pgfscope}%
\end{pgfscope}%
\begin{pgfscope}%
\pgfpathrectangle{\pgfqpoint{2.629775in}{0.720000in}}{\pgfqpoint{0.550375in}{2.664000in}}%
\pgfusepath{clip}%
\pgfsetrectcap%
\pgfsetroundjoin%
\pgfsetlinewidth{0.702625pt}%
\definecolor{currentstroke}{rgb}{0.101961,0.101961,0.101961}%
\pgfsetstrokecolor{currentstroke}%
\pgfsetdash{}{0pt}%
\pgfpathmoveto{\pgfqpoint{3.047547in}{1.069377in}}%
\pgfpathlineto{\pgfqpoint{3.071428in}{1.069377in}}%
\pgfusepath{stroke}%
\end{pgfscope}%
\begin{pgfscope}%
\pgfpathrectangle{\pgfqpoint{2.629775in}{0.720000in}}{\pgfqpoint{0.550375in}{2.664000in}}%
\pgfusepath{clip}%
\pgfsetbuttcap%
\pgfsetroundjoin%
\definecolor{currentfill}{rgb}{0.101961,0.101961,0.101961}%
\pgfsetfillcolor{currentfill}%
\pgfsetlinewidth{0.501875pt}%
\definecolor{currentstroke}{rgb}{0.101961,0.101961,0.101961}%
\pgfsetstrokecolor{currentstroke}%
\pgfsetdash{}{0pt}%
\pgfsys@defobject{currentmarker}{\pgfqpoint{0.000000in}{-0.020833in}}{\pgfqpoint{0.000000in}{0.020833in}}{%
\pgfpathmoveto{\pgfqpoint{0.000000in}{-0.020833in}}%
\pgfpathlineto{\pgfqpoint{0.000000in}{0.020833in}}%
\pgfusepath{stroke,fill}%
}%
\begin{pgfscope}%
\pgfsys@transformshift{3.071428in}{1.069377in}%
\pgfsys@useobject{currentmarker}{}%
\end{pgfscope}%
\begin{pgfscope}%
\pgfsys@transformshift{3.069222in}{1.069377in}%
\pgfsys@useobject{currentmarker}{}%
\end{pgfscope}%
\begin{pgfscope}%
\pgfsys@transformshift{3.048947in}{1.069377in}%
\pgfsys@useobject{currentmarker}{}%
\end{pgfscope}%
\begin{pgfscope}%
\pgfsys@transformshift{3.047547in}{1.069377in}%
\pgfsys@useobject{currentmarker}{}%
\end{pgfscope}%
\end{pgfscope}%
\begin{pgfscope}%
\pgfpathrectangle{\pgfqpoint{2.629775in}{0.720000in}}{\pgfqpoint{0.550375in}{2.664000in}}%
\pgfusepath{clip}%
\pgfsetrectcap%
\pgfsetroundjoin%
\pgfsetlinewidth{0.702625pt}%
\definecolor{currentstroke}{rgb}{0.101961,0.101961,0.101961}%
\pgfsetstrokecolor{currentstroke}%
\pgfsetdash{}{0pt}%
\pgfpathmoveto{\pgfqpoint{3.009630in}{0.851016in}}%
\pgfpathlineto{\pgfqpoint{3.100244in}{0.851016in}}%
\pgfusepath{stroke}%
\end{pgfscope}%
\begin{pgfscope}%
\pgfpathrectangle{\pgfqpoint{2.629775in}{0.720000in}}{\pgfqpoint{0.550375in}{2.664000in}}%
\pgfusepath{clip}%
\pgfsetbuttcap%
\pgfsetroundjoin%
\definecolor{currentfill}{rgb}{0.101961,0.101961,0.101961}%
\pgfsetfillcolor{currentfill}%
\pgfsetlinewidth{0.501875pt}%
\definecolor{currentstroke}{rgb}{0.101961,0.101961,0.101961}%
\pgfsetstrokecolor{currentstroke}%
\pgfsetdash{}{0pt}%
\pgfsys@defobject{currentmarker}{\pgfqpoint{0.000000in}{-0.020833in}}{\pgfqpoint{0.000000in}{0.020833in}}{%
\pgfpathmoveto{\pgfqpoint{0.000000in}{-0.020833in}}%
\pgfpathlineto{\pgfqpoint{0.000000in}{0.020833in}}%
\pgfusepath{stroke,fill}%
}%
\begin{pgfscope}%
\pgfsys@transformshift{3.090485in}{0.851016in}%
\pgfsys@useobject{currentmarker}{}%
\end{pgfscope}%
\begin{pgfscope}%
\pgfsys@transformshift{3.075000in}{0.851016in}%
\pgfsys@useobject{currentmarker}{}%
\end{pgfscope}%
\begin{pgfscope}%
\pgfsys@transformshift{3.009630in}{0.851016in}%
\pgfsys@useobject{currentmarker}{}%
\end{pgfscope}%
\begin{pgfscope}%
\pgfsys@transformshift{3.100244in}{0.851016in}%
\pgfsys@useobject{currentmarker}{}%
\end{pgfscope}%
\end{pgfscope}%
\begin{pgfscope}%
\pgfsetrectcap%
\pgfsetmiterjoin%
\pgfsetlinewidth{0.803000pt}%
\definecolor{currentstroke}{rgb}{0.101961,0.101961,0.101961}%
\pgfsetstrokecolor{currentstroke}%
\pgfsetdash{}{0pt}%
\pgfpathmoveto{\pgfqpoint{2.629775in}{0.720000in}}%
\pgfpathlineto{\pgfqpoint{2.629775in}{3.384000in}}%
\pgfusepath{stroke}%
\end{pgfscope}%
\begin{pgfscope}%
\pgfsetrectcap%
\pgfsetmiterjoin%
\pgfsetlinewidth{0.803000pt}%
\definecolor{currentstroke}{rgb}{0.101961,0.101961,0.101961}%
\pgfsetstrokecolor{currentstroke}%
\pgfsetdash{}{0pt}%
\pgfpathmoveto{\pgfqpoint{2.629775in}{0.720000in}}%
\pgfpathlineto{\pgfqpoint{3.180150in}{0.720000in}}%
\pgfusepath{stroke}%
\end{pgfscope}%
\begin{pgfscope}%
\definecolor{textcolor}{rgb}{0.549020,0.549020,0.549020}%
\pgfsetstrokecolor{textcolor}%
\pgfsetfillcolor{textcolor}%
\pgftext[x=3.296700in,y=3.546000in,right,top]{\color{textcolor}\sffamily\fontsize{5.500000}{6.600000}\selectfont OLD DATA}%
\end{pgfscope}%
\end{pgfpicture}%
\makeatother%
\endgroup%
}
\caption{[PLACEHOLDER: OLD DATA, COVEX is the rung 3 covex-v6-fair-final stand-in (OLD); the DEV candidates covex-v6-fair-final-cand-tl and -sync are not in F13.csv; arms with blank GPU fields in every standard F13 row (not plotted): lmcache, warm-all-lru; seeds missing per arm: mooncake has 1/4; eco-llm has 3/4] Efficiency of the cloud GPU on the standard grid (16 devices, 1x memory, 3 s robot deadline, seeds 4 to 7): correct answers (robot plus single-shot) per measured GPU-second, GPU busy fraction of the 1200 s per-run budget, and D174 uncached tokens charged per GPU-second. Bars are the mean over the seeds present for the arm; marks show each seed. Arms with no committed GPU metrics (LMCache, Warm-all LRU) are not drawn. Blue is COVEX, dark grey stateful baselines, light grey stateless baselines.}
\label{fig:F13_gpu_efficiency}
\end{figure}
\clearpage
\begin{figure}[t]
\centering
{\fontfamily{LinuxBiolinumT}\selectfont
%% Creator: Matplotlib, PGF backend
%%
%% To include the figure in your LaTeX document, write
%%   \input{<filename>.pgf}
%%
%% Make sure the required packages are loaded in your preamble
%%   \usepackage{pgf}
%%
%% Also ensure that all the required font packages are loaded; for instance,
%% the lmodern package is sometimes necessary when using math font.
%%   \usepackage{lmodern}
%%
%% Figures using additional raster images can only be included by \input if
%% they are in the same directory as the main LaTeX file. For loading figures
%% from other directories you can use the `import` package
%%   \usepackage{import}
%%
%% and then include the figures with
%%   \import{<path to file>}{<filename>.pgf}
%%
%% Matplotlib used the following preamble
%%
\begingroup%
\makeatletter%
\begin{pgfpicture}%
\pgfpathrectangle{\pgfpointorigin}{\pgfqpoint{3.330000in}{2.600000in}}%
\pgfusepath{use as bounding box, clip}%
\begin{pgfscope}%
\pgfsetbuttcap%
\pgfsetmiterjoin%
\definecolor{currentfill}{rgb}{1.000000,1.000000,1.000000}%
\pgfsetfillcolor{currentfill}%
\pgfsetlinewidth{0.000000pt}%
\definecolor{currentstroke}{rgb}{1.000000,1.000000,1.000000}%
\pgfsetstrokecolor{currentstroke}%
\pgfsetdash{}{0pt}%
\pgfpathmoveto{\pgfqpoint{0.000000in}{0.000000in}}%
\pgfpathlineto{\pgfqpoint{3.330000in}{0.000000in}}%
\pgfpathlineto{\pgfqpoint{3.330000in}{2.600000in}}%
\pgfpathlineto{\pgfqpoint{0.000000in}{2.600000in}}%
\pgfpathlineto{\pgfqpoint{0.000000in}{0.000000in}}%
\pgfpathclose%
\pgfusepath{fill}%
\end{pgfscope}%
\begin{pgfscope}%
\pgfsetbuttcap%
\pgfsetmiterjoin%
\definecolor{currentfill}{rgb}{1.000000,1.000000,1.000000}%
\pgfsetfillcolor{currentfill}%
\pgfsetlinewidth{0.000000pt}%
\definecolor{currentstroke}{rgb}{0.000000,0.000000,0.000000}%
\pgfsetstrokecolor{currentstroke}%
\pgfsetstrokeopacity{0.000000}%
\pgfsetdash{}{0pt}%
\pgfpathmoveto{\pgfqpoint{0.466200in}{0.442000in}}%
\pgfpathlineto{\pgfqpoint{3.230100in}{0.442000in}}%
\pgfpathlineto{\pgfqpoint{3.230100in}{2.210000in}}%
\pgfpathlineto{\pgfqpoint{0.466200in}{2.210000in}}%
\pgfpathlineto{\pgfqpoint{0.466200in}{0.442000in}}%
\pgfpathclose%
\pgfusepath{fill}%
\end{pgfscope}%
\begin{pgfscope}%
\definecolor{textcolor}{rgb}{0.101961,0.101961,0.101961}%
\pgfsetstrokecolor{textcolor}%
\pgfsetfillcolor{textcolor}%
\pgftext[x=0.591832in,y=0.393389in,,top]{\color{textcolor}\sffamily\fontsize{7.000000}{8.400000}\selectfont 8}%
\end{pgfscope}%
\begin{pgfscope}%
\definecolor{textcolor}{rgb}{0.101961,0.101961,0.101961}%
\pgfsetstrokecolor{textcolor}%
\pgfsetfillcolor{textcolor}%
\pgftext[x=1.010605in,y=0.393389in,,top]{\color{textcolor}\sffamily\fontsize{7.000000}{8.400000}\selectfont 12}%
\end{pgfscope}%
\begin{pgfscope}%
\definecolor{textcolor}{rgb}{0.101961,0.101961,0.101961}%
\pgfsetstrokecolor{textcolor}%
\pgfsetfillcolor{textcolor}%
\pgftext[x=1.429377in,y=0.393389in,,top]{\color{textcolor}\sffamily\fontsize{7.000000}{8.400000}\selectfont 16}%
\end{pgfscope}%
\begin{pgfscope}%
\definecolor{textcolor}{rgb}{0.101961,0.101961,0.101961}%
\pgfsetstrokecolor{textcolor}%
\pgfsetfillcolor{textcolor}%
\pgftext[x=1.848150in,y=0.393389in,,top]{\color{textcolor}\sffamily\fontsize{7.000000}{8.400000}\selectfont 20}%
\end{pgfscope}%
\begin{pgfscope}%
\definecolor{textcolor}{rgb}{0.101961,0.101961,0.101961}%
\pgfsetstrokecolor{textcolor}%
\pgfsetfillcolor{textcolor}%
\pgftext[x=2.266923in,y=0.393389in,,top]{\color{textcolor}\sffamily\fontsize{7.000000}{8.400000}\selectfont 24}%
\end{pgfscope}%
\begin{pgfscope}%
\definecolor{textcolor}{rgb}{0.101961,0.101961,0.101961}%
\pgfsetstrokecolor{textcolor}%
\pgfsetfillcolor{textcolor}%
\pgftext[x=3.104468in,y=0.393389in,,top]{\color{textcolor}\sffamily\fontsize{7.000000}{8.400000}\selectfont 32}%
\end{pgfscope}%
\begin{pgfscope}%
\definecolor{textcolor}{rgb}{0.000000,0.000000,0.000000}%
\pgfsetstrokecolor{textcolor}%
\pgfsetfillcolor{textcolor}%
\pgftext[x=1.848150in,y=0.251414in,,top]{\color{textcolor}\sffamily\fontsize{7.500000}{9.000000}\selectfont devices in the fleet}%
\end{pgfscope}%
\begin{pgfscope}%
\pgfpathrectangle{\pgfqpoint{0.466200in}{0.442000in}}{\pgfqpoint{2.763900in}{1.768000in}}%
\pgfusepath{clip}%
\pgfsetrectcap%
\pgfsetroundjoin%
\pgfsetlinewidth{0.401500pt}%
\definecolor{currentstroke}{rgb}{0.901961,0.901961,0.901961}%
\pgfsetstrokecolor{currentstroke}%
\pgfsetdash{}{0pt}%
\pgfpathmoveto{\pgfqpoint{0.466200in}{0.442000in}}%
\pgfpathlineto{\pgfqpoint{3.230100in}{0.442000in}}%
\pgfusepath{stroke}%
\end{pgfscope}%
\begin{pgfscope}%
\definecolor{textcolor}{rgb}{0.101961,0.101961,0.101961}%
\pgfsetstrokecolor{textcolor}%
\pgfsetfillcolor{textcolor}%
\pgftext[x=0.365939in, y=0.408242in, left, base]{\color{textcolor}\sffamily\fontsize{7.000000}{8.400000}\selectfont 0}%
\end{pgfscope}%
\begin{pgfscope}%
\pgfpathrectangle{\pgfqpoint{0.466200in}{0.442000in}}{\pgfqpoint{2.763900in}{1.768000in}}%
\pgfusepath{clip}%
\pgfsetrectcap%
\pgfsetroundjoin%
\pgfsetlinewidth{0.401500pt}%
\definecolor{currentstroke}{rgb}{0.901961,0.901961,0.901961}%
\pgfsetstrokecolor{currentstroke}%
\pgfsetdash{}{0pt}%
\pgfpathmoveto{\pgfqpoint{0.466200in}{0.680919in}}%
\pgfpathlineto{\pgfqpoint{3.230100in}{0.680919in}}%
\pgfusepath{stroke}%
\end{pgfscope}%
\begin{pgfscope}%
\definecolor{textcolor}{rgb}{0.101961,0.101961,0.101961}%
\pgfsetstrokecolor{textcolor}%
\pgfsetfillcolor{textcolor}%
\pgftext[x=0.365939in, y=0.647161in, left, base]{\color{textcolor}\sffamily\fontsize{7.000000}{8.400000}\selectfont 1}%
\end{pgfscope}%
\begin{pgfscope}%
\pgfpathrectangle{\pgfqpoint{0.466200in}{0.442000in}}{\pgfqpoint{2.763900in}{1.768000in}}%
\pgfusepath{clip}%
\pgfsetrectcap%
\pgfsetroundjoin%
\pgfsetlinewidth{0.401500pt}%
\definecolor{currentstroke}{rgb}{0.901961,0.901961,0.901961}%
\pgfsetstrokecolor{currentstroke}%
\pgfsetdash{}{0pt}%
\pgfpathmoveto{\pgfqpoint{0.466200in}{0.919838in}}%
\pgfpathlineto{\pgfqpoint{3.230100in}{0.919838in}}%
\pgfusepath{stroke}%
\end{pgfscope}%
\begin{pgfscope}%
\definecolor{textcolor}{rgb}{0.101961,0.101961,0.101961}%
\pgfsetstrokecolor{textcolor}%
\pgfsetfillcolor{textcolor}%
\pgftext[x=0.365939in, y=0.886080in, left, base]{\color{textcolor}\sffamily\fontsize{7.000000}{8.400000}\selectfont 2}%
\end{pgfscope}%
\begin{pgfscope}%
\pgfpathrectangle{\pgfqpoint{0.466200in}{0.442000in}}{\pgfqpoint{2.763900in}{1.768000in}}%
\pgfusepath{clip}%
\pgfsetrectcap%
\pgfsetroundjoin%
\pgfsetlinewidth{0.401500pt}%
\definecolor{currentstroke}{rgb}{0.901961,0.901961,0.901961}%
\pgfsetstrokecolor{currentstroke}%
\pgfsetdash{}{0pt}%
\pgfpathmoveto{\pgfqpoint{0.466200in}{1.158757in}}%
\pgfpathlineto{\pgfqpoint{3.230100in}{1.158757in}}%
\pgfusepath{stroke}%
\end{pgfscope}%
\begin{pgfscope}%
\definecolor{textcolor}{rgb}{0.101961,0.101961,0.101961}%
\pgfsetstrokecolor{textcolor}%
\pgfsetfillcolor{textcolor}%
\pgftext[x=0.365939in, y=1.124999in, left, base]{\color{textcolor}\sffamily\fontsize{7.000000}{8.400000}\selectfont 3}%
\end{pgfscope}%
\begin{pgfscope}%
\pgfpathrectangle{\pgfqpoint{0.466200in}{0.442000in}}{\pgfqpoint{2.763900in}{1.768000in}}%
\pgfusepath{clip}%
\pgfsetrectcap%
\pgfsetroundjoin%
\pgfsetlinewidth{0.401500pt}%
\definecolor{currentstroke}{rgb}{0.901961,0.901961,0.901961}%
\pgfsetstrokecolor{currentstroke}%
\pgfsetdash{}{0pt}%
\pgfpathmoveto{\pgfqpoint{0.466200in}{1.397676in}}%
\pgfpathlineto{\pgfqpoint{3.230100in}{1.397676in}}%
\pgfusepath{stroke}%
\end{pgfscope}%
\begin{pgfscope}%
\definecolor{textcolor}{rgb}{0.101961,0.101961,0.101961}%
\pgfsetstrokecolor{textcolor}%
\pgfsetfillcolor{textcolor}%
\pgftext[x=0.365939in, y=1.363918in, left, base]{\color{textcolor}\sffamily\fontsize{7.000000}{8.400000}\selectfont 4}%
\end{pgfscope}%
\begin{pgfscope}%
\pgfpathrectangle{\pgfqpoint{0.466200in}{0.442000in}}{\pgfqpoint{2.763900in}{1.768000in}}%
\pgfusepath{clip}%
\pgfsetrectcap%
\pgfsetroundjoin%
\pgfsetlinewidth{0.401500pt}%
\definecolor{currentstroke}{rgb}{0.901961,0.901961,0.901961}%
\pgfsetstrokecolor{currentstroke}%
\pgfsetdash{}{0pt}%
\pgfpathmoveto{\pgfqpoint{0.466200in}{1.636595in}}%
\pgfpathlineto{\pgfqpoint{3.230100in}{1.636595in}}%
\pgfusepath{stroke}%
\end{pgfscope}%
\begin{pgfscope}%
\definecolor{textcolor}{rgb}{0.101961,0.101961,0.101961}%
\pgfsetstrokecolor{textcolor}%
\pgfsetfillcolor{textcolor}%
\pgftext[x=0.365939in, y=1.602837in, left, base]{\color{textcolor}\sffamily\fontsize{7.000000}{8.400000}\selectfont 5}%
\end{pgfscope}%
\begin{pgfscope}%
\definecolor{textcolor}{rgb}{0.000000,0.000000,0.000000}%
\pgfsetstrokecolor{textcolor}%
\pgfsetfillcolor{textcolor}%
\pgftext[x=0.310383in,y=1.326000in,,bottom,rotate=90.000000]{\color{textcolor}\sffamily\fontsize{7.500000}{9.000000}\selectfont decision latency, wall (ms)}%
\end{pgfscope}%
\begin{pgfscope}%
\pgfpathrectangle{\pgfqpoint{0.466200in}{0.442000in}}{\pgfqpoint{2.763900in}{1.768000in}}%
\pgfusepath{clip}%
\pgfsetrectcap%
\pgfsetroundjoin%
\pgfsetlinewidth{1.003750pt}%
\definecolor{currentstroke}{rgb}{0.156863,0.356863,0.568627}%
\pgfsetstrokecolor{currentstroke}%
\pgfsetdash{}{0pt}%
\pgfpathmoveto{\pgfqpoint{1.429377in}{0.930476in}}%
\pgfpathlineto{\pgfqpoint{1.848150in}{0.815550in}}%
\pgfpathlineto{\pgfqpoint{2.266923in}{0.810724in}}%
\pgfpathlineto{\pgfqpoint{3.104468in}{0.777896in}}%
\pgfusepath{stroke}%
\end{pgfscope}%
\begin{pgfscope}%
\pgfpathrectangle{\pgfqpoint{0.466200in}{0.442000in}}{\pgfqpoint{2.763900in}{1.768000in}}%
\pgfusepath{clip}%
\pgfsetbuttcap%
\pgfsetroundjoin%
\definecolor{currentfill}{rgb}{0.156863,0.356863,0.568627}%
\pgfsetfillcolor{currentfill}%
\pgfsetlinewidth{1.003750pt}%
\definecolor{currentstroke}{rgb}{0.156863,0.356863,0.568627}%
\pgfsetstrokecolor{currentstroke}%
\pgfsetdash{}{0pt}%
\pgfsys@defobject{currentmarker}{\pgfqpoint{-0.022222in}{-0.022222in}}{\pgfqpoint{0.022222in}{0.022222in}}{%
\pgfpathmoveto{\pgfqpoint{0.000000in}{-0.022222in}}%
\pgfpathcurveto{\pgfqpoint{0.005893in}{-0.022222in}}{\pgfqpoint{0.011546in}{-0.019881in}}{\pgfqpoint{0.015713in}{-0.015713in}}%
\pgfpathcurveto{\pgfqpoint{0.019881in}{-0.011546in}}{\pgfqpoint{0.022222in}{-0.005893in}}{\pgfqpoint{0.022222in}{0.000000in}}%
\pgfpathcurveto{\pgfqpoint{0.022222in}{0.005893in}}{\pgfqpoint{0.019881in}{0.011546in}}{\pgfqpoint{0.015713in}{0.015713in}}%
\pgfpathcurveto{\pgfqpoint{0.011546in}{0.019881in}}{\pgfqpoint{0.005893in}{0.022222in}}{\pgfqpoint{0.000000in}{0.022222in}}%
\pgfpathcurveto{\pgfqpoint{-0.005893in}{0.022222in}}{\pgfqpoint{-0.011546in}{0.019881in}}{\pgfqpoint{-0.015713in}{0.015713in}}%
\pgfpathcurveto{\pgfqpoint{-0.019881in}{0.011546in}}{\pgfqpoint{-0.022222in}{0.005893in}}{\pgfqpoint{-0.022222in}{0.000000in}}%
\pgfpathcurveto{\pgfqpoint{-0.022222in}{-0.005893in}}{\pgfqpoint{-0.019881in}{-0.011546in}}{\pgfqpoint{-0.015713in}{-0.015713in}}%
\pgfpathcurveto{\pgfqpoint{-0.011546in}{-0.019881in}}{\pgfqpoint{-0.005893in}{-0.022222in}}{\pgfqpoint{0.000000in}{-0.022222in}}%
\pgfpathlineto{\pgfqpoint{0.000000in}{-0.022222in}}%
\pgfpathclose%
\pgfusepath{stroke,fill}%
}%
\begin{pgfscope}%
\pgfsys@transformshift{1.429377in}{0.930476in}%
\pgfsys@useobject{currentmarker}{}%
\end{pgfscope}%
\begin{pgfscope}%
\pgfsys@transformshift{1.848150in}{0.815550in}%
\pgfsys@useobject{currentmarker}{}%
\end{pgfscope}%
\begin{pgfscope}%
\pgfsys@transformshift{2.266923in}{0.810724in}%
\pgfsys@useobject{currentmarker}{}%
\end{pgfscope}%
\begin{pgfscope}%
\pgfsys@transformshift{3.104468in}{0.777896in}%
\pgfsys@useobject{currentmarker}{}%
\end{pgfscope}%
\end{pgfscope}%
\begin{pgfscope}%
\pgfpathrectangle{\pgfqpoint{0.466200in}{0.442000in}}{\pgfqpoint{2.763900in}{1.768000in}}%
\pgfusepath{clip}%
\pgfsetbuttcap%
\pgfsetroundjoin%
\pgfsetlinewidth{0.803000pt}%
\definecolor{currentstroke}{rgb}{0.156863,0.356863,0.568627}%
\pgfsetstrokecolor{currentstroke}%
\pgfsetdash{{2.000000pt}{1.200000pt}}{0.000000pt}%
\pgfpathmoveto{\pgfqpoint{1.429377in}{1.303242in}}%
\pgfpathlineto{\pgfqpoint{1.848150in}{1.486417in}}%
\pgfpathlineto{\pgfqpoint{2.266923in}{1.503501in}}%
\pgfpathlineto{\pgfqpoint{3.104468in}{1.529967in}}%
\pgfusepath{stroke}%
\end{pgfscope}%
\begin{pgfscope}%
\pgfpathrectangle{\pgfqpoint{0.466200in}{0.442000in}}{\pgfqpoint{2.763900in}{1.768000in}}%
\pgfusepath{clip}%
\pgfsetbuttcap%
\pgfsetroundjoin%
\definecolor{currentfill}{rgb}{1.000000,1.000000,1.000000}%
\pgfsetfillcolor{currentfill}%
\pgfsetlinewidth{0.702625pt}%
\definecolor{currentstroke}{rgb}{0.156863,0.356863,0.568627}%
\pgfsetstrokecolor{currentstroke}%
\pgfsetdash{}{0pt}%
\pgfsys@defobject{currentmarker}{\pgfqpoint{-0.022222in}{-0.022222in}}{\pgfqpoint{0.022222in}{0.022222in}}{%
\pgfpathmoveto{\pgfqpoint{0.000000in}{-0.022222in}}%
\pgfpathcurveto{\pgfqpoint{0.005893in}{-0.022222in}}{\pgfqpoint{0.011546in}{-0.019881in}}{\pgfqpoint{0.015713in}{-0.015713in}}%
\pgfpathcurveto{\pgfqpoint{0.019881in}{-0.011546in}}{\pgfqpoint{0.022222in}{-0.005893in}}{\pgfqpoint{0.022222in}{0.000000in}}%
\pgfpathcurveto{\pgfqpoint{0.022222in}{0.005893in}}{\pgfqpoint{0.019881in}{0.011546in}}{\pgfqpoint{0.015713in}{0.015713in}}%
\pgfpathcurveto{\pgfqpoint{0.011546in}{0.019881in}}{\pgfqpoint{0.005893in}{0.022222in}}{\pgfqpoint{0.000000in}{0.022222in}}%
\pgfpathcurveto{\pgfqpoint{-0.005893in}{0.022222in}}{\pgfqpoint{-0.011546in}{0.019881in}}{\pgfqpoint{-0.015713in}{0.015713in}}%
\pgfpathcurveto{\pgfqpoint{-0.019881in}{0.011546in}}{\pgfqpoint{-0.022222in}{0.005893in}}{\pgfqpoint{-0.022222in}{0.000000in}}%
\pgfpathcurveto{\pgfqpoint{-0.022222in}{-0.005893in}}{\pgfqpoint{-0.019881in}{-0.011546in}}{\pgfqpoint{-0.015713in}{-0.015713in}}%
\pgfpathcurveto{\pgfqpoint{-0.011546in}{-0.019881in}}{\pgfqpoint{-0.005893in}{-0.022222in}}{\pgfqpoint{0.000000in}{-0.022222in}}%
\pgfpathlineto{\pgfqpoint{0.000000in}{-0.022222in}}%
\pgfpathclose%
\pgfusepath{stroke,fill}%
}%
\begin{pgfscope}%
\pgfsys@transformshift{1.429377in}{1.303242in}%
\pgfsys@useobject{currentmarker}{}%
\end{pgfscope}%
\begin{pgfscope}%
\pgfsys@transformshift{1.848150in}{1.486417in}%
\pgfsys@useobject{currentmarker}{}%
\end{pgfscope}%
\begin{pgfscope}%
\pgfsys@transformshift{2.266923in}{1.503501in}%
\pgfsys@useobject{currentmarker}{}%
\end{pgfscope}%
\begin{pgfscope}%
\pgfsys@transformshift{3.104468in}{1.529967in}%
\pgfsys@useobject{currentmarker}{}%
\end{pgfscope}%
\end{pgfscope}%
\begin{pgfscope}%
\pgfpathrectangle{\pgfqpoint{0.466200in}{0.442000in}}{\pgfqpoint{2.763900in}{1.768000in}}%
\pgfusepath{clip}%
\pgfsetrectcap%
\pgfsetroundjoin%
\pgfsetlinewidth{1.505625pt}%
\definecolor{currentstroke}{rgb}{0.156863,0.356863,0.568627}%
\pgfsetstrokecolor{currentstroke}%
\pgfsetdash{}{0pt}%
\pgfpathmoveto{\pgfqpoint{1.848150in}{0.815550in}}%
\pgfusepath{stroke}%
\end{pgfscope}%
\begin{pgfscope}%
\pgfpathrectangle{\pgfqpoint{0.466200in}{0.442000in}}{\pgfqpoint{2.763900in}{1.768000in}}%
\pgfusepath{clip}%
\pgfsetbuttcap%
\pgfsetroundjoin%
\definecolor{currentfill}{rgb}{1.000000,1.000000,1.000000}%
\pgfsetfillcolor{currentfill}%
\pgfsetlinewidth{0.702625pt}%
\definecolor{currentstroke}{rgb}{0.156863,0.356863,0.568627}%
\pgfsetstrokecolor{currentstroke}%
\pgfsetdash{}{0pt}%
\pgfsys@defobject{currentmarker}{\pgfqpoint{-0.022222in}{-0.022222in}}{\pgfqpoint{0.022222in}{0.022222in}}{%
\pgfpathmoveto{\pgfqpoint{0.000000in}{-0.022222in}}%
\pgfpathcurveto{\pgfqpoint{0.005893in}{-0.022222in}}{\pgfqpoint{0.011546in}{-0.019881in}}{\pgfqpoint{0.015713in}{-0.015713in}}%
\pgfpathcurveto{\pgfqpoint{0.019881in}{-0.011546in}}{\pgfqpoint{0.022222in}{-0.005893in}}{\pgfqpoint{0.022222in}{0.000000in}}%
\pgfpathcurveto{\pgfqpoint{0.022222in}{0.005893in}}{\pgfqpoint{0.019881in}{0.011546in}}{\pgfqpoint{0.015713in}{0.015713in}}%
\pgfpathcurveto{\pgfqpoint{0.011546in}{0.019881in}}{\pgfqpoint{0.005893in}{0.022222in}}{\pgfqpoint{0.000000in}{0.022222in}}%
\pgfpathcurveto{\pgfqpoint{-0.005893in}{0.022222in}}{\pgfqpoint{-0.011546in}{0.019881in}}{\pgfqpoint{-0.015713in}{0.015713in}}%
\pgfpathcurveto{\pgfqpoint{-0.019881in}{0.011546in}}{\pgfqpoint{-0.022222in}{0.005893in}}{\pgfqpoint{-0.022222in}{0.000000in}}%
\pgfpathcurveto{\pgfqpoint{-0.022222in}{-0.005893in}}{\pgfqpoint{-0.019881in}{-0.011546in}}{\pgfqpoint{-0.015713in}{-0.015713in}}%
\pgfpathcurveto{\pgfqpoint{-0.011546in}{-0.019881in}}{\pgfqpoint{-0.005893in}{-0.022222in}}{\pgfqpoint{0.000000in}{-0.022222in}}%
\pgfpathlineto{\pgfqpoint{0.000000in}{-0.022222in}}%
\pgfpathclose%
\pgfusepath{stroke,fill}%
}%
\begin{pgfscope}%
\pgfsys@transformshift{1.848150in}{0.815550in}%
\pgfsys@useobject{currentmarker}{}%
\end{pgfscope}%
\end{pgfscope}%
\begin{pgfscope}%
\pgfpathrectangle{\pgfqpoint{0.466200in}{0.442000in}}{\pgfqpoint{2.763900in}{1.768000in}}%
\pgfusepath{clip}%
\pgfsetrectcap%
\pgfsetroundjoin%
\pgfsetlinewidth{1.505625pt}%
\definecolor{currentstroke}{rgb}{0.156863,0.356863,0.568627}%
\pgfsetstrokecolor{currentstroke}%
\pgfsetdash{}{0pt}%
\pgfpathmoveto{\pgfqpoint{2.266923in}{0.810724in}}%
\pgfusepath{stroke}%
\end{pgfscope}%
\begin{pgfscope}%
\pgfpathrectangle{\pgfqpoint{0.466200in}{0.442000in}}{\pgfqpoint{2.763900in}{1.768000in}}%
\pgfusepath{clip}%
\pgfsetbuttcap%
\pgfsetroundjoin%
\definecolor{currentfill}{rgb}{1.000000,1.000000,1.000000}%
\pgfsetfillcolor{currentfill}%
\pgfsetlinewidth{0.702625pt}%
\definecolor{currentstroke}{rgb}{0.156863,0.356863,0.568627}%
\pgfsetstrokecolor{currentstroke}%
\pgfsetdash{}{0pt}%
\pgfsys@defobject{currentmarker}{\pgfqpoint{-0.022222in}{-0.022222in}}{\pgfqpoint{0.022222in}{0.022222in}}{%
\pgfpathmoveto{\pgfqpoint{0.000000in}{-0.022222in}}%
\pgfpathcurveto{\pgfqpoint{0.005893in}{-0.022222in}}{\pgfqpoint{0.011546in}{-0.019881in}}{\pgfqpoint{0.015713in}{-0.015713in}}%
\pgfpathcurveto{\pgfqpoint{0.019881in}{-0.011546in}}{\pgfqpoint{0.022222in}{-0.005893in}}{\pgfqpoint{0.022222in}{0.000000in}}%
\pgfpathcurveto{\pgfqpoint{0.022222in}{0.005893in}}{\pgfqpoint{0.019881in}{0.011546in}}{\pgfqpoint{0.015713in}{0.015713in}}%
\pgfpathcurveto{\pgfqpoint{0.011546in}{0.019881in}}{\pgfqpoint{0.005893in}{0.022222in}}{\pgfqpoint{0.000000in}{0.022222in}}%
\pgfpathcurveto{\pgfqpoint{-0.005893in}{0.022222in}}{\pgfqpoint{-0.011546in}{0.019881in}}{\pgfqpoint{-0.015713in}{0.015713in}}%
\pgfpathcurveto{\pgfqpoint{-0.019881in}{0.011546in}}{\pgfqpoint{-0.022222in}{0.005893in}}{\pgfqpoint{-0.022222in}{0.000000in}}%
\pgfpathcurveto{\pgfqpoint{-0.022222in}{-0.005893in}}{\pgfqpoint{-0.019881in}{-0.011546in}}{\pgfqpoint{-0.015713in}{-0.015713in}}%
\pgfpathcurveto{\pgfqpoint{-0.011546in}{-0.019881in}}{\pgfqpoint{-0.005893in}{-0.022222in}}{\pgfqpoint{0.000000in}{-0.022222in}}%
\pgfpathlineto{\pgfqpoint{0.000000in}{-0.022222in}}%
\pgfpathclose%
\pgfusepath{stroke,fill}%
}%
\begin{pgfscope}%
\pgfsys@transformshift{2.266923in}{0.810724in}%
\pgfsys@useobject{currentmarker}{}%
\end{pgfscope}%
\end{pgfscope}%
\begin{pgfscope}%
\pgfpathrectangle{\pgfqpoint{0.466200in}{0.442000in}}{\pgfqpoint{2.763900in}{1.768000in}}%
\pgfusepath{clip}%
\pgfsetrectcap%
\pgfsetroundjoin%
\pgfsetlinewidth{1.505625pt}%
\definecolor{currentstroke}{rgb}{0.156863,0.356863,0.568627}%
\pgfsetstrokecolor{currentstroke}%
\pgfsetdash{}{0pt}%
\pgfpathmoveto{\pgfqpoint{3.104468in}{0.777896in}}%
\pgfusepath{stroke}%
\end{pgfscope}%
\begin{pgfscope}%
\pgfpathrectangle{\pgfqpoint{0.466200in}{0.442000in}}{\pgfqpoint{2.763900in}{1.768000in}}%
\pgfusepath{clip}%
\pgfsetbuttcap%
\pgfsetroundjoin%
\definecolor{currentfill}{rgb}{1.000000,1.000000,1.000000}%
\pgfsetfillcolor{currentfill}%
\pgfsetlinewidth{0.702625pt}%
\definecolor{currentstroke}{rgb}{0.156863,0.356863,0.568627}%
\pgfsetstrokecolor{currentstroke}%
\pgfsetdash{}{0pt}%
\pgfsys@defobject{currentmarker}{\pgfqpoint{-0.022222in}{-0.022222in}}{\pgfqpoint{0.022222in}{0.022222in}}{%
\pgfpathmoveto{\pgfqpoint{0.000000in}{-0.022222in}}%
\pgfpathcurveto{\pgfqpoint{0.005893in}{-0.022222in}}{\pgfqpoint{0.011546in}{-0.019881in}}{\pgfqpoint{0.015713in}{-0.015713in}}%
\pgfpathcurveto{\pgfqpoint{0.019881in}{-0.011546in}}{\pgfqpoint{0.022222in}{-0.005893in}}{\pgfqpoint{0.022222in}{0.000000in}}%
\pgfpathcurveto{\pgfqpoint{0.022222in}{0.005893in}}{\pgfqpoint{0.019881in}{0.011546in}}{\pgfqpoint{0.015713in}{0.015713in}}%
\pgfpathcurveto{\pgfqpoint{0.011546in}{0.019881in}}{\pgfqpoint{0.005893in}{0.022222in}}{\pgfqpoint{0.000000in}{0.022222in}}%
\pgfpathcurveto{\pgfqpoint{-0.005893in}{0.022222in}}{\pgfqpoint{-0.011546in}{0.019881in}}{\pgfqpoint{-0.015713in}{0.015713in}}%
\pgfpathcurveto{\pgfqpoint{-0.019881in}{0.011546in}}{\pgfqpoint{-0.022222in}{0.005893in}}{\pgfqpoint{-0.022222in}{0.000000in}}%
\pgfpathcurveto{\pgfqpoint{-0.022222in}{-0.005893in}}{\pgfqpoint{-0.019881in}{-0.011546in}}{\pgfqpoint{-0.015713in}{-0.015713in}}%
\pgfpathcurveto{\pgfqpoint{-0.011546in}{-0.019881in}}{\pgfqpoint{-0.005893in}{-0.022222in}}{\pgfqpoint{0.000000in}{-0.022222in}}%
\pgfpathlineto{\pgfqpoint{0.000000in}{-0.022222in}}%
\pgfpathclose%
\pgfusepath{stroke,fill}%
}%
\begin{pgfscope}%
\pgfsys@transformshift{3.104468in}{0.777896in}%
\pgfsys@useobject{currentmarker}{}%
\end{pgfscope}%
\end{pgfscope}%
\begin{pgfscope}%
\pgfpathrectangle{\pgfqpoint{0.466200in}{0.442000in}}{\pgfqpoint{2.763900in}{1.768000in}}%
\pgfusepath{clip}%
\pgfsetrectcap%
\pgfsetroundjoin%
\pgfsetlinewidth{1.003750pt}%
\definecolor{currentstroke}{rgb}{0.478431,0.478431,0.478431}%
\pgfsetstrokecolor{currentstroke}%
\pgfsetdash{}{0pt}%
\pgfpathmoveto{\pgfqpoint{0.591832in}{0.653491in}}%
\pgfpathlineto{\pgfqpoint{1.429377in}{0.574266in}}%
\pgfusepath{stroke}%
\end{pgfscope}%
\begin{pgfscope}%
\pgfpathrectangle{\pgfqpoint{0.466200in}{0.442000in}}{\pgfqpoint{2.763900in}{1.768000in}}%
\pgfusepath{clip}%
\pgfsetbuttcap%
\pgfsetmiterjoin%
\definecolor{currentfill}{rgb}{0.478431,0.478431,0.478431}%
\pgfsetfillcolor{currentfill}%
\pgfsetlinewidth{1.003750pt}%
\definecolor{currentstroke}{rgb}{0.478431,0.478431,0.478431}%
\pgfsetstrokecolor{currentstroke}%
\pgfsetdash{}{0pt}%
\pgfsys@defobject{currentmarker}{\pgfqpoint{-0.022222in}{-0.022222in}}{\pgfqpoint{0.022222in}{0.022222in}}{%
\pgfpathmoveto{\pgfqpoint{-0.022222in}{-0.022222in}}%
\pgfpathlineto{\pgfqpoint{0.022222in}{-0.022222in}}%
\pgfpathlineto{\pgfqpoint{0.022222in}{0.022222in}}%
\pgfpathlineto{\pgfqpoint{-0.022222in}{0.022222in}}%
\pgfpathlineto{\pgfqpoint{-0.022222in}{-0.022222in}}%
\pgfpathclose%
\pgfusepath{stroke,fill}%
}%
\begin{pgfscope}%
\pgfsys@transformshift{0.591832in}{0.653491in}%
\pgfsys@useobject{currentmarker}{}%
\end{pgfscope}%
\begin{pgfscope}%
\pgfsys@transformshift{1.429377in}{0.574266in}%
\pgfsys@useobject{currentmarker}{}%
\end{pgfscope}%
\end{pgfscope}%
\begin{pgfscope}%
\pgfpathrectangle{\pgfqpoint{0.466200in}{0.442000in}}{\pgfqpoint{2.763900in}{1.768000in}}%
\pgfusepath{clip}%
\pgfsetbuttcap%
\pgfsetroundjoin%
\pgfsetlinewidth{0.803000pt}%
\definecolor{currentstroke}{rgb}{0.478431,0.478431,0.478431}%
\pgfsetstrokecolor{currentstroke}%
\pgfsetdash{{2.000000pt}{1.200000pt}}{0.000000pt}%
\pgfpathmoveto{\pgfqpoint{0.591832in}{0.694064in}}%
\pgfpathlineto{\pgfqpoint{1.429377in}{0.705162in}}%
\pgfusepath{stroke}%
\end{pgfscope}%
\begin{pgfscope}%
\pgfpathrectangle{\pgfqpoint{0.466200in}{0.442000in}}{\pgfqpoint{2.763900in}{1.768000in}}%
\pgfusepath{clip}%
\pgfsetbuttcap%
\pgfsetmiterjoin%
\definecolor{currentfill}{rgb}{1.000000,1.000000,1.000000}%
\pgfsetfillcolor{currentfill}%
\pgfsetlinewidth{0.702625pt}%
\definecolor{currentstroke}{rgb}{0.478431,0.478431,0.478431}%
\pgfsetstrokecolor{currentstroke}%
\pgfsetdash{}{0pt}%
\pgfsys@defobject{currentmarker}{\pgfqpoint{-0.022222in}{-0.022222in}}{\pgfqpoint{0.022222in}{0.022222in}}{%
\pgfpathmoveto{\pgfqpoint{-0.022222in}{-0.022222in}}%
\pgfpathlineto{\pgfqpoint{0.022222in}{-0.022222in}}%
\pgfpathlineto{\pgfqpoint{0.022222in}{0.022222in}}%
\pgfpathlineto{\pgfqpoint{-0.022222in}{0.022222in}}%
\pgfpathlineto{\pgfqpoint{-0.022222in}{-0.022222in}}%
\pgfpathclose%
\pgfusepath{stroke,fill}%
}%
\begin{pgfscope}%
\pgfsys@transformshift{0.591832in}{0.694064in}%
\pgfsys@useobject{currentmarker}{}%
\end{pgfscope}%
\begin{pgfscope}%
\pgfsys@transformshift{1.429377in}{0.705162in}%
\pgfsys@useobject{currentmarker}{}%
\end{pgfscope}%
\end{pgfscope}%
\begin{pgfscope}%
\pgfpathrectangle{\pgfqpoint{0.466200in}{0.442000in}}{\pgfqpoint{2.763900in}{1.768000in}}%
\pgfusepath{clip}%
\pgfsetrectcap%
\pgfsetroundjoin%
\pgfsetlinewidth{1.003750pt}%
\definecolor{currentstroke}{rgb}{0.603922,0.603922,0.603922}%
\pgfsetstrokecolor{currentstroke}%
\pgfsetdash{}{0pt}%
\pgfpathmoveto{\pgfqpoint{0.591832in}{0.667420in}}%
\pgfpathlineto{\pgfqpoint{1.010605in}{0.709494in}}%
\pgfpathlineto{\pgfqpoint{1.429377in}{0.628930in}}%
\pgfusepath{stroke}%
\end{pgfscope}%
\begin{pgfscope}%
\pgfpathrectangle{\pgfqpoint{0.466200in}{0.442000in}}{\pgfqpoint{2.763900in}{1.768000in}}%
\pgfusepath{clip}%
\pgfsetbuttcap%
\pgfsetmiterjoin%
\definecolor{currentfill}{rgb}{0.603922,0.603922,0.603922}%
\pgfsetfillcolor{currentfill}%
\pgfsetlinewidth{1.003750pt}%
\definecolor{currentstroke}{rgb}{0.603922,0.603922,0.603922}%
\pgfsetstrokecolor{currentstroke}%
\pgfsetdash{}{0pt}%
\pgfsys@defobject{currentmarker}{\pgfqpoint{-0.022222in}{-0.022222in}}{\pgfqpoint{0.022222in}{0.022222in}}{%
\pgfpathmoveto{\pgfqpoint{0.000000in}{0.022222in}}%
\pgfpathlineto{\pgfqpoint{-0.022222in}{-0.022222in}}%
\pgfpathlineto{\pgfqpoint{0.022222in}{-0.022222in}}%
\pgfpathlineto{\pgfqpoint{0.000000in}{0.022222in}}%
\pgfpathclose%
\pgfusepath{stroke,fill}%
}%
\begin{pgfscope}%
\pgfsys@transformshift{0.591832in}{0.667420in}%
\pgfsys@useobject{currentmarker}{}%
\end{pgfscope}%
\begin{pgfscope}%
\pgfsys@transformshift{1.010605in}{0.709494in}%
\pgfsys@useobject{currentmarker}{}%
\end{pgfscope}%
\begin{pgfscope}%
\pgfsys@transformshift{1.429377in}{0.628930in}%
\pgfsys@useobject{currentmarker}{}%
\end{pgfscope}%
\end{pgfscope}%
\begin{pgfscope}%
\pgfpathrectangle{\pgfqpoint{0.466200in}{0.442000in}}{\pgfqpoint{2.763900in}{1.768000in}}%
\pgfusepath{clip}%
\pgfsetbuttcap%
\pgfsetroundjoin%
\pgfsetlinewidth{0.803000pt}%
\definecolor{currentstroke}{rgb}{0.603922,0.603922,0.603922}%
\pgfsetstrokecolor{currentstroke}%
\pgfsetdash{{2.000000pt}{1.200000pt}}{0.000000pt}%
\pgfpathmoveto{\pgfqpoint{0.591832in}{0.714583in}}%
\pgfpathlineto{\pgfqpoint{1.010605in}{0.935272in}}%
\pgfpathlineto{\pgfqpoint{1.429377in}{0.761544in}}%
\pgfusepath{stroke}%
\end{pgfscope}%
\begin{pgfscope}%
\pgfpathrectangle{\pgfqpoint{0.466200in}{0.442000in}}{\pgfqpoint{2.763900in}{1.768000in}}%
\pgfusepath{clip}%
\pgfsetbuttcap%
\pgfsetmiterjoin%
\definecolor{currentfill}{rgb}{1.000000,1.000000,1.000000}%
\pgfsetfillcolor{currentfill}%
\pgfsetlinewidth{0.702625pt}%
\definecolor{currentstroke}{rgb}{0.603922,0.603922,0.603922}%
\pgfsetstrokecolor{currentstroke}%
\pgfsetdash{}{0pt}%
\pgfsys@defobject{currentmarker}{\pgfqpoint{-0.022222in}{-0.022222in}}{\pgfqpoint{0.022222in}{0.022222in}}{%
\pgfpathmoveto{\pgfqpoint{0.000000in}{0.022222in}}%
\pgfpathlineto{\pgfqpoint{-0.022222in}{-0.022222in}}%
\pgfpathlineto{\pgfqpoint{0.022222in}{-0.022222in}}%
\pgfpathlineto{\pgfqpoint{0.000000in}{0.022222in}}%
\pgfpathclose%
\pgfusepath{stroke,fill}%
}%
\begin{pgfscope}%
\pgfsys@transformshift{0.591832in}{0.714583in}%
\pgfsys@useobject{currentmarker}{}%
\end{pgfscope}%
\begin{pgfscope}%
\pgfsys@transformshift{1.010605in}{0.935272in}%
\pgfsys@useobject{currentmarker}{}%
\end{pgfscope}%
\begin{pgfscope}%
\pgfsys@transformshift{1.429377in}{0.761544in}%
\pgfsys@useobject{currentmarker}{}%
\end{pgfscope}%
\end{pgfscope}%
\begin{pgfscope}%
\pgfpathrectangle{\pgfqpoint{0.466200in}{0.442000in}}{\pgfqpoint{2.763900in}{1.768000in}}%
\pgfusepath{clip}%
\pgfsetrectcap%
\pgfsetroundjoin%
\pgfsetlinewidth{1.003750pt}%
\definecolor{currentstroke}{rgb}{0.603922,0.603922,0.603922}%
\pgfsetstrokecolor{currentstroke}%
\pgfsetdash{}{0pt}%
\pgfpathmoveto{\pgfqpoint{0.591832in}{0.651890in}}%
\pgfpathlineto{\pgfqpoint{1.010605in}{0.771206in}}%
\pgfpathlineto{\pgfqpoint{1.429377in}{0.717826in}}%
\pgfusepath{stroke}%
\end{pgfscope}%
\begin{pgfscope}%
\pgfpathrectangle{\pgfqpoint{0.466200in}{0.442000in}}{\pgfqpoint{2.763900in}{1.768000in}}%
\pgfusepath{clip}%
\pgfsetbuttcap%
\pgfsetmiterjoin%
\definecolor{currentfill}{rgb}{0.603922,0.603922,0.603922}%
\pgfsetfillcolor{currentfill}%
\pgfsetlinewidth{1.003750pt}%
\definecolor{currentstroke}{rgb}{0.603922,0.603922,0.603922}%
\pgfsetstrokecolor{currentstroke}%
\pgfsetdash{}{0pt}%
\pgfsys@defobject{currentmarker}{\pgfqpoint{-0.031427in}{-0.031427in}}{\pgfqpoint{0.031427in}{0.031427in}}{%
\pgfpathmoveto{\pgfqpoint{-0.000000in}{-0.031427in}}%
\pgfpathlineto{\pgfqpoint{0.031427in}{0.000000in}}%
\pgfpathlineto{\pgfqpoint{0.000000in}{0.031427in}}%
\pgfpathlineto{\pgfqpoint{-0.031427in}{0.000000in}}%
\pgfpathlineto{\pgfqpoint{-0.000000in}{-0.031427in}}%
\pgfpathclose%
\pgfusepath{stroke,fill}%
}%
\begin{pgfscope}%
\pgfsys@transformshift{0.591832in}{0.651890in}%
\pgfsys@useobject{currentmarker}{}%
\end{pgfscope}%
\begin{pgfscope}%
\pgfsys@transformshift{1.010605in}{0.771206in}%
\pgfsys@useobject{currentmarker}{}%
\end{pgfscope}%
\begin{pgfscope}%
\pgfsys@transformshift{1.429377in}{0.717826in}%
\pgfsys@useobject{currentmarker}{}%
\end{pgfscope}%
\end{pgfscope}%
\begin{pgfscope}%
\pgfpathrectangle{\pgfqpoint{0.466200in}{0.442000in}}{\pgfqpoint{2.763900in}{1.768000in}}%
\pgfusepath{clip}%
\pgfsetbuttcap%
\pgfsetroundjoin%
\pgfsetlinewidth{0.803000pt}%
\definecolor{currentstroke}{rgb}{0.603922,0.603922,0.603922}%
\pgfsetstrokecolor{currentstroke}%
\pgfsetdash{{2.000000pt}{1.200000pt}}{0.000000pt}%
\pgfpathmoveto{\pgfqpoint{0.591832in}{0.687513in}}%
\pgfpathlineto{\pgfqpoint{1.010605in}{0.824680in}}%
\pgfpathlineto{\pgfqpoint{1.429377in}{0.956849in}}%
\pgfusepath{stroke}%
\end{pgfscope}%
\begin{pgfscope}%
\pgfpathrectangle{\pgfqpoint{0.466200in}{0.442000in}}{\pgfqpoint{2.763900in}{1.768000in}}%
\pgfusepath{clip}%
\pgfsetbuttcap%
\pgfsetmiterjoin%
\definecolor{currentfill}{rgb}{1.000000,1.000000,1.000000}%
\pgfsetfillcolor{currentfill}%
\pgfsetlinewidth{0.702625pt}%
\definecolor{currentstroke}{rgb}{0.603922,0.603922,0.603922}%
\pgfsetstrokecolor{currentstroke}%
\pgfsetdash{}{0pt}%
\pgfsys@defobject{currentmarker}{\pgfqpoint{-0.031427in}{-0.031427in}}{\pgfqpoint{0.031427in}{0.031427in}}{%
\pgfpathmoveto{\pgfqpoint{-0.000000in}{-0.031427in}}%
\pgfpathlineto{\pgfqpoint{0.031427in}{0.000000in}}%
\pgfpathlineto{\pgfqpoint{0.000000in}{0.031427in}}%
\pgfpathlineto{\pgfqpoint{-0.031427in}{0.000000in}}%
\pgfpathlineto{\pgfqpoint{-0.000000in}{-0.031427in}}%
\pgfpathclose%
\pgfusepath{stroke,fill}%
}%
\begin{pgfscope}%
\pgfsys@transformshift{0.591832in}{0.687513in}%
\pgfsys@useobject{currentmarker}{}%
\end{pgfscope}%
\begin{pgfscope}%
\pgfsys@transformshift{1.010605in}{0.824680in}%
\pgfsys@useobject{currentmarker}{}%
\end{pgfscope}%
\begin{pgfscope}%
\pgfsys@transformshift{1.429377in}{0.956849in}%
\pgfsys@useobject{currentmarker}{}%
\end{pgfscope}%
\end{pgfscope}%
\begin{pgfscope}%
\pgfpathrectangle{\pgfqpoint{0.466200in}{0.442000in}}{\pgfqpoint{2.763900in}{1.768000in}}%
\pgfusepath{clip}%
\pgfsetrectcap%
\pgfsetroundjoin%
\pgfsetlinewidth{1.003750pt}%
\definecolor{currentstroke}{rgb}{0.603922,0.603922,0.603922}%
\pgfsetstrokecolor{currentstroke}%
\pgfsetdash{}{0pt}%
\pgfpathmoveto{\pgfqpoint{0.591832in}{0.657242in}}%
\pgfpathlineto{\pgfqpoint{1.429377in}{0.819731in}}%
\pgfusepath{stroke}%
\end{pgfscope}%
\begin{pgfscope}%
\pgfpathrectangle{\pgfqpoint{0.466200in}{0.442000in}}{\pgfqpoint{2.763900in}{1.768000in}}%
\pgfusepath{clip}%
\pgfsetbuttcap%
\pgfsetmiterjoin%
\definecolor{currentfill}{rgb}{0.603922,0.603922,0.603922}%
\pgfsetfillcolor{currentfill}%
\pgfsetlinewidth{1.003750pt}%
\definecolor{currentstroke}{rgb}{0.603922,0.603922,0.603922}%
\pgfsetstrokecolor{currentstroke}%
\pgfsetdash{}{0pt}%
\pgfsys@defobject{currentmarker}{\pgfqpoint{-0.022222in}{-0.022222in}}{\pgfqpoint{0.022222in}{0.022222in}}{%
\pgfpathmoveto{\pgfqpoint{-0.000000in}{-0.022222in}}%
\pgfpathlineto{\pgfqpoint{0.022222in}{0.022222in}}%
\pgfpathlineto{\pgfqpoint{-0.022222in}{0.022222in}}%
\pgfpathlineto{\pgfqpoint{-0.000000in}{-0.022222in}}%
\pgfpathclose%
\pgfusepath{stroke,fill}%
}%
\begin{pgfscope}%
\pgfsys@transformshift{0.591832in}{0.657242in}%
\pgfsys@useobject{currentmarker}{}%
\end{pgfscope}%
\begin{pgfscope}%
\pgfsys@transformshift{1.429377in}{0.819731in}%
\pgfsys@useobject{currentmarker}{}%
\end{pgfscope}%
\end{pgfscope}%
\begin{pgfscope}%
\pgfpathrectangle{\pgfqpoint{0.466200in}{0.442000in}}{\pgfqpoint{2.763900in}{1.768000in}}%
\pgfusepath{clip}%
\pgfsetbuttcap%
\pgfsetroundjoin%
\pgfsetlinewidth{0.803000pt}%
\definecolor{currentstroke}{rgb}{0.603922,0.603922,0.603922}%
\pgfsetstrokecolor{currentstroke}%
\pgfsetdash{{2.000000pt}{1.200000pt}}{0.000000pt}%
\pgfpathmoveto{\pgfqpoint{0.591832in}{0.699008in}}%
\pgfpathlineto{\pgfqpoint{1.429377in}{0.977560in}}%
\pgfusepath{stroke}%
\end{pgfscope}%
\begin{pgfscope}%
\pgfpathrectangle{\pgfqpoint{0.466200in}{0.442000in}}{\pgfqpoint{2.763900in}{1.768000in}}%
\pgfusepath{clip}%
\pgfsetbuttcap%
\pgfsetmiterjoin%
\definecolor{currentfill}{rgb}{1.000000,1.000000,1.000000}%
\pgfsetfillcolor{currentfill}%
\pgfsetlinewidth{0.702625pt}%
\definecolor{currentstroke}{rgb}{0.603922,0.603922,0.603922}%
\pgfsetstrokecolor{currentstroke}%
\pgfsetdash{}{0pt}%
\pgfsys@defobject{currentmarker}{\pgfqpoint{-0.022222in}{-0.022222in}}{\pgfqpoint{0.022222in}{0.022222in}}{%
\pgfpathmoveto{\pgfqpoint{-0.000000in}{-0.022222in}}%
\pgfpathlineto{\pgfqpoint{0.022222in}{0.022222in}}%
\pgfpathlineto{\pgfqpoint{-0.022222in}{0.022222in}}%
\pgfpathlineto{\pgfqpoint{-0.000000in}{-0.022222in}}%
\pgfpathclose%
\pgfusepath{stroke,fill}%
}%
\begin{pgfscope}%
\pgfsys@transformshift{0.591832in}{0.699008in}%
\pgfsys@useobject{currentmarker}{}%
\end{pgfscope}%
\begin{pgfscope}%
\pgfsys@transformshift{1.429377in}{0.977560in}%
\pgfsys@useobject{currentmarker}{}%
\end{pgfscope}%
\end{pgfscope}%
\begin{pgfscope}%
\pgfsetrectcap%
\pgfsetmiterjoin%
\pgfsetlinewidth{0.803000pt}%
\definecolor{currentstroke}{rgb}{0.101961,0.101961,0.101961}%
\pgfsetstrokecolor{currentstroke}%
\pgfsetdash{}{0pt}%
\pgfpathmoveto{\pgfqpoint{0.466200in}{0.442000in}}%
\pgfpathlineto{\pgfqpoint{0.466200in}{2.210000in}}%
\pgfusepath{stroke}%
\end{pgfscope}%
\begin{pgfscope}%
\pgfsetrectcap%
\pgfsetmiterjoin%
\pgfsetlinewidth{0.803000pt}%
\definecolor{currentstroke}{rgb}{0.101961,0.101961,0.101961}%
\pgfsetstrokecolor{currentstroke}%
\pgfsetdash{}{0pt}%
\pgfpathmoveto{\pgfqpoint{0.466200in}{0.442000in}}%
\pgfpathlineto{\pgfqpoint{3.230100in}{0.442000in}}%
\pgfusepath{stroke}%
\end{pgfscope}%
\begin{pgfscope}%
\pgfsetrectcap%
\pgfsetroundjoin%
\pgfsetlinewidth{1.003750pt}%
\definecolor{currentstroke}{rgb}{0.156863,0.356863,0.568627}%
\pgfsetstrokecolor{currentstroke}%
\pgfsetdash{}{0pt}%
\pgfpathmoveto{\pgfqpoint{0.511339in}{2.133264in}}%
\pgfpathlineto{\pgfqpoint{0.592589in}{2.133264in}}%
\pgfpathlineto{\pgfqpoint{0.673839in}{2.133264in}}%
\pgfusepath{stroke}%
\end{pgfscope}%
\begin{pgfscope}%
\pgfsetbuttcap%
\pgfsetroundjoin%
\definecolor{currentfill}{rgb}{0.156863,0.356863,0.568627}%
\pgfsetfillcolor{currentfill}%
\pgfsetlinewidth{1.003750pt}%
\definecolor{currentstroke}{rgb}{0.156863,0.356863,0.568627}%
\pgfsetstrokecolor{currentstroke}%
\pgfsetdash{}{0pt}%
\pgfsys@defobject{currentmarker}{\pgfqpoint{-0.022222in}{-0.022222in}}{\pgfqpoint{0.022222in}{0.022222in}}{%
\pgfpathmoveto{\pgfqpoint{0.000000in}{-0.022222in}}%
\pgfpathcurveto{\pgfqpoint{0.005893in}{-0.022222in}}{\pgfqpoint{0.011546in}{-0.019881in}}{\pgfqpoint{0.015713in}{-0.015713in}}%
\pgfpathcurveto{\pgfqpoint{0.019881in}{-0.011546in}}{\pgfqpoint{0.022222in}{-0.005893in}}{\pgfqpoint{0.022222in}{0.000000in}}%
\pgfpathcurveto{\pgfqpoint{0.022222in}{0.005893in}}{\pgfqpoint{0.019881in}{0.011546in}}{\pgfqpoint{0.015713in}{0.015713in}}%
\pgfpathcurveto{\pgfqpoint{0.011546in}{0.019881in}}{\pgfqpoint{0.005893in}{0.022222in}}{\pgfqpoint{0.000000in}{0.022222in}}%
\pgfpathcurveto{\pgfqpoint{-0.005893in}{0.022222in}}{\pgfqpoint{-0.011546in}{0.019881in}}{\pgfqpoint{-0.015713in}{0.015713in}}%
\pgfpathcurveto{\pgfqpoint{-0.019881in}{0.011546in}}{\pgfqpoint{-0.022222in}{0.005893in}}{\pgfqpoint{-0.022222in}{0.000000in}}%
\pgfpathcurveto{\pgfqpoint{-0.022222in}{-0.005893in}}{\pgfqpoint{-0.019881in}{-0.011546in}}{\pgfqpoint{-0.015713in}{-0.015713in}}%
\pgfpathcurveto{\pgfqpoint{-0.011546in}{-0.019881in}}{\pgfqpoint{-0.005893in}{-0.022222in}}{\pgfqpoint{0.000000in}{-0.022222in}}%
\pgfpathlineto{\pgfqpoint{0.000000in}{-0.022222in}}%
\pgfpathclose%
\pgfusepath{stroke,fill}%
}%
\begin{pgfscope}%
\pgfsys@transformshift{0.592589in}{2.133264in}%
\pgfsys@useobject{currentmarker}{}%
\end{pgfscope}%
\end{pgfscope}%
\begin{pgfscope}%
\definecolor{textcolor}{rgb}{0.101961,0.101961,0.101961}%
\pgfsetstrokecolor{textcolor}%
\pgfsetfillcolor{textcolor}%
\pgftext[x=0.746061in,y=2.101667in,left,base]{\color{textcolor}\sffamily\fontsize{6.500000}{7.800000}\selectfont COVEX (old: rung 3)}%
\end{pgfscope}%
\begin{pgfscope}%
\pgfsetrectcap%
\pgfsetroundjoin%
\pgfsetlinewidth{1.003750pt}%
\definecolor{currentstroke}{rgb}{0.478431,0.478431,0.478431}%
\pgfsetstrokecolor{currentstroke}%
\pgfsetdash{}{0pt}%
\pgfpathmoveto{\pgfqpoint{0.511339in}{2.026667in}}%
\pgfpathlineto{\pgfqpoint{0.592589in}{2.026667in}}%
\pgfpathlineto{\pgfqpoint{0.673839in}{2.026667in}}%
\pgfusepath{stroke}%
\end{pgfscope}%
\begin{pgfscope}%
\pgfsetbuttcap%
\pgfsetmiterjoin%
\definecolor{currentfill}{rgb}{0.478431,0.478431,0.478431}%
\pgfsetfillcolor{currentfill}%
\pgfsetlinewidth{1.003750pt}%
\definecolor{currentstroke}{rgb}{0.478431,0.478431,0.478431}%
\pgfsetstrokecolor{currentstroke}%
\pgfsetdash{}{0pt}%
\pgfsys@defobject{currentmarker}{\pgfqpoint{-0.022222in}{-0.022222in}}{\pgfqpoint{0.022222in}{0.022222in}}{%
\pgfpathmoveto{\pgfqpoint{-0.022222in}{-0.022222in}}%
\pgfpathlineto{\pgfqpoint{0.022222in}{-0.022222in}}%
\pgfpathlineto{\pgfqpoint{0.022222in}{0.022222in}}%
\pgfpathlineto{\pgfqpoint{-0.022222in}{0.022222in}}%
\pgfpathlineto{\pgfqpoint{-0.022222in}{-0.022222in}}%
\pgfpathclose%
\pgfusepath{stroke,fill}%
}%
\begin{pgfscope}%
\pgfsys@transformshift{0.592589in}{2.026667in}%
\pgfsys@useobject{currentmarker}{}%
\end{pgfscope}%
\end{pgfscope}%
\begin{pgfscope}%
\definecolor{textcolor}{rgb}{0.101961,0.101961,0.101961}%
\pgfsetstrokecolor{textcolor}%
\pgfsetfillcolor{textcolor}%
\pgftext[x=0.746061in,y=1.995069in,left,base]{\color{textcolor}\sffamily\fontsize{6.500000}{7.800000}\selectfont Mooncake}%
\end{pgfscope}%
\begin{pgfscope}%
\pgfsetrectcap%
\pgfsetroundjoin%
\pgfsetlinewidth{1.003750pt}%
\definecolor{currentstroke}{rgb}{0.603922,0.603922,0.603922}%
\pgfsetstrokecolor{currentstroke}%
\pgfsetdash{}{0pt}%
\pgfpathmoveto{\pgfqpoint{0.511339in}{1.924699in}}%
\pgfpathlineto{\pgfqpoint{0.592589in}{1.924699in}}%
\pgfpathlineto{\pgfqpoint{0.673839in}{1.924699in}}%
\pgfusepath{stroke}%
\end{pgfscope}%
\begin{pgfscope}%
\pgfsetbuttcap%
\pgfsetmiterjoin%
\definecolor{currentfill}{rgb}{0.603922,0.603922,0.603922}%
\pgfsetfillcolor{currentfill}%
\pgfsetlinewidth{1.003750pt}%
\definecolor{currentstroke}{rgb}{0.603922,0.603922,0.603922}%
\pgfsetstrokecolor{currentstroke}%
\pgfsetdash{}{0pt}%
\pgfsys@defobject{currentmarker}{\pgfqpoint{-0.022222in}{-0.022222in}}{\pgfqpoint{0.022222in}{0.022222in}}{%
\pgfpathmoveto{\pgfqpoint{0.000000in}{0.022222in}}%
\pgfpathlineto{\pgfqpoint{-0.022222in}{-0.022222in}}%
\pgfpathlineto{\pgfqpoint{0.022222in}{-0.022222in}}%
\pgfpathlineto{\pgfqpoint{0.000000in}{0.022222in}}%
\pgfpathclose%
\pgfusepath{stroke,fill}%
}%
\begin{pgfscope}%
\pgfsys@transformshift{0.592589in}{1.924699in}%
\pgfsys@useobject{currentmarker}{}%
\end{pgfscope}%
\end{pgfscope}%
\begin{pgfscope}%
\definecolor{textcolor}{rgb}{0.101961,0.101961,0.101961}%
\pgfsetstrokecolor{textcolor}%
\pgfsetfillcolor{textcolor}%
\pgftext[x=0.746061in,y=1.893102in,left,base]{\color{textcolor}\sffamily\fontsize{6.500000}{7.800000}\selectfont LLF}%
\end{pgfscope}%
\begin{pgfscope}%
\pgfsetrectcap%
\pgfsetroundjoin%
\pgfsetlinewidth{1.003750pt}%
\definecolor{currentstroke}{rgb}{0.603922,0.603922,0.603922}%
\pgfsetstrokecolor{currentstroke}%
\pgfsetdash{}{0pt}%
\pgfpathmoveto{\pgfqpoint{0.511339in}{1.822732in}}%
\pgfpathlineto{\pgfqpoint{0.592589in}{1.822732in}}%
\pgfpathlineto{\pgfqpoint{0.673839in}{1.822732in}}%
\pgfusepath{stroke}%
\end{pgfscope}%
\begin{pgfscope}%
\pgfsetbuttcap%
\pgfsetmiterjoin%
\definecolor{currentfill}{rgb}{0.603922,0.603922,0.603922}%
\pgfsetfillcolor{currentfill}%
\pgfsetlinewidth{1.003750pt}%
\definecolor{currentstroke}{rgb}{0.603922,0.603922,0.603922}%
\pgfsetstrokecolor{currentstroke}%
\pgfsetdash{}{0pt}%
\pgfsys@defobject{currentmarker}{\pgfqpoint{-0.031427in}{-0.031427in}}{\pgfqpoint{0.031427in}{0.031427in}}{%
\pgfpathmoveto{\pgfqpoint{-0.000000in}{-0.031427in}}%
\pgfpathlineto{\pgfqpoint{0.031427in}{0.000000in}}%
\pgfpathlineto{\pgfqpoint{0.000000in}{0.031427in}}%
\pgfpathlineto{\pgfqpoint{-0.031427in}{0.000000in}}%
\pgfpathlineto{\pgfqpoint{-0.000000in}{-0.031427in}}%
\pgfpathclose%
\pgfusepath{stroke,fill}%
}%
\begin{pgfscope}%
\pgfsys@transformshift{0.592589in}{1.822732in}%
\pgfsys@useobject{currentmarker}{}%
\end{pgfscope}%
\end{pgfscope}%
\begin{pgfscope}%
\definecolor{textcolor}{rgb}{0.101961,0.101961,0.101961}%
\pgfsetstrokecolor{textcolor}%
\pgfsetfillcolor{textcolor}%
\pgftext[x=0.746061in,y=1.791134in,left,base]{\color{textcolor}\sffamily\fontsize{6.500000}{7.800000}\selectfont SimpleStream}%
\end{pgfscope}%
\begin{pgfscope}%
\pgfsetrectcap%
\pgfsetroundjoin%
\pgfsetlinewidth{1.003750pt}%
\definecolor{currentstroke}{rgb}{0.603922,0.603922,0.603922}%
\pgfsetstrokecolor{currentstroke}%
\pgfsetdash{}{0pt}%
\pgfpathmoveto{\pgfqpoint{1.631141in}{2.133264in}}%
\pgfpathlineto{\pgfqpoint{1.712391in}{2.133264in}}%
\pgfpathlineto{\pgfqpoint{1.793641in}{2.133264in}}%
\pgfusepath{stroke}%
\end{pgfscope}%
\begin{pgfscope}%
\pgfsetbuttcap%
\pgfsetmiterjoin%
\definecolor{currentfill}{rgb}{0.603922,0.603922,0.603922}%
\pgfsetfillcolor{currentfill}%
\pgfsetlinewidth{1.003750pt}%
\definecolor{currentstroke}{rgb}{0.603922,0.603922,0.603922}%
\pgfsetstrokecolor{currentstroke}%
\pgfsetdash{}{0pt}%
\pgfsys@defobject{currentmarker}{\pgfqpoint{-0.022222in}{-0.022222in}}{\pgfqpoint{0.022222in}{0.022222in}}{%
\pgfpathmoveto{\pgfqpoint{-0.000000in}{-0.022222in}}%
\pgfpathlineto{\pgfqpoint{0.022222in}{0.022222in}}%
\pgfpathlineto{\pgfqpoint{-0.022222in}{0.022222in}}%
\pgfpathlineto{\pgfqpoint{-0.000000in}{-0.022222in}}%
\pgfpathclose%
\pgfusepath{stroke,fill}%
}%
\begin{pgfscope}%
\pgfsys@transformshift{1.712391in}{2.133264in}%
\pgfsys@useobject{currentmarker}{}%
\end{pgfscope}%
\end{pgfscope}%
\begin{pgfscope}%
\definecolor{textcolor}{rgb}{0.101961,0.101961,0.101961}%
\pgfsetstrokecolor{textcolor}%
\pgfsetfillcolor{textcolor}%
\pgftext[x=1.865863in,y=2.101667in,left,base]{\color{textcolor}\sffamily\fontsize{6.500000}{7.800000}\selectfont ECO-LLM}%
\end{pgfscope}%
\begin{pgfscope}%
\pgfsetrectcap%
\pgfsetroundjoin%
\pgfsetlinewidth{1.003750pt}%
\definecolor{currentstroke}{rgb}{0.333333,0.333333,0.333333}%
\pgfsetstrokecolor{currentstroke}%
\pgfsetdash{}{0pt}%
\pgfpathmoveto{\pgfqpoint{1.631141in}{2.031296in}}%
\pgfpathlineto{\pgfqpoint{1.712391in}{2.031296in}}%
\pgfpathlineto{\pgfqpoint{1.793641in}{2.031296in}}%
\pgfusepath{stroke}%
\end{pgfscope}%
\begin{pgfscope}%
\definecolor{textcolor}{rgb}{0.101961,0.101961,0.101961}%
\pgfsetstrokecolor{textcolor}%
\pgfsetfillcolor{textcolor}%
\pgftext[x=1.865863in,y=1.999699in,left,base]{\color{textcolor}\sffamily\fontsize{6.500000}{7.800000}\selectfont p50}%
\end{pgfscope}%
\begin{pgfscope}%
\pgfsetbuttcap%
\pgfsetroundjoin%
\pgfsetlinewidth{0.803000pt}%
\definecolor{currentstroke}{rgb}{0.333333,0.333333,0.333333}%
\pgfsetstrokecolor{currentstroke}%
\pgfsetdash{{2.000000pt}{1.200000pt}}{0.000000pt}%
\pgfpathmoveto{\pgfqpoint{1.631141in}{1.929329in}}%
\pgfpathlineto{\pgfqpoint{1.712391in}{1.929329in}}%
\pgfpathlineto{\pgfqpoint{1.793641in}{1.929329in}}%
\pgfusepath{stroke}%
\end{pgfscope}%
\begin{pgfscope}%
\definecolor{textcolor}{rgb}{0.101961,0.101961,0.101961}%
\pgfsetstrokecolor{textcolor}%
\pgfsetfillcolor{textcolor}%
\pgftext[x=1.865863in,y=1.897732in,left,base]{\color{textcolor}\sffamily\fontsize{6.500000}{7.800000}\selectfont p99}%
\end{pgfscope}%
\begin{pgfscope}%
\definecolor{textcolor}{rgb}{0.549020,0.549020,0.549020}%
\pgfsetstrokecolor{textcolor}%
\pgfsetfillcolor{textcolor}%
\pgftext[x=3.296700in,y=2.561000in,right,top]{\color{textcolor}\sffamily\fontsize{5.500000}{6.600000}\selectfont OLD DATA}%
\end{pgfscope}%
\end{pgfpicture}%
\makeatother%
\endgroup%
}
\caption{[PLACEHOLDER: OLD DATA, COVEX is the rung 3 covex-v6-fair-final stand-in (OLD); DEV candidates are not in F14.csv; dry latency at 16/24/32/48 devices is not committed] Decision latency of the policy per decision versus fleet size, live runs, 1x memory, 3 s robot deadline. Clock plotted: wall clock (decision\_wall\_ms), not CPU time. Solid lines are p50, dashed lines p99; each point is the median over the available seeds of the per-run percentile. COVEX points above 16 devices (hollow) are single seed 4 runs with the 3:11:2 mix flag, the only committed larger fleets; baselines have no cells above 16 devices. Each run logs about 12,000 decisions.}
\label{fig:F14_decision_latency}
\end{figure}
\clearpage
\begin{figure}[t]
\centering
{\fontfamily{LinuxBiolinumT}\selectfont
%% Creator: Matplotlib, PGF backend
%%
%% To include the figure in your LaTeX document, write
%%   \input{<filename>.pgf}
%%
%% Make sure the required packages are loaded in your preamble
%%   \usepackage{pgf}
%%
%% Also ensure that all the required font packages are loaded; for instance,
%% the lmodern package is sometimes necessary when using math font.
%%   \usepackage{lmodern}
%%
%% Figures using additional raster images can only be included by \input if
%% they are in the same directory as the main LaTeX file. For loading figures
%% from other directories you can use the `import` package
%%   \usepackage{import}
%%
%% and then include the figures with
%%   \import{<path to file>}{<filename>.pgf}
%%
%% Matplotlib used the following preamble
%%
\begingroup%
\makeatletter%
\begin{pgfpicture}%
\pgfpathrectangle{\pgfpointorigin}{\pgfqpoint{3.330000in}{2.500000in}}%
\pgfusepath{use as bounding box, clip}%
\begin{pgfscope}%
\pgfsetbuttcap%
\pgfsetmiterjoin%
\definecolor{currentfill}{rgb}{1.000000,1.000000,1.000000}%
\pgfsetfillcolor{currentfill}%
\pgfsetlinewidth{0.000000pt}%
\definecolor{currentstroke}{rgb}{1.000000,1.000000,1.000000}%
\pgfsetstrokecolor{currentstroke}%
\pgfsetdash{}{0pt}%
\pgfpathmoveto{\pgfqpoint{0.000000in}{0.000000in}}%
\pgfpathlineto{\pgfqpoint{3.330000in}{0.000000in}}%
\pgfpathlineto{\pgfqpoint{3.330000in}{2.500000in}}%
\pgfpathlineto{\pgfqpoint{0.000000in}{2.500000in}}%
\pgfpathlineto{\pgfqpoint{0.000000in}{0.000000in}}%
\pgfpathclose%
\pgfusepath{fill}%
\end{pgfscope}%
\begin{pgfscope}%
\pgfsetbuttcap%
\pgfsetmiterjoin%
\definecolor{currentfill}{rgb}{1.000000,1.000000,1.000000}%
\pgfsetfillcolor{currentfill}%
\pgfsetlinewidth{0.000000pt}%
\definecolor{currentstroke}{rgb}{0.000000,0.000000,0.000000}%
\pgfsetstrokecolor{currentstroke}%
\pgfsetstrokeopacity{0.000000}%
\pgfsetdash{}{0pt}%
\pgfpathmoveto{\pgfqpoint{0.466200in}{0.450000in}}%
\pgfpathlineto{\pgfqpoint{3.230100in}{0.450000in}}%
\pgfpathlineto{\pgfqpoint{3.230100in}{2.350000in}}%
\pgfpathlineto{\pgfqpoint{0.466200in}{2.350000in}}%
\pgfpathlineto{\pgfqpoint{0.466200in}{0.450000in}}%
\pgfpathclose%
\pgfusepath{fill}%
\end{pgfscope}%
\begin{pgfscope}%
\definecolor{textcolor}{rgb}{0.101961,0.101961,0.101961}%
\pgfsetstrokecolor{textcolor}%
\pgfsetfillcolor{textcolor}%
\pgftext[x=0.591832in,y=0.401389in,,top]{\color{textcolor}\sffamily\fontsize{7.000000}{8.400000}\selectfont 0.5}%
\end{pgfscope}%
\begin{pgfscope}%
\definecolor{textcolor}{rgb}{0.101961,0.101961,0.101961}%
\pgfsetstrokecolor{textcolor}%
\pgfsetfillcolor{textcolor}%
\pgftext[x=1.094359in,y=0.401389in,,top]{\color{textcolor}\sffamily\fontsize{7.000000}{8.400000}\selectfont 1.0}%
\end{pgfscope}%
\begin{pgfscope}%
\definecolor{textcolor}{rgb}{0.101961,0.101961,0.101961}%
\pgfsetstrokecolor{textcolor}%
\pgfsetfillcolor{textcolor}%
\pgftext[x=1.596886in,y=0.401389in,,top]{\color{textcolor}\sffamily\fontsize{7.000000}{8.400000}\selectfont 1.5}%
\end{pgfscope}%
\begin{pgfscope}%
\definecolor{textcolor}{rgb}{0.101961,0.101961,0.101961}%
\pgfsetstrokecolor{textcolor}%
\pgfsetfillcolor{textcolor}%
\pgftext[x=2.099414in,y=0.401389in,,top]{\color{textcolor}\sffamily\fontsize{7.000000}{8.400000}\selectfont 2.0}%
\end{pgfscope}%
\begin{pgfscope}%
\definecolor{textcolor}{rgb}{0.101961,0.101961,0.101961}%
\pgfsetstrokecolor{textcolor}%
\pgfsetfillcolor{textcolor}%
\pgftext[x=2.601941in,y=0.401389in,,top]{\color{textcolor}\sffamily\fontsize{7.000000}{8.400000}\selectfont 2.5}%
\end{pgfscope}%
\begin{pgfscope}%
\definecolor{textcolor}{rgb}{0.101961,0.101961,0.101961}%
\pgfsetstrokecolor{textcolor}%
\pgfsetfillcolor{textcolor}%
\pgftext[x=3.104468in,y=0.401389in,,top]{\color{textcolor}\sffamily\fontsize{7.000000}{8.400000}\selectfont 3.0}%
\end{pgfscope}%
\begin{pgfscope}%
\definecolor{textcolor}{rgb}{0.000000,0.000000,0.000000}%
\pgfsetstrokecolor{textcolor}%
\pgfsetfillcolor{textcolor}%
\pgftext[x=1.848150in,y=0.259414in,,top]{\color{textcolor}\sffamily\fontsize{7.500000}{9.000000}\selectfont time to event (s)}%
\end{pgfscope}%
\begin{pgfscope}%
\pgfpathrectangle{\pgfqpoint{0.466200in}{0.450000in}}{\pgfqpoint{2.763900in}{1.900000in}}%
\pgfusepath{clip}%
\pgfsetrectcap%
\pgfsetroundjoin%
\pgfsetlinewidth{0.401500pt}%
\definecolor{currentstroke}{rgb}{0.901961,0.901961,0.901961}%
\pgfsetstrokecolor{currentstroke}%
\pgfsetdash{}{0pt}%
\pgfpathmoveto{\pgfqpoint{0.466200in}{0.450000in}}%
\pgfpathlineto{\pgfqpoint{3.230100in}{0.450000in}}%
\pgfusepath{stroke}%
\end{pgfscope}%
\begin{pgfscope}%
\definecolor{textcolor}{rgb}{0.101961,0.101961,0.101961}%
\pgfsetstrokecolor{textcolor}%
\pgfsetfillcolor{textcolor}%
\pgftext[x=0.285595in, y=0.416242in, left, base]{\color{textcolor}\sffamily\fontsize{7.000000}{8.400000}\selectfont 0.0}%
\end{pgfscope}%
\begin{pgfscope}%
\pgfpathrectangle{\pgfqpoint{0.466200in}{0.450000in}}{\pgfqpoint{2.763900in}{1.900000in}}%
\pgfusepath{clip}%
\pgfsetrectcap%
\pgfsetroundjoin%
\pgfsetlinewidth{0.401500pt}%
\definecolor{currentstroke}{rgb}{0.901961,0.901961,0.901961}%
\pgfsetstrokecolor{currentstroke}%
\pgfsetdash{}{0pt}%
\pgfpathmoveto{\pgfqpoint{0.466200in}{0.661111in}}%
\pgfpathlineto{\pgfqpoint{3.230100in}{0.661111in}}%
\pgfusepath{stroke}%
\end{pgfscope}%
\begin{pgfscope}%
\definecolor{textcolor}{rgb}{0.101961,0.101961,0.101961}%
\pgfsetstrokecolor{textcolor}%
\pgfsetfillcolor{textcolor}%
\pgftext[x=0.285595in, y=0.627353in, left, base]{\color{textcolor}\sffamily\fontsize{7.000000}{8.400000}\selectfont 0.1}%
\end{pgfscope}%
\begin{pgfscope}%
\pgfpathrectangle{\pgfqpoint{0.466200in}{0.450000in}}{\pgfqpoint{2.763900in}{1.900000in}}%
\pgfusepath{clip}%
\pgfsetrectcap%
\pgfsetroundjoin%
\pgfsetlinewidth{0.401500pt}%
\definecolor{currentstroke}{rgb}{0.901961,0.901961,0.901961}%
\pgfsetstrokecolor{currentstroke}%
\pgfsetdash{}{0pt}%
\pgfpathmoveto{\pgfqpoint{0.466200in}{0.872222in}}%
\pgfpathlineto{\pgfqpoint{3.230100in}{0.872222in}}%
\pgfusepath{stroke}%
\end{pgfscope}%
\begin{pgfscope}%
\definecolor{textcolor}{rgb}{0.101961,0.101961,0.101961}%
\pgfsetstrokecolor{textcolor}%
\pgfsetfillcolor{textcolor}%
\pgftext[x=0.285595in, y=0.838465in, left, base]{\color{textcolor}\sffamily\fontsize{7.000000}{8.400000}\selectfont 0.2}%
\end{pgfscope}%
\begin{pgfscope}%
\pgfpathrectangle{\pgfqpoint{0.466200in}{0.450000in}}{\pgfqpoint{2.763900in}{1.900000in}}%
\pgfusepath{clip}%
\pgfsetrectcap%
\pgfsetroundjoin%
\pgfsetlinewidth{0.401500pt}%
\definecolor{currentstroke}{rgb}{0.901961,0.901961,0.901961}%
\pgfsetstrokecolor{currentstroke}%
\pgfsetdash{}{0pt}%
\pgfpathmoveto{\pgfqpoint{0.466200in}{1.083333in}}%
\pgfpathlineto{\pgfqpoint{3.230100in}{1.083333in}}%
\pgfusepath{stroke}%
\end{pgfscope}%
\begin{pgfscope}%
\definecolor{textcolor}{rgb}{0.101961,0.101961,0.101961}%
\pgfsetstrokecolor{textcolor}%
\pgfsetfillcolor{textcolor}%
\pgftext[x=0.285595in, y=1.049576in, left, base]{\color{textcolor}\sffamily\fontsize{7.000000}{8.400000}\selectfont 0.3}%
\end{pgfscope}%
\begin{pgfscope}%
\pgfpathrectangle{\pgfqpoint{0.466200in}{0.450000in}}{\pgfqpoint{2.763900in}{1.900000in}}%
\pgfusepath{clip}%
\pgfsetrectcap%
\pgfsetroundjoin%
\pgfsetlinewidth{0.401500pt}%
\definecolor{currentstroke}{rgb}{0.901961,0.901961,0.901961}%
\pgfsetstrokecolor{currentstroke}%
\pgfsetdash{}{0pt}%
\pgfpathmoveto{\pgfqpoint{0.466200in}{1.294444in}}%
\pgfpathlineto{\pgfqpoint{3.230100in}{1.294444in}}%
\pgfusepath{stroke}%
\end{pgfscope}%
\begin{pgfscope}%
\definecolor{textcolor}{rgb}{0.101961,0.101961,0.101961}%
\pgfsetstrokecolor{textcolor}%
\pgfsetfillcolor{textcolor}%
\pgftext[x=0.285595in, y=1.260687in, left, base]{\color{textcolor}\sffamily\fontsize{7.000000}{8.400000}\selectfont 0.4}%
\end{pgfscope}%
\begin{pgfscope}%
\pgfpathrectangle{\pgfqpoint{0.466200in}{0.450000in}}{\pgfqpoint{2.763900in}{1.900000in}}%
\pgfusepath{clip}%
\pgfsetrectcap%
\pgfsetroundjoin%
\pgfsetlinewidth{0.401500pt}%
\definecolor{currentstroke}{rgb}{0.901961,0.901961,0.901961}%
\pgfsetstrokecolor{currentstroke}%
\pgfsetdash{}{0pt}%
\pgfpathmoveto{\pgfqpoint{0.466200in}{1.505556in}}%
\pgfpathlineto{\pgfqpoint{3.230100in}{1.505556in}}%
\pgfusepath{stroke}%
\end{pgfscope}%
\begin{pgfscope}%
\definecolor{textcolor}{rgb}{0.101961,0.101961,0.101961}%
\pgfsetstrokecolor{textcolor}%
\pgfsetfillcolor{textcolor}%
\pgftext[x=0.285595in, y=1.471798in, left, base]{\color{textcolor}\sffamily\fontsize{7.000000}{8.400000}\selectfont 0.5}%
\end{pgfscope}%
\begin{pgfscope}%
\pgfpathrectangle{\pgfqpoint{0.466200in}{0.450000in}}{\pgfqpoint{2.763900in}{1.900000in}}%
\pgfusepath{clip}%
\pgfsetrectcap%
\pgfsetroundjoin%
\pgfsetlinewidth{0.401500pt}%
\definecolor{currentstroke}{rgb}{0.901961,0.901961,0.901961}%
\pgfsetstrokecolor{currentstroke}%
\pgfsetdash{}{0pt}%
\pgfpathmoveto{\pgfqpoint{0.466200in}{1.716667in}}%
\pgfpathlineto{\pgfqpoint{3.230100in}{1.716667in}}%
\pgfusepath{stroke}%
\end{pgfscope}%
\begin{pgfscope}%
\definecolor{textcolor}{rgb}{0.101961,0.101961,0.101961}%
\pgfsetstrokecolor{textcolor}%
\pgfsetfillcolor{textcolor}%
\pgftext[x=0.285595in, y=1.682909in, left, base]{\color{textcolor}\sffamily\fontsize{7.000000}{8.400000}\selectfont 0.6}%
\end{pgfscope}%
\begin{pgfscope}%
\pgfpathrectangle{\pgfqpoint{0.466200in}{0.450000in}}{\pgfqpoint{2.763900in}{1.900000in}}%
\pgfusepath{clip}%
\pgfsetrectcap%
\pgfsetroundjoin%
\pgfsetlinewidth{0.401500pt}%
\definecolor{currentstroke}{rgb}{0.901961,0.901961,0.901961}%
\pgfsetstrokecolor{currentstroke}%
\pgfsetdash{}{0pt}%
\pgfpathmoveto{\pgfqpoint{0.466200in}{1.927778in}}%
\pgfpathlineto{\pgfqpoint{3.230100in}{1.927778in}}%
\pgfusepath{stroke}%
\end{pgfscope}%
\begin{pgfscope}%
\definecolor{textcolor}{rgb}{0.101961,0.101961,0.101961}%
\pgfsetstrokecolor{textcolor}%
\pgfsetfillcolor{textcolor}%
\pgftext[x=0.285595in, y=1.894020in, left, base]{\color{textcolor}\sffamily\fontsize{7.000000}{8.400000}\selectfont 0.7}%
\end{pgfscope}%
\begin{pgfscope}%
\pgfpathrectangle{\pgfqpoint{0.466200in}{0.450000in}}{\pgfqpoint{2.763900in}{1.900000in}}%
\pgfusepath{clip}%
\pgfsetrectcap%
\pgfsetroundjoin%
\pgfsetlinewidth{0.401500pt}%
\definecolor{currentstroke}{rgb}{0.901961,0.901961,0.901961}%
\pgfsetstrokecolor{currentstroke}%
\pgfsetdash{}{0pt}%
\pgfpathmoveto{\pgfqpoint{0.466200in}{2.138889in}}%
\pgfpathlineto{\pgfqpoint{3.230100in}{2.138889in}}%
\pgfusepath{stroke}%
\end{pgfscope}%
\begin{pgfscope}%
\definecolor{textcolor}{rgb}{0.101961,0.101961,0.101961}%
\pgfsetstrokecolor{textcolor}%
\pgfsetfillcolor{textcolor}%
\pgftext[x=0.285595in, y=2.105131in, left, base]{\color{textcolor}\sffamily\fontsize{7.000000}{8.400000}\selectfont 0.8}%
\end{pgfscope}%
\begin{pgfscope}%
\pgfpathrectangle{\pgfqpoint{0.466200in}{0.450000in}}{\pgfqpoint{2.763900in}{1.900000in}}%
\pgfusepath{clip}%
\pgfsetrectcap%
\pgfsetroundjoin%
\pgfsetlinewidth{0.401500pt}%
\definecolor{currentstroke}{rgb}{0.901961,0.901961,0.901961}%
\pgfsetstrokecolor{currentstroke}%
\pgfsetdash{}{0pt}%
\pgfpathmoveto{\pgfqpoint{0.466200in}{2.350000in}}%
\pgfpathlineto{\pgfqpoint{3.230100in}{2.350000in}}%
\pgfusepath{stroke}%
\end{pgfscope}%
\begin{pgfscope}%
\definecolor{textcolor}{rgb}{0.101961,0.101961,0.101961}%
\pgfsetstrokecolor{textcolor}%
\pgfsetfillcolor{textcolor}%
\pgftext[x=0.285595in, y=2.316242in, left, base]{\color{textcolor}\sffamily\fontsize{7.000000}{8.400000}\selectfont 0.9}%
\end{pgfscope}%
\begin{pgfscope}%
\definecolor{textcolor}{rgb}{0.000000,0.000000,0.000000}%
\pgfsetstrokecolor{textcolor}%
\pgfsetfillcolor{textcolor}%
\pgftext[x=0.230039in,y=1.400000in,,bottom,rotate=90.000000]{\color{textcolor}\sffamily\fontsize{7.500000}{9.000000}\selectfont Nexar AP}%
\end{pgfscope}%
\begin{pgfscope}%
\pgfpathrectangle{\pgfqpoint{0.466200in}{0.450000in}}{\pgfqpoint{2.763900in}{1.900000in}}%
\pgfusepath{clip}%
\pgfsetbuttcap%
\pgfsetroundjoin%
\definecolor{currentfill}{rgb}{0.478431,0.478431,0.478431}%
\pgfsetfillcolor{currentfill}%
\pgfsetfillopacity{0.150000}%
\pgfsetlinewidth{0.000000pt}%
\definecolor{currentstroke}{rgb}{0.478431,0.478431,0.478431}%
\pgfsetstrokecolor{currentstroke}%
\pgfsetstrokeopacity{0.150000}%
\pgfsetdash{}{0pt}%
\pgfsys@defobject{currentmarker}{\pgfqpoint{0.591832in}{1.558370in}}{\pgfqpoint{3.104468in}{2.102160in}}{%
\pgfpathmoveto{\pgfqpoint{0.591832in}{2.102160in}}%
\pgfpathlineto{\pgfqpoint{0.591832in}{1.838345in}}%
\pgfpathlineto{\pgfqpoint{1.094359in}{1.683947in}}%
\pgfpathlineto{\pgfqpoint{1.596886in}{1.622483in}}%
\pgfpathlineto{\pgfqpoint{2.099414in}{1.610799in}}%
\pgfpathlineto{\pgfqpoint{2.601941in}{1.588647in}}%
\pgfpathlineto{\pgfqpoint{3.104468in}{1.558370in}}%
\pgfpathlineto{\pgfqpoint{3.104468in}{1.833386in}}%
\pgfpathlineto{\pgfqpoint{3.104468in}{1.833386in}}%
\pgfpathlineto{\pgfqpoint{2.601941in}{1.868229in}}%
\pgfpathlineto{\pgfqpoint{2.099414in}{1.879475in}}%
\pgfpathlineto{\pgfqpoint{1.596886in}{1.889580in}}%
\pgfpathlineto{\pgfqpoint{1.094359in}{1.958237in}}%
\pgfpathlineto{\pgfqpoint{0.591832in}{2.102160in}}%
\pgfpathlineto{\pgfqpoint{0.591832in}{2.102160in}}%
\pgfpathclose%
\pgfusepath{fill}%
}%
\begin{pgfscope}%
\pgfsys@transformshift{0.000000in}{0.000000in}%
\pgfsys@useobject{currentmarker}{}%
\end{pgfscope}%
\end{pgfscope}%
\begin{pgfscope}%
\pgfpathrectangle{\pgfqpoint{0.466200in}{0.450000in}}{\pgfqpoint{2.763900in}{1.900000in}}%
\pgfusepath{clip}%
\pgfsetbuttcap%
\pgfsetroundjoin%
\definecolor{currentfill}{rgb}{0.478431,0.305882,0.611765}%
\pgfsetfillcolor{currentfill}%
\pgfsetfillopacity{0.150000}%
\pgfsetlinewidth{0.000000pt}%
\definecolor{currentstroke}{rgb}{0.478431,0.305882,0.611765}%
\pgfsetstrokecolor{currentstroke}%
\pgfsetstrokeopacity{0.150000}%
\pgfsetdash{}{0pt}%
\pgfsys@defobject{currentmarker}{\pgfqpoint{0.591832in}{1.450232in}}{\pgfqpoint{3.104468in}{1.903725in}}{%
\pgfpathmoveto{\pgfqpoint{0.591832in}{1.903725in}}%
\pgfpathlineto{\pgfqpoint{0.591832in}{1.636848in}}%
\pgfpathlineto{\pgfqpoint{1.094359in}{1.557568in}}%
\pgfpathlineto{\pgfqpoint{1.596886in}{1.480742in}}%
\pgfpathlineto{\pgfqpoint{2.099414in}{1.467083in}}%
\pgfpathlineto{\pgfqpoint{2.601941in}{1.450232in}}%
\pgfpathlineto{\pgfqpoint{3.104468in}{1.455534in}}%
\pgfpathlineto{\pgfqpoint{3.104468in}{1.693863in}}%
\pgfpathlineto{\pgfqpoint{3.104468in}{1.693863in}}%
\pgfpathlineto{\pgfqpoint{2.601941in}{1.699886in}}%
\pgfpathlineto{\pgfqpoint{2.099414in}{1.712425in}}%
\pgfpathlineto{\pgfqpoint{1.596886in}{1.725660in}}%
\pgfpathlineto{\pgfqpoint{1.094359in}{1.822131in}}%
\pgfpathlineto{\pgfqpoint{0.591832in}{1.903725in}}%
\pgfpathlineto{\pgfqpoint{0.591832in}{1.903725in}}%
\pgfpathclose%
\pgfusepath{fill}%
}%
\begin{pgfscope}%
\pgfsys@transformshift{0.000000in}{0.000000in}%
\pgfsys@useobject{currentmarker}{}%
\end{pgfscope}%
\end{pgfscope}%
\begin{pgfscope}%
\pgfpathrectangle{\pgfqpoint{0.466200in}{0.450000in}}{\pgfqpoint{2.763900in}{1.900000in}}%
\pgfusepath{clip}%
\pgfsetrectcap%
\pgfsetroundjoin%
\pgfsetlinewidth{1.104125pt}%
\definecolor{currentstroke}{rgb}{0.478431,0.478431,0.478431}%
\pgfsetstrokecolor{currentstroke}%
\pgfsetdash{}{0pt}%
\pgfpathmoveto{\pgfqpoint{0.591832in}{1.969209in}}%
\pgfpathlineto{\pgfqpoint{1.094359in}{1.815117in}}%
\pgfpathlineto{\pgfqpoint{1.596886in}{1.743240in}}%
\pgfpathlineto{\pgfqpoint{2.099414in}{1.733133in}}%
\pgfpathlineto{\pgfqpoint{2.601941in}{1.712836in}}%
\pgfpathlineto{\pgfqpoint{3.104468in}{1.684173in}}%
\pgfusepath{stroke}%
\end{pgfscope}%
\begin{pgfscope}%
\pgfpathrectangle{\pgfqpoint{0.466200in}{0.450000in}}{\pgfqpoint{2.763900in}{1.900000in}}%
\pgfusepath{clip}%
\pgfsetbuttcap%
\pgfsetroundjoin%
\definecolor{currentfill}{rgb}{0.478431,0.478431,0.478431}%
\pgfsetfillcolor{currentfill}%
\pgfsetlinewidth{0.702625pt}%
\definecolor{currentstroke}{rgb}{0.478431,0.478431,0.478431}%
\pgfsetstrokecolor{currentstroke}%
\pgfsetdash{}{0pt}%
\pgfsys@defobject{currentmarker}{\pgfqpoint{-0.020833in}{-0.020833in}}{\pgfqpoint{0.020833in}{0.020833in}}{%
\pgfpathmoveto{\pgfqpoint{0.000000in}{-0.020833in}}%
\pgfpathcurveto{\pgfqpoint{0.005525in}{-0.020833in}}{\pgfqpoint{0.010825in}{-0.018638in}}{\pgfqpoint{0.014731in}{-0.014731in}}%
\pgfpathcurveto{\pgfqpoint{0.018638in}{-0.010825in}}{\pgfqpoint{0.020833in}{-0.005525in}}{\pgfqpoint{0.020833in}{0.000000in}}%
\pgfpathcurveto{\pgfqpoint{0.020833in}{0.005525in}}{\pgfqpoint{0.018638in}{0.010825in}}{\pgfqpoint{0.014731in}{0.014731in}}%
\pgfpathcurveto{\pgfqpoint{0.010825in}{0.018638in}}{\pgfqpoint{0.005525in}{0.020833in}}{\pgfqpoint{0.000000in}{0.020833in}}%
\pgfpathcurveto{\pgfqpoint{-0.005525in}{0.020833in}}{\pgfqpoint{-0.010825in}{0.018638in}}{\pgfqpoint{-0.014731in}{0.014731in}}%
\pgfpathcurveto{\pgfqpoint{-0.018638in}{0.010825in}}{\pgfqpoint{-0.020833in}{0.005525in}}{\pgfqpoint{-0.020833in}{0.000000in}}%
\pgfpathcurveto{\pgfqpoint{-0.020833in}{-0.005525in}}{\pgfqpoint{-0.018638in}{-0.010825in}}{\pgfqpoint{-0.014731in}{-0.014731in}}%
\pgfpathcurveto{\pgfqpoint{-0.010825in}{-0.018638in}}{\pgfqpoint{-0.005525in}{-0.020833in}}{\pgfqpoint{0.000000in}{-0.020833in}}%
\pgfpathlineto{\pgfqpoint{0.000000in}{-0.020833in}}%
\pgfpathclose%
\pgfusepath{stroke,fill}%
}%
\begin{pgfscope}%
\pgfsys@transformshift{0.591832in}{1.969209in}%
\pgfsys@useobject{currentmarker}{}%
\end{pgfscope}%
\begin{pgfscope}%
\pgfsys@transformshift{1.094359in}{1.815117in}%
\pgfsys@useobject{currentmarker}{}%
\end{pgfscope}%
\begin{pgfscope}%
\pgfsys@transformshift{1.596886in}{1.743240in}%
\pgfsys@useobject{currentmarker}{}%
\end{pgfscope}%
\begin{pgfscope}%
\pgfsys@transformshift{2.099414in}{1.733133in}%
\pgfsys@useobject{currentmarker}{}%
\end{pgfscope}%
\begin{pgfscope}%
\pgfsys@transformshift{2.601941in}{1.712836in}%
\pgfsys@useobject{currentmarker}{}%
\end{pgfscope}%
\begin{pgfscope}%
\pgfsys@transformshift{3.104468in}{1.684173in}%
\pgfsys@useobject{currentmarker}{}%
\end{pgfscope}%
\end{pgfscope}%
\begin{pgfscope}%
\pgfpathrectangle{\pgfqpoint{0.466200in}{0.450000in}}{\pgfqpoint{2.763900in}{1.900000in}}%
\pgfusepath{clip}%
\pgfsetrectcap%
\pgfsetroundjoin%
\pgfsetlinewidth{1.104125pt}%
\definecolor{currentstroke}{rgb}{0.478431,0.305882,0.611765}%
\pgfsetstrokecolor{currentstroke}%
\pgfsetdash{}{0pt}%
\pgfpathmoveto{\pgfqpoint{0.591832in}{1.765487in}}%
\pgfpathlineto{\pgfqpoint{1.094359in}{1.677246in}}%
\pgfpathlineto{\pgfqpoint{1.596886in}{1.584101in}}%
\pgfpathlineto{\pgfqpoint{2.099414in}{1.568506in}}%
\pgfpathlineto{\pgfqpoint{2.601941in}{1.555412in}}%
\pgfpathlineto{\pgfqpoint{3.104468in}{1.558755in}}%
\pgfusepath{stroke}%
\end{pgfscope}%
\begin{pgfscope}%
\pgfpathrectangle{\pgfqpoint{0.466200in}{0.450000in}}{\pgfqpoint{2.763900in}{1.900000in}}%
\pgfusepath{clip}%
\pgfsetbuttcap%
\pgfsetroundjoin%
\definecolor{currentfill}{rgb}{0.478431,0.305882,0.611765}%
\pgfsetfillcolor{currentfill}%
\pgfsetlinewidth{0.702625pt}%
\definecolor{currentstroke}{rgb}{0.478431,0.305882,0.611765}%
\pgfsetstrokecolor{currentstroke}%
\pgfsetdash{}{0pt}%
\pgfsys@defobject{currentmarker}{\pgfqpoint{-0.020833in}{-0.020833in}}{\pgfqpoint{0.020833in}{0.020833in}}{%
\pgfpathmoveto{\pgfqpoint{0.000000in}{-0.020833in}}%
\pgfpathcurveto{\pgfqpoint{0.005525in}{-0.020833in}}{\pgfqpoint{0.010825in}{-0.018638in}}{\pgfqpoint{0.014731in}{-0.014731in}}%
\pgfpathcurveto{\pgfqpoint{0.018638in}{-0.010825in}}{\pgfqpoint{0.020833in}{-0.005525in}}{\pgfqpoint{0.020833in}{0.000000in}}%
\pgfpathcurveto{\pgfqpoint{0.020833in}{0.005525in}}{\pgfqpoint{0.018638in}{0.010825in}}{\pgfqpoint{0.014731in}{0.014731in}}%
\pgfpathcurveto{\pgfqpoint{0.010825in}{0.018638in}}{\pgfqpoint{0.005525in}{0.020833in}}{\pgfqpoint{0.000000in}{0.020833in}}%
\pgfpathcurveto{\pgfqpoint{-0.005525in}{0.020833in}}{\pgfqpoint{-0.010825in}{0.018638in}}{\pgfqpoint{-0.014731in}{0.014731in}}%
\pgfpathcurveto{\pgfqpoint{-0.018638in}{0.010825in}}{\pgfqpoint{-0.020833in}{0.005525in}}{\pgfqpoint{-0.020833in}{0.000000in}}%
\pgfpathcurveto{\pgfqpoint{-0.020833in}{-0.005525in}}{\pgfqpoint{-0.018638in}{-0.010825in}}{\pgfqpoint{-0.014731in}{-0.014731in}}%
\pgfpathcurveto{\pgfqpoint{-0.010825in}{-0.018638in}}{\pgfqpoint{-0.005525in}{-0.020833in}}{\pgfqpoint{0.000000in}{-0.020833in}}%
\pgfpathlineto{\pgfqpoint{0.000000in}{-0.020833in}}%
\pgfpathclose%
\pgfusepath{stroke,fill}%
}%
\begin{pgfscope}%
\pgfsys@transformshift{0.591832in}{1.765487in}%
\pgfsys@useobject{currentmarker}{}%
\end{pgfscope}%
\begin{pgfscope}%
\pgfsys@transformshift{1.094359in}{1.677246in}%
\pgfsys@useobject{currentmarker}{}%
\end{pgfscope}%
\begin{pgfscope}%
\pgfsys@transformshift{1.596886in}{1.584101in}%
\pgfsys@useobject{currentmarker}{}%
\end{pgfscope}%
\begin{pgfscope}%
\pgfsys@transformshift{2.099414in}{1.568506in}%
\pgfsys@useobject{currentmarker}{}%
\end{pgfscope}%
\begin{pgfscope}%
\pgfsys@transformshift{2.601941in}{1.555412in}%
\pgfsys@useobject{currentmarker}{}%
\end{pgfscope}%
\begin{pgfscope}%
\pgfsys@transformshift{3.104468in}{1.558755in}%
\pgfsys@useobject{currentmarker}{}%
\end{pgfscope}%
\end{pgfscope}%
\begin{pgfscope}%
\pgfpathrectangle{\pgfqpoint{0.466200in}{0.450000in}}{\pgfqpoint{2.763900in}{1.900000in}}%
\pgfusepath{clip}%
\pgfsetbuttcap%
\pgfsetroundjoin%
\pgfsetlinewidth{0.803000pt}%
\definecolor{currentstroke}{rgb}{0.478431,0.478431,0.478431}%
\pgfsetstrokecolor{currentstroke}%
\pgfsetdash{{2.400000pt}{1.200000pt}}{0.000000pt}%
\pgfpathmoveto{\pgfqpoint{0.591832in}{2.099584in}}%
\pgfpathlineto{\pgfqpoint{1.094359in}{1.790487in}}%
\pgfpathlineto{\pgfqpoint{1.596886in}{1.492719in}}%
\pgfpathlineto{\pgfqpoint{2.099414in}{1.050540in}}%
\pgfpathlineto{\pgfqpoint{2.601941in}{0.736034in}}%
\pgfpathlineto{\pgfqpoint{3.104468in}{0.670413in}}%
\pgfusepath{stroke}%
\end{pgfscope}%
\begin{pgfscope}%
\pgfpathrectangle{\pgfqpoint{0.466200in}{0.450000in}}{\pgfqpoint{2.763900in}{1.900000in}}%
\pgfusepath{clip}%
\pgfsetbuttcap%
\pgfsetmiterjoin%
\definecolor{currentfill}{rgb}{1.000000,1.000000,1.000000}%
\pgfsetfillcolor{currentfill}%
\pgfsetlinewidth{0.702625pt}%
\definecolor{currentstroke}{rgb}{0.478431,0.478431,0.478431}%
\pgfsetstrokecolor{currentstroke}%
\pgfsetdash{}{0pt}%
\pgfsys@defobject{currentmarker}{\pgfqpoint{-0.020833in}{-0.020833in}}{\pgfqpoint{0.020833in}{0.020833in}}{%
\pgfpathmoveto{\pgfqpoint{-0.020833in}{-0.020833in}}%
\pgfpathlineto{\pgfqpoint{0.020833in}{-0.020833in}}%
\pgfpathlineto{\pgfqpoint{0.020833in}{0.020833in}}%
\pgfpathlineto{\pgfqpoint{-0.020833in}{0.020833in}}%
\pgfpathlineto{\pgfqpoint{-0.020833in}{-0.020833in}}%
\pgfpathclose%
\pgfusepath{stroke,fill}%
}%
\begin{pgfscope}%
\pgfsys@transformshift{0.591832in}{2.099584in}%
\pgfsys@useobject{currentmarker}{}%
\end{pgfscope}%
\begin{pgfscope}%
\pgfsys@transformshift{1.094359in}{1.790487in}%
\pgfsys@useobject{currentmarker}{}%
\end{pgfscope}%
\begin{pgfscope}%
\pgfsys@transformshift{1.596886in}{1.492719in}%
\pgfsys@useobject{currentmarker}{}%
\end{pgfscope}%
\begin{pgfscope}%
\pgfsys@transformshift{2.099414in}{1.050540in}%
\pgfsys@useobject{currentmarker}{}%
\end{pgfscope}%
\begin{pgfscope}%
\pgfsys@transformshift{2.601941in}{0.736034in}%
\pgfsys@useobject{currentmarker}{}%
\end{pgfscope}%
\begin{pgfscope}%
\pgfsys@transformshift{3.104468in}{0.670413in}%
\pgfsys@useobject{currentmarker}{}%
\end{pgfscope}%
\end{pgfscope}%
\begin{pgfscope}%
\pgfpathrectangle{\pgfqpoint{0.466200in}{0.450000in}}{\pgfqpoint{2.763900in}{1.900000in}}%
\pgfusepath{clip}%
\pgfsetbuttcap%
\pgfsetroundjoin%
\pgfsetlinewidth{0.803000pt}%
\definecolor{currentstroke}{rgb}{0.478431,0.305882,0.611765}%
\pgfsetstrokecolor{currentstroke}%
\pgfsetdash{{2.400000pt}{1.200000pt}}{0.000000pt}%
\pgfpathmoveto{\pgfqpoint{0.591832in}{1.870072in}}%
\pgfpathlineto{\pgfqpoint{1.094359in}{1.606699in}}%
\pgfpathlineto{\pgfqpoint{1.596886in}{1.181228in}}%
\pgfpathlineto{\pgfqpoint{2.099414in}{0.876356in}}%
\pgfpathlineto{\pgfqpoint{2.601941in}{0.682850in}}%
\pgfpathlineto{\pgfqpoint{3.104468in}{0.570634in}}%
\pgfusepath{stroke}%
\end{pgfscope}%
\begin{pgfscope}%
\pgfpathrectangle{\pgfqpoint{0.466200in}{0.450000in}}{\pgfqpoint{2.763900in}{1.900000in}}%
\pgfusepath{clip}%
\pgfsetbuttcap%
\pgfsetmiterjoin%
\definecolor{currentfill}{rgb}{1.000000,1.000000,1.000000}%
\pgfsetfillcolor{currentfill}%
\pgfsetlinewidth{0.702625pt}%
\definecolor{currentstroke}{rgb}{0.478431,0.305882,0.611765}%
\pgfsetstrokecolor{currentstroke}%
\pgfsetdash{}{0pt}%
\pgfsys@defobject{currentmarker}{\pgfqpoint{-0.020833in}{-0.020833in}}{\pgfqpoint{0.020833in}{0.020833in}}{%
\pgfpathmoveto{\pgfqpoint{-0.020833in}{-0.020833in}}%
\pgfpathlineto{\pgfqpoint{0.020833in}{-0.020833in}}%
\pgfpathlineto{\pgfqpoint{0.020833in}{0.020833in}}%
\pgfpathlineto{\pgfqpoint{-0.020833in}{0.020833in}}%
\pgfpathlineto{\pgfqpoint{-0.020833in}{-0.020833in}}%
\pgfpathclose%
\pgfusepath{stroke,fill}%
}%
\begin{pgfscope}%
\pgfsys@transformshift{0.591832in}{1.870072in}%
\pgfsys@useobject{currentmarker}{}%
\end{pgfscope}%
\begin{pgfscope}%
\pgfsys@transformshift{1.094359in}{1.606699in}%
\pgfsys@useobject{currentmarker}{}%
\end{pgfscope}%
\begin{pgfscope}%
\pgfsys@transformshift{1.596886in}{1.181228in}%
\pgfsys@useobject{currentmarker}{}%
\end{pgfscope}%
\begin{pgfscope}%
\pgfsys@transformshift{2.099414in}{0.876356in}%
\pgfsys@useobject{currentmarker}{}%
\end{pgfscope}%
\begin{pgfscope}%
\pgfsys@transformshift{2.601941in}{0.682850in}%
\pgfsys@useobject{currentmarker}{}%
\end{pgfscope}%
\begin{pgfscope}%
\pgfsys@transformshift{3.104468in}{0.570634in}%
\pgfsys@useobject{currentmarker}{}%
\end{pgfscope}%
\end{pgfscope}%
\begin{pgfscope}%
\pgfpathrectangle{\pgfqpoint{0.466200in}{0.450000in}}{\pgfqpoint{2.763900in}{1.900000in}}%
\pgfusepath{clip}%
\pgfsetbuttcap%
\pgfsetroundjoin%
\pgfsetlinewidth{0.501875pt}%
\definecolor{currentstroke}{rgb}{0.549020,0.549020,0.549020}%
\pgfsetstrokecolor{currentstroke}%
\pgfsetdash{{0.500000pt}{0.825000pt}}{0.000000pt}%
\pgfpathmoveto{\pgfqpoint{0.466200in}{1.505556in}}%
\pgfpathlineto{\pgfqpoint{3.230100in}{1.505556in}}%
\pgfusepath{stroke}%
\end{pgfscope}%
\begin{pgfscope}%
\pgfsetrectcap%
\pgfsetmiterjoin%
\pgfsetlinewidth{0.803000pt}%
\definecolor{currentstroke}{rgb}{0.101961,0.101961,0.101961}%
\pgfsetstrokecolor{currentstroke}%
\pgfsetdash{}{0pt}%
\pgfpathmoveto{\pgfqpoint{0.466200in}{0.450000in}}%
\pgfpathlineto{\pgfqpoint{0.466200in}{2.350000in}}%
\pgfusepath{stroke}%
\end{pgfscope}%
\begin{pgfscope}%
\pgfsetrectcap%
\pgfsetmiterjoin%
\pgfsetlinewidth{0.803000pt}%
\definecolor{currentstroke}{rgb}{0.101961,0.101961,0.101961}%
\pgfsetstrokecolor{currentstroke}%
\pgfsetdash{}{0pt}%
\pgfpathmoveto{\pgfqpoint{0.466200in}{0.450000in}}%
\pgfpathlineto{\pgfqpoint{3.230100in}{0.450000in}}%
\pgfusepath{stroke}%
\end{pgfscope}%
\begin{pgfscope}%
\pgfsetrectcap%
\pgfsetroundjoin%
\pgfsetlinewidth{1.104125pt}%
\definecolor{currentstroke}{rgb}{0.478431,0.478431,0.478431}%
\pgfsetstrokecolor{currentstroke}%
\pgfsetdash{}{0pt}%
\pgfpathmoveto{\pgfqpoint{1.968688in}{2.237153in}}%
\pgfpathlineto{\pgfqpoint{2.067994in}{2.237153in}}%
\pgfpathlineto{\pgfqpoint{2.167299in}{2.237153in}}%
\pgfusepath{stroke}%
\end{pgfscope}%
\begin{pgfscope}%
\pgfsetbuttcap%
\pgfsetroundjoin%
\definecolor{currentfill}{rgb}{0.478431,0.478431,0.478431}%
\pgfsetfillcolor{currentfill}%
\pgfsetlinewidth{0.702625pt}%
\definecolor{currentstroke}{rgb}{0.478431,0.478431,0.478431}%
\pgfsetstrokecolor{currentstroke}%
\pgfsetdash{}{0pt}%
\pgfsys@defobject{currentmarker}{\pgfqpoint{-0.020833in}{-0.020833in}}{\pgfqpoint{0.020833in}{0.020833in}}{%
\pgfpathmoveto{\pgfqpoint{0.000000in}{-0.020833in}}%
\pgfpathcurveto{\pgfqpoint{0.005525in}{-0.020833in}}{\pgfqpoint{0.010825in}{-0.018638in}}{\pgfqpoint{0.014731in}{-0.014731in}}%
\pgfpathcurveto{\pgfqpoint{0.018638in}{-0.010825in}}{\pgfqpoint{0.020833in}{-0.005525in}}{\pgfqpoint{0.020833in}{0.000000in}}%
\pgfpathcurveto{\pgfqpoint{0.020833in}{0.005525in}}{\pgfqpoint{0.018638in}{0.010825in}}{\pgfqpoint{0.014731in}{0.014731in}}%
\pgfpathcurveto{\pgfqpoint{0.010825in}{0.018638in}}{\pgfqpoint{0.005525in}{0.020833in}}{\pgfqpoint{0.000000in}{0.020833in}}%
\pgfpathcurveto{\pgfqpoint{-0.005525in}{0.020833in}}{\pgfqpoint{-0.010825in}{0.018638in}}{\pgfqpoint{-0.014731in}{0.014731in}}%
\pgfpathcurveto{\pgfqpoint{-0.018638in}{0.010825in}}{\pgfqpoint{-0.020833in}{0.005525in}}{\pgfqpoint{-0.020833in}{0.000000in}}%
\pgfpathcurveto{\pgfqpoint{-0.020833in}{-0.005525in}}{\pgfqpoint{-0.018638in}{-0.010825in}}{\pgfqpoint{-0.014731in}{-0.014731in}}%
\pgfpathcurveto{\pgfqpoint{-0.010825in}{-0.018638in}}{\pgfqpoint{-0.005525in}{-0.020833in}}{\pgfqpoint{0.000000in}{-0.020833in}}%
\pgfpathlineto{\pgfqpoint{0.000000in}{-0.020833in}}%
\pgfpathclose%
\pgfusepath{stroke,fill}%
}%
\begin{pgfscope}%
\pgfsys@transformshift{2.067994in}{2.237153in}%
\pgfsys@useobject{currentmarker}{}%
\end{pgfscope}%
\end{pgfscope}%
\begin{pgfscope}%
\definecolor{textcolor}{rgb}{0.101961,0.101961,0.101961}%
\pgfsetstrokecolor{textcolor}%
\pgfsetfillcolor{textcolor}%
\pgftext[x=2.239522in,y=2.205556in,left,base]{\color{textcolor}\sffamily\fontsize{6.500000}{7.800000}\selectfont cloud 32B, 2 fps episode}%
\end{pgfscope}%
\begin{pgfscope}%
\pgfsetrectcap%
\pgfsetroundjoin%
\pgfsetlinewidth{1.104125pt}%
\definecolor{currentstroke}{rgb}{0.478431,0.305882,0.611765}%
\pgfsetstrokecolor{currentstroke}%
\pgfsetdash{}{0pt}%
\pgfpathmoveto{\pgfqpoint{1.968688in}{2.135185in}}%
\pgfpathlineto{\pgfqpoint{2.067994in}{2.135185in}}%
\pgfpathlineto{\pgfqpoint{2.167299in}{2.135185in}}%
\pgfusepath{stroke}%
\end{pgfscope}%
\begin{pgfscope}%
\pgfsetbuttcap%
\pgfsetroundjoin%
\definecolor{currentfill}{rgb}{0.478431,0.305882,0.611765}%
\pgfsetfillcolor{currentfill}%
\pgfsetlinewidth{0.702625pt}%
\definecolor{currentstroke}{rgb}{0.478431,0.305882,0.611765}%
\pgfsetstrokecolor{currentstroke}%
\pgfsetdash{}{0pt}%
\pgfsys@defobject{currentmarker}{\pgfqpoint{-0.020833in}{-0.020833in}}{\pgfqpoint{0.020833in}{0.020833in}}{%
\pgfpathmoveto{\pgfqpoint{0.000000in}{-0.020833in}}%
\pgfpathcurveto{\pgfqpoint{0.005525in}{-0.020833in}}{\pgfqpoint{0.010825in}{-0.018638in}}{\pgfqpoint{0.014731in}{-0.014731in}}%
\pgfpathcurveto{\pgfqpoint{0.018638in}{-0.010825in}}{\pgfqpoint{0.020833in}{-0.005525in}}{\pgfqpoint{0.020833in}{0.000000in}}%
\pgfpathcurveto{\pgfqpoint{0.020833in}{0.005525in}}{\pgfqpoint{0.018638in}{0.010825in}}{\pgfqpoint{0.014731in}{0.014731in}}%
\pgfpathcurveto{\pgfqpoint{0.010825in}{0.018638in}}{\pgfqpoint{0.005525in}{0.020833in}}{\pgfqpoint{0.000000in}{0.020833in}}%
\pgfpathcurveto{\pgfqpoint{-0.005525in}{0.020833in}}{\pgfqpoint{-0.010825in}{0.018638in}}{\pgfqpoint{-0.014731in}{0.014731in}}%
\pgfpathcurveto{\pgfqpoint{-0.018638in}{0.010825in}}{\pgfqpoint{-0.020833in}{0.005525in}}{\pgfqpoint{-0.020833in}{0.000000in}}%
\pgfpathcurveto{\pgfqpoint{-0.020833in}{-0.005525in}}{\pgfqpoint{-0.018638in}{-0.010825in}}{\pgfqpoint{-0.014731in}{-0.014731in}}%
\pgfpathcurveto{\pgfqpoint{-0.010825in}{-0.018638in}}{\pgfqpoint{-0.005525in}{-0.020833in}}{\pgfqpoint{0.000000in}{-0.020833in}}%
\pgfpathlineto{\pgfqpoint{0.000000in}{-0.020833in}}%
\pgfpathclose%
\pgfusepath{stroke,fill}%
}%
\begin{pgfscope}%
\pgfsys@transformshift{2.067994in}{2.135185in}%
\pgfsys@useobject{currentmarker}{}%
\end{pgfscope}%
\end{pgfscope}%
\begin{pgfscope}%
\definecolor{textcolor}{rgb}{0.101961,0.101961,0.101961}%
\pgfsetstrokecolor{textcolor}%
\pgfsetfillcolor{textcolor}%
\pgftext[x=2.239522in,y=2.103588in,left,base]{\color{textcolor}\sffamily\fontsize{6.500000}{7.800000}\selectfont device 4B, 2 fps episode}%
\end{pgfscope}%
\begin{pgfscope}%
\pgfsetbuttcap%
\pgfsetroundjoin%
\pgfsetlinewidth{0.803000pt}%
\definecolor{currentstroke}{rgb}{0.478431,0.478431,0.478431}%
\pgfsetstrokecolor{currentstroke}%
\pgfsetdash{{2.400000pt}{1.200000pt}}{0.000000pt}%
\pgfpathmoveto{\pgfqpoint{1.968688in}{2.033218in}}%
\pgfpathlineto{\pgfqpoint{2.067994in}{2.033218in}}%
\pgfpathlineto{\pgfqpoint{2.167299in}{2.033218in}}%
\pgfusepath{stroke}%
\end{pgfscope}%
\begin{pgfscope}%
\pgfsetbuttcap%
\pgfsetmiterjoin%
\definecolor{currentfill}{rgb}{1.000000,1.000000,1.000000}%
\pgfsetfillcolor{currentfill}%
\pgfsetlinewidth{0.702625pt}%
\definecolor{currentstroke}{rgb}{0.478431,0.478431,0.478431}%
\pgfsetstrokecolor{currentstroke}%
\pgfsetdash{}{0pt}%
\pgfsys@defobject{currentmarker}{\pgfqpoint{-0.020833in}{-0.020833in}}{\pgfqpoint{0.020833in}{0.020833in}}{%
\pgfpathmoveto{\pgfqpoint{-0.020833in}{-0.020833in}}%
\pgfpathlineto{\pgfqpoint{0.020833in}{-0.020833in}}%
\pgfpathlineto{\pgfqpoint{0.020833in}{0.020833in}}%
\pgfpathlineto{\pgfqpoint{-0.020833in}{0.020833in}}%
\pgfpathlineto{\pgfqpoint{-0.020833in}{-0.020833in}}%
\pgfpathclose%
\pgfusepath{stroke,fill}%
}%
\begin{pgfscope}%
\pgfsys@transformshift{2.067994in}{2.033218in}%
\pgfsys@useobject{currentmarker}{}%
\end{pgfscope}%
\end{pgfscope}%
\begin{pgfscope}%
\definecolor{textcolor}{rgb}{0.101961,0.101961,0.101961}%
\pgfsetstrokecolor{textcolor}%
\pgfsetfillcolor{textcolor}%
\pgftext[x=2.239522in,y=2.001621in,left,base]{\color{textcolor}\sffamily\fontsize{6.500000}{7.800000}\selectfont cloud 32B, one frame}%
\end{pgfscope}%
\begin{pgfscope}%
\pgfsetbuttcap%
\pgfsetroundjoin%
\pgfsetlinewidth{0.803000pt}%
\definecolor{currentstroke}{rgb}{0.478431,0.305882,0.611765}%
\pgfsetstrokecolor{currentstroke}%
\pgfsetdash{{2.400000pt}{1.200000pt}}{0.000000pt}%
\pgfpathmoveto{\pgfqpoint{1.968688in}{1.931250in}}%
\pgfpathlineto{\pgfqpoint{2.067994in}{1.931250in}}%
\pgfpathlineto{\pgfqpoint{2.167299in}{1.931250in}}%
\pgfusepath{stroke}%
\end{pgfscope}%
\begin{pgfscope}%
\pgfsetbuttcap%
\pgfsetmiterjoin%
\definecolor{currentfill}{rgb}{1.000000,1.000000,1.000000}%
\pgfsetfillcolor{currentfill}%
\pgfsetlinewidth{0.702625pt}%
\definecolor{currentstroke}{rgb}{0.478431,0.305882,0.611765}%
\pgfsetstrokecolor{currentstroke}%
\pgfsetdash{}{0pt}%
\pgfsys@defobject{currentmarker}{\pgfqpoint{-0.020833in}{-0.020833in}}{\pgfqpoint{0.020833in}{0.020833in}}{%
\pgfpathmoveto{\pgfqpoint{-0.020833in}{-0.020833in}}%
\pgfpathlineto{\pgfqpoint{0.020833in}{-0.020833in}}%
\pgfpathlineto{\pgfqpoint{0.020833in}{0.020833in}}%
\pgfpathlineto{\pgfqpoint{-0.020833in}{0.020833in}}%
\pgfpathlineto{\pgfqpoint{-0.020833in}{-0.020833in}}%
\pgfpathclose%
\pgfusepath{stroke,fill}%
}%
\begin{pgfscope}%
\pgfsys@transformshift{2.067994in}{1.931250in}%
\pgfsys@useobject{currentmarker}{}%
\end{pgfscope}%
\end{pgfscope}%
\begin{pgfscope}%
\definecolor{textcolor}{rgb}{0.101961,0.101961,0.101961}%
\pgfsetstrokecolor{textcolor}%
\pgfsetfillcolor{textcolor}%
\pgftext[x=2.239522in,y=1.899653in,left,base]{\color{textcolor}\sffamily\fontsize{6.500000}{7.800000}\selectfont device 4B, one frame}%
\end{pgfscope}%
\end{pgfpicture}%
\makeatother%
\endgroup%
}
\caption{[PLACEHOLDER: committed model-level measurement, StudyEARLY1; no COVEX data, not FINAL] Crash-anticipation average precision (Nexar, 300 validation clips) versus time to event for the cloud model (Qwen3-VL-32B) and the on-device model (Qwen3-VL-4B), model-level measurement from StudyEARLY1 with 95 percent bootstrap bands over source videos for the 2 fps episodic format; the one-frame format is dashed. Dotted line is the chance level (0.5 for the 2 fps format). Cloud delivery latency (median 0.69 s) is not charged here.}
\label{fig:F15_ap_vs_time_to_event}
\end{figure}
\clearpage
\begin{figure*}[t]
\centering
{\fontfamily{LinuxBiolinumT}\selectfont
%% Creator: Matplotlib, PGF backend
%%
%% To include the figure in your LaTeX document, write
%%   \input{<filename>.pgf}
%%
%% Make sure the required packages are loaded in your preamble
%%   \usepackage{pgf}
%%
%% Also ensure that all the required font packages are loaded; for instance,
%% the lmodern package is sometimes necessary when using math font.
%%   \usepackage{lmodern}
%%
%% Figures using additional raster images can only be included by \input if
%% they are in the same directory as the main LaTeX file. For loading figures
%% from other directories you can use the `import` package
%%   \usepackage{import}
%%
%% and then include the figures with
%%   \import{<path to file>}{<filename>.pgf}
%%
%% Matplotlib used the following preamble
%%
\begingroup%
\makeatletter%
\begin{pgfpicture}%
\pgfpathrectangle{\pgfpointorigin}{\pgfqpoint{7.000000in}{2.500000in}}%
\pgfusepath{use as bounding box, clip}%
\begin{pgfscope}%
\pgfsetbuttcap%
\pgfsetmiterjoin%
\definecolor{currentfill}{rgb}{1.000000,1.000000,1.000000}%
\pgfsetfillcolor{currentfill}%
\pgfsetlinewidth{0.000000pt}%
\definecolor{currentstroke}{rgb}{1.000000,1.000000,1.000000}%
\pgfsetstrokecolor{currentstroke}%
\pgfsetdash{}{0pt}%
\pgfpathmoveto{\pgfqpoint{0.000000in}{0.000000in}}%
\pgfpathlineto{\pgfqpoint{7.000000in}{0.000000in}}%
\pgfpathlineto{\pgfqpoint{7.000000in}{2.500000in}}%
\pgfpathlineto{\pgfqpoint{0.000000in}{2.500000in}}%
\pgfpathlineto{\pgfqpoint{0.000000in}{0.000000in}}%
\pgfpathclose%
\pgfusepath{fill}%
\end{pgfscope}%
\begin{pgfscope}%
\pgfsetbuttcap%
\pgfsetmiterjoin%
\definecolor{currentfill}{rgb}{1.000000,1.000000,1.000000}%
\pgfsetfillcolor{currentfill}%
\pgfsetlinewidth{0.000000pt}%
\definecolor{currentstroke}{rgb}{0.000000,0.000000,0.000000}%
\pgfsetstrokecolor{currentstroke}%
\pgfsetstrokeopacity{0.000000}%
\pgfsetdash{}{0pt}%
\pgfpathmoveto{\pgfqpoint{0.595000in}{0.500000in}}%
\pgfpathlineto{\pgfqpoint{3.115000in}{0.500000in}}%
\pgfpathlineto{\pgfqpoint{3.115000in}{2.350000in}}%
\pgfpathlineto{\pgfqpoint{0.595000in}{2.350000in}}%
\pgfpathlineto{\pgfqpoint{0.595000in}{0.500000in}}%
\pgfpathclose%
\pgfusepath{fill}%
\end{pgfscope}%
\begin{pgfscope}%
\definecolor{textcolor}{rgb}{0.101961,0.101961,0.101961}%
\pgfsetstrokecolor{textcolor}%
\pgfsetfillcolor{textcolor}%
\pgftext[x=1.282273in,y=0.451389in,,top]{\color{textcolor}\sffamily\fontsize{7.000000}{8.400000}\selectfont device 4B}%
\end{pgfscope}%
\begin{pgfscope}%
\definecolor{textcolor}{rgb}{0.101961,0.101961,0.101961}%
\pgfsetstrokecolor{textcolor}%
\pgfsetfillcolor{textcolor}%
\pgftext[x=2.427727in,y=0.451389in,,top]{\color{textcolor}\sffamily\fontsize{7.000000}{8.400000}\selectfont cloud 32B}%
\end{pgfscope}%
\begin{pgfscope}%
\pgfpathrectangle{\pgfqpoint{0.595000in}{0.500000in}}{\pgfqpoint{2.520000in}{1.850000in}}%
\pgfusepath{clip}%
\pgfsetrectcap%
\pgfsetroundjoin%
\pgfsetlinewidth{0.401500pt}%
\definecolor{currentstroke}{rgb}{0.901961,0.901961,0.901961}%
\pgfsetstrokecolor{currentstroke}%
\pgfsetdash{}{0pt}%
\pgfpathmoveto{\pgfqpoint{0.595000in}{0.500000in}}%
\pgfpathlineto{\pgfqpoint{3.115000in}{0.500000in}}%
\pgfusepath{stroke}%
\end{pgfscope}%
\begin{pgfscope}%
\definecolor{textcolor}{rgb}{0.101961,0.101961,0.101961}%
\pgfsetstrokecolor{textcolor}%
\pgfsetfillcolor{textcolor}%
\pgftext[x=0.414395in, y=0.466242in, left, base]{\color{textcolor}\sffamily\fontsize{7.000000}{8.400000}\selectfont 0.0}%
\end{pgfscope}%
\begin{pgfscope}%
\pgfpathrectangle{\pgfqpoint{0.595000in}{0.500000in}}{\pgfqpoint{2.520000in}{1.850000in}}%
\pgfusepath{clip}%
\pgfsetrectcap%
\pgfsetroundjoin%
\pgfsetlinewidth{0.401500pt}%
\definecolor{currentstroke}{rgb}{0.901961,0.901961,0.901961}%
\pgfsetstrokecolor{currentstroke}%
\pgfsetdash{}{0pt}%
\pgfpathmoveto{\pgfqpoint{0.595000in}{0.731250in}}%
\pgfpathlineto{\pgfqpoint{3.115000in}{0.731250in}}%
\pgfusepath{stroke}%
\end{pgfscope}%
\begin{pgfscope}%
\definecolor{textcolor}{rgb}{0.101961,0.101961,0.101961}%
\pgfsetstrokecolor{textcolor}%
\pgfsetfillcolor{textcolor}%
\pgftext[x=0.414395in, y=0.697492in, left, base]{\color{textcolor}\sffamily\fontsize{7.000000}{8.400000}\selectfont 0.1}%
\end{pgfscope}%
\begin{pgfscope}%
\pgfpathrectangle{\pgfqpoint{0.595000in}{0.500000in}}{\pgfqpoint{2.520000in}{1.850000in}}%
\pgfusepath{clip}%
\pgfsetrectcap%
\pgfsetroundjoin%
\pgfsetlinewidth{0.401500pt}%
\definecolor{currentstroke}{rgb}{0.901961,0.901961,0.901961}%
\pgfsetstrokecolor{currentstroke}%
\pgfsetdash{}{0pt}%
\pgfpathmoveto{\pgfqpoint{0.595000in}{0.962500in}}%
\pgfpathlineto{\pgfqpoint{3.115000in}{0.962500in}}%
\pgfusepath{stroke}%
\end{pgfscope}%
\begin{pgfscope}%
\definecolor{textcolor}{rgb}{0.101961,0.101961,0.101961}%
\pgfsetstrokecolor{textcolor}%
\pgfsetfillcolor{textcolor}%
\pgftext[x=0.414395in, y=0.928742in, left, base]{\color{textcolor}\sffamily\fontsize{7.000000}{8.400000}\selectfont 0.2}%
\end{pgfscope}%
\begin{pgfscope}%
\pgfpathrectangle{\pgfqpoint{0.595000in}{0.500000in}}{\pgfqpoint{2.520000in}{1.850000in}}%
\pgfusepath{clip}%
\pgfsetrectcap%
\pgfsetroundjoin%
\pgfsetlinewidth{0.401500pt}%
\definecolor{currentstroke}{rgb}{0.901961,0.901961,0.901961}%
\pgfsetstrokecolor{currentstroke}%
\pgfsetdash{}{0pt}%
\pgfpathmoveto{\pgfqpoint{0.595000in}{1.193750in}}%
\pgfpathlineto{\pgfqpoint{3.115000in}{1.193750in}}%
\pgfusepath{stroke}%
\end{pgfscope}%
\begin{pgfscope}%
\definecolor{textcolor}{rgb}{0.101961,0.101961,0.101961}%
\pgfsetstrokecolor{textcolor}%
\pgfsetfillcolor{textcolor}%
\pgftext[x=0.414395in, y=1.159992in, left, base]{\color{textcolor}\sffamily\fontsize{7.000000}{8.400000}\selectfont 0.3}%
\end{pgfscope}%
\begin{pgfscope}%
\pgfpathrectangle{\pgfqpoint{0.595000in}{0.500000in}}{\pgfqpoint{2.520000in}{1.850000in}}%
\pgfusepath{clip}%
\pgfsetrectcap%
\pgfsetroundjoin%
\pgfsetlinewidth{0.401500pt}%
\definecolor{currentstroke}{rgb}{0.901961,0.901961,0.901961}%
\pgfsetstrokecolor{currentstroke}%
\pgfsetdash{}{0pt}%
\pgfpathmoveto{\pgfqpoint{0.595000in}{1.425000in}}%
\pgfpathlineto{\pgfqpoint{3.115000in}{1.425000in}}%
\pgfusepath{stroke}%
\end{pgfscope}%
\begin{pgfscope}%
\definecolor{textcolor}{rgb}{0.101961,0.101961,0.101961}%
\pgfsetstrokecolor{textcolor}%
\pgfsetfillcolor{textcolor}%
\pgftext[x=0.414395in, y=1.391242in, left, base]{\color{textcolor}\sffamily\fontsize{7.000000}{8.400000}\selectfont 0.4}%
\end{pgfscope}%
\begin{pgfscope}%
\pgfpathrectangle{\pgfqpoint{0.595000in}{0.500000in}}{\pgfqpoint{2.520000in}{1.850000in}}%
\pgfusepath{clip}%
\pgfsetrectcap%
\pgfsetroundjoin%
\pgfsetlinewidth{0.401500pt}%
\definecolor{currentstroke}{rgb}{0.901961,0.901961,0.901961}%
\pgfsetstrokecolor{currentstroke}%
\pgfsetdash{}{0pt}%
\pgfpathmoveto{\pgfqpoint{0.595000in}{1.656250in}}%
\pgfpathlineto{\pgfqpoint{3.115000in}{1.656250in}}%
\pgfusepath{stroke}%
\end{pgfscope}%
\begin{pgfscope}%
\definecolor{textcolor}{rgb}{0.101961,0.101961,0.101961}%
\pgfsetstrokecolor{textcolor}%
\pgfsetfillcolor{textcolor}%
\pgftext[x=0.414395in, y=1.622492in, left, base]{\color{textcolor}\sffamily\fontsize{7.000000}{8.400000}\selectfont 0.5}%
\end{pgfscope}%
\begin{pgfscope}%
\pgfpathrectangle{\pgfqpoint{0.595000in}{0.500000in}}{\pgfqpoint{2.520000in}{1.850000in}}%
\pgfusepath{clip}%
\pgfsetrectcap%
\pgfsetroundjoin%
\pgfsetlinewidth{0.401500pt}%
\definecolor{currentstroke}{rgb}{0.901961,0.901961,0.901961}%
\pgfsetstrokecolor{currentstroke}%
\pgfsetdash{}{0pt}%
\pgfpathmoveto{\pgfqpoint{0.595000in}{1.887500in}}%
\pgfpathlineto{\pgfqpoint{3.115000in}{1.887500in}}%
\pgfusepath{stroke}%
\end{pgfscope}%
\begin{pgfscope}%
\definecolor{textcolor}{rgb}{0.101961,0.101961,0.101961}%
\pgfsetstrokecolor{textcolor}%
\pgfsetfillcolor{textcolor}%
\pgftext[x=0.414395in, y=1.853742in, left, base]{\color{textcolor}\sffamily\fontsize{7.000000}{8.400000}\selectfont 0.6}%
\end{pgfscope}%
\begin{pgfscope}%
\pgfpathrectangle{\pgfqpoint{0.595000in}{0.500000in}}{\pgfqpoint{2.520000in}{1.850000in}}%
\pgfusepath{clip}%
\pgfsetrectcap%
\pgfsetroundjoin%
\pgfsetlinewidth{0.401500pt}%
\definecolor{currentstroke}{rgb}{0.901961,0.901961,0.901961}%
\pgfsetstrokecolor{currentstroke}%
\pgfsetdash{}{0pt}%
\pgfpathmoveto{\pgfqpoint{0.595000in}{2.118750in}}%
\pgfpathlineto{\pgfqpoint{3.115000in}{2.118750in}}%
\pgfusepath{stroke}%
\end{pgfscope}%
\begin{pgfscope}%
\definecolor{textcolor}{rgb}{0.101961,0.101961,0.101961}%
\pgfsetstrokecolor{textcolor}%
\pgfsetfillcolor{textcolor}%
\pgftext[x=0.414395in, y=2.084992in, left, base]{\color{textcolor}\sffamily\fontsize{7.000000}{8.400000}\selectfont 0.7}%
\end{pgfscope}%
\begin{pgfscope}%
\pgfpathrectangle{\pgfqpoint{0.595000in}{0.500000in}}{\pgfqpoint{2.520000in}{1.850000in}}%
\pgfusepath{clip}%
\pgfsetrectcap%
\pgfsetroundjoin%
\pgfsetlinewidth{0.401500pt}%
\definecolor{currentstroke}{rgb}{0.901961,0.901961,0.901961}%
\pgfsetstrokecolor{currentstroke}%
\pgfsetdash{}{0pt}%
\pgfpathmoveto{\pgfqpoint{0.595000in}{2.350000in}}%
\pgfpathlineto{\pgfqpoint{3.115000in}{2.350000in}}%
\pgfusepath{stroke}%
\end{pgfscope}%
\begin{pgfscope}%
\definecolor{textcolor}{rgb}{0.101961,0.101961,0.101961}%
\pgfsetstrokecolor{textcolor}%
\pgfsetfillcolor{textcolor}%
\pgftext[x=0.414395in, y=2.316242in, left, base]{\color{textcolor}\sffamily\fontsize{7.000000}{8.400000}\selectfont 0.8}%
\end{pgfscope}%
\begin{pgfscope}%
\definecolor{textcolor}{rgb}{0.000000,0.000000,0.000000}%
\pgfsetstrokecolor{textcolor}%
\pgfsetfillcolor{textcolor}%
\pgftext[x=0.358839in,y=1.425000in,,bottom,rotate=90.000000]{\color{textcolor}\sffamily\fontsize{7.500000}{9.000000}\selectfont accuracy, 69 long-history questions}%
\end{pgfscope}%
\begin{pgfscope}%
\pgfpathrectangle{\pgfqpoint{0.595000in}{0.500000in}}{\pgfqpoint{2.520000in}{1.850000in}}%
\pgfusepath{clip}%
\pgfsetbuttcap%
\pgfsetmiterjoin%
\definecolor{currentfill}{rgb}{0.784314,0.784314,0.784314}%
\pgfsetfillcolor{currentfill}%
\pgfsetlinewidth{0.000000pt}%
\definecolor{currentstroke}{rgb}{0.000000,0.000000,0.000000}%
\pgfsetstrokecolor{currentstroke}%
\pgfsetstrokeopacity{0.000000}%
\pgfsetdash{}{0pt}%
\pgfpathmoveto{\pgfqpoint{0.869909in}{0.500000in}}%
\pgfpathlineto{\pgfqpoint{1.259364in}{0.500000in}}%
\pgfpathlineto{\pgfqpoint{1.259364in}{0.868659in}}%
\pgfpathlineto{\pgfqpoint{0.869909in}{0.868659in}}%
\pgfpathlineto{\pgfqpoint{0.869909in}{0.500000in}}%
\pgfpathclose%
\pgfusepath{fill}%
\end{pgfscope}%
\begin{pgfscope}%
\pgfpathrectangle{\pgfqpoint{0.595000in}{0.500000in}}{\pgfqpoint{2.520000in}{1.850000in}}%
\pgfusepath{clip}%
\pgfsetbuttcap%
\pgfsetmiterjoin%
\definecolor{currentfill}{rgb}{0.156863,0.356863,0.568627}%
\pgfsetfillcolor{currentfill}%
\pgfsetlinewidth{0.000000pt}%
\definecolor{currentstroke}{rgb}{0.000000,0.000000,0.000000}%
\pgfsetstrokecolor{currentstroke}%
\pgfsetstrokeopacity{0.000000}%
\pgfsetdash{}{0pt}%
\pgfpathmoveto{\pgfqpoint{1.305182in}{0.500000in}}%
\pgfpathlineto{\pgfqpoint{1.694636in}{0.500000in}}%
\pgfpathlineto{\pgfqpoint{1.694636in}{1.446023in}}%
\pgfpathlineto{\pgfqpoint{1.305182in}{1.446023in}}%
\pgfpathlineto{\pgfqpoint{1.305182in}{0.500000in}}%
\pgfpathclose%
\pgfusepath{fill}%
\end{pgfscope}%
\begin{pgfscope}%
\pgfpathrectangle{\pgfqpoint{0.595000in}{0.500000in}}{\pgfqpoint{2.520000in}{1.850000in}}%
\pgfusepath{clip}%
\pgfsetbuttcap%
\pgfsetmiterjoin%
\definecolor{currentfill}{rgb}{0.784314,0.784314,0.784314}%
\pgfsetfillcolor{currentfill}%
\pgfsetlinewidth{0.000000pt}%
\definecolor{currentstroke}{rgb}{0.000000,0.000000,0.000000}%
\pgfsetstrokecolor{currentstroke}%
\pgfsetstrokeopacity{0.000000}%
\pgfsetdash{}{0pt}%
\pgfpathmoveto{\pgfqpoint{2.015364in}{0.500000in}}%
\pgfpathlineto{\pgfqpoint{2.404818in}{0.500000in}}%
\pgfpathlineto{\pgfqpoint{2.404818in}{1.103261in}}%
\pgfpathlineto{\pgfqpoint{2.015364in}{1.103261in}}%
\pgfpathlineto{\pgfqpoint{2.015364in}{0.500000in}}%
\pgfpathclose%
\pgfusepath{fill}%
\end{pgfscope}%
\begin{pgfscope}%
\pgfpathrectangle{\pgfqpoint{0.595000in}{0.500000in}}{\pgfqpoint{2.520000in}{1.850000in}}%
\pgfusepath{clip}%
\pgfsetbuttcap%
\pgfsetmiterjoin%
\definecolor{currentfill}{rgb}{0.156863,0.356863,0.568627}%
\pgfsetfillcolor{currentfill}%
\pgfsetlinewidth{0.000000pt}%
\definecolor{currentstroke}{rgb}{0.000000,0.000000,0.000000}%
\pgfsetstrokecolor{currentstroke}%
\pgfsetstrokeopacity{0.000000}%
\pgfsetdash{}{0pt}%
\pgfpathmoveto{\pgfqpoint{2.450636in}{0.500000in}}%
\pgfpathlineto{\pgfqpoint{2.840091in}{0.500000in}}%
\pgfpathlineto{\pgfqpoint{2.840091in}{1.740036in}}%
\pgfpathlineto{\pgfqpoint{2.450636in}{1.740036in}}%
\pgfpathlineto{\pgfqpoint{2.450636in}{0.500000in}}%
\pgfpathclose%
\pgfusepath{fill}%
\end{pgfscope}%
\begin{pgfscope}%
\pgfpathrectangle{\pgfqpoint{0.595000in}{0.500000in}}{\pgfqpoint{2.520000in}{1.850000in}}%
\pgfusepath{clip}%
\pgfsetrectcap%
\pgfsetroundjoin%
\pgfsetlinewidth{0.803000pt}%
\definecolor{currentstroke}{rgb}{0.101961,0.101961,0.101961}%
\pgfsetstrokecolor{currentstroke}%
\pgfsetdash{}{0pt}%
\pgfpathmoveto{\pgfqpoint{1.064636in}{0.711405in}}%
\pgfpathlineto{\pgfqpoint{1.064636in}{1.108985in}}%
\pgfusepath{stroke}%
\end{pgfscope}%
\begin{pgfscope}%
\pgfpathrectangle{\pgfqpoint{0.595000in}{0.500000in}}{\pgfqpoint{2.520000in}{1.850000in}}%
\pgfusepath{clip}%
\pgfsetrectcap%
\pgfsetroundjoin%
\pgfsetlinewidth{0.803000pt}%
\definecolor{currentstroke}{rgb}{0.101961,0.101961,0.101961}%
\pgfsetstrokecolor{currentstroke}%
\pgfsetdash{}{0pt}%
\pgfpathmoveto{\pgfqpoint{1.499909in}{1.190684in}}%
\pgfpathlineto{\pgfqpoint{1.499909in}{1.724487in}}%
\pgfusepath{stroke}%
\end{pgfscope}%
\begin{pgfscope}%
\pgfpathrectangle{\pgfqpoint{0.595000in}{0.500000in}}{\pgfqpoint{2.520000in}{1.850000in}}%
\pgfusepath{clip}%
\pgfsetrectcap%
\pgfsetroundjoin%
\pgfsetlinewidth{0.803000pt}%
\definecolor{currentstroke}{rgb}{0.101961,0.101961,0.101961}%
\pgfsetstrokecolor{currentstroke}%
\pgfsetdash{}{0pt}%
\pgfpathmoveto{\pgfqpoint{2.210091in}{0.897416in}}%
\pgfpathlineto{\pgfqpoint{2.210091in}{1.367432in}}%
\pgfusepath{stroke}%
\end{pgfscope}%
\begin{pgfscope}%
\pgfpathrectangle{\pgfqpoint{0.595000in}{0.500000in}}{\pgfqpoint{2.520000in}{1.850000in}}%
\pgfusepath{clip}%
\pgfsetrectcap%
\pgfsetroundjoin%
\pgfsetlinewidth{0.803000pt}%
\definecolor{currentstroke}{rgb}{0.101961,0.101961,0.101961}%
\pgfsetstrokecolor{currentstroke}%
\pgfsetdash{}{0pt}%
\pgfpathmoveto{\pgfqpoint{2.645364in}{1.470751in}}%
\pgfpathlineto{\pgfqpoint{2.645364in}{2.000484in}}%
\pgfusepath{stroke}%
\end{pgfscope}%
\begin{pgfscope}%
\pgfsetrectcap%
\pgfsetmiterjoin%
\pgfsetlinewidth{0.803000pt}%
\definecolor{currentstroke}{rgb}{0.101961,0.101961,0.101961}%
\pgfsetstrokecolor{currentstroke}%
\pgfsetdash{}{0pt}%
\pgfpathmoveto{\pgfqpoint{0.595000in}{0.500000in}}%
\pgfpathlineto{\pgfqpoint{0.595000in}{2.350000in}}%
\pgfusepath{stroke}%
\end{pgfscope}%
\begin{pgfscope}%
\pgfsetrectcap%
\pgfsetmiterjoin%
\pgfsetlinewidth{0.803000pt}%
\definecolor{currentstroke}{rgb}{0.101961,0.101961,0.101961}%
\pgfsetstrokecolor{currentstroke}%
\pgfsetdash{}{0pt}%
\pgfpathmoveto{\pgfqpoint{0.595000in}{0.500000in}}%
\pgfpathlineto{\pgfqpoint{3.115000in}{0.500000in}}%
\pgfusepath{stroke}%
\end{pgfscope}%
\begin{pgfscope}%
\definecolor{textcolor}{rgb}{0.101961,0.101961,0.101961}%
\pgfsetstrokecolor{textcolor}%
\pgfsetfillcolor{textcolor}%
\pgftext[x=1.064636in,y=1.143672in,,bottom]{\color{textcolor}\sffamily\fontsize{7.000000}{8.400000}\selectfont 16}%
\end{pgfscope}%
\begin{pgfscope}%
\definecolor{textcolor}{rgb}{0.101961,0.101961,0.101961}%
\pgfsetstrokecolor{textcolor}%
\pgfsetfillcolor{textcolor}%
\pgftext[x=1.499909in,y=1.759175in,,bottom]{\color{textcolor}\sffamily\fontsize{7.000000}{8.400000}\selectfont 41}%
\end{pgfscope}%
\begin{pgfscope}%
\definecolor{textcolor}{rgb}{0.101961,0.101961,0.101961}%
\pgfsetstrokecolor{textcolor}%
\pgfsetfillcolor{textcolor}%
\pgftext[x=2.210091in,y=1.402120in,,bottom]{\color{textcolor}\sffamily\fontsize{7.000000}{8.400000}\selectfont 26}%
\end{pgfscope}%
\begin{pgfscope}%
\definecolor{textcolor}{rgb}{0.101961,0.101961,0.101961}%
\pgfsetstrokecolor{textcolor}%
\pgfsetfillcolor{textcolor}%
\pgftext[x=2.645364in,y=2.035172in,,bottom]{\color{textcolor}\sffamily\fontsize{7.000000}{8.400000}\selectfont 54}%
\end{pgfscope}%
\begin{pgfscope}%
\pgfsetbuttcap%
\pgfsetmiterjoin%
\definecolor{currentfill}{rgb}{0.784314,0.784314,0.784314}%
\pgfsetfillcolor{currentfill}%
\pgfsetlinewidth{0.000000pt}%
\definecolor{currentstroke}{rgb}{0.000000,0.000000,0.000000}%
\pgfsetstrokecolor{currentstroke}%
\pgfsetstrokeopacity{0.000000}%
\pgfsetdash{}{0pt}%
\pgfpathmoveto{\pgfqpoint{0.643611in}{2.228472in}}%
\pgfpathlineto{\pgfqpoint{0.740833in}{2.228472in}}%
\pgfpathlineto{\pgfqpoint{0.740833in}{2.296528in}}%
\pgfpathlineto{\pgfqpoint{0.643611in}{2.296528in}}%
\pgfpathlineto{\pgfqpoint{0.643611in}{2.228472in}}%
\pgfpathclose%
\pgfusepath{fill}%
\end{pgfscope}%
\begin{pgfscope}%
\definecolor{textcolor}{rgb}{0.101961,0.101961,0.101961}%
\pgfsetstrokecolor{textcolor}%
\pgfsetfillcolor{textcolor}%
\pgftext[x=0.818611in,y=2.228472in,left,base]{\color{textcolor}\sffamily\fontsize{7.000000}{8.400000}\selectfont recent window (no history)}%
\end{pgfscope}%
\begin{pgfscope}%
\pgfsetbuttcap%
\pgfsetmiterjoin%
\definecolor{currentfill}{rgb}{0.156863,0.356863,0.568627}%
\pgfsetfillcolor{currentfill}%
\pgfsetlinewidth{0.000000pt}%
\definecolor{currentstroke}{rgb}{0.000000,0.000000,0.000000}%
\pgfsetstrokecolor{currentstroke}%
\pgfsetstrokeopacity{0.000000}%
\pgfsetdash{}{0pt}%
\pgfpathmoveto{\pgfqpoint{0.643611in}{2.087500in}}%
\pgfpathlineto{\pgfqpoint{0.740833in}{2.087500in}}%
\pgfpathlineto{\pgfqpoint{0.740833in}{2.155556in}}%
\pgfpathlineto{\pgfqpoint{0.643611in}{2.155556in}}%
\pgfpathlineto{\pgfqpoint{0.643611in}{2.087500in}}%
\pgfpathclose%
\pgfusepath{fill}%
\end{pgfscope}%
\begin{pgfscope}%
\definecolor{textcolor}{rgb}{0.101961,0.101961,0.101961}%
\pgfsetstrokecolor{textcolor}%
\pgfsetfillcolor{textcolor}%
\pgftext[x=0.818611in,y=2.087500in,left,base]{\color{textcolor}\sffamily\fontsize{7.000000}{8.400000}\selectfont full history}%
\end{pgfscope}%
\begin{pgfscope}%
\pgfsetbuttcap%
\pgfsetmiterjoin%
\definecolor{currentfill}{rgb}{1.000000,1.000000,1.000000}%
\pgfsetfillcolor{currentfill}%
\pgfsetlinewidth{0.000000pt}%
\definecolor{currentstroke}{rgb}{0.000000,0.000000,0.000000}%
\pgfsetstrokecolor{currentstroke}%
\pgfsetstrokeopacity{0.000000}%
\pgfsetdash{}{0pt}%
\pgfpathmoveto{\pgfqpoint{4.060000in}{0.500000in}}%
\pgfpathlineto{\pgfqpoint{6.860000in}{0.500000in}}%
\pgfpathlineto{\pgfqpoint{6.860000in}{2.350000in}}%
\pgfpathlineto{\pgfqpoint{4.060000in}{2.350000in}}%
\pgfpathlineto{\pgfqpoint{4.060000in}{0.500000in}}%
\pgfpathclose%
\pgfusepath{fill}%
\end{pgfscope}%
\begin{pgfscope}%
\definecolor{textcolor}{rgb}{0.101961,0.101961,0.101961}%
\pgfsetstrokecolor{textcolor}%
\pgfsetfillcolor{textcolor}%
\pgftext[x=4.823636in,y=0.451389in,,top]{\color{textcolor}\sffamily\fontsize{7.000000}{8.400000}\selectfont 11k tokens}%
\end{pgfscope}%
\begin{pgfscope}%
\definecolor{textcolor}{rgb}{0.101961,0.101961,0.101961}%
\pgfsetstrokecolor{textcolor}%
\pgfsetfillcolor{textcolor}%
\pgftext[x=6.096364in,y=0.451389in,,top]{\color{textcolor}\sffamily\fontsize{7.000000}{8.400000}\selectfont 28k tokens}%
\end{pgfscope}%
\begin{pgfscope}%
\pgfpathrectangle{\pgfqpoint{4.060000in}{0.500000in}}{\pgfqpoint{2.800000in}{1.850000in}}%
\pgfusepath{clip}%
\pgfsetrectcap%
\pgfsetroundjoin%
\pgfsetlinewidth{0.401500pt}%
\definecolor{currentstroke}{rgb}{0.901961,0.901961,0.901961}%
\pgfsetstrokecolor{currentstroke}%
\pgfsetdash{}{0pt}%
\pgfpathmoveto{\pgfqpoint{4.060000in}{0.662884in}}%
\pgfpathlineto{\pgfqpoint{6.860000in}{0.662884in}}%
\pgfusepath{stroke}%
\end{pgfscope}%
\begin{pgfscope}%
\definecolor{textcolor}{rgb}{0.101961,0.101961,0.101961}%
\pgfsetstrokecolor{textcolor}%
\pgfsetfillcolor{textcolor}%
\pgftext[x=3.879395in, y=0.629127in, left, base]{\color{textcolor}\sffamily\fontsize{7.000000}{8.400000}\selectfont 0.3}%
\end{pgfscope}%
\begin{pgfscope}%
\pgfpathrectangle{\pgfqpoint{4.060000in}{0.500000in}}{\pgfqpoint{2.800000in}{1.850000in}}%
\pgfusepath{clip}%
\pgfsetrectcap%
\pgfsetroundjoin%
\pgfsetlinewidth{0.401500pt}%
\definecolor{currentstroke}{rgb}{0.901961,0.901961,0.901961}%
\pgfsetstrokecolor{currentstroke}%
\pgfsetdash{}{0pt}%
\pgfpathmoveto{\pgfqpoint{4.060000in}{1.146547in}}%
\pgfpathlineto{\pgfqpoint{6.860000in}{1.146547in}}%
\pgfusepath{stroke}%
\end{pgfscope}%
\begin{pgfscope}%
\definecolor{textcolor}{rgb}{0.101961,0.101961,0.101961}%
\pgfsetstrokecolor{textcolor}%
\pgfsetfillcolor{textcolor}%
\pgftext[x=3.959739in, y=1.112790in, left, base]{\color{textcolor}\sffamily\fontsize{7.000000}{8.400000}\selectfont 1}%
\end{pgfscope}%
\begin{pgfscope}%
\pgfpathrectangle{\pgfqpoint{4.060000in}{0.500000in}}{\pgfqpoint{2.800000in}{1.850000in}}%
\pgfusepath{clip}%
\pgfsetrectcap%
\pgfsetroundjoin%
\pgfsetlinewidth{0.401500pt}%
\definecolor{currentstroke}{rgb}{0.901961,0.901961,0.901961}%
\pgfsetstrokecolor{currentstroke}%
\pgfsetdash{}{0pt}%
\pgfpathmoveto{\pgfqpoint{4.060000in}{1.587884in}}%
\pgfpathlineto{\pgfqpoint{6.860000in}{1.587884in}}%
\pgfusepath{stroke}%
\end{pgfscope}%
\begin{pgfscope}%
\definecolor{textcolor}{rgb}{0.101961,0.101961,0.101961}%
\pgfsetstrokecolor{textcolor}%
\pgfsetfillcolor{textcolor}%
\pgftext[x=3.959739in, y=1.554127in, left, base]{\color{textcolor}\sffamily\fontsize{7.000000}{8.400000}\selectfont 3}%
\end{pgfscope}%
\begin{pgfscope}%
\pgfpathrectangle{\pgfqpoint{4.060000in}{0.500000in}}{\pgfqpoint{2.800000in}{1.850000in}}%
\pgfusepath{clip}%
\pgfsetrectcap%
\pgfsetroundjoin%
\pgfsetlinewidth{0.401500pt}%
\definecolor{currentstroke}{rgb}{0.901961,0.901961,0.901961}%
\pgfsetstrokecolor{currentstroke}%
\pgfsetdash{}{0pt}%
\pgfpathmoveto{\pgfqpoint{4.060000in}{2.071547in}}%
\pgfpathlineto{\pgfqpoint{6.860000in}{2.071547in}}%
\pgfusepath{stroke}%
\end{pgfscope}%
\begin{pgfscope}%
\definecolor{textcolor}{rgb}{0.101961,0.101961,0.101961}%
\pgfsetstrokecolor{textcolor}%
\pgfsetfillcolor{textcolor}%
\pgftext[x=3.908089in, y=2.037790in, left, base]{\color{textcolor}\sffamily\fontsize{7.000000}{8.400000}\selectfont 10}%
\end{pgfscope}%
\begin{pgfscope}%
\pgfsetbuttcap%
\pgfsetroundjoin%
\definecolor{currentfill}{rgb}{0.000000,0.000000,0.000000}%
\pgfsetfillcolor{currentfill}%
\pgfsetlinewidth{0.602250pt}%
\definecolor{currentstroke}{rgb}{0.000000,0.000000,0.000000}%
\pgfsetstrokecolor{currentstroke}%
\pgfsetdash{}{0pt}%
\pgfsys@defobject{currentmarker}{\pgfqpoint{-0.027778in}{0.000000in}}{\pgfqpoint{-0.000000in}{0.000000in}}{%
\pgfpathmoveto{\pgfqpoint{-0.000000in}{0.000000in}}%
\pgfpathlineto{\pgfqpoint{-0.027778in}{0.000000in}}%
\pgfusepath{stroke,fill}%
}%
\begin{pgfscope}%
\pgfsys@transformshift{4.060000in}{0.500000in}%
\pgfsys@useobject{currentmarker}{}%
\end{pgfscope}%
\end{pgfscope}%
\begin{pgfscope}%
\pgfsetbuttcap%
\pgfsetroundjoin%
\definecolor{currentfill}{rgb}{0.000000,0.000000,0.000000}%
\pgfsetfillcolor{currentfill}%
\pgfsetlinewidth{0.602250pt}%
\definecolor{currentstroke}{rgb}{0.000000,0.000000,0.000000}%
\pgfsetstrokecolor{currentstroke}%
\pgfsetdash{}{0pt}%
\pgfsys@defobject{currentmarker}{\pgfqpoint{-0.027778in}{0.000000in}}{\pgfqpoint{-0.000000in}{0.000000in}}{%
\pgfpathmoveto{\pgfqpoint{-0.000000in}{0.000000in}}%
\pgfpathlineto{\pgfqpoint{-0.027778in}{0.000000in}}%
\pgfusepath{stroke,fill}%
}%
\begin{pgfscope}%
\pgfsys@transformshift{4.060000in}{0.778453in}%
\pgfsys@useobject{currentmarker}{}%
\end{pgfscope}%
\end{pgfscope}%
\begin{pgfscope}%
\pgfsetbuttcap%
\pgfsetroundjoin%
\definecolor{currentfill}{rgb}{0.000000,0.000000,0.000000}%
\pgfsetfillcolor{currentfill}%
\pgfsetlinewidth{0.602250pt}%
\definecolor{currentstroke}{rgb}{0.000000,0.000000,0.000000}%
\pgfsetstrokecolor{currentstroke}%
\pgfsetdash{}{0pt}%
\pgfsys@defobject{currentmarker}{\pgfqpoint{-0.027778in}{0.000000in}}{\pgfqpoint{-0.000000in}{0.000000in}}{%
\pgfpathmoveto{\pgfqpoint{-0.000000in}{0.000000in}}%
\pgfpathlineto{\pgfqpoint{-0.027778in}{0.000000in}}%
\pgfusepath{stroke,fill}%
}%
\begin{pgfscope}%
\pgfsys@transformshift{4.060000in}{0.868095in}%
\pgfsys@useobject{currentmarker}{}%
\end{pgfscope}%
\end{pgfscope}%
\begin{pgfscope}%
\pgfsetbuttcap%
\pgfsetroundjoin%
\definecolor{currentfill}{rgb}{0.000000,0.000000,0.000000}%
\pgfsetfillcolor{currentfill}%
\pgfsetlinewidth{0.602250pt}%
\definecolor{currentstroke}{rgb}{0.000000,0.000000,0.000000}%
\pgfsetstrokecolor{currentstroke}%
\pgfsetdash{}{0pt}%
\pgfsys@defobject{currentmarker}{\pgfqpoint{-0.027778in}{0.000000in}}{\pgfqpoint{-0.000000in}{0.000000in}}{%
\pgfpathmoveto{\pgfqpoint{-0.000000in}{0.000000in}}%
\pgfpathlineto{\pgfqpoint{-0.027778in}{0.000000in}}%
\pgfusepath{stroke,fill}%
}%
\begin{pgfscope}%
\pgfsys@transformshift{4.060000in}{0.941337in}%
\pgfsys@useobject{currentmarker}{}%
\end{pgfscope}%
\end{pgfscope}%
\begin{pgfscope}%
\pgfsetbuttcap%
\pgfsetroundjoin%
\definecolor{currentfill}{rgb}{0.000000,0.000000,0.000000}%
\pgfsetfillcolor{currentfill}%
\pgfsetlinewidth{0.602250pt}%
\definecolor{currentstroke}{rgb}{0.000000,0.000000,0.000000}%
\pgfsetstrokecolor{currentstroke}%
\pgfsetdash{}{0pt}%
\pgfsys@defobject{currentmarker}{\pgfqpoint{-0.027778in}{0.000000in}}{\pgfqpoint{-0.000000in}{0.000000in}}{%
\pgfpathmoveto{\pgfqpoint{-0.000000in}{0.000000in}}%
\pgfpathlineto{\pgfqpoint{-0.027778in}{0.000000in}}%
\pgfusepath{stroke,fill}%
}%
\begin{pgfscope}%
\pgfsys@transformshift{4.060000in}{1.003263in}%
\pgfsys@useobject{currentmarker}{}%
\end{pgfscope}%
\end{pgfscope}%
\begin{pgfscope}%
\pgfsetbuttcap%
\pgfsetroundjoin%
\definecolor{currentfill}{rgb}{0.000000,0.000000,0.000000}%
\pgfsetfillcolor{currentfill}%
\pgfsetlinewidth{0.602250pt}%
\definecolor{currentstroke}{rgb}{0.000000,0.000000,0.000000}%
\pgfsetstrokecolor{currentstroke}%
\pgfsetdash{}{0pt}%
\pgfsys@defobject{currentmarker}{\pgfqpoint{-0.027778in}{0.000000in}}{\pgfqpoint{-0.000000in}{0.000000in}}{%
\pgfpathmoveto{\pgfqpoint{-0.000000in}{0.000000in}}%
\pgfpathlineto{\pgfqpoint{-0.027778in}{0.000000in}}%
\pgfusepath{stroke,fill}%
}%
\begin{pgfscope}%
\pgfsys@transformshift{4.060000in}{1.056905in}%
\pgfsys@useobject{currentmarker}{}%
\end{pgfscope}%
\end{pgfscope}%
\begin{pgfscope}%
\pgfsetbuttcap%
\pgfsetroundjoin%
\definecolor{currentfill}{rgb}{0.000000,0.000000,0.000000}%
\pgfsetfillcolor{currentfill}%
\pgfsetlinewidth{0.602250pt}%
\definecolor{currentstroke}{rgb}{0.000000,0.000000,0.000000}%
\pgfsetstrokecolor{currentstroke}%
\pgfsetdash{}{0pt}%
\pgfsys@defobject{currentmarker}{\pgfqpoint{-0.027778in}{0.000000in}}{\pgfqpoint{-0.000000in}{0.000000in}}{%
\pgfpathmoveto{\pgfqpoint{-0.000000in}{0.000000in}}%
\pgfpathlineto{\pgfqpoint{-0.027778in}{0.000000in}}%
\pgfusepath{stroke,fill}%
}%
\begin{pgfscope}%
\pgfsys@transformshift{4.060000in}{1.104222in}%
\pgfsys@useobject{currentmarker}{}%
\end{pgfscope}%
\end{pgfscope}%
\begin{pgfscope}%
\pgfsetbuttcap%
\pgfsetroundjoin%
\definecolor{currentfill}{rgb}{0.000000,0.000000,0.000000}%
\pgfsetfillcolor{currentfill}%
\pgfsetlinewidth{0.602250pt}%
\definecolor{currentstroke}{rgb}{0.000000,0.000000,0.000000}%
\pgfsetstrokecolor{currentstroke}%
\pgfsetdash{}{0pt}%
\pgfsys@defobject{currentmarker}{\pgfqpoint{-0.027778in}{0.000000in}}{\pgfqpoint{-0.000000in}{0.000000in}}{%
\pgfpathmoveto{\pgfqpoint{-0.000000in}{0.000000in}}%
\pgfpathlineto{\pgfqpoint{-0.027778in}{0.000000in}}%
\pgfusepath{stroke,fill}%
}%
\begin{pgfscope}%
\pgfsys@transformshift{4.060000in}{1.425000in}%
\pgfsys@useobject{currentmarker}{}%
\end{pgfscope}%
\end{pgfscope}%
\begin{pgfscope}%
\pgfsetbuttcap%
\pgfsetroundjoin%
\definecolor{currentfill}{rgb}{0.000000,0.000000,0.000000}%
\pgfsetfillcolor{currentfill}%
\pgfsetlinewidth{0.602250pt}%
\definecolor{currentstroke}{rgb}{0.000000,0.000000,0.000000}%
\pgfsetstrokecolor{currentstroke}%
\pgfsetdash{}{0pt}%
\pgfsys@defobject{currentmarker}{\pgfqpoint{-0.027778in}{0.000000in}}{\pgfqpoint{-0.000000in}{0.000000in}}{%
\pgfpathmoveto{\pgfqpoint{-0.000000in}{0.000000in}}%
\pgfpathlineto{\pgfqpoint{-0.027778in}{0.000000in}}%
\pgfusepath{stroke,fill}%
}%
\begin{pgfscope}%
\pgfsys@transformshift{4.060000in}{1.703453in}%
\pgfsys@useobject{currentmarker}{}%
\end{pgfscope}%
\end{pgfscope}%
\begin{pgfscope}%
\pgfsetbuttcap%
\pgfsetroundjoin%
\definecolor{currentfill}{rgb}{0.000000,0.000000,0.000000}%
\pgfsetfillcolor{currentfill}%
\pgfsetlinewidth{0.602250pt}%
\definecolor{currentstroke}{rgb}{0.000000,0.000000,0.000000}%
\pgfsetstrokecolor{currentstroke}%
\pgfsetdash{}{0pt}%
\pgfsys@defobject{currentmarker}{\pgfqpoint{-0.027778in}{0.000000in}}{\pgfqpoint{-0.000000in}{0.000000in}}{%
\pgfpathmoveto{\pgfqpoint{-0.000000in}{0.000000in}}%
\pgfpathlineto{\pgfqpoint{-0.027778in}{0.000000in}}%
\pgfusepath{stroke,fill}%
}%
\begin{pgfscope}%
\pgfsys@transformshift{4.060000in}{1.793095in}%
\pgfsys@useobject{currentmarker}{}%
\end{pgfscope}%
\end{pgfscope}%
\begin{pgfscope}%
\pgfsetbuttcap%
\pgfsetroundjoin%
\definecolor{currentfill}{rgb}{0.000000,0.000000,0.000000}%
\pgfsetfillcolor{currentfill}%
\pgfsetlinewidth{0.602250pt}%
\definecolor{currentstroke}{rgb}{0.000000,0.000000,0.000000}%
\pgfsetstrokecolor{currentstroke}%
\pgfsetdash{}{0pt}%
\pgfsys@defobject{currentmarker}{\pgfqpoint{-0.027778in}{0.000000in}}{\pgfqpoint{-0.000000in}{0.000000in}}{%
\pgfpathmoveto{\pgfqpoint{-0.000000in}{0.000000in}}%
\pgfpathlineto{\pgfqpoint{-0.027778in}{0.000000in}}%
\pgfusepath{stroke,fill}%
}%
\begin{pgfscope}%
\pgfsys@transformshift{4.060000in}{1.866337in}%
\pgfsys@useobject{currentmarker}{}%
\end{pgfscope}%
\end{pgfscope}%
\begin{pgfscope}%
\pgfsetbuttcap%
\pgfsetroundjoin%
\definecolor{currentfill}{rgb}{0.000000,0.000000,0.000000}%
\pgfsetfillcolor{currentfill}%
\pgfsetlinewidth{0.602250pt}%
\definecolor{currentstroke}{rgb}{0.000000,0.000000,0.000000}%
\pgfsetstrokecolor{currentstroke}%
\pgfsetdash{}{0pt}%
\pgfsys@defobject{currentmarker}{\pgfqpoint{-0.027778in}{0.000000in}}{\pgfqpoint{-0.000000in}{0.000000in}}{%
\pgfpathmoveto{\pgfqpoint{-0.000000in}{0.000000in}}%
\pgfpathlineto{\pgfqpoint{-0.027778in}{0.000000in}}%
\pgfusepath{stroke,fill}%
}%
\begin{pgfscope}%
\pgfsys@transformshift{4.060000in}{1.928263in}%
\pgfsys@useobject{currentmarker}{}%
\end{pgfscope}%
\end{pgfscope}%
\begin{pgfscope}%
\pgfsetbuttcap%
\pgfsetroundjoin%
\definecolor{currentfill}{rgb}{0.000000,0.000000,0.000000}%
\pgfsetfillcolor{currentfill}%
\pgfsetlinewidth{0.602250pt}%
\definecolor{currentstroke}{rgb}{0.000000,0.000000,0.000000}%
\pgfsetstrokecolor{currentstroke}%
\pgfsetdash{}{0pt}%
\pgfsys@defobject{currentmarker}{\pgfqpoint{-0.027778in}{0.000000in}}{\pgfqpoint{-0.000000in}{0.000000in}}{%
\pgfpathmoveto{\pgfqpoint{-0.000000in}{0.000000in}}%
\pgfpathlineto{\pgfqpoint{-0.027778in}{0.000000in}}%
\pgfusepath{stroke,fill}%
}%
\begin{pgfscope}%
\pgfsys@transformshift{4.060000in}{1.981905in}%
\pgfsys@useobject{currentmarker}{}%
\end{pgfscope}%
\end{pgfscope}%
\begin{pgfscope}%
\pgfsetbuttcap%
\pgfsetroundjoin%
\definecolor{currentfill}{rgb}{0.000000,0.000000,0.000000}%
\pgfsetfillcolor{currentfill}%
\pgfsetlinewidth{0.602250pt}%
\definecolor{currentstroke}{rgb}{0.000000,0.000000,0.000000}%
\pgfsetstrokecolor{currentstroke}%
\pgfsetdash{}{0pt}%
\pgfsys@defobject{currentmarker}{\pgfqpoint{-0.027778in}{0.000000in}}{\pgfqpoint{-0.000000in}{0.000000in}}{%
\pgfpathmoveto{\pgfqpoint{-0.000000in}{0.000000in}}%
\pgfpathlineto{\pgfqpoint{-0.027778in}{0.000000in}}%
\pgfusepath{stroke,fill}%
}%
\begin{pgfscope}%
\pgfsys@transformshift{4.060000in}{2.029222in}%
\pgfsys@useobject{currentmarker}{}%
\end{pgfscope}%
\end{pgfscope}%
\begin{pgfscope}%
\pgfsetbuttcap%
\pgfsetroundjoin%
\definecolor{currentfill}{rgb}{0.000000,0.000000,0.000000}%
\pgfsetfillcolor{currentfill}%
\pgfsetlinewidth{0.602250pt}%
\definecolor{currentstroke}{rgb}{0.000000,0.000000,0.000000}%
\pgfsetstrokecolor{currentstroke}%
\pgfsetdash{}{0pt}%
\pgfsys@defobject{currentmarker}{\pgfqpoint{-0.027778in}{0.000000in}}{\pgfqpoint{-0.000000in}{0.000000in}}{%
\pgfpathmoveto{\pgfqpoint{-0.000000in}{0.000000in}}%
\pgfpathlineto{\pgfqpoint{-0.027778in}{0.000000in}}%
\pgfusepath{stroke,fill}%
}%
\begin{pgfscope}%
\pgfsys@transformshift{4.060000in}{2.350000in}%
\pgfsys@useobject{currentmarker}{}%
\end{pgfscope}%
\end{pgfscope}%
\begin{pgfscope}%
\definecolor{textcolor}{rgb}{0.000000,0.000000,0.000000}%
\pgfsetstrokecolor{textcolor}%
\pgfsetfillcolor{textcolor}%
\pgftext[x=3.823839in,y=1.425000in,,bottom,rotate=90.000000]{\color{textcolor}\sffamily\fontsize{7.500000}{9.000000}\selectfont first-token time (s)}%
\end{pgfscope}%
\begin{pgfscope}%
\pgfpathrectangle{\pgfqpoint{4.060000in}{0.500000in}}{\pgfqpoint{2.800000in}{1.850000in}}%
\pgfusepath{clip}%
\pgfsetbuttcap%
\pgfsetmiterjoin%
\definecolor{currentfill}{rgb}{0.784314,0.784314,0.784314}%
\pgfsetfillcolor{currentfill}%
\pgfsetlinewidth{0.000000pt}%
\definecolor{currentstroke}{rgb}{0.000000,0.000000,0.000000}%
\pgfsetstrokecolor{currentstroke}%
\pgfsetstrokeopacity{0.000000}%
\pgfsetdash{}{0pt}%
\pgfpathmoveto{\pgfqpoint{4.639091in}{-226.701573in}}%
\pgfpathlineto{\pgfqpoint{4.639091in}{1.505623in}}%
\pgfpathlineto{\pgfqpoint{4.320909in}{1.505623in}}%
\pgfpathlineto{\pgfqpoint{4.320909in}{-226.701573in}}%
\pgfusepath{fill}%
\end{pgfscope}%
\begin{pgfscope}%
\pgfpathrectangle{\pgfqpoint{4.060000in}{0.500000in}}{\pgfqpoint{2.800000in}{1.850000in}}%
\pgfusepath{clip}%
\pgfsetbuttcap%
\pgfsetmiterjoin%
\definecolor{currentfill}{rgb}{0.478431,0.478431,0.478431}%
\pgfsetfillcolor{currentfill}%
\pgfsetlinewidth{0.000000pt}%
\definecolor{currentstroke}{rgb}{0.000000,0.000000,0.000000}%
\pgfsetstrokecolor{currentstroke}%
\pgfsetstrokeopacity{0.000000}%
\pgfsetdash{}{0pt}%
\pgfpathmoveto{\pgfqpoint{4.982727in}{-226.701573in}}%
\pgfpathlineto{\pgfqpoint{4.982727in}{0.833286in}}%
\pgfpathlineto{\pgfqpoint{4.664545in}{0.833286in}}%
\pgfpathlineto{\pgfqpoint{4.664545in}{-226.701573in}}%
\pgfusepath{fill}%
\end{pgfscope}%
\begin{pgfscope}%
\pgfpathrectangle{\pgfqpoint{4.060000in}{0.500000in}}{\pgfqpoint{2.800000in}{1.850000in}}%
\pgfusepath{clip}%
\pgfsetbuttcap%
\pgfsetmiterjoin%
\definecolor{currentfill}{rgb}{0.156863,0.356863,0.568627}%
\pgfsetfillcolor{currentfill}%
\pgfsetlinewidth{0.000000pt}%
\definecolor{currentstroke}{rgb}{0.000000,0.000000,0.000000}%
\pgfsetstrokecolor{currentstroke}%
\pgfsetstrokeopacity{0.000000}%
\pgfsetdash{}{0pt}%
\pgfpathmoveto{\pgfqpoint{5.326364in}{-226.701573in}}%
\pgfpathlineto{\pgfqpoint{5.326364in}{0.793242in}}%
\pgfpathlineto{\pgfqpoint{5.008182in}{0.793242in}}%
\pgfpathlineto{\pgfqpoint{5.008182in}{-226.701573in}}%
\pgfusepath{fill}%
\end{pgfscope}%
\begin{pgfscope}%
\pgfpathrectangle{\pgfqpoint{4.060000in}{0.500000in}}{\pgfqpoint{2.800000in}{1.850000in}}%
\pgfusepath{clip}%
\pgfsetbuttcap%
\pgfsetmiterjoin%
\definecolor{currentfill}{rgb}{0.784314,0.784314,0.784314}%
\pgfsetfillcolor{currentfill}%
\pgfsetlinewidth{0.000000pt}%
\definecolor{currentstroke}{rgb}{0.000000,0.000000,0.000000}%
\pgfsetstrokecolor{currentstroke}%
\pgfsetstrokeopacity{0.000000}%
\pgfsetdash{}{0pt}%
\pgfpathmoveto{\pgfqpoint{5.911818in}{-226.701573in}}%
\pgfpathlineto{\pgfqpoint{5.911818in}{1.948818in}}%
\pgfpathlineto{\pgfqpoint{5.593636in}{1.948818in}}%
\pgfpathlineto{\pgfqpoint{5.593636in}{-226.701573in}}%
\pgfusepath{fill}%
\end{pgfscope}%
\begin{pgfscope}%
\pgfpathrectangle{\pgfqpoint{4.060000in}{0.500000in}}{\pgfqpoint{2.800000in}{1.850000in}}%
\pgfusepath{clip}%
\pgfsetbuttcap%
\pgfsetmiterjoin%
\definecolor{currentfill}{rgb}{0.478431,0.478431,0.478431}%
\pgfsetfillcolor{currentfill}%
\pgfsetlinewidth{0.000000pt}%
\definecolor{currentstroke}{rgb}{0.000000,0.000000,0.000000}%
\pgfsetstrokecolor{currentstroke}%
\pgfsetstrokeopacity{0.000000}%
\pgfsetdash{}{0pt}%
\pgfpathmoveto{\pgfqpoint{6.255455in}{-226.701573in}}%
\pgfpathlineto{\pgfqpoint{6.255455in}{1.146949in}}%
\pgfpathlineto{\pgfqpoint{5.937273in}{1.146949in}}%
\pgfpathlineto{\pgfqpoint{5.937273in}{-226.701573in}}%
\pgfusepath{fill}%
\end{pgfscope}%
\begin{pgfscope}%
\pgfpathrectangle{\pgfqpoint{4.060000in}{0.500000in}}{\pgfqpoint{2.800000in}{1.850000in}}%
\pgfusepath{clip}%
\pgfsetbuttcap%
\pgfsetmiterjoin%
\definecolor{currentfill}{rgb}{0.156863,0.356863,0.568627}%
\pgfsetfillcolor{currentfill}%
\pgfsetlinewidth{0.000000pt}%
\definecolor{currentstroke}{rgb}{0.000000,0.000000,0.000000}%
\pgfsetstrokecolor{currentstroke}%
\pgfsetstrokeopacity{0.000000}%
\pgfsetdash{}{0pt}%
\pgfpathmoveto{\pgfqpoint{6.599091in}{-226.701573in}}%
\pgfpathlineto{\pgfqpoint{6.599091in}{1.092215in}}%
\pgfpathlineto{\pgfqpoint{6.280909in}{1.092215in}}%
\pgfpathlineto{\pgfqpoint{6.280909in}{-226.701573in}}%
\pgfusepath{fill}%
\end{pgfscope}%
\begin{pgfscope}%
\pgfpathrectangle{\pgfqpoint{4.060000in}{0.500000in}}{\pgfqpoint{2.800000in}{1.850000in}}%
\pgfusepath{clip}%
\pgfsetrectcap%
\pgfsetroundjoin%
\pgfsetlinewidth{0.803000pt}%
\definecolor{currentstroke}{rgb}{0.101961,0.101961,0.101961}%
\pgfsetstrokecolor{currentstroke}%
\pgfsetdash{}{0pt}%
\pgfpathmoveto{\pgfqpoint{4.480000in}{1.497069in}}%
\pgfpathlineto{\pgfqpoint{4.480000in}{1.521177in}}%
\pgfusepath{stroke}%
\end{pgfscope}%
\begin{pgfscope}%
\pgfpathrectangle{\pgfqpoint{4.060000in}{0.500000in}}{\pgfqpoint{2.800000in}{1.850000in}}%
\pgfusepath{clip}%
\pgfsetrectcap%
\pgfsetroundjoin%
\pgfsetlinewidth{0.803000pt}%
\definecolor{currentstroke}{rgb}{0.101961,0.101961,0.101961}%
\pgfsetstrokecolor{currentstroke}%
\pgfsetdash{}{0pt}%
\pgfpathmoveto{\pgfqpoint{4.823636in}{0.777447in}}%
\pgfpathlineto{\pgfqpoint{4.823636in}{0.897521in}}%
\pgfusepath{stroke}%
\end{pgfscope}%
\begin{pgfscope}%
\pgfpathrectangle{\pgfqpoint{4.060000in}{0.500000in}}{\pgfqpoint{2.800000in}{1.850000in}}%
\pgfusepath{clip}%
\pgfsetrectcap%
\pgfsetroundjoin%
\pgfsetlinewidth{0.803000pt}%
\definecolor{currentstroke}{rgb}{0.101961,0.101961,0.101961}%
\pgfsetstrokecolor{currentstroke}%
\pgfsetdash{}{0pt}%
\pgfpathmoveto{\pgfqpoint{5.167273in}{0.740566in}}%
\pgfpathlineto{\pgfqpoint{5.167273in}{0.869698in}}%
\pgfusepath{stroke}%
\end{pgfscope}%
\begin{pgfscope}%
\pgfpathrectangle{\pgfqpoint{4.060000in}{0.500000in}}{\pgfqpoint{2.800000in}{1.850000in}}%
\pgfusepath{clip}%
\pgfsetrectcap%
\pgfsetroundjoin%
\pgfsetlinewidth{0.803000pt}%
\definecolor{currentstroke}{rgb}{0.101961,0.101961,0.101961}%
\pgfsetstrokecolor{currentstroke}%
\pgfsetdash{}{0pt}%
\pgfpathmoveto{\pgfqpoint{5.752727in}{1.945836in}}%
\pgfpathlineto{\pgfqpoint{5.752727in}{1.951671in}}%
\pgfusepath{stroke}%
\end{pgfscope}%
\begin{pgfscope}%
\pgfpathrectangle{\pgfqpoint{4.060000in}{0.500000in}}{\pgfqpoint{2.800000in}{1.850000in}}%
\pgfusepath{clip}%
\pgfsetrectcap%
\pgfsetroundjoin%
\pgfsetlinewidth{0.803000pt}%
\definecolor{currentstroke}{rgb}{0.101961,0.101961,0.101961}%
\pgfsetstrokecolor{currentstroke}%
\pgfsetdash{}{0pt}%
\pgfpathmoveto{\pgfqpoint{6.096364in}{1.134725in}}%
\pgfpathlineto{\pgfqpoint{6.096364in}{1.148950in}}%
\pgfusepath{stroke}%
\end{pgfscope}%
\begin{pgfscope}%
\pgfpathrectangle{\pgfqpoint{4.060000in}{0.500000in}}{\pgfqpoint{2.800000in}{1.850000in}}%
\pgfusepath{clip}%
\pgfsetrectcap%
\pgfsetroundjoin%
\pgfsetlinewidth{0.803000pt}%
\definecolor{currentstroke}{rgb}{0.101961,0.101961,0.101961}%
\pgfsetstrokecolor{currentstroke}%
\pgfsetdash{}{0pt}%
\pgfpathmoveto{\pgfqpoint{6.440000in}{1.058409in}}%
\pgfpathlineto{\pgfqpoint{6.440000in}{1.128471in}}%
\pgfusepath{stroke}%
\end{pgfscope}%
\begin{pgfscope}%
\pgfsetrectcap%
\pgfsetmiterjoin%
\pgfsetlinewidth{0.803000pt}%
\definecolor{currentstroke}{rgb}{0.101961,0.101961,0.101961}%
\pgfsetstrokecolor{currentstroke}%
\pgfsetdash{}{0pt}%
\pgfpathmoveto{\pgfqpoint{4.060000in}{0.500000in}}%
\pgfpathlineto{\pgfqpoint{4.060000in}{2.350000in}}%
\pgfusepath{stroke}%
\end{pgfscope}%
\begin{pgfscope}%
\pgfsetrectcap%
\pgfsetmiterjoin%
\pgfsetlinewidth{0.803000pt}%
\definecolor{currentstroke}{rgb}{0.101961,0.101961,0.101961}%
\pgfsetstrokecolor{currentstroke}%
\pgfsetdash{}{0pt}%
\pgfpathmoveto{\pgfqpoint{4.060000in}{0.500000in}}%
\pgfpathlineto{\pgfqpoint{6.860000in}{0.500000in}}%
\pgfusepath{stroke}%
\end{pgfscope}%
\begin{pgfscope}%
\definecolor{textcolor}{rgb}{0.101961,0.101961,0.101961}%
\pgfsetstrokecolor{textcolor}%
\pgfsetfillcolor{textcolor}%
\pgftext[x=4.480000in,y=1.566703in,,bottom]{\color{textcolor}\sffamily\fontsize{6.500000}{7.800000}\selectfont 2.44}%
\end{pgfscope}%
\begin{pgfscope}%
\definecolor{textcolor}{rgb}{0.101961,0.101961,0.101961}%
\pgfsetstrokecolor{textcolor}%
\pgfsetfillcolor{textcolor}%
\pgftext[x=4.823636in,y=0.943048in,,bottom]{\color{textcolor}\sffamily\fontsize{6.500000}{7.800000}\selectfont 0.46}%
\end{pgfscope}%
\begin{pgfscope}%
\definecolor{textcolor}{rgb}{0.101961,0.101961,0.101961}%
\pgfsetstrokecolor{textcolor}%
\pgfsetfillcolor{textcolor}%
\pgftext[x=5.167273in,y=0.915225in,,bottom]{\color{textcolor}\sffamily\fontsize{6.500000}{7.800000}\selectfont 0.42}%
\end{pgfscope}%
\begin{pgfscope}%
\definecolor{textcolor}{rgb}{0.101961,0.101961,0.101961}%
\pgfsetstrokecolor{textcolor}%
\pgfsetfillcolor{textcolor}%
\pgftext[x=5.752727in,y=1.997198in,,bottom]{\color{textcolor}\sffamily\fontsize{6.500000}{7.800000}\selectfont 7.37}%
\end{pgfscope}%
\begin{pgfscope}%
\definecolor{textcolor}{rgb}{0.101961,0.101961,0.101961}%
\pgfsetstrokecolor{textcolor}%
\pgfsetfillcolor{textcolor}%
\pgftext[x=6.096364in,y=1.194477in,,bottom]{\color{textcolor}\sffamily\fontsize{6.500000}{7.800000}\selectfont 1.00}%
\end{pgfscope}%
\begin{pgfscope}%
\definecolor{textcolor}{rgb}{0.101961,0.101961,0.101961}%
\pgfsetstrokecolor{textcolor}%
\pgfsetfillcolor{textcolor}%
\pgftext[x=6.440000in,y=1.173997in,,bottom]{\color{textcolor}\sffamily\fontsize{6.500000}{7.800000}\selectfont 0.87}%
\end{pgfscope}%
\begin{pgfscope}%
\pgfsetbuttcap%
\pgfsetmiterjoin%
\definecolor{currentfill}{rgb}{0.784314,0.784314,0.784314}%
\pgfsetfillcolor{currentfill}%
\pgfsetlinewidth{0.000000pt}%
\definecolor{currentstroke}{rgb}{0.000000,0.000000,0.000000}%
\pgfsetstrokecolor{currentstroke}%
\pgfsetstrokeopacity{0.000000}%
\pgfsetdash{}{0pt}%
\pgfpathmoveto{\pgfqpoint{4.108611in}{2.233333in}}%
\pgfpathlineto{\pgfqpoint{4.205833in}{2.233333in}}%
\pgfpathlineto{\pgfqpoint{4.205833in}{2.301389in}}%
\pgfpathlineto{\pgfqpoint{4.108611in}{2.301389in}}%
\pgfpathlineto{\pgfqpoint{4.108611in}{2.233333in}}%
\pgfpathclose%
\pgfusepath{fill}%
\end{pgfscope}%
\begin{pgfscope}%
\definecolor{textcolor}{rgb}{0.101961,0.101961,0.101961}%
\pgfsetstrokecolor{textcolor}%
\pgfsetfillcolor{textcolor}%
\pgftext[x=4.283611in,y=2.233333in,left,base]{\color{textcolor}\sffamily\fontsize{7.000000}{8.400000}\selectfont COLD}%
\end{pgfscope}%
\begin{pgfscope}%
\pgfsetbuttcap%
\pgfsetmiterjoin%
\definecolor{currentfill}{rgb}{0.478431,0.478431,0.478431}%
\pgfsetfillcolor{currentfill}%
\pgfsetlinewidth{0.000000pt}%
\definecolor{currentstroke}{rgb}{0.000000,0.000000,0.000000}%
\pgfsetstrokecolor{currentstroke}%
\pgfsetstrokeopacity{0.000000}%
\pgfsetdash{}{0pt}%
\pgfpathmoveto{\pgfqpoint{4.634493in}{2.233333in}}%
\pgfpathlineto{\pgfqpoint{4.731715in}{2.233333in}}%
\pgfpathlineto{\pgfqpoint{4.731715in}{2.301389in}}%
\pgfpathlineto{\pgfqpoint{4.634493in}{2.301389in}}%
\pgfpathlineto{\pgfqpoint{4.634493in}{2.233333in}}%
\pgfpathclose%
\pgfusepath{fill}%
\end{pgfscope}%
\begin{pgfscope}%
\definecolor{textcolor}{rgb}{0.101961,0.101961,0.101961}%
\pgfsetstrokecolor{textcolor}%
\pgfsetfillcolor{textcolor}%
\pgftext[x=4.809493in,y=2.233333in,left,base]{\color{textcolor}\sffamily\fontsize{7.000000}{8.400000}\selectfont HOST-WARM}%
\end{pgfscope}%
\begin{pgfscope}%
\pgfsetbuttcap%
\pgfsetmiterjoin%
\definecolor{currentfill}{rgb}{0.156863,0.356863,0.568627}%
\pgfsetfillcolor{currentfill}%
\pgfsetlinewidth{0.000000pt}%
\definecolor{currentstroke}{rgb}{0.000000,0.000000,0.000000}%
\pgfsetstrokecolor{currentstroke}%
\pgfsetstrokeopacity{0.000000}%
\pgfsetdash{}{0pt}%
\pgfpathmoveto{\pgfqpoint{5.513653in}{2.233333in}}%
\pgfpathlineto{\pgfqpoint{5.610875in}{2.233333in}}%
\pgfpathlineto{\pgfqpoint{5.610875in}{2.301389in}}%
\pgfpathlineto{\pgfqpoint{5.513653in}{2.301389in}}%
\pgfpathlineto{\pgfqpoint{5.513653in}{2.233333in}}%
\pgfpathclose%
\pgfusepath{fill}%
\end{pgfscope}%
\begin{pgfscope}%
\definecolor{textcolor}{rgb}{0.101961,0.101961,0.101961}%
\pgfsetstrokecolor{textcolor}%
\pgfsetfillcolor{textcolor}%
\pgftext[x=5.688653in,y=2.233333in,left,base]{\color{textcolor}\sffamily\fontsize{7.000000}{8.400000}\selectfont GPU-WARM}%
\end{pgfscope}%
\begin{pgfscope}%
\definecolor{textcolor}{rgb}{0.101961,0.101961,0.101961}%
\pgfsetstrokecolor{textcolor}%
\pgfsetfillcolor{textcolor}%
\pgftext[x=0.021000in,y=2.462500in,left,top]{\color{textcolor}\sffamily\fontsize{7.500000}{9.000000}\bfseries\selectfont (a)}%
\end{pgfscope}%
\begin{pgfscope}%
\definecolor{textcolor}{rgb}{0.101961,0.101961,0.101961}%
\pgfsetstrokecolor{textcolor}%
\pgfsetfillcolor{textcolor}%
\pgftext[x=3.500000in,y=2.462500in,left,top]{\color{textcolor}\sffamily\fontsize{7.500000}{9.000000}\bfseries\selectfont (b)}%
\end{pgfscope}%
\end{pgfpicture}%
\makeatother%
\endgroup%
}
\caption{[PLACEHOLDER: committed accuracy and cost-profile measurements, no COVEX data, not FINAL] (a) Accuracy on the 69 EVAL long-history-required questions with the recent window only versus the full history, for the on-device 4B model and the cloud 32B model (Wilson 95 percent bands, each question one trial). (b) First-token time of an isolated 32B escalation when the prefix is cold, in host memory, or on the GPU, at about 11k and 28k prompt tokens (bars are medians, whiskers min to max over 12 and 6 requests). The cold series comes from the GPU-only server and the two warm series from the host-offload server of the same study. Both panels are committed cost-profile and accuracy measurements, not COVEX system results.}
\label{fig:F16_history_and_warm}
\end{figure*}
\clearpage
\begin{figure}[t]
\centering
{\fontfamily{LinuxBiolinumT}\selectfont
%% Creator: Matplotlib, PGF backend
%%
%% To include the figure in your LaTeX document, write
%%   \input{<filename>.pgf}
%%
%% Make sure the required packages are loaded in your preamble
%%   \usepackage{pgf}
%%
%% Also ensure that all the required font packages are loaded; for instance,
%% the lmodern package is sometimes necessary when using math font.
%%   \usepackage{lmodern}
%%
%% Figures using additional raster images can only be included by \input if
%% they are in the same directory as the main LaTeX file. For loading figures
%% from other directories you can use the `import` package
%%   \usepackage{import}
%%
%% and then include the figures with
%%   \import{<path to file>}{<filename>.pgf}
%%
%% Matplotlib used the following preamble
%%
\begingroup%
\makeatletter%
\begin{pgfpicture}%
\pgfpathrectangle{\pgfpointorigin}{\pgfqpoint{3.330000in}{2.400000in}}%
\pgfusepath{use as bounding box, clip}%
\begin{pgfscope}%
\pgfsetbuttcap%
\pgfsetmiterjoin%
\definecolor{currentfill}{rgb}{1.000000,1.000000,1.000000}%
\pgfsetfillcolor{currentfill}%
\pgfsetlinewidth{0.000000pt}%
\definecolor{currentstroke}{rgb}{1.000000,1.000000,1.000000}%
\pgfsetstrokecolor{currentstroke}%
\pgfsetdash{}{0pt}%
\pgfpathmoveto{\pgfqpoint{0.000000in}{0.000000in}}%
\pgfpathlineto{\pgfqpoint{3.330000in}{0.000000in}}%
\pgfpathlineto{\pgfqpoint{3.330000in}{2.400000in}}%
\pgfpathlineto{\pgfqpoint{0.000000in}{2.400000in}}%
\pgfpathlineto{\pgfqpoint{0.000000in}{0.000000in}}%
\pgfpathclose%
\pgfusepath{fill}%
\end{pgfscope}%
\begin{pgfscope}%
\pgfsetbuttcap%
\pgfsetmiterjoin%
\definecolor{currentfill}{rgb}{1.000000,1.000000,1.000000}%
\pgfsetfillcolor{currentfill}%
\pgfsetlinewidth{0.000000pt}%
\definecolor{currentstroke}{rgb}{0.000000,0.000000,0.000000}%
\pgfsetstrokecolor{currentstroke}%
\pgfsetstrokeopacity{0.000000}%
\pgfsetdash{}{0pt}%
\pgfpathmoveto{\pgfqpoint{0.466200in}{0.408000in}}%
\pgfpathlineto{\pgfqpoint{3.230100in}{0.408000in}}%
\pgfpathlineto{\pgfqpoint{3.230100in}{2.088000in}}%
\pgfpathlineto{\pgfqpoint{0.466200in}{2.088000in}}%
\pgfpathlineto{\pgfqpoint{0.466200in}{0.408000in}}%
\pgfpathclose%
\pgfusepath{fill}%
\end{pgfscope}%
\begin{pgfscope}%
\definecolor{textcolor}{rgb}{0.000000,0.000000,0.000000}%
\pgfsetstrokecolor{textcolor}%
\pgfsetfillcolor{textcolor}%
\pgftext[x=1.848150in,y=0.303833in,,top]{\color{textcolor}\sffamily\fontsize{7.500000}{9.000000}\selectfont deployment time (provisional axis)}%
\end{pgfscope}%
\begin{pgfscope}%
\pgfpathrectangle{\pgfqpoint{0.466200in}{0.408000in}}{\pgfqpoint{2.763900in}{1.680000in}}%
\pgfusepath{clip}%
\pgfsetrectcap%
\pgfsetroundjoin%
\pgfsetlinewidth{0.401500pt}%
\definecolor{currentstroke}{rgb}{0.901961,0.901961,0.901961}%
\pgfsetstrokecolor{currentstroke}%
\pgfsetdash{}{0pt}%
\pgfpathmoveto{\pgfqpoint{0.466200in}{0.408000in}}%
\pgfpathlineto{\pgfqpoint{3.230100in}{0.408000in}}%
\pgfusepath{stroke}%
\end{pgfscope}%
\begin{pgfscope}%
\definecolor{textcolor}{rgb}{0.101961,0.101961,0.101961}%
\pgfsetstrokecolor{textcolor}%
\pgfsetfillcolor{textcolor}%
\pgftext[x=0.285595in, y=0.374242in, left, base]{\color{textcolor}\sffamily\fontsize{7.000000}{8.400000}\selectfont 0.0}%
\end{pgfscope}%
\begin{pgfscope}%
\pgfpathrectangle{\pgfqpoint{0.466200in}{0.408000in}}{\pgfqpoint{2.763900in}{1.680000in}}%
\pgfusepath{clip}%
\pgfsetrectcap%
\pgfsetroundjoin%
\pgfsetlinewidth{0.401500pt}%
\definecolor{currentstroke}{rgb}{0.901961,0.901961,0.901961}%
\pgfsetstrokecolor{currentstroke}%
\pgfsetdash{}{0pt}%
\pgfpathmoveto{\pgfqpoint{0.466200in}{0.744000in}}%
\pgfpathlineto{\pgfqpoint{3.230100in}{0.744000in}}%
\pgfusepath{stroke}%
\end{pgfscope}%
\begin{pgfscope}%
\definecolor{textcolor}{rgb}{0.101961,0.101961,0.101961}%
\pgfsetstrokecolor{textcolor}%
\pgfsetfillcolor{textcolor}%
\pgftext[x=0.285595in, y=0.710242in, left, base]{\color{textcolor}\sffamily\fontsize{7.000000}{8.400000}\selectfont 0.2}%
\end{pgfscope}%
\begin{pgfscope}%
\pgfpathrectangle{\pgfqpoint{0.466200in}{0.408000in}}{\pgfqpoint{2.763900in}{1.680000in}}%
\pgfusepath{clip}%
\pgfsetrectcap%
\pgfsetroundjoin%
\pgfsetlinewidth{0.401500pt}%
\definecolor{currentstroke}{rgb}{0.901961,0.901961,0.901961}%
\pgfsetstrokecolor{currentstroke}%
\pgfsetdash{}{0pt}%
\pgfpathmoveto{\pgfqpoint{0.466200in}{1.080000in}}%
\pgfpathlineto{\pgfqpoint{3.230100in}{1.080000in}}%
\pgfusepath{stroke}%
\end{pgfscope}%
\begin{pgfscope}%
\definecolor{textcolor}{rgb}{0.101961,0.101961,0.101961}%
\pgfsetstrokecolor{textcolor}%
\pgfsetfillcolor{textcolor}%
\pgftext[x=0.285595in, y=1.046242in, left, base]{\color{textcolor}\sffamily\fontsize{7.000000}{8.400000}\selectfont 0.4}%
\end{pgfscope}%
\begin{pgfscope}%
\pgfpathrectangle{\pgfqpoint{0.466200in}{0.408000in}}{\pgfqpoint{2.763900in}{1.680000in}}%
\pgfusepath{clip}%
\pgfsetrectcap%
\pgfsetroundjoin%
\pgfsetlinewidth{0.401500pt}%
\definecolor{currentstroke}{rgb}{0.901961,0.901961,0.901961}%
\pgfsetstrokecolor{currentstroke}%
\pgfsetdash{}{0pt}%
\pgfpathmoveto{\pgfqpoint{0.466200in}{1.416000in}}%
\pgfpathlineto{\pgfqpoint{3.230100in}{1.416000in}}%
\pgfusepath{stroke}%
\end{pgfscope}%
\begin{pgfscope}%
\definecolor{textcolor}{rgb}{0.101961,0.101961,0.101961}%
\pgfsetstrokecolor{textcolor}%
\pgfsetfillcolor{textcolor}%
\pgftext[x=0.285595in, y=1.382242in, left, base]{\color{textcolor}\sffamily\fontsize{7.000000}{8.400000}\selectfont 0.6}%
\end{pgfscope}%
\begin{pgfscope}%
\pgfpathrectangle{\pgfqpoint{0.466200in}{0.408000in}}{\pgfqpoint{2.763900in}{1.680000in}}%
\pgfusepath{clip}%
\pgfsetrectcap%
\pgfsetroundjoin%
\pgfsetlinewidth{0.401500pt}%
\definecolor{currentstroke}{rgb}{0.901961,0.901961,0.901961}%
\pgfsetstrokecolor{currentstroke}%
\pgfsetdash{}{0pt}%
\pgfpathmoveto{\pgfqpoint{0.466200in}{1.752000in}}%
\pgfpathlineto{\pgfqpoint{3.230100in}{1.752000in}}%
\pgfusepath{stroke}%
\end{pgfscope}%
\begin{pgfscope}%
\definecolor{textcolor}{rgb}{0.101961,0.101961,0.101961}%
\pgfsetstrokecolor{textcolor}%
\pgfsetfillcolor{textcolor}%
\pgftext[x=0.285595in, y=1.718242in, left, base]{\color{textcolor}\sffamily\fontsize{7.000000}{8.400000}\selectfont 0.8}%
\end{pgfscope}%
\begin{pgfscope}%
\pgfpathrectangle{\pgfqpoint{0.466200in}{0.408000in}}{\pgfqpoint{2.763900in}{1.680000in}}%
\pgfusepath{clip}%
\pgfsetrectcap%
\pgfsetroundjoin%
\pgfsetlinewidth{0.401500pt}%
\definecolor{currentstroke}{rgb}{0.901961,0.901961,0.901961}%
\pgfsetstrokecolor{currentstroke}%
\pgfsetdash{}{0pt}%
\pgfpathmoveto{\pgfqpoint{0.466200in}{2.088000in}}%
\pgfpathlineto{\pgfqpoint{3.230100in}{2.088000in}}%
\pgfusepath{stroke}%
\end{pgfscope}%
\begin{pgfscope}%
\definecolor{textcolor}{rgb}{0.101961,0.101961,0.101961}%
\pgfsetstrokecolor{textcolor}%
\pgfsetfillcolor{textcolor}%
\pgftext[x=0.285595in, y=2.054242in, left, base]{\color{textcolor}\sffamily\fontsize{7.000000}{8.400000}\selectfont 1.0}%
\end{pgfscope}%
\begin{pgfscope}%
\definecolor{textcolor}{rgb}{0.000000,0.000000,0.000000}%
\pgfsetstrokecolor{textcolor}%
\pgfsetfillcolor{textcolor}%
\pgftext[x=0.230039in,y=1.248000in,,bottom,rotate=90.000000]{\color{textcolor}\sffamily\fontsize{7.500000}{9.000000}\selectfont on-time rate (provisional axis)}%
\end{pgfscope}%
\begin{pgfscope}%
\pgfsetrectcap%
\pgfsetmiterjoin%
\pgfsetlinewidth{0.803000pt}%
\definecolor{currentstroke}{rgb}{0.101961,0.101961,0.101961}%
\pgfsetstrokecolor{currentstroke}%
\pgfsetdash{}{0pt}%
\pgfpathmoveto{\pgfqpoint{0.466200in}{0.408000in}}%
\pgfpathlineto{\pgfqpoint{0.466200in}{2.088000in}}%
\pgfusepath{stroke}%
\end{pgfscope}%
\begin{pgfscope}%
\pgfsetrectcap%
\pgfsetmiterjoin%
\pgfsetlinewidth{0.803000pt}%
\definecolor{currentstroke}{rgb}{0.101961,0.101961,0.101961}%
\pgfsetstrokecolor{currentstroke}%
\pgfsetdash{}{0pt}%
\pgfpathmoveto{\pgfqpoint{0.466200in}{0.408000in}}%
\pgfpathlineto{\pgfqpoint{3.230100in}{0.408000in}}%
\pgfusepath{stroke}%
\end{pgfscope}%
\begin{pgfscope}%
\pgfsetrectcap%
\pgfsetroundjoin%
\pgfsetlinewidth{1.204500pt}%
\definecolor{currentstroke}{rgb}{0.156863,0.356863,0.568627}%
\pgfsetstrokecolor{currentstroke}%
\pgfsetdash{}{0pt}%
\pgfpathmoveto{\pgfqpoint{0.514811in}{2.005361in}}%
\pgfpathlineto{\pgfqpoint{0.582867in}{2.005361in}}%
\pgfpathlineto{\pgfqpoint{0.650922in}{2.005361in}}%
\pgfusepath{stroke}%
\end{pgfscope}%
\begin{pgfscope}%
\definecolor{textcolor}{rgb}{0.101961,0.101961,0.101961}%
\pgfsetstrokecolor{textcolor}%
\pgfsetfillcolor{textcolor}%
\pgftext[x=0.728700in,y=1.971333in,left,base]{\color{textcolor}\sffamily\fontsize{7.000000}{8.400000}\selectfont COVEX}%
\end{pgfscope}%
\begin{pgfscope}%
\pgfsetrectcap%
\pgfsetroundjoin%
\pgfsetlinewidth{1.204500pt}%
\definecolor{currentstroke}{rgb}{0.478431,0.478431,0.478431}%
\pgfsetstrokecolor{currentstroke}%
\pgfsetdash{}{0pt}%
\pgfpathmoveto{\pgfqpoint{0.514811in}{1.869790in}}%
\pgfpathlineto{\pgfqpoint{0.582867in}{1.869790in}}%
\pgfpathlineto{\pgfqpoint{0.650922in}{1.869790in}}%
\pgfusepath{stroke}%
\end{pgfscope}%
\begin{pgfscope}%
\definecolor{textcolor}{rgb}{0.101961,0.101961,0.101961}%
\pgfsetstrokecolor{textcolor}%
\pgfsetfillcolor{textcolor}%
\pgftext[x=0.728700in,y=1.835762in,left,base]{\color{textcolor}\sffamily\fontsize{7.000000}{8.400000}\selectfont stateful baseline}%
\end{pgfscope}%
\begin{pgfscope}%
\pgfsetrectcap%
\pgfsetroundjoin%
\pgfsetlinewidth{1.204500pt}%
\definecolor{currentstroke}{rgb}{0.784314,0.784314,0.784314}%
\pgfsetstrokecolor{currentstroke}%
\pgfsetdash{}{0pt}%
\pgfpathmoveto{\pgfqpoint{0.514811in}{1.734219in}}%
\pgfpathlineto{\pgfqpoint{0.582867in}{1.734219in}}%
\pgfpathlineto{\pgfqpoint{0.650922in}{1.734219in}}%
\pgfusepath{stroke}%
\end{pgfscope}%
\begin{pgfscope}%
\definecolor{textcolor}{rgb}{0.101961,0.101961,0.101961}%
\pgfsetstrokecolor{textcolor}%
\pgfsetfillcolor{textcolor}%
\pgftext[x=0.728700in,y=1.700192in,left,base]{\color{textcolor}\sffamily\fontsize{7.000000}{8.400000}\selectfont stateless baseline}%
\end{pgfscope}%
\begin{pgfscope}%
\definecolor{textcolor}{rgb}{0.549020,0.549020,0.549020}%
\pgfsetstrokecolor{textcolor}%
\pgfsetfillcolor{textcolor}%
\pgftext[x=1.848150in,y=1.200000in,,]{\color{textcolor}\sffamily\fontsize{7.000000}{8.400000}\selectfont NO DATA YET: real-deployment data}%
\end{pgfscope}%
\end{pgfpicture}%
\makeatother%
\endgroup%
}
\caption{[PLACEHOLDER: NO DATA, no real-deployment data exists] Real deployment. No real-deployment data exists yet; the panel shows empty axes and legend only.}
\label{fig:F17_real_deployment}
\end{figure}
\clearpage
\begin{figure*}[t]
\centering
{\fontfamily{LinuxBiolinumT}\selectfont
%% Creator: Matplotlib, PGF backend
%%
%% To include the figure in your LaTeX document, write
%%   \input{<filename>.pgf}
%%
%% Make sure the required packages are loaded in your preamble
%%   \usepackage{pgf}
%%
%% Also ensure that all the required font packages are loaded; for instance,
%% the lmodern package is sometimes necessary when using math font.
%%   \usepackage{lmodern}
%%
%% Figures using additional raster images can only be included by \input if
%% they are in the same directory as the main LaTeX file. For loading figures
%% from other directories you can use the `import` package
%%   \usepackage{import}
%%
%% and then include the figures with
%%   \import{<path to file>}{<filename>.pgf}
%%
%% Matplotlib used the following preamble
%%
\begingroup%
\makeatletter%
\begin{pgfpicture}%
\pgfpathrectangle{\pgfpointorigin}{\pgfqpoint{7.000000in}{3.500000in}}%
\pgfusepath{use as bounding box, clip}%
\begin{pgfscope}%
\pgfsetbuttcap%
\pgfsetmiterjoin%
\definecolor{currentfill}{rgb}{1.000000,1.000000,1.000000}%
\pgfsetfillcolor{currentfill}%
\pgfsetlinewidth{0.000000pt}%
\definecolor{currentstroke}{rgb}{1.000000,1.000000,1.000000}%
\pgfsetstrokecolor{currentstroke}%
\pgfsetdash{}{0pt}%
\pgfpathmoveto{\pgfqpoint{0.000000in}{0.000000in}}%
\pgfpathlineto{\pgfqpoint{7.000000in}{0.000000in}}%
\pgfpathlineto{\pgfqpoint{7.000000in}{3.500000in}}%
\pgfpathlineto{\pgfqpoint{0.000000in}{3.500000in}}%
\pgfpathlineto{\pgfqpoint{0.000000in}{0.000000in}}%
\pgfpathclose%
\pgfusepath{fill}%
\end{pgfscope}%
\begin{pgfscope}%
\pgfsetbuttcap%
\pgfsetmiterjoin%
\definecolor{currentfill}{rgb}{1.000000,1.000000,1.000000}%
\pgfsetfillcolor{currentfill}%
\pgfsetlinewidth{0.000000pt}%
\definecolor{currentstroke}{rgb}{0.000000,0.000000,0.000000}%
\pgfsetstrokecolor{currentstroke}%
\pgfsetstrokeopacity{0.000000}%
\pgfsetdash{}{0pt}%
\pgfpathmoveto{\pgfqpoint{1.505000in}{0.525000in}}%
\pgfpathlineto{\pgfqpoint{2.798250in}{0.525000in}}%
\pgfpathlineto{\pgfqpoint{2.798250in}{3.255000in}}%
\pgfpathlineto{\pgfqpoint{1.505000in}{3.255000in}}%
\pgfpathlineto{\pgfqpoint{1.505000in}{0.525000in}}%
\pgfpathclose%
\pgfusepath{fill}%
\end{pgfscope}%
\begin{pgfscope}%
\pgfpathrectangle{\pgfqpoint{1.505000in}{0.525000in}}{\pgfqpoint{1.293250in}{2.730000in}}%
\pgfusepath{clip}%
\pgfsetrectcap%
\pgfsetroundjoin%
\pgfsetlinewidth{0.401500pt}%
\definecolor{currentstroke}{rgb}{0.901961,0.901961,0.901961}%
\pgfsetstrokecolor{currentstroke}%
\pgfsetdash{}{0pt}%
\pgfpathmoveto{\pgfqpoint{1.505000in}{0.525000in}}%
\pgfpathlineto{\pgfqpoint{1.505000in}{3.255000in}}%
\pgfusepath{stroke}%
\end{pgfscope}%
\begin{pgfscope}%
\pgfsetbuttcap%
\pgfsetroundjoin%
\definecolor{currentfill}{rgb}{0.101961,0.101961,0.101961}%
\pgfsetfillcolor{currentfill}%
\pgfsetlinewidth{0.803000pt}%
\definecolor{currentstroke}{rgb}{0.101961,0.101961,0.101961}%
\pgfsetstrokecolor{currentstroke}%
\pgfsetdash{}{0pt}%
\pgfsys@defobject{currentmarker}{\pgfqpoint{0.000000in}{-0.027778in}}{\pgfqpoint{0.000000in}{0.000000in}}{%
\pgfpathmoveto{\pgfqpoint{0.000000in}{0.000000in}}%
\pgfpathlineto{\pgfqpoint{0.000000in}{-0.027778in}}%
\pgfusepath{stroke,fill}%
}%
\begin{pgfscope}%
\pgfsys@transformshift{1.505000in}{0.525000in}%
\pgfsys@useobject{currentmarker}{}%
\end{pgfscope}%
\end{pgfscope}%
\begin{pgfscope}%
\definecolor{textcolor}{rgb}{0.101961,0.101961,0.101961}%
\pgfsetstrokecolor{textcolor}%
\pgfsetfillcolor{textcolor}%
\pgftext[x=1.505000in,y=0.448611in,,top]{\color{textcolor}\sffamily\fontsize{7.000000}{8.400000}\selectfont 0}%
\end{pgfscope}%
\begin{pgfscope}%
\pgfpathrectangle{\pgfqpoint{1.505000in}{0.525000in}}{\pgfqpoint{1.293250in}{2.730000in}}%
\pgfusepath{clip}%
\pgfsetrectcap%
\pgfsetroundjoin%
\pgfsetlinewidth{0.401500pt}%
\definecolor{currentstroke}{rgb}{0.901961,0.901961,0.901961}%
\pgfsetstrokecolor{currentstroke}%
\pgfsetdash{}{0pt}%
\pgfpathmoveto{\pgfqpoint{2.022300in}{0.525000in}}%
\pgfpathlineto{\pgfqpoint{2.022300in}{3.255000in}}%
\pgfusepath{stroke}%
\end{pgfscope}%
\begin{pgfscope}%
\pgfsetbuttcap%
\pgfsetroundjoin%
\definecolor{currentfill}{rgb}{0.101961,0.101961,0.101961}%
\pgfsetfillcolor{currentfill}%
\pgfsetlinewidth{0.803000pt}%
\definecolor{currentstroke}{rgb}{0.101961,0.101961,0.101961}%
\pgfsetstrokecolor{currentstroke}%
\pgfsetdash{}{0pt}%
\pgfsys@defobject{currentmarker}{\pgfqpoint{0.000000in}{-0.027778in}}{\pgfqpoint{0.000000in}{0.000000in}}{%
\pgfpathmoveto{\pgfqpoint{0.000000in}{0.000000in}}%
\pgfpathlineto{\pgfqpoint{0.000000in}{-0.027778in}}%
\pgfusepath{stroke,fill}%
}%
\begin{pgfscope}%
\pgfsys@transformshift{2.022300in}{0.525000in}%
\pgfsys@useobject{currentmarker}{}%
\end{pgfscope}%
\end{pgfscope}%
\begin{pgfscope}%
\definecolor{textcolor}{rgb}{0.101961,0.101961,0.101961}%
\pgfsetstrokecolor{textcolor}%
\pgfsetfillcolor{textcolor}%
\pgftext[x=2.022300in,y=0.448611in,,top]{\color{textcolor}\sffamily\fontsize{7.000000}{8.400000}\selectfont 0.5}%
\end{pgfscope}%
\begin{pgfscope}%
\pgfpathrectangle{\pgfqpoint{1.505000in}{0.525000in}}{\pgfqpoint{1.293250in}{2.730000in}}%
\pgfusepath{clip}%
\pgfsetrectcap%
\pgfsetroundjoin%
\pgfsetlinewidth{0.401500pt}%
\definecolor{currentstroke}{rgb}{0.901961,0.901961,0.901961}%
\pgfsetstrokecolor{currentstroke}%
\pgfsetdash{}{0pt}%
\pgfpathmoveto{\pgfqpoint{2.539600in}{0.525000in}}%
\pgfpathlineto{\pgfqpoint{2.539600in}{3.255000in}}%
\pgfusepath{stroke}%
\end{pgfscope}%
\begin{pgfscope}%
\pgfsetbuttcap%
\pgfsetroundjoin%
\definecolor{currentfill}{rgb}{0.101961,0.101961,0.101961}%
\pgfsetfillcolor{currentfill}%
\pgfsetlinewidth{0.803000pt}%
\definecolor{currentstroke}{rgb}{0.101961,0.101961,0.101961}%
\pgfsetstrokecolor{currentstroke}%
\pgfsetdash{}{0pt}%
\pgfsys@defobject{currentmarker}{\pgfqpoint{0.000000in}{-0.027778in}}{\pgfqpoint{0.000000in}{0.000000in}}{%
\pgfpathmoveto{\pgfqpoint{0.000000in}{0.000000in}}%
\pgfpathlineto{\pgfqpoint{0.000000in}{-0.027778in}}%
\pgfusepath{stroke,fill}%
}%
\begin{pgfscope}%
\pgfsys@transformshift{2.539600in}{0.525000in}%
\pgfsys@useobject{currentmarker}{}%
\end{pgfscope}%
\end{pgfscope}%
\begin{pgfscope}%
\definecolor{textcolor}{rgb}{0.101961,0.101961,0.101961}%
\pgfsetstrokecolor{textcolor}%
\pgfsetfillcolor{textcolor}%
\pgftext[x=2.539600in,y=0.448611in,,top]{\color{textcolor}\sffamily\fontsize{7.000000}{8.400000}\selectfont 1}%
\end{pgfscope}%
\begin{pgfscope}%
\definecolor{textcolor}{rgb}{0.101961,0.101961,0.101961}%
\pgfsetstrokecolor{textcolor}%
\pgfsetfillcolor{textcolor}%
\pgftext[x=2.151625in,y=0.334414in,,top]{\color{textcolor}\sffamily\fontsize{7.500000}{9.000000}\selectfont Robot}%
\end{pgfscope}%
\begin{pgfscope}%
\pgfpathrectangle{\pgfqpoint{1.505000in}{0.525000in}}{\pgfqpoint{1.293250in}{2.730000in}}%
\pgfusepath{clip}%
\pgfsetbuttcap%
\pgfsetmiterjoin%
\definecolor{currentfill}{rgb}{0.156863,0.356863,0.568627}%
\pgfsetfillcolor{currentfill}%
\pgfsetlinewidth{0.000000pt}%
\definecolor{currentstroke}{rgb}{0.000000,0.000000,0.000000}%
\pgfsetstrokecolor{currentstroke}%
\pgfsetstrokeopacity{0.000000}%
\pgfsetdash{}{0pt}%
\pgfpathmoveto{\pgfqpoint{1.505000in}{3.101648in}}%
\pgfpathlineto{\pgfqpoint{2.038299in}{3.101648in}}%
\pgfpathlineto{\pgfqpoint{2.038299in}{3.206130in}}%
\pgfpathlineto{\pgfqpoint{1.505000in}{3.206130in}}%
\pgfpathlineto{\pgfqpoint{1.505000in}{3.101648in}}%
\pgfpathclose%
\pgfusepath{fill}%
\end{pgfscope}%
\begin{pgfscope}%
\pgfpathrectangle{\pgfqpoint{1.505000in}{0.525000in}}{\pgfqpoint{1.293250in}{2.730000in}}%
\pgfusepath{clip}%
\pgfsetbuttcap%
\pgfsetmiterjoin%
\definecolor{currentfill}{rgb}{0.784314,0.784314,0.784314}%
\pgfsetfillcolor{currentfill}%
\pgfsetlinewidth{0.000000pt}%
\definecolor{currentstroke}{rgb}{0.000000,0.000000,0.000000}%
\pgfsetstrokecolor{currentstroke}%
\pgfsetstrokeopacity{0.000000}%
\pgfsetdash{}{0pt}%
\pgfpathmoveto{\pgfqpoint{1.505000in}{2.933130in}}%
\pgfpathlineto{\pgfqpoint{1.731119in}{2.933130in}}%
\pgfpathlineto{\pgfqpoint{1.731119in}{3.037611in}}%
\pgfpathlineto{\pgfqpoint{1.505000in}{3.037611in}}%
\pgfpathlineto{\pgfqpoint{1.505000in}{2.933130in}}%
\pgfpathclose%
\pgfusepath{fill}%
\end{pgfscope}%
\begin{pgfscope}%
\pgfpathrectangle{\pgfqpoint{1.505000in}{0.525000in}}{\pgfqpoint{1.293250in}{2.730000in}}%
\pgfusepath{clip}%
\pgfsetbuttcap%
\pgfsetmiterjoin%
\definecolor{currentfill}{rgb}{0.784314,0.784314,0.784314}%
\pgfsetfillcolor{currentfill}%
\pgfsetlinewidth{0.000000pt}%
\definecolor{currentstroke}{rgb}{0.000000,0.000000,0.000000}%
\pgfsetstrokecolor{currentstroke}%
\pgfsetstrokeopacity{0.000000}%
\pgfsetdash{}{0pt}%
\pgfpathmoveto{\pgfqpoint{1.505000in}{2.764611in}}%
\pgfpathlineto{\pgfqpoint{1.856977in}{2.764611in}}%
\pgfpathlineto{\pgfqpoint{1.856977in}{2.869093in}}%
\pgfpathlineto{\pgfqpoint{1.505000in}{2.869093in}}%
\pgfpathlineto{\pgfqpoint{1.505000in}{2.764611in}}%
\pgfpathclose%
\pgfusepath{fill}%
\end{pgfscope}%
\begin{pgfscope}%
\pgfpathrectangle{\pgfqpoint{1.505000in}{0.525000in}}{\pgfqpoint{1.293250in}{2.730000in}}%
\pgfusepath{clip}%
\pgfsetbuttcap%
\pgfsetmiterjoin%
\definecolor{currentfill}{rgb}{0.784314,0.784314,0.784314}%
\pgfsetfillcolor{currentfill}%
\pgfsetlinewidth{0.000000pt}%
\definecolor{currentstroke}{rgb}{0.000000,0.000000,0.000000}%
\pgfsetstrokecolor{currentstroke}%
\pgfsetstrokeopacity{0.000000}%
\pgfsetdash{}{0pt}%
\pgfpathmoveto{\pgfqpoint{1.505000in}{2.596093in}}%
\pgfpathlineto{\pgfqpoint{1.720453in}{2.596093in}}%
\pgfpathlineto{\pgfqpoint{1.720453in}{2.700574in}}%
\pgfpathlineto{\pgfqpoint{1.505000in}{2.700574in}}%
\pgfpathlineto{\pgfqpoint{1.505000in}{2.596093in}}%
\pgfpathclose%
\pgfusepath{fill}%
\end{pgfscope}%
\begin{pgfscope}%
\pgfpathrectangle{\pgfqpoint{1.505000in}{0.525000in}}{\pgfqpoint{1.293250in}{2.730000in}}%
\pgfusepath{clip}%
\pgfsetbuttcap%
\pgfsetmiterjoin%
\definecolor{currentfill}{rgb}{0.784314,0.784314,0.784314}%
\pgfsetfillcolor{currentfill}%
\pgfsetlinewidth{0.000000pt}%
\definecolor{currentstroke}{rgb}{0.000000,0.000000,0.000000}%
\pgfsetstrokecolor{currentstroke}%
\pgfsetstrokeopacity{0.000000}%
\pgfsetdash{}{0pt}%
\pgfpathmoveto{\pgfqpoint{1.505000in}{2.427574in}}%
\pgfpathlineto{\pgfqpoint{1.792981in}{2.427574in}}%
\pgfpathlineto{\pgfqpoint{1.792981in}{2.532056in}}%
\pgfpathlineto{\pgfqpoint{1.505000in}{2.532056in}}%
\pgfpathlineto{\pgfqpoint{1.505000in}{2.427574in}}%
\pgfpathclose%
\pgfusepath{fill}%
\end{pgfscope}%
\begin{pgfscope}%
\pgfpathrectangle{\pgfqpoint{1.505000in}{0.525000in}}{\pgfqpoint{1.293250in}{2.730000in}}%
\pgfusepath{clip}%
\pgfsetbuttcap%
\pgfsetmiterjoin%
\definecolor{currentfill}{rgb}{0.784314,0.784314,0.784314}%
\pgfsetfillcolor{currentfill}%
\pgfsetlinewidth{0.000000pt}%
\definecolor{currentstroke}{rgb}{0.000000,0.000000,0.000000}%
\pgfsetstrokecolor{currentstroke}%
\pgfsetstrokeopacity{0.000000}%
\pgfsetdash{}{0pt}%
\pgfpathmoveto{\pgfqpoint{1.505000in}{2.259056in}}%
\pgfpathlineto{\pgfqpoint{1.797248in}{2.259056in}}%
\pgfpathlineto{\pgfqpoint{1.797248in}{2.363537in}}%
\pgfpathlineto{\pgfqpoint{1.505000in}{2.363537in}}%
\pgfpathlineto{\pgfqpoint{1.505000in}{2.259056in}}%
\pgfpathclose%
\pgfusepath{fill}%
\end{pgfscope}%
\begin{pgfscope}%
\pgfpathrectangle{\pgfqpoint{1.505000in}{0.525000in}}{\pgfqpoint{1.293250in}{2.730000in}}%
\pgfusepath{clip}%
\pgfsetbuttcap%
\pgfsetmiterjoin%
\definecolor{currentfill}{rgb}{0.478431,0.478431,0.478431}%
\pgfsetfillcolor{currentfill}%
\pgfsetlinewidth{0.000000pt}%
\definecolor{currentstroke}{rgb}{0.000000,0.000000,0.000000}%
\pgfsetstrokecolor{currentstroke}%
\pgfsetstrokeopacity{0.000000}%
\pgfsetdash{}{0pt}%
\pgfpathmoveto{\pgfqpoint{1.505000in}{2.090537in}}%
\pgfpathlineto{\pgfqpoint{1.790848in}{2.090537in}}%
\pgfpathlineto{\pgfqpoint{1.790848in}{2.195019in}}%
\pgfpathlineto{\pgfqpoint{1.505000in}{2.195019in}}%
\pgfpathlineto{\pgfqpoint{1.505000in}{2.090537in}}%
\pgfpathclose%
\pgfusepath{fill}%
\end{pgfscope}%
\begin{pgfscope}%
\pgfpathrectangle{\pgfqpoint{1.505000in}{0.525000in}}{\pgfqpoint{1.293250in}{2.730000in}}%
\pgfusepath{clip}%
\pgfsetbuttcap%
\pgfsetmiterjoin%
\definecolor{currentfill}{rgb}{0.784314,0.784314,0.784314}%
\pgfsetfillcolor{currentfill}%
\pgfsetlinewidth{0.000000pt}%
\definecolor{currentstroke}{rgb}{0.000000,0.000000,0.000000}%
\pgfsetstrokecolor{currentstroke}%
\pgfsetstrokeopacity{0.000000}%
\pgfsetdash{}{0pt}%
\pgfpathmoveto{\pgfqpoint{1.505000in}{1.922019in}}%
\pgfpathlineto{\pgfqpoint{1.788715in}{1.922019in}}%
\pgfpathlineto{\pgfqpoint{1.788715in}{2.026500in}}%
\pgfpathlineto{\pgfqpoint{1.505000in}{2.026500in}}%
\pgfpathlineto{\pgfqpoint{1.505000in}{1.922019in}}%
\pgfpathclose%
\pgfusepath{fill}%
\end{pgfscope}%
\begin{pgfscope}%
\pgfpathrectangle{\pgfqpoint{1.505000in}{0.525000in}}{\pgfqpoint{1.293250in}{2.730000in}}%
\pgfusepath{clip}%
\pgfsetbuttcap%
\pgfsetmiterjoin%
\definecolor{currentfill}{rgb}{0.784314,0.784314,0.784314}%
\pgfsetfillcolor{currentfill}%
\pgfsetlinewidth{0.000000pt}%
\definecolor{currentstroke}{rgb}{0.000000,0.000000,0.000000}%
\pgfsetstrokecolor{currentstroke}%
\pgfsetstrokeopacity{0.000000}%
\pgfsetdash{}{0pt}%
\pgfpathmoveto{\pgfqpoint{1.505000in}{1.753500in}}%
\pgfpathlineto{\pgfqpoint{1.784449in}{1.753500in}}%
\pgfpathlineto{\pgfqpoint{1.784449in}{1.857981in}}%
\pgfpathlineto{\pgfqpoint{1.505000in}{1.857981in}}%
\pgfpathlineto{\pgfqpoint{1.505000in}{1.753500in}}%
\pgfpathclose%
\pgfusepath{fill}%
\end{pgfscope}%
\begin{pgfscope}%
\pgfpathrectangle{\pgfqpoint{1.505000in}{0.525000in}}{\pgfqpoint{1.293250in}{2.730000in}}%
\pgfusepath{clip}%
\pgfsetbuttcap%
\pgfsetmiterjoin%
\definecolor{currentfill}{rgb}{0.478431,0.478431,0.478431}%
\pgfsetfillcolor{currentfill}%
\pgfsetlinewidth{0.000000pt}%
\definecolor{currentstroke}{rgb}{0.000000,0.000000,0.000000}%
\pgfsetstrokecolor{currentstroke}%
\pgfsetstrokeopacity{0.000000}%
\pgfsetdash{}{0pt}%
\pgfpathmoveto{\pgfqpoint{1.505000in}{1.584981in}}%
\pgfpathlineto{\pgfqpoint{1.784449in}{1.584981in}}%
\pgfpathlineto{\pgfqpoint{1.784449in}{1.689463in}}%
\pgfpathlineto{\pgfqpoint{1.505000in}{1.689463in}}%
\pgfpathlineto{\pgfqpoint{1.505000in}{1.584981in}}%
\pgfpathclose%
\pgfusepath{fill}%
\end{pgfscope}%
\begin{pgfscope}%
\pgfpathrectangle{\pgfqpoint{1.505000in}{0.525000in}}{\pgfqpoint{1.293250in}{2.730000in}}%
\pgfusepath{clip}%
\pgfsetbuttcap%
\pgfsetmiterjoin%
\definecolor{currentfill}{rgb}{0.784314,0.784314,0.784314}%
\pgfsetfillcolor{currentfill}%
\pgfsetlinewidth{0.000000pt}%
\definecolor{currentstroke}{rgb}{0.000000,0.000000,0.000000}%
\pgfsetstrokecolor{currentstroke}%
\pgfsetstrokeopacity{0.000000}%
\pgfsetdash{}{0pt}%
\pgfpathmoveto{\pgfqpoint{1.505000in}{1.416463in}}%
\pgfpathlineto{\pgfqpoint{1.782315in}{1.416463in}}%
\pgfpathlineto{\pgfqpoint{1.782315in}{1.520944in}}%
\pgfpathlineto{\pgfqpoint{1.505000in}{1.520944in}}%
\pgfpathlineto{\pgfqpoint{1.505000in}{1.416463in}}%
\pgfpathclose%
\pgfusepath{fill}%
\end{pgfscope}%
\begin{pgfscope}%
\pgfpathrectangle{\pgfqpoint{1.505000in}{0.525000in}}{\pgfqpoint{1.293250in}{2.730000in}}%
\pgfusepath{clip}%
\pgfsetbuttcap%
\pgfsetmiterjoin%
\definecolor{currentfill}{rgb}{0.478431,0.478431,0.478431}%
\pgfsetfillcolor{currentfill}%
\pgfsetlinewidth{0.000000pt}%
\definecolor{currentstroke}{rgb}{0.000000,0.000000,0.000000}%
\pgfsetstrokecolor{currentstroke}%
\pgfsetstrokeopacity{0.000000}%
\pgfsetdash{}{0pt}%
\pgfpathmoveto{\pgfqpoint{1.505000in}{1.247944in}}%
\pgfpathlineto{\pgfqpoint{1.778049in}{1.247944in}}%
\pgfpathlineto{\pgfqpoint{1.778049in}{1.352426in}}%
\pgfpathlineto{\pgfqpoint{1.505000in}{1.352426in}}%
\pgfpathlineto{\pgfqpoint{1.505000in}{1.247944in}}%
\pgfpathclose%
\pgfusepath{fill}%
\end{pgfscope}%
\begin{pgfscope}%
\pgfpathrectangle{\pgfqpoint{1.505000in}{0.525000in}}{\pgfqpoint{1.293250in}{2.730000in}}%
\pgfusepath{clip}%
\pgfsetbuttcap%
\pgfsetmiterjoin%
\definecolor{currentfill}{rgb}{0.784314,0.784314,0.784314}%
\pgfsetfillcolor{currentfill}%
\pgfsetlinewidth{0.000000pt}%
\definecolor{currentstroke}{rgb}{0.000000,0.000000,0.000000}%
\pgfsetstrokecolor{currentstroke}%
\pgfsetstrokeopacity{0.000000}%
\pgfsetdash{}{0pt}%
\pgfpathmoveto{\pgfqpoint{1.505000in}{1.079426in}}%
\pgfpathlineto{\pgfqpoint{1.778049in}{1.079426in}}%
\pgfpathlineto{\pgfqpoint{1.778049in}{1.183907in}}%
\pgfpathlineto{\pgfqpoint{1.505000in}{1.183907in}}%
\pgfpathlineto{\pgfqpoint{1.505000in}{1.079426in}}%
\pgfpathclose%
\pgfusepath{fill}%
\end{pgfscope}%
\begin{pgfscope}%
\pgfpathrectangle{\pgfqpoint{1.505000in}{0.525000in}}{\pgfqpoint{1.293250in}{2.730000in}}%
\pgfusepath{clip}%
\pgfsetbuttcap%
\pgfsetmiterjoin%
\definecolor{currentfill}{rgb}{0.478431,0.478431,0.478431}%
\pgfsetfillcolor{currentfill}%
\pgfsetlinewidth{0.000000pt}%
\definecolor{currentstroke}{rgb}{0.000000,0.000000,0.000000}%
\pgfsetstrokecolor{currentstroke}%
\pgfsetstrokeopacity{0.000000}%
\pgfsetdash{}{0pt}%
\pgfpathmoveto{\pgfqpoint{1.505000in}{0.910907in}}%
\pgfpathlineto{\pgfqpoint{1.782315in}{0.910907in}}%
\pgfpathlineto{\pgfqpoint{1.782315in}{1.015389in}}%
\pgfpathlineto{\pgfqpoint{1.505000in}{1.015389in}}%
\pgfpathlineto{\pgfqpoint{1.505000in}{0.910907in}}%
\pgfpathclose%
\pgfusepath{fill}%
\end{pgfscope}%
\begin{pgfscope}%
\pgfpathrectangle{\pgfqpoint{1.505000in}{0.525000in}}{\pgfqpoint{1.293250in}{2.730000in}}%
\pgfusepath{clip}%
\pgfsetbuttcap%
\pgfsetmiterjoin%
\definecolor{currentfill}{rgb}{0.478431,0.478431,0.478431}%
\pgfsetfillcolor{currentfill}%
\pgfsetlinewidth{0.000000pt}%
\definecolor{currentstroke}{rgb}{0.000000,0.000000,0.000000}%
\pgfsetstrokecolor{currentstroke}%
\pgfsetstrokeopacity{0.000000}%
\pgfsetdash{}{0pt}%
\pgfpathmoveto{\pgfqpoint{1.505000in}{0.742389in}}%
\pgfpathlineto{\pgfqpoint{1.778049in}{0.742389in}}%
\pgfpathlineto{\pgfqpoint{1.778049in}{0.846870in}}%
\pgfpathlineto{\pgfqpoint{1.505000in}{0.846870in}}%
\pgfpathlineto{\pgfqpoint{1.505000in}{0.742389in}}%
\pgfpathclose%
\pgfusepath{fill}%
\end{pgfscope}%
\begin{pgfscope}%
\pgfpathrectangle{\pgfqpoint{1.505000in}{0.525000in}}{\pgfqpoint{1.293250in}{2.730000in}}%
\pgfusepath{clip}%
\pgfsetbuttcap%
\pgfsetmiterjoin%
\definecolor{currentfill}{rgb}{0.478431,0.305882,0.611765}%
\pgfsetfillcolor{currentfill}%
\pgfsetlinewidth{0.000000pt}%
\definecolor{currentstroke}{rgb}{0.000000,0.000000,0.000000}%
\pgfsetstrokecolor{currentstroke}%
\pgfsetstrokeopacity{0.000000}%
\pgfsetdash{}{0pt}%
\pgfpathmoveto{\pgfqpoint{1.505000in}{0.573870in}}%
\pgfpathlineto{\pgfqpoint{1.771649in}{0.573870in}}%
\pgfpathlineto{\pgfqpoint{1.771649in}{0.678352in}}%
\pgfpathlineto{\pgfqpoint{1.505000in}{0.678352in}}%
\pgfpathlineto{\pgfqpoint{1.505000in}{0.573870in}}%
\pgfpathclose%
\pgfusepath{fill}%
\end{pgfscope}%
\begin{pgfscope}%
\pgfsetrectcap%
\pgfsetmiterjoin%
\pgfsetlinewidth{0.803000pt}%
\definecolor{currentstroke}{rgb}{0.101961,0.101961,0.101961}%
\pgfsetstrokecolor{currentstroke}%
\pgfsetdash{}{0pt}%
\pgfpathmoveto{\pgfqpoint{1.505000in}{0.525000in}}%
\pgfpathlineto{\pgfqpoint{1.505000in}{3.255000in}}%
\pgfusepath{stroke}%
\end{pgfscope}%
\begin{pgfscope}%
\pgfsetrectcap%
\pgfsetmiterjoin%
\pgfsetlinewidth{0.803000pt}%
\definecolor{currentstroke}{rgb}{0.101961,0.101961,0.101961}%
\pgfsetstrokecolor{currentstroke}%
\pgfsetdash{}{0pt}%
\pgfpathmoveto{\pgfqpoint{1.505000in}{0.525000in}}%
\pgfpathlineto{\pgfqpoint{2.798250in}{0.525000in}}%
\pgfusepath{stroke}%
\end{pgfscope}%
\begin{pgfscope}%
\pgfpathrectangle{\pgfqpoint{1.505000in}{0.525000in}}{\pgfqpoint{1.293250in}{2.730000in}}%
\pgfusepath{clip}%
\pgfsetbuttcap%
\pgfsetroundjoin%
\definecolor{currentfill}{rgb}{0.101961,0.101961,0.101961}%
\pgfsetfillcolor{currentfill}%
\pgfsetlinewidth{0.301125pt}%
\definecolor{currentstroke}{rgb}{1.000000,1.000000,1.000000}%
\pgfsetstrokecolor{currentstroke}%
\pgfsetdash{}{0pt}%
\pgfsys@defobject{currentmarker}{\pgfqpoint{-0.013889in}{-0.013889in}}{\pgfqpoint{0.013889in}{0.013889in}}{%
\pgfpathmoveto{\pgfqpoint{0.000000in}{-0.013889in}}%
\pgfpathcurveto{\pgfqpoint{0.003683in}{-0.013889in}}{\pgfqpoint{0.007216in}{-0.012425in}}{\pgfqpoint{0.009821in}{-0.009821in}}%
\pgfpathcurveto{\pgfqpoint{0.012425in}{-0.007216in}}{\pgfqpoint{0.013889in}{-0.003683in}}{\pgfqpoint{0.013889in}{0.000000in}}%
\pgfpathcurveto{\pgfqpoint{0.013889in}{0.003683in}}{\pgfqpoint{0.012425in}{0.007216in}}{\pgfqpoint{0.009821in}{0.009821in}}%
\pgfpathcurveto{\pgfqpoint{0.007216in}{0.012425in}}{\pgfqpoint{0.003683in}{0.013889in}}{\pgfqpoint{0.000000in}{0.013889in}}%
\pgfpathcurveto{\pgfqpoint{-0.003683in}{0.013889in}}{\pgfqpoint{-0.007216in}{0.012425in}}{\pgfqpoint{-0.009821in}{0.009821in}}%
\pgfpathcurveto{\pgfqpoint{-0.012425in}{0.007216in}}{\pgfqpoint{-0.013889in}{0.003683in}}{\pgfqpoint{-0.013889in}{0.000000in}}%
\pgfpathcurveto{\pgfqpoint{-0.013889in}{-0.003683in}}{\pgfqpoint{-0.012425in}{-0.007216in}}{\pgfqpoint{-0.009821in}{-0.009821in}}%
\pgfpathcurveto{\pgfqpoint{-0.007216in}{-0.012425in}}{\pgfqpoint{-0.003683in}{-0.013889in}}{\pgfqpoint{0.000000in}{-0.013889in}}%
\pgfpathlineto{\pgfqpoint{0.000000in}{-0.013889in}}%
\pgfpathclose%
\pgfusepath{stroke,fill}%
}%
\begin{pgfscope}%
\pgfsys@transformshift{2.155544in}{3.153889in}%
\pgfsys@useobject{currentmarker}{}%
\end{pgfscope}%
\begin{pgfscope}%
\pgfsys@transformshift{1.961202in}{3.153889in}%
\pgfsys@useobject{currentmarker}{}%
\end{pgfscope}%
\begin{pgfscope}%
\pgfsys@transformshift{2.026575in}{3.153889in}%
\pgfsys@useobject{currentmarker}{}%
\end{pgfscope}%
\begin{pgfscope}%
\pgfsys@transformshift{1.997667in}{3.153889in}%
\pgfsys@useobject{currentmarker}{}%
\end{pgfscope}%
\end{pgfscope}%
\begin{pgfscope}%
\definecolor{textcolor}{rgb}{0.101961,0.101961,0.101961}%
\pgfsetstrokecolor{textcolor}%
\pgfsetfillcolor{textcolor}%
\pgftext[x=2.580984in,y=3.153889in,left,]{\color{textcolor}\sffamily\fontsize{6.000000}{7.200000}\bfseries\selectfont 0.52}%
\end{pgfscope}%
\begin{pgfscope}%
\pgfpathrectangle{\pgfqpoint{1.505000in}{0.525000in}}{\pgfqpoint{1.293250in}{2.730000in}}%
\pgfusepath{clip}%
\pgfsetbuttcap%
\pgfsetroundjoin%
\definecolor{currentfill}{rgb}{0.101961,0.101961,0.101961}%
\pgfsetfillcolor{currentfill}%
\pgfsetlinewidth{0.301125pt}%
\definecolor{currentstroke}{rgb}{1.000000,1.000000,1.000000}%
\pgfsetstrokecolor{currentstroke}%
\pgfsetdash{}{0pt}%
\pgfsys@defobject{currentmarker}{\pgfqpoint{-0.013889in}{-0.013889in}}{\pgfqpoint{0.013889in}{0.013889in}}{%
\pgfpathmoveto{\pgfqpoint{0.000000in}{-0.013889in}}%
\pgfpathcurveto{\pgfqpoint{0.003683in}{-0.013889in}}{\pgfqpoint{0.007216in}{-0.012425in}}{\pgfqpoint{0.009821in}{-0.009821in}}%
\pgfpathcurveto{\pgfqpoint{0.012425in}{-0.007216in}}{\pgfqpoint{0.013889in}{-0.003683in}}{\pgfqpoint{0.013889in}{0.000000in}}%
\pgfpathcurveto{\pgfqpoint{0.013889in}{0.003683in}}{\pgfqpoint{0.012425in}{0.007216in}}{\pgfqpoint{0.009821in}{0.009821in}}%
\pgfpathcurveto{\pgfqpoint{0.007216in}{0.012425in}}{\pgfqpoint{0.003683in}{0.013889in}}{\pgfqpoint{0.000000in}{0.013889in}}%
\pgfpathcurveto{\pgfqpoint{-0.003683in}{0.013889in}}{\pgfqpoint{-0.007216in}{0.012425in}}{\pgfqpoint{-0.009821in}{0.009821in}}%
\pgfpathcurveto{\pgfqpoint{-0.012425in}{0.007216in}}{\pgfqpoint{-0.013889in}{0.003683in}}{\pgfqpoint{-0.013889in}{0.000000in}}%
\pgfpathcurveto{\pgfqpoint{-0.013889in}{-0.003683in}}{\pgfqpoint{-0.012425in}{-0.007216in}}{\pgfqpoint{-0.009821in}{-0.009821in}}%
\pgfpathcurveto{\pgfqpoint{-0.007216in}{-0.012425in}}{\pgfqpoint{-0.003683in}{-0.013889in}}{\pgfqpoint{0.000000in}{-0.013889in}}%
\pgfpathlineto{\pgfqpoint{0.000000in}{-0.013889in}}%
\pgfpathclose%
\pgfusepath{stroke,fill}%
}%
\begin{pgfscope}%
\pgfsys@transformshift{1.646082in}{2.985370in}%
\pgfsys@useobject{currentmarker}{}%
\end{pgfscope}%
\begin{pgfscope}%
\pgfsys@transformshift{1.814565in}{2.985370in}%
\pgfsys@useobject{currentmarker}{}%
\end{pgfscope}%
\begin{pgfscope}%
\pgfsys@transformshift{1.676008in}{2.985370in}%
\pgfsys@useobject{currentmarker}{}%
\end{pgfscope}%
\begin{pgfscope}%
\pgfsys@transformshift{1.800600in}{2.985370in}%
\pgfsys@useobject{currentmarker}{}%
\end{pgfscope}%
\end{pgfscope}%
\begin{pgfscope}%
\definecolor{textcolor}{rgb}{0.101961,0.101961,0.101961}%
\pgfsetstrokecolor{textcolor}%
\pgfsetfillcolor{textcolor}%
\pgftext[x=2.580984in,y=2.985370in,left,]{\color{textcolor}\sffamily\fontsize{6.000000}{7.200000}\selectfont 0.22}%
\end{pgfscope}%
\begin{pgfscope}%
\pgfpathrectangle{\pgfqpoint{1.505000in}{0.525000in}}{\pgfqpoint{1.293250in}{2.730000in}}%
\pgfusepath{clip}%
\pgfsetbuttcap%
\pgfsetroundjoin%
\definecolor{currentfill}{rgb}{0.101961,0.101961,0.101961}%
\pgfsetfillcolor{currentfill}%
\pgfsetlinewidth{0.301125pt}%
\definecolor{currentstroke}{rgb}{1.000000,1.000000,1.000000}%
\pgfsetstrokecolor{currentstroke}%
\pgfsetdash{}{0pt}%
\pgfsys@defobject{currentmarker}{\pgfqpoint{-0.013889in}{-0.013889in}}{\pgfqpoint{0.013889in}{0.013889in}}{%
\pgfpathmoveto{\pgfqpoint{0.000000in}{-0.013889in}}%
\pgfpathcurveto{\pgfqpoint{0.003683in}{-0.013889in}}{\pgfqpoint{0.007216in}{-0.012425in}}{\pgfqpoint{0.009821in}{-0.009821in}}%
\pgfpathcurveto{\pgfqpoint{0.012425in}{-0.007216in}}{\pgfqpoint{0.013889in}{-0.003683in}}{\pgfqpoint{0.013889in}{0.000000in}}%
\pgfpathcurveto{\pgfqpoint{0.013889in}{0.003683in}}{\pgfqpoint{0.012425in}{0.007216in}}{\pgfqpoint{0.009821in}{0.009821in}}%
\pgfpathcurveto{\pgfqpoint{0.007216in}{0.012425in}}{\pgfqpoint{0.003683in}{0.013889in}}{\pgfqpoint{0.000000in}{0.013889in}}%
\pgfpathcurveto{\pgfqpoint{-0.003683in}{0.013889in}}{\pgfqpoint{-0.007216in}{0.012425in}}{\pgfqpoint{-0.009821in}{0.009821in}}%
\pgfpathcurveto{\pgfqpoint{-0.012425in}{0.007216in}}{\pgfqpoint{-0.013889in}{0.003683in}}{\pgfqpoint{-0.013889in}{0.000000in}}%
\pgfpathcurveto{\pgfqpoint{-0.013889in}{-0.003683in}}{\pgfqpoint{-0.012425in}{-0.007216in}}{\pgfqpoint{-0.009821in}{-0.009821in}}%
\pgfpathcurveto{\pgfqpoint{-0.007216in}{-0.012425in}}{\pgfqpoint{-0.003683in}{-0.013889in}}{\pgfqpoint{0.000000in}{-0.013889in}}%
\pgfpathlineto{\pgfqpoint{0.000000in}{-0.013889in}}%
\pgfpathclose%
\pgfusepath{stroke,fill}%
}%
\begin{pgfscope}%
\pgfsys@transformshift{1.998786in}{2.816852in}%
\pgfsys@useobject{currentmarker}{}%
\end{pgfscope}%
\begin{pgfscope}%
\pgfsys@transformshift{1.798272in}{2.816852in}%
\pgfsys@useobject{currentmarker}{}%
\end{pgfscope}%
\begin{pgfscope}%
\pgfsys@transformshift{1.710210in}{2.816852in}%
\pgfsys@useobject{currentmarker}{}%
\end{pgfscope}%
\begin{pgfscope}%
\pgfsys@transformshift{1.918840in}{2.816852in}%
\pgfsys@useobject{currentmarker}{}%
\end{pgfscope}%
\end{pgfscope}%
\begin{pgfscope}%
\definecolor{textcolor}{rgb}{0.101961,0.101961,0.101961}%
\pgfsetstrokecolor{textcolor}%
\pgfsetfillcolor{textcolor}%
\pgftext[x=2.580984in,y=2.816852in,left,]{\color{textcolor}\sffamily\fontsize{6.000000}{7.200000}\selectfont 0.34}%
\end{pgfscope}%
\begin{pgfscope}%
\pgfpathrectangle{\pgfqpoint{1.505000in}{0.525000in}}{\pgfqpoint{1.293250in}{2.730000in}}%
\pgfusepath{clip}%
\pgfsetbuttcap%
\pgfsetroundjoin%
\definecolor{currentfill}{rgb}{0.101961,0.101961,0.101961}%
\pgfsetfillcolor{currentfill}%
\pgfsetlinewidth{0.301125pt}%
\definecolor{currentstroke}{rgb}{1.000000,1.000000,1.000000}%
\pgfsetstrokecolor{currentstroke}%
\pgfsetdash{}{0pt}%
\pgfsys@defobject{currentmarker}{\pgfqpoint{-0.013889in}{-0.013889in}}{\pgfqpoint{0.013889in}{0.013889in}}{%
\pgfpathmoveto{\pgfqpoint{0.000000in}{-0.013889in}}%
\pgfpathcurveto{\pgfqpoint{0.003683in}{-0.013889in}}{\pgfqpoint{0.007216in}{-0.012425in}}{\pgfqpoint{0.009821in}{-0.009821in}}%
\pgfpathcurveto{\pgfqpoint{0.012425in}{-0.007216in}}{\pgfqpoint{0.013889in}{-0.003683in}}{\pgfqpoint{0.013889in}{0.000000in}}%
\pgfpathcurveto{\pgfqpoint{0.013889in}{0.003683in}}{\pgfqpoint{0.012425in}{0.007216in}}{\pgfqpoint{0.009821in}{0.009821in}}%
\pgfpathcurveto{\pgfqpoint{0.007216in}{0.012425in}}{\pgfqpoint{0.003683in}{0.013889in}}{\pgfqpoint{0.000000in}{0.013889in}}%
\pgfpathcurveto{\pgfqpoint{-0.003683in}{0.013889in}}{\pgfqpoint{-0.007216in}{0.012425in}}{\pgfqpoint{-0.009821in}{0.009821in}}%
\pgfpathcurveto{\pgfqpoint{-0.012425in}{0.007216in}}{\pgfqpoint{-0.013889in}{0.003683in}}{\pgfqpoint{-0.013889in}{0.000000in}}%
\pgfpathcurveto{\pgfqpoint{-0.013889in}{-0.003683in}}{\pgfqpoint{-0.012425in}{-0.007216in}}{\pgfqpoint{-0.009821in}{-0.009821in}}%
\pgfpathcurveto{\pgfqpoint{-0.007216in}{-0.012425in}}{\pgfqpoint{-0.003683in}{-0.013889in}}{\pgfqpoint{0.000000in}{-0.013889in}}%
\pgfpathlineto{\pgfqpoint{0.000000in}{-0.013889in}}%
\pgfpathclose%
\pgfusepath{stroke,fill}%
}%
\begin{pgfscope}%
\pgfsys@transformshift{1.685271in}{2.648333in}%
\pgfsys@useobject{currentmarker}{}%
\end{pgfscope}%
\begin{pgfscope}%
\pgfsys@transformshift{1.814565in}{2.648333in}%
\pgfsys@useobject{currentmarker}{}%
\end{pgfscope}%
\begin{pgfscope}%
\pgfsys@transformshift{1.641807in}{2.648333in}%
\pgfsys@useobject{currentmarker}{}%
\end{pgfscope}%
\begin{pgfscope}%
\pgfsys@transformshift{1.741480in}{2.648333in}%
\pgfsys@useobject{currentmarker}{}%
\end{pgfscope}%
\end{pgfscope}%
\begin{pgfscope}%
\definecolor{textcolor}{rgb}{0.101961,0.101961,0.101961}%
\pgfsetstrokecolor{textcolor}%
\pgfsetfillcolor{textcolor}%
\pgftext[x=2.580984in,y=2.648333in,left,]{\color{textcolor}\sffamily\fontsize{6.000000}{7.200000}\selectfont 0.21}%
\end{pgfscope}%
\begin{pgfscope}%
\pgfpathrectangle{\pgfqpoint{1.505000in}{0.525000in}}{\pgfqpoint{1.293250in}{2.730000in}}%
\pgfusepath{clip}%
\pgfsetbuttcap%
\pgfsetroundjoin%
\definecolor{currentfill}{rgb}{0.101961,0.101961,0.101961}%
\pgfsetfillcolor{currentfill}%
\pgfsetlinewidth{0.301125pt}%
\definecolor{currentstroke}{rgb}{1.000000,1.000000,1.000000}%
\pgfsetstrokecolor{currentstroke}%
\pgfsetdash{}{0pt}%
\pgfsys@defobject{currentmarker}{\pgfqpoint{-0.013889in}{-0.013889in}}{\pgfqpoint{0.013889in}{0.013889in}}{%
\pgfpathmoveto{\pgfqpoint{0.000000in}{-0.013889in}}%
\pgfpathcurveto{\pgfqpoint{0.003683in}{-0.013889in}}{\pgfqpoint{0.007216in}{-0.012425in}}{\pgfqpoint{0.009821in}{-0.009821in}}%
\pgfpathcurveto{\pgfqpoint{0.012425in}{-0.007216in}}{\pgfqpoint{0.013889in}{-0.003683in}}{\pgfqpoint{0.013889in}{0.000000in}}%
\pgfpathcurveto{\pgfqpoint{0.013889in}{0.003683in}}{\pgfqpoint{0.012425in}{0.007216in}}{\pgfqpoint{0.009821in}{0.009821in}}%
\pgfpathcurveto{\pgfqpoint{0.007216in}{0.012425in}}{\pgfqpoint{0.003683in}{0.013889in}}{\pgfqpoint{0.000000in}{0.013889in}}%
\pgfpathcurveto{\pgfqpoint{-0.003683in}{0.013889in}}{\pgfqpoint{-0.007216in}{0.012425in}}{\pgfqpoint{-0.009821in}{0.009821in}}%
\pgfpathcurveto{\pgfqpoint{-0.012425in}{0.007216in}}{\pgfqpoint{-0.013889in}{0.003683in}}{\pgfqpoint{-0.013889in}{0.000000in}}%
\pgfpathcurveto{\pgfqpoint{-0.013889in}{-0.003683in}}{\pgfqpoint{-0.012425in}{-0.007216in}}{\pgfqpoint{-0.009821in}{-0.009821in}}%
\pgfpathcurveto{\pgfqpoint{-0.007216in}{-0.012425in}}{\pgfqpoint{-0.003683in}{-0.013889in}}{\pgfqpoint{0.000000in}{-0.013889in}}%
\pgfpathlineto{\pgfqpoint{0.000000in}{-0.013889in}}%
\pgfpathclose%
\pgfusepath{stroke,fill}%
}%
\begin{pgfscope}%
\pgfsys@transformshift{1.912570in}{2.479815in}%
\pgfsys@useobject{currentmarker}{}%
\end{pgfscope}%
\begin{pgfscope}%
\pgfsys@transformshift{1.814565in}{2.479815in}%
\pgfsys@useobject{currentmarker}{}%
\end{pgfscope}%
\begin{pgfscope}%
\pgfsys@transformshift{1.607605in}{2.479815in}%
\pgfsys@useobject{currentmarker}{}%
\end{pgfscope}%
\begin{pgfscope}%
\pgfsys@transformshift{1.830160in}{2.479815in}%
\pgfsys@useobject{currentmarker}{}%
\end{pgfscope}%
\end{pgfscope}%
\begin{pgfscope}%
\definecolor{textcolor}{rgb}{0.101961,0.101961,0.101961}%
\pgfsetstrokecolor{textcolor}%
\pgfsetfillcolor{textcolor}%
\pgftext[x=2.580984in,y=2.479815in,left,]{\color{textcolor}\sffamily\fontsize{6.000000}{7.200000}\selectfont 0.28}%
\end{pgfscope}%
\begin{pgfscope}%
\pgfpathrectangle{\pgfqpoint{1.505000in}{0.525000in}}{\pgfqpoint{1.293250in}{2.730000in}}%
\pgfusepath{clip}%
\pgfsetbuttcap%
\pgfsetroundjoin%
\definecolor{currentfill}{rgb}{0.101961,0.101961,0.101961}%
\pgfsetfillcolor{currentfill}%
\pgfsetlinewidth{0.301125pt}%
\definecolor{currentstroke}{rgb}{1.000000,1.000000,1.000000}%
\pgfsetstrokecolor{currentstroke}%
\pgfsetdash{}{0pt}%
\pgfsys@defobject{currentmarker}{\pgfqpoint{-0.013889in}{-0.013889in}}{\pgfqpoint{0.013889in}{0.013889in}}{%
\pgfpathmoveto{\pgfqpoint{0.000000in}{-0.013889in}}%
\pgfpathcurveto{\pgfqpoint{0.003683in}{-0.013889in}}{\pgfqpoint{0.007216in}{-0.012425in}}{\pgfqpoint{0.009821in}{-0.009821in}}%
\pgfpathcurveto{\pgfqpoint{0.012425in}{-0.007216in}}{\pgfqpoint{0.013889in}{-0.003683in}}{\pgfqpoint{0.013889in}{0.000000in}}%
\pgfpathcurveto{\pgfqpoint{0.013889in}{0.003683in}}{\pgfqpoint{0.012425in}{0.007216in}}{\pgfqpoint{0.009821in}{0.009821in}}%
\pgfpathcurveto{\pgfqpoint{0.007216in}{0.012425in}}{\pgfqpoint{0.003683in}{0.013889in}}{\pgfqpoint{0.000000in}{0.013889in}}%
\pgfpathcurveto{\pgfqpoint{-0.003683in}{0.013889in}}{\pgfqpoint{-0.007216in}{0.012425in}}{\pgfqpoint{-0.009821in}{0.009821in}}%
\pgfpathcurveto{\pgfqpoint{-0.012425in}{0.007216in}}{\pgfqpoint{-0.013889in}{0.003683in}}{\pgfqpoint{-0.013889in}{0.000000in}}%
\pgfpathcurveto{\pgfqpoint{-0.013889in}{-0.003683in}}{\pgfqpoint{-0.012425in}{-0.007216in}}{\pgfqpoint{-0.009821in}{-0.009821in}}%
\pgfpathcurveto{\pgfqpoint{-0.007216in}{-0.012425in}}{\pgfqpoint{-0.003683in}{-0.013889in}}{\pgfqpoint{0.000000in}{-0.013889in}}%
\pgfpathlineto{\pgfqpoint{0.000000in}{-0.013889in}}%
\pgfpathclose%
\pgfusepath{stroke,fill}%
}%
\begin{pgfscope}%
\pgfsys@transformshift{1.889056in}{2.311296in}%
\pgfsys@useobject{currentmarker}{}%
\end{pgfscope}%
\begin{pgfscope}%
\pgfsys@transformshift{1.822712in}{2.311296in}%
\pgfsys@useobject{currentmarker}{}%
\end{pgfscope}%
\begin{pgfscope}%
\pgfsys@transformshift{1.650357in}{2.311296in}%
\pgfsys@useobject{currentmarker}{}%
\end{pgfscope}%
\begin{pgfscope}%
\pgfsys@transformshift{1.820307in}{2.311296in}%
\pgfsys@useobject{currentmarker}{}%
\end{pgfscope}%
\end{pgfscope}%
\begin{pgfscope}%
\definecolor{textcolor}{rgb}{0.101961,0.101961,0.101961}%
\pgfsetstrokecolor{textcolor}%
\pgfsetfillcolor{textcolor}%
\pgftext[x=2.580984in,y=2.311296in,left,]{\color{textcolor}\sffamily\fontsize{6.000000}{7.200000}\selectfont 0.28}%
\end{pgfscope}%
\begin{pgfscope}%
\pgfpathrectangle{\pgfqpoint{1.505000in}{0.525000in}}{\pgfqpoint{1.293250in}{2.730000in}}%
\pgfusepath{clip}%
\pgfsetbuttcap%
\pgfsetroundjoin%
\definecolor{currentfill}{rgb}{0.101961,0.101961,0.101961}%
\pgfsetfillcolor{currentfill}%
\pgfsetlinewidth{0.301125pt}%
\definecolor{currentstroke}{rgb}{1.000000,1.000000,1.000000}%
\pgfsetstrokecolor{currentstroke}%
\pgfsetdash{}{0pt}%
\pgfsys@defobject{currentmarker}{\pgfqpoint{-0.013889in}{-0.013889in}}{\pgfqpoint{0.013889in}{0.013889in}}{%
\pgfpathmoveto{\pgfqpoint{0.000000in}{-0.013889in}}%
\pgfpathcurveto{\pgfqpoint{0.003683in}{-0.013889in}}{\pgfqpoint{0.007216in}{-0.012425in}}{\pgfqpoint{0.009821in}{-0.009821in}}%
\pgfpathcurveto{\pgfqpoint{0.012425in}{-0.007216in}}{\pgfqpoint{0.013889in}{-0.003683in}}{\pgfqpoint{0.013889in}{0.000000in}}%
\pgfpathcurveto{\pgfqpoint{0.013889in}{0.003683in}}{\pgfqpoint{0.012425in}{0.007216in}}{\pgfqpoint{0.009821in}{0.009821in}}%
\pgfpathcurveto{\pgfqpoint{0.007216in}{0.012425in}}{\pgfqpoint{0.003683in}{0.013889in}}{\pgfqpoint{0.000000in}{0.013889in}}%
\pgfpathcurveto{\pgfqpoint{-0.003683in}{0.013889in}}{\pgfqpoint{-0.007216in}{0.012425in}}{\pgfqpoint{-0.009821in}{0.009821in}}%
\pgfpathcurveto{\pgfqpoint{-0.012425in}{0.007216in}}{\pgfqpoint{-0.013889in}{0.003683in}}{\pgfqpoint{-0.013889in}{0.000000in}}%
\pgfpathcurveto{\pgfqpoint{-0.013889in}{-0.003683in}}{\pgfqpoint{-0.012425in}{-0.007216in}}{\pgfqpoint{-0.009821in}{-0.009821in}}%
\pgfpathcurveto{\pgfqpoint{-0.007216in}{-0.012425in}}{\pgfqpoint{-0.003683in}{-0.013889in}}{\pgfqpoint{0.000000in}{-0.013889in}}%
\pgfpathlineto{\pgfqpoint{0.000000in}{-0.013889in}}%
\pgfpathclose%
\pgfusepath{stroke,fill}%
}%
\begin{pgfscope}%
\pgfsys@transformshift{1.889056in}{2.142778in}%
\pgfsys@useobject{currentmarker}{}%
\end{pgfscope}%
\begin{pgfscope}%
\pgfsys@transformshift{1.822712in}{2.142778in}%
\pgfsys@useobject{currentmarker}{}%
\end{pgfscope}%
\begin{pgfscope}%
\pgfsys@transformshift{1.599055in}{2.142778in}%
\pgfsys@useobject{currentmarker}{}%
\end{pgfscope}%
\begin{pgfscope}%
\pgfsys@transformshift{1.849867in}{2.142778in}%
\pgfsys@useobject{currentmarker}{}%
\end{pgfscope}%
\end{pgfscope}%
\begin{pgfscope}%
\definecolor{textcolor}{rgb}{0.101961,0.101961,0.101961}%
\pgfsetstrokecolor{textcolor}%
\pgfsetfillcolor{textcolor}%
\pgftext[x=2.580984in,y=2.142778in,left,]{\color{textcolor}\sffamily\fontsize{6.000000}{7.200000}\selectfont 0.28}%
\end{pgfscope}%
\begin{pgfscope}%
\pgfpathrectangle{\pgfqpoint{1.505000in}{0.525000in}}{\pgfqpoint{1.293250in}{2.730000in}}%
\pgfusepath{clip}%
\pgfsetbuttcap%
\pgfsetroundjoin%
\definecolor{currentfill}{rgb}{0.101961,0.101961,0.101961}%
\pgfsetfillcolor{currentfill}%
\pgfsetlinewidth{0.301125pt}%
\definecolor{currentstroke}{rgb}{1.000000,1.000000,1.000000}%
\pgfsetstrokecolor{currentstroke}%
\pgfsetdash{}{0pt}%
\pgfsys@defobject{currentmarker}{\pgfqpoint{-0.013889in}{-0.013889in}}{\pgfqpoint{0.013889in}{0.013889in}}{%
\pgfpathmoveto{\pgfqpoint{0.000000in}{-0.013889in}}%
\pgfpathcurveto{\pgfqpoint{0.003683in}{-0.013889in}}{\pgfqpoint{0.007216in}{-0.012425in}}{\pgfqpoint{0.009821in}{-0.009821in}}%
\pgfpathcurveto{\pgfqpoint{0.012425in}{-0.007216in}}{\pgfqpoint{0.013889in}{-0.003683in}}{\pgfqpoint{0.013889in}{0.000000in}}%
\pgfpathcurveto{\pgfqpoint{0.013889in}{0.003683in}}{\pgfqpoint{0.012425in}{0.007216in}}{\pgfqpoint{0.009821in}{0.009821in}}%
\pgfpathcurveto{\pgfqpoint{0.007216in}{0.012425in}}{\pgfqpoint{0.003683in}{0.013889in}}{\pgfqpoint{0.000000in}{0.013889in}}%
\pgfpathcurveto{\pgfqpoint{-0.003683in}{0.013889in}}{\pgfqpoint{-0.007216in}{0.012425in}}{\pgfqpoint{-0.009821in}{0.009821in}}%
\pgfpathcurveto{\pgfqpoint{-0.012425in}{0.007216in}}{\pgfqpoint{-0.013889in}{0.003683in}}{\pgfqpoint{-0.013889in}{0.000000in}}%
\pgfpathcurveto{\pgfqpoint{-0.013889in}{-0.003683in}}{\pgfqpoint{-0.012425in}{-0.007216in}}{\pgfqpoint{-0.009821in}{-0.009821in}}%
\pgfpathcurveto{\pgfqpoint{-0.007216in}{-0.012425in}}{\pgfqpoint{-0.003683in}{-0.013889in}}{\pgfqpoint{0.000000in}{-0.013889in}}%
\pgfpathlineto{\pgfqpoint{0.000000in}{-0.013889in}}%
\pgfpathclose%
\pgfusepath{stroke,fill}%
}%
\begin{pgfscope}%
\pgfsys@transformshift{1.889056in}{1.974259in}%
\pgfsys@useobject{currentmarker}{}%
\end{pgfscope}%
\begin{pgfscope}%
\pgfsys@transformshift{1.822712in}{1.974259in}%
\pgfsys@useobject{currentmarker}{}%
\end{pgfscope}%
\begin{pgfscope}%
\pgfsys@transformshift{1.599055in}{1.974259in}%
\pgfsys@useobject{currentmarker}{}%
\end{pgfscope}%
\begin{pgfscope}%
\pgfsys@transformshift{1.840013in}{1.974259in}%
\pgfsys@useobject{currentmarker}{}%
\end{pgfscope}%
\end{pgfscope}%
\begin{pgfscope}%
\definecolor{textcolor}{rgb}{0.101961,0.101961,0.101961}%
\pgfsetstrokecolor{textcolor}%
\pgfsetfillcolor{textcolor}%
\pgftext[x=2.580984in,y=1.974259in,left,]{\color{textcolor}\sffamily\fontsize{6.000000}{7.200000}\selectfont 0.27}%
\end{pgfscope}%
\begin{pgfscope}%
\pgfpathrectangle{\pgfqpoint{1.505000in}{0.525000in}}{\pgfqpoint{1.293250in}{2.730000in}}%
\pgfusepath{clip}%
\pgfsetbuttcap%
\pgfsetroundjoin%
\definecolor{currentfill}{rgb}{0.101961,0.101961,0.101961}%
\pgfsetfillcolor{currentfill}%
\pgfsetlinewidth{0.301125pt}%
\definecolor{currentstroke}{rgb}{1.000000,1.000000,1.000000}%
\pgfsetstrokecolor{currentstroke}%
\pgfsetdash{}{0pt}%
\pgfsys@defobject{currentmarker}{\pgfqpoint{-0.013889in}{-0.013889in}}{\pgfqpoint{0.013889in}{0.013889in}}{%
\pgfpathmoveto{\pgfqpoint{0.000000in}{-0.013889in}}%
\pgfpathcurveto{\pgfqpoint{0.003683in}{-0.013889in}}{\pgfqpoint{0.007216in}{-0.012425in}}{\pgfqpoint{0.009821in}{-0.009821in}}%
\pgfpathcurveto{\pgfqpoint{0.012425in}{-0.007216in}}{\pgfqpoint{0.013889in}{-0.003683in}}{\pgfqpoint{0.013889in}{0.000000in}}%
\pgfpathcurveto{\pgfqpoint{0.013889in}{0.003683in}}{\pgfqpoint{0.012425in}{0.007216in}}{\pgfqpoint{0.009821in}{0.009821in}}%
\pgfpathcurveto{\pgfqpoint{0.007216in}{0.012425in}}{\pgfqpoint{0.003683in}{0.013889in}}{\pgfqpoint{0.000000in}{0.013889in}}%
\pgfpathcurveto{\pgfqpoint{-0.003683in}{0.013889in}}{\pgfqpoint{-0.007216in}{0.012425in}}{\pgfqpoint{-0.009821in}{0.009821in}}%
\pgfpathcurveto{\pgfqpoint{-0.012425in}{0.007216in}}{\pgfqpoint{-0.013889in}{0.003683in}}{\pgfqpoint{-0.013889in}{0.000000in}}%
\pgfpathcurveto{\pgfqpoint{-0.013889in}{-0.003683in}}{\pgfqpoint{-0.012425in}{-0.007216in}}{\pgfqpoint{-0.009821in}{-0.009821in}}%
\pgfpathcurveto{\pgfqpoint{-0.007216in}{-0.012425in}}{\pgfqpoint{-0.003683in}{-0.013889in}}{\pgfqpoint{0.000000in}{-0.013889in}}%
\pgfpathlineto{\pgfqpoint{0.000000in}{-0.013889in}}%
\pgfpathclose%
\pgfusepath{stroke,fill}%
}%
\begin{pgfscope}%
\pgfsys@transformshift{1.904732in}{1.805741in}%
\pgfsys@useobject{currentmarker}{}%
\end{pgfscope}%
\begin{pgfscope}%
\pgfsys@transformshift{1.806419in}{1.805741in}%
\pgfsys@useobject{currentmarker}{}%
\end{pgfscope}%
\begin{pgfscope}%
\pgfsys@transformshift{1.607605in}{1.805741in}%
\pgfsys@useobject{currentmarker}{}%
\end{pgfscope}%
\begin{pgfscope}%
\pgfsys@transformshift{1.810453in}{1.805741in}%
\pgfsys@useobject{currentmarker}{}%
\end{pgfscope}%
\end{pgfscope}%
\begin{pgfscope}%
\definecolor{textcolor}{rgb}{0.101961,0.101961,0.101961}%
\pgfsetstrokecolor{textcolor}%
\pgfsetfillcolor{textcolor}%
\pgftext[x=2.580984in,y=1.805741in,left,]{\color{textcolor}\sffamily\fontsize{6.000000}{7.200000}\selectfont 0.27}%
\end{pgfscope}%
\begin{pgfscope}%
\pgfpathrectangle{\pgfqpoint{1.505000in}{0.525000in}}{\pgfqpoint{1.293250in}{2.730000in}}%
\pgfusepath{clip}%
\pgfsetbuttcap%
\pgfsetroundjoin%
\definecolor{currentfill}{rgb}{0.101961,0.101961,0.101961}%
\pgfsetfillcolor{currentfill}%
\pgfsetlinewidth{0.301125pt}%
\definecolor{currentstroke}{rgb}{1.000000,1.000000,1.000000}%
\pgfsetstrokecolor{currentstroke}%
\pgfsetdash{}{0pt}%
\pgfsys@defobject{currentmarker}{\pgfqpoint{-0.013889in}{-0.013889in}}{\pgfqpoint{0.013889in}{0.013889in}}{%
\pgfpathmoveto{\pgfqpoint{0.000000in}{-0.013889in}}%
\pgfpathcurveto{\pgfqpoint{0.003683in}{-0.013889in}}{\pgfqpoint{0.007216in}{-0.012425in}}{\pgfqpoint{0.009821in}{-0.009821in}}%
\pgfpathcurveto{\pgfqpoint{0.012425in}{-0.007216in}}{\pgfqpoint{0.013889in}{-0.003683in}}{\pgfqpoint{0.013889in}{0.000000in}}%
\pgfpathcurveto{\pgfqpoint{0.013889in}{0.003683in}}{\pgfqpoint{0.012425in}{0.007216in}}{\pgfqpoint{0.009821in}{0.009821in}}%
\pgfpathcurveto{\pgfqpoint{0.007216in}{0.012425in}}{\pgfqpoint{0.003683in}{0.013889in}}{\pgfqpoint{0.000000in}{0.013889in}}%
\pgfpathcurveto{\pgfqpoint{-0.003683in}{0.013889in}}{\pgfqpoint{-0.007216in}{0.012425in}}{\pgfqpoint{-0.009821in}{0.009821in}}%
\pgfpathcurveto{\pgfqpoint{-0.012425in}{0.007216in}}{\pgfqpoint{-0.013889in}{0.003683in}}{\pgfqpoint{-0.013889in}{0.000000in}}%
\pgfpathcurveto{\pgfqpoint{-0.013889in}{-0.003683in}}{\pgfqpoint{-0.012425in}{-0.007216in}}{\pgfqpoint{-0.009821in}{-0.009821in}}%
\pgfpathcurveto{\pgfqpoint{-0.007216in}{-0.012425in}}{\pgfqpoint{-0.003683in}{-0.013889in}}{\pgfqpoint{0.000000in}{-0.013889in}}%
\pgfpathlineto{\pgfqpoint{0.000000in}{-0.013889in}}%
\pgfpathclose%
\pgfusepath{stroke,fill}%
}%
\begin{pgfscope}%
\pgfsys@transformshift{1.889056in}{1.637222in}%
\pgfsys@useobject{currentmarker}{}%
\end{pgfscope}%
\begin{pgfscope}%
\pgfsys@transformshift{1.830858in}{1.637222in}%
\pgfsys@useobject{currentmarker}{}%
\end{pgfscope}%
\begin{pgfscope}%
\pgfsys@transformshift{1.599055in}{1.637222in}%
\pgfsys@useobject{currentmarker}{}%
\end{pgfscope}%
\begin{pgfscope}%
\pgfsys@transformshift{1.810453in}{1.637222in}%
\pgfsys@useobject{currentmarker}{}%
\end{pgfscope}%
\end{pgfscope}%
\begin{pgfscope}%
\definecolor{textcolor}{rgb}{0.101961,0.101961,0.101961}%
\pgfsetstrokecolor{textcolor}%
\pgfsetfillcolor{textcolor}%
\pgftext[x=2.580984in,y=1.637222in,left,]{\color{textcolor}\sffamily\fontsize{6.000000}{7.200000}\selectfont 0.27}%
\end{pgfscope}%
\begin{pgfscope}%
\pgfpathrectangle{\pgfqpoint{1.505000in}{0.525000in}}{\pgfqpoint{1.293250in}{2.730000in}}%
\pgfusepath{clip}%
\pgfsetbuttcap%
\pgfsetroundjoin%
\definecolor{currentfill}{rgb}{0.101961,0.101961,0.101961}%
\pgfsetfillcolor{currentfill}%
\pgfsetlinewidth{0.301125pt}%
\definecolor{currentstroke}{rgb}{1.000000,1.000000,1.000000}%
\pgfsetstrokecolor{currentstroke}%
\pgfsetdash{}{0pt}%
\pgfsys@defobject{currentmarker}{\pgfqpoint{-0.013889in}{-0.013889in}}{\pgfqpoint{0.013889in}{0.013889in}}{%
\pgfpathmoveto{\pgfqpoint{0.000000in}{-0.013889in}}%
\pgfpathcurveto{\pgfqpoint{0.003683in}{-0.013889in}}{\pgfqpoint{0.007216in}{-0.012425in}}{\pgfqpoint{0.009821in}{-0.009821in}}%
\pgfpathcurveto{\pgfqpoint{0.012425in}{-0.007216in}}{\pgfqpoint{0.013889in}{-0.003683in}}{\pgfqpoint{0.013889in}{0.000000in}}%
\pgfpathcurveto{\pgfqpoint{0.013889in}{0.003683in}}{\pgfqpoint{0.012425in}{0.007216in}}{\pgfqpoint{0.009821in}{0.009821in}}%
\pgfpathcurveto{\pgfqpoint{0.007216in}{0.012425in}}{\pgfqpoint{0.003683in}{0.013889in}}{\pgfqpoint{0.000000in}{0.013889in}}%
\pgfpathcurveto{\pgfqpoint{-0.003683in}{0.013889in}}{\pgfqpoint{-0.007216in}{0.012425in}}{\pgfqpoint{-0.009821in}{0.009821in}}%
\pgfpathcurveto{\pgfqpoint{-0.012425in}{0.007216in}}{\pgfqpoint{-0.013889in}{0.003683in}}{\pgfqpoint{-0.013889in}{0.000000in}}%
\pgfpathcurveto{\pgfqpoint{-0.013889in}{-0.003683in}}{\pgfqpoint{-0.012425in}{-0.007216in}}{\pgfqpoint{-0.009821in}{-0.009821in}}%
\pgfpathcurveto{\pgfqpoint{-0.007216in}{-0.012425in}}{\pgfqpoint{-0.003683in}{-0.013889in}}{\pgfqpoint{0.000000in}{-0.013889in}}%
\pgfpathlineto{\pgfqpoint{0.000000in}{-0.013889in}}%
\pgfpathclose%
\pgfusepath{stroke,fill}%
}%
\begin{pgfscope}%
\pgfsys@transformshift{1.889056in}{1.468704in}%
\pgfsys@useobject{currentmarker}{}%
\end{pgfscope}%
\begin{pgfscope}%
\pgfsys@transformshift{1.830858in}{1.468704in}%
\pgfsys@useobject{currentmarker}{}%
\end{pgfscope}%
\begin{pgfscope}%
\pgfsys@transformshift{1.590504in}{1.468704in}%
\pgfsys@useobject{currentmarker}{}%
\end{pgfscope}%
\begin{pgfscope}%
\pgfsys@transformshift{1.810453in}{1.468704in}%
\pgfsys@useobject{currentmarker}{}%
\end{pgfscope}%
\end{pgfscope}%
\begin{pgfscope}%
\definecolor{textcolor}{rgb}{0.101961,0.101961,0.101961}%
\pgfsetstrokecolor{textcolor}%
\pgfsetfillcolor{textcolor}%
\pgftext[x=2.580984in,y=1.468704in,left,]{\color{textcolor}\sffamily\fontsize{6.000000}{7.200000}\selectfont 0.27}%
\end{pgfscope}%
\begin{pgfscope}%
\pgfpathrectangle{\pgfqpoint{1.505000in}{0.525000in}}{\pgfqpoint{1.293250in}{2.730000in}}%
\pgfusepath{clip}%
\pgfsetbuttcap%
\pgfsetroundjoin%
\definecolor{currentfill}{rgb}{0.101961,0.101961,0.101961}%
\pgfsetfillcolor{currentfill}%
\pgfsetlinewidth{0.301125pt}%
\definecolor{currentstroke}{rgb}{1.000000,1.000000,1.000000}%
\pgfsetstrokecolor{currentstroke}%
\pgfsetdash{}{0pt}%
\pgfsys@defobject{currentmarker}{\pgfqpoint{-0.013889in}{-0.013889in}}{\pgfqpoint{0.013889in}{0.013889in}}{%
\pgfpathmoveto{\pgfqpoint{0.000000in}{-0.013889in}}%
\pgfpathcurveto{\pgfqpoint{0.003683in}{-0.013889in}}{\pgfqpoint{0.007216in}{-0.012425in}}{\pgfqpoint{0.009821in}{-0.009821in}}%
\pgfpathcurveto{\pgfqpoint{0.012425in}{-0.007216in}}{\pgfqpoint{0.013889in}{-0.003683in}}{\pgfqpoint{0.013889in}{0.000000in}}%
\pgfpathcurveto{\pgfqpoint{0.013889in}{0.003683in}}{\pgfqpoint{0.012425in}{0.007216in}}{\pgfqpoint{0.009821in}{0.009821in}}%
\pgfpathcurveto{\pgfqpoint{0.007216in}{0.012425in}}{\pgfqpoint{0.003683in}{0.013889in}}{\pgfqpoint{0.000000in}{0.013889in}}%
\pgfpathcurveto{\pgfqpoint{-0.003683in}{0.013889in}}{\pgfqpoint{-0.007216in}{0.012425in}}{\pgfqpoint{-0.009821in}{0.009821in}}%
\pgfpathcurveto{\pgfqpoint{-0.012425in}{0.007216in}}{\pgfqpoint{-0.013889in}{0.003683in}}{\pgfqpoint{-0.013889in}{0.000000in}}%
\pgfpathcurveto{\pgfqpoint{-0.013889in}{-0.003683in}}{\pgfqpoint{-0.012425in}{-0.007216in}}{\pgfqpoint{-0.009821in}{-0.009821in}}%
\pgfpathcurveto{\pgfqpoint{-0.007216in}{-0.012425in}}{\pgfqpoint{-0.003683in}{-0.013889in}}{\pgfqpoint{0.000000in}{-0.013889in}}%
\pgfpathlineto{\pgfqpoint{0.000000in}{-0.013889in}}%
\pgfpathclose%
\pgfusepath{stroke,fill}%
}%
\begin{pgfscope}%
\pgfsys@transformshift{1.889056in}{1.300185in}%
\pgfsys@useobject{currentmarker}{}%
\end{pgfscope}%
\begin{pgfscope}%
\pgfsys@transformshift{1.822712in}{1.300185in}%
\pgfsys@useobject{currentmarker}{}%
\end{pgfscope}%
\begin{pgfscope}%
\pgfsys@transformshift{1.581954in}{1.300185in}%
\pgfsys@useobject{currentmarker}{}%
\end{pgfscope}%
\begin{pgfscope}%
\pgfsys@transformshift{1.810453in}{1.300185in}%
\pgfsys@useobject{currentmarker}{}%
\end{pgfscope}%
\end{pgfscope}%
\begin{pgfscope}%
\definecolor{textcolor}{rgb}{0.101961,0.101961,0.101961}%
\pgfsetstrokecolor{textcolor}%
\pgfsetfillcolor{textcolor}%
\pgftext[x=2.580984in,y=1.300185in,left,]{\color{textcolor}\sffamily\fontsize{6.000000}{7.200000}\selectfont 0.26}%
\end{pgfscope}%
\begin{pgfscope}%
\pgfpathrectangle{\pgfqpoint{1.505000in}{0.525000in}}{\pgfqpoint{1.293250in}{2.730000in}}%
\pgfusepath{clip}%
\pgfsetbuttcap%
\pgfsetroundjoin%
\definecolor{currentfill}{rgb}{0.101961,0.101961,0.101961}%
\pgfsetfillcolor{currentfill}%
\pgfsetlinewidth{0.301125pt}%
\definecolor{currentstroke}{rgb}{1.000000,1.000000,1.000000}%
\pgfsetstrokecolor{currentstroke}%
\pgfsetdash{}{0pt}%
\pgfsys@defobject{currentmarker}{\pgfqpoint{-0.013889in}{-0.013889in}}{\pgfqpoint{0.013889in}{0.013889in}}{%
\pgfpathmoveto{\pgfqpoint{0.000000in}{-0.013889in}}%
\pgfpathcurveto{\pgfqpoint{0.003683in}{-0.013889in}}{\pgfqpoint{0.007216in}{-0.012425in}}{\pgfqpoint{0.009821in}{-0.009821in}}%
\pgfpathcurveto{\pgfqpoint{0.012425in}{-0.007216in}}{\pgfqpoint{0.013889in}{-0.003683in}}{\pgfqpoint{0.013889in}{0.000000in}}%
\pgfpathcurveto{\pgfqpoint{0.013889in}{0.003683in}}{\pgfqpoint{0.012425in}{0.007216in}}{\pgfqpoint{0.009821in}{0.009821in}}%
\pgfpathcurveto{\pgfqpoint{0.007216in}{0.012425in}}{\pgfqpoint{0.003683in}{0.013889in}}{\pgfqpoint{0.000000in}{0.013889in}}%
\pgfpathcurveto{\pgfqpoint{-0.003683in}{0.013889in}}{\pgfqpoint{-0.007216in}{0.012425in}}{\pgfqpoint{-0.009821in}{0.009821in}}%
\pgfpathcurveto{\pgfqpoint{-0.012425in}{0.007216in}}{\pgfqpoint{-0.013889in}{0.003683in}}{\pgfqpoint{-0.013889in}{0.000000in}}%
\pgfpathcurveto{\pgfqpoint{-0.013889in}{-0.003683in}}{\pgfqpoint{-0.012425in}{-0.007216in}}{\pgfqpoint{-0.009821in}{-0.009821in}}%
\pgfpathcurveto{\pgfqpoint{-0.007216in}{-0.012425in}}{\pgfqpoint{-0.003683in}{-0.013889in}}{\pgfqpoint{0.000000in}{-0.013889in}}%
\pgfpathlineto{\pgfqpoint{0.000000in}{-0.013889in}}%
\pgfpathclose%
\pgfusepath{stroke,fill}%
}%
\begin{pgfscope}%
\pgfsys@transformshift{1.889056in}{1.131667in}%
\pgfsys@useobject{currentmarker}{}%
\end{pgfscope}%
\begin{pgfscope}%
\pgfsys@transformshift{1.814565in}{1.131667in}%
\pgfsys@useobject{currentmarker}{}%
\end{pgfscope}%
\begin{pgfscope}%
\pgfsys@transformshift{1.590504in}{1.131667in}%
\pgfsys@useobject{currentmarker}{}%
\end{pgfscope}%
\begin{pgfscope}%
\pgfsys@transformshift{1.810453in}{1.131667in}%
\pgfsys@useobject{currentmarker}{}%
\end{pgfscope}%
\end{pgfscope}%
\begin{pgfscope}%
\definecolor{textcolor}{rgb}{0.101961,0.101961,0.101961}%
\pgfsetstrokecolor{textcolor}%
\pgfsetfillcolor{textcolor}%
\pgftext[x=2.580984in,y=1.131667in,left,]{\color{textcolor}\sffamily\fontsize{6.000000}{7.200000}\selectfont 0.26}%
\end{pgfscope}%
\begin{pgfscope}%
\pgfpathrectangle{\pgfqpoint{1.505000in}{0.525000in}}{\pgfqpoint{1.293250in}{2.730000in}}%
\pgfusepath{clip}%
\pgfsetbuttcap%
\pgfsetroundjoin%
\definecolor{currentfill}{rgb}{0.101961,0.101961,0.101961}%
\pgfsetfillcolor{currentfill}%
\pgfsetlinewidth{0.301125pt}%
\definecolor{currentstroke}{rgb}{1.000000,1.000000,1.000000}%
\pgfsetstrokecolor{currentstroke}%
\pgfsetdash{}{0pt}%
\pgfsys@defobject{currentmarker}{\pgfqpoint{-0.013889in}{-0.013889in}}{\pgfqpoint{0.013889in}{0.013889in}}{%
\pgfpathmoveto{\pgfqpoint{0.000000in}{-0.013889in}}%
\pgfpathcurveto{\pgfqpoint{0.003683in}{-0.013889in}}{\pgfqpoint{0.007216in}{-0.012425in}}{\pgfqpoint{0.009821in}{-0.009821in}}%
\pgfpathcurveto{\pgfqpoint{0.012425in}{-0.007216in}}{\pgfqpoint{0.013889in}{-0.003683in}}{\pgfqpoint{0.013889in}{0.000000in}}%
\pgfpathcurveto{\pgfqpoint{0.013889in}{0.003683in}}{\pgfqpoint{0.012425in}{0.007216in}}{\pgfqpoint{0.009821in}{0.009821in}}%
\pgfpathcurveto{\pgfqpoint{0.007216in}{0.012425in}}{\pgfqpoint{0.003683in}{0.013889in}}{\pgfqpoint{0.000000in}{0.013889in}}%
\pgfpathcurveto{\pgfqpoint{-0.003683in}{0.013889in}}{\pgfqpoint{-0.007216in}{0.012425in}}{\pgfqpoint{-0.009821in}{0.009821in}}%
\pgfpathcurveto{\pgfqpoint{-0.012425in}{0.007216in}}{\pgfqpoint{-0.013889in}{0.003683in}}{\pgfqpoint{-0.013889in}{0.000000in}}%
\pgfpathcurveto{\pgfqpoint{-0.013889in}{-0.003683in}}{\pgfqpoint{-0.012425in}{-0.007216in}}{\pgfqpoint{-0.009821in}{-0.009821in}}%
\pgfpathcurveto{\pgfqpoint{-0.007216in}{-0.012425in}}{\pgfqpoint{-0.003683in}{-0.013889in}}{\pgfqpoint{0.000000in}{-0.013889in}}%
\pgfpathlineto{\pgfqpoint{0.000000in}{-0.013889in}}%
\pgfpathclose%
\pgfusepath{stroke,fill}%
}%
\begin{pgfscope}%
\pgfsys@transformshift{1.889056in}{0.963148in}%
\pgfsys@useobject{currentmarker}{}%
\end{pgfscope}%
\begin{pgfscope}%
\pgfsys@transformshift{1.830858in}{0.963148in}%
\pgfsys@useobject{currentmarker}{}%
\end{pgfscope}%
\begin{pgfscope}%
\pgfsys@transformshift{1.590504in}{0.963148in}%
\pgfsys@useobject{currentmarker}{}%
\end{pgfscope}%
\begin{pgfscope}%
\pgfsys@transformshift{1.810453in}{0.963148in}%
\pgfsys@useobject{currentmarker}{}%
\end{pgfscope}%
\end{pgfscope}%
\begin{pgfscope}%
\definecolor{textcolor}{rgb}{0.101961,0.101961,0.101961}%
\pgfsetstrokecolor{textcolor}%
\pgfsetfillcolor{textcolor}%
\pgftext[x=2.580984in,y=0.963148in,left,]{\color{textcolor}\sffamily\fontsize{6.000000}{7.200000}\selectfont 0.27}%
\end{pgfscope}%
\begin{pgfscope}%
\pgfpathrectangle{\pgfqpoint{1.505000in}{0.525000in}}{\pgfqpoint{1.293250in}{2.730000in}}%
\pgfusepath{clip}%
\pgfsetbuttcap%
\pgfsetroundjoin%
\definecolor{currentfill}{rgb}{0.101961,0.101961,0.101961}%
\pgfsetfillcolor{currentfill}%
\pgfsetlinewidth{0.301125pt}%
\definecolor{currentstroke}{rgb}{1.000000,1.000000,1.000000}%
\pgfsetstrokecolor{currentstroke}%
\pgfsetdash{}{0pt}%
\pgfsys@defobject{currentmarker}{\pgfqpoint{-0.013889in}{-0.013889in}}{\pgfqpoint{0.013889in}{0.013889in}}{%
\pgfpathmoveto{\pgfqpoint{0.000000in}{-0.013889in}}%
\pgfpathcurveto{\pgfqpoint{0.003683in}{-0.013889in}}{\pgfqpoint{0.007216in}{-0.012425in}}{\pgfqpoint{0.009821in}{-0.009821in}}%
\pgfpathcurveto{\pgfqpoint{0.012425in}{-0.007216in}}{\pgfqpoint{0.013889in}{-0.003683in}}{\pgfqpoint{0.013889in}{0.000000in}}%
\pgfpathcurveto{\pgfqpoint{0.013889in}{0.003683in}}{\pgfqpoint{0.012425in}{0.007216in}}{\pgfqpoint{0.009821in}{0.009821in}}%
\pgfpathcurveto{\pgfqpoint{0.007216in}{0.012425in}}{\pgfqpoint{0.003683in}{0.013889in}}{\pgfqpoint{0.000000in}{0.013889in}}%
\pgfpathcurveto{\pgfqpoint{-0.003683in}{0.013889in}}{\pgfqpoint{-0.007216in}{0.012425in}}{\pgfqpoint{-0.009821in}{0.009821in}}%
\pgfpathcurveto{\pgfqpoint{-0.012425in}{0.007216in}}{\pgfqpoint{-0.013889in}{0.003683in}}{\pgfqpoint{-0.013889in}{0.000000in}}%
\pgfpathcurveto{\pgfqpoint{-0.013889in}{-0.003683in}}{\pgfqpoint{-0.012425in}{-0.007216in}}{\pgfqpoint{-0.009821in}{-0.009821in}}%
\pgfpathcurveto{\pgfqpoint{-0.007216in}{-0.012425in}}{\pgfqpoint{-0.003683in}{-0.013889in}}{\pgfqpoint{0.000000in}{-0.013889in}}%
\pgfpathlineto{\pgfqpoint{0.000000in}{-0.013889in}}%
\pgfpathclose%
\pgfusepath{stroke,fill}%
}%
\begin{pgfscope}%
\pgfsys@transformshift{1.889056in}{0.794630in}%
\pgfsys@useobject{currentmarker}{}%
\end{pgfscope}%
\begin{pgfscope}%
\pgfsys@transformshift{1.814565in}{0.794630in}%
\pgfsys@useobject{currentmarker}{}%
\end{pgfscope}%
\begin{pgfscope}%
\pgfsys@transformshift{1.590504in}{0.794630in}%
\pgfsys@useobject{currentmarker}{}%
\end{pgfscope}%
\begin{pgfscope}%
\pgfsys@transformshift{1.810453in}{0.794630in}%
\pgfsys@useobject{currentmarker}{}%
\end{pgfscope}%
\end{pgfscope}%
\begin{pgfscope}%
\definecolor{textcolor}{rgb}{0.101961,0.101961,0.101961}%
\pgfsetstrokecolor{textcolor}%
\pgfsetfillcolor{textcolor}%
\pgftext[x=2.580984in,y=0.794630in,left,]{\color{textcolor}\sffamily\fontsize{6.000000}{7.200000}\selectfont 0.26}%
\end{pgfscope}%
\begin{pgfscope}%
\pgfpathrectangle{\pgfqpoint{1.505000in}{0.525000in}}{\pgfqpoint{1.293250in}{2.730000in}}%
\pgfusepath{clip}%
\pgfsetbuttcap%
\pgfsetroundjoin%
\definecolor{currentfill}{rgb}{0.101961,0.101961,0.101961}%
\pgfsetfillcolor{currentfill}%
\pgfsetlinewidth{0.301125pt}%
\definecolor{currentstroke}{rgb}{1.000000,1.000000,1.000000}%
\pgfsetstrokecolor{currentstroke}%
\pgfsetdash{}{0pt}%
\pgfsys@defobject{currentmarker}{\pgfqpoint{-0.013889in}{-0.013889in}}{\pgfqpoint{0.013889in}{0.013889in}}{%
\pgfpathmoveto{\pgfqpoint{0.000000in}{-0.013889in}}%
\pgfpathcurveto{\pgfqpoint{0.003683in}{-0.013889in}}{\pgfqpoint{0.007216in}{-0.012425in}}{\pgfqpoint{0.009821in}{-0.009821in}}%
\pgfpathcurveto{\pgfqpoint{0.012425in}{-0.007216in}}{\pgfqpoint{0.013889in}{-0.003683in}}{\pgfqpoint{0.013889in}{0.000000in}}%
\pgfpathcurveto{\pgfqpoint{0.013889in}{0.003683in}}{\pgfqpoint{0.012425in}{0.007216in}}{\pgfqpoint{0.009821in}{0.009821in}}%
\pgfpathcurveto{\pgfqpoint{0.007216in}{0.012425in}}{\pgfqpoint{0.003683in}{0.013889in}}{\pgfqpoint{0.000000in}{0.013889in}}%
\pgfpathcurveto{\pgfqpoint{-0.003683in}{0.013889in}}{\pgfqpoint{-0.007216in}{0.012425in}}{\pgfqpoint{-0.009821in}{0.009821in}}%
\pgfpathcurveto{\pgfqpoint{-0.012425in}{0.007216in}}{\pgfqpoint{-0.013889in}{0.003683in}}{\pgfqpoint{-0.013889in}{0.000000in}}%
\pgfpathcurveto{\pgfqpoint{-0.013889in}{-0.003683in}}{\pgfqpoint{-0.012425in}{-0.007216in}}{\pgfqpoint{-0.009821in}{-0.009821in}}%
\pgfpathcurveto{\pgfqpoint{-0.007216in}{-0.012425in}}{\pgfqpoint{-0.003683in}{-0.013889in}}{\pgfqpoint{0.000000in}{-0.013889in}}%
\pgfpathlineto{\pgfqpoint{0.000000in}{-0.013889in}}%
\pgfpathclose%
\pgfusepath{stroke,fill}%
}%
\begin{pgfscope}%
\pgfsys@transformshift{1.889056in}{0.626111in}%
\pgfsys@useobject{currentmarker}{}%
\end{pgfscope}%
\begin{pgfscope}%
\pgfsys@transformshift{1.814565in}{0.626111in}%
\pgfsys@useobject{currentmarker}{}%
\end{pgfscope}%
\begin{pgfscope}%
\pgfsys@transformshift{1.564853in}{0.626111in}%
\pgfsys@useobject{currentmarker}{}%
\end{pgfscope}%
\begin{pgfscope}%
\pgfsys@transformshift{1.810453in}{0.626111in}%
\pgfsys@useobject{currentmarker}{}%
\end{pgfscope}%
\end{pgfscope}%
\begin{pgfscope}%
\definecolor{textcolor}{rgb}{0.101961,0.101961,0.101961}%
\pgfsetstrokecolor{textcolor}%
\pgfsetfillcolor{textcolor}%
\pgftext[x=2.580984in,y=0.626111in,left,]{\color{textcolor}\sffamily\fontsize{6.000000}{7.200000}\selectfont 0.26}%
\end{pgfscope}%
\begin{pgfscope}%
\definecolor{textcolor}{rgb}{0.101961,0.101961,0.101961}%
\pgfsetstrokecolor{textcolor}%
\pgfsetfillcolor{textcolor}%
\pgftext[x=1.466203in,y=3.153889in,left,]{\makebox[0pt][r]{\color{textcolor}\sffamily\fontsize{7.000000}{8.400000}\bfseries\selectfont \textbf{COVEX (old: rung 3)}}}%
\end{pgfscope}%
\begin{pgfscope}%
\definecolor{textcolor}{rgb}{0.101961,0.101961,0.101961}%
\pgfsetstrokecolor{textcolor}%
\pgfsetfillcolor{textcolor}%
\pgftext[x=1.466203in,y=2.985370in,left,]{\makebox[0pt][r]{\color{textcolor}\sffamily\fontsize{7.000000}{8.400000}\selectfont StreamScout~\cite{streamscout2026}}}%
\end{pgfscope}%
\begin{pgfscope}%
\definecolor{textcolor}{rgb}{0.101961,0.101961,0.101961}%
\pgfsetstrokecolor{textcolor}%
\pgfsetfillcolor{textcolor}%
\pgftext[x=1.466203in,y=2.816852in,left,]{\makebox[0pt][r]{\color{textcolor}\sffamily\fontsize{7.000000}{8.400000}\selectfont Confidence cascade~\cite{chen2024frugalgpt}}}%
\end{pgfscope}%
\begin{pgfscope}%
\definecolor{textcolor}{rgb}{0.101961,0.101961,0.101961}%
\pgfsetstrokecolor{textcolor}%
\pgfsetfillcolor{textcolor}%
\pgftext[x=1.466203in,y=2.648333in,left,]{\makebox[0pt][r]{\color{textcolor}\sffamily\fontsize{7.000000}{8.400000}\selectfont SimpleStream~\cite{simplestream2026}}}%
\end{pgfscope}%
\begin{pgfscope}%
\definecolor{textcolor}{rgb}{0.101961,0.101961,0.101961}%
\pgfsetstrokecolor{textcolor}%
\pgfsetfillcolor{textcolor}%
\pgftext[x=1.466203in,y=2.479815in,left,]{\makebox[0pt][r]{\color{textcolor}\sffamily\fontsize{7.000000}{8.400000}\selectfont LLF~\cite{mok1983llf}}}%
\end{pgfscope}%
\begin{pgfscope}%
\definecolor{textcolor}{rgb}{0.101961,0.101961,0.101961}%
\pgfsetstrokecolor{textcolor}%
\pgfsetfillcolor{textcolor}%
\pgftext[x=1.466203in,y=2.311296in,left,]{\makebox[0pt][r]{\color{textcolor}\sffamily\fontsize{7.000000}{8.400000}\selectfont ECO-LLM~\cite{ecollm2025}}}%
\end{pgfscope}%
\begin{pgfscope}%
\definecolor{textcolor}{rgb}{0.101961,0.101961,0.101961}%
\pgfsetstrokecolor{textcolor}%
\pgfsetfillcolor{textcolor}%
\pgftext[x=1.466203in,y=2.142778in,left,]{\makebox[0pt][r]{\color{textcolor}\sffamily\fontsize{7.000000}{8.400000}\selectfont Mooncake~\cite{mooncake2025}}}%
\end{pgfscope}%
\begin{pgfscope}%
\definecolor{textcolor}{rgb}{0.101961,0.101961,0.101961}%
\pgfsetstrokecolor{textcolor}%
\pgfsetfillcolor{textcolor}%
\pgftext[x=1.466203in,y=1.974259in,left,]{\makebox[0pt][r]{\color{textcolor}\sffamily\fontsize{7.000000}{8.400000}\selectfont EDF + feasibility~\cite{liu1973edf}}}%
\end{pgfscope}%
\begin{pgfscope}%
\definecolor{textcolor}{rgb}{0.101961,0.101961,0.101961}%
\pgfsetstrokecolor{textcolor}%
\pgfsetfillcolor{textcolor}%
\pgftext[x=1.466203in,y=1.805741in,left,]{\makebox[0pt][r]{\color{textcolor}\sffamily\fontsize{7.000000}{8.400000}\selectfont EDF~\cite{liu1973edf}}}%
\end{pgfscope}%
\begin{pgfscope}%
\definecolor{textcolor}{rgb}{0.101961,0.101961,0.101961}%
\pgfsetstrokecolor{textcolor}%
\pgfsetfillcolor{textcolor}%
\pgftext[x=1.466203in,y=1.637222in,left,]{\makebox[0pt][r]{\color{textcolor}\sffamily\fontsize{7.000000}{8.400000}\selectfont KVFlow~\cite{kvflow2025}}}%
\end{pgfscope}%
\begin{pgfscope}%
\definecolor{textcolor}{rgb}{0.101961,0.101961,0.101961}%
\pgfsetstrokecolor{textcolor}%
\pgfsetfillcolor{textcolor}%
\pgftext[x=1.466203in,y=1.468704in,left,]{\makebox[0pt][r]{\color{textcolor}\sffamily\fontsize{7.000000}{8.400000}\selectfont Always-cloud FCFS}}%
\end{pgfscope}%
\begin{pgfscope}%
\definecolor{textcolor}{rgb}{0.101961,0.101961,0.101961}%
\pgfsetstrokecolor{textcolor}%
\pgfsetfillcolor{textcolor}%
\pgftext[x=1.466203in,y=1.300185in,left,]{\makebox[0pt][r]{\color{textcolor}\sffamily\fontsize{7.000000}{8.400000}\selectfont Warm-all LRU}}%
\end{pgfscope}%
\begin{pgfscope}%
\definecolor{textcolor}{rgb}{0.101961,0.101961,0.101961}%
\pgfsetstrokecolor{textcolor}%
\pgfsetfillcolor{textcolor}%
\pgftext[x=1.466203in,y=1.131667in,left,]{\makebox[0pt][r]{\color{textcolor}\sffamily\fontsize{7.000000}{8.400000}\selectfont QueryStream~\cite{querystream2026}}}%
\end{pgfscope}%
\begin{pgfscope}%
\definecolor{textcolor}{rgb}{0.101961,0.101961,0.101961}%
\pgfsetstrokecolor{textcolor}%
\pgfsetfillcolor{textcolor}%
\pgftext[x=1.466203in,y=0.963148in,left,]{\makebox[0pt][r]{\color{textcolor}\sffamily\fontsize{7.000000}{8.400000}\selectfont AttentionStore~\cite{cachedattention2024}}}%
\end{pgfscope}%
\begin{pgfscope}%
\definecolor{textcolor}{rgb}{0.101961,0.101961,0.101961}%
\pgfsetstrokecolor{textcolor}%
\pgfsetfillcolor{textcolor}%
\pgftext[x=1.466203in,y=0.794630in,left,]{\makebox[0pt][r]{\color{textcolor}\sffamily\fontsize{7.000000}{8.400000}\selectfont LMCache~\cite{lmcache2025}}}%
\end{pgfscope}%
\begin{pgfscope}%
\definecolor{textcolor}{rgb}{0.101961,0.101961,0.101961}%
\pgfsetstrokecolor{textcolor}%
\pgfsetfillcolor{textcolor}%
\pgftext[x=1.466203in,y=0.626111in,left,]{\makebox[0pt][r]{\color{textcolor}\sffamily\fontsize{7.000000}{8.400000}\selectfont On-device only}}%
\end{pgfscope}%
\begin{pgfscope}%
\pgfsetbuttcap%
\pgfsetmiterjoin%
\definecolor{currentfill}{rgb}{1.000000,1.000000,1.000000}%
\pgfsetfillcolor{currentfill}%
\pgfsetlinewidth{0.000000pt}%
\definecolor{currentstroke}{rgb}{0.000000,0.000000,0.000000}%
\pgfsetstrokecolor{currentstroke}%
\pgfsetstrokeopacity{0.000000}%
\pgfsetdash{}{0pt}%
\pgfpathmoveto{\pgfqpoint{2.882250in}{0.525000in}}%
\pgfpathlineto{\pgfqpoint{4.175500in}{0.525000in}}%
\pgfpathlineto{\pgfqpoint{4.175500in}{3.255000in}}%
\pgfpathlineto{\pgfqpoint{2.882250in}{3.255000in}}%
\pgfpathlineto{\pgfqpoint{2.882250in}{0.525000in}}%
\pgfpathclose%
\pgfusepath{fill}%
\end{pgfscope}%
\begin{pgfscope}%
\pgfpathrectangle{\pgfqpoint{2.882250in}{0.525000in}}{\pgfqpoint{1.293250in}{2.730000in}}%
\pgfusepath{clip}%
\pgfsetrectcap%
\pgfsetroundjoin%
\pgfsetlinewidth{0.401500pt}%
\definecolor{currentstroke}{rgb}{0.901961,0.901961,0.901961}%
\pgfsetstrokecolor{currentstroke}%
\pgfsetdash{}{0pt}%
\pgfpathmoveto{\pgfqpoint{2.882250in}{0.525000in}}%
\pgfpathlineto{\pgfqpoint{2.882250in}{3.255000in}}%
\pgfusepath{stroke}%
\end{pgfscope}%
\begin{pgfscope}%
\pgfsetbuttcap%
\pgfsetroundjoin%
\definecolor{currentfill}{rgb}{0.101961,0.101961,0.101961}%
\pgfsetfillcolor{currentfill}%
\pgfsetlinewidth{0.803000pt}%
\definecolor{currentstroke}{rgb}{0.101961,0.101961,0.101961}%
\pgfsetstrokecolor{currentstroke}%
\pgfsetdash{}{0pt}%
\pgfsys@defobject{currentmarker}{\pgfqpoint{0.000000in}{-0.027778in}}{\pgfqpoint{0.000000in}{0.000000in}}{%
\pgfpathmoveto{\pgfqpoint{0.000000in}{0.000000in}}%
\pgfpathlineto{\pgfqpoint{0.000000in}{-0.027778in}}%
\pgfusepath{stroke,fill}%
}%
\begin{pgfscope}%
\pgfsys@transformshift{2.882250in}{0.525000in}%
\pgfsys@useobject{currentmarker}{}%
\end{pgfscope}%
\end{pgfscope}%
\begin{pgfscope}%
\definecolor{textcolor}{rgb}{0.101961,0.101961,0.101961}%
\pgfsetstrokecolor{textcolor}%
\pgfsetfillcolor{textcolor}%
\pgftext[x=2.882250in,y=0.448611in,,top]{\color{textcolor}\sffamily\fontsize{7.000000}{8.400000}\selectfont 0}%
\end{pgfscope}%
\begin{pgfscope}%
\pgfpathrectangle{\pgfqpoint{2.882250in}{0.525000in}}{\pgfqpoint{1.293250in}{2.730000in}}%
\pgfusepath{clip}%
\pgfsetrectcap%
\pgfsetroundjoin%
\pgfsetlinewidth{0.401500pt}%
\definecolor{currentstroke}{rgb}{0.901961,0.901961,0.901961}%
\pgfsetstrokecolor{currentstroke}%
\pgfsetdash{}{0pt}%
\pgfpathmoveto{\pgfqpoint{3.399550in}{0.525000in}}%
\pgfpathlineto{\pgfqpoint{3.399550in}{3.255000in}}%
\pgfusepath{stroke}%
\end{pgfscope}%
\begin{pgfscope}%
\pgfsetbuttcap%
\pgfsetroundjoin%
\definecolor{currentfill}{rgb}{0.101961,0.101961,0.101961}%
\pgfsetfillcolor{currentfill}%
\pgfsetlinewidth{0.803000pt}%
\definecolor{currentstroke}{rgb}{0.101961,0.101961,0.101961}%
\pgfsetstrokecolor{currentstroke}%
\pgfsetdash{}{0pt}%
\pgfsys@defobject{currentmarker}{\pgfqpoint{0.000000in}{-0.027778in}}{\pgfqpoint{0.000000in}{0.000000in}}{%
\pgfpathmoveto{\pgfqpoint{0.000000in}{0.000000in}}%
\pgfpathlineto{\pgfqpoint{0.000000in}{-0.027778in}}%
\pgfusepath{stroke,fill}%
}%
\begin{pgfscope}%
\pgfsys@transformshift{3.399550in}{0.525000in}%
\pgfsys@useobject{currentmarker}{}%
\end{pgfscope}%
\end{pgfscope}%
\begin{pgfscope}%
\definecolor{textcolor}{rgb}{0.101961,0.101961,0.101961}%
\pgfsetstrokecolor{textcolor}%
\pgfsetfillcolor{textcolor}%
\pgftext[x=3.399550in,y=0.448611in,,top]{\color{textcolor}\sffamily\fontsize{7.000000}{8.400000}\selectfont 0.5}%
\end{pgfscope}%
\begin{pgfscope}%
\pgfpathrectangle{\pgfqpoint{2.882250in}{0.525000in}}{\pgfqpoint{1.293250in}{2.730000in}}%
\pgfusepath{clip}%
\pgfsetrectcap%
\pgfsetroundjoin%
\pgfsetlinewidth{0.401500pt}%
\definecolor{currentstroke}{rgb}{0.901961,0.901961,0.901961}%
\pgfsetstrokecolor{currentstroke}%
\pgfsetdash{}{0pt}%
\pgfpathmoveto{\pgfqpoint{3.916850in}{0.525000in}}%
\pgfpathlineto{\pgfqpoint{3.916850in}{3.255000in}}%
\pgfusepath{stroke}%
\end{pgfscope}%
\begin{pgfscope}%
\pgfsetbuttcap%
\pgfsetroundjoin%
\definecolor{currentfill}{rgb}{0.101961,0.101961,0.101961}%
\pgfsetfillcolor{currentfill}%
\pgfsetlinewidth{0.803000pt}%
\definecolor{currentstroke}{rgb}{0.101961,0.101961,0.101961}%
\pgfsetstrokecolor{currentstroke}%
\pgfsetdash{}{0pt}%
\pgfsys@defobject{currentmarker}{\pgfqpoint{0.000000in}{-0.027778in}}{\pgfqpoint{0.000000in}{0.000000in}}{%
\pgfpathmoveto{\pgfqpoint{0.000000in}{0.000000in}}%
\pgfpathlineto{\pgfqpoint{0.000000in}{-0.027778in}}%
\pgfusepath{stroke,fill}%
}%
\begin{pgfscope}%
\pgfsys@transformshift{3.916850in}{0.525000in}%
\pgfsys@useobject{currentmarker}{}%
\end{pgfscope}%
\end{pgfscope}%
\begin{pgfscope}%
\definecolor{textcolor}{rgb}{0.101961,0.101961,0.101961}%
\pgfsetstrokecolor{textcolor}%
\pgfsetfillcolor{textcolor}%
\pgftext[x=3.916850in,y=0.448611in,,top]{\color{textcolor}\sffamily\fontsize{7.000000}{8.400000}\selectfont 1}%
\end{pgfscope}%
\begin{pgfscope}%
\definecolor{textcolor}{rgb}{0.101961,0.101961,0.101961}%
\pgfsetstrokecolor{textcolor}%
\pgfsetfillcolor{textcolor}%
\pgftext[x=3.528875in,y=0.334414in,,top]{\color{textcolor}\sffamily\fontsize{7.500000}{9.000000}\selectfont Single-shot}%
\end{pgfscope}%
\begin{pgfscope}%
\pgfpathrectangle{\pgfqpoint{2.882250in}{0.525000in}}{\pgfqpoint{1.293250in}{2.730000in}}%
\pgfusepath{clip}%
\pgfsetbuttcap%
\pgfsetmiterjoin%
\definecolor{currentfill}{rgb}{0.156863,0.356863,0.568627}%
\pgfsetfillcolor{currentfill}%
\pgfsetlinewidth{0.000000pt}%
\definecolor{currentstroke}{rgb}{0.000000,0.000000,0.000000}%
\pgfsetstrokecolor{currentstroke}%
\pgfsetstrokeopacity{0.000000}%
\pgfsetdash{}{0pt}%
\pgfpathmoveto{\pgfqpoint{2.882250in}{3.101648in}}%
\pgfpathlineto{\pgfqpoint{3.501996in}{3.101648in}}%
\pgfpathlineto{\pgfqpoint{3.501996in}{3.206130in}}%
\pgfpathlineto{\pgfqpoint{2.882250in}{3.206130in}}%
\pgfpathlineto{\pgfqpoint{2.882250in}{3.101648in}}%
\pgfpathclose%
\pgfusepath{fill}%
\end{pgfscope}%
\begin{pgfscope}%
\pgfpathrectangle{\pgfqpoint{2.882250in}{0.525000in}}{\pgfqpoint{1.293250in}{2.730000in}}%
\pgfusepath{clip}%
\pgfsetbuttcap%
\pgfsetmiterjoin%
\definecolor{currentfill}{rgb}{0.784314,0.784314,0.784314}%
\pgfsetfillcolor{currentfill}%
\pgfsetlinewidth{0.000000pt}%
\definecolor{currentstroke}{rgb}{0.000000,0.000000,0.000000}%
\pgfsetstrokecolor{currentstroke}%
\pgfsetstrokeopacity{0.000000}%
\pgfsetdash{}{0pt}%
\pgfpathmoveto{\pgfqpoint{2.882250in}{2.933130in}}%
\pgfpathlineto{\pgfqpoint{3.503444in}{2.933130in}}%
\pgfpathlineto{\pgfqpoint{3.503444in}{3.037611in}}%
\pgfpathlineto{\pgfqpoint{2.882250in}{3.037611in}}%
\pgfpathlineto{\pgfqpoint{2.882250in}{2.933130in}}%
\pgfpathclose%
\pgfusepath{fill}%
\end{pgfscope}%
\begin{pgfscope}%
\pgfpathrectangle{\pgfqpoint{2.882250in}{0.525000in}}{\pgfqpoint{1.293250in}{2.730000in}}%
\pgfusepath{clip}%
\pgfsetbuttcap%
\pgfsetmiterjoin%
\definecolor{currentfill}{rgb}{0.784314,0.784314,0.784314}%
\pgfsetfillcolor{currentfill}%
\pgfsetlinewidth{0.000000pt}%
\definecolor{currentstroke}{rgb}{0.000000,0.000000,0.000000}%
\pgfsetstrokecolor{currentstroke}%
\pgfsetstrokeopacity{0.000000}%
\pgfsetdash{}{0pt}%
\pgfpathmoveto{\pgfqpoint{2.882250in}{2.764611in}}%
\pgfpathlineto{\pgfqpoint{3.460728in}{2.764611in}}%
\pgfpathlineto{\pgfqpoint{3.460728in}{2.869093in}}%
\pgfpathlineto{\pgfqpoint{2.882250in}{2.869093in}}%
\pgfpathlineto{\pgfqpoint{2.882250in}{2.764611in}}%
\pgfpathclose%
\pgfusepath{fill}%
\end{pgfscope}%
\begin{pgfscope}%
\pgfpathrectangle{\pgfqpoint{2.882250in}{0.525000in}}{\pgfqpoint{1.293250in}{2.730000in}}%
\pgfusepath{clip}%
\pgfsetbuttcap%
\pgfsetmiterjoin%
\definecolor{currentfill}{rgb}{0.784314,0.784314,0.784314}%
\pgfsetfillcolor{currentfill}%
\pgfsetlinewidth{0.000000pt}%
\definecolor{currentstroke}{rgb}{0.000000,0.000000,0.000000}%
\pgfsetstrokecolor{currentstroke}%
\pgfsetstrokeopacity{0.000000}%
\pgfsetdash{}{0pt}%
\pgfpathmoveto{\pgfqpoint{2.882250in}{2.596093in}}%
\pgfpathlineto{\pgfqpoint{3.503444in}{2.596093in}}%
\pgfpathlineto{\pgfqpoint{3.503444in}{2.700574in}}%
\pgfpathlineto{\pgfqpoint{2.882250in}{2.700574in}}%
\pgfpathlineto{\pgfqpoint{2.882250in}{2.596093in}}%
\pgfpathclose%
\pgfusepath{fill}%
\end{pgfscope}%
\begin{pgfscope}%
\pgfpathrectangle{\pgfqpoint{2.882250in}{0.525000in}}{\pgfqpoint{1.293250in}{2.730000in}}%
\pgfusepath{clip}%
\pgfsetbuttcap%
\pgfsetmiterjoin%
\definecolor{currentfill}{rgb}{0.784314,0.784314,0.784314}%
\pgfsetfillcolor{currentfill}%
\pgfsetlinewidth{0.000000pt}%
\definecolor{currentstroke}{rgb}{0.000000,0.000000,0.000000}%
\pgfsetstrokecolor{currentstroke}%
\pgfsetstrokeopacity{0.000000}%
\pgfsetdash{}{0pt}%
\pgfpathmoveto{\pgfqpoint{2.882250in}{2.427574in}}%
\pgfpathlineto{\pgfqpoint{3.462176in}{2.427574in}}%
\pgfpathlineto{\pgfqpoint{3.462176in}{2.532056in}}%
\pgfpathlineto{\pgfqpoint{2.882250in}{2.532056in}}%
\pgfpathlineto{\pgfqpoint{2.882250in}{2.427574in}}%
\pgfpathclose%
\pgfusepath{fill}%
\end{pgfscope}%
\begin{pgfscope}%
\pgfpathrectangle{\pgfqpoint{2.882250in}{0.525000in}}{\pgfqpoint{1.293250in}{2.730000in}}%
\pgfusepath{clip}%
\pgfsetbuttcap%
\pgfsetmiterjoin%
\definecolor{currentfill}{rgb}{0.784314,0.784314,0.784314}%
\pgfsetfillcolor{currentfill}%
\pgfsetlinewidth{0.000000pt}%
\definecolor{currentstroke}{rgb}{0.000000,0.000000,0.000000}%
\pgfsetstrokecolor{currentstroke}%
\pgfsetstrokeopacity{0.000000}%
\pgfsetdash{}{0pt}%
\pgfpathmoveto{\pgfqpoint{2.882250in}{2.259056in}}%
\pgfpathlineto{\pgfqpoint{3.459280in}{2.259056in}}%
\pgfpathlineto{\pgfqpoint{3.459280in}{2.363537in}}%
\pgfpathlineto{\pgfqpoint{2.882250in}{2.363537in}}%
\pgfpathlineto{\pgfqpoint{2.882250in}{2.259056in}}%
\pgfpathclose%
\pgfusepath{fill}%
\end{pgfscope}%
\begin{pgfscope}%
\pgfpathrectangle{\pgfqpoint{2.882250in}{0.525000in}}{\pgfqpoint{1.293250in}{2.730000in}}%
\pgfusepath{clip}%
\pgfsetbuttcap%
\pgfsetmiterjoin%
\definecolor{currentfill}{rgb}{0.478431,0.478431,0.478431}%
\pgfsetfillcolor{currentfill}%
\pgfsetlinewidth{0.000000pt}%
\definecolor{currentstroke}{rgb}{0.000000,0.000000,0.000000}%
\pgfsetstrokecolor{currentstroke}%
\pgfsetstrokeopacity{0.000000}%
\pgfsetdash{}{0pt}%
\pgfpathmoveto{\pgfqpoint{2.882250in}{2.090537in}}%
\pgfpathlineto{\pgfqpoint{3.460004in}{2.090537in}}%
\pgfpathlineto{\pgfqpoint{3.460004in}{2.195019in}}%
\pgfpathlineto{\pgfqpoint{2.882250in}{2.195019in}}%
\pgfpathlineto{\pgfqpoint{2.882250in}{2.090537in}}%
\pgfpathclose%
\pgfusepath{fill}%
\end{pgfscope}%
\begin{pgfscope}%
\pgfpathrectangle{\pgfqpoint{2.882250in}{0.525000in}}{\pgfqpoint{1.293250in}{2.730000in}}%
\pgfusepath{clip}%
\pgfsetbuttcap%
\pgfsetmiterjoin%
\definecolor{currentfill}{rgb}{0.784314,0.784314,0.784314}%
\pgfsetfillcolor{currentfill}%
\pgfsetlinewidth{0.000000pt}%
\definecolor{currentstroke}{rgb}{0.000000,0.000000,0.000000}%
\pgfsetstrokecolor{currentstroke}%
\pgfsetstrokeopacity{0.000000}%
\pgfsetdash{}{0pt}%
\pgfpathmoveto{\pgfqpoint{2.882250in}{1.922019in}}%
\pgfpathlineto{\pgfqpoint{3.457832in}{1.922019in}}%
\pgfpathlineto{\pgfqpoint{3.457832in}{2.026500in}}%
\pgfpathlineto{\pgfqpoint{2.882250in}{2.026500in}}%
\pgfpathlineto{\pgfqpoint{2.882250in}{1.922019in}}%
\pgfpathclose%
\pgfusepath{fill}%
\end{pgfscope}%
\begin{pgfscope}%
\pgfpathrectangle{\pgfqpoint{2.882250in}{0.525000in}}{\pgfqpoint{1.293250in}{2.730000in}}%
\pgfusepath{clip}%
\pgfsetbuttcap%
\pgfsetmiterjoin%
\definecolor{currentfill}{rgb}{0.784314,0.784314,0.784314}%
\pgfsetfillcolor{currentfill}%
\pgfsetlinewidth{0.000000pt}%
\definecolor{currentstroke}{rgb}{0.000000,0.000000,0.000000}%
\pgfsetstrokecolor{currentstroke}%
\pgfsetstrokeopacity{0.000000}%
\pgfsetdash{}{0pt}%
\pgfpathmoveto{\pgfqpoint{2.882250in}{1.753500in}}%
\pgfpathlineto{\pgfqpoint{3.441904in}{1.753500in}}%
\pgfpathlineto{\pgfqpoint{3.441904in}{1.857981in}}%
\pgfpathlineto{\pgfqpoint{2.882250in}{1.857981in}}%
\pgfpathlineto{\pgfqpoint{2.882250in}{1.753500in}}%
\pgfpathclose%
\pgfusepath{fill}%
\end{pgfscope}%
\begin{pgfscope}%
\pgfpathrectangle{\pgfqpoint{2.882250in}{0.525000in}}{\pgfqpoint{1.293250in}{2.730000in}}%
\pgfusepath{clip}%
\pgfsetbuttcap%
\pgfsetmiterjoin%
\definecolor{currentfill}{rgb}{0.478431,0.478431,0.478431}%
\pgfsetfillcolor{currentfill}%
\pgfsetlinewidth{0.000000pt}%
\definecolor{currentstroke}{rgb}{0.000000,0.000000,0.000000}%
\pgfsetstrokecolor{currentstroke}%
\pgfsetstrokeopacity{0.000000}%
\pgfsetdash{}{0pt}%
\pgfpathmoveto{\pgfqpoint{2.882250in}{1.584981in}}%
\pgfpathlineto{\pgfqpoint{3.436112in}{1.584981in}}%
\pgfpathlineto{\pgfqpoint{3.436112in}{1.689463in}}%
\pgfpathlineto{\pgfqpoint{2.882250in}{1.689463in}}%
\pgfpathlineto{\pgfqpoint{2.882250in}{1.584981in}}%
\pgfpathclose%
\pgfusepath{fill}%
\end{pgfscope}%
\begin{pgfscope}%
\pgfpathrectangle{\pgfqpoint{2.882250in}{0.525000in}}{\pgfqpoint{1.293250in}{2.730000in}}%
\pgfusepath{clip}%
\pgfsetbuttcap%
\pgfsetmiterjoin%
\definecolor{currentfill}{rgb}{0.784314,0.784314,0.784314}%
\pgfsetfillcolor{currentfill}%
\pgfsetlinewidth{0.000000pt}%
\definecolor{currentstroke}{rgb}{0.000000,0.000000,0.000000}%
\pgfsetstrokecolor{currentstroke}%
\pgfsetstrokeopacity{0.000000}%
\pgfsetdash{}{0pt}%
\pgfpathmoveto{\pgfqpoint{2.882250in}{1.416463in}}%
\pgfpathlineto{\pgfqpoint{3.436112in}{1.416463in}}%
\pgfpathlineto{\pgfqpoint{3.436112in}{1.520944in}}%
\pgfpathlineto{\pgfqpoint{2.882250in}{1.520944in}}%
\pgfpathlineto{\pgfqpoint{2.882250in}{1.416463in}}%
\pgfpathclose%
\pgfusepath{fill}%
\end{pgfscope}%
\begin{pgfscope}%
\pgfpathrectangle{\pgfqpoint{2.882250in}{0.525000in}}{\pgfqpoint{1.293250in}{2.730000in}}%
\pgfusepath{clip}%
\pgfsetbuttcap%
\pgfsetmiterjoin%
\definecolor{currentfill}{rgb}{0.478431,0.478431,0.478431}%
\pgfsetfillcolor{currentfill}%
\pgfsetlinewidth{0.000000pt}%
\definecolor{currentstroke}{rgb}{0.000000,0.000000,0.000000}%
\pgfsetstrokecolor{currentstroke}%
\pgfsetstrokeopacity{0.000000}%
\pgfsetdash{}{0pt}%
\pgfpathmoveto{\pgfqpoint{2.882250in}{1.247944in}}%
\pgfpathlineto{\pgfqpoint{3.435388in}{1.247944in}}%
\pgfpathlineto{\pgfqpoint{3.435388in}{1.352426in}}%
\pgfpathlineto{\pgfqpoint{2.882250in}{1.352426in}}%
\pgfpathlineto{\pgfqpoint{2.882250in}{1.247944in}}%
\pgfpathclose%
\pgfusepath{fill}%
\end{pgfscope}%
\begin{pgfscope}%
\pgfpathrectangle{\pgfqpoint{2.882250in}{0.525000in}}{\pgfqpoint{1.293250in}{2.730000in}}%
\pgfusepath{clip}%
\pgfsetbuttcap%
\pgfsetmiterjoin%
\definecolor{currentfill}{rgb}{0.784314,0.784314,0.784314}%
\pgfsetfillcolor{currentfill}%
\pgfsetlinewidth{0.000000pt}%
\definecolor{currentstroke}{rgb}{0.000000,0.000000,0.000000}%
\pgfsetstrokecolor{currentstroke}%
\pgfsetstrokeopacity{0.000000}%
\pgfsetdash{}{0pt}%
\pgfpathmoveto{\pgfqpoint{2.882250in}{1.079426in}}%
\pgfpathlineto{\pgfqpoint{3.433216in}{1.079426in}}%
\pgfpathlineto{\pgfqpoint{3.433216in}{1.183907in}}%
\pgfpathlineto{\pgfqpoint{2.882250in}{1.183907in}}%
\pgfpathlineto{\pgfqpoint{2.882250in}{1.079426in}}%
\pgfpathclose%
\pgfusepath{fill}%
\end{pgfscope}%
\begin{pgfscope}%
\pgfpathrectangle{\pgfqpoint{2.882250in}{0.525000in}}{\pgfqpoint{1.293250in}{2.730000in}}%
\pgfusepath{clip}%
\pgfsetbuttcap%
\pgfsetmiterjoin%
\definecolor{currentfill}{rgb}{0.478431,0.478431,0.478431}%
\pgfsetfillcolor{currentfill}%
\pgfsetlinewidth{0.000000pt}%
\definecolor{currentstroke}{rgb}{0.000000,0.000000,0.000000}%
\pgfsetstrokecolor{currentstroke}%
\pgfsetstrokeopacity{0.000000}%
\pgfsetdash{}{0pt}%
\pgfpathmoveto{\pgfqpoint{2.882250in}{0.910907in}}%
\pgfpathlineto{\pgfqpoint{3.431768in}{0.910907in}}%
\pgfpathlineto{\pgfqpoint{3.431768in}{1.015389in}}%
\pgfpathlineto{\pgfqpoint{2.882250in}{1.015389in}}%
\pgfpathlineto{\pgfqpoint{2.882250in}{0.910907in}}%
\pgfpathclose%
\pgfusepath{fill}%
\end{pgfscope}%
\begin{pgfscope}%
\pgfpathrectangle{\pgfqpoint{2.882250in}{0.525000in}}{\pgfqpoint{1.293250in}{2.730000in}}%
\pgfusepath{clip}%
\pgfsetbuttcap%
\pgfsetmiterjoin%
\definecolor{currentfill}{rgb}{0.478431,0.478431,0.478431}%
\pgfsetfillcolor{currentfill}%
\pgfsetlinewidth{0.000000pt}%
\definecolor{currentstroke}{rgb}{0.000000,0.000000,0.000000}%
\pgfsetstrokecolor{currentstroke}%
\pgfsetstrokeopacity{0.000000}%
\pgfsetdash{}{0pt}%
\pgfpathmoveto{\pgfqpoint{2.882250in}{0.742389in}}%
\pgfpathlineto{\pgfqpoint{3.431768in}{0.742389in}}%
\pgfpathlineto{\pgfqpoint{3.431768in}{0.846870in}}%
\pgfpathlineto{\pgfqpoint{2.882250in}{0.846870in}}%
\pgfpathlineto{\pgfqpoint{2.882250in}{0.742389in}}%
\pgfpathclose%
\pgfusepath{fill}%
\end{pgfscope}%
\begin{pgfscope}%
\pgfpathrectangle{\pgfqpoint{2.882250in}{0.525000in}}{\pgfqpoint{1.293250in}{2.730000in}}%
\pgfusepath{clip}%
\pgfsetbuttcap%
\pgfsetmiterjoin%
\definecolor{currentfill}{rgb}{0.478431,0.305882,0.611765}%
\pgfsetfillcolor{currentfill}%
\pgfsetlinewidth{0.000000pt}%
\definecolor{currentstroke}{rgb}{0.000000,0.000000,0.000000}%
\pgfsetstrokecolor{currentstroke}%
\pgfsetstrokeopacity{0.000000}%
\pgfsetdash{}{0pt}%
\pgfpathmoveto{\pgfqpoint{2.882250in}{0.573870in}}%
\pgfpathlineto{\pgfqpoint{3.438284in}{0.573870in}}%
\pgfpathlineto{\pgfqpoint{3.438284in}{0.678352in}}%
\pgfpathlineto{\pgfqpoint{2.882250in}{0.678352in}}%
\pgfpathlineto{\pgfqpoint{2.882250in}{0.573870in}}%
\pgfpathclose%
\pgfusepath{fill}%
\end{pgfscope}%
\begin{pgfscope}%
\pgfsetrectcap%
\pgfsetmiterjoin%
\pgfsetlinewidth{0.803000pt}%
\definecolor{currentstroke}{rgb}{0.101961,0.101961,0.101961}%
\pgfsetstrokecolor{currentstroke}%
\pgfsetdash{}{0pt}%
\pgfpathmoveto{\pgfqpoint{2.882250in}{0.525000in}}%
\pgfpathlineto{\pgfqpoint{2.882250in}{3.255000in}}%
\pgfusepath{stroke}%
\end{pgfscope}%
\begin{pgfscope}%
\pgfsetrectcap%
\pgfsetmiterjoin%
\pgfsetlinewidth{0.803000pt}%
\definecolor{currentstroke}{rgb}{0.101961,0.101961,0.101961}%
\pgfsetstrokecolor{currentstroke}%
\pgfsetdash{}{0pt}%
\pgfpathmoveto{\pgfqpoint{2.882250in}{0.525000in}}%
\pgfpathlineto{\pgfqpoint{4.175500in}{0.525000in}}%
\pgfusepath{stroke}%
\end{pgfscope}%
\begin{pgfscope}%
\pgfpathrectangle{\pgfqpoint{2.882250in}{0.525000in}}{\pgfqpoint{1.293250in}{2.730000in}}%
\pgfusepath{clip}%
\pgfsetbuttcap%
\pgfsetroundjoin%
\definecolor{currentfill}{rgb}{0.101961,0.101961,0.101961}%
\pgfsetfillcolor{currentfill}%
\pgfsetlinewidth{0.301125pt}%
\definecolor{currentstroke}{rgb}{1.000000,1.000000,1.000000}%
\pgfsetstrokecolor{currentstroke}%
\pgfsetdash{}{0pt}%
\pgfsys@defobject{currentmarker}{\pgfqpoint{-0.013889in}{-0.013889in}}{\pgfqpoint{0.013889in}{0.013889in}}{%
\pgfpathmoveto{\pgfqpoint{0.000000in}{-0.013889in}}%
\pgfpathcurveto{\pgfqpoint{0.003683in}{-0.013889in}}{\pgfqpoint{0.007216in}{-0.012425in}}{\pgfqpoint{0.009821in}{-0.009821in}}%
\pgfpathcurveto{\pgfqpoint{0.012425in}{-0.007216in}}{\pgfqpoint{0.013889in}{-0.003683in}}{\pgfqpoint{0.013889in}{0.000000in}}%
\pgfpathcurveto{\pgfqpoint{0.013889in}{0.003683in}}{\pgfqpoint{0.012425in}{0.007216in}}{\pgfqpoint{0.009821in}{0.009821in}}%
\pgfpathcurveto{\pgfqpoint{0.007216in}{0.012425in}}{\pgfqpoint{0.003683in}{0.013889in}}{\pgfqpoint{0.000000in}{0.013889in}}%
\pgfpathcurveto{\pgfqpoint{-0.003683in}{0.013889in}}{\pgfqpoint{-0.007216in}{0.012425in}}{\pgfqpoint{-0.009821in}{0.009821in}}%
\pgfpathcurveto{\pgfqpoint{-0.012425in}{0.007216in}}{\pgfqpoint{-0.013889in}{0.003683in}}{\pgfqpoint{-0.013889in}{0.000000in}}%
\pgfpathcurveto{\pgfqpoint{-0.013889in}{-0.003683in}}{\pgfqpoint{-0.012425in}{-0.007216in}}{\pgfqpoint{-0.009821in}{-0.009821in}}%
\pgfpathcurveto{\pgfqpoint{-0.007216in}{-0.012425in}}{\pgfqpoint{-0.003683in}{-0.013889in}}{\pgfqpoint{0.000000in}{-0.013889in}}%
\pgfpathlineto{\pgfqpoint{0.000000in}{-0.013889in}}%
\pgfpathclose%
\pgfusepath{stroke,fill}%
}%
\begin{pgfscope}%
\pgfsys@transformshift{3.499532in}{3.153889in}%
\pgfsys@useobject{currentmarker}{}%
\end{pgfscope}%
\begin{pgfscope}%
\pgfsys@transformshift{3.498360in}{3.153889in}%
\pgfsys@useobject{currentmarker}{}%
\end{pgfscope}%
\begin{pgfscope}%
\pgfsys@transformshift{3.505328in}{3.153889in}%
\pgfsys@useobject{currentmarker}{}%
\end{pgfscope}%
\begin{pgfscope}%
\pgfsys@transformshift{3.504739in}{3.153889in}%
\pgfsys@useobject{currentmarker}{}%
\end{pgfscope}%
\end{pgfscope}%
\begin{pgfscope}%
\definecolor{textcolor}{rgb}{0.101961,0.101961,0.101961}%
\pgfsetstrokecolor{textcolor}%
\pgfsetfillcolor{textcolor}%
\pgftext[x=3.958234in,y=3.153889in,left,]{\color{textcolor}\sffamily\fontsize{6.000000}{7.200000}\bfseries\selectfont 0.60}%
\end{pgfscope}%
\begin{pgfscope}%
\pgfpathrectangle{\pgfqpoint{2.882250in}{0.525000in}}{\pgfqpoint{1.293250in}{2.730000in}}%
\pgfusepath{clip}%
\pgfsetbuttcap%
\pgfsetroundjoin%
\definecolor{currentfill}{rgb}{0.101961,0.101961,0.101961}%
\pgfsetfillcolor{currentfill}%
\pgfsetlinewidth{0.301125pt}%
\definecolor{currentstroke}{rgb}{1.000000,1.000000,1.000000}%
\pgfsetstrokecolor{currentstroke}%
\pgfsetdash{}{0pt}%
\pgfsys@defobject{currentmarker}{\pgfqpoint{-0.013889in}{-0.013889in}}{\pgfqpoint{0.013889in}{0.013889in}}{%
\pgfpathmoveto{\pgfqpoint{0.000000in}{-0.013889in}}%
\pgfpathcurveto{\pgfqpoint{0.003683in}{-0.013889in}}{\pgfqpoint{0.007216in}{-0.012425in}}{\pgfqpoint{0.009821in}{-0.009821in}}%
\pgfpathcurveto{\pgfqpoint{0.012425in}{-0.007216in}}{\pgfqpoint{0.013889in}{-0.003683in}}{\pgfqpoint{0.013889in}{0.000000in}}%
\pgfpathcurveto{\pgfqpoint{0.013889in}{0.003683in}}{\pgfqpoint{0.012425in}{0.007216in}}{\pgfqpoint{0.009821in}{0.009821in}}%
\pgfpathcurveto{\pgfqpoint{0.007216in}{0.012425in}}{\pgfqpoint{0.003683in}{0.013889in}}{\pgfqpoint{0.000000in}{0.013889in}}%
\pgfpathcurveto{\pgfqpoint{-0.003683in}{0.013889in}}{\pgfqpoint{-0.007216in}{0.012425in}}{\pgfqpoint{-0.009821in}{0.009821in}}%
\pgfpathcurveto{\pgfqpoint{-0.012425in}{0.007216in}}{\pgfqpoint{-0.013889in}{0.003683in}}{\pgfqpoint{-0.013889in}{0.000000in}}%
\pgfpathcurveto{\pgfqpoint{-0.013889in}{-0.003683in}}{\pgfqpoint{-0.012425in}{-0.007216in}}{\pgfqpoint{-0.009821in}{-0.009821in}}%
\pgfpathcurveto{\pgfqpoint{-0.007216in}{-0.012425in}}{\pgfqpoint{-0.003683in}{-0.013889in}}{\pgfqpoint{0.000000in}{-0.013889in}}%
\pgfpathlineto{\pgfqpoint{0.000000in}{-0.013889in}}%
\pgfpathclose%
\pgfusepath{stroke,fill}%
}%
\begin{pgfscope}%
\pgfsys@transformshift{3.502430in}{2.985370in}%
\pgfsys@useobject{currentmarker}{}%
\end{pgfscope}%
\begin{pgfscope}%
\pgfsys@transformshift{3.504172in}{2.985370in}%
\pgfsys@useobject{currentmarker}{}%
\end{pgfscope}%
\begin{pgfscope}%
\pgfsys@transformshift{3.502430in}{2.985370in}%
\pgfsys@useobject{currentmarker}{}%
\end{pgfscope}%
\begin{pgfscope}%
\pgfsys@transformshift{3.504739in}{2.985370in}%
\pgfsys@useobject{currentmarker}{}%
\end{pgfscope}%
\end{pgfscope}%
\begin{pgfscope}%
\definecolor{textcolor}{rgb}{0.101961,0.101961,0.101961}%
\pgfsetstrokecolor{textcolor}%
\pgfsetfillcolor{textcolor}%
\pgftext[x=3.958234in,y=2.985370in,left,]{\color{textcolor}\sffamily\fontsize{6.000000}{7.200000}\selectfont 0.60}%
\end{pgfscope}%
\begin{pgfscope}%
\pgfpathrectangle{\pgfqpoint{2.882250in}{0.525000in}}{\pgfqpoint{1.293250in}{2.730000in}}%
\pgfusepath{clip}%
\pgfsetbuttcap%
\pgfsetroundjoin%
\definecolor{currentfill}{rgb}{0.101961,0.101961,0.101961}%
\pgfsetfillcolor{currentfill}%
\pgfsetlinewidth{0.301125pt}%
\definecolor{currentstroke}{rgb}{1.000000,1.000000,1.000000}%
\pgfsetstrokecolor{currentstroke}%
\pgfsetdash{}{0pt}%
\pgfsys@defobject{currentmarker}{\pgfqpoint{-0.013889in}{-0.013889in}}{\pgfqpoint{0.013889in}{0.013889in}}{%
\pgfpathmoveto{\pgfqpoint{0.000000in}{-0.013889in}}%
\pgfpathcurveto{\pgfqpoint{0.003683in}{-0.013889in}}{\pgfqpoint{0.007216in}{-0.012425in}}{\pgfqpoint{0.009821in}{-0.009821in}}%
\pgfpathcurveto{\pgfqpoint{0.012425in}{-0.007216in}}{\pgfqpoint{0.013889in}{-0.003683in}}{\pgfqpoint{0.013889in}{0.000000in}}%
\pgfpathcurveto{\pgfqpoint{0.013889in}{0.003683in}}{\pgfqpoint{0.012425in}{0.007216in}}{\pgfqpoint{0.009821in}{0.009821in}}%
\pgfpathcurveto{\pgfqpoint{0.007216in}{0.012425in}}{\pgfqpoint{0.003683in}{0.013889in}}{\pgfqpoint{0.000000in}{0.013889in}}%
\pgfpathcurveto{\pgfqpoint{-0.003683in}{0.013889in}}{\pgfqpoint{-0.007216in}{0.012425in}}{\pgfqpoint{-0.009821in}{0.009821in}}%
\pgfpathcurveto{\pgfqpoint{-0.012425in}{0.007216in}}{\pgfqpoint{-0.013889in}{0.003683in}}{\pgfqpoint{-0.013889in}{0.000000in}}%
\pgfpathcurveto{\pgfqpoint{-0.013889in}{-0.003683in}}{\pgfqpoint{-0.012425in}{-0.007216in}}{\pgfqpoint{-0.009821in}{-0.009821in}}%
\pgfpathcurveto{\pgfqpoint{-0.007216in}{-0.012425in}}{\pgfqpoint{-0.003683in}{-0.013889in}}{\pgfqpoint{0.000000in}{-0.013889in}}%
\pgfpathlineto{\pgfqpoint{0.000000in}{-0.013889in}}%
\pgfpathclose%
\pgfusepath{stroke,fill}%
}%
\begin{pgfscope}%
\pgfsys@transformshift{3.461858in}{2.816852in}%
\pgfsys@useobject{currentmarker}{}%
\end{pgfscope}%
\begin{pgfscope}%
\pgfsys@transformshift{3.466392in}{2.816852in}%
\pgfsys@useobject{currentmarker}{}%
\end{pgfscope}%
\begin{pgfscope}%
\pgfsys@transformshift{3.456062in}{2.816852in}%
\pgfsys@useobject{currentmarker}{}%
\end{pgfscope}%
\begin{pgfscope}%
\pgfsys@transformshift{3.458629in}{2.816852in}%
\pgfsys@useobject{currentmarker}{}%
\end{pgfscope}%
\end{pgfscope}%
\begin{pgfscope}%
\definecolor{textcolor}{rgb}{0.101961,0.101961,0.101961}%
\pgfsetstrokecolor{textcolor}%
\pgfsetfillcolor{textcolor}%
\pgftext[x=3.958234in,y=2.816852in,left,]{\color{textcolor}\sffamily\fontsize{6.000000}{7.200000}\selectfont 0.56}%
\end{pgfscope}%
\begin{pgfscope}%
\pgfpathrectangle{\pgfqpoint{2.882250in}{0.525000in}}{\pgfqpoint{1.293250in}{2.730000in}}%
\pgfusepath{clip}%
\pgfsetbuttcap%
\pgfsetroundjoin%
\definecolor{currentfill}{rgb}{0.101961,0.101961,0.101961}%
\pgfsetfillcolor{currentfill}%
\pgfsetlinewidth{0.301125pt}%
\definecolor{currentstroke}{rgb}{1.000000,1.000000,1.000000}%
\pgfsetstrokecolor{currentstroke}%
\pgfsetdash{}{0pt}%
\pgfsys@defobject{currentmarker}{\pgfqpoint{-0.013889in}{-0.013889in}}{\pgfqpoint{0.013889in}{0.013889in}}{%
\pgfpathmoveto{\pgfqpoint{0.000000in}{-0.013889in}}%
\pgfpathcurveto{\pgfqpoint{0.003683in}{-0.013889in}}{\pgfqpoint{0.007216in}{-0.012425in}}{\pgfqpoint{0.009821in}{-0.009821in}}%
\pgfpathcurveto{\pgfqpoint{0.012425in}{-0.007216in}}{\pgfqpoint{0.013889in}{-0.003683in}}{\pgfqpoint{0.013889in}{0.000000in}}%
\pgfpathcurveto{\pgfqpoint{0.013889in}{0.003683in}}{\pgfqpoint{0.012425in}{0.007216in}}{\pgfqpoint{0.009821in}{0.009821in}}%
\pgfpathcurveto{\pgfqpoint{0.007216in}{0.012425in}}{\pgfqpoint{0.003683in}{0.013889in}}{\pgfqpoint{0.000000in}{0.013889in}}%
\pgfpathcurveto{\pgfqpoint{-0.003683in}{0.013889in}}{\pgfqpoint{-0.007216in}{0.012425in}}{\pgfqpoint{-0.009821in}{0.009821in}}%
\pgfpathcurveto{\pgfqpoint{-0.012425in}{0.007216in}}{\pgfqpoint{-0.013889in}{0.003683in}}{\pgfqpoint{-0.013889in}{0.000000in}}%
\pgfpathcurveto{\pgfqpoint{-0.013889in}{-0.003683in}}{\pgfqpoint{-0.012425in}{-0.007216in}}{\pgfqpoint{-0.009821in}{-0.009821in}}%
\pgfpathcurveto{\pgfqpoint{-0.007216in}{-0.012425in}}{\pgfqpoint{-0.003683in}{-0.013889in}}{\pgfqpoint{0.000000in}{-0.013889in}}%
\pgfpathlineto{\pgfqpoint{0.000000in}{-0.013889in}}%
\pgfpathclose%
\pgfusepath{stroke,fill}%
}%
\begin{pgfscope}%
\pgfsys@transformshift{3.502430in}{2.648333in}%
\pgfsys@useobject{currentmarker}{}%
\end{pgfscope}%
\begin{pgfscope}%
\pgfsys@transformshift{3.504172in}{2.648333in}%
\pgfsys@useobject{currentmarker}{}%
\end{pgfscope}%
\begin{pgfscope}%
\pgfsys@transformshift{3.502430in}{2.648333in}%
\pgfsys@useobject{currentmarker}{}%
\end{pgfscope}%
\begin{pgfscope}%
\pgfsys@transformshift{3.504739in}{2.648333in}%
\pgfsys@useobject{currentmarker}{}%
\end{pgfscope}%
\end{pgfscope}%
\begin{pgfscope}%
\definecolor{textcolor}{rgb}{0.101961,0.101961,0.101961}%
\pgfsetstrokecolor{textcolor}%
\pgfsetfillcolor{textcolor}%
\pgftext[x=3.958234in,y=2.648333in,left,]{\color{textcolor}\sffamily\fontsize{6.000000}{7.200000}\selectfont 0.60}%
\end{pgfscope}%
\begin{pgfscope}%
\pgfpathrectangle{\pgfqpoint{2.882250in}{0.525000in}}{\pgfqpoint{1.293250in}{2.730000in}}%
\pgfusepath{clip}%
\pgfsetbuttcap%
\pgfsetroundjoin%
\definecolor{currentfill}{rgb}{0.101961,0.101961,0.101961}%
\pgfsetfillcolor{currentfill}%
\pgfsetlinewidth{0.301125pt}%
\definecolor{currentstroke}{rgb}{1.000000,1.000000,1.000000}%
\pgfsetstrokecolor{currentstroke}%
\pgfsetdash{}{0pt}%
\pgfsys@defobject{currentmarker}{\pgfqpoint{-0.013889in}{-0.013889in}}{\pgfqpoint{0.013889in}{0.013889in}}{%
\pgfpathmoveto{\pgfqpoint{0.000000in}{-0.013889in}}%
\pgfpathcurveto{\pgfqpoint{0.003683in}{-0.013889in}}{\pgfqpoint{0.007216in}{-0.012425in}}{\pgfqpoint{0.009821in}{-0.009821in}}%
\pgfpathcurveto{\pgfqpoint{0.012425in}{-0.007216in}}{\pgfqpoint{0.013889in}{-0.003683in}}{\pgfqpoint{0.013889in}{0.000000in}}%
\pgfpathcurveto{\pgfqpoint{0.013889in}{0.003683in}}{\pgfqpoint{0.012425in}{0.007216in}}{\pgfqpoint{0.009821in}{0.009821in}}%
\pgfpathcurveto{\pgfqpoint{0.007216in}{0.012425in}}{\pgfqpoint{0.003683in}{0.013889in}}{\pgfqpoint{0.000000in}{0.013889in}}%
\pgfpathcurveto{\pgfqpoint{-0.003683in}{0.013889in}}{\pgfqpoint{-0.007216in}{0.012425in}}{\pgfqpoint{-0.009821in}{0.009821in}}%
\pgfpathcurveto{\pgfqpoint{-0.012425in}{0.007216in}}{\pgfqpoint{-0.013889in}{0.003683in}}{\pgfqpoint{-0.013889in}{0.000000in}}%
\pgfpathcurveto{\pgfqpoint{-0.013889in}{-0.003683in}}{\pgfqpoint{-0.012425in}{-0.007216in}}{\pgfqpoint{-0.009821in}{-0.009821in}}%
\pgfpathcurveto{\pgfqpoint{-0.007216in}{-0.012425in}}{\pgfqpoint{-0.003683in}{-0.013889in}}{\pgfqpoint{0.000000in}{-0.013889in}}%
\pgfpathlineto{\pgfqpoint{0.000000in}{-0.013889in}}%
\pgfpathclose%
\pgfusepath{stroke,fill}%
}%
\begin{pgfscope}%
\pgfsys@transformshift{3.461858in}{2.479815in}%
\pgfsys@useobject{currentmarker}{}%
\end{pgfscope}%
\begin{pgfscope}%
\pgfsys@transformshift{3.478017in}{2.479815in}%
\pgfsys@useobject{currentmarker}{}%
\end{pgfscope}%
\begin{pgfscope}%
\pgfsys@transformshift{3.464756in}{2.479815in}%
\pgfsys@useobject{currentmarker}{}%
\end{pgfscope}%
\begin{pgfscope}%
\pgfsys@transformshift{3.444219in}{2.479815in}%
\pgfsys@useobject{currentmarker}{}%
\end{pgfscope}%
\end{pgfscope}%
\begin{pgfscope}%
\definecolor{textcolor}{rgb}{0.101961,0.101961,0.101961}%
\pgfsetstrokecolor{textcolor}%
\pgfsetfillcolor{textcolor}%
\pgftext[x=3.958234in,y=2.479815in,left,]{\color{textcolor}\sffamily\fontsize{6.000000}{7.200000}\selectfont 0.56}%
\end{pgfscope}%
\begin{pgfscope}%
\pgfpathrectangle{\pgfqpoint{2.882250in}{0.525000in}}{\pgfqpoint{1.293250in}{2.730000in}}%
\pgfusepath{clip}%
\pgfsetbuttcap%
\pgfsetroundjoin%
\definecolor{currentfill}{rgb}{0.101961,0.101961,0.101961}%
\pgfsetfillcolor{currentfill}%
\pgfsetlinewidth{0.301125pt}%
\definecolor{currentstroke}{rgb}{1.000000,1.000000,1.000000}%
\pgfsetstrokecolor{currentstroke}%
\pgfsetdash{}{0pt}%
\pgfsys@defobject{currentmarker}{\pgfqpoint{-0.013889in}{-0.013889in}}{\pgfqpoint{0.013889in}{0.013889in}}{%
\pgfpathmoveto{\pgfqpoint{0.000000in}{-0.013889in}}%
\pgfpathcurveto{\pgfqpoint{0.003683in}{-0.013889in}}{\pgfqpoint{0.007216in}{-0.012425in}}{\pgfqpoint{0.009821in}{-0.009821in}}%
\pgfpathcurveto{\pgfqpoint{0.012425in}{-0.007216in}}{\pgfqpoint{0.013889in}{-0.003683in}}{\pgfqpoint{0.013889in}{0.000000in}}%
\pgfpathcurveto{\pgfqpoint{0.013889in}{0.003683in}}{\pgfqpoint{0.012425in}{0.007216in}}{\pgfqpoint{0.009821in}{0.009821in}}%
\pgfpathcurveto{\pgfqpoint{0.007216in}{0.012425in}}{\pgfqpoint{0.003683in}{0.013889in}}{\pgfqpoint{0.000000in}{0.013889in}}%
\pgfpathcurveto{\pgfqpoint{-0.003683in}{0.013889in}}{\pgfqpoint{-0.007216in}{0.012425in}}{\pgfqpoint{-0.009821in}{0.009821in}}%
\pgfpathcurveto{\pgfqpoint{-0.012425in}{0.007216in}}{\pgfqpoint{-0.013889in}{0.003683in}}{\pgfqpoint{-0.013889in}{0.000000in}}%
\pgfpathcurveto{\pgfqpoint{-0.013889in}{-0.003683in}}{\pgfqpoint{-0.012425in}{-0.007216in}}{\pgfqpoint{-0.009821in}{-0.009821in}}%
\pgfpathcurveto{\pgfqpoint{-0.007216in}{-0.012425in}}{\pgfqpoint{-0.003683in}{-0.013889in}}{\pgfqpoint{0.000000in}{-0.013889in}}%
\pgfpathlineto{\pgfqpoint{0.000000in}{-0.013889in}}%
\pgfpathclose%
\pgfusepath{stroke,fill}%
}%
\begin{pgfscope}%
\pgfsys@transformshift{3.438674in}{2.311296in}%
\pgfsys@useobject{currentmarker}{}%
\end{pgfscope}%
\begin{pgfscope}%
\pgfsys@transformshift{3.437330in}{2.311296in}%
\pgfsys@useobject{currentmarker}{}%
\end{pgfscope}%
\begin{pgfscope}%
\pgfsys@transformshift{3.470552in}{2.311296in}%
\pgfsys@useobject{currentmarker}{}%
\end{pgfscope}%
\begin{pgfscope}%
\pgfsys@transformshift{3.490330in}{2.311296in}%
\pgfsys@useobject{currentmarker}{}%
\end{pgfscope}%
\end{pgfscope}%
\begin{pgfscope}%
\definecolor{textcolor}{rgb}{0.101961,0.101961,0.101961}%
\pgfsetstrokecolor{textcolor}%
\pgfsetfillcolor{textcolor}%
\pgftext[x=3.958234in,y=2.311296in,left,]{\color{textcolor}\sffamily\fontsize{6.000000}{7.200000}\selectfont 0.56}%
\end{pgfscope}%
\begin{pgfscope}%
\pgfpathrectangle{\pgfqpoint{2.882250in}{0.525000in}}{\pgfqpoint{1.293250in}{2.730000in}}%
\pgfusepath{clip}%
\pgfsetbuttcap%
\pgfsetroundjoin%
\definecolor{currentfill}{rgb}{0.101961,0.101961,0.101961}%
\pgfsetfillcolor{currentfill}%
\pgfsetlinewidth{0.301125pt}%
\definecolor{currentstroke}{rgb}{1.000000,1.000000,1.000000}%
\pgfsetstrokecolor{currentstroke}%
\pgfsetdash{}{0pt}%
\pgfsys@defobject{currentmarker}{\pgfqpoint{-0.013889in}{-0.013889in}}{\pgfqpoint{0.013889in}{0.013889in}}{%
\pgfpathmoveto{\pgfqpoint{0.000000in}{-0.013889in}}%
\pgfpathcurveto{\pgfqpoint{0.003683in}{-0.013889in}}{\pgfqpoint{0.007216in}{-0.012425in}}{\pgfqpoint{0.009821in}{-0.009821in}}%
\pgfpathcurveto{\pgfqpoint{0.012425in}{-0.007216in}}{\pgfqpoint{0.013889in}{-0.003683in}}{\pgfqpoint{0.013889in}{0.000000in}}%
\pgfpathcurveto{\pgfqpoint{0.013889in}{0.003683in}}{\pgfqpoint{0.012425in}{0.007216in}}{\pgfqpoint{0.009821in}{0.009821in}}%
\pgfpathcurveto{\pgfqpoint{0.007216in}{0.012425in}}{\pgfqpoint{0.003683in}{0.013889in}}{\pgfqpoint{0.000000in}{0.013889in}}%
\pgfpathcurveto{\pgfqpoint{-0.003683in}{0.013889in}}{\pgfqpoint{-0.007216in}{0.012425in}}{\pgfqpoint{-0.009821in}{0.009821in}}%
\pgfpathcurveto{\pgfqpoint{-0.012425in}{0.007216in}}{\pgfqpoint{-0.013889in}{0.003683in}}{\pgfqpoint{-0.013889in}{0.000000in}}%
\pgfpathcurveto{\pgfqpoint{-0.013889in}{-0.003683in}}{\pgfqpoint{-0.012425in}{-0.007216in}}{\pgfqpoint{-0.009821in}{-0.009821in}}%
\pgfpathcurveto{\pgfqpoint{-0.007216in}{-0.012425in}}{\pgfqpoint{-0.003683in}{-0.013889in}}{\pgfqpoint{0.000000in}{-0.013889in}}%
\pgfpathlineto{\pgfqpoint{0.000000in}{-0.013889in}}%
\pgfpathclose%
\pgfusepath{stroke,fill}%
}%
\begin{pgfscope}%
\pgfsys@transformshift{3.438674in}{2.142778in}%
\pgfsys@useobject{currentmarker}{}%
\end{pgfscope}%
\begin{pgfscope}%
\pgfsys@transformshift{3.440237in}{2.142778in}%
\pgfsys@useobject{currentmarker}{}%
\end{pgfscope}%
\begin{pgfscope}%
\pgfsys@transformshift{3.467654in}{2.142778in}%
\pgfsys@useobject{currentmarker}{}%
\end{pgfscope}%
\begin{pgfscope}%
\pgfsys@transformshift{3.493212in}{2.142778in}%
\pgfsys@useobject{currentmarker}{}%
\end{pgfscope}%
\end{pgfscope}%
\begin{pgfscope}%
\definecolor{textcolor}{rgb}{0.101961,0.101961,0.101961}%
\pgfsetstrokecolor{textcolor}%
\pgfsetfillcolor{textcolor}%
\pgftext[x=3.958234in,y=2.142778in,left,]{\color{textcolor}\sffamily\fontsize{6.000000}{7.200000}\selectfont 0.56}%
\end{pgfscope}%
\begin{pgfscope}%
\pgfpathrectangle{\pgfqpoint{2.882250in}{0.525000in}}{\pgfqpoint{1.293250in}{2.730000in}}%
\pgfusepath{clip}%
\pgfsetbuttcap%
\pgfsetroundjoin%
\definecolor{currentfill}{rgb}{0.101961,0.101961,0.101961}%
\pgfsetfillcolor{currentfill}%
\pgfsetlinewidth{0.301125pt}%
\definecolor{currentstroke}{rgb}{1.000000,1.000000,1.000000}%
\pgfsetstrokecolor{currentstroke}%
\pgfsetdash{}{0pt}%
\pgfsys@defobject{currentmarker}{\pgfqpoint{-0.013889in}{-0.013889in}}{\pgfqpoint{0.013889in}{0.013889in}}{%
\pgfpathmoveto{\pgfqpoint{0.000000in}{-0.013889in}}%
\pgfpathcurveto{\pgfqpoint{0.003683in}{-0.013889in}}{\pgfqpoint{0.007216in}{-0.012425in}}{\pgfqpoint{0.009821in}{-0.009821in}}%
\pgfpathcurveto{\pgfqpoint{0.012425in}{-0.007216in}}{\pgfqpoint{0.013889in}{-0.003683in}}{\pgfqpoint{0.013889in}{0.000000in}}%
\pgfpathcurveto{\pgfqpoint{0.013889in}{0.003683in}}{\pgfqpoint{0.012425in}{0.007216in}}{\pgfqpoint{0.009821in}{0.009821in}}%
\pgfpathcurveto{\pgfqpoint{0.007216in}{0.012425in}}{\pgfqpoint{0.003683in}{0.013889in}}{\pgfqpoint{0.000000in}{0.013889in}}%
\pgfpathcurveto{\pgfqpoint{-0.003683in}{0.013889in}}{\pgfqpoint{-0.007216in}{0.012425in}}{\pgfqpoint{-0.009821in}{0.009821in}}%
\pgfpathcurveto{\pgfqpoint{-0.012425in}{0.007216in}}{\pgfqpoint{-0.013889in}{0.003683in}}{\pgfqpoint{-0.013889in}{0.000000in}}%
\pgfpathcurveto{\pgfqpoint{-0.013889in}{-0.003683in}}{\pgfqpoint{-0.012425in}{-0.007216in}}{\pgfqpoint{-0.009821in}{-0.009821in}}%
\pgfpathcurveto{\pgfqpoint{-0.007216in}{-0.012425in}}{\pgfqpoint{-0.003683in}{-0.013889in}}{\pgfqpoint{0.000000in}{-0.013889in}}%
\pgfpathlineto{\pgfqpoint{0.000000in}{-0.013889in}}%
\pgfpathclose%
\pgfusepath{stroke,fill}%
}%
\begin{pgfscope}%
\pgfsys@transformshift{3.438674in}{1.974259in}%
\pgfsys@useobject{currentmarker}{}%
\end{pgfscope}%
\begin{pgfscope}%
\pgfsys@transformshift{3.434424in}{1.974259in}%
\pgfsys@useobject{currentmarker}{}%
\end{pgfscope}%
\begin{pgfscope}%
\pgfsys@transformshift{3.467654in}{1.974259in}%
\pgfsys@useobject{currentmarker}{}%
\end{pgfscope}%
\begin{pgfscope}%
\pgfsys@transformshift{3.490330in}{1.974259in}%
\pgfsys@useobject{currentmarker}{}%
\end{pgfscope}%
\end{pgfscope}%
\begin{pgfscope}%
\definecolor{textcolor}{rgb}{0.101961,0.101961,0.101961}%
\pgfsetstrokecolor{textcolor}%
\pgfsetfillcolor{textcolor}%
\pgftext[x=3.958234in,y=1.974259in,left,]{\color{textcolor}\sffamily\fontsize{6.000000}{7.200000}\selectfont 0.56}%
\end{pgfscope}%
\begin{pgfscope}%
\pgfpathrectangle{\pgfqpoint{2.882250in}{0.525000in}}{\pgfqpoint{1.293250in}{2.730000in}}%
\pgfusepath{clip}%
\pgfsetbuttcap%
\pgfsetroundjoin%
\definecolor{currentfill}{rgb}{0.101961,0.101961,0.101961}%
\pgfsetfillcolor{currentfill}%
\pgfsetlinewidth{0.301125pt}%
\definecolor{currentstroke}{rgb}{1.000000,1.000000,1.000000}%
\pgfsetstrokecolor{currentstroke}%
\pgfsetdash{}{0pt}%
\pgfsys@defobject{currentmarker}{\pgfqpoint{-0.013889in}{-0.013889in}}{\pgfqpoint{0.013889in}{0.013889in}}{%
\pgfpathmoveto{\pgfqpoint{0.000000in}{-0.013889in}}%
\pgfpathcurveto{\pgfqpoint{0.003683in}{-0.013889in}}{\pgfqpoint{0.007216in}{-0.012425in}}{\pgfqpoint{0.009821in}{-0.009821in}}%
\pgfpathcurveto{\pgfqpoint{0.012425in}{-0.007216in}}{\pgfqpoint{0.013889in}{-0.003683in}}{\pgfqpoint{0.013889in}{0.000000in}}%
\pgfpathcurveto{\pgfqpoint{0.013889in}{0.003683in}}{\pgfqpoint{0.012425in}{0.007216in}}{\pgfqpoint{0.009821in}{0.009821in}}%
\pgfpathcurveto{\pgfqpoint{0.007216in}{0.012425in}}{\pgfqpoint{0.003683in}{0.013889in}}{\pgfqpoint{0.000000in}{0.013889in}}%
\pgfpathcurveto{\pgfqpoint{-0.003683in}{0.013889in}}{\pgfqpoint{-0.007216in}{0.012425in}}{\pgfqpoint{-0.009821in}{0.009821in}}%
\pgfpathcurveto{\pgfqpoint{-0.012425in}{0.007216in}}{\pgfqpoint{-0.013889in}{0.003683in}}{\pgfqpoint{-0.013889in}{0.000000in}}%
\pgfpathcurveto{\pgfqpoint{-0.013889in}{-0.003683in}}{\pgfqpoint{-0.012425in}{-0.007216in}}{\pgfqpoint{-0.009821in}{-0.009821in}}%
\pgfpathcurveto{\pgfqpoint{-0.007216in}{-0.012425in}}{\pgfqpoint{-0.003683in}{-0.013889in}}{\pgfqpoint{0.000000in}{-0.013889in}}%
\pgfpathlineto{\pgfqpoint{0.000000in}{-0.013889in}}%
\pgfpathclose%
\pgfusepath{stroke,fill}%
}%
\begin{pgfscope}%
\pgfsys@transformshift{3.435775in}{1.805741in}%
\pgfsys@useobject{currentmarker}{}%
\end{pgfscope}%
\begin{pgfscope}%
\pgfsys@transformshift{3.443143in}{1.805741in}%
\pgfsys@useobject{currentmarker}{}%
\end{pgfscope}%
\begin{pgfscope}%
\pgfsys@transformshift{3.444470in}{1.805741in}%
\pgfsys@useobject{currentmarker}{}%
\end{pgfscope}%
\begin{pgfscope}%
\pgfsys@transformshift{3.444219in}{1.805741in}%
\pgfsys@useobject{currentmarker}{}%
\end{pgfscope}%
\end{pgfscope}%
\begin{pgfscope}%
\definecolor{textcolor}{rgb}{0.101961,0.101961,0.101961}%
\pgfsetstrokecolor{textcolor}%
\pgfsetfillcolor{textcolor}%
\pgftext[x=3.958234in,y=1.805741in,left,]{\color{textcolor}\sffamily\fontsize{6.000000}{7.200000}\selectfont 0.54}%
\end{pgfscope}%
\begin{pgfscope}%
\pgfpathrectangle{\pgfqpoint{2.882250in}{0.525000in}}{\pgfqpoint{1.293250in}{2.730000in}}%
\pgfusepath{clip}%
\pgfsetbuttcap%
\pgfsetroundjoin%
\definecolor{currentfill}{rgb}{0.101961,0.101961,0.101961}%
\pgfsetfillcolor{currentfill}%
\pgfsetlinewidth{0.301125pt}%
\definecolor{currentstroke}{rgb}{1.000000,1.000000,1.000000}%
\pgfsetstrokecolor{currentstroke}%
\pgfsetdash{}{0pt}%
\pgfsys@defobject{currentmarker}{\pgfqpoint{-0.013889in}{-0.013889in}}{\pgfqpoint{0.013889in}{0.013889in}}{%
\pgfpathmoveto{\pgfqpoint{0.000000in}{-0.013889in}}%
\pgfpathcurveto{\pgfqpoint{0.003683in}{-0.013889in}}{\pgfqpoint{0.007216in}{-0.012425in}}{\pgfqpoint{0.009821in}{-0.009821in}}%
\pgfpathcurveto{\pgfqpoint{0.012425in}{-0.007216in}}{\pgfqpoint{0.013889in}{-0.003683in}}{\pgfqpoint{0.013889in}{0.000000in}}%
\pgfpathcurveto{\pgfqpoint{0.013889in}{0.003683in}}{\pgfqpoint{0.012425in}{0.007216in}}{\pgfqpoint{0.009821in}{0.009821in}}%
\pgfpathcurveto{\pgfqpoint{0.007216in}{0.012425in}}{\pgfqpoint{0.003683in}{0.013889in}}{\pgfqpoint{0.000000in}{0.013889in}}%
\pgfpathcurveto{\pgfqpoint{-0.003683in}{0.013889in}}{\pgfqpoint{-0.007216in}{0.012425in}}{\pgfqpoint{-0.009821in}{0.009821in}}%
\pgfpathcurveto{\pgfqpoint{-0.012425in}{0.007216in}}{\pgfqpoint{-0.013889in}{0.003683in}}{\pgfqpoint{-0.013889in}{0.000000in}}%
\pgfpathcurveto{\pgfqpoint{-0.013889in}{-0.003683in}}{\pgfqpoint{-0.012425in}{-0.007216in}}{\pgfqpoint{-0.009821in}{-0.009821in}}%
\pgfpathcurveto{\pgfqpoint{-0.007216in}{-0.012425in}}{\pgfqpoint{-0.003683in}{-0.013889in}}{\pgfqpoint{0.000000in}{-0.013889in}}%
\pgfpathlineto{\pgfqpoint{0.000000in}{-0.013889in}}%
\pgfpathclose%
\pgfusepath{stroke,fill}%
}%
\begin{pgfscope}%
\pgfsys@transformshift{3.438674in}{1.637222in}%
\pgfsys@useobject{currentmarker}{}%
\end{pgfscope}%
\begin{pgfscope}%
\pgfsys@transformshift{3.434424in}{1.637222in}%
\pgfsys@useobject{currentmarker}{}%
\end{pgfscope}%
\begin{pgfscope}%
\pgfsys@transformshift{3.438674in}{1.637222in}%
\pgfsys@useobject{currentmarker}{}%
\end{pgfscope}%
\begin{pgfscope}%
\pgfsys@transformshift{3.432692in}{1.637222in}%
\pgfsys@useobject{currentmarker}{}%
\end{pgfscope}%
\end{pgfscope}%
\begin{pgfscope}%
\definecolor{textcolor}{rgb}{0.101961,0.101961,0.101961}%
\pgfsetstrokecolor{textcolor}%
\pgfsetfillcolor{textcolor}%
\pgftext[x=3.958234in,y=1.637222in,left,]{\color{textcolor}\sffamily\fontsize{6.000000}{7.200000}\selectfont 0.54}%
\end{pgfscope}%
\begin{pgfscope}%
\pgfpathrectangle{\pgfqpoint{2.882250in}{0.525000in}}{\pgfqpoint{1.293250in}{2.730000in}}%
\pgfusepath{clip}%
\pgfsetbuttcap%
\pgfsetroundjoin%
\definecolor{currentfill}{rgb}{0.101961,0.101961,0.101961}%
\pgfsetfillcolor{currentfill}%
\pgfsetlinewidth{0.301125pt}%
\definecolor{currentstroke}{rgb}{1.000000,1.000000,1.000000}%
\pgfsetstrokecolor{currentstroke}%
\pgfsetdash{}{0pt}%
\pgfsys@defobject{currentmarker}{\pgfqpoint{-0.013889in}{-0.013889in}}{\pgfqpoint{0.013889in}{0.013889in}}{%
\pgfpathmoveto{\pgfqpoint{0.000000in}{-0.013889in}}%
\pgfpathcurveto{\pgfqpoint{0.003683in}{-0.013889in}}{\pgfqpoint{0.007216in}{-0.012425in}}{\pgfqpoint{0.009821in}{-0.009821in}}%
\pgfpathcurveto{\pgfqpoint{0.012425in}{-0.007216in}}{\pgfqpoint{0.013889in}{-0.003683in}}{\pgfqpoint{0.013889in}{0.000000in}}%
\pgfpathcurveto{\pgfqpoint{0.013889in}{0.003683in}}{\pgfqpoint{0.012425in}{0.007216in}}{\pgfqpoint{0.009821in}{0.009821in}}%
\pgfpathcurveto{\pgfqpoint{0.007216in}{0.012425in}}{\pgfqpoint{0.003683in}{0.013889in}}{\pgfqpoint{0.000000in}{0.013889in}}%
\pgfpathcurveto{\pgfqpoint{-0.003683in}{0.013889in}}{\pgfqpoint{-0.007216in}{0.012425in}}{\pgfqpoint{-0.009821in}{0.009821in}}%
\pgfpathcurveto{\pgfqpoint{-0.012425in}{0.007216in}}{\pgfqpoint{-0.013889in}{0.003683in}}{\pgfqpoint{-0.013889in}{0.000000in}}%
\pgfpathcurveto{\pgfqpoint{-0.013889in}{-0.003683in}}{\pgfqpoint{-0.012425in}{-0.007216in}}{\pgfqpoint{-0.009821in}{-0.009821in}}%
\pgfpathcurveto{\pgfqpoint{-0.007216in}{-0.012425in}}{\pgfqpoint{-0.003683in}{-0.013889in}}{\pgfqpoint{0.000000in}{-0.013889in}}%
\pgfpathlineto{\pgfqpoint{0.000000in}{-0.013889in}}%
\pgfpathclose%
\pgfusepath{stroke,fill}%
}%
\begin{pgfscope}%
\pgfsys@transformshift{3.438674in}{1.468704in}%
\pgfsys@useobject{currentmarker}{}%
\end{pgfscope}%
\begin{pgfscope}%
\pgfsys@transformshift{3.437330in}{1.468704in}%
\pgfsys@useobject{currentmarker}{}%
\end{pgfscope}%
\begin{pgfscope}%
\pgfsys@transformshift{3.438674in}{1.468704in}%
\pgfsys@useobject{currentmarker}{}%
\end{pgfscope}%
\begin{pgfscope}%
\pgfsys@transformshift{3.429810in}{1.468704in}%
\pgfsys@useobject{currentmarker}{}%
\end{pgfscope}%
\end{pgfscope}%
\begin{pgfscope}%
\definecolor{textcolor}{rgb}{0.101961,0.101961,0.101961}%
\pgfsetstrokecolor{textcolor}%
\pgfsetfillcolor{textcolor}%
\pgftext[x=3.958234in,y=1.468704in,left,]{\color{textcolor}\sffamily\fontsize{6.000000}{7.200000}\selectfont 0.54}%
\end{pgfscope}%
\begin{pgfscope}%
\pgfpathrectangle{\pgfqpoint{2.882250in}{0.525000in}}{\pgfqpoint{1.293250in}{2.730000in}}%
\pgfusepath{clip}%
\pgfsetbuttcap%
\pgfsetroundjoin%
\definecolor{currentfill}{rgb}{0.101961,0.101961,0.101961}%
\pgfsetfillcolor{currentfill}%
\pgfsetlinewidth{0.301125pt}%
\definecolor{currentstroke}{rgb}{1.000000,1.000000,1.000000}%
\pgfsetstrokecolor{currentstroke}%
\pgfsetdash{}{0pt}%
\pgfsys@defobject{currentmarker}{\pgfqpoint{-0.013889in}{-0.013889in}}{\pgfqpoint{0.013889in}{0.013889in}}{%
\pgfpathmoveto{\pgfqpoint{0.000000in}{-0.013889in}}%
\pgfpathcurveto{\pgfqpoint{0.003683in}{-0.013889in}}{\pgfqpoint{0.007216in}{-0.012425in}}{\pgfqpoint{0.009821in}{-0.009821in}}%
\pgfpathcurveto{\pgfqpoint{0.012425in}{-0.007216in}}{\pgfqpoint{0.013889in}{-0.003683in}}{\pgfqpoint{0.013889in}{0.000000in}}%
\pgfpathcurveto{\pgfqpoint{0.013889in}{0.003683in}}{\pgfqpoint{0.012425in}{0.007216in}}{\pgfqpoint{0.009821in}{0.009821in}}%
\pgfpathcurveto{\pgfqpoint{0.007216in}{0.012425in}}{\pgfqpoint{0.003683in}{0.013889in}}{\pgfqpoint{0.000000in}{0.013889in}}%
\pgfpathcurveto{\pgfqpoint{-0.003683in}{0.013889in}}{\pgfqpoint{-0.007216in}{0.012425in}}{\pgfqpoint{-0.009821in}{0.009821in}}%
\pgfpathcurveto{\pgfqpoint{-0.012425in}{0.007216in}}{\pgfqpoint{-0.013889in}{0.003683in}}{\pgfqpoint{-0.013889in}{0.000000in}}%
\pgfpathcurveto{\pgfqpoint{-0.013889in}{-0.003683in}}{\pgfqpoint{-0.012425in}{-0.007216in}}{\pgfqpoint{-0.009821in}{-0.009821in}}%
\pgfpathcurveto{\pgfqpoint{-0.007216in}{-0.012425in}}{\pgfqpoint{-0.003683in}{-0.013889in}}{\pgfqpoint{0.000000in}{-0.013889in}}%
\pgfpathlineto{\pgfqpoint{0.000000in}{-0.013889in}}%
\pgfpathclose%
\pgfusepath{stroke,fill}%
}%
\begin{pgfscope}%
\pgfsys@transformshift{3.438674in}{1.300185in}%
\pgfsys@useobject{currentmarker}{}%
\end{pgfscope}%
\begin{pgfscope}%
\pgfsys@transformshift{3.431518in}{1.300185in}%
\pgfsys@useobject{currentmarker}{}%
\end{pgfscope}%
\begin{pgfscope}%
\pgfsys@transformshift{3.432877in}{1.300185in}%
\pgfsys@useobject{currentmarker}{}%
\end{pgfscope}%
\begin{pgfscope}%
\pgfsys@transformshift{3.438456in}{1.300185in}%
\pgfsys@useobject{currentmarker}{}%
\end{pgfscope}%
\end{pgfscope}%
\begin{pgfscope}%
\definecolor{textcolor}{rgb}{0.101961,0.101961,0.101961}%
\pgfsetstrokecolor{textcolor}%
\pgfsetfillcolor{textcolor}%
\pgftext[x=3.958234in,y=1.300185in,left,]{\color{textcolor}\sffamily\fontsize{6.000000}{7.200000}\selectfont 0.53}%
\end{pgfscope}%
\begin{pgfscope}%
\pgfpathrectangle{\pgfqpoint{2.882250in}{0.525000in}}{\pgfqpoint{1.293250in}{2.730000in}}%
\pgfusepath{clip}%
\pgfsetbuttcap%
\pgfsetroundjoin%
\definecolor{currentfill}{rgb}{0.101961,0.101961,0.101961}%
\pgfsetfillcolor{currentfill}%
\pgfsetlinewidth{0.301125pt}%
\definecolor{currentstroke}{rgb}{1.000000,1.000000,1.000000}%
\pgfsetstrokecolor{currentstroke}%
\pgfsetdash{}{0pt}%
\pgfsys@defobject{currentmarker}{\pgfqpoint{-0.013889in}{-0.013889in}}{\pgfqpoint{0.013889in}{0.013889in}}{%
\pgfpathmoveto{\pgfqpoint{0.000000in}{-0.013889in}}%
\pgfpathcurveto{\pgfqpoint{0.003683in}{-0.013889in}}{\pgfqpoint{0.007216in}{-0.012425in}}{\pgfqpoint{0.009821in}{-0.009821in}}%
\pgfpathcurveto{\pgfqpoint{0.012425in}{-0.007216in}}{\pgfqpoint{0.013889in}{-0.003683in}}{\pgfqpoint{0.013889in}{0.000000in}}%
\pgfpathcurveto{\pgfqpoint{0.013889in}{0.003683in}}{\pgfqpoint{0.012425in}{0.007216in}}{\pgfqpoint{0.009821in}{0.009821in}}%
\pgfpathcurveto{\pgfqpoint{0.007216in}{0.012425in}}{\pgfqpoint{0.003683in}{0.013889in}}{\pgfqpoint{0.000000in}{0.013889in}}%
\pgfpathcurveto{\pgfqpoint{-0.003683in}{0.013889in}}{\pgfqpoint{-0.007216in}{0.012425in}}{\pgfqpoint{-0.009821in}{0.009821in}}%
\pgfpathcurveto{\pgfqpoint{-0.012425in}{0.007216in}}{\pgfqpoint{-0.013889in}{0.003683in}}{\pgfqpoint{-0.013889in}{0.000000in}}%
\pgfpathcurveto{\pgfqpoint{-0.013889in}{-0.003683in}}{\pgfqpoint{-0.012425in}{-0.007216in}}{\pgfqpoint{-0.009821in}{-0.009821in}}%
\pgfpathcurveto{\pgfqpoint{-0.007216in}{-0.012425in}}{\pgfqpoint{-0.003683in}{-0.013889in}}{\pgfqpoint{0.000000in}{-0.013889in}}%
\pgfpathlineto{\pgfqpoint{0.000000in}{-0.013889in}}%
\pgfpathclose%
\pgfusepath{stroke,fill}%
}%
\begin{pgfscope}%
\pgfsys@transformshift{3.435775in}{1.131667in}%
\pgfsys@useobject{currentmarker}{}%
\end{pgfscope}%
\begin{pgfscope}%
\pgfsys@transformshift{3.431518in}{1.131667in}%
\pgfsys@useobject{currentmarker}{}%
\end{pgfscope}%
\begin{pgfscope}%
\pgfsys@transformshift{3.432877in}{1.131667in}%
\pgfsys@useobject{currentmarker}{}%
\end{pgfscope}%
\begin{pgfscope}%
\pgfsys@transformshift{3.432692in}{1.131667in}%
\pgfsys@useobject{currentmarker}{}%
\end{pgfscope}%
\end{pgfscope}%
\begin{pgfscope}%
\definecolor{textcolor}{rgb}{0.101961,0.101961,0.101961}%
\pgfsetstrokecolor{textcolor}%
\pgfsetfillcolor{textcolor}%
\pgftext[x=3.958234in,y=1.131667in,left,]{\color{textcolor}\sffamily\fontsize{6.000000}{7.200000}\selectfont 0.53}%
\end{pgfscope}%
\begin{pgfscope}%
\pgfpathrectangle{\pgfqpoint{2.882250in}{0.525000in}}{\pgfqpoint{1.293250in}{2.730000in}}%
\pgfusepath{clip}%
\pgfsetbuttcap%
\pgfsetroundjoin%
\definecolor{currentfill}{rgb}{0.101961,0.101961,0.101961}%
\pgfsetfillcolor{currentfill}%
\pgfsetlinewidth{0.301125pt}%
\definecolor{currentstroke}{rgb}{1.000000,1.000000,1.000000}%
\pgfsetstrokecolor{currentstroke}%
\pgfsetdash{}{0pt}%
\pgfsys@defobject{currentmarker}{\pgfqpoint{-0.013889in}{-0.013889in}}{\pgfqpoint{0.013889in}{0.013889in}}{%
\pgfpathmoveto{\pgfqpoint{0.000000in}{-0.013889in}}%
\pgfpathcurveto{\pgfqpoint{0.003683in}{-0.013889in}}{\pgfqpoint{0.007216in}{-0.012425in}}{\pgfqpoint{0.009821in}{-0.009821in}}%
\pgfpathcurveto{\pgfqpoint{0.012425in}{-0.007216in}}{\pgfqpoint{0.013889in}{-0.003683in}}{\pgfqpoint{0.013889in}{0.000000in}}%
\pgfpathcurveto{\pgfqpoint{0.013889in}{0.003683in}}{\pgfqpoint{0.012425in}{0.007216in}}{\pgfqpoint{0.009821in}{0.009821in}}%
\pgfpathcurveto{\pgfqpoint{0.007216in}{0.012425in}}{\pgfqpoint{0.003683in}{0.013889in}}{\pgfqpoint{0.000000in}{0.013889in}}%
\pgfpathcurveto{\pgfqpoint{-0.003683in}{0.013889in}}{\pgfqpoint{-0.007216in}{0.012425in}}{\pgfqpoint{-0.009821in}{0.009821in}}%
\pgfpathcurveto{\pgfqpoint{-0.012425in}{0.007216in}}{\pgfqpoint{-0.013889in}{0.003683in}}{\pgfqpoint{-0.013889in}{0.000000in}}%
\pgfpathcurveto{\pgfqpoint{-0.013889in}{-0.003683in}}{\pgfqpoint{-0.012425in}{-0.007216in}}{\pgfqpoint{-0.009821in}{-0.009821in}}%
\pgfpathcurveto{\pgfqpoint{-0.007216in}{-0.012425in}}{\pgfqpoint{-0.003683in}{-0.013889in}}{\pgfqpoint{0.000000in}{-0.013889in}}%
\pgfpathlineto{\pgfqpoint{0.000000in}{-0.013889in}}%
\pgfpathclose%
\pgfusepath{stroke,fill}%
}%
\begin{pgfscope}%
\pgfsys@transformshift{3.432877in}{0.963148in}%
\pgfsys@useobject{currentmarker}{}%
\end{pgfscope}%
\begin{pgfscope}%
\pgfsys@transformshift{3.431518in}{0.963148in}%
\pgfsys@useobject{currentmarker}{}%
\end{pgfscope}%
\begin{pgfscope}%
\pgfsys@transformshift{3.429979in}{0.963148in}%
\pgfsys@useobject{currentmarker}{}%
\end{pgfscope}%
\begin{pgfscope}%
\pgfsys@transformshift{3.432692in}{0.963148in}%
\pgfsys@useobject{currentmarker}{}%
\end{pgfscope}%
\end{pgfscope}%
\begin{pgfscope}%
\definecolor{textcolor}{rgb}{0.101961,0.101961,0.101961}%
\pgfsetstrokecolor{textcolor}%
\pgfsetfillcolor{textcolor}%
\pgftext[x=3.958234in,y=0.963148in,left,]{\color{textcolor}\sffamily\fontsize{6.000000}{7.200000}\selectfont 0.53}%
\end{pgfscope}%
\begin{pgfscope}%
\pgfpathrectangle{\pgfqpoint{2.882250in}{0.525000in}}{\pgfqpoint{1.293250in}{2.730000in}}%
\pgfusepath{clip}%
\pgfsetbuttcap%
\pgfsetroundjoin%
\definecolor{currentfill}{rgb}{0.101961,0.101961,0.101961}%
\pgfsetfillcolor{currentfill}%
\pgfsetlinewidth{0.301125pt}%
\definecolor{currentstroke}{rgb}{1.000000,1.000000,1.000000}%
\pgfsetstrokecolor{currentstroke}%
\pgfsetdash{}{0pt}%
\pgfsys@defobject{currentmarker}{\pgfqpoint{-0.013889in}{-0.013889in}}{\pgfqpoint{0.013889in}{0.013889in}}{%
\pgfpathmoveto{\pgfqpoint{0.000000in}{-0.013889in}}%
\pgfpathcurveto{\pgfqpoint{0.003683in}{-0.013889in}}{\pgfqpoint{0.007216in}{-0.012425in}}{\pgfqpoint{0.009821in}{-0.009821in}}%
\pgfpathcurveto{\pgfqpoint{0.012425in}{-0.007216in}}{\pgfqpoint{0.013889in}{-0.003683in}}{\pgfqpoint{0.013889in}{0.000000in}}%
\pgfpathcurveto{\pgfqpoint{0.013889in}{0.003683in}}{\pgfqpoint{0.012425in}{0.007216in}}{\pgfqpoint{0.009821in}{0.009821in}}%
\pgfpathcurveto{\pgfqpoint{0.007216in}{0.012425in}}{\pgfqpoint{0.003683in}{0.013889in}}{\pgfqpoint{0.000000in}{0.013889in}}%
\pgfpathcurveto{\pgfqpoint{-0.003683in}{0.013889in}}{\pgfqpoint{-0.007216in}{0.012425in}}{\pgfqpoint{-0.009821in}{0.009821in}}%
\pgfpathcurveto{\pgfqpoint{-0.012425in}{0.007216in}}{\pgfqpoint{-0.013889in}{0.003683in}}{\pgfqpoint{-0.013889in}{0.000000in}}%
\pgfpathcurveto{\pgfqpoint{-0.013889in}{-0.003683in}}{\pgfqpoint{-0.012425in}{-0.007216in}}{\pgfqpoint{-0.009821in}{-0.009821in}}%
\pgfpathcurveto{\pgfqpoint{-0.007216in}{-0.012425in}}{\pgfqpoint{-0.003683in}{-0.013889in}}{\pgfqpoint{0.000000in}{-0.013889in}}%
\pgfpathlineto{\pgfqpoint{0.000000in}{-0.013889in}}%
\pgfpathclose%
\pgfusepath{stroke,fill}%
}%
\begin{pgfscope}%
\pgfsys@transformshift{3.432877in}{0.794630in}%
\pgfsys@useobject{currentmarker}{}%
\end{pgfscope}%
\begin{pgfscope}%
\pgfsys@transformshift{3.431518in}{0.794630in}%
\pgfsys@useobject{currentmarker}{}%
\end{pgfscope}%
\begin{pgfscope}%
\pgfsys@transformshift{3.432877in}{0.794630in}%
\pgfsys@useobject{currentmarker}{}%
\end{pgfscope}%
\begin{pgfscope}%
\pgfsys@transformshift{3.429810in}{0.794630in}%
\pgfsys@useobject{currentmarker}{}%
\end{pgfscope}%
\end{pgfscope}%
\begin{pgfscope}%
\definecolor{textcolor}{rgb}{0.101961,0.101961,0.101961}%
\pgfsetstrokecolor{textcolor}%
\pgfsetfillcolor{textcolor}%
\pgftext[x=3.958234in,y=0.794630in,left,]{\color{textcolor}\sffamily\fontsize{6.000000}{7.200000}\selectfont 0.53}%
\end{pgfscope}%
\begin{pgfscope}%
\pgfpathrectangle{\pgfqpoint{2.882250in}{0.525000in}}{\pgfqpoint{1.293250in}{2.730000in}}%
\pgfusepath{clip}%
\pgfsetbuttcap%
\pgfsetroundjoin%
\definecolor{currentfill}{rgb}{0.101961,0.101961,0.101961}%
\pgfsetfillcolor{currentfill}%
\pgfsetlinewidth{0.301125pt}%
\definecolor{currentstroke}{rgb}{1.000000,1.000000,1.000000}%
\pgfsetstrokecolor{currentstroke}%
\pgfsetdash{}{0pt}%
\pgfsys@defobject{currentmarker}{\pgfqpoint{-0.013889in}{-0.013889in}}{\pgfqpoint{0.013889in}{0.013889in}}{%
\pgfpathmoveto{\pgfqpoint{0.000000in}{-0.013889in}}%
\pgfpathcurveto{\pgfqpoint{0.003683in}{-0.013889in}}{\pgfqpoint{0.007216in}{-0.012425in}}{\pgfqpoint{0.009821in}{-0.009821in}}%
\pgfpathcurveto{\pgfqpoint{0.012425in}{-0.007216in}}{\pgfqpoint{0.013889in}{-0.003683in}}{\pgfqpoint{0.013889in}{0.000000in}}%
\pgfpathcurveto{\pgfqpoint{0.013889in}{0.003683in}}{\pgfqpoint{0.012425in}{0.007216in}}{\pgfqpoint{0.009821in}{0.009821in}}%
\pgfpathcurveto{\pgfqpoint{0.007216in}{0.012425in}}{\pgfqpoint{0.003683in}{0.013889in}}{\pgfqpoint{0.000000in}{0.013889in}}%
\pgfpathcurveto{\pgfqpoint{-0.003683in}{0.013889in}}{\pgfqpoint{-0.007216in}{0.012425in}}{\pgfqpoint{-0.009821in}{0.009821in}}%
\pgfpathcurveto{\pgfqpoint{-0.012425in}{0.007216in}}{\pgfqpoint{-0.013889in}{0.003683in}}{\pgfqpoint{-0.013889in}{0.000000in}}%
\pgfpathcurveto{\pgfqpoint{-0.013889in}{-0.003683in}}{\pgfqpoint{-0.012425in}{-0.007216in}}{\pgfqpoint{-0.009821in}{-0.009821in}}%
\pgfpathcurveto{\pgfqpoint{-0.007216in}{-0.012425in}}{\pgfqpoint{-0.003683in}{-0.013889in}}{\pgfqpoint{0.000000in}{-0.013889in}}%
\pgfpathlineto{\pgfqpoint{0.000000in}{-0.013889in}}%
\pgfpathclose%
\pgfusepath{stroke,fill}%
}%
\begin{pgfscope}%
\pgfsys@transformshift{3.438674in}{0.626111in}%
\pgfsys@useobject{currentmarker}{}%
\end{pgfscope}%
\begin{pgfscope}%
\pgfsys@transformshift{3.440237in}{0.626111in}%
\pgfsys@useobject{currentmarker}{}%
\end{pgfscope}%
\begin{pgfscope}%
\pgfsys@transformshift{3.435775in}{0.626111in}%
\pgfsys@useobject{currentmarker}{}%
\end{pgfscope}%
\begin{pgfscope}%
\pgfsys@transformshift{3.438456in}{0.626111in}%
\pgfsys@useobject{currentmarker}{}%
\end{pgfscope}%
\end{pgfscope}%
\begin{pgfscope}%
\definecolor{textcolor}{rgb}{0.101961,0.101961,0.101961}%
\pgfsetstrokecolor{textcolor}%
\pgfsetfillcolor{textcolor}%
\pgftext[x=3.958234in,y=0.626111in,left,]{\color{textcolor}\sffamily\fontsize{6.000000}{7.200000}\selectfont 0.54}%
\end{pgfscope}%
\begin{pgfscope}%
\pgfsetbuttcap%
\pgfsetmiterjoin%
\definecolor{currentfill}{rgb}{1.000000,1.000000,1.000000}%
\pgfsetfillcolor{currentfill}%
\pgfsetlinewidth{0.000000pt}%
\definecolor{currentstroke}{rgb}{0.000000,0.000000,0.000000}%
\pgfsetstrokecolor{currentstroke}%
\pgfsetstrokeopacity{0.000000}%
\pgfsetdash{}{0pt}%
\pgfpathmoveto{\pgfqpoint{4.259500in}{0.525000in}}%
\pgfpathlineto{\pgfqpoint{5.552750in}{0.525000in}}%
\pgfpathlineto{\pgfqpoint{5.552750in}{3.255000in}}%
\pgfpathlineto{\pgfqpoint{4.259500in}{3.255000in}}%
\pgfpathlineto{\pgfqpoint{4.259500in}{0.525000in}}%
\pgfpathclose%
\pgfusepath{fill}%
\end{pgfscope}%
\begin{pgfscope}%
\pgfpathrectangle{\pgfqpoint{4.259500in}{0.525000in}}{\pgfqpoint{1.293250in}{2.730000in}}%
\pgfusepath{clip}%
\pgfsetrectcap%
\pgfsetroundjoin%
\pgfsetlinewidth{0.401500pt}%
\definecolor{currentstroke}{rgb}{0.901961,0.901961,0.901961}%
\pgfsetstrokecolor{currentstroke}%
\pgfsetdash{}{0pt}%
\pgfpathmoveto{\pgfqpoint{4.259500in}{0.525000in}}%
\pgfpathlineto{\pgfqpoint{4.259500in}{3.255000in}}%
\pgfusepath{stroke}%
\end{pgfscope}%
\begin{pgfscope}%
\pgfsetbuttcap%
\pgfsetroundjoin%
\definecolor{currentfill}{rgb}{0.101961,0.101961,0.101961}%
\pgfsetfillcolor{currentfill}%
\pgfsetlinewidth{0.803000pt}%
\definecolor{currentstroke}{rgb}{0.101961,0.101961,0.101961}%
\pgfsetstrokecolor{currentstroke}%
\pgfsetdash{}{0pt}%
\pgfsys@defobject{currentmarker}{\pgfqpoint{0.000000in}{-0.027778in}}{\pgfqpoint{0.000000in}{0.000000in}}{%
\pgfpathmoveto{\pgfqpoint{0.000000in}{0.000000in}}%
\pgfpathlineto{\pgfqpoint{0.000000in}{-0.027778in}}%
\pgfusepath{stroke,fill}%
}%
\begin{pgfscope}%
\pgfsys@transformshift{4.259500in}{0.525000in}%
\pgfsys@useobject{currentmarker}{}%
\end{pgfscope}%
\end{pgfscope}%
\begin{pgfscope}%
\definecolor{textcolor}{rgb}{0.101961,0.101961,0.101961}%
\pgfsetstrokecolor{textcolor}%
\pgfsetfillcolor{textcolor}%
\pgftext[x=4.259500in,y=0.448611in,,top]{\color{textcolor}\sffamily\fontsize{7.000000}{8.400000}\selectfont 0}%
\end{pgfscope}%
\begin{pgfscope}%
\pgfpathrectangle{\pgfqpoint{4.259500in}{0.525000in}}{\pgfqpoint{1.293250in}{2.730000in}}%
\pgfusepath{clip}%
\pgfsetrectcap%
\pgfsetroundjoin%
\pgfsetlinewidth{0.401500pt}%
\definecolor{currentstroke}{rgb}{0.901961,0.901961,0.901961}%
\pgfsetstrokecolor{currentstroke}%
\pgfsetdash{}{0pt}%
\pgfpathmoveto{\pgfqpoint{4.776800in}{0.525000in}}%
\pgfpathlineto{\pgfqpoint{4.776800in}{3.255000in}}%
\pgfusepath{stroke}%
\end{pgfscope}%
\begin{pgfscope}%
\pgfsetbuttcap%
\pgfsetroundjoin%
\definecolor{currentfill}{rgb}{0.101961,0.101961,0.101961}%
\pgfsetfillcolor{currentfill}%
\pgfsetlinewidth{0.803000pt}%
\definecolor{currentstroke}{rgb}{0.101961,0.101961,0.101961}%
\pgfsetstrokecolor{currentstroke}%
\pgfsetdash{}{0pt}%
\pgfsys@defobject{currentmarker}{\pgfqpoint{0.000000in}{-0.027778in}}{\pgfqpoint{0.000000in}{0.000000in}}{%
\pgfpathmoveto{\pgfqpoint{0.000000in}{0.000000in}}%
\pgfpathlineto{\pgfqpoint{0.000000in}{-0.027778in}}%
\pgfusepath{stroke,fill}%
}%
\begin{pgfscope}%
\pgfsys@transformshift{4.776800in}{0.525000in}%
\pgfsys@useobject{currentmarker}{}%
\end{pgfscope}%
\end{pgfscope}%
\begin{pgfscope}%
\definecolor{textcolor}{rgb}{0.101961,0.101961,0.101961}%
\pgfsetstrokecolor{textcolor}%
\pgfsetfillcolor{textcolor}%
\pgftext[x=4.776800in,y=0.448611in,,top]{\color{textcolor}\sffamily\fontsize{7.000000}{8.400000}\selectfont 0.5}%
\end{pgfscope}%
\begin{pgfscope}%
\pgfpathrectangle{\pgfqpoint{4.259500in}{0.525000in}}{\pgfqpoint{1.293250in}{2.730000in}}%
\pgfusepath{clip}%
\pgfsetrectcap%
\pgfsetroundjoin%
\pgfsetlinewidth{0.401500pt}%
\definecolor{currentstroke}{rgb}{0.901961,0.901961,0.901961}%
\pgfsetstrokecolor{currentstroke}%
\pgfsetdash{}{0pt}%
\pgfpathmoveto{\pgfqpoint{5.294100in}{0.525000in}}%
\pgfpathlineto{\pgfqpoint{5.294100in}{3.255000in}}%
\pgfusepath{stroke}%
\end{pgfscope}%
\begin{pgfscope}%
\pgfsetbuttcap%
\pgfsetroundjoin%
\definecolor{currentfill}{rgb}{0.101961,0.101961,0.101961}%
\pgfsetfillcolor{currentfill}%
\pgfsetlinewidth{0.803000pt}%
\definecolor{currentstroke}{rgb}{0.101961,0.101961,0.101961}%
\pgfsetstrokecolor{currentstroke}%
\pgfsetdash{}{0pt}%
\pgfsys@defobject{currentmarker}{\pgfqpoint{0.000000in}{-0.027778in}}{\pgfqpoint{0.000000in}{0.000000in}}{%
\pgfpathmoveto{\pgfqpoint{0.000000in}{0.000000in}}%
\pgfpathlineto{\pgfqpoint{0.000000in}{-0.027778in}}%
\pgfusepath{stroke,fill}%
}%
\begin{pgfscope}%
\pgfsys@transformshift{5.294100in}{0.525000in}%
\pgfsys@useobject{currentmarker}{}%
\end{pgfscope}%
\end{pgfscope}%
\begin{pgfscope}%
\definecolor{textcolor}{rgb}{0.101961,0.101961,0.101961}%
\pgfsetstrokecolor{textcolor}%
\pgfsetfillcolor{textcolor}%
\pgftext[x=5.294100in,y=0.448611in,,top]{\color{textcolor}\sffamily\fontsize{7.000000}{8.400000}\selectfont 1}%
\end{pgfscope}%
\begin{pgfscope}%
\definecolor{textcolor}{rgb}{0.101961,0.101961,0.101961}%
\pgfsetstrokecolor{textcolor}%
\pgfsetfillcolor{textcolor}%
\pgftext[x=4.906125in,y=0.334414in,,top]{\color{textcolor}\sffamily\fontsize{7.500000}{9.000000}\selectfont Pooled}%
\end{pgfscope}%
\begin{pgfscope}%
\pgfpathrectangle{\pgfqpoint{4.259500in}{0.525000in}}{\pgfqpoint{1.293250in}{2.730000in}}%
\pgfusepath{clip}%
\pgfsetbuttcap%
\pgfsetmiterjoin%
\definecolor{currentfill}{rgb}{0.156863,0.356863,0.568627}%
\pgfsetfillcolor{currentfill}%
\pgfsetlinewidth{0.000000pt}%
\definecolor{currentstroke}{rgb}{0.000000,0.000000,0.000000}%
\pgfsetstrokecolor{currentstroke}%
\pgfsetstrokeopacity{0.000000}%
\pgfsetdash{}{0pt}%
\pgfpathmoveto{\pgfqpoint{4.259500in}{3.101648in}}%
\pgfpathlineto{\pgfqpoint{4.857341in}{3.101648in}}%
\pgfpathlineto{\pgfqpoint{4.857341in}{3.206130in}}%
\pgfpathlineto{\pgfqpoint{4.259500in}{3.206130in}}%
\pgfpathlineto{\pgfqpoint{4.259500in}{3.101648in}}%
\pgfpathclose%
\pgfusepath{fill}%
\end{pgfscope}%
\begin{pgfscope}%
\pgfpathrectangle{\pgfqpoint{4.259500in}{0.525000in}}{\pgfqpoint{1.293250in}{2.730000in}}%
\pgfusepath{clip}%
\pgfsetbuttcap%
\pgfsetmiterjoin%
\definecolor{currentfill}{rgb}{0.784314,0.784314,0.784314}%
\pgfsetfillcolor{currentfill}%
\pgfsetlinewidth{0.000000pt}%
\definecolor{currentstroke}{rgb}{0.000000,0.000000,0.000000}%
\pgfsetstrokecolor{currentstroke}%
\pgfsetstrokeopacity{0.000000}%
\pgfsetdash{}{0pt}%
\pgfpathmoveto{\pgfqpoint{4.259500in}{2.933130in}}%
\pgfpathlineto{\pgfqpoint{4.780584in}{2.933130in}}%
\pgfpathlineto{\pgfqpoint{4.780584in}{3.037611in}}%
\pgfpathlineto{\pgfqpoint{4.259500in}{3.037611in}}%
\pgfpathlineto{\pgfqpoint{4.259500in}{2.933130in}}%
\pgfpathclose%
\pgfusepath{fill}%
\end{pgfscope}%
\begin{pgfscope}%
\pgfpathrectangle{\pgfqpoint{4.259500in}{0.525000in}}{\pgfqpoint{1.293250in}{2.730000in}}%
\pgfusepath{clip}%
\pgfsetbuttcap%
\pgfsetmiterjoin%
\definecolor{currentfill}{rgb}{0.784314,0.784314,0.784314}%
\pgfsetfillcolor{currentfill}%
\pgfsetlinewidth{0.000000pt}%
\definecolor{currentstroke}{rgb}{0.000000,0.000000,0.000000}%
\pgfsetstrokecolor{currentstroke}%
\pgfsetstrokeopacity{0.000000}%
\pgfsetdash{}{0pt}%
\pgfpathmoveto{\pgfqpoint{4.259500in}{2.764611in}}%
\pgfpathlineto{\pgfqpoint{4.780584in}{2.764611in}}%
\pgfpathlineto{\pgfqpoint{4.780584in}{2.869093in}}%
\pgfpathlineto{\pgfqpoint{4.259500in}{2.869093in}}%
\pgfpathlineto{\pgfqpoint{4.259500in}{2.764611in}}%
\pgfpathclose%
\pgfusepath{fill}%
\end{pgfscope}%
\begin{pgfscope}%
\pgfpathrectangle{\pgfqpoint{4.259500in}{0.525000in}}{\pgfqpoint{1.293250in}{2.730000in}}%
\pgfusepath{clip}%
\pgfsetbuttcap%
\pgfsetmiterjoin%
\definecolor{currentfill}{rgb}{0.784314,0.784314,0.784314}%
\pgfsetfillcolor{currentfill}%
\pgfsetlinewidth{0.000000pt}%
\definecolor{currentstroke}{rgb}{0.000000,0.000000,0.000000}%
\pgfsetstrokecolor{currentstroke}%
\pgfsetstrokeopacity{0.000000}%
\pgfsetdash{}{0pt}%
\pgfpathmoveto{\pgfqpoint{4.259500in}{2.596093in}}%
\pgfpathlineto{\pgfqpoint{4.777881in}{2.596093in}}%
\pgfpathlineto{\pgfqpoint{4.777881in}{2.700574in}}%
\pgfpathlineto{\pgfqpoint{4.259500in}{2.700574in}}%
\pgfpathlineto{\pgfqpoint{4.259500in}{2.596093in}}%
\pgfpathclose%
\pgfusepath{fill}%
\end{pgfscope}%
\begin{pgfscope}%
\pgfpathrectangle{\pgfqpoint{4.259500in}{0.525000in}}{\pgfqpoint{1.293250in}{2.730000in}}%
\pgfusepath{clip}%
\pgfsetbuttcap%
\pgfsetmiterjoin%
\definecolor{currentfill}{rgb}{0.784314,0.784314,0.784314}%
\pgfsetfillcolor{currentfill}%
\pgfsetlinewidth{0.000000pt}%
\definecolor{currentstroke}{rgb}{0.000000,0.000000,0.000000}%
\pgfsetstrokecolor{currentstroke}%
\pgfsetstrokeopacity{0.000000}%
\pgfsetdash{}{0pt}%
\pgfpathmoveto{\pgfqpoint{4.259500in}{2.427574in}}%
\pgfpathlineto{\pgfqpoint{4.765449in}{2.427574in}}%
\pgfpathlineto{\pgfqpoint{4.765449in}{2.532056in}}%
\pgfpathlineto{\pgfqpoint{4.259500in}{2.532056in}}%
\pgfpathlineto{\pgfqpoint{4.259500in}{2.427574in}}%
\pgfpathclose%
\pgfusepath{fill}%
\end{pgfscope}%
\begin{pgfscope}%
\pgfpathrectangle{\pgfqpoint{4.259500in}{0.525000in}}{\pgfqpoint{1.293250in}{2.730000in}}%
\pgfusepath{clip}%
\pgfsetbuttcap%
\pgfsetmiterjoin%
\definecolor{currentfill}{rgb}{0.784314,0.784314,0.784314}%
\pgfsetfillcolor{currentfill}%
\pgfsetlinewidth{0.000000pt}%
\definecolor{currentstroke}{rgb}{0.000000,0.000000,0.000000}%
\pgfsetstrokecolor{currentstroke}%
\pgfsetstrokeopacity{0.000000}%
\pgfsetdash{}{0pt}%
\pgfpathmoveto{\pgfqpoint{4.259500in}{2.259056in}}%
\pgfpathlineto{\pgfqpoint{4.764368in}{2.259056in}}%
\pgfpathlineto{\pgfqpoint{4.764368in}{2.363537in}}%
\pgfpathlineto{\pgfqpoint{4.259500in}{2.363537in}}%
\pgfpathlineto{\pgfqpoint{4.259500in}{2.259056in}}%
\pgfpathclose%
\pgfusepath{fill}%
\end{pgfscope}%
\begin{pgfscope}%
\pgfpathrectangle{\pgfqpoint{4.259500in}{0.525000in}}{\pgfqpoint{1.293250in}{2.730000in}}%
\pgfusepath{clip}%
\pgfsetbuttcap%
\pgfsetmiterjoin%
\definecolor{currentfill}{rgb}{0.478431,0.478431,0.478431}%
\pgfsetfillcolor{currentfill}%
\pgfsetlinewidth{0.000000pt}%
\definecolor{currentstroke}{rgb}{0.000000,0.000000,0.000000}%
\pgfsetstrokecolor{currentstroke}%
\pgfsetstrokeopacity{0.000000}%
\pgfsetdash{}{0pt}%
\pgfpathmoveto{\pgfqpoint{4.259500in}{2.090537in}}%
\pgfpathlineto{\pgfqpoint{4.763286in}{2.090537in}}%
\pgfpathlineto{\pgfqpoint{4.763286in}{2.195019in}}%
\pgfpathlineto{\pgfqpoint{4.259500in}{2.195019in}}%
\pgfpathlineto{\pgfqpoint{4.259500in}{2.090537in}}%
\pgfpathclose%
\pgfusepath{fill}%
\end{pgfscope}%
\begin{pgfscope}%
\pgfpathrectangle{\pgfqpoint{4.259500in}{0.525000in}}{\pgfqpoint{1.293250in}{2.730000in}}%
\pgfusepath{clip}%
\pgfsetbuttcap%
\pgfsetmiterjoin%
\definecolor{currentfill}{rgb}{0.784314,0.784314,0.784314}%
\pgfsetfillcolor{currentfill}%
\pgfsetlinewidth{0.000000pt}%
\definecolor{currentstroke}{rgb}{0.000000,0.000000,0.000000}%
\pgfsetstrokecolor{currentstroke}%
\pgfsetstrokeopacity{0.000000}%
\pgfsetdash{}{0pt}%
\pgfpathmoveto{\pgfqpoint{4.259500in}{1.922019in}}%
\pgfpathlineto{\pgfqpoint{4.761124in}{1.922019in}}%
\pgfpathlineto{\pgfqpoint{4.761124in}{2.026500in}}%
\pgfpathlineto{\pgfqpoint{4.259500in}{2.026500in}}%
\pgfpathlineto{\pgfqpoint{4.259500in}{1.922019in}}%
\pgfpathclose%
\pgfusepath{fill}%
\end{pgfscope}%
\begin{pgfscope}%
\pgfpathrectangle{\pgfqpoint{4.259500in}{0.525000in}}{\pgfqpoint{1.293250in}{2.730000in}}%
\pgfusepath{clip}%
\pgfsetbuttcap%
\pgfsetmiterjoin%
\definecolor{currentfill}{rgb}{0.784314,0.784314,0.784314}%
\pgfsetfillcolor{currentfill}%
\pgfsetlinewidth{0.000000pt}%
\definecolor{currentstroke}{rgb}{0.000000,0.000000,0.000000}%
\pgfsetstrokecolor{currentstroke}%
\pgfsetstrokeopacity{0.000000}%
\pgfsetdash{}{0pt}%
\pgfpathmoveto{\pgfqpoint{4.259500in}{1.753500in}}%
\pgfpathlineto{\pgfqpoint{4.748151in}{1.753500in}}%
\pgfpathlineto{\pgfqpoint{4.748151in}{1.857981in}}%
\pgfpathlineto{\pgfqpoint{4.259500in}{1.857981in}}%
\pgfpathlineto{\pgfqpoint{4.259500in}{1.753500in}}%
\pgfpathclose%
\pgfusepath{fill}%
\end{pgfscope}%
\begin{pgfscope}%
\pgfpathrectangle{\pgfqpoint{4.259500in}{0.525000in}}{\pgfqpoint{1.293250in}{2.730000in}}%
\pgfusepath{clip}%
\pgfsetbuttcap%
\pgfsetmiterjoin%
\definecolor{currentfill}{rgb}{0.478431,0.478431,0.478431}%
\pgfsetfillcolor{currentfill}%
\pgfsetlinewidth{0.000000pt}%
\definecolor{currentstroke}{rgb}{0.000000,0.000000,0.000000}%
\pgfsetstrokecolor{currentstroke}%
\pgfsetstrokeopacity{0.000000}%
\pgfsetdash{}{0pt}%
\pgfpathmoveto{\pgfqpoint{4.259500in}{1.584981in}}%
\pgfpathlineto{\pgfqpoint{4.743827in}{1.584981in}}%
\pgfpathlineto{\pgfqpoint{4.743827in}{1.689463in}}%
\pgfpathlineto{\pgfqpoint{4.259500in}{1.689463in}}%
\pgfpathlineto{\pgfqpoint{4.259500in}{1.584981in}}%
\pgfpathclose%
\pgfusepath{fill}%
\end{pgfscope}%
\begin{pgfscope}%
\pgfpathrectangle{\pgfqpoint{4.259500in}{0.525000in}}{\pgfqpoint{1.293250in}{2.730000in}}%
\pgfusepath{clip}%
\pgfsetbuttcap%
\pgfsetmiterjoin%
\definecolor{currentfill}{rgb}{0.784314,0.784314,0.784314}%
\pgfsetfillcolor{currentfill}%
\pgfsetlinewidth{0.000000pt}%
\definecolor{currentstroke}{rgb}{0.000000,0.000000,0.000000}%
\pgfsetstrokecolor{currentstroke}%
\pgfsetstrokeopacity{0.000000}%
\pgfsetdash{}{0pt}%
\pgfpathmoveto{\pgfqpoint{4.259500in}{1.416463in}}%
\pgfpathlineto{\pgfqpoint{4.743286in}{1.416463in}}%
\pgfpathlineto{\pgfqpoint{4.743286in}{1.520944in}}%
\pgfpathlineto{\pgfqpoint{4.259500in}{1.520944in}}%
\pgfpathlineto{\pgfqpoint{4.259500in}{1.416463in}}%
\pgfpathclose%
\pgfusepath{fill}%
\end{pgfscope}%
\begin{pgfscope}%
\pgfpathrectangle{\pgfqpoint{4.259500in}{0.525000in}}{\pgfqpoint{1.293250in}{2.730000in}}%
\pgfusepath{clip}%
\pgfsetbuttcap%
\pgfsetmiterjoin%
\definecolor{currentfill}{rgb}{0.478431,0.478431,0.478431}%
\pgfsetfillcolor{currentfill}%
\pgfsetlinewidth{0.000000pt}%
\definecolor{currentstroke}{rgb}{0.000000,0.000000,0.000000}%
\pgfsetstrokecolor{currentstroke}%
\pgfsetstrokeopacity{0.000000}%
\pgfsetdash{}{0pt}%
\pgfpathmoveto{\pgfqpoint{4.259500in}{1.247944in}}%
\pgfpathlineto{\pgfqpoint{4.741665in}{1.247944in}}%
\pgfpathlineto{\pgfqpoint{4.741665in}{1.352426in}}%
\pgfpathlineto{\pgfqpoint{4.259500in}{1.352426in}}%
\pgfpathlineto{\pgfqpoint{4.259500in}{1.247944in}}%
\pgfpathclose%
\pgfusepath{fill}%
\end{pgfscope}%
\begin{pgfscope}%
\pgfpathrectangle{\pgfqpoint{4.259500in}{0.525000in}}{\pgfqpoint{1.293250in}{2.730000in}}%
\pgfusepath{clip}%
\pgfsetbuttcap%
\pgfsetmiterjoin%
\definecolor{currentfill}{rgb}{0.784314,0.784314,0.784314}%
\pgfsetfillcolor{currentfill}%
\pgfsetlinewidth{0.000000pt}%
\definecolor{currentstroke}{rgb}{0.000000,0.000000,0.000000}%
\pgfsetstrokecolor{currentstroke}%
\pgfsetstrokeopacity{0.000000}%
\pgfsetdash{}{0pt}%
\pgfpathmoveto{\pgfqpoint{4.259500in}{1.079426in}}%
\pgfpathlineto{\pgfqpoint{4.740043in}{1.079426in}}%
\pgfpathlineto{\pgfqpoint{4.740043in}{1.183907in}}%
\pgfpathlineto{\pgfqpoint{4.259500in}{1.183907in}}%
\pgfpathlineto{\pgfqpoint{4.259500in}{1.079426in}}%
\pgfpathclose%
\pgfusepath{fill}%
\end{pgfscope}%
\begin{pgfscope}%
\pgfpathrectangle{\pgfqpoint{4.259500in}{0.525000in}}{\pgfqpoint{1.293250in}{2.730000in}}%
\pgfusepath{clip}%
\pgfsetbuttcap%
\pgfsetmiterjoin%
\definecolor{currentfill}{rgb}{0.478431,0.478431,0.478431}%
\pgfsetfillcolor{currentfill}%
\pgfsetlinewidth{0.000000pt}%
\definecolor{currentstroke}{rgb}{0.000000,0.000000,0.000000}%
\pgfsetstrokecolor{currentstroke}%
\pgfsetstrokeopacity{0.000000}%
\pgfsetdash{}{0pt}%
\pgfpathmoveto{\pgfqpoint{4.259500in}{0.910907in}}%
\pgfpathlineto{\pgfqpoint{4.740043in}{0.910907in}}%
\pgfpathlineto{\pgfqpoint{4.740043in}{1.015389in}}%
\pgfpathlineto{\pgfqpoint{4.259500in}{1.015389in}}%
\pgfpathlineto{\pgfqpoint{4.259500in}{0.910907in}}%
\pgfpathclose%
\pgfusepath{fill}%
\end{pgfscope}%
\begin{pgfscope}%
\pgfpathrectangle{\pgfqpoint{4.259500in}{0.525000in}}{\pgfqpoint{1.293250in}{2.730000in}}%
\pgfusepath{clip}%
\pgfsetbuttcap%
\pgfsetmiterjoin%
\definecolor{currentfill}{rgb}{0.478431,0.478431,0.478431}%
\pgfsetfillcolor{currentfill}%
\pgfsetlinewidth{0.000000pt}%
\definecolor{currentstroke}{rgb}{0.000000,0.000000,0.000000}%
\pgfsetstrokecolor{currentstroke}%
\pgfsetstrokeopacity{0.000000}%
\pgfsetdash{}{0pt}%
\pgfpathmoveto{\pgfqpoint{4.259500in}{0.742389in}}%
\pgfpathlineto{\pgfqpoint{4.738962in}{0.742389in}}%
\pgfpathlineto{\pgfqpoint{4.738962in}{0.846870in}}%
\pgfpathlineto{\pgfqpoint{4.259500in}{0.846870in}}%
\pgfpathlineto{\pgfqpoint{4.259500in}{0.742389in}}%
\pgfpathclose%
\pgfusepath{fill}%
\end{pgfscope}%
\begin{pgfscope}%
\pgfpathrectangle{\pgfqpoint{4.259500in}{0.525000in}}{\pgfqpoint{1.293250in}{2.730000in}}%
\pgfusepath{clip}%
\pgfsetbuttcap%
\pgfsetmiterjoin%
\definecolor{currentfill}{rgb}{0.478431,0.305882,0.611765}%
\pgfsetfillcolor{currentfill}%
\pgfsetlinewidth{0.000000pt}%
\definecolor{currentstroke}{rgb}{0.000000,0.000000,0.000000}%
\pgfsetstrokecolor{currentstroke}%
\pgfsetstrokeopacity{0.000000}%
\pgfsetdash{}{0pt}%
\pgfpathmoveto{\pgfqpoint{4.259500in}{0.573870in}}%
\pgfpathlineto{\pgfqpoint{4.742205in}{0.573870in}}%
\pgfpathlineto{\pgfqpoint{4.742205in}{0.678352in}}%
\pgfpathlineto{\pgfqpoint{4.259500in}{0.678352in}}%
\pgfpathlineto{\pgfqpoint{4.259500in}{0.573870in}}%
\pgfpathclose%
\pgfusepath{fill}%
\end{pgfscope}%
\begin{pgfscope}%
\pgfsetrectcap%
\pgfsetmiterjoin%
\pgfsetlinewidth{0.803000pt}%
\definecolor{currentstroke}{rgb}{0.101961,0.101961,0.101961}%
\pgfsetstrokecolor{currentstroke}%
\pgfsetdash{}{0pt}%
\pgfpathmoveto{\pgfqpoint{4.259500in}{0.525000in}}%
\pgfpathlineto{\pgfqpoint{4.259500in}{3.255000in}}%
\pgfusepath{stroke}%
\end{pgfscope}%
\begin{pgfscope}%
\pgfsetrectcap%
\pgfsetmiterjoin%
\pgfsetlinewidth{0.803000pt}%
\definecolor{currentstroke}{rgb}{0.101961,0.101961,0.101961}%
\pgfsetstrokecolor{currentstroke}%
\pgfsetdash{}{0pt}%
\pgfpathmoveto{\pgfqpoint{4.259500in}{0.525000in}}%
\pgfpathlineto{\pgfqpoint{5.552750in}{0.525000in}}%
\pgfusepath{stroke}%
\end{pgfscope}%
\begin{pgfscope}%
\pgfpathrectangle{\pgfqpoint{4.259500in}{0.525000in}}{\pgfqpoint{1.293250in}{2.730000in}}%
\pgfusepath{clip}%
\pgfsetbuttcap%
\pgfsetroundjoin%
\definecolor{currentfill}{rgb}{0.101961,0.101961,0.101961}%
\pgfsetfillcolor{currentfill}%
\pgfsetlinewidth{0.301125pt}%
\definecolor{currentstroke}{rgb}{1.000000,1.000000,1.000000}%
\pgfsetstrokecolor{currentstroke}%
\pgfsetdash{}{0pt}%
\pgfsys@defobject{currentmarker}{\pgfqpoint{-0.013889in}{-0.013889in}}{\pgfqpoint{0.013889in}{0.013889in}}{%
\pgfpathmoveto{\pgfqpoint{0.000000in}{-0.013889in}}%
\pgfpathcurveto{\pgfqpoint{0.003683in}{-0.013889in}}{\pgfqpoint{0.007216in}{-0.012425in}}{\pgfqpoint{0.009821in}{-0.009821in}}%
\pgfpathcurveto{\pgfqpoint{0.012425in}{-0.007216in}}{\pgfqpoint{0.013889in}{-0.003683in}}{\pgfqpoint{0.013889in}{0.000000in}}%
\pgfpathcurveto{\pgfqpoint{0.013889in}{0.003683in}}{\pgfqpoint{0.012425in}{0.007216in}}{\pgfqpoint{0.009821in}{0.009821in}}%
\pgfpathcurveto{\pgfqpoint{0.007216in}{0.012425in}}{\pgfqpoint{0.003683in}{0.013889in}}{\pgfqpoint{0.000000in}{0.013889in}}%
\pgfpathcurveto{\pgfqpoint{-0.003683in}{0.013889in}}{\pgfqpoint{-0.007216in}{0.012425in}}{\pgfqpoint{-0.009821in}{0.009821in}}%
\pgfpathcurveto{\pgfqpoint{-0.012425in}{0.007216in}}{\pgfqpoint{-0.013889in}{0.003683in}}{\pgfqpoint{-0.013889in}{0.000000in}}%
\pgfpathcurveto{\pgfqpoint{-0.013889in}{-0.003683in}}{\pgfqpoint{-0.012425in}{-0.007216in}}{\pgfqpoint{-0.009821in}{-0.009821in}}%
\pgfpathcurveto{\pgfqpoint{-0.007216in}{-0.012425in}}{\pgfqpoint{-0.003683in}{-0.013889in}}{\pgfqpoint{0.000000in}{-0.013889in}}%
\pgfpathlineto{\pgfqpoint{0.000000in}{-0.013889in}}%
\pgfpathclose%
\pgfusepath{stroke,fill}%
}%
\begin{pgfscope}%
\pgfsys@transformshift{4.885761in}{3.153889in}%
\pgfsys@useobject{currentmarker}{}%
\end{pgfscope}%
\begin{pgfscope}%
\pgfsys@transformshift{4.833564in}{3.153889in}%
\pgfsys@useobject{currentmarker}{}%
\end{pgfscope}%
\begin{pgfscope}%
\pgfsys@transformshift{4.856884in}{3.153889in}%
\pgfsys@useobject{currentmarker}{}%
\end{pgfscope}%
\begin{pgfscope}%
\pgfsys@transformshift{4.852611in}{3.153889in}%
\pgfsys@useobject{currentmarker}{}%
\end{pgfscope}%
\end{pgfscope}%
\begin{pgfscope}%
\definecolor{textcolor}{rgb}{0.101961,0.101961,0.101961}%
\pgfsetstrokecolor{textcolor}%
\pgfsetfillcolor{textcolor}%
\pgftext[x=5.335484in,y=3.153889in,left,]{\color{textcolor}\sffamily\fontsize{6.000000}{7.200000}\bfseries\selectfont 0.58}%
\end{pgfscope}%
\begin{pgfscope}%
\pgfpathrectangle{\pgfqpoint{4.259500in}{0.525000in}}{\pgfqpoint{1.293250in}{2.730000in}}%
\pgfusepath{clip}%
\pgfsetbuttcap%
\pgfsetroundjoin%
\definecolor{currentfill}{rgb}{0.101961,0.101961,0.101961}%
\pgfsetfillcolor{currentfill}%
\pgfsetlinewidth{0.301125pt}%
\definecolor{currentstroke}{rgb}{1.000000,1.000000,1.000000}%
\pgfsetstrokecolor{currentstroke}%
\pgfsetdash{}{0pt}%
\pgfsys@defobject{currentmarker}{\pgfqpoint{-0.013889in}{-0.013889in}}{\pgfqpoint{0.013889in}{0.013889in}}{%
\pgfpathmoveto{\pgfqpoint{0.000000in}{-0.013889in}}%
\pgfpathcurveto{\pgfqpoint{0.003683in}{-0.013889in}}{\pgfqpoint{0.007216in}{-0.012425in}}{\pgfqpoint{0.009821in}{-0.009821in}}%
\pgfpathcurveto{\pgfqpoint{0.012425in}{-0.007216in}}{\pgfqpoint{0.013889in}{-0.003683in}}{\pgfqpoint{0.013889in}{0.000000in}}%
\pgfpathcurveto{\pgfqpoint{0.013889in}{0.003683in}}{\pgfqpoint{0.012425in}{0.007216in}}{\pgfqpoint{0.009821in}{0.009821in}}%
\pgfpathcurveto{\pgfqpoint{0.007216in}{0.012425in}}{\pgfqpoint{0.003683in}{0.013889in}}{\pgfqpoint{0.000000in}{0.013889in}}%
\pgfpathcurveto{\pgfqpoint{-0.003683in}{0.013889in}}{\pgfqpoint{-0.007216in}{0.012425in}}{\pgfqpoint{-0.009821in}{0.009821in}}%
\pgfpathcurveto{\pgfqpoint{-0.012425in}{0.007216in}}{\pgfqpoint{-0.013889in}{0.003683in}}{\pgfqpoint{-0.013889in}{0.000000in}}%
\pgfpathcurveto{\pgfqpoint{-0.013889in}{-0.003683in}}{\pgfqpoint{-0.012425in}{-0.007216in}}{\pgfqpoint{-0.009821in}{-0.009821in}}%
\pgfpathcurveto{\pgfqpoint{-0.007216in}{-0.012425in}}{\pgfqpoint{-0.003683in}{-0.013889in}}{\pgfqpoint{0.000000in}{-0.013889in}}%
\pgfpathlineto{\pgfqpoint{0.000000in}{-0.013889in}}%
\pgfpathclose%
\pgfusepath{stroke,fill}%
}%
\begin{pgfscope}%
\pgfsys@transformshift{4.750353in}{2.985370in}%
\pgfsys@useobject{currentmarker}{}%
\end{pgfscope}%
\begin{pgfscope}%
\pgfsys@transformshift{4.799291in}{2.985370in}%
\pgfsys@useobject{currentmarker}{}%
\end{pgfscope}%
\begin{pgfscope}%
\pgfsys@transformshift{4.765978in}{2.985370in}%
\pgfsys@useobject{currentmarker}{}%
\end{pgfscope}%
\begin{pgfscope}%
\pgfsys@transformshift{4.808016in}{2.985370in}%
\pgfsys@useobject{currentmarker}{}%
\end{pgfscope}%
\end{pgfscope}%
\begin{pgfscope}%
\definecolor{textcolor}{rgb}{0.101961,0.101961,0.101961}%
\pgfsetstrokecolor{textcolor}%
\pgfsetfillcolor{textcolor}%
\pgftext[x=5.335484in,y=2.985370in,left,]{\color{textcolor}\sffamily\fontsize{6.000000}{7.200000}\selectfont 0.50}%
\end{pgfscope}%
\begin{pgfscope}%
\pgfpathrectangle{\pgfqpoint{4.259500in}{0.525000in}}{\pgfqpoint{1.293250in}{2.730000in}}%
\pgfusepath{clip}%
\pgfsetbuttcap%
\pgfsetroundjoin%
\definecolor{currentfill}{rgb}{0.101961,0.101961,0.101961}%
\pgfsetfillcolor{currentfill}%
\pgfsetlinewidth{0.301125pt}%
\definecolor{currentstroke}{rgb}{1.000000,1.000000,1.000000}%
\pgfsetstrokecolor{currentstroke}%
\pgfsetdash{}{0pt}%
\pgfsys@defobject{currentmarker}{\pgfqpoint{-0.013889in}{-0.013889in}}{\pgfqpoint{0.013889in}{0.013889in}}{%
\pgfpathmoveto{\pgfqpoint{0.000000in}{-0.013889in}}%
\pgfpathcurveto{\pgfqpoint{0.003683in}{-0.013889in}}{\pgfqpoint{0.007216in}{-0.012425in}}{\pgfqpoint{0.009821in}{-0.009821in}}%
\pgfpathcurveto{\pgfqpoint{0.012425in}{-0.007216in}}{\pgfqpoint{0.013889in}{-0.003683in}}{\pgfqpoint{0.013889in}{0.000000in}}%
\pgfpathcurveto{\pgfqpoint{0.013889in}{0.003683in}}{\pgfqpoint{0.012425in}{0.007216in}}{\pgfqpoint{0.009821in}{0.009821in}}%
\pgfpathcurveto{\pgfqpoint{0.007216in}{0.012425in}}{\pgfqpoint{0.003683in}{0.013889in}}{\pgfqpoint{0.000000in}{0.013889in}}%
\pgfpathcurveto{\pgfqpoint{-0.003683in}{0.013889in}}{\pgfqpoint{-0.007216in}{0.012425in}}{\pgfqpoint{-0.009821in}{0.009821in}}%
\pgfpathcurveto{\pgfqpoint{-0.012425in}{0.007216in}}{\pgfqpoint{-0.013889in}{0.003683in}}{\pgfqpoint{-0.013889in}{0.000000in}}%
\pgfpathcurveto{\pgfqpoint{-0.013889in}{-0.003683in}}{\pgfqpoint{-0.012425in}{-0.007216in}}{\pgfqpoint{-0.009821in}{-0.009821in}}%
\pgfpathcurveto{\pgfqpoint{-0.007216in}{-0.012425in}}{\pgfqpoint{-0.003683in}{-0.013889in}}{\pgfqpoint{0.000000in}{-0.013889in}}%
\pgfpathlineto{\pgfqpoint{0.000000in}{-0.013889in}}%
\pgfpathclose%
\pgfusepath{stroke,fill}%
}%
\begin{pgfscope}%
\pgfsys@transformshift{4.815941in}{2.816852in}%
\pgfsys@useobject{currentmarker}{}%
\end{pgfscope}%
\begin{pgfscope}%
\pgfsys@transformshift{4.767161in}{2.816852in}%
\pgfsys@useobject{currentmarker}{}%
\end{pgfscope}%
\begin{pgfscope}%
\pgfsys@transformshift{4.740005in}{2.816852in}%
\pgfsys@useobject{currentmarker}{}%
\end{pgfscope}%
\begin{pgfscope}%
\pgfsys@transformshift{4.799097in}{2.816852in}%
\pgfsys@useobject{currentmarker}{}%
\end{pgfscope}%
\end{pgfscope}%
\begin{pgfscope}%
\definecolor{textcolor}{rgb}{0.101961,0.101961,0.101961}%
\pgfsetstrokecolor{textcolor}%
\pgfsetfillcolor{textcolor}%
\pgftext[x=5.335484in,y=2.816852in,left,]{\color{textcolor}\sffamily\fontsize{6.000000}{7.200000}\selectfont 0.50}%
\end{pgfscope}%
\begin{pgfscope}%
\pgfpathrectangle{\pgfqpoint{4.259500in}{0.525000in}}{\pgfqpoint{1.293250in}{2.730000in}}%
\pgfusepath{clip}%
\pgfsetbuttcap%
\pgfsetroundjoin%
\definecolor{currentfill}{rgb}{0.101961,0.101961,0.101961}%
\pgfsetfillcolor{currentfill}%
\pgfsetlinewidth{0.301125pt}%
\definecolor{currentstroke}{rgb}{1.000000,1.000000,1.000000}%
\pgfsetstrokecolor{currentstroke}%
\pgfsetdash{}{0pt}%
\pgfsys@defobject{currentmarker}{\pgfqpoint{-0.013889in}{-0.013889in}}{\pgfqpoint{0.013889in}{0.013889in}}{%
\pgfpathmoveto{\pgfqpoint{0.000000in}{-0.013889in}}%
\pgfpathcurveto{\pgfqpoint{0.003683in}{-0.013889in}}{\pgfqpoint{0.007216in}{-0.012425in}}{\pgfqpoint{0.009821in}{-0.009821in}}%
\pgfpathcurveto{\pgfqpoint{0.012425in}{-0.007216in}}{\pgfqpoint{0.013889in}{-0.003683in}}{\pgfqpoint{0.013889in}{0.000000in}}%
\pgfpathcurveto{\pgfqpoint{0.013889in}{0.003683in}}{\pgfqpoint{0.012425in}{0.007216in}}{\pgfqpoint{0.009821in}{0.009821in}}%
\pgfpathcurveto{\pgfqpoint{0.007216in}{0.012425in}}{\pgfqpoint{0.003683in}{0.013889in}}{\pgfqpoint{0.000000in}{0.013889in}}%
\pgfpathcurveto{\pgfqpoint{-0.003683in}{0.013889in}}{\pgfqpoint{-0.007216in}{0.012425in}}{\pgfqpoint{-0.009821in}{0.009821in}}%
\pgfpathcurveto{\pgfqpoint{-0.012425in}{0.007216in}}{\pgfqpoint{-0.013889in}{0.003683in}}{\pgfqpoint{-0.013889in}{0.000000in}}%
\pgfpathcurveto{\pgfqpoint{-0.013889in}{-0.003683in}}{\pgfqpoint{-0.012425in}{-0.007216in}}{\pgfqpoint{-0.009821in}{-0.009821in}}%
\pgfpathcurveto{\pgfqpoint{-0.007216in}{-0.012425in}}{\pgfqpoint{-0.003683in}{-0.013889in}}{\pgfqpoint{0.000000in}{-0.013889in}}%
\pgfpathlineto{\pgfqpoint{0.000000in}{-0.013889in}}%
\pgfpathclose%
\pgfusepath{stroke,fill}%
}%
\begin{pgfscope}%
\pgfsys@transformshift{4.760932in}{2.648333in}%
\pgfsys@useobject{currentmarker}{}%
\end{pgfscope}%
\begin{pgfscope}%
\pgfsys@transformshift{4.799291in}{2.648333in}%
\pgfsys@useobject{currentmarker}{}%
\end{pgfscope}%
\begin{pgfscope}%
\pgfsys@transformshift{4.757320in}{2.648333in}%
\pgfsys@useobject{currentmarker}{}%
\end{pgfscope}%
\begin{pgfscope}%
\pgfsys@transformshift{4.794638in}{2.648333in}%
\pgfsys@useobject{currentmarker}{}%
\end{pgfscope}%
\end{pgfscope}%
\begin{pgfscope}%
\definecolor{textcolor}{rgb}{0.101961,0.101961,0.101961}%
\pgfsetstrokecolor{textcolor}%
\pgfsetfillcolor{textcolor}%
\pgftext[x=5.335484in,y=2.648333in,left,]{\color{textcolor}\sffamily\fontsize{6.000000}{7.200000}\selectfont 0.50}%
\end{pgfscope}%
\begin{pgfscope}%
\pgfpathrectangle{\pgfqpoint{4.259500in}{0.525000in}}{\pgfqpoint{1.293250in}{2.730000in}}%
\pgfusepath{clip}%
\pgfsetbuttcap%
\pgfsetroundjoin%
\definecolor{currentfill}{rgb}{0.101961,0.101961,0.101961}%
\pgfsetfillcolor{currentfill}%
\pgfsetlinewidth{0.301125pt}%
\definecolor{currentstroke}{rgb}{1.000000,1.000000,1.000000}%
\pgfsetstrokecolor{currentstroke}%
\pgfsetdash{}{0pt}%
\pgfsys@defobject{currentmarker}{\pgfqpoint{-0.013889in}{-0.013889in}}{\pgfqpoint{0.013889in}{0.013889in}}{%
\pgfpathmoveto{\pgfqpoint{0.000000in}{-0.013889in}}%
\pgfpathcurveto{\pgfqpoint{0.003683in}{-0.013889in}}{\pgfqpoint{0.007216in}{-0.012425in}}{\pgfqpoint{0.009821in}{-0.009821in}}%
\pgfpathcurveto{\pgfqpoint{0.012425in}{-0.007216in}}{\pgfqpoint{0.013889in}{-0.003683in}}{\pgfqpoint{0.013889in}{0.000000in}}%
\pgfpathcurveto{\pgfqpoint{0.013889in}{0.003683in}}{\pgfqpoint{0.012425in}{0.007216in}}{\pgfqpoint{0.009821in}{0.009821in}}%
\pgfpathcurveto{\pgfqpoint{0.007216in}{0.012425in}}{\pgfqpoint{0.003683in}{0.013889in}}{\pgfqpoint{0.000000in}{0.013889in}}%
\pgfpathcurveto{\pgfqpoint{-0.003683in}{0.013889in}}{\pgfqpoint{-0.007216in}{0.012425in}}{\pgfqpoint{-0.009821in}{0.009821in}}%
\pgfpathcurveto{\pgfqpoint{-0.012425in}{0.007216in}}{\pgfqpoint{-0.013889in}{0.003683in}}{\pgfqpoint{-0.013889in}{0.000000in}}%
\pgfpathcurveto{\pgfqpoint{-0.013889in}{-0.003683in}}{\pgfqpoint{-0.012425in}{-0.007216in}}{\pgfqpoint{-0.009821in}{-0.009821in}}%
\pgfpathcurveto{\pgfqpoint{-0.007216in}{-0.012425in}}{\pgfqpoint{-0.003683in}{-0.013889in}}{\pgfqpoint{0.000000in}{-0.013889in}}%
\pgfpathlineto{\pgfqpoint{0.000000in}{-0.013889in}}%
\pgfpathclose%
\pgfusepath{stroke,fill}%
}%
\begin{pgfscope}%
\pgfsys@transformshift{4.792668in}{2.479815in}%
\pgfsys@useobject{currentmarker}{}%
\end{pgfscope}%
\begin{pgfscope}%
\pgfsys@transformshift{4.780013in}{2.479815in}%
\pgfsys@useobject{currentmarker}{}%
\end{pgfscope}%
\begin{pgfscope}%
\pgfsys@transformshift{4.720525in}{2.479815in}%
\pgfsys@useobject{currentmarker}{}%
\end{pgfscope}%
\begin{pgfscope}%
\pgfsys@transformshift{4.767881in}{2.479815in}%
\pgfsys@useobject{currentmarker}{}%
\end{pgfscope}%
\end{pgfscope}%
\begin{pgfscope}%
\definecolor{textcolor}{rgb}{0.101961,0.101961,0.101961}%
\pgfsetstrokecolor{textcolor}%
\pgfsetfillcolor{textcolor}%
\pgftext[x=5.335484in,y=2.479815in,left,]{\color{textcolor}\sffamily\fontsize{6.000000}{7.200000}\selectfont 0.49}%
\end{pgfscope}%
\begin{pgfscope}%
\pgfpathrectangle{\pgfqpoint{4.259500in}{0.525000in}}{\pgfqpoint{1.293250in}{2.730000in}}%
\pgfusepath{clip}%
\pgfsetbuttcap%
\pgfsetroundjoin%
\definecolor{currentfill}{rgb}{0.101961,0.101961,0.101961}%
\pgfsetfillcolor{currentfill}%
\pgfsetlinewidth{0.301125pt}%
\definecolor{currentstroke}{rgb}{1.000000,1.000000,1.000000}%
\pgfsetstrokecolor{currentstroke}%
\pgfsetdash{}{0pt}%
\pgfsys@defobject{currentmarker}{\pgfqpoint{-0.013889in}{-0.013889in}}{\pgfqpoint{0.013889in}{0.013889in}}{%
\pgfpathmoveto{\pgfqpoint{0.000000in}{-0.013889in}}%
\pgfpathcurveto{\pgfqpoint{0.003683in}{-0.013889in}}{\pgfqpoint{0.007216in}{-0.012425in}}{\pgfqpoint{0.009821in}{-0.009821in}}%
\pgfpathcurveto{\pgfqpoint{0.012425in}{-0.007216in}}{\pgfqpoint{0.013889in}{-0.003683in}}{\pgfqpoint{0.013889in}{0.000000in}}%
\pgfpathcurveto{\pgfqpoint{0.013889in}{0.003683in}}{\pgfqpoint{0.012425in}{0.007216in}}{\pgfqpoint{0.009821in}{0.009821in}}%
\pgfpathcurveto{\pgfqpoint{0.007216in}{0.012425in}}{\pgfqpoint{0.003683in}{0.013889in}}{\pgfqpoint{0.000000in}{0.013889in}}%
\pgfpathcurveto{\pgfqpoint{-0.003683in}{0.013889in}}{\pgfqpoint{-0.007216in}{0.012425in}}{\pgfqpoint{-0.009821in}{0.009821in}}%
\pgfpathcurveto{\pgfqpoint{-0.012425in}{0.007216in}}{\pgfqpoint{-0.013889in}{0.003683in}}{\pgfqpoint{-0.013889in}{0.000000in}}%
\pgfpathcurveto{\pgfqpoint{-0.013889in}{-0.003683in}}{\pgfqpoint{-0.012425in}{-0.007216in}}{\pgfqpoint{-0.009821in}{-0.009821in}}%
\pgfpathcurveto{\pgfqpoint{-0.007216in}{-0.012425in}}{\pgfqpoint{-0.003683in}{-0.013889in}}{\pgfqpoint{0.000000in}{-0.013889in}}%
\pgfpathlineto{\pgfqpoint{0.000000in}{-0.013889in}}%
\pgfpathclose%
\pgfusepath{stroke,fill}%
}%
\begin{pgfscope}%
\pgfsys@transformshift{4.769395in}{2.311296in}%
\pgfsys@useobject{currentmarker}{}%
\end{pgfscope}%
\begin{pgfscope}%
\pgfsys@transformshift{4.752167in}{2.311296in}%
\pgfsys@useobject{currentmarker}{}%
\end{pgfscope}%
\begin{pgfscope}%
\pgfsys@transformshift{4.735676in}{2.311296in}%
\pgfsys@useobject{currentmarker}{}%
\end{pgfscope}%
\begin{pgfscope}%
\pgfsys@transformshift{4.801327in}{2.311296in}%
\pgfsys@useobject{currentmarker}{}%
\end{pgfscope}%
\end{pgfscope}%
\begin{pgfscope}%
\definecolor{textcolor}{rgb}{0.101961,0.101961,0.101961}%
\pgfsetstrokecolor{textcolor}%
\pgfsetfillcolor{textcolor}%
\pgftext[x=5.335484in,y=2.311296in,left,]{\color{textcolor}\sffamily\fontsize{6.000000}{7.200000}\selectfont 0.49}%
\end{pgfscope}%
\begin{pgfscope}%
\pgfpathrectangle{\pgfqpoint{4.259500in}{0.525000in}}{\pgfqpoint{1.293250in}{2.730000in}}%
\pgfusepath{clip}%
\pgfsetbuttcap%
\pgfsetroundjoin%
\definecolor{currentfill}{rgb}{0.101961,0.101961,0.101961}%
\pgfsetfillcolor{currentfill}%
\pgfsetlinewidth{0.301125pt}%
\definecolor{currentstroke}{rgb}{1.000000,1.000000,1.000000}%
\pgfsetstrokecolor{currentstroke}%
\pgfsetdash{}{0pt}%
\pgfsys@defobject{currentmarker}{\pgfqpoint{-0.013889in}{-0.013889in}}{\pgfqpoint{0.013889in}{0.013889in}}{%
\pgfpathmoveto{\pgfqpoint{0.000000in}{-0.013889in}}%
\pgfpathcurveto{\pgfqpoint{0.003683in}{-0.013889in}}{\pgfqpoint{0.007216in}{-0.012425in}}{\pgfqpoint{0.009821in}{-0.009821in}}%
\pgfpathcurveto{\pgfqpoint{0.012425in}{-0.007216in}}{\pgfqpoint{0.013889in}{-0.003683in}}{\pgfqpoint{0.013889in}{0.000000in}}%
\pgfpathcurveto{\pgfqpoint{0.013889in}{0.003683in}}{\pgfqpoint{0.012425in}{0.007216in}}{\pgfqpoint{0.009821in}{0.009821in}}%
\pgfpathcurveto{\pgfqpoint{0.007216in}{0.012425in}}{\pgfqpoint{0.003683in}{0.013889in}}{\pgfqpoint{0.000000in}{0.013889in}}%
\pgfpathcurveto{\pgfqpoint{-0.003683in}{0.013889in}}{\pgfqpoint{-0.007216in}{0.012425in}}{\pgfqpoint{-0.009821in}{0.009821in}}%
\pgfpathcurveto{\pgfqpoint{-0.012425in}{0.007216in}}{\pgfqpoint{-0.013889in}{0.003683in}}{\pgfqpoint{-0.013889in}{0.000000in}}%
\pgfpathcurveto{\pgfqpoint{-0.013889in}{-0.003683in}}{\pgfqpoint{-0.012425in}{-0.007216in}}{\pgfqpoint{-0.009821in}{-0.009821in}}%
\pgfpathcurveto{\pgfqpoint{-0.007216in}{-0.012425in}}{\pgfqpoint{-0.003683in}{-0.013889in}}{\pgfqpoint{0.000000in}{-0.013889in}}%
\pgfpathlineto{\pgfqpoint{0.000000in}{-0.013889in}}%
\pgfpathclose%
\pgfusepath{stroke,fill}%
}%
\begin{pgfscope}%
\pgfsys@transformshift{4.769395in}{2.142778in}%
\pgfsys@useobject{currentmarker}{}%
\end{pgfscope}%
\begin{pgfscope}%
\pgfsys@transformshift{4.754309in}{2.142778in}%
\pgfsys@useobject{currentmarker}{}%
\end{pgfscope}%
\begin{pgfscope}%
\pgfsys@transformshift{4.720525in}{2.142778in}%
\pgfsys@useobject{currentmarker}{}%
\end{pgfscope}%
\begin{pgfscope}%
\pgfsys@transformshift{4.810246in}{2.142778in}%
\pgfsys@useobject{currentmarker}{}%
\end{pgfscope}%
\end{pgfscope}%
\begin{pgfscope}%
\definecolor{textcolor}{rgb}{0.101961,0.101961,0.101961}%
\pgfsetstrokecolor{textcolor}%
\pgfsetfillcolor{textcolor}%
\pgftext[x=5.335484in,y=2.142778in,left,]{\color{textcolor}\sffamily\fontsize{6.000000}{7.200000}\selectfont 0.49}%
\end{pgfscope}%
\begin{pgfscope}%
\pgfpathrectangle{\pgfqpoint{4.259500in}{0.525000in}}{\pgfqpoint{1.293250in}{2.730000in}}%
\pgfusepath{clip}%
\pgfsetbuttcap%
\pgfsetroundjoin%
\definecolor{currentfill}{rgb}{0.101961,0.101961,0.101961}%
\pgfsetfillcolor{currentfill}%
\pgfsetlinewidth{0.301125pt}%
\definecolor{currentstroke}{rgb}{1.000000,1.000000,1.000000}%
\pgfsetstrokecolor{currentstroke}%
\pgfsetdash{}{0pt}%
\pgfsys@defobject{currentmarker}{\pgfqpoint{-0.013889in}{-0.013889in}}{\pgfqpoint{0.013889in}{0.013889in}}{%
\pgfpathmoveto{\pgfqpoint{0.000000in}{-0.013889in}}%
\pgfpathcurveto{\pgfqpoint{0.003683in}{-0.013889in}}{\pgfqpoint{0.007216in}{-0.012425in}}{\pgfqpoint{0.009821in}{-0.009821in}}%
\pgfpathcurveto{\pgfqpoint{0.012425in}{-0.007216in}}{\pgfqpoint{0.013889in}{-0.003683in}}{\pgfqpoint{0.013889in}{0.000000in}}%
\pgfpathcurveto{\pgfqpoint{0.013889in}{0.003683in}}{\pgfqpoint{0.012425in}{0.007216in}}{\pgfqpoint{0.009821in}{0.009821in}}%
\pgfpathcurveto{\pgfqpoint{0.007216in}{0.012425in}}{\pgfqpoint{0.003683in}{0.013889in}}{\pgfqpoint{0.000000in}{0.013889in}}%
\pgfpathcurveto{\pgfqpoint{-0.003683in}{0.013889in}}{\pgfqpoint{-0.007216in}{0.012425in}}{\pgfqpoint{-0.009821in}{0.009821in}}%
\pgfpathcurveto{\pgfqpoint{-0.012425in}{0.007216in}}{\pgfqpoint{-0.013889in}{0.003683in}}{\pgfqpoint{-0.013889in}{0.000000in}}%
\pgfpathcurveto{\pgfqpoint{-0.013889in}{-0.003683in}}{\pgfqpoint{-0.012425in}{-0.007216in}}{\pgfqpoint{-0.009821in}{-0.009821in}}%
\pgfpathcurveto{\pgfqpoint{-0.007216in}{-0.012425in}}{\pgfqpoint{-0.003683in}{-0.013889in}}{\pgfqpoint{0.000000in}{-0.013889in}}%
\pgfpathlineto{\pgfqpoint{0.000000in}{-0.013889in}}%
\pgfpathclose%
\pgfusepath{stroke,fill}%
}%
\begin{pgfscope}%
\pgfsys@transformshift{4.769395in}{1.974259in}%
\pgfsys@useobject{currentmarker}{}%
\end{pgfscope}%
\begin{pgfscope}%
\pgfsys@transformshift{4.750025in}{1.974259in}%
\pgfsys@useobject{currentmarker}{}%
\end{pgfscope}%
\begin{pgfscope}%
\pgfsys@transformshift{4.720525in}{1.974259in}%
\pgfsys@useobject{currentmarker}{}%
\end{pgfscope}%
\begin{pgfscope}%
\pgfsys@transformshift{4.805787in}{1.974259in}%
\pgfsys@useobject{currentmarker}{}%
\end{pgfscope}%
\end{pgfscope}%
\begin{pgfscope}%
\definecolor{textcolor}{rgb}{0.101961,0.101961,0.101961}%
\pgfsetstrokecolor{textcolor}%
\pgfsetfillcolor{textcolor}%
\pgftext[x=5.335484in,y=1.974259in,left,]{\color{textcolor}\sffamily\fontsize{6.000000}{7.200000}\selectfont 0.48}%
\end{pgfscope}%
\begin{pgfscope}%
\pgfpathrectangle{\pgfqpoint{4.259500in}{0.525000in}}{\pgfqpoint{1.293250in}{2.730000in}}%
\pgfusepath{clip}%
\pgfsetbuttcap%
\pgfsetroundjoin%
\definecolor{currentfill}{rgb}{0.101961,0.101961,0.101961}%
\pgfsetfillcolor{currentfill}%
\pgfsetlinewidth{0.301125pt}%
\definecolor{currentstroke}{rgb}{1.000000,1.000000,1.000000}%
\pgfsetstrokecolor{currentstroke}%
\pgfsetdash{}{0pt}%
\pgfsys@defobject{currentmarker}{\pgfqpoint{-0.013889in}{-0.013889in}}{\pgfqpoint{0.013889in}{0.013889in}}{%
\pgfpathmoveto{\pgfqpoint{0.000000in}{-0.013889in}}%
\pgfpathcurveto{\pgfqpoint{0.003683in}{-0.013889in}}{\pgfqpoint{0.007216in}{-0.012425in}}{\pgfqpoint{0.009821in}{-0.009821in}}%
\pgfpathcurveto{\pgfqpoint{0.012425in}{-0.007216in}}{\pgfqpoint{0.013889in}{-0.003683in}}{\pgfqpoint{0.013889in}{0.000000in}}%
\pgfpathcurveto{\pgfqpoint{0.013889in}{0.003683in}}{\pgfqpoint{0.012425in}{0.007216in}}{\pgfqpoint{0.009821in}{0.009821in}}%
\pgfpathcurveto{\pgfqpoint{0.007216in}{0.012425in}}{\pgfqpoint{0.003683in}{0.013889in}}{\pgfqpoint{0.000000in}{0.013889in}}%
\pgfpathcurveto{\pgfqpoint{-0.003683in}{0.013889in}}{\pgfqpoint{-0.007216in}{0.012425in}}{\pgfqpoint{-0.009821in}{0.009821in}}%
\pgfpathcurveto{\pgfqpoint{-0.012425in}{0.007216in}}{\pgfqpoint{-0.013889in}{0.003683in}}{\pgfqpoint{-0.013889in}{0.000000in}}%
\pgfpathcurveto{\pgfqpoint{-0.013889in}{-0.003683in}}{\pgfqpoint{-0.012425in}{-0.007216in}}{\pgfqpoint{-0.009821in}{-0.009821in}}%
\pgfpathcurveto{\pgfqpoint{-0.007216in}{-0.012425in}}{\pgfqpoint{-0.003683in}{-0.013889in}}{\pgfqpoint{0.000000in}{-0.013889in}}%
\pgfpathlineto{\pgfqpoint{0.000000in}{-0.013889in}}%
\pgfpathclose%
\pgfusepath{stroke,fill}%
}%
\begin{pgfscope}%
\pgfsys@transformshift{4.771511in}{1.805741in}%
\pgfsys@useobject{currentmarker}{}%
\end{pgfscope}%
\begin{pgfscope}%
\pgfsys@transformshift{4.752167in}{1.805741in}%
\pgfsys@useobject{currentmarker}{}%
\end{pgfscope}%
\begin{pgfscope}%
\pgfsys@transformshift{4.705374in}{1.805741in}%
\pgfsys@useobject{currentmarker}{}%
\end{pgfscope}%
\begin{pgfscope}%
\pgfsys@transformshift{4.763422in}{1.805741in}%
\pgfsys@useobject{currentmarker}{}%
\end{pgfscope}%
\end{pgfscope}%
\begin{pgfscope}%
\definecolor{textcolor}{rgb}{0.101961,0.101961,0.101961}%
\pgfsetstrokecolor{textcolor}%
\pgfsetfillcolor{textcolor}%
\pgftext[x=5.335484in,y=1.805741in,left,]{\color{textcolor}\sffamily\fontsize{6.000000}{7.200000}\selectfont 0.47}%
\end{pgfscope}%
\begin{pgfscope}%
\pgfpathrectangle{\pgfqpoint{4.259500in}{0.525000in}}{\pgfqpoint{1.293250in}{2.730000in}}%
\pgfusepath{clip}%
\pgfsetbuttcap%
\pgfsetroundjoin%
\definecolor{currentfill}{rgb}{0.101961,0.101961,0.101961}%
\pgfsetfillcolor{currentfill}%
\pgfsetlinewidth{0.301125pt}%
\definecolor{currentstroke}{rgb}{1.000000,1.000000,1.000000}%
\pgfsetstrokecolor{currentstroke}%
\pgfsetdash{}{0pt}%
\pgfsys@defobject{currentmarker}{\pgfqpoint{-0.013889in}{-0.013889in}}{\pgfqpoint{0.013889in}{0.013889in}}{%
\pgfpathmoveto{\pgfqpoint{0.000000in}{-0.013889in}}%
\pgfpathcurveto{\pgfqpoint{0.003683in}{-0.013889in}}{\pgfqpoint{0.007216in}{-0.012425in}}{\pgfqpoint{0.009821in}{-0.009821in}}%
\pgfpathcurveto{\pgfqpoint{0.012425in}{-0.007216in}}{\pgfqpoint{0.013889in}{-0.003683in}}{\pgfqpoint{0.013889in}{0.000000in}}%
\pgfpathcurveto{\pgfqpoint{0.013889in}{0.003683in}}{\pgfqpoint{0.012425in}{0.007216in}}{\pgfqpoint{0.009821in}{0.009821in}}%
\pgfpathcurveto{\pgfqpoint{0.007216in}{0.012425in}}{\pgfqpoint{0.003683in}{0.013889in}}{\pgfqpoint{0.000000in}{0.013889in}}%
\pgfpathcurveto{\pgfqpoint{-0.003683in}{0.013889in}}{\pgfqpoint{-0.007216in}{0.012425in}}{\pgfqpoint{-0.009821in}{0.009821in}}%
\pgfpathcurveto{\pgfqpoint{-0.012425in}{0.007216in}}{\pgfqpoint{-0.013889in}{0.003683in}}{\pgfqpoint{-0.013889in}{0.000000in}}%
\pgfpathcurveto{\pgfqpoint{-0.013889in}{-0.003683in}}{\pgfqpoint{-0.012425in}{-0.007216in}}{\pgfqpoint{-0.009821in}{-0.009821in}}%
\pgfpathcurveto{\pgfqpoint{-0.007216in}{-0.012425in}}{\pgfqpoint{-0.003683in}{-0.013889in}}{\pgfqpoint{0.000000in}{-0.013889in}}%
\pgfpathlineto{\pgfqpoint{0.000000in}{-0.013889in}}%
\pgfpathclose%
\pgfusepath{stroke,fill}%
}%
\begin{pgfscope}%
\pgfsys@transformshift{4.769395in}{1.637222in}%
\pgfsys@useobject{currentmarker}{}%
\end{pgfscope}%
\begin{pgfscope}%
\pgfsys@transformshift{4.752167in}{1.637222in}%
\pgfsys@useobject{currentmarker}{}%
\end{pgfscope}%
\begin{pgfscope}%
\pgfsys@transformshift{4.698880in}{1.637222in}%
\pgfsys@useobject{currentmarker}{}%
\end{pgfscope}%
\begin{pgfscope}%
\pgfsys@transformshift{4.754503in}{1.637222in}%
\pgfsys@useobject{currentmarker}{}%
\end{pgfscope}%
\end{pgfscope}%
\begin{pgfscope}%
\definecolor{textcolor}{rgb}{0.101961,0.101961,0.101961}%
\pgfsetstrokecolor{textcolor}%
\pgfsetfillcolor{textcolor}%
\pgftext[x=5.335484in,y=1.637222in,left,]{\color{textcolor}\sffamily\fontsize{6.000000}{7.200000}\selectfont 0.47}%
\end{pgfscope}%
\begin{pgfscope}%
\pgfpathrectangle{\pgfqpoint{4.259500in}{0.525000in}}{\pgfqpoint{1.293250in}{2.730000in}}%
\pgfusepath{clip}%
\pgfsetbuttcap%
\pgfsetroundjoin%
\definecolor{currentfill}{rgb}{0.101961,0.101961,0.101961}%
\pgfsetfillcolor{currentfill}%
\pgfsetlinewidth{0.301125pt}%
\definecolor{currentstroke}{rgb}{1.000000,1.000000,1.000000}%
\pgfsetstrokecolor{currentstroke}%
\pgfsetdash{}{0pt}%
\pgfsys@defobject{currentmarker}{\pgfqpoint{-0.013889in}{-0.013889in}}{\pgfqpoint{0.013889in}{0.013889in}}{%
\pgfpathmoveto{\pgfqpoint{0.000000in}{-0.013889in}}%
\pgfpathcurveto{\pgfqpoint{0.003683in}{-0.013889in}}{\pgfqpoint{0.007216in}{-0.012425in}}{\pgfqpoint{0.009821in}{-0.009821in}}%
\pgfpathcurveto{\pgfqpoint{0.012425in}{-0.007216in}}{\pgfqpoint{0.013889in}{-0.003683in}}{\pgfqpoint{0.013889in}{0.000000in}}%
\pgfpathcurveto{\pgfqpoint{0.013889in}{0.003683in}}{\pgfqpoint{0.012425in}{0.007216in}}{\pgfqpoint{0.009821in}{0.009821in}}%
\pgfpathcurveto{\pgfqpoint{0.007216in}{0.012425in}}{\pgfqpoint{0.003683in}{0.013889in}}{\pgfqpoint{0.000000in}{0.013889in}}%
\pgfpathcurveto{\pgfqpoint{-0.003683in}{0.013889in}}{\pgfqpoint{-0.007216in}{0.012425in}}{\pgfqpoint{-0.009821in}{0.009821in}}%
\pgfpathcurveto{\pgfqpoint{-0.012425in}{0.007216in}}{\pgfqpoint{-0.013889in}{0.003683in}}{\pgfqpoint{-0.013889in}{0.000000in}}%
\pgfpathcurveto{\pgfqpoint{-0.013889in}{-0.003683in}}{\pgfqpoint{-0.012425in}{-0.007216in}}{\pgfqpoint{-0.009821in}{-0.009821in}}%
\pgfpathcurveto{\pgfqpoint{-0.007216in}{-0.012425in}}{\pgfqpoint{-0.003683in}{-0.013889in}}{\pgfqpoint{0.000000in}{-0.013889in}}%
\pgfpathlineto{\pgfqpoint{0.000000in}{-0.013889in}}%
\pgfpathclose%
\pgfusepath{stroke,fill}%
}%
\begin{pgfscope}%
\pgfsys@transformshift{4.769395in}{1.468704in}%
\pgfsys@useobject{currentmarker}{}%
\end{pgfscope}%
\begin{pgfscope}%
\pgfsys@transformshift{4.754309in}{1.468704in}%
\pgfsys@useobject{currentmarker}{}%
\end{pgfscope}%
\begin{pgfscope}%
\pgfsys@transformshift{4.696716in}{1.468704in}%
\pgfsys@useobject{currentmarker}{}%
\end{pgfscope}%
\begin{pgfscope}%
\pgfsys@transformshift{4.752273in}{1.468704in}%
\pgfsys@useobject{currentmarker}{}%
\end{pgfscope}%
\end{pgfscope}%
\begin{pgfscope}%
\definecolor{textcolor}{rgb}{0.101961,0.101961,0.101961}%
\pgfsetstrokecolor{textcolor}%
\pgfsetfillcolor{textcolor}%
\pgftext[x=5.335484in,y=1.468704in,left,]{\color{textcolor}\sffamily\fontsize{6.000000}{7.200000}\selectfont 0.47}%
\end{pgfscope}%
\begin{pgfscope}%
\pgfpathrectangle{\pgfqpoint{4.259500in}{0.525000in}}{\pgfqpoint{1.293250in}{2.730000in}}%
\pgfusepath{clip}%
\pgfsetbuttcap%
\pgfsetroundjoin%
\definecolor{currentfill}{rgb}{0.101961,0.101961,0.101961}%
\pgfsetfillcolor{currentfill}%
\pgfsetlinewidth{0.301125pt}%
\definecolor{currentstroke}{rgb}{1.000000,1.000000,1.000000}%
\pgfsetstrokecolor{currentstroke}%
\pgfsetdash{}{0pt}%
\pgfsys@defobject{currentmarker}{\pgfqpoint{-0.013889in}{-0.013889in}}{\pgfqpoint{0.013889in}{0.013889in}}{%
\pgfpathmoveto{\pgfqpoint{0.000000in}{-0.013889in}}%
\pgfpathcurveto{\pgfqpoint{0.003683in}{-0.013889in}}{\pgfqpoint{0.007216in}{-0.012425in}}{\pgfqpoint{0.009821in}{-0.009821in}}%
\pgfpathcurveto{\pgfqpoint{0.012425in}{-0.007216in}}{\pgfqpoint{0.013889in}{-0.003683in}}{\pgfqpoint{0.013889in}{0.000000in}}%
\pgfpathcurveto{\pgfqpoint{0.013889in}{0.003683in}}{\pgfqpoint{0.012425in}{0.007216in}}{\pgfqpoint{0.009821in}{0.009821in}}%
\pgfpathcurveto{\pgfqpoint{0.007216in}{0.012425in}}{\pgfqpoint{0.003683in}{0.013889in}}{\pgfqpoint{0.000000in}{0.013889in}}%
\pgfpathcurveto{\pgfqpoint{-0.003683in}{0.013889in}}{\pgfqpoint{-0.007216in}{0.012425in}}{\pgfqpoint{-0.009821in}{0.009821in}}%
\pgfpathcurveto{\pgfqpoint{-0.012425in}{0.007216in}}{\pgfqpoint{-0.013889in}{0.003683in}}{\pgfqpoint{-0.013889in}{0.000000in}}%
\pgfpathcurveto{\pgfqpoint{-0.013889in}{-0.003683in}}{\pgfqpoint{-0.012425in}{-0.007216in}}{\pgfqpoint{-0.009821in}{-0.009821in}}%
\pgfpathcurveto{\pgfqpoint{-0.007216in}{-0.012425in}}{\pgfqpoint{-0.003683in}{-0.013889in}}{\pgfqpoint{0.000000in}{-0.013889in}}%
\pgfpathlineto{\pgfqpoint{0.000000in}{-0.013889in}}%
\pgfpathclose%
\pgfusepath{stroke,fill}%
}%
\begin{pgfscope}%
\pgfsys@transformshift{4.769395in}{1.300185in}%
\pgfsys@useobject{currentmarker}{}%
\end{pgfscope}%
\begin{pgfscope}%
\pgfsys@transformshift{4.747883in}{1.300185in}%
\pgfsys@useobject{currentmarker}{}%
\end{pgfscope}%
\begin{pgfscope}%
\pgfsys@transformshift{4.690223in}{1.300185in}%
\pgfsys@useobject{currentmarker}{}%
\end{pgfscope}%
\begin{pgfscope}%
\pgfsys@transformshift{4.758962in}{1.300185in}%
\pgfsys@useobject{currentmarker}{}%
\end{pgfscope}%
\end{pgfscope}%
\begin{pgfscope}%
\definecolor{textcolor}{rgb}{0.101961,0.101961,0.101961}%
\pgfsetstrokecolor{textcolor}%
\pgfsetfillcolor{textcolor}%
\pgftext[x=5.335484in,y=1.300185in,left,]{\color{textcolor}\sffamily\fontsize{6.000000}{7.200000}\selectfont 0.47}%
\end{pgfscope}%
\begin{pgfscope}%
\pgfpathrectangle{\pgfqpoint{4.259500in}{0.525000in}}{\pgfqpoint{1.293250in}{2.730000in}}%
\pgfusepath{clip}%
\pgfsetbuttcap%
\pgfsetroundjoin%
\definecolor{currentfill}{rgb}{0.101961,0.101961,0.101961}%
\pgfsetfillcolor{currentfill}%
\pgfsetlinewidth{0.301125pt}%
\definecolor{currentstroke}{rgb}{1.000000,1.000000,1.000000}%
\pgfsetstrokecolor{currentstroke}%
\pgfsetdash{}{0pt}%
\pgfsys@defobject{currentmarker}{\pgfqpoint{-0.013889in}{-0.013889in}}{\pgfqpoint{0.013889in}{0.013889in}}{%
\pgfpathmoveto{\pgfqpoint{0.000000in}{-0.013889in}}%
\pgfpathcurveto{\pgfqpoint{0.003683in}{-0.013889in}}{\pgfqpoint{0.007216in}{-0.012425in}}{\pgfqpoint{0.009821in}{-0.009821in}}%
\pgfpathcurveto{\pgfqpoint{0.012425in}{-0.007216in}}{\pgfqpoint{0.013889in}{-0.003683in}}{\pgfqpoint{0.013889in}{0.000000in}}%
\pgfpathcurveto{\pgfqpoint{0.013889in}{0.003683in}}{\pgfqpoint{0.012425in}{0.007216in}}{\pgfqpoint{0.009821in}{0.009821in}}%
\pgfpathcurveto{\pgfqpoint{0.007216in}{0.012425in}}{\pgfqpoint{0.003683in}{0.013889in}}{\pgfqpoint{0.000000in}{0.013889in}}%
\pgfpathcurveto{\pgfqpoint{-0.003683in}{0.013889in}}{\pgfqpoint{-0.007216in}{0.012425in}}{\pgfqpoint{-0.009821in}{0.009821in}}%
\pgfpathcurveto{\pgfqpoint{-0.012425in}{0.007216in}}{\pgfqpoint{-0.013889in}{0.003683in}}{\pgfqpoint{-0.013889in}{0.000000in}}%
\pgfpathcurveto{\pgfqpoint{-0.013889in}{-0.003683in}}{\pgfqpoint{-0.012425in}{-0.007216in}}{\pgfqpoint{-0.009821in}{-0.009821in}}%
\pgfpathcurveto{\pgfqpoint{-0.007216in}{-0.012425in}}{\pgfqpoint{-0.003683in}{-0.013889in}}{\pgfqpoint{0.000000in}{-0.013889in}}%
\pgfpathlineto{\pgfqpoint{0.000000in}{-0.013889in}}%
\pgfpathclose%
\pgfusepath{stroke,fill}%
}%
\begin{pgfscope}%
\pgfsys@transformshift{4.767279in}{1.131667in}%
\pgfsys@useobject{currentmarker}{}%
\end{pgfscope}%
\begin{pgfscope}%
\pgfsys@transformshift{4.745741in}{1.131667in}%
\pgfsys@useobject{currentmarker}{}%
\end{pgfscope}%
\begin{pgfscope}%
\pgfsys@transformshift{4.692387in}{1.131667in}%
\pgfsys@useobject{currentmarker}{}%
\end{pgfscope}%
\begin{pgfscope}%
\pgfsys@transformshift{4.754503in}{1.131667in}%
\pgfsys@useobject{currentmarker}{}%
\end{pgfscope}%
\end{pgfscope}%
\begin{pgfscope}%
\definecolor{textcolor}{rgb}{0.101961,0.101961,0.101961}%
\pgfsetstrokecolor{textcolor}%
\pgfsetfillcolor{textcolor}%
\pgftext[x=5.335484in,y=1.131667in,left,]{\color{textcolor}\sffamily\fontsize{6.000000}{7.200000}\selectfont 0.46}%
\end{pgfscope}%
\begin{pgfscope}%
\pgfpathrectangle{\pgfqpoint{4.259500in}{0.525000in}}{\pgfqpoint{1.293250in}{2.730000in}}%
\pgfusepath{clip}%
\pgfsetbuttcap%
\pgfsetroundjoin%
\definecolor{currentfill}{rgb}{0.101961,0.101961,0.101961}%
\pgfsetfillcolor{currentfill}%
\pgfsetlinewidth{0.301125pt}%
\definecolor{currentstroke}{rgb}{1.000000,1.000000,1.000000}%
\pgfsetstrokecolor{currentstroke}%
\pgfsetdash{}{0pt}%
\pgfsys@defobject{currentmarker}{\pgfqpoint{-0.013889in}{-0.013889in}}{\pgfqpoint{0.013889in}{0.013889in}}{%
\pgfpathmoveto{\pgfqpoint{0.000000in}{-0.013889in}}%
\pgfpathcurveto{\pgfqpoint{0.003683in}{-0.013889in}}{\pgfqpoint{0.007216in}{-0.012425in}}{\pgfqpoint{0.009821in}{-0.009821in}}%
\pgfpathcurveto{\pgfqpoint{0.012425in}{-0.007216in}}{\pgfqpoint{0.013889in}{-0.003683in}}{\pgfqpoint{0.013889in}{0.000000in}}%
\pgfpathcurveto{\pgfqpoint{0.013889in}{0.003683in}}{\pgfqpoint{0.012425in}{0.007216in}}{\pgfqpoint{0.009821in}{0.009821in}}%
\pgfpathcurveto{\pgfqpoint{0.007216in}{0.012425in}}{\pgfqpoint{0.003683in}{0.013889in}}{\pgfqpoint{0.000000in}{0.013889in}}%
\pgfpathcurveto{\pgfqpoint{-0.003683in}{0.013889in}}{\pgfqpoint{-0.007216in}{0.012425in}}{\pgfqpoint{-0.009821in}{0.009821in}}%
\pgfpathcurveto{\pgfqpoint{-0.012425in}{0.007216in}}{\pgfqpoint{-0.013889in}{0.003683in}}{\pgfqpoint{-0.013889in}{0.000000in}}%
\pgfpathcurveto{\pgfqpoint{-0.013889in}{-0.003683in}}{\pgfqpoint{-0.012425in}{-0.007216in}}{\pgfqpoint{-0.009821in}{-0.009821in}}%
\pgfpathcurveto{\pgfqpoint{-0.007216in}{-0.012425in}}{\pgfqpoint{-0.003683in}{-0.013889in}}{\pgfqpoint{0.000000in}{-0.013889in}}%
\pgfpathlineto{\pgfqpoint{0.000000in}{-0.013889in}}%
\pgfpathclose%
\pgfusepath{stroke,fill}%
}%
\begin{pgfscope}%
\pgfsys@transformshift{4.765163in}{0.963148in}%
\pgfsys@useobject{currentmarker}{}%
\end{pgfscope}%
\begin{pgfscope}%
\pgfsys@transformshift{4.750025in}{0.963148in}%
\pgfsys@useobject{currentmarker}{}%
\end{pgfscope}%
\begin{pgfscope}%
\pgfsys@transformshift{4.690223in}{0.963148in}%
\pgfsys@useobject{currentmarker}{}%
\end{pgfscope}%
\begin{pgfscope}%
\pgfsys@transformshift{4.754503in}{0.963148in}%
\pgfsys@useobject{currentmarker}{}%
\end{pgfscope}%
\end{pgfscope}%
\begin{pgfscope}%
\definecolor{textcolor}{rgb}{0.101961,0.101961,0.101961}%
\pgfsetstrokecolor{textcolor}%
\pgfsetfillcolor{textcolor}%
\pgftext[x=5.335484in,y=0.963148in,left,]{\color{textcolor}\sffamily\fontsize{6.000000}{7.200000}\selectfont 0.46}%
\end{pgfscope}%
\begin{pgfscope}%
\pgfpathrectangle{\pgfqpoint{4.259500in}{0.525000in}}{\pgfqpoint{1.293250in}{2.730000in}}%
\pgfusepath{clip}%
\pgfsetbuttcap%
\pgfsetroundjoin%
\definecolor{currentfill}{rgb}{0.101961,0.101961,0.101961}%
\pgfsetfillcolor{currentfill}%
\pgfsetlinewidth{0.301125pt}%
\definecolor{currentstroke}{rgb}{1.000000,1.000000,1.000000}%
\pgfsetstrokecolor{currentstroke}%
\pgfsetdash{}{0pt}%
\pgfsys@defobject{currentmarker}{\pgfqpoint{-0.013889in}{-0.013889in}}{\pgfqpoint{0.013889in}{0.013889in}}{%
\pgfpathmoveto{\pgfqpoint{0.000000in}{-0.013889in}}%
\pgfpathcurveto{\pgfqpoint{0.003683in}{-0.013889in}}{\pgfqpoint{0.007216in}{-0.012425in}}{\pgfqpoint{0.009821in}{-0.009821in}}%
\pgfpathcurveto{\pgfqpoint{0.012425in}{-0.007216in}}{\pgfqpoint{0.013889in}{-0.003683in}}{\pgfqpoint{0.013889in}{0.000000in}}%
\pgfpathcurveto{\pgfqpoint{0.013889in}{0.003683in}}{\pgfqpoint{0.012425in}{0.007216in}}{\pgfqpoint{0.009821in}{0.009821in}}%
\pgfpathcurveto{\pgfqpoint{0.007216in}{0.012425in}}{\pgfqpoint{0.003683in}{0.013889in}}{\pgfqpoint{0.000000in}{0.013889in}}%
\pgfpathcurveto{\pgfqpoint{-0.003683in}{0.013889in}}{\pgfqpoint{-0.007216in}{0.012425in}}{\pgfqpoint{-0.009821in}{0.009821in}}%
\pgfpathcurveto{\pgfqpoint{-0.012425in}{0.007216in}}{\pgfqpoint{-0.013889in}{0.003683in}}{\pgfqpoint{-0.013889in}{0.000000in}}%
\pgfpathcurveto{\pgfqpoint{-0.013889in}{-0.003683in}}{\pgfqpoint{-0.012425in}{-0.007216in}}{\pgfqpoint{-0.009821in}{-0.009821in}}%
\pgfpathcurveto{\pgfqpoint{-0.007216in}{-0.012425in}}{\pgfqpoint{-0.003683in}{-0.013889in}}{\pgfqpoint{0.000000in}{-0.013889in}}%
\pgfpathlineto{\pgfqpoint{0.000000in}{-0.013889in}}%
\pgfpathclose%
\pgfusepath{stroke,fill}%
}%
\begin{pgfscope}%
\pgfsys@transformshift{4.765163in}{0.794630in}%
\pgfsys@useobject{currentmarker}{}%
\end{pgfscope}%
\begin{pgfscope}%
\pgfsys@transformshift{4.745741in}{0.794630in}%
\pgfsys@useobject{currentmarker}{}%
\end{pgfscope}%
\begin{pgfscope}%
\pgfsys@transformshift{4.692387in}{0.794630in}%
\pgfsys@useobject{currentmarker}{}%
\end{pgfscope}%
\begin{pgfscope}%
\pgfsys@transformshift{4.752273in}{0.794630in}%
\pgfsys@useobject{currentmarker}{}%
\end{pgfscope}%
\end{pgfscope}%
\begin{pgfscope}%
\definecolor{textcolor}{rgb}{0.101961,0.101961,0.101961}%
\pgfsetstrokecolor{textcolor}%
\pgfsetfillcolor{textcolor}%
\pgftext[x=5.335484in,y=0.794630in,left,]{\color{textcolor}\sffamily\fontsize{6.000000}{7.200000}\selectfont 0.46}%
\end{pgfscope}%
\begin{pgfscope}%
\pgfpathrectangle{\pgfqpoint{4.259500in}{0.525000in}}{\pgfqpoint{1.293250in}{2.730000in}}%
\pgfusepath{clip}%
\pgfsetbuttcap%
\pgfsetroundjoin%
\definecolor{currentfill}{rgb}{0.101961,0.101961,0.101961}%
\pgfsetfillcolor{currentfill}%
\pgfsetlinewidth{0.301125pt}%
\definecolor{currentstroke}{rgb}{1.000000,1.000000,1.000000}%
\pgfsetstrokecolor{currentstroke}%
\pgfsetdash{}{0pt}%
\pgfsys@defobject{currentmarker}{\pgfqpoint{-0.013889in}{-0.013889in}}{\pgfqpoint{0.013889in}{0.013889in}}{%
\pgfpathmoveto{\pgfqpoint{0.000000in}{-0.013889in}}%
\pgfpathcurveto{\pgfqpoint{0.003683in}{-0.013889in}}{\pgfqpoint{0.007216in}{-0.012425in}}{\pgfqpoint{0.009821in}{-0.009821in}}%
\pgfpathcurveto{\pgfqpoint{0.012425in}{-0.007216in}}{\pgfqpoint{0.013889in}{-0.003683in}}{\pgfqpoint{0.013889in}{0.000000in}}%
\pgfpathcurveto{\pgfqpoint{0.013889in}{0.003683in}}{\pgfqpoint{0.012425in}{0.007216in}}{\pgfqpoint{0.009821in}{0.009821in}}%
\pgfpathcurveto{\pgfqpoint{0.007216in}{0.012425in}}{\pgfqpoint{0.003683in}{0.013889in}}{\pgfqpoint{0.000000in}{0.013889in}}%
\pgfpathcurveto{\pgfqpoint{-0.003683in}{0.013889in}}{\pgfqpoint{-0.007216in}{0.012425in}}{\pgfqpoint{-0.009821in}{0.009821in}}%
\pgfpathcurveto{\pgfqpoint{-0.012425in}{0.007216in}}{\pgfqpoint{-0.013889in}{0.003683in}}{\pgfqpoint{-0.013889in}{0.000000in}}%
\pgfpathcurveto{\pgfqpoint{-0.013889in}{-0.003683in}}{\pgfqpoint{-0.012425in}{-0.007216in}}{\pgfqpoint{-0.009821in}{-0.009821in}}%
\pgfpathcurveto{\pgfqpoint{-0.007216in}{-0.012425in}}{\pgfqpoint{-0.003683in}{-0.013889in}}{\pgfqpoint{0.000000in}{-0.013889in}}%
\pgfpathlineto{\pgfqpoint{0.000000in}{-0.013889in}}%
\pgfpathclose%
\pgfusepath{stroke,fill}%
}%
\begin{pgfscope}%
\pgfsys@transformshift{4.769395in}{0.626111in}%
\pgfsys@useobject{currentmarker}{}%
\end{pgfscope}%
\begin{pgfscope}%
\pgfsys@transformshift{4.752167in}{0.626111in}%
\pgfsys@useobject{currentmarker}{}%
\end{pgfscope}%
\begin{pgfscope}%
\pgfsys@transformshift{4.688058in}{0.626111in}%
\pgfsys@useobject{currentmarker}{}%
\end{pgfscope}%
\begin{pgfscope}%
\pgfsys@transformshift{4.758962in}{0.626111in}%
\pgfsys@useobject{currentmarker}{}%
\end{pgfscope}%
\end{pgfscope}%
\begin{pgfscope}%
\definecolor{textcolor}{rgb}{0.101961,0.101961,0.101961}%
\pgfsetstrokecolor{textcolor}%
\pgfsetfillcolor{textcolor}%
\pgftext[x=5.335484in,y=0.626111in,left,]{\color{textcolor}\sffamily\fontsize{6.000000}{7.200000}\selectfont 0.47}%
\end{pgfscope}%
\begin{pgfscope}%
\pgfsetbuttcap%
\pgfsetmiterjoin%
\definecolor{currentfill}{rgb}{1.000000,1.000000,1.000000}%
\pgfsetfillcolor{currentfill}%
\pgfsetlinewidth{0.000000pt}%
\definecolor{currentstroke}{rgb}{0.000000,0.000000,0.000000}%
\pgfsetstrokecolor{currentstroke}%
\pgfsetstrokeopacity{0.000000}%
\pgfsetdash{}{0pt}%
\pgfpathmoveto{\pgfqpoint{5.636750in}{0.525000in}}%
\pgfpathlineto{\pgfqpoint{6.930000in}{0.525000in}}%
\pgfpathlineto{\pgfqpoint{6.930000in}{3.255000in}}%
\pgfpathlineto{\pgfqpoint{5.636750in}{3.255000in}}%
\pgfpathlineto{\pgfqpoint{5.636750in}{0.525000in}}%
\pgfpathclose%
\pgfusepath{fill}%
\end{pgfscope}%
\begin{pgfscope}%
\pgfpathrectangle{\pgfqpoint{5.636750in}{0.525000in}}{\pgfqpoint{1.293250in}{2.730000in}}%
\pgfusepath{clip}%
\pgfsetrectcap%
\pgfsetroundjoin%
\pgfsetlinewidth{0.401500pt}%
\definecolor{currentstroke}{rgb}{0.901961,0.901961,0.901961}%
\pgfsetstrokecolor{currentstroke}%
\pgfsetdash{}{0pt}%
\pgfpathmoveto{\pgfqpoint{5.636750in}{0.525000in}}%
\pgfpathlineto{\pgfqpoint{5.636750in}{3.255000in}}%
\pgfusepath{stroke}%
\end{pgfscope}%
\begin{pgfscope}%
\pgfsetbuttcap%
\pgfsetroundjoin%
\definecolor{currentfill}{rgb}{0.101961,0.101961,0.101961}%
\pgfsetfillcolor{currentfill}%
\pgfsetlinewidth{0.803000pt}%
\definecolor{currentstroke}{rgb}{0.101961,0.101961,0.101961}%
\pgfsetstrokecolor{currentstroke}%
\pgfsetdash{}{0pt}%
\pgfsys@defobject{currentmarker}{\pgfqpoint{0.000000in}{-0.027778in}}{\pgfqpoint{0.000000in}{0.000000in}}{%
\pgfpathmoveto{\pgfqpoint{0.000000in}{0.000000in}}%
\pgfpathlineto{\pgfqpoint{0.000000in}{-0.027778in}}%
\pgfusepath{stroke,fill}%
}%
\begin{pgfscope}%
\pgfsys@transformshift{5.636750in}{0.525000in}%
\pgfsys@useobject{currentmarker}{}%
\end{pgfscope}%
\end{pgfscope}%
\begin{pgfscope}%
\definecolor{textcolor}{rgb}{0.101961,0.101961,0.101961}%
\pgfsetstrokecolor{textcolor}%
\pgfsetfillcolor{textcolor}%
\pgftext[x=5.636750in,y=0.448611in,,top]{\color{textcolor}\sffamily\fontsize{7.000000}{8.400000}\selectfont 0}%
\end{pgfscope}%
\begin{pgfscope}%
\pgfpathrectangle{\pgfqpoint{5.636750in}{0.525000in}}{\pgfqpoint{1.293250in}{2.730000in}}%
\pgfusepath{clip}%
\pgfsetrectcap%
\pgfsetroundjoin%
\pgfsetlinewidth{0.401500pt}%
\definecolor{currentstroke}{rgb}{0.901961,0.901961,0.901961}%
\pgfsetstrokecolor{currentstroke}%
\pgfsetdash{}{0pt}%
\pgfpathmoveto{\pgfqpoint{6.154050in}{0.525000in}}%
\pgfpathlineto{\pgfqpoint{6.154050in}{3.255000in}}%
\pgfusepath{stroke}%
\end{pgfscope}%
\begin{pgfscope}%
\pgfsetbuttcap%
\pgfsetroundjoin%
\definecolor{currentfill}{rgb}{0.101961,0.101961,0.101961}%
\pgfsetfillcolor{currentfill}%
\pgfsetlinewidth{0.803000pt}%
\definecolor{currentstroke}{rgb}{0.101961,0.101961,0.101961}%
\pgfsetstrokecolor{currentstroke}%
\pgfsetdash{}{0pt}%
\pgfsys@defobject{currentmarker}{\pgfqpoint{0.000000in}{-0.027778in}}{\pgfqpoint{0.000000in}{0.000000in}}{%
\pgfpathmoveto{\pgfqpoint{0.000000in}{0.000000in}}%
\pgfpathlineto{\pgfqpoint{0.000000in}{-0.027778in}}%
\pgfusepath{stroke,fill}%
}%
\begin{pgfscope}%
\pgfsys@transformshift{6.154050in}{0.525000in}%
\pgfsys@useobject{currentmarker}{}%
\end{pgfscope}%
\end{pgfscope}%
\begin{pgfscope}%
\definecolor{textcolor}{rgb}{0.101961,0.101961,0.101961}%
\pgfsetstrokecolor{textcolor}%
\pgfsetfillcolor{textcolor}%
\pgftext[x=6.154050in,y=0.448611in,,top]{\color{textcolor}\sffamily\fontsize{7.000000}{8.400000}\selectfont 0.5}%
\end{pgfscope}%
\begin{pgfscope}%
\pgfpathrectangle{\pgfqpoint{5.636750in}{0.525000in}}{\pgfqpoint{1.293250in}{2.730000in}}%
\pgfusepath{clip}%
\pgfsetrectcap%
\pgfsetroundjoin%
\pgfsetlinewidth{0.401500pt}%
\definecolor{currentstroke}{rgb}{0.901961,0.901961,0.901961}%
\pgfsetstrokecolor{currentstroke}%
\pgfsetdash{}{0pt}%
\pgfpathmoveto{\pgfqpoint{6.671350in}{0.525000in}}%
\pgfpathlineto{\pgfqpoint{6.671350in}{3.255000in}}%
\pgfusepath{stroke}%
\end{pgfscope}%
\begin{pgfscope}%
\pgfsetbuttcap%
\pgfsetroundjoin%
\definecolor{currentfill}{rgb}{0.101961,0.101961,0.101961}%
\pgfsetfillcolor{currentfill}%
\pgfsetlinewidth{0.803000pt}%
\definecolor{currentstroke}{rgb}{0.101961,0.101961,0.101961}%
\pgfsetstrokecolor{currentstroke}%
\pgfsetdash{}{0pt}%
\pgfsys@defobject{currentmarker}{\pgfqpoint{0.000000in}{-0.027778in}}{\pgfqpoint{0.000000in}{0.000000in}}{%
\pgfpathmoveto{\pgfqpoint{0.000000in}{0.000000in}}%
\pgfpathlineto{\pgfqpoint{0.000000in}{-0.027778in}}%
\pgfusepath{stroke,fill}%
}%
\begin{pgfscope}%
\pgfsys@transformshift{6.671350in}{0.525000in}%
\pgfsys@useobject{currentmarker}{}%
\end{pgfscope}%
\end{pgfscope}%
\begin{pgfscope}%
\definecolor{textcolor}{rgb}{0.101961,0.101961,0.101961}%
\pgfsetstrokecolor{textcolor}%
\pgfsetfillcolor{textcolor}%
\pgftext[x=6.671350in,y=0.448611in,,top]{\color{textcolor}\sffamily\fontsize{7.000000}{8.400000}\selectfont 1}%
\end{pgfscope}%
\begin{pgfscope}%
\definecolor{textcolor}{rgb}{0.101961,0.101961,0.101961}%
\pgfsetstrokecolor{textcolor}%
\pgfsetfillcolor{textcolor}%
\pgftext[x=6.283375in,y=0.334414in,,top]{\color{textcolor}\sffamily\fontsize{7.500000}{9.000000}\selectfont Dashcam}%
\end{pgfscope}%
\begin{pgfscope}%
\pgfpathrectangle{\pgfqpoint{5.636750in}{0.525000in}}{\pgfqpoint{1.293250in}{2.730000in}}%
\pgfusepath{clip}%
\pgfsetbuttcap%
\pgfsetmiterjoin%
\definecolor{currentfill}{rgb}{0.156863,0.356863,0.568627}%
\pgfsetfillcolor{currentfill}%
\pgfsetlinewidth{0.000000pt}%
\definecolor{currentstroke}{rgb}{0.000000,0.000000,0.000000}%
\pgfsetstrokecolor{currentstroke}%
\pgfsetstrokeopacity{0.000000}%
\pgfsetdash{}{0pt}%
\pgfpathmoveto{\pgfqpoint{5.636750in}{3.101648in}}%
\pgfpathlineto{\pgfqpoint{6.379633in}{3.101648in}}%
\pgfpathlineto{\pgfqpoint{6.379633in}{3.206130in}}%
\pgfpathlineto{\pgfqpoint{5.636750in}{3.206130in}}%
\pgfpathlineto{\pgfqpoint{5.636750in}{3.101648in}}%
\pgfpathclose%
\pgfusepath{fill}%
\end{pgfscope}%
\begin{pgfscope}%
\pgfpathrectangle{\pgfqpoint{5.636750in}{0.525000in}}{\pgfqpoint{1.293250in}{2.730000in}}%
\pgfusepath{clip}%
\pgfsetbuttcap%
\pgfsetmiterjoin%
\definecolor{currentfill}{rgb}{0.784314,0.784314,0.784314}%
\pgfsetfillcolor{currentfill}%
\pgfsetlinewidth{0.000000pt}%
\definecolor{currentstroke}{rgb}{0.000000,0.000000,0.000000}%
\pgfsetstrokecolor{currentstroke}%
\pgfsetstrokeopacity{0.000000}%
\pgfsetdash{}{0pt}%
\pgfpathmoveto{\pgfqpoint{5.636750in}{2.933130in}}%
\pgfpathlineto{\pgfqpoint{6.438082in}{2.933130in}}%
\pgfpathlineto{\pgfqpoint{6.438082in}{3.037611in}}%
\pgfpathlineto{\pgfqpoint{5.636750in}{3.037611in}}%
\pgfpathlineto{\pgfqpoint{5.636750in}{2.933130in}}%
\pgfpathclose%
\pgfusepath{fill}%
\end{pgfscope}%
\begin{pgfscope}%
\pgfpathrectangle{\pgfqpoint{5.636750in}{0.525000in}}{\pgfqpoint{1.293250in}{2.730000in}}%
\pgfusepath{clip}%
\pgfsetbuttcap%
\pgfsetmiterjoin%
\definecolor{currentfill}{rgb}{0.784314,0.784314,0.784314}%
\pgfsetfillcolor{currentfill}%
\pgfsetlinewidth{0.000000pt}%
\definecolor{currentstroke}{rgb}{0.000000,0.000000,0.000000}%
\pgfsetstrokecolor{currentstroke}%
\pgfsetstrokeopacity{0.000000}%
\pgfsetdash{}{0pt}%
\pgfpathmoveto{\pgfqpoint{5.636750in}{2.596093in}}%
\pgfpathlineto{\pgfqpoint{6.618714in}{2.596093in}}%
\pgfpathlineto{\pgfqpoint{6.618714in}{2.700574in}}%
\pgfpathlineto{\pgfqpoint{5.636750in}{2.700574in}}%
\pgfpathlineto{\pgfqpoint{5.636750in}{2.596093in}}%
\pgfpathclose%
\pgfusepath{fill}%
\end{pgfscope}%
\begin{pgfscope}%
\pgfpathrectangle{\pgfqpoint{5.636750in}{0.525000in}}{\pgfqpoint{1.293250in}{2.730000in}}%
\pgfusepath{clip}%
\pgfsetbuttcap%
\pgfsetmiterjoin%
\definecolor{currentfill}{rgb}{0.784314,0.784314,0.784314}%
\pgfsetfillcolor{currentfill}%
\pgfsetlinewidth{0.000000pt}%
\definecolor{currentstroke}{rgb}{0.000000,0.000000,0.000000}%
\pgfsetstrokecolor{currentstroke}%
\pgfsetstrokeopacity{0.000000}%
\pgfsetdash{}{0pt}%
\pgfpathmoveto{\pgfqpoint{5.636750in}{2.427574in}}%
\pgfpathlineto{\pgfqpoint{6.279596in}{2.427574in}}%
\pgfpathlineto{\pgfqpoint{6.279596in}{2.532056in}}%
\pgfpathlineto{\pgfqpoint{5.636750in}{2.532056in}}%
\pgfpathlineto{\pgfqpoint{5.636750in}{2.427574in}}%
\pgfpathclose%
\pgfusepath{fill}%
\end{pgfscope}%
\begin{pgfscope}%
\pgfpathrectangle{\pgfqpoint{5.636750in}{0.525000in}}{\pgfqpoint{1.293250in}{2.730000in}}%
\pgfusepath{clip}%
\pgfsetbuttcap%
\pgfsetmiterjoin%
\definecolor{currentfill}{rgb}{0.784314,0.784314,0.784314}%
\pgfsetfillcolor{currentfill}%
\pgfsetlinewidth{0.000000pt}%
\definecolor{currentstroke}{rgb}{0.000000,0.000000,0.000000}%
\pgfsetstrokecolor{currentstroke}%
\pgfsetstrokeopacity{0.000000}%
\pgfsetdash{}{0pt}%
\pgfpathmoveto{\pgfqpoint{5.636750in}{2.259056in}}%
\pgfpathlineto{\pgfqpoint{6.061780in}{2.259056in}}%
\pgfpathlineto{\pgfqpoint{6.061780in}{2.363537in}}%
\pgfpathlineto{\pgfqpoint{5.636750in}{2.363537in}}%
\pgfpathlineto{\pgfqpoint{5.636750in}{2.259056in}}%
\pgfpathclose%
\pgfusepath{fill}%
\end{pgfscope}%
\begin{pgfscope}%
\pgfpathrectangle{\pgfqpoint{5.636750in}{0.525000in}}{\pgfqpoint{1.293250in}{2.730000in}}%
\pgfusepath{clip}%
\pgfsetbuttcap%
\pgfsetmiterjoin%
\definecolor{currentfill}{rgb}{0.478431,0.478431,0.478431}%
\pgfsetfillcolor{currentfill}%
\pgfsetlinewidth{0.000000pt}%
\definecolor{currentstroke}{rgb}{0.000000,0.000000,0.000000}%
\pgfsetstrokecolor{currentstroke}%
\pgfsetstrokeopacity{0.000000}%
\pgfsetdash{}{0pt}%
\pgfpathmoveto{\pgfqpoint{5.636750in}{2.090537in}}%
\pgfpathlineto{\pgfqpoint{5.974429in}{2.090537in}}%
\pgfpathlineto{\pgfqpoint{5.974429in}{2.195019in}}%
\pgfpathlineto{\pgfqpoint{5.636750in}{2.195019in}}%
\pgfpathlineto{\pgfqpoint{5.636750in}{2.090537in}}%
\pgfpathclose%
\pgfusepath{fill}%
\end{pgfscope}%
\begin{pgfscope}%
\pgfpathrectangle{\pgfqpoint{5.636750in}{0.525000in}}{\pgfqpoint{1.293250in}{2.730000in}}%
\pgfusepath{clip}%
\pgfsetbuttcap%
\pgfsetmiterjoin%
\definecolor{currentfill}{rgb}{0.784314,0.784314,0.784314}%
\pgfsetfillcolor{currentfill}%
\pgfsetlinewidth{0.000000pt}%
\definecolor{currentstroke}{rgb}{0.000000,0.000000,0.000000}%
\pgfsetstrokecolor{currentstroke}%
\pgfsetstrokeopacity{0.000000}%
\pgfsetdash{}{0pt}%
\pgfpathmoveto{\pgfqpoint{5.636750in}{1.922019in}}%
\pgfpathlineto{\pgfqpoint{5.972392in}{1.922019in}}%
\pgfpathlineto{\pgfqpoint{5.972392in}{2.026500in}}%
\pgfpathlineto{\pgfqpoint{5.636750in}{2.026500in}}%
\pgfpathlineto{\pgfqpoint{5.636750in}{1.922019in}}%
\pgfpathclose%
\pgfusepath{fill}%
\end{pgfscope}%
\begin{pgfscope}%
\pgfpathrectangle{\pgfqpoint{5.636750in}{0.525000in}}{\pgfqpoint{1.293250in}{2.730000in}}%
\pgfusepath{clip}%
\pgfsetbuttcap%
\pgfsetmiterjoin%
\definecolor{currentfill}{rgb}{0.784314,0.784314,0.784314}%
\pgfsetfillcolor{currentfill}%
\pgfsetlinewidth{0.000000pt}%
\definecolor{currentstroke}{rgb}{0.000000,0.000000,0.000000}%
\pgfsetstrokecolor{currentstroke}%
\pgfsetstrokeopacity{0.000000}%
\pgfsetdash{}{0pt}%
\pgfpathmoveto{\pgfqpoint{5.636750in}{1.753500in}}%
\pgfpathlineto{\pgfqpoint{5.918212in}{1.753500in}}%
\pgfpathlineto{\pgfqpoint{5.918212in}{1.857981in}}%
\pgfpathlineto{\pgfqpoint{5.636750in}{1.857981in}}%
\pgfpathlineto{\pgfqpoint{5.636750in}{1.753500in}}%
\pgfpathclose%
\pgfusepath{fill}%
\end{pgfscope}%
\begin{pgfscope}%
\pgfpathrectangle{\pgfqpoint{5.636750in}{0.525000in}}{\pgfqpoint{1.293250in}{2.730000in}}%
\pgfusepath{clip}%
\pgfsetbuttcap%
\pgfsetmiterjoin%
\definecolor{currentfill}{rgb}{0.478431,0.478431,0.478431}%
\pgfsetfillcolor{currentfill}%
\pgfsetlinewidth{0.000000pt}%
\definecolor{currentstroke}{rgb}{0.000000,0.000000,0.000000}%
\pgfsetstrokecolor{currentstroke}%
\pgfsetstrokeopacity{0.000000}%
\pgfsetdash{}{0pt}%
\pgfpathmoveto{\pgfqpoint{5.636750in}{1.584981in}}%
\pgfpathlineto{\pgfqpoint{5.718895in}{1.584981in}}%
\pgfpathlineto{\pgfqpoint{5.718895in}{1.689463in}}%
\pgfpathlineto{\pgfqpoint{5.636750in}{1.689463in}}%
\pgfpathlineto{\pgfqpoint{5.636750in}{1.584981in}}%
\pgfpathclose%
\pgfusepath{fill}%
\end{pgfscope}%
\begin{pgfscope}%
\pgfpathrectangle{\pgfqpoint{5.636750in}{0.525000in}}{\pgfqpoint{1.293250in}{2.730000in}}%
\pgfusepath{clip}%
\pgfsetbuttcap%
\pgfsetmiterjoin%
\definecolor{currentfill}{rgb}{0.784314,0.784314,0.784314}%
\pgfsetfillcolor{currentfill}%
\pgfsetlinewidth{0.000000pt}%
\definecolor{currentstroke}{rgb}{0.000000,0.000000,0.000000}%
\pgfsetstrokecolor{currentstroke}%
\pgfsetstrokeopacity{0.000000}%
\pgfsetdash{}{0pt}%
\pgfpathmoveto{\pgfqpoint{5.636750in}{1.416463in}}%
\pgfpathlineto{\pgfqpoint{5.719070in}{1.416463in}}%
\pgfpathlineto{\pgfqpoint{5.719070in}{1.520944in}}%
\pgfpathlineto{\pgfqpoint{5.636750in}{1.520944in}}%
\pgfpathlineto{\pgfqpoint{5.636750in}{1.416463in}}%
\pgfpathclose%
\pgfusepath{fill}%
\end{pgfscope}%
\begin{pgfscope}%
\pgfpathrectangle{\pgfqpoint{5.636750in}{0.525000in}}{\pgfqpoint{1.293250in}{2.730000in}}%
\pgfusepath{clip}%
\pgfsetbuttcap%
\pgfsetmiterjoin%
\definecolor{currentfill}{rgb}{0.478431,0.478431,0.478431}%
\pgfsetfillcolor{currentfill}%
\pgfsetlinewidth{0.000000pt}%
\definecolor{currentstroke}{rgb}{0.000000,0.000000,0.000000}%
\pgfsetstrokecolor{currentstroke}%
\pgfsetstrokeopacity{0.000000}%
\pgfsetdash{}{0pt}%
\pgfpathmoveto{\pgfqpoint{5.636750in}{1.247944in}}%
\pgfpathlineto{\pgfqpoint{5.693467in}{1.247944in}}%
\pgfpathlineto{\pgfqpoint{5.693467in}{1.352426in}}%
\pgfpathlineto{\pgfqpoint{5.636750in}{1.352426in}}%
\pgfpathlineto{\pgfqpoint{5.636750in}{1.247944in}}%
\pgfpathclose%
\pgfusepath{fill}%
\end{pgfscope}%
\begin{pgfscope}%
\pgfpathrectangle{\pgfqpoint{5.636750in}{0.525000in}}{\pgfqpoint{1.293250in}{2.730000in}}%
\pgfusepath{clip}%
\pgfsetbuttcap%
\pgfsetmiterjoin%
\definecolor{currentfill}{rgb}{0.784314,0.784314,0.784314}%
\pgfsetfillcolor{currentfill}%
\pgfsetlinewidth{0.000000pt}%
\definecolor{currentstroke}{rgb}{0.000000,0.000000,0.000000}%
\pgfsetstrokecolor{currentstroke}%
\pgfsetstrokeopacity{0.000000}%
\pgfsetdash{}{0pt}%
\pgfpathmoveto{\pgfqpoint{5.636750in}{1.079426in}}%
\pgfpathlineto{\pgfqpoint{5.716471in}{1.079426in}}%
\pgfpathlineto{\pgfqpoint{5.716471in}{1.183907in}}%
\pgfpathlineto{\pgfqpoint{5.636750in}{1.183907in}}%
\pgfpathlineto{\pgfqpoint{5.636750in}{1.079426in}}%
\pgfpathclose%
\pgfusepath{fill}%
\end{pgfscope}%
\begin{pgfscope}%
\pgfpathrectangle{\pgfqpoint{5.636750in}{0.525000in}}{\pgfqpoint{1.293250in}{2.730000in}}%
\pgfusepath{clip}%
\pgfsetbuttcap%
\pgfsetmiterjoin%
\definecolor{currentfill}{rgb}{0.478431,0.478431,0.478431}%
\pgfsetfillcolor{currentfill}%
\pgfsetlinewidth{0.000000pt}%
\definecolor{currentstroke}{rgb}{0.000000,0.000000,0.000000}%
\pgfsetstrokecolor{currentstroke}%
\pgfsetstrokeopacity{0.000000}%
\pgfsetdash{}{0pt}%
\pgfpathmoveto{\pgfqpoint{5.636750in}{0.910907in}}%
\pgfpathlineto{\pgfqpoint{5.717654in}{0.910907in}}%
\pgfpathlineto{\pgfqpoint{5.717654in}{1.015389in}}%
\pgfpathlineto{\pgfqpoint{5.636750in}{1.015389in}}%
\pgfpathlineto{\pgfqpoint{5.636750in}{0.910907in}}%
\pgfpathclose%
\pgfusepath{fill}%
\end{pgfscope}%
\begin{pgfscope}%
\pgfpathrectangle{\pgfqpoint{5.636750in}{0.525000in}}{\pgfqpoint{1.293250in}{2.730000in}}%
\pgfusepath{clip}%
\pgfsetbuttcap%
\pgfsetmiterjoin%
\definecolor{currentfill}{rgb}{0.478431,0.478431,0.478431}%
\pgfsetfillcolor{currentfill}%
\pgfsetlinewidth{0.000000pt}%
\definecolor{currentstroke}{rgb}{0.000000,0.000000,0.000000}%
\pgfsetstrokecolor{currentstroke}%
\pgfsetstrokeopacity{0.000000}%
\pgfsetdash{}{0pt}%
\pgfpathmoveto{\pgfqpoint{5.636750in}{0.742389in}}%
\pgfpathlineto{\pgfqpoint{5.718604in}{0.742389in}}%
\pgfpathlineto{\pgfqpoint{5.718604in}{0.846870in}}%
\pgfpathlineto{\pgfqpoint{5.636750in}{0.846870in}}%
\pgfpathlineto{\pgfqpoint{5.636750in}{0.742389in}}%
\pgfpathclose%
\pgfusepath{fill}%
\end{pgfscope}%
\begin{pgfscope}%
\pgfpathrectangle{\pgfqpoint{5.636750in}{0.525000in}}{\pgfqpoint{1.293250in}{2.730000in}}%
\pgfusepath{clip}%
\pgfsetbuttcap%
\pgfsetmiterjoin%
\definecolor{currentfill}{rgb}{0.478431,0.305882,0.611765}%
\pgfsetfillcolor{currentfill}%
\pgfsetlinewidth{0.000000pt}%
\definecolor{currentstroke}{rgb}{0.000000,0.000000,0.000000}%
\pgfsetstrokecolor{currentstroke}%
\pgfsetstrokeopacity{0.000000}%
\pgfsetdash{}{0pt}%
\pgfpathmoveto{\pgfqpoint{5.636750in}{0.573870in}}%
\pgfpathlineto{\pgfqpoint{5.636750in}{0.573870in}}%
\pgfpathlineto{\pgfqpoint{5.636750in}{0.678352in}}%
\pgfpathlineto{\pgfqpoint{5.636750in}{0.678352in}}%
\pgfpathlineto{\pgfqpoint{5.636750in}{0.573870in}}%
\pgfpathclose%
\pgfusepath{fill}%
\end{pgfscope}%
\begin{pgfscope}%
\pgfsetrectcap%
\pgfsetmiterjoin%
\pgfsetlinewidth{0.803000pt}%
\definecolor{currentstroke}{rgb}{0.101961,0.101961,0.101961}%
\pgfsetstrokecolor{currentstroke}%
\pgfsetdash{}{0pt}%
\pgfpathmoveto{\pgfqpoint{5.636750in}{0.525000in}}%
\pgfpathlineto{\pgfqpoint{5.636750in}{3.255000in}}%
\pgfusepath{stroke}%
\end{pgfscope}%
\begin{pgfscope}%
\pgfsetrectcap%
\pgfsetmiterjoin%
\pgfsetlinewidth{0.803000pt}%
\definecolor{currentstroke}{rgb}{0.101961,0.101961,0.101961}%
\pgfsetstrokecolor{currentstroke}%
\pgfsetdash{}{0pt}%
\pgfpathmoveto{\pgfqpoint{5.636750in}{0.525000in}}%
\pgfpathlineto{\pgfqpoint{6.930000in}{0.525000in}}%
\pgfusepath{stroke}%
\end{pgfscope}%
\begin{pgfscope}%
\pgfpathrectangle{\pgfqpoint{5.636750in}{0.525000in}}{\pgfqpoint{1.293250in}{2.730000in}}%
\pgfusepath{clip}%
\pgfsetbuttcap%
\pgfsetroundjoin%
\definecolor{currentfill}{rgb}{0.101961,0.101961,0.101961}%
\pgfsetfillcolor{currentfill}%
\pgfsetlinewidth{0.301125pt}%
\definecolor{currentstroke}{rgb}{1.000000,1.000000,1.000000}%
\pgfsetstrokecolor{currentstroke}%
\pgfsetdash{}{0pt}%
\pgfsys@defobject{currentmarker}{\pgfqpoint{-0.013889in}{-0.013889in}}{\pgfqpoint{0.013889in}{0.013889in}}{%
\pgfpathmoveto{\pgfqpoint{0.000000in}{-0.013889in}}%
\pgfpathcurveto{\pgfqpoint{0.003683in}{-0.013889in}}{\pgfqpoint{0.007216in}{-0.012425in}}{\pgfqpoint{0.009821in}{-0.009821in}}%
\pgfpathcurveto{\pgfqpoint{0.012425in}{-0.007216in}}{\pgfqpoint{0.013889in}{-0.003683in}}{\pgfqpoint{0.013889in}{0.000000in}}%
\pgfpathcurveto{\pgfqpoint{0.013889in}{0.003683in}}{\pgfqpoint{0.012425in}{0.007216in}}{\pgfqpoint{0.009821in}{0.009821in}}%
\pgfpathcurveto{\pgfqpoint{0.007216in}{0.012425in}}{\pgfqpoint{0.003683in}{0.013889in}}{\pgfqpoint{0.000000in}{0.013889in}}%
\pgfpathcurveto{\pgfqpoint{-0.003683in}{0.013889in}}{\pgfqpoint{-0.007216in}{0.012425in}}{\pgfqpoint{-0.009821in}{0.009821in}}%
\pgfpathcurveto{\pgfqpoint{-0.012425in}{0.007216in}}{\pgfqpoint{-0.013889in}{0.003683in}}{\pgfqpoint{-0.013889in}{0.000000in}}%
\pgfpathcurveto{\pgfqpoint{-0.013889in}{-0.003683in}}{\pgfqpoint{-0.012425in}{-0.007216in}}{\pgfqpoint{-0.009821in}{-0.009821in}}%
\pgfpathcurveto{\pgfqpoint{-0.007216in}{-0.012425in}}{\pgfqpoint{-0.003683in}{-0.013889in}}{\pgfqpoint{0.000000in}{-0.013889in}}%
\pgfpathlineto{\pgfqpoint{0.000000in}{-0.013889in}}%
\pgfpathclose%
\pgfusepath{stroke,fill}%
}%
\begin{pgfscope}%
\pgfsys@transformshift{6.372492in}{3.153889in}%
\pgfsys@useobject{currentmarker}{}%
\end{pgfscope}%
\begin{pgfscope}%
\pgfsys@transformshift{6.357838in}{3.153889in}%
\pgfsys@useobject{currentmarker}{}%
\end{pgfscope}%
\begin{pgfscope}%
\pgfsys@transformshift{6.386661in}{3.153889in}%
\pgfsys@useobject{currentmarker}{}%
\end{pgfscope}%
\begin{pgfscope}%
\pgfsys@transformshift{6.401541in}{3.153889in}%
\pgfsys@useobject{currentmarker}{}%
\end{pgfscope}%
\end{pgfscope}%
\begin{pgfscope}%
\definecolor{textcolor}{rgb}{0.101961,0.101961,0.101961}%
\pgfsetstrokecolor{textcolor}%
\pgfsetfillcolor{textcolor}%
\pgftext[x=6.712734in,y=3.153889in,left,]{\color{textcolor}\sffamily\fontsize{6.000000}{7.200000}\bfseries\selectfont 0.72}%
\end{pgfscope}%
\begin{pgfscope}%
\pgfpathrectangle{\pgfqpoint{5.636750in}{0.525000in}}{\pgfqpoint{1.293250in}{2.730000in}}%
\pgfusepath{clip}%
\pgfsetbuttcap%
\pgfsetroundjoin%
\definecolor{currentfill}{rgb}{0.101961,0.101961,0.101961}%
\pgfsetfillcolor{currentfill}%
\pgfsetlinewidth{0.301125pt}%
\definecolor{currentstroke}{rgb}{1.000000,1.000000,1.000000}%
\pgfsetstrokecolor{currentstroke}%
\pgfsetdash{}{0pt}%
\pgfsys@defobject{currentmarker}{\pgfqpoint{-0.013889in}{-0.013889in}}{\pgfqpoint{0.013889in}{0.013889in}}{%
\pgfpathmoveto{\pgfqpoint{0.000000in}{-0.013889in}}%
\pgfpathcurveto{\pgfqpoint{0.003683in}{-0.013889in}}{\pgfqpoint{0.007216in}{-0.012425in}}{\pgfqpoint{0.009821in}{-0.009821in}}%
\pgfpathcurveto{\pgfqpoint{0.012425in}{-0.007216in}}{\pgfqpoint{0.013889in}{-0.003683in}}{\pgfqpoint{0.013889in}{0.000000in}}%
\pgfpathcurveto{\pgfqpoint{0.013889in}{0.003683in}}{\pgfqpoint{0.012425in}{0.007216in}}{\pgfqpoint{0.009821in}{0.009821in}}%
\pgfpathcurveto{\pgfqpoint{0.007216in}{0.012425in}}{\pgfqpoint{0.003683in}{0.013889in}}{\pgfqpoint{0.000000in}{0.013889in}}%
\pgfpathcurveto{\pgfqpoint{-0.003683in}{0.013889in}}{\pgfqpoint{-0.007216in}{0.012425in}}{\pgfqpoint{-0.009821in}{0.009821in}}%
\pgfpathcurveto{\pgfqpoint{-0.012425in}{0.007216in}}{\pgfqpoint{-0.013889in}{0.003683in}}{\pgfqpoint{-0.013889in}{0.000000in}}%
\pgfpathcurveto{\pgfqpoint{-0.013889in}{-0.003683in}}{\pgfqpoint{-0.012425in}{-0.007216in}}{\pgfqpoint{-0.009821in}{-0.009821in}}%
\pgfpathcurveto{\pgfqpoint{-0.007216in}{-0.012425in}}{\pgfqpoint{-0.003683in}{-0.013889in}}{\pgfqpoint{0.000000in}{-0.013889in}}%
\pgfpathlineto{\pgfqpoint{0.000000in}{-0.013889in}}%
\pgfpathclose%
\pgfusepath{stroke,fill}%
}%
\begin{pgfscope}%
\pgfsys@transformshift{6.436212in}{2.985370in}%
\pgfsys@useobject{currentmarker}{}%
\end{pgfscope}%
\begin{pgfscope}%
\pgfsys@transformshift{6.370925in}{2.985370in}%
\pgfsys@useobject{currentmarker}{}%
\end{pgfscope}%
\begin{pgfscope}%
\pgfsys@transformshift{6.465201in}{2.985370in}%
\pgfsys@useobject{currentmarker}{}%
\end{pgfscope}%
\begin{pgfscope}%
\pgfsys@transformshift{6.479991in}{2.985370in}%
\pgfsys@useobject{currentmarker}{}%
\end{pgfscope}%
\end{pgfscope}%
\begin{pgfscope}%
\definecolor{textcolor}{rgb}{0.101961,0.101961,0.101961}%
\pgfsetstrokecolor{textcolor}%
\pgfsetfillcolor{textcolor}%
\pgftext[x=6.712734in,y=2.985370in,left,]{\color{textcolor}\sffamily\fontsize{6.000000}{7.200000}\selectfont 0.77}%
\end{pgfscope}%
\begin{pgfscope}%
\definecolor{textcolor}{rgb}{0.549020,0.549020,0.549020}%
\pgfsetstrokecolor{textcolor}%
\pgfsetfillcolor{textcolor}%
\pgftext[x=5.657442in,y=2.816852in,left,]{\color{textcolor}\sffamily\fontsize{6.000000}{7.200000}\selectfont n/a}%
\end{pgfscope}%
\begin{pgfscope}%
\pgfpathrectangle{\pgfqpoint{5.636750in}{0.525000in}}{\pgfqpoint{1.293250in}{2.730000in}}%
\pgfusepath{clip}%
\pgfsetbuttcap%
\pgfsetroundjoin%
\definecolor{currentfill}{rgb}{0.101961,0.101961,0.101961}%
\pgfsetfillcolor{currentfill}%
\pgfsetlinewidth{0.301125pt}%
\definecolor{currentstroke}{rgb}{1.000000,1.000000,1.000000}%
\pgfsetstrokecolor{currentstroke}%
\pgfsetdash{}{0pt}%
\pgfsys@defobject{currentmarker}{\pgfqpoint{-0.013889in}{-0.013889in}}{\pgfqpoint{0.013889in}{0.013889in}}{%
\pgfpathmoveto{\pgfqpoint{0.000000in}{-0.013889in}}%
\pgfpathcurveto{\pgfqpoint{0.003683in}{-0.013889in}}{\pgfqpoint{0.007216in}{-0.012425in}}{\pgfqpoint{0.009821in}{-0.009821in}}%
\pgfpathcurveto{\pgfqpoint{0.012425in}{-0.007216in}}{\pgfqpoint{0.013889in}{-0.003683in}}{\pgfqpoint{0.013889in}{0.000000in}}%
\pgfpathcurveto{\pgfqpoint{0.013889in}{0.003683in}}{\pgfqpoint{0.012425in}{0.007216in}}{\pgfqpoint{0.009821in}{0.009821in}}%
\pgfpathcurveto{\pgfqpoint{0.007216in}{0.012425in}}{\pgfqpoint{0.003683in}{0.013889in}}{\pgfqpoint{0.000000in}{0.013889in}}%
\pgfpathcurveto{\pgfqpoint{-0.003683in}{0.013889in}}{\pgfqpoint{-0.007216in}{0.012425in}}{\pgfqpoint{-0.009821in}{0.009821in}}%
\pgfpathcurveto{\pgfqpoint{-0.012425in}{0.007216in}}{\pgfqpoint{-0.013889in}{0.003683in}}{\pgfqpoint{-0.013889in}{0.000000in}}%
\pgfpathcurveto{\pgfqpoint{-0.013889in}{-0.003683in}}{\pgfqpoint{-0.012425in}{-0.007216in}}{\pgfqpoint{-0.009821in}{-0.009821in}}%
\pgfpathcurveto{\pgfqpoint{-0.007216in}{-0.012425in}}{\pgfqpoint{-0.003683in}{-0.013889in}}{\pgfqpoint{0.000000in}{-0.013889in}}%
\pgfpathlineto{\pgfqpoint{0.000000in}{-0.013889in}}%
\pgfpathclose%
\pgfusepath{stroke,fill}%
}%
\begin{pgfscope}%
\pgfsys@transformshift{6.606524in}{2.648333in}%
\pgfsys@useobject{currentmarker}{}%
\end{pgfscope}%
\begin{pgfscope}%
\pgfsys@transformshift{6.577570in}{2.648333in}%
\pgfsys@useobject{currentmarker}{}%
\end{pgfscope}%
\begin{pgfscope}%
\pgfsys@transformshift{6.649096in}{2.648333in}%
\pgfsys@useobject{currentmarker}{}%
\end{pgfscope}%
\begin{pgfscope}%
\pgfsys@transformshift{6.641667in}{2.648333in}%
\pgfsys@useobject{currentmarker}{}%
\end{pgfscope}%
\end{pgfscope}%
\begin{pgfscope}%
\definecolor{textcolor}{rgb}{0.101961,0.101961,0.101961}%
\pgfsetstrokecolor{textcolor}%
\pgfsetfillcolor{textcolor}%
\pgftext[x=6.712734in,y=2.648333in,left,]{\color{textcolor}\sffamily\fontsize{6.000000}{7.200000}\selectfont 0.95}%
\end{pgfscope}%
\begin{pgfscope}%
\pgfpathrectangle{\pgfqpoint{5.636750in}{0.525000in}}{\pgfqpoint{1.293250in}{2.730000in}}%
\pgfusepath{clip}%
\pgfsetbuttcap%
\pgfsetroundjoin%
\definecolor{currentfill}{rgb}{0.101961,0.101961,0.101961}%
\pgfsetfillcolor{currentfill}%
\pgfsetlinewidth{0.301125pt}%
\definecolor{currentstroke}{rgb}{1.000000,1.000000,1.000000}%
\pgfsetstrokecolor{currentstroke}%
\pgfsetdash{}{0pt}%
\pgfsys@defobject{currentmarker}{\pgfqpoint{-0.013889in}{-0.013889in}}{\pgfqpoint{0.013889in}{0.013889in}}{%
\pgfpathmoveto{\pgfqpoint{0.000000in}{-0.013889in}}%
\pgfpathcurveto{\pgfqpoint{0.003683in}{-0.013889in}}{\pgfqpoint{0.007216in}{-0.012425in}}{\pgfqpoint{0.009821in}{-0.009821in}}%
\pgfpathcurveto{\pgfqpoint{0.012425in}{-0.007216in}}{\pgfqpoint{0.013889in}{-0.003683in}}{\pgfqpoint{0.013889in}{0.000000in}}%
\pgfpathcurveto{\pgfqpoint{0.013889in}{0.003683in}}{\pgfqpoint{0.012425in}{0.007216in}}{\pgfqpoint{0.009821in}{0.009821in}}%
\pgfpathcurveto{\pgfqpoint{0.007216in}{0.012425in}}{\pgfqpoint{0.003683in}{0.013889in}}{\pgfqpoint{0.000000in}{0.013889in}}%
\pgfpathcurveto{\pgfqpoint{-0.003683in}{0.013889in}}{\pgfqpoint{-0.007216in}{0.012425in}}{\pgfqpoint{-0.009821in}{0.009821in}}%
\pgfpathcurveto{\pgfqpoint{-0.012425in}{0.007216in}}{\pgfqpoint{-0.013889in}{0.003683in}}{\pgfqpoint{-0.013889in}{0.000000in}}%
\pgfpathcurveto{\pgfqpoint{-0.013889in}{-0.003683in}}{\pgfqpoint{-0.012425in}{-0.007216in}}{\pgfqpoint{-0.009821in}{-0.009821in}}%
\pgfpathcurveto{\pgfqpoint{-0.007216in}{-0.012425in}}{\pgfqpoint{-0.003683in}{-0.013889in}}{\pgfqpoint{0.000000in}{-0.013889in}}%
\pgfpathlineto{\pgfqpoint{0.000000in}{-0.013889in}}%
\pgfpathclose%
\pgfusepath{stroke,fill}%
}%
\begin{pgfscope}%
\pgfsys@transformshift{6.219573in}{2.479815in}%
\pgfsys@useobject{currentmarker}{}%
\end{pgfscope}%
\begin{pgfscope}%
\pgfsys@transformshift{6.320519in}{2.479815in}%
\pgfsys@useobject{currentmarker}{}%
\end{pgfscope}%
\begin{pgfscope}%
\pgfsys@transformshift{6.259004in}{2.479815in}%
\pgfsys@useobject{currentmarker}{}%
\end{pgfscope}%
\begin{pgfscope}%
\pgfsys@transformshift{6.319287in}{2.479815in}%
\pgfsys@useobject{currentmarker}{}%
\end{pgfscope}%
\end{pgfscope}%
\begin{pgfscope}%
\definecolor{textcolor}{rgb}{0.101961,0.101961,0.101961}%
\pgfsetstrokecolor{textcolor}%
\pgfsetfillcolor{textcolor}%
\pgftext[x=6.712734in,y=2.479815in,left,]{\color{textcolor}\sffamily\fontsize{6.000000}{7.200000}\selectfont 0.62}%
\end{pgfscope}%
\begin{pgfscope}%
\pgfpathrectangle{\pgfqpoint{5.636750in}{0.525000in}}{\pgfqpoint{1.293250in}{2.730000in}}%
\pgfusepath{clip}%
\pgfsetbuttcap%
\pgfsetroundjoin%
\definecolor{currentfill}{rgb}{0.101961,0.101961,0.101961}%
\pgfsetfillcolor{currentfill}%
\pgfsetlinewidth{0.301125pt}%
\definecolor{currentstroke}{rgb}{1.000000,1.000000,1.000000}%
\pgfsetstrokecolor{currentstroke}%
\pgfsetdash{}{0pt}%
\pgfsys@defobject{currentmarker}{\pgfqpoint{-0.013889in}{-0.013889in}}{\pgfqpoint{0.013889in}{0.013889in}}{%
\pgfpathmoveto{\pgfqpoint{0.000000in}{-0.013889in}}%
\pgfpathcurveto{\pgfqpoint{0.003683in}{-0.013889in}}{\pgfqpoint{0.007216in}{-0.012425in}}{\pgfqpoint{0.009821in}{-0.009821in}}%
\pgfpathcurveto{\pgfqpoint{0.012425in}{-0.007216in}}{\pgfqpoint{0.013889in}{-0.003683in}}{\pgfqpoint{0.013889in}{0.000000in}}%
\pgfpathcurveto{\pgfqpoint{0.013889in}{0.003683in}}{\pgfqpoint{0.012425in}{0.007216in}}{\pgfqpoint{0.009821in}{0.009821in}}%
\pgfpathcurveto{\pgfqpoint{0.007216in}{0.012425in}}{\pgfqpoint{0.003683in}{0.013889in}}{\pgfqpoint{0.000000in}{0.013889in}}%
\pgfpathcurveto{\pgfqpoint{-0.003683in}{0.013889in}}{\pgfqpoint{-0.007216in}{0.012425in}}{\pgfqpoint{-0.009821in}{0.009821in}}%
\pgfpathcurveto{\pgfqpoint{-0.012425in}{0.007216in}}{\pgfqpoint{-0.013889in}{0.003683in}}{\pgfqpoint{-0.013889in}{0.000000in}}%
\pgfpathcurveto{\pgfqpoint{-0.013889in}{-0.003683in}}{\pgfqpoint{-0.012425in}{-0.007216in}}{\pgfqpoint{-0.009821in}{-0.009821in}}%
\pgfpathcurveto{\pgfqpoint{-0.007216in}{-0.012425in}}{\pgfqpoint{-0.003683in}{-0.013889in}}{\pgfqpoint{0.000000in}{-0.013889in}}%
\pgfpathlineto{\pgfqpoint{0.000000in}{-0.013889in}}%
\pgfpathclose%
\pgfusepath{stroke,fill}%
}%
\begin{pgfscope}%
\pgfsys@transformshift{5.706134in}{2.311296in}%
\pgfsys@useobject{currentmarker}{}%
\end{pgfscope}%
\begin{pgfscope}%
\pgfsys@transformshift{6.059230in}{2.311296in}%
\pgfsys@useobject{currentmarker}{}%
\end{pgfscope}%
\begin{pgfscope}%
\pgfsys@transformshift{6.419974in}{2.311296in}%
\pgfsys@useobject{currentmarker}{}%
\end{pgfscope}%
\end{pgfscope}%
\begin{pgfscope}%
\definecolor{textcolor}{rgb}{0.101961,0.101961,0.101961}%
\pgfsetstrokecolor{textcolor}%
\pgfsetfillcolor{textcolor}%
\pgftext[x=6.712734in,y=2.311296in,left,]{\color{textcolor}\sffamily\fontsize{6.000000}{7.200000}\selectfont 0.41}%
\end{pgfscope}%
\begin{pgfscope}%
\pgfpathrectangle{\pgfqpoint{5.636750in}{0.525000in}}{\pgfqpoint{1.293250in}{2.730000in}}%
\pgfusepath{clip}%
\pgfsetbuttcap%
\pgfsetroundjoin%
\definecolor{currentfill}{rgb}{0.101961,0.101961,0.101961}%
\pgfsetfillcolor{currentfill}%
\pgfsetlinewidth{0.301125pt}%
\definecolor{currentstroke}{rgb}{1.000000,1.000000,1.000000}%
\pgfsetstrokecolor{currentstroke}%
\pgfsetdash{}{0pt}%
\pgfsys@defobject{currentmarker}{\pgfqpoint{-0.013889in}{-0.013889in}}{\pgfqpoint{0.013889in}{0.013889in}}{%
\pgfpathmoveto{\pgfqpoint{0.000000in}{-0.013889in}}%
\pgfpathcurveto{\pgfqpoint{0.003683in}{-0.013889in}}{\pgfqpoint{0.007216in}{-0.012425in}}{\pgfqpoint{0.009821in}{-0.009821in}}%
\pgfpathcurveto{\pgfqpoint{0.012425in}{-0.007216in}}{\pgfqpoint{0.013889in}{-0.003683in}}{\pgfqpoint{0.013889in}{0.000000in}}%
\pgfpathcurveto{\pgfqpoint{0.013889in}{0.003683in}}{\pgfqpoint{0.012425in}{0.007216in}}{\pgfqpoint{0.009821in}{0.009821in}}%
\pgfpathcurveto{\pgfqpoint{0.007216in}{0.012425in}}{\pgfqpoint{0.003683in}{0.013889in}}{\pgfqpoint{0.000000in}{0.013889in}}%
\pgfpathcurveto{\pgfqpoint{-0.003683in}{0.013889in}}{\pgfqpoint{-0.007216in}{0.012425in}}{\pgfqpoint{-0.009821in}{0.009821in}}%
\pgfpathcurveto{\pgfqpoint{-0.012425in}{0.007216in}}{\pgfqpoint{-0.013889in}{0.003683in}}{\pgfqpoint{-0.013889in}{0.000000in}}%
\pgfpathcurveto{\pgfqpoint{-0.013889in}{-0.003683in}}{\pgfqpoint{-0.012425in}{-0.007216in}}{\pgfqpoint{-0.009821in}{-0.009821in}}%
\pgfpathcurveto{\pgfqpoint{-0.007216in}{-0.012425in}}{\pgfqpoint{-0.003683in}{-0.013889in}}{\pgfqpoint{0.000000in}{-0.013889in}}%
\pgfpathlineto{\pgfqpoint{0.000000in}{-0.013889in}}%
\pgfpathclose%
\pgfusepath{stroke,fill}%
}%
\begin{pgfscope}%
\pgfsys@transformshift{5.706966in}{2.142778in}%
\pgfsys@useobject{currentmarker}{}%
\end{pgfscope}%
\begin{pgfscope}%
\pgfsys@transformshift{5.693644in}{2.142778in}%
\pgfsys@useobject{currentmarker}{}%
\end{pgfscope}%
\begin{pgfscope}%
\pgfsys@transformshift{6.055073in}{2.142778in}%
\pgfsys@useobject{currentmarker}{}%
\end{pgfscope}%
\begin{pgfscope}%
\pgfsys@transformshift{6.442032in}{2.142778in}%
\pgfsys@useobject{currentmarker}{}%
\end{pgfscope}%
\end{pgfscope}%
\begin{pgfscope}%
\definecolor{textcolor}{rgb}{0.101961,0.101961,0.101961}%
\pgfsetstrokecolor{textcolor}%
\pgfsetfillcolor{textcolor}%
\pgftext[x=6.712734in,y=2.142778in,left,]{\color{textcolor}\sffamily\fontsize{6.000000}{7.200000}\selectfont 0.33}%
\end{pgfscope}%
\begin{pgfscope}%
\pgfpathrectangle{\pgfqpoint{5.636750in}{0.525000in}}{\pgfqpoint{1.293250in}{2.730000in}}%
\pgfusepath{clip}%
\pgfsetbuttcap%
\pgfsetroundjoin%
\definecolor{currentfill}{rgb}{0.101961,0.101961,0.101961}%
\pgfsetfillcolor{currentfill}%
\pgfsetlinewidth{0.301125pt}%
\definecolor{currentstroke}{rgb}{1.000000,1.000000,1.000000}%
\pgfsetstrokecolor{currentstroke}%
\pgfsetdash{}{0pt}%
\pgfsys@defobject{currentmarker}{\pgfqpoint{-0.013889in}{-0.013889in}}{\pgfqpoint{0.013889in}{0.013889in}}{%
\pgfpathmoveto{\pgfqpoint{0.000000in}{-0.013889in}}%
\pgfpathcurveto{\pgfqpoint{0.003683in}{-0.013889in}}{\pgfqpoint{0.007216in}{-0.012425in}}{\pgfqpoint{0.009821in}{-0.009821in}}%
\pgfpathcurveto{\pgfqpoint{0.012425in}{-0.007216in}}{\pgfqpoint{0.013889in}{-0.003683in}}{\pgfqpoint{0.013889in}{0.000000in}}%
\pgfpathcurveto{\pgfqpoint{0.013889in}{0.003683in}}{\pgfqpoint{0.012425in}{0.007216in}}{\pgfqpoint{0.009821in}{0.009821in}}%
\pgfpathcurveto{\pgfqpoint{0.007216in}{0.012425in}}{\pgfqpoint{0.003683in}{0.013889in}}{\pgfqpoint{0.000000in}{0.013889in}}%
\pgfpathcurveto{\pgfqpoint{-0.003683in}{0.013889in}}{\pgfqpoint{-0.007216in}{0.012425in}}{\pgfqpoint{-0.009821in}{0.009821in}}%
\pgfpathcurveto{\pgfqpoint{-0.012425in}{0.007216in}}{\pgfqpoint{-0.013889in}{0.003683in}}{\pgfqpoint{-0.013889in}{0.000000in}}%
\pgfpathcurveto{\pgfqpoint{-0.013889in}{-0.003683in}}{\pgfqpoint{-0.012425in}{-0.007216in}}{\pgfqpoint{-0.009821in}{-0.009821in}}%
\pgfpathcurveto{\pgfqpoint{-0.007216in}{-0.012425in}}{\pgfqpoint{-0.003683in}{-0.013889in}}{\pgfqpoint{0.000000in}{-0.013889in}}%
\pgfpathlineto{\pgfqpoint{0.000000in}{-0.013889in}}%
\pgfpathclose%
\pgfusepath{stroke,fill}%
}%
\begin{pgfscope}%
\pgfsys@transformshift{5.708055in}{1.974259in}%
\pgfsys@useobject{currentmarker}{}%
\end{pgfscope}%
\begin{pgfscope}%
\pgfsys@transformshift{5.693470in}{1.974259in}%
\pgfsys@useobject{currentmarker}{}%
\end{pgfscope}%
\begin{pgfscope}%
\pgfsys@transformshift{6.063026in}{1.974259in}%
\pgfsys@useobject{currentmarker}{}%
\end{pgfscope}%
\begin{pgfscope}%
\pgfsys@transformshift{6.425016in}{1.974259in}%
\pgfsys@useobject{currentmarker}{}%
\end{pgfscope}%
\end{pgfscope}%
\begin{pgfscope}%
\definecolor{textcolor}{rgb}{0.101961,0.101961,0.101961}%
\pgfsetstrokecolor{textcolor}%
\pgfsetfillcolor{textcolor}%
\pgftext[x=6.712734in,y=1.974259in,left,]{\color{textcolor}\sffamily\fontsize{6.000000}{7.200000}\selectfont 0.32}%
\end{pgfscope}%
\begin{pgfscope}%
\pgfpathrectangle{\pgfqpoint{5.636750in}{0.525000in}}{\pgfqpoint{1.293250in}{2.730000in}}%
\pgfusepath{clip}%
\pgfsetbuttcap%
\pgfsetroundjoin%
\definecolor{currentfill}{rgb}{0.101961,0.101961,0.101961}%
\pgfsetfillcolor{currentfill}%
\pgfsetlinewidth{0.301125pt}%
\definecolor{currentstroke}{rgb}{1.000000,1.000000,1.000000}%
\pgfsetstrokecolor{currentstroke}%
\pgfsetdash{}{0pt}%
\pgfsys@defobject{currentmarker}{\pgfqpoint{-0.013889in}{-0.013889in}}{\pgfqpoint{0.013889in}{0.013889in}}{%
\pgfpathmoveto{\pgfqpoint{0.000000in}{-0.013889in}}%
\pgfpathcurveto{\pgfqpoint{0.003683in}{-0.013889in}}{\pgfqpoint{0.007216in}{-0.012425in}}{\pgfqpoint{0.009821in}{-0.009821in}}%
\pgfpathcurveto{\pgfqpoint{0.012425in}{-0.007216in}}{\pgfqpoint{0.013889in}{-0.003683in}}{\pgfqpoint{0.013889in}{0.000000in}}%
\pgfpathcurveto{\pgfqpoint{0.013889in}{0.003683in}}{\pgfqpoint{0.012425in}{0.007216in}}{\pgfqpoint{0.009821in}{0.009821in}}%
\pgfpathcurveto{\pgfqpoint{0.007216in}{0.012425in}}{\pgfqpoint{0.003683in}{0.013889in}}{\pgfqpoint{0.000000in}{0.013889in}}%
\pgfpathcurveto{\pgfqpoint{-0.003683in}{0.013889in}}{\pgfqpoint{-0.007216in}{0.012425in}}{\pgfqpoint{-0.009821in}{0.009821in}}%
\pgfpathcurveto{\pgfqpoint{-0.012425in}{0.007216in}}{\pgfqpoint{-0.013889in}{0.003683in}}{\pgfqpoint{-0.013889in}{0.000000in}}%
\pgfpathcurveto{\pgfqpoint{-0.013889in}{-0.003683in}}{\pgfqpoint{-0.012425in}{-0.007216in}}{\pgfqpoint{-0.009821in}{-0.009821in}}%
\pgfpathcurveto{\pgfqpoint{-0.007216in}{-0.012425in}}{\pgfqpoint{-0.003683in}{-0.013889in}}{\pgfqpoint{0.000000in}{-0.013889in}}%
\pgfpathlineto{\pgfqpoint{0.000000in}{-0.013889in}}%
\pgfpathclose%
\pgfusepath{stroke,fill}%
}%
\begin{pgfscope}%
\pgfsys@transformshift{5.897136in}{1.805741in}%
\pgfsys@useobject{currentmarker}{}%
\end{pgfscope}%
\begin{pgfscope}%
\pgfsys@transformshift{5.902777in}{1.805741in}%
\pgfsys@useobject{currentmarker}{}%
\end{pgfscope}%
\begin{pgfscope}%
\pgfsys@transformshift{5.899935in}{1.805741in}%
\pgfsys@useobject{currentmarker}{}%
\end{pgfscope}%
\begin{pgfscope}%
\pgfsys@transformshift{5.973000in}{1.805741in}%
\pgfsys@useobject{currentmarker}{}%
\end{pgfscope}%
\end{pgfscope}%
\begin{pgfscope}%
\definecolor{textcolor}{rgb}{0.101961,0.101961,0.101961}%
\pgfsetstrokecolor{textcolor}%
\pgfsetfillcolor{textcolor}%
\pgftext[x=6.712734in,y=1.805741in,left,]{\color{textcolor}\sffamily\fontsize{6.000000}{7.200000}\selectfont 0.27}%
\end{pgfscope}%
\begin{pgfscope}%
\pgfpathrectangle{\pgfqpoint{5.636750in}{0.525000in}}{\pgfqpoint{1.293250in}{2.730000in}}%
\pgfusepath{clip}%
\pgfsetbuttcap%
\pgfsetroundjoin%
\definecolor{currentfill}{rgb}{0.101961,0.101961,0.101961}%
\pgfsetfillcolor{currentfill}%
\pgfsetlinewidth{0.301125pt}%
\definecolor{currentstroke}{rgb}{1.000000,1.000000,1.000000}%
\pgfsetstrokecolor{currentstroke}%
\pgfsetdash{}{0pt}%
\pgfsys@defobject{currentmarker}{\pgfqpoint{-0.013889in}{-0.013889in}}{\pgfqpoint{0.013889in}{0.013889in}}{%
\pgfpathmoveto{\pgfqpoint{0.000000in}{-0.013889in}}%
\pgfpathcurveto{\pgfqpoint{0.003683in}{-0.013889in}}{\pgfqpoint{0.007216in}{-0.012425in}}{\pgfqpoint{0.009821in}{-0.009821in}}%
\pgfpathcurveto{\pgfqpoint{0.012425in}{-0.007216in}}{\pgfqpoint{0.013889in}{-0.003683in}}{\pgfqpoint{0.013889in}{0.000000in}}%
\pgfpathcurveto{\pgfqpoint{0.013889in}{0.003683in}}{\pgfqpoint{0.012425in}{0.007216in}}{\pgfqpoint{0.009821in}{0.009821in}}%
\pgfpathcurveto{\pgfqpoint{0.007216in}{0.012425in}}{\pgfqpoint{0.003683in}{0.013889in}}{\pgfqpoint{0.000000in}{0.013889in}}%
\pgfpathcurveto{\pgfqpoint{-0.003683in}{0.013889in}}{\pgfqpoint{-0.007216in}{0.012425in}}{\pgfqpoint{-0.009821in}{0.009821in}}%
\pgfpathcurveto{\pgfqpoint{-0.012425in}{0.007216in}}{\pgfqpoint{-0.013889in}{0.003683in}}{\pgfqpoint{-0.013889in}{0.000000in}}%
\pgfpathcurveto{\pgfqpoint{-0.013889in}{-0.003683in}}{\pgfqpoint{-0.012425in}{-0.007216in}}{\pgfqpoint{-0.009821in}{-0.009821in}}%
\pgfpathcurveto{\pgfqpoint{-0.007216in}{-0.012425in}}{\pgfqpoint{-0.003683in}{-0.013889in}}{\pgfqpoint{0.000000in}{-0.013889in}}%
\pgfpathlineto{\pgfqpoint{0.000000in}{-0.013889in}}%
\pgfpathclose%
\pgfusepath{stroke,fill}%
}%
\begin{pgfscope}%
\pgfsys@transformshift{5.706469in}{1.637222in}%
\pgfsys@useobject{currentmarker}{}%
\end{pgfscope}%
\begin{pgfscope}%
\pgfsys@transformshift{5.686415in}{1.637222in}%
\pgfsys@useobject{currentmarker}{}%
\end{pgfscope}%
\begin{pgfscope}%
\pgfsys@transformshift{5.763478in}{1.637222in}%
\pgfsys@useobject{currentmarker}{}%
\end{pgfscope}%
\begin{pgfscope}%
\pgfsys@transformshift{5.719219in}{1.637222in}%
\pgfsys@useobject{currentmarker}{}%
\end{pgfscope}%
\end{pgfscope}%
\begin{pgfscope}%
\definecolor{textcolor}{rgb}{0.101961,0.101961,0.101961}%
\pgfsetstrokecolor{textcolor}%
\pgfsetfillcolor{textcolor}%
\pgftext[x=6.712734in,y=1.637222in,left,]{\color{textcolor}\sffamily\fontsize{6.000000}{7.200000}\selectfont 0.08}%
\end{pgfscope}%
\begin{pgfscope}%
\pgfpathrectangle{\pgfqpoint{5.636750in}{0.525000in}}{\pgfqpoint{1.293250in}{2.730000in}}%
\pgfusepath{clip}%
\pgfsetbuttcap%
\pgfsetroundjoin%
\definecolor{currentfill}{rgb}{0.101961,0.101961,0.101961}%
\pgfsetfillcolor{currentfill}%
\pgfsetlinewidth{0.301125pt}%
\definecolor{currentstroke}{rgb}{1.000000,1.000000,1.000000}%
\pgfsetstrokecolor{currentstroke}%
\pgfsetdash{}{0pt}%
\pgfsys@defobject{currentmarker}{\pgfqpoint{-0.013889in}{-0.013889in}}{\pgfqpoint{0.013889in}{0.013889in}}{%
\pgfpathmoveto{\pgfqpoint{0.000000in}{-0.013889in}}%
\pgfpathcurveto{\pgfqpoint{0.003683in}{-0.013889in}}{\pgfqpoint{0.007216in}{-0.012425in}}{\pgfqpoint{0.009821in}{-0.009821in}}%
\pgfpathcurveto{\pgfqpoint{0.012425in}{-0.007216in}}{\pgfqpoint{0.013889in}{-0.003683in}}{\pgfqpoint{0.013889in}{0.000000in}}%
\pgfpathcurveto{\pgfqpoint{0.013889in}{0.003683in}}{\pgfqpoint{0.012425in}{0.007216in}}{\pgfqpoint{0.009821in}{0.009821in}}%
\pgfpathcurveto{\pgfqpoint{0.007216in}{0.012425in}}{\pgfqpoint{0.003683in}{0.013889in}}{\pgfqpoint{0.000000in}{0.013889in}}%
\pgfpathcurveto{\pgfqpoint{-0.003683in}{0.013889in}}{\pgfqpoint{-0.007216in}{0.012425in}}{\pgfqpoint{-0.009821in}{0.009821in}}%
\pgfpathcurveto{\pgfqpoint{-0.012425in}{0.007216in}}{\pgfqpoint{-0.013889in}{0.003683in}}{\pgfqpoint{-0.013889in}{0.000000in}}%
\pgfpathcurveto{\pgfqpoint{-0.013889in}{-0.003683in}}{\pgfqpoint{-0.012425in}{-0.007216in}}{\pgfqpoint{-0.009821in}{-0.009821in}}%
\pgfpathcurveto{\pgfqpoint{-0.007216in}{-0.012425in}}{\pgfqpoint{-0.003683in}{-0.013889in}}{\pgfqpoint{0.000000in}{-0.013889in}}%
\pgfpathlineto{\pgfqpoint{0.000000in}{-0.013889in}}%
\pgfpathclose%
\pgfusepath{stroke,fill}%
}%
\begin{pgfscope}%
\pgfsys@transformshift{5.708754in}{1.468704in}%
\pgfsys@useobject{currentmarker}{}%
\end{pgfscope}%
\begin{pgfscope}%
\pgfsys@transformshift{5.690888in}{1.468704in}%
\pgfsys@useobject{currentmarker}{}%
\end{pgfscope}%
\begin{pgfscope}%
\pgfsys@transformshift{5.757446in}{1.468704in}%
\pgfsys@useobject{currentmarker}{}%
\end{pgfscope}%
\begin{pgfscope}%
\pgfsys@transformshift{5.719190in}{1.468704in}%
\pgfsys@useobject{currentmarker}{}%
\end{pgfscope}%
\end{pgfscope}%
\begin{pgfscope}%
\definecolor{textcolor}{rgb}{0.101961,0.101961,0.101961}%
\pgfsetstrokecolor{textcolor}%
\pgfsetfillcolor{textcolor}%
\pgftext[x=6.712734in,y=1.468704in,left,]{\color{textcolor}\sffamily\fontsize{6.000000}{7.200000}\selectfont 0.08}%
\end{pgfscope}%
\begin{pgfscope}%
\pgfpathrectangle{\pgfqpoint{5.636750in}{0.525000in}}{\pgfqpoint{1.293250in}{2.730000in}}%
\pgfusepath{clip}%
\pgfsetbuttcap%
\pgfsetroundjoin%
\definecolor{currentfill}{rgb}{0.101961,0.101961,0.101961}%
\pgfsetfillcolor{currentfill}%
\pgfsetlinewidth{0.301125pt}%
\definecolor{currentstroke}{rgb}{1.000000,1.000000,1.000000}%
\pgfsetstrokecolor{currentstroke}%
\pgfsetdash{}{0pt}%
\pgfsys@defobject{currentmarker}{\pgfqpoint{-0.013889in}{-0.013889in}}{\pgfqpoint{0.013889in}{0.013889in}}{%
\pgfpathmoveto{\pgfqpoint{0.000000in}{-0.013889in}}%
\pgfpathcurveto{\pgfqpoint{0.003683in}{-0.013889in}}{\pgfqpoint{0.007216in}{-0.012425in}}{\pgfqpoint{0.009821in}{-0.009821in}}%
\pgfpathcurveto{\pgfqpoint{0.012425in}{-0.007216in}}{\pgfqpoint{0.013889in}{-0.003683in}}{\pgfqpoint{0.013889in}{0.000000in}}%
\pgfpathcurveto{\pgfqpoint{0.013889in}{0.003683in}}{\pgfqpoint{0.012425in}{0.007216in}}{\pgfqpoint{0.009821in}{0.009821in}}%
\pgfpathcurveto{\pgfqpoint{0.007216in}{0.012425in}}{\pgfqpoint{0.003683in}{0.013889in}}{\pgfqpoint{0.000000in}{0.013889in}}%
\pgfpathcurveto{\pgfqpoint{-0.003683in}{0.013889in}}{\pgfqpoint{-0.007216in}{0.012425in}}{\pgfqpoint{-0.009821in}{0.009821in}}%
\pgfpathcurveto{\pgfqpoint{-0.012425in}{0.007216in}}{\pgfqpoint{-0.013889in}{0.003683in}}{\pgfqpoint{-0.013889in}{0.000000in}}%
\pgfpathcurveto{\pgfqpoint{-0.013889in}{-0.003683in}}{\pgfqpoint{-0.012425in}{-0.007216in}}{\pgfqpoint{-0.009821in}{-0.009821in}}%
\pgfpathcurveto{\pgfqpoint{-0.007216in}{-0.012425in}}{\pgfqpoint{-0.003683in}{-0.013889in}}{\pgfqpoint{0.000000in}{-0.013889in}}%
\pgfpathlineto{\pgfqpoint{0.000000in}{-0.013889in}}%
\pgfpathclose%
\pgfusepath{stroke,fill}%
}%
\begin{pgfscope}%
\pgfsys@transformshift{5.686217in}{1.300185in}%
\pgfsys@useobject{currentmarker}{}%
\end{pgfscope}%
\begin{pgfscope}%
\pgfsys@transformshift{5.680426in}{1.300185in}%
\pgfsys@useobject{currentmarker}{}%
\end{pgfscope}%
\begin{pgfscope}%
\pgfsys@transformshift{5.711399in}{1.300185in}%
\pgfsys@useobject{currentmarker}{}%
\end{pgfscope}%
\begin{pgfscope}%
\pgfsys@transformshift{5.695827in}{1.300185in}%
\pgfsys@useobject{currentmarker}{}%
\end{pgfscope}%
\end{pgfscope}%
\begin{pgfscope}%
\definecolor{textcolor}{rgb}{0.101961,0.101961,0.101961}%
\pgfsetstrokecolor{textcolor}%
\pgfsetfillcolor{textcolor}%
\pgftext[x=6.712734in,y=1.300185in,left,]{\color{textcolor}\sffamily\fontsize{6.000000}{7.200000}\selectfont 0.05}%
\end{pgfscope}%
\begin{pgfscope}%
\pgfpathrectangle{\pgfqpoint{5.636750in}{0.525000in}}{\pgfqpoint{1.293250in}{2.730000in}}%
\pgfusepath{clip}%
\pgfsetbuttcap%
\pgfsetroundjoin%
\definecolor{currentfill}{rgb}{0.101961,0.101961,0.101961}%
\pgfsetfillcolor{currentfill}%
\pgfsetlinewidth{0.301125pt}%
\definecolor{currentstroke}{rgb}{1.000000,1.000000,1.000000}%
\pgfsetstrokecolor{currentstroke}%
\pgfsetdash{}{0pt}%
\pgfsys@defobject{currentmarker}{\pgfqpoint{-0.013889in}{-0.013889in}}{\pgfqpoint{0.013889in}{0.013889in}}{%
\pgfpathmoveto{\pgfqpoint{0.000000in}{-0.013889in}}%
\pgfpathcurveto{\pgfqpoint{0.003683in}{-0.013889in}}{\pgfqpoint{0.007216in}{-0.012425in}}{\pgfqpoint{0.009821in}{-0.009821in}}%
\pgfpathcurveto{\pgfqpoint{0.012425in}{-0.007216in}}{\pgfqpoint{0.013889in}{-0.003683in}}{\pgfqpoint{0.013889in}{0.000000in}}%
\pgfpathcurveto{\pgfqpoint{0.013889in}{0.003683in}}{\pgfqpoint{0.012425in}{0.007216in}}{\pgfqpoint{0.009821in}{0.009821in}}%
\pgfpathcurveto{\pgfqpoint{0.007216in}{0.012425in}}{\pgfqpoint{0.003683in}{0.013889in}}{\pgfqpoint{0.000000in}{0.013889in}}%
\pgfpathcurveto{\pgfqpoint{-0.003683in}{0.013889in}}{\pgfqpoint{-0.007216in}{0.012425in}}{\pgfqpoint{-0.009821in}{0.009821in}}%
\pgfpathcurveto{\pgfqpoint{-0.012425in}{0.007216in}}{\pgfqpoint{-0.013889in}{0.003683in}}{\pgfqpoint{-0.013889in}{0.000000in}}%
\pgfpathcurveto{\pgfqpoint{-0.013889in}{-0.003683in}}{\pgfqpoint{-0.012425in}{-0.007216in}}{\pgfqpoint{-0.009821in}{-0.009821in}}%
\pgfpathcurveto{\pgfqpoint{-0.007216in}{-0.012425in}}{\pgfqpoint{-0.003683in}{-0.013889in}}{\pgfqpoint{0.000000in}{-0.013889in}}%
\pgfpathlineto{\pgfqpoint{0.000000in}{-0.013889in}}%
\pgfpathclose%
\pgfusepath{stroke,fill}%
}%
\begin{pgfscope}%
\pgfsys@transformshift{5.703498in}{1.131667in}%
\pgfsys@useobject{currentmarker}{}%
\end{pgfscope}%
\begin{pgfscope}%
\pgfsys@transformshift{5.683940in}{1.131667in}%
\pgfsys@useobject{currentmarker}{}%
\end{pgfscope}%
\begin{pgfscope}%
\pgfsys@transformshift{5.758169in}{1.131667in}%
\pgfsys@useobject{currentmarker}{}%
\end{pgfscope}%
\begin{pgfscope}%
\pgfsys@transformshift{5.720278in}{1.131667in}%
\pgfsys@useobject{currentmarker}{}%
\end{pgfscope}%
\end{pgfscope}%
\begin{pgfscope}%
\definecolor{textcolor}{rgb}{0.101961,0.101961,0.101961}%
\pgfsetstrokecolor{textcolor}%
\pgfsetfillcolor{textcolor}%
\pgftext[x=6.712734in,y=1.131667in,left,]{\color{textcolor}\sffamily\fontsize{6.000000}{7.200000}\selectfont 0.08}%
\end{pgfscope}%
\begin{pgfscope}%
\pgfpathrectangle{\pgfqpoint{5.636750in}{0.525000in}}{\pgfqpoint{1.293250in}{2.730000in}}%
\pgfusepath{clip}%
\pgfsetbuttcap%
\pgfsetroundjoin%
\definecolor{currentfill}{rgb}{0.101961,0.101961,0.101961}%
\pgfsetfillcolor{currentfill}%
\pgfsetlinewidth{0.301125pt}%
\definecolor{currentstroke}{rgb}{1.000000,1.000000,1.000000}%
\pgfsetstrokecolor{currentstroke}%
\pgfsetdash{}{0pt}%
\pgfsys@defobject{currentmarker}{\pgfqpoint{-0.013889in}{-0.013889in}}{\pgfqpoint{0.013889in}{0.013889in}}{%
\pgfpathmoveto{\pgfqpoint{0.000000in}{-0.013889in}}%
\pgfpathcurveto{\pgfqpoint{0.003683in}{-0.013889in}}{\pgfqpoint{0.007216in}{-0.012425in}}{\pgfqpoint{0.009821in}{-0.009821in}}%
\pgfpathcurveto{\pgfqpoint{0.012425in}{-0.007216in}}{\pgfqpoint{0.013889in}{-0.003683in}}{\pgfqpoint{0.013889in}{0.000000in}}%
\pgfpathcurveto{\pgfqpoint{0.013889in}{0.003683in}}{\pgfqpoint{0.012425in}{0.007216in}}{\pgfqpoint{0.009821in}{0.009821in}}%
\pgfpathcurveto{\pgfqpoint{0.007216in}{0.012425in}}{\pgfqpoint{0.003683in}{0.013889in}}{\pgfqpoint{0.000000in}{0.013889in}}%
\pgfpathcurveto{\pgfqpoint{-0.003683in}{0.013889in}}{\pgfqpoint{-0.007216in}{0.012425in}}{\pgfqpoint{-0.009821in}{0.009821in}}%
\pgfpathcurveto{\pgfqpoint{-0.012425in}{0.007216in}}{\pgfqpoint{-0.013889in}{0.003683in}}{\pgfqpoint{-0.013889in}{0.000000in}}%
\pgfpathcurveto{\pgfqpoint{-0.013889in}{-0.003683in}}{\pgfqpoint{-0.012425in}{-0.007216in}}{\pgfqpoint{-0.009821in}{-0.009821in}}%
\pgfpathcurveto{\pgfqpoint{-0.007216in}{-0.012425in}}{\pgfqpoint{-0.003683in}{-0.013889in}}{\pgfqpoint{0.000000in}{-0.013889in}}%
\pgfpathlineto{\pgfqpoint{0.000000in}{-0.013889in}}%
\pgfpathclose%
\pgfusepath{stroke,fill}%
}%
\begin{pgfscope}%
\pgfsys@transformshift{5.706245in}{0.963148in}%
\pgfsys@useobject{currentmarker}{}%
\end{pgfscope}%
\begin{pgfscope}%
\pgfsys@transformshift{5.688594in}{0.963148in}%
\pgfsys@useobject{currentmarker}{}%
\end{pgfscope}%
\begin{pgfscope}%
\pgfsys@transformshift{5.758749in}{0.963148in}%
\pgfsys@useobject{currentmarker}{}%
\end{pgfscope}%
\begin{pgfscope}%
\pgfsys@transformshift{5.717029in}{0.963148in}%
\pgfsys@useobject{currentmarker}{}%
\end{pgfscope}%
\end{pgfscope}%
\begin{pgfscope}%
\definecolor{textcolor}{rgb}{0.101961,0.101961,0.101961}%
\pgfsetstrokecolor{textcolor}%
\pgfsetfillcolor{textcolor}%
\pgftext[x=6.712734in,y=0.963148in,left,]{\color{textcolor}\sffamily\fontsize{6.000000}{7.200000}\selectfont 0.08}%
\end{pgfscope}%
\begin{pgfscope}%
\pgfpathrectangle{\pgfqpoint{5.636750in}{0.525000in}}{\pgfqpoint{1.293250in}{2.730000in}}%
\pgfusepath{clip}%
\pgfsetbuttcap%
\pgfsetroundjoin%
\definecolor{currentfill}{rgb}{0.101961,0.101961,0.101961}%
\pgfsetfillcolor{currentfill}%
\pgfsetlinewidth{0.301125pt}%
\definecolor{currentstroke}{rgb}{1.000000,1.000000,1.000000}%
\pgfsetstrokecolor{currentstroke}%
\pgfsetdash{}{0pt}%
\pgfsys@defobject{currentmarker}{\pgfqpoint{-0.013889in}{-0.013889in}}{\pgfqpoint{0.013889in}{0.013889in}}{%
\pgfpathmoveto{\pgfqpoint{0.000000in}{-0.013889in}}%
\pgfpathcurveto{\pgfqpoint{0.003683in}{-0.013889in}}{\pgfqpoint{0.007216in}{-0.012425in}}{\pgfqpoint{0.009821in}{-0.009821in}}%
\pgfpathcurveto{\pgfqpoint{0.012425in}{-0.007216in}}{\pgfqpoint{0.013889in}{-0.003683in}}{\pgfqpoint{0.013889in}{0.000000in}}%
\pgfpathcurveto{\pgfqpoint{0.013889in}{0.003683in}}{\pgfqpoint{0.012425in}{0.007216in}}{\pgfqpoint{0.009821in}{0.009821in}}%
\pgfpathcurveto{\pgfqpoint{0.007216in}{0.012425in}}{\pgfqpoint{0.003683in}{0.013889in}}{\pgfqpoint{0.000000in}{0.013889in}}%
\pgfpathcurveto{\pgfqpoint{-0.003683in}{0.013889in}}{\pgfqpoint{-0.007216in}{0.012425in}}{\pgfqpoint{-0.009821in}{0.009821in}}%
\pgfpathcurveto{\pgfqpoint{-0.012425in}{0.007216in}}{\pgfqpoint{-0.013889in}{0.003683in}}{\pgfqpoint{-0.013889in}{0.000000in}}%
\pgfpathcurveto{\pgfqpoint{-0.013889in}{-0.003683in}}{\pgfqpoint{-0.012425in}{-0.007216in}}{\pgfqpoint{-0.009821in}{-0.009821in}}%
\pgfpathcurveto{\pgfqpoint{-0.007216in}{-0.012425in}}{\pgfqpoint{-0.003683in}{-0.013889in}}{\pgfqpoint{0.000000in}{-0.013889in}}%
\pgfpathlineto{\pgfqpoint{0.000000in}{-0.013889in}}%
\pgfpathclose%
\pgfusepath{stroke,fill}%
}%
\begin{pgfscope}%
\pgfsys@transformshift{5.707309in}{0.794630in}%
\pgfsys@useobject{currentmarker}{}%
\end{pgfscope}%
\begin{pgfscope}%
\pgfsys@transformshift{5.689154in}{0.794630in}%
\pgfsys@useobject{currentmarker}{}%
\end{pgfscope}%
\begin{pgfscope}%
\pgfsys@transformshift{5.758660in}{0.794630in}%
\pgfsys@useobject{currentmarker}{}%
\end{pgfscope}%
\begin{pgfscope}%
\pgfsys@transformshift{5.719296in}{0.794630in}%
\pgfsys@useobject{currentmarker}{}%
\end{pgfscope}%
\end{pgfscope}%
\begin{pgfscope}%
\definecolor{textcolor}{rgb}{0.101961,0.101961,0.101961}%
\pgfsetstrokecolor{textcolor}%
\pgfsetfillcolor{textcolor}%
\pgftext[x=6.712734in,y=0.794630in,left,]{\color{textcolor}\sffamily\fontsize{6.000000}{7.200000}\selectfont 0.08}%
\end{pgfscope}%
\begin{pgfscope}%
\pgfpathrectangle{\pgfqpoint{5.636750in}{0.525000in}}{\pgfqpoint{1.293250in}{2.730000in}}%
\pgfusepath{clip}%
\pgfsetbuttcap%
\pgfsetroundjoin%
\definecolor{currentfill}{rgb}{0.101961,0.101961,0.101961}%
\pgfsetfillcolor{currentfill}%
\pgfsetlinewidth{0.301125pt}%
\definecolor{currentstroke}{rgb}{1.000000,1.000000,1.000000}%
\pgfsetstrokecolor{currentstroke}%
\pgfsetdash{}{0pt}%
\pgfsys@defobject{currentmarker}{\pgfqpoint{-0.013889in}{-0.013889in}}{\pgfqpoint{0.013889in}{0.013889in}}{%
\pgfpathmoveto{\pgfqpoint{0.000000in}{-0.013889in}}%
\pgfpathcurveto{\pgfqpoint{0.003683in}{-0.013889in}}{\pgfqpoint{0.007216in}{-0.012425in}}{\pgfqpoint{0.009821in}{-0.009821in}}%
\pgfpathcurveto{\pgfqpoint{0.012425in}{-0.007216in}}{\pgfqpoint{0.013889in}{-0.003683in}}{\pgfqpoint{0.013889in}{0.000000in}}%
\pgfpathcurveto{\pgfqpoint{0.013889in}{0.003683in}}{\pgfqpoint{0.012425in}{0.007216in}}{\pgfqpoint{0.009821in}{0.009821in}}%
\pgfpathcurveto{\pgfqpoint{0.007216in}{0.012425in}}{\pgfqpoint{0.003683in}{0.013889in}}{\pgfqpoint{0.000000in}{0.013889in}}%
\pgfpathcurveto{\pgfqpoint{-0.003683in}{0.013889in}}{\pgfqpoint{-0.007216in}{0.012425in}}{\pgfqpoint{-0.009821in}{0.009821in}}%
\pgfpathcurveto{\pgfqpoint{-0.012425in}{0.007216in}}{\pgfqpoint{-0.013889in}{0.003683in}}{\pgfqpoint{-0.013889in}{0.000000in}}%
\pgfpathcurveto{\pgfqpoint{-0.013889in}{-0.003683in}}{\pgfqpoint{-0.012425in}{-0.007216in}}{\pgfqpoint{-0.009821in}{-0.009821in}}%
\pgfpathcurveto{\pgfqpoint{-0.007216in}{-0.012425in}}{\pgfqpoint{-0.003683in}{-0.013889in}}{\pgfqpoint{0.000000in}{-0.013889in}}%
\pgfpathlineto{\pgfqpoint{0.000000in}{-0.013889in}}%
\pgfpathclose%
\pgfusepath{stroke,fill}%
}%
\begin{pgfscope}%
\pgfsys@transformshift{5.636750in}{0.626111in}%
\pgfsys@useobject{currentmarker}{}%
\end{pgfscope}%
\begin{pgfscope}%
\pgfsys@transformshift{5.636750in}{0.626111in}%
\pgfsys@useobject{currentmarker}{}%
\end{pgfscope}%
\begin{pgfscope}%
\pgfsys@transformshift{5.636750in}{0.626111in}%
\pgfsys@useobject{currentmarker}{}%
\end{pgfscope}%
\begin{pgfscope}%
\pgfsys@transformshift{5.636750in}{0.626111in}%
\pgfsys@useobject{currentmarker}{}%
\end{pgfscope}%
\end{pgfscope}%
\begin{pgfscope}%
\definecolor{textcolor}{rgb}{0.101961,0.101961,0.101961}%
\pgfsetstrokecolor{textcolor}%
\pgfsetfillcolor{textcolor}%
\pgftext[x=6.712734in,y=0.626111in,left,]{\color{textcolor}\sffamily\fontsize{6.000000}{7.200000}\selectfont 0.00}%
\end{pgfscope}%
\begin{pgfscope}%
\definecolor{textcolor}{rgb}{0.549020,0.549020,0.549020}%
\pgfsetstrokecolor{textcolor}%
\pgfsetfillcolor{textcolor}%
\pgftext[x=6.930000in,y=3.447500in,right,top]{\color{textcolor}\sffamily\fontsize{5.500000}{6.600000}\selectfont OLD DATA}%
\end{pgfscope}%
\begin{pgfscope}%
\pgfsetbuttcap%
\pgfsetmiterjoin%
\definecolor{currentfill}{rgb}{0.156863,0.356863,0.568627}%
\pgfsetfillcolor{currentfill}%
\pgfsetlinewidth{0.000000pt}%
\definecolor{currentstroke}{rgb}{0.000000,0.000000,0.000000}%
\pgfsetstrokecolor{currentstroke}%
\pgfsetstrokeopacity{0.000000}%
\pgfsetdash{}{0pt}%
\pgfpathmoveto{\pgfqpoint{2.114210in}{0.094306in}}%
\pgfpathlineto{\pgfqpoint{2.211433in}{0.094306in}}%
\pgfpathlineto{\pgfqpoint{2.211433in}{0.162361in}}%
\pgfpathlineto{\pgfqpoint{2.114210in}{0.162361in}}%
\pgfpathlineto{\pgfqpoint{2.114210in}{0.094306in}}%
\pgfpathclose%
\pgfusepath{fill}%
\end{pgfscope}%
\begin{pgfscope}%
\definecolor{textcolor}{rgb}{0.101961,0.101961,0.101961}%
\pgfsetstrokecolor{textcolor}%
\pgfsetfillcolor{textcolor}%
\pgftext[x=2.289210in,y=0.094306in,left,base]{\color{textcolor}\sffamily\fontsize{7.000000}{8.400000}\selectfont COVEX (old: rung 3)}%
\end{pgfscope}%
\begin{pgfscope}%
\pgfsetbuttcap%
\pgfsetmiterjoin%
\definecolor{currentfill}{rgb}{0.478431,0.478431,0.478431}%
\pgfsetfillcolor{currentfill}%
\pgfsetlinewidth{0.000000pt}%
\definecolor{currentstroke}{rgb}{0.000000,0.000000,0.000000}%
\pgfsetstrokecolor{currentstroke}%
\pgfsetstrokeopacity{0.000000}%
\pgfsetdash{}{0pt}%
\pgfpathmoveto{\pgfqpoint{3.372845in}{0.094306in}}%
\pgfpathlineto{\pgfqpoint{3.470067in}{0.094306in}}%
\pgfpathlineto{\pgfqpoint{3.470067in}{0.162361in}}%
\pgfpathlineto{\pgfqpoint{3.372845in}{0.162361in}}%
\pgfpathlineto{\pgfqpoint{3.372845in}{0.094306in}}%
\pgfpathclose%
\pgfusepath{fill}%
\end{pgfscope}%
\begin{pgfscope}%
\definecolor{textcolor}{rgb}{0.101961,0.101961,0.101961}%
\pgfsetstrokecolor{textcolor}%
\pgfsetfillcolor{textcolor}%
\pgftext[x=3.547845in,y=0.094306in,left,base]{\color{textcolor}\sffamily\fontsize{7.000000}{8.400000}\selectfont Stateful baseline}%
\end{pgfscope}%
\begin{pgfscope}%
\pgfsetbuttcap%
\pgfsetmiterjoin%
\definecolor{currentfill}{rgb}{0.784314,0.784314,0.784314}%
\pgfsetfillcolor{currentfill}%
\pgfsetlinewidth{0.000000pt}%
\definecolor{currentstroke}{rgb}{0.000000,0.000000,0.000000}%
\pgfsetstrokecolor{currentstroke}%
\pgfsetstrokeopacity{0.000000}%
\pgfsetdash{}{0pt}%
\pgfpathmoveto{\pgfqpoint{4.402193in}{0.094306in}}%
\pgfpathlineto{\pgfqpoint{4.499416in}{0.094306in}}%
\pgfpathlineto{\pgfqpoint{4.499416in}{0.162361in}}%
\pgfpathlineto{\pgfqpoint{4.402193in}{0.162361in}}%
\pgfpathlineto{\pgfqpoint{4.402193in}{0.094306in}}%
\pgfpathclose%
\pgfusepath{fill}%
\end{pgfscope}%
\begin{pgfscope}%
\definecolor{textcolor}{rgb}{0.101961,0.101961,0.101961}%
\pgfsetstrokecolor{textcolor}%
\pgfsetfillcolor{textcolor}%
\pgftext[x=4.577193in,y=0.094306in,left,base]{\color{textcolor}\sffamily\fontsize{7.000000}{8.400000}\selectfont Stateless baseline}%
\end{pgfscope}%
\begin{pgfscope}%
\pgfsetbuttcap%
\pgfsetmiterjoin%
\definecolor{currentfill}{rgb}{0.478431,0.305882,0.611765}%
\pgfsetfillcolor{currentfill}%
\pgfsetlinewidth{0.000000pt}%
\definecolor{currentstroke}{rgb}{0.000000,0.000000,0.000000}%
\pgfsetstrokecolor{currentstroke}%
\pgfsetstrokeopacity{0.000000}%
\pgfsetdash{}{0pt}%
\pgfpathmoveto{\pgfqpoint{5.471748in}{0.094306in}}%
\pgfpathlineto{\pgfqpoint{5.568970in}{0.094306in}}%
\pgfpathlineto{\pgfqpoint{5.568970in}{0.162361in}}%
\pgfpathlineto{\pgfqpoint{5.471748in}{0.162361in}}%
\pgfpathlineto{\pgfqpoint{5.471748in}{0.094306in}}%
\pgfpathclose%
\pgfusepath{fill}%
\end{pgfscope}%
\begin{pgfscope}%
\definecolor{textcolor}{rgb}{0.101961,0.101961,0.101961}%
\pgfsetstrokecolor{textcolor}%
\pgfsetfillcolor{textcolor}%
\pgftext[x=5.646748in,y=0.094306in,left,base]{\color{textcolor}\sffamily\fontsize{7.000000}{8.400000}\selectfont On-device only}%
\end{pgfscope}%
\end{pgfpicture}%
\makeatother%
\endgroup%
}
\caption{[PLACEHOLDER: OLD DATA, rung 3 stand-in for COVEX, baselines on dev seeds 4-7] Correct-by-deadline fraction on robot, single-shot and pooled (robot plus single-shot) questions, and per-episode on-time attainment on dashcam streams, for COVEX, the baselines and the on-device model (standard condition: 1$\times$ memory, 16 devices, 3\,s robot deadline). Bars pool seeds 4 to 7 (dashcam: mean over seeds); dots are single seeds. Confidence cascade has no dashcam cells; ECO-LLM lacks seed 5.}
\label{fig:F1_headline}
\end{figure*}
\clearpage
\begin{figure}[t]
\centering
{\fontfamily{LinuxBiolinumT}\selectfont
%% Creator: Matplotlib, PGF backend
%%
%% To include the figure in your LaTeX document, write
%%   \input{<filename>.pgf}
%%
%% Make sure the required packages are loaded in your preamble
%%   \usepackage{pgf}
%%
%% Also ensure that all the required font packages are loaded; for instance,
%% the lmodern package is sometimes necessary when using math font.
%%   \usepackage{lmodern}
%%
%% Figures using additional raster images can only be included by \input if
%% they are in the same directory as the main LaTeX file. For loading figures
%% from other directories you can use the `import` package
%%   \usepackage{import}
%%
%% and then include the figures with
%%   \import{<path to file>}{<filename>.pgf}
%%
%% Matplotlib used the following preamble
%%
\begingroup%
\makeatletter%
\begin{pgfpicture}%
\pgfpathrectangle{\pgfpointorigin}{\pgfqpoint{3.330000in}{3.000000in}}%
\pgfusepath{use as bounding box, clip}%
\begin{pgfscope}%
\pgfsetbuttcap%
\pgfsetmiterjoin%
\definecolor{currentfill}{rgb}{1.000000,1.000000,1.000000}%
\pgfsetfillcolor{currentfill}%
\pgfsetlinewidth{0.000000pt}%
\definecolor{currentstroke}{rgb}{1.000000,1.000000,1.000000}%
\pgfsetstrokecolor{currentstroke}%
\pgfsetdash{}{0pt}%
\pgfpathmoveto{\pgfqpoint{0.000000in}{0.000000in}}%
\pgfpathlineto{\pgfqpoint{3.330000in}{0.000000in}}%
\pgfpathlineto{\pgfqpoint{3.330000in}{3.000000in}}%
\pgfpathlineto{\pgfqpoint{0.000000in}{3.000000in}}%
\pgfpathlineto{\pgfqpoint{0.000000in}{0.000000in}}%
\pgfpathclose%
\pgfusepath{fill}%
\end{pgfscope}%
\begin{pgfscope}%
\pgfsetbuttcap%
\pgfsetmiterjoin%
\definecolor{currentfill}{rgb}{1.000000,1.000000,1.000000}%
\pgfsetfillcolor{currentfill}%
\pgfsetlinewidth{0.000000pt}%
\definecolor{currentstroke}{rgb}{0.000000,0.000000,0.000000}%
\pgfsetstrokecolor{currentstroke}%
\pgfsetstrokeopacity{0.000000}%
\pgfsetdash{}{0pt}%
\pgfpathmoveto{\pgfqpoint{0.499500in}{0.990000in}}%
\pgfpathlineto{\pgfqpoint{3.230100in}{0.990000in}}%
\pgfpathlineto{\pgfqpoint{3.230100in}{2.850000in}}%
\pgfpathlineto{\pgfqpoint{0.499500in}{2.850000in}}%
\pgfpathlineto{\pgfqpoint{0.499500in}{0.990000in}}%
\pgfpathclose%
\pgfusepath{fill}%
\end{pgfscope}%
\begin{pgfscope}%
\pgfpathrectangle{\pgfqpoint{0.499500in}{0.990000in}}{\pgfqpoint{2.730600in}{1.860000in}}%
\pgfusepath{clip}%
\pgfsetrectcap%
\pgfsetroundjoin%
\pgfsetlinewidth{0.401500pt}%
\definecolor{currentstroke}{rgb}{0.901961,0.901961,0.901961}%
\pgfsetstrokecolor{currentstroke}%
\pgfsetdash{}{0pt}%
\pgfpathmoveto{\pgfqpoint{0.576781in}{0.990000in}}%
\pgfpathlineto{\pgfqpoint{0.576781in}{2.850000in}}%
\pgfusepath{stroke}%
\end{pgfscope}%
\begin{pgfscope}%
\definecolor{textcolor}{rgb}{0.101961,0.101961,0.101961}%
\pgfsetstrokecolor{textcolor}%
\pgfsetfillcolor{textcolor}%
\pgftext[x=0.576781in,y=0.941389in,,top]{\color{textcolor}\sffamily\fontsize{7.000000}{8.400000}\selectfont 0.0}%
\end{pgfscope}%
\begin{pgfscope}%
\pgfpathrectangle{\pgfqpoint{0.499500in}{0.990000in}}{\pgfqpoint{2.730600in}{1.860000in}}%
\pgfusepath{clip}%
\pgfsetrectcap%
\pgfsetroundjoin%
\pgfsetlinewidth{0.401500pt}%
\definecolor{currentstroke}{rgb}{0.901961,0.901961,0.901961}%
\pgfsetstrokecolor{currentstroke}%
\pgfsetdash{}{0pt}%
\pgfpathmoveto{\pgfqpoint{1.091989in}{0.990000in}}%
\pgfpathlineto{\pgfqpoint{1.091989in}{2.850000in}}%
\pgfusepath{stroke}%
\end{pgfscope}%
\begin{pgfscope}%
\definecolor{textcolor}{rgb}{0.101961,0.101961,0.101961}%
\pgfsetstrokecolor{textcolor}%
\pgfsetfillcolor{textcolor}%
\pgftext[x=1.091989in,y=0.941389in,,top]{\color{textcolor}\sffamily\fontsize{7.000000}{8.400000}\selectfont 0.2}%
\end{pgfscope}%
\begin{pgfscope}%
\pgfpathrectangle{\pgfqpoint{0.499500in}{0.990000in}}{\pgfqpoint{2.730600in}{1.860000in}}%
\pgfusepath{clip}%
\pgfsetrectcap%
\pgfsetroundjoin%
\pgfsetlinewidth{0.401500pt}%
\definecolor{currentstroke}{rgb}{0.901961,0.901961,0.901961}%
\pgfsetstrokecolor{currentstroke}%
\pgfsetdash{}{0pt}%
\pgfpathmoveto{\pgfqpoint{1.607196in}{0.990000in}}%
\pgfpathlineto{\pgfqpoint{1.607196in}{2.850000in}}%
\pgfusepath{stroke}%
\end{pgfscope}%
\begin{pgfscope}%
\definecolor{textcolor}{rgb}{0.101961,0.101961,0.101961}%
\pgfsetstrokecolor{textcolor}%
\pgfsetfillcolor{textcolor}%
\pgftext[x=1.607196in,y=0.941389in,,top]{\color{textcolor}\sffamily\fontsize{7.000000}{8.400000}\selectfont 0.4}%
\end{pgfscope}%
\begin{pgfscope}%
\pgfpathrectangle{\pgfqpoint{0.499500in}{0.990000in}}{\pgfqpoint{2.730600in}{1.860000in}}%
\pgfusepath{clip}%
\pgfsetrectcap%
\pgfsetroundjoin%
\pgfsetlinewidth{0.401500pt}%
\definecolor{currentstroke}{rgb}{0.901961,0.901961,0.901961}%
\pgfsetstrokecolor{currentstroke}%
\pgfsetdash{}{0pt}%
\pgfpathmoveto{\pgfqpoint{2.122404in}{0.990000in}}%
\pgfpathlineto{\pgfqpoint{2.122404in}{2.850000in}}%
\pgfusepath{stroke}%
\end{pgfscope}%
\begin{pgfscope}%
\definecolor{textcolor}{rgb}{0.101961,0.101961,0.101961}%
\pgfsetstrokecolor{textcolor}%
\pgfsetfillcolor{textcolor}%
\pgftext[x=2.122404in,y=0.941389in,,top]{\color{textcolor}\sffamily\fontsize{7.000000}{8.400000}\selectfont 0.6}%
\end{pgfscope}%
\begin{pgfscope}%
\pgfpathrectangle{\pgfqpoint{0.499500in}{0.990000in}}{\pgfqpoint{2.730600in}{1.860000in}}%
\pgfusepath{clip}%
\pgfsetrectcap%
\pgfsetroundjoin%
\pgfsetlinewidth{0.401500pt}%
\definecolor{currentstroke}{rgb}{0.901961,0.901961,0.901961}%
\pgfsetstrokecolor{currentstroke}%
\pgfsetdash{}{0pt}%
\pgfpathmoveto{\pgfqpoint{2.637611in}{0.990000in}}%
\pgfpathlineto{\pgfqpoint{2.637611in}{2.850000in}}%
\pgfusepath{stroke}%
\end{pgfscope}%
\begin{pgfscope}%
\definecolor{textcolor}{rgb}{0.101961,0.101961,0.101961}%
\pgfsetstrokecolor{textcolor}%
\pgfsetfillcolor{textcolor}%
\pgftext[x=2.637611in,y=0.941389in,,top]{\color{textcolor}\sffamily\fontsize{7.000000}{8.400000}\selectfont 0.8}%
\end{pgfscope}%
\begin{pgfscope}%
\pgfpathrectangle{\pgfqpoint{0.499500in}{0.990000in}}{\pgfqpoint{2.730600in}{1.860000in}}%
\pgfusepath{clip}%
\pgfsetrectcap%
\pgfsetroundjoin%
\pgfsetlinewidth{0.401500pt}%
\definecolor{currentstroke}{rgb}{0.901961,0.901961,0.901961}%
\pgfsetstrokecolor{currentstroke}%
\pgfsetdash{}{0pt}%
\pgfpathmoveto{\pgfqpoint{3.152819in}{0.990000in}}%
\pgfpathlineto{\pgfqpoint{3.152819in}{2.850000in}}%
\pgfusepath{stroke}%
\end{pgfscope}%
\begin{pgfscope}%
\definecolor{textcolor}{rgb}{0.101961,0.101961,0.101961}%
\pgfsetstrokecolor{textcolor}%
\pgfsetfillcolor{textcolor}%
\pgftext[x=3.152819in,y=0.941389in,,top]{\color{textcolor}\sffamily\fontsize{7.000000}{8.400000}\selectfont 1.0}%
\end{pgfscope}%
\begin{pgfscope}%
\definecolor{textcolor}{rgb}{0.101961,0.101961,0.101961}%
\pgfsetstrokecolor{textcolor}%
\pgfsetfillcolor{textcolor}%
\pgftext[x=1.864800in,y=0.827192in,,top]{\color{textcolor}\sffamily\fontsize{7.500000}{9.000000}\selectfont On-time rate}%
\end{pgfscope}%
\begin{pgfscope}%
\pgfpathrectangle{\pgfqpoint{0.499500in}{0.990000in}}{\pgfqpoint{2.730600in}{1.860000in}}%
\pgfusepath{clip}%
\pgfsetrectcap%
\pgfsetroundjoin%
\pgfsetlinewidth{0.401500pt}%
\definecolor{currentstroke}{rgb}{0.901961,0.901961,0.901961}%
\pgfsetstrokecolor{currentstroke}%
\pgfsetdash{}{0pt}%
\pgfpathmoveto{\pgfqpoint{0.499500in}{1.053409in}}%
\pgfpathlineto{\pgfqpoint{3.230100in}{1.053409in}}%
\pgfusepath{stroke}%
\end{pgfscope}%
\begin{pgfscope}%
\definecolor{textcolor}{rgb}{0.101961,0.101961,0.101961}%
\pgfsetstrokecolor{textcolor}%
\pgfsetfillcolor{textcolor}%
\pgftext[x=0.318895in, y=1.019651in, left, base]{\color{textcolor}\sffamily\fontsize{7.000000}{8.400000}\selectfont 0.0}%
\end{pgfscope}%
\begin{pgfscope}%
\pgfpathrectangle{\pgfqpoint{0.499500in}{0.990000in}}{\pgfqpoint{2.730600in}{1.860000in}}%
\pgfusepath{clip}%
\pgfsetrectcap%
\pgfsetroundjoin%
\pgfsetlinewidth{0.401500pt}%
\definecolor{currentstroke}{rgb}{0.901961,0.901961,0.901961}%
\pgfsetstrokecolor{currentstroke}%
\pgfsetdash{}{0pt}%
\pgfpathmoveto{\pgfqpoint{0.499500in}{1.264773in}}%
\pgfpathlineto{\pgfqpoint{3.230100in}{1.264773in}}%
\pgfusepath{stroke}%
\end{pgfscope}%
\begin{pgfscope}%
\definecolor{textcolor}{rgb}{0.101961,0.101961,0.101961}%
\pgfsetstrokecolor{textcolor}%
\pgfsetfillcolor{textcolor}%
\pgftext[x=0.318895in, y=1.231015in, left, base]{\color{textcolor}\sffamily\fontsize{7.000000}{8.400000}\selectfont 0.1}%
\end{pgfscope}%
\begin{pgfscope}%
\pgfpathrectangle{\pgfqpoint{0.499500in}{0.990000in}}{\pgfqpoint{2.730600in}{1.860000in}}%
\pgfusepath{clip}%
\pgfsetrectcap%
\pgfsetroundjoin%
\pgfsetlinewidth{0.401500pt}%
\definecolor{currentstroke}{rgb}{0.901961,0.901961,0.901961}%
\pgfsetstrokecolor{currentstroke}%
\pgfsetdash{}{0pt}%
\pgfpathmoveto{\pgfqpoint{0.499500in}{1.476136in}}%
\pgfpathlineto{\pgfqpoint{3.230100in}{1.476136in}}%
\pgfusepath{stroke}%
\end{pgfscope}%
\begin{pgfscope}%
\definecolor{textcolor}{rgb}{0.101961,0.101961,0.101961}%
\pgfsetstrokecolor{textcolor}%
\pgfsetfillcolor{textcolor}%
\pgftext[x=0.318895in, y=1.442379in, left, base]{\color{textcolor}\sffamily\fontsize{7.000000}{8.400000}\selectfont 0.2}%
\end{pgfscope}%
\begin{pgfscope}%
\pgfpathrectangle{\pgfqpoint{0.499500in}{0.990000in}}{\pgfqpoint{2.730600in}{1.860000in}}%
\pgfusepath{clip}%
\pgfsetrectcap%
\pgfsetroundjoin%
\pgfsetlinewidth{0.401500pt}%
\definecolor{currentstroke}{rgb}{0.901961,0.901961,0.901961}%
\pgfsetstrokecolor{currentstroke}%
\pgfsetdash{}{0pt}%
\pgfpathmoveto{\pgfqpoint{0.499500in}{1.687500in}}%
\pgfpathlineto{\pgfqpoint{3.230100in}{1.687500in}}%
\pgfusepath{stroke}%
\end{pgfscope}%
\begin{pgfscope}%
\definecolor{textcolor}{rgb}{0.101961,0.101961,0.101961}%
\pgfsetstrokecolor{textcolor}%
\pgfsetfillcolor{textcolor}%
\pgftext[x=0.318895in, y=1.653742in, left, base]{\color{textcolor}\sffamily\fontsize{7.000000}{8.400000}\selectfont 0.3}%
\end{pgfscope}%
\begin{pgfscope}%
\pgfpathrectangle{\pgfqpoint{0.499500in}{0.990000in}}{\pgfqpoint{2.730600in}{1.860000in}}%
\pgfusepath{clip}%
\pgfsetrectcap%
\pgfsetroundjoin%
\pgfsetlinewidth{0.401500pt}%
\definecolor{currentstroke}{rgb}{0.901961,0.901961,0.901961}%
\pgfsetstrokecolor{currentstroke}%
\pgfsetdash{}{0pt}%
\pgfpathmoveto{\pgfqpoint{0.499500in}{1.898864in}}%
\pgfpathlineto{\pgfqpoint{3.230100in}{1.898864in}}%
\pgfusepath{stroke}%
\end{pgfscope}%
\begin{pgfscope}%
\definecolor{textcolor}{rgb}{0.101961,0.101961,0.101961}%
\pgfsetstrokecolor{textcolor}%
\pgfsetfillcolor{textcolor}%
\pgftext[x=0.318895in, y=1.865106in, left, base]{\color{textcolor}\sffamily\fontsize{7.000000}{8.400000}\selectfont 0.4}%
\end{pgfscope}%
\begin{pgfscope}%
\pgfpathrectangle{\pgfqpoint{0.499500in}{0.990000in}}{\pgfqpoint{2.730600in}{1.860000in}}%
\pgfusepath{clip}%
\pgfsetrectcap%
\pgfsetroundjoin%
\pgfsetlinewidth{0.401500pt}%
\definecolor{currentstroke}{rgb}{0.901961,0.901961,0.901961}%
\pgfsetstrokecolor{currentstroke}%
\pgfsetdash{}{0pt}%
\pgfpathmoveto{\pgfqpoint{0.499500in}{2.110227in}}%
\pgfpathlineto{\pgfqpoint{3.230100in}{2.110227in}}%
\pgfusepath{stroke}%
\end{pgfscope}%
\begin{pgfscope}%
\definecolor{textcolor}{rgb}{0.101961,0.101961,0.101961}%
\pgfsetstrokecolor{textcolor}%
\pgfsetfillcolor{textcolor}%
\pgftext[x=0.318895in, y=2.076470in, left, base]{\color{textcolor}\sffamily\fontsize{7.000000}{8.400000}\selectfont 0.5}%
\end{pgfscope}%
\begin{pgfscope}%
\pgfpathrectangle{\pgfqpoint{0.499500in}{0.990000in}}{\pgfqpoint{2.730600in}{1.860000in}}%
\pgfusepath{clip}%
\pgfsetrectcap%
\pgfsetroundjoin%
\pgfsetlinewidth{0.401500pt}%
\definecolor{currentstroke}{rgb}{0.901961,0.901961,0.901961}%
\pgfsetstrokecolor{currentstroke}%
\pgfsetdash{}{0pt}%
\pgfpathmoveto{\pgfqpoint{0.499500in}{2.321591in}}%
\pgfpathlineto{\pgfqpoint{3.230100in}{2.321591in}}%
\pgfusepath{stroke}%
\end{pgfscope}%
\begin{pgfscope}%
\definecolor{textcolor}{rgb}{0.101961,0.101961,0.101961}%
\pgfsetstrokecolor{textcolor}%
\pgfsetfillcolor{textcolor}%
\pgftext[x=0.318895in, y=2.287833in, left, base]{\color{textcolor}\sffamily\fontsize{7.000000}{8.400000}\selectfont 0.6}%
\end{pgfscope}%
\begin{pgfscope}%
\pgfpathrectangle{\pgfqpoint{0.499500in}{0.990000in}}{\pgfqpoint{2.730600in}{1.860000in}}%
\pgfusepath{clip}%
\pgfsetrectcap%
\pgfsetroundjoin%
\pgfsetlinewidth{0.401500pt}%
\definecolor{currentstroke}{rgb}{0.901961,0.901961,0.901961}%
\pgfsetstrokecolor{currentstroke}%
\pgfsetdash{}{0pt}%
\pgfpathmoveto{\pgfqpoint{0.499500in}{2.532955in}}%
\pgfpathlineto{\pgfqpoint{3.230100in}{2.532955in}}%
\pgfusepath{stroke}%
\end{pgfscope}%
\begin{pgfscope}%
\definecolor{textcolor}{rgb}{0.101961,0.101961,0.101961}%
\pgfsetstrokecolor{textcolor}%
\pgfsetfillcolor{textcolor}%
\pgftext[x=0.318895in, y=2.499197in, left, base]{\color{textcolor}\sffamily\fontsize{7.000000}{8.400000}\selectfont 0.7}%
\end{pgfscope}%
\begin{pgfscope}%
\pgfpathrectangle{\pgfqpoint{0.499500in}{0.990000in}}{\pgfqpoint{2.730600in}{1.860000in}}%
\pgfusepath{clip}%
\pgfsetrectcap%
\pgfsetroundjoin%
\pgfsetlinewidth{0.401500pt}%
\definecolor{currentstroke}{rgb}{0.901961,0.901961,0.901961}%
\pgfsetstrokecolor{currentstroke}%
\pgfsetdash{}{0pt}%
\pgfpathmoveto{\pgfqpoint{0.499500in}{2.744318in}}%
\pgfpathlineto{\pgfqpoint{3.230100in}{2.744318in}}%
\pgfusepath{stroke}%
\end{pgfscope}%
\begin{pgfscope}%
\definecolor{textcolor}{rgb}{0.101961,0.101961,0.101961}%
\pgfsetstrokecolor{textcolor}%
\pgfsetfillcolor{textcolor}%
\pgftext[x=0.318895in, y=2.710561in, left, base]{\color{textcolor}\sffamily\fontsize{7.000000}{8.400000}\selectfont 0.8}%
\end{pgfscope}%
\begin{pgfscope}%
\definecolor{textcolor}{rgb}{0.101961,0.101961,0.101961}%
\pgfsetstrokecolor{textcolor}%
\pgfsetfillcolor{textcolor}%
\pgftext[x=0.291117in,y=1.920000in,,bottom,rotate=90.000000]{\color{textcolor}\sffamily\fontsize{7.500000}{9.000000}\selectfont Correct by deadline}%
\end{pgfscope}%
\begin{pgfscope}%
\pgfpathrectangle{\pgfqpoint{0.499500in}{0.990000in}}{\pgfqpoint{2.730600in}{1.860000in}}%
\pgfusepath{clip}%
\pgfsetrectcap%
\pgfsetroundjoin%
\pgfsetlinewidth{0.602250pt}%
\definecolor{currentstroke}{rgb}{0.901961,0.901961,0.901961}%
\pgfsetstrokecolor{currentstroke}%
\pgfsetdash{}{0pt}%
\pgfpathmoveto{\pgfqpoint{0.576781in}{1.053409in}}%
\pgfpathlineto{\pgfqpoint{2.778601in}{2.860000in}}%
\pgfusepath{stroke}%
\end{pgfscope}%
\begin{pgfscope}%
\pgfsetrectcap%
\pgfsetmiterjoin%
\pgfsetlinewidth{0.803000pt}%
\definecolor{currentstroke}{rgb}{0.101961,0.101961,0.101961}%
\pgfsetstrokecolor{currentstroke}%
\pgfsetdash{}{0pt}%
\pgfpathmoveto{\pgfqpoint{0.499500in}{0.990000in}}%
\pgfpathlineto{\pgfqpoint{0.499500in}{2.850000in}}%
\pgfusepath{stroke}%
\end{pgfscope}%
\begin{pgfscope}%
\pgfsetrectcap%
\pgfsetmiterjoin%
\pgfsetlinewidth{0.803000pt}%
\definecolor{currentstroke}{rgb}{0.101961,0.101961,0.101961}%
\pgfsetstrokecolor{currentstroke}%
\pgfsetdash{}{0pt}%
\pgfpathmoveto{\pgfqpoint{0.499500in}{0.990000in}}%
\pgfpathlineto{\pgfqpoint{3.230100in}{0.990000in}}%
\pgfusepath{stroke}%
\end{pgfscope}%
\begin{pgfscope}%
\pgfpathrectangle{\pgfqpoint{0.499500in}{0.990000in}}{\pgfqpoint{2.730600in}{1.860000in}}%
\pgfusepath{clip}%
\pgfsetbuttcap%
\pgfsetroundjoin%
\definecolor{currentfill}{rgb}{0.784314,0.784314,0.784314}%
\pgfsetfillcolor{currentfill}%
\pgfsetlinewidth{0.301125pt}%
\definecolor{currentstroke}{rgb}{1.000000,1.000000,1.000000}%
\pgfsetstrokecolor{currentstroke}%
\pgfsetdash{}{0pt}%
\pgfsys@defobject{currentmarker}{\pgfqpoint{-0.027778in}{-0.027778in}}{\pgfqpoint{0.027778in}{0.027778in}}{%
\pgfpathmoveto{\pgfqpoint{0.000000in}{-0.027778in}}%
\pgfpathcurveto{\pgfqpoint{0.007367in}{-0.027778in}}{\pgfqpoint{0.014433in}{-0.024851in}}{\pgfqpoint{0.019642in}{-0.019642in}}%
\pgfpathcurveto{\pgfqpoint{0.024851in}{-0.014433in}}{\pgfqpoint{0.027778in}{-0.007367in}}{\pgfqpoint{0.027778in}{0.000000in}}%
\pgfpathcurveto{\pgfqpoint{0.027778in}{0.007367in}}{\pgfqpoint{0.024851in}{0.014433in}}{\pgfqpoint{0.019642in}{0.019642in}}%
\pgfpathcurveto{\pgfqpoint{0.014433in}{0.024851in}}{\pgfqpoint{0.007367in}{0.027778in}}{\pgfqpoint{0.000000in}{0.027778in}}%
\pgfpathcurveto{\pgfqpoint{-0.007367in}{0.027778in}}{\pgfqpoint{-0.014433in}{0.024851in}}{\pgfqpoint{-0.019642in}{0.019642in}}%
\pgfpathcurveto{\pgfqpoint{-0.024851in}{0.014433in}}{\pgfqpoint{-0.027778in}{0.007367in}}{\pgfqpoint{-0.027778in}{0.000000in}}%
\pgfpathcurveto{\pgfqpoint{-0.027778in}{-0.007367in}}{\pgfqpoint{-0.024851in}{-0.014433in}}{\pgfqpoint{-0.019642in}{-0.019642in}}%
\pgfpathcurveto{\pgfqpoint{-0.014433in}{-0.024851in}}{\pgfqpoint{-0.007367in}{-0.027778in}}{\pgfqpoint{0.000000in}{-0.027778in}}%
\pgfpathlineto{\pgfqpoint{0.000000in}{-0.027778in}}%
\pgfpathclose%
\pgfusepath{stroke,fill}%
}%
\begin{pgfscope}%
\pgfsys@transformshift{3.037346in}{1.521886in}%
\pgfsys@useobject{currentmarker}{}%
\end{pgfscope}%
\end{pgfscope}%
\begin{pgfscope}%
\pgfpathrectangle{\pgfqpoint{0.499500in}{0.990000in}}{\pgfqpoint{2.730600in}{1.860000in}}%
\pgfusepath{clip}%
\pgfsetbuttcap%
\pgfsetroundjoin%
\definecolor{currentfill}{rgb}{0.784314,0.784314,0.784314}%
\pgfsetfillcolor{currentfill}%
\pgfsetlinewidth{0.301125pt}%
\definecolor{currentstroke}{rgb}{1.000000,1.000000,1.000000}%
\pgfsetstrokecolor{currentstroke}%
\pgfsetdash{}{0pt}%
\pgfsys@defobject{currentmarker}{\pgfqpoint{-0.027778in}{-0.027778in}}{\pgfqpoint{0.027778in}{0.027778in}}{%
\pgfpathmoveto{\pgfqpoint{-0.027778in}{-0.027778in}}%
\pgfpathlineto{\pgfqpoint{0.027778in}{-0.027778in}}%
\pgfpathlineto{\pgfqpoint{0.027778in}{0.027778in}}%
\pgfpathlineto{\pgfqpoint{-0.027778in}{0.027778in}}%
\pgfpathlineto{\pgfqpoint{-0.027778in}{-0.027778in}}%
\pgfpathclose%
\pgfusepath{stroke,fill}%
}%
\begin{pgfscope}%
\pgfsys@transformshift{3.149206in}{2.322476in}%
\pgfsys@useobject{currentmarker}{}%
\end{pgfscope}%
\end{pgfscope}%
\begin{pgfscope}%
\pgfpathrectangle{\pgfqpoint{0.499500in}{0.990000in}}{\pgfqpoint{2.730600in}{1.860000in}}%
\pgfusepath{clip}%
\pgfsetbuttcap%
\pgfsetroundjoin%
\definecolor{currentfill}{rgb}{0.784314,0.784314,0.784314}%
\pgfsetfillcolor{currentfill}%
\pgfsetlinewidth{0.301125pt}%
\definecolor{currentstroke}{rgb}{1.000000,1.000000,1.000000}%
\pgfsetstrokecolor{currentstroke}%
\pgfsetdash{}{0pt}%
\pgfsys@defobject{currentmarker}{\pgfqpoint{-0.027778in}{-0.027778in}}{\pgfqpoint{0.027778in}{0.027778in}}{%
\pgfpathmoveto{\pgfqpoint{0.000000in}{0.027778in}}%
\pgfpathlineto{\pgfqpoint{-0.027778in}{-0.027778in}}%
\pgfpathlineto{\pgfqpoint{0.027778in}{-0.027778in}}%
\pgfpathlineto{\pgfqpoint{0.000000in}{0.027778in}}%
\pgfpathclose%
\pgfusepath{stroke,fill}%
}%
\begin{pgfscope}%
\pgfsys@transformshift{2.576013in}{2.690491in}%
\pgfsys@useobject{currentmarker}{}%
\end{pgfscope}%
\end{pgfscope}%
\begin{pgfscope}%
\pgfpathrectangle{\pgfqpoint{0.499500in}{0.990000in}}{\pgfqpoint{2.730600in}{1.860000in}}%
\pgfusepath{clip}%
\pgfsetbuttcap%
\pgfsetroundjoin%
\definecolor{currentfill}{rgb}{0.784314,0.784314,0.784314}%
\pgfsetfillcolor{currentfill}%
\pgfsetlinewidth{0.301125pt}%
\definecolor{currentstroke}{rgb}{1.000000,1.000000,1.000000}%
\pgfsetstrokecolor{currentstroke}%
\pgfsetdash{}{0pt}%
\pgfsys@defobject{currentmarker}{\pgfqpoint{-0.027778in}{-0.027778in}}{\pgfqpoint{0.027778in}{0.027778in}}{%
\pgfpathmoveto{\pgfqpoint{0.000000in}{-0.027778in}}%
\pgfpathcurveto{\pgfqpoint{0.007367in}{-0.027778in}}{\pgfqpoint{0.014433in}{-0.024851in}}{\pgfqpoint{0.019642in}{-0.019642in}}%
\pgfpathcurveto{\pgfqpoint{0.024851in}{-0.014433in}}{\pgfqpoint{0.027778in}{-0.007367in}}{\pgfqpoint{0.027778in}{0.000000in}}%
\pgfpathcurveto{\pgfqpoint{0.027778in}{0.007367in}}{\pgfqpoint{0.024851in}{0.014433in}}{\pgfqpoint{0.019642in}{0.019642in}}%
\pgfpathcurveto{\pgfqpoint{0.014433in}{0.024851in}}{\pgfqpoint{0.007367in}{0.027778in}}{\pgfqpoint{0.000000in}{0.027778in}}%
\pgfpathcurveto{\pgfqpoint{-0.007367in}{0.027778in}}{\pgfqpoint{-0.014433in}{0.024851in}}{\pgfqpoint{-0.019642in}{0.019642in}}%
\pgfpathcurveto{\pgfqpoint{-0.024851in}{0.014433in}}{\pgfqpoint{-0.027778in}{0.007367in}}{\pgfqpoint{-0.027778in}{0.000000in}}%
\pgfpathcurveto{\pgfqpoint{-0.027778in}{-0.007367in}}{\pgfqpoint{-0.024851in}{-0.014433in}}{\pgfqpoint{-0.019642in}{-0.019642in}}%
\pgfpathcurveto{\pgfqpoint{-0.014433in}{-0.024851in}}{\pgfqpoint{-0.007367in}{-0.027778in}}{\pgfqpoint{0.000000in}{-0.027778in}}%
\pgfpathlineto{\pgfqpoint{0.000000in}{-0.027778in}}%
\pgfpathclose%
\pgfusepath{stroke,fill}%
}%
\begin{pgfscope}%
\pgfsys@transformshift{3.147748in}{1.494238in}%
\pgfsys@useobject{currentmarker}{}%
\end{pgfscope}%
\end{pgfscope}%
\begin{pgfscope}%
\pgfpathrectangle{\pgfqpoint{0.499500in}{0.990000in}}{\pgfqpoint{2.730600in}{1.860000in}}%
\pgfusepath{clip}%
\pgfsetbuttcap%
\pgfsetroundjoin%
\definecolor{currentfill}{rgb}{0.784314,0.784314,0.784314}%
\pgfsetfillcolor{currentfill}%
\pgfsetlinewidth{0.301125pt}%
\definecolor{currentstroke}{rgb}{1.000000,1.000000,1.000000}%
\pgfsetstrokecolor{currentstroke}%
\pgfsetdash{}{0pt}%
\pgfsys@defobject{currentmarker}{\pgfqpoint{-0.027778in}{-0.027778in}}{\pgfqpoint{0.027778in}{0.027778in}}{%
\pgfpathmoveto{\pgfqpoint{-0.027778in}{-0.027778in}}%
\pgfpathlineto{\pgfqpoint{0.027778in}{-0.027778in}}%
\pgfpathlineto{\pgfqpoint{0.027778in}{0.027778in}}%
\pgfpathlineto{\pgfqpoint{-0.027778in}{0.027778in}}%
\pgfpathlineto{\pgfqpoint{-0.027778in}{-0.027778in}}%
\pgfpathclose%
\pgfusepath{stroke,fill}%
}%
\begin{pgfscope}%
\pgfsys@transformshift{3.151015in}{2.322476in}%
\pgfsys@useobject{currentmarker}{}%
\end{pgfscope}%
\end{pgfscope}%
\begin{pgfscope}%
\pgfpathrectangle{\pgfqpoint{0.499500in}{0.990000in}}{\pgfqpoint{2.730600in}{1.860000in}}%
\pgfusepath{clip}%
\pgfsetbuttcap%
\pgfsetroundjoin%
\definecolor{currentfill}{rgb}{0.784314,0.784314,0.784314}%
\pgfsetfillcolor{currentfill}%
\pgfsetlinewidth{0.301125pt}%
\definecolor{currentstroke}{rgb}{1.000000,1.000000,1.000000}%
\pgfsetstrokecolor{currentstroke}%
\pgfsetdash{}{0pt}%
\pgfsys@defobject{currentmarker}{\pgfqpoint{-0.027778in}{-0.027778in}}{\pgfqpoint{0.027778in}{0.027778in}}{%
\pgfpathmoveto{\pgfqpoint{0.000000in}{0.027778in}}%
\pgfpathlineto{\pgfqpoint{-0.027778in}{-0.027778in}}%
\pgfpathlineto{\pgfqpoint{0.027778in}{-0.027778in}}%
\pgfpathlineto{\pgfqpoint{0.000000in}{0.027778in}}%
\pgfpathclose%
\pgfusepath{stroke,fill}%
}%
\begin{pgfscope}%
\pgfsys@transformshift{3.029933in}{3.059513in}%
\pgfsys@useobject{currentmarker}{}%
\end{pgfscope}%
\end{pgfscope}%
\begin{pgfscope}%
\pgfpathrectangle{\pgfqpoint{0.499500in}{0.990000in}}{\pgfqpoint{2.730600in}{1.860000in}}%
\pgfusepath{clip}%
\pgfsetbuttcap%
\pgfsetroundjoin%
\definecolor{currentfill}{rgb}{0.784314,0.784314,0.784314}%
\pgfsetfillcolor{currentfill}%
\pgfsetlinewidth{0.301125pt}%
\definecolor{currentstroke}{rgb}{1.000000,1.000000,1.000000}%
\pgfsetstrokecolor{currentstroke}%
\pgfsetdash{}{0pt}%
\pgfsys@defobject{currentmarker}{\pgfqpoint{-0.027778in}{-0.027778in}}{\pgfqpoint{0.027778in}{0.027778in}}{%
\pgfpathmoveto{\pgfqpoint{0.000000in}{-0.027778in}}%
\pgfpathcurveto{\pgfqpoint{0.007367in}{-0.027778in}}{\pgfqpoint{0.014433in}{-0.024851in}}{\pgfqpoint{0.019642in}{-0.019642in}}%
\pgfpathcurveto{\pgfqpoint{0.024851in}{-0.014433in}}{\pgfqpoint{0.027778in}{-0.007367in}}{\pgfqpoint{0.027778in}{0.000000in}}%
\pgfpathcurveto{\pgfqpoint{0.027778in}{0.007367in}}{\pgfqpoint{0.024851in}{0.014433in}}{\pgfqpoint{0.019642in}{0.019642in}}%
\pgfpathcurveto{\pgfqpoint{0.014433in}{0.024851in}}{\pgfqpoint{0.007367in}{0.027778in}}{\pgfqpoint{0.000000in}{0.027778in}}%
\pgfpathcurveto{\pgfqpoint{-0.007367in}{0.027778in}}{\pgfqpoint{-0.014433in}{0.024851in}}{\pgfqpoint{-0.019642in}{0.019642in}}%
\pgfpathcurveto{\pgfqpoint{-0.024851in}{0.014433in}}{\pgfqpoint{-0.027778in}{0.007367in}}{\pgfqpoint{-0.027778in}{0.000000in}}%
\pgfpathcurveto{\pgfqpoint{-0.027778in}{-0.007367in}}{\pgfqpoint{-0.024851in}{-0.014433in}}{\pgfqpoint{-0.019642in}{-0.019642in}}%
\pgfpathcurveto{\pgfqpoint{-0.014433in}{-0.024851in}}{\pgfqpoint{-0.007367in}{-0.027778in}}{\pgfqpoint{0.000000in}{-0.027778in}}%
\pgfpathlineto{\pgfqpoint{0.000000in}{-0.027778in}}%
\pgfpathclose%
\pgfusepath{stroke,fill}%
}%
\begin{pgfscope}%
\pgfsys@transformshift{0.708723in}{1.638153in}%
\pgfsys@useobject{currentmarker}{}%
\end{pgfscope}%
\end{pgfscope}%
\begin{pgfscope}%
\pgfpathrectangle{\pgfqpoint{0.499500in}{0.990000in}}{\pgfqpoint{2.730600in}{1.860000in}}%
\pgfusepath{clip}%
\pgfsetbuttcap%
\pgfsetroundjoin%
\definecolor{currentfill}{rgb}{0.784314,0.784314,0.784314}%
\pgfsetfillcolor{currentfill}%
\pgfsetlinewidth{0.301125pt}%
\definecolor{currentstroke}{rgb}{1.000000,1.000000,1.000000}%
\pgfsetstrokecolor{currentstroke}%
\pgfsetdash{}{0pt}%
\pgfsys@defobject{currentmarker}{\pgfqpoint{-0.027778in}{-0.027778in}}{\pgfqpoint{0.027778in}{0.027778in}}{%
\pgfpathmoveto{\pgfqpoint{-0.027778in}{-0.027778in}}%
\pgfpathlineto{\pgfqpoint{0.027778in}{-0.027778in}}%
\pgfpathlineto{\pgfqpoint{0.027778in}{0.027778in}}%
\pgfpathlineto{\pgfqpoint{-0.027778in}{0.027778in}}%
\pgfpathlineto{\pgfqpoint{-0.027778in}{-0.027778in}}%
\pgfpathclose%
\pgfusepath{stroke,fill}%
}%
\begin{pgfscope}%
\pgfsys@transformshift{1.687079in}{2.238244in}%
\pgfsys@useobject{currentmarker}{}%
\end{pgfscope}%
\end{pgfscope}%
\begin{pgfscope}%
\pgfpathrectangle{\pgfqpoint{0.499500in}{0.990000in}}{\pgfqpoint{2.730600in}{1.860000in}}%
\pgfusepath{clip}%
\pgfsetbuttcap%
\pgfsetroundjoin%
\definecolor{currentfill}{rgb}{0.784314,0.784314,0.784314}%
\pgfsetfillcolor{currentfill}%
\pgfsetlinewidth{0.301125pt}%
\definecolor{currentstroke}{rgb}{1.000000,1.000000,1.000000}%
\pgfsetstrokecolor{currentstroke}%
\pgfsetdash{}{0pt}%
\pgfsys@defobject{currentmarker}{\pgfqpoint{-0.027778in}{-0.027778in}}{\pgfqpoint{0.027778in}{0.027778in}}{%
\pgfpathmoveto{\pgfqpoint{0.000000in}{0.027778in}}%
\pgfpathlineto{\pgfqpoint{-0.027778in}{-0.027778in}}%
\pgfpathlineto{\pgfqpoint{0.027778in}{-0.027778in}}%
\pgfpathlineto{\pgfqpoint{0.000000in}{0.027778in}}%
\pgfpathclose%
\pgfusepath{stroke,fill}%
}%
\begin{pgfscope}%
\pgfsys@transformshift{2.187318in}{2.366711in}%
\pgfsys@useobject{currentmarker}{}%
\end{pgfscope}%
\end{pgfscope}%
\begin{pgfscope}%
\pgfpathrectangle{\pgfqpoint{0.499500in}{0.990000in}}{\pgfqpoint{2.730600in}{1.860000in}}%
\pgfusepath{clip}%
\pgfsetbuttcap%
\pgfsetroundjoin%
\definecolor{currentfill}{rgb}{0.784314,0.784314,0.784314}%
\pgfsetfillcolor{currentfill}%
\pgfsetlinewidth{0.301125pt}%
\definecolor{currentstroke}{rgb}{1.000000,1.000000,1.000000}%
\pgfsetstrokecolor{currentstroke}%
\pgfsetdash{}{0pt}%
\pgfsys@defobject{currentmarker}{\pgfqpoint{-0.027778in}{-0.027778in}}{\pgfqpoint{0.027778in}{0.027778in}}{%
\pgfpathmoveto{\pgfqpoint{0.000000in}{-0.027778in}}%
\pgfpathcurveto{\pgfqpoint{0.007367in}{-0.027778in}}{\pgfqpoint{0.014433in}{-0.024851in}}{\pgfqpoint{0.019642in}{-0.019642in}}%
\pgfpathcurveto{\pgfqpoint{0.024851in}{-0.014433in}}{\pgfqpoint{0.027778in}{-0.007367in}}{\pgfqpoint{0.027778in}{0.000000in}}%
\pgfpathcurveto{\pgfqpoint{0.027778in}{0.007367in}}{\pgfqpoint{0.024851in}{0.014433in}}{\pgfqpoint{0.019642in}{0.019642in}}%
\pgfpathcurveto{\pgfqpoint{0.014433in}{0.024851in}}{\pgfqpoint{0.007367in}{0.027778in}}{\pgfqpoint{0.000000in}{0.027778in}}%
\pgfpathcurveto{\pgfqpoint{-0.007367in}{0.027778in}}{\pgfqpoint{-0.014433in}{0.024851in}}{\pgfqpoint{-0.019642in}{0.019642in}}%
\pgfpathcurveto{\pgfqpoint{-0.024851in}{0.014433in}}{\pgfqpoint{-0.027778in}{0.007367in}}{\pgfqpoint{-0.027778in}{0.000000in}}%
\pgfpathcurveto{\pgfqpoint{-0.027778in}{-0.007367in}}{\pgfqpoint{-0.024851in}{-0.014433in}}{\pgfqpoint{-0.019642in}{-0.019642in}}%
\pgfpathcurveto{\pgfqpoint{-0.014433in}{-0.024851in}}{\pgfqpoint{-0.007367in}{-0.027778in}}{\pgfqpoint{0.000000in}{-0.027778in}}%
\pgfpathlineto{\pgfqpoint{0.000000in}{-0.027778in}}%
\pgfpathclose%
\pgfusepath{stroke,fill}%
}%
\begin{pgfscope}%
\pgfsys@transformshift{1.483531in}{1.628649in}%
\pgfsys@useobject{currentmarker}{}%
\end{pgfscope}%
\end{pgfscope}%
\begin{pgfscope}%
\pgfpathrectangle{\pgfqpoint{0.499500in}{0.990000in}}{\pgfqpoint{2.730600in}{1.860000in}}%
\pgfusepath{clip}%
\pgfsetbuttcap%
\pgfsetroundjoin%
\definecolor{currentfill}{rgb}{0.784314,0.784314,0.784314}%
\pgfsetfillcolor{currentfill}%
\pgfsetlinewidth{0.301125pt}%
\definecolor{currentstroke}{rgb}{1.000000,1.000000,1.000000}%
\pgfsetstrokecolor{currentstroke}%
\pgfsetdash{}{0pt}%
\pgfsys@defobject{currentmarker}{\pgfqpoint{-0.027778in}{-0.027778in}}{\pgfqpoint{0.027778in}{0.027778in}}{%
\pgfpathmoveto{\pgfqpoint{-0.027778in}{-0.027778in}}%
\pgfpathlineto{\pgfqpoint{0.027778in}{-0.027778in}}%
\pgfpathlineto{\pgfqpoint{0.027778in}{0.027778in}}%
\pgfpathlineto{\pgfqpoint{-0.027778in}{0.027778in}}%
\pgfpathlineto{\pgfqpoint{-0.027778in}{-0.027778in}}%
\pgfpathclose%
\pgfusepath{stroke,fill}%
}%
\begin{pgfscope}%
\pgfsys@transformshift{1.794515in}{2.247040in}%
\pgfsys@useobject{currentmarker}{}%
\end{pgfscope}%
\end{pgfscope}%
\begin{pgfscope}%
\pgfpathrectangle{\pgfqpoint{0.499500in}{0.990000in}}{\pgfqpoint{2.730600in}{1.860000in}}%
\pgfusepath{clip}%
\pgfsetbuttcap%
\pgfsetroundjoin%
\definecolor{currentfill}{rgb}{0.784314,0.784314,0.784314}%
\pgfsetfillcolor{currentfill}%
\pgfsetlinewidth{0.301125pt}%
\definecolor{currentstroke}{rgb}{1.000000,1.000000,1.000000}%
\pgfsetstrokecolor{currentstroke}%
\pgfsetdash{}{0pt}%
\pgfsys@defobject{currentmarker}{\pgfqpoint{-0.027778in}{-0.027778in}}{\pgfqpoint{0.027778in}{0.027778in}}{%
\pgfpathmoveto{\pgfqpoint{0.000000in}{0.027778in}}%
\pgfpathlineto{\pgfqpoint{-0.027778in}{-0.027778in}}%
\pgfpathlineto{\pgfqpoint{0.027778in}{-0.027778in}}%
\pgfpathlineto{\pgfqpoint{0.000000in}{0.027778in}}%
\pgfpathclose%
\pgfusepath{stroke,fill}%
}%
\begin{pgfscope}%
\pgfsys@transformshift{1.647460in}{1.921723in}%
\pgfsys@useobject{currentmarker}{}%
\end{pgfscope}%
\end{pgfscope}%
\begin{pgfscope}%
\pgfpathrectangle{\pgfqpoint{0.499500in}{0.990000in}}{\pgfqpoint{2.730600in}{1.860000in}}%
\pgfusepath{clip}%
\pgfsetbuttcap%
\pgfsetroundjoin%
\definecolor{currentfill}{rgb}{0.478431,0.478431,0.478431}%
\pgfsetfillcolor{currentfill}%
\pgfsetlinewidth{0.301125pt}%
\definecolor{currentstroke}{rgb}{1.000000,1.000000,1.000000}%
\pgfsetstrokecolor{currentstroke}%
\pgfsetdash{}{0pt}%
\pgfsys@defobject{currentmarker}{\pgfqpoint{-0.027778in}{-0.027778in}}{\pgfqpoint{0.027778in}{0.027778in}}{%
\pgfpathmoveto{\pgfqpoint{0.000000in}{-0.027778in}}%
\pgfpathcurveto{\pgfqpoint{0.007367in}{-0.027778in}}{\pgfqpoint{0.014433in}{-0.024851in}}{\pgfqpoint{0.019642in}{-0.019642in}}%
\pgfpathcurveto{\pgfqpoint{0.024851in}{-0.014433in}}{\pgfqpoint{0.027778in}{-0.007367in}}{\pgfqpoint{0.027778in}{0.000000in}}%
\pgfpathcurveto{\pgfqpoint{0.027778in}{0.007367in}}{\pgfqpoint{0.024851in}{0.014433in}}{\pgfqpoint{0.019642in}{0.019642in}}%
\pgfpathcurveto{\pgfqpoint{0.014433in}{0.024851in}}{\pgfqpoint{0.007367in}{0.027778in}}{\pgfqpoint{0.000000in}{0.027778in}}%
\pgfpathcurveto{\pgfqpoint{-0.007367in}{0.027778in}}{\pgfqpoint{-0.014433in}{0.024851in}}{\pgfqpoint{-0.019642in}{0.019642in}}%
\pgfpathcurveto{\pgfqpoint{-0.024851in}{0.014433in}}{\pgfqpoint{-0.027778in}{0.007367in}}{\pgfqpoint{-0.027778in}{0.000000in}}%
\pgfpathcurveto{\pgfqpoint{-0.027778in}{-0.007367in}}{\pgfqpoint{-0.024851in}{-0.014433in}}{\pgfqpoint{-0.019642in}{-0.019642in}}%
\pgfpathcurveto{\pgfqpoint{-0.014433in}{-0.024851in}}{\pgfqpoint{-0.007367in}{-0.027778in}}{\pgfqpoint{0.000000in}{-0.027778in}}%
\pgfpathlineto{\pgfqpoint{0.000000in}{-0.027778in}}%
\pgfpathclose%
\pgfusepath{stroke,fill}%
}%
\begin{pgfscope}%
\pgfsys@transformshift{0.712854in}{1.636002in}%
\pgfsys@useobject{currentmarker}{}%
\end{pgfscope}%
\end{pgfscope}%
\begin{pgfscope}%
\pgfpathrectangle{\pgfqpoint{0.499500in}{0.990000in}}{\pgfqpoint{2.730600in}{1.860000in}}%
\pgfusepath{clip}%
\pgfsetbuttcap%
\pgfsetroundjoin%
\definecolor{currentfill}{rgb}{0.478431,0.478431,0.478431}%
\pgfsetfillcolor{currentfill}%
\pgfsetlinewidth{0.301125pt}%
\definecolor{currentstroke}{rgb}{1.000000,1.000000,1.000000}%
\pgfsetstrokecolor{currentstroke}%
\pgfsetdash{}{0pt}%
\pgfsys@defobject{currentmarker}{\pgfqpoint{-0.027778in}{-0.027778in}}{\pgfqpoint{0.027778in}{0.027778in}}{%
\pgfpathmoveto{\pgfqpoint{-0.027778in}{-0.027778in}}%
\pgfpathlineto{\pgfqpoint{0.027778in}{-0.027778in}}%
\pgfpathlineto{\pgfqpoint{0.027778in}{0.027778in}}%
\pgfpathlineto{\pgfqpoint{-0.027778in}{0.027778in}}%
\pgfpathlineto{\pgfqpoint{-0.027778in}{-0.027778in}}%
\pgfpathclose%
\pgfusepath{stroke,fill}%
}%
\begin{pgfscope}%
\pgfsys@transformshift{1.538925in}{2.233609in}%
\pgfsys@useobject{currentmarker}{}%
\end{pgfscope}%
\end{pgfscope}%
\begin{pgfscope}%
\pgfpathrectangle{\pgfqpoint{0.499500in}{0.990000in}}{\pgfqpoint{2.730600in}{1.860000in}}%
\pgfusepath{clip}%
\pgfsetbuttcap%
\pgfsetroundjoin%
\definecolor{currentfill}{rgb}{0.478431,0.478431,0.478431}%
\pgfsetfillcolor{currentfill}%
\pgfsetlinewidth{0.301125pt}%
\definecolor{currentstroke}{rgb}{1.000000,1.000000,1.000000}%
\pgfsetstrokecolor{currentstroke}%
\pgfsetdash{}{0pt}%
\pgfsys@defobject{currentmarker}{\pgfqpoint{-0.027778in}{-0.027778in}}{\pgfqpoint{0.027778in}{0.027778in}}{%
\pgfpathmoveto{\pgfqpoint{0.000000in}{0.027778in}}%
\pgfpathlineto{\pgfqpoint{-0.027778in}{-0.027778in}}%
\pgfpathlineto{\pgfqpoint{0.027778in}{-0.027778in}}%
\pgfpathlineto{\pgfqpoint{0.000000in}{0.027778in}}%
\pgfpathclose%
\pgfusepath{stroke,fill}%
}%
\begin{pgfscope}%
\pgfsys@transformshift{1.429085in}{1.743270in}%
\pgfsys@useobject{currentmarker}{}%
\end{pgfscope}%
\end{pgfscope}%
\begin{pgfscope}%
\pgfpathrectangle{\pgfqpoint{0.499500in}{0.990000in}}{\pgfqpoint{2.730600in}{1.860000in}}%
\pgfusepath{clip}%
\pgfsetbuttcap%
\pgfsetroundjoin%
\definecolor{currentfill}{rgb}{0.784314,0.784314,0.784314}%
\pgfsetfillcolor{currentfill}%
\pgfsetlinewidth{0.301125pt}%
\definecolor{currentstroke}{rgb}{1.000000,1.000000,1.000000}%
\pgfsetstrokecolor{currentstroke}%
\pgfsetdash{}{0pt}%
\pgfsys@defobject{currentmarker}{\pgfqpoint{-0.027778in}{-0.027778in}}{\pgfqpoint{0.027778in}{0.027778in}}{%
\pgfpathmoveto{\pgfqpoint{0.000000in}{-0.027778in}}%
\pgfpathcurveto{\pgfqpoint{0.007367in}{-0.027778in}}{\pgfqpoint{0.014433in}{-0.024851in}}{\pgfqpoint{0.019642in}{-0.019642in}}%
\pgfpathcurveto{\pgfqpoint{0.024851in}{-0.014433in}}{\pgfqpoint{0.027778in}{-0.007367in}}{\pgfqpoint{0.027778in}{0.000000in}}%
\pgfpathcurveto{\pgfqpoint{0.027778in}{0.007367in}}{\pgfqpoint{0.024851in}{0.014433in}}{\pgfqpoint{0.019642in}{0.019642in}}%
\pgfpathcurveto{\pgfqpoint{0.014433in}{0.024851in}}{\pgfqpoint{0.007367in}{0.027778in}}{\pgfqpoint{0.000000in}{0.027778in}}%
\pgfpathcurveto{\pgfqpoint{-0.007367in}{0.027778in}}{\pgfqpoint{-0.014433in}{0.024851in}}{\pgfqpoint{-0.019642in}{0.019642in}}%
\pgfpathcurveto{\pgfqpoint{-0.024851in}{0.014433in}}{\pgfqpoint{-0.027778in}{0.007367in}}{\pgfqpoint{-0.027778in}{0.000000in}}%
\pgfpathcurveto{\pgfqpoint{-0.027778in}{-0.007367in}}{\pgfqpoint{-0.024851in}{-0.014433in}}{\pgfqpoint{-0.019642in}{-0.019642in}}%
\pgfpathcurveto{\pgfqpoint{-0.014433in}{-0.024851in}}{\pgfqpoint{-0.007367in}{-0.027778in}}{\pgfqpoint{0.000000in}{-0.027778in}}%
\pgfpathlineto{\pgfqpoint{0.000000in}{-0.027778in}}%
\pgfpathclose%
\pgfusepath{stroke,fill}%
}%
\begin{pgfscope}%
\pgfsys@transformshift{0.712854in}{1.630970in}%
\pgfsys@useobject{currentmarker}{}%
\end{pgfscope}%
\end{pgfscope}%
\begin{pgfscope}%
\pgfpathrectangle{\pgfqpoint{0.499500in}{0.990000in}}{\pgfqpoint{2.730600in}{1.860000in}}%
\pgfusepath{clip}%
\pgfsetbuttcap%
\pgfsetroundjoin%
\definecolor{currentfill}{rgb}{0.784314,0.784314,0.784314}%
\pgfsetfillcolor{currentfill}%
\pgfsetlinewidth{0.301125pt}%
\definecolor{currentstroke}{rgb}{1.000000,1.000000,1.000000}%
\pgfsetstrokecolor{currentstroke}%
\pgfsetdash{}{0pt}%
\pgfsys@defobject{currentmarker}{\pgfqpoint{-0.027778in}{-0.027778in}}{\pgfqpoint{0.027778in}{0.027778in}}{%
\pgfpathmoveto{\pgfqpoint{-0.027778in}{-0.027778in}}%
\pgfpathlineto{\pgfqpoint{0.027778in}{-0.027778in}}%
\pgfpathlineto{\pgfqpoint{0.027778in}{0.027778in}}%
\pgfpathlineto{\pgfqpoint{-0.027778in}{0.027778in}}%
\pgfpathlineto{\pgfqpoint{-0.027778in}{-0.027778in}}%
\pgfpathclose%
\pgfusepath{stroke,fill}%
}%
\begin{pgfscope}%
\pgfsys@transformshift{1.522765in}{2.229168in}%
\pgfsys@useobject{currentmarker}{}%
\end{pgfscope}%
\end{pgfscope}%
\begin{pgfscope}%
\pgfpathrectangle{\pgfqpoint{0.499500in}{0.990000in}}{\pgfqpoint{2.730600in}{1.860000in}}%
\pgfusepath{clip}%
\pgfsetbuttcap%
\pgfsetroundjoin%
\definecolor{currentfill}{rgb}{0.784314,0.784314,0.784314}%
\pgfsetfillcolor{currentfill}%
\pgfsetlinewidth{0.301125pt}%
\definecolor{currentstroke}{rgb}{1.000000,1.000000,1.000000}%
\pgfsetstrokecolor{currentstroke}%
\pgfsetdash{}{0pt}%
\pgfsys@defobject{currentmarker}{\pgfqpoint{-0.027778in}{-0.027778in}}{\pgfqpoint{0.027778in}{0.027778in}}{%
\pgfpathmoveto{\pgfqpoint{0.000000in}{0.027778in}}%
\pgfpathlineto{\pgfqpoint{-0.027778in}{-0.027778in}}%
\pgfpathlineto{\pgfqpoint{0.027778in}{-0.027778in}}%
\pgfpathlineto{\pgfqpoint{0.000000in}{0.027778in}}%
\pgfpathclose%
\pgfusepath{stroke,fill}%
}%
\begin{pgfscope}%
\pgfsys@transformshift{1.424381in}{1.739108in}%
\pgfsys@useobject{currentmarker}{}%
\end{pgfscope}%
\end{pgfscope}%
\begin{pgfscope}%
\pgfpathrectangle{\pgfqpoint{0.499500in}{0.990000in}}{\pgfqpoint{2.730600in}{1.860000in}}%
\pgfusepath{clip}%
\pgfsetbuttcap%
\pgfsetroundjoin%
\definecolor{currentfill}{rgb}{0.784314,0.784314,0.784314}%
\pgfsetfillcolor{currentfill}%
\pgfsetlinewidth{0.301125pt}%
\definecolor{currentstroke}{rgb}{1.000000,1.000000,1.000000}%
\pgfsetstrokecolor{currentstroke}%
\pgfsetdash{}{0pt}%
\pgfsys@defobject{currentmarker}{\pgfqpoint{-0.027778in}{-0.027778in}}{\pgfqpoint{0.027778in}{0.027778in}}{%
\pgfpathmoveto{\pgfqpoint{0.000000in}{-0.027778in}}%
\pgfpathcurveto{\pgfqpoint{0.007367in}{-0.027778in}}{\pgfqpoint{0.014433in}{-0.024851in}}{\pgfqpoint{0.019642in}{-0.019642in}}%
\pgfpathcurveto{\pgfqpoint{0.024851in}{-0.014433in}}{\pgfqpoint{0.027778in}{-0.007367in}}{\pgfqpoint{0.027778in}{0.000000in}}%
\pgfpathcurveto{\pgfqpoint{0.027778in}{0.007367in}}{\pgfqpoint{0.024851in}{0.014433in}}{\pgfqpoint{0.019642in}{0.019642in}}%
\pgfpathcurveto{\pgfqpoint{0.014433in}{0.024851in}}{\pgfqpoint{0.007367in}{0.027778in}}{\pgfqpoint{0.000000in}{0.027778in}}%
\pgfpathcurveto{\pgfqpoint{-0.007367in}{0.027778in}}{\pgfqpoint{-0.014433in}{0.024851in}}{\pgfqpoint{-0.019642in}{0.019642in}}%
\pgfpathcurveto{\pgfqpoint{-0.024851in}{0.014433in}}{\pgfqpoint{-0.027778in}{0.007367in}}{\pgfqpoint{-0.027778in}{0.000000in}}%
\pgfpathcurveto{\pgfqpoint{-0.027778in}{-0.007367in}}{\pgfqpoint{-0.024851in}{-0.014433in}}{\pgfqpoint{-0.019642in}{-0.019642in}}%
\pgfpathcurveto{\pgfqpoint{-0.014433in}{-0.024851in}}{\pgfqpoint{-0.007367in}{-0.027778in}}{\pgfqpoint{0.000000in}{-0.027778in}}%
\pgfpathlineto{\pgfqpoint{0.000000in}{-0.027778in}}%
\pgfpathclose%
\pgfusepath{stroke,fill}%
}%
\begin{pgfscope}%
\pgfsys@transformshift{0.654205in}{1.619924in}%
\pgfsys@useobject{currentmarker}{}%
\end{pgfscope}%
\end{pgfscope}%
\begin{pgfscope}%
\pgfpathrectangle{\pgfqpoint{0.499500in}{0.990000in}}{\pgfqpoint{2.730600in}{1.860000in}}%
\pgfusepath{clip}%
\pgfsetbuttcap%
\pgfsetroundjoin%
\definecolor{currentfill}{rgb}{0.784314,0.784314,0.784314}%
\pgfsetfillcolor{currentfill}%
\pgfsetlinewidth{0.301125pt}%
\definecolor{currentstroke}{rgb}{1.000000,1.000000,1.000000}%
\pgfsetstrokecolor{currentstroke}%
\pgfsetdash{}{0pt}%
\pgfsys@defobject{currentmarker}{\pgfqpoint{-0.027778in}{-0.027778in}}{\pgfqpoint{0.027778in}{0.027778in}}{%
\pgfpathmoveto{\pgfqpoint{-0.027778in}{-0.027778in}}%
\pgfpathlineto{\pgfqpoint{0.027778in}{-0.027778in}}%
\pgfpathlineto{\pgfqpoint{0.027778in}{0.027778in}}%
\pgfpathlineto{\pgfqpoint{-0.027778in}{0.027778in}}%
\pgfpathlineto{\pgfqpoint{-0.027778in}{-0.027778in}}%
\pgfpathclose%
\pgfusepath{stroke,fill}%
}%
\begin{pgfscope}%
\pgfsys@transformshift{0.975078in}{2.196750in}%
\pgfsys@useobject{currentmarker}{}%
\end{pgfscope}%
\end{pgfscope}%
\begin{pgfscope}%
\pgfpathrectangle{\pgfqpoint{0.499500in}{0.990000in}}{\pgfqpoint{2.730600in}{1.860000in}}%
\pgfusepath{clip}%
\pgfsetbuttcap%
\pgfsetroundjoin%
\definecolor{currentfill}{rgb}{0.784314,0.784314,0.784314}%
\pgfsetfillcolor{currentfill}%
\pgfsetlinewidth{0.301125pt}%
\definecolor{currentstroke}{rgb}{1.000000,1.000000,1.000000}%
\pgfsetstrokecolor{currentstroke}%
\pgfsetdash{}{0pt}%
\pgfsys@defobject{currentmarker}{\pgfqpoint{-0.027778in}{-0.027778in}}{\pgfqpoint{0.027778in}{0.027778in}}{%
\pgfpathmoveto{\pgfqpoint{0.000000in}{0.027778in}}%
\pgfpathlineto{\pgfqpoint{-0.027778in}{-0.027778in}}%
\pgfpathlineto{\pgfqpoint{0.027778in}{-0.027778in}}%
\pgfpathlineto{\pgfqpoint{0.000000in}{0.027778in}}%
\pgfpathclose%
\pgfusepath{stroke,fill}%
}%
\begin{pgfscope}%
\pgfsys@transformshift{1.285723in}{1.628422in}%
\pgfsys@useobject{currentmarker}{}%
\end{pgfscope}%
\end{pgfscope}%
\begin{pgfscope}%
\pgfpathrectangle{\pgfqpoint{0.499500in}{0.990000in}}{\pgfqpoint{2.730600in}{1.860000in}}%
\pgfusepath{clip}%
\pgfsetbuttcap%
\pgfsetroundjoin%
\definecolor{currentfill}{rgb}{0.478431,0.478431,0.478431}%
\pgfsetfillcolor{currentfill}%
\pgfsetlinewidth{0.301125pt}%
\definecolor{currentstroke}{rgb}{1.000000,1.000000,1.000000}%
\pgfsetstrokecolor{currentstroke}%
\pgfsetdash{}{0pt}%
\pgfsys@defobject{currentmarker}{\pgfqpoint{-0.027778in}{-0.027778in}}{\pgfqpoint{0.027778in}{0.027778in}}{%
\pgfpathmoveto{\pgfqpoint{0.000000in}{-0.027778in}}%
\pgfpathcurveto{\pgfqpoint{0.007367in}{-0.027778in}}{\pgfqpoint{0.014433in}{-0.024851in}}{\pgfqpoint{0.019642in}{-0.019642in}}%
\pgfpathcurveto{\pgfqpoint{0.024851in}{-0.014433in}}{\pgfqpoint{0.027778in}{-0.007367in}}{\pgfqpoint{0.027778in}{0.000000in}}%
\pgfpathcurveto{\pgfqpoint{0.027778in}{0.007367in}}{\pgfqpoint{0.024851in}{0.014433in}}{\pgfqpoint{0.019642in}{0.019642in}}%
\pgfpathcurveto{\pgfqpoint{0.014433in}{0.024851in}}{\pgfqpoint{0.007367in}{0.027778in}}{\pgfqpoint{0.000000in}{0.027778in}}%
\pgfpathcurveto{\pgfqpoint{-0.007367in}{0.027778in}}{\pgfqpoint{-0.014433in}{0.024851in}}{\pgfqpoint{-0.019642in}{0.019642in}}%
\pgfpathcurveto{\pgfqpoint{-0.024851in}{0.014433in}}{\pgfqpoint{-0.027778in}{0.007367in}}{\pgfqpoint{-0.027778in}{0.000000in}}%
\pgfpathcurveto{\pgfqpoint{-0.027778in}{-0.007367in}}{\pgfqpoint{-0.024851in}{-0.014433in}}{\pgfqpoint{-0.019642in}{-0.019642in}}%
\pgfpathcurveto{\pgfqpoint{-0.014433in}{-0.024851in}}{\pgfqpoint{-0.007367in}{-0.027778in}}{\pgfqpoint{0.000000in}{-0.027778in}}%
\pgfpathlineto{\pgfqpoint{0.000000in}{-0.027778in}}%
\pgfpathclose%
\pgfusepath{stroke,fill}%
}%
\begin{pgfscope}%
\pgfsys@transformshift{0.657770in}{1.620033in}%
\pgfsys@useobject{currentmarker}{}%
\end{pgfscope}%
\end{pgfscope}%
\begin{pgfscope}%
\pgfpathrectangle{\pgfqpoint{0.499500in}{0.990000in}}{\pgfqpoint{2.730600in}{1.860000in}}%
\pgfusepath{clip}%
\pgfsetbuttcap%
\pgfsetroundjoin%
\definecolor{currentfill}{rgb}{0.478431,0.478431,0.478431}%
\pgfsetfillcolor{currentfill}%
\pgfsetlinewidth{0.301125pt}%
\definecolor{currentstroke}{rgb}{1.000000,1.000000,1.000000}%
\pgfsetstrokecolor{currentstroke}%
\pgfsetdash{}{0pt}%
\pgfsys@defobject{currentmarker}{\pgfqpoint{-0.027778in}{-0.027778in}}{\pgfqpoint{0.027778in}{0.027778in}}{%
\pgfpathmoveto{\pgfqpoint{-0.027778in}{-0.027778in}}%
\pgfpathlineto{\pgfqpoint{0.027778in}{-0.027778in}}%
\pgfpathlineto{\pgfqpoint{0.027778in}{0.027778in}}%
\pgfpathlineto{\pgfqpoint{-0.027778in}{0.027778in}}%
\pgfpathlineto{\pgfqpoint{-0.027778in}{-0.027778in}}%
\pgfpathclose%
\pgfusepath{stroke,fill}%
}%
\begin{pgfscope}%
\pgfsys@transformshift{0.829109in}{2.184929in}%
\pgfsys@useobject{currentmarker}{}%
\end{pgfscope}%
\end{pgfscope}%
\begin{pgfscope}%
\pgfpathrectangle{\pgfqpoint{0.499500in}{0.990000in}}{\pgfqpoint{2.730600in}{1.860000in}}%
\pgfusepath{clip}%
\pgfsetbuttcap%
\pgfsetroundjoin%
\definecolor{currentfill}{rgb}{0.478431,0.478431,0.478431}%
\pgfsetfillcolor{currentfill}%
\pgfsetlinewidth{0.301125pt}%
\definecolor{currentstroke}{rgb}{1.000000,1.000000,1.000000}%
\pgfsetstrokecolor{currentstroke}%
\pgfsetdash{}{0pt}%
\pgfsys@defobject{currentmarker}{\pgfqpoint{-0.027778in}{-0.027778in}}{\pgfqpoint{0.027778in}{0.027778in}}{%
\pgfpathmoveto{\pgfqpoint{0.000000in}{0.027778in}}%
\pgfpathlineto{\pgfqpoint{-0.027778in}{-0.027778in}}%
\pgfpathlineto{\pgfqpoint{0.027778in}{-0.027778in}}%
\pgfpathlineto{\pgfqpoint{0.000000in}{0.027778in}}%
\pgfpathclose%
\pgfusepath{stroke,fill}%
}%
\begin{pgfscope}%
\pgfsys@transformshift{0.786165in}{1.221228in}%
\pgfsys@useobject{currentmarker}{}%
\end{pgfscope}%
\end{pgfscope}%
\begin{pgfscope}%
\pgfpathrectangle{\pgfqpoint{0.499500in}{0.990000in}}{\pgfqpoint{2.730600in}{1.860000in}}%
\pgfusepath{clip}%
\pgfsetbuttcap%
\pgfsetroundjoin%
\definecolor{currentfill}{rgb}{0.784314,0.784314,0.784314}%
\pgfsetfillcolor{currentfill}%
\pgfsetlinewidth{0.301125pt}%
\definecolor{currentstroke}{rgb}{1.000000,1.000000,1.000000}%
\pgfsetstrokecolor{currentstroke}%
\pgfsetdash{}{0pt}%
\pgfsys@defobject{currentmarker}{\pgfqpoint{-0.027778in}{-0.027778in}}{\pgfqpoint{0.027778in}{0.027778in}}{%
\pgfpathmoveto{\pgfqpoint{0.000000in}{-0.027778in}}%
\pgfpathcurveto{\pgfqpoint{0.007367in}{-0.027778in}}{\pgfqpoint{0.014433in}{-0.024851in}}{\pgfqpoint{0.019642in}{-0.019642in}}%
\pgfpathcurveto{\pgfqpoint{0.024851in}{-0.014433in}}{\pgfqpoint{0.027778in}{-0.007367in}}{\pgfqpoint{0.027778in}{0.000000in}}%
\pgfpathcurveto{\pgfqpoint{0.027778in}{0.007367in}}{\pgfqpoint{0.024851in}{0.014433in}}{\pgfqpoint{0.019642in}{0.019642in}}%
\pgfpathcurveto{\pgfqpoint{0.014433in}{0.024851in}}{\pgfqpoint{0.007367in}{0.027778in}}{\pgfqpoint{0.000000in}{0.027778in}}%
\pgfpathcurveto{\pgfqpoint{-0.007367in}{0.027778in}}{\pgfqpoint{-0.014433in}{0.024851in}}{\pgfqpoint{-0.019642in}{0.019642in}}%
\pgfpathcurveto{\pgfqpoint{-0.024851in}{0.014433in}}{\pgfqpoint{-0.027778in}{0.007367in}}{\pgfqpoint{-0.027778in}{0.000000in}}%
\pgfpathcurveto{\pgfqpoint{-0.027778in}{-0.007367in}}{\pgfqpoint{-0.024851in}{-0.014433in}}{\pgfqpoint{-0.019642in}{-0.019642in}}%
\pgfpathcurveto{\pgfqpoint{-0.014433in}{-0.024851in}}{\pgfqpoint{-0.007367in}{-0.027778in}}{\pgfqpoint{0.000000in}{-0.027778in}}%
\pgfpathlineto{\pgfqpoint{0.000000in}{-0.027778in}}%
\pgfpathclose%
\pgfusepath{stroke,fill}%
}%
\begin{pgfscope}%
\pgfsys@transformshift{0.653702in}{1.615666in}%
\pgfsys@useobject{currentmarker}{}%
\end{pgfscope}%
\end{pgfscope}%
\begin{pgfscope}%
\pgfpathrectangle{\pgfqpoint{0.499500in}{0.990000in}}{\pgfqpoint{2.730600in}{1.860000in}}%
\pgfusepath{clip}%
\pgfsetbuttcap%
\pgfsetroundjoin%
\definecolor{currentfill}{rgb}{0.784314,0.784314,0.784314}%
\pgfsetfillcolor{currentfill}%
\pgfsetlinewidth{0.301125pt}%
\definecolor{currentstroke}{rgb}{1.000000,1.000000,1.000000}%
\pgfsetstrokecolor{currentstroke}%
\pgfsetdash{}{0pt}%
\pgfsys@defobject{currentmarker}{\pgfqpoint{-0.027778in}{-0.027778in}}{\pgfqpoint{0.027778in}{0.027778in}}{%
\pgfpathmoveto{\pgfqpoint{-0.027778in}{-0.027778in}}%
\pgfpathlineto{\pgfqpoint{0.027778in}{-0.027778in}}%
\pgfpathlineto{\pgfqpoint{0.027778in}{0.027778in}}%
\pgfpathlineto{\pgfqpoint{-0.027778in}{0.027778in}}%
\pgfpathlineto{\pgfqpoint{-0.027778in}{-0.027778in}}%
\pgfpathclose%
\pgfusepath{stroke,fill}%
}%
\begin{pgfscope}%
\pgfsys@transformshift{0.834547in}{2.184942in}%
\pgfsys@useobject{currentmarker}{}%
\end{pgfscope}%
\end{pgfscope}%
\begin{pgfscope}%
\pgfpathrectangle{\pgfqpoint{0.499500in}{0.990000in}}{\pgfqpoint{2.730600in}{1.860000in}}%
\pgfusepath{clip}%
\pgfsetbuttcap%
\pgfsetroundjoin%
\definecolor{currentfill}{rgb}{0.784314,0.784314,0.784314}%
\pgfsetfillcolor{currentfill}%
\pgfsetlinewidth{0.301125pt}%
\definecolor{currentstroke}{rgb}{1.000000,1.000000,1.000000}%
\pgfsetstrokecolor{currentstroke}%
\pgfsetdash{}{0pt}%
\pgfsys@defobject{currentmarker}{\pgfqpoint{-0.027778in}{-0.027778in}}{\pgfqpoint{0.027778in}{0.027778in}}{%
\pgfpathmoveto{\pgfqpoint{0.000000in}{0.027778in}}%
\pgfpathlineto{\pgfqpoint{-0.027778in}{-0.027778in}}%
\pgfpathlineto{\pgfqpoint{0.027778in}{-0.027778in}}%
\pgfpathlineto{\pgfqpoint{0.000000in}{0.027778in}}%
\pgfpathclose%
\pgfusepath{stroke,fill}%
}%
\begin{pgfscope}%
\pgfsys@transformshift{0.787592in}{1.221584in}%
\pgfsys@useobject{currentmarker}{}%
\end{pgfscope}%
\end{pgfscope}%
\begin{pgfscope}%
\pgfpathrectangle{\pgfqpoint{0.499500in}{0.990000in}}{\pgfqpoint{2.730600in}{1.860000in}}%
\pgfusepath{clip}%
\pgfsetbuttcap%
\pgfsetroundjoin%
\definecolor{currentfill}{rgb}{0.784314,0.784314,0.784314}%
\pgfsetfillcolor{currentfill}%
\pgfsetlinewidth{0.301125pt}%
\definecolor{currentstroke}{rgb}{1.000000,1.000000,1.000000}%
\pgfsetstrokecolor{currentstroke}%
\pgfsetdash{}{0pt}%
\pgfsys@defobject{currentmarker}{\pgfqpoint{-0.027778in}{-0.027778in}}{\pgfqpoint{0.027778in}{0.027778in}}{%
\pgfpathmoveto{\pgfqpoint{0.000000in}{-0.027778in}}%
\pgfpathcurveto{\pgfqpoint{0.007367in}{-0.027778in}}{\pgfqpoint{0.014433in}{-0.024851in}}{\pgfqpoint{0.019642in}{-0.019642in}}%
\pgfpathcurveto{\pgfqpoint{0.024851in}{-0.014433in}}{\pgfqpoint{0.027778in}{-0.007367in}}{\pgfqpoint{0.027778in}{0.000000in}}%
\pgfpathcurveto{\pgfqpoint{0.027778in}{0.007367in}}{\pgfqpoint{0.024851in}{0.014433in}}{\pgfqpoint{0.019642in}{0.019642in}}%
\pgfpathcurveto{\pgfqpoint{0.014433in}{0.024851in}}{\pgfqpoint{0.007367in}{0.027778in}}{\pgfqpoint{0.000000in}{0.027778in}}%
\pgfpathcurveto{\pgfqpoint{-0.007367in}{0.027778in}}{\pgfqpoint{-0.014433in}{0.024851in}}{\pgfqpoint{-0.019642in}{0.019642in}}%
\pgfpathcurveto{\pgfqpoint{-0.024851in}{0.014433in}}{\pgfqpoint{-0.027778in}{0.007367in}}{\pgfqpoint{-0.027778in}{0.000000in}}%
\pgfpathcurveto{\pgfqpoint{-0.027778in}{-0.007367in}}{\pgfqpoint{-0.024851in}{-0.014433in}}{\pgfqpoint{-0.019642in}{-0.019642in}}%
\pgfpathcurveto{\pgfqpoint{-0.014433in}{-0.024851in}}{\pgfqpoint{-0.007367in}{-0.027778in}}{\pgfqpoint{0.000000in}{-0.027778in}}%
\pgfpathlineto{\pgfqpoint{0.000000in}{-0.027778in}}%
\pgfpathclose%
\pgfusepath{stroke,fill}%
}%
\begin{pgfscope}%
\pgfsys@transformshift{0.632355in}{1.607345in}%
\pgfsys@useobject{currentmarker}{}%
\end{pgfscope}%
\end{pgfscope}%
\begin{pgfscope}%
\pgfpathrectangle{\pgfqpoint{0.499500in}{0.990000in}}{\pgfqpoint{2.730600in}{1.860000in}}%
\pgfusepath{clip}%
\pgfsetbuttcap%
\pgfsetroundjoin%
\definecolor{currentfill}{rgb}{0.784314,0.784314,0.784314}%
\pgfsetfillcolor{currentfill}%
\pgfsetlinewidth{0.301125pt}%
\definecolor{currentstroke}{rgb}{1.000000,1.000000,1.000000}%
\pgfsetstrokecolor{currentstroke}%
\pgfsetdash{}{0pt}%
\pgfsys@defobject{currentmarker}{\pgfqpoint{-0.027778in}{-0.027778in}}{\pgfqpoint{0.027778in}{0.027778in}}{%
\pgfpathmoveto{\pgfqpoint{-0.027778in}{-0.027778in}}%
\pgfpathlineto{\pgfqpoint{0.027778in}{-0.027778in}}%
\pgfpathlineto{\pgfqpoint{0.027778in}{0.027778in}}%
\pgfpathlineto{\pgfqpoint{-0.027778in}{0.027778in}}%
\pgfpathlineto{\pgfqpoint{-0.027778in}{-0.027778in}}%
\pgfpathclose%
\pgfusepath{stroke,fill}%
}%
\begin{pgfscope}%
\pgfsys@transformshift{0.809245in}{2.179005in}%
\pgfsys@useobject{currentmarker}{}%
\end{pgfscope}%
\end{pgfscope}%
\begin{pgfscope}%
\pgfpathrectangle{\pgfqpoint{0.499500in}{0.990000in}}{\pgfqpoint{2.730600in}{1.860000in}}%
\pgfusepath{clip}%
\pgfsetbuttcap%
\pgfsetroundjoin%
\definecolor{currentfill}{rgb}{0.784314,0.784314,0.784314}%
\pgfsetfillcolor{currentfill}%
\pgfsetlinewidth{0.301125pt}%
\definecolor{currentstroke}{rgb}{1.000000,1.000000,1.000000}%
\pgfsetstrokecolor{currentstroke}%
\pgfsetdash{}{0pt}%
\pgfsys@defobject{currentmarker}{\pgfqpoint{-0.027778in}{-0.027778in}}{\pgfqpoint{0.027778in}{0.027778in}}{%
\pgfpathmoveto{\pgfqpoint{0.000000in}{0.027778in}}%
\pgfpathlineto{\pgfqpoint{-0.027778in}{-0.027778in}}%
\pgfpathlineto{\pgfqpoint{0.027778in}{-0.027778in}}%
\pgfpathlineto{\pgfqpoint{0.000000in}{0.027778in}}%
\pgfpathclose%
\pgfusepath{stroke,fill}%
}%
\begin{pgfscope}%
\pgfsys@transformshift{0.780585in}{1.216276in}%
\pgfsys@useobject{currentmarker}{}%
\end{pgfscope}%
\end{pgfscope}%
\begin{pgfscope}%
\pgfpathrectangle{\pgfqpoint{0.499500in}{0.990000in}}{\pgfqpoint{2.730600in}{1.860000in}}%
\pgfusepath{clip}%
\pgfsetbuttcap%
\pgfsetroundjoin%
\definecolor{currentfill}{rgb}{0.478431,0.478431,0.478431}%
\pgfsetfillcolor{currentfill}%
\pgfsetlinewidth{0.301125pt}%
\definecolor{currentstroke}{rgb}{1.000000,1.000000,1.000000}%
\pgfsetstrokecolor{currentstroke}%
\pgfsetdash{}{0pt}%
\pgfsys@defobject{currentmarker}{\pgfqpoint{-0.027778in}{-0.027778in}}{\pgfqpoint{0.027778in}{0.027778in}}{%
\pgfpathmoveto{\pgfqpoint{0.000000in}{-0.027778in}}%
\pgfpathcurveto{\pgfqpoint{0.007367in}{-0.027778in}}{\pgfqpoint{0.014433in}{-0.024851in}}{\pgfqpoint{0.019642in}{-0.019642in}}%
\pgfpathcurveto{\pgfqpoint{0.024851in}{-0.014433in}}{\pgfqpoint{0.027778in}{-0.007367in}}{\pgfqpoint{0.027778in}{0.000000in}}%
\pgfpathcurveto{\pgfqpoint{0.027778in}{0.007367in}}{\pgfqpoint{0.024851in}{0.014433in}}{\pgfqpoint{0.019642in}{0.019642in}}%
\pgfpathcurveto{\pgfqpoint{0.014433in}{0.024851in}}{\pgfqpoint{0.007367in}{0.027778in}}{\pgfqpoint{0.000000in}{0.027778in}}%
\pgfpathcurveto{\pgfqpoint{-0.007367in}{0.027778in}}{\pgfqpoint{-0.014433in}{0.024851in}}{\pgfqpoint{-0.019642in}{0.019642in}}%
\pgfpathcurveto{\pgfqpoint{-0.024851in}{0.014433in}}{\pgfqpoint{-0.027778in}{0.007367in}}{\pgfqpoint{-0.027778in}{0.000000in}}%
\pgfpathcurveto{\pgfqpoint{-0.027778in}{-0.007367in}}{\pgfqpoint{-0.024851in}{-0.014433in}}{\pgfqpoint{-0.019642in}{-0.019642in}}%
\pgfpathcurveto{\pgfqpoint{-0.014433in}{-0.024851in}}{\pgfqpoint{-0.007367in}{-0.027778in}}{\pgfqpoint{0.000000in}{-0.027778in}}%
\pgfpathlineto{\pgfqpoint{0.000000in}{-0.027778in}}%
\pgfpathclose%
\pgfusepath{stroke,fill}%
}%
\begin{pgfscope}%
\pgfsys@transformshift{0.642689in}{1.615666in}%
\pgfsys@useobject{currentmarker}{}%
\end{pgfscope}%
\end{pgfscope}%
\begin{pgfscope}%
\pgfpathrectangle{\pgfqpoint{0.499500in}{0.990000in}}{\pgfqpoint{2.730600in}{1.860000in}}%
\pgfusepath{clip}%
\pgfsetbuttcap%
\pgfsetroundjoin%
\definecolor{currentfill}{rgb}{0.478431,0.478431,0.478431}%
\pgfsetfillcolor{currentfill}%
\pgfsetlinewidth{0.301125pt}%
\definecolor{currentstroke}{rgb}{1.000000,1.000000,1.000000}%
\pgfsetstrokecolor{currentstroke}%
\pgfsetdash{}{0pt}%
\pgfsys@defobject{currentmarker}{\pgfqpoint{-0.027778in}{-0.027778in}}{\pgfqpoint{0.027778in}{0.027778in}}{%
\pgfpathmoveto{\pgfqpoint{-0.027778in}{-0.027778in}}%
\pgfpathlineto{\pgfqpoint{0.027778in}{-0.027778in}}%
\pgfpathlineto{\pgfqpoint{0.027778in}{0.027778in}}%
\pgfpathlineto{\pgfqpoint{-0.027778in}{0.027778in}}%
\pgfpathlineto{\pgfqpoint{-0.027778in}{-0.027778in}}%
\pgfpathclose%
\pgfusepath{stroke,fill}%
}%
\begin{pgfscope}%
\pgfsys@transformshift{0.780427in}{2.176044in}%
\pgfsys@useobject{currentmarker}{}%
\end{pgfscope}%
\end{pgfscope}%
\begin{pgfscope}%
\pgfpathrectangle{\pgfqpoint{0.499500in}{0.990000in}}{\pgfqpoint{2.730600in}{1.860000in}}%
\pgfusepath{clip}%
\pgfsetbuttcap%
\pgfsetroundjoin%
\definecolor{currentfill}{rgb}{0.478431,0.478431,0.478431}%
\pgfsetfillcolor{currentfill}%
\pgfsetlinewidth{0.301125pt}%
\definecolor{currentstroke}{rgb}{1.000000,1.000000,1.000000}%
\pgfsetstrokecolor{currentstroke}%
\pgfsetdash{}{0pt}%
\pgfsys@defobject{currentmarker}{\pgfqpoint{-0.027778in}{-0.027778in}}{\pgfqpoint{0.027778in}{0.027778in}}{%
\pgfpathmoveto{\pgfqpoint{0.000000in}{0.027778in}}%
\pgfpathlineto{\pgfqpoint{-0.027778in}{-0.027778in}}%
\pgfpathlineto{\pgfqpoint{0.027778in}{-0.027778in}}%
\pgfpathlineto{\pgfqpoint{0.000000in}{0.027778in}}%
\pgfpathclose%
\pgfusepath{stroke,fill}%
}%
\begin{pgfscope}%
\pgfsys@transformshift{0.783379in}{1.218692in}%
\pgfsys@useobject{currentmarker}{}%
\end{pgfscope}%
\end{pgfscope}%
\begin{pgfscope}%
\pgfpathrectangle{\pgfqpoint{0.499500in}{0.990000in}}{\pgfqpoint{2.730600in}{1.860000in}}%
\pgfusepath{clip}%
\pgfsetbuttcap%
\pgfsetroundjoin%
\definecolor{currentfill}{rgb}{0.478431,0.478431,0.478431}%
\pgfsetfillcolor{currentfill}%
\pgfsetlinewidth{0.301125pt}%
\definecolor{currentstroke}{rgb}{1.000000,1.000000,1.000000}%
\pgfsetstrokecolor{currentstroke}%
\pgfsetdash{}{0pt}%
\pgfsys@defobject{currentmarker}{\pgfqpoint{-0.027778in}{-0.027778in}}{\pgfqpoint{0.027778in}{0.027778in}}{%
\pgfpathmoveto{\pgfqpoint{0.000000in}{-0.027778in}}%
\pgfpathcurveto{\pgfqpoint{0.007367in}{-0.027778in}}{\pgfqpoint{0.014433in}{-0.024851in}}{\pgfqpoint{0.019642in}{-0.019642in}}%
\pgfpathcurveto{\pgfqpoint{0.024851in}{-0.014433in}}{\pgfqpoint{0.027778in}{-0.007367in}}{\pgfqpoint{0.027778in}{0.000000in}}%
\pgfpathcurveto{\pgfqpoint{0.027778in}{0.007367in}}{\pgfqpoint{0.024851in}{0.014433in}}{\pgfqpoint{0.019642in}{0.019642in}}%
\pgfpathcurveto{\pgfqpoint{0.014433in}{0.024851in}}{\pgfqpoint{0.007367in}{0.027778in}}{\pgfqpoint{0.000000in}{0.027778in}}%
\pgfpathcurveto{\pgfqpoint{-0.007367in}{0.027778in}}{\pgfqpoint{-0.014433in}{0.024851in}}{\pgfqpoint{-0.019642in}{0.019642in}}%
\pgfpathcurveto{\pgfqpoint{-0.024851in}{0.014433in}}{\pgfqpoint{-0.027778in}{0.007367in}}{\pgfqpoint{-0.027778in}{0.000000in}}%
\pgfpathcurveto{\pgfqpoint{-0.027778in}{-0.007367in}}{\pgfqpoint{-0.024851in}{-0.014433in}}{\pgfqpoint{-0.019642in}{-0.019642in}}%
\pgfpathcurveto{\pgfqpoint{-0.014433in}{-0.024851in}}{\pgfqpoint{-0.007367in}{-0.027778in}}{\pgfqpoint{0.000000in}{-0.027778in}}%
\pgfpathlineto{\pgfqpoint{0.000000in}{-0.027778in}}%
\pgfpathclose%
\pgfusepath{stroke,fill}%
}%
\begin{pgfscope}%
\pgfsys@transformshift{0.627668in}{1.607345in}%
\pgfsys@useobject{currentmarker}{}%
\end{pgfscope}%
\end{pgfscope}%
\begin{pgfscope}%
\pgfpathrectangle{\pgfqpoint{0.499500in}{0.990000in}}{\pgfqpoint{2.730600in}{1.860000in}}%
\pgfusepath{clip}%
\pgfsetbuttcap%
\pgfsetroundjoin%
\definecolor{currentfill}{rgb}{0.478431,0.478431,0.478431}%
\pgfsetfillcolor{currentfill}%
\pgfsetlinewidth{0.301125pt}%
\definecolor{currentstroke}{rgb}{1.000000,1.000000,1.000000}%
\pgfsetstrokecolor{currentstroke}%
\pgfsetdash{}{0pt}%
\pgfsys@defobject{currentmarker}{\pgfqpoint{-0.027778in}{-0.027778in}}{\pgfqpoint{0.027778in}{0.027778in}}{%
\pgfpathmoveto{\pgfqpoint{-0.027778in}{-0.027778in}}%
\pgfpathlineto{\pgfqpoint{0.027778in}{-0.027778in}}%
\pgfpathlineto{\pgfqpoint{0.027778in}{0.027778in}}%
\pgfpathlineto{\pgfqpoint{-0.027778in}{0.027778in}}%
\pgfpathlineto{\pgfqpoint{-0.027778in}{-0.027778in}}%
\pgfpathclose%
\pgfusepath{stroke,fill}%
}%
\begin{pgfscope}%
\pgfsys@transformshift{0.789412in}{2.176053in}%
\pgfsys@useobject{currentmarker}{}%
\end{pgfscope}%
\end{pgfscope}%
\begin{pgfscope}%
\pgfpathrectangle{\pgfqpoint{0.499500in}{0.990000in}}{\pgfqpoint{2.730600in}{1.860000in}}%
\pgfusepath{clip}%
\pgfsetbuttcap%
\pgfsetroundjoin%
\definecolor{currentfill}{rgb}{0.478431,0.478431,0.478431}%
\pgfsetfillcolor{currentfill}%
\pgfsetlinewidth{0.301125pt}%
\definecolor{currentstroke}{rgb}{1.000000,1.000000,1.000000}%
\pgfsetstrokecolor{currentstroke}%
\pgfsetdash{}{0pt}%
\pgfsys@defobject{currentmarker}{\pgfqpoint{-0.027778in}{-0.027778in}}{\pgfqpoint{0.027778in}{0.027778in}}{%
\pgfpathmoveto{\pgfqpoint{0.000000in}{0.027778in}}%
\pgfpathlineto{\pgfqpoint{-0.027778in}{-0.027778in}}%
\pgfpathlineto{\pgfqpoint{0.027778in}{-0.027778in}}%
\pgfpathlineto{\pgfqpoint{0.000000in}{0.027778in}}%
\pgfpathclose%
\pgfusepath{stroke,fill}%
}%
\begin{pgfscope}%
\pgfsys@transformshift{0.786059in}{1.220634in}%
\pgfsys@useobject{currentmarker}{}%
\end{pgfscope}%
\end{pgfscope}%
\begin{pgfscope}%
\pgfpathrectangle{\pgfqpoint{0.499500in}{0.990000in}}{\pgfqpoint{2.730600in}{1.860000in}}%
\pgfusepath{clip}%
\pgfsetbuttcap%
\pgfsetroundjoin%
\definecolor{currentfill}{rgb}{0.478431,0.478431,0.478431}%
\pgfsetfillcolor{currentfill}%
\pgfsetlinewidth{0.301125pt}%
\definecolor{currentstroke}{rgb}{1.000000,1.000000,1.000000}%
\pgfsetstrokecolor{currentstroke}%
\pgfsetdash{}{0pt}%
\pgfsys@defobject{currentmarker}{\pgfqpoint{-0.027778in}{-0.027778in}}{\pgfqpoint{0.027778in}{0.027778in}}{%
\pgfpathmoveto{\pgfqpoint{0.000000in}{-0.027778in}}%
\pgfpathcurveto{\pgfqpoint{0.007367in}{-0.027778in}}{\pgfqpoint{0.014433in}{-0.024851in}}{\pgfqpoint{0.019642in}{-0.019642in}}%
\pgfpathcurveto{\pgfqpoint{0.024851in}{-0.014433in}}{\pgfqpoint{0.027778in}{-0.007367in}}{\pgfqpoint{0.027778in}{0.000000in}}%
\pgfpathcurveto{\pgfqpoint{0.027778in}{0.007367in}}{\pgfqpoint{0.024851in}{0.014433in}}{\pgfqpoint{0.019642in}{0.019642in}}%
\pgfpathcurveto{\pgfqpoint{0.014433in}{0.024851in}}{\pgfqpoint{0.007367in}{0.027778in}}{\pgfqpoint{0.000000in}{0.027778in}}%
\pgfpathcurveto{\pgfqpoint{-0.007367in}{0.027778in}}{\pgfqpoint{-0.014433in}{0.024851in}}{\pgfqpoint{-0.019642in}{0.019642in}}%
\pgfpathcurveto{\pgfqpoint{-0.024851in}{0.014433in}}{\pgfqpoint{-0.027778in}{0.007367in}}{\pgfqpoint{-0.027778in}{0.000000in}}%
\pgfpathcurveto{\pgfqpoint{-0.027778in}{-0.007367in}}{\pgfqpoint{-0.024851in}{-0.014433in}}{\pgfqpoint{-0.019642in}{-0.019642in}}%
\pgfpathcurveto{\pgfqpoint{-0.014433in}{-0.024851in}}{\pgfqpoint{-0.007367in}{-0.027778in}}{\pgfqpoint{0.000000in}{-0.027778in}}%
\pgfpathlineto{\pgfqpoint{0.000000in}{-0.027778in}}%
\pgfpathclose%
\pgfusepath{stroke,fill}%
}%
\begin{pgfscope}%
\pgfsys@transformshift{0.628480in}{1.607138in}%
\pgfsys@useobject{currentmarker}{}%
\end{pgfscope}%
\end{pgfscope}%
\begin{pgfscope}%
\pgfpathrectangle{\pgfqpoint{0.499500in}{0.990000in}}{\pgfqpoint{2.730600in}{1.860000in}}%
\pgfusepath{clip}%
\pgfsetbuttcap%
\pgfsetroundjoin%
\definecolor{currentfill}{rgb}{0.478431,0.478431,0.478431}%
\pgfsetfillcolor{currentfill}%
\pgfsetlinewidth{0.301125pt}%
\definecolor{currentstroke}{rgb}{1.000000,1.000000,1.000000}%
\pgfsetstrokecolor{currentstroke}%
\pgfsetdash{}{0pt}%
\pgfsys@defobject{currentmarker}{\pgfqpoint{-0.027778in}{-0.027778in}}{\pgfqpoint{0.027778in}{0.027778in}}{%
\pgfpathmoveto{\pgfqpoint{-0.027778in}{-0.027778in}}%
\pgfpathlineto{\pgfqpoint{0.027778in}{-0.027778in}}%
\pgfpathlineto{\pgfqpoint{0.027778in}{0.027778in}}%
\pgfpathlineto{\pgfqpoint{-0.027778in}{0.027778in}}%
\pgfpathlineto{\pgfqpoint{-0.027778in}{-0.027778in}}%
\pgfpathclose%
\pgfusepath{stroke,fill}%
}%
\begin{pgfscope}%
\pgfsys@transformshift{0.756991in}{2.183428in}%
\pgfsys@useobject{currentmarker}{}%
\end{pgfscope}%
\end{pgfscope}%
\begin{pgfscope}%
\pgfpathrectangle{\pgfqpoint{0.499500in}{0.990000in}}{\pgfqpoint{2.730600in}{1.860000in}}%
\pgfusepath{clip}%
\pgfsetbuttcap%
\pgfsetroundjoin%
\definecolor{currentfill}{rgb}{0.478431,0.478431,0.478431}%
\pgfsetfillcolor{currentfill}%
\pgfsetlinewidth{0.301125pt}%
\definecolor{currentstroke}{rgb}{1.000000,1.000000,1.000000}%
\pgfsetstrokecolor{currentstroke}%
\pgfsetdash{}{0pt}%
\pgfsys@defobject{currentmarker}{\pgfqpoint{-0.027778in}{-0.027778in}}{\pgfqpoint{0.027778in}{0.027778in}}{%
\pgfpathmoveto{\pgfqpoint{0.000000in}{0.027778in}}%
\pgfpathlineto{\pgfqpoint{-0.027778in}{-0.027778in}}%
\pgfpathlineto{\pgfqpoint{0.027778in}{-0.027778in}}%
\pgfpathlineto{\pgfqpoint{0.000000in}{0.027778in}}%
\pgfpathclose%
\pgfusepath{stroke,fill}%
}%
\begin{pgfscope}%
\pgfsys@transformshift{0.724656in}{1.169280in}%
\pgfsys@useobject{currentmarker}{}%
\end{pgfscope}%
\end{pgfscope}%
\begin{pgfscope}%
\pgfpathrectangle{\pgfqpoint{0.499500in}{0.990000in}}{\pgfqpoint{2.730600in}{1.860000in}}%
\pgfusepath{clip}%
\pgfsetbuttcap%
\pgfsetroundjoin%
\definecolor{currentfill}{rgb}{0.478431,0.305882,0.611765}%
\pgfsetfillcolor{currentfill}%
\pgfsetlinewidth{0.301125pt}%
\definecolor{currentstroke}{rgb}{1.000000,1.000000,1.000000}%
\pgfsetstrokecolor{currentstroke}%
\pgfsetdash{}{0pt}%
\pgfsys@defobject{currentmarker}{\pgfqpoint{-0.027778in}{-0.027778in}}{\pgfqpoint{0.027778in}{0.027778in}}{%
\pgfpathmoveto{\pgfqpoint{0.000000in}{0.027778in}}%
\pgfpathlineto{\pgfqpoint{-0.027778in}{-0.027778in}}%
\pgfpathlineto{\pgfqpoint{0.027778in}{-0.027778in}}%
\pgfpathlineto{\pgfqpoint{0.000000in}{0.027778in}}%
\pgfpathclose%
\pgfusepath{stroke,fill}%
}%
\begin{pgfscope}%
\pgfsys@transformshift{0.576781in}{1.053409in}%
\pgfsys@useobject{currentmarker}{}%
\end{pgfscope}%
\end{pgfscope}%
\begin{pgfscope}%
\pgfpathrectangle{\pgfqpoint{0.499500in}{0.990000in}}{\pgfqpoint{2.730600in}{1.860000in}}%
\pgfusepath{clip}%
\pgfsetbuttcap%
\pgfsetroundjoin%
\definecolor{currentfill}{rgb}{0.156863,0.356863,0.568627}%
\pgfsetfillcolor{currentfill}%
\pgfsetlinewidth{0.602250pt}%
\definecolor{currentstroke}{rgb}{0.101961,0.101961,0.101961}%
\pgfsetstrokecolor{currentstroke}%
\pgfsetdash{}{0pt}%
\pgfsys@defobject{currentmarker}{\pgfqpoint{-0.040493in}{-0.040493in}}{\pgfqpoint{0.040493in}{0.040493in}}{%
\pgfpathmoveto{\pgfqpoint{0.000000in}{-0.040493in}}%
\pgfpathcurveto{\pgfqpoint{0.010739in}{-0.040493in}}{\pgfqpoint{0.021039in}{-0.036226in}}{\pgfqpoint{0.028633in}{-0.028633in}}%
\pgfpathcurveto{\pgfqpoint{0.036226in}{-0.021039in}}{\pgfqpoint{0.040493in}{-0.010739in}}{\pgfqpoint{0.040493in}{0.000000in}}%
\pgfpathcurveto{\pgfqpoint{0.040493in}{0.010739in}}{\pgfqpoint{0.036226in}{0.021039in}}{\pgfqpoint{0.028633in}{0.028633in}}%
\pgfpathcurveto{\pgfqpoint{0.021039in}{0.036226in}}{\pgfqpoint{0.010739in}{0.040493in}}{\pgfqpoint{0.000000in}{0.040493in}}%
\pgfpathcurveto{\pgfqpoint{-0.010739in}{0.040493in}}{\pgfqpoint{-0.021039in}{0.036226in}}{\pgfqpoint{-0.028633in}{0.028633in}}%
\pgfpathcurveto{\pgfqpoint{-0.036226in}{0.021039in}}{\pgfqpoint{-0.040493in}{0.010739in}}{\pgfqpoint{-0.040493in}{0.000000in}}%
\pgfpathcurveto{\pgfqpoint{-0.040493in}{-0.010739in}}{\pgfqpoint{-0.036226in}{-0.021039in}}{\pgfqpoint{-0.028633in}{-0.028633in}}%
\pgfpathcurveto{\pgfqpoint{-0.021039in}{-0.036226in}}{\pgfqpoint{-0.010739in}{-0.040493in}}{\pgfqpoint{0.000000in}{-0.040493in}}%
\pgfpathlineto{\pgfqpoint{0.000000in}{-0.040493in}}%
\pgfpathclose%
\pgfusepath{stroke,fill}%
}%
\begin{pgfscope}%
\pgfsys@transformshift{2.911094in}{2.136677in}%
\pgfsys@useobject{currentmarker}{}%
\end{pgfscope}%
\end{pgfscope}%
\begin{pgfscope}%
\pgfpathrectangle{\pgfqpoint{0.499500in}{0.990000in}}{\pgfqpoint{2.730600in}{1.860000in}}%
\pgfusepath{clip}%
\pgfsetbuttcap%
\pgfsetroundjoin%
\definecolor{currentfill}{rgb}{0.156863,0.356863,0.568627}%
\pgfsetfillcolor{currentfill}%
\pgfsetlinewidth{0.602250pt}%
\definecolor{currentstroke}{rgb}{0.101961,0.101961,0.101961}%
\pgfsetstrokecolor{currentstroke}%
\pgfsetdash{}{0pt}%
\pgfsys@defobject{currentmarker}{\pgfqpoint{-0.040493in}{-0.040493in}}{\pgfqpoint{0.040493in}{0.040493in}}{%
\pgfpathmoveto{\pgfqpoint{-0.040493in}{-0.040493in}}%
\pgfpathlineto{\pgfqpoint{0.040493in}{-0.040493in}}%
\pgfpathlineto{\pgfqpoint{0.040493in}{0.040493in}}%
\pgfpathlineto{\pgfqpoint{-0.040493in}{0.040493in}}%
\pgfpathlineto{\pgfqpoint{-0.040493in}{-0.040493in}}%
\pgfpathclose%
\pgfusepath{stroke,fill}%
}%
\begin{pgfscope}%
\pgfsys@transformshift{3.122146in}{2.319507in}%
\pgfsys@useobject{currentmarker}{}%
\end{pgfscope}%
\end{pgfscope}%
\begin{pgfscope}%
\pgfpathrectangle{\pgfqpoint{0.499500in}{0.990000in}}{\pgfqpoint{2.730600in}{1.860000in}}%
\pgfusepath{clip}%
\pgfsetbuttcap%
\pgfsetroundjoin%
\definecolor{currentfill}{rgb}{0.156863,0.356863,0.568627}%
\pgfsetfillcolor{currentfill}%
\pgfsetlinewidth{0.602250pt}%
\definecolor{currentstroke}{rgb}{0.101961,0.101961,0.101961}%
\pgfsetstrokecolor{currentstroke}%
\pgfsetdash{}{0pt}%
\pgfsys@defobject{currentmarker}{\pgfqpoint{-0.040493in}{-0.040493in}}{\pgfqpoint{0.040493in}{0.040493in}}{%
\pgfpathmoveto{\pgfqpoint{0.000000in}{0.040493in}}%
\pgfpathlineto{\pgfqpoint{-0.040493in}{-0.040493in}}%
\pgfpathlineto{\pgfqpoint{0.040493in}{-0.040493in}}%
\pgfpathlineto{\pgfqpoint{0.000000in}{0.040493in}}%
\pgfpathclose%
\pgfusepath{stroke,fill}%
}%
\begin{pgfscope}%
\pgfsys@transformshift{2.438043in}{2.571082in}%
\pgfsys@useobject{currentmarker}{}%
\end{pgfscope}%
\end{pgfscope}%
\begin{pgfscope}%
\definecolor{textcolor}{rgb}{0.549020,0.549020,0.549020}%
\pgfsetstrokecolor{textcolor}%
\pgfsetfillcolor{textcolor}%
\pgftext[x=3.296700in,y=2.955000in,right,top]{\color{textcolor}\sffamily\fontsize{5.500000}{6.600000}\selectfont OLD DATA}%
\end{pgfscope}%
\begin{pgfscope}%
\pgfsetbuttcap%
\pgfsetroundjoin%
\definecolor{currentfill}{rgb}{0.478431,0.478431,0.478431}%
\pgfsetfillcolor{currentfill}%
\pgfsetlinewidth{1.003750pt}%
\definecolor{currentstroke}{rgb}{0.478431,0.478431,0.478431}%
\pgfsetstrokecolor{currentstroke}%
\pgfsetdash{}{0pt}%
\pgfsys@defobject{currentmarker}{\pgfqpoint{-0.027778in}{-0.027778in}}{\pgfqpoint{0.027778in}{0.027778in}}{%
\pgfpathmoveto{\pgfqpoint{0.000000in}{-0.027778in}}%
\pgfpathcurveto{\pgfqpoint{0.007367in}{-0.027778in}}{\pgfqpoint{0.014433in}{-0.024851in}}{\pgfqpoint{0.019642in}{-0.019642in}}%
\pgfpathcurveto{\pgfqpoint{0.024851in}{-0.014433in}}{\pgfqpoint{0.027778in}{-0.007367in}}{\pgfqpoint{0.027778in}{0.000000in}}%
\pgfpathcurveto{\pgfqpoint{0.027778in}{0.007367in}}{\pgfqpoint{0.024851in}{0.014433in}}{\pgfqpoint{0.019642in}{0.019642in}}%
\pgfpathcurveto{\pgfqpoint{0.014433in}{0.024851in}}{\pgfqpoint{0.007367in}{0.027778in}}{\pgfqpoint{0.000000in}{0.027778in}}%
\pgfpathcurveto{\pgfqpoint{-0.007367in}{0.027778in}}{\pgfqpoint{-0.014433in}{0.024851in}}{\pgfqpoint{-0.019642in}{0.019642in}}%
\pgfpathcurveto{\pgfqpoint{-0.024851in}{0.014433in}}{\pgfqpoint{-0.027778in}{0.007367in}}{\pgfqpoint{-0.027778in}{0.000000in}}%
\pgfpathcurveto{\pgfqpoint{-0.027778in}{-0.007367in}}{\pgfqpoint{-0.024851in}{-0.014433in}}{\pgfqpoint{-0.019642in}{-0.019642in}}%
\pgfpathcurveto{\pgfqpoint{-0.014433in}{-0.024851in}}{\pgfqpoint{-0.007367in}{-0.027778in}}{\pgfqpoint{0.000000in}{-0.027778in}}%
\pgfpathlineto{\pgfqpoint{0.000000in}{-0.027778in}}%
\pgfpathclose%
\pgfusepath{stroke,fill}%
}%
\begin{pgfscope}%
\pgfsys@transformshift{0.251322in}{0.372685in}%
\pgfsys@useobject{currentmarker}{}%
\end{pgfscope}%
\end{pgfscope}%
\begin{pgfscope}%
\definecolor{textcolor}{rgb}{0.101961,0.101961,0.101961}%
\pgfsetstrokecolor{textcolor}%
\pgfsetfillcolor{textcolor}%
\pgftext[x=0.367989in,y=0.338657in,left,base]{\color{textcolor}\sffamily\fontsize{7.000000}{8.400000}\selectfont Robot}%
\end{pgfscope}%
\begin{pgfscope}%
\pgfsetbuttcap%
\pgfsetmiterjoin%
\definecolor{currentfill}{rgb}{0.478431,0.478431,0.478431}%
\pgfsetfillcolor{currentfill}%
\pgfsetlinewidth{1.003750pt}%
\definecolor{currentstroke}{rgb}{0.478431,0.478431,0.478431}%
\pgfsetstrokecolor{currentstroke}%
\pgfsetdash{}{0pt}%
\pgfsys@defobject{currentmarker}{\pgfqpoint{-0.027778in}{-0.027778in}}{\pgfqpoint{0.027778in}{0.027778in}}{%
\pgfpathmoveto{\pgfqpoint{-0.027778in}{-0.027778in}}%
\pgfpathlineto{\pgfqpoint{0.027778in}{-0.027778in}}%
\pgfpathlineto{\pgfqpoint{0.027778in}{0.027778in}}%
\pgfpathlineto{\pgfqpoint{-0.027778in}{0.027778in}}%
\pgfpathlineto{\pgfqpoint{-0.027778in}{-0.027778in}}%
\pgfpathclose%
\pgfusepath{stroke,fill}%
}%
\begin{pgfscope}%
\pgfsys@transformshift{0.251322in}{0.256558in}%
\pgfsys@useobject{currentmarker}{}%
\end{pgfscope}%
\end{pgfscope}%
\begin{pgfscope}%
\definecolor{textcolor}{rgb}{0.101961,0.101961,0.101961}%
\pgfsetstrokecolor{textcolor}%
\pgfsetfillcolor{textcolor}%
\pgftext[x=0.367989in,y=0.222531in,left,base]{\color{textcolor}\sffamily\fontsize{7.000000}{8.400000}\selectfont Single-shot}%
\end{pgfscope}%
\begin{pgfscope}%
\pgfsetbuttcap%
\pgfsetmiterjoin%
\definecolor{currentfill}{rgb}{0.478431,0.478431,0.478431}%
\pgfsetfillcolor{currentfill}%
\pgfsetlinewidth{1.003750pt}%
\definecolor{currentstroke}{rgb}{0.478431,0.478431,0.478431}%
\pgfsetstrokecolor{currentstroke}%
\pgfsetdash{}{0pt}%
\pgfsys@defobject{currentmarker}{\pgfqpoint{-0.027778in}{-0.027778in}}{\pgfqpoint{0.027778in}{0.027778in}}{%
\pgfpathmoveto{\pgfqpoint{0.000000in}{0.027778in}}%
\pgfpathlineto{\pgfqpoint{-0.027778in}{-0.027778in}}%
\pgfpathlineto{\pgfqpoint{0.027778in}{-0.027778in}}%
\pgfpathlineto{\pgfqpoint{0.000000in}{0.027778in}}%
\pgfpathclose%
\pgfusepath{stroke,fill}%
}%
\begin{pgfscope}%
\pgfsys@transformshift{0.251322in}{0.140432in}%
\pgfsys@useobject{currentmarker}{}%
\end{pgfscope}%
\end{pgfscope}%
\begin{pgfscope}%
\definecolor{textcolor}{rgb}{0.101961,0.101961,0.101961}%
\pgfsetstrokecolor{textcolor}%
\pgfsetfillcolor{textcolor}%
\pgftext[x=0.367989in,y=0.106404in,left,base]{\color{textcolor}\sffamily\fontsize{7.000000}{8.400000}\selectfont Dashcam}%
\end{pgfscope}%
\begin{pgfscope}%
\pgfsetbuttcap%
\pgfsetroundjoin%
\definecolor{currentfill}{rgb}{0.156863,0.356863,0.568627}%
\pgfsetfillcolor{currentfill}%
\pgfsetlinewidth{1.003750pt}%
\definecolor{currentstroke}{rgb}{0.156863,0.356863,0.568627}%
\pgfsetstrokecolor{currentstroke}%
\pgfsetdash{}{0pt}%
\pgfsys@defobject{currentmarker}{\pgfqpoint{-0.027778in}{-0.027778in}}{\pgfqpoint{0.027778in}{0.027778in}}{%
\pgfpathmoveto{\pgfqpoint{0.000000in}{-0.027778in}}%
\pgfpathcurveto{\pgfqpoint{0.007367in}{-0.027778in}}{\pgfqpoint{0.014433in}{-0.024851in}}{\pgfqpoint{0.019642in}{-0.019642in}}%
\pgfpathcurveto{\pgfqpoint{0.024851in}{-0.014433in}}{\pgfqpoint{0.027778in}{-0.007367in}}{\pgfqpoint{0.027778in}{0.000000in}}%
\pgfpathcurveto{\pgfqpoint{0.027778in}{0.007367in}}{\pgfqpoint{0.024851in}{0.014433in}}{\pgfqpoint{0.019642in}{0.019642in}}%
\pgfpathcurveto{\pgfqpoint{0.014433in}{0.024851in}}{\pgfqpoint{0.007367in}{0.027778in}}{\pgfqpoint{0.000000in}{0.027778in}}%
\pgfpathcurveto{\pgfqpoint{-0.007367in}{0.027778in}}{\pgfqpoint{-0.014433in}{0.024851in}}{\pgfqpoint{-0.019642in}{0.019642in}}%
\pgfpathcurveto{\pgfqpoint{-0.024851in}{0.014433in}}{\pgfqpoint{-0.027778in}{0.007367in}}{\pgfqpoint{-0.027778in}{0.000000in}}%
\pgfpathcurveto{\pgfqpoint{-0.027778in}{-0.007367in}}{\pgfqpoint{-0.024851in}{-0.014433in}}{\pgfqpoint{-0.019642in}{-0.019642in}}%
\pgfpathcurveto{\pgfqpoint{-0.014433in}{-0.024851in}}{\pgfqpoint{-0.007367in}{-0.027778in}}{\pgfqpoint{0.000000in}{-0.027778in}}%
\pgfpathlineto{\pgfqpoint{0.000000in}{-0.027778in}}%
\pgfpathclose%
\pgfusepath{stroke,fill}%
}%
\begin{pgfscope}%
\pgfsys@transformshift{1.383522in}{0.494213in}%
\pgfsys@useobject{currentmarker}{}%
\end{pgfscope}%
\end{pgfscope}%
\begin{pgfscope}%
\definecolor{textcolor}{rgb}{0.101961,0.101961,0.101961}%
\pgfsetstrokecolor{textcolor}%
\pgfsetfillcolor{textcolor}%
\pgftext[x=1.500189in,y=0.460185in,left,base]{\color{textcolor}\sffamily\fontsize{7.000000}{8.400000}\selectfont COVEX (old: rung 3)}%
\end{pgfscope}%
\begin{pgfscope}%
\pgfsetbuttcap%
\pgfsetroundjoin%
\definecolor{currentfill}{rgb}{0.478431,0.478431,0.478431}%
\pgfsetfillcolor{currentfill}%
\pgfsetlinewidth{1.003750pt}%
\definecolor{currentstroke}{rgb}{0.478431,0.478431,0.478431}%
\pgfsetstrokecolor{currentstroke}%
\pgfsetdash{}{0pt}%
\pgfsys@defobject{currentmarker}{\pgfqpoint{-0.027778in}{-0.027778in}}{\pgfqpoint{0.027778in}{0.027778in}}{%
\pgfpathmoveto{\pgfqpoint{0.000000in}{-0.027778in}}%
\pgfpathcurveto{\pgfqpoint{0.007367in}{-0.027778in}}{\pgfqpoint{0.014433in}{-0.024851in}}{\pgfqpoint{0.019642in}{-0.019642in}}%
\pgfpathcurveto{\pgfqpoint{0.024851in}{-0.014433in}}{\pgfqpoint{0.027778in}{-0.007367in}}{\pgfqpoint{0.027778in}{0.000000in}}%
\pgfpathcurveto{\pgfqpoint{0.027778in}{0.007367in}}{\pgfqpoint{0.024851in}{0.014433in}}{\pgfqpoint{0.019642in}{0.019642in}}%
\pgfpathcurveto{\pgfqpoint{0.014433in}{0.024851in}}{\pgfqpoint{0.007367in}{0.027778in}}{\pgfqpoint{0.000000in}{0.027778in}}%
\pgfpathcurveto{\pgfqpoint{-0.007367in}{0.027778in}}{\pgfqpoint{-0.014433in}{0.024851in}}{\pgfqpoint{-0.019642in}{0.019642in}}%
\pgfpathcurveto{\pgfqpoint{-0.024851in}{0.014433in}}{\pgfqpoint{-0.027778in}{0.007367in}}{\pgfqpoint{-0.027778in}{0.000000in}}%
\pgfpathcurveto{\pgfqpoint{-0.027778in}{-0.007367in}}{\pgfqpoint{-0.024851in}{-0.014433in}}{\pgfqpoint{-0.019642in}{-0.019642in}}%
\pgfpathcurveto{\pgfqpoint{-0.014433in}{-0.024851in}}{\pgfqpoint{-0.007367in}{-0.027778in}}{\pgfqpoint{0.000000in}{-0.027778in}}%
\pgfpathlineto{\pgfqpoint{0.000000in}{-0.027778in}}%
\pgfpathclose%
\pgfusepath{stroke,fill}%
}%
\begin{pgfscope}%
\pgfsys@transformshift{1.383522in}{0.372685in}%
\pgfsys@useobject{currentmarker}{}%
\end{pgfscope}%
\end{pgfscope}%
\begin{pgfscope}%
\definecolor{textcolor}{rgb}{0.101961,0.101961,0.101961}%
\pgfsetstrokecolor{textcolor}%
\pgfsetfillcolor{textcolor}%
\pgftext[x=1.500189in,y=0.338657in,left,base]{\color{textcolor}\sffamily\fontsize{7.000000}{8.400000}\selectfont Stateful}%
\end{pgfscope}%
\begin{pgfscope}%
\pgfsetbuttcap%
\pgfsetroundjoin%
\definecolor{currentfill}{rgb}{0.784314,0.784314,0.784314}%
\pgfsetfillcolor{currentfill}%
\pgfsetlinewidth{1.003750pt}%
\definecolor{currentstroke}{rgb}{0.784314,0.784314,0.784314}%
\pgfsetstrokecolor{currentstroke}%
\pgfsetdash{}{0pt}%
\pgfsys@defobject{currentmarker}{\pgfqpoint{-0.027778in}{-0.027778in}}{\pgfqpoint{0.027778in}{0.027778in}}{%
\pgfpathmoveto{\pgfqpoint{0.000000in}{-0.027778in}}%
\pgfpathcurveto{\pgfqpoint{0.007367in}{-0.027778in}}{\pgfqpoint{0.014433in}{-0.024851in}}{\pgfqpoint{0.019642in}{-0.019642in}}%
\pgfpathcurveto{\pgfqpoint{0.024851in}{-0.014433in}}{\pgfqpoint{0.027778in}{-0.007367in}}{\pgfqpoint{0.027778in}{0.000000in}}%
\pgfpathcurveto{\pgfqpoint{0.027778in}{0.007367in}}{\pgfqpoint{0.024851in}{0.014433in}}{\pgfqpoint{0.019642in}{0.019642in}}%
\pgfpathcurveto{\pgfqpoint{0.014433in}{0.024851in}}{\pgfqpoint{0.007367in}{0.027778in}}{\pgfqpoint{0.000000in}{0.027778in}}%
\pgfpathcurveto{\pgfqpoint{-0.007367in}{0.027778in}}{\pgfqpoint{-0.014433in}{0.024851in}}{\pgfqpoint{-0.019642in}{0.019642in}}%
\pgfpathcurveto{\pgfqpoint{-0.024851in}{0.014433in}}{\pgfqpoint{-0.027778in}{0.007367in}}{\pgfqpoint{-0.027778in}{0.000000in}}%
\pgfpathcurveto{\pgfqpoint{-0.027778in}{-0.007367in}}{\pgfqpoint{-0.024851in}{-0.014433in}}{\pgfqpoint{-0.019642in}{-0.019642in}}%
\pgfpathcurveto{\pgfqpoint{-0.014433in}{-0.024851in}}{\pgfqpoint{-0.007367in}{-0.027778in}}{\pgfqpoint{0.000000in}{-0.027778in}}%
\pgfpathlineto{\pgfqpoint{0.000000in}{-0.027778in}}%
\pgfpathclose%
\pgfusepath{stroke,fill}%
}%
\begin{pgfscope}%
\pgfsys@transformshift{1.383522in}{0.256558in}%
\pgfsys@useobject{currentmarker}{}%
\end{pgfscope}%
\end{pgfscope}%
\begin{pgfscope}%
\definecolor{textcolor}{rgb}{0.101961,0.101961,0.101961}%
\pgfsetstrokecolor{textcolor}%
\pgfsetfillcolor{textcolor}%
\pgftext[x=1.500189in,y=0.222531in,left,base]{\color{textcolor}\sffamily\fontsize{7.000000}{8.400000}\selectfont Stateless}%
\end{pgfscope}%
\begin{pgfscope}%
\pgfsetbuttcap%
\pgfsetroundjoin%
\definecolor{currentfill}{rgb}{0.478431,0.305882,0.611765}%
\pgfsetfillcolor{currentfill}%
\pgfsetlinewidth{1.003750pt}%
\definecolor{currentstroke}{rgb}{0.478431,0.305882,0.611765}%
\pgfsetstrokecolor{currentstroke}%
\pgfsetdash{}{0pt}%
\pgfsys@defobject{currentmarker}{\pgfqpoint{-0.027778in}{-0.027778in}}{\pgfqpoint{0.027778in}{0.027778in}}{%
\pgfpathmoveto{\pgfqpoint{0.000000in}{-0.027778in}}%
\pgfpathcurveto{\pgfqpoint{0.007367in}{-0.027778in}}{\pgfqpoint{0.014433in}{-0.024851in}}{\pgfqpoint{0.019642in}{-0.019642in}}%
\pgfpathcurveto{\pgfqpoint{0.024851in}{-0.014433in}}{\pgfqpoint{0.027778in}{-0.007367in}}{\pgfqpoint{0.027778in}{0.000000in}}%
\pgfpathcurveto{\pgfqpoint{0.027778in}{0.007367in}}{\pgfqpoint{0.024851in}{0.014433in}}{\pgfqpoint{0.019642in}{0.019642in}}%
\pgfpathcurveto{\pgfqpoint{0.014433in}{0.024851in}}{\pgfqpoint{0.007367in}{0.027778in}}{\pgfqpoint{0.000000in}{0.027778in}}%
\pgfpathcurveto{\pgfqpoint{-0.007367in}{0.027778in}}{\pgfqpoint{-0.014433in}{0.024851in}}{\pgfqpoint{-0.019642in}{0.019642in}}%
\pgfpathcurveto{\pgfqpoint{-0.024851in}{0.014433in}}{\pgfqpoint{-0.027778in}{0.007367in}}{\pgfqpoint{-0.027778in}{0.000000in}}%
\pgfpathcurveto{\pgfqpoint{-0.027778in}{-0.007367in}}{\pgfqpoint{-0.024851in}{-0.014433in}}{\pgfqpoint{-0.019642in}{-0.019642in}}%
\pgfpathcurveto{\pgfqpoint{-0.014433in}{-0.024851in}}{\pgfqpoint{-0.007367in}{-0.027778in}}{\pgfqpoint{0.000000in}{-0.027778in}}%
\pgfpathlineto{\pgfqpoint{0.000000in}{-0.027778in}}%
\pgfpathclose%
\pgfusepath{stroke,fill}%
}%
\begin{pgfscope}%
\pgfsys@transformshift{1.383522in}{0.140432in}%
\pgfsys@useobject{currentmarker}{}%
\end{pgfscope}%
\end{pgfscope}%
\begin{pgfscope}%
\definecolor{textcolor}{rgb}{0.101961,0.101961,0.101961}%
\pgfsetstrokecolor{textcolor}%
\pgfsetfillcolor{textcolor}%
\pgftext[x=1.500189in,y=0.106404in,left,base]{\color{textcolor}\sffamily\fontsize{7.000000}{8.400000}\selectfont On-device only}%
\end{pgfscope}%
\end{pgfpicture}%
\makeatother%
\endgroup%
}
\caption{[PLACEHOLDER: OLD DATA, rung 3 stand-in for COVEX, baselines on dev seeds 4-7] On-time rate against correct-by-deadline per arm (mean over seeds 4 to 7, standard condition); one marker per arm per question type. Dashcam correctness is per-episode attainment. The on-device model has no on-time rate for robot and single-shot and appears for dashcam only; confidence cascade has no cells.}
\label{fig:F2_ontime_vs_correct}
\end{figure}
\clearpage
\begin{figure*}[t]
\centering
{\fontfamily{LinuxBiolinumT}\selectfont
%% Creator: Matplotlib, PGF backend
%%
%% To include the figure in your LaTeX document, write
%%   \input{<filename>.pgf}
%%
%% Make sure the required packages are loaded in your preamble
%%   \usepackage{pgf}
%%
%% Also ensure that all the required font packages are loaded; for instance,
%% the lmodern package is sometimes necessary when using math font.
%%   \usepackage{lmodern}
%%
%% Figures using additional raster images can only be included by \input if
%% they are in the same directory as the main LaTeX file. For loading figures
%% from other directories you can use the `import` package
%%   \usepackage{import}
%%
%% and then include the figures with
%%   \import{<path to file>}{<filename>.pgf}
%%
%% Matplotlib used the following preamble
%%
\begingroup%
\makeatletter%
\begin{pgfpicture}%
\pgfpathrectangle{\pgfpointorigin}{\pgfqpoint{7.000000in}{2.300000in}}%
\pgfusepath{use as bounding box, clip}%
\begin{pgfscope}%
\pgfsetbuttcap%
\pgfsetmiterjoin%
\definecolor{currentfill}{rgb}{1.000000,1.000000,1.000000}%
\pgfsetfillcolor{currentfill}%
\pgfsetlinewidth{0.000000pt}%
\definecolor{currentstroke}{rgb}{1.000000,1.000000,1.000000}%
\pgfsetstrokecolor{currentstroke}%
\pgfsetdash{}{0pt}%
\pgfpathmoveto{\pgfqpoint{0.000000in}{0.000000in}}%
\pgfpathlineto{\pgfqpoint{7.000000in}{0.000000in}}%
\pgfpathlineto{\pgfqpoint{7.000000in}{2.300000in}}%
\pgfpathlineto{\pgfqpoint{0.000000in}{2.300000in}}%
\pgfpathlineto{\pgfqpoint{0.000000in}{0.000000in}}%
\pgfpathclose%
\pgfusepath{fill}%
\end{pgfscope}%
\begin{pgfscope}%
\pgfsetbuttcap%
\pgfsetmiterjoin%
\definecolor{currentfill}{rgb}{1.000000,1.000000,1.000000}%
\pgfsetfillcolor{currentfill}%
\pgfsetlinewidth{0.000000pt}%
\definecolor{currentstroke}{rgb}{0.000000,0.000000,0.000000}%
\pgfsetstrokecolor{currentstroke}%
\pgfsetstrokeopacity{0.000000}%
\pgfsetdash{}{0pt}%
\pgfpathmoveto{\pgfqpoint{0.840000in}{0.414000in}}%
\pgfpathlineto{\pgfqpoint{6.720000in}{0.414000in}}%
\pgfpathlineto{\pgfqpoint{6.720000in}{2.070000in}}%
\pgfpathlineto{\pgfqpoint{0.840000in}{2.070000in}}%
\pgfpathlineto{\pgfqpoint{0.840000in}{0.414000in}}%
\pgfpathclose%
\pgfusepath{fill}%
\end{pgfscope}%
\begin{pgfscope}%
\pgfsetbuttcap%
\pgfsetmiterjoin%
\definecolor{currentfill}{rgb}{1.000000,1.000000,1.000000}%
\pgfsetfillcolor{currentfill}%
\pgfsetlinewidth{0.000000pt}%
\definecolor{currentstroke}{rgb}{0.000000,0.000000,0.000000}%
\pgfsetstrokecolor{currentstroke}%
\pgfsetstrokeopacity{0.000000}%
\pgfsetdash{}{0pt}%
\pgfpathmoveto{\pgfqpoint{0.490000in}{0.690000in}}%
\pgfpathlineto{\pgfqpoint{2.450000in}{0.690000in}}%
\pgfpathlineto{\pgfqpoint{2.450000in}{2.116000in}}%
\pgfpathlineto{\pgfqpoint{0.490000in}{2.116000in}}%
\pgfpathlineto{\pgfqpoint{0.490000in}{0.690000in}}%
\pgfpathclose%
\pgfusepath{fill}%
\end{pgfscope}%
\begin{pgfscope}%
\definecolor{textcolor}{rgb}{0.101961,0.101961,0.101961}%
\pgfsetstrokecolor{textcolor}%
\pgfsetfillcolor{textcolor}%
\pgftext[x=1.470000in,y=0.662222in,,top]{\color{textcolor}\sffamily\fontsize{7.500000}{9.000000}\selectfont Robot: answer latency (s)}%
\end{pgfscope}%
\begin{pgfscope}%
\pgfpathrectangle{\pgfqpoint{0.490000in}{0.690000in}}{\pgfqpoint{1.960000in}{1.426000in}}%
\pgfusepath{clip}%
\pgfsetrectcap%
\pgfsetroundjoin%
\pgfsetlinewidth{0.401500pt}%
\definecolor{currentstroke}{rgb}{0.901961,0.901961,0.901961}%
\pgfsetstrokecolor{currentstroke}%
\pgfsetdash{}{0pt}%
\pgfpathmoveto{\pgfqpoint{0.490000in}{0.690000in}}%
\pgfpathlineto{\pgfqpoint{2.450000in}{0.690000in}}%
\pgfusepath{stroke}%
\end{pgfscope}%
\begin{pgfscope}%
\definecolor{textcolor}{rgb}{0.101961,0.101961,0.101961}%
\pgfsetstrokecolor{textcolor}%
\pgfsetfillcolor{textcolor}%
\pgftext[x=0.389739in, y=0.656242in, left, base]{\color{textcolor}\sffamily\fontsize{7.000000}{8.400000}\selectfont 0}%
\end{pgfscope}%
\begin{pgfscope}%
\pgfpathrectangle{\pgfqpoint{0.490000in}{0.690000in}}{\pgfqpoint{1.960000in}{1.426000in}}%
\pgfusepath{clip}%
\pgfsetrectcap%
\pgfsetroundjoin%
\pgfsetlinewidth{0.401500pt}%
\definecolor{currentstroke}{rgb}{0.901961,0.901961,0.901961}%
\pgfsetstrokecolor{currentstroke}%
\pgfsetdash{}{0pt}%
\pgfpathmoveto{\pgfqpoint{0.490000in}{1.403000in}}%
\pgfpathlineto{\pgfqpoint{2.450000in}{1.403000in}}%
\pgfusepath{stroke}%
\end{pgfscope}%
\begin{pgfscope}%
\definecolor{textcolor}{rgb}{0.101961,0.101961,0.101961}%
\pgfsetstrokecolor{textcolor}%
\pgfsetfillcolor{textcolor}%
\pgftext[x=0.309395in, y=1.369242in, left, base]{\color{textcolor}\sffamily\fontsize{7.000000}{8.400000}\selectfont 0.5}%
\end{pgfscope}%
\begin{pgfscope}%
\pgfpathrectangle{\pgfqpoint{0.490000in}{0.690000in}}{\pgfqpoint{1.960000in}{1.426000in}}%
\pgfusepath{clip}%
\pgfsetrectcap%
\pgfsetroundjoin%
\pgfsetlinewidth{0.401500pt}%
\definecolor{currentstroke}{rgb}{0.901961,0.901961,0.901961}%
\pgfsetstrokecolor{currentstroke}%
\pgfsetdash{}{0pt}%
\pgfpathmoveto{\pgfqpoint{0.490000in}{2.116000in}}%
\pgfpathlineto{\pgfqpoint{2.450000in}{2.116000in}}%
\pgfusepath{stroke}%
\end{pgfscope}%
\begin{pgfscope}%
\definecolor{textcolor}{rgb}{0.101961,0.101961,0.101961}%
\pgfsetstrokecolor{textcolor}%
\pgfsetfillcolor{textcolor}%
\pgftext[x=0.389739in, y=2.082242in, left, base]{\color{textcolor}\sffamily\fontsize{7.000000}{8.400000}\selectfont 1}%
\end{pgfscope}%
\begin{pgfscope}%
\definecolor{textcolor}{rgb}{0.101961,0.101961,0.101961}%
\pgfsetstrokecolor{textcolor}%
\pgfsetfillcolor{textcolor}%
\pgftext[x=0.281617in,y=1.403000in,,bottom,rotate=90.000000]{\color{textcolor}\sffamily\fontsize{7.500000}{9.000000}\selectfont Fraction of requests}%
\end{pgfscope}%
\begin{pgfscope}%
\pgfsetrectcap%
\pgfsetmiterjoin%
\pgfsetlinewidth{0.803000pt}%
\definecolor{currentstroke}{rgb}{0.101961,0.101961,0.101961}%
\pgfsetstrokecolor{currentstroke}%
\pgfsetdash{}{0pt}%
\pgfpathmoveto{\pgfqpoint{0.490000in}{0.690000in}}%
\pgfpathlineto{\pgfqpoint{0.490000in}{2.116000in}}%
\pgfusepath{stroke}%
\end{pgfscope}%
\begin{pgfscope}%
\pgfsetrectcap%
\pgfsetmiterjoin%
\pgfsetlinewidth{0.803000pt}%
\definecolor{currentstroke}{rgb}{0.101961,0.101961,0.101961}%
\pgfsetstrokecolor{currentstroke}%
\pgfsetdash{}{0pt}%
\pgfpathmoveto{\pgfqpoint{0.490000in}{0.690000in}}%
\pgfpathlineto{\pgfqpoint{2.450000in}{0.690000in}}%
\pgfusepath{stroke}%
\end{pgfscope}%
\begin{pgfscope}%
\pgfsetbuttcap%
\pgfsetmiterjoin%
\definecolor{currentfill}{rgb}{1.000000,1.000000,1.000000}%
\pgfsetfillcolor{currentfill}%
\pgfsetlinewidth{0.000000pt}%
\definecolor{currentstroke}{rgb}{0.000000,0.000000,0.000000}%
\pgfsetstrokecolor{currentstroke}%
\pgfsetstrokeopacity{0.000000}%
\pgfsetdash{}{0pt}%
\pgfpathmoveto{\pgfqpoint{2.660000in}{0.690000in}}%
\pgfpathlineto{\pgfqpoint{4.620000in}{0.690000in}}%
\pgfpathlineto{\pgfqpoint{4.620000in}{2.116000in}}%
\pgfpathlineto{\pgfqpoint{2.660000in}{2.116000in}}%
\pgfpathlineto{\pgfqpoint{2.660000in}{0.690000in}}%
\pgfpathclose%
\pgfusepath{fill}%
\end{pgfscope}%
\begin{pgfscope}%
\definecolor{textcolor}{rgb}{0.101961,0.101961,0.101961}%
\pgfsetstrokecolor{textcolor}%
\pgfsetfillcolor{textcolor}%
\pgftext[x=3.640000in,y=0.662222in,,top]{\color{textcolor}\sffamily\fontsize{7.500000}{9.000000}\selectfont Single-shot: answer latency (s)}%
\end{pgfscope}%
\begin{pgfscope}%
\pgfpathrectangle{\pgfqpoint{2.660000in}{0.690000in}}{\pgfqpoint{1.960000in}{1.426000in}}%
\pgfusepath{clip}%
\pgfsetrectcap%
\pgfsetroundjoin%
\pgfsetlinewidth{0.401500pt}%
\definecolor{currentstroke}{rgb}{0.901961,0.901961,0.901961}%
\pgfsetstrokecolor{currentstroke}%
\pgfsetdash{}{0pt}%
\pgfpathmoveto{\pgfqpoint{2.660000in}{0.690000in}}%
\pgfpathlineto{\pgfqpoint{4.620000in}{0.690000in}}%
\pgfusepath{stroke}%
\end{pgfscope}%
\begin{pgfscope}%
\pgfpathrectangle{\pgfqpoint{2.660000in}{0.690000in}}{\pgfqpoint{1.960000in}{1.426000in}}%
\pgfusepath{clip}%
\pgfsetrectcap%
\pgfsetroundjoin%
\pgfsetlinewidth{0.401500pt}%
\definecolor{currentstroke}{rgb}{0.901961,0.901961,0.901961}%
\pgfsetstrokecolor{currentstroke}%
\pgfsetdash{}{0pt}%
\pgfpathmoveto{\pgfqpoint{2.660000in}{1.403000in}}%
\pgfpathlineto{\pgfqpoint{4.620000in}{1.403000in}}%
\pgfusepath{stroke}%
\end{pgfscope}%
\begin{pgfscope}%
\pgfpathrectangle{\pgfqpoint{2.660000in}{0.690000in}}{\pgfqpoint{1.960000in}{1.426000in}}%
\pgfusepath{clip}%
\pgfsetrectcap%
\pgfsetroundjoin%
\pgfsetlinewidth{0.401500pt}%
\definecolor{currentstroke}{rgb}{0.901961,0.901961,0.901961}%
\pgfsetstrokecolor{currentstroke}%
\pgfsetdash{}{0pt}%
\pgfpathmoveto{\pgfqpoint{2.660000in}{2.116000in}}%
\pgfpathlineto{\pgfqpoint{4.620000in}{2.116000in}}%
\pgfusepath{stroke}%
\end{pgfscope}%
\begin{pgfscope}%
\pgfsetrectcap%
\pgfsetmiterjoin%
\pgfsetlinewidth{0.803000pt}%
\definecolor{currentstroke}{rgb}{0.101961,0.101961,0.101961}%
\pgfsetstrokecolor{currentstroke}%
\pgfsetdash{}{0pt}%
\pgfpathmoveto{\pgfqpoint{2.660000in}{0.690000in}}%
\pgfpathlineto{\pgfqpoint{2.660000in}{2.116000in}}%
\pgfusepath{stroke}%
\end{pgfscope}%
\begin{pgfscope}%
\pgfsetrectcap%
\pgfsetmiterjoin%
\pgfsetlinewidth{0.803000pt}%
\definecolor{currentstroke}{rgb}{0.101961,0.101961,0.101961}%
\pgfsetstrokecolor{currentstroke}%
\pgfsetdash{}{0pt}%
\pgfpathmoveto{\pgfqpoint{2.660000in}{0.690000in}}%
\pgfpathlineto{\pgfqpoint{4.620000in}{0.690000in}}%
\pgfusepath{stroke}%
\end{pgfscope}%
\begin{pgfscope}%
\pgfsetbuttcap%
\pgfsetmiterjoin%
\definecolor{currentfill}{rgb}{1.000000,1.000000,1.000000}%
\pgfsetfillcolor{currentfill}%
\pgfsetlinewidth{0.000000pt}%
\definecolor{currentstroke}{rgb}{0.000000,0.000000,0.000000}%
\pgfsetstrokecolor{currentstroke}%
\pgfsetstrokeopacity{0.000000}%
\pgfsetdash{}{0pt}%
\pgfpathmoveto{\pgfqpoint{4.830000in}{0.690000in}}%
\pgfpathlineto{\pgfqpoint{6.790000in}{0.690000in}}%
\pgfpathlineto{\pgfqpoint{6.790000in}{2.116000in}}%
\pgfpathlineto{\pgfqpoint{4.830000in}{2.116000in}}%
\pgfpathlineto{\pgfqpoint{4.830000in}{0.690000in}}%
\pgfpathclose%
\pgfusepath{fill}%
\end{pgfscope}%
\begin{pgfscope}%
\definecolor{textcolor}{rgb}{0.101961,0.101961,0.101961}%
\pgfsetstrokecolor{textcolor}%
\pgfsetfillcolor{textcolor}%
\pgftext[x=5.810000in,y=0.662222in,,top]{\color{textcolor}\sffamily\fontsize{7.500000}{9.000000}\selectfont Dashcam: answer latency (s)}%
\end{pgfscope}%
\begin{pgfscope}%
\pgfpathrectangle{\pgfqpoint{4.830000in}{0.690000in}}{\pgfqpoint{1.960000in}{1.426000in}}%
\pgfusepath{clip}%
\pgfsetrectcap%
\pgfsetroundjoin%
\pgfsetlinewidth{0.401500pt}%
\definecolor{currentstroke}{rgb}{0.901961,0.901961,0.901961}%
\pgfsetstrokecolor{currentstroke}%
\pgfsetdash{}{0pt}%
\pgfpathmoveto{\pgfqpoint{4.830000in}{0.690000in}}%
\pgfpathlineto{\pgfqpoint{6.790000in}{0.690000in}}%
\pgfusepath{stroke}%
\end{pgfscope}%
\begin{pgfscope}%
\pgfpathrectangle{\pgfqpoint{4.830000in}{0.690000in}}{\pgfqpoint{1.960000in}{1.426000in}}%
\pgfusepath{clip}%
\pgfsetrectcap%
\pgfsetroundjoin%
\pgfsetlinewidth{0.401500pt}%
\definecolor{currentstroke}{rgb}{0.901961,0.901961,0.901961}%
\pgfsetstrokecolor{currentstroke}%
\pgfsetdash{}{0pt}%
\pgfpathmoveto{\pgfqpoint{4.830000in}{1.403000in}}%
\pgfpathlineto{\pgfqpoint{6.790000in}{1.403000in}}%
\pgfusepath{stroke}%
\end{pgfscope}%
\begin{pgfscope}%
\pgfpathrectangle{\pgfqpoint{4.830000in}{0.690000in}}{\pgfqpoint{1.960000in}{1.426000in}}%
\pgfusepath{clip}%
\pgfsetrectcap%
\pgfsetroundjoin%
\pgfsetlinewidth{0.401500pt}%
\definecolor{currentstroke}{rgb}{0.901961,0.901961,0.901961}%
\pgfsetstrokecolor{currentstroke}%
\pgfsetdash{}{0pt}%
\pgfpathmoveto{\pgfqpoint{4.830000in}{2.116000in}}%
\pgfpathlineto{\pgfqpoint{6.790000in}{2.116000in}}%
\pgfusepath{stroke}%
\end{pgfscope}%
\begin{pgfscope}%
\pgfsetrectcap%
\pgfsetmiterjoin%
\pgfsetlinewidth{0.803000pt}%
\definecolor{currentstroke}{rgb}{0.101961,0.101961,0.101961}%
\pgfsetstrokecolor{currentstroke}%
\pgfsetdash{}{0pt}%
\pgfpathmoveto{\pgfqpoint{4.830000in}{0.690000in}}%
\pgfpathlineto{\pgfqpoint{4.830000in}{2.116000in}}%
\pgfusepath{stroke}%
\end{pgfscope}%
\begin{pgfscope}%
\pgfsetrectcap%
\pgfsetmiterjoin%
\pgfsetlinewidth{0.803000pt}%
\definecolor{currentstroke}{rgb}{0.101961,0.101961,0.101961}%
\pgfsetstrokecolor{currentstroke}%
\pgfsetdash{}{0pt}%
\pgfpathmoveto{\pgfqpoint{4.830000in}{0.690000in}}%
\pgfpathlineto{\pgfqpoint{6.790000in}{0.690000in}}%
\pgfusepath{stroke}%
\end{pgfscope}%
\begin{pgfscope}%
\definecolor{textcolor}{rgb}{0.549020,0.549020,0.549020}%
\pgfsetstrokecolor{textcolor}%
\pgfsetfillcolor{textcolor}%
\pgftext[x=3.500000in,y=1.150000in,,]{\color{textcolor}\sffamily\fontsize{7.000000}{8.400000}\selectfont NO DATA YET: delivered-answer latency field and device-only rows in F3}%
\end{pgfscope}%
\begin{pgfscope}%
\pgfsetrectcap%
\pgfsetroundjoin%
\pgfsetlinewidth{1.505625pt}%
\definecolor{currentstroke}{rgb}{0.156863,0.356863,0.568627}%
\pgfsetstrokecolor{currentstroke}%
\pgfsetdash{}{0pt}%
\pgfpathmoveto{\pgfqpoint{1.296109in}{0.134333in}}%
\pgfpathlineto{\pgfqpoint{1.373887in}{0.134333in}}%
\pgfpathlineto{\pgfqpoint{1.451664in}{0.134333in}}%
\pgfusepath{stroke}%
\end{pgfscope}%
\begin{pgfscope}%
\definecolor{textcolor}{rgb}{0.101961,0.101961,0.101961}%
\pgfsetstrokecolor{textcolor}%
\pgfsetfillcolor{textcolor}%
\pgftext[x=1.529442in,y=0.100306in,left,base]{\color{textcolor}\sffamily\fontsize{7.000000}{8.400000}\selectfont COVEX (old: rung 3)}%
\end{pgfscope}%
\begin{pgfscope}%
\pgfsetrectcap%
\pgfsetroundjoin%
\pgfsetlinewidth{1.505625pt}%
\definecolor{currentstroke}{rgb}{0.478431,0.305882,0.611765}%
\pgfsetstrokecolor{currentstroke}%
\pgfsetdash{}{0pt}%
\pgfpathmoveto{\pgfqpoint{2.613077in}{0.134333in}}%
\pgfpathlineto{\pgfqpoint{2.690854in}{0.134333in}}%
\pgfpathlineto{\pgfqpoint{2.768632in}{0.134333in}}%
\pgfusepath{stroke}%
\end{pgfscope}%
\begin{pgfscope}%
\definecolor{textcolor}{rgb}{0.101961,0.101961,0.101961}%
\pgfsetstrokecolor{textcolor}%
\pgfsetfillcolor{textcolor}%
\pgftext[x=2.846410in,y=0.100306in,left,base]{\color{textcolor}\sffamily\fontsize{7.000000}{8.400000}\selectfont On-device only}%
\end{pgfscope}%
\begin{pgfscope}%
\pgfsetrectcap%
\pgfsetroundjoin%
\pgfsetlinewidth{1.505625pt}%
\definecolor{currentstroke}{rgb}{0.478431,0.478431,0.478431}%
\pgfsetstrokecolor{currentstroke}%
\pgfsetdash{}{0pt}%
\pgfpathmoveto{\pgfqpoint{3.631285in}{0.134333in}}%
\pgfpathlineto{\pgfqpoint{3.709063in}{0.134333in}}%
\pgfpathlineto{\pgfqpoint{3.786840in}{0.134333in}}%
\pgfusepath{stroke}%
\end{pgfscope}%
\begin{pgfscope}%
\definecolor{textcolor}{rgb}{0.101961,0.101961,0.101961}%
\pgfsetstrokecolor{textcolor}%
\pgfsetfillcolor{textcolor}%
\pgftext[x=3.864618in,y=0.100306in,left,base]{\color{textcolor}\sffamily\fontsize{7.000000}{8.400000}\selectfont Three strongest baselines}%
\end{pgfscope}%
\begin{pgfscope}%
\pgfsetbuttcap%
\pgfsetroundjoin%
\pgfsetlinewidth{0.803000pt}%
\definecolor{currentstroke}{rgb}{0.101961,0.101961,0.101961}%
\pgfsetstrokecolor{currentstroke}%
\pgfsetdash{{2.960000pt}{1.280000pt}}{0.000000pt}%
\pgfpathmoveto{\pgfqpoint{5.098543in}{0.134333in}}%
\pgfpathlineto{\pgfqpoint{5.176321in}{0.134333in}}%
\pgfpathlineto{\pgfqpoint{5.254099in}{0.134333in}}%
\pgfusepath{stroke}%
\end{pgfscope}%
\begin{pgfscope}%
\definecolor{textcolor}{rgb}{0.101961,0.101961,0.101961}%
\pgfsetstrokecolor{textcolor}%
\pgfsetfillcolor{textcolor}%
\pgftext[x=5.331877in,y=0.100306in,left,base]{\color{textcolor}\sffamily\fontsize{7.000000}{8.400000}\selectfont Deadline}%
\end{pgfscope}%
\end{pgfpicture}%
\makeatother%
\endgroup%
}
\caption{[PLACEHOLDER: NO DATA, NO DATA YET: delivered-answer latency field and device-only rows in F3] Latency CDFs per question type with the deadline marked, for COVEX, the on-device model and the three strongest baselines. Not drawn: the committed latency table has no delivered-answer latency field and no on-device rows.}
\label{fig:F3_latency_cdf}
\end{figure*}
\clearpage
\begin{figure}[t]
\centering
{\fontfamily{LinuxBiolinumT}\selectfont
%% Creator: Matplotlib, PGF backend
%%
%% To include the figure in your LaTeX document, write
%%   \input{<filename>.pgf}
%%
%% Make sure the required packages are loaded in your preamble
%%   \usepackage{pgf}
%%
%% Also ensure that all the required font packages are loaded; for instance,
%% the lmodern package is sometimes necessary when using math font.
%%   \usepackage{lmodern}
%%
%% Figures using additional raster images can only be included by \input if
%% they are in the same directory as the main LaTeX file. For loading figures
%% from other directories you can use the `import` package
%%   \usepackage{import}
%%
%% and then include the figures with
%%   \import{<path to file>}{<filename>.pgf}
%%
%% Matplotlib used the following preamble
%%
\begingroup%
\makeatletter%
\begin{pgfpicture}%
\pgfpathrectangle{\pgfpointorigin}{\pgfqpoint{3.330000in}{2.200000in}}%
\pgfusepath{use as bounding box, clip}%
\begin{pgfscope}%
\pgfsetbuttcap%
\pgfsetmiterjoin%
\definecolor{currentfill}{rgb}{1.000000,1.000000,1.000000}%
\pgfsetfillcolor{currentfill}%
\pgfsetlinewidth{0.000000pt}%
\definecolor{currentstroke}{rgb}{1.000000,1.000000,1.000000}%
\pgfsetstrokecolor{currentstroke}%
\pgfsetdash{}{0pt}%
\pgfpathmoveto{\pgfqpoint{0.000000in}{0.000000in}}%
\pgfpathlineto{\pgfqpoint{3.330000in}{0.000000in}}%
\pgfpathlineto{\pgfqpoint{3.330000in}{2.200000in}}%
\pgfpathlineto{\pgfqpoint{0.000000in}{2.200000in}}%
\pgfpathlineto{\pgfqpoint{0.000000in}{0.000000in}}%
\pgfpathclose%
\pgfusepath{fill}%
\end{pgfscope}%
\begin{pgfscope}%
\pgfsetbuttcap%
\pgfsetmiterjoin%
\definecolor{currentfill}{rgb}{1.000000,1.000000,1.000000}%
\pgfsetfillcolor{currentfill}%
\pgfsetlinewidth{0.000000pt}%
\definecolor{currentstroke}{rgb}{0.000000,0.000000,0.000000}%
\pgfsetstrokecolor{currentstroke}%
\pgfsetstrokeopacity{0.000000}%
\pgfsetdash{}{0pt}%
\pgfpathmoveto{\pgfqpoint{1.248750in}{0.440000in}}%
\pgfpathlineto{\pgfqpoint{3.230100in}{0.440000in}}%
\pgfpathlineto{\pgfqpoint{3.230100in}{2.068000in}}%
\pgfpathlineto{\pgfqpoint{1.248750in}{2.068000in}}%
\pgfpathlineto{\pgfqpoint{1.248750in}{0.440000in}}%
\pgfpathclose%
\pgfusepath{fill}%
\end{pgfscope}%
\begin{pgfscope}%
\pgfpathrectangle{\pgfqpoint{1.248750in}{0.440000in}}{\pgfqpoint{1.981350in}{1.628000in}}%
\pgfusepath{clip}%
\pgfsetrectcap%
\pgfsetroundjoin%
\pgfsetlinewidth{0.401500pt}%
\definecolor{currentstroke}{rgb}{0.901961,0.901961,0.901961}%
\pgfsetstrokecolor{currentstroke}%
\pgfsetdash{}{0pt}%
\pgfpathmoveto{\pgfqpoint{1.457313in}{0.440000in}}%
\pgfpathlineto{\pgfqpoint{1.457313in}{2.068000in}}%
\pgfusepath{stroke}%
\end{pgfscope}%
\begin{pgfscope}%
\pgfsetbuttcap%
\pgfsetroundjoin%
\definecolor{currentfill}{rgb}{0.101961,0.101961,0.101961}%
\pgfsetfillcolor{currentfill}%
\pgfsetlinewidth{0.803000pt}%
\definecolor{currentstroke}{rgb}{0.101961,0.101961,0.101961}%
\pgfsetstrokecolor{currentstroke}%
\pgfsetdash{}{0pt}%
\pgfsys@defobject{currentmarker}{\pgfqpoint{0.000000in}{-0.027778in}}{\pgfqpoint{0.000000in}{0.000000in}}{%
\pgfpathmoveto{\pgfqpoint{0.000000in}{0.000000in}}%
\pgfpathlineto{\pgfqpoint{0.000000in}{-0.027778in}}%
\pgfusepath{stroke,fill}%
}%
\begin{pgfscope}%
\pgfsys@transformshift{1.457313in}{0.440000in}%
\pgfsys@useobject{currentmarker}{}%
\end{pgfscope}%
\end{pgfscope}%
\begin{pgfscope}%
\definecolor{textcolor}{rgb}{0.101961,0.101961,0.101961}%
\pgfsetstrokecolor{textcolor}%
\pgfsetfillcolor{textcolor}%
\pgftext[x=1.457313in,y=0.363611in,,top]{\color{textcolor}\sffamily\fontsize{7.000000}{8.400000}\selectfont 0}%
\end{pgfscope}%
\begin{pgfscope}%
\pgfpathrectangle{\pgfqpoint{1.248750in}{0.440000in}}{\pgfqpoint{1.981350in}{1.628000in}}%
\pgfusepath{clip}%
\pgfsetrectcap%
\pgfsetroundjoin%
\pgfsetlinewidth{0.401500pt}%
\definecolor{currentstroke}{rgb}{0.901961,0.901961,0.901961}%
\pgfsetstrokecolor{currentstroke}%
\pgfsetdash{}{0pt}%
\pgfpathmoveto{\pgfqpoint{1.891820in}{0.440000in}}%
\pgfpathlineto{\pgfqpoint{1.891820in}{2.068000in}}%
\pgfusepath{stroke}%
\end{pgfscope}%
\begin{pgfscope}%
\pgfsetbuttcap%
\pgfsetroundjoin%
\definecolor{currentfill}{rgb}{0.101961,0.101961,0.101961}%
\pgfsetfillcolor{currentfill}%
\pgfsetlinewidth{0.803000pt}%
\definecolor{currentstroke}{rgb}{0.101961,0.101961,0.101961}%
\pgfsetstrokecolor{currentstroke}%
\pgfsetdash{}{0pt}%
\pgfsys@defobject{currentmarker}{\pgfqpoint{0.000000in}{-0.027778in}}{\pgfqpoint{0.000000in}{0.000000in}}{%
\pgfpathmoveto{\pgfqpoint{0.000000in}{0.000000in}}%
\pgfpathlineto{\pgfqpoint{0.000000in}{-0.027778in}}%
\pgfusepath{stroke,fill}%
}%
\begin{pgfscope}%
\pgfsys@transformshift{1.891820in}{0.440000in}%
\pgfsys@useobject{currentmarker}{}%
\end{pgfscope}%
\end{pgfscope}%
\begin{pgfscope}%
\definecolor{textcolor}{rgb}{0.101961,0.101961,0.101961}%
\pgfsetstrokecolor{textcolor}%
\pgfsetfillcolor{textcolor}%
\pgftext[x=1.891820in,y=0.363611in,,top]{\color{textcolor}\sffamily\fontsize{7.000000}{8.400000}\selectfont .25}%
\end{pgfscope}%
\begin{pgfscope}%
\pgfpathrectangle{\pgfqpoint{1.248750in}{0.440000in}}{\pgfqpoint{1.981350in}{1.628000in}}%
\pgfusepath{clip}%
\pgfsetrectcap%
\pgfsetroundjoin%
\pgfsetlinewidth{0.401500pt}%
\definecolor{currentstroke}{rgb}{0.901961,0.901961,0.901961}%
\pgfsetstrokecolor{currentstroke}%
\pgfsetdash{}{0pt}%
\pgfpathmoveto{\pgfqpoint{2.326326in}{0.440000in}}%
\pgfpathlineto{\pgfqpoint{2.326326in}{2.068000in}}%
\pgfusepath{stroke}%
\end{pgfscope}%
\begin{pgfscope}%
\pgfsetbuttcap%
\pgfsetroundjoin%
\definecolor{currentfill}{rgb}{0.101961,0.101961,0.101961}%
\pgfsetfillcolor{currentfill}%
\pgfsetlinewidth{0.803000pt}%
\definecolor{currentstroke}{rgb}{0.101961,0.101961,0.101961}%
\pgfsetstrokecolor{currentstroke}%
\pgfsetdash{}{0pt}%
\pgfsys@defobject{currentmarker}{\pgfqpoint{0.000000in}{-0.027778in}}{\pgfqpoint{0.000000in}{0.000000in}}{%
\pgfpathmoveto{\pgfqpoint{0.000000in}{0.000000in}}%
\pgfpathlineto{\pgfqpoint{0.000000in}{-0.027778in}}%
\pgfusepath{stroke,fill}%
}%
\begin{pgfscope}%
\pgfsys@transformshift{2.326326in}{0.440000in}%
\pgfsys@useobject{currentmarker}{}%
\end{pgfscope}%
\end{pgfscope}%
\begin{pgfscope}%
\definecolor{textcolor}{rgb}{0.101961,0.101961,0.101961}%
\pgfsetstrokecolor{textcolor}%
\pgfsetfillcolor{textcolor}%
\pgftext[x=2.326326in,y=0.363611in,,top]{\color{textcolor}\sffamily\fontsize{7.000000}{8.400000}\selectfont .5}%
\end{pgfscope}%
\begin{pgfscope}%
\pgfpathrectangle{\pgfqpoint{1.248750in}{0.440000in}}{\pgfqpoint{1.981350in}{1.628000in}}%
\pgfusepath{clip}%
\pgfsetrectcap%
\pgfsetroundjoin%
\pgfsetlinewidth{0.401500pt}%
\definecolor{currentstroke}{rgb}{0.901961,0.901961,0.901961}%
\pgfsetstrokecolor{currentstroke}%
\pgfsetdash{}{0pt}%
\pgfpathmoveto{\pgfqpoint{2.760833in}{0.440000in}}%
\pgfpathlineto{\pgfqpoint{2.760833in}{2.068000in}}%
\pgfusepath{stroke}%
\end{pgfscope}%
\begin{pgfscope}%
\pgfsetbuttcap%
\pgfsetroundjoin%
\definecolor{currentfill}{rgb}{0.101961,0.101961,0.101961}%
\pgfsetfillcolor{currentfill}%
\pgfsetlinewidth{0.803000pt}%
\definecolor{currentstroke}{rgb}{0.101961,0.101961,0.101961}%
\pgfsetstrokecolor{currentstroke}%
\pgfsetdash{}{0pt}%
\pgfsys@defobject{currentmarker}{\pgfqpoint{0.000000in}{-0.027778in}}{\pgfqpoint{0.000000in}{0.000000in}}{%
\pgfpathmoveto{\pgfqpoint{0.000000in}{0.000000in}}%
\pgfpathlineto{\pgfqpoint{0.000000in}{-0.027778in}}%
\pgfusepath{stroke,fill}%
}%
\begin{pgfscope}%
\pgfsys@transformshift{2.760833in}{0.440000in}%
\pgfsys@useobject{currentmarker}{}%
\end{pgfscope}%
\end{pgfscope}%
\begin{pgfscope}%
\definecolor{textcolor}{rgb}{0.101961,0.101961,0.101961}%
\pgfsetstrokecolor{textcolor}%
\pgfsetfillcolor{textcolor}%
\pgftext[x=2.760833in,y=0.363611in,,top]{\color{textcolor}\sffamily\fontsize{7.000000}{8.400000}\selectfont .75}%
\end{pgfscope}%
\begin{pgfscope}%
\pgfpathrectangle{\pgfqpoint{1.248750in}{0.440000in}}{\pgfqpoint{1.981350in}{1.628000in}}%
\pgfusepath{clip}%
\pgfsetrectcap%
\pgfsetroundjoin%
\pgfsetlinewidth{0.401500pt}%
\definecolor{currentstroke}{rgb}{0.901961,0.901961,0.901961}%
\pgfsetstrokecolor{currentstroke}%
\pgfsetdash{}{0pt}%
\pgfpathmoveto{\pgfqpoint{3.195339in}{0.440000in}}%
\pgfpathlineto{\pgfqpoint{3.195339in}{2.068000in}}%
\pgfusepath{stroke}%
\end{pgfscope}%
\begin{pgfscope}%
\pgfsetbuttcap%
\pgfsetroundjoin%
\definecolor{currentfill}{rgb}{0.101961,0.101961,0.101961}%
\pgfsetfillcolor{currentfill}%
\pgfsetlinewidth{0.803000pt}%
\definecolor{currentstroke}{rgb}{0.101961,0.101961,0.101961}%
\pgfsetstrokecolor{currentstroke}%
\pgfsetdash{}{0pt}%
\pgfsys@defobject{currentmarker}{\pgfqpoint{0.000000in}{-0.027778in}}{\pgfqpoint{0.000000in}{0.000000in}}{%
\pgfpathmoveto{\pgfqpoint{0.000000in}{0.000000in}}%
\pgfpathlineto{\pgfqpoint{0.000000in}{-0.027778in}}%
\pgfusepath{stroke,fill}%
}%
\begin{pgfscope}%
\pgfsys@transformshift{3.195339in}{0.440000in}%
\pgfsys@useobject{currentmarker}{}%
\end{pgfscope}%
\end{pgfscope}%
\begin{pgfscope}%
\definecolor{textcolor}{rgb}{0.101961,0.101961,0.101961}%
\pgfsetstrokecolor{textcolor}%
\pgfsetfillcolor{textcolor}%
\pgftext[x=3.195339in,y=0.363611in,,top]{\color{textcolor}\sffamily\fontsize{7.000000}{8.400000}\selectfont 1}%
\end{pgfscope}%
\begin{pgfscope}%
\definecolor{textcolor}{rgb}{0.101961,0.101961,0.101961}%
\pgfsetstrokecolor{textcolor}%
\pgfsetfillcolor{textcolor}%
\pgftext[x=1.500916in, y=0.181899in, left, base]{\color{textcolor}\sffamily\fontsize{7.500000}{9.000000}\selectfont Dashcam on-time attainment gain}%
\end{pgfscope}%
\begin{pgfscope}%
\definecolor{textcolor}{rgb}{0.101961,0.101961,0.101961}%
\pgfsetstrokecolor{textcolor}%
\pgfsetfillcolor{textcolor}%
\pgftext[x=1.422731in, y=0.075494in, left, base]{\color{textcolor}\sffamily\fontsize{7.500000}{9.000000}\selectfont over the on-device model (per device)}%
\end{pgfscope}%
\begin{pgfscope}%
\pgfpathrectangle{\pgfqpoint{1.248750in}{0.440000in}}{\pgfqpoint{1.981350in}{1.628000in}}%
\pgfusepath{clip}%
\pgfsetbuttcap%
\pgfsetmiterjoin%
\pgfsetlinewidth{0.803000pt}%
\definecolor{currentstroke}{rgb}{0.156863,0.356863,0.568627}%
\pgfsetstrokecolor{currentstroke}%
\pgfsetdash{}{0pt}%
\pgfpathmoveto{\pgfqpoint{2.670336in}{1.795623in}}%
\pgfpathlineto{\pgfqpoint{2.772813in}{1.795623in}}%
\pgfpathlineto{\pgfqpoint{2.772813in}{1.964685in}}%
\pgfpathlineto{\pgfqpoint{2.670336in}{1.964685in}}%
\pgfpathlineto{\pgfqpoint{2.670336in}{1.795623in}}%
\pgfpathclose%
\pgfusepath{stroke}%
\end{pgfscope}%
\begin{pgfscope}%
\pgfpathrectangle{\pgfqpoint{1.248750in}{0.440000in}}{\pgfqpoint{1.981350in}{1.628000in}}%
\pgfusepath{clip}%
\pgfsetbuttcap%
\pgfsetmiterjoin%
\pgfsetlinewidth{0.803000pt}%
\definecolor{currentstroke}{rgb}{0.784314,0.784314,0.784314}%
\pgfsetstrokecolor{currentstroke}%
\pgfsetdash{}{0pt}%
\pgfpathmoveto{\pgfqpoint{3.073647in}{1.482546in}}%
\pgfpathlineto{\pgfqpoint{3.154105in}{1.482546in}}%
\pgfpathlineto{\pgfqpoint{3.154105in}{1.651608in}}%
\pgfpathlineto{\pgfqpoint{3.073647in}{1.651608in}}%
\pgfpathlineto{\pgfqpoint{3.073647in}{1.482546in}}%
\pgfpathclose%
\pgfusepath{stroke}%
\end{pgfscope}%
\begin{pgfscope}%
\pgfpathrectangle{\pgfqpoint{1.248750in}{0.440000in}}{\pgfqpoint{1.981350in}{1.628000in}}%
\pgfusepath{clip}%
\pgfsetbuttcap%
\pgfsetmiterjoin%
\pgfsetlinewidth{0.803000pt}%
\definecolor{currentstroke}{rgb}{0.784314,0.784314,0.784314}%
\pgfsetstrokecolor{currentstroke}%
\pgfsetdash{}{0pt}%
\pgfpathmoveto{\pgfqpoint{2.763807in}{1.169469in}}%
\pgfpathlineto{\pgfqpoint{2.863890in}{1.169469in}}%
\pgfpathlineto{\pgfqpoint{2.863890in}{1.338531in}}%
\pgfpathlineto{\pgfqpoint{2.763807in}{1.338531in}}%
\pgfpathlineto{\pgfqpoint{2.763807in}{1.169469in}}%
\pgfpathclose%
\pgfusepath{stroke}%
\end{pgfscope}%
\begin{pgfscope}%
\pgfpathrectangle{\pgfqpoint{1.248750in}{0.440000in}}{\pgfqpoint{1.981350in}{1.628000in}}%
\pgfusepath{clip}%
\pgfsetbuttcap%
\pgfsetmiterjoin%
\pgfsetlinewidth{0.803000pt}%
\definecolor{currentstroke}{rgb}{0.784314,0.784314,0.784314}%
\pgfsetstrokecolor{currentstroke}%
\pgfsetdash{}{0pt}%
\pgfpathmoveto{\pgfqpoint{2.491024in}{0.856392in}}%
\pgfpathlineto{\pgfqpoint{2.603840in}{0.856392in}}%
\pgfpathlineto{\pgfqpoint{2.603840in}{1.025454in}}%
\pgfpathlineto{\pgfqpoint{2.491024in}{1.025454in}}%
\pgfpathlineto{\pgfqpoint{2.491024in}{0.856392in}}%
\pgfpathclose%
\pgfusepath{stroke}%
\end{pgfscope}%
\begin{pgfscope}%
\pgfpathrectangle{\pgfqpoint{1.248750in}{0.440000in}}{\pgfqpoint{1.981350in}{1.628000in}}%
\pgfusepath{clip}%
\pgfsetbuttcap%
\pgfsetmiterjoin%
\pgfsetlinewidth{0.803000pt}%
\definecolor{currentstroke}{rgb}{0.784314,0.784314,0.784314}%
\pgfsetstrokecolor{currentstroke}%
\pgfsetdash{}{0pt}%
\pgfpathmoveto{\pgfqpoint{1.584040in}{0.543315in}}%
\pgfpathlineto{\pgfqpoint{2.766524in}{0.543315in}}%
\pgfpathlineto{\pgfqpoint{2.766524in}{0.712377in}}%
\pgfpathlineto{\pgfqpoint{1.584040in}{0.712377in}}%
\pgfpathlineto{\pgfqpoint{1.584040in}{0.543315in}}%
\pgfpathclose%
\pgfusepath{stroke}%
\end{pgfscope}%
\begin{pgfscope}%
\pgfsetrectcap%
\pgfsetmiterjoin%
\pgfsetlinewidth{0.803000pt}%
\definecolor{currentstroke}{rgb}{0.101961,0.101961,0.101961}%
\pgfsetstrokecolor{currentstroke}%
\pgfsetdash{}{0pt}%
\pgfpathmoveto{\pgfqpoint{1.248750in}{0.440000in}}%
\pgfpathlineto{\pgfqpoint{1.248750in}{2.068000in}}%
\pgfusepath{stroke}%
\end{pgfscope}%
\begin{pgfscope}%
\pgfsetrectcap%
\pgfsetmiterjoin%
\pgfsetlinewidth{0.803000pt}%
\definecolor{currentstroke}{rgb}{0.101961,0.101961,0.101961}%
\pgfsetstrokecolor{currentstroke}%
\pgfsetdash{}{0pt}%
\pgfpathmoveto{\pgfqpoint{1.248750in}{0.440000in}}%
\pgfpathlineto{\pgfqpoint{3.230100in}{0.440000in}}%
\pgfusepath{stroke}%
\end{pgfscope}%
\begin{pgfscope}%
\pgfpathrectangle{\pgfqpoint{1.248750in}{0.440000in}}{\pgfqpoint{1.981350in}{1.628000in}}%
\pgfusepath{clip}%
\pgfsetrectcap%
\pgfsetroundjoin%
\pgfsetlinewidth{1.003750pt}%
\definecolor{currentstroke}{rgb}{0.101961,0.101961,0.101961}%
\pgfsetstrokecolor{currentstroke}%
\pgfsetdash{}{0pt}%
\pgfpathmoveto{\pgfqpoint{2.686531in}{1.795623in}}%
\pgfpathlineto{\pgfqpoint{2.686531in}{1.964685in}}%
\pgfusepath{stroke}%
\end{pgfscope}%
\begin{pgfscope}%
\pgfpathrectangle{\pgfqpoint{1.248750in}{0.440000in}}{\pgfqpoint{1.981350in}{1.628000in}}%
\pgfusepath{clip}%
\pgfsetbuttcap%
\pgfsetroundjoin%
\definecolor{currentfill}{rgb}{0.156863,0.356863,0.568627}%
\pgfsetfillcolor{currentfill}%
\pgfsetfillopacity{0.900000}%
\pgfsetlinewidth{0.200750pt}%
\definecolor{currentstroke}{rgb}{0.101961,0.101961,0.101961}%
\pgfsetstrokecolor{currentstroke}%
\pgfsetstrokeopacity{0.900000}%
\pgfsetdash{}{0pt}%
\pgfsys@defobject{currentmarker}{\pgfqpoint{-0.015528in}{-0.015528in}}{\pgfqpoint{0.015528in}{0.015528in}}{%
\pgfpathmoveto{\pgfqpoint{0.000000in}{-0.015528in}}%
\pgfpathcurveto{\pgfqpoint{0.004118in}{-0.015528in}}{\pgfqpoint{0.008068in}{-0.013892in}}{\pgfqpoint{0.010980in}{-0.010980in}}%
\pgfpathcurveto{\pgfqpoint{0.013892in}{-0.008068in}}{\pgfqpoint{0.015528in}{-0.004118in}}{\pgfqpoint{0.015528in}{0.000000in}}%
\pgfpathcurveto{\pgfqpoint{0.015528in}{0.004118in}}{\pgfqpoint{0.013892in}{0.008068in}}{\pgfqpoint{0.010980in}{0.010980in}}%
\pgfpathcurveto{\pgfqpoint{0.008068in}{0.013892in}}{\pgfqpoint{0.004118in}{0.015528in}}{\pgfqpoint{0.000000in}{0.015528in}}%
\pgfpathcurveto{\pgfqpoint{-0.004118in}{0.015528in}}{\pgfqpoint{-0.008068in}{0.013892in}}{\pgfqpoint{-0.010980in}{0.010980in}}%
\pgfpathcurveto{\pgfqpoint{-0.013892in}{0.008068in}}{\pgfqpoint{-0.015528in}{0.004118in}}{\pgfqpoint{-0.015528in}{0.000000in}}%
\pgfpathcurveto{\pgfqpoint{-0.015528in}{-0.004118in}}{\pgfqpoint{-0.013892in}{-0.008068in}}{\pgfqpoint{-0.010980in}{-0.010980in}}%
\pgfpathcurveto{\pgfqpoint{-0.008068in}{-0.013892in}}{\pgfqpoint{-0.004118in}{-0.015528in}}{\pgfqpoint{0.000000in}{-0.015528in}}%
\pgfpathlineto{\pgfqpoint{0.000000in}{-0.015528in}}%
\pgfpathclose%
\pgfusepath{stroke,fill}%
}%
\begin{pgfscope}%
\pgfsys@transformshift{2.661390in}{1.881412in}%
\pgfsys@useobject{currentmarker}{}%
\end{pgfscope}%
\begin{pgfscope}%
\pgfsys@transformshift{2.644861in}{1.928104in}%
\pgfsys@useobject{currentmarker}{}%
\end{pgfscope}%
\begin{pgfscope}%
\pgfsys@transformshift{2.638497in}{1.842276in}%
\pgfsys@useobject{currentmarker}{}%
\end{pgfscope}%
\begin{pgfscope}%
\pgfsys@transformshift{2.681450in}{1.927911in}%
\pgfsys@useobject{currentmarker}{}%
\end{pgfscope}%
\begin{pgfscope}%
\pgfsys@transformshift{2.764995in}{1.860124in}%
\pgfsys@useobject{currentmarker}{}%
\end{pgfscope}%
\begin{pgfscope}%
\pgfsys@transformshift{2.777751in}{1.871992in}%
\pgfsys@useobject{currentmarker}{}%
\end{pgfscope}%
\begin{pgfscope}%
\pgfsys@transformshift{2.680925in}{1.915037in}%
\pgfsys@useobject{currentmarker}{}%
\end{pgfscope}%
\begin{pgfscope}%
\pgfsys@transformshift{2.684363in}{1.870488in}%
\pgfsys@useobject{currentmarker}{}%
\end{pgfscope}%
\begin{pgfscope}%
\pgfsys@transformshift{2.738807in}{1.885433in}%
\pgfsys@useobject{currentmarker}{}%
\end{pgfscope}%
\begin{pgfscope}%
\pgfsys@transformshift{2.673318in}{1.829864in}%
\pgfsys@useobject{currentmarker}{}%
\end{pgfscope}%
\begin{pgfscope}%
\pgfsys@transformshift{2.645737in}{1.907139in}%
\pgfsys@useobject{currentmarker}{}%
\end{pgfscope}%
\begin{pgfscope}%
\pgfsys@transformshift{2.684938in}{1.884214in}%
\pgfsys@useobject{currentmarker}{}%
\end{pgfscope}%
\begin{pgfscope}%
\pgfsys@transformshift{2.679363in}{1.862029in}%
\pgfsys@useobject{currentmarker}{}%
\end{pgfscope}%
\begin{pgfscope}%
\pgfsys@transformshift{2.713685in}{1.910856in}%
\pgfsys@useobject{currentmarker}{}%
\end{pgfscope}%
\begin{pgfscope}%
\pgfsys@transformshift{2.625727in}{1.859205in}%
\pgfsys@useobject{currentmarker}{}%
\end{pgfscope}%
\begin{pgfscope}%
\pgfsys@transformshift{2.780634in}{1.875204in}%
\pgfsys@useobject{currentmarker}{}%
\end{pgfscope}%
\begin{pgfscope}%
\pgfsys@transformshift{2.718044in}{1.841199in}%
\pgfsys@useobject{currentmarker}{}%
\end{pgfscope}%
\begin{pgfscope}%
\pgfsys@transformshift{2.790007in}{1.869841in}%
\pgfsys@useobject{currentmarker}{}%
\end{pgfscope}%
\begin{pgfscope}%
\pgfsys@transformshift{2.792345in}{1.848588in}%
\pgfsys@useobject{currentmarker}{}%
\end{pgfscope}%
\begin{pgfscope}%
\pgfsys@transformshift{2.688123in}{1.854853in}%
\pgfsys@useobject{currentmarker}{}%
\end{pgfscope}%
\begin{pgfscope}%
\pgfsys@transformshift{2.783678in}{1.906804in}%
\pgfsys@useobject{currentmarker}{}%
\end{pgfscope}%
\begin{pgfscope}%
\pgfsys@transformshift{2.836295in}{1.856779in}%
\pgfsys@useobject{currentmarker}{}%
\end{pgfscope}%
\begin{pgfscope}%
\pgfsys@transformshift{2.771167in}{1.878577in}%
\pgfsys@useobject{currentmarker}{}%
\end{pgfscope}%
\begin{pgfscope}%
\pgfsys@transformshift{2.659544in}{1.931326in}%
\pgfsys@useobject{currentmarker}{}%
\end{pgfscope}%
\end{pgfscope}%
\begin{pgfscope}%
\definecolor{textcolor}{rgb}{0.101961,0.101961,0.101961}%
\pgfsetstrokecolor{textcolor}%
\pgfsetfillcolor{textcolor}%
\pgftext[x=2.582277in,y=1.880154in,left,]{\makebox[0pt][r]{\color{textcolor}\sffamily\fontsize{5.500000}{6.600000}\selectfont 0.67}}%
\end{pgfscope}%
\begin{pgfscope}%
\pgfpathrectangle{\pgfqpoint{1.248750in}{0.440000in}}{\pgfqpoint{1.981350in}{1.628000in}}%
\pgfusepath{clip}%
\pgfsetrectcap%
\pgfsetroundjoin%
\pgfsetlinewidth{1.003750pt}%
\definecolor{currentstroke}{rgb}{0.101961,0.101961,0.101961}%
\pgfsetstrokecolor{currentstroke}%
\pgfsetdash{}{0pt}%
\pgfpathmoveto{\pgfqpoint{3.124032in}{1.482546in}}%
\pgfpathlineto{\pgfqpoint{3.124032in}{1.651608in}}%
\pgfusepath{stroke}%
\end{pgfscope}%
\begin{pgfscope}%
\pgfpathrectangle{\pgfqpoint{1.248750in}{0.440000in}}{\pgfqpoint{1.981350in}{1.628000in}}%
\pgfusepath{clip}%
\pgfsetbuttcap%
\pgfsetroundjoin%
\definecolor{currentfill}{rgb}{0.784314,0.784314,0.784314}%
\pgfsetfillcolor{currentfill}%
\pgfsetfillopacity{0.900000}%
\pgfsetlinewidth{0.200750pt}%
\definecolor{currentstroke}{rgb}{0.101961,0.101961,0.101961}%
\pgfsetstrokecolor{currentstroke}%
\pgfsetstrokeopacity{0.900000}%
\pgfsetdash{}{0pt}%
\pgfsys@defobject{currentmarker}{\pgfqpoint{-0.015528in}{-0.015528in}}{\pgfqpoint{0.015528in}{0.015528in}}{%
\pgfpathmoveto{\pgfqpoint{0.000000in}{-0.015528in}}%
\pgfpathcurveto{\pgfqpoint{0.004118in}{-0.015528in}}{\pgfqpoint{0.008068in}{-0.013892in}}{\pgfqpoint{0.010980in}{-0.010980in}}%
\pgfpathcurveto{\pgfqpoint{0.013892in}{-0.008068in}}{\pgfqpoint{0.015528in}{-0.004118in}}{\pgfqpoint{0.015528in}{0.000000in}}%
\pgfpathcurveto{\pgfqpoint{0.015528in}{0.004118in}}{\pgfqpoint{0.013892in}{0.008068in}}{\pgfqpoint{0.010980in}{0.010980in}}%
\pgfpathcurveto{\pgfqpoint{0.008068in}{0.013892in}}{\pgfqpoint{0.004118in}{0.015528in}}{\pgfqpoint{0.000000in}{0.015528in}}%
\pgfpathcurveto{\pgfqpoint{-0.004118in}{0.015528in}}{\pgfqpoint{-0.008068in}{0.013892in}}{\pgfqpoint{-0.010980in}{0.010980in}}%
\pgfpathcurveto{\pgfqpoint{-0.013892in}{0.008068in}}{\pgfqpoint{-0.015528in}{0.004118in}}{\pgfqpoint{-0.015528in}{0.000000in}}%
\pgfpathcurveto{\pgfqpoint{-0.015528in}{-0.004118in}}{\pgfqpoint{-0.013892in}{-0.008068in}}{\pgfqpoint{-0.010980in}{-0.010980in}}%
\pgfpathcurveto{\pgfqpoint{-0.008068in}{-0.013892in}}{\pgfqpoint{-0.004118in}{-0.015528in}}{\pgfqpoint{0.000000in}{-0.015528in}}%
\pgfpathlineto{\pgfqpoint{0.000000in}{-0.015528in}}%
\pgfpathclose%
\pgfusepath{stroke,fill}%
}%
\begin{pgfscope}%
\pgfsys@transformshift{3.107543in}{1.568335in}%
\pgfsys@useobject{currentmarker}{}%
\end{pgfscope}%
\begin{pgfscope}%
\pgfsys@transformshift{3.089308in}{1.615027in}%
\pgfsys@useobject{currentmarker}{}%
\end{pgfscope}%
\begin{pgfscope}%
\pgfsys@transformshift{3.078998in}{1.529199in}%
\pgfsys@useobject{currentmarker}{}%
\end{pgfscope}%
\begin{pgfscope}%
\pgfsys@transformshift{3.094508in}{1.614834in}%
\pgfsys@useobject{currentmarker}{}%
\end{pgfscope}%
\begin{pgfscope}%
\pgfsys@transformshift{3.098928in}{1.547047in}%
\pgfsys@useobject{currentmarker}{}%
\end{pgfscope}%
\begin{pgfscope}%
\pgfsys@transformshift{3.092945in}{1.558915in}%
\pgfsys@useobject{currentmarker}{}%
\end{pgfscope}%
\begin{pgfscope}%
\pgfsys@transformshift{3.049250in}{1.601960in}%
\pgfsys@useobject{currentmarker}{}%
\end{pgfscope}%
\begin{pgfscope}%
\pgfsys@transformshift{3.047705in}{1.557412in}%
\pgfsys@useobject{currentmarker}{}%
\end{pgfscope}%
\begin{pgfscope}%
\pgfsys@transformshift{3.043161in}{1.572356in}%
\pgfsys@useobject{currentmarker}{}%
\end{pgfscope}%
\begin{pgfscope}%
\pgfsys@transformshift{3.057596in}{1.516787in}%
\pgfsys@useobject{currentmarker}{}%
\end{pgfscope}%
\begin{pgfscope}%
\pgfsys@transformshift{3.043664in}{1.594062in}%
\pgfsys@useobject{currentmarker}{}%
\end{pgfscope}%
\begin{pgfscope}%
\pgfsys@transformshift{3.038225in}{1.571137in}%
\pgfsys@useobject{currentmarker}{}%
\end{pgfscope}%
\begin{pgfscope}%
\pgfsys@transformshift{3.169059in}{1.548953in}%
\pgfsys@useobject{currentmarker}{}%
\end{pgfscope}%
\begin{pgfscope}%
\pgfsys@transformshift{3.155497in}{1.597779in}%
\pgfsys@useobject{currentmarker}{}%
\end{pgfscope}%
\begin{pgfscope}%
\pgfsys@transformshift{3.159960in}{1.546128in}%
\pgfsys@useobject{currentmarker}{}%
\end{pgfscope}%
\begin{pgfscope}%
\pgfsys@transformshift{3.157963in}{1.562127in}%
\pgfsys@useobject{currentmarker}{}%
\end{pgfscope}%
\begin{pgfscope}%
\pgfsys@transformshift{3.167107in}{1.528122in}%
\pgfsys@useobject{currentmarker}{}%
\end{pgfscope}%
\begin{pgfscope}%
\pgfsys@transformshift{3.149016in}{1.556764in}%
\pgfsys@useobject{currentmarker}{}%
\end{pgfscope}%
\begin{pgfscope}%
\pgfsys@transformshift{3.144744in}{1.535511in}%
\pgfsys@useobject{currentmarker}{}%
\end{pgfscope}%
\begin{pgfscope}%
\pgfsys@transformshift{3.153290in}{1.541776in}%
\pgfsys@useobject{currentmarker}{}%
\end{pgfscope}%
\begin{pgfscope}%
\pgfsys@transformshift{3.153641in}{1.593727in}%
\pgfsys@useobject{currentmarker}{}%
\end{pgfscope}%
\begin{pgfscope}%
\pgfsys@transformshift{3.160139in}{1.543702in}%
\pgfsys@useobject{currentmarker}{}%
\end{pgfscope}%
\begin{pgfscope}%
\pgfsys@transformshift{3.140520in}{1.565501in}%
\pgfsys@useobject{currentmarker}{}%
\end{pgfscope}%
\begin{pgfscope}%
\pgfsys@transformshift{3.145025in}{1.618250in}%
\pgfsys@useobject{currentmarker}{}%
\end{pgfscope}%
\end{pgfscope}%
\begin{pgfscope}%
\definecolor{textcolor}{rgb}{0.101961,0.101961,0.101961}%
\pgfsetstrokecolor{textcolor}%
\pgfsetfillcolor{textcolor}%
\pgftext[x=2.994775in,y=1.567077in,left,]{\makebox[0pt][r]{\color{textcolor}\sffamily\fontsize{5.500000}{6.600000}\selectfont 0.91}}%
\end{pgfscope}%
\begin{pgfscope}%
\pgfpathrectangle{\pgfqpoint{1.248750in}{0.440000in}}{\pgfqpoint{1.981350in}{1.628000in}}%
\pgfusepath{clip}%
\pgfsetrectcap%
\pgfsetroundjoin%
\pgfsetlinewidth{1.003750pt}%
\definecolor{currentstroke}{rgb}{0.101961,0.101961,0.101961}%
\pgfsetstrokecolor{currentstroke}%
\pgfsetdash{}{0pt}%
\pgfpathmoveto{\pgfqpoint{2.826823in}{1.169469in}}%
\pgfpathlineto{\pgfqpoint{2.826823in}{1.338531in}}%
\pgfusepath{stroke}%
\end{pgfscope}%
\begin{pgfscope}%
\pgfpathrectangle{\pgfqpoint{1.248750in}{0.440000in}}{\pgfqpoint{1.981350in}{1.628000in}}%
\pgfusepath{clip}%
\pgfsetbuttcap%
\pgfsetroundjoin%
\definecolor{currentfill}{rgb}{0.784314,0.784314,0.784314}%
\pgfsetfillcolor{currentfill}%
\pgfsetfillopacity{0.900000}%
\pgfsetlinewidth{0.200750pt}%
\definecolor{currentstroke}{rgb}{0.101961,0.101961,0.101961}%
\pgfsetstrokecolor{currentstroke}%
\pgfsetstrokeopacity{0.900000}%
\pgfsetdash{}{0pt}%
\pgfsys@defobject{currentmarker}{\pgfqpoint{-0.015528in}{-0.015528in}}{\pgfqpoint{0.015528in}{0.015528in}}{%
\pgfpathmoveto{\pgfqpoint{0.000000in}{-0.015528in}}%
\pgfpathcurveto{\pgfqpoint{0.004118in}{-0.015528in}}{\pgfqpoint{0.008068in}{-0.013892in}}{\pgfqpoint{0.010980in}{-0.010980in}}%
\pgfpathcurveto{\pgfqpoint{0.013892in}{-0.008068in}}{\pgfqpoint{0.015528in}{-0.004118in}}{\pgfqpoint{0.015528in}{0.000000in}}%
\pgfpathcurveto{\pgfqpoint{0.015528in}{0.004118in}}{\pgfqpoint{0.013892in}{0.008068in}}{\pgfqpoint{0.010980in}{0.010980in}}%
\pgfpathcurveto{\pgfqpoint{0.008068in}{0.013892in}}{\pgfqpoint{0.004118in}{0.015528in}}{\pgfqpoint{0.000000in}{0.015528in}}%
\pgfpathcurveto{\pgfqpoint{-0.004118in}{0.015528in}}{\pgfqpoint{-0.008068in}{0.013892in}}{\pgfqpoint{-0.010980in}{0.010980in}}%
\pgfpathcurveto{\pgfqpoint{-0.013892in}{0.008068in}}{\pgfqpoint{-0.015528in}{0.004118in}}{\pgfqpoint{-0.015528in}{0.000000in}}%
\pgfpathcurveto{\pgfqpoint{-0.015528in}{-0.004118in}}{\pgfqpoint{-0.013892in}{-0.008068in}}{\pgfqpoint{-0.010980in}{-0.010980in}}%
\pgfpathcurveto{\pgfqpoint{-0.008068in}{-0.013892in}}{\pgfqpoint{-0.004118in}{-0.015528in}}{\pgfqpoint{0.000000in}{-0.015528in}}%
\pgfpathlineto{\pgfqpoint{0.000000in}{-0.015528in}}%
\pgfpathclose%
\pgfusepath{stroke,fill}%
}%
\begin{pgfscope}%
\pgfsys@transformshift{2.824441in}{1.255258in}%
\pgfsys@useobject{currentmarker}{}%
\end{pgfscope}%
\begin{pgfscope}%
\pgfsys@transformshift{2.832182in}{1.301950in}%
\pgfsys@useobject{currentmarker}{}%
\end{pgfscope}%
\begin{pgfscope}%
\pgfsys@transformshift{2.797467in}{1.216122in}%
\pgfsys@useobject{currentmarker}{}%
\end{pgfscope}%
\begin{pgfscope}%
\pgfsys@transformshift{2.778743in}{1.301757in}%
\pgfsys@useobject{currentmarker}{}%
\end{pgfscope}%
\begin{pgfscope}%
\pgfsys@transformshift{2.796549in}{1.233970in}%
\pgfsys@useobject{currentmarker}{}%
\end{pgfscope}%
\begin{pgfscope}%
\pgfsys@transformshift{2.798228in}{1.245838in}%
\pgfsys@useobject{currentmarker}{}%
\end{pgfscope}%
\begin{pgfscope}%
\pgfsys@transformshift{2.718999in}{1.288883in}%
\pgfsys@useobject{currentmarker}{}%
\end{pgfscope}%
\begin{pgfscope}%
\pgfsys@transformshift{2.709116in}{1.244335in}%
\pgfsys@useobject{currentmarker}{}%
\end{pgfscope}%
\begin{pgfscope}%
\pgfsys@transformshift{2.691640in}{1.259279in}%
\pgfsys@useobject{currentmarker}{}%
\end{pgfscope}%
\begin{pgfscope}%
\pgfsys@transformshift{2.707535in}{1.203710in}%
\pgfsys@useobject{currentmarker}{}%
\end{pgfscope}%
\begin{pgfscope}%
\pgfsys@transformshift{2.645737in}{1.280985in}%
\pgfsys@useobject{currentmarker}{}%
\end{pgfscope}%
\begin{pgfscope}%
\pgfsys@transformshift{2.702691in}{1.258060in}%
\pgfsys@useobject{currentmarker}{}%
\end{pgfscope}%
\begin{pgfscope}%
\pgfsys@transformshift{2.875591in}{1.235876in}%
\pgfsys@useobject{currentmarker}{}%
\end{pgfscope}%
\begin{pgfscope}%
\pgfsys@transformshift{2.861546in}{1.284702in}%
\pgfsys@useobject{currentmarker}{}%
\end{pgfscope}%
\begin{pgfscope}%
\pgfsys@transformshift{2.836235in}{1.233051in}%
\pgfsys@useobject{currentmarker}{}%
\end{pgfscope}%
\begin{pgfscope}%
\pgfsys@transformshift{2.856278in}{1.249050in}%
\pgfsys@useobject{currentmarker}{}%
\end{pgfscope}%
\begin{pgfscope}%
\pgfsys@transformshift{2.815974in}{1.215045in}%
\pgfsys@useobject{currentmarker}{}%
\end{pgfscope}%
\begin{pgfscope}%
\pgfsys@transformshift{2.829205in}{1.243687in}%
\pgfsys@useobject{currentmarker}{}%
\end{pgfscope}%
\begin{pgfscope}%
\pgfsys@transformshift{2.896201in}{1.222434in}%
\pgfsys@useobject{currentmarker}{}%
\end{pgfscope}%
\begin{pgfscope}%
\pgfsys@transformshift{2.891360in}{1.228699in}%
\pgfsys@useobject{currentmarker}{}%
\end{pgfscope}%
\begin{pgfscope}%
\pgfsys@transformshift{2.871510in}{1.280650in}%
\pgfsys@useobject{currentmarker}{}%
\end{pgfscope}%
\begin{pgfscope}%
\pgfsys@transformshift{2.884695in}{1.230625in}%
\pgfsys@useobject{currentmarker}{}%
\end{pgfscope}%
\begin{pgfscope}%
\pgfsys@transformshift{2.870919in}{1.252424in}%
\pgfsys@useobject{currentmarker}{}%
\end{pgfscope}%
\begin{pgfscope}%
\pgfsys@transformshift{2.854619in}{1.305173in}%
\pgfsys@useobject{currentmarker}{}%
\end{pgfscope}%
\end{pgfscope}%
\begin{pgfscope}%
\definecolor{textcolor}{rgb}{0.101961,0.101961,0.101961}%
\pgfsetstrokecolor{textcolor}%
\pgfsetfillcolor{textcolor}%
\pgftext[x=2.602287in,y=1.254000in,left,]{\makebox[0pt][r]{\color{textcolor}\sffamily\fontsize{5.500000}{6.600000}\selectfont 0.68}}%
\end{pgfscope}%
\begin{pgfscope}%
\pgfpathrectangle{\pgfqpoint{1.248750in}{0.440000in}}{\pgfqpoint{1.981350in}{1.628000in}}%
\pgfusepath{clip}%
\pgfsetrectcap%
\pgfsetroundjoin%
\pgfsetlinewidth{1.003750pt}%
\definecolor{currentstroke}{rgb}{0.101961,0.101961,0.101961}%
\pgfsetstrokecolor{currentstroke}%
\pgfsetdash{}{0pt}%
\pgfpathmoveto{\pgfqpoint{2.547488in}{0.856392in}}%
\pgfpathlineto{\pgfqpoint{2.547488in}{1.025454in}}%
\pgfusepath{stroke}%
\end{pgfscope}%
\begin{pgfscope}%
\pgfpathrectangle{\pgfqpoint{1.248750in}{0.440000in}}{\pgfqpoint{1.981350in}{1.628000in}}%
\pgfusepath{clip}%
\pgfsetbuttcap%
\pgfsetroundjoin%
\definecolor{currentfill}{rgb}{0.784314,0.784314,0.784314}%
\pgfsetfillcolor{currentfill}%
\pgfsetfillopacity{0.900000}%
\pgfsetlinewidth{0.200750pt}%
\definecolor{currentstroke}{rgb}{0.101961,0.101961,0.101961}%
\pgfsetstrokecolor{currentstroke}%
\pgfsetstrokeopacity{0.900000}%
\pgfsetdash{}{0pt}%
\pgfsys@defobject{currentmarker}{\pgfqpoint{-0.015528in}{-0.015528in}}{\pgfqpoint{0.015528in}{0.015528in}}{%
\pgfpathmoveto{\pgfqpoint{0.000000in}{-0.015528in}}%
\pgfpathcurveto{\pgfqpoint{0.004118in}{-0.015528in}}{\pgfqpoint{0.008068in}{-0.013892in}}{\pgfqpoint{0.010980in}{-0.010980in}}%
\pgfpathcurveto{\pgfqpoint{0.013892in}{-0.008068in}}{\pgfqpoint{0.015528in}{-0.004118in}}{\pgfqpoint{0.015528in}{0.000000in}}%
\pgfpathcurveto{\pgfqpoint{0.015528in}{0.004118in}}{\pgfqpoint{0.013892in}{0.008068in}}{\pgfqpoint{0.010980in}{0.010980in}}%
\pgfpathcurveto{\pgfqpoint{0.008068in}{0.013892in}}{\pgfqpoint{0.004118in}{0.015528in}}{\pgfqpoint{0.000000in}{0.015528in}}%
\pgfpathcurveto{\pgfqpoint{-0.004118in}{0.015528in}}{\pgfqpoint{-0.008068in}{0.013892in}}{\pgfqpoint{-0.010980in}{0.010980in}}%
\pgfpathcurveto{\pgfqpoint{-0.013892in}{0.008068in}}{\pgfqpoint{-0.015528in}{0.004118in}}{\pgfqpoint{-0.015528in}{0.000000in}}%
\pgfpathcurveto{\pgfqpoint{-0.015528in}{-0.004118in}}{\pgfqpoint{-0.013892in}{-0.008068in}}{\pgfqpoint{-0.010980in}{-0.010980in}}%
\pgfpathcurveto{\pgfqpoint{-0.008068in}{-0.013892in}}{\pgfqpoint{-0.004118in}{-0.015528in}}{\pgfqpoint{0.000000in}{-0.015528in}}%
\pgfpathlineto{\pgfqpoint{0.000000in}{-0.015528in}}%
\pgfpathclose%
\pgfusepath{stroke,fill}%
}%
\begin{pgfscope}%
\pgfsys@transformshift{2.473252in}{0.942181in}%
\pgfsys@useobject{currentmarker}{}%
\end{pgfscope}%
\begin{pgfscope}%
\pgfsys@transformshift{2.479629in}{0.988873in}%
\pgfsys@useobject{currentmarker}{}%
\end{pgfscope}%
\begin{pgfscope}%
\pgfsys@transformshift{2.435120in}{0.903045in}%
\pgfsys@useobject{currentmarker}{}%
\end{pgfscope}%
\begin{pgfscope}%
\pgfsys@transformshift{2.434677in}{0.988680in}%
\pgfsys@useobject{currentmarker}{}%
\end{pgfscope}%
\begin{pgfscope}%
\pgfsys@transformshift{2.450347in}{0.920893in}%
\pgfsys@useobject{currentmarker}{}%
\end{pgfscope}%
\begin{pgfscope}%
\pgfsys@transformshift{2.391325in}{0.932761in}%
\pgfsys@useobject{currentmarker}{}%
\end{pgfscope}%
\begin{pgfscope}%
\pgfsys@transformshift{2.638427in}{0.975806in}%
\pgfsys@useobject{currentmarker}{}%
\end{pgfscope}%
\begin{pgfscope}%
\pgfsys@transformshift{2.600380in}{0.931258in}%
\pgfsys@useobject{currentmarker}{}%
\end{pgfscope}%
\begin{pgfscope}%
\pgfsys@transformshift{2.614218in}{0.946202in}%
\pgfsys@useobject{currentmarker}{}%
\end{pgfscope}%
\begin{pgfscope}%
\pgfsys@transformshift{2.617168in}{0.890634in}%
\pgfsys@useobject{currentmarker}{}%
\end{pgfscope}%
\begin{pgfscope}%
\pgfsys@transformshift{2.567224in}{0.967909in}%
\pgfsys@useobject{currentmarker}{}%
\end{pgfscope}%
\begin{pgfscope}%
\pgfsys@transformshift{2.594398in}{0.944983in}%
\pgfsys@useobject{currentmarker}{}%
\end{pgfscope}%
\begin{pgfscope}%
\pgfsys@transformshift{2.506787in}{0.922799in}%
\pgfsys@useobject{currentmarker}{}%
\end{pgfscope}%
\begin{pgfscope}%
\pgfsys@transformshift{2.527753in}{0.971625in}%
\pgfsys@useobject{currentmarker}{}%
\end{pgfscope}%
\begin{pgfscope}%
\pgfsys@transformshift{2.494823in}{0.919974in}%
\pgfsys@useobject{currentmarker}{}%
\end{pgfscope}%
\begin{pgfscope}%
\pgfsys@transformshift{2.494967in}{0.935973in}%
\pgfsys@useobject{currentmarker}{}%
\end{pgfscope}%
\begin{pgfscope}%
\pgfsys@transformshift{2.515126in}{0.901968in}%
\pgfsys@useobject{currentmarker}{}%
\end{pgfscope}%
\begin{pgfscope}%
\pgfsys@transformshift{2.520083in}{0.930610in}%
\pgfsys@useobject{currentmarker}{}%
\end{pgfscope}%
\begin{pgfscope}%
\pgfsys@transformshift{2.652095in}{0.909357in}%
\pgfsys@useobject{currentmarker}{}%
\end{pgfscope}%
\begin{pgfscope}%
\pgfsys@transformshift{2.628553in}{0.915622in}%
\pgfsys@useobject{currentmarker}{}%
\end{pgfscope}%
\begin{pgfscope}%
\pgfsys@transformshift{2.592931in}{0.967573in}%
\pgfsys@useobject{currentmarker}{}%
\end{pgfscope}%
\begin{pgfscope}%
\pgfsys@transformshift{2.632130in}{0.917548in}%
\pgfsys@useobject{currentmarker}{}%
\end{pgfscope}%
\begin{pgfscope}%
\pgfsys@transformshift{2.597724in}{0.939347in}%
\pgfsys@useobject{currentmarker}{}%
\end{pgfscope}%
\begin{pgfscope}%
\pgfsys@transformshift{2.594224in}{0.992096in}%
\pgfsys@useobject{currentmarker}{}%
\end{pgfscope}%
\end{pgfscope}%
\begin{pgfscope}%
\definecolor{textcolor}{rgb}{0.101961,0.101961,0.101961}%
\pgfsetstrokecolor{textcolor}%
\pgfsetfillcolor{textcolor}%
\pgftext[x=2.347874in,y=0.940923in,left,]{\makebox[0pt][r]{\color{textcolor}\sffamily\fontsize{5.500000}{6.600000}\selectfont 0.54}}%
\end{pgfscope}%
\begin{pgfscope}%
\pgfpathrectangle{\pgfqpoint{1.248750in}{0.440000in}}{\pgfqpoint{1.981350in}{1.628000in}}%
\pgfusepath{clip}%
\pgfsetrectcap%
\pgfsetroundjoin%
\pgfsetlinewidth{1.003750pt}%
\definecolor{currentstroke}{rgb}{0.101961,0.101961,0.101961}%
\pgfsetstrokecolor{currentstroke}%
\pgfsetdash{}{0pt}%
\pgfpathmoveto{\pgfqpoint{2.179469in}{0.543315in}}%
\pgfpathlineto{\pgfqpoint{2.179469in}{0.712377in}}%
\pgfusepath{stroke}%
\end{pgfscope}%
\begin{pgfscope}%
\pgfpathrectangle{\pgfqpoint{1.248750in}{0.440000in}}{\pgfqpoint{1.981350in}{1.628000in}}%
\pgfusepath{clip}%
\pgfsetbuttcap%
\pgfsetroundjoin%
\definecolor{currentfill}{rgb}{0.784314,0.784314,0.784314}%
\pgfsetfillcolor{currentfill}%
\pgfsetfillopacity{0.900000}%
\pgfsetlinewidth{0.200750pt}%
\definecolor{currentstroke}{rgb}{0.101961,0.101961,0.101961}%
\pgfsetstrokecolor{currentstroke}%
\pgfsetstrokeopacity{0.900000}%
\pgfsetdash{}{0pt}%
\pgfsys@defobject{currentmarker}{\pgfqpoint{-0.015528in}{-0.015528in}}{\pgfqpoint{0.015528in}{0.015528in}}{%
\pgfpathmoveto{\pgfqpoint{0.000000in}{-0.015528in}}%
\pgfpathcurveto{\pgfqpoint{0.004118in}{-0.015528in}}{\pgfqpoint{0.008068in}{-0.013892in}}{\pgfqpoint{0.010980in}{-0.010980in}}%
\pgfpathcurveto{\pgfqpoint{0.013892in}{-0.008068in}}{\pgfqpoint{0.015528in}{-0.004118in}}{\pgfqpoint{0.015528in}{0.000000in}}%
\pgfpathcurveto{\pgfqpoint{0.015528in}{0.004118in}}{\pgfqpoint{0.013892in}{0.008068in}}{\pgfqpoint{0.010980in}{0.010980in}}%
\pgfpathcurveto{\pgfqpoint{0.008068in}{0.013892in}}{\pgfqpoint{0.004118in}{0.015528in}}{\pgfqpoint{0.000000in}{0.015528in}}%
\pgfpathcurveto{\pgfqpoint{-0.004118in}{0.015528in}}{\pgfqpoint{-0.008068in}{0.013892in}}{\pgfqpoint{-0.010980in}{0.010980in}}%
\pgfpathcurveto{\pgfqpoint{-0.013892in}{0.008068in}}{\pgfqpoint{-0.015528in}{0.004118in}}{\pgfqpoint{-0.015528in}{0.000000in}}%
\pgfpathcurveto{\pgfqpoint{-0.015528in}{-0.004118in}}{\pgfqpoint{-0.013892in}{-0.008068in}}{\pgfqpoint{-0.010980in}{-0.010980in}}%
\pgfpathcurveto{\pgfqpoint{-0.008068in}{-0.013892in}}{\pgfqpoint{-0.004118in}{-0.015528in}}{\pgfqpoint{0.000000in}{-0.015528in}}%
\pgfpathlineto{\pgfqpoint{0.000000in}{-0.015528in}}%
\pgfpathclose%
\pgfusepath{stroke,fill}%
}%
\begin{pgfscope}%
\pgfsys@transformshift{1.566612in}{0.629105in}%
\pgfsys@useobject{currentmarker}{}%
\end{pgfscope}%
\begin{pgfscope}%
\pgfsys@transformshift{1.572179in}{0.675796in}%
\pgfsys@useobject{currentmarker}{}%
\end{pgfscope}%
\begin{pgfscope}%
\pgfsys@transformshift{1.578096in}{0.589968in}%
\pgfsys@useobject{currentmarker}{}%
\end{pgfscope}%
\begin{pgfscope}%
\pgfsys@transformshift{1.588218in}{0.675603in}%
\pgfsys@useobject{currentmarker}{}%
\end{pgfscope}%
\begin{pgfscope}%
\pgfsys@transformshift{1.582647in}{0.607816in}%
\pgfsys@useobject{currentmarker}{}%
\end{pgfscope}%
\begin{pgfscope}%
\pgfsys@transformshift{1.559707in}{0.619685in}%
\pgfsys@useobject{currentmarker}{}%
\end{pgfscope}%
\begin{pgfscope}%
\pgfsys@transformshift{2.179155in}{0.662729in}%
\pgfsys@useobject{currentmarker}{}%
\end{pgfscope}%
\begin{pgfscope}%
\pgfsys@transformshift{2.192189in}{0.618181in}%
\pgfsys@useobject{currentmarker}{}%
\end{pgfscope}%
\begin{pgfscope}%
\pgfsys@transformshift{2.156062in}{0.633125in}%
\pgfsys@useobject{currentmarker}{}%
\end{pgfscope}%
\begin{pgfscope}%
\pgfsys@transformshift{2.160355in}{0.577557in}%
\pgfsys@useobject{currentmarker}{}%
\end{pgfscope}%
\begin{pgfscope}%
\pgfsys@transformshift{2.182520in}{0.654832in}%
\pgfsys@useobject{currentmarker}{}%
\end{pgfscope}%
\begin{pgfscope}%
\pgfsys@transformshift{2.179783in}{0.631906in}%
\pgfsys@useobject{currentmarker}{}%
\end{pgfscope}%
\begin{pgfscope}%
\pgfsys@transformshift{2.804772in}{0.609722in}%
\pgfsys@useobject{currentmarker}{}%
\end{pgfscope}%
\begin{pgfscope}%
\pgfsys@transformshift{2.793245in}{0.658548in}%
\pgfsys@useobject{currentmarker}{}%
\end{pgfscope}%
\begin{pgfscope}%
\pgfsys@transformshift{2.755289in}{0.606897in}%
\pgfsys@useobject{currentmarker}{}%
\end{pgfscope}%
\begin{pgfscope}%
\pgfsys@transformshift{2.794054in}{0.622896in}%
\pgfsys@useobject{currentmarker}{}%
\end{pgfscope}%
\begin{pgfscope}%
\pgfsys@transformshift{2.770269in}{0.588891in}%
\pgfsys@useobject{currentmarker}{}%
\end{pgfscope}%
\begin{pgfscope}%
\pgfsys@transformshift{2.818429in}{0.617533in}%
\pgfsys@useobject{currentmarker}{}%
\end{pgfscope}%
\end{pgfscope}%
\begin{pgfscope}%
\definecolor{textcolor}{rgb}{0.101961,0.101961,0.101961}%
\pgfsetstrokecolor{textcolor}%
\pgfsetfillcolor{textcolor}%
\pgftext[x=1.516257in,y=0.627846in,left,]{\makebox[0pt][r]{\color{textcolor}\sffamily\fontsize{5.500000}{6.600000}\selectfont 0.06}}%
\end{pgfscope}%
\begin{pgfscope}%
\definecolor{textcolor}{rgb}{0.101961,0.101961,0.101961}%
\pgfsetstrokecolor{textcolor}%
\pgfsetfillcolor{textcolor}%
\pgftext[x=1.189310in,y=1.880154in,left,]{\makebox[0pt][r]{\color{textcolor}\sffamily\fontsize{7.000000}{8.400000}\bfseries\selectfont \textbf{COVEX (old: rung 3)}}}%
\end{pgfscope}%
\begin{pgfscope}%
\definecolor{textcolor}{rgb}{0.101961,0.101961,0.101961}%
\pgfsetstrokecolor{textcolor}%
\pgfsetfillcolor{textcolor}%
\pgftext[x=1.189310in,y=1.567077in,left,]{\makebox[0pt][r]{\color{textcolor}\sffamily\fontsize{7.000000}{8.400000}\selectfont SimpleStream~\cite{simplestream2026}}}%
\end{pgfscope}%
\begin{pgfscope}%
\definecolor{textcolor}{rgb}{0.101961,0.101961,0.101961}%
\pgfsetstrokecolor{textcolor}%
\pgfsetfillcolor{textcolor}%
\pgftext[x=1.189310in,y=1.254000in,left,]{\makebox[0pt][r]{\color{textcolor}\sffamily\fontsize{7.000000}{8.400000}\selectfont StreamScout~\cite{streamscout2026}}}%
\end{pgfscope}%
\begin{pgfscope}%
\definecolor{textcolor}{rgb}{0.101961,0.101961,0.101961}%
\pgfsetstrokecolor{textcolor}%
\pgfsetfillcolor{textcolor}%
\pgftext[x=1.189310in,y=0.940923in,left,]{\makebox[0pt][r]{\color{textcolor}\sffamily\fontsize{7.000000}{8.400000}\selectfont LLF~\cite{mok1983llf}}}%
\end{pgfscope}%
\begin{pgfscope}%
\definecolor{textcolor}{rgb}{0.101961,0.101961,0.101961}%
\pgfsetstrokecolor{textcolor}%
\pgfsetfillcolor{textcolor}%
\pgftext[x=1.189310in,y=0.627846in,left,]{\makebox[0pt][r]{\color{textcolor}\sffamily\fontsize{7.000000}{8.400000}\selectfont ECO-LLM~\cite{ecollm2025}}}%
\end{pgfscope}%
\begin{pgfscope}%
\pgfpathrectangle{\pgfqpoint{1.248750in}{0.440000in}}{\pgfqpoint{1.981350in}{1.628000in}}%
\pgfusepath{clip}%
\pgfsetbuttcap%
\pgfsetroundjoin%
\pgfsetlinewidth{0.803000pt}%
\definecolor{currentstroke}{rgb}{0.101961,0.101961,0.101961}%
\pgfsetstrokecolor{currentstroke}%
\pgfsetdash{}{0pt}%
\pgfpathmoveto{\pgfqpoint{2.625727in}{1.823795in}}%
\pgfpathcurveto{\pgfqpoint{2.635118in}{1.823795in}}{\pgfqpoint{2.644126in}{1.827526in}}{\pgfqpoint{2.650766in}{1.834166in}}%
\pgfpathcurveto{\pgfqpoint{2.657406in}{1.840806in}}{\pgfqpoint{2.661137in}{1.849814in}}{\pgfqpoint{2.661137in}{1.859205in}}%
\pgfpathcurveto{\pgfqpoint{2.661137in}{1.868595in}}{\pgfqpoint{2.657406in}{1.877603in}}{\pgfqpoint{2.650766in}{1.884243in}}%
\pgfpathcurveto{\pgfqpoint{2.644126in}{1.890884in}}{\pgfqpoint{2.635118in}{1.894615in}}{\pgfqpoint{2.625727in}{1.894615in}}%
\pgfpathcurveto{\pgfqpoint{2.616337in}{1.894615in}}{\pgfqpoint{2.607329in}{1.890884in}}{\pgfqpoint{2.600689in}{1.884243in}}%
\pgfpathcurveto{\pgfqpoint{2.594049in}{1.877603in}}{\pgfqpoint{2.590318in}{1.868595in}}{\pgfqpoint{2.590318in}{1.859205in}}%
\pgfpathcurveto{\pgfqpoint{2.590318in}{1.849814in}}{\pgfqpoint{2.594049in}{1.840806in}}{\pgfqpoint{2.600689in}{1.834166in}}%
\pgfpathcurveto{\pgfqpoint{2.607329in}{1.827526in}}{\pgfqpoint{2.616337in}{1.823795in}}{\pgfqpoint{2.625727in}{1.823795in}}%
\pgfpathlineto{\pgfqpoint{2.625727in}{1.823795in}}%
\pgfpathclose%
\pgfusepath{stroke}%
\end{pgfscope}%
\begin{pgfscope}%
\pgfpathrectangle{\pgfqpoint{1.248750in}{0.440000in}}{\pgfqpoint{1.981350in}{1.628000in}}%
\pgfusepath{clip}%
\pgfsetbuttcap%
\pgfsetroundjoin%
\pgfsetlinewidth{0.803000pt}%
\definecolor{currentstroke}{rgb}{0.101961,0.101961,0.101961}%
\pgfsetstrokecolor{currentstroke}%
\pgfsetdash{}{0pt}%
\pgfpathmoveto{\pgfqpoint{3.038225in}{1.535727in}}%
\pgfpathcurveto{\pgfqpoint{3.047616in}{1.535727in}}{\pgfqpoint{3.056624in}{1.539458in}}{\pgfqpoint{3.063264in}{1.546099in}}%
\pgfpathcurveto{\pgfqpoint{3.069904in}{1.552739in}}{\pgfqpoint{3.073635in}{1.561746in}}{\pgfqpoint{3.073635in}{1.571137in}}%
\pgfpathcurveto{\pgfqpoint{3.073635in}{1.580528in}}{\pgfqpoint{3.069904in}{1.589535in}}{\pgfqpoint{3.063264in}{1.596176in}}%
\pgfpathcurveto{\pgfqpoint{3.056624in}{1.602816in}}{\pgfqpoint{3.047616in}{1.606547in}}{\pgfqpoint{3.038225in}{1.606547in}}%
\pgfpathcurveto{\pgfqpoint{3.028835in}{1.606547in}}{\pgfqpoint{3.019827in}{1.602816in}}{\pgfqpoint{3.013187in}{1.596176in}}%
\pgfpathcurveto{\pgfqpoint{3.006547in}{1.589535in}}{\pgfqpoint{3.002816in}{1.580528in}}{\pgfqpoint{3.002816in}{1.571137in}}%
\pgfpathcurveto{\pgfqpoint{3.002816in}{1.561746in}}{\pgfqpoint{3.006547in}{1.552739in}}{\pgfqpoint{3.013187in}{1.546099in}}%
\pgfpathcurveto{\pgfqpoint{3.019827in}{1.539458in}}{\pgfqpoint{3.028835in}{1.535727in}}{\pgfqpoint{3.038225in}{1.535727in}}%
\pgfpathlineto{\pgfqpoint{3.038225in}{1.535727in}}%
\pgfpathclose%
\pgfusepath{stroke}%
\end{pgfscope}%
\begin{pgfscope}%
\pgfpathrectangle{\pgfqpoint{1.248750in}{0.440000in}}{\pgfqpoint{1.981350in}{1.628000in}}%
\pgfusepath{clip}%
\pgfsetbuttcap%
\pgfsetroundjoin%
\pgfsetlinewidth{0.803000pt}%
\definecolor{currentstroke}{rgb}{0.101961,0.101961,0.101961}%
\pgfsetstrokecolor{currentstroke}%
\pgfsetdash{}{0pt}%
\pgfpathmoveto{\pgfqpoint{2.645737in}{1.245576in}}%
\pgfpathcurveto{\pgfqpoint{2.655128in}{1.245576in}}{\pgfqpoint{2.664136in}{1.249307in}}{\pgfqpoint{2.670776in}{1.255947in}}%
\pgfpathcurveto{\pgfqpoint{2.677416in}{1.262587in}}{\pgfqpoint{2.681147in}{1.271595in}}{\pgfqpoint{2.681147in}{1.280985in}}%
\pgfpathcurveto{\pgfqpoint{2.681147in}{1.290376in}}{\pgfqpoint{2.677416in}{1.299384in}}{\pgfqpoint{2.670776in}{1.306024in}}%
\pgfpathcurveto{\pgfqpoint{2.664136in}{1.312664in}}{\pgfqpoint{2.655128in}{1.316395in}}{\pgfqpoint{2.645737in}{1.316395in}}%
\pgfpathcurveto{\pgfqpoint{2.636347in}{1.316395in}}{\pgfqpoint{2.627339in}{1.312664in}}{\pgfqpoint{2.620699in}{1.306024in}}%
\pgfpathcurveto{\pgfqpoint{2.614058in}{1.299384in}}{\pgfqpoint{2.610327in}{1.290376in}}{\pgfqpoint{2.610327in}{1.280985in}}%
\pgfpathcurveto{\pgfqpoint{2.610327in}{1.271595in}}{\pgfqpoint{2.614058in}{1.262587in}}{\pgfqpoint{2.620699in}{1.255947in}}%
\pgfpathcurveto{\pgfqpoint{2.627339in}{1.249307in}}{\pgfqpoint{2.636347in}{1.245576in}}{\pgfqpoint{2.645737in}{1.245576in}}%
\pgfpathlineto{\pgfqpoint{2.645737in}{1.245576in}}%
\pgfpathclose%
\pgfusepath{stroke}%
\end{pgfscope}%
\begin{pgfscope}%
\pgfpathrectangle{\pgfqpoint{1.248750in}{0.440000in}}{\pgfqpoint{1.981350in}{1.628000in}}%
\pgfusepath{clip}%
\pgfsetbuttcap%
\pgfsetroundjoin%
\pgfsetlinewidth{0.803000pt}%
\definecolor{currentstroke}{rgb}{0.101961,0.101961,0.101961}%
\pgfsetstrokecolor{currentstroke}%
\pgfsetdash{}{0pt}%
\pgfpathmoveto{\pgfqpoint{2.391325in}{0.897352in}}%
\pgfpathcurveto{\pgfqpoint{2.400716in}{0.897352in}}{\pgfqpoint{2.409723in}{0.901083in}}{\pgfqpoint{2.416364in}{0.907723in}}%
\pgfpathcurveto{\pgfqpoint{2.423004in}{0.914363in}}{\pgfqpoint{2.426735in}{0.923371in}}{\pgfqpoint{2.426735in}{0.932761in}}%
\pgfpathcurveto{\pgfqpoint{2.426735in}{0.942152in}}{\pgfqpoint{2.423004in}{0.951160in}}{\pgfqpoint{2.416364in}{0.957800in}}%
\pgfpathcurveto{\pgfqpoint{2.409723in}{0.964440in}}{\pgfqpoint{2.400716in}{0.968171in}}{\pgfqpoint{2.391325in}{0.968171in}}%
\pgfpathcurveto{\pgfqpoint{2.381934in}{0.968171in}}{\pgfqpoint{2.372927in}{0.964440in}}{\pgfqpoint{2.366286in}{0.957800in}}%
\pgfpathcurveto{\pgfqpoint{2.359646in}{0.951160in}}{\pgfqpoint{2.355915in}{0.942152in}}{\pgfqpoint{2.355915in}{0.932761in}}%
\pgfpathcurveto{\pgfqpoint{2.355915in}{0.923371in}}{\pgfqpoint{2.359646in}{0.914363in}}{\pgfqpoint{2.366286in}{0.907723in}}%
\pgfpathcurveto{\pgfqpoint{2.372927in}{0.901083in}}{\pgfqpoint{2.381934in}{0.897352in}}{\pgfqpoint{2.391325in}{0.897352in}}%
\pgfpathlineto{\pgfqpoint{2.391325in}{0.897352in}}%
\pgfpathclose%
\pgfusepath{stroke}%
\end{pgfscope}%
\begin{pgfscope}%
\pgfpathrectangle{\pgfqpoint{1.248750in}{0.440000in}}{\pgfqpoint{1.981350in}{1.628000in}}%
\pgfusepath{clip}%
\pgfsetbuttcap%
\pgfsetroundjoin%
\pgfsetlinewidth{0.803000pt}%
\definecolor{currentstroke}{rgb}{0.101961,0.101961,0.101961}%
\pgfsetstrokecolor{currentstroke}%
\pgfsetdash{}{0pt}%
\pgfpathmoveto{\pgfqpoint{1.559707in}{0.584275in}}%
\pgfpathcurveto{\pgfqpoint{1.569098in}{0.584275in}}{\pgfqpoint{1.578105in}{0.588006in}}{\pgfqpoint{1.584746in}{0.594646in}}%
\pgfpathcurveto{\pgfqpoint{1.591386in}{0.601286in}}{\pgfqpoint{1.595117in}{0.610294in}}{\pgfqpoint{1.595117in}{0.619685in}}%
\pgfpathcurveto{\pgfqpoint{1.595117in}{0.629075in}}{\pgfqpoint{1.591386in}{0.638083in}}{\pgfqpoint{1.584746in}{0.644723in}}%
\pgfpathcurveto{\pgfqpoint{1.578105in}{0.651363in}}{\pgfqpoint{1.569098in}{0.655094in}}{\pgfqpoint{1.559707in}{0.655094in}}%
\pgfpathcurveto{\pgfqpoint{1.550316in}{0.655094in}}{\pgfqpoint{1.541309in}{0.651363in}}{\pgfqpoint{1.534669in}{0.644723in}}%
\pgfpathcurveto{\pgfqpoint{1.528028in}{0.638083in}}{\pgfqpoint{1.524297in}{0.629075in}}{\pgfqpoint{1.524297in}{0.619685in}}%
\pgfpathcurveto{\pgfqpoint{1.524297in}{0.610294in}}{\pgfqpoint{1.528028in}{0.601286in}}{\pgfqpoint{1.534669in}{0.594646in}}%
\pgfpathcurveto{\pgfqpoint{1.541309in}{0.588006in}}{\pgfqpoint{1.550316in}{0.584275in}}{\pgfqpoint{1.559707in}{0.584275in}}%
\pgfpathlineto{\pgfqpoint{1.559707in}{0.584275in}}%
\pgfpathclose%
\pgfusepath{stroke}%
\end{pgfscope}%
\begin{pgfscope}%
\definecolor{textcolor}{rgb}{0.549020,0.549020,0.549020}%
\pgfsetstrokecolor{textcolor}%
\pgfsetfillcolor{textcolor}%
\pgftext[x=3.296700in,y=2.167000in,right,top]{\color{textcolor}\sffamily\fontsize{5.500000}{6.600000}\selectfont OLD DATA}%
\end{pgfscope}%
\end{pgfpicture}%
\makeatother%
\endgroup%
}
\caption{[PLACEHOLDER: OLD DATA, rung 3 stand-in for COVEX, baselines on dev seeds 4-7; dashcam only] Per-device gain in dashcam on-time attainment over the on-device model (which meets none: attainment 0 in every seed), for COVEX and the four strongest baselines. One dot per device and seed (seeds 4 to 7, standard condition); boxes show quartiles, the bar the median, and the ringed dot with its value the worst device. Robot and single-shot are not shown: the committed per-device table has no on-device reference there.}
\label{fig:F4_per_device}
\end{figure}
\clearpage
\bibliographystyle{ACM-Reference-Format}
\bibliography{refs_stub}
\end{document}
